\IfFileExists{stacks-project-book.cls}{%
\documentclass{stacks-project-book}
}{%
\documentclass{amsbook}
}

\usepackage{verbatim}
\newenvironment{reference}{\comment}{\endcomment}
\newenvironment{slogan}{\comment}{\endcomment}
\newenvironment{history}{\comment}{\endcomment}

\usepackage[all]{xy}

\xyoption{2cell}
\UseAllTwocells

\usepackage{multicol}

\usepackage{lmodern}
\usepackage[T1]{fontenc}
\usepackage[french]{babel}
\AtBeginDocument{\shorthandoff{?}}


\usepackage{hyperref}


\theoremstyle{plain}
\newtheorem{theorem}[subsection]{Théorème}
\newtheorem{proposition}[subsection]{Proposition}
\newtheorem{lemma}[subsection]{Lemme}

\theoremstyle{definition}
\newtheorem{definition}[subsection]{Définition}
\newtheorem{example}[subsection]{Exemple}
\newtheorem{exercise}[subsection]{Exercice}
\newtheorem{situation}[subsection]{Situation}

\theoremstyle{remark}
\newtheorem{remark}[subsection]{Remarque}
\newtheorem{remarks}[subsection]{Remarques}

\renewcommand{\proofname}{Démonstration}

\numberwithin{equation}{subsection}

\def\lim{\mathop{\mathrm{lim}}\nolimits}
\def\colim{\mathop{\mathrm{colim}}\nolimits}
\def\Spec{\mathop{\mathrm{Spec}}}
\def\Hom{\mathop{\mathrm{Hom}}\nolimits}
\def\Ext{\mathop{\mathrm{Ext}}\nolimits}
\def\SheafHom{\mathop{\mathcal{H}\!\mathit{om}}\nolimits}
\def\SheafExt{\mathop{\mathcal{E}\!\mathit{xt}}\nolimits}
\def\Sch{\mathit{Sch}}
\def\Mor{\mathop{\mathrm{Mor}}\nolimits}
\def\Ob{\mathop{\mathrm{Ob}}\nolimits}
\def\Sh{\mathop{\mathit{Sh}}\nolimits}
\def\NL{\mathop{N\!L}\nolimits}
\def\CH{\mathop{\mathrm{CH}}\nolimits}
\def\proetale{{pro\text{-}\acute{e}tale}}
\def\etale{{\acute{e}tale}}
\def\QCoh{\mathit{QCoh}}
\def\Ker{\mathop{\mathrm{Ker}}}
\def\Im{\mathop{\mathrm{Im}}}
\def\Coker{\mathop{\mathrm{Coker}}}
\def\Coim{\mathop{\mathrm{Coim}}}

\DeclareMathSymbol{\boxtimes}{\mathbin}{AMSa}{"02}

\def\QCohstack{\mathcal{QC}\!\mathit{oh}}
\def\Cohstack{\mathcal{C}\!\mathit{oh}}
\def\Spacesstack{\mathcal{S}\!\mathit{paces}}
\def\Quotfunctor{\mathrm{Quot}}
\def\Hilbfunctor{\mathrm{Hilb}}
\def\Curvesstack{\mathcal{C}\!\mathit{urves}}
\def\Polarizedstack{\mathcal{P}\!\mathit{olarized}}
\def\Complexesstack{\mathcal{C}\!\mathit{omplexes}}
\def\Pic{\mathop{\mathrm{Pic}}\nolimits}
\def\Picardstack{\mathcal{P}\!\mathit{ic}}
\def\Picardfunctor{\mathrm{Pic}}
\def\Deformationcategory{\mathcal{D}\!\mathit{ef}}
\begin{document}
\title{Le projet Stacks --- édition française, partie 12}
\maketitle
\tableofcontents

% OK, start here.
%

\chapter{Limites de champs algébriques}



\phantomsection
\label{stacks-limits-section-phantom}


\section{Introduction}
\label{stacks-limits-section-introduction}

\noindent
Dans ce chapitre, nous rassemblons des résultats relatifs aux limites de
champs algébriques. De nombreux résultats sur les limites de champs
algébriques et d'espaces algébriques ont été obtenus par David Rydh dans
\cite{rydh_approx}.


\section{Conventions}
\label{stacks-limits-section-conventions}

\noindent
Nous continuons d'employer les conventions et l'abus de langage introduits
dans Propriétés des champs, section
\ref{stacks-properties-section-conventions}.







\section{Morphismes de présentation finie}
\label{stacks-limits-section-finite-presentation}

\noindent
Cette section est l'analogue de Limites d'espaces, section
\ref{spaces-limits-section-finite-presentation}. Nous y avons défini ce que
signifie, pour une transformation de foncteurs sur $\Sch$, le fait de commuter
aux limites (nous suggérons de consulter la caractérisation de Limites
d'espaces, lemme
\ref{spaces-limits-lemma-characterize-relative-limit-preserving}). Dans
Critères de représentabilité, section \ref{criteria-section-limit-preserving},
nous avons défini la notion de « commutation aux limites sur les objets ».
Rappelons enfin que, dans Axiomes d'Artin, section \ref{artin-section-limits},
nous avons défini ce que signifie, pour une catégorie fibrée en groupoïdes sur
$\Sch$, le fait de commuter aux limites. En combinant ces notions, nous
obtenons la définition suivante.

\begin{definition}
\label{stacks-limits-definition-limit-preserving}
Soit $S$ un schéma. Soit $f : \mathcal{X} \to \mathcal{Y}$ un
$1$-morphisme de catégories fibrées en groupoïdes sur $(\Sch/S)_{fppf}$.
Nous disons que $f$ {\it commute aux limites} si, pour toute limite projective
filtrante $U = \lim U_i$ de schémas affines sur $S$, le diagramme
$$
\xymatrix{
\colim \mathcal{X}_{U_i} \ar[r] \ar[d]_f & \mathcal{X}_U \ar[d]^f \\
\colim \mathcal{Y}_{U_i} \ar[r] & \mathcal{Y}_U
}
$$
de catégories fibres est $2$-cartésien.
\end{definition}

\begin{lemma}
\label{stacks-limits-lemma-limit-preserving-objects}
Soit $S$ un schéma. Soit $f : \mathcal{X} \to \mathcal{Y}$ un
$1$-morphisme de catégories fibrées en groupoïdes sur $(\Sch/S)_{fppf}$.
Si $f$ commute aux limites (définition \ref{stacks-limits-definition-limit-preserving}),
alors $f$ commute aux limites sur les objets (Critères de représentabilité,
section \ref{criteria-section-limit-preserving}).
\end{lemma}

\begin{proof}
Si, pour toute limite projective filtrante $U = \lim U_i$ de schémas affines
sur $U$, le foncteur
$$
\colim \mathcal{X}_{U_i} \longrightarrow
(\colim \mathcal{Y}_{U_i}) \times_{\mathcal{Y}_U} \mathcal{X}_U
$$
est essentiellement surjectif, alors $f$ commute aux limites sur les objets.
\end{proof}

\begin{lemma}
\label{stacks-limits-lemma-base-change-limit-preserving}
Soient $p : \mathcal{X} \to \mathcal{Y}$ et
$q : \mathcal{Z} \to \mathcal{Y}$ des $1$-morphismes de catégories fibrées
en groupoïdes sur $(\Sch/S)_{fppf}$. Si
$p : \mathcal{X} \to \mathcal{Y}$ commute aux limites, il en va de même du
changement de base
$p' : \mathcal{X} \times_\mathcal{Y} \mathcal{Z} \to \mathcal{Z}$
de $p$ suivant $q$.
\end{lemma}

\begin{proof}
C'est formel. Soit $U = \lim_{i \in I} U_i$ la limite projective filtrante
de schémas affines $U_i$ sur $S$. Pour tout $i$, nous avons
$$
(\mathcal{X} \times_\mathcal{Y} \mathcal{Z})_{U_i} =
\mathcal{X}_{U_i} \times_{\mathcal{Y}_{U_i}} \mathcal{Z}_{U_i}
$$
Les colimites filtrantes commutent aux $2$-produits fibrés de catégories
(nous omettons les détails). Puisque $p$ commute aux limites, il vient
\begin{align*}
\colim (\mathcal{X} \times_\mathcal{Y} \mathcal{Z})_{U_i}
& =
\colim \mathcal{X}_{U_i} \times_{\colim \mathcal{Y}_{U_i}}
\colim \mathcal{Z}_{U_i} \\
& =
\mathcal{X}_U \times_{\mathcal{Y}_U} \colim \mathcal{Y}_{U_i}
\times_{\colim \mathcal{Y}_{U_i}}
\colim \mathcal{Z}_{U_i} \\
& =
\mathcal{X}_U \times_{\mathcal{Y}_U} \colim \mathcal{Z}_{U_i} \\
& =
\mathcal{X}_U \times_{\mathcal{Y}_U} \mathcal{Z}_U \times_{\mathcal{Z}_U}
\colim \mathcal{Z}_{U_i} \\
& =
(\mathcal{X} \times_\mathcal{Y} \mathcal{Z})_U \times_{\mathcal{Z}_U}
\colim \mathcal{Z}_{U_i}
\end{align*}
comme voulu.
\end{proof}

\begin{lemma}
\label{stacks-limits-lemma-composition-limit-preserving}
Soient $p : \mathcal{X} \to \mathcal{Y}$ et
$q : \mathcal{Y} \to \mathcal{Z}$ des $1$-morphismes de catégories fibrées
en groupoïdes sur $(\Sch/S)_{fppf}$. Si $p$ et $q$ commutent aux limites,
leur composé $q \circ p$ y commute également.
\end{lemma}

\begin{proof}
C'est formel. Soit $U = \lim_{i \in I} U_i$ la limite projective filtrante
de schémas affines $U_i$ sur $S$. Puisque $p$ et $q$ commutent aux limites,
nous obtenons
\begin{align*}
\colim \mathcal{X}_{U_i}
& =
\mathcal{X}_U \times_{\mathcal{Y}_U} \colim \mathcal{Y}_{U_i} \\
& =
\mathcal{X}_U \times_{\mathcal{Y}_U} \mathcal{Y}_U
\times_{\mathcal{Z}_U} \colim \mathcal{Z}_{U_i} \\
& =
\mathcal{X}_U \times_{\mathcal{Z}_U} \colim \mathcal{Z}_{U_i}
\end{align*}
comme voulu.
\end{proof}

\begin{lemma}
\label{stacks-limits-lemma-representable-by-spaces-limit-preserving}
Soit $p : \mathcal{X} \to \mathcal{Y}$ un $1$-morphisme de catégories
fibrées en groupoïdes sur $(\Sch/S)_{fppf}$. Si $p$ est représentable par
des espaces algébriques, les conditions suivantes sont équivalentes :
\begin{enumerate}
\item $p$ commute aux limites,
\item $p$ commute aux limites sur les objets, et
\item $p$ est localement de présentation finie (voir Champs algébriques,
définition \ref{algebraic-definition-relative-representable-property}).
\end{enumerate}
\end{lemma}

\begin{proof}
Dans Critères de représentabilité, lemme
\ref{criteria-lemma-representable-by-spaces-limit-preserving}, nous avons vu
que (2) et (3) sont équivalentes. Il suffit donc de montrer l'équivalence de
(1) et (2). Une implication a été établie dans le lemme
\ref{stacks-limits-lemma-limit-preserving-objects}. Pour la réciproque, soit
$U = \lim_{i \in I} U_i$ la limite projective filtrante de schémas affines
$U_i$ sur $S$. Nous devons montrer que
$$
\colim \mathcal{X}_{U_i} \longrightarrow
\mathcal{X}_U \times_{\mathcal{Y}_U} \colim \mathcal{Y}_{U_i}
$$
est une équivalence. L'hypothèse (2) assure qu'il est essentiellement
surjectif ; il reste donc à prouver qu'il est pleinement fidèle. Comme $p$
est fidèle sur les catégories fibres (Champs algébriques, lemme
\ref{algebraic-lemma-criterion-map-representable-spaces-fibred-in-groupoids}),
ce foncteur est fidèle. Soient $x_i$ et $x'_i$ des objets de la catégorie
fibre de $\mathcal{X}$ sur $U_i$. Le foncteur ci-dessus envoie $x_i$ sur
$(x_i|_U, p(x_i), can)$, où $can$ est l'isomorphisme canonique
$p(x_i|_U) \to p(x_i)|_U$. Supposons donc donné un morphisme
$$
(\alpha, \beta_i) : (x_i|_U, p(x_i), can) \longrightarrow
(x'_i|_U, p(x'_i), can)
$$
dans la catégorie qui figure au membre de droite de la première flèche
affichée de cette démonstration. Il s'agit de construire un $i' \geq i$ et un
morphisme $x_i|_{U_{i'}} \to x'_i|_{U_{i'}}$ dont l'image soit
$(\alpha, \beta_i|_{U_{i'}})$.

\medskip\noindent
Posons $y_i = p(x_i)$ et $y'_i = p(x'_i)$. D'après Champs algébriques,
lemme
\ref{algebraic-lemma-criterion-map-representable-spaces-fibred-in-groupoids},
le foncteur
$$
X_{y_i} : (\Sch/U_i)^{opp} \to \textit{Ensembles},\quad
V/U_i \mapsto
\{(x, \phi) \mid x \in \Ob(\mathcal{X}_V), \phi : f(x) \to y_i|V\}/\cong
$$
est un espace algébrique sur $U_i$, et il en va de même du foncteur
$X_{y'_i}$ défini de façon analogue. Puisque (2) équivaut à (3), nous voyons
que $X_{y'_i}$ est localement de présentation finie sur $U_i$. Remarquons que
$(x_i, \text{id})$ et $(x'_i, \text{id})$ définissent des points à valeurs
dans $U_i$ de $X_{y_i}$ et de $X_{y'_i}$. Il existe une transformation de
foncteurs
$$
\beta_i : X_{y_i} \to X_{y'_i},\quad
(x/V, \phi) \mapsto (x/V, \beta_i|_V \circ \phi)
$$
autrement dit, un morphisme d'espaces algébriques sur $U_i$. Nous affirmons
que le diagramme
$$
\xymatrix{
U \ar[d] \ar[rr] & & U_i \ar[d]^{(x'_i, \text{id})} \\
U_i \ar[r]^{(x_i, \text{id})} & X_{y_i} \ar[r]^{\beta_i} & X_{y'_i}
}
$$
est commutatif. En effet, cela équivaut à dire que les couples
$(x_i|_U, \beta_i|_U)$ et $(x'_i|_U, \text{id})$ intervenant dans la
définition du foncteur $X_{y'_i}$ sont isomorphes. Or le morphisme
$\alpha : x_i|_U \to x'_i|_U$ fournit précisément un tel isomorphisme. En
remontant cet argument, on voit qu'il suffit de trouver un $i' \geq i$ pour
lequel le diagramme
$$
\xymatrix{
U_{i'} \ar[d] \ar[rr] & & U_i \ar[d]^{(x'_i, \text{id})} \\
U_i \ar[r]^{(x_i, \text{id})} & X_{y_i} \ar[r]^{\beta_i} & X_{y'_i}
}
$$
soit commutatif : on obtient alors un isomorphisme
$x_i|_{U_{i'}} \to x'_i|_{U_{i'}}$ qui résout le problème posé au paragraphe
précédent. Or le morphisme diagonal
$$
\Delta : X_{y'_i} \to X_{y'_i} \times_{U_i} X_{y'_i}
$$
est localement de présentation finie (Morphismes d'espaces, lemme
\ref{spaces-morphisms-lemma-diagonal-morphism-finite-type}). Puisque
$U \to U_i$ égalise les deux morphismes vers $X_{y'_i}$, il existe donc un
$i' \geq i$ tel que $U_{i'} \to U_i$ les égalise également ; voir Limites
d'espaces, proposition
\ref{spaces-limits-proposition-characterize-locally-finite-presentation}.
\end{proof}

\begin{lemma}
\label{stacks-limits-lemma-limit-preserving-diagonal}
Soit $p : \mathcal{X} \to \mathcal{Y}$ un $1$-morphisme de catégories
fibrées en groupoïdes sur $(\Sch/S)_{fppf}$. Les conditions suivantes sont
équivalentes :
\begin{enumerate}
\item la diagonale
$\Delta : \mathcal{X} \to \mathcal{X} \times_\mathcal{Y} \mathcal{X}$
commute aux limites, et
\item pour toute limite projective filtrante $U = \lim U_i$ de schémas
affines sur $S$, le foncteur
$$
\colim \mathcal{X}_{U_i} \longrightarrow
\mathcal{X}_U \times_{\mathcal{Y}_U} \colim \mathcal{Y}_{U_i}
$$
est pleinement fidèle.
\end{enumerate}
En particulier, si $p$ commute aux limites, il en va de même de $\Delta$.
\end{lemma}

\begin{proof}
Soit $U = \lim U_i$ une limite projective filtrante de schémas affines sur
$S$. Nous affirmons que le foncteur
$$
\colim \mathcal{X}_{U_i} \longrightarrow
\mathcal{X}_U \times_{\mathcal{Y}_U} \colim \mathcal{Y}_{U_i}
$$
est pleinement fidèle si et seulement si le foncteur
$$
\colim \mathcal{X}_{U_i} \longrightarrow
\mathcal{X}_U \times_{(\mathcal{X} \times_\mathcal{Y} \mathcal{X})_U}
\colim (\mathcal{X} \times_\mathcal{Y} \mathcal{X})_{U_i}
$$
est une équivalence. Cela démontrera le lemme. Comme
$(\mathcal{X} \times_\mathcal{Y} \mathcal{X})_U =
\mathcal{X}_U \times_{\mathcal{Y}_U} \mathcal{X}_U$
et
$(\mathcal{X} \times_\mathcal{Y} \mathcal{X})_{U_i} =
\mathcal{X}_{U_i} \times_{\mathcal{Y}_{U_i}} \mathcal{X}_{U_i}$
il s'agit d'une assertion purement catégorique, que nous examinons au
paragraphe suivant.

\medskip\noindent
Soit $\mathcal{I}$ une catégorie filtrante d'indices. Soient
$(\mathcal{C}_i)$ et $(\mathcal{D}_i)$ des systèmes de groupoïdes sur
$\mathcal{I}$. Soit
$p : (\mathcal{C}_i) \to (\mathcal{D}_i)$ un morphisme de systèmes de
groupoïdes sur $\mathcal{I}$. Supposons donnés un foncteur de groupoïdes
$p : \mathcal{C} \to \mathcal{D}$ et des foncteurs
$f : \colim \mathcal{C}_i \to \mathcal{C}$ et
$g : \colim \mathcal{D}_i \to \mathcal{D}$
qui s'insèrent dans un diagramme commutatif
$$
\xymatrix{
\colim \mathcal{C}_i \ar[d]_p \ar[r]_f & \mathcal{C} \ar[d]^p \\
\colim \mathcal{D}_i \ar[r]^g & \mathcal{D}
}
$$
Nous affirmons alors que
$$
A : \colim \mathcal{C}_i \longrightarrow
\mathcal{C} \times_\mathcal{D} \colim \mathcal{D}_i
$$
est pleinement fidèle si et seulement si le foncteur
$$
B : \colim \mathcal{C}_i \longrightarrow
\mathcal{C}
\times_{\Delta, \mathcal{C} \times_\mathcal{D} \mathcal{C}, f \times_g f}
\colim (\mathcal{C}_i \times_{\mathcal{D}_i} \mathcal{C}_i)
$$
est une équivalence. Posons $\mathcal{C}' = \colim \mathcal{C}_i$ et
$\mathcal{D}' = \colim \mathcal{D}_i$. Puisque les $2$-produits fibrés
commutent aux colimites filtrantes, $A$ et $B$ deviennent les foncteurs
$$
A' : \mathcal{C}' \to \mathcal{C} \times_\mathcal{D} \mathcal{D}'
\quad\text{et}\quad
B' : \mathcal{C}' \longrightarrow
\mathcal{C}
\times_{\Delta, \mathcal{C} \times_\mathcal{D} \mathcal{C}, f \times_g f}
(\mathcal{C}' \times_{\mathcal{D}'} \mathcal{C}')
$$
Il suffit donc de prouver que, si
$$
\xymatrix{
\mathcal{C}' \ar[d]_p \ar[r]_f & \mathcal{C} \ar[d]^p \\
\mathcal{D}' \ar[r]^g & \mathcal{D}
}
$$
est un diagramme commutatif de groupoïdes, alors $A'$ est pleinement fidèle
si et seulement si $B'$ est une équivalence. Cela résulte de Catégories,
lemme \ref{categories-lemma-fully-faithful-diagonal-equivalence} (avec une
catégorie de base triviale, c'est-à-dire ponctuelle), car
$$
\mathcal{C}
\times_{\Delta, \mathcal{C} \times_\mathcal{D} \mathcal{C}, f \times_g f}
(\mathcal{C}' \times_{\mathcal{D}'} \mathcal{C}') =
\mathcal{C}'
\times_{A', \mathcal{C} \times_\mathcal{D} \mathcal{D}', A'}
\mathcal{C}'
$$
La démonstration est achevée.
\end{proof}

\begin{lemma}
\label{stacks-limits-lemma-locally-finite-presentation-limit-preserving}
Soit $S$ un schéma. Soit $\mathcal{X}$ un champ algébrique sur $S$. Si
$\mathcal{X} \to S$ est localement de présentation finie, alors
$\mathcal{X}$ commute aux limites au sens de Axiomes d'Artin, définition
\ref{artin-definition-limit-preserving} (de manière équivalente, le morphisme
$\mathcal{X} \to S$ commute aux limites).
\end{lemma}

\begin{proof}
Choisissons un morphisme lisse surjectif $U \to \mathcal{X}$, où $U$ est un
schéma. Alors $U \to S$ est localement de présentation finie ; voir Morphismes
de champs, section \ref{stacks-morphisms-section-finite-presentation}. Nous
pouvons écrire $\mathcal{X} = [U/R]$ pour un groupoïde lisse en espaces
algébriques $(U, R, s, t, c)$ ; voir Champs algébriques, lemme
\ref{algebraic-lemma-stack-presentation}. Puisque $U$ est localement de présentation
finie sur $S$, l'espace algébrique $R$ l'est lui aussi sur $S$. Rappelons que
$[U/R]$ est le champ en groupoïdes sur $(\Sch/S)_{fppf}$ obtenu en passant au
champ associé à la catégorie fibrée en groupoïdes dont la catégorie fibre sur
$T$ est le groupoïde $(U(T), R(T), s, t, c)$. Comme les foncteurs $U$ et $R$
commutent aux limites (Limites d'espaces, proposition
\ref{spaces-limits-proposition-characterize-locally-finite-presentation}),
cette catégorie fibrée en groupoïdes commute aux limites. Il suffit donc de
montrer que le passage au champ fppf associé conserve cette propriété. C'est
bien le cas (indication : utiliser Topologies, lemme
\ref{topologies-lemma-limit-fppf-topology}). Nous donnons néanmoins ci-dessous
une démonstration directe, en utilisant que nous connaissons ici explicitement
le passage au champ associé.

\medskip\noindent
Soit $T = \lim T_\lambda$ une limite projective filtrante de schémas affines
sur $S$. Nous devons montrer que le foncteur
$$
\colim [U/R]_{T_\lambda} \longrightarrow [U/R]_T
$$
est une équivalence de catégories. Montrons d'abord qu'il est essentiellement
surjectif. Soit $x \in \Ob([U/R]_T)$. Groupoïdes en espaces, lemme
\ref{spaces-groupoids-lemma-quotient-stack-objects}, décrit la catégorie
$[U/R]_T$. En particulier, $x$ correspond à un recouvrement fppf
$\{T_i \to T\}_{i \in I}$ et à une donnée de descente dans $[U/R]$
$(u_i, r_{ij})$ relative à ce recouvrement. Après avoir raffiné celui-ci, nous
pouvons supposer qu'il s'agit d'un recouvrement fppf standard du schéma affine
$T$. Par Topologies, lemme \ref{topologies-lemma-limit-fppf-topology}, nous
pouvons choisir un $\lambda$ et un recouvrement fppf standard
$\{T_{\lambda, i} \to T_\lambda\}_{i \in I}$ dont le changement de base à
$T$ soit égal à $\{T_i \to T\}_{i \in I}$. Pour tout $i$, quitte à augmenter
$\lambda$, il existe un morphisme $u_{\lambda, i} : T_{\lambda, i} \to U$
dont le composé avec $T_i \to T_{\lambda, i}$ est le morphisme donné $u_i$
(c'est ici que nous utilisons la commutation de $U$ aux limites). De même,
pour tous $i, j$, quitte à augmenter $\lambda$, il existe un morphisme
$r_{\lambda, ij} : T_{\lambda, i} \times_{T_\lambda} T_{\lambda, j} \to R$
dont le composé avec $T_{ij} \to T_{\lambda, ij}$ est le morphisme donné
$r_{ij}$ (c'est ici que nous utilisons la commutation de $R$ aux limites).
Quitte à augmenter encore $\lambda$, nous pouvons supposer que
$$
s \circ r_{\lambda, ij} = u_{\lambda, i} \circ \text{pr}_0
\quad\text{et}\quad
t \circ r_{\lambda, ij} = u_{\lambda, j} \circ \text{pr}_1,
$$
et
$$
c \circ (r_{\lambda, jk} \circ \text{pr}_{12}, r_{\lambda, ij}
\circ \text{pr}_{01}) = r_{\lambda, ik} \circ \text{pr}_{02}.
$$
Autrement dit, nous pouvons supposer que
$(u_{\lambda, i}, r_{\lambda, ij})$ est une donnée de descente dans $[U/R]$
relative au recouvrement
$\{T_{\lambda, i} \to T_\lambda\}_{i \in I}$. Nous obtenons alors un objet
correspondant de $[U/R]$ sur $T_\lambda$ dont l'image réciproque sur $T$ est
isomorphe à $x$, comme voulu. La pleine fidélité se démontre exactement de la
même façon, à l'aide de la description des morphismes dans les catégories
fibres de $[U/T]$ donnée dans Groupoïdes en espaces, lemme
\ref{spaces-groupoids-lemma-quotient-stack-objects}.
\end{proof}

\begin{proposition}
\label{stacks-limits-proposition-characterize-locally-finite-presentation}
\begin{reference}
C'est un cas particulier de \cite[Lemme 2.3.15]{Emerton-Gee}.
\end{reference}
Soit $f : \mathcal{X} \to \mathcal{Y}$ un morphisme de champs algébriques.
Les conditions suivantes sont équivalentes :
\begin{enumerate}
\item $f$ commute aux limites,
\item $f$ commute aux limites sur les objets, et
\item $f$ est localement de présentation finie.
\end{enumerate}
\end{proposition}

\begin{proof}
Supposons (3). Soit $T = \lim T_i$ une limite projective filtrante de schémas
affines. Considérons le foncteur
$$
\colim \mathcal{X}_{T_i} \longrightarrow
\mathcal{X}_T \times_{\mathcal{Y}_T} \colim \mathcal{Y}_{T_i}
$$
Soit $(x, y_i, \beta)$ un objet du membre de droite, c'est-à-dire que
$x \in \Ob(\mathcal{X}_T)$, $y_i \in \Ob(\mathcal{Y}_{T_i})$ et que
$\beta : f(x) \to y_i|_T$ est un morphisme de $\mathcal{Y}_T$. Nous pouvons
alors considérer $(x, y_i, \beta)$ comme un objet du champ algébrique
$\mathcal{X}_{y_i} = \mathcal{X} \times_{\mathcal{Y}, y_i} T_i$ au-dessus de
$T$. Comme
$\mathcal{X}_{y_i} \to T_i$ est localement de présentation finie (c'est un
changement de base de $f$), il commute aux limites par le lemme
\ref{stacks-limits-lemma-locally-finite-presentation-limit-preserving}. Ainsi,
$(x, y_i, \beta)$ provient d'un objet sur $T_{i'}$ pour un certain
$i' \geq i$. En explicitant les définitions, nous voyons que
$(x, y_i, \beta)$ appartient à l'image essentielle du foncteur affiché.
Celui-ci est donc essentiellement surjectif. Autrement dit, $f$ commute aux
limites sur les objets. Appliquons maintenant ce résultat à la diagonale
$\Delta$ de $f$. D'après Morphismes de champs, lemme
\ref{stacks-morphisms-lemma-diagonal-morphism-finite-type}, le morphisme
$\Delta$ est localement de présentation finie. L'argument précédent montre
donc que $\Delta$ commute aux limites sur les objets. Par le lemme
\ref{stacks-limits-lemma-representable-by-spaces-limit-preserving}, $\Delta$ commute aux limites.
Le lemme \ref{stacks-limits-lemma-limit-preserving-diagonal} entraîne alors que le foncteur
affiché ci-dessus est pleinement fidèle. C'est donc une équivalence, puisque
nous avons déjà prouvé sa surjectivité essentielle, et (1) est établi.

\medskip\noindent
L'implication (1) $\Rightarrow$ (2) est immédiate. Supposons (2). Choisissons
un schéma $V$ et un morphisme lisse surjectif $V \to \mathcal{Y}$. Par
Critères de représentabilité, lemme
\ref{criteria-lemma-base-change-limit-preserving}, le changement de base
$\mathcal{X} \times_\mathcal{Y} V \to V$ commute aux limites sur les objets.
Choisissons un schéma $U$ et un morphisme lisse surjectif
$U \to \mathcal{X} \times_\mathcal{Y} V$. Un morphisme lisse étant localement
de présentation finie, $U \to \mathcal{X} \times_\mathcal{Y} V$ commute aux
limites, d'après la première partie de la démonstration. Par Critères de
représentabilité, lemme \ref{criteria-lemma-composition-limit-preserving}, le
composé $U \to V$ commute aux limites sur les objets. Nous en concluons que
$U \to V$ est localement de présentation finie ; voir Critères de
représentabilité, lemme
\ref{criteria-lemma-representable-by-spaces-limit-preserving}. C'est exactement
la condition disant que $f$ est localement de présentation finie ; voir
Morphismes de champs, définition
\ref{stacks-morphisms-definition-locally-finite-presentation}.
\end{proof}




\section{Descente de propriétés}
\label{stacks-limits-section-descent}

\noindent
Cette section est l'analogue de Limites, section \ref{limits-section-descent}.

\begin{situation}
\label{stacks-limits-situation-descent}
Soit $Y = \lim_{i \in I} Y_i$ la limite projective d'un système filtrant
d'espaces algébriques dont les morphismes de transition sont affines. Nous
supposons que $X_i$ est quasi-compact et quasi-séparé pour tout $i \in I$.
Nous choisissons en outre un élément $0 \in I$.
\end{situation}

\begin{lemma}
\label{stacks-limits-lemma-eventually-separated}
Dans la situation \ref{stacks-limits-situation-descent}, supposons que
$\mathcal{X}_0 \to Y_0$ soit un morphisme d'un champ algébrique vers $Y_0$.
Supposons $\mathcal{X}_0$ quasi-compact et quasi-séparé. Si
$Y \times_{Y_0} \mathcal{X}_0 \to Y$ est séparé, alors
$Y_i \times_{Y_0} \mathcal{X}_0 \to Y_i$ est séparé pour tout $i \in I$
suffisamment grand.
\end{lemma}

\begin{proof}
Écrivons $\mathcal{X} = Y \times_{Y_0} \mathcal{X}_0$ et
$\mathcal{X}_i = Y_i \times_{Y_0} \mathcal{X}_0$. Choisissons un schéma
affine $U_0$ et un morphisme lisse surjectif
$U_0 \to \mathcal{X}_0$. Posons $U = Y \times_{Y_0} U_0$ et
$U_i = Y_i \times_{Y_0} U_0$. Alors $U$ et $U_i$ sont affines, tandis que
$U \to \mathcal{X}$ et $U_i \to \mathcal{X}_i$ sont lisses et surjectifs.
Posons $R_0 = U_0 \times_{\mathcal{X}_0} U_0$, puis
$R = Y \times_{Y_0} R_0$ et $R_i = Y_i \times_{Y_0} R_0$. Nous avons alors
$R = U \times_\mathcal{X} U$ et
$R_i = U_i \times_{\mathcal{X}_i} U_i$.

\medskip\noindent
Avec ces notations, remarquons que, si $\mathcal{X} \to Y$ est séparé, alors
$R \to U \times_Y U$ est propre : c'est le changement de base de
$\mathcal{X} \to \mathcal{X} \times_Y \mathcal{X}$ suivant
$U \times_Y U \to \mathcal{X} \times_Y \mathcal{X}$. Réciproquement,
$\mathcal{X}_i \to Y_i$ est séparé dès que
$R_i \to U_i \times_{Y_i} U_i$ est propre, car
$U_i \times_{Y_i} U_i \to \mathcal{X}_i \times_{Y_i} \mathcal{X}_i$ est
lisse et surjectif ; voir Propriétés des champs, lemme
\ref{stacks-properties-lemma-check-property-covering}. Remarquons que
$R_0 \to U_0 \times_{Y_0} U_0$ est localement de type fini et que $R_0$ est
quasi-compact et quasi-séparé. Par Limites d'espaces, lemme
\ref{spaces-limits-lemma-eventually-proper}, le morphisme
$R_i \to U_i \times_{Y_i} U_i$ est propre pour $i$ assez grand, ce qui achève
la démonstration.
\end{proof}







\section{Descente d'objets relatifs}
\label{stacks-limits-section-descending-relative}

\noindent
Cette section est l'analogue de Limites d'espaces, section
\ref{spaces-limits-section-descending-relative}.

\begin{lemma}
\label{stacks-limits-lemma-descend-a-stack-down}
Soit $I$ un ensemble filtrant. Soit $(X_i, f_{ii'})$ un système projectif
d'espaces algébriques indexé par $I$. Supposons que
\begin{enumerate}
\item les morphismes $f_{ii'} : X_i \to X_{i'}$ sont affines,
\item les espaces $X_i$ sont quasi-compacts et quasi-séparés.
\end{enumerate}
Soit $X = \lim X_i$. Si $\mathcal{X}$ est un champ algébrique de présentation
finie sur $X$, il existe un $i \in I$ et un champ algébrique $\mathcal{X}_i$
de présentation finie sur $X_i$ tels que
$\mathcal{X} \cong \mathcal{X}_i \times_{X_i} X$ comme champs algébriques sur
$X$.
\end{lemma}

\begin{proof}
Par Morphismes de champs, définition
\ref{stacks-morphisms-definition-locally-finite-presentation}, le morphisme
$\mathcal{X} \to X$ est quasi-compact, localement de présentation finie et
quasi-séparé. Comme $X$ et $\mathcal{X} \to X$ sont quasi-compacts, le champ
$\mathcal{X}$ est quasi-compact (Morphismes de champs, définition
\ref{stacks-morphisms-definition-quasi-compact}). Il existe donc un schéma
affine $U$ et un morphisme lisse surjectif $U \to \mathcal{X}$ (Propriétés
des champs, lemme \ref{stacks-properties-lemma-quasi-compact-stack}). Posons
$R = U \times_\mathcal{X} U$. Nous obtenons un groupoïde lisse en espaces
algébriques $(U, R, s, t, c)$ sur $X$ tel que $\mathcal{X} = [U/R]$ ; voir
Champs algébriques, lemme \ref{algebraic-lemma-stack-presentation}. Puisque
$\mathcal{X} \to X$ et $X$ sont quasi-séparés, $\mathcal{X}$ est
quasi-séparé (Morphismes de champs, lemme
\ref{stacks-morphisms-lemma-composition-separated}). Ainsi,
$R \to U \times U$ est quasi-compact et quasi-séparé (Morphismes de champs,
lemme \ref{stacks-morphisms-lemma-fibre-product-after-map}) ; par conséquent,
$R$ est un espace algébrique quasi-compact et quasi-séparé. D'autre part,
$U \to X$ est localement de présentation finie, et il en va donc de même de
$R \to X$ (car $s : R \to U$ est lisse, donc localement de présentation
finie). Le quintuplet $(U, R, s, t, c)$ est ainsi un objet groupoïde dans la
catégorie des espaces algébriques de présentation finie sur $X$. Par Limites
d'espaces, lemme \ref{spaces-limits-lemma-descend-finite-presentation}, il
existe un $i$ et un groupoïde en espaces algébriques
$(U_i, R_i, s_i, t_i, c_i)$ sur $X_i$ dont le changement de base à $X$ est
isomorphe à $(U, R, s, t, c)$. Quitte à augmenter $i$, nous pouvons supposer
$s_i$ et $t_i$ lisses ; voir Limites d'espaces, lemme
\ref{spaces-limits-lemma-descend-smooth}. Le champ quotient
$\mathcal{X}_i = [U_i/R_i]$ est algébrique (Champs algébriques, théorème
\ref{algebraic-theorem-smooth-groupoid-gives-algebraic-stack}).

\medskip\noindent
Il existe un morphisme $[U/R] \to [U_i/R_i]$ ; voir Groupoïdes en espaces,
lemme \ref{spaces-groupoids-lemma-quotient-stack-functorial}. Nous affirmons
que, joint aux morphismes $[U/R] \to X$ et $[U_i/R_i] \to X_i$ (Groupoïdes
en espaces, lemme \ref{spaces-groupoids-lemma-quotient-stack-arrows}), il
fournit un isomorphisme, c'est-à-dire une équivalence,
$$
[U/R] \longrightarrow [U_i/R_i] \times_{X_i} X
$$
Le morphisme correspondant
$$
[U/_{\!p}R] \longrightarrow [U_i/_{\!p}R_i] \times_{X_i} X
$$
au niveau des « préfaisceaux de groupoïdes » de Groupoïdes en espaces,
équation (\ref{spaces-groupoids-equation-quotient-stack}), est un
isomorphisme. L'affirmation résulte donc du fait que le passage au champ
associé commute aux produits fibrés ; voir Champs, lemme
\ref{stacks-lemma-stackification-fibre-product-fibred-categories}.
\end{proof}




\section[Immersion fermée et présentation finie]{Immersion fermée d'un morphisme de type fini dans un morphisme de présentation finie}
\label{stacks-limits-section-finite-type-closed-in-finite-presentation}

\noindent
Cette section est l'analogue de Limites d'espaces, section
\ref{spaces-limits-section-finite-type-closed-in-finite-presentation}.

\begin{lemma}
\label{stacks-limits-lemma-finite-type-closed-in-finite-presentation}
Soit $f : \mathcal{X} \to Y$ un morphisme d'un champ algébrique vers un espace
algébrique. Supposons que :
\begin{enumerate}
\item $f$ est de type fini et quasi-séparé ;
\item $Y$ est quasi-compact et quasi-séparé.
\end{enumerate}
Il existe alors un morphisme de présentation finie
$f' : \mathcal{X}' \to Y$ et une immersion fermée
$\mathcal{X} \to \mathcal{X}'$ de champs algébriques sur $Y$.
\end{lemma}

\begin{proof}
Écrivons $Y = \lim_{i \in I} Y_i$ comme limite projective d'espaces algébriques
indexés par un ensemble filtrant $I$, dont les morphismes de transition sont
affines et les $Y_i$ noethériens ; voir Limites d'espaces, proposition
\ref{spaces-limits-proposition-approximate}. Nous utiliserons les résultats de
Limites d'espaces, section
\ref{spaces-limits-section-finite-type-quasi-separated}.

\medskip\noindent
Choisissons une présentation $\mathcal{X} = [U/R]$. Notons
$(U, R, s, t, c, e, i)$ le groupoïde correspondant en espaces algébriques sur
$Y$. Nous pouvons et allons supposer $U$ affine. Alors $U$, $R$ et
$R \times_{s, U, t} R$ sont des espaces algébriques quasi-séparés et de type
fini sur $Y$. Nous avons deux morphismes $s, t : R \to U$, trois morphismes
$c : R \times_{s, U, t} R \to R$,
$\text{pr}_1 : R \times_{s, U, t} R \to R$,
$\text{pr}_2 : R \times_{s, U, t} R \to R$,
un morphisme $e : U \to R$ et, enfin, un morphisme $i : R \to R$. Ces
morphismes satisfont une liste d'axiomes explicitée dans Groupoïdes, section
\ref{groupoids-section-groupoids}.

\medskip\noindent
D'après Limites d'espaces, remarque
\ref{spaces-limits-remark-finite-type-gives-well-defined-system}
nous pouvons trouver un $i_0 \in I$ et des systèmes projectifs
\begin{enumerate}
\item $(U_i)_{i \geq i_0}$,
\item $(R_i)_{i \geq i_0}$,
\item $(T_i)_{i \geq i_0}$
\end{enumerate}
sur $(Y_i)_{i \geq i_0}$ tels que
$U = \lim_{i \geq i_0} U_i$,
$R = \lim_{i \geq i_0} R_i$ et
$R \times_{s, U, t} R = \lim_{i \geq i_0} T_i$
et tels qu'il existe des morphismes de systèmes projectifs
\begin{enumerate}
\item $(s_i)_{i \geq i_0} : (R_i)_{i \geq i_0} \to (U_i)_{i \geq i_0}$,
\item $(t_i)_{i \geq i_0} : (R_i)_{i \geq i_0} \to (U_i)_{i \geq i_0}$,
\item $(c_i)_{i \geq i_0} : (T_i)_{i \geq i_0} \to (R_i)_{i \geq i_0}$,
\item $(p_i)_{i \geq i_0} : (T_i)_{i \geq i_0} \to (R_i)_{i \geq i_0}$,
\item $(q_i)_{i \geq i_0} : (T_i)_{i \geq i_0} \to (R_i)_{i \geq i_0}$,
\item $(e_i)_{i \geq i_0} : (U_i)_{i \geq i_0} \to (R_i)_{i \geq i_0}$,
\item $(i_i)_{i \geq i_0} : (R_i)_{i \geq i_0} \to (R_i)_{i \geq i_0}$
\end{enumerate}
vérifiant
$s = \lim_{i \geq i_0} s_i$,
$t = \lim_{i \geq i_0} t_i$,
$c = \lim_{i \geq i_0} c_i$,
$\text{pr}_1 = \lim_{i \geq i_0} p_i$,
$\text{pr}_2 = \lim_{i \geq i_0} q_i$,
$e = \lim_{i \geq i_0} e_i$ et
$i = \lim_{i \geq i_0} i_i$.
Par Limites d'espaces, lemme
\ref{spaces-limits-lemma-morphism-good-diagram-smooth}
nous pouvons supposer $s_i$ et $t_i$ lisses, quitte à augmenter $i_0$.
Par Limites d'espaces, lemme
\ref{spaces-limits-lemma-morphism-good-diagram-flat}
nous pouvons supposer que les applications
$R \to U \times_{U_i, s_i} R_i$ données par $s$ et $R \to R_i$, ainsi que
$R \to U \times_{U_i, t_i} R_i$ données par $t$ et $R \to R_i$, sont des
isomorphismes pour tout $i \geq i_0$. Par Limites d'espaces, lemme
\ref{spaces-limits-lemma-good-diagram-fibre-product}, nous pouvons supposer les
diagrammes
$$
\xymatrix{
T_i \ar[r]_{q_i} \ar[d]_{p_i} & R_i \ar[d]^{t_i} \\
R_i \ar[r]^{s_i} & U_i
}
$$
cartésiens. L'assertion d'unicité de Limites d'espaces, lemme
\ref{spaces-limits-lemma-morphism-good-diagram}, garantit alors que, pour un
$i$ assez grand, les relations mentionnées plus haut entre les morphismes
$s, t, c, e, i$ sont satisfaites par $s_i, t_i, c_i, e_i, i_i$. Fixons un tel
$i$.

\medskip\noindent
Il s'ensuit que $(U_i, R_i, s_i, t_i, c_i, e_i, i_i)$ est un groupoïde lisse en
espaces algébriques sur $Y_i$. Ainsi, $\mathcal{X}_i = [U_i/R_i]$ est un champ
algébrique (Champs algébriques, théorème
\ref{algebraic-theorem-smooth-groupoid-gives-algebraic-stack}).
Le morphisme de groupoïdes
$$
(U, R, s, t, c, e, i) \to (U_i, R_i, s_i, t_i, c_i, e_i, i_i)
$$
au-dessus de $Y \to Y_i$ détermine un diagramme commutatif
$$
\xymatrix{
\mathcal{X} \ar[d] \ar[r] & \mathcal{X}_i \ar[d] \\
Y \ar[r] & Y_i
}
$$
(Groupoïdes en espaces, lemme
\ref{spaces-groupoids-lemma-quotient-stack-functorial}).
Nous affirmons que le morphisme $\mathcal{X} \to Y \times_{Y_i} \mathcal{X}_i$
est une immersion fermée. Cette affirmation achève la démonstration, car le
champ algébrique $\mathcal{X}_i \to Y_i$ est de présentation finie par
construction. Pour la prouver, remarquons que le diagramme de gauche
$$
\xymatrix{
U \ar[d] \ar[r] & U_i \ar[d] \\
\mathcal{X} \ar[r] & \mathcal{X}_i
}
\quad\quad
\xymatrix{
U \ar[d] \ar[r] & Y \times_{Y_i} U_i \ar[d] \\
\mathcal{X} \ar[r] & Y \times_{Y_i} \mathcal{X}_i
}
$$
est cartésien par Groupoïdes en espaces, lemme
\ref{spaces-groupoids-lemma-criterion-fibre-product}
et les résultats rappelés ci-dessus. Le diagramme commutatif de droite est donc
lui aussi cartésien. Le résultat voulu découle alors du fait que
$U \to Y \times_{Y_i} U_i$ est une immersion fermée, par construction du
système projectif $(U_i)$ dans Limites d'espaces, lemme
\ref{spaces-limits-lemma-limit-from-good-diagram}, du fait que
$Y \times_{Y_i} U_i \to Y \times_{Y_i} \mathcal{X}_i$ est lisse et surjectif,
et de Propriétés des champs, lemme
\ref{stacks-properties-lemma-check-immersion-covering}.
\end{proof}

\noindent
Il existe une variante pour les champs algébriques séparés.

\begin{lemma}
\label{stacks-limits-lemma-separated-closed-in-finite-presentation}
Soit $f : \mathcal{X} \to Y$ un morphisme d'un champ algébrique vers un espace
algébrique. Supposons que :
\begin{enumerate}
\item $f$ est de type fini et séparé ;
\item $Y$ est quasi-compact et quasi-séparé.
\end{enumerate}
Il existe alors un morphisme séparé de présentation finie
$f' : \mathcal{X}' \to Y$ et une immersion fermée
$\mathcal{X} \to \mathcal{X}'$ de champs algébriques sur $Y$.
\end{lemma}

\begin{proof}
Nous appliquons d'abord exactement le même procédé que dans la démonstration du
lemme \ref{stacks-limits-lemma-finite-type-closed-in-finite-presentation}, dont nous reprenons
les notations, pour construire l'immersion $\mathcal{X} \to \mathcal{X}'$ sous
la forme d'un morphisme
$\mathcal{X} \to \mathcal{X}' = Y \times_{Y_i} \mathcal{X}_i$, où
$\mathcal{X}_i = [U_i/R_i]$. Il suffit donc de montrer que
$\mathcal{X}_i \to Y_i$ est séparé pour $i$ assez grand. Autrement dit, il
suffit de montrer que
$\mathcal{X}_i \to \mathcal{X}_i \times_{Y_i} \mathcal{X}_i$ est propre pour
$i$ assez grand. Puisque le morphisme
$U_i \times_{Y_i} U_i \to \mathcal{X}_i \times_{Y_i} \mathcal{X}_i$ est
surjectif et lisse, et puisque
$R_i = \mathcal{X}_i
\times_{\mathcal{X}_i \times_{Y_i} \mathcal{X}_i} U_i \times_{Y_i} U_i$
il suffit de montrer que le morphisme
$(s_i, t_i) : R_i \to U_i \times_{Y_i} U_i$
est propre pour $i$ assez grand ; voir Propriétés des champs, lemme
\ref{stacks-properties-lemma-check-property-covering}.
Nous le démontrons au paragraphe suivant.

\medskip\noindent
Observons que $U \times_Y U \to Y$ est quasi-séparé et de type fini. Nous
pouvons donc appliquer la construction de Limites d'espaces, remarque
\ref{spaces-limits-remark-finite-type-gives-well-defined-system}
pour trouver un $i_1 \in I$ et un système projectif $(V_i)_{i \geq i_1}$ tels
que $U \times_Y U = \lim_{i \geq i_1} V_i$. Par Limites d'espaces, lemme
\ref{spaces-limits-lemma-good-diagram-fibre-product}, pour $i$ assez grand, la
fonctorialité de la construction appliquée aux projections
$U \times_Y U \to U$ donne des immersions fermées
$$
V_i \to U_i \times_{Y_i} U_i
$$
(Il existe ici un léger décalage : il faudrait en réalité remplacer $Y_i$ par
l'image schématique de $Y \to Y_i$, mais cela ne change manifestement pas le
produit fibré.) D'autre part, par Limites d'espaces, lemme
\ref{spaces-limits-lemma-morphism-good-diagram-proper}
la fonctorialité appliquée au morphisme propre
$(s, t) : R \to U \times_Y U$ (nous utilisons ici que $\mathcal{X}$ est séparé)
fournit des morphismes $R_i \to V_i$ propres pour $i$ assez grand. En les
composant, nous obtenons un morphisme propre
$R_i \to U_i \times_{Y_i} U_i$ pour tout $i$ assez grand. La fonctorialité de
la construction de Limites d'espaces, remarque
\ref{spaces-limits-remark-finite-type-gives-well-defined-system}
montre que ce morphisme coïncide avec $(s_i, t_i)$ pour $i$ assez grand, ce qui
achève la démonstration.
\end{proof}






\section{Morphismes universellement fermés}
\label{stacks-limits-section-universally-closed}

\noindent
Cette section est l'analogue de Limites d'espaces, section
\ref{spaces-limits-section-universally-closed}.

\begin{lemma}
\label{stacks-limits-lemma-separate}
Soit $g : Z \to Y$ un morphisme de schémas affines. Soit
$f : \mathcal{X} \to Y$ un morphisme quasi-compact de champs algébriques. Soit
$z \in Z$, et soit $T \subset |\mathcal{X} \times_Y Z|$ une partie fermée telle
que $z \not \in \Im(T \to |Z|)$. Si $\mathcal{X}$ est quasi-compact, il existe
un voisinage ouvert $V \subset Z$ de $z$, un diagramme commutatif
$$
\xymatrix{
V \ar[d] \ar[r]_a & Z' \ar[d]^b \\
Z \ar[r]^g & Y,
}
$$
et une partie fermée $T' \subset |X \times_Y Z'|$ tels que :
\begin{enumerate}
\item $Z'$ est un schéma affine de présentation finie sur $Y$ ;
\item en posant $z' = a(z)$, on a $z' \not \in \Im(T' \to |Z'|)$ ;
\item l'image réciproque de $T$ dans $|\mathcal{X} \times_Y V|$ s'envoie dans
$T'$ par $|\mathcal{X} \times_Y V| \to |\mathcal{X} \times_Y Z'|$.
\end{enumerate}
\end{lemma}

\begin{proof}
Nous allons déduire l'assertion du résultat correspondant pour les morphismes
de schémas. Comme $\mathcal{X}$ est quasi-compact, choisissons un schéma affine
$W$ et un morphisme lisse surjectif $W \to \mathcal{X}$. Soit
$T_W \subset |W \times_Y Z|$ l'image réciproque de $T$. Alors $z$ n'appartient
pas à l'image de $T_W$. Le cas des schémas (Limites, lemme
\ref{limits-lemma-separate}) fournit un voisinage ouvert $V \subset Z$ de $z$,
un diagramme commutatif de schémas
$$
\xymatrix{
V \ar[d] \ar[r]_a & Z' \ar[d]^b \\
Z \ar[r]^g & Y,
}
$$
et une partie fermée $T' \subset |W \times_Y Z'|$ tels que :
\begin{enumerate}
\item $Z'$ est un schéma affine de présentation finie sur $Y$ ;
\item en posant $z' = a(z)$, on a $z' \not \in \Im(T' \to |Z'|)$ ;
\item $T_1 = T_W \cap |W \times_Y V|$ s'envoie dans $T'$ par
$|W \times_Y V| \to |W \times_Y Z'|$.
\end{enumerate}
Le diagramme commutatif
$$
\xymatrix{
W \times_Y Z \ar[d] &
W \times_Y V \ar[l] \ar[rr]_{a_1} \ar[d]_c & &
W \times_Y Z' \ar[d]^q \\
\mathcal{X} \times_Y Z &
\mathcal{X} \times_Y V \ar[l] \ar[rr]^{a_2} & &
\mathcal{X} \times_Y Z'
}
$$
est formé de carrés cartésiens, et ses flèches verticales sont surjectives,
lisses et, a fortiori, ouvertes. Le carré de gauche montre que
$T_1 = T_W \cap |W \times_Y V|$ est l'image réciproque de
$T_2 = T \cap |\mathcal{X} \times_Y V|$ par $c$. Par Propriétés des champs, lemme
\ref{stacks-properties-lemma-points-cartesian}, on obtient
$a_1(T_1) = q^{-1}(a_2(T_2))$.
Par Topologie, lemme \ref{topology-lemma-open-morphism-quotient-topology}, on a
$$
q^{-1}\left(\overline{a_2(T_2)}\right) =
\overline{q^{-1}(a_2(T_2))} =
\overline{a_1(T_1)} \subset T'
$$
Comme $q$ est surjectif, l'image de $\overline{a_2(T_2)} \to |Z'|$ ne contient
pas $z'$, puisque cela vaut pour $T'$. Le diagramme ci-dessus, avec
$Z', V, a, b$, et la partie fermée
$\overline{a_2(T_2)} \subset |\mathcal{X} \times_Y Z'|$ répondent donc au
problème posé dans le lemme.
\end{proof}

\begin{lemma}
\label{stacks-limits-lemma-test-universally-closed}
Soit $f : \mathcal{X} \to \mathcal{Y}$ un morphisme quasi-compact de champs
algébriques. Les conditions suivantes sont équivalentes :
\begin{enumerate}
\item $f$ est universellement fermé ;
\item pour tout morphisme $Z \to \mathcal{Y}$ localement de présentation finie,
où $Z$ est un schéma affine, l'application
$|\mathcal{X} \times_Y Z| \to |Z|$ est fermée ;
\item il existe un schéma $V$ et un morphisme lisse surjectif
$V \to \mathcal{Y}$ tels que l'application
$|\mathbf{A}^n \times (\mathcal{X} \times_\mathcal{Y} V)|
\to |\mathbf{A}^n \times V|$ soit fermée pour tout $n \geq 0$.
\end{enumerate}
\end{lemma}

\begin{proof}
Il est clair que (1) implique (2).

\medskip\noindent
Supposons (2). Choisissons un schéma $V$ qui soit réunion disjointe de schémas
affines, ainsi qu'un morphisme lisse surjectif $V \to \mathcal{Y}$. Pour montrer
que $f$ est universellement fermé, il suffit de montrer que le changement de
base $\mathcal{X} \times_\mathcal{Y} V \to V$ de $f$ est universellement fermé ;
voir Morphismes de champs, lemme
\ref{stacks-morphisms-lemma-universally-closed-local}.
Remarquons que la condition (2) vaut pour ce changement de base. Afin de prouver
que (2) implique (1), nous pouvons donc supposer que $Y = \mathcal{Y}$ est un
schéma affine.

\medskip\noindent
Supposons (2), et supposons que $\mathcal{Y} = Y$ est un schéma affine. Si $f$
n'est pas universellement fermé, il existe un schéma affine $Z$ sur $Y$ tel que
$|\mathcal{X} \times_Y Z| \to |Z|$ ne soit pas fermé ; voir Morphismes de
champs, lemme
\ref{stacks-morphisms-lemma-universally-closed-local}.
Il existe donc une partie fermée $T \subset |\mathcal{X} \times_Y Z|$ telle que
$\Im(T \to |Z|)$ ne soit pas fermée. Choisissons $z \in |Z|$ dans l'adhérence
de l'image de $T$, mais non dans cette image. Appliquons le lemme
\ref{stacks-limits-lemma-separate}. Nous obtenons un voisinage ouvert $V \subset Z$, un
diagramme commutatif
$$
\xymatrix{
V \ar[d] \ar[r]_a & Z' \ar[d]^b \\
Z \ar[r]^g & Y,
}
$$
et une partie fermée $T' \subset |\mathcal{X} \times_Y Z'|$ tels que :
\begin{enumerate}
\item $Z'$ est un schéma affine de présentation finie sur $Y$ ;
\item en posant $z' = a(z)$, on a $z' \not \in \Im(T' \to |Z'|)$ ;
\item l'image réciproque de $T$ dans $|\mathcal{X} \times_Y V|$ s'envoie dans
$T'$ par
$|\mathcal{X} \times_Y V| \to |\mathcal{X} \times_Y Z'|$.
\end{enumerate}
Nous affirmons que $z'$ appartient à l'adhérence de $\Im(T' \to |Z'|)$. Il en
résulte que $|\mathcal{X} \times_Y Z'| \to |Z'|$ n'est pas fermée, contrairement
à (2). Autrement dit, cette affirmation montre que (2) implique (1). Pour la
vérifier, considérons le diagramme commutatif suivant :
$$
\xymatrix{
\mathcal{X} \times_Y Z \ar[d] &
\mathcal{X} \times_Y V \ar[l] \ar[d] \ar[r] &
\mathcal{X} \times_Y Z' \ar[d] \\
Z & V \ar[l] \ar[r]^a & Z'
}
$$
Soit $T_V \subset |\mathcal{X} \times_Y V|$ l'image réciproque de $T$. Par
Propriétés des champs, lemme \ref{stacks-properties-lemma-points-cartesian},
l'image de $T_V$ dans $|V|$ est l'image réciproque de l'image de $T$ dans
$|Z|$. Puisque $z$ appartient à l'adhérence de l'image de $T \to |Z|$ et que
$|V| \to |Z|$ est ouverte, $z$ appartient donc à l'adhérence de l'image de
$T_V \to |V|$. Comme l'image de $T_V$ dans
$|\mathcal{X} \times_Y Z'|$ est contenue dans $|T'|$, il s'ensuit immédiatement
que $z' = a(z)$ appartient à l'adhérence de l'image de $T'$.

\medskip\noindent
Il est clair que (1) implique (3). Soit $V \to \mathcal{Y}$ comme dans (3). Si
nous montrons que $\mathcal{X} \times_Y V \to V$ est universellement fermé,
alors $f$ est universellement fermé par Morphismes de champs, lemme
\ref{stacks-morphisms-lemma-universally-closed-local}.
Il suffit donc de montrer que $f : \mathcal{X} \to \mathcal{Y}$ satisfait (2)
lorsque $f$ est un morphisme quasi-compact de champs algébriques, que
$\mathcal{Y} = Y$ est un schéma et que
$|\mathbf{A}^n \times \mathcal{X}| \to |\mathbf{A}^n \times Y|$
soit fermée pour tout $n$. Soit $Z \to Y$ localement de présentation finie,
où $Z$ est un schéma affine. Nous devons montrer que l'application
$|\mathcal{X} \times_Y Z| \to |Z|$ est fermée. Comme $Y$ est un schéma, que
$Z$ est affine et que $Z \to Y$ est localement de présentation finie, il
existe une immersion $Z \to \mathbf{A}^n \times Y$ ; voir Morphismes, lemme
\ref{morphisms-lemma-quasi-affine-finite-type-over-S}. Considérons le diagramme
cartésien
$$
\vcenter{
\xymatrix{
\mathcal{X} \times_Y Z \ar[d] \ar[r] &
\mathbf{A}^n \times \mathcal{X} \ar[d] \\
Z \ar[r] & \mathbf{A}^n \times Y
}
}
\quad
\begin{matrix}
\text{qui induit le} \\
\text{carré cartésien}
\end{matrix}
\quad
\vcenter{
\xymatrix{
|\mathcal{X} \times_Y Z| \ar[d] \ar[r] &
|\mathbf{A}^n \times \mathcal{X}| \ar[d] \\
|Z| \ar[r] & |\mathbf{A}^n \times Y|
}
}
$$
des espaces topologiques, dont les flèches horizontales sont des
homéomorphismes sur des parties localement fermées (Propriétés des champs, lemme
\ref{stacks-properties-lemma-subspace-induced-topology}).
Toute partie fermée $T$ de $|X \times_Y Z|$ est donc l'image réciproque d'une
partie fermée $T'$ de $|\mathbf{A}^n \times Y|$. Comme, par hypothèse, l'image
de $T'$ dans $|\mathbf{A}^n \times X|$ est fermée, nous en concluons que l'image
de $T$ dans $|Z|$ est fermée, comme voulu.
\end{proof}












% OK, start here.
%

\chapter{Cohomologie des champs algébriques}


\phantomsection
\label{stacks-cohomology-section-phantom}





\section{Introduction}
\label{stacks-cohomology-section-introduction}

\noindent
Dans ce chapitre, nous traitons de la cohomologie des champs algébriques.
Il s'agit notamment de la cohomologie des faisceaux quasi-cohérents, c'est-à-dire que
nous démontrons des analogues des résultats des chapitres intitulés
``Cohomologie des schémas'' et ``Cohomologie des espaces algébriques''.
Les résultats de ce chapitre diffèrent de ceux
de \cite{LM-B} principalement parce que nous utilisons systématiquement les
« grands sites ». Avant de lire ce chapitre, nous invitons le lecteur à parcourir
le chapitre ``Faisceaux sur les champs algébriques'' afin de se
familiariser avec la terminologie qui y est introduite ; voir
Faisceaux sur les champs, section \ref{stacks-sheaves-section-introduction}.



\section{Conventions et abus de langage}
\label{stacks-cohomology-section-conventions}

\noindent
Nous continuons d'employer les conventions et abus de langage
introduits dans
Propriétés des champs, section \ref{stacks-properties-section-conventions}.











\section{Notations}
\label{stacks-cohomology-section-notation}

\noindent
Différentes topologies. Si nous désignons un champ algébrique par une lettre
calligraphique, telle que $\mathcal{X}, \mathcal{Y}, \mathcal{Z}$, alors les notations
$\mathcal{X}_{Zar}, \mathcal{X}_\etale, \mathcal{X}_{smooth},
\mathcal{X}_{syntomic}, \mathcal{X}_{fppf}$ désignent les sites introduits
dans
Faisceaux sur les champs, définition
\ref{stacks-sheaves-definition-inherited-topologies}.
(Il faut penser à un « grand site ».) Par conséquent, le faisceau structural de
$\mathcal{X}$ est un faisceau sur $\mathcal{X}_{fppf}$.
En revanche, les espaces algébriques et les schémas
sont habituellement désignés par des majuscules romaines, telles que $X, Y, Z$, et, dans ce cas,
$X_\etale$ désigne le petit site étale de $X$ (tel qu'il est
défini dans
Topologies, définition
\ref{topologies-definition-big-small-etale}
ou
Propriétés des espaces, définition
\ref{spaces-properties-definition-etale-site}).
Cette distinction paraît suffisamment claire.

\medskip\noindent
La topologie sous-entendue est la topologie fppf. Ainsi, nous dirons parfois
« faisceau sur $\mathcal{X}$ » ou « faisceau de $\mathcal{O}_\mathcal{X}$-modules »
pour désigner un faisceau sur $\mathcal{X}_{fppf}$ ou un objet de
$\textit{Mod}(\mathcal{X}_{fppf}, \mathcal{O}_\mathcal{X})$.

\medskip\noindent
Si $f : \mathcal{X} \to \mathcal{Y}$ est un morphisme de champs
algébriques, alors les foncteurs $f_*$ et $f^{-1}$ définis sur les préfaisceaux
préservent les faisceaux pour chacune des topologies mentionnées ci-dessus. En particulier,
lorsque nous parlons de l'image directe ou de l'image inverse d'un faisceau, il n'est pas nécessaire
de préciser la topologie considérée. Il n'en va pas de même
lorsque nous calculons les groupes de cohomologie ou les images directes supérieures. Dans ce
cas, nous préciserons toujours la topologie considérée.

\medskip\noindent
Supposons que $f : X \to \mathcal{Y}$ soit un morphisme d'un espace
algébrique $X$ vers un champ algébrique $\mathcal{Y}$. Soit $\mathcal{G}$
un faisceau sur $\mathcal{Y}_\tau$ pour une topologie $\tau$. Alors
$f^{-1}\mathcal{G}$ est un faisceau pour la topologie $\tau$ sur $\mathcal{S}_X$
(le champ algébrique associé à $X$), car, suivant nos conventions, $f$
est en fait un $1$-morphisme $f : \mathcal{S}_X \to \mathcal{Y}$.
Si $\tau = \etale$ ou si la topologie est plus fine, nous écrivons
$f^{-1}\mathcal{G}|_{X_\etale}$
pour désigner la restriction au site étale de $X$ ; voir
Faisceaux sur les champs, section \ref{stacks-sheaves-section-compare}.
Si $\mathcal{G}$ est un $\mathcal{O}_\mathcal{X}$-module, nous écrivons parfois
plutôt
$f^*\mathcal{G}$ et $f^*\mathcal{G}|_{X_\etale}$.




\section{Image inverse des modules quasi-cohérents}
\label{stacks-cohomology-section-pullback}

\noindent
Soit $f : \mathcal{X} \to \mathcal{Y}$ un morphisme de champs algébriques.
Il est très général que les modules quasi-cohérents sur les topos annelés
soient compatibles aux images inverses. En particulier, l'image inverse $f^*$ préserve
les modules quasi-cohérents et nous obtenons un foncteur
$$
f^* :
\QCoh(\mathcal{O}_\mathcal{Y})
\longrightarrow
\QCoh(\mathcal{O}_\mathcal{X}),
$$
voir Faisceaux sur les champs, lemme
\ref{stacks-sheaves-lemma-pullback-quasi-coherent}.
En général, ce foncteur n'est pas exact, mais il l'est si $f$ est plat.

\begin{lemma}
\label{stacks-cohomology-lemma-flat-pullback-quasi-coherent}
Si $f : \mathcal{X} \to \mathcal{Y}$ est un morphisme plat de champs algébriques,
alors $f^* : \QCoh(\mathcal{O}_\mathcal{Y}) \to
\QCoh(\mathcal{O}_\mathcal{X})$ est un foncteur exact.
\end{lemma}

\begin{proof}
Choisissons un schéma $V$ et un morphisme lisse surjectif $V \to \mathcal{Y}$.
Choisissons un schéma $U$ et un morphisme lisse surjectif
$U \to V \times_\mathcal{Y} \mathcal{X}$. Alors $U \to \mathcal{X}$ est
encore lisse et surjectif, comme composé de deux tels morphismes.
Du diagramme commutatif
$$
\xymatrix{
U \ar[d] \ar[r]_{f'} & V \ar[d] \\
\mathcal{X} \ar[r]^f & \mathcal{Y}
}
$$
nous déduisons un diagramme commutatif
$$
\xymatrix{
\QCoh(\mathcal{O}_U) & \QCoh(\mathcal{O}_V) \ar[l] \\
\QCoh(\mathcal{O}_\mathcal{X}) \ar[u] &
\QCoh(\mathcal{O}_\mathcal{Y}) \ar[l] \ar[u]
}
$$
de catégories abéliennes. La démonstration du fait que les deux catégories inférieures de ce
diagramme sont abéliennes a montré que les foncteurs verticaux sont des foncteurs
exacts et fidèles (voir la démonstration de
Faisceaux sur les champs, lemme
\ref{stacks-sheaves-lemma-quasi-coherent-algebraic-stack}).
Puisque $f'$ est un morphisme plat de schémas (d'après notre définition des
morphismes plats de champs algébriques), nous voyons que $(f')^*$ est un
foncteur exact sur les faisceaux quasi-cohérents de $V$. Cela conclut.
\end{proof}

\begin{lemma}
\label{stacks-cohomology-lemma-quasi-coherent-check-exact}
Soit $\mathcal{X}$ un champ algébrique. Soit $I$ un ensemble et,
pour $i \in I$, soit $x_i : U_i \to \mathcal{X}$ un objet
de $\mathcal{X}$. Supposons que $x_i$ soit plat et que
$\coprod x_i : \coprod U_i \to \mathcal{X}$ soit surjectif.
Soit $\varphi : \mathcal{F} \to \mathcal{G}$ une flèche de
$\QCoh(\mathcal{O}_\mathcal{X})$. Notons $\varphi_i$
la restriction de $\varphi$ à $(U_i)_\etale$.
Alors $\varphi$ est injective, resp.\ surjective, resp.\ un isomorphisme
si et seulement si chaque $\varphi_i$ l'est.
\end{lemma}

\begin{proof}
Choisissons un schéma $U$ et un morphisme lisse surjectif
$x : U \to \mathcal{X}$. Nous pouvons considérer, et considérons, $x$
comme un objet de $\mathcal{X}$.
On obtient ainsi une présentation $\mathcal{X} = [U/R]$ par un groupoïde
en espaces $(U, R, s, t, c)$ et, par conséquent, une équivalence
$$
\QCoh(\mathcal{O}_\mathcal{X}) = \QCoh(U, R, s, t, c)
$$
Voir la discussion dans Faisceaux sur les champs, section
\ref{stacks-sheaves-section-quasi-coherent-algebraic-stacks}.
La structure de catégorie abélienne du membre de droite est telle que $\varphi$
est injective, resp.\ surjective, resp.\ un isomorphisme
si et seulement si la restriction $\varphi|_{U_\etale}$ l'est ; voir
Groupoïdes dans les espaces, lemme \ref{spaces-groupoids-lemma-abelian}.

\medskip\noindent
Pour chaque $i$, choisissons un recouvrement étale
$\{W_{i, j} \to V \times_\mathcal{X} U_i\}_{j \in J_i}$
par des schémas. Notons $g_{i, j} : W_{i, j} \to V$ et
$h_{i, j} : W_{i, j} \to U_i$ les flèches évidentes.
Chacun des morphismes de schémas $g_{i, j} : W_{i, j} \to U$ est plat,
et ils sont conjointement surjectifs.
De même, pour chaque $i$ fixé, les morphismes de schémas
$h_{i, j} : W_{i, j} \to U_i$ sont plats et conjointement surjectifs.
D'après Faisceaux sur les champs, lemme \ref{stacks-sheaves-lemma-quasi-coherent},
l'image inverse par $(g_{i, j})_{small}$ de la restriction
$\varphi|_{U_\etale}$ est la restriction $\varphi|_{(W_{i, j})_\etale}$,
et l'image inverse par $(h_{i, j})_{small}$ de la restriction
$\varphi|_{(U_i)_\etale}$ est la restriction $\varphi|_{(W_{i, j})_\etale}$.
L'image inverse des modules quasi-cohérents par un morphisme plat de schémas est exacte,
et l'image inverse par une famille conjointement surjective de morphismes plats de schémas
reflète les applications injectives, resp.\ surjectives, resp.\ bijectives
de modules quasi-cohérents (en fait, cela vaut pour tous les modules,
car on peut vérifier l'exactitude sur les germes). Ainsi,
$$
\varphi|_{U_\etale} \text{ injective}
\Leftrightarrow
\varphi|_{(W_{i, j})_\etale} \text{ injective pour tous }i, j
\Leftrightarrow
\varphi|_{(U_i)_\etale} \text{ injective pour tout }i
$$
Ceci achève la démonstration.
\end{proof}




\section{Images directes supérieures de certaines classes de modules}
\label{stacks-cohomology-section-key}

\noindent
Le lemme suivant fonde notre étude des
images directes supérieures de certaines classes de faisceaux de modules.
Il en existe deux versions : l'une pour la topologie étale et
l'autre pour la topologie fppf.

\begin{lemma}
\label{stacks-cohomology-lemma-general-pushforward}
Soit $\mathcal{M}$ une règle qui associe à tout champ algébrique
$\mathcal{X}$ une sous-catégorie $\mathcal{M}_\mathcal{X}$ de
$\textit{Mod}(\mathcal{X}_\etale, \mathcal{O}_\mathcal{X})$
telle que
\begin{enumerate}
\item $\mathcal{M}_\mathcal{X}$ soit une sous-catégorie faible de Serre
de $\textit{Mod}(\mathcal{X}_\etale, \mathcal{O}_\mathcal{X})$
(voir Homologie, définition \ref{homology-definition-serre-subcategory})
pour tout champ algébrique $\mathcal{X}$,
\item pour tout morphisme lisse de champs algébriques
$f : \mathcal{Y} \to \mathcal{X}$, le foncteur $f^*$ envoie
$\mathcal{M}_\mathcal{X}$ dans $\mathcal{M}_\mathcal{Y}$,
\item si $f_i : \mathcal{X}_i \to \mathcal{X}$ est une famille de morphismes
lisses de champs algébriques telle que
$|\mathcal{X}| = \bigcup |f_i|(|\mathcal{X}_i|)$, alors un objet
$\mathcal{F}$ de
$\textit{Mod}(\mathcal{X}_\etale, \mathcal{O}_\mathcal{X})$
appartient à $\mathcal{M}_\mathcal{X}$ si et seulement si
$f_i^*\mathcal{F}$ appartient à $\mathcal{M}_{\mathcal{X}_i}$ pour tout $i$, et
\item si $f : \mathcal{Y} \to \mathcal{X}$ est un morphisme de champs
algébriques tel que $\mathcal{X}$ et $\mathcal{Y}$ soient représentables
par des schémas affines, alors $R^if_*$ envoie $\mathcal{M}_\mathcal{Y}$
dans $\mathcal{M}_\mathcal{X}$.
\end{enumerate}
Alors, pour tout morphisme quasi-compact et quasi-séparé
$f : \mathcal{Y} \to \mathcal{X}$ de champs algébriques,
$R^if_*$ envoie $\mathcal{M}_\mathcal{Y}$
dans $\mathcal{M}_\mathcal{X}$. (Les images directes supérieures sont calculées pour la topologie
étale.)
\end{lemma}

\begin{proof}
Soit $f : \mathcal{Y} \to \mathcal{X}$ un morphisme quasi-compact et quasi-séparé
de champs algébriques, et soit $\mathcal{F}$ un objet de
$\mathcal{M}_\mathcal{Y}$. Choisissons un morphisme lisse surjectif
$\mathcal{U} \to \mathcal{X}$, où $\mathcal{U}$ est représentable par
un schéma. D'après
Faisceaux sur les champs, lemme
\ref{stacks-sheaves-lemma-base-change-higher-direct-images}
la prise des images directes supérieures commute au changement de base.
L'hypothèse (2) montre que l'image inverse de $\mathcal{F}$
sur $\mathcal{U} \times_\mathcal{X} \mathcal{Y}$ appartient à
$\mathcal{M}_{\mathcal{U} \times_\mathcal{X} \mathcal{Y}}$
car la projection
$\mathcal{U} \times_\mathcal{X} \mathcal{Y} \to \mathcal{Y}$
est lisse comme changement de base d'un morphisme lisse. Ainsi, (3) montre que nous pouvons
remplacer $\mathcal{Y} \to \mathcal{X}$ par la projection
$\mathcal{U} \times_\mathcal{X} \mathcal{Y} \to \mathcal{U}$.
Autrement dit, nous pouvons supposer que $\mathcal{X}$
est représentable par un schéma.
En appliquant encore (3), nous voyons que la question est locale sur
$\mathcal{X}$ pour la topologie de Zariski ; nous pouvons donc supposer que $\mathcal{X}$ est représentable par
un schéma affine. Comme $f$ est quasi-compact, il s'ensuit que
$\mathcal{Y}$ est également quasi-compacte. Nous pouvons donc choisir un morphisme lisse surjectif
$g : \mathcal{V} \to \mathcal{Y}$, où $\mathcal{V}$ est représentable
par un schéma affine.

\medskip\noindent
Dans cette situation, nous disposons de la suite spectrale
$$
E_2^{p, q} = R^q(f \circ g_p)_*g_p^*\mathcal{F}
\Rightarrow
R^{p + q}f_*\mathcal{F}
$$
de
Faisceaux sur les champs, proposition
\ref{stacks-sheaves-proposition-smooth-covering-compute-direct-image}.
Rappelons qu'il s'agit d'une suite spectrale du premier quadrant ; nous pouvons donc
appliquer la dernière partie de Homologie, lemme \ref{homology-lemma-first-quadrant-ss}.
Notons que les morphismes
$$
g_p : \mathcal{V}_p =
\mathcal{V} \times_\mathcal{Y} \ldots \times_\mathcal{Y} \mathcal{V}
\longrightarrow
\mathcal{Y}
$$
sont lisses comme composés de changements de base du morphisme lisse $g$.
Ainsi, les faisceaux $g_p^*\mathcal{F}$ appartiennent à
$\mathcal{M}_{\mathcal{V}_p}$ d'après (2). Il suffit donc de démontrer que les
images directes supérieures des objets de $\mathcal{M}_{\mathcal{V}_p}$ par
les morphismes
$$
\mathcal{V}_p =
\mathcal{V} \times_\mathcal{Y} \ldots \times_\mathcal{Y} \mathcal{V}
\longrightarrow
\mathcal{X}
$$
appartiennent à $\mathcal{M}_\mathcal{X}$. Les champs algébriques $\mathcal{V}_p$
sont quasi-compacts et quasi-séparés d'après
Morphismes de champs, lemme
\ref{stacks-morphisms-lemma-quasi-compact-quasi-separated-permanence}.
Bien entendu, chaque $\mathcal{V}_p$ est représentable par un espace algébrique
(la diagonale du champ algébrique $\mathcal{Y}$ est représentable
par des espaces algébriques). Nous sommes donc ramenés au cas où
$\mathcal{Y}$ est représentable par un espace algébrique et $\mathcal{X}$
par un schéma affine.

\medskip\noindent
Lorsque $\mathcal{Y}$ est représentable par un espace
algébrique et $\mathcal{X}$ par un schéma affine, choisissons
de nouveau un morphisme lisse surjectif $\mathcal{V} \to \mathcal{Y}$, où
$\mathcal{V}$ est représentable par un schéma affine. En reprenant
l'argument précédent, nous sommes de nouveau ramenés aux morphismes
$\mathcal{V}_p \to \mathcal{X}$. Mais, dans la situation présente, les champs
algébriques $\mathcal{V}_p$ sont représentables par des schémas quasi-compacts et quasi-séparés
(car la diagonale d'un espace algébrique est représentable par des
schémas).

\medskip\noindent
Nous pouvons donc supposer que $\mathcal{Y}$ est représentable par un schéma et
$\mathcal{X}$ par un schéma affine. Choisissons (encore)
un morphisme lisse surjectif $\mathcal{V} \to \mathcal{Y}$, où
$\mathcal{V}$ est représentable par un schéma affine. Dans ce cas, tous
les champs algébriques $\mathcal{V}_p$ sont représentables par des schémas
séparés (car la diagonale d'un schéma est séparée).

\medskip\noindent
Nous pouvons donc supposer que $\mathcal{Y}$ est représentable par un schéma séparé et
$\mathcal{X}$ par un schéma affine. Choisissons (une fois encore)
un morphisme lisse surjectif $\mathcal{V} \to \mathcal{Y}$, où
$\mathcal{V}$ est représentable par un schéma affine. Dans ce cas, tous
les champs algébriques $\mathcal{V}_p$ sont représentables par des schémas affines
(car la diagonale d'un schéma séparé est une immersion fermée, donc affine),
et ce cas est couvert par l'hypothèse (4).
Ceci achève la démonstration.
\end{proof}

\noindent
Voici la version pour la topologie fppf.

\begin{lemma}
\label{stacks-cohomology-lemma-general-pushforward-fppf}
Soit $\mathcal{M}$ une règle qui associe à tout champ algébrique
$\mathcal{X}$ une sous-catégorie $\mathcal{M}_\mathcal{X}$ de
$\textit{Mod}(\mathcal{O}_\mathcal{X})$
telle que
\begin{enumerate}
\item $\mathcal{O}_\mathcal{X}$ soit une sous-catégorie faible de Serre
de $\textit{Mod}(\mathcal{O}_\mathcal{X})$
pour tout champ algébrique $\mathcal{X}$,
\item pour tout morphisme lisse de champs algébriques
$f : \mathcal{Y} \to \mathcal{X}$, le foncteur $f^*$ envoie
$\mathcal{M}_\mathcal{X}$ dans $\mathcal{M}_\mathcal{Y}$,
\item si $f_i : \mathcal{X}_i \to \mathcal{X}$ est une famille de morphismes
lisses de champs algébriques telle que
$|\mathcal{X}| = \bigcup |f_i|(|\mathcal{X}_i|)$, alors un objet
$\mathcal{F}$ de $\textit{Mod}(\mathcal{O}_\mathcal{X})$
appartient à $\mathcal{M}_\mathcal{X}$ si et seulement si
$f_i^*\mathcal{F}$ appartient à $\mathcal{M}_{\mathcal{X}_i}$ pour tout $i$, et
\item si $f : \mathcal{Y} \to \mathcal{X}$ est un morphisme de champs
algébriques tel que $\mathcal{X}$ et $\mathcal{Y}$ soient représentables
par des schémas affines, alors $R^if_*$ envoie $\mathcal{M}_\mathcal{Y}$
dans $\mathcal{M}_\mathcal{X}$.
\end{enumerate}
Alors, pour tout morphisme quasi-compact et quasi-séparé
$f : \mathcal{Y} \to \mathcal{X}$ de champs algébriques,
$R^if_*$ envoie $\mathcal{M}_\mathcal{Y}$
dans $\mathcal{M}_\mathcal{X}$. (Les images directes supérieures sont calculées pour la topologie
fppf.)
\end{lemma}

\begin{proof}
Identique à la démonstration du lemme \ref{stacks-cohomology-lemma-general-pushforward}.
\end{proof}


\section{Modules localement quasi-cohérents}
\label{stacks-cohomology-section-locally-quasi-coherent}

\noindent
Soit $\mathcal{X}$ un champ algébrique. Soit $\mathcal{F}$ un préfaisceau
de $\mathcal{O}_\mathcal{X}$-modules. On peut se demander si $\mathcal{F}$
est {\it localement quasi-cohérent} ; voir
Faisceaux sur les champs, définition
\ref{stacks-sheaves-definition-locally-quasi-coherent}.
En bref, cela signifie que $\mathcal{F}$ est un $\mathcal{O}_\mathcal{X}$-module
pour la topologie étale tel que, pour tout morphisme $f : U \to \mathcal{X}$,
la restriction $f^*\mathcal{F}|_{U_\etale}$ soit quasi-cohérente
sur $U_\etale$. (La définition proprement dite est légèrement différente, mais
équivalente.) Un fait utile est que
$$
\textit{LQCoh}(\mathcal{O}_\mathcal{X}) \subset
\textit{Mod}(\mathcal{X}_\etale, \mathcal{O}_\mathcal{X})
$$
est une sous-catégorie faible de Serre ; voir
Faisceaux sur les champs, lemme \ref{stacks-sheaves-lemma-lqc-colimits}.

\begin{lemma}
\label{stacks-cohomology-lemma-check-lqc-on-etale-covering}
Soit $\mathcal{X}$ un champ algébrique. Soit
$f_j : \mathcal{X}_j \to \mathcal{X}$ une famille de morphismes
lisses de champs algébriques telle que
$|\mathcal{X}| =\bigcup |f_j|(|\mathcal{X}_j|)$.
Soit $\mathcal{F}$ un faisceau de $\mathcal{O}_\mathcal{X}$-modules
sur $\mathcal{X}_\etale$. Si chaque $f_j^{-1}\mathcal{F}$
est localement quasi-cohérent, alors $\mathcal{F}$ l'est aussi.
\end{lemma}

\begin{proof}
Nous pouvons remplacer chacun des champs algébriques $\mathcal{X}_j$ par
un schéma $U_j$ (en utilisant le fait que tout champ algébrique possède un recouvrement lisse par
un schéma et que les composés de morphismes lisses sont lisses ; voir
Morphismes de champs, lemme \ref{stacks-morphisms-lemma-composition-smooth}).
L'image inverse de $\mathcal{F}$ sur $(\Sch/U_j)_\etale$ est encore
localement quasi-cohérente ; voir
Faisceaux sur les champs, lemme \ref{stacks-sheaves-lemma-pullback-lqc}.
Alors $f = \coprod f_j : U = \coprod U_j \to \mathcal{X}$ est un morphisme
lisse surjectif. Soit $x$ un objet de $\mathcal{X}$. D'après
Faisceaux sur les champs, lemme
\ref{stacks-sheaves-lemma-surjective-flat-locally-finite-presentation}
il existe un recouvrement étale $\{x_i \to x\}_{i \in I}$
tel que chaque $x_i$ se relève en un objet $u_i$ de $(\Sch/U)_\etale$.
Cela signifie simplement que $x$, $x_i$ sont au-dessus de schémas $V$, $V_i$, que
$\{V_i \to V\}$ est un recouvrement étale et que $x_i$ provient
d'un morphisme $u_i : V_i \to U$. La restriction
$x_i^*\mathcal{F}|_{V_{i, \etale}}$ est égale à la restriction
de $f^*\mathcal{F}$ à $V_{i, \etale}$ ; voir
Faisceaux sur les champs, lemme \ref{stacks-sheaves-lemma-comparison}.
Ainsi, $x^*\mathcal{F}|_{V_\etale}$
est un faisceau sur le petit site étale de $V$ dont la restriction
à $V_{i, \etale}$ est quasi-cohérente pour tout $i$.
Il est donc quasi-cohérent, comme souhaité, par exemple d'après
Propriétés des espaces, lemme
\ref{spaces-properties-lemma-characterize-quasi-coherent}.
\end{proof}

\begin{lemma}
\label{stacks-cohomology-lemma-pushforward-locally-quasi-coherent}
Soit $f : \mathcal{X} \to \mathcal{Y}$ un morphisme quasi-compact et
quasi-séparé de champs algébriques. Soit
$\mathcal{F}$ un $\mathcal{O}_\mathcal{X}$-module localement quasi-cohérent
sur $\mathcal{X}_\etale$.
Alors $R^if_*\mathcal{F}$ (calculé pour la topologie étale) est
localement quasi-cohérent sur $\mathcal{Y}_\etale$.
\end{lemma}

\begin{proof}
Nous allons utiliser le
lemme \ref{stacks-cohomology-lemma-general-pushforward}
pour le démontrer. Vérifions ses hypothèses (1) -- (4).
Les parties (1) et (2) résultent de
Faisceaux sur les champs, lemme \ref{stacks-sheaves-lemma-lqc-colimits}.
La partie (3) résulte du
lemme \ref{stacks-cohomology-lemma-check-lqc-on-etale-covering}.
Il suffit donc de démontrer (4).

\medskip\noindent
Supposons que $f : \mathcal{X} \to \mathcal{Y}$ soit un morphisme de champs algébriques
tel que $\mathcal{X}$ et $\mathcal{Y}$ soient représentables par des schémas
affines $X$ et $Y$. Choisissons un objet quelconque $y$ de $\mathcal{Y}$ au-dessus d'un
schéma $V$. Pour plus de clarté, notons $\mathcal{V} = (\Sch/V)_{fppf}$ le
champ algébrique correspondant à $V$. Considérons le diagramme cartésien
$$
\xymatrix{
\mathcal{Z} \ar[d] \ar[r]_g \ar[d]_{f'} & \mathcal{X} \ar[d]^f \\
\mathcal{V} \ar[r]^y & \mathcal{Y}
}
$$
Ainsi, $\mathcal{Z}$ est représentable par le schéma $Z = V \times_Y X$
et $f'$ est quasi-compact et séparé (même affine). D'après
Faisceaux sur les champs, lemme
\ref{stacks-sheaves-lemma-compare-representable-morphism-cohomology}
nous avons
$$
R^if_*\mathcal{F}|_{V_\etale} =
R^if'_{small, *}\big(g^*\mathcal{F}|_{Z_\etale}\big)
$$
Le membre de droite est un faisceau quasi-cohérent sur $V_\etale$ d'après
Cohomologie des espaces, lemme
\ref{spaces-cohomology-lemma-higher-direct-image}.
Il s'ensuit que le membre de gauche est quasi-cohérent, ce qu'il
fallait démontrer.
\end{proof}

\begin{lemma}
\label{stacks-cohomology-lemma-check-lqc-on-flat-covering}
Soit $\mathcal{X}$ un champ algébrique. Soit
$f_j : \mathcal{X}_j \to \mathcal{X}$ une famille de morphismes plats
et localement de présentation finie de champs algébriques telle que
$|\mathcal{X}| =\bigcup |f_j|(|\mathcal{X}_j|)$.
Soit $\mathcal{F}$ un faisceau de $\mathcal{O}_\mathcal{X}$-modules
sur $\mathcal{X}_{fppf}$. Si chaque $f_j^{-1}\mathcal{F}$
est localement quasi-cohérent, alors $\mathcal{F}$ l'est aussi.
\end{lemma}

\begin{proof}
Supposons d'abord qu'il existe un morphisme $a : \mathcal{U} \to \mathcal{X}$
surjectif, plat, localement de présentation finie, quasi-compact
et quasi-séparé, tel que $a^*\mathcal{F}$ soit localement quasi-cohérent.
Il existe alors une suite exacte
$$
0 \to \mathcal{F} \to a_*a^*\mathcal{F} \to b_*b^*\mathcal{F}
$$
où $b$ est le morphisme
$b : \mathcal{U} \times_\mathcal{X} \mathcal{U} \to \mathcal{X}$ ; voir
Faisceaux sur les champs, proposition
\ref{stacks-sheaves-proposition-exactness-cech-complex} et
lemme \ref{stacks-sheaves-lemma-surjective-flat-locally-finite-presentation}.
De plus, l'image inverse $b^*\mathcal{F}$ est l'image inverse de $a^*\mathcal{F}$
par l'une des projections ; elle est donc localement quasi-cohérente
(Faisceaux sur les champs, lemme \ref{stacks-sheaves-lemma-pullback-lqc}).
Les modules $a_*a^*\mathcal{F}$ et $b_*b^*\mathcal{F}$ sont localement
quasi-cohérents d'après le lemme \ref{stacks-cohomology-lemma-pushforward-locally-quasi-coherent}.
(Notons que $a_*$ et $b_*$ ne dépendent pas de la topologie
utilisée pour les calculer.)
Nous en concluons que $\mathcal{F}$ est localement quasi-cohérent ; voir
Faisceaux sur les champs, lemme \ref{stacks-sheaves-lemma-lqc-colimits}.

\medskip\noindent
Nous allons ramener la démonstration du cas général à la
situation du premier paragraphe. Soit $x$ un objet de $\mathcal{X}$
au-dessus du schéma $U$. Nous devons montrer que
$\mathcal{F}|_{U_\etale}$ est un $\mathcal{O}_U$-module quasi-cohérent.
Il suffit de le faire localement sur $U$ pour la topologie de Zariski ; nous pouvons donc
supposer $U$ affine. D'après
Morphismes de champs, lemme
\ref{stacks-morphisms-lemma-surjective-family-flat-locally-finite-presentation}
il existe un recouvrement fppf $\{a_i : U_i \to U\}$ tel que
chaque $x \circ a_i$ se factorise par un certain $f_j$. Ainsi, $a_i^*\mathcal{F}$
est localement quasi-cohérent sur $(\Sch/U_i)_{fppf}$. En raffinant
le recouvrement, nous pouvons supposer que $\{U_i \to U\}_{i = 1, \ldots, n}$
est un recouvrement fppf standard. Alors $x^*\mathcal{F}$ est un
module fppf sur $(\Sch/U)_{fppf}$ dont l'image inverse par le morphisme
$a : U_1 \amalg \ldots \amalg U_n \to U$ est localement quasi-cohérente.
Le premier paragraphe montre donc que $x^*\mathcal{F}$ est localement
quasi-cohérent, ce qui implique assurément que $\mathcal{F}|_{U_\etale}$
est quasi-cohérent.
\end{proof}






\section{Morphismes de comparaison plats}
\label{stacks-cohomology-section-flat-comparison}

\noindent
Soit $\mathcal{X}$ un champ algébrique et soit $\mathcal{F}$ un objet
de $\textit{Mod}(\mathcal{X}_\etale, \mathcal{O}_\mathcal{X})$.
Étant donné un objet $x$ de $\mathcal{X}$ au-dessus du schéma $U$, la
restriction $\mathcal{F}|_{U_\etale}$ est la restriction de
$x^{-1}\mathcal{F}$ au petit site étale de $U$ ; voir
Faisceaux sur les champs, définition \ref{stacks-sheaves-definition-pullback}.
Soit ensuite $\varphi : x \to x'$ un morphisme de $\mathcal{X}$ au-dessus
d'un morphisme de schémas $f : U \to U'$. Nous avons donc un diagramme $2$-commutatif
$$
\xymatrix{
U \ar[rd]_x \ar[rr]_f & & U' \ar[ld]^{x'} \\
& \mathcal{X}
}
$$
À $\varphi$ est associé un morphisme de comparaison entre les restrictions
\begin{equation}
\label{stacks-cohomology-equation-comparison-modules}
c_\varphi :
f_{small}^*(\mathcal{F}|_{U'_\etale})
\longrightarrow
\mathcal{F}|_{U_\etale}
\end{equation}
voir Faisceaux sur les champs, équation
(\ref{stacks-sheaves-equation-comparison-modules}).
Dans cette situation, nous pouvons considérer la propriété suivante
de $\mathcal{F}$.

\begin{definition}
\label{stacks-cohomology-definition-flat-base-change}
Soit $\mathcal{X}$ un champ algébrique et soit $\mathcal{F}$ dans
$\textit{Mod}(\mathcal{X}_\etale, \mathcal{O}_\mathcal{X})$.
Nous disons que $\mathcal{F}$ possède la {\it propriété de changement de base plat}\footnote{Cette
terminologie n'est peut-être pas standard.}
si et seulement si $c_\varphi$ est un isomorphisme dès que $f$ est plat.
\end{definition}

\noindent
Voici un lemme donnant quelques propriétés de cette notion.

\begin{lemma}
\label{stacks-cohomology-lemma-check-flat-comparison-on-etale-covering}
Soit $\mathcal{X}$ un champ algébrique. Soit $\mathcal{F}$
un $\mathcal{O}_\mathcal{X}$-module sur $\mathcal{X}_\etale$.
\begin{enumerate}
\item Si $\mathcal{F}$ possède la propriété de changement de base plat, alors, pour tout morphisme
$g : \mathcal{Y} \to \mathcal{X}$ de champs algébriques, son
image inverse $g^*\mathcal{F}$ la possède aussi.
\item La sous-catégorie pleine de
$\textit{Mod}(\mathcal{X}_\etale, \mathcal{O}_\mathcal{X})$
formée des modules possédant la propriété de changement de base plat
est une sous-catégorie faible de Serre.
\item Soit $f_i : \mathcal{X}_i \to \mathcal{X}$ une famille de morphismes
lisses de champs algébriques telle que
$|\mathcal{X}| = \bigcup_i |f_i|(|\mathcal{X}_i|)$. Si chaque
$f_i^*\mathcal{F}$ possède la propriété de changement de base plat, alors
$\mathcal{F}$ la possède aussi.
\item La catégorie des $\mathcal{O}_\mathcal{X}$-modules
sur $\mathcal{X}_\etale$ possédant la propriété de changement de base plat
admet des limites inductives, qui coïncident avec celles de
$\textit{Mod}(\mathcal{X}_\etale, \mathcal{O}_\mathcal{X})$.
\item Si $\mathcal{F}$ et $\mathcal{G}$ sont dans
$\textit{Mod}(\mathcal{X}_\etale, \mathcal{O}_\mathcal{X})$
et possèdent la propriété de changement de base plat, alors le produit tensoriel
$\mathcal{F} \otimes_{\mathcal{O}_\mathcal{X}} \mathcal{G}$
possède cette propriété.
\item Si $\mathcal{F}$ et $\mathcal{G}$ sont dans
$\textit{Mod}(\mathcal{X}_\etale, \mathcal{O}_\mathcal{X})$
avec $\mathcal{F}$ de présentation finie et $\mathcal{G}$ possédant
la propriété de changement de base plat, alors le faisceau
$\SheafHom_{\mathcal{O}_\mathcal{X}}(\mathcal{F}, \mathcal{G})$
possède la propriété de changement de base plat.
\end{enumerate}
\end{lemma}

\begin{proof}
Soit $g : \mathcal{Y} \to \mathcal{X}$ comme en (1).
Soit $y$ un objet de $\mathcal{Y}$ au-dessus d'un schéma $V$. D'après
Faisceaux sur les champs, lemme \ref{stacks-sheaves-lemma-comparison}
nous avons
$(g^*\mathcal{F})|_{V_\etale} = \mathcal{F}|_{V_\etale}$.
De plus, un morphisme de comparaison pour le faisceau $g^*\mathcal{F}$ sur $\mathcal{Y}$
est un cas particulier d'un morphisme de comparaison pour le faisceau $\mathcal{F}$ sur
$\mathcal{X}$ ; voir
Faisceaux sur les champs, lemme \ref{stacks-sheaves-lemma-comparison}.
Ainsi, (1) est clair.

\medskip\noindent
Démonstration de (2). Nous utilisons la caractérisation des sous-catégories faibles de Serre de
Homologie, lemme \ref{homology-lemma-characterize-weak-serre-subcategory}.
Les noyaux et conoyaux des
morphismes entre faisceaux possédant la propriété de changement de base plat
possèdent eux aussi cette propriété. C'est clair, car
$f_{small}^*$ est exact pour un morphisme plat de schémas et les
foncteurs de restriction $(-)|_{U_\etale}$ sont exacts (puisque nous
travaillons avec la topologie étale). Enfin, si
$0 \to \mathcal{F}_1 \to \mathcal{F}_2 \to \mathcal{F}_3 \to 0$
est une suite exacte courte dans
$\textit{Mod}(\mathcal{X}_\etale, \mathcal{O}_\mathcal{X})$
et que les deux faisceaux extrêmes possèdent la propriété de changement de base plat, alors
celui du milieu la possède aussi, toujours par l'exactitude de
$f_{small}^*$ et des foncteurs de restriction (ainsi que par le lemme des cinq).

\medskip\noindent
Démonstration de (3).
Soit $f_i : \mathcal{X}_i \to \mathcal{X}$ une famille conjointement surjective de
morphismes lisses de champs algébriques et supposons que chaque $f_i^*\mathcal{F}$
possède la propriété de changement de base plat. D'après la partie (1), la définition
d'un champ algébrique et le fait que les composés de morphismes lisses
sont lisses (voir
Morphismes de champs, lemme \ref{stacks-morphisms-lemma-composition-smooth})
nous pouvons supposer que chaque $\mathcal{X}_i$ est représentable par un schéma.
Soit $\varphi : x \to x'$ un morphisme de $\mathcal{X}$ au-dessus
d'un morphisme plat de schémas $a : U \to U'$. D'après
Faisceaux sur les champs, lemme
\ref{stacks-sheaves-lemma-surjective-flat-locally-finite-presentation}
il existe une famille conjointement surjective de morphismes étales
$U'_i \to U'$ telle que $U'_i \to U' \to \mathcal{X}$ se factorise par
$\mathcal{X}_i$. Nous obtenons ainsi des diagrammes commutatifs
$$
\xymatrix{
U_i = U \times_{U'} U_i' \ar[r]_-{a_i} \ar[d] &
U_i' \ar[r]_{x_i'} \ar[d] & \mathcal{X}_i \ar[d]^{f_i} \\
U \ar[r]^a & U' \ar[r]^{x'} & \mathcal{X}
}
$$
Notons que chaque $a_i$ est un morphisme plat de schémas, comme changement de base de $a$.
Notons $\psi_i : x_i \to x'_i$ le morphisme de $\mathcal{X}_i$ au-dessus de
$a_i$ et de but $x_i'$. Par hypothèse, le morphisme de comparaison
$c_{\psi_i} :
(a_i)_{small}^*\big(f_i^*\mathcal{F}|_{(U'_i)_\etale}\big)
\to f_i^*\mathcal{F}|_{(U_i)_\etale}$ est un isomorphisme.
Comme les flèches verticales $U_i' \to U'$ et $U_i \to U$ sont étales,
les faisceaux $f_i^*\mathcal{F}|_{(U_i')_\etale}$ et
$f_i^*\mathcal{F}|_{(U_i)_\etale}$ sont les restrictions de
$\mathcal{F}|_{U'_\etale}$ et $\mathcal{F}|_{U_\etale}$,
et le morphisme $c_{\psi_i}$ est la restriction de $c_\varphi$ à
$(U_i)_\etale$ ; voir
Faisceaux sur les champs, lemme \ref{stacks-sheaves-lemma-comparison}.
Puisque $\{U_i \to U\}$ est un recouvrement étale, il s'ensuit
que le morphisme de comparaison $c_\varphi$ est un isomorphisme, ce qu'il
fallait démontrer.

\medskip\noindent
Démonstration de (4). Soit $\mathcal{I} \to
\textit{Mod}(\mathcal{X}_\etale, \mathcal{O}_\mathcal{X})$,
$i \mapsto \mathcal{F}_i$ un diagramme, et supposons que chaque $\mathcal{F}_i$
possède la propriété de changement de base plat. Soit $\varphi : x \to x'$ un morphisme
de $\mathcal{X}$ au-dessus du morphisme plat de schémas $f : U \to U'$.
Rappelons que $\colim_i \mathcal{F}_i$ est le faisceau associé à la limite inductive des préfaisceaux.
Comme nous utilisons la topologie étale, il est clair que
$$
(\colim_i \mathcal{F}_i)|_{U_\etale} =
\colim_i {\mathcal{F}_i}|_{U_\etale}
$$
et de même pour la restriction à $U'_\etale$. Ainsi,
\begin{align*}
f_{small}^*((\colim_i \mathcal{F}_i)|_{U'_\etale})
& =
f_{small}^*(\colim_i {\mathcal{F}_i}|_{U'_\etale}) \\
& =
\colim_i f_{small}^*({\mathcal{F}_i}|_{U'_\etale}) \\
& \xrightarrow{\colim c_\varphi}
\colim_i \mathcal{F}_i|_{U_\etale} \\
& =
(\colim_i \mathcal{F}_i)|_{U_\etale}
\end{align*}
Pour la deuxième égalité, nous avons utilisé que $f_{small}^*$ commute aux limites inductives
(en tant qu'adjoint à gauche). La flèche est un isomorphisme puisque chaque
$\mathcal{F}_i$ possède la propriété de changement de base plat. La limite inductive
possède donc cette propriété, et (4) est vraie.

\medskip\noindent
La partie (5) vaut parce que les produits tensoriels commutent aux images inverses ; voir
Modules sur les sites, lemme \ref{sites-modules-lemma-tensor-product-pullback}.
Nous omettons les détails.

\medskip\noindent
Soient $\mathcal{F}$ et $\mathcal{G}$ comme en (6). Puisque $\mathcal{F}$
est quasi-cohérent, il possède la propriété de changement de base plat d'après
Faisceaux sur les champs, lemme \ref{stacks-sheaves-lemma-quasi-coherent}.
Soit $\varphi : x \to x'$ un morphisme de $\mathcal{X}$ au-dessus
du morphisme plat de schémas $f : U \to U'$.
Comme nous utilisons la topologie étale, nous avons
$$
\SheafHom_{\mathcal{O}_\mathcal{X}}(\mathcal{F}, \mathcal{G})|_{U_\etale} =
\SheafHom_{\mathcal{O}_U}(\mathcal{F}|_{U_\etale}, \mathcal{G}|_{U_\etale})
$$
et de même pour la restriction à $U'_\etale$ (nous omettons les détails). Ainsi,
\begin{align*}
f_{small}^*(
\SheafHom_{\mathcal{O}_\mathcal{X}}(\mathcal{F}, \mathcal{G})|_{U'_\etale})
& =
f_{small}^*(
\SheafHom_{\mathcal{O}_{U'}}(
\mathcal{F}|_{U'_\etale}, \mathcal{G}|_{U'_\etale})) \\
& =
\SheafHom_{\mathcal{O}_{U'}}(
f_{small}^*(\mathcal{F}|_{U'_\etale}),
f_{small}^*(\mathcal{G}|_{U'_\etale})) \\
& \xrightarrow{c_\varphi}
\SheafHom_{\mathcal{O}_U}(\mathcal{F}|_{U_\etale}, \mathcal{G}|_{U_\etale}) \\
& =
\SheafHom_{\mathcal{O}_\mathcal{X}}(\mathcal{F}, \mathcal{G})|_{U_\etale}
\end{align*}
Ici, la deuxième égalité est donnée par
Modules sur les sites, lemme \ref{sites-modules-lemma-pullback-internal-hom}
qui utilise que $f : U \to U'$ est plat et que, par conséquent, le morphisme
de sites annelés $f_{small}$ est plat lui aussi.
La flèche est un isomorphisme puisque $\mathcal{F}$ et $\mathcal{G}$
possèdent toutes deux la propriété de changement de base plat. Ainsi, notre $\SheafHom$
possède lui aussi cette propriété, comme souhaité.
\end{proof}

\begin{lemma}
\label{stacks-cohomology-lemma-flat-comparison}
Soit $f : \mathcal{X} \to \mathcal{Y}$ un morphisme quasi-compact et
quasi-séparé de champs algébriques. Soit
$\mathcal{F}$ un objet de
$\textit{Mod}(\mathcal{X}_\etale, \mathcal{O}_\mathcal{X})$
qui soit localement quasi-cohérent et possède la propriété de changement de base plat.
Alors chaque $R^if_*\mathcal{F}$ (calculé pour la topologie étale)
possède la propriété de changement de base plat.
\end{lemma}

\begin{proof}
Nous allons utiliser le
lemme \ref{stacks-cohomology-lemma-general-pushforward}
pour le démontrer. Pour tout champ algébrique $\mathcal{X}$, notons
$\textit{LQCoh}^{fbc}(\mathcal{O}_\mathcal{X})$ la sous-catégorie pleine de
$\textit{Mod}(\mathcal{X}_\etale, \mathcal{O}_\mathcal{X})$
formée des faisceaux localement quasi-cohérents possédant la propriété de
changement de base plat. Une fois vérifiées les conditions (1) -- (4) du
lemme \ref{stacks-cohomology-lemma-general-pushforward}
le lemme en résultera. Les propriétés (1), (2) et (3) résultent de
Faisceaux sur les champs, lemmes \ref{stacks-sheaves-lemma-pullback-lqc} et
\ref{stacks-sheaves-lemma-lqc-colimits}
ainsi que des
lemmes \ref{stacks-cohomology-lemma-check-lqc-on-etale-covering} et
\ref{stacks-cohomology-lemma-check-flat-comparison-on-etale-covering}.
Il suffit donc de démontrer la partie (4).

\medskip\noindent
Supposons que $f : \mathcal{X} \to \mathcal{Y}$ soit un morphisme de champs algébriques
tel que $\mathcal{X}$ et $\mathcal{Y}$ soient représentables par des schémas
affines $X$ et $Y$. Dans ce cas, supposons que
$\psi : y \to y'$ soit un morphisme de $\mathcal{Y}$ au-dessus
d'un morphisme plat de schémas $b : V \to V'$. Pour plus de clarté, notons
$\mathcal{V} = (\Sch/V)_{fppf}$ et $\mathcal{V}' = (\Sch/V')_{fppf}$
les champs algébriques correspondants. Considérons le diagramme
de champs algébriques
$$
\xymatrix{
\mathcal{Z} \ar[d]_{f''} \ar[r]_a &
\mathcal{Z}' \ar[r]_{x'} \ar[d]_{f'} & \mathcal{X} \ar[d]^f \\
\mathcal{V} \ar[r]^b & \mathcal{V}' \ar[r]^{y'} & \mathcal{Y}
}
$$
dont les deux carrés sont cartésiens. Comme $f$ est représentable par des schémas
(et quasi-compact et séparé -- même affine), nous voyons que $\mathcal{Z}$ et
$\mathcal{Z}'$ sont représentables par des schémas $Z$ et $Z'$ et qu'en
fait $Z = V \times_{V'} Z'$. Puisque $\mathcal{F}$ possède la propriété de
changement de base plat, nous voyons que
$$
a_{small}^*\big(\mathcal{F}|_{Z'_\etale}\big)
\longrightarrow
\mathcal{F}|_{Z_\etale}
$$
est un isomorphisme. De plus,
$$
R^if_*\mathcal{F}|_{V'_\etale} =
R^i(f')_{small, *}\big(\mathcal{F}|_{Z'_\etale}\big)
$$
et
$$
R^if_*\mathcal{F}|_{V_\etale} =
R^i(f'')_{small, *}\big(\mathcal{F}|_{Z_\etale}\big)
$$
d'après
Faisceaux sur les champs, lemme
\ref{stacks-sheaves-lemma-compare-representable-morphism-cohomology}.
Nous voyons donc que le morphisme de comparaison
$$
c_\psi :
b_{small}^*(R^if_*\mathcal{F}|_{V'_\etale})
\longrightarrow
R^if_*\mathcal{F}|_{V_\etale}
$$
est un isomorphisme d'après
Cohomologie des espaces, lemme
\ref{spaces-cohomology-lemma-flat-base-change-cohomology}.
Ainsi, $R^if_*\mathcal{F}$ possède la propriété de changement de base plat.
Puisque $R^if_*\mathcal{F}$ est localement quasi-cohérent d'après le
lemme \ref{stacks-cohomology-lemma-pushforward-locally-quasi-coherent}
nous avons terminé.
\end{proof}





\section[Modules quasi-cohérents et changement de base plat]{Modules localement quasi-cohérents possédant la propriété de changement de base plat}
\label{stacks-cohomology-section-loc-qcoh-flat-base-change}

\noindent
Soit $\mathcal{X}$ un champ algébrique. Nous noterons\footnote{Nous nous excusons
de cette notation affreuse.}
$$
\textit{LQCoh}^{fbc}(\mathcal{O}_\mathcal{X})
\subset
\textit{Mod}(\mathcal{X}_\etale, \mathcal{O}_\mathcal{X})
$$
la sous-catégorie pleine dont les objets sont les
$\mathcal{O}_\mathcal{X}$-modules étales $\mathcal{F}$
qui sont à la fois localement quasi-cohérents
(section \ref{stacks-cohomology-section-locally-quasi-coherent})
et possèdent la propriété de changement de base plat
(section \ref{stacks-cohomology-section-flat-comparison}).
Nous avons
$$
\QCoh(\mathcal{O}_\mathcal{X}) \subset
\textit{LQCoh}^{fbc}(\mathcal{O}_\mathcal{X})
$$
d'après Faisceaux sur les champs, lemme \ref{stacks-sheaves-lemma-quasi-coherent}.

\begin{proposition}
\label{stacks-cohomology-proposition-loc-qcoh-flat-base-change}
Résumé des résultats sur les modules localement quasi-cohérents possédant la propriété de
changement de base plat.
\begin{enumerate}
\item Soit $\mathcal{X}$ un champ algébrique.
Si $\mathcal{F}$ appartient à $\textit{LQCoh}^{fbc}(\mathcal{O}_\mathcal{X})$,
alors $\mathcal{F}$ est un faisceau pour la topologie fppf, c'est-à-dire
un objet de $\textit{Mod}(\mathcal{O}_\mathcal{X})$.
\item La catégorie $\textit{LQCoh}^{fbc}(\mathcal{O}_\mathcal{X})$
est une sous-catégorie faible de Serre à la fois de $\textit{Mod}(\mathcal{O}_\mathcal{X})$
et de $\textit{Mod}(\mathcal{X}_\etale, \mathcal{O}_\mathcal{X})$.
\item L'image inverse $f^*$ par tout morphisme de champs algébriques
$f : \mathcal{X} \to \mathcal{Y}$ induit un foncteur
$f^* : \textit{LQCoh}^{fbc}(\mathcal{O}_\mathcal{Y}) \to
\textit{LQCoh}^{fbc}(\mathcal{O}_\mathcal{X})$.
\item Si $f : \mathcal{X} \to \mathcal{Y}$ est un morphisme
quasi-compact et quasi-séparé de champs algébriques
et si $\mathcal{F}$ est un objet de
$\textit{LQCoh}^{fbc}(\mathcal{O}_\mathcal{X})$, alors
\begin{enumerate}
\item l'image directe totale $Rf_*\mathcal{F}$ et les images directes
supérieures $R^if_*\mathcal{F}$ peuvent être calculées indifféremment pour la topologie étale ou
la topologie fppf, avec le même résultat, et
\item chaque $R^if_*\mathcal{F}$ est un objet de
$\textit{LQCoh}^{fbc}(\mathcal{O}_\mathcal{Y})$.
\end{enumerate}
\item La catégorie $\textit{LQCoh}^{fbc}(\mathcal{O}_\mathcal{X})$ admet des
limites inductives, qui coïncident avec celles de
$\textit{Mod}(\mathcal{X}_\etale, \mathcal{O}_\mathcal{X})$
ainsi qu'avec celles de $\textit{Mod}(\mathcal{O}_\mathcal{X})$.
\item Si $\mathcal{F}$ et $\mathcal{G}$ appartiennent à
$\textit{LQCoh}^{fbc}(\mathcal{O}_\mathcal{X})$, alors le produit tensoriel
$\mathcal{F} \otimes_{\mathcal{O}_\mathcal{X}} \mathcal{G}$
appartient à $\textit{LQCoh}^{fbc}(\mathcal{O}_\mathcal{X})$.
\item Si $\mathcal{F}$ est de présentation finie et si $\mathcal{G}$ appartient à
$\textit{LQCoh}^{fbc}(\mathcal{O}_\mathcal{X})$, alors
$\SheafHom_{\mathcal{O}_\mathcal{X}}(\mathcal{F}, \mathcal{G})$
appartient à $\textit{LQCoh}^{fbc}(\mathcal{O}_\mathcal{X})$.
\end{enumerate}
\end{proposition}

\begin{proof}
La partie (1) est
Faisceaux sur les champs, lemme
\ref{stacks-sheaves-lemma-lqc-flat-base-change-fppf-sheaf}.

\medskip\noindent
La partie (2) pour l'inclusion
$\textit{LQCoh}^{fbc}(\mathcal{O}_\mathcal{X}) \subset
\textit{Mod}(\mathcal{X}_\etale, \mathcal{O}_\mathcal{X})$
a été établie dans la démonstration du
lemme \ref{stacks-cohomology-lemma-flat-comparison}.
Démontrons (2) pour l'inclusion
$\textit{LQCoh}^{fbc}(\mathcal{O}_\mathcal{X}) \subset
\textit{Mod}(\mathcal{O}_\mathcal{X})$.
Soit $\varphi : \mathcal{F} \to \mathcal{G}$ un morphisme entre des
objets de $\textit{LQCoh}^{fbc}(\mathcal{O}_\mathcal{X})$. Puisque
$\Ker(\varphi)$
est le même qu'il soit calculé pour la topologie étale ou pour la topologie
fppf, nous voyons que $\Ker(\varphi)$ appartient à
$\textit{LQCoh}^{fbc}(\mathcal{O}_\mathcal{X})$
d'après le cas étale. D'autre part,
le conoyau calculé pour la topologie fppf est le faisceau fppf associé
au conoyau calculé pour la topologie étale. Or ce
conoyau étale appartient à $\textit{LQCoh}^{fbc}(\mathcal{O}_\mathcal{X})$,
donc il est un faisceau fppf
d'après (1), et nous voyons que le conoyau appartient à
$\textit{LQCoh}^{fbc}(\mathcal{O}_\mathcal{X})$.
Enfin, supposons que
$$
0 \to \mathcal{F}_1 \to \mathcal{F}_2 \to \mathcal{F}_3 \to 0
$$
soit une suite exacte dans $\textit{Mod}(\mathcal{O}_\mathcal{X})$
(c'est-à-dire pour la topologie fppf), avec $\mathcal{F}_1$, $\mathcal{F}_2$
dans $\textit{LQCoh}^{fbc}(\mathcal{O}_\mathcal{X})$. Pour montrer que
$\mathcal{F}_2$ est un objet de $\textit{LQCoh}^{fbc}(\mathcal{O}_\mathcal{X})$,
il suffit de montrer que
la suite est aussi exacte pour la topologie étale. Pour cela, il
suffit de montrer que tout élément de $H^1_{fppf}(x, \mathcal{F}_1)$
devient nul sur les membres d'un recouvrement étale de $x$ (pour tout
objet $x$ de $\mathcal{X}$). Cela est vrai parce que
$H^1_{fppf}(x, \mathcal{F}_1) = H^1_\etale(x, \mathcal{F}_1)$ d'après
Faisceaux sur les champs, lemme \ref{stacks-sheaves-lemma-compare-fppf-etale}
et par la localité de la cohomologie ; voir
Cohomologie sur les sites, lemme
\ref{sites-cohomology-lemma-kill-cohomology-class-on-covering}.
Ceci démontre (2).

\medskip\noindent
La partie (3) résulte du
lemme \ref{stacks-cohomology-lemma-check-flat-comparison-on-etale-covering}
et de
Faisceaux sur les champs, lemme \ref{stacks-sheaves-lemma-pullback-lqc}.

\medskip\noindent
La partie (4)(b), pour $R^if_*\mathcal{F}$ calculé en cohomologie étale,
résulte du lemme \ref{stacks-cohomology-lemma-flat-comparison}.
La partie (4)(a) résulte alors de
Faisceaux sur les champs, lemme \ref{stacks-sheaves-lemma-compare-fppf-etale}
combiné avec (1) ci-dessus.

\medskip\noindent
La partie (5), pour la topologie étale, résulte de
Faisceaux sur les champs, lemme \ref{stacks-sheaves-lemma-lqc-colimits} et du
lemme \ref{stacks-cohomology-lemma-check-flat-comparison-on-etale-covering}.
La version fppf en résulte, puisque la limite inductive pour la topologie
étale est déjà un faisceau fppf d'après la partie (1).

\medskip\noindent
Les parties (6) et (7) résultent des parties correspondantes du
lemme \ref{stacks-cohomology-lemma-check-flat-comparison-on-etale-covering} et de
Faisceaux sur les champs, lemme \ref{stacks-sheaves-lemma-lqc-colimits}.
\end{proof}

\begin{lemma}
\label{stacks-cohomology-lemma-check-lqc-fbc-on-covering}
Soit $\mathcal{X}$ un champ algébrique.
\begin{enumerate}
\item Soit $f_j : \mathcal{X}_j \to \mathcal{X}$ une famille de morphismes
lisses de champs algébriques telle que
$|\mathcal{X}| =\bigcup |f_j|(|\mathcal{X}_j|)$.
Soit $\mathcal{F}$ un faisceau de $\mathcal{O}_\mathcal{X}$-modules
sur $\mathcal{X}_\etale$. Si chaque $f_j^{-1}\mathcal{F}$
appartient à $\textit{LQCoh}^{fpc}(\mathcal{O}_{\mathcal{X}_i})$, alors
$\mathcal{F}$ appartient à
$\textit{LQCoh}^{fbc}(\mathcal{O}_\mathcal{X})$.
\item Soit $f_j : \mathcal{X}_j \to \mathcal{X}$ une famille de morphismes plats
et localement de présentation finie de champs algébriques telle que
$|\mathcal{X}| =\bigcup |f_j|(|\mathcal{X}_j|)$.
Soit $\mathcal{F}$ un faisceau de $\mathcal{O}_\mathcal{X}$-modules
sur $\mathcal{X}_{fppf}$. Si chaque $f_j^{-1}\mathcal{F}$
appartient à $\textit{LQCoh}^{fbc}(\mathcal{O}_{\mathcal{X}_i})$, alors
$\mathcal{F}$ appartient à $\textit{LQCoh}^{fbc}(\mathcal{O}_\mathcal{X})$.
\end{enumerate}
\end{lemma}

\begin{proof}
La partie (1) résulte de la combinaison des
lemmes \ref{stacks-cohomology-lemma-check-lqc-on-etale-covering} et
\ref{stacks-cohomology-lemma-check-flat-comparison-on-etale-covering}.
La démonstration de (2) est analogue à celle du
lemme \ref{stacks-cohomology-lemma-check-lqc-on-flat-covering}.
Soit $\mathcal{F}$ un faisceau de $\mathcal{O}_\mathcal{X}$-modules
sur $\mathcal{X}_{fppf}$.

\medskip\noindent
Supposons d'abord qu'il existe un morphisme $a : \mathcal{U} \to \mathcal{X}$
surjectif, plat, localement de présentation finie, quasi-compact
et quasi-séparé, tel que $a^*\mathcal{F}$ soit localement quasi-cohérent
et possède la propriété de changement de base plat.
Il existe alors une suite exacte
$$
0 \to \mathcal{F} \to a_*a^*\mathcal{F} \to b_*b^*\mathcal{F}
$$
où $b$ est le morphisme
$b : \mathcal{U} \times_\mathcal{X} \mathcal{U} \to \mathcal{X}$ ; voir
Faisceaux sur les champs, proposition
\ref{stacks-sheaves-proposition-exactness-cech-complex} et
lemme \ref{stacks-sheaves-lemma-surjective-flat-locally-finite-presentation}.
De plus, l'image inverse $b^*\mathcal{F}$ est l'image inverse de $a^*\mathcal{F}$
par l'une des projections ; elle est donc localement quasi-cohérente
et possède la propriété de changement de base plat ; voir
proposition \ref{stacks-cohomology-proposition-loc-qcoh-flat-base-change}.
Les modules $a_*a^*\mathcal{F}$ et $b_*b^*\mathcal{F}$ sont localement
quasi-cohérents et possèdent la propriété de changement de base plat d'après la
proposition \ref{stacks-cohomology-proposition-loc-qcoh-flat-base-change}.
Nous en concluons que $\mathcal{F}$ est localement quasi-cohérent et
possède la propriété de changement de base plat d'après la
proposition \ref{stacks-cohomology-proposition-loc-qcoh-flat-base-change}.

\medskip\noindent
Choisissons un schéma $U$ et un morphisme lisse surjectif $x : U \to \mathcal{X}$.
D'après la partie (1), il suffit de montrer que $x^*\mathcal{F}$ est localement
quasi-cohérent et possède la propriété de changement de base plat.
Toujours d'après la partie (1), il suffit de le faire localement sur $U$ pour la topologie de Zariski ;
nous pouvons donc supposer que $U$ est affine. D'après
Morphismes de champs, lemme
\ref{stacks-morphisms-lemma-surjective-family-flat-locally-finite-presentation}
il existe un recouvrement fppf $\{a_i : U_i \to U\}$ tel que
chaque $x \circ a_i$ se factorise par un certain $f_j$. Ainsi, le module
$a_i^*\mathcal{F}$ sur $(\Sch/U_i)_{fppf}$
est localement quasi-cohérent et possède la propriété de changement de base plat.
Après avoir raffiné le recouvrement, nous pouvons supposer que $\{U_i \to U\}_{i = 1, \ldots, n}$
est un recouvrement fppf standard. Alors $x^*\mathcal{F}$ est un
module fppf sur $(\Sch/U)_{fppf}$ dont l'image inverse par le morphisme
$a : U_1 \amalg \ldots \amalg U_n \to U$ est localement quasi-cohérente
et possède la propriété de changement de base plat.
Le paragraphe précédent montre donc que $x^*\mathcal{F}$ est localement
quasi-cohérent et possède la propriété de changement de base plat, comme souhaité.
\end{proof}

\begin{lemma}
\label{stacks-cohomology-lemma-loc-qcoh-fbc-compute-higher-direct-image}
Soit $f : \mathcal{X} \to \mathcal{Y}$ un morphisme de champs algébriques
qui soit quasi-compact, quasi-séparé et
représentable par des espaces algébriques. Soit $\mathcal{F}$ dans
$\textit{LQCoh}^{fbc}(\mathcal{O}_\mathcal{X})$.
Alors, pour tout objet $y : V \to \mathcal{Y}$ de $\mathcal{Y}$, nous avons
$$
(R^if_*\mathcal{F})|_{V_\etale} = R^if'_{small, *}(\mathcal{F}|_{U_\etale})
$$
où $f' : U = V \times_\mathcal{Y} \mathcal{X} \to V$ est le changement
de base de $f$.
\end{lemma}

\begin{proof}
D'après Faisceaux sur les champs, lemme
\ref{stacks-sheaves-lemma-base-change-higher-direct-images}
nous pouvons nous ramener au cas où $\mathcal{X}$
est représenté par $U$ et $\mathcal{Y}$
par $V$. Cela utilise bien entendu aussi que l'image inverse de
$\mathcal{F}$ sur $U$ appartient à $\textit{LQCoh}^{fbc}(\mathcal{O}_U)$
d'après la proposition \ref{stacks-cohomology-proposition-loc-qcoh-flat-base-change}.
Le résultat découle alors de
Faisceaux sur les champs, lemme
\ref{stacks-sheaves-lemma-compare-morphism-cohomology}
et du fait que $R^if_*$ peut être calculé pour la
topologie étale d'après la
proposition \ref{stacks-cohomology-proposition-loc-qcoh-flat-base-change}.
\end{proof}

\begin{lemma}
\label{stacks-cohomology-lemma-loc-qcoh-fbc-affine-direct-image}
Soit $f : \mathcal{X} \to \mathcal{Y}$ un morphisme affine de champs
algébriques. Le foncteur $f_* : \textit{LQCoh}^{fbc}(\mathcal{O}_\mathcal{X}) \to
\textit{LQCoh}^{fbc}(\mathcal{O}_\mathcal{Y})$ est exact et commute
aux sommes directes. Les foncteurs $R^if_*$ pour $i > 0$ s'annulent sur
$\textit{LQCoh}^{fbc}(\mathcal{O}_\mathcal{X})$.
\end{lemma}

\begin{proof}
Les foncteurs existent d'après la proposition \ref{stacks-cohomology-proposition-loc-qcoh-flat-base-change}.
D'après le lemme \ref{stacks-cohomology-lemma-loc-qcoh-fbc-compute-higher-direct-image},
on se ramène au cas d'un morphisme affine d'espaces algébriques,
les images directes supérieures étant prises dans le cadre des modules
quasi-cohérents sur les espaces algébriques. D'après la discussion de
Cohomologie des espaces, section
\ref{spaces-cohomology-section-higher-direct-image}
nous nous ramenons au cas d'un morphisme affine de schémas.
Pour les morphismes affines de schémas, l'annulation des
images directes supérieures sur les modules quasi-cohérents résulte de
Cohomologie des schémas, lemme \ref{coherent-lemma-relative-affine-vanishing}.
L'annulation de $R^1f_*$ implique l'exactitude de $f_*$.
La commutation aux sommes directes résulte de
Morphismes, lemme \ref{morphisms-lemma-affine-equivalence-modules}
par exemple.
\end{proof}










\section{Modules parasites}
\label{stacks-cohomology-section-parasitic}

\noindent
La définition suivante est compatible avec
Descente, définition \ref{descent-definition-parasitic}.

\begin{definition}
\label{stacks-cohomology-definition-parasitic}
Soit $\mathcal{X}$ un champ algébrique.
Un préfaisceau de $\mathcal{O}_\mathcal{X}$-modules $\mathcal{F}$ est
{\it parasite} si $\mathcal{F}(x) = 0$ pour tout objet $x$
de $\mathcal{X}$ au-dessus d'un schéma $U$ tel que le morphisme
correspondant $x : U \to \mathcal{X}$ soit plat.
\end{definition}

\noindent
Voici un lemme donnant quelques propriétés de cette notion.

\begin{lemma}
\label{stacks-cohomology-lemma-parasitic}
Soit $\mathcal{X}$ un champ algébrique. Soit $\mathcal{F}$
un préfaisceau de $\mathcal{O}_\mathcal{X}$-modules.
\begin{enumerate}
\item Si $\mathcal{F}$ est parasite et si
$g : \mathcal{Y} \to \mathcal{X}$ est un morphisme plat de champs algébriques,
alors $g^*\mathcal{F}$ est parasite.
\item Pour $\tau \in \{Zariski, \etale, smooth, syntomic, fppf\}$,
nous avons :
\begin{enumerate}
\item le faisceau $\tau$ associé à un préfaisceau parasite de modules est
parasite, et
\item la sous-catégorie pleine de
$\textit{Mod}(\mathcal{X}_\tau, \mathcal{O}_\mathcal{X})$
formée des modules parasites est une sous-catégorie de Serre.
\end{enumerate}
\item Supposons que $\mathcal{F}$ soit un faisceau pour la topologie étale.
Soit $f_i : \mathcal{X}_i \to \mathcal{X}$ une famille de
morphismes lisses de champs algébriques telle que
$|\mathcal{X}| = \bigcup_i |f_i|(|\mathcal{X}_i|)$. Si chaque
$f_i^*\mathcal{F}$ est parasite, alors $\mathcal{F}$ l'est aussi.
\item Supposons que $\mathcal{F}$ soit un faisceau pour la topologie fppf.
Soit $f_i : \mathcal{X}_i \to \mathcal{X}$ une famille de
morphismes plats et localement de présentation finie de champs algébriques telle que
$|\mathcal{X}| = \bigcup_i |f_i|(|\mathcal{X}_i|)$. Si chaque
$f_i^*\mathcal{F}$ est parasite, alors $\mathcal{F}$ l'est aussi.
\end{enumerate}
\end{lemma}

\begin{proof}
Pour démontrer la partie (1), soit $y$ un objet de $\mathcal{Y}$ au-dessus
d'un schéma $V$ tel que le morphisme correspondant $y : V \to \mathcal{Y}$
soit plat. Alors $g(y) : V \to \mathcal{Y} \to \mathcal{X}$ est plat
comme composé de morphismes plats (voir
Morphismes de champs, lemme \ref{stacks-morphisms-lemma-composition-flat})
et donc $\mathcal{F}(g(y))$ est nul par hypothèse. Puisque
$g^*\mathcal{F} = g^{-1}\mathcal{F}(y) = \mathcal{F}(g(y))$, nous en concluons que
$g^*\mathcal{F}$ est parasite.

\medskip\noindent
Pour démontrer la partie (2)(a), notons que si $\{x_i \to x\}$ est un recouvrement pour la topologie $\tau$
de $\mathcal{X}$, alors chacun des morphismes $x_i \to x$ est au-dessus
d'un morphisme plat de schémas. Ainsi, si $x$ est au-dessus d'un schéma
$U$ tel que $x : U \to \mathcal{X}$ soit plat, il en va de même de tous les
objets $x_i$. Par conséquent, le préfaisceau $\mathcal{F}^+$ (voir
Sites, section \ref{sites-section-sheafification})
est parasite si le préfaisceau $\mathcal{F}$ est
parasite. Cela démontre (2)(a), puisque le faisceau associé à $\mathcal{F}$
est $(\mathcal{F}^+)^+$.

\medskip\noindent
Soit $\mathcal{F}$ un $\tau$-module parasite. Il résulte immédiatement des
définitions que tout sous-module de $\mathcal{F}$ est parasite. D'autre
part, si $\mathcal{F}' \subset \mathcal{F}$ est un sous-module, il est
également clair que le préfaisceau
$x \mapsto \mathcal{F}(x)/\mathcal{F}'(x)$
est parasite. Ainsi, le quotient $\mathcal{F}/\mathcal{F}'$ est un module
parasite d'après (2)(a). Enfin, il faut montrer que, pour toute suite exacte courte
$0 \to \mathcal{F}_1 \to \mathcal{F}_2 \to \mathcal{F}_3 \to 0$
dont $\mathcal{F}_1$ et $\mathcal{F}_3$ sont parasites, $\mathcal{F}_2$
est parasite. Cela résulte immédiatement de l'évaluation en $x$, lorsque cet objet est
au-dessus d'un schéma plat sur $\mathcal{X}$. Ceci démontre (2)(b) ; voir
Homologie, lemme \ref{homology-lemma-characterize-serre-subcategory}.

\medskip\noindent
Soit $f_i : \mathcal{X}_i \to \mathcal{X}$ une famille conjointement surjective de
morphismes lisses de champs algébriques, et supposons que chaque $f_i^*\mathcal{F}$
soit parasite. Soit $x$ un objet de $\mathcal{X}$ au-dessus d'un
schéma $U$ tel que $x : U \to \mathcal{X}$ soit plat. Considérons un recouvrement
lisse surjectif $W_i \to U \times_{x, \mathcal{X}} \mathcal{X}_i$.
Notons $y_i : W_i \to \mathcal{X}_i$ la projection. Il s'ensuit
que $\{f_i(y_i) \to x\}$ est un recouvrement pour la topologie lisse
sur $\mathcal{X}$. Comme un composé de morphismes plats est plat, nous voyons que
$f_i^*\mathcal{F}(y_i) = 0$. D'autre part, comme nous l'avons vu dans la démonstration
de (1), nous avons $f_i^*\mathcal{F}(y_i) = \mathcal{F}(f_i(y_i))$.
Ainsi, pour un certain recouvrement lisse $\{x_i \to x\}_{i \in I}$
de $\mathcal{X}$, nous avons $\mathcal{F}(x_i) = 0$. Cela implique
$\mathcal{F}(x) = 0$, car la topologie lisse est la même
que la topologie étale ; voir
Compléments sur les morphismes, lemme \ref{more-morphisms-lemma-etale-dominates-smooth}.
En effet, $\{x_i \to x\}_{i \in I}$ est au-dessus d'un recouvrement lisse
$\{U_i \to U\}_{i \in I}$ de schémas. D'après le lemme qui vient d'être cité,
il existe un recouvrement étale $\{V_j \to U\}_{j \in J}$ qui
raffine $\{U_i \to U\}_{i \in I}$. Notons $x'_j = x|_{V_j}$.
Alors $\{x'_j \to x\}$ est un recouvrement étale dans $\mathcal{X}$
raffinant $\{x_i \to x\}_{i \in I}$. Cela signifie que l'application
$\mathcal{F}(x) \to \prod_{j \in J} \mathcal{F}(x'_j)$, qui est
injective puisque $\mathcal{F}$ est un faisceau pour la topologie étale,
se factorise par $\mathcal{F}(x) \to \prod_{i \in I} \mathcal{F}(x_i)$,
qui est nulle. Ainsi, $\mathcal{F}(x) = 0$, comme souhaité.

\medskip\noindent
Démonstration de (4) : omise. Indication : elle est semblable, mais plus simple, à celle de (3).
\end{proof}

\noindent
Les modules parasites sont préservés par toute image directe, sans aucune restriction.

\begin{lemma}
\label{stacks-cohomology-lemma-pushforward-parasitic}
Soit $\tau \in \{\etale, fppf\}$.
Soit $\mathcal{X}$ un champ algébrique.
Soit $\mathcal{F}$ un objet parasite de
$\textit{Mod}(\mathcal{X}_\tau, \mathcal{O}_\mathcal{X})$.
\begin{enumerate}
\item $H^i_\tau(\mathcal{X}, \mathcal{F}) = 0$ pour tout $i$.
\item Soit $f : \mathcal{X} \to \mathcal{Y}$ un morphisme de champs algébriques.
Alors $R^if_*\mathcal{F}$ (calculé pour la topologie $\tau$) est un
objet parasite de $\textit{Mod}(\mathcal{Y}_\tau, \mathcal{O}_\mathcal{Y})$.
\end{enumerate}
\end{lemma}

\begin{proof}
Nous ramenons d'abord (2) à (1).
D'après Faisceaux sur les champs, lemme \ref{stacks-sheaves-lemma-pushforward-restriction},
nous voyons que $R^if_*\mathcal{F}$ est le faisceau associé au préfaisceau
$$
y \longmapsto
H^i_\tau\Big(V \times_{y, \mathcal{Y}} \mathcal{X},
\ \text{pr}^{-1}\mathcal{F}\Big)
$$
Ici, $y$ est un objet type de $\mathcal{Y}$ au-dessus du schéma $V$.
D'après le lemme \ref{stacks-cohomology-lemma-parasitic}, il suffit de montrer que
ces groupes de cohomologie sont nuls lorsque $y : V \to \mathcal{Y}$ est plat.
Notons que $\text{pr} : V \times_{y, \mathcal{Y}} \mathcal{X} \to \mathcal{X}$
est plat comme changement de base de $y$. Ainsi, d'après le
lemme \ref{stacks-cohomology-lemma-parasitic}, nous voyons que $\text{pr}^{-1}\mathcal{F}$
est parasite. Il suffit donc de démontrer (1).

\medskip\noindent
Pour démontrer (1), nous pouvons utiliser la suite spectrale de
Faisceaux sur les champs, proposition
\ref{stacks-sheaves-proposition-smooth-covering-compute-cohomology}
pour nous ramener au cas où $\mathcal{X}$
est un champ algébrique représentable par un espace algébrique.
Notons que, dans la suite spectrale, chaque
$f_p^{-1}\mathcal{F} = f_p^*\mathcal{F}$ est un module parasite d'après le
lemme \ref{stacks-cohomology-lemma-parasitic}, car les morphismes
$f_p : \mathcal{U}_p =
\mathcal{U} \times_\mathcal{X} \ldots
\times_\mathcal{X} \mathcal{U} \to \mathcal{X}$ sont plats.
En réutilisant encore une fois cette suite spectrale (comme dans la
démonstration du lemme \ref{stacks-cohomology-lemma-general-pushforward}),
nous nous ramenons au cas où le
champ algébrique $\mathcal{X}$ est représentable par un schéma $X$.
Alors $H^i_\tau(\mathcal{X}, \mathcal{F}) = H^i((\Sch/X)_\tau, \mathcal{F})$.
Dans ce cas, l'annulation résulte facilement d'un argument
avec des recouvrements de {\v C}ech ; voir
Descente, lemme \ref{descent-lemma-cohomology-parasitic}.
\end{proof}

\noindent
Le lemme suivant donne l'une des principales raisons de nous intéresser aux
modules parasites. Pour en comprendre l'énoncé, rappelons que les
foncteurs
$\QCoh(\mathcal{O}_\mathcal{X}) \to
\textit{Mod}(\mathcal{X}_\etale, \mathcal{O}_\mathcal{X})$
et
$\QCoh(\mathcal{O}_\mathcal{X}) \to
\textit{Mod}(\mathcal{O}_\mathcal{X})$
ne sont pas exacts en général.

\begin{lemma}
\label{stacks-cohomology-lemma-exact-sequence-quasi-coherent-parasitic-cohomology}
Soit $\mathcal{X}$ un champ algébrique. Soient
$\alpha : \mathcal{F} \to \mathcal{G}$ et
$\beta : \mathcal{G} \to \mathcal{H}$
des morphismes dans $\QCoh(\mathcal{O}_\mathcal{X})$ tels que
$\beta \circ \alpha = 0$. Les assertions suivantes sont équivalentes :
\begin{enumerate}
\item dans la catégorie abélienne $\QCoh(\mathcal{O}_\mathcal{X})$,
le complexe $\mathcal{F} \to \mathcal{G} \to \mathcal{H}$
est exact en $\mathcal{G}$ ;
\item $\Ker(\beta)/\Im(\alpha)$, calculé soit dans
$\textit{Mod}(\mathcal{X}_\etale, \mathcal{O}_\mathcal{X})$, soit dans
$\textit{Mod}(\mathcal{X}_{fppf}, \mathcal{O}_\mathcal{X})$
est parasite.
\end{enumerate}
\end{lemma}

\begin{proof}
Nous avons $\QCoh(\mathcal{O}_\mathcal{X}) \subset
\textit{LQCoh}^{fbc}(\mathcal{O}_\mathcal{X})$ ; voir
section \ref{stacks-cohomology-section-loc-qcoh-flat-base-change}.
Ainsi, les objets $\Ker(\beta)/\Im(\alpha)$ calculés dans
$\textit{Mod}(\mathcal{X}_\etale, \mathcal{O}_\mathcal{X})$ ou dans
$\textit{Mod}(\mathcal{X}_{fppf}, \mathcal{O}_\mathcal{X})$ coïncident ; voir
proposition \ref{stacks-cohomology-proposition-loc-qcoh-flat-base-change}.
Désormais, nous utiliserons la topologie étale sur $\mathcal{X}$.

\medskip\noindent
Soit $\mathcal{E}$ la cohomologie de
$\mathcal{F} \to \mathcal{G} \to \mathcal{H}$
calculée dans la catégorie abélienne $\QCoh(\mathcal{O}_\mathcal{X})$.
Soit $x : U \to \mathcal{X}$ un morphisme plat, où $U$ est un schéma.
Comme nous utilisons la topologie étale, le foncteur de restriction
$\textit{Mod}(\mathcal{X}_\etale, \mathcal{O}_\mathcal{X})
\to \textit{Mod}(U_\etale, \mathcal{O}_U)$ est exact.
D'autre part, d'après le lemme \ref{stacks-cohomology-lemma-flat-pullback-quasi-coherent} et
Faisceaux sur les champs, lemme \ref{stacks-sheaves-lemma-compare-quasi-coherent}
le foncteur de restriction
$$
\QCoh(\mathcal{O}_\mathcal{X}) \xrightarrow{x^*}
\QCoh((\Sch/U)_\etale, \mathcal{O}) \xrightarrow{{-}|_{U_\etale}}
\QCoh(U_\etale, \mathcal{O}_U)
$$
est lui aussi exact. Nous en concluons que
$\mathcal{E}|_{U_\etale} = (\Ker(\beta)/\Im(\alpha))|_{U_\etale}$.

\medskip\noindent
Si (1) vaut, alors $\mathcal{E} = 0$ ; ainsi, $\Ker(\beta)/\Im(\alpha)$
se restreint à zéro sur $U_\etale$ pour tout $U$ plat sur $\mathcal{X}$, et c'est
la définition d'un module parasite. Si (2) vaut, alors
$\Ker(\beta)/\Im(\alpha)$ se restreint à zéro sur $U_\etale$ pour tout $U$
plat sur $\mathcal{X}$ ; ainsi, $\mathcal{E}$ se restreint à zéro sur
$U_\etale$ pour tout $U$ plat sur $\mathcal{X}$. Cela implique certainement
que le module quasi-cohérent $\mathcal{E}$ est nul ; on peut par exemple appliquer le
lemme \ref{stacks-cohomology-lemma-quasi-coherent-check-exact}
au morphisme $0 \to \mathcal{E}$.
\end{proof}





\section{Modules quasi-cohérents}
\label{stacks-cohomology-section-quasi-coherent}

\noindent
Nous avons vu que la catégorie des modules quasi-cohérents sur un champ
algébrique est équivalente à la catégorie des modules quasi-cohérents sur une
présentation ; voir
Faisceaux sur les champs, section
\ref{stacks-sheaves-section-quasi-coherent-algebraic-stacks}.
Ce fait est à la base de ce qui suit.

\begin{lemma}
\label{stacks-cohomology-lemma-adjoint}
Soit $\mathcal{X}$ un champ algébrique. Soit
$\textit{LQCoh}^{fbc}(\mathcal{O}_\mathcal{X})$
la catégorie des modules localement quasi-cohérents possédant la
propriété de changement de base plat ; voir
section \ref{stacks-cohomology-section-loc-qcoh-flat-base-change}.
Le foncteur d'inclusion
$i : \QCoh(\mathcal{O}_\mathcal{X}) \to
\textit{LQCoh}^{fbc}(\mathcal{O}_\mathcal{X})$
admet un adjoint à droite
$$
Q : \textit{LQCoh}^{fbc}(\mathcal{O}_\mathcal{X}) \to
\QCoh(\mathcal{O}_\mathcal{X})
$$
tel que $Q \circ i$ soit le foncteur identité.
\end{lemma}

\begin{proof}
Choisissons un schéma $U$ et un morphisme lisse surjectif $f : U \to \mathcal{X}$.
Posons $R = U \times_\mathcal{X} U$ ; nous obtenons ainsi un groupoïde lisse
$(U, R, s, t, c)$ en espaces algébriques tel que
$\mathcal{X} = [U/R]$ ; voir
Champs algébriques, lemme \ref{algebraic-lemma-stack-presentation}.
Nous pouvons remplacer, et remplaçons, $\mathcal{X}$ par $[U/R]$. D'après
Faisceaux sur les champs, proposition \ref{stacks-sheaves-proposition-quasi-coherent}
il existe une équivalence
$$
q_1 :
\QCoh(U, R, s, t, c)
\longrightarrow
\QCoh(\mathcal{O}_\mathcal{X})
$$
Construisons un foncteur
$$
q_2 :
\textit{LQCoh}^{fbc}(\mathcal{O}_\mathcal{X})
\longrightarrow
\QCoh(U, R, s, t, c)
$$
par la règle suivante : si $\mathcal{F}$ est un objet de
$\textit{LQCoh}^{fbc}(\mathcal{O}_\mathcal{X})$, nous posons
$$
q_2(\mathcal{F}) = (f^*\mathcal{F}|_{U_\etale}, \alpha)
$$
où $\alpha$ est l'isomorphisme
$$
t_{small}^*(f^*\mathcal{F}|_{U_\etale})
\to
t^*f^*\mathcal{F}|_{R_\etale} \to
s^*f^*\mathcal{F}|_{R_\etale} \to
s_{small}^*(f^*\mathcal{F}|_{U_\etale})
$$
où les deux morphismes extrêmes sont les morphismes de comparaison. Notons que
$q_2(\mathcal{F})$ est quasi-cohérent précisément parce que $\mathcal{F}$ est
localement quasi-cohérent, et que nous avons utilisé (et eu besoin de)
la propriété de changement de base plat dans la construction de
la donnée de descente $\alpha$. Nous omettons la
vérification de la condition de cocycle (voir
Groupoïdes dans les espaces, définition
\ref{spaces-groupoids-definition-groupoid-module}).
En examinant la démonstration de
Faisceaux sur les champs, proposition \ref{stacks-sheaves-proposition-quasi-coherent}
nous voyons que $q_2 \circ i$ est un quasi-inverse de $q_1$.
Définissons $Q = q_1 \circ q_2$.
Soit $\mathcal{F}$ un objet de
$\textit{LQCoh}^{fbc}(\mathcal{O}_\mathcal{X})$ et
soit $\mathcal{G}$ un objet de $\QCoh(\mathcal{O}_\mathcal{X})$.
Nous avons
\begin{align*}
\Mor_{\textit{LQCoh}^{fbc}(\mathcal{O}_\mathcal{X})}
(i(\mathcal{G}), \mathcal{F})
& =
\Mor_{\QCoh(U, R, s, t, c)}(q_2(i(\mathcal{G})), q_2(\mathcal{F})) \\
& =
\Mor_{\QCoh(\mathcal{O}_\mathcal{X})}(\mathcal{G}, Q(\mathcal{F}))
\end{align*}
où la première égalité résulte de
Faisceaux sur les champs, lemme \ref{stacks-sheaves-lemma-map-from-quasi-coherent}
et la deuxième vaut parce que $q_1 \circ i$ et $q_2$ sont des équivalences de catégories
quasi-inverses. L'assertion $Q \circ i \cong \text{id}$
est une conséquence formelle du fait que $i$ est pleinement fidèle.
\end{proof}

\begin{lemma}
\label{stacks-cohomology-lemma-adjoint-kernel-parasitic}
Soit $\mathcal{X}$ un champ algébrique.
Soit $Q : \textit{LQCoh}^{fbc}(\mathcal{O}_\mathcal{X}) \to
\QCoh(\mathcal{O}_\mathcal{X})$
le foncteur construit dans le lemme \ref{stacks-cohomology-lemma-adjoint}.
\begin{enumerate}
\item Le noyau de $Q$ est exactement la collection des objets parasites
de $\textit{LQCoh}^{fbc}(\mathcal{O}_\mathcal{X})$.
\item Pour tout objet $\mathcal{F}$
de $\textit{LQCoh}^{fbc}(\mathcal{O}_\mathcal{X})$, le noyau et le
conoyau du
morphisme d'adjonction $Q(\mathcal{F}) \to \mathcal{F}$ sont parasites.
\item Le foncteur $Q$ est exact et commute à toutes les limites projectives et inductives.
\end{enumerate}
\end{lemma}

\begin{proof}
Écrivons $\mathcal{X} = [U/R]$ comme dans la démonstration du lemme \ref{stacks-cohomology-lemma-adjoint}.
Soit $\mathcal{F}$ un objet de
$\textit{LQCoh}^{fbc}(\mathcal{O}_\mathcal{X})$.
Il ressort de la démonstration du lemme \ref{stacks-cohomology-lemma-adjoint}
que $\mathcal{F}$ appartient au noyau de $Q$ si et seulement si
$\mathcal{F}|_{U_\etale} = 0$.
En particulier, si $\mathcal{F}$ est parasite, alors $\mathcal{F}$ appartient au
noyau. Soit ensuite $x : V \to \mathcal{X}$ un morphisme plat, où
$V$ est un schéma. Posons $W = V \times_\mathcal{X} U$ et considérons le diagramme
$$
\xymatrix{
W \ar[d]_p \ar[r]_q & V \ar[d] \\
U \ar[r] & \mathcal{X}
}
$$
Notons que la projection $p : W \to U$ est plate et que la projection
$q : W \to V$ est lisse et surjective. Il s'ensuit que $q_{small}^*$
est un foncteur fidèle sur les modules quasi-cohérents. Par hypothèse, $\mathcal{F}$
possède la propriété de changement de base plat, de sorte que nous obtenons
$p_{small}^*\mathcal{F}|_{U_\etale} \cong
q_{small}^*\mathcal{F}|_{V_\etale}$. Ainsi, si $\mathcal{F}$
appartient au noyau de $Q$, alors $\mathcal{F}|_{V_\etale} = 0$,
ce qui achève la démonstration de (1).

\medskip\noindent
La partie (2) résulte de la discussion précédente et du fait
que le morphisme $Q(\mathcal{F}) \to \mathcal{F}$ devient un isomorphisme après
restriction à $U_\etale$.

\medskip\noindent
Pour démontrer la partie (3), notons que $Q$ est exact à gauche comme adjoint à droite.
Soit $0 \to \mathcal{F} \to \mathcal{G} \to \mathcal{H} \to 0$
une suite exacte courte dans $\textit{LQCoh}^{fbc}(\mathcal{O}_\mathcal{X})$.
Considérons le diagramme commutatif suivant
$$
\xymatrix{
0 \ar[r] &
Q(\mathcal{F}) \ar[r] \ar[d]_a &
Q(\mathcal{G}) \ar[r] \ar[d]_b &
Q(\mathcal{H}) \ar[r] \ar[d]_c & 0 \\
0 \ar[r] &
\mathcal{F} \ar[r] &
\mathcal{G} \ar[r] &
\mathcal{H} \ar[r] & 0
}
$$
Puisque les noyaux et conoyaux de $a$, $b$ et $c$ sont parasites d'après
la partie (2), et puisque la ligne inférieure est une suite exacte courte, nous voyons que
la ligne supérieure, considérée comme complexe de $\mathcal{O}_\mathcal{X}$-modules, a des faisceaux de
cohomologie parasites (nous omettons les détails ; on utilise ici que la catégorie des
modules parasites est une sous-catégorie de Serre de la catégorie de tous les modules).
Par l'exactitude à gauche de $Q$, seule l'exactitude en $Q(\mathcal{H})$
reste à établir. Or le conoyau $\mathcal{Q}$ de
$Q(\mathcal{G}) \to Q(\mathcal{H}))$
peut être calculé soit dans $\textit{Mod}(\mathcal{O}_\mathcal{X})$,
soit dans $\QCoh(\mathcal{O}_\mathcal{X})$, avec le même résultat, car
le foncteur d'inclusion $\QCoh(\mathcal{O}_\mathcal{X}) \to
\textit{LQCoh}^{fbc}(\mathcal{O}_\mathcal{X})$
est un adjoint à gauche, donc est exact à droite. Ainsi, $\mathcal{Q} = Q(\mathcal{Q})$
est à la fois quasi-cohérent et parasite, donc $0$ d'après la partie (1), comme souhaité.

\medskip\noindent
Comme adjoint à droite, $Q$ commute à toutes les limites projectives. Puisque $Q$ est exact, pour
montrer que $Q$ commute à toutes les limites inductives, il suffit de montrer que $Q$
commute aux sommes directes ; voir
Catégories, lemme \ref{categories-lemma-colimits-coproducts-coequalizers}.
Soit $\mathcal{F}_i$, $i \in I$, une famille d'objets de
$\textit{LQCoh}^{fbc}(\mathcal{O}_\mathcal{X})$. Pour voir que
$Q(\bigoplus \mathcal{F}_i)$ est égal à $\bigoplus Q(\mathcal{F}_i)$,
considérons la construction de $Q$ dans la démonstration du
lemme \ref{stacks-cohomology-lemma-adjoint}. Elle utilise une présentation $\mathcal{X} = [U/R]$,
où $U$ est un schéma. Alors $Q(\mathcal{F})$
se calcule en prenant d'abord la paire
$(\mathcal{F}|_{U_\etale}, \alpha)$ dans $\QCoh(U, R, s, t, c)$,
puis en utilisant l'équivalence
$\QCoh(U, R, s, t, c) \cong \QCoh(\mathcal{O}_\mathcal{X})$.
Puisque le foncteur de restriction $\textit{Mod}(\mathcal{O}_\mathcal{X}) \to
\textit{Mod}(\mathcal{O}_{U_\etale})$,
$\mathcal{F} \mapsto \mathcal{F}|_{U_\etale}$ commute aux
sommes directes, l'égalité voulue est claire.
\end{proof}

\begin{lemma}
\label{stacks-cohomology-lemma-flat-morphism-and-Q}
Soit $f : \mathcal{X} \to \mathcal{Y}$ un morphisme plat de champs algébriques.
Alors $Q_\mathcal{X} \circ f^* = f^* \circ Q_\mathcal{Y}$, où
$Q_\mathcal{X}$ et $Q_\mathcal{Y}$ sont comme dans le lemme \ref{stacks-cohomology-lemma-adjoint}.
\end{lemma}

\begin{proof}
Observons que $f^*$ préserve à la fois $\QCoh$ et $\textit{LQCoh}^{fbc}$ ;
voir Faisceaux sur les champs, lemme
\ref{stacks-sheaves-lemma-pullback-quasi-coherent} et
proposition \ref{stacks-cohomology-proposition-loc-qcoh-flat-base-change}.
Si $\mathcal{F}$ appartient à $\textit{LQCoh}^{fbc}(\mathcal{O}_\mathcal{Y})$,
alors $Q_\mathcal{Y}(\mathcal{F}) \to \mathcal{F}$
a un noyau et un conoyau parasites
d'après le lemme \ref{stacks-cohomology-lemma-adjoint-kernel-parasitic}.
Comme $f$ est plat, nous obtenons que $f^*Q_\mathcal{Y}(\mathcal{F}) \to f^*\mathcal{F}$
a un noyau et un conoyau parasites d'après le
lemme \ref{stacks-cohomology-lemma-parasitic}.
Ainsi, le morphisme induit
$f^*Q_\mathcal{Y}(\mathcal{F}) \to Q_\mathcal{X}(f^*\mathcal{F})$
a un noyau et un conoyau parasites et est donc un isomorphisme,
par exemple d'après le
lemme \ref{stacks-cohomology-lemma-exact-sequence-quasi-coherent-parasitic-cohomology}.
\end{proof}

\begin{lemma}
\label{stacks-cohomology-lemma-flat-object-and-Q}
Soit $\mathcal{X}$ un champ algébrique. Soit $x$ un objet de $\mathcal{X}$
au-dessus du schéma $U$ tel que $x : U \to \mathcal{X}$ soit plat.
Alors, pour $\mathcal{F}$ dans $\QCoh^{fbc}(\mathcal{O}_\mathcal{X})$,
nous avons $Q(\mathcal{F})|_{U_\etale} = \mathcal{F}|_{U_\etale}$.
\end{lemma}

\begin{proof}
C'est vrai parce que le noyau et le conoyau de $Q(\mathcal{F}) \to \mathcal{F}$
sont parasites ; voir le lemme \ref{stacks-cohomology-lemma-adjoint-kernel-parasitic}.
\end{proof}

\begin{remark}
\label{stacks-cohomology-remark-QCoh-abelian}
Soit $\mathcal{X}$ un champ algébrique. La catégorie
$\QCoh(\mathcal{O}_\mathcal{X})$ est abélienne, le foncteur d'inclusion
$\QCoh(\mathcal{O}_\mathcal{X}) \to \textit{Mod}(\mathcal{O}_\mathcal{X})$
est exact à droite, mais n'est pas exact en général ; voir Faisceaux sur les champs, lemme
\ref{stacks-sheaves-lemma-quasi-coherent-algebraic-stack}.
Nous pouvons utiliser le foncteur $Q$ des
lemmes \ref{stacks-cohomology-lemma-adjoint} et \ref{stacks-cohomology-lemma-adjoint-kernel-parasitic}
pour comprendre ceci. Plus précisément, soit $\varphi : \mathcal{F} \to \mathcal{G}$
un morphisme de $\mathcal{O}_\mathcal{X}$-modules quasi-cohérents. Alors
\begin{enumerate}
\item le conoyau $\Coker(\varphi)$ calculé dans
$\textit{Mod}(\mathcal{O}_\mathcal{X})$ est quasi-cohérent et
est le conoyau de $\varphi$ dans $\QCoh(\mathcal{O}_\mathcal{X})$ ;
\item l'image $\Im(\varphi)$ calculée dans
$\textit{Mod}(\mathcal{O}_\mathcal{X})$ est quasi-cohérente et
est l'image de $\varphi$ dans $\QCoh(\mathcal{O}_\mathcal{X})$ ; et
\item le noyau $\Ker(\varphi)$ calculé dans
$\textit{Mod}(\mathcal{O}_\mathcal{X})$
appartient à $\textit{LQCoh}^{fbc}(\mathcal{O}_\mathcal{X})$
d'après la proposition \ref{stacks-cohomology-proposition-loc-qcoh-flat-base-change}, et
$Q(\Ker(\varphi))$ est le noyau dans $\QCoh(\mathcal{O}_\mathcal{X})$.
\end{enumerate}
Cela résulte des références indiquées.
\end{remark}

\begin{remark}
\label{stacks-cohomology-remark-QCoh-tensor}
Soit $\mathcal{X}$ un champ algébrique. Étant donnés deux
$\mathcal{O}_\mathcal{X}$-modules quasi-cohérents $\mathcal{F}$ et $\mathcal{G}$,
le module produit tensoriel
$\mathcal{F} \otimes_{\mathcal{O}_\mathcal{X}} \mathcal{G}$
est quasi-cohérent ; voir Faisceaux sur les champs, lemme
\ref{stacks-sheaves-lemma-quasi-coherent-algebraic-stack}, partie (5).
De même, étant donnés deux modules localement quasi-cohérents possédant
la propriété de changement de base plat, leur produit tensoriel possède la
même propriété ; voir la proposition \ref{stacks-cohomology-proposition-loc-qcoh-flat-base-change}.
Ainsi, les foncteurs d'inclusion
$$
\QCoh(\mathcal{O}_\mathcal{X}) \to
\textit{LQCoh}^{fbc}(\mathcal{O}_\mathcal{X}) \to
\textit{Mod}(\mathcal{O}_\mathcal{X})
$$
sont des foncteurs de catégories monoïdales symétriques. Plus intéressant encore,
le foncteur
$$
Q :
\textit{LQCoh}^{fbc}(\mathcal{O}_\mathcal{X})
\longrightarrow
\QCoh(\mathcal{O}_\mathcal{X})
$$
est lui aussi un foncteur de catégories monoïdales symétriques. En effet, étant donnés
$\mathcal{F}$ et $\mathcal{G}$ dans
$\textit{LQCoh}^{fbc}(\mathcal{O}_\mathcal{X})$, nous obtenons
$$
\xymatrix{
Q(\mathcal{F})
\otimes_{\mathcal{O}_\mathcal{X}}
Q(\mathcal{G}) \ar[rr] \ar[rd] & &
\mathcal{F}
\otimes_{\mathcal{O}_\mathcal{X}}
\mathcal{G} \\
&
Q(\mathcal{F} \otimes_{\mathcal{O}_\mathcal{X}} \mathcal{G}) \ar[ru]
}
$$
où la flèche sud-ouest provient de la propriété universelle
de la flèche nord-ouest (et du fait déjà mentionné que
l'objet du coin supérieur gauche est quasi-cohérent).
Si nous restreignons ce diagramme à $U_\etale$ pour $U \to \mathcal{X}$
plat, alors les trois flèches deviennent des isomorphismes (voir les
lemmes \ref{stacks-cohomology-lemma-adjoint} et \ref{stacks-cohomology-lemma-adjoint-kernel-parasitic}
et la définition \ref{stacks-cohomology-definition-parasitic}).
Ainsi, $Q(\mathcal{F}) \otimes_{\mathcal{O}_\mathcal{X}}
Q(\mathcal{G}) \to
Q(\mathcal{F} \otimes_{\mathcal{O}_\mathcal{X}} \mathcal{G})$
est un isomorphisme ; voir par exemple le
lemme \ref{stacks-cohomology-lemma-quasi-coherent-check-exact}.
\end{remark}

\begin{remark}
\label{stacks-cohomology-remark-bousfield-colocalization}
Soit $\mathcal{X}$ un champ algébrique. Soit
$\textit{Parasitic}(\mathcal{O}_\mathcal{X}) \subset
\textit{Mod}(\mathcal{O}_\mathcal{X})$ la
sous-catégorie pleine formée des modules parasites. Les résultats des
lemmes \ref{stacks-cohomology-lemma-adjoint} et \ref{stacks-cohomology-lemma-adjoint-kernel-parasitic}
impliquent que
$$
\QCoh(\mathcal{O}_\mathcal{X}) =
\textit{LQCoh}^{fbc}(\mathcal{O}_\mathcal{X}) /
\textit{Parasitic}(\mathcal{O}_\mathcal{X})
\cap \textit{LQCoh}^{fbc}(\mathcal{O}_\mathcal{X})
$$
autrement dit : la catégorie des modules quasi-cohérents est la catégorie
des modules localement quasi-cohérents possédant la propriété de changement de base plat,
quotientée par la sous-catégorie de Serre formée des objets parasites.
Voir Homologie, lemme \ref{homology-lemma-serre-subcategory-is-kernel}.
L'existence du foncteur d'inclusion
$i : \QCoh(\mathcal{O}_\mathcal{X}) \to
\textit{LQCoh}^{fbc}(\mathcal{O}_\mathcal{X})$
qui est adjoint à gauche du foncteur quotient est un trait essentiel de
la situation. Dans Catégories dérivées de champs, section
\ref{stacks-perfect-section-derived}, et en particulier dans le
lemme \ref{stacks-perfect-lemma-bousfield-colocalization}
nous démontrons qu'un résultat analogue vaut au niveau des catégories dérivées.
\end{remark}

\begin{lemma}
\label{stacks-cohomology-lemma-internal-hom-fp-into-qcoh}
Soit $\mathcal{X}$ un champ algébrique. Soit $\mathcal{F}$ un
$\mathcal{O}_\mathcal{X}$-module de présentation finie et soit
$\mathcal{G}$ un $\mathcal{O}_\mathcal{X}$-module quasi-cohérent.
Les Hom internes
$\SheafHom_{\mathcal{O}_\mathcal{X}}(\mathcal{F}, \mathcal{G})$
calculés dans
$\textit{Mod}(\mathcal{X}_\etale, \mathcal{O}_\mathcal{X})$ ou dans
$\textit{Mod}(\mathcal{O}_\mathcal{X})$ coïncident, et leur valeur commune
est un objet de $\textit{LQCoh}^{fbc}(\mathcal{O}_\mathcal{X})$.
Le module quasi-cohérent
$
hom(\mathcal{F}, \mathcal{G}) =
Q(\SheafHom_{\mathcal{O}_\mathcal{X}}(\mathcal{F}, \mathcal{G}))
$
possède la propriété universelle suivante
$$
\Hom_\mathcal{X}(\mathcal{H}, hom(\mathcal{F}, \mathcal{G})) =
\Hom_\mathcal{X}(\mathcal{H} \otimes_{\mathcal{O}_\mathcal{X}} \mathcal{F},
\mathcal{G})
$$
pour $\mathcal{H}$ dans $\QCoh(\mathcal{O}_\mathcal{X})$.
\end{lemma}

\begin{proof}
La construction de
$\SheafHom_{\mathcal{O}_\mathcal{X}}(\mathcal{F}, \mathcal{G})$ dans
Modules sur les sites, section \ref{sites-modules-section-internal-hom}
ne dépend que de $\mathcal{F}$ et $\mathcal{G}$ comme préfaisceaux
de modules ; le résultat $\SheafHom$ est un faisceau pour la topologie
fppf parce que $\mathcal{F}$ et $\mathcal{G}$ sont supposés être des
faisceaux pour la topologie fppf ; voir
Modules sur les sites, lemme \ref{sites-modules-lemma-internal-hom}.
D'après Faisceaux sur les champs, lemme \ref{stacks-sheaves-lemma-lqc-colimits},
nous voyons que $\SheafHom_{\mathcal{O}_\mathcal{X}}(\mathcal{F}, \mathcal{G})$
est localement quasi-cohérent.
D'après le lemme \ref{stacks-cohomology-lemma-check-flat-comparison-on-etale-covering},
nous voyons que $\SheafHom_{\mathcal{O}_\mathcal{X}}(\mathcal{F}, \mathcal{G})$
possède la propriété de changement de base plat. Ainsi,
$\SheafHom_{\mathcal{O}_\mathcal{X}}(\mathcal{F}, \mathcal{G})$
est un objet de $\textit{LQCoh}^{fbc}(\mathcal{O}_\mathcal{X})$
et il est légitime de lui appliquer le foncteur $Q$ du lemme \ref{stacks-cohomology-lemma-adjoint}.
Par la propriété universelle de $Q$, nous avons
$$
\Hom_\mathcal{X}(\mathcal{H},
Q(\SheafHom_{\mathcal{O}_\mathcal{X}}(\mathcal{F}, \mathcal{G}))) =
\Hom_\mathcal{X}(\mathcal{H},
\SheafHom_{\mathcal{O}_\mathcal{X}}(\mathcal{F}, \mathcal{G}))
$$
pour $\mathcal{H}$ quasi-cohérent ; la formule affichée
du lemme résulte donc de Modules sur les sites, lemme
\ref{sites-modules-lemma-internal-hom-adjoint-tensor}.
\end{proof}

\begin{lemma}
\label{stacks-cohomology-lemma-flat-morphism-and-hom}
Soit $f : \mathcal{X} \to \mathcal{Y}$ un morphisme plat de champs algébriques.
Soit $\mathcal{F}$ un
$\mathcal{O}_\mathcal{Y}$-module de présentation finie et soit
$\mathcal{G}$ un $\mathcal{O}_\mathcal{Y}$-module quasi-cohérent.
Alors $f^*hom(\mathcal{F}, \mathcal{G}) = hom(f^*\mathcal{F}, f^*\mathcal{G})$
avec les notations du lemme \ref{stacks-cohomology-lemma-internal-hom-fp-into-qcoh}.
\end{lemma}

\begin{proof}
Nous avons $f^*\SheafHom_{\mathcal{O}_\mathcal{Y}}(\mathcal{F}, \mathcal{G}) =
\SheafHom_{\mathcal{O}_\mathcal{X}}(f^*\mathcal{F}, f^*\mathcal{G})$
d'après Modules sur les sites, lemme \ref{sites-modules-lemma-pullback-internal-hom}.
(Observons que ce n'est pas à cette étape qu'est utilisée la platitude de $f$,
car le morphisme de topos annelés associé à $f$ est toujours plat ; voir
Faisceaux sur les champs, remarque \ref{stacks-sheaves-remark-flat}.)
Appliquons alors le lemme \ref{stacks-cohomology-lemma-flat-morphism-and-Q} (et c'est ici que nous utilisons
la platitude de $f$).
\end{proof}








\section{Image directe des modules quasi-cohérents}
\label{stacks-cohomology-section-pushforward-quasi-coherent}

\noindent
Soit $f : \mathcal{X} \to \mathcal{Y}$ un morphisme de champs algébriques.
Considérons l'image directe
$$
f_* :
\textit{Mod}(\mathcal{O}_\mathcal{X})
\longrightarrow
\textit{Mod}(\mathcal{O}_\mathcal{Y})
$$
Il s'avère que ce foncteur ne préserve presque jamais les sous-catégories
de faisceaux quasi-cohérents. Considérons par exemple le morphisme de schémas
$$
j : X = \mathbf{A}^2_k \setminus \{0\} \longrightarrow \mathbf{A}^2_k = Y.
$$
Il lui est associé le morphisme correspondant de champs algébriques
$$
f = j_{big} : \mathcal{X} = (\Sch/X)_{fppf} \to
(\Sch/Y)_{fppf} = \mathcal{Y}
$$
L'image directe $f_*\mathcal{O}_\mathcal{X}$ du faisceau structural a pour
sections globales $k[x, y]$. Ainsi, si $f_*\mathcal{O}_\mathcal{X}$ était
quasi-cohérent sur $\mathcal{Y}$, nous aurions
$f_*\mathcal{O}_\mathcal{X} = \mathcal{O}_\mathcal{Y}$. Cependant,
considérons $T = \Spec(k) \to \mathbf{A}^2_k = Y$ envoyé sur $0$.
Alors $\Gamma(T, f_*\mathcal{O}_\mathcal{X}) = 0$ parce que
$X \times_Y T = \emptyset$, tandis que $\Gamma(T, \mathcal{O}_\mathcal{Y}) = k$.
En revanche, pour tout morphisme plat $T \to Y$, nous avons l'égalité
$\Gamma(T, f_*\mathcal{O}_\mathcal{X}) = \Gamma(T, \mathcal{O}_\mathcal{Y})$
qui résulte de
Cohomologie des schémas, lemme \ref{coherent-lemma-flat-base-change-cohomology}
en utilisant que $j$ est quasi-compact et quasi-séparé.

\medskip\noindent
Soit $f : \mathcal{X} \to \mathcal{Y}$ un morphisme quasi-compact et
quasi-séparé de champs algébriques. Nous contournons le problème
mentionné ci-dessus à l'aide des trois observations suivantes :
\begin{enumerate}
\item $f_*$ préserve bien les modules localement quasi-cohérents
(lemme \ref{stacks-cohomology-lemma-pushforward-locally-quasi-coherent}) ;
\item $f_*$ transforme un faisceau quasi-cohérent en un faisceau localement quasi-cohérent
dont les morphismes de comparaison plats sont des isomorphismes
(lemme \ref{stacks-cohomology-lemma-flat-comparison}) ; et
\item les $\mathcal{O}_\mathcal{Y}$-modules localement quasi-cohérents
possédant la propriété de changement de base plat donnent des modules quasi-cohérents
sur une présentation de $\mathcal{Y}$, et donc des modules quasi-cohérents
sur $\mathcal{Y}$ ; voir
Faisceaux sur les champs, section
\ref{stacks-sheaves-section-quasi-coherent-algebraic-stacks}.
\end{enumerate}
Nous obtenons ainsi un foncteur
$$
f_{\QCoh, *} :
\QCoh(\mathcal{O}_\mathcal{X})
\longrightarrow
\QCoh(\mathcal{O}_\mathcal{Y})
$$
qui est adjoint à droite de
$f^* : \QCoh(\mathcal{O}_\mathcal{Y}) \to
\QCoh(\mathcal{O}_\mathcal{X})$
et tel que, de plus,
$$
\Gamma(y, f_*\mathcal{F}) = \Gamma(y, f_{\QCoh, *}\mathcal{F})
$$
pour tout $y \in \Ob(\mathcal{Y})$ tel que le $1$-morphisme
associé $y : V \to \mathcal{Y}$ soit plat ; voir
lemme \ref{stacks-cohomology-lemma-direct-image-quasi-coherent-over-flat}.
De plus, une construction semblable produira des foncteurs
$R^if_{\QCoh, *}$.
Cependant, ces résultats ne suffiront pas à produire un foncteur
image directe totale (de complexes à faisceaux de cohomologie
quasi-cohérents).

\begin{proposition}
\label{stacks-cohomology-proposition-direct-image-quasi-coherent}
Soit $f : \mathcal{X} \to \mathcal{Y}$ un morphisme quasi-compact et quasi-séparé
de champs algébriques. Le foncteur
$f^* : \QCoh(\mathcal{O}_\mathcal{Y}) \to
\QCoh(\mathcal{O}_\mathcal{X})$
admet un adjoint à droite
$$
f_{\QCoh, *} :
\QCoh(\mathcal{O}_\mathcal{X})
\longrightarrow
\QCoh(\mathcal{O}_\mathcal{Y})
$$
qui peut être défini comme le composé
$$
\QCoh(\mathcal{O}_\mathcal{X}) \to
\textit{LQCoh}^{fbc}(\mathcal{O}_\mathcal{X})
\xrightarrow{f_*} \textit{LQCoh}^{fbc}(\mathcal{O}_\mathcal{Y})
\xrightarrow{Q} \QCoh(\mathcal{O}_\mathcal{Y})
$$
où les foncteurs $f_*$ et $Q$ sont ceux de la
proposition \ref{stacks-cohomology-proposition-loc-qcoh-flat-base-change}
et du
lemme \ref{stacks-cohomology-lemma-adjoint}.
De plus, si nous définissons $R^if_{\QCoh, *}$ comme le composé
$$
\QCoh(\mathcal{O}_\mathcal{X}) \to
\textit{LQCoh}^{fbc}(\mathcal{O}_\mathcal{X})
\xrightarrow{R^if_*} \textit{LQCoh}^{fbc}(\mathcal{O}_\mathcal{Y})
\xrightarrow{Q} \QCoh(\mathcal{O}_\mathcal{Y})
$$
alors la suite de foncteurs $\{R^if_{\QCoh, *}\}_{i \geq 0}$
forme un $\delta$-foncteur cohomologique.
\end{proposition}

\begin{proof}
C'est une combinaison des résultats mentionnés dans l'énoncé.
L'adjonction se démontre comme suit. Soit $\mathcal{F}$
un $\mathcal{O}_\mathcal{X}$-module quasi-cohérent et soit
$\mathcal{G}$ un $\mathcal{O}_\mathcal{Y}$-module quasi-cohérent.
Alors nous avons
\begin{align*}
\Mor_{\QCoh(\mathcal{O}_\mathcal{X})}(f^*\mathcal{G}, \mathcal{F})
& =
\Mor_{\textit{LQCoh}^{fbc}(\mathcal{O}_\mathcal{Y})}
(\mathcal{G}, f_*\mathcal{F}) \\
& =
\Mor_{\QCoh(\mathcal{O}_\mathcal{Y})}(\mathcal{G}, Q(f_*\mathcal{F}))
\\
& =
\Mor_{\QCoh(\mathcal{O}_\mathcal{Y})}(\mathcal{G},
f_{\QCoh, *}\mathcal{F})
\end{align*}
la première égalité résultant de l'adjonction de $f_*$ et $f^*$ (pour des faisceaux
de modules arbitraires). D'après la
proposition \ref{stacks-cohomology-proposition-loc-qcoh-flat-base-change}
nous voyons que $f_*\mathcal{F}$ est un objet de
$\textit{LQCoh}^{fbc}(\mathcal{O}_\mathcal{Y})$
(et peut être calculé pour la topologie fppf comme pour la topologie étale), et nous
obtenons la deuxième égalité d'après le lemme \ref{stacks-cohomology-lemma-adjoint}. La troisième
égalité est la définition de $f_{\QCoh, *}$.

\medskip\noindent
Pour voir que $\{R^if_{\QCoh, *}\}_{i \geq 0}$ est un
$\delta$-foncteur cohomologique au sens de
Homologie, définition \ref{homology-definition-cohomological-delta-functor}
soit
$$
0 \to \mathcal{F}_1 \to \mathcal{F}_2 \to \mathcal{F}_3 \to 0
$$
une suite exacte courte de $\QCoh(\mathcal{O}_\mathcal{X})$.
Cette suite n'est pas nécessairement exacte dans
$\textit{Mod}(\mathcal{O}_\mathcal{X})$, mais nous savons qu'elle l'est
à des modules parasites près ; voir
lemme \ref{stacks-cohomology-lemma-exact-sequence-quasi-coherent-parasitic-cohomology}.
Nous pouvons donc décomposer la suite en suites exactes courtes
$$
\begin{matrix}
0 \to \mathcal{P}_1 \to \mathcal{F}_1 \to \mathcal{I}_2 \to 0 \\
0 \to \mathcal{I}_2 \to \mathcal{F}_2 \to \mathcal{Q}_2 \to 0 \\
0 \to \mathcal{P}_2 \to \mathcal{Q}_2 \to \mathcal{I}_3 \to 0 \\
0 \to \mathcal{I}_3 \to \mathcal{F}_3 \to \mathcal{P}_3 \to 0
\end{matrix}
$$
de $\textit{Mod}(\mathcal{O}_\mathcal{X})$, les $\mathcal{P}_i$ étant parasites.
Notons que chacun des faisceaux
$\mathcal{P}_j$, $\mathcal{I}_j$, $\mathcal{Q}_j$ est un objet de
$\textit{LQCoh}^{fbc}(\mathcal{O}_\mathcal{X})$ ; voir
proposition \ref{stacks-cohomology-proposition-loc-qcoh-flat-base-change}.
En appliquant $R^if_*$, nous obtenons des suites exactes longues
$$
\begin{matrix}
0 \to f_*\mathcal{P}_1 \to f_*\mathcal{F}_1 \to f_*\mathcal{I}_2 \to
R^1f_*\mathcal{P}_1 \to \ldots \\
0 \to f_*\mathcal{I}_2 \to f_*\mathcal{F}_2 \to f_*\mathcal{Q}_2 \to
R^1f_*\mathcal{I}_2 \to \ldots \\
0 \to f_*\mathcal{P}_2 \to f_*\mathcal{Q}_2 \to f_*\mathcal{I}_3 \to
R^1f_*\mathcal{P}_2 \to \ldots \\
0 \to f_*\mathcal{I}_3 \to f_*\mathcal{F}_3 \to f_*\mathcal{P}_3 \to
R^1f_*\mathcal{I}_3 \to \ldots
\end{matrix}
$$
où tous les termes sont des objets de
$\textit{LQCoh}^{fbc}(\mathcal{O}_\mathcal{Y})$ d'après la
proposition \ref{stacks-cohomology-proposition-loc-qcoh-flat-base-change}.
D'après le
lemme \ref{stacks-cohomology-lemma-pushforward-parasitic},
les faisceaux $R^if_*\mathcal{P}_j$ sont parasites et s'annulent donc après application
du foncteur $Q$ ; voir le
lemme \ref{stacks-cohomology-lemma-adjoint-kernel-parasitic}.
Puisque $Q$ est exact, les morphismes
$$
Q(R^if_*\mathcal{F}_3)
\cong
Q(R^if_*\mathcal{I}_3)
\cong
Q(R^if_*\mathcal{Q}_2)
\rightarrow
Q(R^{i + 1}f_*\mathcal{I}_2)
\cong
Q(R^{i + 1}f_*\mathcal{F}_1)
$$
peuvent servir de morphisme de connexion qui fait de la famille de foncteurs
$\{R^if_{\QCoh, *}\}_{i \geq 0}$
un $\delta$-foncteur cohomologique.
\end{proof}

\begin{lemma}
\label{stacks-cohomology-lemma-direct-image-quasi-coherent-over-flat}
Soit $f : \mathcal{X} \to \mathcal{Y}$ un morphisme quasi-compact et quasi-séparé
de champs algébriques. Soit $y : V \to \mathcal{Y}$ dans $\Ob(\mathcal{Y})$,
avec $y$ plat. Soit $\mathcal{F}$ dans
$\QCoh(\mathcal{O}_\mathcal{X})$.
Alors $(f_*\mathcal{F})(y) = (f_{\QCoh, *}\mathcal{F})(y)$
et $(R^if_*\mathcal{F})(y) = (R^if_{\QCoh, *}\mathcal{F})(y)$
pour tout $i \in \mathbf{Z}$.
\end{lemma}

\begin{proof}
Cela résulte de la construction des foncteurs $R^if_{\QCoh, *}$ dans la
proposition \ref{stacks-cohomology-proposition-direct-image-quasi-coherent},
de la définition des modules parasites dans définition \ref{stacks-cohomology-definition-parasitic}
et du lemme \ref{stacks-cohomology-lemma-adjoint-kernel-parasitic}, partie (2).
\end{proof}

\begin{remark}
\label{stacks-cohomology-remark-direct-image-quasi-coherent-tensor}
Soit $f : \mathcal{X} \to \mathcal{Y}$ un morphisme quasi-compact et quasi-séparé
de champs algébriques. Soient $\mathcal{F}$ et $\mathcal{G}$ dans
$\QCoh(\mathcal{O}_\mathcal{X})$. Il existe alors un diagramme commutatif
canonique
$$
\xymatrix{
f_{\QCoh, *}\mathcal{F}
\otimes_{\mathcal{O}_\mathcal{Y}}
f_{\QCoh, *}\mathcal{G} \ar[r] \ar[d] &
f_*\mathcal{F}
\otimes_{\mathcal{O}_\mathcal{Y}}
f_*\mathcal{G} \ar[d]^c \\
f_{\QCoh, *}(\mathcal{F}
\otimes_{\mathcal{O}_\mathcal{X}}
\mathcal{G}) \ar[r] &
f_*(\mathcal{F}
\otimes_{\mathcal{O}_\mathcal{X}}
\mathcal{G})
}
$$
La flèche verticale $c$ de droite est le cup-produit relatif naïf
(en degré $0$) ;
voir Cohomologie sur les sites, section \ref{sites-cohomology-section-cup-product}.
La source et le but de $c$ appartiennent à
$\textit{LQCoh}^{fbc}(\mathcal{O}_\mathcal{X})$ ; voir la
proposition \ref{stacks-cohomology-proposition-loc-qcoh-flat-base-change}.
En appliquant $Q$ à $c$, nous obtenons la flèche verticale de gauche,
car $Q$ commute aux produits tensoriels ; voir la
remarque \ref{stacks-cohomology-remark-QCoh-tensor}.
Cette construction est fonctorielle en $\mathcal{F}$ et $\mathcal{G}$.
\end{remark}

\begin{lemma}
\label{stacks-cohomology-lemma-leray}
Soit $f : \mathcal{X} \to \mathcal{Y}$
un morphisme quasi-compact et quasi-séparé de champs algébriques.
Soit $\mathcal{F}$ un faisceau quasi-cohérent sur $\mathcal{X}$. Alors
il existe une suite spectrale dont la page $E_2$ est
$$
E_2^{p, q} = H^p(\mathcal{Y}, R^qf_{\QCoh, *}\mathcal{F})
$$
et qui converge vers $H^{p + q}(\mathcal{X}, \mathcal{F})$.
\end{lemma}

\begin{proof}
D'après Cohomologie sur les sites, lemme \ref{sites-cohomology-lemma-Leray},
la suite spectrale de Leray dont la page est
$$
E_2^{p, q} = H^p(\mathcal{Y}, R^qf_*\mathcal{F})
$$
converge vers $H^{p + q}(\mathcal{X}, \mathcal{F})$. Le noyau et le conoyau
du morphisme d'adjonction
$$
R^qf_{\QCoh, *}\mathcal{F} \longrightarrow R^qf_*\mathcal{F}
$$
sont des modules parasites sur $\mathcal{Y}$
(lemme \ref{stacks-cohomology-lemma-adjoint-kernel-parasitic})
et ont donc une cohomologie nulle
(lemme \ref{stacks-cohomology-lemma-pushforward-parasitic}).
Il en résulte formellement que
$H^p(\mathcal{Y}, R^qf_{\QCoh, *}\mathcal{F}) =
H^p(\mathcal{Y}, R^qf_*\mathcal{F})$, ce qui conclut.
\end{proof}

\begin{lemma}
\label{stacks-cohomology-lemma-relative-leray}
Soient $f : \mathcal{X} \to \mathcal{Y}$ et $g : \mathcal{Y} \to \mathcal{Z}$
des morphismes quasi-compacts et quasi-séparés de champs algébriques.
Soit $\mathcal{F}$ un faisceau quasi-cohérent sur $\mathcal{X}$. Alors
il existe une suite spectrale dont la page $E_2$ est
$$
E_2^{p, q} = R^pg_{\QCoh, *}(R^qf_{\QCoh, *}\mathcal{F})
$$
et qui converge vers $R^{p + q}(g \circ f)_{\QCoh, *}\mathcal{F}$.
\end{lemma}

\begin{proof}
D'après Cohomologie sur les sites, lemme \ref{sites-cohomology-lemma-relative-Leray},
la suite spectrale de Leray dont la page est
$$
E_2^{p, q} = R^pg_*(R^qf_*\mathcal{F})
$$
converge vers $R^{p + q}(g \circ f)_*\mathcal{F}$. D'après les résultats de la
proposition \ref{stacks-cohomology-proposition-loc-qcoh-flat-base-change}
les termes de cette suite spectrale sont tous des objets de
$\textit{LQCoh}^{fbc}(\mathcal{O}_\mathcal{Z})$. En appliquant le foncteur exact
$Q_\mathcal{Z} : \textit{LQCoh}^{fbc}(\mathcal{O}_\mathcal{Z}) \to
\QCoh(\mathcal{O}_\mathcal{Z})$, nous obtenons une suite spectrale dans
$\QCoh(\mathcal{O}_\mathcal{Z})$ convergeant vers
$R^{p + q}(g \circ f)_{\QCoh, *}\mathcal{F}$. Le résultat
en découle donc si nous pouvons montrer que
$$
Q_\mathcal{Z}(R^pg_*(R^qf_*\mathcal{F})) =
Q_\mathcal{Z}(R^pg_*(Q_\mathcal{X}(R^qf_*\mathcal{F}))
$$
Cela résulte du fait que le noyau et le conoyau du morphisme
$$
Q_\mathcal{X}(R^qf_*\mathcal{F}) \longrightarrow R^qf_*\mathcal{F}
$$
sont parasites (lemme \ref{stacks-cohomology-lemma-adjoint-kernel-parasitic}) et que
$R^pg_*$ transforme les modules parasites en modules parasites
(lemme \ref{stacks-cohomology-lemma-pushforward-parasitic}).
\end{proof}

\noindent
Pour terminer cette section, explicitons les suites spectrales
associées à un recouvrement lisse par un schéma. Comparer avec
Faisceaux sur les champs, sections \ref{stacks-sheaves-section-cohomology} et
\ref{stacks-sheaves-section-higher-direct-images}.

\begin{proposition}
\label{stacks-cohomology-proposition-smooth-covering-compute-cohomology}
Soit $f : \mathcal{U} \to \mathcal{X}$ un morphisme de champs algébriques.
Supposons que $f$ soit représentable par des espaces algébriques, surjectif, plat et
localement de présentation finie. Soit $\mathcal{F}$ un
$\mathcal{O}_\mathcal{X}$-module quasi-cohérent. Il existe alors une suite spectrale
$$
E_2^{p, q} = H^q(\mathcal{U}_p, f_p^*\mathcal{F})
\Rightarrow
H^{p + q}(\mathcal{X}, \mathcal{F})
$$
où $f_p$ est le morphisme
$\mathcal{U} \times_\mathcal{X} \ldots \times_\mathcal{X} \mathcal{U} \to
\mathcal{X}$ (avec $p + 1$ facteurs).
\end{proposition}

\begin{proof}
C'est un cas particulier de
Faisceaux sur les champs, proposition
\ref{stacks-sheaves-proposition-smooth-covering-compute-cohomology}.
\end{proof}

\begin{proposition}
\label{stacks-cohomology-proposition-smooth-covering-compute-direct-image}
Soient $f : \mathcal{U} \to \mathcal{X}$ et $g : \mathcal{X} \to \mathcal{Y}$
des morphismes composables de champs algébriques.
Supposons que
\begin{enumerate}
\item $f$ soit représentable par des espaces algébriques, surjectif,
plat, localement de présentation finie, quasi-compact et quasi-séparé, et
\item $g$ soit quasi-compact et quasi-séparé.
\end{enumerate}
Si $\mathcal{F}$ appartient à $\QCoh(\mathcal{O}_\mathcal{X})$, alors
il existe une suite spectrale
$$
E_2^{p, q} = R^q(g \circ f_p)_{\QCoh, *}f_p^*\mathcal{F}
\Rightarrow
R^{p + q}g_{\QCoh, *}\mathcal{F}
$$
dans $\QCoh(\mathcal{O}_\mathcal{Y})$.
\end{proposition}

\begin{proof}
Notons que chacun des morphismes
$f_p : \mathcal{U} \times_\mathcal{X} \ldots \times_\mathcal{X} \mathcal{U} \to
\mathcal{X}$ est quasi-compact et quasi-séparé ; ainsi, $g \circ f_p$
est quasi-compact et quasi-séparé, et l'assertion a donc un sens
(c'est-à-dire que les foncteurs $R^q(g \circ f_p)_{\QCoh, *}$ sont définis).
Il existe une suite spectrale
$$
E_2^{p, q} = R^q(g \circ f_p)_*f_p^{-1}\mathcal{F}
\Rightarrow
R^{p + q}g_*\mathcal{F}
$$
d'après Faisceaux sur les champs, proposition
\ref{stacks-sheaves-proposition-smooth-covering-compute-direct-image}.
En appliquant le foncteur exact
$Q_\mathcal{Y} : \textit{LQCoh}^{fbc}(\mathcal{O}_\mathcal{Y}) \to
\QCoh(\mathcal{O}_\mathcal{Y})$, on obtient la suite spectrale voulue dans
$\QCoh(\mathcal{O}_\mathcal{Y})$.
\end{proof}







\section{Remarques supplémentaires sur les modules quasi-cohérents}
\label{stacks-cohomology-section-further-remarks}

\noindent
Dans cette section, nous réunissons quelques résultats qui aident à comprendre comment
utiliser les modules quasi-cohérents sur les champs algébriques.

\medskip\noindent
Soit $f : \mathcal{U} \to \mathcal{X}$ un morphisme de champs algébriques.
Supposons que $\mathcal{U}$ soit représenté par l'espace algébrique $U$.
Considérons le foncteur
$$
a :
\textit{Mod}(\mathcal{X}_\etale, \mathcal{O}_\mathcal{X})
\longrightarrow
\textit{Mod}(U_\etale, \mathcal{O}_U),\quad
\mathcal{F}
\longmapsto
f^*\mathcal{F}|_{U_\etale}
$$
obtenu en faisant suivre l'image inverse (Faisceaux sur les champs, section
\ref{stacks-sheaves-section-modules}) de la restriction
(Faisceaux sur les champs, section
\ref{stacks-sheaves-section-restriction-algebraic-spaces}).
En appliquant ce foncteur aux modules localement quasi-cohérents, on obtient un foncteur
$$
b : \textit{LQCoh}(\mathcal{O}_\mathcal{X})
\longrightarrow
\QCoh(U_\etale, \mathcal{O}_U)
$$
Voir Faisceaux sur les champs, lemmes \ref{stacks-sheaves-lemma-pullback-lqc} et
\ref{stacks-sheaves-lemma-compare-locally-quasi-coherent}.
On peut encore restreindre notre foncteur à des sous-catégories plus petites et obtenir
$$
c :
\textit{LQCoh}^{fbc}(\mathcal{O}_\mathcal{X})
\longrightarrow
\QCoh(U_\etale, \mathcal{O}_U)
$$
ainsi que
$$
d :
\QCoh(\mathcal{O}_\mathcal{X})
\longrightarrow
\QCoh(U_\etale, \mathcal{O}_U)
$$
Ces foncteurs possèdent les propriétés suivantes :\footnote{Nous suggérons
d'établir pourquoi ces assertions sont vraies sur un coin de table plutôt que de
suivre les références indiquées.}
\begin{enumerate}
\item Le foncteur $a$ est exact. En effet, l'image inverse $f^* = f^{-1}$ est exacte
(Faisceaux sur les champs, section \ref{stacks-sheaves-section-modules}),
et la restriction à $U_\etale$ est exacte ; voir
Faisceaux sur les champs, équation (\ref{stacks-sheaves-equation-restrict}).
\item Le foncteur $b$ est exact. En effet, d'après
Faisceaux sur les champs, lemme \ref{stacks-sheaves-lemma-lqc-colimits}
le foncteur d'inclusion $\textit{LQCoh}(\mathcal{O}_\mathcal{X}) \to
\textit{Mod}(\mathcal{X}_\etale, \mathcal{O}_\mathcal{X})$ est exact.
\item Le foncteur $c$ est exact. En effet, d'après la
proposition \ref{stacks-cohomology-proposition-loc-qcoh-flat-base-change}
le foncteur d'inclusion $\textit{LQCoh}^{fbc}(\mathcal{O}_\mathcal{X}) \to
\textit{Mod}(\mathcal{X}_\etale, \mathcal{O}_\mathcal{X})$ est exact.
\item Le foncteur $d$ est exact à droite, mais n'est pas exact en général.
En effet, d'après Faisceaux sur les champs, lemme
\ref{stacks-sheaves-lemma-qc-colimits}
le foncteur d'inclusion $\QCoh(\mathcal{O}_\mathcal{X}) \to
\textit{Mod}(\mathcal{X}_\etale, \mathcal{O}_\mathcal{X})$ est exact à
droite. Nous omettons de donner un exemple montrant qu'il n'est pas exact.
\item Si $f$ est plat, alors $d$ est exact. Cela résulte de la combinaison du
Lemme \ref{stacks-cohomology-lemma-flat-pullback-quasi-coherent} et de
Faisceaux sur les champs, lemme \ref{stacks-sheaves-lemma-compare-quasi-coherent}.
\item Si $f$ est plat, alors $c$ annule les objets parasites.
En effet, $f^*$ préserve les objets parasites d'après le Lemme \ref{stacks-cohomology-lemma-parasitic}.
Pour tout schéma $V$ étale sur $U$, et donc plat sur $\mathcal{X}$,
on voit alors que
$0 = f^*\mathcal{F}|_{V_\etale} = c(\mathcal{F})|_{V_\etale}$
par compatibilité de la restriction à la localisation étale
(Faisceaux sur les champs, remarque \ref{stacks-sheaves-remark-compare}).
Ainsi, manifestement, $c(\mathcal{F}) = 0$.
\item Si $f$ est plat, alors $c = d \circ Q$. En effet, le noyau et le
conoyau de $Q(\mathcal{F}) \to \mathcal{F}$ sont parasites d'après le
Lemme \ref{stacks-cohomology-lemma-adjoint-kernel-parasitic}. Puisque
$c$ est exact (3) et annule les objets parasites (6),
on voit que l'image par $c$ du morphisme $Q(\mathcal{F}) \to \mathcal{F}$
est un isomorphisme.
\item Les foncteurs $a, b, c, d$ commutent aux limites inductives
et aux sommes directes arbitraires. Cela vaut pour $f^*$ et la restriction,
qui sont des adjoints à gauche, et vaut donc pour $a$. Le résultat pour
$b$, $c$, $d$ découle alors des références données ci-dessus.
\item Les foncteurs $a, b, c, d$ commutent aux produits tensoriels.
\item Si $f$ est plat et surjectif, si $\mathcal{F}$ appartient à
$\textit{LQCoh}^{fbc}(\mathcal{O}_\mathcal{X})$ et si $c(\mathcal{F}) = 0$,
alors $\mathcal{F}$ est parasite. En effet, (7) donne $d(Q(\mathcal{F})) = 0$.
On peut supposer que $U$ est un schéma, par compatibilité de la restriction
avec la localisation étale (voir la référence ci-dessus). Le
Lemme \ref{stacks-cohomology-lemma-quasi-coherent-check-exact}, appliqué à $0 \to Q(\mathcal{F})$
et au morphisme $f : U \to \mathcal{X}$, montre alors que $Q(\mathcal{F}) = 0$.
Ainsi, $\mathcal{F}$ est parasite d'après le Lemme \ref{stacks-cohomology-lemma-adjoint-kernel-parasitic}.
\item Si $f$ est plat et surjectif, alors le foncteur $d$
reflète l'exactitude. Plus précisément, soit $\mathcal{F}^\bullet$ un complexe de
$\QCoh(\mathcal{O}_\mathcal{X})$. Alors $\mathcal{F}^\bullet$ est exact
dans $\QCoh(\mathcal{O}_\mathcal{X})$ si et seulement si $d(\mathcal{F}^\bullet)$
est exact. Nous avons vu une implication en (5). Réciproquement,
supposons que $H^i(d(\mathcal{F}^\bullet)) = 0$. Alors
$\mathcal{G} = H^i(\mathcal{F}^\bullet)$ est un objet de
$\QCoh(\mathcal{O}_\mathcal{X})$ tel que $d(\mathcal{G}) = 0$.
Ainsi, $\mathcal{G}$ est à la fois quasi-cohérent et parasite d'après (10),
donc égal à $0$, par exemple d'après la Remarque \ref{stacks-cohomology-remark-bousfield-colocalization}.
\item
\label{stacks-cohomology-item-hom-restriction}
Si $f$ est plat,
$\mathcal{F}, \mathcal{G} \in \Ob(\QCoh(\mathcal{O}_\mathcal{X}))$,
et si $\mathcal{F}$ est de présentation finie,
alors on a
$$
d(hom(\mathcal{F}, \mathcal{G})) =
\SheafHom_{\mathcal{O}_U}(d(\mathcal{F}), d(\mathcal{G}))
$$
avec les notations du Lemme \ref{stacks-cohomology-lemma-internal-hom-fp-into-qcoh}.
La manière la plus simple de le voir est peut-être la suivante :
\begin{align*}
d(hom(\mathcal{F}, \mathcal{G}))
& =
d(Q(\SheafHom_{\mathcal{O}_\mathcal{X}}(\mathcal{F}, \mathcal{G}))) \\
& =
c(\SheafHom_{\mathcal{O}_\mathcal{X}}(\mathcal{F}, \mathcal{G})) \\
& =
f^*\SheafHom_{\mathcal{O}_\mathcal{X}}(\mathcal{F}, \mathcal{G})|_{U_\etale} \\
& =
\SheafHom_{\mathcal{O}_\mathcal{U}}(f^*\mathcal{F},
f^*\mathcal{G})|_{U_\etale} \\
& =
\SheafHom_{\mathcal{O}_U}(f^*\mathcal{F}|_{U_\etale},
f^*\mathcal{G}|_{U_\etale})
\end{align*}
La première égalité résulte de la construction de $hom$. La deuxième
résulte de (7). La troisième résulte de la définition de $c$. La quatrième résulte de
Modules sur les sites, lemme \ref{sites-modules-lemma-pullback-internal-hom}.
La dernière égalité résulte de la même référence appliquée au morphisme plat
de topos annelés $i_U (U_\etale, \mathcal{O}_U) \to
(\mathcal{U}_\etale, \mathcal{O}_\mathcal{U})$ de
Faisceaux sur les champs, lemme \ref{stacks-sheaves-lemma-compare}.
\item Ajouter ici d'autres éléments.
\end{enumerate}






\section{Limites inductives et cohomologie}
\label{stacks-cohomology-section-colimits}

\noindent
Le lemme suivant s'applique notamment aux diagrammes de faisceaux
quasi-cohérents.

\begin{lemma}
\label{stacks-cohomology-lemma-colimits}
Soit $\mathcal{X}$ un champ algébrique quasi-compact et quasi-séparé.
Alors
$$
\colim_i H^p(\mathcal{X}, \mathcal{F}_i)
\longrightarrow
H^p(\mathcal{X}, \colim_i \mathcal{F}_i)
$$
est un isomorphisme pour tout diagramme filtrant de faisceaux abéliens sur
$\mathcal{X}$. Il en va de même pour les faisceaux abéliens sur $\mathcal{X}_\etale$
lorsque la cohomologie est prise pour la topologie étale.
\end{lemma}

\begin{proof}
Soit $\tau = fppf$, resp.\ $\tau = \etale$. Le lemme résulte de
Cohomologie sur les sites, lemme \ref{sites-cohomology-lemma-colim-global},
appliqué au site $\mathcal{X}_\tau$. Pour en vérifier les hypothèses, nous
utilisons Cohomologie sur les sites, remarque \ref{sites-cohomology-remark-colim-global}.
Soit en effet $\mathcal{B} \subset \Ob(\mathcal{X}_\tau)$ l'ensemble
des objets au-dessus de schémas affines. Autrement dit, un élément
de $\mathcal{B}$ est un morphisme $x : U \to \mathcal{X}$ où $U$ est affine.
Vérifions successivement chacune des conditions (1) -- (4) de la remarque :
\begin{enumerate}
\item Puisque $\mathcal{X}$ est quasi-compact, il existe un morphisme
lisse et surjectif $x : U \to \mathcal{X}$ où $U$ est affine
(Propriétés des champs, lemme
\ref{stacks-properties-lemma-quasi-compact-stack}).
Alors $h_x^\# \to *$ est un morphisme surjectif de faisceaux sur $\mathcal{X}_\tau$.
\item Puisque les recouvrements dans $\mathcal{X}_\tau$ sont des recouvrements
fppf, resp.\ étales, tout recouvrement de $U \in \mathcal{B}$ est raffiné
par un recouvrement fppf affine fini ; voir
Topologies, lemme \ref{topologies-lemma-fppf-affine},
resp.\ lemme \ref{topologies-lemma-etale-affine}.
\item Soient $x : U \to \mathcal{X}$ et $x' : U' \to \mathcal{X}$ dans
$\mathcal{B}$. Le produit $h_x^\# \times h_{x'}^\#$ dans
$\Sh(\mathcal{X}_\tau)$ est égal au faisceau sur $\mathcal{X}_\tau$
déterminé par l'espace algébrique $W = U \times_{x, \mathcal{X}, x'} U'$ au-dessus de
$\mathcal{X}$ : pour un objet $y : V \to \mathcal{X}$ de
$\mathcal{X}_\tau$, on a
$(h_x^\# \times h_{x'}^\#)(y) = \{f : V \to W \mid y =
x \circ \text{pr}_1 \circ f = x' \circ \text{pr}_2 \circ f\}$.
L'espace algébrique $W$ est quasi-compact, car $\mathcal{X}$
est quasi-séparé ; voir, par exemple, Morphismes de champs, lemme
\ref{stacks-morphisms-lemma-quasi-compact-quasi-separated-permanence}
On peut donc choisir un schéma affine $U''$ et
un morphisme étale surjectif $U'' \to W$. Notons $x'' : U'' \to \mathcal{X}$
le composé de $U'' \to W$ et $W \to \mathcal{X}$. Alors
$h_{x''}^\# \to h_x^\# \times h_{x'}^\#$ est surjectif, comme voulu.
\item Soient $x : U \to \mathcal{X}$ et $x' : U' \to \mathcal{X}$ dans
$\mathcal{B}$. Soient $a, b : U \to U'$ des morphismes au-dessus de $\mathcal{X}$,
c'est-à-dire $a, b  : x \to x'$ des morphismes dans $\mathcal{X}_\tau$.
Alors l'égalisateur de $h_a$ et $h_b$ est représenté par
l'égalisateur de $a, b : U \to U'$, qui est un
schéma affine au-dessus de $\mathcal{X}$ et appartient donc à $\mathcal{B}$.
\end{enumerate}
Cela achève la démonstration.
\end{proof}

\begin{lemma}
\label{stacks-cohomology-lemma-colimit-cohomology}
Soit $f : \mathcal{X} \to \mathcal{Y}$ un morphisme quasi-compact et quasi-séparé
de champs algébriques. Soit $\mathcal{F} = \colim \mathcal{F}_i$
une limite inductive filtrante de faisceaux abéliens sur $\mathcal{X}$.
Alors, pour tout $p \geq 0$, on a
$$
R^pf_*\mathcal{F} = \colim R^pf_*\mathcal{F}_i.
$$
Il en va de même pour les faisceaux abéliens sur $\mathcal{X}_\etale$
lorsque les images directes supérieures sont prises pour la topologie étale.
\end{lemma}

\begin{proof}
Nous le démontrons pour la topologie fppf ; la démonstration pour la topologie
étale est la même. Rappelons que $R^if_*\mathcal{F}$ est le faisceau sur
$\mathcal{Y}_{fppf}$ associé au préfaisceau
$$
(y : V \to \mathcal{Y})
\longmapsto
H^i(V \times_{y, \mathcal{Y}} \mathcal{X}, \text{pr}^{-1}\mathcal{F})
$$
Voir
Faisceaux sur les champs, lemme \ref{stacks-sheaves-lemma-pushforward-restriction}.
Rappelons que la limite inductive est le faisceau associé à la limite inductive des préfaisceaux.
Lorsque $V$ est affine, le produit fibré $V \times_\mathcal{Y} \mathcal{X}$
est quasi-compact et quasi-séparé. On peut donc appliquer le
Lemme \ref{stacks-cohomology-lemma-colimits} à $H^p(V \times_\mathcal{Y} \mathcal{X}, -)$
lorsque $V$ est affine. Puisque tout $V$ possède un recouvrement fppf par des
objets affines, cela démontre le lemme. Nous omettons certains détails.
\end{proof}

\begin{lemma}
\label{stacks-cohomology-lemma-quasi-coherent-pushforward-direct-sums}
Soit $f : \mathcal{X} \to \mathcal{Y}$ un morphisme quasi-compact et quasi-séparé
de champs algébriques. Le foncteur $f_{\QCoh, *}$
et les foncteurs $R^if_{\QCoh, *}$ commutent aux sommes directes
et aux limites inductives filtrantes.
\end{lemma}

\begin{proof}
Les foncteurs $f_*$ et $R^if_*$ commutent aux sommes directes et aux limites
inductives filtrantes sur tous les modules d'après le Lemme \ref{stacks-cohomology-lemma-colimit-cohomology}.
Le lemme en résulte, car $f_{\QCoh, *} = Q \circ f_*$ et
$R^if_{\QCoh, *} = Q \circ R^if_*$, et $Q$ commute à
toutes les limites inductives ; voir le Lemme \ref{stacks-cohomology-lemma-adjoint-kernel-parasitic}.
\end{proof}

\begin{lemma}
\label{stacks-cohomology-lemma-quasi-coherent-pushforward-affine}
Soit $f : \mathcal{X} \to \mathcal{Y}$ un morphisme affine de champs
algébriques. Les foncteurs $R^if_{\QCoh, *}$, $i > 0$, sont nuls, et le foncteur
$f_{\QCoh, *}$ est exact et commute aux sommes directes et à toutes les limites inductives.
\end{lemma}

\begin{proof}
Puisque $R^if_{\QCoh, *} = Q \circ R^if_*$, l'annulation résulte du
Lemme \ref{stacks-cohomology-lemma-loc-qcoh-fbc-affine-direct-image}.
Cette annulation entraîne que $f_{\QCoh, *}$ est exact, car
$\{R^if_{\QCoh, *}\}_{i \geq 0}$ forme un $\delta$-foncteur ; voir la
proposition \ref{stacks-cohomology-proposition-direct-image-quasi-coherent}.
Puis $f_{\QCoh, *}$ commute aux sommes directes, par exemple d'après le
Lemme \ref{stacks-cohomology-lemma-quasi-coherent-pushforward-direct-sums}.
Un foncteur exact qui commute aux sommes directes commute
à toutes les limites inductives.
\end{proof}

\noindent
Le lemme suivant affirme que les modules de présentation finie se comportent
comme prévu sur les champs algébriques quasi-compacts et quasi-séparés.

\begin{lemma}
\label{stacks-cohomology-lemma-finite-presentation-quasi-compact-colimit}
Soit $\mathcal{X}$ un champ algébrique quasi-compact et quasi-séparé.
Soit $I$ un ensemble ordonné filtrant et soit $(\mathcal{F}_i, \varphi_{ii'})$ un
système indexé par $I$ de $\mathcal{O}_\mathcal{X}$-modules. Soit $\mathcal{G}$ un
$\mathcal{O}_\mathcal{X}$-module de présentation finie. Alors on a
$$
\colim_i \Hom_\mathcal{X}(\mathcal{G}, \mathcal{F}_i)
=
\Hom_\mathcal{X}(\mathcal{G}, \colim_i \mathcal{F}_i).
$$
En particulier, $\Hom_\mathcal{X}(\mathcal{G}, -)$ commute aux limites
inductives filtrantes dans $\QCoh(\mathcal{O}_\mathcal{X})$.
\end{lemma}

\begin{proof}
L'égalité affichée est un cas particulier de Modules sur les sites, lemme
\ref{sites-modules-lemma-finite-presentation-quasi-compact-colimit}.
Pour l'appliquer, il faut vérifier les hypothèses de la partie (4) de
Sites, lemme \ref{sites-lemma-directed-colimits-global-sections},
pour le site $\mathcal{X}_{fppf}$. À cette fin, nous vérifions
les hypothèses (2)(a), (2)(b), (2)(c) de
Sites, remarque \ref{sites-remark-stronger-conditions}.
Soit en effet $\mathcal{B} \subset \Ob(\mathcal{X}_{fppf})$ l'ensemble
des objets au-dessus de schémas affines. Autrement dit, un élément
de $\mathcal{B}$ est un morphisme $x : U \to \mathcal{X}$ où $U$ est affine.
Vérifions successivement chacune des conditions (2)(a), (2)(b) et (2)(c)
de la remarque :
\begin{enumerate}
\item Puisque $\mathcal{X}$ est quasi-compact, il existe un morphisme
lisse et surjectif $x : U \to \mathcal{X}$ où $U$ est affine
(Propriétés des champs, lemme
\ref{stacks-properties-lemma-quasi-compact-stack}).
Alors $h_x^\# \to *$ est un morphisme surjectif de faisceaux sur $\mathcal{X}_{fppf}$.
\item Puisque les recouvrements dans $\mathcal{X}_{fppf}$ sont des recouvrements
fppf, tout recouvrement de $U \in \mathcal{B}$ est raffiné
par un recouvrement fppf affine fini ; voir
Topologies, lemme \ref{topologies-lemma-fppf-affine}.
\item Soient $x : U \to \mathcal{X}$ et $x' : U' \to \mathcal{X}$ dans
$\mathcal{B}$. Le produit $h_x^\# \times h_{x'}^\#$ dans
$\Sh(\mathcal{X}_{fppf})$ est égal au faisceau sur $\mathcal{X}_{fppf}$
déterminé par l'espace algébrique $W = U \times_{x, \mathcal{X}, x'} U'$ au-dessus de
$\mathcal{X}$ : pour un objet $y : V \to \mathcal{X}$ de
$\mathcal{X}_{fppf}$, on a
$(h_x^\# \times h_{x'}^\#)(y) = \{f : V \to W \mid y =
x \circ \text{pr}_1 \circ f = x' \circ \text{pr}_2 \circ f\}$.
L'espace algébrique $W$ est quasi-compact, car $\mathcal{X}$
est quasi-séparé ; voir, par exemple, Morphismes de champs, lemme
\ref{stacks-morphisms-lemma-quasi-compact-quasi-separated-permanence}
On peut donc choisir un schéma affine $U''$ et
un morphisme étale surjectif $U'' \to W$. Notons $x'' : U'' \to \mathcal{X}$
le composé de $U'' \to W$ et $W \to \mathcal{X}$. Alors
$h_{x''}^\# \to h_x^\# \times h_{x'}^\#$ est surjectif, comme voulu.
\end{enumerate}
Pour la dernière assertion, observons que le foncteur d'inclusion
$\QCoh(\mathcal{O}_X) \to \textit{Mod}(\mathcal{O}_X)$
commute aux limites inductives et que les modules de présentation finie
sont quasi-cohérents. Voir Faisceaux sur les champs, lemme
\ref{stacks-sheaves-lemma-quasi-coherent-algebraic-stack}.
\end{proof}



\section{Les sites lisse-étale et plat-fppf}
\label{stacks-cohomology-section-lisse-etale}

\noindent
Dans l'ouvrage \cite{LM-B}, beaucoup des résultats ci-dessus sont démontrés à l'aide du
site lisse-étale d'un champ algébrique. Nous définissons ici ce site.
Dans Exemples, section \ref{examples-section-lisse-etale-not-functorial},
nous montrons que le site lisse-étale n'est pas fonctoriel.
Nous définissons aussi son analogue, le site plat-fppf, qui est mieux adapté
au développement des champs algébriques adopté dans le projet Champs
(car nous prenons la topologie fppf comme topologie de base). Bien entendu, le
site plat-fppf n'est pas non plus fonctoriel.

\begin{definition}
\label{stacks-cohomology-definition-lisse-etale}
Soit $\mathcal{X}$ un champ algébrique.
\begin{enumerate}
\item Le {\it site lisse-étale} de $\mathcal{X}$ est la sous-catégorie pleine
$\mathcal{X}_{lisse,\etale}$\footnote{Dans la littérature, le
site est noté $\text{Lis-\'et}(\mathcal{X})$ ou
$\text{Lis-Et}(\mathcal{X})$, et le topos associé est noté
$\mathcal{X}_{\text{lis-\'e}t}$ ou $\mathcal{X}_{\text{lis-et}}$.
Dans le projet Champs, notre convention consiste à nommer le site et à
noter le topos correspondant $\Sh(\mathcal{C})$.} de $\mathcal{X}$
dont les objets sont les $x \in \Ob(\mathcal{X})$ au-dessus d'un schéma $U$
tels que $x : U \to \mathcal{X}$ soit lisse. Un recouvrement de
$\mathcal{X}_{lisse,\etale}$ est une famille de morphismes
$\{x_i \to x\}_{i \in I}$ de $\mathcal{X}_{lisse,\etale}$
qui forme un recouvrement de $\mathcal{X}_\etale$.
\item Le {\it site plat-fppf} de $\mathcal{X}$ est la sous-catégorie pleine
$\mathcal{X}_{flat,fppf}$ de $\mathcal{X}$
dont les objets sont les $x \in \Ob(\mathcal{X})$ au-dessus d'un schéma $U$
tels que $x : U \to \mathcal{X}$ soit plat. Un recouvrement de
$\mathcal{X}_{flat,fppf}$ est une famille de morphismes
$\{x_i \to x\}_{i \in I}$ de $\mathcal{X}_{flat,fppf}$
qui forme un recouvrement de $\mathcal{X}_{fppf}$.
\end{enumerate}
\end{definition}

\noindent
Nous notons $\mathcal{O}_{\mathcal{X}_{lisse,\etale}}$
la restriction de $\mathcal{O}_\mathcal{X}$ au site lisse-étale,
et de même pour $\mathcal{O}_{\mathcal{X}_{flat,fppf}}$.
La relation entre le site lisse-étale et le site étale est
la suivante (nous nous en tenons surtout aux propriétés « topologiques » dans ce lemme).

\begin{lemma}
\label{stacks-cohomology-lemma-lisse-etale}
Soit $\mathcal{X}$ un champ algébrique.
\begin{enumerate}
\item Le foncteur d'inclusion
$\mathcal{X}_{lisse,\etale} \to \mathcal{X}_\etale$
est pleinement fidèle, continu et cocontinu. Il s'ensuit que
\begin{enumerate}
\item il existe un morphisme de topos
$$
g :
\Sh(\mathcal{X}_{lisse,\etale})
\longrightarrow
\Sh(\mathcal{X}_\etale)
$$
où $g^{-1}$ est donné par la restriction,
\item le foncteur $g^{-1}$ admet un adjoint à gauche $g_!^{Sh}$ sur les faisceaux d'ensembles,
\item les morphismes d'adjonction $g^{-1}g_* \to \text{id}$ et
$\text{id} \to g^{-1}g_!^{Sh}$ sont des isomorphismes,
\item le foncteur $g^{-1}$ admet un adjoint à gauche $g_!$ sur les faisceaux abéliens,
\item le morphisme d'adjonction $\text{id} \to g^{-1}g_!$ est un isomorphisme, et
\item on a $g^{-1}\mathcal{O}_\mathcal{X} =
\mathcal{O}_{\mathcal{X}_{lisse,\etale}}$ ; ainsi, $g$ induit un morphisme plat
de topos annelés tel que $g^{-1} = g^*$.
\end{enumerate}
\item Le foncteur d'inclusion
$\mathcal{X}_{flat,fppf} \to \mathcal{X}_{fppf}$
est pleinement fidèle, continu et cocontinu. Il s'ensuit que
\begin{enumerate}
\item il existe un morphisme de topos
$$
g :
\Sh(\mathcal{X}_{flat,fppf})
\longrightarrow
\Sh(\mathcal{X}_{fppf})
$$
où $g^{-1}$ est donné par la restriction,
\item le foncteur $g^{-1}$ admet un adjoint à gauche $g_!^{Sh}$ sur les faisceaux d'ensembles,
\item les morphismes d'adjonction $g^{-1}g_* \to \text{id}$ et
$\text{id} \to g^{-1}g_!^{Sh}$ sont des isomorphismes,
\item le foncteur $g^{-1}$ admet un adjoint à gauche $g_!$ sur les faisceaux abéliens,
\item le morphisme d'adjonction $\text{id} \to g^{-1}g_!$ est un isomorphisme, et
\item on a $g^{-1}\mathcal{O}_\mathcal{X} =
\mathcal{O}_{\mathcal{X}_{flat,fppf}}$ ; ainsi, $g$ induit un morphisme plat
de topos annelés tel que $g^{-1} = g^*$.
\end{enumerate}
\end{enumerate}
\end{lemma}

\begin{proof}
Dans les deux cas, il est immédiat que le foncteur est pleinement fidèle,
continu et cocontinu (voir
Sites, définitions \ref{sites-definition-continuous} et
\ref{sites-definition-cocontinuous}).
Les propriétés (a), (b), (c) résultent donc de
Sites, lemmes \ref{sites-lemma-when-shriek} et
\ref{sites-lemma-back-and-forth}.
Les parties (d), (e) résultent de
Modules sur les sites, lemmes \ref{sites-modules-lemma-g-shriek-adjoint} et
\ref{sites-modules-lemma-back-and-forth}.
La partie (f) est immédiate.
\end{proof}

\begin{lemma}
\label{stacks-cohomology-lemma-lisse-etale-cohomology}
Soit $\mathcal{X}$ un champ algébrique. On reprend les notations du
Lemme \ref{stacks-cohomology-lemma-lisse-etale}.
\begin{enumerate}
\item Pour un faisceau abélien $\mathcal{F}$ sur $\mathcal{X}_\etale$, on a
\begin{enumerate}
\item $H^p(\mathcal{X}_\etale, \mathcal{F}) =
H^p(\mathcal{X}_{lisse,\etale}, g^{-1}\mathcal{F})$, et
\item $H^p(x, \mathcal{F}) =
H^p(\mathcal{X}_{lisse,\etale}/x, g^{-1}\mathcal{F})$
pour tout objet $x$ de $\mathcal{X}_{lisse,\etale}$.
\end{enumerate}
Il en va de même pour les faisceaux de modules.
\item Pour un faisceau abélien $\mathcal{F}$ sur $\mathcal{X}_{fppf}$, on a
\begin{enumerate}
\item $H^p(\mathcal{X}_{fppf}, \mathcal{F}) =
H^p(\mathcal{X}_{flat,fppf}, g^{-1}\mathcal{F})$, et
\item $H^p(x, \mathcal{F}) =
H^p(\mathcal{X}_{flat,fppf}/x, g^{-1}\mathcal{F})$
pour tout objet $x$ de $\mathcal{X}_{flat,fppf}$.
\end{enumerate}
Il en va de même pour les faisceaux de modules.
\end{enumerate}
\end{lemma}

\begin{proof}
La partie (1)(a) résulte de Faisceaux sur les champs, lemme
\ref{stacks-sheaves-lemma-cohomology-on-subcategory}, appliqué au foncteur
d'inclusion $\mathcal{X}_{lisse,\etale} \to \mathcal{X}_\etale$.
La partie (1)(b) résulte de la partie (1)(a). En effet, si $x$ est au-dessus du
schéma $U$, alors le site $\mathcal{X}_\etale/x$ est équivalent
à $(\Sch/U)_\etale$, et $\mathcal{X}_{lisse,\etale}$ est équivalent
à $U_{lisse,\etale}$. La partie (2) se démontre de la même manière.
\end{proof}

\begin{lemma}
\label{stacks-cohomology-lemma-lisse-etale-modules}
Soit $\mathcal{X}$ un champ algébrique. On reprend les notations du
Lemme \ref{stacks-cohomology-lemma-lisse-etale}.
\begin{enumerate}
\item Il existe un foncteur
$$
g_! :
\textit{Mod}(\mathcal{X}_{lisse,\etale},
\mathcal{O}_{\mathcal{X}_{lisse,\etale}})
\longrightarrow
\textit{Mod}(\mathcal{X}_\etale, \mathcal{O}_{\mathcal{X}})
$$
qui est adjoint à gauche de $g^*$. En outre, il coïncide avec le foncteur $g_!$
sur les faisceaux abéliens, et $g^*g_! = \text{id}$.
\item Il existe un foncteur
$$
g_! :
\textit{Mod}(\mathcal{X}_{flat,fppf},
\mathcal{O}_{\mathcal{X}_{flat,fppf}})
\longrightarrow
\textit{Mod}(\mathcal{X}_{fppf}, \mathcal{O}_{\mathcal{X}})
$$
qui est adjoint à gauche de $g^*$. En outre, il coïncide avec le foncteur $g_!$
sur les faisceaux abéliens, et $g^*g_! = \text{id}$.
\end{enumerate}
\end{lemma}

\begin{proof}
Dans les deux cas, l'existence du foncteur $g_!$ résulte de
Modules sur les sites, lemme \ref{sites-modules-lemma-lower-shriek-modules}.
Pour voir que $g_!$ coïncide avec le foncteur sur les faisceaux abéliens, nous
montrons que les morphismes de Modules sur les sites, équation
(\ref{sites-modules-equation-compare-on-localizations})
sont des isomorphismes.

\medskip\noindent
Cas lisse-étale. Soit $x \in \Ob(\mathcal{X}_{lisse,\etale})$
au-dessus d'un schéma $U$, avec $x : U \to \mathcal{X}$ lisse.
Considérons le foncteur pleinement fidèle induit
$$
g' :
\mathcal{X}_{lisse,\etale}/x
\longrightarrow
\mathcal{X}_\etale/x
$$
Le membre de droite s'identifie à $(\Sch/U)_\etale$, et celui de
gauche à la sous-catégorie pleine des schémas $U'/U$ tels que
le composé $U' \to U \to \mathcal{X}$ soit lisse. Ainsi,
Cohomologie étale, lemme
\ref{etale-cohomology-lemma-compare-structure-sheaves}
s'applique.

\medskip\noindent
Cas plat-fppf. Soit $x \in \Ob(\mathcal{X}_{flat,fppf})$
au-dessus d'un schéma $U$, avec $x : U \to \mathcal{X}$ plat.
Considérons le foncteur pleinement fidèle induit
$$
g' :
\mathcal{X}_{flat,fppf}/x
\longrightarrow
\mathcal{X}_{fppf}/x
$$
Le membre de droite s'identifie à $(\Sch/U)_{fppf}$, et celui de
gauche à la sous-catégorie pleine des schémas $U'/U$ tels que
le composé $U' \to U \to \mathcal{X}$ soit plat. Ainsi,
Cohomologie étale, lemme
\ref{etale-cohomology-lemma-compare-structure-sheaves}
s'applique.

\medskip\noindent
Dans les deux cas, l'égalité $g^*g_! = \text{id}$ résulte de
$g^* = g^{-1}$ et de
l'égalité pour les faisceaux abéliens du Lemme \ref{stacks-cohomology-lemma-lisse-etale}.
\end{proof}

\begin{lemma}
\label{stacks-cohomology-lemma-lisse-etale-structure-sheaf}
Soit $\mathcal{X}$ un champ algébrique. On reprend les notations des
lemmes \ref{stacks-cohomology-lemma-lisse-etale} et \ref{stacks-cohomology-lemma-lisse-etale-modules}.
\begin{enumerate}
\item On a $g_!\mathcal{O}_{\mathcal{X}_{lisse,\etale}} =
\mathcal{O}_\mathcal{X}$.
\item On a $g_!\mathcal{O}_{\mathcal{X}_{flat, fppf}} =
\mathcal{O}_\mathcal{X}$.
\end{enumerate}
\end{lemma}

\begin{proof}
Dans cette démonstration, nous posons
$\mathcal{C} = \mathcal{X}_\etale$
(resp.\ $\mathcal{C} = \mathcal{X}_{fppf}$)
et nous notons
$\mathcal{C}' = \mathcal{X}_{lisse,\etale}$
(resp.\ $\mathcal{C}' = \mathcal{X}_{flat, fppf}$).
Alors $\mathcal{C}'$ est une sous-catégorie pleine de $\mathcal{C}$.
Dans cette démonstration, nous regardons les objets $V$ de $\mathcal{C}$
comme des schémas au-dessus de $\mathcal{X}$ et les objets $U$ de $\mathcal{C}'$
comme des schémas lisses (resp.\ plats) au-dessus de $\mathcal{X}$.
Enfin, nous posons $\mathcal{O} = \mathcal{O}_\mathcal{X}$
et $\mathcal{O}' = \mathcal{O}_{\mathcal{X}_{lisse,\etale}}$
(resp.\ $\mathcal{O}' = \mathcal{O}_{\mathcal{X}_{flat,fppf}}$).
Avec ces notations, on a $\mathcal{O}(V) = \Gamma(V, \mathcal{O}_V)$
et $\mathcal{O}'(U) = \Gamma(U, \mathcal{O}_U)$.
Considérons l'homomorphisme de $\mathcal{O}$-modules
$g_!\mathcal{O}' \to \mathcal{O}$
adjoint à l'identification $\mathcal{O}' = g^{-1}\mathcal{O}$.

\medskip\noindent
Rappelons que $g_!\mathcal{O}'$ est le faisceau associé au préfaisceau
$g_{p!}\mathcal{O}'$ défini par la règle
$$
V \longmapsto \colim_{V \to U} \mathcal{O}'(U)
$$
où la limite inductive est prise dans la catégorie des groupes abéliens
(Modules sur les sites, définition \ref{sites-modules-definition-g-shriek}).
Nous utiliserons fréquemment ci-dessous le fait que, si
$$
V \to U \to U'
$$
sont des morphismes et si $f' \in \mathcal{O}'(U')$ a pour restriction
$f \in \mathcal{O}'(U)$, alors $(V \to U, f)$ et $(V \to U', f')$
définissent le même élément de la limite inductive. De plus,
$g_!\mathcal{O}' \to \mathcal{O}$ envoie simplement l'élément
$(V \to U, f)$ sur l'image réciproque de $f$ dans $V$.

\medskip\noindent
Démontrons que $g_!\mathcal{O}' \to \mathcal{O}$ est surjectif.
Soit $h \in \mathcal{O}(V)$ pour un objet $V$ de $\mathcal{C}$.
Il suffit de montrer que $h$ appartient localement à l'image. Choisissons un objet
$U$ de $\mathcal{C}'$ correspondant à un morphisme lisse surjectif
$U \to \mathcal{X}$. Puisque $U \times_\mathcal{X} V \to V$ est lisse et
surjectif, après avoir remplacé $V$ par les membres d'un recouvrement étale de $V$,
nous pouvons supposer qu'il existe un morphisme $V \to U$ ; voir
Topologies sur les espaces algébriques, lemme
\ref{spaces-topologies-lemma-etale-dominates-smooth}. À l'aide de $h$, nous obtenons
un morphisme $V \to U \times \mathbf{A}^1$ tel que, si l'on écrit
$\mathbf{A}^1 = \Spec(\mathbf{Z}[t])$, l'élément
$t \in \mathcal{O}(U \times \mathbf{A}^1)$ ait pour image réciproque $h$.
Comme $U \times \mathbf{A}^1$ est un objet de $\mathcal{C}'$,
$(V \to U \times \mathbf{A}^1, t)$ est un élément de la limite inductive
ci-dessus qui s'envoie sur $h \in \mathcal{O}(V)$,
comme voulu.

\medskip\noindent
Supposons que $s \in g_!\mathcal{O}'(V)$ soit une section
dont l'image dans $\mathcal{O}(V)$ est nulle. Pour achever la démonstration,
il reste à montrer que $s$ est nul. Après avoir remplacé $V$ par les membres
d'un recouvrement, nous pouvons supposer que $s$ est un élément de la limite inductive
$$
\colim_{V \to U} \mathcal{O}'(U)
$$
Écrivons $s = \sum (\varphi_i, s_i)$ comme une somme finie, où
$\varphi_i : V \to U_i$, les $U_i$ sont lisses (resp.\ plats) sur $\mathcal{X}$ et
$s_i \in \Gamma(U_i, \mathcal{O}_{U_i})$. Choisissons un schéma $W$ étale
et surjectif sur l'espace algébrique
$U = U_1 \times_\mathcal{X} \ldots \times_\mathcal{X} U_n$.
Notons que $W$ est encore lisse (resp.\ plat) sur $\mathcal{X}$, c'est-à-dire
qu'il définit un objet de $\mathcal{C}'$. Le produit fibré
$$
V' = V \times_{(\varphi_1, \ldots, \varphi_n), U} W
$$
est étale et surjectif sur $V$ ; il suffit donc de montrer que l'image de $s$
est nulle dans $g_!\mathcal{O}'(V')$. La restriction
$\sum (\varphi_i, s_i)|_{V'}$ correspond à la somme des images réciproques
des fonctions $s_i$ sur $W$. Autrement dit, nous sommes ramenés au cas
d'un couple $(\varphi, s)$, où $\varphi : V \to U$ est un morphisme avec $U$
dans $\mathcal{C}'$ et où la restriction de $s \in \mathcal{O}'(U)$ dans
$\mathcal{O}(V)$ est nulle. D'après le diagramme commutatif
$$
\xymatrix{
V \ar[rr]_-{(\varphi, 0)} \ar[rrd]_\varphi & & U \times \mathbf{A}^1 \\
& & U \ar[u]_{(\text{id}, 0)}
}
$$
on voit que
$((\varphi, 0) : V \to U \times \mathbf{A}^1, \text{pr}_2^*x)$
représente zéro dans la limite inductive ci-dessus. Nous pouvons donc
remplacer $U$ par $U \times \mathbf{A}^1$, $\varphi$ par $(\varphi, 0)$
et $s$ par $\text{pr}_1^*s + \text{pr}_2^*x$. Nous pouvons ainsi supposer que
le lieu $Z : s = 0$ dans $U$ défini par $s$ est lisse (resp.\ plat)
sur $\mathcal{X}$. On voit alors que $(V \to Z, 0)$ et $(\varphi, s)$
ont la même valeur dans la limite inductive, c'est-à-dire que l'élément $s$
est nul, comme voulu.
\end{proof}

\noindent
Les sites lisse-étale et plat-fppf permettent de caractériser
les modules parasites comme suit.

\begin{lemma}
\label{stacks-cohomology-lemma-parasitic-in-terms-flat-fppf}
Soit $\mathcal{X}$ un champ algébrique.
\begin{enumerate}
\item Soit $\mathcal{F}$ un $\mathcal{O}_\mathcal{X}$-module
sur $\mathcal{X}_\etale$ ayant la propriété de changement de base plat.
Les conditions suivantes sont équivalentes :
\begin{enumerate}
\item $\mathcal{F}$ est parasite ;
\item $g^*\mathcal{F} = 0$, où
$g : \Sh(\mathcal{X}_{lisse,\etale}) \to
\Sh(\mathcal{X}_\etale)$ est défini comme dans le lemme \ref{stacks-cohomology-lemma-lisse-etale}.
\end{enumerate}
\item Soit $\mathcal{F}$ un $\mathcal{O}_\mathcal{X}$-module sur
$\mathcal{X}_{fppf}$. Les conditions suivantes sont équivalentes :
\begin{enumerate}
\item $\mathcal{F}$ est parasite ;
\item $g^*\mathcal{F} = 0$, où
$g :  \Sh(\mathcal{X}_{flat,fppf}) \to \Sh(\mathcal{X}_{fppf})$
est défini comme dans le lemme \ref{stacks-cohomology-lemma-lisse-etale}.
\end{enumerate}
\end{enumerate}
\end{lemma}

\begin{proof}
La partie (2) résulte immédiatement des définitions (c'est l'un des avantages
du site plat-fppf sur le site lisse-étale). L'implication
(1)(a) $\Rightarrow$ (1)(b) est elle aussi immédiate. Pour établir (1)(b)
$\Rightarrow$ (1)(a), soient $U$ un schéma et $x : U \to \mathcal{X}$
un morphisme lisse surjectif. Alors $x$ est un objet du
site lisse-étale de $\mathcal{X}$. La condition (1)(b)
implique donc que $\mathcal{F}|_{U_\etale} = 0$. Soit $V \to \mathcal{X}$
un morphisme plat, où $V$ est un schéma. Posons $W = U \times_\mathcal{X} V$
et considérons le diagramme
$$
\xymatrix{
W \ar[d]_p \ar[r]_q & V \ar[d] \\
U \ar[r] & \mathcal{X}
}
$$
La projection $p : W \to U$ est plate, tandis que la projection
$q : W \to V$ est lisse et surjective. Il en résulte que $q_{small}^*$
est un foncteur fidèle sur les modules quasi-cohérents. Par hypothèse,
$\mathcal{F}$ possède la propriété de changement de base plat, d'où un isomorphisme
$p_{small}^*\mathcal{F}|_{U_\etale} \cong
q_{small}^*\mathcal{F}|_{V_\etale}$. Ainsi, si $\mathcal{F}$
appartient au noyau de $g^*$, alors $\mathcal{F}|_{V_\etale} = 0$,
comme voulu.
\end{proof}






\section{Fonctorialité des sites lisse-étale et plat-fppf}
\label{stacks-cohomology-section-lisse-etale-functorial}

\noindent
Le site lisse-étale est fonctoriel pour les morphismes lisses de champs
algébriques, et le site plat-fppf l'est pour les morphismes plats de champs
algébriques. Avertissons le lecteur que les topos lisse-étale et plat-fppf ne
sont pas fonctoriels par rapport à tous les morphismes de champs algébriques ; voir
Exemples, section \ref{examples-section-lisse-etale-not-functorial}.

\begin{lemma}
\label{stacks-cohomology-lemma-lisse-etale-functorial}
Soit $f : \mathcal{X} \to \mathcal{Y}$ un morphisme de champs algébriques.
\begin{enumerate}
\item Si $f$ est lisse, alors $f$ se restreint en un foncteur continu et
cocontinu
$\mathcal{X}_{lisse,\etale} \to \mathcal{Y}_{lisse,\etale}$
qui donne un morphisme de topos annelés s'insérant dans le diagramme
commutatif suivant
$$
\xymatrix{
\Sh(\mathcal{X}_{lisse,\etale}) \ar[r]_{g'} \ar[d]_{f'} &
\Sh(\mathcal{X}_\etale) \ar[d]^f \\
\Sh(\mathcal{Y}_{lisse,\etale}) \ar[r]^g &
\Sh(\mathcal{Y}_\etale)
}
$$
On a $f'_*(g')^{-1} = g^{-1}f_*$ et $g'_!(f')^{-1} = f^{-1}g_!$.
\item Si $f$ est plat, alors $f$ se restreint en un foncteur continu et
cocontinu
$\mathcal{X}_{flat,fppf} \to \mathcal{Y}_{flat,fppf}$
qui donne un morphisme de topos annelés s'insérant dans le diagramme
commutatif suivant
$$
\xymatrix{
\Sh(\mathcal{X}_{flat,fppf}) \ar[r]_{g'} \ar[d]_{f'} &
\Sh(\mathcal{X}_{fppf}) \ar[d]^f \\
\Sh(\mathcal{Y}_{flat,fppf}) \ar[r]^g &
\Sh(\mathcal{Y}_{fppf})
}
$$
On a $f'_*(g')^{-1} = g^{-1}f_*$ et $g'_!(f')^{-1} = f^{-1}g_!$.
\end{enumerate}
\end{lemma}

\begin{proof}
La première assertion résulte du fait que, si $x \in \Ob(\mathcal{X})$
est au-dessus d'un schéma $U$, si $x : U \to \mathcal{X}$ est lisse
(resp.\ plat) et si $f$ est lisse (resp.\ plat), alors
$f(x) : U \to \mathcal{Y}$ est lisse (resp.\ plat) ; voir
Morphismes de champs, lemmes \ref{stacks-morphisms-lemma-composition-smooth} et
\ref{stacks-morphisms-lemma-composition-flat}. Le foncteur induit
$\mathcal{X}_{lisse,\etale} \to \mathcal{Y}_{lisse,\etale}$
(resp.\ $\mathcal{X}_{flat,fppf} \to \mathcal{Y}_{flat,fppf}$) est
continu et cocontinu par notre définition des recouvrements dans ces
catégories. Enfin, la commutativité du diagramme résulte du fait que les
morphismes horizontaux sont donnés par les foncteurs d'inclusion (voir le
lemme \ref{stacks-cohomology-lemma-lisse-etale}) et du
lemme \ref{sites-lemma-composition-cocontinuous} de Sites.

\medskip\noindent
Pour montrer que $f'_*(g')^{-1} = g^{-1}f_*$, soit $\mathcal{F}$ un faisceau
sur $\mathcal{X}_\etale$ (resp.\ $\mathcal{X}_{fppf}$).
Il existe un morphisme canonique d'image inverse
$$
g^{-1}f_*\mathcal{F} \longrightarrow f'_*(g')^{-1}\mathcal{F}
$$
voir Sites, section \ref{sites-section-pullback}.
Nous affirmons que ce morphisme est un isomorphisme.
Pour le montrer, prenons un objet $y$ de $\mathcal{Y}_{lisse,\etale}$
(resp.\ $\mathcal{Y}_{flat,fppf}$). Supposons que $y$ soit au-dessus du schéma
$V$ et que $y : V \to \mathcal{Y}$ soit lisse (resp.\ plat). Comme
$g^{-1}$ est le foncteur de restriction, on obtient
$$
\left(g^{-1}f_*\mathcal{F}\right)(y) =
\Gamma(V \times_{y, \mathcal{Y}} \mathcal{X},\ \text{pr}^{-1}\mathcal{F})
$$
d'après Faisceaux sur les champs, équation
(\ref{stacks-sheaves-equation-pushforward}). Soit
$(V \times_{y, \mathcal{Y}} \mathcal{X})' \subset
V \times_{y, \mathcal{Y}} \mathcal{X}$
la sous-catégorie pleine formée des objets
$z : W \to V \times_{y, \mathcal{Y}} \mathcal{X}$ tels que le morphisme
induit $W \to \mathcal{X}$ soit lisse (resp.\ plat). Notons
$$
\text{pr}' :
(V \times_{y, \mathcal{Y}} \mathcal{X})'
\longrightarrow
\mathcal{X}_{lisse,\etale}
\ (\text{resp. }\mathcal{X}_{flat,fppf})
$$
la restriction du foncteur $\text{pr}$ utilisé dans la formule ci-dessus.
Exactement le même argument que celui qui démontre
Faisceaux sur les champs, équation (\ref{stacks-sheaves-equation-pushforward}),
montre que, pour tout faisceau $\mathcal{H}$ sur $\mathcal{X}_{lisse,\etale}$
(resp.\ $\mathcal{X}_{flat,fppf}$), on a
\begin{equation}
\label{stacks-cohomology-equation-pushforward-lisse-etale}
f'_*\mathcal{H}(y) =
\Gamma((V \times_{y, \mathcal{Y}} \mathcal{X})',
\ (\text{pr}')^{-1}\mathcal{H})
\end{equation}
Comme $(g')^{-1}$ est le foncteur de restriction, on voit que
$$
\left(f'_*(g')^{-1}\mathcal{F}\right)(y) =
\Gamma((V \times_{y, \mathcal{Y}} \mathcal{X})',
\ \text{pr}^{-1}\mathcal{F}|_{(V \times_{y, \mathcal{Y}} \mathcal{X})'})
$$
D'après le
lemme \ref{stacks-sheaves-lemma-cohomology-on-subcategory} de Faisceaux sur les champs,
on voit que
$$
\Gamma((V \times_{y, \mathcal{Y}} \mathcal{X})',
\ \text{pr}^{-1}\mathcal{F}|_{(V \times_{y, \mathcal{Y}} \mathcal{X})'})
=
\Gamma(V \times_{y, \mathcal{Y}} \mathcal{X},\ \text{pr}^{-1}\mathcal{F})
$$
sont égaux, comme voulu ; bien que nous omettions la vérification des hypothèses
du lemme, notons que le fait que $V \to \mathcal{Y}$ soit lisse
(resp.\ plat) sert à vérifier la seconde condition.

\medskip\noindent
Enfin, l'égalité $g'_!(f')^{-1} = f^{-1}g_!$ découle formellement de
l'égalité $f'_*(g')^{-1} = g^{-1}f_*$ par les adjonctions entre
$f^{-1}$ et $f_*$, entre $g_!$ et $g^{-1}$, et par leurs versions
« primées ».
\end{proof}

\begin{lemma}
\label{stacks-cohomology-lemma-lisse-etale-functorial-pushforward}
Sous les hypothèses et avec les notations du lemme
\ref{stacks-cohomology-lemma-lisse-etale-functorial}, soit $\mathcal{H}$ un faisceau abélien sur
$\mathcal{X}_{lisse,\etale}$ (resp.\ $\mathcal{X}_{flat,fppf}$). Alors
\begin{equation}
\label{stacks-cohomology-equation-higher-direct-image-lisse-etale}
R^pf'_*\mathcal{H} =
\text{faisceau associé au préfaisceau }y \longmapsto
H^p((V \times_{y, \mathcal{Y}} \mathcal{X})', (\text{pr}')^{-1}\mathcal{H})
\end{equation}
Ici, $y$ est un objet de $\mathcal{Y}_{lisse,\etale}$
(resp.\ $\mathcal{Y}_{flat,fppf}$) au-dessus du schéma $V$, et les notations
$(V \times_{y, \mathcal{Y}} \mathcal{X})'$ et $\text{pr}'$ sont
expliquées dans la démonstration.
\end{lemma}

\begin{proof}
Comme dans la démonstration du lemme \ref{stacks-cohomology-lemma-lisse-etale-functorial},
soit $(V \times_{y, \mathcal{Y}} \mathcal{X})' \subset
V \times_{y, \mathcal{Y}} \mathcal{X}$ la sous-catégorie pleine
formée des objets $(x, \varphi)$, où $x$ est un objet de
$\mathcal{X}_{lisse,\etale}$ (resp.\ $\mathcal{X}_{flat,fppf}$)
et $\varphi : f(x) \to y$ est un morphisme dans $\mathcal{Y}$.
D'après l'équation (\ref{stacks-cohomology-equation-pushforward-lisse-etale}),
on a
$$
f'_*\mathcal{H}(y) =
\Gamma((V \times_{y, \mathcal{Y}} \mathcal{X})',
\ (\text{pr}')^{-1}\mathcal{H})
$$
où $\text{pr}'$ est la projection. Pour un objet $(x, \varphi)$
de $(V \times_{y, \mathcal{Y}} \mathcal{X})'$,
on peut considérer $\varphi$ comme une section de $(f')^{-1}h_y$ au-dessus de $x$.
Ainsi, $(V \times_\mathcal{Y} \mathcal{X})'$ est la localisation
du site $\mathcal{X}_{lisse,\etale}$
(resp. $\mathcal{X}_{flat,fppf}$) en le faisceau d'ensembles $(f')^{-1}h_y$ ; voir
Sites, lemme \ref{sites-lemma-localize-topos-site}. Le morphisme
$$
\text{pr}' : (V \times_{y, \mathcal{Y}} \mathcal{X})'
\to \mathcal{X}_{lisse,\etale}
\ (\text{resp. }
\text{pr}' : (V \times_{y, \mathcal{Y}} \mathcal{X})'
\to \mathcal{X}_{flat,fppf})
$$
est le morphisme de localisation.
En particulier, le foncteur image inverse $(\text{pr}')^{-1}$ préserve
les faisceaux abéliens injectifs ; voir
Cohomologie sur les sites, lemme
\ref{sites-cohomology-lemma-cohomology-on-sheaf-sets}.

\medskip\noindent
Choisissons une résolution injective $\mathcal{H} \to \mathcal{I}^\bullet$
sur $\mathcal{X}_{lisse,\etale}$ (resp.\ $\mathcal{X}_{flat,fppf}$).
La formule de l'image directe montre que $R^if'_*\mathcal{H}$ est le
faisceau associé au préfaisceau qui, à $y$, associe la cohomologie
du complexe
$$
\begin{matrix}
\Gamma\Big((V \times_{y, \mathcal{Y}} \mathcal{X})',
(\text{pr}')^{-1}\mathcal{I}^{i - 1}\Big) \\
\downarrow \\
\Gamma\Big((V \times_{y, \mathcal{Y}} \mathcal{X})',
(\text{pr}')^{-1}\mathcal{I}^i\Big) \\
\downarrow \\
\Gamma\Big((V \times_{y, \mathcal{Y}} \mathcal{X})',
(\text{pr}')^{-1}\mathcal{I}^{i + 1}\Big)
\end{matrix}
$$
Comme $(\text{pr}')^{-1}$ est exact et préserve les objets injectifs, le
complexe $(\text{pr}')^{-1}\mathcal{I}^\bullet$ est une
résolution injective de $(\text{pr}')^{-1}\mathcal{H}$,
ce qui achève la démonstration.
\end{proof}

\begin{lemma}
\label{stacks-cohomology-lemma-lisse-etale-functorial-cohomology}
Sous les hypothèses et avec les notations du lemme
\ref{stacks-cohomology-lemma-lisse-etale-functorial}, le morphisme canonique de changement de base
$$
g^{-1}Rf_*\mathcal{F} \longrightarrow Rf'_*(g')^{-1}\mathcal{F}
$$
est un isomorphisme pour tout faisceau abélien $\mathcal{F}$
sur $\mathcal{X}_\etale$ (resp.\ $\mathcal{X}_{fppf}$).
\end{lemma}

\begin{proof}
En comparant les formules pour $g^{-1}R^pf_*\mathcal{F}$
et $R^pf'_*(g')^{-1}\mathcal{F}$ données par le
lemme \ref{stacks-sheaves-lemma-pushforward-restriction} de Faisceaux sur les champs
et le lemme \ref{stacks-cohomology-lemma-lisse-etale-functorial-pushforward},
on voit qu'il suffit de montrer
$$
H^p((V \times_{y, \mathcal{Y}} \mathcal{X})',
\ \text{pr}^{-1}\mathcal{F}|_{(V \times_{y, \mathcal{Y}} \mathcal{X})'})
=
H^p_\tau(V \times_{y, \mathcal{Y}} \mathcal{X},\ \text{pr}^{-1}\mathcal{F})
$$
où $\tau = \etale$ (resp.\ $\tau = fppf$). Ici, $y$ est un objet
de $\mathcal{Y}$ au-dessus d'un schéma $V$ tel que le morphisme
$y : V \to \mathcal{Y}$ soit lisse (resp.\ plat).
Cette égalité résulte du
lemme \ref{stacks-sheaves-lemma-cohomology-on-subcategory} de Faisceaux sur les champs.
Bien que nous omettions la vérification des hypothèses
du lemme, notons que le fait que $V \to \mathcal{Y}$ soit lisse
(resp.\ plat) sert à vérifier la seconde condition.
\end{proof}






\section{Modules quasi-cohérents et sites lisse-étale et plat-fppf}
\label{stacks-cohomology-section-quasi-coherent-modules-II}

\noindent
Dans cette section, nous expliquons comment décrire les modules quasi-cohérents
sur un champ algébrique au moyen de son site lisse-étale ou plat-fppf.

\begin{lemma}
\label{stacks-cohomology-lemma-check-qc-on-etale-covering}
Soit $\mathcal{X}$ un champ algébrique.
\begin{enumerate}
\item Soit $f_j : \mathcal{X}_j \to \mathcal{X}$ une famille de morphismes
lisses de champs algébriques telle que
$|\mathcal{X}| =\bigcup |f_j|(|\mathcal{X}_j|)$.
Soit $\mathcal{F}$ un faisceau de $\mathcal{O}_\mathcal{X}$-modules
sur $\mathcal{X}_\etale$. Si chaque $f_j^{-1}\mathcal{F}$ est
quasi-cohérent, alors $\mathcal{F}$ l'est aussi.
\item Soit $f_j : \mathcal{X}_j \to \mathcal{X}$ une famille de morphismes
plats et localement de présentation finie de champs algébriques telle que
$|\mathcal{X}| =\bigcup |f_j|(|\mathcal{X}_j|)$.
Soit $\mathcal{F}$ un faisceau de $\mathcal{O}_\mathcal{X}$-modules
sur $\mathcal{X}_{fppf}$. Si chaque $f_j^{-1}\mathcal{F}$ est
quasi-cohérent, alors $\mathcal{F}$ l'est aussi.
\end{enumerate}
\end{lemma}

\begin{proof}
Preuve de (1). Nous pouvons remplacer chacun des champs algébriques $\mathcal{X}_j$
par un schéma $U_j$ (puisque tout champ algébrique admet un recouvrement lisse par
un schéma et que les composés de morphismes lisses sont lisses ; voir
Morphismes de champs, lemme \ref{stacks-morphisms-lemma-composition-smooth}).
L'image inverse de $\mathcal{F}$ sur $(\Sch/U_j)_\etale$ est encore
quasi-cohérente ; voir
Modules sur les sites, lemme \ref{sites-modules-lemma-local-pullback}.
Alors $f = \coprod f_j : U = \coprod U_j \to \mathcal{X}$ est un morphisme lisse
surjectif. Soit $x : V \to \mathcal{X}$ un objet de $\mathcal{X}$. D'après
Faisceaux sur les champs, lemme
\ref{stacks-sheaves-lemma-surjective-flat-locally-finite-presentation}
il existe un recouvrement étale $\{x_i \to x\}_{i \in I}$ tel que chaque
$x_i$ se relève en un objet $u_i$ de $(\Sch/U)_\etale$.
Cela signifie simplement que $x_i$ est au-dessus d'un schéma $V_i$, que
$\{V_i \to V\}$ est un recouvrement étale et que $x_i$ provient
d'un morphisme $u_i : V_i \to U$. Alors
$x_i^*\mathcal{F} = u_i^*f^*\mathcal{F}$ est quasi-cohérent.
Il s'ensuit que $x^*\mathcal{F}$ sur $(\Sch/V)_\etale$
est quasi-cohérent, par exemple d'après
Modules sur les sites, lemme \ref{sites-modules-lemma-local-final-object}.
D'après Faisceaux sur les champs, lemme
\ref{stacks-sheaves-lemma-characterize-quasi-coherent-bis}
on voit que $x^*\mathcal{F}$ est un faisceau fppf et, puisque $x$
était arbitraire, que $\mathcal{F}$ est un faisceau pour la
topologie fppf. En appliquant Faisceaux sur les champs, lemme
\ref{stacks-sheaves-lemma-characterize-quasi-coherent}
on voit que $\mathcal{F}$ est quasi-cohérent.

\medskip\noindent
Preuve de (2). La démonstration emploie exactement le même argument, que nous
détaillons ici. Nous pouvons remplacer chacun des champs algébriques $\mathcal{X}_j$
par un schéma $U_j$ (puisque tout champ algébrique admet un recouvrement lisse par
un schéma et que les morphismes plats et localement de présentation finie sont
stables par composition ; voir Morphismes de champs, lemmes
\ref{stacks-morphisms-lemma-composition-flat} et
\ref{stacks-morphisms-lemma-composition-finite-presentation}).
L'image inverse de $\mathcal{F}$ sur $(\Sch/U_j)_\etale$ est encore
localement quasi-cohérente ; voir
Faisceaux sur les champs, lemme \ref{stacks-sheaves-lemma-pullback-quasi-coherent}.
Alors $f = \coprod f_j : U = \coprod U_j \to \mathcal{X}$ est un morphisme
surjectif, plat et localement de présentation finie. Soit
$x : V \to \mathcal{X}$ un objet de $\mathcal{X}$. D'après
Faisceaux sur les champs, lemme
\ref{stacks-sheaves-lemma-surjective-flat-locally-finite-presentation}
il existe un recouvrement fppf $\{x_i \to x\}_{i \in I}$ tel que chaque
$x_i$ se relève en un objet $u_i$ de $(\Sch/U)_\etale$.
Cela signifie simplement que $x_i$ est au-dessus d'un schéma $V_i$, que
$\{V_i \to V\}$ est un recouvrement fppf et que $x_i$ provient
d'un morphisme $u_i : V_i \to U$. Alors
$x_i^*\mathcal{F} = u_i^*f^*\mathcal{F}$ est quasi-cohérent.
Il s'ensuit que $x^*\mathcal{F}$ sur $(\Sch/V)_\etale$
est quasi-cohérent, par exemple d'après
Modules sur les sites, lemme \ref{sites-modules-lemma-local-final-object}.
D'après Faisceaux sur les champs, lemme
\ref{stacks-sheaves-lemma-characterize-quasi-coherent}
on voit que $\mathcal{F}$ est quasi-cohérent.
\end{proof}

\noindent
Rappelons que nous avons défini la notion de module quasi-cohérent sur
un topos annelé quelconque dans
Modules sur les sites, section \ref{sites-modules-section-local}.

\begin{lemma}
\label{stacks-cohomology-lemma-shriek-quasi-coherent}
Soit $\mathcal{X}$ un champ algébrique. Les notations sont celles du
lemme \ref{stacks-cohomology-lemma-lisse-etale}.
\begin{enumerate}
\item Soit $\mathcal{H}$ un
$\mathcal{O}_{\mathcal{X}_{lisse,\etale}}$-module quasi-cohérent
sur le site lisse-étale de $\mathcal{X}$. Alors $g_!\mathcal{H}$ est un
module quasi-cohérent sur $\mathcal{X}$.
\item Soit $\mathcal{H}$ un
$\mathcal{O}_{\mathcal{X}_{flat,fppf}}$-module quasi-cohérent
sur le site plat-fppf de $\mathcal{X}$. Alors $g_!\mathcal{H}$ est un
module quasi-cohérent sur $\mathcal{X}$.
\end{enumerate}
\end{lemma}

\begin{proof}
Choisissons un schéma $U$ et un morphisme lisse surjectif $x : U \to \mathcal{X}$.
D'après
Modules sur les sites, définition \ref{sites-modules-definition-site-local},
il existe un recouvrement étale (resp.\ fppf)
$\{U_i \to U\}_{i \in I}$ tel que chaque image inverse $f_i^{-1}\mathcal{H}$
admette une présentation globale (voir
Modules sur les sites, définition \ref{sites-modules-definition-global}).
Ici, $f_i : U_i \to \mathcal{X}$ est le composé
$U_i \to U \to \mathcal{X}$, qui est un morphisme de champs algébriques.
(Rappelons que l'image inverse « est » la restriction à $\mathcal{X}/f_i$ ; voir
Faisceaux sur les champs, définition \ref{stacks-sheaves-definition-pullback}
et la discussion qui suit.) Comme chaque $f_i$ est lisse (resp.\ plat), le
lemme \ref{stacks-cohomology-lemma-lisse-etale-functorial}
donne $f_i^{-1}g_!\mathcal{H} = g_{i, !}(f'_i)^{-1}\mathcal{H}$.
Le lemme \ref{stacks-cohomology-lemma-check-qc-on-etale-covering}
ramène l'énoncé au cas où $\mathcal{H}$
admet une présentation globale. Écrivons
$$
\bigoplus\nolimits_{j \in J} \mathcal{O} \longrightarrow
\bigoplus\nolimits_{i \in I} \mathcal{O} \longrightarrow
\mathcal{H} \longrightarrow 0
$$
de $\mathcal{O}$-modules, où
$\mathcal{O} = \mathcal{O}_{\mathcal{X}_{lisse,\etale}}$
(resp.\ $\mathcal{O} = \mathcal{O}_{\mathcal{X}_{flat,fppf}}$).
Comme $g_!$ commute aux limites inductives arbitraires (en tant que foncteur
adjoint à gauche ; voir le lemme \ref{stacks-cohomology-lemma-lisse-etale-modules} et
Catégories, lemme \ref{categories-lemma-adjoint-exact}),
on en déduit l'existence d'une suite exacte
$$
\bigoplus\nolimits_{j \in J} g_!\mathcal{O} \longrightarrow
\bigoplus\nolimits_{i \in I} g_!\mathcal{O} \longrightarrow
g_!\mathcal{H} \longrightarrow 0
$$
Le lemme \ref{stacks-cohomology-lemma-lisse-etale-structure-sheaf}
montre que $g_!\mathcal{O} = \mathcal{O}_\mathcal{X}$.
Dans le cas (2), la démonstration est terminée. Dans le cas (1), on applique
Faisceaux sur les champs, lemme
\ref{stacks-sheaves-lemma-characterize-quasi-coherent-bis}
pour conclure.
\end{proof}

\begin{lemma}
\label{stacks-cohomology-lemma-quasi-coherent}
Soit $\mathcal{X}$ un champ algébrique.
\begin{enumerate}
\item Pour le site lisse-étale, avec $g$ comme dans le
lemme \ref{stacks-cohomology-lemma-lisse-etale}, on a :
\begin{enumerate}
\item les foncteurs $g^{-1}$ et $g_!$ définissent des foncteurs quasi-inverses
$$
\xymatrix{
\QCoh(\mathcal{O}_\mathcal{X}) \ar@<1ex>[r]^-{g^{-1}} &
\QCoh(\mathcal{X}_{lisse,\etale},
\mathcal{O}_{\mathcal{X}_{lisse,\etale}}) \ar@<1ex>[l]^-{g_!}
}
$$
\item si $\mathcal{F}$ appartient à $\textit{LQCoh}^{fbc}(\mathcal{O}_\mathcal{X})$,
alors $g^{-1}\mathcal{F}$ appartient à
$\QCoh(\mathcal{O}_{\mathcal{X}_{lisse,\etale}})$ ;
\item $Q(\mathcal{F}) = g_!g^{-1}\mathcal{F}$, où $Q$ est celui du
lemme \ref{stacks-cohomology-lemma-adjoint}.
\end{enumerate}
\item Pour le site plat-fppf, avec $g$ comme dans le
lemme \ref{stacks-cohomology-lemma-lisse-etale}, on a :
\begin{enumerate}
\item les foncteurs $g^{-1}$ et $g_!$ définissent des foncteurs quasi-inverses
$$
\xymatrix{
\QCoh(\mathcal{O}_\mathcal{X}) \ar@<1ex>[r]^-{g^{-1}} &
\QCoh(\mathcal{X}_{flat,fppf},
\mathcal{O}_{\mathcal{X}_{flat,fppf}}) \ar@<1ex>[l]^-{g_!}
}
$$
\item si $\mathcal{F}$ appartient à $\textit{LQCoh}^{fbc}(\mathcal{O}_\mathcal{X})$,
alors $g^{-1}\mathcal{F}$ appartient à
$\QCoh(\mathcal{O}_{\mathcal{X}_{flat,fppf}})$
et
\item $Q(\mathcal{F}) = g_!g^{-1}\mathcal{F}$, où $Q$ est celui du
lemme \ref{stacks-cohomology-lemma-adjoint}.
\end{enumerate}
\end{enumerate}
\end{lemma}

\begin{proof}
L'image inverse par tout morphisme de topos annelés préserve les catégories de
modules quasi-cohérents ; voir
Modules sur les sites, lemme \ref{sites-modules-lemma-local-pullback}.
Par conséquent, $g^{-1}$ préserve les catégories de modules quasi-cohérents ;
nous utilisons ici l'égalité
$\QCoh(\mathcal{O}_\mathcal{X}) =
\QCoh(\mathcal{X}_\etale, \mathcal{O}_\mathcal{X})$
fournie par Faisceaux sur les champs, lemme
\ref{stacks-sheaves-lemma-characterize-quasi-coherent-bis}.
Il en va de même pour $g_!$ d'après le
lemme \ref{stacks-cohomology-lemma-shriek-quasi-coherent}.
Le lemme \ref{stacks-cohomology-lemma-lisse-etale} montre que
$\mathcal{H} \to g^{-1}g_!\mathcal{H}$ est un isomorphisme.
Réciproquement, si $\mathcal{F}$ appartient à $\QCoh(\mathcal{O}_\mathcal{X})$,
alors le morphisme $g_!g^{-1}\mathcal{F} \to \mathcal{F}$ est un morphisme de
modules quasi-cohérents sur $\mathcal{X}$ dont la restriction à tout schéma lisse sur
$\mathcal{X}$ est un isomorphisme. La discussion de
Faisceaux sur les champs, sections
\ref{stacks-sheaves-section-quasi-coherent-presentation} et
\ref{stacks-sheaves-section-quasi-coherent-algebraic-stacks}
(par comparaison avec les modules quasi-cohérents sur des présentations),
montre donc qu'il s'agit d'un isomorphisme. Cela démontre (1)(a) et (2)(a).

\medskip\noindent
Soit $\mathcal{F}$ un objet de
$\textit{LQCoh}^{fbc}(\mathcal{O}_\mathcal{X})$. D'après le
lemme \ref{stacks-cohomology-lemma-adjoint-kernel-parasitic},
le noyau et le conoyau du morphisme
$Q(\mathcal{F}) \to \mathcal{F}$ sont parasites. Par conséquent, d'après le
lemme \ref{stacks-cohomology-lemma-parasitic-in-terms-flat-fppf}
et puisque $g^* = g^{-1}$ est exact, on conclut que
$g^*Q(\mathcal{F}) \to g^*\mathcal{F}$ est un isomorphisme. Ainsi,
$g^*\mathcal{F}$ est quasi-cohérent. Cela démontre (1)(b) et (2)(b).
Enfin, (1)(c) et (2)(c) résultent du fait que
$g_!g^*Q(\mathcal{F}) \to Q(\mathcal{F})$ est un isomorphisme d'après
les arguments qui précèdent.
\end{proof}

\begin{lemma}
\label{stacks-cohomology-lemma-quasi-coherent-weak-serre}
Soit $\mathcal{X}$ un champ algébrique.
\begin{enumerate}
\item $\QCoh(\mathcal{O}_{\mathcal{X}_{lisse,\etale}})$
est une sous-catégorie faible de Serre de
$\textit{Mod}(\mathcal{O}_{\mathcal{X}_{lisse,\etale}})$.
\item $\QCoh(\mathcal{O}_{\mathcal{X}_{flat,fppf}})$
est une sous-catégorie faible de Serre de
$\textit{Mod}(\mathcal{O}_{\mathcal{X}_{flat,fppf}})$.
\end{enumerate}
\end{lemma}

\begin{proof}
Nous allons vérifier les conditions (1), (2), (3) et (4) du
lemme \ref{homology-lemma-characterize-weak-serre-subcategory} d'Homologie.

\medskip\noindent
Comme $0$ est un module quasi-cohérent sur tout site annelé, la condition (1)
est satisfaite.

\medskip\noindent
Par définition, $\QCoh(\mathcal{O})$
est une sous-catégorie strictement pleine de $\textit{Mod}(\mathcal{O})$ ; donc (2) est satisfaite.

\medskip\noindent
Soit $\varphi : \mathcal{G} \to \mathcal{F}$ un morphisme de modules
quasi-cohérents sur $\mathcal{X}_{lisse,\etale}$ ou $\mathcal{X}_{flat,fppf}$.
On a $g^*g_!\mathcal{F} = \mathcal{F}$ et de même pour
$\mathcal{G}$ et $\varphi$ ; voir le lemme \ref{stacks-cohomology-lemma-lisse-etale-modules}.
D'après le lemme \ref{stacks-cohomology-lemma-shriek-quasi-coherent},
$g_!\mathcal{F}$ et $g_!\mathcal{G}$ sont des
$\mathcal{O}_\mathcal{X}$-modules quasi-cohérents. D'après Faisceaux sur les champs, lemme
\ref{stacks-sheaves-lemma-quasi-coherent-algebraic-stack}
le module $\Coker(g_!\varphi)$ est quasi-cohérent
sur $\mathcal{X}$ (et c'est le conoyau dans la catégorie
des modules quasi-cohérents sur $\mathcal{X}$).
Comme $g^*$ est exact (voir le lemme \ref{stacks-cohomology-lemma-lisse-etale}),
$g^*\Coker(g_!\varphi) = \Coker(g^*g_!\varphi) = \Coker(\varphi)$
est également quasi-cohérent (voir le lemme \ref{stacks-cohomology-lemma-quasi-coherent}).
D'après la proposition \ref{stacks-cohomology-proposition-loc-qcoh-flat-base-change},
le noyau $\Ker(g_!\varphi)$ appartient à
$\textit{LQCoh}^{fbc}(\mathcal{O}_\mathcal{X})$.
Comme $g^*$ est exact, on a
$g^*\Ker(g_!\varphi) = \Ker(g^*g_!\varphi) = \Ker(\varphi)$.
Puisque $g^*$ envoie les objets de
$\textit{LQCoh}^{fbc}(\mathcal{O}_\mathcal{X})$ sur des modules
quasi-cohérents d'après le lemme \ref{stacks-cohomology-lemma-quasi-coherent}, on en déduit que
$\Ker(\varphi)$ est également quasi-cohérent. Cela démontre (3).

\medskip\noindent
Enfin, supposons que
$$
0 \to \mathcal{F} \to \mathcal{E} \to \mathcal{G} \to 0
$$
soit une extension de $\mathcal{O}_{\mathcal{X}_{lisse,\etale}}$-modules
(resp.\ de $\mathcal{O}_{\mathcal{X}_{flat,fppf}}$-modules), où $\mathcal{F}$
et $\mathcal{G}$ sont quasi-cohérents. Pour démontrer (4) et achever la preuve,
il nous faut montrer que
$\mathcal{E}$ est quasi-cohérent sur $\mathcal{X}_{lisse,\etale}$
(resp.\ $\mathcal{X}_{flat,fppf}$). Soit $U$
un objet de $\mathcal{X}_{lisse,\etale}$
(resp.\ $\mathcal{X}_{flat,fppf}$) ; nous considérons $U$ comme un schéma lisse
(resp.\ plat) sur $\mathcal{X}$. Il faut montrer que la restriction
de $\mathcal{E}$ à $U_{lisse,\etale}$ (resp.\ $=U_{flat,fppf}$)
est quasi-cohérente. Nous pouvons donc supposer
que $\mathcal{X} = U$ est un schéma. Puisque $\mathcal{G}$
est quasi-cohérent sur $U_{lisse,\etale}$ (resp.\ $U_{flat,fppf}$),
nous pouvons supposer, après avoir remplacé $U$ par les membres
d'un recouvrement étale (resp.\ fppf), que $\mathcal{G}$ admet une
présentation
$$
\bigoplus\nolimits_{j \in J} \mathcal{O} \longrightarrow
\bigoplus\nolimits_{i \in I} \mathcal{O} \longrightarrow
\mathcal{G} \longrightarrow 0
$$
sur $U_{lisse,\etale}$ (resp.\ $U_{flat,fppf}$), où $\mathcal{O}$
est le faisceau structural du site. Nous pouvons aussi supposer $U$ affine.
Comme $\mathcal{F}$ est quasi-cohérent, on a
$$
H^1(U_{lisse,\etale}, \mathcal{F}) = 0,
\quad\text{resp.}\quad
H^1(U_{flat,fppf}, \mathcal{F}) = 0
$$
En effet, $\mathcal{F}$ est l'image inverse d'un module quasi-cohérent
$\mathcal{F}'$ sur le grand site de $U$
(d'après le lemme \ref{stacks-cohomology-lemma-quasi-coherent}), les cohomologies
de $\mathcal{F}$ et $\mathcal{F}'$ coïncident (d'après le
lemme \ref{stacks-cohomology-lemma-lisse-etale-cohomology}), et l'on sait que
la cohomologie de $\mathcal{F}'$ sur le grand site du schéma affine $U$
est nulle (pour l'obtenir dans la situation présente, il faut
combiner Descentee, propositions
\ref{descent-proposition-equivalence-quasi-coherent} et
\ref{descent-proposition-same-cohomology-quasi-coherent} avec
Cohomologie des schémas, lemme
\ref{coherent-lemma-quasi-coherent-affine-cohomology-zero}).
Nous pouvons donc relever le morphisme
$\bigoplus_{i \in I} \mathcal{O} \to \mathcal{G}$
à $\mathcal{E}$. Une chasse au diagramme donne alors
une suite exacte
$$
\bigoplus\nolimits_{j \in J} \mathcal{O} \to
\mathcal{F} \oplus \bigoplus\nolimits_{i \in I} \mathcal{O}
\to
\mathcal{E} \to 0
$$
D'après (3), démontré ci-dessus, on en conclut que $\mathcal{E}$ est
quasi-cohérent, comme voulu.
\end{proof}









\section{Faisceaux cohérents sur les champs localement noethériens}
\sectionmark{Faisceaux cohérents sur les champs noethériens}
\label{stacks-cohomology-section-coherent-sheaves}

\noindent
Cette section est l'analogue de
Cohomologie des espaces, section
\ref{spaces-cohomology-section-coherent}.
Nous avons défini la notion de module cohérent sur tout topos annelé dans
Modules sur les sites, section \ref{sites-modules-section-local}.
Cependant, pour tout champ algébrique $\mathcal{X}$, la catégorie des
$\mathcal{O}_\mathcal{X}$-modules cohérents est nulle, essentiellement parce que
le site $\mathcal{X}$ contient trop d'objets non noethériens (même si
$\mathcal{X}$ est lui-même localement noethérien). Nous définirons plutôt les
modules cohérents au moyen du lemme suivant.

\begin{lemma}
\label{stacks-cohomology-lemma-coherent-Noetherian}
Soit $\mathcal{X}$ un champ algébrique localement noethérien.
Soit $\mathcal{F}$ un $\mathcal{O}_\mathcal{X}$-module.
Les conditions suivantes sont équivalentes :
\begin{enumerate}
\item $\mathcal{F}$ est un $\mathcal{O}_\mathcal{X}$-module
quasi-cohérent de type fini ;
\item $\mathcal{F}$ est un $\mathcal{O}_\mathcal{X}$-module
de présentation finie ;
\item $\mathcal{F}$ est quasi-cohérent et, pour tout morphisme
$f : U \to \mathcal{X}$ où $U$ est un espace algébrique localement noethérien,
l'image inverse $f^*\mathcal{F}|_{U_\etale}$ est cohérente ;
\item $\mathcal{F}$ est quasi-cohérent et il existe un espace algébrique
$U$ et un morphisme $f : U \to \mathcal{X}$ localement de type fini,
plat et surjectif, tels que l'image inverse $f^*\mathcal{F}|_{U_\etale}$
soit cohérente.
\end{enumerate}
\end{lemma}

\begin{proof}
Soit $f : U \to \mathcal{X}$ comme dans (4).
Alors $U$ est localement noethérien (Morphismes de champs, lemme
\ref{stacks-morphisms-lemma-locally-finite-type-locally-noetherian})
et l'on voit que l'énoncé du lemme a un sens. De plus,
$f$ est localement de présentation finie d'après Morphismes de champs, lemme
\ref{stacks-morphisms-lemma-noetherian-finite-type-finite-presentation}.
Soit $x$ un objet de $\mathcal{X}$ au-dessus du schéma $V$.
Pour démontrer (2), il faut montrer qu'après avoir remplacé $V$
par les membres d'un recouvrement fppf de $V$, la restriction $x^*\mathcal{F}$
admet une présentation globale finie sur $\mathcal{X}/x \cong (\Sch/V)_{fppf}$.
La projection $W = U \times_\mathcal{X} V \to V$ est
localement de présentation finie, plate et surjective.
On peut donc remplacer $V$ par les membres d'un recouvrement étale de $W$
par des schémas et supposer qu'il existe un morphisme $h : V \to U$ tel que $f \circ h = x$.
Comme $\mathcal{F}$ est quasi-cohérent, on voit que
la restriction $x^*\mathcal{F}$
est l'image inverse de $h_{small}^*(f^*\mathcal{F})|_{U_\etale}$
par $\pi_V$ ; voir Faisceaux sur les champs, lemme
\ref{stacks-sheaves-lemma-compare-quasi-coherent}.
Puisque, par hypothèse, $f^*\mathcal{F}|_{U_\etale}$ admet localement pour la
topologie étale une présentation finie, on en conclut
(4) $\Rightarrow$ (2).

\medskip\noindent
La condition (2) implique (1) pour tout topos annelé (cela résulte aussitôt de la définition).
Les propriétés « de type fini » et « quasi-cohérent »
sont préservées par l'image inverse de tout morphisme de topos annelés ; voir
Modules sur les sites, lemme \ref{sites-modules-lemma-local-pullback}.
Par conséquent, (1) implique (3) ; voir Cohomologie des espaces, lemme
\ref{spaces-cohomology-lemma-coherent-Noetherian}.
Enfin, (3) implique trivialement (4).
\end{proof}

\begin{definition}
\label{stacks-cohomology-definition-coherent}
Soit $\mathcal{X}$ un champ algébrique localement noethérien.
Un $\mathcal{O}_\mathcal{X}$-module $\mathcal{F}$ est dit {\it cohérent}
si $\mathcal{F}$ satisfait l'une (et donc chacune) des conditions
équivalentes du lemme \ref{stacks-cohomology-lemma-coherent-Noetherian}.
La catégorie des $\mathcal{O}_\mathcal{X}$-modules cohérents est
notée $\textit{Coh}(\mathcal{O}_\mathcal{X})$.
\end{definition}

\begin{lemma}
\label{stacks-cohomology-lemma-elementary-coherent}
Soit $\mathcal{X}$ un champ algébrique localement noethérien.
Le module $\mathcal{O}_\mathcal{X}$ est cohérent, tout
$\mathcal{O}_\mathcal{X}$-module inversible est cohérent et, plus généralement, tout
$\mathcal{O}_\mathcal{X}$-module localement libre de type fini est cohérent.
\end{lemma}

\begin{proof}
Cela résulte de la définition et de
Cohomologie des espaces, lemme
\ref{spaces-cohomology-lemma-coherent-Noetherian}.
\end{proof}

\begin{lemma}
\label{stacks-cohomology-lemma-pullback-coherent}
Soit $f : \mathcal{X} \to \mathcal{Y}$ un morphisme de champs algébriques
localement noethériens. Alors $f^*$ envoie les modules cohérents
sur $\mathcal{Y}$ sur des modules cohérents sur $\mathcal{X}$.
\end{lemma}

\begin{proof}
Cela résulte aussitôt de la définition et du fait que l'image inverse par tout
morphisme de topos annelés préserve les modules de présentation finie ; voir
Modules sur les sites, lemme \ref{sites-modules-lemma-local-pullback}.
\end{proof}

\begin{lemma}
\label{stacks-cohomology-lemma-coherent-abelian-Noetherian}
Soit $\mathcal{X}$ un champ algébrique localement noethérien.
La catégorie des $\mathcal{O}_\mathcal{X}$-modules cohérents est abélienne.
Si $\varphi : \mathcal{F} \to \mathcal{G}$ est un morphisme
de $\mathcal{O}_\mathcal{X}$-modules cohérents, alors :
\begin{enumerate}
\item le conoyau $\Coker(\varphi)$ calculé dans
$\textit{Mod}(\mathcal{O}_\mathcal{X})$ est un
$\mathcal{O}_\mathcal{X}$-module cohérent ;
\item l'image $\Im(\varphi)$ calculée dans
$\textit{Mod}(\mathcal{O}_\mathcal{X})$ est un
$\mathcal{O}_\mathcal{X}$-module cohérent ;
\item le noyau $\Ker(\varphi)$ calculé dans
$\textit{Mod}(\mathcal{O}_\mathcal{X})$
peut ne pas être cohérent, mais il appartient à
$\textit{LQCoh}^{fbc}(\mathcal{O}_\mathcal{X})$ et $Q(\Ker(\varphi))$
est cohérent et constitue le noyau de $\varphi$ dans
$\textit{Coh}(\mathcal{O}_\mathcal{X})$.
\end{enumerate}
Le foncteur d'inclusion $\textit{Coh}(\mathcal{O}_\mathcal{X}) \to
\QCoh(\mathcal{O}_\mathcal{X})$ est exact.
\end{lemma}

\begin{proof}
Les règles données pour former les noyaux, images et conoyaux dans
$\textit{Coh}(\mathcal{O}_\mathcal{X})$ coïncident avec la prescription
pour les modules quasi-cohérents de la remarque \ref{stacks-cohomology-remark-QCoh-abelian}.
Le lemme résultera donc du fait que les modules quasi-cohérents
$\Coker(\varphi)$, $\Im(\varphi)$ et $Q(\Ker(\varphi))$ sont cohérents.
D'après le lemme \ref{stacks-cohomology-lemma-coherent-Noetherian}, il suffit de le montrer après
restriction à $U_\etale$ pour un morphisme lisse surjectif
$f : U \to \mathcal{X}$. Le foncteur
$\mathcal{F} \mapsto f^*\mathcal{F}|_{U_\etale}$
est exact. Par conséquent, $f^*\Coker(\varphi)$ et $f^*\Im(\varphi)$ sont
le conoyau et l'image d'un morphisme entre $\mathcal{O}_U$-modules cohérents,
donc sont cohérents comme voulu. Le foncteur
$\mathcal{F} \mapsto f^*\mathcal{F}|_{U_\etale}$
annule les modules parasites d'après le lemme \ref{stacks-cohomology-lemma-parasitic}.
Ainsi, $f^*Q(\Ker(\varphi))|_{U_\etale} = f^*\Ker(\varphi)|_{U_\etale}$
d'après la partie (2) du lemme \ref{stacks-cohomology-lemma-adjoint-kernel-parasitic}.
On conclut donc de même que $Q(\Ker(\varphi))$ est cohérent.
\end{proof}

\begin{lemma}
\label{stacks-cohomology-lemma-extension-coherent}
Soit $\mathcal{X}$ un champ algébrique localement noethérien.
Étant donnée une suite exacte courte
$0 \to \mathcal{F}_1 \to \mathcal{F}_2 \to \mathcal{F}_3 \to 0$
dans $\textit{Mod}(\mathcal{O}_\mathcal{X})$,
si $\mathcal{F}_1$ et $\mathcal{F}_3$ sont cohérents, alors
$\mathcal{F}_2$ est cohérent.
\end{lemma}

\begin{proof}
D'après Faisceaux sur les champs, lemme
\ref{stacks-sheaves-lemma-quasi-coherent-algebraic-stack}, partie (7),
$\mathcal{F}_2$ est quasi-cohérent. On peut alors vérifier que
$\mathcal{F}_2$ est cohérent par restriction à $U_\etale$
pour un morphisme $U \to \mathcal{X}$ lisse et surjectif.
Cela résulte de Cohomologie des espaces, lemme
\ref{spaces-cohomology-lemma-coherent-abelian-Noetherian}.
Quelques détails sont omis.
\end{proof}

\noindent
Les modules cohérents forment une sous-catégorie de Serre de la
catégorie des $\mathcal{O}_\mathcal{X}$-modules quasi-cohérents.
Cela n'est pas vrai pour les modules sur un topos annelé général.

\begin{lemma}
\label{stacks-cohomology-lemma-coherent-Noetherian-quasi-coherent-sub-quotient}
Soit $\mathcal{X}$ un champ algébrique localement noethérien.
Alors $\textit{Coh}(\mathcal{O}_\mathcal{X})$ est une sous-catégorie de Serre de
$\QCoh(\mathcal{O}_\mathcal{X})$. Soit $\varphi : \mathcal{F} \to \mathcal{G}$
un morphisme de $\mathcal{O}_\mathcal{X}$-modules quasi-cohérents. On a :
\begin{enumerate}
\item si $\mathcal{F}$ est cohérent et $\varphi$ est surjectif,
alors $\mathcal{G}$ est cohérent ;
\item si $\mathcal{F}$ est cohérent, alors $\Im(\varphi)$ est cohérente ;
\item si $\mathcal{G}$ est cohérent et $\Ker(\varphi)$ est parasite, alors
$\mathcal{F}$ est cohérent.
\end{enumerate}
\end{lemma}

\begin{proof}
Choisissons un schéma $U$ et un morphisme lisse surjectif $f : U \to \mathcal{X}$.
Alors le foncteur
$f^* : \QCoh(\mathcal{O}_\mathcal{X}) \to \QCoh(\mathcal{O}_U)$
est exact (lemme \ref{stacks-cohomology-lemma-flat-pullback-quasi-coherent})
et, en outre, $\textit{Coh}(\mathcal{O}_\mathcal{X})$ est par définition
la sous-catégorie pleine de $\QCoh(\mathcal{O}_\mathcal{X})$ formée
des objets $\mathcal{F}$ tels que $f^*\mathcal{F}$ appartienne à
$\textit{Coh}(\mathcal{O}_U)$. Le fait que
$\textit{Coh}(\mathcal{O}_\mathcal{X})$ soit une sous-catégorie de Serre de
$\QCoh(\mathcal{O}_\mathcal{X})$ résulte aussitôt de ceci
et du fait correspondant pour $U$ ; voir
Cohomologie des espaces, lemmes
\ref{spaces-cohomology-lemma-coherent-abelian-Noetherian} et
\ref{spaces-cohomology-lemma-coherent-Noetherian-quasi-coherent-sub-quotient}.
Nous omettons la preuve de (1), (2) et (3). Indication : comparer
avec la démonstration du lemme \ref{stacks-cohomology-lemma-coherent-abelian-Noetherian}.
\end{proof}

\noindent
Soit $\mathcal{X}$ un champ algébrique localement noethérien.
Soit $U$ un espace algébrique et soit $f : U \to \mathcal{X}$
un morphisme surjectif, localement de présentation finie et plat.
Observons que $U$ est localement noethérien
(Morphismes de champs, lemme
\ref{stacks-morphisms-lemma-locally-finite-type-locally-noetherian}).
Soit $(U, R, s, t, c)$ le groupoïde en espaces algébriques
et $f_{can} : [U/R] \to \mathcal{X}$ l'isomorphisme
construit dans
Champs algébriques, lemme \ref{algebraic-lemma-map-space-into-stack} et
remarque \ref{algebraic-remark-flat-fp-presentation}.
Comme dans Faisceaux sur les champs, section
\ref{stacks-sheaves-section-quasi-coherent-algebraic-stacks}
on obtient des équivalences
$$
\QCoh(\mathcal{O}_\mathcal{X})
\cong
\QCoh(\mathcal{O}_{[U/R]})
\cong
\QCoh(U, R, s, t, c)
$$
où la seconde équivalence est celle de
Faisceaux sur les champs, proposition \ref{stacks-sheaves-proposition-quasi-coherent}.
Rappelons que, dans Groupoïdes dans les espaces, section
\ref{spaces-groupoids-section-colimits}
nous avons défini la sous-catégorie pleine
$$
\textit{Coh}(U, R, s, t, c) \subset \QCoh(U, R, s, t, c)
$$
des {\it modules cohérents} comme celle des $(\mathcal{G}, \alpha)$ tels que
$\mathcal{G}$ soit un $\mathcal{O}_U$-module cohérent.

\begin{lemma}
\label{stacks-cohomology-lemma-coherent-presentation}
Dans la situation décrite ci-dessus, l'équivalence
$\QCoh(\mathcal{O}_\mathcal{X}) \cong \QCoh(U, R, s, t, c)$
envoie les modules cohérents sur les modules cohérents et réciproquement, c'est-à-dire
qu'elle induit une équivalence
$\textit{Coh}(\mathcal{O}_\mathcal{X}) \cong \textit{Coh}(U, R, s, t, c)$.
\end{lemma}

\begin{proof}
Cela résulte aussitôt de la définition des
$\mathcal{O}_\mathcal{X}$-modules cohérents.
Pour expliciter les références, ce qui précède utilise
Morphismes de champs, lemme
\ref{stacks-morphisms-lemma-locally-finite-type-locally-noetherian},
Champs algébriques, lemme \ref{algebraic-lemma-map-space-into-stack} et
remarque \ref{algebraic-remark-flat-fp-presentation},
Faisceaux sur les champs, section
\ref{stacks-sheaves-section-quasi-coherent-algebraic-stacks},
Faisceaux sur les champs, proposition \ref{stacks-sheaves-proposition-quasi-coherent},
et Groupoïdes dans les espaces, section
\ref{spaces-groupoids-section-colimits}.
\end{proof}

\begin{lemma}
\label{stacks-cohomology-lemma-coherent-hom}
Soit $\mathcal{X}$ un champ algébrique localement noethérien. Soient $\mathcal{F}$
et $\mathcal{G}$ des $\mathcal{O}_\mathcal{X}$-modules cohérents. Alors
le Hom interne $hom(\mathcal{F}, \mathcal{G})$
construit dans le lemme \ref{stacks-cohomology-lemma-internal-hom-fp-into-qcoh}
est un $\mathcal{O}_\mathcal{X}$-module cohérent.
\end{lemma}

\begin{proof}
Soit $U \to \mathcal{X}$ un morphisme lisse surjectif issu d'un schéma.
D'après le point (\ref{stacks-cohomology-item-hom-restriction}) de la section \ref{stacks-cohomology-section-further-remarks},
la restriction de
$hom(\mathcal{F}, \mathcal{G})$ à $U$ est le faisceau Hom des restrictions.
Ce lemme résulte donc du cas des espaces algébriques ; voir
Cohomologie des espaces, lemme \ref{spaces-cohomology-lemma-tensor-hom-coherent}.
\end{proof}






\section{Faisceaux cohérents sur les champs noethériens}
\label{stacks-cohomology-section-coherent-on-noetherian}

\noindent
Cette section est l'analogue de
Cohomologie des espaces, section
\ref{spaces-cohomology-section-coherent-quasi-compact}.

\begin{lemma}
\label{stacks-cohomology-lemma-directed-colimit-coherent}
Soit $\mathcal{X}$ un champ algébrique noethérien. Tout
$\mathcal{O}_\mathcal{X}$-module quasi-cohérent est la colimite filtrante
de ses sous-modules cohérents.
\end{lemma}

\begin{proof}
Soit $\mathcal{F}$ un $\mathcal{O}_\mathcal{X}$-module quasi-cohérent.
Si $\mathcal{G}, \mathcal{H} \subset \mathcal{F}$ sont des sous-modules
$\mathcal{O}_\mathcal{X}$-cohérents, alors l'image de
$\mathcal{G} \oplus \mathcal{H} \to \mathcal{F}$ est encore un
sous-module $\mathcal{O}_\mathcal{X}$-cohérent qui les contient tous deux,
voir le lemme \ref{stacks-cohomology-lemma-coherent-Noetherian-quasi-coherent-sub-quotient}.
Le système est donc filtrant.
Il suffit dès lors de montrer que $\mathcal{F}$ s'écrit comme
colimite filtrante de modules cohérents : on pourra alors prendre les
images de ces modules dans $\mathcal{F}$ et conclure qu'il y en a
suffisamment.

\medskip\noindent
Soient $U$ un schéma affine et $U \to \mathcal{X}$
un morphisme lisse surjectif (Propriétés des champs, lemme
\ref{stacks-properties-lemma-quasi-compact-stack}).
Posons $R = U \times_\mathcal{X} U$, de sorte que $\mathcal{X} = [U/R]$
comme dans Champs algébriques, lemme \ref{algebraic-lemma-stack-presentation}.
D'après le lemme \ref{stacks-cohomology-lemma-coherent-presentation}, on a
$\QCoh(\mathcal{O}_X) = \QCoh(U, R, s, t, c)$ et
$\textit{Coh}(\mathcal{O}_X) = \textit{Coh}(U, R, s, t, c)$.
On est ainsi ramené à établir l'énoncé correspondant pour
$\QCoh(U, R, s, t, c)$. C'est le contenu de
Groupoïdes dans les espaces, lemme \ref{spaces-groupoids-lemma-colimit-coherent} ;
nous en vérifions les hypothèses au paragraphe suivant.

\medskip\noindent
Nous invitons le lecteur à passer le reste de la démonstration.
Le schéma affine $U$ est noethérien ; cela résulte de notre définition
d'un champ $\mathcal{X}$ localement noethérien, voir
Propriétés des champs, définition
\ref{stacks-properties-definition-type-property} et
remarque \ref{stacks-properties-remark-list-properties-local-smooth-topology}.
Les morphismes de projection $s, t : R \to U$ sont lisses (voir la référence
ci-dessus),
quasi-séparés et quasi-compacts (Morphismes de champs, lemme
\ref{stacks-morphisms-lemma-quasi-compact-quasi-separated-permanence}).
En particulier, $R$ est un espace algébrique quasi-compact et quasi-séparé,
lisse sur $U$, donc noethérien
(Morphismes d'espaces, lemme
\ref{spaces-morphisms-lemma-finite-presentation-noetherian}).
\end{proof}











% OK, start here.
%

\chapter{Catégories dérivées des champs}


\phantomsection
\label{stacks-perfect-section-phantom}





\section{Introduction}
\label{stacks-perfect-section-introduction}

\noindent
Dans ce chapitre, nous étudions les catégories dérivées associées aux champs
algébriques. Il s'agit notamment des catégories dérivées de faisceaux
quasi-cohérents : nous démontrons des analogues des résultats pour les schémas
(voir Catégories dérivées des schémas, section
\ref{perfect-section-introduction}) et pour les espaces algébriques (voir
Catégories dérivées des espaces, section
\ref{spaces-perfect-section-introduction}). Les résultats du présent chapitre
diffèrent de ceux de \cite{LM-B} principalement parce que nous utilisons
systématiquement les « gros sites ». Avant d'aborder ce chapitre, le lecteur
est invité à parcourir rapidement les chapitres « Faisceaux sur les champs
algébriques » et « Cohomologie des champs algébriques », où est introduite la
terminologie employée ici.



\section{Conventions, notations et abus de langage}
\label{stacks-perfect-section-conventions}

\noindent
Nous continuons d'utiliser les conventions et les abus de langage introduits
dans Propriétés des champs, section
\ref{stacks-properties-section-conventions}. Nous adoptons les notations
expliquées dans Cohomologie des champs, section
\ref{stacks-cohomology-section-notation}.












\section{Les sites lisse-\'etale et plat-fppf}
\label{stacks-perfect-section-lisse-etale}

\noindent
Cette section est l'analogue, pour les catégories dérivées, de Cohomologie des
champs, section \ref{stacks-cohomology-section-lisse-etale}.

\begin{lemma}
\label{stacks-perfect-lemma-shriek-derived}
Soit $\mathcal{X}$ un champ algébrique. Nous employons les notations de
Cohomologie des champs, lemmes
\ref{stacks-cohomology-lemma-lisse-etale} et
\ref{stacks-cohomology-lemma-lisse-etale-modules}.
\begin{enumerate}
\item Le foncteur
$g_! : \textit{Ab}(\mathcal{X}_{lisse,\etale}) \to
\textit{Ab}(\mathcal{X}_\etale)$
admet un foncteur dérivé gauche
$$
Lg_! :
D(\mathcal{X}_{lisse,\etale})
\longrightarrow
D(\mathcal{X}_\etale)
$$
qui est adjoint à gauche de $g^{-1}$ et tel que
$g^{-1}Lg_! = \text{id}$.
\item Le foncteur $g_! :
\textit{Mod}(\mathcal{X}_{lisse,\etale},
\mathcal{O}_{\mathcal{X}_{lisse,\etale}}) \to
\textit{Mod}(\mathcal{X}_\etale, \mathcal{O}_{\mathcal{X}})$
admet un foncteur dérivé gauche
$$
Lg_! :
D(\mathcal{O}_{\mathcal{X}_{lisse,\etale}})
\longrightarrow
D(\mathcal{X}_\etale, \mathcal{O}_\mathcal{X})
$$
qui est adjoint à gauche de $g^*$ et tel que $g^*Lg_! = \text{id}$.
\item Le foncteur $g_! : \textit{Ab}(\mathcal{X}_{flat,fppf}) \to
\textit{Ab}(\mathcal{X}_{fppf})$
admet un foncteur dérivé gauche
$$
Lg_! :
D(\mathcal{X}_{flat, fppf})
\longrightarrow
D(\mathcal{X}_{fppf})
$$
qui est adjoint à gauche de $g^{-1}$ et tel que
$g^{-1}Lg_! = \text{id}$.
\item Le foncteur $g_! :
\textit{Mod}(\mathcal{X}_{flat,fppf},
\mathcal{O}_{\mathcal{X}_{flat,fppf}}) \to
\textit{Mod}(\mathcal{X}_{fppf}, \mathcal{O}_{\mathcal{X}})$
admet un foncteur dérivé gauche
$$
Lg_! :
D(\mathcal{O}_{\mathcal{X}_{flat, fppf}})
\longrightarrow
D(\mathcal{O}_\mathcal{X})
$$
qui est adjoint à gauche de $g^*$ et tel que $g^*Lg_! = \text{id}$.
\end{enumerate}
Attention : il n'est pas clair a priori que $Lg_!$ sur les modules coïncide avec
$Lg_!$ sur les faisceaux abéliens ; voir Cohomologie sur les sites, remarque
\ref{sites-cohomology-remark-when-derived-shriek-equal}.
\end{lemma}

\begin{proof}
L'existence du foncteur $Lg_!$ et son adjonction à $g^*$ résultent de
Cohomologie sur les sites, lemme
\ref{sites-cohomology-lemma-existence-derived-lower-shriek}. (Dans le cas des
faisceaux abéliens, on prend le faisceau constant $\mathbf{Z}$ comme faisceau
structural.) De plus, sa valeur sur un complexe $\mathcal{H}^\bullet$ se calcule
en choisissant une résolution gauche convenable
$\mathcal{K}^\bullet \to \mathcal{H}^\bullet$, puis en appliquant le foncteur
$g_!$ à $\mathcal{K}^\bullet$. Comme
$g^{-1}g_!\mathcal{K}^\bullet = \mathcal{K}^\bullet$ d'après Cohomologie des
champs, lemmes \ref{stacks-cohomology-lemma-lisse-etale-modules} et
\ref{stacks-cohomology-lemma-lisse-etale}, la dernière assertion est vérifiée
dans chacun des cas.
\end{proof}

\begin{lemma}
\label{stacks-perfect-lemma-lisse-etale-functorial-derived}
Sous les hypothèses et avec les notations de Cohomologie des champs, lemme
\ref{stacks-cohomology-lemma-lisse-etale-functorial}, nous avons
$$
g^{-1} \circ Rf_* = Rf'_* \circ (g')^{-1}
\quad\text{et}\quad
L(g')_! \circ (f')^{-1} = f^{-1} \circ Lg_!
$$
sur les catégories dérivées non bornées, aussi bien pour les modules que pour
les faisceaux abéliens.
\end{lemma}

\begin{proof}
Soit $\tau = \etale$ (resp.\ $\tau = fppf$), et soit $\mathcal{F}$ un faisceau
abélien sur $\mathcal{X}_\tau$. D'après Cohomologie des champs, lemme
\ref{stacks-cohomology-lemma-lisse-etale-functorial-cohomology}
l'application canonique de changement de base
$$
g^{-1}Rf_*\mathcal{F} \longrightarrow Rf'_*(g')^{-1}\mathcal{F}
$$
est un isomorphisme. La suite de la démonstration est formelle. Puisque la
cohomologie des groupes abéliens coïncide avec celle des faisceaux de modules,
nous en déduisons également que
$g^{-1} Rf_*\mathcal{F} = Rf'_* (g')^{-1}\mathcal{F}$ lorsque $\mathcal{F}$ est
un faisceau de modules sur $\mathcal{X}_\tau$.

\medskip\noindent
Montrons ensuite que, pour $\mathcal{G}$, faisceau de modules ou faisceau de
groupes abéliens sur $\mathcal{Y}_{lisse,\etale}$ (resp.\
$\mathcal{Y}_{flat,fppf}$), l'application canonique
$$
L(g')_!(f')^{-1}\mathcal{G} \to f^{-1}Lg_!\mathcal{G}
$$
est un isomorphisme. Il suffit pour cela de démontrer que, pour tout faisceau
injectif $\mathcal{I}$ sur $\mathcal{X}_\tau$, l'application induite
$$
\Hom(L(g')_!(f')^{-1}\mathcal{G}, \mathcal{I}[n])
\leftarrow
\Hom(f^{-1}Lg_!\mathcal{G}, \mathcal{I}[n])
$$
est un isomorphisme pour tout $n \in \mathbf{Z}$ (les Hom étant pris dans les
catégories dérivées appropriées). Par les adjonctions entre $f^{-1}$ et $Rf_*$,
entre $Lg_!$ et $g^{-1}$, et par leurs versions « primées », cela résulte de
l'isomorphisme
$g^{-1} Rf_*\mathcal{I} \to Rf'_* (g')^{-1}\mathcal{I}$ démontré ci-dessus.

\medskip\noindent
Dans le cas d'un complexe borné $\mathcal{G}^\bullet$, de modules ou de groupes
abéliens, sur $\mathcal{Y}_{lisse,\etale}$ (resp.\
$\mathcal{Y}_{fppf}$), l'application canonique
\begin{equation}
\label{stacks-perfect-equation-to-show}
L(g')_!(f')^{-1}\mathcal{G}^\bullet \to f^{-1}Lg_!\mathcal{G}^\bullet
\end{equation}
est un isomorphisme. Cela découle du cas d'un faisceau par les arguments usuels
de troncature et du fait que les foncteurs $L(g')_!(f')^{-1}$ et
$f^{-1}Lg_!$ sont des foncteurs exacts de catégories triangulées.

\medskip\noindent
Supposons que $\mathcal{G}^\bullet$ soit un complexe borné supérieurement, de
modules ou de groupes abéliens, sur $\mathcal{Y}_{lisse,\etale}$ (resp.\
$\mathcal{Y}_{fppf}$). L'application canonique
(\ref{stacks-perfect-equation-to-show}) est un isomorphisme, car les troncatures bêtes
$\sigma_{\geq -n}$ (voir Homologie, section
\ref{homology-section-truncations}) permettent d'écrire
$\mathcal{G}^\bullet$ comme une colimite
$\mathcal{G}^\bullet = \colim \mathcal{G}_n^\bullet$
de complexes bornés. On obtient ainsi un triangle distingué
$$
\bigoplus\nolimits_{n \geq 1} \mathcal{G}_n^\bullet \to
\bigoplus\nolimits_{n \geq 1} \mathcal{G}_n^\bullet \to
\mathcal{G}^\bullet \to \ldots
$$
et chacun des foncteurs $L(g')_!$, $(f')^{-1}$, $f^{-1}$ et $Lg_!$ commute aux
sommes directes de complexes.

\medskip\noindent
Si $\mathcal{G}^\bullet$ est un complexe arbitraire, de modules ou de groupes
abéliens, sur $\mathcal{Y}_{lisse,\etale}$ (resp.\
$\mathcal{Y}_{fppf}$), nous employons les troncatures canoniques
$\tau_{\leq n}$ (voir Homologie, section
\ref{homology-section-truncations}) pour écrire $\mathcal{G}^\bullet$ comme une
colimite de complexes bornés supérieurement, puis nous reprenons l'argument du
paragraphe précédent.

\medskip\noindent
Enfin, par les adjonctions entre $f^{-1}$ et $Rf_*$, entre $Lg_!$ et $g^{-1}$,
et par leurs versions « primées », nous concluons en toute généralité que la
première identité du lemme résulte de la seconde.
\end{proof}

\begin{lemma}
\label{stacks-perfect-lemma-higher-shriek-quasi-coherent}
Soit $\mathcal{X}$ un champ algébrique. Nous employons les notations de
Cohomologie des champs, lemme
\ref{stacks-cohomology-lemma-lisse-etale}.
\begin{enumerate}
\item Soit $\mathcal{H}$ un
$\mathcal{O}_{\mathcal{X}_{lisse,\etale}}$-module quasi-cohérent sur le site
lisse-\'etale de $\mathcal{X}$. Pour tout $p \in \mathbf{Z}$, le faisceau
$H^p(Lg_!\mathcal{H})$ est un module localement quasi-cohérent sur
$\mathcal{X}$ et possède la propriété de changement de base plat.
\item Soit $\mathcal{H}$ un
$\mathcal{O}_{\mathcal{X}_{flat,fppf}}$-module quasi-cohérent sur le site
plat-fppf de $\mathcal{X}$. Pour tout $p \in \mathbf{Z}$, le faisceau
$H^p(Lg_!\mathcal{H})$ est un module localement quasi-cohérent sur
$\mathcal{X}$ et possède la propriété de changement de base plat.
\end{enumerate}
\end{lemma}

\begin{proof}
Choisissons un schéma $U$ et un morphisme lisse surjectif
$x : U \to \mathcal{X}$. D'après Modules sur les sites, définition
\ref{sites-modules-definition-site-local}, il existe un recouvrement étale
(resp.\ fppf) $\{U_i \to U\}_{i \in I}$ tel que chaque image inverse
$f_i^{-1}\mathcal{H}$ possède une présentation globale (voir Modules sur les
sites, définition \ref{sites-modules-definition-global}). Ici,
$f_i : U_i \to \mathcal{X}$ est le composé $U_i \to U \to \mathcal{X}$, qui est
un morphisme de champs algébriques. (Rappelons que l'image inverse « est » la
restriction à $\mathcal{X}/f_i$ ; voir Faisceaux sur les champs, définition
\ref{stacks-sheaves-definition-pullback} et la discussion qui suit.) Après
avoir raffiné le recouvrement, nous pouvons supposer que chaque $U_i$ est un
schéma affine. Puisque chaque $f_i$ est lisse (resp.\ plat), le lemme
\ref{stacks-perfect-lemma-lisse-etale-functorial-derived} donne
$f_i^{-1}Lg_!\mathcal{H} = Lg_{i, !}(f'_i)^{-1}\mathcal{H}$. À l'aide de
Cohomologie des champs, lemme
\ref{stacks-cohomology-lemma-check-lqc-fbc-on-covering}, nous ramenons l'énoncé
du lemme au cas où $\mathcal{H}$ possède une présentation globale et où
$\mathcal{X} = (\Sch/X)_{fppf}$ pour un schéma affine $X = \Spec(A)$.

\medskip\noindent
Écrivons cette présentation sous la forme
$$
\bigoplus\nolimits_{j \in J} \mathcal{O} \longrightarrow
\bigoplus\nolimits_{i \in I} \mathcal{O} \longrightarrow
\mathcal{H} \longrightarrow 0
$$
où $\mathcal{O} = \mathcal{O}_{\mathcal{X}_{lisse,\etale}}$ (resp.\
$\mathcal{O} = \mathcal{O}_{\mathcal{X}_{flat,fppf}}$). Remarquons que le site
$\mathcal{X}_{lisse,\etale}$ (resp.\ $\mathcal{X}_{flat,fppf}$) possède un
objet final, à savoir $X/X$, qui est quasi-compact (voir Cohomologie sur les
sites, section \ref{sites-cohomology-section-limits}). Nous avons donc
$$
\Gamma(\bigoplus\nolimits_{i \in I} \mathcal{O}) =
\bigoplus\nolimits_{i \in I} A
$$
d'après Sites, lemme \ref{sites-lemma-directed-colimits-sections}. Par
conséquent, l'application de la présentation correspond à une présentation
analogue
$$
\bigoplus\nolimits_{j \in J} A \longrightarrow
\bigoplus\nolimits_{i \in I} A \longrightarrow
M \longrightarrow 0
$$
d'un $A$-module $M$. De plus, $\mathcal{H}$ est égal à la restriction au site
lisse-\'etale (resp.\ plat-fppf) du faisceau quasi-cohérent $M^a$ associé à
$M$. Choisissons une résolution
$$
\ldots \to F_2 \to F_1 \to F_0 \to M \to 0
$$
par des $A$-modules libres. Le complexe
$$
\ldots \mathcal{O} \otimes_A F_2 \to \mathcal{O} \otimes_A F_1 \to
\mathcal{O} \otimes_A F_0 \to \mathcal{H} \to 0
$$
est une résolution de $\mathcal{H}$ par des $\mathcal{O}$-modules libres, car,
pour tout objet $U/X$ de $\mathcal{X}_{lisse,\etale}$ (resp.\
$\mathcal{X}_{flat,fppf}$), le morphisme structural $U \to X$ est plat. Par
construction, la valeur de $Lg_!\mathcal{H}$ est donc
$$
\ldots \to
\mathcal{O}_\mathcal{X} \otimes_A F_2 \to
\mathcal{O}_\mathcal{X} \otimes_A F_1 \to
\mathcal{O}_\mathcal{X} \otimes_A F_0 \to 0 \to \ldots
$$
Comme il s'agit d'un complexe de modules quasi-cohérents sur
$\mathcal{X}_\etale$ (resp.\ $\mathcal{X}_{fppf}$), il résulte de Cohomologie
des champs, proposition
\ref{stacks-cohomology-proposition-loc-qcoh-flat-base-change}, que
$H^p(Lg_!\mathcal{H})$ est quasi-cohérent.
\end{proof}




\section{Cohomologie et sites lisse-\'etale et plat-fppf}
\label{stacks-perfect-section-compare-unbounded-lisse-etale}

\noindent
Nous avons déjà vu que la cohomologie d'un faisceau sur un champ algébrique
$\mathcal{X}$ peut se calculer sur le site plat-fppf. Dans cette section, nous
démontrons qu'il en va de même pour les objets, éventuellement non bornés, de
la catégorie dérivée de $\mathcal{X}$.

\begin{lemma}
\label{stacks-perfect-lemma-higher-shriek-Z}
Soit $\mathcal{X}$ un champ algébrique. Nous avons
$Lg_!\mathbf{Z} = \mathbf{Z}$, qu'il s'agisse du $Lg_!$ de la partie (1) du
lemme \ref{stacks-perfect-lemma-shriek-derived} ou du $Lg_!$ de la partie (3) du lemme
\ref{stacks-perfect-lemma-shriek-derived}.
\end{lemma}

\begin{proof}
Nous le démontrons pour la comparaison du site plat-fppf au site fppf ; le cas
du site lisse-\'etale est exactement analogue. Il faut montrer que
$H^i(Lg_!\mathbf{Z})$ est $0$ pour $i \not = 0$ et que l'application canonique
$H^0(Lg_!\mathbf{Z}) \to \mathbf{Z}$ est un isomorphisme. Soit
$f : \mathcal{U} \to \mathcal{X}$ un morphisme plat et surjectif, où
$\mathcal{U}$ est un schéma, et supposons en outre $f$ localement de
présentation finie. (On peut, par exemple, choisir une présentation
$U \to \mathcal{X}$ et prendre pour $\mathcal{U}$ le champ algébrique
correspondant à $U$.) D'après Faisceaux sur les champs, lemmes
\ref{stacks-sheaves-lemma-check-exactness-covering} et
\ref{stacks-sheaves-lemma-surjective-flat-locally-finite-presentation}
il suffit de montrer que l'image inverse $f^{-1}H^i(Lg_!\mathbf{Z})$ est $0$
pour $i \not = 0$ et que l'image inverse
$H^0(Lg_!\mathbf{Z}) \to f^{-1}\mathbf{Z}$ est un isomorphisme. Le lemme
\ref{stacks-perfect-lemma-lisse-etale-functorial-derived} donne
$f^{-1}Lg_!\mathbf{Z} = L(g')_!\mathbf{Z}$, où
$g' : \Sh(\mathcal{U}_{flat, fppf}) \to \Sh(\mathcal{U}_{fppf})$
est le morphisme de comparaison correspondant pour $\mathcal{U}$. Nous sommes
ainsi ramenés au cas étudié au paragraphe suivant.

\medskip\noindent
Supposons $\mathcal{X} = (\Sch/X)_{fppf}$ pour un schéma $X$. La catégorie
$\mathcal{X}_{flat, fppf}$ possède alors un objet final $e$, à savoir $X/X$ ;
de plus, le foncteur
$u : \mathcal{X}_{flat, fppf} \to \mathcal{X}_{fppf}$ envoie $e$ sur l'objet
final. Puisque $\mathbf{Z}$ est le faisceau abélien libre sur l'objet final,
lorsque celui-ci existe, nous obtenons $Lg_!\mathbf{Z} = \mathbf{Z}$ par la
construction même de $Lg_!$ dans Cohomologie sur les sites, lemme
\ref{sites-cohomology-lemma-existence-derived-lower-shriek}.
\end{proof}

\begin{lemma}
\label{stacks-perfect-lemma-lisse-etale-cohomology}
Soit $\mathcal{X}$ un champ algébrique. Nous employons les notations du lemme
\ref{stacks-perfect-lemma-shriek-derived}.
\begin{enumerate}
\item Pour $K$ dans $D(\mathcal{X}_\etale)$, nous avons
\begin{enumerate}
\item $R\Gamma(\mathcal{X}_\etale, K) =
R\Gamma(\mathcal{X}_{lisse,\etale}, g^{-1}K)$, et
\item $R\Gamma(x, K) =
R\Gamma(\mathcal{X}_{lisse,\etale}/x, g^{-1}K)$
pour tout objet $x$ de $\mathcal{X}_{lisse,\etale}$.
\end{enumerate}
\item Pour $K$ dans $D(\mathcal{X}_{fppf})$, nous avons
\begin{enumerate}
\item $R\Gamma(\mathcal{X}_{fppf}, K) =
R\Gamma(\mathcal{X}_{flat,fppf}, g^{-1}K)$, et
\item $H^p(x, K) =
R\Gamma(\mathcal{X}_{flat,fppf}/x, g^{-1}K)$
pour tout objet $x$ de $\mathcal{X}_{flat,fppf}$.
\end{enumerate}
\end{enumerate}
Dans les deux cas, les mêmes assertions valent pour les modules, puisque
$g^{-1} = g^*$ et que le calcul de la cohomologie ne présente aucune différence
d'après Cohomologie sur les sites, lemme
\ref{sites-cohomology-lemma-modules-abelian-unbounded}.
\end{lemma}

\begin{proof}
Nous le démontrons pour la comparaison du site plat-fppf au site fppf ; le cas
du site lisse-\'etale est exactement analogue. D'après le lemme
\ref{stacks-perfect-lemma-higher-shriek-Z}, nous avons $Lg_!\mathbf{Z} = \mathbf{Z}$. Il vient
\begin{align*}
R\Gamma(\mathcal{X}_{fppf}, K)
& =
R\Hom(\mathbf{Z}, K) \\
& =
R\Hom(Lg_!\mathbf{Z}, K) \\
& =
R\Hom(\mathbf{Z}, g^{-1}K) \\
& =
R\Gamma(\mathcal{X}_{lisse,\etale}, g^{-1}K)
\end{align*}
Ceci démontre (1)(a). La partie (1)(b) résulte de la partie (1)(a). En effet,
si $x$ est au-dessus du schéma $U$, alors le site $\mathcal{X}_\etale/x$ est
équivalent à $(\Sch/U)_\etale$, et $\mathcal{X}_{lisse,\etale}$ est équivalent
à $U_{lisse, \etale}$.
\end{proof}










\section{Catégories dérivées de modules quasi-cohérents}
\label{stacks-perfect-section-derived}

\noindent
Soit $\mathcal{X}$ un champ algébrique. Comme le foncteur d'inclusion
$\QCoh(\mathcal{O}_\mathcal{X}) \to
\textit{Mod}(\mathcal{O}_\mathcal{X})$ n'est pas exact, nous ne pouvons pas
définir $D_\QCoh(\mathcal{O}_\mathcal{X})$ comme la sous-catégorie pleine de
$D(\mathcal{O}_\mathcal{X})$ formée des complexes dont les faisceaux de
cohomologie sont quasi-cohérents. Nous définissons plutôt la catégorie dérivée
des modules quasi-cohérents comme un quotient, par analogie avec Cohomologie
des champs, remarque
\ref{stacks-cohomology-remark-bousfield-colocalization}.

\medskip\noindent
Rappelons que $\textit{LQCoh}^{fbc}(\mathcal{O}_\mathcal{X}) \subset
\textit{Mod}(\mathcal{O}_\mathcal{X})$ désigne la sous-catégorie pleine des
$\mathcal{O}_\mathcal{X}$-modules localement quasi-cohérents possédant la
propriété de changement de base plat ; voir Cohomologie des champs, section
\ref{stacks-cohomology-section-loc-qcoh-flat-base-change}.
Nous abrégerons
$$
D_{\textit{LQCoh}^{fbc}}(\mathcal{O}_\mathcal{X}) =
D_{\textit{LQCoh}^{fbc}(\mathcal{O}_\mathcal{X})}(\mathcal{O}_\mathcal{X})
$$
Du lemme \ref{derived-lemma-cohomology-in-serre-subcategory} de Catégories
dérivées et de la partie (2) de Cohomologie des champs, proposition
\ref{stacks-cohomology-proposition-loc-qcoh-flat-base-change}, nous déduisons
que $D_{\textit{LQCoh}^{fbc}}(\mathcal{O}_\mathcal{X})$ est une sous-catégorie
triangulée strictement pleine et saturée de
$D(\mathcal{O}_\mathcal{X})$.

\medskip\noindent
Notons $\textit{Parasitic}(\mathcal{O}_\mathcal{X}) \subset
\textit{Mod}(\mathcal{O}_\mathcal{X})$ la sous-catégorie pleine des
$\mathcal{O}_\mathcal{X}$-modules parasites ; voir Cohomologie des champs,
section
\ref{stacks-cohomology-section-parasitic}.
Abrégeons de même
$$
D_{\textit{Parasitic}}(\mathcal{O}_\mathcal{X}) =
D_{\textit{Parasitic}(\mathcal{O}_\mathcal{X})}(\mathcal{O}_\mathcal{X})
$$
Comme précédemment, c'est une sous-catégorie triangulée strictement pleine et
saturée de $D(\mathcal{O}_\mathcal{X})$, puisque
$\textit{Parasitic}(\mathcal{O}_\mathcal{X})$ est une sous-catégorie de Serre
de $\textit{Mod}(\mathcal{O}_\mathcal{X})$ ; voir Cohomologie des champs, lemme
\ref{stacks-cohomology-lemma-parasitic}.

\medskip\noindent
L'intersection des sous-catégories faiblement de Serre
$\textit{Parasitic}(\mathcal{O}_\mathcal{X}) \cap
\textit{LQCoh}^{fbc}(\mathcal{O}_\mathcal{X})$
de $\textit{Mod}(\mathcal{O}_\mathcal{X})$ est encore une sous-catégorie
faiblement de Serre. Nous adoptons de même l'abréviation
\begin{align*}
D_{\textit{Parasitic} \cap \textit{LQCoh}^{fbc}}(\mathcal{O}_\mathcal{X})
& =
D_{\textit{Parasitic}(\mathcal{O}_\mathcal{X}) \cap
\textit{LQCoh}^{fbc}(\mathcal{O}_\mathcal{X})}(\mathcal{O}_\mathcal{X}) \\
& =
D_{\textit{Parasitic}}(\mathcal{O}_\mathcal{X})
\cap
D_{\textit{LQCoh}^{fbc}}(\mathcal{O}_\mathcal{X})
\end{align*}
Comme précédemment, c'est une sous-catégorie triangulée strictement pleine et
saturée de $D(\mathcal{O}_\mathcal{X})$. A fortiori, c'est donc une
sous-catégorie triangulée strictement pleine et saturée de
$D_{\textit{LQCoh}^{fbc}}(\mathcal{O}_\mathcal{X})$.

\begin{definition}
\label{stacks-perfect-definition-derived}
Soit $\mathcal{X}$ un champ algébrique. Avec les notations ci-dessus, nous
définissons la {\it catégorie dérivée des $\mathcal{O}_\mathcal{X}$-modules à
faisceaux de cohomologie quasi-cohérents} comme le quotient de
Verdier\footnote{Cette définition diffère de celle de la littérature ; voir
\cite[6.3]{olsson_sheaves}. Elle coïncide toutefois avec cette dernière d'après
le lemme \ref{stacks-perfect-lemma-derived-quasi-coherent}.}
$$
D_\QCoh(\mathcal{O}_\mathcal{X}) =
D_{\textit{LQCoh}^{fbc}}(\mathcal{O}_\mathcal{X})/
D_{\textit{Parasitic} \cap \textit{LQCoh}^{fbc}}(\mathcal{O}_\mathcal{X})
$$
\end{definition}

\noindent
Le quotient de Verdier est défini dans Catégories dérivées, section
\ref{derived-section-quotients}. Un morphisme $a : E \to E'$ de
$D_{\textit{LQCoh}^{fbc}}(\mathcal{O}_\mathcal{X})$ devient un isomorphisme
dans $D_\QCoh(\mathcal{O}_\mathcal{X})$ si et seulement si le cône $C(a)$ a des
faisceaux de cohomologie parasites ; voir Catégories dérivées, lemme
\ref{derived-lemma-operations}.

\medskip\noindent
Considérons les foncteurs
$$
D_{\textit{LQCoh}^{fbc}}(\mathcal{O}_\mathcal{X})
\xrightarrow{H^i}
\textit{LQCoh}^{fbc}(\mathcal{O}_\mathcal{X})
\xrightarrow{Q}
\QCoh(\mathcal{O}_\mathcal{X})
$$
Remarquons que $Q$ annule la sous-catégorie
$\textit{Parasitic}(\mathcal{O}_\mathcal{X}) \cap
\textit{LQCoh}^{fbc}(\mathcal{O}_\mathcal{X})$ ; voir Cohomologie des champs,
lemme \ref{stacks-cohomology-lemma-adjoint-kernel-parasitic}. D'après
Catégories dérivées, lemme
\ref{derived-lemma-universal-property-quotient}, nous obtenons un foncteur
cohomologique
\begin{equation}
\label{stacks-perfect-equation-Hi-quasi-coherent}
H^i :
D_\QCoh(\mathcal{O}_\mathcal{X})
\longrightarrow
\QCoh(\mathcal{O}_\mathcal{X})
\end{equation}
De plus, remarquons que $E \in D_\QCoh(\mathcal{O}_\mathcal{X})$ est nul si et
seulement si $H^i(E) = 0$ pour tout $i \in \mathbf{Z}$, puisque le noyau de
$Q$ est exactement
$\textit{Parasitic}(\mathcal{O}_\mathcal{X}) \cap
\textit{LQCoh}^{fbc}(\mathcal{O}_\mathcal{X})$ d'après Cohomologie des champs,
lemme \ref{stacks-cohomology-lemma-adjoint-kernel-parasitic}.

\medskip\noindent
Remarquons que les catégories
$\textit{Parasitic}(\mathcal{O}_\mathcal{X}) \cap
\textit{LQCoh}^{fbc}(\mathcal{O}_\mathcal{X})$ et
$\textit{LQCoh}^{fbc}(\mathcal{O}_\mathcal{X})$
constituent aussi des sous-catégories faiblement de Serre de la catégorie
abélienne
$\textit{Mod}(\mathcal{X}_\etale, \mathcal{O}_\mathcal{X})$
des modules pour la topologie étale ; voir Cohomologie des champs, proposition
\ref{stacks-cohomology-proposition-loc-qcoh-flat-base-change} et lemme
\ref{stacks-cohomology-lemma-parasitic}. L'énoncé du lemme suivant a donc un
sens.

\begin{lemma}
\label{stacks-perfect-lemma-compare-etale-fppf-QCoh}
Soit $\mathcal{X}$ un champ algébrique. Posons, par abréviation,
$\mathcal{P}_\mathcal{X} = \textit{Parasitic}(\mathcal{O}_\mathcal{X}) \cap
\textit{LQCoh}^{fbc}(\mathcal{O}_\mathcal{X})$.
Le morphisme de comparaison
$\epsilon : \mathcal{X}_{fppf} \to \mathcal{X}_\etale$
induit un diagramme commutatif
$$
\xymatrix{
D_{\textit{Parasitic} \cap \textit{LQCoh}^{fbc}}(\mathcal{O}_\mathcal{X})
\ar[r] &
D_{\textit{LQCoh}^{fbc}}(\mathcal{O}_\mathcal{X}) \ar[r] &
D(\mathcal{O}_\mathcal{X}) \\
D_{\mathcal{P}_\mathcal{X}}(\mathcal{X}_\etale, \mathcal{O}_\mathcal{X})
\ar[r] \ar[u]^{\epsilon^*} &
D_{\textit{LQCoh}^{fbc}(\mathcal{O}_\mathcal{X})}(
\mathcal{X}_\etale, \mathcal{O}_\mathcal{X})
\ar[r] \ar[u]^{\epsilon^*} &
D(\mathcal{X}_\etale, \mathcal{O}_\mathcal{X})
\ar[u]^{\epsilon^*}
}
$$
De plus, les deux flèches verticales de gauche sont des équivalences de
catégories triangulées ; nous obtenons donc aussi une équivalence
$$
D_{\textit{LQCoh}^{fbc}(\mathcal{O}_\mathcal{X})}
(\mathcal{X}_\etale, \mathcal{O}_\mathcal{X})
/
D_{\mathcal{P}_\mathcal{X}}(\mathcal{X}_\etale, \mathcal{O}_\mathcal{X})
\longrightarrow
D_\QCoh(\mathcal{O}_\mathcal{X})
$$
\end{lemma}

\begin{proof}
Puisque $\epsilon^*$ est exact, il est clair que nous obtenons le diagramme de
l'énoncé. Montrons que la flèche verticale médiane est une équivalence en
appliquant Cohomologie sur les sites, lemme
\ref{sites-cohomology-lemma-compare-topologies-derived-adequate-modules}
à la situation suivante :
$\mathcal{C} = \mathcal{X}$,
$\tau = fppf$,
$\tau' = \etale$,
$\mathcal{O} = \mathcal{O}_\mathcal{X}$,
$\mathcal{A} = \textit{LQCoh}^{fbc}(\mathcal{O}_\mathcal{X})$, et
$\mathcal{B}$ est l'ensemble des objets de $\mathcal{X}$ situés au-dessus de
schémas affines. Pour appliquer le lemme, il faut vérifier les conditions (1),
(2), (3) et (4). Les conditions (1) et (2) résultent clairement de la
discussion précédente ; plus précisément, elles découlent de Cohomologie des
champs, proposition
\ref{stacks-cohomology-proposition-loc-qcoh-flat-base-change}. La condition
(3) est satisfaite parce que tout schéma possède un recouvrement ouvert de
Zariski par des schémas affines. La condition (4) résulte de Descente, lemme
\ref{descent-lemma-quasi-coherent-and-flat-base-change}.

\medskip\noindent
Nous omettons de vérifier que l'équivalence de catégories $\epsilon^* :
D_{\textit{LQCoh}^{fbc}(\mathcal{O}_\mathcal{X})}(
\mathcal{X}_\etale, \mathcal{O}_\mathcal{X})
\to
D_{\textit{LQCoh}^{fbc}}(\mathcal{O}_\mathcal{X})$
induit une équivalence entre les sous-catégories de complexes à faisceaux de
cohomologie parasites.
\end{proof}

\noindent
Soit $\mathcal{X}$ un champ algébrique. D'après Cohomologie des champs, lemme
\ref{stacks-cohomology-lemma-quasi-coherent-weak-serre}
la catégorie des modules quasi-cohérents
$\QCoh(\mathcal{O}_{\mathcal{X}_{lisse,\etale}})$
forme une sous-catégorie faiblement de Serre de
$\textit{Mod}(\mathcal{O}_{\mathcal{X}_{lisse,\etale}})$
et la catégorie des modules quasi-cohérents
$\QCoh(\mathcal{O}_{\mathcal{X}_{flat,fppf}})$
forme une sous-catégorie faiblement de Serre de
$\textit{Mod}(\mathcal{O}_{\mathcal{X}_{flat,fppf}})$.
Nous pouvons donc considérer
$$
D_\QCoh(\mathcal{O}_{\mathcal{X}_{lisse,\etale}}) =
D_{\QCoh(\mathcal{O}_{\mathcal{X}_{lisse,\etale}})}(
\mathcal{O}_{\mathcal{X}_{lisse,\etale}})
\subset
D(\mathcal{O}_{\mathcal{X}_{lisse,\etale}})
$$
et, de même,
$$
D_\QCoh(\mathcal{O}_{\mathcal{X}_{flat,fppf}}) =
D_{\QCoh(\mathcal{O}_{\mathcal{X}_{flat,fppf}})}(
\mathcal{O}_{\mathcal{X}_{flat,fppf}})
\subset
D(\mathcal{O}_{\mathcal{X}_{flat,fppf}})
$$
Comme ci-dessus, ce sont des sous-catégories triangulées strictement pleines et
saturées. Il se trouve que $D_\QCoh(\mathcal{O}_\mathcal{X})$ est équivalente à
chacune d'elles.

\begin{lemma}
\label{stacks-perfect-lemma-derived-quasi-coherent}
Soit $\mathcal{X}$ un champ algébrique. Posons
$\mathcal{P}_\mathcal{X} = \textit{Parasitic}(\mathcal{O}_\mathcal{X}) \cap
\textit{LQCoh}^{fbc}(\mathcal{O}_\mathcal{X})$.
\begin{enumerate}
\item
Soit $\mathcal{F}^\bullet$ un objet de
$D_{\textit{LQCoh}^{fbc}(\mathcal{O}_\mathcal{X})}
(\mathcal{X}_\etale, \mathcal{O}_\mathcal{X})$.
Pour le site lisse-\'etale, prenons $g$ comme dans Cohomologie des champs,
lemme \ref{stacks-cohomology-lemma-lisse-etale}. Nous avons
\begin{enumerate}
\item $g^*\mathcal{F}^\bullet$ appartient à
$D_\QCoh(\mathcal{O}_{\mathcal{X}_{lisse,\etale}})$,
\item $g^*\mathcal{F}^\bullet = 0$ si et seulement si
$\mathcal{F}^\bullet$ appartient à
$D_{\mathcal{P}_\mathcal{X}}(\mathcal{X}_\etale, \mathcal{O}_\mathcal{X})$,
\item $Lg_!\mathcal{H}^\bullet$ appartient à
$D_{\textit{LQCoh}^{fbc}(\mathcal{O}_\mathcal{X})}(
\mathcal{X}_\etale, \mathcal{O}_\mathcal{X})$
pour $\mathcal{H}^\bullet$ appartenant à
$D_\QCoh(\mathcal{O}_{\mathcal{X}_{lisse,\etale}})$, et
\item les foncteurs $g^*$ et $Lg_!$ définissent des foncteurs mutuellement
inverses
$$
\xymatrix{
D_\QCoh(\mathcal{O}_\mathcal{X}) \ar@<1ex>[r]^-{g^*} &
D_\QCoh(\mathcal{O}_{\mathcal{X}_{lisse,\etale}})
\ar@<1ex>[l]^-{Lg_!}
}
$$
\end{enumerate}
\item
Soit $\mathcal{F}^\bullet$ un objet de
$D_{\textit{LQCoh}^{fbc}}(\mathcal{O}_\mathcal{X})$. Pour le site plat-fppf,
prenons $g$ comme dans Cohomologie des champs, lemme
\ref{stacks-cohomology-lemma-lisse-etale}. Nous avons
\begin{enumerate}
\item $g^*\mathcal{F}^\bullet$ appartient à
$D_\QCoh(\mathcal{O}_{\mathcal{X}_{flat, fppf}})$,
\item $g^*\mathcal{F}^\bullet = 0$ si et seulement si
$\mathcal{F}^\bullet$ appartient à
$D_{\mathcal{P}_\mathcal{X}}(\mathcal{O}_\mathcal{X})$,
\item $Lg_!\mathcal{H}^\bullet$ appartient à
$D_{\textit{LQCoh}^{fbc}}(\mathcal{O}_\mathcal{X})$
pour $\mathcal{H}^\bullet$ appartenant à
$D_\QCoh(\mathcal{O}_{\mathcal{X}_{flat,fppf}})$, et
\item les foncteurs $g^*$ et $Lg_!$ définissent des foncteurs mutuellement
inverses
$$
\xymatrix{
D_\QCoh(\mathcal{O}_\mathcal{X}) \ar@<1ex>[r]^-{g^*} &
D_\QCoh(\mathcal{O}_{\mathcal{X}_{flat,fppf}}) \ar@<1ex>[l]^-{Lg_!}
}
$$
\end{enumerate}
\end{enumerate}
\end{lemma}

\begin{proof}
Le foncteur $g^* = g^{-1}$ est exact ; (1)(a), (2)(a), (1)(b) et (2)(b)
résultent donc de Cohomologie des champs, lemmes
\ref{stacks-cohomology-lemma-quasi-coherent} et
\ref{stacks-cohomology-lemma-parasitic-in-terms-flat-fppf}.

\medskip\noindent
Démonstration de (1)(c) et (2)(c). La construction de $Lg_!$ au lemme
\ref{stacks-perfect-lemma-shriek-derived} (par l'intermédiaire de Cohomologie sur les sites,
lemme \ref{sites-cohomology-lemma-existence-derived-lower-shriek}, qui utilise
à son tour Catégories dérivées, proposition
\ref{derived-proposition-left-derived-exists}) montre que $Lg_!$, sur tout objet
$\mathcal{H}^\bullet$ de
$D(\mathcal{O}_{\mathcal{X}_{lisse,\etale}})$, se calcule par
$$
Lg_!\mathcal{H}^\bullet = \colim g_!\mathcal{K}_n^\bullet =
g_! \colim \mathcal{K}_n^\bullet
$$
(colimites terme à terme), où le quasi-isomorphisme
$\colim \mathcal{K}_n^\bullet \to \mathcal{H}^\bullet$ induit des
quasi-isomorphismes
$\mathcal{K}_n^\bullet \to \tau_{\leq n} \mathcal{H}^\bullet$. Comme les
foncteurs d'inclusion
$$
\textit{LQCoh}^{fbc}(\mathcal{O}_\mathcal{X}) \subset
\textit{Mod}(\mathcal{X}_\etale, \mathcal{O}_\mathcal{X})
\quad\text{et}\quad
\textit{LQCoh}^{fbc}(\mathcal{O}_\mathcal{X}) \subset
\textit{Mod}(\mathcal{O}_\mathcal{X})
$$
sont compatibles aux colimites filtrantes, il suffit de démontrer (c) pour les
complexes bornés supérieurement $\mathcal{H}^\bullet$ appartenant à
$D_\QCoh(\mathcal{O}_{\mathcal{X}_{lisse,\etale}})$ et à
$D_\QCoh(\mathcal{O}_{\mathcal{X}_{flat,fppf}})$. Dans ce cas, pour montrer que
$H^n(Lg_!\mathcal{H}^\bullet)$ appartient à
$\textit{LQCoh}^{fbc}(\mathcal{O}_\mathcal{X})$, nous pouvons raisonner par
récurrence sur l'entier $m$ tel que $\mathcal{H}^i = 0$ pour $i > m$. Si
$m < n$, alors $H^n(Lg_!\mathcal{H}^\bullet) = 0$ et le résultat est établi.
Dans le cas général, considérons le triangle distingué
$$
\tau_{\leq m - 1}\mathcal{H}^\bullet \to \mathcal{H}^\bullet \to
H^m(\mathcal{H}^\bullet)[-m] \to \ldots
$$
(Catégories dérivées, remarque
\ref{derived-remark-truncation-distinguished-triangle})
et appliquons le foncteur $Lg_!$. Comme
$\textit{LQCoh}^{fbc}(\mathcal{O}_\mathcal{X})$
est une sous-catégorie faiblement de Serre de la catégorie des modules, il
suffit de démontrer (c) pour deux termes sur trois. Le résultat vaut pour
$Lg_!\tau_{\leq m - 1}\mathcal{H}^\bullet$ par récurrence et pour
$Lg_!H^m(\mathcal{H}^\bullet)[-m]$ d'après le lemme
\ref{stacks-perfect-lemma-higher-shriek-quasi-coherent}. Par conséquent, (c) est vérifié.

\medskip\noindent
Démontrons (2)(d). D'après (2)(a) et (2)(b), le foncteur $g^{-1} = g^*$ induit
un foncteur
$$
c :
D_\QCoh(\mathcal{O}_\mathcal{X})
\longrightarrow
D_\QCoh(\mathcal{O}_{\mathcal{X}_{flat, fppf}})
$$
voir Catégories dérivées, lemme
\ref{derived-lemma-universal-property-quotient}. Nous avons donc le diagramme
suivant de catégories triangulées
$$
\xymatrix{
D_{\textit{LQCoh}^{fbc}}(\mathcal{O}_\mathcal{X})
\ar[rd]^{g^{-1}} \ar[rr]_q & &
D_\QCoh(\mathcal{O}_\mathcal{X}) \ar[ld]^c \\
& D_\QCoh(\mathcal{O}_{\mathcal{X}_{flat, fppf}})
\ar@<1ex>[lu]^{Lg_!}
}
$$
où $q$ est le foncteur quotient, le triangle intérieur est commutatif et
$g^{-1}Lg_! = \text{id}$. Pour tout objet $E$ de
$D_{\textit{LQCoh}^{fbc}}(\mathcal{O}_\mathcal{X})$, l'image de l'application
$a : Lg_!g^{-1}E \to E$ est un quasi-isomorphisme dans
$D(\mathcal{O}_{\mathcal{X}_{flat, fppf}})$. Le cône de $a$ est donc envoyé
sur zéro par $g^{-1}$ et, d'après (2)(b), $q(a)$ est un isomorphisme. Ainsi,
$q \circ Lg_!$ est un quasi-inverse de $c$.

\medskip\noindent
Dans le cas du site lisse-\'etale, le même argument démontre que
$$
D_{\textit{LQCoh}^{fbc}(\mathcal{O}_\mathcal{X})}(
\mathcal{X}_\etale, \mathcal{O}_\mathcal{X})
/
D_{\mathcal{P}_\mathcal{X}}(
\mathcal{X}_\etale, \mathcal{O}_\mathcal{X})
$$
est équivalent à
$D_\QCoh(\mathcal{O}_{\mathcal{X}_{lisse,\etale}})$. L'application de la
dernière équivalence du lemme \ref{stacks-perfect-lemma-compare-etale-fppf-QCoh} achève la
démonstration.
\end{proof}

\noindent
Le lemme suivant montre que le foncteur quotient
$D_{\textit{LQCoh}^{fbc}}(\mathcal{O}_\mathcal{X}) \to
D_\QCoh(\mathcal{O}_\mathcal{X})$ admet un adjoint à gauche. Voir la remarque
\ref{stacks-perfect-remark-QCoh-admissible}.

\begin{lemma}
\label{stacks-perfect-lemma-bousfield-colocalization}
Soit $\mathcal{X}$ un champ algébrique. Soit $E$ un objet de
$D_{\textit{LQCoh}^{fbc}}(\mathcal{O}_\mathcal{X})$. Il existe un triangle
distingué canonique
$$
E' \to E \to P \to E'[1]
$$
dans $D_{\textit{LQCoh}^{fbc}}(\mathcal{O}_\mathcal{X})$, tel que $P$
appartienne à $D_{\textit{Parasitic} \cap \textit{LQCoh}^{fbc}}
(\mathcal{O}_\mathcal{X})$
et que
$$
\Hom_{D(\mathcal{O}_\mathcal{X})}(E', P') = 0
$$
pour tout $P'$ appartenant à
$D_{\textit{Parasitic} \cap \textit{LQCoh}^{fbc}}(\mathcal{O}_\mathcal{X})$.
\end{lemma}

\begin{proof}
Considérons le morphisme de topos annelés
$g : \Sh(\mathcal{X}_{flat, fppf}) \longrightarrow \Sh(\mathcal{X}_{fppf})$
étudié dans Cohomologie des champs, section
\ref{stacks-cohomology-section-lisse-etale}.
Posons $E' = Lg_!g^*E$ et soit $P$ le cône de l'application d'adjonction
$E' \to E$ ; voir la partie (4) du lemme \ref{stacks-perfect-lemma-shriek-derived}. D'après
les parties (2)(a) et (2)(c) du lemme
\ref{stacks-perfect-lemma-derived-quasi-coherent}, $E'$ appartient à
$D_{\textit{LQCoh}^{fbc}}(\mathcal{O}_\mathcal{X})$. Par conséquent, $P$
appartient aussi à $D_{\textit{LQCoh}^{fbc}}(\mathcal{O}_\mathcal{X})$.
L'application $g^*E' \to g^*E$ est un isomorphisme, car
$g^*Lg_! = \text{id}$ d'après la partie (4) du lemme
\ref{stacks-perfect-lemma-shriek-derived}. Ainsi, $g^*P = 0$, de sorte que $P$ est un objet de
$D_{\textit{Parasitic} \cap \textit{LQCoh}^{fbc}}(\mathcal{O}_\mathcal{X})$
d'après la partie (2)(b) du lemme \ref{stacks-perfect-lemma-derived-quasi-coherent}. Enfin,
pour $P'$ appartenant à
$D_{\textit{Parasitic} \cap \textit{LQCoh}^{fbc}}(\mathcal{O}_\mathcal{X})$
nous avons
$$
\Hom(E', P') = \Hom(Lg_!g^*E, P') = \Hom(g^*E, g^*P') = 0
$$
car $g^*P' = 0$ d'après la partie (2)(b) du lemme
\ref{stacks-perfect-lemma-derived-quasi-coherent}. Le triangle distingué
$E' \to E \to P \to E'[1]$ est canonique — plus précisément, il est unique à
un isomorphisme de triangles induisant l'identité sur $E$ près — d'après la
discussion de Catégories dérivées, section \ref{derived-section-admissible}.
\end{proof}

\begin{remark}
\label{stacks-perfect-remark-QCoh-admissible}
Le lemme \ref{stacks-perfect-lemma-bousfield-colocalization} montre que
$$
D_{\textit{Parasitic} \cap \textit{LQCoh}^{fbc}}(\mathcal{O}_\mathcal{X})
\subset
D_{\textit{LQCoh}^{fbc}}(\mathcal{O}_\mathcal{X})
$$
est une sous-catégorie admissible à gauche ; voir Catégories dérivées, section
\ref{derived-section-admissible}. En particulier, si
$\mathcal{A} \subset D_{\textit{LQCoh}^{fbc}}(\mathcal{O}_\mathcal{X})$
désigne son orthogonal gauche, alors Catégories dérivées, proposition
\ref{derived-proposition-summarize-admissible}, implique que $\mathcal{A}$ est
admissible à droite dans
$D_{\textit{LQCoh}^{fbc}}(\mathcal{O}_\mathcal{X})$ et que le composé
$$
\mathcal{A} \longrightarrow
D_{\textit{LQCoh}^{fbc}}(\mathcal{O}_\mathcal{X}) \longrightarrow
D_\QCoh(\mathcal{O}_\mathcal{X})
$$
est une équivalence. Cela signifie que nous pouvons considérer
$D_\QCoh(\mathcal{O}_\mathcal{X})$ comme une sous-catégorie triangulée
strictement pleine et saturée de
$D_{\textit{LQCoh}^{fbc}}(\mathcal{O}_\mathcal{X})$, ainsi que de
$D(\mathcal{X}_{fppf}, \mathcal{O}_\mathcal{X})$.
\end{remark}





\section{Image directe dérivée des modules quasi-cohérents}
\label{stacks-perfect-section-derived-pushforward}

\noindent
Comme première application de ce qui précède, nous construisons l'image directe
dérivée. Dans Exemples, section
\ref{examples-section-derived-push-quasi-coherent}, le lecteur trouvera un
exemple de morphisme quasi-compact et quasi-séparé
$f : \mathcal{X} \to \mathcal{Y}$ de champs algébriques tel que le foncteur
image directe $Rf_*$ n'induise pas de foncteur
$D_\QCoh(\mathcal{O}_\mathcal{X}) \to
D_\QCoh(\mathcal{O}_\mathcal{Y})$. Il est donc nécessaire de se restreindre aux
complexes bornés inférieurement.

\begin{proposition}
\label{stacks-perfect-proposition-derived-direct-image-quasi-coherent}
Soit $f : \mathcal{X} \to \mathcal{Y}$ un morphisme quasi-compact et
quasi-séparé de champs algébriques. Le foncteur $Rf_*$ induit un diagramme
commutatif
$$
\xymatrix{
D^{+}_{\textit{Parasitic} \cap \textit{LQCoh}^{fbc}}(\mathcal{O}_\mathcal{X})
\ar[r] \ar[d]^{Rf_*} &
D^{+}_{\textit{LQCoh}^{fbc}}(\mathcal{O}_\mathcal{X})
\ar[r] \ar[d]^{Rf_*} &
D(\mathcal{O}_\mathcal{X})
\ar[d]^{Rf_*} \\
D^{+}_{\textit{Parasitic} \cap \textit{LQCoh}^{fbc}}(\mathcal{O}_\mathcal{Y})
\ar[r] &
D^{+}_{\textit{LQCoh}^{fbc}}(\mathcal{O}_\mathcal{Y}) \ar[r] &
D(\mathcal{O}_\mathcal{Y})
}
$$
et induit donc un foncteur
$$
Rf_{\QCoh, *} :
D^{+}_\QCoh(\mathcal{O}_\mathcal{X})
\longrightarrow
D^{+}_\QCoh(\mathcal{O}_\mathcal{Y})
$$
sur les catégories quotients. De plus, les foncteurs $R^if_\QCoh$ de
Cohomologie des champs, proposition
\ref{stacks-cohomology-proposition-direct-image-quasi-coherent}, sont égaux à
$H^i \circ Rf_{\QCoh, *}$, où $H^i$ est défini dans
(\ref{stacks-perfect-equation-Hi-quasi-coherent}).
\end{proposition}

\begin{proof}
Il faut montrer que $Rf_*E$ est un objet de
$D^{+}_{\textit{LQCoh}^{fbc}}(\mathcal{O}_\mathcal{Y})$ pour $E$ appartenant à
$D^{+}_{\textit{LQCoh}^{fbc}}(\mathcal{O}_\mathcal{X})$. Cela résulte de
Cohomologie des champs, proposition
\ref{stacks-cohomology-proposition-loc-qcoh-flat-base-change}, et de la suite
spectrale $R^if_*H^j(E) \Rightarrow R^{i + j}f_*E$. Le cas des modules
parasites se traite de la même manière à l'aide de Cohomologie des champs,
lemme \ref{stacks-cohomology-lemma-pushforward-parasitic}. La dernière
assertion découle immédiatement de la définition de $H^i$ dans
(\ref{stacks-perfect-equation-Hi-quasi-coherent}).
\end{proof}




\section{Image inverse dérivée des modules quasi-cohérents}
\label{stacks-perfect-section-derived-pullback}

\noindent
L'image inverse dérivée des complexes à faisceaux de cohomologie
quasi-cohérents existe en toute généralité.

\begin{proposition}
\label{stacks-perfect-proposition-derived-pullback-quasi-coherent}
Soit $f : \mathcal{X} \to \mathcal{Y}$ un morphisme de champs algébriques. Le
foncteur exact $f^*$ induit un diagramme commutatif
$$
\xymatrix{
D_{\textit{LQCoh}^{fbc}}(\mathcal{O}_\mathcal{X}) \ar[r] &
D(\mathcal{O}_\mathcal{X}) \\
D_{\textit{LQCoh}^{fbc}}(\mathcal{O}_\mathcal{Y})
\ar[r] \ar[u]^{f^*} &
D(\mathcal{O}_\mathcal{Y}) \ar[u]^{f^*}
}
$$
Le composé
$$
D_{\textit{LQCoh}^{fbc}}(\mathcal{O}_\mathcal{Y})
\xrightarrow{f^*}
D_{\textit{LQCoh}^{fbc}}(\mathcal{O}_\mathcal{X})
\xrightarrow{q_\mathcal{X}}
D_\QCoh(\mathcal{O}_\mathcal{X})
$$
est dérivable à gauche relativement à la localisation
$D_{\textit{LQCoh}^{fbc}}(\mathcal{O}_\mathcal{Y}) \to
D_\QCoh(\mathcal{O}_\mathcal{Y})$
et nous pouvons définir $Lf^*_\QCoh$ comme son foncteur dérivé gauche
$$
Lf_\QCoh^* :
D_\QCoh(\mathcal{O}_\mathcal{Y})
\longrightarrow
D_\QCoh(\mathcal{O}_\mathcal{X})
$$
(voir Catégories dérivées, définitions
\ref{derived-definition-right-derived-functor-defined} et
\ref{derived-definition-everywhere-defined}). Si $f$ est quasi-compact et
quasi-séparé, alors $Lf^*_\QCoh$ et $Rf_{\QCoh, *}$ satisfont à l'adjonction
suivante :
$$
\Hom_{D_\QCoh(\mathcal{O}_\mathcal{X})}(Lf^*_\QCoh A, B)
=
\Hom_{D_\QCoh(\mathcal{O}_\mathcal{Y})}(A, Rf_{\QCoh, *}B)
$$
pour $A \in D_\QCoh(\mathcal{O}_\mathcal{Y})$ et
$B \in D^{+}_\QCoh(\mathcal{O}_\mathcal{X})$.
\end{proposition}

\begin{proof}
Pour démontrer la première assertion, il faut montrer que $f^*E$ est un objet
de $D_{\textit{LQCoh}^{fbc}}(\mathcal{O}_\mathcal{X})$ pour $E$ appartenant à
$D_{\textit{LQCoh}^{fbc}}(\mathcal{O}_\mathcal{Y})$. Puisque
$f^* = f^{-1}$ est exact, cela résulte immédiatement du fait que $f^*$ envoie
$\textit{LQCoh}^{fbc}(\mathcal{O}_\mathcal{Y})$ dans
$\textit{LQCoh}^{fbc}(\mathcal{O}_\mathcal{X})$ d'après Cohomologie des champs,
proposition
\ref{stacks-cohomology-proposition-loc-qcoh-flat-base-change}.

\medskip\noindent
Posons $\mathcal{D} = D_{\textit{LQCoh}^{fbc}}(\mathcal{O}_\mathcal{Y})$. Soit
$S$ la collection des morphismes de $\mathcal{D}$ dont le cône est un objet de
$D_{\textit{Parasitic} \cap \textit{LQCoh}^{fbc}}(\mathcal{O}_\mathcal{Y})$.
Posons $\mathcal{D}' = D_\QCoh(\mathcal{O}_\mathcal{X})$. Posons
$F = q_\mathcal{X} \circ f^* : \mathcal{D} \to \mathcal{D}'$. Alors
$\mathcal{D}, S, \mathcal{D}', F$ sont comme dans Catégories dérivées,
situation \ref{derived-situation-derived-functor} et définition
\ref{derived-definition-right-derived-functor-defined}. Montrons que $LF(E)$
est défini pour tout objet $E$ de $\mathcal{D}$. Considérons en effet le
triangle
$$
E' \to E \to P \to E'[1]
$$
construit au lemme \ref{stacks-perfect-lemma-bousfield-colocalization}. Remarquons que
$s : E' \to E$ est un élément de $S$. Nous affirmons que $E'$ calcule $LF$.
En effet, supposons que $s' : E'' \to E$ soit un autre élément de $S$, c'est-à-dire
qu'il s'insère dans un triangle $E'' \to E \to P' \to E''[1]$ avec $P'$
appartenant à
$D_{\textit{Parasitic} \cap \textit{LQCoh}^{fbc}}(\mathcal{O}_\mathcal{Y})$.
D'après le lemme \ref{stacks-perfect-lemma-bousfield-colocalization} et sa démonstration,
$E' \to E$ se factorise par $E'' \to E$. Nous voyons ainsi que $E' \to E$ est
cofinal dans le système $S/E$. Il est donc clair que $E'$ calcule $LF$.

\medskip\noindent
Pour établir la dernière assertion, écrivons $B = q_\mathcal{X}(H)$ et
$A = q_\mathcal{Y}(E)$. Choisissons $E' \to E$ comme ci-dessus. Nous
utiliserons d'une part l'égalité
$Rf_{\QCoh, *}(B) = q_\mathcal{Y}(Rf_*H)$
et d'autre part l'égalité
$Lf^*_\QCoh(A) = q_\mathcal{X}(f^*E')$.
\begin{align*}
\Hom_{D_\QCoh(\mathcal{O}_\mathcal{X})}(Lf^*_\QCoh A, B)
& = 
\Hom_{D_\QCoh(\mathcal{O}_\mathcal{X})}(q_\mathcal{X}(f^*E'),
q_\mathcal{X}(H)) \\
& = 
\colim_{H \to H'} \Hom_{D(\mathcal{O}_\mathcal{X})}(f^*E', H') \\
& = \colim_{H \to H'} \Hom_{D(\mathcal{O}_\mathcal{Y})}(E', Rf_*H') \\
& = \Hom_{D(\mathcal{O}_\mathcal{Y})}(E', Rf_*H) \\
& =
\Hom_{D_\QCoh(\mathcal{O}_\mathcal{Y})}(A, Rf_{\QCoh, *}B)
\end{align*}
Ici, la colimite porte sur les morphismes $s : H \to H'$ de
$D^+_{\textit{LQCoh}^{fbc}}(\mathcal{O}_\mathcal{X})$
dont le cône $P(s)$ est un objet de
$D^+_{\textit{Parasitic} \cap \textit{LQCoh}^{fbc}}(\mathcal{O}_\mathcal{X})$.
Nous avons vu la première égalité ci-dessus. La deuxième résulte de la
construction du quotient de Verdier. La troisième découle de Cohomologie sur
les sites, lemme \ref{sites-cohomology-lemma-adjoint}. Comme $Rf_*P(s)$ est un
objet de
$D^+_{\textit{Parasitic} \cap \textit{LQCoh}^{fbc}}(\mathcal{O}_\mathcal{Y})$
d'après la proposition
\ref{stacks-perfect-proposition-derived-direct-image-quasi-coherent}, nous avons
$\Hom_{D(\mathcal{O}_\mathcal{Y})}(E', Rf_*P(s)) = 0$. La quatrième égalité est
donc vérifiée. La dernière résulte de la construction de $E'$.
\end{proof}








\section{Objets quasi-cohérents dans la catégorie dérivée}
\label{stacks-perfect-section-QC}

\noindent
Cette section prolonge Faisceaux sur les champs, section
\ref{stacks-sheaves-section-QC}. Soit $\mathcal{X}$ un champ algébrique. Dans
cette section, nous avons défini une catégorie triangulée
$$
\mathit{QC}(\mathcal{X}) = \mathit{QC}(\mathcal{X}_{affine}, \mathcal{O})
$$
et nous avons démontré que, si $\mathcal{X}$ est représentable par un espace
algébrique $X$, alors $\mathit{QC}(\mathcal{X})$ est équivalente à
$D_\QCoh(\mathcal{O}_X)$. Il se trouve que la théorie développée jusqu'ici
suffit exactement à démontrer le même résultat pour tout champ algébrique.

\begin{lemma}
\label{stacks-perfect-lemma-cohomology-parasitic}
Soit $\mathcal{X}$ un champ algébrique. Soit $K$ un objet de
$D(\mathcal{X}_{fppf})$ dont les faisceaux de cohomologie sont parasites. Alors
$R\Gamma(x, K) = 0$ pour tout objet $x$ de $\mathcal{X}$ situé au-dessus d'un
schéma $U$ tel que $U \to \mathcal{X}$ soit plat.
\end{lemma}

\begin{proof}
Notons $g : \Sh(\mathcal{X}_{flat, fppf}) \to \Sh(\mathcal{X}_{fppf})$ le
morphisme de topos considéré à la section \ref{stacks-perfect-section-lisse-etale}. Soit $x$
un objet de $\mathcal{X}$ situé au-dessus d'un schéma $U$ tel que
$U \to \mathcal{X}$ soit plat, c'est-à-dire un objet $x$ de
$\mathcal{X}_{flat, fppf}$. D'après la partie (2)(b) du lemme
\ref{stacks-perfect-lemma-lisse-etale-cohomology}, nous avons
$R\Gamma(x, K) = R\Gamma(\mathcal{X}_{flat, fppf}/x, g^{-1}K)$.
Or notre hypothèse signifie que les faisceaux de cohomologie de l'objet
$g^{-1}K$ de $D(\mathcal{X}_{flat, fppf})$ sont nuls ; voir Cohomologie des
champs, définition \ref{stacks-cohomology-definition-parasitic}. Ainsi,
$g^{-1}K = 0$, ce qui conclut.
\end{proof}

\begin{lemma}
\label{stacks-perfect-lemma-cohomology-parasitic-complex}
Soit $\mathcal{X}$ un champ algébrique. Soit $K$ un objet de
$D(\mathcal{X}_{fppf})$ tel que $R\Gamma(x, K) = 0$ pour tout objet $x$ de
$\mathcal{X}$ situé au-dessus d'un schéma affine $U$ tel que
$U \to \mathcal{X}$ soit plat. Alors $H^i(\mathcal{X}, K) = 0$ pour tout $i$.
\end{lemma}

\begin{proof}
Notons $g : \Sh(\mathcal{X}_{flat, fppf}) \to \Sh(\mathcal{X}_{fppf})$ le
morphisme de topos considéré à la section \ref{stacks-perfect-section-lisse-etale}. D'après la
partie (2)(b) du lemme \ref{stacks-perfect-lemma-lisse-etale-cohomology}, notre hypothèse
signifie que $g^{-1}K$ a une cohomologie nulle sur tout objet de
$\mathcal{X}_{flat, fppf}$ situé au-dessus d'un schéma affine. Comme tout objet
$x$ de $\mathcal{X}_{flat, fppf}$ possède un recouvrement par de tels objets,
nous en déduisons que les faisceaux de cohomologie de $g^{-1}K$ sont nuls,
c'est-à-dire que $g^{-1}K = 0$. Bien entendu,
$R\Gamma(\mathcal{X}_{flat, fppf}, g^{-1}K) = 0$, ce qui donne le résultat voulu
d'après la partie (2)(a) du lemme \ref{stacks-perfect-lemma-lisse-etale-cohomology}.
\end{proof}

\begin{lemma}
\label{stacks-perfect-lemma-higher-shriek-QC}
Soit $\mathcal{X}$ un champ algébrique. Soit $K$ un objet de
$D_\QCoh(\mathcal{O}_{\mathcal{X}_{flat, fppf}})$. Alors $Lg_!K$ possède la
propriété suivante : pour tout morphisme $x \to x'$ de
$\mathcal{X}_{affine}$, l'application
$$
R\Gamma(x', Lg_!K) \otimes_{\mathcal{O}(x')}^\mathbf{L} \mathcal{O}(x)
\longrightarrow
R\Gamma(x, Lg_!K)
$$
est un quasi-isomorphisme.
\end{lemma}

\begin{proof}
D'après la partie (2)(c) du lemme \ref{stacks-perfect-lemma-derived-quasi-coherent}, l'objet
$Lg_!K$ appartient à
$D_{\textit{LQCoh}^{fbc}}(\mathcal{O}_\mathcal{X})$. Il en résulte aussitôt que
l'application affichée dans le lemme est un isomorphisme lorsque
$\mathcal{O}(x') \to \mathcal{O}(x)$ est un homomorphisme d'anneaux plat ; nous
omettons les détails.

\medskip\noindent
Montrons dans ce paragraphe que la question est locale pour la topologie
étale. Soit $x \to x'$ un morphisme arbitraire de $\mathcal{X}_{affine}$. Soit
$\{x'_i \to x'\}$ un recouvrement dans
$\mathcal{X}_{affine, \etale}$. Posons $x_i = x \times_{x'} x'_i$, de sorte que
$\{x_i \to x\}$ soit aussi un recouvrement de
$\mathcal{X}_{affine, \etale}$. Alors
$\mathcal{O}(x') \to \prod \mathcal{O}(x'_i)$ est un homomorphisme d'anneaux
étale fidèlement plat et
$$
\prod \mathcal{O}(x_i) =
\mathcal{O}(x) \otimes_{\mathcal{O}(x')} \left(\prod \mathcal{O}(x'_i)\right)
$$
Un argument élémentaire d'algèbre, que nous omettons, montre donc qu'il suffit
d'établir le résultat de l'énoncé pour chacun des morphismes
$x_i \to x'_i$ de $\mathcal{X}_{affine}$. Autrement dit, le problème est local
pour la topologie étale.

\medskip\noindent
Choisissons un schéma $X$ et un morphisme lisse surjectif
$f : X \to \mathcal{X}$. Nous pouvons considérer $f$ comme un objet de
$\mathcal{X}$, conformément à notre abus de notation ; alors
$(\Sch/X)_{fppf} = \mathcal{X}/f$ ; voir Faisceaux sur les champs, section
\ref{stacks-sheaves-section-restriction}. D'après Faisceaux sur les champs,
lemme
\ref{stacks-sheaves-lemma-surjective-flat-locally-finite-presentation}
par exemple, il existe un recouvrement étale $\{x'_i \to x'\}$ tel que
$x'_i : U'_i = p(x'_i) \to \mathcal{X}$ se factorise par $f$. D'après le
résultat du paragraphe précédent, nous pouvons supposer que $x \to x'$ est
l'image d'un morphisme $U \to U'$ de $(\textit{Aff}/X)_{fppf}$ par le foncteur
$(\Sch/X)_{fppf} \to \mathcal{X}$. À ce stade, nous utilisons le fait que la
restriction à $(\Sch/X)_{fppf}$ de $Lg_!K$ est égale à
$f^*Lg_!K = L(g')_!(f')^*K$ d'après le lemme
\ref{stacks-perfect-lemma-lisse-etale-functorial-derived}. Nous sommes ainsi ramenés au cas
étudié au paragraphe suivant.

\medskip\noindent
Supposons $\mathcal{X} = (\Sch/X)_{fppf}$ et que $x \to x'$ corresponde au
morphisme de schémas affines $U \to U'$. Nous pouvons encore travailler
localement pour la topologie étale, ou de Zariski, sur $U'$ ; nous pouvons donc
supposer que $U' \to X$ se factorise par un ouvert affine de $X$. Nous sommes
ainsi ramenés au cas du paragraphe suivant.

\medskip\noindent
Supposons $\mathcal{X} = (\Sch/X)_{fppf}$, où $X = \Spec(R)$ est un schéma
affine, et que $x \to x'$ corresponde au morphisme de schémas affines
$U \to U'$. Soit $M^\bullet$ un complexe de $R$-modules représentant
$R\Gamma(X, K)$. La construction de Compléments d'algèbre, lemme
\ref{more-algebra-lemma-K-flat-resolution}, permet de supposer que
$M^\bullet = \colim P_n^\bullet$, où chaque $P_n^\bullet$ est un complexe borné
supérieurement de $R$-modules libres. Nous omettons les détails ; voir aussi
Compléments d'algèbre, remarque \ref{more-algebra-remark-P-resolution}.
Considérons le complexe de modules $M^\bullet_{flat, fppf}$ sur
$X_{flat, fppf} = (\Sch/X)_{flat, fppf}$ défini par la règle
$$
U \longmapsto \Gamma(U, M^\bullet \otimes_R \mathcal{O}_U)
$$
C'est un complexe de faisceaux d'après la discussion de Descente, section
\ref{descent-section-quasi-coherent-sheaves}. Il existe une application
canonique $M^\bullet_{flat, fppf} \to K$ qui, d'après les remarques faites au
début de la démonstration, induit un isomorphisme sur les sections au-dessus des
objets affines de $X_{flat, fppf}$. Comme tout objet de $X_{flat, fppf}$ possède
un recouvrement par des objets affines, nous voyons que
$M^\bullet_{flat, fppf}$ coïncide avec $K$.

\medskip\noindent
Soit $M^\bullet_{fppf}$ le complexe de modules sur $X_{fppf}$ donné par la
même formule que ci-dessus. Rappelons que
$Lg_!\mathcal{O} = g_!\mathcal{O} = \mathcal{O}$. Puisque $Lg_!$ est le
foncteur dérivé gauche de $g_!$, nous en déduisons que
$Lg_!P_{n, flat, fppf}^\bullet = P_{n, fppf}^\bullet$. Comme le foncteur $Lg_!$
commute aux colimites homotopiques, ou directement par sa construction dans
Cohomologie sur les sites, lemme
\ref{sites-cohomology-lemma-existence-derived-lower-shriek}, et comme
$M^\bullet = \colim P_n^\bullet$, nous concluons que
$Lg_!M^\bullet_{flat, fppf} = M^\bullet_{fppf}$. Écrivons
$U = \Spec(A)$, $U' = \Spec(A')$ et supposons que $U \to U'$ corresponde à
l'homomorphisme d'anneaux $A' \to A$. D'après ce qui précède,
$$
R\Gamma(U, Lg_!K) = M^\bullet \otimes_R A
\quad\text{et}\quad
R\Gamma(U', Lg_!K) = M^\bullet \otimes_R A'
$$
Puisque $M^\bullet$ est un complexe K-plat de $R$-modules, la transitivité du
produit tensoriel montre que
$$
R\Gamma(U', Lg_!K) \otimes_{A'}^\mathbf{L} A
\longrightarrow
R\Gamma(U, Lg_!K)
$$
est le quasi-isomorphisme recherché.
\end{proof}

\begin{proposition}
\label{stacks-perfect-proposition-QC-compare}
Soit $\mathcal{X}$ un champ algébrique. Alors $\mathit{QC}(\mathcal{X})$ est
canoniquement équivalente à $D_\QCoh(\mathcal{O}_\mathcal{X})$.
\end{proposition}

\begin{proof}
D'après Faisceaux sur les champs, lemme
\ref{stacks-sheaves-lemma-QC-compare-fppf}, l'image inverse par le morphisme de
comparaison
$\epsilon : \mathcal{X}_{affine, fppf} \to \mathcal{X}_{affine}$
identifie $\mathit{QC}(\mathcal{X})$ à une sous-catégorie pleine
$Q_\mathcal{X} \subset D(\mathcal{X}_{affine, fppf}, \mathcal{O})$. À l'aide
de l'équivalence de topos annelés de Faisceaux sur les champs, équation
(\ref{stacks-sheaves-equation-alternative-ringed}), nous pouvons et allons
considérer $Q_\mathcal{X}$ comme une sous-catégorie pleine de
$D(\mathcal{X}_{fppf}, \mathcal{O})$.

\medskip\noindent
De même, le lemme \ref{stacks-perfect-lemma-bousfield-colocalization} et la remarque
\ref{stacks-perfect-remark-QCoh-admissible} montrent que
$D_\QCoh(\mathcal{O}_\mathcal{X})$ peut être considérée comme l'orthogonal
gauche $\mathcal{A}$ de la sous-catégorie admissible à gauche
$D_{\textit{Parasitic} \cap \textit{LQCoh}^{fbc}}(\mathcal{O}_\mathcal{X})$
de $D_{\textit{LQCoh}^{fbc}}(\mathcal{O}_\mathcal{X})$.

\medskip\noindent
Pour conclure, nous montrerons que $Q_\mathcal{X}$ est égale à $\mathcal{A}$
comme sous-catégorie de $D(\mathcal{X}_{fppf}, \mathcal{O})$.

\medskip\noindent
Étape 1 : $Q_\mathcal{X}$ est contenue dans
$D_{\textit{LQCoh}^{fbc}}(\mathcal{O}_\mathcal{X})$.
Un objet $K$ de $Q_\mathcal{X}$ est caractérisé par la propriété suivante :
$K$, considéré comme un objet de
$D(\mathcal{X}_{affine, fppf}, \mathcal{O})$, est tel que $R\epsilon_*K$ soit un
objet de $\mathit{QC}(\mathcal{X}_{affine}, \mathcal{O})$. Cela signifie
exactement que, pour tout morphisme $x \to x'$ de
$\mathcal{X}_{affine}$, l'application
$$
R\Gamma(x', K) \otimes_{\mathcal{O}(x')}^\mathbf{L} \mathcal{O}(x)
\longrightarrow
R\Gamma(x, K)
$$
est un isomorphisme ; voir la note de bas de page dans l'énoncé de Cohomologie
sur les sites, lemme
\ref{sites-cohomology-lemma-cartesion-plus-topology}. Maintenant, si
$x' \to x$ est situé au-dessus d'un morphisme plat de schémas affines, cela
signifie que
$$
H^i(x', K) \otimes_{\mathcal{O}(x')} \mathcal{O}(x)
\cong
H^i(x, K)
$$
Cela signifie clairement que $H^i(K)$ est un faisceau pour la topologie étale
(Faisceaux sur les champs, lemme
\ref{stacks-sheaves-lemma-quasi-coherent-alternative}) et qu'il possède la
propriété de changement de base plat ; nous omettons un petit détail.

\medskip\noindent
Étape 2 : $Q_\mathcal{X}$ est contenue dans $\mathcal{A}$. Il suffit pour cela
de montrer que, pour $K$ dans $Q_\mathcal{X}$, nous avons
$\Hom(K, P) = 0$ pour tout $P$ dans
$D_{\textit{Parasitic} \cap \textit{LQCoh}^{fbc}}(\mathcal{O}_\mathcal{X})$.
Considérons l'objet
$$
H = R\SheafHom_{\mathcal{O}_\mathcal{X}}(K, P)
$$
Soit $x$ un objet de $\mathcal{X}$ situé au-dessus d'un schéma affine
$U = p(x)$. D'après Cohomologie sur les sites, lemme
\ref{sites-cohomology-lemma-section-RHom-over-U}, nous avons la première
égalité de
$$
R\Gamma(x, H) =
R\Hom_{\mathcal{O}_\mathcal{X}}(K|_{\mathcal{X}/x}, P|_{\mathcal{X}/x}) =
R\Hom_{\mathcal{O}}(K|_{\mathcal{X}_{affine}/x}, P|_{\mathcal{X}_{affine}/x})
$$
La seconde égalité vient de ce que le topos du site $\mathcal{X}/x$ est
équivalent au topos du site $\mathcal{X}_{affine}/x$ ; voir Faisceaux sur les
champs, équation (\ref{stacks-sheaves-equation-alternative-ringed}). Nous
pouvons écrire $K = \epsilon^*N$ pour un certain $N$ de
$\mathit{QC}(\mathcal{O})$. Alors Cohomologie sur les sites, lemme
\ref{sites-cohomology-lemma-QC-hom-out-of-plus-topology}, donne
$$
R\Gamma(x, H) =
R\Hom_{D(\mathcal{O}(x))}(R\Gamma(x, N), R\Gamma(x, P))
$$
Le lemme \ref{stacks-perfect-lemma-cohomology-parasitic} montre que
$R\Gamma(x, P) = 0$ si $U \to \mathcal{X}$ est plat, et donc que
$R\Gamma(x, H) = 0$ sous la même hypothèse. D'après le lemme
\ref{stacks-perfect-lemma-cohomology-parasitic-complex}, nous concluons que
$R\Gamma(\mathcal{X}, H) = 0$, et par conséquent que $\Hom(K, P) = 0$.

\medskip\noindent
Étape 3 : $\mathcal{A}$ est contenue dans $Q_\mathcal{X}$. Soit $K$ un objet de
$\mathcal{A}$ et soit $x \to x'$ un morphisme de
$\mathcal{X}_{affine}$. Il faut montrer que
$$
R\Gamma(x', K) \otimes_{\mathcal{O}(x')}^\mathbf{L} \mathcal{O}(x)
\longrightarrow
R\Gamma(x, K)
$$
est un quasi-isomorphisme ; voir la note de bas de page dans l'énoncé de
Cohomologie sur les sites, lemme
\ref{sites-cohomology-lemma-cartesion-plus-topology}. La démonstration du lemme
\ref{stacks-perfect-lemma-bousfield-colocalization} et la discussion de la remarque
\ref{stacks-perfect-remark-QCoh-admissible} montrent que $\mathcal{A}$ est l'image de la
restriction de $Lg_!$ à
$D_\QCoh(\mathcal{O}_{\mathcal{X}_{flat, fppf}})$. Nous pouvons donc supposer
$K = Lg_!M$ pour un certain $M$ de
$D_\QCoh(\mathcal{O}_{\mathcal{X}_{flat, fppf}})$. L'égalité recherchée résulte
alors du lemme \ref{stacks-perfect-lemma-higher-shriek-QC}.
\end{proof}






















% OK, start here.
%

\chapter{Introduction aux champs algébriques}


\phantomsection
\label{stacks-introduction-section-phantom}





\section{Pourquoi lire ce chapitre ?}
\label{stacks-introduction-section-introduction}

\noindent
Nous donnons une introduction informelle aux champs algébriques. Le but est
d'introduire rapidement un langage simple qui permette de penser les
propriétés locales et globales de votre problème de modules favori.
Une fois cela fait, vous devriez pouvoir formuler des questions bien posées
sur les problèmes de modules et commencer à les résoudre, en admettant
l'existence d'une théorie générale. Si vous aboutissez à un résultat
intéressant, vous pourrez revenir à la théorie générale développée dans les
autres parties du projet Champs et combler les lacunes selon les besoins.

\medskip\noindent
Le point de vue adopté ici est proche de celui de
\cite{KatzMazur} et \cite{mumford_picard}.






\section{Préliminaires}
\label{stacks-introduction-section-preliminary}

\noindent
Soit $S$ un schéma. Une {\it courbe elliptique} sur $S$ est un triplet
$(E, f, 0)$, où $E$ est un schéma et où $f : E \to S$ et $0 : S \to E$
sont des morphismes de schémas tels que
\begin{enumerate}
\item $f : E \to S$ soit propre et lisse de dimension relative $1$ ;
\item pour tout $s \in S$, la fibre $E_s$ soit une courbe connexe
de genre $1$, c'est-à-dire que $H^0(E_s, \mathcal{O})$ et
$H^1(E_s, \mathcal{O})$ soient tous deux des espaces vectoriels de dimension
$1$ sur $\kappa(s)$ ;
\item $0$ soit une section de $f$.
\end{enumerate}
Étant données des courbes elliptiques $(E, f, 0)/S$ et $(E', f', 0')/S'$, un
{\it morphisme de courbes elliptiques au-dessus de $a : S \to S'$}
est un morphisme $\alpha : E \to E'$ tel que le diagramme
$$
\xymatrix{
E \ar[rr]_\alpha \ar[d]^f & & E' \ar[d]_{f'}  \\
S \ar@/^5ex/[u]^0 \ar[rr]^a & & S' \ar@/_5ex/[u]_{0'}
}
$$
soit commutatif et que le carré intérieur soit cartésien ; autrement dit, le
morphisme $\alpha$ induit un isomorphisme $E \to S \times_{S'} E'$.
Nous allons définir le champ des courbes elliptiques $\mathcal{M}_{1, 1}$.
Dans le reste du projet Champs, nous développons la méthode introduite dans
l'article de Deligne et Mumford \cite{DM}, qui consiste à présenter
$\mathcal{M}_{1, 1}$ comme une catégorie munie d'un foncteur
$$
p : \mathcal{M}_{1, 1} \longrightarrow \Sch, \quad
(E, f, 0)/S \longmapsto S
$$
Cela conduit à travailler avec les catégories fibrées sur la catégorie des
schémas, les topologies, les champs fibrés en groupoïdes, les recouvrements,
etc. Dans ce chapitre, nous laissons tout cet appareil de côté et envisageons
les choses un peu autrement, probablement d'une manière plus proche de celle
dont les fondateurs de la théorie ont commencé à les concevoir.


\section{Le champ de modules des courbes elliptiques}
\label{stacks-introduction-section-moduli-elliptic-curves}

\noindent
Voici la démarche que nous allons suivre :
\begin{enumerate}
\item Partons de votre catégorie de schémas favorite $\Sch$.
\item Ajoutons un nouveau symbole $\mathcal{M}_{1, 1}$.
\item Un morphisme $S \to \mathcal{M}_{1, 1}$ {\bf est} une courbe elliptique
$(E, f, 0)$ sur $S$.
\item Un diagramme
$$
\xymatrix{
S \ar[rr]_a \ar[rd]_{(E, f, 0)} & & S' \ar[ld]^{(E', F', 0')} \\
& \mathcal{M}_{1, 1}
}
$$
{\bf est} commutatif si et seulement s'il existe un morphisme
$\alpha : E \to E'$ de courbes elliptiques au-dessus de $a : S \to S'$.
Nous disons que $\alpha$ {\it atteste} la commutativité du diagramme.
\item Remarquons que les diagrammes commutatifs se composent comme suit :
$$
\xymatrix{
S \ar[rrr]_a \ar[rrrd]_{(E, f, 0)} & & &
S' \ar[d]_{(E', F', 0')} \ar[rrr]_{a'} & & &
S'' \ar[llld]^{(E'', F'', 0'')}
\\
& & & \mathcal{M}_{1, 1}
}
$$
en effet, $\alpha' \circ \alpha$ atteste la commutativité du triangle extérieur
si $\alpha$ et $\alpha'$ attestent respectivement la commutativité des triangles
de gauche et de droite.
\item La composée
$$
S \xrightarrow{a} S' \xrightarrow{(E', f', 0')} \mathcal{M}_{1, 1}
$$
est donnée par $(E' \times_{S'} S, f' \times_{S'} S, 0' \times_{S'} S)$.
\end{enumerate}
Au terme de cette procédure, nous avons agrandi la catégorie $\Sch$ des
schémas en lui adjoignant exactement un objet...

\medskip\noindent
À ceci près que nous n'avons pas défini ce qu'est un morphisme de
$\mathcal{M}_{1, 1}$ vers un schéma $T$. La réponse consiste à retenir la
notion la plus faible possible pour laquelle les compositions aient un sens.
Ainsi, un morphisme $F : \mathcal{M}_{1, 1} \to T$ est une règle qui associe à
toute courbe elliptique $(E, f, 0)/S$ un morphisme
$F(E, f, 0) : S \to T$, de telle sorte que, pour tout diagramme commutatif
$$
\xymatrix{
S \ar[rr]_a \ar[rd]_{(E, f, 0)} & & S' \ar[ld]^{(E', F', 0')} \\
& \mathcal{M}_{1, 1}
}
$$
le diagramme
$$
\xymatrix{
S \ar[rr]_a \ar[rd]_{F(E, f, 0)} & & S' \ar[ld]^{F(E', F', 0')} \\
& T
}
$$
soit lui aussi commutatif. Le $j$-invariant, dont vous avez peut-être entendu
parler, en fournit un exemple :
$$
j : \mathcal{M}_{1, 1} \longrightarrow \mathbf{A}^1_{\mathbf{Z}}
$$
Bien, nous avons donc terminé...

\medskip\noindent
Eh bien non ! Il nous faut encore définir une notion de morphisme
$\mathcal{M}_{1, 1} \to \mathcal{M}_{1, 1}$. Nous procédons exactement comme
précédemment : un morphisme
$F : \mathcal{M}_{1, 1} \to \mathcal{M}_{1, 1}$
est une règle qui associe à toute courbe elliptique $(E, f, 0)/S$ une autre
courbe elliptique $F(E, f, 0)$, tout en préservant la commutativité des
diagrammes précédents. Toutefois, comme je ne connais aucun exemple non
trivial d'un tel foncteur, je définirai pour l'instant l'ensemble des
morphismes de $\mathcal{M}_{1, 1}$ vers lui-même comme réduit à l'identité.

\medskip\noindent
J'espère que vous voyez comment ajouter d'autres objets à cette catégorie
agrandie. Il semble intuitivement clair que, pour tout problème de modules
« bien sage », nous pouvons effectuer la construction précédente et ajouter
un objet à notre catégorie. En fait, une grande partie de la géométrie
algébrique contemporaine se déroule dans un tel univers, où $\Sch$ est
agrandie par une infinité dénombrable de champs de modules construits
explicitement.

\medskip\noindent
Vous pourriez objecter que la structure obtenue n'est pas une catégorie,
car il subsiste un certain « flou » quant aux diagrammes qui commutent et aux
combinaisons de diagrammes qui continuent de commuter : il faut en effet
produire une attestation de la commutativité. Or il se trouve que cette idée,
qui consiste à munir la commutativité d'attestations, constitue une approche
valable des $2$-catégories ! Nous nous y tiendrons donc.






\section{Produits fibrés}
\label{stacks-introduction-section-fibre-products}

\noindent
La question posée ici est de déterminer ce que doit être le produit fibré
$$
\xymatrix{
& ? \ar@{..>}[rd] \ar@{..>}[ld] \\
S \ar[rd]_{(E, f, 0)} & & S' \ar[ld]^{(E', f', 0')} \\
& \mathcal{M}_{1, 1}
}
$$
La réponse est la suivante : un morphisme d'un schéma $T$ vers $?$
doit être un triplet $(a, a', \alpha)$, où
$a : T \to S$ et $a' : T \to S'$ sont des morphismes de schémas et où
$\alpha : E \times_{S, a} T \to E' \times_{S', a'} T$ est un isomorphisme de
courbes elliptiques sur $T$. Cela a un sens en vertu de notre définition
antérieure de la composition et des diagrammes commutatifs.

\begin{lemma}[Fait clé]
\label{stacks-introduction-lemma-key-fact}
Le foncteur $\Sch^{opp} \to \textit{Ensembles}$,
$T \mapsto \{(a, a', \alpha)\text{ comme ci-dessus}\}$
est représentable par un schéma $S \times_{\mathcal{M}_{1, 1}} S'$.
\end{lemma}

\begin{proof}
Idée de la démonstration. Relier ce foncteur à
$$
\mathit{Isom}_{S \times S'}(E \times S', S \times E')
$$
et utiliser la théorie de Grothendieck des schémas de Hilbert.
\end{proof}

\begin{remark}
\label{stacks-introduction-remark-diagonal}
Nous avons la formule
$S \times_{\mathcal{M}_{1, 1}} S' =
(S \times S')
\times_{\mathcal{M}_{1, 1} \times \mathcal{M}_{1, 1}}
\mathcal{M}_{1, 1}$.
Le fait clé est donc une propriété de la diagonale
$\Delta_{\mathcal{M}_{1, 1}}$ de $\mathcal{M}_{1, 1}$.
\end{remark}

\noindent
Quoi qu'il en soit, le fait clé nous permet de poser la définition suivante.

\begin{definition}
\label{stacks-introduction-definition-smooth}
Nous disons qu'un morphisme $S \to \mathcal{M}_{1, 1}$ est {\it lisse} si,
pour tout morphisme $S' \to \mathcal{M}_{1, 1}$, le morphisme de projection
$$
S \times_{\mathcal{M}_{1, 1}} S' \longrightarrow S'
$$
est lisse.
\end{definition}

\noindent
Remarquons que cette définition est compatible avec la notion de morphisme
lisse de schémas, puisque tout changement de base d'un morphisme lisse est
lisse. Il est en outre clair comment étendre cette définition à d'autres
propriétés des morphismes vers $\mathcal{M}_{1, 1}$ (ou vers votre champ de
modules favori). Nous l'utiliserons notamment ci-dessous pour les morphismes
{\it surjectifs}.





\section{La définition}
\label{stacks-introduction-section-definition}

\noindent
Nous allons la formuler comme une définition plutôt que comme un résultat,
car nous attendons du lecteur qu'il essaie d'autres cas, et pas seulement le
champ $\mathcal{M}_{1, 1}$ ni la seule catégorie $\Sch$ de tous les schémas.

\begin{definition}
\label{stacks-introduction-definition-algebraic-stack}
Nous disons que $\mathcal{M}_{1, 1}$ est un {\it champ algébrique} si et
seulement si
\begin{enumerate}
\item les objets vérifient la descente pour la topologie \'etale sur $\Sch$ ;
\item le fait clé est vérifié ;
\item il existe un morphisme surjectif et lisse
$S \to \mathcal{M}_{1, 1}$.
\end{enumerate}
\end{definition}

\noindent
La première condition est une « propriété de faisceau ». Nous allons
l'expliciter, car elle recèle un point technique. Supposons donnés un schéma
$S$, un recouvrement \'etale $\{S_i \to S\}$ et des morphismes
$e_i : S_i \to \mathcal{M}_{1, 1}$ tels que les diagrammes
$$
\xymatrix{
S_i \times_S S_j \ar[rd]_{e_i \circ \text{pr}_1} \ar[rr]_{\text{id}} & &
S_i \times_S S_j \ar[ld]^{e_j \circ \text{pr}_2} \\
& \mathcal{M}_{1, 1}
}
$$
soient commutatifs. La condition de faisceau ne garantit {\it pas}, à elle
seule, l'existence d'un morphisme $e : S \to \mathcal{M}_{1, 1}$ dans cette
situation. En effet, il faut choisir des attestations $\alpha_{ij}$ pour les
diagrammes précédents et imposer que
$$
\text{pr}_{02}^*\alpha_{ik} =
\text{pr}_{12}^*\alpha_{jk} \circ \text{pr}_{01}^*\alpha_{ij}
$$
comme attestations sur $S_i \times_S S_j \times_S S_k$. Le sens de cette
condition devrait être clair... Sinon, je crains qu'il ne vous faille lire
une partie des développements sur les catégories fibrées en groupoïdes, etc.
Quoi qu'il en soit, l'égalité affichée est souvent appelée la
{\it condition de cocycle}. Voici une formulation plus précise de la
« propriété de faisceau » : étant donnés $\{S_i \to S\}$,
$e_i : S_i \to \mathcal{M}_{1, 1}$ et des attestations $\alpha_{ij}$
satisfaisant la condition de cocycle, il existe un unique morphisme
$e : S \to \mathcal{M}_{1, 1}$, à isomorphisme unique près, tel que
$e_i \cong e|_{S_i}$ et qui redonne les $\alpha_{ij}$.

\medskip\noindent
Comme vous le voyez, la formulation même d'un énoncé précis demande déjà un
peu de travail. La démonstration de cette « propriété de faisceau » repose sur
une technique fondamentale de la géométrie algébrique : la théorie de la
descente. Je suggère, dans un premier temps, d'admettre simplement cette
« propriété de faisceau » et d'examiner ce qu'elle implique en pratique.
Une certaine agilité d'esprit est en effet nécessaire pour condenser cette
« propriété de faisceau » en un énoncé maniable tenant sur un coin de table.
La variante la plus simple qui présente déjà quelque intérêt est peut-être la
suivante. Supposons donnée une extension galoisienne finie $L/K$ de corps, de
groupe de Galois $G = \text{Gal}(L/K)$. Posons $T = \Spec(L)$ et
$S = \Spec(K)$. Alors $\{T \to S\}$ est un recouvrement \'etale.
Soit $(E, f, 0)$ une courbe elliptique sur $L$. (Cela signifie simplement que
$E \subset \mathbf{P}^2_L$ est donnée par une équation de Weierstrass et que
$0$ est le point à l'infini usuel.) Notons
$E_\sigma = E \times_{T, \Spec(\sigma)} T$ le changement de base.
(Cela revient à appliquer $\sigma$ aux coefficients de l'équation de
Weierstrass, à moins que ce ne soit $\sigma^{-1}$.) Supposons en outre que,
pour tout $\sigma \in G$, nous soit donné un isomorphisme
$$
\alpha_\sigma : E \longrightarrow E_\sigma
$$
sur $T$. Dans cette situation, la condition de cocycle signifie que
$$
(\alpha_\tau)^\sigma \circ \alpha_\sigma = \alpha_{\tau\sigma}
$$
pour $\sigma, \tau \in G$. Si vous avez déjà pratiqué la cohomologie des
groupes, cette formule vous sera familière. La condition de « recollement »
sur $\mathcal{M}_{1, 1}$ affirme donc que, si ce système d'équations admet une
solution, il existe une courbe elliptique $E'$ sur $S$ telle que
$E \cong E' \times_S T$ (elle en dit un peu plus, puisqu'elle indique aussi
comment retrouver les $\alpha_\sigma$).

\medskip\noindent
Défi : pouvez-vous le démontrer en n'utilisant que des courbes elliptiques
définies par des équations de Weierstrass ?




\section{Un recouvrement lisse}
\label{stacks-introduction-section-smooth}

\noindent
Il nous reste à trouver un recouvrement lisse de $\mathcal{M}_{1, 1}$.
En un certain sens, l'existence d'un recouvrement lisse
{\it implique}\footnote{Il y a là une légère tricherie : pour vérifier la
lissité, il faut démontrer un résultat proche du fait clé, puisque la lissité
est définie au moyen de produits fibrés. L'avantage est qu'il suffit de
démontrer l'existence de ces produits fibrés lorsque, d'un côté, figure le
morphisme dont on cherche précisément à montrer qu'il fournit le recouvrement
lisse.} le fait clé ! Dans le cas des courbes elliptiques, nous en construisons
un au moyen de l'équation de Weierstrass.

\medskip\noindent
Posons
$$
W = \Spec(\mathbf{Z}[a_1, a_2, a_3, a_4, a_6, 1/\Delta])
$$
où $\Delta \in \mathbf{Z}[a_1, a_2, a_3, a_4, a_6]$ est un certain polynôme
(voir ci-dessous). Posons
$$
\mathbf{P}_W^2 \supset E_W :
zy^2 + a_1 xyz + a_3 yz^2 = x^3 + a_2x^2z + a_4xz^2 + a_6z^3.
$$
Notons $f_W : E_W \to W$ la projection. Enfin, notons $0_W : W \to E_W$ la
section de $f_W$ donnée par $(0 : 1 : 0)$. Il existe alors un polynôme homogène
de degré $12$, noté $\Delta$, en les $a_i$, où $\deg(a_i) = i$, tel que
$E_W \to W$ soit lisse. On peut l'obtenir explicitement en calculant les
dérivées partielles de l'équation de Weierstrass ; on peut bien sûr aussi le
chercher dans la littérature. On peut encore demander à pari/gp de le
calculer. Le voici :
\begin{align*}
\Delta & = -a_6a_1^6 + a_4a_3a_1^5 + ((-a_3^2 - 12a_6)a_2 + a_4^2)a_1^4 + \\
& (8a_4a_3a_2 + (a_3^3 + 36a_6a_3))a_1^3 + \\
& ((-8a_3^2 - 48a_6)a_2^2 + 8a_4^2a_2 + (-30a_4a_3^2 + 72a_6a_4))a_1^2 + \\
& (16a_4a_3a_2^2 + (36a_3^3 + 144a_6a_3)a_2 - 96a_4^2a_3)a_1 + \\
& (-16a_3^2 - 64a_6)a_2^3 + 16a_4^2a_2^2 + (72a_4a_3^2 + 288a_6a_4)a_2 + \\
& -27a_3^4 - 216a_6a_3^2  -64a_4^3 - 432a_6^2
\end{align*}
Vous reconnaîtrez peut-être les deux derniers termes du cas
$y^2 = x^3 + Ax + B$, où le discriminant vaut
$-64A^3 - 432B^2 = -16(4A^3 + 27B^2)$.

\begin{lemma}
\label{stacks-introduction-lemma-Weierstrass-smooth-cover}
Le morphisme $W \xrightarrow{(E_W, f_W, 0_W)} \mathcal{M}_{1, 1}$ est lisse
et surjectif.
\end{lemma}

\begin{proof}
La surjectivité résulte du fait que toute courbe elliptique sur un corps admet
une équation de Weierstrass. Esquissons une manière de démontrer la lissité.
Considérons le sous-schéma en groupes
$$
H =
\left\{
\left(
\begin{matrix}
u^2 & s & 0 \\
0 & u^3 & 0 \\
r & t & 1
\end{matrix}
\right)
\middle|
\begin{matrix}
u\text{ inversible} \\
s, r, t\text{ arbitraires}
\end{matrix}
\right\}
\subset
\text{GL}_{3, \mathbf{Z}}
$$
Il existe une action $H \times W \to W$ de $H$ sur le schéma de Weierstrass
$W$. Pour en trouver les équations, il suffit d'expliciter l'effet, sur
l'équation générale de Weierstrass, d'un changement de coordonnées donné par
une matrice de $H$. On constate alors que les assertions suivantes sont vraies :
\begin{enumerate}
\item toute courbe elliptique $(E, f, 0)/S$ admet, localement pour la topologie
de Zariski sur $S$, une équation de Weierstrass ;
\item deux équations de Weierstrass quelconques de $(E, f, 0)$ diffèrent,
localement pour la topologie de Zariski, par un élément de $H$.
\end{enumerate}
En considérant le produit fibré
$S \times_{\mathcal{M}_{1, 1}} W =
\mathit{Isom}_{S \times W}(E \times W, S \times E_W)$,
nous en concluons que le morphisme $W \to \mathcal{M}_{1, 1}$ est un
$H$-torseur. Puisque $H \to \Spec(\mathbf{Z})$ est lisse et que les torseurs
sous des schémas en groupes lisses sont lisses, le résultat s'ensuit.
\end{proof}

\begin{remark}
\label{stacks-introduction-remark-quotient-stack}
L'argument esquissé ci-dessus montre en fait que
$\mathcal{M}_{1, 1} = [W/H]$ est un champ quotient global.
Dans environ 50\% des cas, un argument établissant qu'un champ de modules est
algébrique montre aussi qu'il s'agit d'un champ quotient global.
\end{remark}



\section{Propriétés des champs algébriques}
\label{stacks-introduction-section-properties}

\noindent
Nous savons maintenant que $\mathcal{M}_{1, 1}$ est un champ algébrique.
Que pouvons-nous en faire ? Ce n'est pas tant le fait qu'il soit algébrique
qui nous aide ici que le point de vue selon lequel les propriétés de
$\mathcal{M}_{1, 1}$ doivent être encodées dans celles des morphismes
$S \to \mathcal{M}_{1, 1}$, c'est-à-dire dans les familles de courbes
elliptiques. En voici quelques exemples.

\medskip\noindent
{\bf Propriétés locales :}
$$
\mathcal{M}_{1, 1} \to \Spec(\mathbf{Z})\text{ est lisse}
\Leftrightarrow
W \to \Spec(\mathbf{Z})\text{ est lisse}
$$
{\bf Idée}. Les propriétés locales d'un champ algébrique sont encodées dans
les propriétés locales de son recouvrement lisse.

\medskip\noindent
{\bf Propriétés globales :}
$$
\begin{matrix}
\mathcal{M}_{1, 1}\text{ est quasi-compact} \Leftarrow W\text{ est quasi-compact}
\\
\mathcal{M}_{1, 1}\text{ est irréductible} \Leftarrow W\text{ est irréductible}
\end{matrix}
$$
{\bf Idée}. Certaines propriétés globales d'un champ algébrique se lisent sur
la propriété correspondante d'un recouvrement lisse
{\it convenable}\footnote{Je suppose qu'il peut exister un champ algébrique
irréductible dépourvu de recouvrement lisse irréductible ; mais, si tel est le
cas, il doit être assez pathologique !}.

\medskip\noindent
{\bf Faisceaux quasi-cohérents :}
$$
\QCoh(\mathcal{O}_{\mathcal{M}_{1, 1}}) =
H\text{-modules quasi-cohérents équivariants sur }W
$$
{\bf Idée}. D'une part, un module quasi-cohérent sur
$\mathcal{M}_{1, 1}$ devrait correspondre à un faisceau quasi-cohérent
$\mathcal{F}_{S, e}$ sur $S$ pour chaque morphisme
$e : S \to \mathcal{M}_{1, 1}$. Cela vaut en particulier pour le morphisme
$(E_W, f_W, 0_W) :  W \to \mathcal{M}_{1, 1}$. Comme ce morphisme est
$H$-équivariant, le module quasi-cohérent $\mathcal{F}_W$ ainsi obtenu est
$H$-équivariant. Réciproquement, à partir d'un module $H$-équivariant, on peut
retrouver les faisceaux $\mathcal{F}_{S, e}$ par la théorie de la descente, en
partant de l'observation que $S \times_{e, \mathcal{M}_{1, 1}} W$ est un
$H$-torseur.

\medskip\noindent
{\bf Groupe de Picard :}
$$
\Pic(\mathcal{M}_{1, 1}) = \Pic_H(W) = \mathbf{Z}/12\mathbf{Z}
$$
{\bf Idée}. Nous avons vu ci-dessus la première égalité. Remarquons que
$\Pic(W) = 0$, car l'anneau
$\mathbf{Z}[a_1, a_2, a_3, a_4, a_6, 1/\Delta]$
a un groupe des classes trivial. Il existe une suite exacte
$$
\mathbf{Z}\Delta \to
\Pic_H(\mathbf{A}^5_{\mathbf{Z}}) \to
\Pic_H(W) \to 0
$$
Le groupe du milieu est égal à $\Hom(H, \mathbf{G}_m) = \mathbf{Z}$.
L'image de $\Delta$ est $12$, car $\Delta$ est de degré $12$.
Cet argument est essentiellement correct ; voir \cite{PicM11}.

\medskip\noindent
{\bf Cohomologie \'etale :} soit $\Lambda$ un anneau.
Il existe une suite spectrale du premier quadrant convergeant vers
$H^{p + q}_\etale(\mathcal{M}_{1, 1}, \Lambda)$, dont la page $E_2$ est
$$
E_2^{p, q} = H_\etale^q(W \times H \times \ldots \times H, \Lambda)
\quad(p\text{ facteurs }H)
$$
{\bf Idée}. Remarquons que
$$
W \times_{\mathcal{M}_{1, 1}} W \times_{\mathcal{M}_{1, 1}}
\ldots \times_{\mathcal{M}_{1, 1}} W
= W \times H \times \ldots \times H
$$
parce que $W \to \mathcal{M}_{1, 1}$ est un $H$-torseur. Cette suite
spectrale est la suite spectrale de {\v C}ech vers la cohomologie associée au
recouvrement lisse $\{W \to \mathcal{M}_{1, 1}\}$. Par exemple,
$H^0_\etale(\mathcal{M}_{1, 1}, \Lambda) = \Lambda$ parce que
$W$ est connexe, et $H^1_\etale(\mathcal{M}_{1, 1}, \Lambda) = 0$ parce que
$H^1_\etale(W, \Lambda) = 0$ (cela exige évidemment une démonstration).
Bien entendu, le recouvrement lisse $W \to \mathcal{M}_{1, 1}$ n'est
peut-être pas « optimal » pour calculer la cohomologie \'etale.


















% OK, start here.
%

\chapter{Compléments sur les morphismes de champs}



\phantomsection
\label{stacks-more-morphisms-section-phantom}


\section{Introduction}
\label{stacks-more-morphisms-section-introduction}

\noindent
Dans ce chapitre, nous poursuivons l'étude des propriétés des morphismes de
champs algébriques. Une référence, dans le cas des champs algébriques quasi-séparés
à diagonale représentable, est \cite{LM-B}.




\section{Conventions et abus de langage}
\label{stacks-more-morphisms-section-conventions}

\noindent
Nous continuons d'employer les conventions et les abus de langage
introduits dans
Propriétés des champs, section \ref{stacks-properties-section-conventions}.






\section{Épaississements}
\label{stacks-more-morphisms-section-thickenings}

\noindent
La terminologie suivante n'est peut-être pas tout à fait standard, mais elle est commode.
Si $\mathcal{Y}$ est un sous-champ fermé d'un champ algébrique $\mathcal{X}$,
alors le morphisme $\mathcal{Y} \to \mathcal{X}$ est représentable.

\begin{definition}
\label{stacks-more-morphisms-definition-thickening}
Épaississements.
\begin{enumerate}
\item On dit qu'un champ algébrique $\mathcal{X}'$ est un {\it épaississement}
d'un champ algébrique $\mathcal{X}$ si $\mathcal{X}$ est un sous-champ fermé
de $\mathcal{X}'$ et si les espaces topologiques associés sont égaux.
\item Étant donnés deux épaississements $\mathcal{X} \subset \mathcal{X}'$ et
$\mathcal{Y} \subset \mathcal{Y}'$, un {\it morphisme d'épaississements}
est un morphisme $f' : \mathcal{X}' \to \mathcal{Y}'$ de champs algébriques
tel que $f'|_\mathcal{X}$ se factorise par le
sous-champ fermé $\mathcal{Y}$. Dans cette situation, on pose
$f = f'|_\mathcal{X} : \mathcal{X} \to \mathcal{Y}$ et on dit que
$(f, f') : (\mathcal{X} \subset \mathcal{X}') \to
(\mathcal{Y} \subset \mathcal{Y}')$ est un morphisme d'épaississements.
\item Soit $\mathcal{Z}$ un champ algébrique. On définit de même les
{\it épaississements sur $\mathcal{Z}$} et les
{\it morphismes d'épaississements sur $\mathcal{Z}$}.
Cela signifie que les champs algébriques
$\mathcal{X}'$ et $\mathcal{Y}'$
sont munis d'un
morphisme structural vers $\mathcal{Z}$ et que $f'$ s'insère dans un
diagramme $2$-commutatif convenable de champs algébriques.
\end{enumerate}
\end{definition}

\noindent
Soit $\mathcal{X} \subset \mathcal{X}'$ un épaississement de champs algébriques.
Soit $U'$ un schéma et soit $U' \to \mathcal{X}'$ un morphisme lisse
surjectif. En posant $U = \mathcal{X} \times_{\mathcal{X}'} U'$, on obtient
un morphisme d'épaississements
$$
(U \subset U') \longrightarrow (\mathcal{X} \subset \mathcal{X}')
$$
et $U \to \mathcal{X}$ est un morphisme lisse surjectif. On peut souvent
déduire des propriétés de l'épaississement $\mathcal{X} \subset \mathcal{X}'$
des propriétés correspondantes de l'épaississement $U \subset U'$.
Parfois, par abus de langage, on dit qu'un morphisme
$\mathcal{X} \to \mathcal{X}'$ est un épaississement s'il s'agit d'une immersion
fermée induisant une bijection $|\mathcal{X}| \to |\mathcal{X}'|$.

\begin{lemma}
\label{stacks-more-morphisms-lemma-thickening}
Soit $i : \mathcal{X} \to \mathcal{X}'$ un morphisme de champs algébriques.
Les conditions suivantes sont équivalentes :
\begin{enumerate}
\item $i$ est un épaississement de champs algébriques (au sens de l'abus de langage ci-dessus), et
\item $i$ est représentable par des espaces algébriques et
est un épaississement au sens de Propriétés des champs, section
\ref{stacks-properties-section-properties-morphisms}.
\end{enumerate}
Dans ce cas, $i$ est une immersion fermée et un homéomorphisme universel.
\end{lemma}

\begin{proof}
D'après Compléments sur les morphismes d'espaces, lemmes
\ref{spaces-more-morphisms-lemma-descending-property-thickening} et
\ref{spaces-more-morphisms-lemma-base-change-thickening}
la propriété $P$ qu'un morphisme d'espaces algébriques soit un
épaississement (du premier ordre) est locale pour la topologie fpqc sur la base et stable par
changement de base. La discussion de Propriétés des champs, section
\ref{stacks-properties-section-properties-morphisms} s'applique donc bien.
Cela étant dit, l'équivalence de (1) et (2) résulte du
fait que $P = P_1 + P_2$, où $P_1$ est la propriété d'être
une immersion fermée et $P_2$ celle d'être surjectif.
(À strictement parler, le lecteur devrait également consulter
Compléments sur les morphismes d'espaces, définition
\ref{spaces-more-morphisms-definition-thickening},
Propriétés des champs, définition \ref{stacks-properties-definition-immersion}
et la discussion qui suit, Morphismes d'espaces, lemme
\ref{spaces-morphisms-lemma-surjective-representable},
Propriétés des champs, section \ref{stacks-properties-section-surjective},
pour constater que toutes les notions coïncident.)
La dernière assertion est claire d'après ce qui précède.
\end{proof}

\noindent
Nous utiliserons le lemme sans autre mention. Les mêmes références,
Compléments sur les morphismes d'espaces, lemmes
\ref{spaces-more-morphisms-lemma-descending-property-thickening} et
\ref{spaces-more-morphisms-lemma-base-change-thickening}
permettent de définir comme suit un épaississement du premier ordre.

\begin{definition}
\label{stacks-more-morphisms-definition-first-order-thickening}
On dit qu'un champ algébrique $\mathcal{X}'$ est un {\it épaississement du premier ordre}
d'un champ algébrique $\mathcal{X}$ si $\mathcal{X}$ est un sous-champ fermé
de $\mathcal{X}'$ et si $\mathcal{X} \to \mathcal{X}'$ est un épaississement
du premier ordre au sens de Propriétés des champs, section
\ref{stacks-properties-section-properties-morphisms}.
\end{definition}

\noindent
Si $(U \subset U') \to (\mathcal{X} \subset \mathcal{X}')$ est un
recouvrement lisse par un schéma comme ci-dessus, cela signifie simplement que $U \subset U'$
est un épaississement du premier ordre. Formulons maintenant les lemmes indispensables.

\begin{lemma}
\label{stacks-more-morphisms-lemma-base-change-thickening}
Soit $\mathcal{Y} \subset \mathcal{Y}'$ un épaississement de champs algébriques.
Soit $\mathcal{X}' \to \mathcal{Y}'$ un morphisme de champs algébriques
et posons $\mathcal{X} = \mathcal{Y} \times_{\mathcal{Y}'} \mathcal{X}'$.
Alors
$(\mathcal{X} \subset \mathcal{X}') \to (\mathcal{Y} \subset \mathcal{Y}')$
est un morphisme d'épaississements. Si $\mathcal{Y} \subset \mathcal{Y}'$ est un
épaississement du premier ordre, alors $\mathcal{X} \subset \mathcal{X}'$ est un
épaississement du premier ordre.
\end{lemma}

\begin{proof}
Voir la discussion ci-dessus, Propriétés des champs, section
\ref{stacks-properties-section-properties-morphisms}, et
et Compléments sur les morphismes d'espaces, lemme
\ref{spaces-more-morphisms-lemma-base-change-thickening}.
\end{proof}

\begin{lemma}
\label{stacks-more-morphisms-lemma-composition-thickening}
Si $\mathcal{X} \subset \mathcal{X}'$ et $\mathcal{X}' \subset \mathcal{X}''$
sont des épaississements de champs algébriques, alors
$\mathcal{X} \subset \mathcal{X}''$ en est également un.
\end{lemma}

\begin{proof}
Voir la discussion ci-dessus, Propriétés des champs, section
\ref{stacks-properties-section-properties-morphisms}, et
et Compléments sur les morphismes d'espaces, lemme
\ref{spaces-more-morphisms-lemma-composition-thickening}
\end{proof}

\begin{example}
\label{stacks-more-morphisms-example-reduction-thickening}
Soit $\mathcal{X}'$ un champ algébrique. Alors $\mathcal{X}'$ est un épaississement
de son réduit $\mathcal{X}'_{red}$, voir
Propriétés des champs, définition
\ref{stacks-properties-definition-reduced-induced-stack}.
De plus, si $\mathcal{X} \subset \mathcal{X}'$ est un épaississement
de champs algébriques, alors
$\mathcal{X}'_{red} = \mathcal{X}_{red} \subset \mathcal{X}$.
Autrement dit, $\mathcal{X} = \mathcal{X}'_{red}$ si et seulement
si $\mathcal{X}$ est un champ algébrique réduit.
\end{example}

\begin{lemma}
\label{stacks-more-morphisms-lemma-reduced-diagonal}
Soit $(f, f') : (\mathcal{X} \subset \mathcal{X}') \to
(\mathcal{Y} \subset \mathcal{Y}')$ un morphisme d'épaississements
de champs algébriques. Alors
$\mathcal{X} \times_\mathcal{Y} \mathcal{X} \to
\mathcal{X}' \times_{\mathcal{Y}'} \mathcal{X}'$
est un épaississement, et le diagramme canonique
$$
\xymatrix{
\mathcal{X} \ar[r]_-\Delta \ar[d] &
\mathcal{X} \times_\mathcal{Y} \mathcal{X} \ar[d] \\
\mathcal{X}' \ar[r]^-{\Delta'} &
\mathcal{X}' \times_{\mathcal{Y}'} \mathcal{X}'
}
$$
est cartésien.
\end{lemma}

\begin{proof}
Puisque $\mathcal{X} \to \mathcal{Y}'$ se factorise par le
sous-champ fermé $\mathcal{Y}$, on a
$\mathcal{X} \times_\mathcal{Y} \mathcal{X} =
\mathcal{X} \times_{\mathcal{Y}'} \mathcal{X}$.
Par conséquent,
$\mathcal{X} \times_\mathcal{Y} \mathcal{X} \to
\mathcal{X}' \times_{\mathcal{Y}'} \mathcal{X}'$
est isomorphe au composé
$$
\mathcal{X} \times_{\mathcal{Y}'} \mathcal{X} \to
\mathcal{X} \times_{\mathcal{Y}'} \mathcal{X}' \to
\mathcal{X}' \times_{\mathcal{Y}'} \mathcal{X}'
$$
dont les deux flèches sont des épaississements, puisqu'elles sont obtenues par changement de base
d'épaississements (lemme \ref{stacks-more-morphisms-lemma-base-change-thickening}).
Le composé est donc lui aussi un épaississement
(lemme \ref{stacks-more-morphisms-lemma-composition-thickening}).
Puisque $\mathcal{X} \to \mathcal{X}'$ est un monomorphisme,
la dernière assertion du lemme résulte de
Propriétés des champs, lemme
\ref{stacks-properties-lemma-monomorphism-diagonal}
appliqué à $\mathcal{X} \to \mathcal{X}' \to \mathcal{Y}'$.
\end{proof}

\begin{lemma}
\label{stacks-more-morphisms-lemma-thickening-diagonals}
Soit $(f, f') : (\mathcal{X} \subset \mathcal{X}') \to
(\mathcal{Y} \subset \mathcal{Y}')$ un morphisme d'épaississements
de champs algébriques.
Soient $\Delta : \mathcal{X} \to \mathcal{X} \times_\mathcal{Y} \mathcal{X}$ et
$\Delta' : \mathcal{X}' \to \mathcal{X}' \times_{\mathcal{Y}'} \mathcal{X}'$
les morphismes diagonaux correspondants.
Alors chacune des propriétés de la liste suivante est satisfaite par $\Delta$ si
et seulement si elle est satisfaite par $\Delta'$ :
(a) représentable par des schémas, (b) affine, (c) surjectif, (d) quasi-compact,
(e) universellement fermé, (f) entier, (g) quasi-séparé, (h) séparé,
(i) universellement injectif, (j) universellement ouvert, (k) localement quasi-fini,
(l) fini, (m) non ramifié, (n) monomorphisme, (o) immersion,
(p) immersion fermée et (q) propre.
\end{lemma}

\begin{proof}
Observons que
$$
(\Delta, \Delta') :
(\mathcal{X} \subset \mathcal{X}')
\longrightarrow
(\mathcal{X} \times_\mathcal{Y} \mathcal{X} \subset
\mathcal{X}' \times_{\mathcal{Y}'} \mathcal{X}')
$$
est un morphisme d'épaississements (lemme \ref{stacks-more-morphisms-lemma-reduced-diagonal}).
De plus, $\Delta$ et $\Delta'$ sont
représentables par des espaces algébriques d'après
Morphismes de champs, lemme \ref{stacks-morphisms-lemma-properties-diagonal}.
Par conséquent, via la discussion de
Propriétés des champs, section
\ref{stacks-properties-section-properties-morphisms}
le lemme résulte, dans les cas (a), (b), (c), (d),
(e), (f), (g), (h), (i) et (j), du
Lemme de Compléments sur les morphismes d'espaces
\ref{spaces-more-morphisms-lemma-thicken-property-morphisms}.

\medskip\noindent
Le lemme \ref{stacks-more-morphisms-lemma-reduced-diagonal} nous dit que
$\mathcal{X} = (\mathcal{X} \times_\mathcal{Y} \mathcal{X})
\times_{(\mathcal{X}' \times_{\mathcal{Y}'} \mathcal{X}')} \mathcal{X}'$.
De plus, $\Delta$ et $\Delta'$ sont localement de type fini d'après
le lemme de Morphismes de champs déjà cité,
Morphismes de champs, lemme \ref{stacks-morphisms-lemma-properties-diagonal}.
Le résultat pour les cas (k), (l), (m), (n), (o), (p) et (q) découle donc du
Lemme de Compléments sur les morphismes d'espaces
\ref{spaces-more-morphisms-lemma-properties-that-extend-over-thickenings}.
\end{proof}

\noindent
On obtient comme conséquence le résultat agréable suivant.

\begin{lemma}
\label{stacks-more-morphisms-lemma-thickening-properties}
\begin{reference}
\cite[Theorem 2.2.5]{Conrad-moduli}
\end{reference}
Soit $\mathcal{X} \subset \mathcal{X}'$ un épaississement de
champs algébriques. Alors
\begin{enumerate}
\item $\mathcal{X}$ est un espace algébrique si et seulement si $\mathcal{X}'$
est un espace algébrique,
\item $\mathcal{X}$ est un schéma si et seulement si $\mathcal{X}'$ est un schéma,
\item $\mathcal{X}$ est DM si et seulement si $\mathcal{X}'$ est DM,
\item $\mathcal{X}$ est quasi-DM si et seulement si $\mathcal{X}'$ est quasi-DM,
\item $\mathcal{X}$ est séparé si et seulement si $\mathcal{X}'$ est séparé,
\item $\mathcal{X}$ est quasi-séparé si et seulement si $\mathcal{X}'$ est
quasi-séparé, et
\item en ajouter ici.
\end{enumerate}
\end{lemma}

\begin{proof}
Dans chaque cas, on se ramène à une question portant sur la diagonale, puis
on applique le lemme \ref{stacks-more-morphisms-lemma-thickening-diagonals} au
morphisme d'épaississements
$$
(\mathcal{X} \subset \mathcal{X}') \to
\left(\Spec(\mathbf{Z}) \subset \Spec(\mathbf{Z})\right)
$$
On procède ainsi après avoir considéré
$\mathcal{X} \subset \mathcal{X}'$ comme un épaississement de champs algébriques
sur $\Spec(\mathbf{Z})$ au moyen de
Champs algébriques, définition \ref{algebraic-definition-viewed-as}.

\medskip\noindent
Cas (1). Un champ algébrique est un espace algébrique si et seulement si sa
diagonale est un monomorphisme, voir
Morphismes de champs, lemme \ref{stacks-morphisms-lemma-hierarchy}
(cela résulte aussi immédiatement de Champs algébriques,
proposition \ref{algebraic-proposition-algebraic-stack-no-automorphisms}).

\medskip\noindent
Cas (2). D'après (1), on peut supposer que $\mathcal{X}$ et $\mathcal{X}'$
sont des espaces algébriques, puis appliquer
Compléments sur les morphismes d'espaces, lemme
\ref{spaces-more-morphisms-lemma-thickening-scheme}.

\medskip\noindent
Cas (3) -- (6). Chacun de ces cas correspond à une condition
sur la diagonale, voir Morphismes de champs, définitions
\ref{stacks-morphisms-definition-separated} et
\ref{stacks-morphisms-definition-absolute-separated}.
\end{proof}









\section{Morphismes d'épaississements}
\label{stacks-more-morphisms-section-morphisms-thickenings}

\noindent
Si $(f, f') : (\mathcal{X} \subset \mathcal{X}') \to
(\mathcal{Y} \subset \mathcal{Y}')$ est un morphisme
d'épaississements de champs algébriques, alors les propriétés du morphisme
$f$ se transmettent souvent à $f'$. Il existe plusieurs variantes.

\begin{lemma}
\label{stacks-more-morphisms-lemma-thicken-property-morphisms}
Soit $(f, f') : (\mathcal{X} \subset \mathcal{X}') \to
(\mathcal{Y} \subset \mathcal{Y}')$
un morphisme d'épaississements de champs algébriques. Alors
\begin{enumerate}
\item $f$ est un morphisme affine si et seulement si $f'$ est un morphisme affine,
\item $f$ est un morphisme surjectif si et seulement si $f'$ est un morphisme surjectif,
\item $f$ est quasi-compact si et seulement si $f'$ est quasi-compact,
\item $f$ est universellement fermé si et seulement si $f'$ est universellement fermé,
\item $f$ est entier si et seulement si $f'$ est entier,
\item $f$ est universellement injectif si et seulement si $f'$ est universellement injectif,
\item $f$ est universellement ouvert si et seulement si $f'$ est universellement ouvert,
\item $f$ est quasi-DM si et seulement si $f'$ est quasi-DM,
\item $f$ est DM si et seulement si $f'$ est DM,
\item $f$ est (quasi-)séparé si et seulement si $f'$ est (quasi-)séparé,
\item $f$ est représentable si et seulement si $f'$ est représentable,
\item $f$ est représentable par des espaces algébriques si et seulement si $f'$ est
représentable par des espaces algébriques,
\item en ajouter ici.
\end{enumerate}
\end{lemma}

\begin{proof}
D'après le lemme \ref{stacks-more-morphisms-lemma-thickening}, les morphismes $\mathcal{X} \to \mathcal{X}'$
et $\mathcal{Y} \to \mathcal{Y}'$ sont des homéomorphismes universels. Ainsi, toute
condition portant sur $|f| : |\mathcal{X}| \to |\mathcal{Y}|$ équivaut à
la condition correspondante sur $|f'| : |\mathcal{X}'| \to |\mathcal{Y}'|$,
et il en va de même après un changement de base arbitraire par un morphisme
$\mathcal{Z}' \to \mathcal{Y}'$. Cela démontre les
énoncés (2), (3), (4), (6) et (7).

\medskip\noindent
Dans les cas (8), (9), (10) et (12), on peut traduire les conditions sur
$f$ et $f'$ en conditions sur les diagonales $\Delta$ et $\Delta'$
comme dans le lemme \ref{stacks-more-morphisms-lemma-thickening-diagonals}. Voir
Morphismes de champs, définition \ref{stacks-morphisms-definition-separated} et
lemme \ref{stacks-morphisms-lemma-hierarchy}.
Ces cas résultent donc du lemme \ref{stacks-more-morphisms-lemma-thickening-diagonals}.

\medskip\noindent
Démonstration de (11). Si $f'$ est représentable, alors $f$ l'est aussi, car
pour un schéma $T$ et un morphisme $T \to \mathcal{Y}$, on a
$\mathcal{X} \times_\mathcal{Y} T =
\mathcal{X} \times_{\mathcal{X}'} (\mathcal{X}' \times_{\mathcal{Y}'} T)$
et $\mathcal{X} \to \mathcal{X}'$ est une immersion fermée (donc représentable).
Réciproquement, supposons $f$ représentable et soit $T' \to \mathcal{Y}'$
un morphisme, où $T'$ est un schéma. Alors
$$
\mathcal{X} \times_{\mathcal{Y}}
(\mathcal{Y} \times_{\mathcal{Y}'} T') =
\mathcal{X} \times_{\mathcal{X}'}
(\mathcal{X}' \times_{\mathcal{Y}'} T') \to
\mathcal{X}' \times_{\mathcal{Y}'} T'
$$
est un épaississement (d'après le lemme \ref{stacks-more-morphisms-lemma-base-change-thickening})
et la source est un schéma. La cible est donc un schéma d'après le
lemme \ref{stacks-more-morphisms-lemma-thickening-properties}.

\medskip\noindent
Dans les cas (1) et (5), si $f$ ou $f'$ possède la propriété indiquée,
alors $f$ et $f'$ sont tous deux représentables d'après (11). Dans ce cas,
choisissons un espace algébrique $V'$ et un morphisme lisse surjectif
$V' \to \mathcal{Y}'$. Posons $V = \mathcal{Y} \times_{\mathcal{Y}'} V'$,
$U' = \mathcal{X}' \times_{\mathcal{Y}'} V'$ et
$U = \mathcal{X} \times_{\mathcal{Y}'} V'$. Les résultats cherchés
découlent alors des résultats correspondants pour
le morphisme d'épaississements $(U \subset U') \to (V \subset V')$
d'espaces algébriques, au moyen du principe de Propriétés des champs, lemme
\ref{stacks-properties-lemma-check-property-covering}.
Voir Compléments sur les morphismes d'espaces, lemme
\ref{spaces-more-morphisms-lemma-thicken-property-morphisms}
pour les résultats correspondants dans le cas des espaces algébriques.
\end{proof}





\section{Déformations infinitésimales des champs algébriques}
\label{stacks-more-morphisms-section-deform}

\noindent
Cette section est l'analogue de
Compléments sur les morphismes d'espaces, section
\ref{spaces-more-morphisms-section-deform}.

\begin{lemma}
\label{stacks-more-morphisms-lemma-flatness-morphism-thickenings-fp-over-ft}
Considérons un diagramme commutatif
$$
\xymatrix{
(\mathcal{X} \subset \mathcal{X}') \ar[rr]_{(f, f')} \ar[rd] & &
(\mathcal{Y} \subset \mathcal{Y}') \ar[ld] \\
& (\mathcal{B} \subset \mathcal{B}')
}
$$
d'épaississements de champs algébriques. Supposons que
\begin{enumerate}
\item $\mathcal{Y}' \to \mathcal{B}'$ soit localement de type fini,
\item $\mathcal{X}' \to \mathcal{B}'$ soit
plat et localement de présentation finie,
\item $f$ soit plat, et
\item $\mathcal{X} = \mathcal{B} \times_{\mathcal{B}'} \mathcal{X}'$ et
$\mathcal{Y} = \mathcal{B} \times_{\mathcal{B}'} \mathcal{Y}'$.
\end{enumerate}
Alors $f'$ est plat et, pour tout $y' \in |\mathcal{Y}'|$ appartenant à l'image de $|f'|$,
le morphisme $\mathcal{Y}' \to \mathcal{B}'$ est plat en $y'$.
\end{lemma}

\begin{proof}
Choisissons un espace algébrique $U'$ et un morphisme lisse surjectif
$U' \to \mathcal{B}'$.
Choisissons un espace algébrique $V'$ et un morphisme lisse surjectif
$V' \to U' \times_{\mathcal{B}'} \mathcal{Y}'$.
Choisissons un espace algébrique $W'$ et un morphisme lisse surjectif
$W' \to V' \times_{\mathcal{Y}'} \mathcal{X}'$. Soient $U, V, W$
les changements de base de $U', V', W'$ par $\mathcal{B} \to \mathcal{B}'$.
Alors la platitude de $f'$ équivaut à celle de $W' \to V'$, et
on sait que $W \to V$ est plat. On peut donc appliquer le lemme
dans le cas des espaces algébriques au diagramme
$$
\xymatrix{
(W \subset W') \ar[rr] \ar[rd] & & (V \subset V') \ar[ld] \\
& (U \subset U')
}
$$
d'épaississements d'espaces algébriques. Voir
Compléments sur les morphismes d'espaces, lemme
\ref{spaces-more-morphisms-lemma-flatness-morphism-thickenings-fp-over-ft}.
L'assertion concernant la platitude de $\mathcal{Y}'/\mathcal{B}'$ aux points de
l'image de $|f'|$ s'en déduit de la même manière.
\end{proof}

\begin{lemma}
\label{stacks-more-morphisms-lemma-deform-property-fp-over-ft}
Considérons un diagramme commutatif
$$
\xymatrix{
(\mathcal{X} \subset \mathcal{X}') \ar[rr]_{(f, f')} \ar[rd] & &
(\mathcal{Y} \subset \mathcal{Y}') \ar[ld] \\
& (\mathcal{B} \subset \mathcal{B}')
}
$$
d'épaississements de champs algébriques.
Supposons que $\mathcal{Y}' \to \mathcal{B}'$ soit localement de type fini,
que $\mathcal{X}' \to \mathcal{B}'$ soit plat et localement de présentation finie,
que $\mathcal{X} = \mathcal{B} \times_{\mathcal{B}'} \mathcal{X}'$ et
que $\mathcal{Y} = \mathcal{B} \times_{\mathcal{B}'} \mathcal{Y}'$. Alors
\begin{enumerate}
\item $f$ est plat si et seulement si $f'$ est plat,
\label{stacks-more-morphisms-item-flat-fp-over-ft}
\item $f$ est un isomorphisme si et seulement si $f'$ est un isomorphisme,
\label{stacks-more-morphisms-item-isomorphism-fp-over-ft}
\item $f$ est une immersion ouverte si et seulement si $f'$ est une immersion ouverte,
\label{stacks-more-morphisms-item-open-immersion-fp-over-ft}
\item $f$ est un monomorphisme si et seulement si $f'$ est un monomorphisme,
\label{stacks-more-morphisms-item-monomorphism-fp-over-ft}
\item $f$ est localement quasi-fini si et seulement si $f'$ est localement quasi-fini,
\label{stacks-more-morphisms-item-quasi-finite-fp-over-ft}
\item $f$ est syntomique si et seulement si $f'$ est syntomique,
\label{stacks-more-morphisms-item-syntomic-fp-over-ft}
\item $f$ est lisse si et seulement si $f'$ est lisse,
\label{stacks-more-morphisms-item-smooth-fp-over-ft}
\item $f$ est non ramifié si et seulement si $f'$ est non ramifié,
\label{stacks-more-morphisms-item-unramified-fp-over-ft}
\item $f$ est \'etale si et seulement si $f'$ est \'etale,
\label{stacks-more-morphisms-item-etale-fp-over-ft}
\item $f$ est fini si et seulement si $f'$ est fini, et
\label{stacks-more-morphisms-item-finite-fp-over-ft}
\item en ajouter ici.
\end{enumerate}
\end{lemma}

\begin{proof}
Dans le cas (\ref{stacks-more-morphisms-item-flat-fp-over-ft}), cela résulte du
lemme \ref{stacks-more-morphisms-lemma-flatness-morphism-thickenings-fp-over-ft}.

\medskip\noindent
Dans les cas
(\ref{stacks-more-morphisms-item-syntomic-fp-over-ft}), (\ref{stacks-more-morphisms-item-smooth-fp-over-ft})
cela se démontre par la méthode employée dans la démonstration du
lemme \ref{stacks-more-morphisms-lemma-flatness-morphism-thickenings-fp-over-ft}.
Plus précisément, choisissons un espace algébrique $U'$ et un morphisme lisse surjectif
$U' \to \mathcal{B}'$.
Choisissons un espace algébrique $V'$ et un morphisme lisse surjectif
$V' \to U' \times_{\mathcal{B}'} \mathcal{Y}'$.
Choisissons un espace algébrique $W'$ et un morphisme lisse surjectif
$W' \to V' \times_{\mathcal{Y}'} \mathcal{X}'$. Soient $U, V, W$
les changements de base de $U', V', W'$ par $\mathcal{B} \to \mathcal{B}'$.
Alors la propriété pour $f$, resp.\ $f'$,
équivaut à la propriété pour $W \to V$, resp.\ $W' \to V'$.
On peut donc appliquer le lemme, dans le cas des espaces algébriques, au
diagramme
$$
\xymatrix{
(W \subset W') \ar[rr] \ar[rd] & & (V \subset V') \ar[ld] \\
& (U \subset U')
}
$$
d'épaississements d'espaces algébriques. Voir
Compléments sur les morphismes d'espaces, lemme
\ref{spaces-more-morphisms-lemma-deform-property-fp-over-ft}.

\medskip\noindent
Dans les cas (\ref{stacks-more-morphisms-item-unramified-fp-over-ft}) et (\ref{stacks-more-morphisms-item-etale-fp-over-ft}),
on constate d'abord que l'hypothèse portant sur $f$ ou $f'$ implique que
$f$ et $f'$ sont tous deux des morphismes DM de champs algébriques, voir
le lemme \ref{stacks-more-morphisms-lemma-thicken-property-morphisms}. On peut alors choisir
un espace algébrique $U'$ et un morphisme lisse surjectif
$U' \to \mathcal{B}'$.
Choisissons un espace algébrique $V'$ et un morphisme lisse surjectif
$V' \to U' \times_{\mathcal{B}'} \mathcal{Y}'$.
Choisissons un espace algébrique $W'$ et un morphisme \'etale surjectif (!)
$W' \to V' \times_{\mathcal{Y}'} \mathcal{X}'$. Soient $U, V, W$
les changements de base de $U', V', W'$ par $\mathcal{B} \to \mathcal{B}'$.
Alors $W \to V \times_\mathcal{Y} \mathcal{X}$ est lui aussi
\'etale et surjectif. La propriété pour $f$, resp.\ $f'$,
équivaut donc à la propriété pour $W \to V$, resp.\ $W' \to V'$.
On peut donc appliquer le lemme, dans le cas des espaces algébriques, au
diagramme
$$
\xymatrix{
(W \subset W') \ar[rr] \ar[rd] & & (V \subset V') \ar[ld] \\
& (U \subset U')
}
$$
d'épaississements d'espaces algébriques. Voir
Compléments sur les morphismes d'espaces, lemme
\ref{spaces-more-morphisms-lemma-deform-property-fp-over-ft}.

\medskip\noindent
Dans les cas (\ref{stacks-more-morphisms-item-isomorphism-fp-over-ft}),
(\ref{stacks-more-morphisms-item-open-immersion-fp-over-ft}),
(\ref{stacks-more-morphisms-item-monomorphism-fp-over-ft}),
(\ref{stacks-more-morphisms-item-finite-fp-over-ft})
on déduit d'abord du lemme \ref{stacks-more-morphisms-lemma-thicken-property-morphisms}
que $f$ et $f'$ sont représentables par des espaces algébriques. On peut donc choisir
un espace algébrique $U'$ et un morphisme lisse surjectif
$U' \to \mathcal{B}'$,
un espace algébrique $V'$ et un morphisme lisse surjectif
$V' \to U' \times_{\mathcal{B}'} \mathcal{Y}'$, puis
$W' = V' \times_{\mathcal{Y}'} \mathcal{X}'$ sera un espace algébrique.
Soient $U, V, W$ les changements de base de
$U', V', W'$ par $\mathcal{B} \to \mathcal{B}'$.
On a également $W = V \times_\mathcal{Y} \mathcal{X}$.
Il faut alors montrer que $W' \to V'$ est un
isomorphisme, resp.\ une immersion ouverte, resp.\ un monomorphisme,
resp.\ un morphisme fini, si et seulement si $W \to V$ possède la même propriété.
Voir Propriétés des champs, lemme
\ref{stacks-properties-lemma-check-property-covering}.
On conclut ainsi en appliquant les résultats
précédents relatifs aux espaces algébriques.

\medskip\noindent
Dans le cas (\ref{stacks-more-morphisms-item-quasi-finite-fp-over-ft}), observons d'abord
que $f$ et $f'$ sont localement de type fini d'après
Morphismes de champs, lemme \ref{stacks-morphisms-lemma-finite-type-permanence}.
D'autre part, $f$ est quasi-DM si et seulement si
$f'$ l'est, d'après le
lemme \ref{stacks-more-morphisms-lemma-thicken-property-morphisms}.
La dernière condition à vérifier pour que $f$ ou $f'$ soit localement quasi-fini
(Morphismes de champs, définition
\ref{stacks-morphisms-definition-locally-quasi-finite})
est une condition portant sur les espaces topologiques sous-jacents;
elle est satisfaite par $f$ si et seulement si elle l'est par $f'$ d'après
la discussion du premier paragraphe de la démonstration.
\end{proof}









\section{Relèvement des schémas affines}
\label{stacks-more-morphisms-section-lifting-affines}

\noindent
Considérons un diagramme aux flèches pleines
$$
\xymatrix{
W \ar[d] \ar@{..>}[r] & W' \ar@{..>}[d] \\
\mathcal{X} \ar[r] & \mathcal{X}'
}
$$
où $\mathcal{X} \subset \mathcal{X}'$ est un épaississement de champs algébriques,
$W$ est un schéma affine et $W \to \mathcal{X}$ est lisse. La question
étudiée dans cette section est de savoir si l'on peut trouver $W'$ et les flèches
en pointillé de sorte que le carré soit cartésien et que $W' \to \mathcal{X}'$ soit lisse.
Nous n'en connaissons pas la réponse en général, mais si $\mathcal{X} \subset \mathcal{X}'$
est un épaississement du premier ordre, nous démontrerons que la réponse est affirmative.

\medskip\noindent
Pour étudier ce problème, introduisons la catégorie suivante.

\begin{remark}[Catégorie des relèvements]
\label{stacks-more-morphisms-remark-gerbe-of-lifts}
Considérons un diagramme
$$
\xymatrix{
W \ar[d]_x \\
\mathcal{X} \ar[r] & \mathcal{X}'
}
$$
où $\mathcal{X} \subset \mathcal{X}'$ est un épaississement de champs algébriques,
$W$ est un espace algébrique et $W \to \mathcal{X}$ est lisse.
Nous allons construire une catégorie $\mathcal{C}$ et un foncteur
$$
p : \mathcal{C} \longrightarrow W_{spaces, \etale}
$$
(voir Propriétés des espaces, définition
\ref{spaces-properties-definition-spaces-etale-site} pour la notation)
comme suit. Un objet de $\mathcal{C}$ sera un système
$(U, U', a, i, x', \alpha)$
qui forme un diagramme commutatif
\begin{equation}
\label{stacks-more-morphisms-equation-object}
\vcenter{
\xymatrix{
U \ar[d]_a \ar[r]_i & U' \ar[dd]^{x'} \\
W \ar[d]_x & \\
\mathcal{X} \ar[r] & \mathcal{X}'
}
}
\end{equation}
dont la commutativité est matérialisée par le $2$-morphisme
$\alpha : x \circ a \to x' \circ i$, et tel que
$U$ et $U'$ soient des espaces algébriques,
$a : U \to W$ soit \'etale, $x' : U' \to \mathcal{X}'$ soit lisse
et $U = \mathcal{X} \times_{\mathcal{X}'} U'$.
En particulier, $U \subset U'$ est un épaississement.
Un morphisme
$$
(V, V', b, j, y', \beta) \to (U, U', a, i, x', \alpha)
$$
est la donnée de $(f, f', \gamma)$, où $f : V \to U$ est un morphisme
au-dessus de $W$, $f' : V' \to U'$ est un morphisme dont la restriction
à $V$ donne $f$, et $\gamma : x' \circ f' \to y'$ est un $2$-morphisme
qui matérialise la commutativité dans le triangle de droite du diagramme ci-dessous
\begin{equation}
\label{stacks-more-morphisms-equation-morphism}
\vcenter{
\xymatrix{
& V \ar[ld]_f \ar[ldd]^b \ar[rr]_j & & V' \ar[ld]_{f'} \ar[lddd]^{y'} \\
U \ar[d]_a \ar[rr]_i & & U' \ar[dd]_{x'} \\
W \ar[d]_x & \\
\mathcal{X} \ar[rr] & & \mathcal{X}'
}
}
\end{equation}
Enfin, on exige que $\gamma$ soit compatible avec $\alpha$ et $\beta$ :
dans le calcul des $2$-catégories de Catégories, sections
\ref{categories-section-formal-cat-cat} et
\ref{categories-section-2-categories}, cela s'écrit
$$
\beta = (\gamma \star \text{id}_j) \circ (\alpha \star \text{id}_f)
$$
(plus brièvement : $\beta = j^*\gamma \circ f^*\alpha$).
Une autre formulation consiste à dire que les objets sont les diagrammes commutatifs
(\ref{stacks-more-morphisms-equation-object}) satisfaisant certaines propriétés supplémentaires, et que les
morphismes sont les diagrammes commutatifs
(\ref{stacks-more-morphisms-equation-morphism}) dans la catégorie $\textit{Espaces}/\mathcal{X}'$
introduite dans Propriétés des champs, remarque
\ref{stacks-properties-remark-representable-over}.
Il est alors clair que $\mathcal{C}$ est une catégorie
et que la règle $p : \mathcal{C} \to W_{spaces, \etale}$ qui
envoie $(U, U', a, i, x', \alpha)$ sur $a : U \to W$
est un foncteur.
\end{remark}

\begin{lemma}
\label{stacks-more-morphisms-lemma-morphisms-lifts-etale}
Pour tout morphisme (\ref{stacks-more-morphisms-equation-morphism}), l'application $f' : V' \to U'$ est \'etale.
\end{lemma}

\begin{proof}
En effet, $f : V \to U$ est \'etale en tant que morphisme de $W_{spaces, \etale}$,
et l'on peut appliquer le
lemme \ref{stacks-more-morphisms-lemma-deform-property-fp-over-ft}, puisque $U' \to \mathcal{X}'$
et $V' \to \mathcal{X}'$ sont lisses et que
$U = \mathcal{X} \times_{\mathcal{X}'} U'$ et
$V = \mathcal{X} \times_{\mathcal{X}'} V'$.
\end{proof}

\begin{lemma}
\label{stacks-more-morphisms-lemma-gerbe-of-lifts-fibred}
La catégorie $p : \mathcal{C} \to W_{spaces, \etale}$ construite
dans la remarque \ref{stacks-more-morphisms-remark-gerbe-of-lifts} est fibrée en groupoïdes.
\end{lemma}

\begin{proof}
Nous affirmons que les catégories fibres de $p$ sont des groupoïdes.
Si $(f, f', \gamma')$, comme dans (\ref{stacks-more-morphisms-equation-morphism}),
est un morphisme tel que $f : V \to U$ soit un isomorphisme, alors
$f'$ est un isomorphisme d'après le lemme \ref{stacks-more-morphisms-lemma-deform-property-fp-over-ft},
et $(f, f', \gamma')$ est donc un isomorphisme.

\medskip\noindent
Considérons un morphisme $f : V \to U$ dans $W_{spaces, \etale}$
et un objet $\xi = (U, U', a, i, x', \alpha)$
de $\mathcal{C}$ au-dessus de $U$. Nous allons construire l'« image inverse »
$f^*\xi$ au-dessus de $V$.
Posons $b = a \circ f$. Soit $f' : V' \to U'$ le morphisme \'etale
dont la restriction à $V$ est $f$ (Compléments sur les morphismes d'espaces,
lemme \ref{spaces-more-morphisms-lemma-topological-invariance}).
Notons $j : V \to V'$ l'épaississement correspondant.
Soient $y' = x' \circ f'$ et $\gamma = \text{id} : x' \circ f' \to y'$.
Posons
$$
\beta = \alpha \star \text{id}_f :
x \circ b = x \circ a \circ f \to
x' \circ i \circ f = x' \circ f' \circ j = y' \circ j
$$
Il est clair que $(f, f', \gamma) : (V, V', b, j, y', \beta) \to
(U, U', a, i, x', \alpha)$ est un morphisme comme dans
(\ref{stacks-more-morphisms-equation-morphism}). Les morphismes $(f, f', \gamma)$
ainsi construits sont fortement cartésiens
(Catégories, définition \ref{categories-definition-cartesian-over-C}).
Nous omettons la démonstration détaillée; essentiellement, la raison est que,
étant donné un morphisme
$(g, g', \epsilon) : (Y, Y', c, k, z', \delta) \to (U, U', a, i, x', \alpha)$
dans $\mathcal{C}$ tel que $g$ se factorise sous la forme $g = f \circ h$ pour un certain
$h : Y \to V$, on obtient une unique factorisation $g' = f' \circ h'$
d'après Compléments sur les morphismes d'espaces,
lemme \ref{spaces-more-morphisms-lemma-topological-invariance}
puis on peut produire le $\zeta$ nécessaire tel que
$(h, h', \zeta) :  (Y, Y', c, k, z', \delta) \to
(V, V', b, j, y', \beta)$ soit un morphisme de $\mathcal{C}$
vérifiant $(g, g', \epsilon) = (f, f', \gamma) \circ (h, h', \zeta)$.

\medskip\noindent
Par conséquent, $p : \mathcal{C} \to W_\etale$ est une catégorie
fibrée (Catégories, définition \ref{categories-definition-fibred-category}).
En combinant cela au fait établi ci-dessus que les catégories fibres sont des groupoïdes,
on conclut que $p : \mathcal{C} \to W_\etale$ est
fibrée en groupoïdes d'après Catégories, lemme
\ref{categories-lemma-fibred-groupoids}.
\end{proof}

\begin{lemma}
\label{stacks-more-morphisms-lemma-gerbe-of-lifts-stack}
La catégorie $p : \mathcal{C} \to W_{spaces, \etale}$ construite
dans la remarque \ref{stacks-more-morphisms-remark-gerbe-of-lifts} est un champ en groupoïdes.
\end{lemma}

\begin{proof}
D'après le lemme \ref{stacks-more-morphisms-lemma-gerbe-of-lifts-fibred}, la première condition
de Champs, définition \ref{stacks-definition-stack-in-groupoids}, est satisfaite.
Comme il est d'usage, nous vérifions la descente des objets et laissons au lecteur
le soin de vérifier celle des morphismes. Supposons donc donnés $a : U \to W$
dans $W_{spaces, \etale}$, un recouvrement $\{U_k \to U\}_{k \in K}$ dans
$W_{spaces, \etale}$, des objets $\xi_k = (U_k, U'_k, a_k, i_k, x'_k, \alpha_k)$
de $\mathcal{C}$ au-dessus de $U_k$, et des morphismes
$$
\varphi_{kk'} = (f_{kk'}, f'_{kk'}, \gamma_{kk'}) :
\xi_k|_{U_k \times_U U_{k'}} \to
\xi_{k'}|_{U_k \times_U U_{k'}}
$$
entre les restrictions, satisfaisant à la condition de cocycle. Pour démontrer
l'effectivité, on peut d'abord raffiner le recouvrement. On peut donc supposer que chaque
$U_k$ est un schéma (même un schéma affine si l'on veut). Écrivons
$$
\xi_k|_{U_k \times_U U_{k'}} =
(U_k \times_U U_{k'}, U'_{kk'}, a_{kk'}, x'_{kk'}, \alpha_{kk'})
$$
On obtient alors un morphisme \'etale (d'après le lemme \ref{stacks-more-morphisms-lemma-morphisms-lifts-etale})
$s_{kk'} : U'_{kk'} \to U'_k$
comme seconde composante du morphisme
$\xi_k|_{U_k \times_U U_{k'}} \to \xi_k$ de $\mathcal{C}$.
De même, on obtient un morphisme \'etale $t_{kk'} : U'_{kk'} \to U'_{k'}$
en considérant la seconde composante du composé
$$
\xi_k|_{U_k \times_U U_{k'}} \xrightarrow{\varphi_{kk'}}
\xi_{k'}|_{U_k \times_U U_{k'}} \to \xi_{k'}
$$
Nous affirmons que
$$
j :
\coprod\nolimits_{(k, k') \in K \times K} U'_{kk'}
\xrightarrow{(\coprod s_{kk'}, \coprod t_{kk'})}
(\coprod\nolimits_{k \in K} U'_k) \times (\coprod\nolimits_{k \in K} U'_k)
$$
est une relation d'équivalence \'etale. Premièrement, nous avons déjà vu
que les composantes $s, t$ du morphisme affiché sont \'etales.
Le changement de base du morphisme $j$ par
$(\coprod U_k) \times (\coprod U_k) \to (\coprod U'_k) \times (\coprod U'_k)$
est un monomorphisme, car c'est l'application
$$
\coprod\nolimits_{(k, k') \in K \times K} U_k \times_U U_{k'}
\longrightarrow
(\coprod\nolimits_{k \in K} U_k) \times (\coprod\nolimits_{k \in K} U_k)
$$
Par conséquent, $j$ est un monomorphisme d'après Compléments sur les morphismes, lemme
\ref{more-morphisms-lemma-properties-that-extend-over-thickenings}.
Enfin, la symétrie de la relation $j$ provient du fait que
$\varphi_{kk'}^{-1}$ est le « renversé » de $\varphi_{k'k}$ (voir
Champs, remarques \ref{stacks-remarks-definition-descent-datum}),
et la transitivité résulte de la condition de cocycle (détails omis).
Le quotient de $\coprod U'_k$ par $j$ est donc un espace algébrique $U'$
(Espaces, théorème \ref{spaces-theorem-presentation}).
Nous avons déjà montré ci-dessus qu'il existe un épaississement
$i : U \to U'$, puisque la restriction de $j$ à
$\coprod U_k$ donne $(\coprod U_k) \times_U (\coprod U_k)$.
Enfin, si l'on considère temporairement les $1$-morphismes
$x'_k : U'_k \to \mathcal{X}'$ comme des objets du champ
$\mathcal{X}'$ au-dessus de $U'_k$, on voit qu'ils sont munis d'une
donnée de descente relativement au recouvrement \'etale
$\{U'_k \to U'\}$, donnée par la troisième composante $\gamma_{kk'}$
des morphismes $\varphi_{kk'}$ dans $\mathcal{C}$.
Puisque $\mathcal{X}'$ est un champ,
cette donnée de descente est effective; en revenant à la formulation initiale, on obtient
un $1$-morphisme $x' : U' \to \mathcal{X}'$ tel que les composés
$U'_k \to U' \to \mathcal{X}'$ soient munis d'isomorphismes vers $x'_k$
compatibles avec $\gamma_{kk'}$. Cela signifie que les morphismes
$\alpha_k : x \circ a_k \to x'_k \circ i_k$ se recollent en un morphisme
$\alpha : x \circ a \to x' \circ i$. Alors $\xi = (U, U', a, i, x', \alpha)$
est l'objet cherché au-dessus de $U$.
\end{proof}

\begin{lemma}
\label{stacks-more-morphisms-lemma-etale-local-lifts}
Soit $\mathcal{X} \subset \mathcal{X}'$ un épaississement de champs algébriques.
Soit $W$ un espace algébrique et soit $W \to \mathcal{X}$ un morphisme lisse.
Il existe un recouvrement \'etale $\{W_i \to W\}_{i \in I}$ et, pour tout $i$,
un diagramme cartésien
$$
\xymatrix{
W_i \ar[r] \ar[d] & W_i' \ar[d] \\
\mathcal{X} \ar[r] & \mathcal{X}'
}
$$
où $W_i' \to \mathcal{X}'$ est lisse.
\end{lemma}

\begin{proof}
Choisissons un schéma $U'$ et un morphisme lisse surjectif $U' \to \mathcal{X}'$.
Comme d'habitude, posons $U = \mathcal{X} \times_{\mathcal{X}'} U'$. Alors
$U \to \mathcal{X}$ est un morphisme lisse surjectif. Par conséquent, le changement de base
$$
V = W \times_{\mathcal{X}} U \longrightarrow W
$$
est un morphisme lisse surjectif d'espaces algébriques.
D'après Topologies sur les espaces, lemme
\ref{spaces-topologies-lemma-etale-dominates-smooth}
on peut trouver un recouvrement \'etale $\{W_i \to W\}$ tel
que $W_i \to W$ se factorise par $V \to W$.
Après avoir recouvert $W_i$ par des ouverts affines (Propriétés des espaces, lemme
\ref{spaces-properties-lemma-cover-by-union-affines})
on peut supposer chaque $W_i$ affine. On peut, et l'on va, remplacer $W$ par $W_i$,
ce qui ramène à la situation étudiée au paragraphe suivant.

\medskip\noindent
Supposons $W$ affine et le morphisme donné $W \to \mathcal{X}$ factorisé
par $U$. On a le diagramme
$$
W \xrightarrow{i} U \to \mathcal{X}
$$
Puisque $W$ et $U$ sont lisses sur $\mathcal{X}$, on voit que
$i$ est localement de type fini (Morphismes de champs, lemme
\ref{stacks-morphisms-lemma-finite-type-permanence}).
Après avoir remplacé $U$ par $\mathbf{A}^n_U$, on peut supposer
que $i$ est une immersion, voir
Morphismes, lemme \ref{morphisms-lemma-quasi-affine-finite-type-over-S}.
D'après Morphismes de champs,
lemme \ref{stacks-morphisms-lemma-lci-permanence}
le morphisme $i$ est d'intersection complète locale.
Par conséquent, $i$ est une immersion Koszul-régulière (au sens de
Diviseurs, définition \ref{divisors-definition-regular-immersion}),
d'après Compléments sur les morphismes, lemme \ref{more-morphisms-lemma-lci}.

\medskip\noindent
On peut encore remplacer $W$ par un recouvrement ouvert affine.
Pour tout point $w \in W$, on peut choisir un ouvert affine
$U'_w \subset U'$ tel que, si $U_w \subset U$ est
l'ouvert affine correspondant, alors $w \in i^{-1}(U_w)$ et
$i^{-1}(U_w) \to U_w$ est une immersion fermée définie
par une suite Koszul-régulière
$f_1, \ldots, f_r \in \Gamma(U_w, \mathcal{O}_{U_w})$.
Cela résulte de la définition des immersions Koszul-régulières
et de Diviseurs, lemme \ref{divisors-lemma-regular-ideal-sheaf-scheme}.
Posons $W_w = i^{-1}(U_w)$; c'est un voisinage ouvert affine
de $w \in W$.
Choisissons des relèvements $f'_1, \ldots, f'_r \in \Gamma(U'_w, \mathcal{O}_{U'_w})$
de $f_1, \ldots, f_r$. Cela est possible puisque $U_w \to U'_w$
est une immersion fermée de schémas affines.
Soit $W'_w \subset U'_w$ le sous-schéma fermé défini par
$f'_1, \ldots, f'_r$.
Nous affirmons que $W'_w \to \mathcal{X}'$ est lisse.
Cette affirmation achève la démonstration, car
$W_w = \mathcal{X} \times_{\mathcal{X}'} W'_w$
par construction.

\medskip\noindent
Pour vérifier l'affirmation, il suffit de vérifier que le changement de base
$W'_w \times_{\mathcal{X}'} X' \to X'$ est lisse pour tout
schéma affine $X'$ lisse sur $\mathcal{X}'$. Choisissons un
morphisme \'etale
$$
Y' \to U'_w \times_{\mathcal{X}'} X'
$$
avec $Y'$ affine. Puisque $U'_w \times_{\mathcal{X}'} X'$ est recouvert
par les images de tels morphismes, il suffit de montrer que le sous-schéma fermé
$Z'$ de $Y'$ défini par $f'_1, \ldots, f'_r$ est lisse sur $X'$.
Considérons le diagramme
$$
\xymatrix{
Z' \ar[r] \ar[d] & Y' \ar[d] \\
W'_w \times_{\mathcal{X}'} X' \ar[d] \ar[r] &
U'_w \times_{\mathcal{X}'} X' \ar[d] \ar[r] & X' \\
W'_w = V(f'_1, \ldots, f'_r) \ar[r] & U'_w
}
$$
Posons $X = \mathcal{X} \times_{\mathcal{X}'} X'$,
$Y = X \times_{X'} Y' = \mathcal{X} \times_{\mathcal{X}'} Y'$ et
$Z = Y \times_{Y'} Z' = X \times_{X'} Z' =
\mathcal{X} \times_{\mathcal{X}'} Z'$.
Alors $(Z \subset Z') \to (Y \subset Y') \subset (X \subset X')$
sont des morphismes (cartésiens) d'épaississements de schémas affines, et
on sait que $Z \to X$ et $Y' \to X'$ sont lisses.
Enfin, la suite de fonctions $f'_1, \ldots, f'_r$
a pour image une suite Koszul-régulière dans $\Gamma(Y', \mathcal{O}_{Y'})$ d'après
Compléments d'algèbre, lemme \ref{more-algebra-lemma-koszul-regular-flat-base-change},
car $Y' \to U'_w$ est lisse, donc plat.
D'après Compléments d'algèbre, lemme \ref{more-algebra-lemma-cut-by-koszul}
(et le fait que les suites Koszul-régulières sont des suites quasi-régulières
d'après Compléments d'algèbre, lemmes \ref{more-algebra-lemma-regular-koszul-regular},
\ref{more-algebra-lemma-koszul-regular-H1-regular} et
\ref{more-algebra-lemma-H1-regular-quasi-regular})
on conclut que $Z' \to X'$ est lisse, comme voulu.
\end{proof}

\begin{lemma}
\label{stacks-more-morphisms-lemma-etale-local-lifts-isomorphic}
Soit $\mathcal{X} \subset \mathcal{X}'$ un épaississement de champs algébriques.
Considérons un diagramme commutatif
$$
\xymatrix{
W'' \ar[d]_{x''} & W \ar[l] \ar[r] \ar[d]_x & W' \ar[d]^{x'} \\
\mathcal{X}' & \mathcal{X} \ar[l] \ar[r] & \mathcal{X}'
}
$$
dont les carrés sont cartésiens, où $W', W, W''$ sont des espaces algébriques et
les flèches verticales sont lisses. Alors il existe
\begin{enumerate}
\item un recouvrement \'etale $\{f'_k : W'_k \to W'\}_{k \in K}$,
\item des morphismes \'etales $f''_k : W'_k \to W''$, et
\item des $2$-morphismes $\gamma_k : x' \circ f'_k \to x'' \circ f''_k$
\end{enumerate}
tels que (a) $(f'_k)^{-1}(W) = (f''_k)^{-1}(W)$, (b)
$f'_k|_{(f'_k)^{-1}(W)} = f''_k|_{(f''_k)^{-1}(W)}$ et
(c) l'image inverse de $\gamma_k$ sur le sous-schéma fermé de (a)
coïncide avec le $2$-morphisme donné par la commutativité du
diagramme initial au-dessus de $W$.
\end{lemma}

\begin{proof}
Notons $i : W \to W'$ et $i'' : W \to W''$ les épaississements donnés.
La commutativité du diagramme de l'énoncé du lemme
signifie qu'il existe un $2$-morphisme $\delta : x' \circ i \to x'' \circ i''$.
C'est le $2$-morphisme auquel renvoie la partie (c) de l'énoncé.
Considérons l'espace algébrique
$$
I' = W' \times_{x', \mathcal{X}', x''} W''
$$
muni des projections $p' : I' \to W'$ et $q' : I' \to W''$.
Observons qu'il existe un $2$-morphisme « universel »
$\gamma : x' \circ p' \to x'' \circ q'$ (nous l'utiliserons plus loin).
Le choix de $\delta$ définit un morphisme
$$
\xymatrix{
W \ar[rr]_\delta & & I' \ar[ld]^{p'} \ar[rd]_{q'} \\
& W' & & W''
}
$$
tel que les composés $W \to I' \to W'$ et $W \to I' \to W''$
soient respectivement $i : W \to W'$ et $i'' : W \to W''$.
Puisque $x''$ est lisse, le morphisme $p' : I' \to W'$ est lisse
en tant que changement de base de $x''$.

\medskip\noindent
Supposons que l'on puisse trouver un recouvrement \'etale $\{f'_k : W'_k \to W'\}$
et des morphismes $\delta_k : W'_k \to I'$ tels que la restriction
de $\delta_k$ à $W_k = (f'_k)^{-1}(W)$ soit égale à $\delta \circ f_k$,
où $f_k = f'_k|_{W_k}$. On a le diagramme
$$
\xymatrix{
W_k \ar[r]^{f_k} \ar[d] & W \ar[r]^\delta & I' \ar[d]^{p'} \\
W'_k \ar[rr]^{f'_k} \ar[rru]^{\delta_k} & & W'
}
$$
Autrement dit, on veut pouvoir prolonger la section donnée
$\delta : W \to I'$ de $p'$ en une section au-dessus de $W'$, après avoir éventuellement
remplacé $W'$ par un recouvrement \'etale.

\medskip\noindent
Si cela est possible, on peut poser $f''_k = q' \circ \delta_k$
et $\gamma_k = \gamma \star \text{id}_{\delta_k}$ (plus brièvement,
$\gamma_k = \delta_k^*\gamma$). Il reste seulement à montrer
que le morphisme $f''_k$ est \'etale.
Par construction, le morphisme $x' \circ p'$ est $2$-isomorphe
à $x'' \circ q'$. Par conséquent, $x'' \circ f''_k$ est $2$-isomorphe
à $x' \circ f'_k$. On en déduit que le composé
$$
W'_k \xrightarrow{f''_k} W'' \xrightarrow{x''} \mathcal{X}'
$$
est lisse, puisque $x' \circ f'_k$ l'est.
Puisque $f_k$ est \'etale, on conclut que $f''_k$ est \'etale
d'après le lemme \ref{stacks-more-morphisms-lemma-deform-property-fp-over-ft}.

\medskip\noindent
Si l'épaississement est du premier ordre, on peut
choisir un recouvrement \'etale quelconque $\{W'_k \to W'\}$ avec $W'_k$ affine.
En effet, puisque $p'$ est lisse, on voit que $p'$ est formellement lisse d'après le
critère infinitésimal de relèvement (Compléments sur les morphismes d'espaces, lemme
\ref{spaces-more-morphisms-lemma-smooth-formally-smooth}).
Comme $W_k$ est affine et que $W_k \to W'_k$ est un épaississement du premier ordre
(en tant que changement de base de $\mathcal{X} \to \mathcal{X}'$, voir
le lemme \ref{stacks-more-morphisms-lemma-base-change-thickening}), on obtient le $\delta_k$
cherché.

\medskip\noindent
Dans le cas général, l'existence du recouvrement et des morphismes
$\delta_k$ résulte de Compléments sur les morphismes d'espaces, lemme
\ref{spaces-more-morphisms-lemma-smooth-strong-lift}.
\end{proof}

\begin{lemma}
\label{stacks-more-morphisms-lemma-gerbe-of-lifts}
La catégorie $p : \mathcal{C} \to W_{spaces, \etale}$ construite
dans la remarque \ref{stacks-more-morphisms-remark-gerbe-of-lifts} est une gerbe.
\end{lemma}

\begin{proof}
Dans le lemme \ref{stacks-more-morphisms-lemma-gerbe-of-lifts-stack},
nous avons vu qu'il s'agit d'un champ en groupoïdes.
Il reste donc à vérifier les conditions (2) et (3) de
Champs, définition \ref{stacks-definition-gerbe}.
La condition (2) résulte du
lemme \ref{stacks-more-morphisms-lemma-etale-local-lifts}.
La condition (3) résulte du
lemme \ref{stacks-more-morphisms-lemma-etale-local-lifts-isomorphic}.
\end{proof}

\begin{lemma}
\label{stacks-more-morphisms-lemma-gerbe-of-lifts-first-order}
Dans la remarque \ref{stacks-more-morphisms-remark-gerbe-of-lifts}, supposons que $\mathcal{X} \subset \mathcal{X}'$
soit un épaississement du premier ordre. Alors
\begin{enumerate}
\item les faisceaux d'automorphismes des objets de la gerbe
$p : \mathcal{C} \to W_{spaces, \etale}$ construite
dans la remarque \ref{stacks-more-morphisms-remark-gerbe-of-lifts} sont abéliens, et
\item le faisceau de groupes $\mathcal{G}$ construit dans
Champs, lemme \ref{stacks-lemma-gerbe-abelian-auts},
est un $\mathcal{O}_W$-module quasi-cohérent.
\end{enumerate}
\end{lemma}

\begin{proof}
Nous allons démontrer simultanément les deux assertions. Plus précisément, étant donné
un objet $\xi = (U, U', a, i, x', \alpha)$, nous munirons
$\mathit{Aut}(\xi)$ d'une structure de
$\mathcal{O}_U$-module quasi-cohérent sur $U_{spaces, \etale}$,
et nous montrerons que cette structure est compatible aux images inverses.
Cela suffira par recollement des faisceaux
(Sites, section \ref{sites-section-glueing-sheaves}),
et par la construction de $\mathcal{G}$ dans la démonstration de
Champs, lemme \ref{stacks-lemma-gerbe-abelian-auts},
comme recollement des faisceaux d'automorphismes $\mathit{Aut}(\xi)$,
ainsi que par le fait qu'il suffit, pour vérifier qu'un module est
quasi-cohérent, de passer à un recouvrement \'etale
(Propriétés des espaces, lemme
\ref{spaces-properties-lemma-characterize-quasi-coherent}).

\medskip\noindent
Nous allons décrire le faisceau $\mathit{Aut}(\xi)$ par la
même méthode que dans la démonstration du
lemme \ref{stacks-more-morphisms-lemma-etale-local-lifts-isomorphic}.
Considérons l'espace algébrique
$$
I' = U' \times_{x', \mathcal{X}', x'} U'
$$
muni des projections $p' : I' \to U'$ et $q' : I' \to U'$.
Sur $I'$, il existe un $2$-morphisme universel
$\gamma : x' \circ p' \to x' \circ q'$.
L'identité $x' \to x'$ définit un morphisme diagonal
$$
\xymatrix{
U' \ar[rr]_{\Delta'} & & I' \ar[ld]^{p'} \ar[rd]_{q'} \\
& U' & & U'
}
$$
tel que les composés $U' \to I' \to U'$ et $U' \to I' \to U'$
soient les morphismes identité. Nous noterons les changements de base de
$U', I', p', q', \Delta'$ à $\mathcal{X}$ par $U, I, p, q, \Delta$.
Puisque $U' \to \mathcal{X}'$ est lisse, on voit que $p' : I' \to U'$
est lisse par changement de base.

\medskip\noindent
Une section de $\mathit{Aut}(\xi)$ au-dessus de $U$ est un morphisme $\delta' : U' \to I'$
tel que $\delta'|_U = \Delta$ et
$p' \circ \delta' = \text{id}_{U'}$. Explicitement,
$(\text{id}_U, q' \circ \delta', (\delta')^*\gamma) : \xi \to \xi$
est une formule pour l'automorphisme correspondant.
Plus généralement, si
$f : V \to U$ est un morphisme \'etale, alors il existe un épaississement
$j : V \to V'$ et un morphisme \'etale
$f' : V' \to U'$ dont la restriction à $V$ est $f$, et $f^*\xi$
correspond à $(V, V', a \circ f, j, x' \circ f', f^*\alpha)$, voir la démonstration du
lemme \ref{stacks-more-morphisms-lemma-gerbe-of-lifts-fibred}.
Une section de $\mathit{Aut}(\xi)$ au-dessus de $V$ est un morphisme
$\delta' : V' \to I'$
tel que $\delta'|_V = \Delta \circ f$
et $p' \circ \delta' = f'$\footnote{Une formule pour l'automorphisme
correspondant est $(\text{id}_V, h', (\delta')^*\gamma)$.
Ici, $h' : V' \to V'$ est l'unique (iso)morphisme tel que
$h'|_V = \text{id}_V$ et que le diagramme
$$
\xymatrix{
V' \ar[r]_{h'} \ar[rd]_{q' \circ \delta'} & V' \ar[d]^{f'} \\
& U'
}
$$
commute. L'unicité et l'existence de $h'$ résultent de l'invariance topologique
du site \'etale, voir Compléments sur les morphismes d'espaces,
théorème \ref{spaces-more-morphisms-theorem-topological-invariance}.
Le lecteur pourrait estimer qu'il faudrait plutôt considérer les morphismes
$\delta'' : V' \to V' \times_{\mathcal{X}'} V'$ tels que
$\delta'' \circ j = \Delta_{V'/\mathcal{X}'}$ et
$\text{pr}_1 \circ \delta'' = \text{id}_{V'}$.
Cela conviendrait également : puisque $V' \times_{\mathcal{X}'} V' \to I'$ est \'etale,
la même invariance topologique montre que l'application qui envoie
$\delta''$ sur $\delta' = (V' \times_{\mathcal{X}'} V' \to I') \circ \delta''$
est une bijection entre les deux
ensembles de morphismes.}.

\medskip\noindent
On conclut que le faisceau d'ensembles $\mathit{Aut}(\xi)$ coïncide avec
le faisceau défini dans
Compléments sur les morphismes d'espaces, remarque
\ref{spaces-more-morphisms-remark-another-special-case}
pour les épaississements $(U \subset U')$ et $(I \subset I')$ sur
$(U \subset U')$, via $\text{id}_{U'}$ et $p'$.
La diagonale $\Delta'$ est une section de ce faisceau et, en faisant
agir sur cette section Compléments sur les morphismes d'espaces, lemme
\ref{spaces-more-morphisms-lemma-action-sheaf}
on obtient un isomorphisme
\begin{equation}
\label{stacks-more-morphisms-equation-isomorphism}
\SheafHom_{\mathcal{O}_U}(\Delta^*\Omega_{I/U}, \mathcal{C}_{U/U'})
\longrightarrow
\mathit{Aut}(\xi)
\end{equation}
sur $U_{spaces, \etale}$. Il reste trois choses à vérifier :
\begin{enumerate}
\item la construction de (\ref{stacks-more-morphisms-equation-isomorphism})
commute à la localisation \'etale,
\item $\SheafHom_{\mathcal{O}_U}(\Delta^*\Omega_{I/U}, \mathcal{C}_{U/U'})$
est un module quasi-cohérent sur $U$,
\item la composition dans $\mathit{Aut}(\xi)$ correspond
à l'addition des sections de ce module quasi-cohérent.
\end{enumerate}
Nous allons les vérifier dans cet ordre.

\medskip\noindent
Pour démontrer (1), il faut montrer que, si $f : V \to U$ est \'etale,
alors (\ref{stacks-more-morphisms-equation-isomorphism}), construit à partir de $\xi$ sur $U$,
se restreint à l'application (\ref{stacks-more-morphisms-equation-isomorphism})
$$
\SheafHom_{\mathcal{O}_V}(
\Delta_V^*\Omega_{V \times_\mathcal{X} V/V}, \mathcal{C}_{V/V'}) \to
\mathit{Aut}(\xi|_V)
$$
construite à partir de $\xi|_V$ sur $V$, dans $V_{spaces, \etale}$.
Cela résulte de la discussion dans la note ci-dessus
et de Compléments sur les morphismes d'espaces, lemme
\ref{spaces-more-morphisms-lemma-action-by-derivations-etale-localization}.

\medskip\noindent
Démonstration de (2). Puisque $p'$ est lisse, le morphisme $I \to U$ est lisse,
et le module des différentielles relatives $\Omega_{I/U}$ est donc
localement libre de type fini (Compléments sur les morphismes d'espaces, lemme
\ref{spaces-more-morphisms-lemma-smooth-omega-finite-locally-free}).
D'autre part, $\mathcal{C}_{U/U'}$ est quasi-cohérent
(Compléments sur les morphismes d'espaces, définition
\ref{spaces-more-morphisms-definition-conormal-sheaf}).
D'après Propriétés des espaces, lemme
\ref{spaces-properties-lemma-properties-quasi-coherent}
on conclut.

\medskip\noindent
Démonstration de (3). Il existe un morphisme $c' : I' \times_{p', U', q'} I' \to I'$
tel que $(U', I', p', q', c')$ soit un groupoïde en espaces algébriques
d'identité $\Delta'$. Voir par exemple
Champs algébriques, lemme \ref{algebraic-lemma-map-space-into-stack}.
La composition dans $\mathit{Aut}(\xi)$ est induite par le morphisme
$c'$ comme suit. Supposons donnés deux morphismes
$$
\delta'_1, \delta'_2 : U' \longrightarrow I'
$$
correspondant, comme ci-dessus, à des sections de $\mathit{Aut}(\xi)$ au-dessus de $U$;
autrement dit, on a $\delta'_i|_U = \Delta_U$ et
$p' \circ \delta'_i = \text{id}_{U'}$. Alors la composition
dans $\mathit{Aut}(\xi)$ est
$$
\delta'_1 \circ \delta'_2 = c'(\delta'_1 \circ q' \circ \delta'_2, \delta'_2)
$$
Nous omettons la vérification détaillée\footnote{Le lecteur voit immédiatement
qu'il est nécessaire de précomposer
$\delta'_1$ par $q' \circ \delta'_2$ pour obtenir un point à valeurs dans $U'$
bien défini du produit fibré $I' \times_{p', U', q'} I'$.}.
Nous sommes donc dans la situation décrite dans
Compléments sur les groupoïdes en espaces, section
\ref{spaces-more-groupoids-section-groupoid-sections}
et le résultat cherché découle du
Lemme de Compléments sur les groupoïdes en espaces
\ref{spaces-more-groupoids-lemma-composition-is-addition}.
\end{proof}

\begin{proposition}[Emerton]
\label{stacks-more-morphisms-proposition-affine-smooth-lift-to-first-order}
\begin{reference}
Courriel de Matthew Emerton daté du 27 avril 2016.
\end{reference}
Soit $\mathcal{X} \subset \mathcal{X}'$ un épaississement du premier ordre
de champs algébriques. Soit $W$ un schéma affine et soit
$W \to \mathcal{X}$ un morphisme lisse. Alors il existe
un diagramme cartésien
$$
\xymatrix{
W \ar[d] \ar[r] & W' \ar[d] \\
\mathcal{X} \ar[r] & \mathcal{X}'
}
$$
où $W' \to \mathcal{X}'$ est lisse et $W'$ affine.
\end{proposition}

\begin{proof}
Considérons la catégorie $p : \mathcal{C} \to W_{spaces, \etale}$
introduite dans la remarque \ref{stacks-more-morphisms-remark-gerbe-of-lifts}.
La proposition affirme qu'il existe un objet de $\mathcal{C}$
au-dessus de $W$. En effet, si l'on dispose d'un tel objet
$(W, W', a, i, y', \alpha)$, alors $W = \mathcal{X} \times_{\mathcal{X}'} W'$.
Par conséquent, $W \to W'$ est un épaississement d'espaces algébriques, donc
$W'$ est affine d'après
Compléments sur les morphismes d'espaces, lemme
\ref{spaces-more-morphisms-lemma-thickening-scheme}
et Compléments sur les morphismes, lemme
\ref{more-morphisms-lemma-thickening-affine-scheme}.

\medskip\noindent
Le lemme \ref{stacks-more-morphisms-lemma-gerbe-of-lifts} nous dit que $\mathcal{C}$ est une gerbe sur
$W_{spaces, \etale}$. Cela signifie que l'on peut trouver des solutions localement pour la topologie \'etale,
et que ces solutions locales sont localement isomorphes pour la topologie \'etale;
cette partie ne requiert pas que l'épaississement soit du premier ordre.
D'après le lemme \ref{stacks-more-morphisms-lemma-gerbe-of-lifts-first-order},
les faisceaux d'automorphismes des objets de notre gerbe sont abéliens et
se recollent pour former un module quasi-cohérent $\mathcal{G}$
sur $W_{spaces, \etale}$. Nous allons vérifier les conditions (1) et (2)
de Cohomologie sur les sites, lemme \ref{sites-cohomology-lemma-existence},
afin de conclure à l'existence d'un objet de $\mathcal{C}$ au-dessus de $W$.
La condition (1) est satisfaite : les recouvrements \'etales $\{W_i \to W\}$
dont chaque $W_i$ est affine sont cofinaux dans l'ensemble de tous les recouvrements.
Pour un tel recouvrement, $W_i$ et $W_i \times_W W_j$ sont affines,
et $H^1(W_i, \mathcal{G})$ ainsi que $H^1(W_i \times_W W_j, \mathcal{G})$
sont nuls : la cohomologie d'un module quasi-cohérent sur un
espace algébrique affine est nulle, par exemple d'après Cohomologie des espaces, proposition
\ref{spaces-cohomology-proposition-vanishing}.
Enfin, la condition (2) est que $H^2(W, \mathcal{G}) = 0$
pour notre faisceau quasi-cohérent $\mathcal{G}$, ce qui résulte encore
de Cohomologie des espaces, proposition
\ref{spaces-cohomology-proposition-vanishing}.
Cela achève la démonstration.
\end{proof}






\section{Déformations infinitésimales}
\label{stacks-more-morphisms-section-inf}

\noindent
Nous poursuivons la discussion de
Les axiomes d'Artin, section \ref{artin-section-inf}.

\begin{lemma}
\label{stacks-more-morphisms-lemma-inf-quasi-coherent}
Soit $\mathcal{X}$ un champ algébrique sur un schéma $S$.
Supposons que $\mathcal{I}_\mathcal{X} \to \mathcal{X}$ soit localement
de présentation finie. Soit $A \to B$ un homomorphisme plat de $S$-algèbres.
Soit $x$ un objet de $\mathcal{X}$ sur $A$, et posons $y = x|_B$.
Alors $\text{Inf}_x(M) \otimes_A B = \text{Inf}_y(M \otimes_A B)$.
\end{lemma}

\begin{proof}
Rappelons que $\text{Inf}_x(M)$ est l'ensemble des automorphismes de la
déformation triviale de $x$ à $A[M]$ qui induisent l'automorphisme
identique de $x$ sur $A$. La déformation triviale est
l'image inverse de $x$ sur $\Spec(A[M])$ par $\Spec(A[M]) \to \Spec(A)$.
Soit $G \to \Spec(A)$ l'espace algébrique groupe des automorphismes de $x$
(celui-ci existe puisque $\mathcal{X}$ est un champ algébrique).
Soit $e : \Spec(A) \to G$ l'élément neutre.
La discussion de Compléments sur les morphismes d'espaces, section
\ref{spaces-more-morphisms-section-action-by-derivations}
donne
$$
\text{Inf}_x(M) = \Hom_A(e^*\Omega_{G/A}, M)
$$
De même,
$$
\text{Inf}_y(M \otimes_A B) = \Hom_B(e_B^*\Omega_{G_B/B}, M \otimes_A B)
$$
Puisque $G \to \Spec(A)$ est localement de présentation finie par
hypothèse, on voit que $\Omega_{G/A}$ est localement de
présentation finie, voir
Compléments sur les morphismes d'espaces, lemme
\ref{spaces-more-morphisms-lemma-finite-presentation-differentials}.
Par conséquent, $e^*\Omega_{G/A}$ est un $A$-module de présentation finie.
De plus, $\Omega_{G_B/B}$ est l'image inverse de $\Omega_{G/A}$ d'après
Compléments sur les morphismes d'espaces, lemme
\ref{spaces-more-morphisms-lemma-base-change-differentials}.
Ainsi, $e_B^*\Omega_{G_B/B} = e^*\Omega_{G/A} \otimes_A B$.
On conclut par Compléments d'algèbre, lemme
\ref{more-algebra-lemma-pseudo-coherence-and-base-change-ext}.
\end{proof}

\begin{lemma}
\label{stacks-more-morphisms-lemma-sheaf-of-infinitesimal-lifts}
Soit $\mathcal{X}$ un champ algébrique sur un schéma de base $S$.
Supposons que $\mathcal{I}_\mathcal{X} \to \mathcal{X}$ soit localement
de présentation finie. Soit $(A' \to A, x)$ une situation de déformation.
Alors le foncteur
$$
F : B' \longmapsto
\{\text{relèvements de }x|_{B' \otimes_{A'} A}\text{ à } B'\}/\text{isomorphismes}
$$
est un faisceau sur le site $(\textit{Aff}/\Spec(A'))_{fppf}$ de
Topologies, définition \ref{topologies-definition-big-small-fppf}.
\end{lemma}

\begin{proof}
Soit $\{T'_i \to T'\}_{i = 1, \ldots, n}$ un recouvrement fppf standard
de schémas affines sur $A'$. Écrivons $T' = \Spec(B')$. Comme d'habitude, notons
$$
T'_{i_0 \ldots i_p} =
T'_{i_0} \times_{T'} \ldots \times_{T'} T'_{i_p} = \Spec(B'_{i_0 \ldots i_p})
$$
où l'anneau est un produit tensoriel convenable.
Posons $B = B' \otimes_{A'} A$ et
$B_{i_0 \ldots i_p} = B'_{i_0 \ldots i_p} \otimes_{A'} A$.
Notons $y = x|_B$ et $y_{i_0 \ldots i_p} = x|_{B_{i_0 \ldots i_p}}$.
Soient $\gamma_i \in F(B'_i)$ tels que $\gamma_{i_0}$ et $\gamma_{i_1}$
s'envoient sur le même élément de $F(B'_{i_0i_1})$.
Il faut trouver un unique $\gamma \in F(B')$ dont l'image soit
$\gamma_i$ dans $F(B'_i)$.

\medskip\noindent
Choisissons un objet véritable $y'_i$ de $\textit{Lift}(y_i, B'_i)$
dans la classe d'isomorphisme $\gamma_i$.
Choisissons des isomorphismes
$\varphi_{i_0i_1} : y'_{i_0}|_{B'_{i_0i_1}} \to y'_{i_1}|_{B'_{i_0i_1}}$
dans la catégorie $\textit{Lift}(y_{i_0i_1}, B'_{i_0i_1})$.
Si les applications $\varphi_{i_0i_1}$ satisfont à la condition de cocycle,
on obtient l'objet $\gamma$ cherché, puisque $\mathcal{X}$ est
un champ pour la topologie fppf. La condition de cocycle affirme que le
composé
$$
y'_{i_0}|_{B'_{i_0i_1i_2}}
\xrightarrow{\varphi_{i_0i_1}|_{B'_{i_0i_1i_2}}}
y'_{i_1}|_{B'_{i_0i_1i_2}}
\xrightarrow{\varphi_{i_1i_2}|_{B'_{i_0i_1i_2}}}
y'_{i_2}|_{B'_{i_0i_1i_2}}
\xrightarrow{\varphi_{i_2i_0}|_{B'_{i_0i_1i_2}}}
y'_{i_0}|_{B'_{i_0i_1i_2}}
$$
est l'identité. Sinon, ces applications donnent des éléments
$$
\delta_{i_0i_1i_2} \in
\text{Inf}_{y_{i_0i_1i_2}}(J_{i_0i_1i_2}) =
\text{Inf}_y(J) \otimes_B B_{i_0i_1i_2}
$$
Ici, $J = \Ker(B' \to B)$ et
$J_{i_0 \ldots i_p} = \Ker(B'_{i_0 \ldots i_p} \to B_{i_0 \ldots i_p})$.
L'égalité de la formule affichée résulte du
lemme \ref{stacks-more-morphisms-lemma-inf-quasi-coherent} appliqué à
$B' \to B'_{i_0 \ldots i_p}$, à $y$ et à $y_{i_0 \ldots i_p}$,
ainsi que de la platitude des applications $B' \to B'_{i_0 \ldots i_p}$,
qui garantit également que
$J_{i_0 \ldots i_p} = J \otimes_{B'} B'_{i_0 \ldots i_p}$.
Un calcul (omis) montre que $\delta_{i_0i_1i_2}$ définit
un $2$-cocycle dans le complexe de {\v C}ech
$$
\prod \text{Inf}_y(J) \otimes_B B_{i_0} \to
\prod \text{Inf}_y(J) \otimes_B B_{i_0i_1} \to
\prod \text{Inf}_y(J) \otimes_B B_{i_0i_1i_2} \to \ldots
$$
D'après Descente, lemme \ref{descent-lemma-standard-covering-Cech-quasi-coherent},
ce complexe est acyclique en degrés positifs et vérifie $H^0 = \text{Inf}_y(J)$.
Puisque $\text{Inf}_{y_{i_0i_1}}(J_{i_0i_1})$ agit sur les
morphismes (Les axiomes d'Artin, remarque \ref{artin-remark-automorphisms}),
cela signifie que l'on peut modifier le choix de $\varphi_{i_0i_1}$
pour se ramener au cas où $\delta_{i_0i_1i_2} = 0$.

\medskip\noindent
Unicité. Il reste à montrer qu'il existe au plus un $\gamma$ dont la restriction
soit $\gamma_i$ pour tout $i$. Supposons donnés des objets $y', z'$ de
$\textit{Lift}(y, B')$ et des isomorphismes $\psi_i : y'|_{B'_i} \to z'|_{B'_i}$
dans $\textit{Lift}(y_i, B'_i)$. On peut alors considérer
$$
\psi_{i_1}^{-1} \circ \psi_{i_0} \in
\text{Inf}_{y_{i_0i_1}}(J_{i_0i_1}) =
\text{Inf}_y(J) \otimes_B B_{i_0i_1}
$$
Comme précédemment, l'obstruction à l'existence d'un isomorphisme entre
$y'$ et $z'$ sur $B'$ est un élément du $H^1$ du
complexe de {\v C}ech affiché ci-dessus, lequel est nul.
\end{proof}

\begin{lemma}
\label{stacks-more-morphisms-lemma-T-quasi-coherent}
Soit $\mathcal{X}$ un champ algébrique sur un schéma $S$, dont le
morphisme structural $\mathcal{X} \to S$ est localement de présentation finie.
Soit $A \to B$ un homomorphisme plat de $S$-algèbres.
Soit $x$ un objet de $\mathcal{X}$ sur $A$.
Alors $T_x(M) \otimes_A B = T_y(M \otimes_A B)$.
\end{lemma}

\begin{proof}
Choisissons un schéma $U$ et un morphisme lisse surjectif $U \to \mathcal{X}$.
Réduisons d'abord le lemme au cas où $x$ se relève à $U$.
Rappelons que $T_x(M)$ est l'ensemble des classes d'isomorphisme des relèvements de $x$
à $A[M]$. Par conséquent, le
lemme \ref{stacks-more-morphisms-lemma-sheaf-of-infinitesimal-lifts}\footnote{Ce lemme s'applique :
$\Delta : \mathcal{X} \to \mathcal{X} \times_S \mathcal{X}$
est localement de présentation finie d'après
Morphismes de champs, lemme
\ref{stacks-morphisms-lemma-finite-presentation-permanence}
et l'hypothèse que $\mathcal{X} \to S$ est localement de présentation finie.
Par conséquent, $\mathcal{I}_\mathcal{X} \to \mathcal{X}$ est localement
de présentation finie en tant que changement de base de $\Delta$.}
affirme que la règle
$$
A_1 \mapsto T_{x|_{A_1}}(M \otimes_A A_1)
$$
est un faisceau sur le petit site \'etale de $\Spec(A)$; le produit tensoriel
est nécessaire pour que $A[M] \to A_1[M \otimes_A A_1]$ soit un homomorphisme plat.
On peut choisir un homomorphisme d'anneaux \'etale fidèlement plat $A \to A_1$
tel que $x|_{A_1}$ se relève en un morphisme $u_1 : \Spec(A_1) \to U$, voir
par exemple Faisceaux sur les champs, lemme
\ref{stacks-sheaves-lemma-surjective-flat-locally-finite-presentation}.
Écrivons $A_2 = A_1 \otimes_A A_1$ et posons $B_1 = B \otimes_A A_1$
et $B_2 = B \otimes_A A_2$. Considérons le diagramme
$$
\xymatrix{
0 \ar[r] &
T_y(M \otimes_A B) \ar[r] &
T_{y|_{B_1}}(M \otimes_A B_1) \ar[r] &
T_{y|_{B_2}}(M \otimes_A B_2) \\
0 \ar[r] &
T_x(M) \ar[r] \ar[u] &
T_{x|_{A_1}}(M \otimes_A A_1) \ar[r] \ar[u] &
T_{x|_{A_2}}(M \otimes_A A_2) \ar[u]
}
$$
Les lignes sont exactes par la condition de faisceau. On a
$M \otimes_A B_i = (M \otimes_A A_i) \otimes_{A_i} B_i$.
Ainsi, si l'on démontre le résultat pour les flèches verticales du milieu et de droite,
le résultat en découle. On est donc ramené au cas étudié au
paragraphe suivant.

\medskip\noindent
Supposons que $x$ soit l'image d'un morphisme $u : \Spec(A) \to U$.
Observons que $T_u(M) \to T_x(M)$ est surjectif, puisque
$U \to \mathcal{X}$ est lisse et représentable par des
espaces algébriques, voir Critères de représentabilité, lemme
\ref{criteria-lemma-representable-by-spaces-formally-smooth}
(voir la discussion qui précède pour une explication) et
Compléments sur les morphismes d'espaces, lemme
\ref{spaces-more-morphisms-lemma-smooth-formally-smooth}.
Posons $R = U \times_\mathcal{X} U$. Rappelons que l'on obtient
un groupoïde $(U, R, s, t, c, e, i)$ en espaces algébriques avec
$\mathcal{X} = [U/R]$. D'après
Les axiomes d'Artin, lemme \ref{artin-lemma-ses-inf-and-T},
on a une suite exacte
$$
T_{e \circ u}(M) \to T_u(M) \oplus T_u(M) \to T_x(M) \to 0
$$
où la nullité du terme de droite a été démontrée ci-dessus. Une suite analogue
vaut après changement de base à $B$. Le résultat cherché découle donc de
celui du lemme pour
$T_u(M)$ et $T_{e \circ u}(M)$. On est ainsi ramené au cas étudié
dans le paragraphe suivant.

\medskip\noindent
Supposons que $\mathcal{X} = X$ soit un espace algébrique localement de
présentation finie sur $S$. Alors
$$
T_x(M) = \Hom_A(x^*\Omega_{X/S}, M)
$$
d'après la discussion de Compléments sur les morphismes d'espaces, section
\ref{spaces-more-morphisms-section-action-by-derivations}.
De même,
$$
T_y(M \otimes_A B) = \Hom_B(y^*\Omega_{X/S}, M \otimes_A B)
$$
Puisque $X \to S$ est localement de présentation finie, on voit que
$\Omega_{X/S}$ est localement de présentation finie, voir
Compléments sur les morphismes d'espaces, lemme
\ref{spaces-more-morphisms-lemma-finite-presentation-differentials}.
Ainsi, $x^*\Omega_{X/S}$ est un $A$-module de présentation finie.
Évidemment, on a $y^*\Omega_{X/S} = x^*\Omega_{X/S} \otimes_A B$.
On conclut par Compléments d'algèbre, lemme
\ref{more-algebra-lemma-pseudo-coherence-and-base-change-ext}.
\end{proof}

\begin{lemma}
\label{stacks-more-morphisms-lemma-local-lift-enough}
Soit $\mathcal{X}$ un champ algébrique sur un schéma $S$, dont le
morphisme structural $\mathcal{X} \to S$ est localement de présentation finie.
Soit $(A' \to A, x)$ une situation de déformation. S'il existe une
$A'$-algèbre $B'$ fidèlement plate et de présentation finie, ainsi qu'un
objet $y'$ de $\mathcal{X}$ sur $B'$ relevant $x|_{B' \otimes_{A'} A}$,
alors il existe un objet $x'$ sur $A'$ relevant $x$.
\end{lemma}

\begin{proof}
Soit $I = \Ker(A' \to A)$. Posons $B'_1 = B' \otimes_{A'} B'$ et
$B'_2 = B' \otimes_{A'} B' \otimes_{A'} B'$. Soient
$J = IB'$, $J_1 = IB'_1$, $J_2 = IB'_2$ et
$B = B'/J$, $B_1 = B'_1/J_1$, $B_2 = B'_2/J_2$.
Posons $y = x|_B$, $y_1 = x|_{B_1}$, $y_2 = x|_{B_2}$.
Soit $F$ le faisceau fppf du
lemme \ref{stacks-more-morphisms-lemma-sheaf-of-infinitesimal-lifts}
(qui s'applique, voir la note dans la démonstration du
lemme \ref{stacks-more-morphisms-lemma-T-quasi-coherent}).
On a donc un diagramme d'égalisateur
$$
\xymatrix{
F(A') \ar[r] &
F(B') \ar@<1ex>[r] \ar@<-1ex>[r] &
F(B'_1)
}
$$
D'autre part, on a $F(B') = \text{Lift}(y, B')$,
$F(B'_1) = \text{Lift}(y_1, B'_1)$ et $F(B'_2) = \text{Lift}(y_2, B'_2)$,
dans la terminologie de Les axiomes d'Artin, section \ref{artin-section-inf}.
Ces ensembles sont non vides et sont (canoniquement) des espaces principaux homogènes sous
$T_y(J)$, $T_{y_1}(J_1)$ et $T_{y_2}(J_2)$, voir
Les axiomes d'Artin, lemme \ref{artin-lemma-properties-lift-RS-star}.
La différence des deux images de $y'$ dans
$F(B'_1)$ est donc un élément
$$
\delta_1 \in T_{y_1}(J_1) = T_x(I) \otimes_A B_1
$$
L'égalité de la formule affichée résulte du
lemme \ref{stacks-more-morphisms-lemma-T-quasi-coherent} appliqué à
$A' \to B'_1$, à $x$ et à $y_1$, ainsi que de la platitude de $A' \to B'_1$,
qui garantit également que $J_1 = I \otimes_{A'} B'_1$. On a des
égalités analogues pour $B'$ et $B'_2$.
Un calcul (omis) montre que $\delta_1$ définit
un $1$-cocycle dans le complexe de {\v C}ech
$$
T_x(I) \otimes_A B \to
T_x(I) \otimes_A B_1 \to
T_x(I) \otimes_A B_2 \to \ldots
$$
D'après Descente, lemme \ref{descent-lemma-standard-covering-Cech-quasi-coherent},
ce complexe est acyclique en degrés positifs et vérifie $H^0 = T_x(I)$.
On peut donc choisir un élément de $T_x(I) \otimes_A B = T_y(J)$
dont le cobord est $\delta_1$. En remplaçant $y'$ par le résultat
de l'action de cet élément, on obtient un nouveau choix $y'$ tel que $\delta_1 = 0$.
Ainsi, $y'$ s'envoie sur le même élément par les deux applications
$F(B') \to F(B'_1)$, et on obtient un élément de $F(A')$ par
la condition de faisceau.
\end{proof}












\section{Morphismes formellement lisses}
\label{stacks-more-morphisms-section-formally-smooth}

\noindent
Dans cette section, nous introduisons la notion de morphisme formellement lisse
$\mathcal{X} \to \mathcal{Y}$ de champs algébriques. Un tel morphisme est
caractérisé par la propriété que les $T$-points de $\mathcal{X}$ se relèvent
aux épaississements infinitésimaux de $T$, pourvu que $T$ soit affine.
Le résultat principal affirme qu'un morphisme formellement lisse et
localement de présentation finie est lisse, voir le
lemme \ref{stacks-more-morphisms-lemma-smooth-formally-smooth}.
Il s'avère que ce critère est souvent plus facile à utiliser que le
critère jacobien.

\begin{definition}
\label{stacks-more-morphisms-definition-formally-smooth}
Un morphisme $f : \mathcal{X} \to \mathcal{Y}$ de champs algébriques est dit
{\it formellement lisse} s'il est formellement lisse sur les objets en tant que
$1$-morphisme de catégories fibrées en groupoïdes, comme expliqué dans
Critères de représentabilité, section \ref{criteria-section-formally-smooth}.
\end{definition}

\noindent
Traduisons la condition de la définition dans le langage que nous employons
actuellement (voir
Propriétés des champs, section \ref{stacks-properties-section-conventions}).
Soit $f : \mathcal{X} \to \mathcal{Y}$ un morphisme de champs algébriques.
Considérons un diagramme $2$-commutatif aux flèches pleines
\begin{equation}
\label{stacks-more-morphisms-equation-diagram}
\vcenter{
\xymatrix{
T \ar[r]_-x \ar[d]_i & \mathcal{X} \ar[d]^f \\
T' \ar[r]^-y \ar@{..>}[ru] & \mathcal{Y}
}
}
\end{equation}
où $i : T \to T'$ est un épaississement du premier ordre de schémas affines.
Soit
$$
\gamma : y \circ i \longrightarrow f \circ x
$$
un $2$-morphisme matérialisant la $2$-commutativité du diagramme.
(Notation comme dans Catégories, sections \ref{categories-section-formal-cat-cat}
et \ref{categories-section-2-categories}.)
Étant donnés (\ref{stacks-more-morphisms-equation-diagram}) et $\gamma$,
une {\it flèche pointillée} est un triplet $(x', \alpha, \beta)$ formé d'un
morphisme $x' : T' \to \mathcal{X}$ et de $2$-flèches
$\alpha : x' \circ i \to x$, $\beta : y \to f \circ x'$
tels que
$\gamma = (\text{id}_f \star \alpha) \circ (\beta \star \text{id}_i)$,
autrement dit tel que
$$
\xymatrix{
& f \circ x' \circ i \ar[rd]^{\text{id}_f \star \alpha} \\
y \circ i \ar[ru]^{\beta \star \text{id}_i} \ar[rr]^\gamma & &
f \circ x
}
$$
soit commutatif. Un {\it morphisme de flèches pointillées}
$(x'_1, \alpha_1, \beta_1) \to (x'_2, \alpha_2, \beta_2)$ est une
$2$-flèche $\theta : x'_1 \to x'_2$ telle que
$\alpha_1 = \alpha_2 \circ (\theta \star \text{id}_i)$ et
$\beta_2 = (\text{id}_f \star \theta) \circ \beta_1$.

\medskip\noindent
La catégorie des flèches pointillées que l'on vient de décrire est un cas particulier de
Catégories, définition \ref{categories-definition-dotted-arrows}.

\begin{lemma}
\label{stacks-more-morphisms-lemma-reformulate-formal-smoothness}
Un morphisme $f : \mathcal{X} \to \mathcal{Y}$ de champs algébriques est
formellement lisse (définition \ref{stacks-more-morphisms-definition-formally-smooth})
si et seulement si, pour tout diagramme (\ref{stacks-more-morphisms-equation-diagram}) et tout $\gamma$,
la catégorie des flèches pointillées est non vide.
\end{lemma}

\begin{proof}
La traduction entre les différents langages est omise.
\end{proof}

\begin{lemma}
\label{stacks-more-morphisms-lemma-base-change-formal-smoothness}
Le changement de base d'un morphisme formellement lisse de champs algébriques
par un morphisme quelconque de champs algébriques est formellement lisse.
\end{lemma}

\begin{proof}
Cela résulte de
Catégories, lemme \ref{categories-lemma-cat-dotted-arrows-base-change},
et de la définition.
\end{proof}

\begin{lemma}
\label{stacks-more-morphisms-lemma-composition-formal-smoothness}
Le composé de morphismes formellement lisses de champs algébriques
est formellement lisse.
\end{lemma}

\begin{proof}
Cela résulte de
Catégories, lemme \ref{categories-lemma-cat-dotted-arrows-composition},
et de la définition.
\end{proof}

\begin{lemma}
\label{stacks-more-morphisms-lemma-formal-smoothness-representable}
Soit $f : \mathcal{X} \to \mathcal{Y}$ un morphisme de champs algébriques
représentable par des espaces algébriques. Les conditions suivantes sont équivalentes :
\begin{enumerate}
\item $f$ est formellement lisse,
\item pour tout schéma $T$ et tout morphisme $T \to \mathcal{Y}$,
le morphisme $\mathcal{X} \times_\mathcal{Y} T \to T$
est formellement lisse comme morphisme d'espaces algébriques.
\end{enumerate}
\end{lemma}

\begin{proof}
Cela résulte de
Catégories, lemme \ref{categories-lemma-cat-dotted-arrows-base-change},
et de la définition.
\end{proof}

\begin{lemma}
\label{stacks-more-morphisms-lemma-lift-to-smooth}
Soit $T \to T'$ un épaississement du premier ordre de schémas affines.
Soit $\mathcal{X}'$ un champ algébrique sur $T'$,
dont le morphisme structural $\mathcal{X}' \to T'$ est lisse.
Soit $x : T \to \mathcal{X}'$ un morphisme
sur $T'$. Alors il existe un morphisme $x' : T' \to \mathcal{X}'$
sur $T'$ tel que $x'|_T = x$.
\end{lemma}

\begin{proof}
On peut appliquer le résultat du lemme \ref{stacks-more-morphisms-lemma-local-lift-enough}.
Il suffit donc de construire un morphisme lisse surjectif
$W' \to T'$, avec $W'$ affine, tel que
$x|_{T \times_{T'} W'}$ se relève à $W'$.
(Nous invitons le lecteur à trouver sa propre démonstration de ce fait
en utilisant le résultat analogue déjà établi
pour les espaces algébriques.) Choisissons un
schéma $U'$ et un morphisme lisse surjectif $U' \to \mathcal{X}'$.
Observons que $U' \to T'$ est lisse et que la projection
$T \times_{\mathcal{X}'} U' \to T$ est lisse et surjective.
Choisissons un schéma affine $W$ et un morphisme \'etale
$W \to T \times_{\mathcal{X}'} U'$ tel que $W \to T$
soit surjectif. Alors $W \to T$ est un morphisme lisse de
schémas affines. Après avoir remplacé $W$ par une réunion disjointe d'ouverts
affines principaux, on peut supposer qu'il existe un
morphisme lisse de schémas affines $W' \to T'$ tel que
$W = T \times_{T'} W'$, voir Algèbre, lemme \ref{algebra-lemma-lift-smooth}.
D'après Compléments sur les morphismes d'espaces, lemme
\ref{spaces-more-morphisms-lemma-smooth-formally-smooth}
on peut trouver un morphisme $W' \to U'$ sur $T'$ relevant
le morphisme donné $W \to U'$. Cela achève la démonstration.
\end{proof}

\noindent
Le lemme suivant est le résultat principal de cette section.
Combiné à
Limites de champs, proposition
\ref{stacks-limits-proposition-characterize-locally-finite-presentation},
il implique que l'on peut reconnaître si un morphisme de champs algébriques
$f : \mathcal{X} \to \mathcal{Y}$ est lisse en termes de
propriétés « simples » du $1$-morphisme de champs en groupoïdes
$\mathcal{X} \to \mathcal{Y}$.

\begin{lemma}[Critère infinitésimal de relèvement]
\label{stacks-more-morphisms-lemma-smooth-formally-smooth}
Soit $f : \mathcal{X} \to \mathcal{Y}$ un morphisme de champs algébriques.
Les conditions suivantes sont équivalentes :
\begin{enumerate}
\item Le morphisme $f$ est lisse.
\item Le morphisme $f$ est localement de présentation finie et
formellement lisse.
\end{enumerate}
\end{lemma}

\begin{proof}
Supposons $f$ lisse. Alors $f$ est localement de présentation finie d'après
Morphismes de champs, lemme
\ref{stacks-morphisms-lemma-smooth-locally-finite-presentation}.
Étant donné un diagramme (\ref{stacks-more-morphisms-equation-diagram})
et un $\gamma : y \circ i \to f \circ x$, il suffit donc de trouver une
flèche pointillée (voir Lemme \ref{stacks-more-morphisms-lemma-reformulate-formal-smoothness}).
En formant les produits fibrés, on obtient
$$
\xymatrix{
T \ar[d] \ar[r] &
T' \times_\mathcal{Y} \mathcal{X} \ar[d] \ar[r] &
\mathcal{X} \ar[d] \\
T' \ar[r] & T' \ar[r] & \mathcal{Y}
}
$$
Il suffit donc de trouver une flèche pointillée dans le
carré de gauche. Puisque $T' \times_\mathcal{Y} \mathcal{X} \to T'$
est lisse
(Morphismes de champs, lemme \ref{stacks-morphisms-lemma-base-change-smooth}),
l'existence d'une flèche pointillée dans le carré de gauche est garantie par le
lemme \ref{stacks-more-morphisms-lemma-lift-to-smooth}.

\medskip\noindent
Réciproquement, supposons que $f$ soit localement de présentation finie et
formellement lisse. Choisissons un schéma $U$ et un morphisme lisse surjectif
$U \to \mathcal{X}$. Alors $a : U \to \mathcal{X}$ et $b : U \to \mathcal{Y}$
sont représentables par des espaces algébriques et localement de présentation finie
(utiliser Morphismes de champs, lemme
\ref{stacks-morphisms-lemma-composition-finite-presentation}
et le fait constaté ci-dessus
qu'un morphisme lisse est localement de présentation finie).
Nous allons appliquer le principe général de
Champs algébriques, lemme
\ref{algebraic-lemma-representable-transformations-property-implication}
en prenant comme donnée l'équivalence de Compléments sur les morphismes d'espaces,
lemme \ref{spaces-more-morphisms-lemma-smooth-formally-smooth},
et utiliser simultanément la traduction du
Lemme de Critères de représentabilité
\ref{criteria-lemma-representable-by-spaces-formally-smooth}.
On applique d'abord ce principe à $a$ pour voir que $a$ est
formellement lisse sur les objets. Ensuite, on utilise que $f$ est
formellement lisse sur les objets par hypothèse
(voir Lemme \ref{stacks-more-morphisms-lemma-reformulate-formal-smoothness})
ainsi que Critères de représentabilité, lemme
\ref{criteria-lemma-composition-formally-smooth}
pour voir que $b = f \circ a$ est formellement lisse sur les objets.
On applique alors le principe une fois encore pour conclure que $b$ est lisse.
Cela signifie que $f$ est lisse par définition de la lissité
des morphismes de champs algébriques, ce qui achève la démonstration.
\end{proof}









\section{Éclatement et platitude}
\label{stacks-more-morphisms-section-blowup-flat}

\noindent
Cette section expose rapidement ce que l'on peut déduire de
Compléments sur les morphismes d'espaces, sections
\ref{spaces-more-morphisms-section-blowup-flat} et
\ref{spaces-more-morphisms-section-applications-flattening-by-blowing-up}
pour les champs algébriques sur des espaces algébriques.

\begin{lemma}
\label{stacks-more-morphisms-lemma-flatten-stack}
Soit $f : \mathcal{X} \to Y$ un morphisme d'un champ algébrique
vers un espace algébrique. Soit $V \subset Y$ un sous-espace ouvert. Supposons que
\begin{enumerate}
\item $Y$ soit quasi-compact et quasi-séparé,
\item $f$ soit de type fini et quasi-séparé,
\item $V$ soit quasi-compact, et
\item $\mathcal{X}_V$ soit plat et localement de présentation finie sur $V$.
\end{enumerate}
Alors il existe un éclatement $V$-admissible $Y' \to Y$
et un sous-champ fermé $\mathcal{X}' \subset \mathcal{X}_{Y'}$
tel que $\mathcal{X}'_V = \mathcal{X}_V$ et que
$\mathcal{X}' \to Y'$ soit plat et de présentation finie.
\end{lemma}

\begin{proof}
Observons que $\mathcal{X}$ est quasi-compact.
Choisissons un schéma affine $U$ et un morphisme lisse
surjectif $U \to \mathcal{X}$.
Soit $R = U \times_\mathcal{X} U$; on obtient ainsi
un groupoïde $(U, R, s, t, c)$ en espaces algébriques sur $Y$, avec
$\mathcal{X} = [U/R]$
(Champs algébriques, lemme \ref{algebraic-lemma-stack-presentation}).
On peut appliquer le
Lemme de Compléments sur les morphismes d'espaces
\ref{spaces-more-morphisms-lemma-flat-after-blowing-up}
à $U \to Y$ et à l'ouvert $V \subset Y$.
On obtient ainsi un éclatement $V$-admissible $Y' \to Y$
tel que la transformée stricte $U' \subset U_{Y'}$
soit plate et de présentation finie sur $Y'$.
Soit $R' \subset R_{Y'}$ la transformée stricte de $R$.
Puisque $s$ et $t$ sont lisses (et en particulier plats),
il résulte de
Diviseurs sur les espaces, lemme
\ref{spaces-divisors-lemma-strict-transform-flat}
que l'on a des diagrammes cartésiens
$$
\vcenter{
\xymatrix{
R' \ar[r] \ar[d] & R_{Y'} \ar[d]^{s_{Y'}} \\
U' \ar[r] & U_{Y'}
}
}
\quad\text{et}\quad
\vcenter{
\xymatrix{
R' \ar[r] \ar[d] & R_{Y'} \ar[d]^{t_{Y'}} \\
U' \ar[r] & U_{Y'}
}
}
$$
Autrement dit, $U'$ est un sous-espace fermé invariant par $R_{Y'}$
de $U_{Y'}$. Ainsi, $U'$ définit un sous-champ fermé
$\mathcal{X}' \subset \mathcal{X}_{Y'}$ d'après
Propriétés des champs, lemme
\ref{stacks-properties-lemma-substacks-presentation}.
Le morphisme $\mathcal{X}' \to Y'$ est plat et
localement de présentation finie, parce que c'est
le cas de $U' \to Y'$. D'autre part,
on sait déjà que $\mathcal{X}' \to Y'$ est quasi-compact et quasi-séparé
(par nos hypothèses sur $f$ et parce que c'est le cas des immersions fermées).
Cela achève la démonstration.
\end{proof}















\section{Lemme de Chow pour les champs algébriques}
\label{stacks-more-morphisms-section-chow}

\noindent
Dans cette section, nous étudions le lemme de Chow pour les champs algébriques.

\begin{lemma}
\label{stacks-more-morphisms-lemma-finite-cover-factor}
Soit $Y$ un espace algébrique quasi-compact et quasi-séparé.
Soit $V \subset Y$ un ouvert quasi-compact. Soit $f : \mathcal{X} \to V$
surjectif, plat et localement de présentation finie.
Alors il existe un morphisme fini surjectif $g : Y' \to Y$ tel que
$V' = g^{-1}(V) \to Y$ se factorise localement pour la topologie de Zariski par $f$.
\end{lemma}

\begin{proof}
Commençons par le démontrer lorsque $Y$ est un schéma.
On peut choisir un schéma $U$ et un morphisme lisse surjectif
$U \to \mathcal{X}$. Alors $\{U \to V\}$ est un recouvrement fppf de schémas.
D'après Compléments sur les morphismes, lemme \ref{more-morphisms-lemma-dominate-fppf-finite},
il existe un morphisme fini surjectif
$V' \to V$ tel que $V' \to V$ se factorise localement pour la topologie de Zariski
par $U$. D'après
Compléments sur les morphismes, lemme
\ref{more-morphisms-lemma-extend-finite-surjective-morphisms}
on peut trouver un morphisme fini surjectif $Y' \to Y$
dont la restriction à $V$ est $V' \to V$, comme souhaité.

\medskip\noindent
Si $Y$ est un espace algébrique, le lemme résulte d'abord d'un
changement de base fini par un morphisme fini surjectif
$Y' \to Y$, où $Y'$ est un schéma. Voir
Limites d'espaces, proposition
\ref{spaces-limits-proposition-there-is-a-scheme-finite-over}.
\end{proof}

\begin{lemma}
\label{stacks-more-morphisms-lemma-make-section}
Soit $f : \mathcal{X} \to Y$ un morphisme d'un champ algébrique
vers un espace algébrique. Soit $V \subset Y$ un sous-espace ouvert.
Supposons que
\begin{enumerate}
\item $f$ soit séparé et de type fini,
\item $Y$ soit quasi-compact et quasi-séparé,
\item $V$ soit quasi-compact, et
\item $\mathcal{X}_V$ soit une gerbe sur $V$.
\end{enumerate}
Alors il existe un diagramme commutatif
$$
\xymatrix{
\overline{Z} \ar[rd]_{\overline{g}} &
Z \ar[l]^j \ar[d]_g \ar[r]_h & \mathcal{X} \ar[ld]^f \\
& Y
}
$$
où $j$ est une immersion ouverte, $\overline{g}$ et $h$ sont propres,
et tel que $|V|$ soit contenu dans l'image de $|g|$.
\end{lemma}

\begin{proof}
Supposons donné un diagramme commutatif
$$
\xymatrix{
\mathcal{X}' \ar[d]_{f'} \ar[r] & \mathcal{X} \ar[d]^f \\
Y' \ar[r] & Y
}
$$
et un ouvert quasi-compact $V' \subset Y'$, tels que
$Y' \to Y$ soit un morphisme propre d'espaces algébriques,
$\mathcal{X}' \to \mathcal{X}$ soit un morphisme propre de champs algébriques,
$V' \subset Y'$ s'envoie surjectivement sur $V$, et
$\mathcal{X}'_{V'}$ soit une gerbe sur $V'$.
Il suffit alors de démontrer le lemme pour la paire
$(f' : \mathcal{X}' \to Y', V')$. Nous omettons certains détails.

\medskip\noindent
Stratégie générale de la démonstration. Nous nous ramènerons
au cas où l'image de $f$ est ouverte et où $f$
admet une section au-dessus de cet ouvert, en appliquant à plusieurs reprises
la remarque précédente. Chaque étape est directe, mais elles sont
assez nombreuses, ce qui rend la démonstration un peu longue.

\medskip\noindent
En utilisant Limites d'espaces, proposition
\ref{spaces-limits-proposition-there-is-a-scheme-finite-over}
on se ramène au cas où $Y$ est un schéma.
(Soit $Y' \to Y$ un morphisme fini surjectif, où $Y'$ est
un schéma. Posons $\mathcal{X}' = \mathcal{X}_{Y'}$ et appliquons
la remarque initiale de la démonstration.)

\medskip\noindent
En utilisant le lemme \ref{stacks-more-morphisms-lemma-flatten-stack}
(et Morphismes de champs, lemme \ref{stacks-morphisms-lemma-gerbe-fppf},
pour voir qu'une gerbe est plate et localement de présentation finie),
on se ramène au cas où $f$ est plat et de présentation finie.

\medskip\noindent
Puisque $f$ est plat et localement de
présentation finie, l'image de $|f|$ est un ouvert $W \subset Y$.
Puisque $\mathcal{X}$ est quasi-compact (car $f$ est de type fini
et $Y$ est quasi-compact), $W$ est quasi-compact.
D'après le lemme \ref{stacks-more-morphisms-lemma-finite-cover-factor},
on peut trouver un morphisme fini surjectif $g : Y' \to Y$
tel que $g^{-1}(W) \to Y$ se factorise localement pour la topologie de Zariski
par $\mathcal{X} \to Y$.
En remplaçant $Y$ par $Y'$ et $\mathcal{X}$ par
$\mathcal{X} \times_Y Y'$, on se ramène à la situation
décrite au paragraphe suivant.

\medskip\noindent
Supposons qu'il existe un entier $n \geq 0$, des ouverts quasi-compacts
$W_i \subset Y$, $i = 1, \ldots, n$, et des
morphismes $x_i : W_i \to \mathcal{X}$ tels que
(a) $f \circ x_i = \text{id}_{W_i}$,
(b) $W = \bigcup_{i = 1, \ldots, n} W_i$ contienne $V$, et
(c) $W$ soit l'image de $|f|$.
Nous raisonnons par récurrence sur $n$. Le cas initial est $n = 0$ : cela
implique $V = \emptyset$ et, dans ce cas, on peut prendre
$\overline{Z} = \emptyset$.
Si $n > 0$, considérons, pour $i = 1, \ldots, n$,
les sous-schémas fermés réduits
$Y_i$ dont l'espace topologique sous-jacent est
$Y \setminus W_i$. Considérons le morphisme fini
$$
Y' = Y \amalg \coprod\nolimits_{i = 1, \ldots, n} Y_i \longrightarrow Y
$$
et l'ouvert quasi-compact
$$
V' = (W_1 \cap \ldots \cap W_n \cap V) \amalg
\coprod_{i = 1, \ldots, n} (V \cap Y_i).
$$
D'après la remarque initiale de la démonstration, si l'on peut démontrer le lemme pour les paires
$$
(\mathcal{X} \to Y, W_1 \cap \ldots \cap W_n \cap V)
\quad\text{et}\quad
(\mathcal{X} \times_Y Y_i \to Y_i, V \cap Y_i),\quad
i = 1, \ldots, n
$$
alors le résultat est vrai. On utilise ici l'égalité ensembliste
$V = (W_1 \cap \ldots \cap W_n \cap V) \cup
\bigcup\nolimits_{i = 1, \ldots, n} (V \cap Y_i)$.
L'hypothèse de récurrence s'applique au second type de
paires ci-dessus. On se ramène donc à la situation décrite au
paragraphe suivant.

\medskip\noindent
Supposons qu'il existe un entier $n \geq 0$, des ouverts quasi-compacts
$W_i \subset Y$, $i = 1, \ldots, n$, et des
morphismes $x_i : W_i \to \mathcal{X}$ tels que
(a) $f \circ x_i = \text{id}_{W_i}$,
(b) $W = \bigcup_{i = 1, \ldots, n} W_i$ contienne $V$,
(c) $W$ soit l'image de $|f|$, et
(d) $V \subset W_1 \cap \ldots \cap W_n$.
Les morphismes
$$
T_{ij} = \mathit{Isom}_\mathcal{X}(x_i|_{W_i \cap W_j \cap V},
x_j|_{W_i \cap W_j \cap V}) \longrightarrow W_i \cap W_j \cap V
$$
sont surjectifs, plats et localement de présentation finie
(Morphismes de champs, lemme \ref{stacks-morphisms-lemma-gerbe-isom-fppf}).
Appliquons le lemme \ref{stacks-more-morphisms-lemma-finite-cover-factor}
à chacun des ouverts quasi-compacts $W_i \cap W_j \cap V$ et aux morphismes
$T_{ij} \to W_i \cap W_j \cap V$; on obtient des morphismes finis surjectifs
$Y'_{ij} \to Y$. En remplaçant $Y$ par le produit fibré de tous
les $Y'_{ij}$ au-dessus de $Y$, on se ramène à la situation décrite
au paragraphe suivant.

\medskip\noindent
Supposons qu'il existe un entier $n \geq 0$, des ouverts quasi-compacts
$W_i \subset Y$, $i = 1, \ldots, n$, et des
morphismes $x_i : W_i \to \mathcal{X}$ tels que
(a) $f \circ x_i = \text{id}_{W_i}$,
(b) $W = \bigcup_{i = 1, \ldots, n} W_i$ contienne $V$,
(c) $W$ soit l'image de $|f|$,
(d) $V \subset W_1 \cap \ldots \cap W_n$, et
(e) $x_i$ et $x_j$ soient localement isomorphes pour la topologie de Zariski au-dessus de $W_i \cap W_j \cap V$.
Soit $y \in V$ arbitraire.
Supposons que l'on puisse trouver un voisinage ouvert quasi-compact
$y \in V_y \subset V$ tel que le lemme soit vrai pour
la paire $(\mathcal{X} \to Y, V_y)$, avec, disons, pour solution
$\overline{Z}_y, Z_y, \overline{g}_y, g_y, h_y$.
Comme $V$ est quasi-compact, on peut trouver un nombre fini de points
$y_1, \ldots, y_m$ tels que $V = V_{y_1} \cup \ldots \cup V_{y_m}$.
Il s'ensuit qu'en posant
$$
\overline{Z} = \coprod \overline{Z}_{y_j},\quad
Z = \coprod Z_{y_j},\quad
\overline{g} = \coprod \overline{g}_{y_j},\quad
g = \coprod g_{y_j},\quad
h = \coprod h_{y_j}
$$
on obtient une solution du lemme. Étant donné $y$, la condition (e)
permet de choisir un voisinage ouvert quasi-compact $y \in V_y \subset V$
et des isomorphismes $\varphi_i : x_1|_{V_y} \to x_i|_{V_y}$ pour
$i = 2, \ldots, n$. Posons $\varphi_{ij} = \varphi_j \circ \varphi_i^{-1}$.
Cela nous conduit à la situation décrite au paragraphe suivant.

\medskip\noindent
Supposons qu'il existe un entier $n \geq 0$, des ouverts quasi-compacts
$W_i \subset Y$, $i = 1, \ldots, n$, et des
morphismes $x_i : W_i \to \mathcal{X}$ tels que
(a) $f \circ x_i = \text{id}_{W_i}$,
(b) $W = \bigcup_{i = 1, \ldots, n} W_i$ contienne $V$,
(c) $W$ soit l'image de $|f|$,
(d) $V \subset W_1 \cap \ldots \cap W_n$, et
(f) il existe des isomorphismes $\varphi_{ij} : x_i|_V \to x_j|_V$
vérifiant $\varphi_{jk} \circ \varphi_{ij} = \varphi_{ik}$.
Les morphismes
$$
I_{ij} = \mathit{Isom}_\mathcal{X}(x_i|_{W_i \cap W_j},
x_j|_{W_i \cap W_j}) \longrightarrow W_i \cap W_j
$$
sont propres parce que $f$ est séparé
(Morphismes de champs, lemme
\ref{stacks-morphisms-lemma-separated-implies-isom}).
Observons que $\varphi_{ij}$ définit une section $V \to I_{ij}$
de $I_{ij} \to W_i \cap W_j$ au-dessus de $V$.
D'après Compléments sur les morphismes d'espaces, lemme
\ref{spaces-more-morphisms-lemma-get-section-after-blowup}
on peut trouver des éclatements $V$-admissibles
$p_{ij} : Y_{ij} \to Y$ tels que $\varphi_{ij}$
se prolonge à $p_{ij}^{-1}(W_i \cap W_j)$.
En remplaçant $Y$ par le produit fibré de tous les $Y_{ij}$
au-dessus de $Y$, on arrive à la situation décrite au paragraphe suivant.

\medskip\noindent
Supposons qu'il existe un entier $n \geq 0$, des ouverts quasi-compacts
$W_i \subset Y$, $i = 1, \ldots, n$, et des
morphismes $x_i : W_i \to \mathcal{X}$ tels que
(a) $f \circ x_i = \text{id}_{W_i}$,
(b) $W = \bigcup_{i = 1, \ldots, n} W_i$ contienne $V$,
(c) $W$ soit l'image de $|f|$,
(d) $V \subset W_1 \cap \ldots \cap W_n$, et
(g) il existe des isomorphismes
$\varphi_{ij} : x_i|_{W_i \cap W_j} \to x_j|_{W_i \cap W_j}$
vérifiant
$$
\varphi_{jk}|_V \circ \varphi_{ij}|_V = \varphi_{ik}|_V.
$$
Après avoir remplacé $Y$, si nécessaire, par un autre éclatement $V$-admissible,
on peut supposer que $V$ est dense et schématiquement dense
dans $Y$, et donc dans tout sous-espace ouvert de $Y$ contenant $V$.
Après un tel remplacement, on conclut que
$$
\varphi_{jk}|_{W_i \cap W_j \cap W_k} \circ
\varphi_{ij}|_{W_i \cap W_j \cap W_k} =
\varphi_{ik}|_{W_i \cap W_j \cap W_k}
$$
en invoquant Morphismes d'espaces, lemme
\ref{spaces-morphisms-lemma-equality-of-morphisms}
et le fait que $I_{ik} \to W_i \cap W_k$ est propre
(donc séparé).
Bien entendu, cela signifie que $(x_i, \varphi_{ij})$
est une donnée de descente, et l'on obtient un morphisme
$x : W \to \mathcal{X}$ qui coïncide avec $x_i$ au-dessus de $W_i$,
car $\mathcal{X}$ est un champ.
Puisque $x$ est une section du morphisme séparé
$\mathcal{X} \to W$, $x$ est propre
(Morphismes de champs, lemme \ref{stacks-morphisms-lemma-section-immersion}).
Le lemme est donc vérifié avec $\overline{Z} = Y$,
$Z = W$, $\overline{g} = \text{id}_Y$, $g = \text{id}_W$,
$h = x$.
\end{proof}

\begin{theorem}[Lemme de Chow]
\label{stacks-more-morphisms-theorem-chow-finite-type}
\begin{reference}
Ce résultat est dû à Ofer Gabber; voir
\cite[Theorem 1.1]{olsson_proper}
\end{reference}
Soit $f : \mathcal{X} \to Y$ un morphisme d'un champ algébrique
vers un espace algébrique. Supposons que
\begin{enumerate}
\item $Y$ soit quasi-compact et quasi-séparé,
\item $f$ soit séparé de type fini.
\end{enumerate}
Alors il existe un diagramme commutatif
$$
\xymatrix{
\mathcal{X} \ar[rd] & X \ar[l] \ar[d] \ar[r] & \overline{X} \ar[ld] \\
& Y
}
$$
où $X \to \mathcal{X}$ est propre surjectif,
$X \to \overline{X}$ est une immersion ouverte, et
$\overline{X} \to Y$ est un morphisme propre d'espaces algébriques.
\end{theorem}

\begin{proof}
L'idée générale consiste à utiliser le fait que $\mathcal{X}$ possède un ouvert dense
qui est une gerbe (Morphismes de champs, proposition
\ref{stacks-morphisms-proposition-open-stratum})
et à invoquer le lemme \ref{stacks-more-morphisms-lemma-make-section}.
Cela ne fonctionne pas directement, car l'ouvert peut ne pas être
quasi-compact et l'on se heurte alors à des difficultés techniques. Nous effectuons donc
d'abord une réduction (standard) au cas noethérien.

\medskip\noindent
Choisissons d'abord une immersion fermée
$\mathcal{X} \to \mathcal{X}'$, où $\mathcal{X}'$
est un champ algébrique séparé et de type fini sur $Y$.
Voir Limites de champs, lemme
\ref{stacks-limits-lemma-separated-closed-in-finite-presentation}.
Il suffit manifestement de démontrer le théorème pour
$\mathcal{X}'$; on peut donc supposer que $\mathcal{X} \to Y$
est séparé et de présentation finie.

\medskip\noindent
Supposons que $\mathcal{X} \to Y$ soit séparé et de présentation finie.
D'après Limites d'espaces, proposition \ref{spaces-limits-proposition-approximate},
on peut écrire $Y = \lim Y_i$ comme la limite projective d'un système filtrant
d'espaces algébriques noethériens à morphismes de transition affines.
D'après Limites de champs, lemme \ref{stacks-limits-lemma-descend-a-stack-down},
il existe un $i$ et un morphisme $\mathcal{X}_i \to Y_i$ de présentation finie,
d'un champ algébrique vers $Y_i$, tels que
$\mathcal{X} = Y \times_{Y_i} \mathcal{X}_i$.
En augmentant $i$, on peut supposer que $\mathcal{X}_i \to Y_i$
est séparé; voir Limites de champs, lemme
\ref{stacks-limits-lemma-eventually-separated}.
Il suffit alors de démontrer le théorème
pour $\mathcal{X}_i \to Y_i$. Cela nous ramène au cas considéré
au paragraphe suivant.

\medskip\noindent
Supposons $Y$ noethérien. On peut remplacer $\mathcal{X}$ par son
réduit (Propriétés des champs, définition
\ref{stacks-properties-definition-reduced-induced-stack}).
Cela nous ramène au cas considéré
au paragraphe suivant.

\medskip\noindent
Supposons que $Y$ soit noethérien et que $\mathcal{X}$ soit réduit.
Puisque $\mathcal{X} \to Y$ est séparé et que $Y$ est quasi-séparé,
le champ algébrique $\mathcal{X}$ est quasi-séparé.
Ainsi, le morphisme d'inertie $\mathcal{I}_\mathcal{X} \to \mathcal{X}$
est quasi-compact. Par Morphismes de champs, proposition
\ref{stacks-morphisms-proposition-open-stratum}
il existe un sous-champ ouvert dense $\mathcal{V} \subset \mathcal{X}$
qui est une gerbe. Soit $\mathcal{V} \to V$ le morphisme
qui présente $\mathcal{V}$ comme une gerbe sur l'espace algébrique $V$.
Voir
Morphismes de champs, lemme \ref{stacks-morphisms-lemma-gerbe-over-iso-classes},
pour une construction de $\mathcal{V} \to V$.
Cette construction montre notamment que le morphisme
$\mathcal{V} \to Y$ se factorise en $\mathcal{V} \to V \to Y$.
On a le diagramme
$$
\xymatrix{
\mathcal{V} \ar[r] \ar[d] & \mathcal{X} \ar[d] \\
V \ar[r] & Y
}
$$
Puisque le morphisme $\mathcal{V} \to V$ est surjectif, plat et
de présentation finie
(Morphismes de champs, lemme \ref{stacks-morphisms-lemma-gerbe-fppf})
et que $\mathcal{V} \to Y$ est localement de présentation finie,
il s'ensuit que $V \to Y$ est localement de présentation finie
(Morphismes de champs, lemme
\ref{stacks-morphisms-lemma-flat-finite-presentation-permanence}).
Notons que $\mathcal{V} \to V$ est un homéomorphisme universel
(Morphismes de champs, lemme
\ref{stacks-morphisms-lemma-gerbe-bijection-points}).
Puisque $\mathcal{V}$ est quasi-compact (voir
Morphismes de champs, lemme
\ref{stacks-morphisms-lemma-locally-closed-in-noetherian}),
$V$ est quasi-compact.
Enfin, puisque $\mathcal{V} \to Y$ est séparé, il en va de même
de $V \to Y$ d'après
Morphismes de champs, lemme
\ref{stacks-morphisms-lemma-check-separated-on-ui-cover}
appliqué à $\mathcal{V} \to V \to Y$
(dont les hypothèses sont satisfaites, comme nous l'avons déjà vu).

\medskip\noindent
Tout ce qui précède signifie que les hypothèses de
Limites d'espaces, lemme \ref{spaces-limits-lemma-embedding-into-affine-over-qs},
s'appliquent au morphisme $V \to Y$. On peut donc trouver un sous-espace
ouvert dense $V' \subset V$ et une immersion $V' \to \mathbf{P}^n_Y$
sur $Y$. On peut manifestement remplacer $V$ par $V'$ et $\mathcal{V}$
par l'image inverse de $V'$ dans $\mathcal{V}$ (rappelons que
$|\mathcal{V}| = |V|$, comme nous l'avons vu plus haut).
On peut donc supposer donné un diagramme
$$
\xymatrix{
\mathcal{V} \ar[rr] \ar[d] & & \mathcal{X} \ar[d] \\
V \ar[r] & \mathbf{P}^n_Y \ar[r] & Y
}
$$
où la flèche $V \to \mathbf{P}^n_Y$ est une immersion.
Soit $\mathcal{X}'$ l'image schématique du morphisme
$$
j : \mathcal{V} \longrightarrow \mathbf{P}^n_Y \times_Y \mathcal{X}
$$
et soit $Y'$ l'image schématique du morphisme
$V \to \mathbf{P}^n_Y$. On obtient un diagramme commutatif
$$
\xymatrix{
\mathcal{V} \ar[r] \ar[d] &
\mathcal{X}' \ar[r] \ar[d] &
\mathbf{P}^n_Y \times_Y \mathcal{X} \ar[d] \ar[r] &
\mathcal{X} \ar[d] \\
V \ar[r] &
Y' \ar[r] &
\mathbf{P}^n_Y \ar[r] &
Y
}
$$
(Voir Morphismes de champs, lemme \ref{stacks-morphisms-lemma-factor-factor}.)
Nous affirmons que $\mathcal{V} = V \times_{Y'} \mathcal{X}'$ et que le
lemme \ref{stacks-more-morphisms-lemma-make-section} s'applique au morphisme $\mathcal{X}' \to Y'$
et au sous-espace ouvert $V \subset Y'$. Si l'affirmation est vraie, on obtient
$$
\xymatrix{
\overline{X} \ar[rd]_{\overline{g}} &
X \ar[l] \ar[d]_g \ar[r]_h & \mathcal{X}' \ar[ld]^f \\
& Y'
}
$$
où $X \to \overline{X}$ est une immersion ouverte, $\overline{g}$ et $h$ sont propres,
et $|V|$ est contenu dans l'image de $|g|$.
La composée $X \to \mathcal{X}' \to \mathcal{X}$ est alors
propre (comme composée de morphismes propres) et son image
contient $|\mathcal{V}|$; cette composée est donc surjective.
De même, $\overline{X} \to Y' \to Y$ est propre en tant que composée
de morphismes propres.

\medskip\noindent
La dernière étape consiste à démontrer l'affirmation.
Observons que $\mathcal{X}' \to Y'$ est séparé et de type fini,
que $Y'$ est quasi-compact et quasi-séparé, et que $V$ est quasi-compact
(nous omettons la vérification complète de tous les détails).
Observons ensuite que
$b : \mathcal{X}' \to \mathcal{X}$ est un isomorphisme au-dessus de
$\mathcal{V}$ d'après Morphismes de champs, lemme
\ref{stacks-morphisms-lemma-scheme-theoretic-image-of-partial-section}.
En particulier, $\mathcal{V}$ s'identifie à un sous-champ ouvert
de $\mathcal{X}'$.
Le morphisme $j$ est quasi-compact
(sa source est quasi-compacte et son but est quasi-séparé); la formation
de l'image schématique de $j$ commute donc au changement de base plat d'après
Morphismes de champs, lemme
\ref{stacks-morphisms-lemma-existence-plus-flat-base-change}.
En particulier, $V \times_{Y'} \mathcal{X}'$ est précisément
l'image schématique de
$\mathcal{V} \to V \times_{Y'} \mathcal{X}'$.
Or, d'après Morphismes de champs, lemme
\ref{stacks-morphisms-lemma-universally-closed-permanence}
l'image de $|\mathcal{V}| \to |V \times_{Y'} \mathcal{X}'|$
est fermée (on utilise que $\mathcal{V} \to V$ est un homéomorphisme universel,
comme nous l'avons vu plus haut, et qu'il est donc universellement fermé).
Cette image est aussi dense (combiner ce que nous venons de dire avec
Morphismes de champs, lemme
\ref{stacks-morphisms-lemma-topology-scheme-theoretic-image})
et l'on conclut que $|\mathcal{V}| = |V \times_{Y'} \mathcal{X}'|$.
Ainsi, $\mathcal{V} \to V \times_{Y'} \mathcal{X}'$ est un
isomorphisme, ce qui achève la démonstration de l'affirmation.
\end{proof}




\section{Critère valuatif noethérien}
\label{stacks-more-morphisms-section-Noetherian-valuative-criterion}

\noindent
Dans cette section, nous étudions des critères valuatifs (raffinés) pour les
morphismes de champs algébriques, en n'utilisant que des anneaux de valuation discrète
dans le cadre noethérien. Il en existe de nombreuses variantes,
et nous en ajouterons ici au fil du temps, selon les besoins.

\medskip\noindent
Soit $f : \mathcal{X} \to \mathcal{Y}$ un morphisme de champs algébriques
(ou d'espaces algébriques, ou encore de schémas). Un critère valuatif {\it raffiné}
est un critère où l'on se donne un morphisme $\mathcal{U} \to \mathcal{X}$
(possédant certaines propriétés) et où l'on ne considère que l'existence ou l'unicité
de flèches pointillées dans des diagrammes aux flèches pleines de la forme
$$
\xymatrix{
\Spec(K) \ar[d] \ar[r] & \mathcal{U} \ar[r] & \mathcal{X} \ar[d] \\
\Spec(A) \ar[rr] \ar@{..>}[rru] & & \mathcal{Y}
}
$$
Nous employons ci-dessous cette terminologie pour décrire les résultats obtenus jusqu'ici.

\medskip\noindent
Critères valuatifs non noethériens pour les morphismes de champs algébriques
\begin{enumerate}
\item Morphismes de champs, section
\ref{stacks-morphisms-section-valuative-second}
(pour la séparation de la diagonale),
\item Morphismes de champs, section
\ref{stacks-morphisms-section-valuative-diagonal}
(pour la séparation),
\item Morphismes de champs, section
\ref{stacks-morphisms-section-valutive-criterion}
(pour le caractère universellement fermé),
\item Morphismes de champs, section
\ref{stacks-morphisms-section-valutive-criterion-properness}
(pour la propreté).
\end{enumerate}
Pour les espaces algébriques, on dispose des critères valuatifs suivants
\begin{enumerate}
\item Morphismes d'espaces, section
\ref{spaces-morphisms-section-valuative-criterion-universally-closed}
(pour le caractère universellement fermé),
\item Morphismes d'espaces, lemme
\ref{spaces-morphisms-lemma-refined-valuative-criterion-universally-closed}
(forme raffinée pour le caractère universellement fermé),
\item Morphismes d'espaces, section
\ref{spaces-morphisms-section-valuative-separatedness}
(pour la séparation),
\item Morphismes d'espaces, section
\ref{spaces-morphisms-section-valuative-criterion-properness}
(pour la propreté),
\item Espaces décents, section
\ref{decent-spaces-section-valuative-criterion-universally-closed}
(pour le caractère universellement fermé des espaces décents),
\item Espaces décents, lemme
\ref{decent-spaces-lemma-re-characterize-universally-closed}
(pour le caractère universellement fermé des morphismes décents entre espaces algébriques),
\item Cohomologie des espaces, section
\ref{spaces-cohomology-section-Noetherian-valuative-criterion}
contient des critères valuatifs noethériens
\begin{enumerate}
\item Cohomologie des espaces, lemme
\ref{spaces-cohomology-lemma-check-separated-dvr}
(pour la séparation au moyen d'anneaux de valuation discrète),
\item Cohomologie des espaces, lemme
\ref{spaces-cohomology-lemma-check-proper-dvr}
(pour la propreté au moyen d'anneaux de valuation discrète),
\item Cohomologie des espaces, remarque
\ref{spaces-cohomology-remark-variant}
(explique comment se ramener à des anneaux de valuation discrète complets),
\end{enumerate}
\item Limites d'espaces, section
\ref{spaces-limits-section-Noetherian-valuative-criterion}
consacrée aux critères valuatifs noethériens
\begin{enumerate}
\item Limites d'espaces, lemme
\ref{spaces-limits-lemma-Noetherian-dvr-valuative-separation}
(pour la séparation au moyen d'anneaux de valuation discrète et de points génériques),
\item Limites d'espaces, lemme
\ref{spaces-limits-lemma-Noetherian-dvr-valuative-proper}
(pour la propreté au moyen d'anneaux de valuation discrète et de points génériques),
\item Limites d'espaces, lemme
\ref{spaces-limits-lemma-check-universally-closed-Noetherian}
(pour le caractère universellement fermé au moyen d'anneaux de valuation discrète).
\end{enumerate}
\item Limites d'espaces, section
\ref{spaces-limits-section-refined-valuative-criteria}
consacrée aux critères valuatifs noethériens raffinés
\begin{enumerate}
\item Limites d'espaces, lemmes
\ref{spaces-limits-lemma-refined-valuative-criterion-proper} et
\ref{spaces-limits-lemma-refined-valuative-criterion-universally-closed}
(forme raffinée pour la propreté au moyen d'anneaux de valuation discrète),
\item Limites d'espaces, lemme
\ref{spaces-limits-lemma-refined-valuative-criterion-separated}
(forme raffinée pour la séparation au moyen d'anneaux de valuation discrète),
\end{enumerate}
\end{enumerate}
Pour les schémas, on dispose des critères valuatifs suivants
\begin{enumerate}
\item Schémas, section
\ref{schemes-section-valuative-criterion-universal-closedness}
(pour le caractère universellement fermé),
\item Schémas, section \ref{schemes-section-valuative-separatedness}
(pour la séparation),
\item Morphismes, section \ref{morphisms-section-valuative-criteria}
(pour la propreté),
\item Morphismes, lemme
\ref{morphisms-lemma-refined-valuative-criterion-universally-closed}
(forme raffinée pour le caractère universellement fermé),
\item Limites, section \ref{limits-section-Noetherian-valuative-criterion}
consacrée aux critères valuatifs noethériens
\begin{enumerate}
\item Limites, lemme \ref{limits-lemma-Noetherian-dvr-valuative-separation}
(pour la séparation au moyen d'anneaux de valuation discrète et de points génériques),
\item Limites, lemme \ref{limits-lemma-Noetherian-dvr-valuative-proper}
(pour la propreté au moyen d'anneaux de valuation discrète et de points génériques),
\item Limites, lemme \ref{limits-lemma-check-universally-closed-Noetherian}
(pour le caractère universellement fermé au moyen d'anneaux de valuation discrète).
\end{enumerate}
\item Limites, section \ref{limits-section-refined-valuative-criteria}
consacrée aux critères valuatifs noethériens raffinés
\begin{enumerate}
\item Limites, lemmes
\ref{limits-lemma-refined-valuative-criterion-proper} et
\ref{limits-lemma-refined-valuative-criterion-universally-closed}
(forme raffinée pour la propreté au moyen d'anneaux de valuation discrète),
\item Limites, lemme \ref{limits-lemma-refined-valuative-criterion-separated}
(forme raffinée pour la séparation au moyen d'anneaux de valuation discrète),
\end{enumerate}
\item Limites, section \ref{limits-section-nagata-valuative}
consacrée aux critères valuatifs sur une base noethérienne, pour lesquels
on peut obtenir des anneaux de valuation discrète essentiellement de type fini
sur la base.
\end{enumerate}
Cela termine notre liste de résultats antérieurs.

\medskip\noindent
Nombre des résultats de cette section peuvent (et devraient peut-être)
être démontrés en invoquant le lemme suivant, même si nous ne l'avons pas
toujours fait.

\begin{lemma}
\label{stacks-more-morphisms-lemma-reach-point-closure-Noetherian}
Soit $f : \mathcal{X} \to \mathcal{Y}$ un morphisme de champs algébriques.
Supposons $f$ de type fini et $\mathcal{Y}$ localement noethérien.
Soit $y \in |\mathcal{Y}|$ un point de l'adhérence de l'image de $|f|$.
Alors il existe un diagramme commutatif
$$
\xymatrix{
\Spec(K) \ar[r] \ar[d] & \mathcal{X} \ar[d]^f \\
\Spec(A) \ar[r] & \mathcal{Y}
}
$$
de champs algébriques, où $A$ est un anneau de valuation discrète et $K$
son corps des fractions, le point fermé de $\Spec(A)$ étant envoyé sur $y$.
\end{lemma}

\begin{proof}
Choisissons un schéma affine $V$, un point $v \in V$ et un morphisme lisse
$V \to \mathcal{Y}$ qui envoie $v$ sur $y$. L'application $|V| \to |\mathcal{Y}|$
est ouverte et, d'après
Propriétés des champs, lemme \ref{stacks-properties-lemma-points-cartesian},
l'image de $|\mathcal{X} \times_\mathcal{Y} V| \to |V|$
est l'image inverse de l'image de $|f|$.
On en déduit que le point $v$ appartient à l'adhérence de
l'image de $|\mathcal{X} \times_\mathcal{Y} V| \to |V|$.
Si l'on démontre le lemme pour
$\mathcal{X} \times_\mathcal{Y} V \to V$ et le point $v$, alors
le lemme en résulte pour $f$ et $y$.
On se ramène ainsi à la situation décrite au
paragraphe suivant.

\medskip\noindent
Supposons donnés $f : \mathcal{X} \to Y$ et $y \in |Y|$ comme dans le lemme, où
$Y$ est un schéma affine noethérien. Puisque $f$ est quasi-compact, on en déduit que
$\mathcal{X}$ est quasi-compact. On peut donc choisir un schéma affine $W$ et
un morphisme lisse surjectif $W \to \mathcal{X}$. L'image de
$|f|$ est alors la même que celle de $|W| \to |Y|$. On se ramène ainsi
au cas des schémas, qui est
Limites, lemme \ref{limits-lemma-reach-point-closure-Noetherian}.
\end{proof}

\begin{lemma}
\label{stacks-more-morphisms-lemma-check-separated-dvr}
Soit $f : \mathcal{X} \to \mathcal{Y}$ un morphisme de champs algébriques. Supposons que
\begin{enumerate}
\item $\mathcal{Y}$ soit localement noethérien,
\item $f$ soit localement de type fini et quasi-séparé,
\item pour tout diagramme commutatif
$$
\xymatrix{
\Spec(K) \ar[r]_x \ar[d]_j & \mathcal{X} \ar[d]^f \\
\Spec(A) \ar[r]^y \ar@{-->}[ru] & \mathcal{Y}
}
$$
où $A$ est un anneau de valuation discrète, $K$ son corps des fractions,
et pour toute $2$-flèche $\gamma : y \circ j \to f \circ x$, la catégorie
des flèches pointillées (Morphismes de champs, définition
\ref{stacks-morphisms-definition-fill-in-diagram})
soit vide, soit un sétoïde ayant exactement une classe d'isomorphisme.
\end{enumerate}
Alors $f$ est séparé.
\end{lemma}

\begin{proof}
Pour démontrer que $f$ est séparé, il faut montrer que
$\Delta : \mathcal{X} \to \mathcal{X} \times_\mathcal{Y} \mathcal{X}$
est propre. On sait déjà que $\Delta$ est représentable par
des espaces algébriques, localement de type fini (Morphismes de champs,
lemme \ref{stacks-morphisms-lemma-properties-diagonal}) et
quasi-compact et quasi-séparé (par définition de la quasi-séparation de $f$).
Choisissons un schéma $U$ et un morphisme lisse surjectif
$U \to \mathcal{X} \times_\mathcal{Y} \mathcal{X}$.
Posons
$$
V = \mathcal{X} \times_{\Delta, \mathcal{X} \times_\mathcal{Y} \mathcal{X}} U
$$
Il suffit de montrer que le morphisme d'espaces algébriques $V \to U$
est propre (Propriétés des champs, lemme
\ref{stacks-properties-lemma-check-property-covering}).
Observons que $U$ est localement noethérien
(utiliser Morphismes de champs, lemme
\ref{stacks-morphisms-lemma-locally-finite-type-locally-noetherian}
et le fait que $U \to \mathcal{Y}$ est localement de type fini),
et que $V \to U$
est de type fini et quasi-séparé (comme changement de base
de $\Delta$, compte tenu des propriétés de $\Delta$ énumérées ci-dessus). D'après
Cohomologie des espaces, lemme \ref{spaces-cohomology-lemma-check-proper-dvr},
il suffit de montrer ceci : étant donné un diagramme commutatif
$$
\xymatrix{
\Spec(K) \ar[r]_v \ar[d]_j &
V \ar[d]^g \ar[r] &
\mathcal{X} \ar[d]^\Delta \\
\Spec(A) \ar[r]^u \ar@{-->}[ru] \ar@{..>}[rru] &
U \ar[r] &
\mathcal{X} \times_\mathcal{Y} \mathcal{X}
}
$$
où $A$ est un anneau de valuation discrète et $K$ son corps des fractions,
il existe une unique flèche en tirets qui rend le diagramme commutatif.
D'après Morphismes de champs, lemme
\ref{stacks-morphisms-lemma-cat-dotted-arrows-base-change}
les catégories des flèches en tirets et des flèches pointillées sont équivalentes.
L'hypothèse (3) entraîne l'existence d'une unique flèche pointillée à
isomorphisme près; voir Morphismes de champs, lemme
\ref{stacks-morphisms-lemma-helper-diagonal}. On en déduit
l'existence de l'unique flèche en tirets voulue.
\end{proof}

\begin{lemma}
\label{stacks-more-morphisms-lemma-refined-valuative-criterion-proper}
Soient $f : \mathcal{X} \to \mathcal{Y}$ et $h : \mathcal{U} \to \mathcal{X}$
des morphismes de champs algébriques. Supposons que $\mathcal{Y}$ soit
localement noethérien, que $f$ et $h$ soient de type fini,
que $f$ soit séparé, et que l'image de
$|h| : |\mathcal{U}| \to |\mathcal{X}|$ soit dense dans $|\mathcal{X}|$.
Si, pour tout diagramme $2$-commutatif
$$
\xymatrix{
\Spec(K) \ar[r]_-u \ar[d]_j & \mathcal{U} \ar[r]_h & \mathcal{X} \ar[d]^f \\
\Spec(A) \ar[rr]^-y & & \mathcal{Y}
}
$$
où $A$ est un anneau de valuation discrète de corps des fractions $K$
et $\gamma : y \circ j \to f \circ h \circ u$, il
existe une extension de corps $K'/K$ et un anneau de valuation $A' \subset K'$
dominant $A$, tels que la catégorie des flèches pointillées du
diagramme induit
$$
\xymatrix{
\Spec(K') \ar[r]_-{x'} \ar[d]_{j'} & \mathcal{X} \ar[d]^f \\
\Spec(A') \ar[r]^-{y'} \ar@{..>}[ru] & \mathcal{Y}
}
$$
muni de la $2$-flèche induite $\gamma' : y' \circ j' \to f \circ x'$
soit non vide (Morphismes de champs, définition
\ref{stacks-morphisms-definition-fill-in-diagram}), alors $f$ est propre.
\end{lemma}

\begin{proof}
Il suffit de démontrer que $f$ est universellement fermé.
Soit $V \to \mathcal{Y}$ un morphisme lisse, où $V$ est un schéma affine.
D'après Propriétés des champs, lemme
\ref{stacks-properties-lemma-points-cartesian}
l'image $I$ de
$|\mathcal{U} \times_\mathcal{Y} V| \to |\mathcal{X} \times_\mathcal{Y} V|$
est l'image inverse de l'image de $|h|$. Puisque
$|\mathcal{X} \times_\mathcal{Y} V| \to |\mathcal{X}|$ est ouverte
(Morphismes de champs, lemme \ref{stacks-morphisms-lemma-fppf-open}),
on en déduit que $I$ est dense dans $|\mathcal{X} \times_\mathcal{Y} V|$.
En outre, comme la catégorie des flèches pointillées se comporte bien par
changement de base (Morphismes de champs, lemme
\ref{stacks-morphisms-lemma-cat-dotted-arrows-base-change})
les morphismes suivants héritent de l'hypothèse d'existence de flèches pointillées
(après extension) :
$\mathcal{U} \times_\mathcal{Y} V \to \mathcal{X} \times_\mathcal{Y} V \to V$.
Les hypothèses du lemme sont donc
satisfaites pour les morphismes
$\mathcal{U} \times_\mathcal{Y} V \to \mathcal{X} \times_\mathcal{Y} V \to V$.
On peut par conséquent supposer que $\mathcal{Y}$ est un schéma affine.

\medskip\noindent
Supposons que $\mathcal{Y} = Y$ soit un schéma affine.
(Désormais, nous n'avons plus à nous soucier des
$2$-flèches $\gamma$ et $\gamma'$; voir Morphismes de champs, lemme
\ref{stacks-morphisms-lemma-cat-dotted-arrows-independent}.)
Alors $\mathcal{U}$ est quasi-compact. Choisissons un schéma affine $U$ et un
morphisme lisse surjectif $U \to \mathcal{U}$.
Nous remplaçons alors $\mathcal{U}$ par $U$.
On peut donc supposer que $\mathcal{U}$ est un schéma affine.

\medskip\noindent
Supposons que $\mathcal{Y} = Y$ et $\mathcal{U} = U$ soient des schémas affines.
D'après le lemme de Chow (théorème \ref{stacks-more-morphisms-theorem-chow-finite-type}),
on peut choisir un morphisme propre surjectif $X \to \mathcal{X}$,
où $X$ est un espace algébrique. Nous utiliserons plus bas que $X \to Y$ est séparé
comme composé de morphismes séparés. Considérons
l'espace algébrique $W = X \times_\mathcal{X} U$. Le morphisme de projection
$W \to X$ est de type fini.
On peut remplacer $X$ par l'image schématique de
$W \to X$ et supposer ainsi que l'image de $|W|$ dans $|X|$
est dense dans $|X|$ (on utilise ici que l'image de $|h|$
est dense dans $|\mathcal{X}|$, de sorte qu'après ce remplacement, le
morphisme $X \to \mathcal{X}$ reste surjectif).
Nous affirmons que, pour tout diagramme commutatif aux flèches pleines
$$
\xymatrix{
\Spec(K) \ar[r] \ar[d] & W \ar[r] & X \ar[d] \\
\Spec(A) \ar[rr] \ar@{..>}[rru] & & Y
}
$$
où $A$ est un anneau de valuation discrète de corps des fractions $K$, il
existe une flèche pointillée qui rend le diagramme commutatif. Tout d'abord, il suffit de
montrer qu'une telle flèche existe après avoir remplacé $K$ par une extension
et $A$ par un anneau de valuation de cette extension dominant $A$; voir
Morphismes d'espaces, lemme \ref{spaces-morphisms-lemma-push-down-solution}.
L'hypothèse du lemme fournit une extension $K'/K$, un anneau de
valuation $A' \subset K'$ dominant $A$, ainsi
qu'une flèche $\Spec(A') \to \mathcal{X}$ qui relève la composée
$\Spec(A') \to \Spec(A) \to Y$ et qui est compatible avec la composée
$\Spec(K') \to \Spec(K) \to W \to X$. Puisque $X \to \mathcal{X}$
est propre, le critère valuatif de propreté
(Morphismes de champs, lemme \ref{stacks-morphisms-lemma-criterion-proper})
donne une extension $K''/K'$, un anneau de valuation $A'' \subset K''$
dominant $A'$, et un morphisme $\Spec(A'') \to X$ qui relève la composée
$\Spec(A'') \to \Spec(A') \to \mathcal{X}$ et qui est compatible avec la composée
$\Spec(K'') \to \Spec(K') \to \Spec(K) \to X$.
L'extension $K''/K$, l'anneau $A'' \subset K''$ et le morphisme $\Spec(A'') \to X$
fournissent alors une solution au problème posé ci-dessus, ce qui démontre l'affirmation.

\medskip\noindent
L'affirmation entraîne que le morphisme $X \to Y$ est propre, par le
cas du lemme relatif aux espaces algébriques
(Limites d'espaces, lemme
\ref{spaces-limits-lemma-refined-valuative-criterion-proper}).
D'après Morphismes de champs, lemme
\ref{stacks-morphisms-lemma-image-proper-is-proper}
on en déduit que $\mathcal{X} \to Y$ est propre, comme souhaité.
\end{proof}

\begin{lemma}
\label{stacks-more-morphisms-lemma-refined-valuative-criterion-separated}
Soient $f : \mathcal{X} \to \mathcal{Y}$ et $h : \mathcal{U} \to \mathcal{X}$
des morphismes de champs algébriques. Supposons que $\mathcal{Y}$ soit
localement noethérien, que $f$ soit localement de type fini et quasi-séparé,
que $h$ soit de type fini, et que l'image de
$|h| : |\mathcal{U}| \to |\mathcal{X}|$ soit dense dans $|\mathcal{X}|$.
Si, pour tout diagramme $2$-commutatif
$$
\xymatrix{
\Spec(K) \ar[r]_-u \ar[d]_j & \mathcal{U} \ar[r]_h & \mathcal{X} \ar[d]^f \\
\Spec(A) \ar[rr]^-y \ar@{..>}[rru] & & \mathcal{Y}
}
$$
où $A$ est un anneau de valuation discrète de corps des fractions $K$
et $\gamma : y \circ j \to f \circ h \circ u$, la catégorie
des flèches pointillées est soit vide, soit un sétoïde ayant exactement
une classe d'isomorphisme, alors $f$ est séparé.
\end{lemma}

\begin{proof}
Il faut démontrer que $\Delta$ est un morphisme propre.
Supposons d'abord que $\Delta$ soit séparé. On peut alors appliquer le
lemme \ref{stacks-more-morphisms-lemma-refined-valuative-criterion-proper}
aux morphismes $\mathcal{U} \to \mathcal{X}$ et
$\Delta : \mathcal{X} \to \mathcal{X} \times_\mathcal{Y} \mathcal{X}$.
Observons que $\Delta$ est quasi-compact puisque $f$ est quasi-séparé.
Bien entendu, $\Delta$ est localement de type fini (cela vaut pour tout
morphisme diagonal; voir Morphismes de champs, lemme
\ref{stacks-morphisms-lemma-properties-diagonal}).
Enfin, supposons donné un diagramme $2$-commutatif
$$
\xymatrix{
\Spec(K) \ar[r]_-u \ar[d]_j &
\mathcal{U} \ar[r]_h &
\mathcal{X} \ar[d]^\Delta \\
\Spec(A) \ar[rr]^-y \ar@{..>}[rru] & &
\mathcal{X} \times_\mathcal{Y} \mathcal{X}
}
$$
où $A$ est un anneau de valuation discrète de corps des fractions $K$
et $\gamma : y \circ j \to \Delta \circ h \circ u$.
D'après Morphismes de champs, lemme \ref{stacks-morphisms-lemma-helper-diagonal},
et l'hypothèse du lemme,
il existe une unique flèche pointillée.
Cela établit la dernière hypothèse du
lemme \ref{stacks-more-morphisms-lemma-refined-valuative-criterion-proper},
et le résultat en découle.

\medskip\noindent
Dans le cas général, il suffit de démontrer que $\Delta$ est séparé,
car on sera alors ramené au cas précédent. Nous affirmons en fait
que les hypothèses du lemme sont satisfaites pour
$$
\mathcal{U} \to \mathcal{X}
\quad\text{et}\quad
\Delta :
\mathcal{X} \to
\mathcal{X} \times_\mathcal{Y} \mathcal{X}
$$
En effet, puisque $\Delta$ est représentable par des espaces algébriques, la
catégorie des flèches pointillées d'un diagramme comme au paragraphe précédent
est un sétoïde (voir par exemple
Morphismes de champs, lemme \ref{stacks-morphisms-lemma-cat-dotted-arrows}).
L'argument du paragraphe précédent montre que ces catégories
sont soit vides, soit munies d'une unique classe d'isomorphisme.
Ainsi, $\Delta$ est séparé.
\end{proof}

\begin{lemma}
\label{stacks-more-morphisms-lemma-valuative-criterion-universally-closed}
Soit $f : \mathcal{X} \to \mathcal{Y}$ un morphisme de champs algébriques.
Supposons que $\mathcal{Y}$ soit localement noethérien et que $f$ soit de type fini.
Si, pour tout diagramme $2$-commutatif
$$
\xymatrix{
\Spec(K) \ar[r]_-x \ar[d]_j & \mathcal{X} \ar[d]^f \\
\Spec(A) \ar[r]^-y & \mathcal{Y}
}
$$
où $A$ est un anneau de valuation discrète de corps des fractions $K$
et $\gamma : y \circ j \to f \circ x$, il existe une extension de corps $K'/K$
et un anneau de valuation $A' \subset K'$ dominant $A$, tels que
la catégorie des flèches pointillées du diagramme induit
$$
\xymatrix{
\Spec(K') \ar[r]_-{x'} \ar[d]_{j'} & \mathcal{X} \ar[d]^f \\
\Spec(A') \ar[r]^-{y'} \ar@{..>}[ru] & \mathcal{Y}
}
$$
muni de la $2$-flèche induite $\gamma' : y' \circ j' \to f \circ x'$
soit non vide (Morphismes de champs, définition
\ref{stacks-morphisms-definition-fill-in-diagram}), alors $f$ est
universellement fermé.
\end{lemma}

\begin{proof}
Soit $V \to \mathcal{Y}$ un morphisme lisse, où $V$ est un schéma affine.
La catégorie des flèches pointillées se comporte bien par changement de base
(Morphismes de champs, lemme
\ref{stacks-morphisms-lemma-cat-dotted-arrows-base-change}).
L'hypothèse d'existence de flèches pointillées (après extension)
est donc héritée par le morphisme $\mathcal{X} \times_\mathcal{Y} V \to V$.
Les hypothèses du lemme sont par conséquent satisfaites pour le morphisme
$\mathcal{X} \times_\mathcal{Y} V \to V$.
On peut donc supposer que $\mathcal{Y}$ est un schéma affine.

\medskip\noindent
Supposons que $\mathcal{Y} = Y$ soit un schéma affine noethérien.
(Désormais, nous n'avons plus à nous soucier des
$2$-flèches $\gamma$ et $\gamma'$; voir Morphismes de champs, lemme
\ref{stacks-morphisms-lemma-cat-dotted-arrows-independent}.)
Pour démontrer que $f$ est universellement fermé, il suffit de montrer que
$|\mathcal{X} \times \mathbf{A}^n| \to |Y \times \mathbf{A}^n|$ est fermée
pour tout $n$, d'après Limites de champs, lemme
\ref{stacks-limits-lemma-test-universally-closed}. Comme l'hypothèse
du lemme est héritée par le morphisme produit
$\mathcal{X} \times \mathbf{A}^n \to Y \times \mathbf{A}^n$ (nous omettons les détails),
on se ramène à démontrer que $|\mathcal{X}| \to |Y|$ est fermée.

\medskip\noindent
Supposons que $Y$ soit un schéma affine noethérien.
Soit $T \subset |\mathcal{X}|$ une partie fermée. Il faut montrer que
l'image de $T$ dans $|Y|$ est fermée. On peut remplacer $\mathcal{X}$
par la structure réduite induite de sous-champ fermé sur $T$; nous
omettons de vérifier que la propriété d'existence de
flèches pointillées est préservée par ce remplacement.
On se ramène ainsi à démontrer que l'image
de $|\mathcal{X}| \to |Y|$ est fermée.

\medskip\noindent
Soit $y \in |Y|$ un point de l'adhérence de l'image de
$|\mathcal{X}| \to |Y|$. D'après le lemme \ref{stacks-more-morphisms-lemma-reach-point-closure-Noetherian},
on peut choisir un diagramme commutatif
$$
\xymatrix{
\Spec(K) \ar[r] \ar[d] & \mathcal{X} \ar[d]^f \\
\Spec(A) \ar[r] & Y
}
$$
où $A$ est un anneau de valuation discrète et $K$ son corps des fractions,
le point fermé de $\Spec(A)$ étant envoyé sur $y$. Il résulte immédiatement
de l'hypothèse du lemme que $y$ appartient à l'image de
$|\mathcal{X}| \to |Y|$, ce qui achève la démonstration.
\end{proof}










\section{Espaces de modules}
\label{stacks-more-morphisms-section-moduli-spaces}

\noindent
Cette section porte sur les morphismes $f : \mathcal{X} \to Y$ de
champs algébriques vers des espaces algébriques. Sous des hypothèses convenables,
$Y$ est appelé un {\it espace de modules} de $\mathcal{X}$. Si
$\mathcal{X} = [U/R]$ est une présentation, on obtient alors un
morphisme $R$-invariant $U \to Y$ et, sous des hypothèses convenables,
$Y$ est un {\it quotient} du groupoïde $(U, R, s, t, c)$.
On trouvera une étude des différents types de quotients
à partir de
Quotients de groupoïdes, section \ref{groupoids-quotients-section-introduction}.

\begin{definition}
\label{stacks-more-morphisms-definition-categorical-quotient}
Soit $\mathcal{X}$ un champ algébrique. Soit
$f : \mathcal{X} \to Y$ un morphisme vers un espace algébrique $Y$.
\begin{enumerate}
\item On dit que $f$ est un {\it espace de modules catégorique} si tout morphisme
$\mathcal{X} \to W$ vers un espace algébrique $W$ se factorise de manière unique par $f$.
\item On dit que $f$ est un {\it espace de modules catégorique uniforme}
si, pour tout morphisme plat $Y' \to Y$ d'espaces algébriques, le changement de base
$f' : Y' \times_Y \mathcal{X} \to Y'$ est un espace de modules catégorique.
\end{enumerate}
Soit $\mathcal{C}$ une sous-catégorie pleine de la catégorie des espaces
algébriques.
\begin{enumerate}
\item[(3)] On dit que $f$ est un {\it espace de modules catégorique dans $\mathcal{C}$}
si $Y \in \Ob(\mathcal{C})$ et si tout morphisme $\mathcal{X} \to W$ avec
$W \in \Ob(\mathcal{C})$ se factorise de manière unique par $f$.
\item[(4)] On dit que $f$ est un {\it espace de modules catégorique uniforme dans $\mathcal{C}$}
si $Y \in \Ob(\mathcal{C})$ et si, pour tout morphisme plat $Y' \to Y$ dans
$\mathcal{C}$, le changement de base $f' : Y' \times_Y \mathcal{X} \to Y'$ est un
espace de modules catégorique dans $\mathcal{C}$.
\end{enumerate}
\end{definition}

\noindent
D'après le lemme de Yoneda, un espace de modules catégorique, s'il existe, est unique.
Rapprochons cette notion du langage introduit pour les quotients.

\begin{lemma}
\label{stacks-more-morphisms-lemma-quotient-compare}
Soit $(U, R, s, t, c)$ un groupoïde en espaces algébriques tel que
$s, t : R \to U$ soient plats et localement de présentation finie.
Considérons le champ algébrique $\mathcal{X} = [U/R]$.
Pour tout espace algébrique $Y$, il existe une correspondance biunivoque ($1$ à $1$) entre les
morphismes $f : \mathcal{X} \to Y$ et les morphismes $R$-invariants
$\phi : U \to Y$.
\end{lemma}

\begin{proof}
Critères de représentabilité, théorème
\ref{criteria-theorem-flat-groupoid-gives-algebraic-stack}
nous dit que $\mathcal{X}$ est un champ algébrique.
Étant donné un morphisme $f : \mathcal{X} \to Y$, soit $\phi : U \to Y$
la composée $U \to \mathcal{X} \to Y$. Puisque $R = U \times_\mathcal{X} U$
(Groupoïdes en espaces, lemme
\ref{spaces-groupoids-lemma-quotient-stack-2-cartesian})
il est immédiat que $\phi$ est $R$-invariant.
Réciproquement, si $\phi : U \to Y$ est un morphisme $R$-invariant vers
un espace algébrique, on obtient un morphisme
$f : \mathcal{X} \to Y$ d'après
Groupoïdes en espaces, lemme
\ref{spaces-groupoids-lemma-quotient-stack-2-coequalizer}.
On peut aussi construire $f$ à partir de $\phi$ en utilisant la description explicite du
champ quotient donnée dans
Groupoïdes en espaces, lemme
\ref{spaces-groupoids-lemma-quotient-stack-objects}.
\end{proof}

\begin{lemma}
\label{stacks-more-morphisms-lemma-categorical-quotient-compare}
Sous les hypothèses et avec les notations du lemme \ref{stacks-more-morphisms-lemma-quotient-compare},
$f$ est un espace de modules catégorique (uniforme)
si et seulement si $\phi$ est un quotient catégorique (uniforme).
Il en va de même des espaces de modules dans une sous-catégorie pleine.
\end{lemma}

\begin{proof}
La correspondance biunivoque ($1$ à $1$) établie dans le
lemme \ref{stacks-more-morphisms-lemma-quotient-compare} montre immédiatement que $f$ est un espace de modules catégorique
si et seulement si $\phi$ est un quotient catégorique
(Quotients de groupoïdes, définition
\ref{groupoids-quotients-definition-categorical}).
Si $Y' \to Y$ est un morphisme, alors
$U' = Y' \times_Y U \to Y' \times_Y \mathcal{X} = \mathcal{X}'$
est un morphisme surjectif, plat et localement de présentation finie
comme changement de base de $U \to \mathcal{X}$
(Critères de représentabilité, lemme
\ref{criteria-lemma-flat-quotient-flat-presentation}).
De plus, $R' = Y' \times_Y R$ est égal à $U' \times_{\mathcal{X}'} U'$
par transitivité des produits fibrés.
Ainsi, $\mathcal{X}' = [U'/R']$; voir
Champs algébriques, remarque \ref{algebraic-remark-flat-fp-presentation}.
Le changement de base de notre situation à $Y'$ est donc une nouvelle situation
du type décrit dans l'énoncé du lemme. Il en
résulte immédiatement que $f$ est un espace de modules catégorique uniforme
si et seulement si $\phi$ est un quotient catégorique uniforme.
\end{proof}

\begin{lemma}
\label{stacks-more-morphisms-lemma-check-uniform-categorical-quotient-on-affines}
Soit $f : \mathcal{X} \to Y$ un morphisme d'un champ algébrique
vers un espace algébrique. Si, pour tout schéma affine $Y'$ et tout
morphisme plat $Y' \to Y$, le changement de base
$f' : Y' \times_Y \mathcal{X} \to Y'$ est un espace de modules catégorique,
alors $f$ est un espace de modules catégorique uniforme.
\end{lemma}

\begin{proof}
Choisissons un recouvrement \'etale $\{Y_i \to Y\}$ où $Y_i$ est un schéma affine.
Pour chaque $i$ et $j$, choisissons un recouvrement par des ouverts affines
$Y_i \times_Y Y_j = \bigcup Y_{ijk}$.
Posons $\mathcal{X}_i = Y_i \times_Y \mathcal{X}$ et
$\mathcal{X}_{ijk} = Y_{ijk} \times_Y \mathcal{X}$.
Soit $g : \mathcal{X} \to W$ un morphisme vers
un espace algébrique. Considérons alors le diagramme
$$
\xymatrix{
\mathcal{X}_i \ar[r] \ar[d] & \mathcal{X} \ar[d] \ar[r]_g & W \\
Y_i \ar[r] \ar@{..>}[rru] & Y
}
$$
L'hypothèse que $\mathcal{X}_i \to Y_i$ est un espace de modules
catégorique fournit une unique flèche pointillée $h_i : Y_i \to W$.
L'hypothèse que $\mathcal{X}_{ijk} \to Y_{ijk}$ est un espace de modules
catégorique entraîne que les restrictions de $h_i$ et $h_j$ à
$Y_{ijk}$ sont égales. Ainsi, $h_i$ et $h_j$ coïncident sur $Y_i \times_Y Y_j$.
Puisque $Y = \coprod Y_i / \coprod Y_i \times_Y Y_j$
(d'après Espaces, section \ref{spaces-section-presentations}), on en déduit
qu'il existe un unique morphisme $Y \to W$ par lequel $g$ se factorise.
Ainsi, $f$ est un espace de modules catégorique. Le même argument s'applique
après un changement de base plat; $f$ est donc un espace de modules catégorique uniforme.
\end{proof}







\section{Le théorème de Keel--Mori}
\label{stacks-more-morphisms-section-Keel-Mori}

\noindent
Dans cette section, nous commençons l'étude du théorème de Keel et Mori
dans le cadre des champs algébriques. Pour une présentation de la
littérature, voir Guide bibliographique, sous-section
\ref{guide-subsection-coarse-moduli-spaces}.

\begin{definition}
\label{stacks-more-morphisms-definition-well-nigh-affine}
Soit $\mathcal{X}$ un champ algébrique. On dit que $\mathcal{X}$
est {\it presque affine} s'il existe un schéma affine $U$
et un morphisme surjectif, plat, fini et de présentation finie
$U \to \mathcal{X}$.
\end{definition}

\noindent
Nous donnons à cette propriété un nom quelque peu ridicule, car nous ne comptons pas
l'utiliser très souvent.

\begin{lemma}
\label{stacks-more-morphisms-lemma-well-nigh-affine}
Soit $\mathcal{X}$ un champ algébrique. Les assertions suivantes sont équivalentes :
\begin{enumerate}
\item $\mathcal{X}$ est presque affine, et
\item il existe un schéma en groupoïdes $(U, R, s, t, c)$ avec $U$ et
$R$ affines et $s, t : R \to U$ finis localement libres, tel que
$\mathcal{X} = [U/R]$.
\end{enumerate}
Lorsque ces assertions sont vraies, $\mathcal{X}$ est quasi-compact, quasi-DM et séparé.
\end{lemma}

\begin{proof}
Supposons $\mathcal{X}$ presque affine. Choisissons un schéma affine $U$
et un morphisme surjectif, plat, fini et de présentation finie
$U \to \mathcal{X}$. Posons $R = U \times_\mathcal{X} U$. On
obtient alors un groupoïde $(U, R, s, t, c)$ en espaces algébriques et un
isomorphisme $[U/R] \to \mathcal{X}$; voir
Champs algébriques, lemme \ref{algebraic-lemma-map-space-into-stack}
et remarque \ref{algebraic-remark-flat-fp-presentation}.
Puisque $s, t : R \to U$ sont des morphismes
plats, finis et de présentation finie
(comme changements de base de $U \to \mathcal{X})$, on voit que
$s, t$ sont finis localement libres
(Morphismes, lemme \ref{morphisms-lemma-finite-flat}).
Il en résulte que $R$ est affine (car les morphismes finis sont affines)
et donc que (2) est vérifiée.

\medskip\noindent
Supposons donné un schéma en groupoïdes $(U, R, s, t, c)$ avec $U$ et
$R$ affines et $s, t : R \to U$ finis localement libres.
Posons $\mathcal{X} = [U/R]$. Alors $\mathcal{X}$ est un champ algébrique
d'après Critères de représentabilité, théorème
\ref{criteria-theorem-flat-groupoid-gives-algebraic-stack} (à strictement parler,
nous n'en avons pas besoin ici, mais on ne saurait trop insister sur ce fait).
Le morphisme $U \to \mathcal{X}$ est surjectif, plat et localement
de présentation finie d'après
Critères de représentabilité, lemme
\ref{criteria-lemma-flat-quotient-flat-presentation}.
On peut donc vérifier que $U \to \mathcal{X}$ est fini en
vérifiant que la projection $U \times_\mathcal{X} U \to U$
possède cette propriété; voir Propriétés des champs, lemme
\ref{stacks-properties-lemma-check-property-covering}.
Puisque $U \times_\mathcal{X} U = R$ d'après
Groupoïdes en espaces, lemme
\ref{spaces-groupoids-lemma-quotient-stack-2-cartesian}
on voit que tel est le cas. Ainsi, $\mathcal{X}$ est presque affine.

\medskip\noindent
Démontrons l'assertion finale. On voit que $\mathcal{X}$
est quasi-compact d'après Propriétés des champs, lemme
\ref{stacks-properties-lemma-quasi-compact-stack}.
On voit que $\mathcal{X} = [U/R]$ est quasi-DM et séparé d'après
Morphismes de champs, lemme
\ref{stacks-morphisms-lemma-properties-diagonal-from-presentation}.
\end{proof}

\begin{lemma}
\label{stacks-more-morphisms-lemma-affine-over-well-nigh-affine}
Soit $\mathcal{X}$ un champ algébrique presque affine.
\begin{enumerate}
\item Si $\mathcal{X}$ est un espace algébrique, alors il est affine.
\item Si $\mathcal{X}' \to \mathcal{X}$ est un morphisme affine
de champs algébriques, alors $\mathcal{X}'$ est presque affine.
\end{enumerate}
\end{lemma}

\begin{proof}
La partie (1) résulte immédiatement de
Limites d'espaces, lemme \ref{spaces-limits-lemma-affine}.
Ce résultat est toutefois bien plus fort qu'il n'est nécessaire, car (1) résulte aussi du
lemme \ref{stacks-more-morphisms-lemma-well-nigh-affine} combiné avec
Groupoïdes, proposition
\ref{groupoids-proposition-finite-flat-equivalence}.

\medskip\noindent
Pour démontrer (2), choisissons un schéma affine $U$ et un
morphisme surjectif, plat, fini et de présentation finie $U \to \mathcal{X}$.
Alors $U' = \mathcal{X}' \times_\mathcal{X} U$ admet un morphisme
affine vers $U$ (Morphismes de champs, lemme
\ref{stacks-morphisms-lemma-base-change-affine}).
Ainsi, $U'$ est un schéma affine. Bien entendu,
$U' \to \mathcal{X}'$ est surjectif, plat, fini et de présentation finie
comme changement de base de $U \to \mathcal{X}$.
\end{proof}

\begin{lemma}
\label{stacks-more-morphisms-lemma-well-nigh-affine-moduli-space}
Soit $\mathcal{X}$ un champ algébrique presque affine. Il existe
un espace de modules catégorique uniforme
$$
f : \mathcal{X} \longrightarrow M
$$
dans la catégorie des schémas affines. De plus,
$f$ est séparé, quasi-compact et un homéomorphisme universel.
\end{lemma}

\begin{proof}
Écrivons $\mathcal{X} = [U/R]$, avec $(U, R, s, t, c)$ comme dans le
lemme \ref{stacks-more-morphisms-lemma-well-nigh-affine}. Soit $C$ l'anneau des
fonctions $R$-invariantes sur $U$; voir
Groupoïdes, section \ref{groupoids-section-finite-flat}.
Posons $M = \Spec(C)$. Le morphisme $R$-invariant
$U \to M$ correspond à un morphisme $f : \mathcal{X} \to M$ d'après le
lemme \ref{stacks-more-morphisms-lemma-quotient-compare}.
La caractérisation des morphismes vers les schémas affines donnée dans
Schémas, lemme \ref{schemes-lemma-morphism-into-affine},
garantit immédiatement que $\phi : U \to M$ est un quotient
catégorique dans la catégorie des schémas affines. Ainsi, $f$ est un
espace de modules catégorique dans la catégorie des schémas affines
(lemme \ref{stacks-more-morphisms-lemma-categorical-quotient-compare}).

\medskip\noindent
Puisque $\mathcal{X}$ est séparé d'après le lemme \ref{stacks-more-morphisms-lemma-well-nigh-affine},
on voit que $f$ est séparé d'après Morphismes de champs, lemme
\ref{stacks-morphisms-lemma-compose-after-separated}.

\medskip\noindent
Puisque $U \to \mathcal{X}$ est surjectif et que $U \to M$ est quasi-compact,
on voit que $f$ est quasi-compact d'après Morphismes de champs, lemme
\ref{stacks-morphisms-lemma-surjection-from-quasi-compact}.

\medskip\noindent
D'après Groupoïdes, lemme \ref{groupoids-lemma-integral-over-invariants},
la composée
$$
U \to \mathcal{X} \to M
$$
est un morphisme entier de schémas affines. En particulier, elle est
universellement fermée
(Morphismes, lemme \ref{morphisms-lemma-integral-universally-closed}).
Puisque $U \to \mathcal{X}$ est surjectif, il s'ensuit que $\mathcal{X} \to M$
est universellement fermé (Morphismes de champs, lemme
\ref{stacks-morphisms-lemma-image-proper-is-proper}).
Pour conclure que $\mathcal{X} \to M$ est un homéomorphisme universel,
il suffit de montrer qu'il est universellement bijectif, c'est-à-dire
surjectif et universellement injectif.

\medskip\noindent
On a $|\mathcal{X}| = |U|/|R|$ d'après
Morphismes de champs, lemme \ref{stacks-morphisms-lemma-points-presentation}.
Ainsi, $|f|$ est surjectif, et même bijectif,
d'après Groupoïdes, lemme \ref{groupoids-lemma-points}.

\medskip\noindent
Soit $C \to C'$ un homomorphisme d'anneaux. Soit $(U', R', s', t', c')$
le changement de base de $(U, R, s, t, c)$ par $M' = \Spec(C') \to M$.
En posant $\mathcal{X}' = [U'/R']$, on observe que
$M' \times_M \mathcal{X} = \mathcal{X}'$ d'après
Quotients de groupoïdes, lemme
\ref{groupoids-quotients-lemma-base-change-quotient-stack}.
Soit $C^1$ l'anneau des fonctions $R'$-invariantes sur $U'$.
Posons $M^1 = \Spec(C^1)$ et considérons le diagramme
$$
\xymatrix{
\mathcal{X}' \ar[d]^{f'} \ar[r] & \mathcal{X} \ar[dd]^f \\
M^1 \ar[d] \\
M' \ar[r] & M
}
$$
D'après Groupoïdes, lemme \ref{groupoids-lemma-invariants-base-change}, et
Algèbre, lemme \ref{algebra-lemma-universally-bijective},
le morphisme $M^1 \to M'$ est un homéomorphisme.
D'autre part, le paragraphe précédent appliqué à
$(U', R', s', t', c')$ montre que $|f'|$ est bijectif.
On en déduit que $f$ induit une bijection sur les points après tout
changement de base par un schéma affine. Ainsi, $f$ est universellement injectif
d'après Morphismes de champs, lemme
\ref{stacks-morphisms-lemma-universally-injective-local}.

\medskip\noindent
Enfin, il reste à montrer que $f$ est un espace de modules catégorique uniforme
dans la catégorie des schémas affines. Cela résulte de la discussion
précédente et du fait que, si
l'homomorphisme d'anneaux $C \to C'$ est plat, alors $C' \to C^1$ est un isomorphisme
d'après Groupoïdes, lemme \ref{groupoids-lemma-invariants-base-change}.
\end{proof}

\begin{lemma}
\label{stacks-more-morphisms-lemma-well-nigh-affine-moduli-space-etale}
Soit $h : \mathcal{X}' \to \mathcal{X}$ un morphisme de champs algébriques.
Supposons que $\mathcal{X}'$ et $\mathcal{X}$ soient presque affines,
que $h$ est \'etale et que $h$ induit des isomorphismes entre les groupes d'automorphismes
(Morphismes de champs, remarque
\ref{stacks-morphisms-remark-identify-automorphism-groups}).
Il existe alors un diagramme cartésien
$$
\xymatrix{
\mathcal{X}' \ar[d] \ar[r] & \mathcal{X} \ar[d] \\
M' \ar[r] & M
}
$$
où $M' \to M$ est \'etale et
les flèches verticales sont les espaces de modules construits dans le
lemme \ref{stacks-more-morphisms-lemma-well-nigh-affine-moduli-space}.
\end{lemma}

\begin{proof}
Remarquons que $h$ est représentable par des espaces algébriques, d'après
Morphismes de champs, lemmes
\ref{stacks-morphisms-lemma-aut-iso-unramified} et
\ref{stacks-morphisms-lemma-stabilizer-preserving}.
Choisissons un schéma affine $U$ et un
morphisme $U \to \mathcal{X}$ surjectif, plat, fini et de présentation finie.
Alors $U' = \mathcal{X}' \times_\mathcal{X} U$ est un espace
algébrique muni d'un morphisme fini (en particulier affine) $U' \to \mathcal{X}'$.
D'après le lemme \ref{stacks-more-morphisms-lemma-affine-over-well-nigh-affine},
nous en déduisons que $U'$ est affine.
En posant $R = U \times_\mathcal{X} U$ et $R' = U' \times_{\mathcal{X}'} U'$,
nous obtenons des groupoïdes $(U, R, s, t, c)$ et $(U', R', s', t', c')$
tels que $\mathcal{X} = [U/R]$ et $\mathcal{X}' = [U'/R']$ ;
voir la démonstration du lemme \ref{stacks-more-morphisms-lemma-well-nigh-affine}.
Nous voyons que les diagrammes
$$
\xymatrix{
R' \ar[d]_{s'} \ar[r]_f & R \ar[d]^s \\
U' \ar[r]^f & U
}
\quad
\quad
\xymatrix{
R' \ar[d]_{t'} \ar[r]_f & R \ar[d]^t \\
U' \ar[r]^f & U
}
\quad
\quad
\xymatrix{
G' \ar[d] \ar[r]_f & G \ar[d] \\
U' \ar[r]^f & U
}
$$
sont cartésiens, où $G$ et $G'$ sont les schémas en groupes stabilisateurs.
Pour les deux premiers, cela résulte de la transitivité des produits fibrés,
et, pour le dernier, du fait qu'il est le changement de base de
l'isomorphisme $\mathcal{I}_{\mathcal{X}'} \to
\mathcal{X}' \times_\mathcal{X} \mathcal{I}_\mathcal{X}$
(d'après le lemme déjà utilisé de Morphismes de champs
\ref{stacks-morphisms-lemma-aut-iso-unramified}).
Rappelons que $M$, resp.\ $M'$, a été construit dans le
lemme \ref{stacks-more-morphisms-lemma-well-nigh-affine-moduli-space}
comme le spectre de l'anneau des fonctions $R$-invariantes sur $U$,
resp.\ de l'anneau des fonctions $R'$-invariantes sur $U'$.
Ainsi, $M' \to M$ est \'etale et $U' = M' \times_M U$
d'après Groupoïdes, lemme \ref{groupoids-lemma-etale}.
Il s'ensuit que $R' = M' \times_M R$ ; autrement dit,
le groupoïde $(U', R', s', t', c')$ est le changement de base de
$(U, R, s, t, c)$ par $M' \to M$.
Cela implique que le diagramme de l'énoncé est
cartésien d'après
Quotients de groupoïdes, lemme
\ref{groupoids-quotients-lemma-base-change-quotient-stack}.
\end{proof}

\begin{lemma}
\label{stacks-more-morphisms-lemma-moduli-space-finite-affine}
Soit $\mathcal{X}$ un champ algébrique presque affine. Le morphisme
$$
f : \mathcal{X} \longrightarrow M
$$
du lemme \ref{stacks-more-morphisms-lemma-well-nigh-affine-moduli-space}
est un espace de modules catégorique uniforme.
\end{lemma}

\begin{proof}
Nous savons déjà que $M$ est un espace de modules catégorique uniforme
dans la catégorie des schémas affines. D'après le
lemme \ref{stacks-more-morphisms-lemma-check-uniform-categorical-quotient-on-affines},
il suffit de montrer que le changement de base
$f' : M' \times_M \mathcal{X} \to M'$
est un espace de modules catégorique pour tout morphisme plat
$M' \to M$ de schémas affines.
Notons que $\mathcal{X}' = M' \times_M \mathcal{X}$ est presque affine d'après le
lemme \ref{stacks-more-morphisms-lemma-affine-over-well-nigh-affine}.
Ainsi, en remplaçant $\mathcal{X}$ par $\mathcal{X}'$
et $M$ par $M'$, nous nous ramenons à prouver que $f$ est un espace de modules
catégorique.

\medskip\noindent
Soit $g : \mathcal{X} \to Y$ un morphisme, où $Y$ est un espace algébrique.
Nous devons montrer que $g = h \circ f$ pour un unique morphisme $h : M \to Y$.

\medskip\noindent
Unicité. Supposons donnés deux morphismes $h_i : M \to Y$ tels que
$g = h_1 \circ f = h_2 \circ f$. Soit $M' \subset M$ l'égalisateur
de $h_1$ et $h_2$. Alors $M' \to M$ est un monomorphisme et
$f : \mathcal{X} \to M$ se factorise par $M'$. Ainsi, $M' \to M$
est un homéomorphisme universel. Nous en concluons que $M'$ est affine
(Morphismes, lemme \ref{morphisms-lemma-universal-homeomorphism}).
Or, puisque $f : \mathcal{X} \to M$
est un espace de modules catégorique dans la catégorie des schémas affines,
nous avons $M' = M$.

\medskip\noindent
Existence. Nous allons montrer ci-dessous que, pour tout $p \in M$, il existe
un carré cartésien
$$
\xymatrix{
\mathcal{X}' \ar[r] \ar[d] & \mathcal{X} \ar[d] \\
M' \ar[r] & M
}
$$
où $M' \to M$ est un morphisme \'etale de schémas affines dont l'image contient $p$, tel que
la composée $\mathcal{X}' \to \mathcal{X} \to Y$ se factorise par $M'$.
Cela signifie que nous pouvons construire le morphisme $h : M \to Y$ localement pour la topologie \'etale sur $M$.
Puisque $Y$ est un faisceau pour la topologie \'etale et en vertu de l'unicité établie
ci-dessus, cela suffit (nous omettons un petit détail).

\medskip\noindent
Soit $y \in |Y|$ l'image de $p$.
Soit $(V, v) \to (Y, y)$ un morphisme \'etale avec $V$ affine.
Considérons $\mathcal{X}' = V \times_Y \mathcal{X}$.
Remarquons que $\mathcal{X}' \to \mathcal{X}$ est séparé et \'etale
comme changement de base de $V \to Y$. De plus, $\mathcal{X}' \to \mathcal{X}$
induit des isomorphismes entre les groupes d'automorphismes
(Morphismes de champs, remarque
\ref{stacks-morphisms-remark-identify-automorphism-groups})
car cela vaut pour
$V \to Y$ ; voir Morphismes de champs, lemme
\ref{stacks-morphisms-lemma-base-change-stabilizer-preserving}.
Choisissons une présentation $\mathcal{X} = [U/R]$
comme dans le lemme \ref{stacks-more-morphisms-lemma-well-nigh-affine}.
Posons $U' = \mathcal{X}' \times_\mathcal{X} U = V \times_Y U$
et choisissons $u' \in U'$ se projetant sur $p$ et $v$ (ce qui est possible
d'après Propriétés des espaces, lemme \ref{spaces-properties-lemma-points-cartesian}).
Puisque $U' \to U$ est séparé et \'etale, nous voyons que
tout ensemble fini de points de $U'$ est contenu dans un ouvert affine ; voir
Compléments sur les morphismes, lemme
\ref{more-morphisms-lemma-separated-locally-quasi-finite-over-affine}.
D'autre part, le morphisme $U' \to \mathcal{X}'$ est
surjectif, fini, plat et localement de présentation finie.
En posant $R' = U' \times_{\mathcal{X}'} U'$, nous voyons
que $s', t' : R' \to U'$ sont finis localement libres.
D'après Groupoïdes, lemme \ref{groupoids-lemma-find-invariant-affine},
il existe un sous-schéma ouvert affine invariant par $R'$, noté $U'' \subset U'$,
contenant $u'$.
Soit $\mathcal{X}'' \subset \mathcal{X}'$ le
sous-champ ouvert correspondant. Alors $\mathcal{X}''$ est
presque affine. D'après le lemme \ref{stacks-more-morphisms-lemma-well-nigh-affine-moduli-space-etale},
nous obtenons un carré cartésien
$$
\xymatrix{
\mathcal{X}'' \ar[r] \ar[d] & \mathcal{X} \ar[d] \\
M'' \ar[r] & M
}
$$
où $M'' \to M$ est \'etale. Puisque $\mathcal{X}'' \to M''$ est
un espace de modules catégorique dans la catégorie des schémas affines,
nous obtenons un morphisme $M'' \to V$ tel que la composée
$\mathcal{X}'' \to \mathcal{X}' \to V$ est égale à la composée
$\mathcal{X}'' \to M'' \to V$. Cela prouve notre assertion et achève
la démonstration.
\end{proof}

\begin{lemma}
\label{stacks-more-morphisms-lemma-etale-separated-over-well-nigh-affine}
Soit $h : \mathcal{X}' \to \mathcal{X}$ un morphisme de champs algébriques.
Supposons que $\mathcal{X}$ soit presque affine, que $h$ soit \'etale, que $h$ soit séparé
et que $h$ induit des isomorphismes entre les groupes d'automorphismes
(Morphismes de champs, remarque
\ref{stacks-morphisms-remark-identify-automorphism-groups}).
Il existe alors un diagramme cartésien
$$
\xymatrix{
\mathcal{X}' \ar[d] \ar[r] & \mathcal{X} \ar[d] \\
M' \ar[r] & M
}
$$
où $M' \to M$ est un morphisme de schémas séparé et \'etale, et où
$\mathcal{X} \to M$ est l'espace de modules construit dans le
lemme \ref{stacks-more-morphisms-lemma-well-nigh-affine-moduli-space}.
\end{lemma}

\begin{proof}
Choisissons un schéma affine $U$ et un morphisme surjectif, plat, fini et
localement de présentation finie $U \to \mathcal{X}$.
Puisque $h$ est représentable par des espaces algébriques
(Morphismes de champs, lemmes
\ref{stacks-morphisms-lemma-aut-iso-unramified} et
\ref{stacks-morphisms-lemma-stabilizer-preserving})
nous voyons que $U' = \mathcal{X}' \times_\mathcal{X} U$ est
un espace algébrique. Puisque $U' \to U$ est séparé et \'etale,
nous voyons que $U'$ est un schéma et que tout ensemble fini de points
de $U'$ est contenu dans un ouvert affine ; voir
Morphismes d'espaces, lemme
\ref{spaces-morphisms-lemma-locally-quasi-finite-separated-representable}
et
Compléments sur les morphismes, lemme
\ref{more-morphisms-lemma-separated-locally-quasi-finite-over-affine}.
En posant $R' = U' \times_{\mathcal{X}'} U'$, nous voyons
que $s', t' : R' \to U'$ sont finis localement libres.
D'après Groupoïdes, lemme \ref{groupoids-lemma-find-invariant-affine},
il existe un recouvrement ouvert $U' = \bigcup U'_i$ par
des sous-schémas ouverts affines invariants par $R'$, $U'_i \subset U'$.
Soient $\mathcal{X}'_i \subset \mathcal{X}'$ les
sous-champs ouverts correspondants. Ils sont presque affines, car $U'_i \to \mathcal{X}'_i$
est surjectif, plat, fini et de présentation finie. D'après le
lemme \ref{stacks-more-morphisms-lemma-well-nigh-affine-moduli-space-etale},
nous obtenons des diagrammes cartésiens
$$
\xymatrix{
\mathcal{X}'_i \ar[r] \ar[d] & \mathcal{X} \ar[d] \\
M'_i \ar[r] & M
}
$$
où $M'_i \to M$ est un morphisme \'etale de schémas affines,
les flèches verticales étant comme dans le
lemme \ref{stacks-more-morphisms-lemma-well-nigh-affine-moduli-space}.
Remarquons que
$\mathcal{X}'_{ij} = \mathcal{X}'_i \cap \mathcal{X}'_j$
est un sous-champ ouvert de $\mathcal{X}'_i$ et de $\mathcal{X}'_j$.
Nous obtenons donc les sous-schémas ouverts correspondants
$V_{ij} \subset M'_i$ et $V_{ji} \subset M'_j$.
En vertu du
lemme \ref{stacks-more-morphisms-lemma-moduli-space-finite-affine},
nous voyons que
$\mathcal{X}'_{ij} \to V_{ij}$ et
$\mathcal{X}'_{ji} \to V_{ji}$ sont tous deux des espaces de modules catégoriques !
Nous obtenons ainsi un unique isomorphisme $\varphi_{ij} : V_{ij} \to V_{ji}$
tel que le diagramme
$$
\xymatrix{
\mathcal{X}'_i \ar[d] & &
\mathcal{X}'_i \cap \mathcal{X}'_j \ar[rr] \ar[ll] \ar[ld] \ar[rd] & &
\mathcal{X}'_j \ar[d] \\
M'_i &
V_{ij} \ar[l] \ar[rr]^{\varphi_{ij}} & &
V_{ji} \ar[r] &
M'_j
}
$$
est commutatif. Ces isomorphismes satisfont à la condition de cocycle de
Schémas, section \ref{schemes-section-glueing-schemes}, comme le montre un calcul
(et une nouvelle application du lemme précédent) que nous omettons.
Nous pouvons donc recoller les schémas affines pour former un schéma $M'$ ; voir
Schémas, lemme \ref{schemes-lemma-glue}.
Identifions les $M'_i$ à leur image dans $M'$.
Nous affirmons qu'il existe un morphisme $\mathcal{X}' \to M'$ s'inscrivant dans des
diagrammes cartésiens
$$
\xymatrix{
\mathcal{X}'_i \ar[r] \ar[d] & \mathcal{X}' \ar[d] \\
M'_i \ar[r] & M'
}
$$
Cela ressort de la description des morphismes vers le schéma recollé $M'$
dans Schémas, lemme \ref{schemes-lemma-glue}, et du fait que donner un morphisme
$\mathcal{X}' \to M'$ revient à donner un morphisme $T \to M'$
pour tout morphisme $T \to \mathcal{X}'$.
De même, il existe un morphisme $M' \to M$ dont les restrictions sont les
morphismes donnés $M'_i \to M$ sur $M'_i$.
Le morphisme $M' \to M$ est \'etale (car il l'est sur les membres d'un
recouvrement \'etale), et la propriété de produit fibré vaut puisqu'elle peut
se vérifier sur les membres du recouvrement ouvert (affine) $M' = \bigcup M'_i$.
Enfin, $M' \to M$ est séparé, car la composée
$U' \to \mathcal{X}' \to M'$ est surjective et universellement fermée,
et nous pouvons appliquer Morphismes, lemme
\ref{morphisms-lemma-image-universally-closed-separated}.
\end{proof}

\begin{lemma}
\label{stacks-more-morphisms-lemma-etale-local-finite-inertia}
Soit $\mathcal{X}$ un champ algébrique. Supposons que
$\mathcal{I}_\mathcal{X} \to \mathcal{X}$ est fini.
Il existe alors un ensemble $I$ et, pour $i \in I$, un morphisme de champs algébriques
$$
g_i : \mathcal{X}_i \longrightarrow \mathcal{X}
$$
ayant les propriétés suivantes :
\begin{enumerate}
\item $|\mathcal{X}| = \bigcup |g_i|(|\mathcal{X}_i|)$,
\item $\mathcal{X}_i$ est presque affine,
\item $\mathcal{I}_{\mathcal{X}_i} \to
\mathcal{X}_i \times_\mathcal{X} \mathcal{I}_\mathcal{X}$
est un isomorphisme, et
\item $g_i : \mathcal{X}_i \to \mathcal{X}$ est représentable
par des espaces algébriques, séparé et \'etale,
\end{enumerate}
\end{lemma}

\begin{proof}
Pour tout $x \in |\mathcal{X}|$, nous pouvons choisir
$g : \mathcal{U} \to \mathcal{X}$, $\mathcal{U} = [U/R]$ et $u$ comme dans
Morphismes de champs,
lemme \ref{stacks-morphisms-lemma-etale-local-quasi-DM-at-x}.
Puis, d'après
Morphismes de champs, lemme
\ref{stacks-morphisms-lemma-stabilizer-preserving-unramified}
nous voyons qu'il existe un sous-champ ouvert
$\mathcal{U}' \subset \mathcal{U}$ contenant $u$
tel que $\mathcal{I}_{\mathcal{U}'} \to
\mathcal{U}' \times_\mathcal{X} \mathcal{I}_\mathcal{X}$
est un isomorphisme.
Soit $U' \subset U$ l'ouvert invariant par $R$ correspondant au
sous-champ ouvert $\mathcal{U}'$.
Soit $u' \in U'$ un point de $U'$ s'envoyant sur $u$.
Remarquons que $t(s^{-1}(\{u'\}))$ est fini, puisque $s : R \to U$ est fini.
D'après Propriétés, lemme \ref{properties-lemma-ample-finite-set-in-affine}
et Groupoïdes, lemme \ref{groupoids-lemma-find-invariant-affine},
nous pouvons trouver un ouvert affine invariant par $R$, noté $U'' \subset U'$,
contenant $u'$. Soit $R''$ la restriction de $R$ à $U''$.
Alors $\mathcal{U}'' = [U''/R'']$ est un sous-champ ouvert de
$\mathcal{U}'$ contenant $u$, est presque affine,
$\mathcal{I}_{\mathcal{U}''} \to
\mathcal{U}'' \times_\mathcal{X} \mathcal{I}_\mathcal{X}$
est un isomorphisme, et $\mathcal{U}'' \to \mathcal{X}$
est représentable par des espaces algébriques et \'etale.
Enfin, $\mathcal{U}'' \to \mathcal{X}$ est séparé, car
$\mathcal{U}''$ est séparé (lemme \ref{stacks-more-morphisms-lemma-well-nigh-affine}),
la diagonale de $\mathcal{X}$ est séparée
(Morphismes de champs, lemme \ref{stacks-morphisms-lemma-diagonal-diagonal}),
et la séparation résulte de Morphismes de champs, lemme
\ref{stacks-morphisms-lemma-compose-after-separated}.
Puisque le point $x \in |\mathcal{X}|$ est arbitraire, la démonstration est achevée.
\end{proof}

\begin{theorem}[Keel-Mori]
\label{stacks-more-morphisms-theorem-keel-mori}
Soit $\mathcal{X}$ un champ algébrique. Supposons que
$\mathcal{I}_\mathcal{X} \to \mathcal{X}$ est fini.
Il existe alors un espace de modules catégorique uniforme
$$
f : \mathcal{X} \longrightarrow M
$$
et $f$ est séparé, quasi-compact et un homéomorphisme universel.
\end{theorem}

\begin{proof}
Choisissons un ensemble $I$\footnote{Le lecteur qui tient encore
compte des questions ensemblistes devra s'assurer que $I$ n'est pas trop grand.}
et, pour $i \in I$, un morphisme de champs algébriques
$g_i : \mathcal{X}_i \to \mathcal{X}$ comme dans
le lemme \ref{stacks-more-morphisms-lemma-etale-local-finite-inertia} ; nous utiliserons sans autre mention toutes
les propriétés énumérées dans ce lemme.
Soit
$$
f_i : \mathcal{X}_i \to M_i
$$
le morphisme donné par le lemme \ref{stacks-more-morphisms-lemma-well-nigh-affine-moduli-space}.
Considérons les champs
$$
\mathcal{X}_{ij} = \mathcal{X}_i \times_{g_i, \mathcal{X}, g_j} \mathcal{X}_j
$$
pour $i, j \in I$. Les projections $\mathcal{X}_{ij} \to \mathcal{X}_i$
et $\mathcal{X}_{ij} \to \mathcal{X}_j$ sont séparées
d'après Morphismes de champs, lemme
\ref{stacks-morphisms-lemma-base-change-separated},
\'etales d'après Morphismes de champs, lemme
\ref{stacks-morphisms-lemma-base-change-etale},
et induisent des isomorphismes sur les groupes d'automorphismes
(comme dans Morphismes de champs, remarque
\ref{stacks-morphisms-remark-identify-automorphism-groups}), d'après
Morphismes de champs, lemme
\ref{stacks-morphisms-lemma-base-change-stabilizer-preserving}.
Nous pouvons donc appliquer le lemme \ref{stacks-more-morphisms-lemma-etale-separated-over-well-nigh-affine}
afin d'obtenir un diagramme commutatif
$$
\xymatrix{
\mathcal{X}_i \ar[d]_{f_i} &
\mathcal{X}_{ij} \ar[d]_{f_{ij}} \ar[l] \ar[r] &
\mathcal{X}_j \ar[d]_{f_j} \\
M_i &
M_{ij} \ar[l] \ar[r] &
M_j
}
$$
dont les carrés sont cartésiens et où $M_{ij} \to M_i$ et $M_{ij} \to M_j$
sont des morphismes de schémas séparés et \'etales ; nous utilisons ici aussi que $f_i$
est un quotient catégorique uniforme d'après le
lemme \ref{stacks-more-morphisms-lemma-moduli-space-finite-affine}.
Affirmation :
$$
\coprod M_{ij} \longrightarrow \coprod M_i \times \coprod M_i
$$
est une relation d'équivalence \'etale.

\medskip\noindent
Démonstration de l'affirmation. Posons $R = \coprod M_{ij}$ et $U = \coprod M_i$.
Nous avons déjà vu que $t : R \to U$ et $s : R \to U$ sont \'etales.
Construisons un morphisme $c : R \times_{s, U, t} R \to R$
compatible avec $\text{pr}_{13} : U \times U \times U \to U \times U$.
Plus précisément, pour $i, j, k \in I$, considérons
$$
\mathcal{X}_{ijk} =
\mathcal{X}_i \times_{g_i, \mathcal{X}, g_j} \mathcal{X}_j
\times_{g_j, \mathcal{X}, g_k} \mathcal{X}_k =
\mathcal{X}_{ij} \times_{\mathcal{X}_j} \mathcal{X}_{jk}
$$
En raisonnant exactement comme au paragraphe précédent,
nous voyons que $M_{ijk} = M_{ij} \times_{M_j} M_{jk}$
est un espace de modules catégorique pour $\mathcal{X}_{ijk}$.
En particulier, il existe un morphisme canonique
$M_{ijk} = M_{ij} \times_{M_j} M_{jk} \to M_{ik}$
provenant de la projection $\mathcal{X}_{ijk} \to \mathcal{X}_{ik}$.
En réunissant ces morphismes, nous obtenons le morphisme $c$.
De même, nous construisons un morphisme $e : U \to R$
compatible avec $\Delta : U \to U \times U$ et
$i : R \to R$ compatible avec la permutation des facteurs $U \times U \to U \times U$.
Soit $k$ un corps algébriquement clos. Alors
$$
\Mor(\Spec(k), \mathcal{X}_i) \to \Mor(\Spec(k), M_i) = M_i(k)
$$
est bijective sur les classes d'isomorphie, et il en reste de même après tout
changement de base par un morphisme $M' \to M_i$. Cela résulte du choix
de $f_i$ et de Morphismes de champs, lemmes
\ref{stacks-morphisms-lemma-universally-injective} et
\ref{stacks-morphisms-lemma-universally-injective-point}.
Par construction des $2$-produits fibrés, le diagramme
$$
\xymatrix{
\Mor(\Spec(k), \mathcal{X}_{ij}) \ar[d] \ar[r] &
\Mor(\Spec(k), \mathcal{X}_j) \ar[d] \\
\Mor(\Spec(k), \mathcal{X}_i) \ar[r] &
\Mor(\Spec(k), \mathcal{X})
}
$$
est un 2-produit fibré de catégories. D'après notre choix de $g_i$, les
foncteurs de ce diagramme induisent des bijections sur les groupes d'automorphismes.
Il s'ensuit que ce diagramme induit un diagramme cartésien
sur les ensembles de classes d'isomorphie ! Nous voyons donc que
$$
R(k) = U(k) \times_{|\Mor(\Spec(k), \mathcal{X})|} U(k)
$$
où $|\Mor(\Spec(k), \mathcal{X})|$ désigne l'ensemble
des classes d'isomorphie.
En particulier, pour tout corps algébriquement clos $k$,
l'application sur les points à valeurs dans $k$ définit une relation d'équivalence.
L'affirmation en résulte par
Groupoïdes, lemme \ref{groupoids-lemma-etale-equivalence-relation}.

\medskip\noindent
Soit $M = U/R$ l'espace algébrique quotient de la
relation d'équivalence \'etale ci-dessus ; voir
Espaces, théorème \ref{spaces-theorem-presentation}.
Il existe un morphisme canonique $f : \mathcal{X} \to M$
s'insérant dans les diagrammes commutatifs
\begin{equation}
\label{stacks-more-morphisms-equation-fundamental-diagram}
\xymatrix{
\mathcal{X}_i \ar[r]_{g_i} \ar[d]_{f_i} & \mathcal{X} \ar[d]^f \\
M_i \ar[r] & M
}
\end{equation}
En effet, un tel morphisme $f$ est défini par un foncteur
$$
f : \Mor(T, \mathcal{X}) \longrightarrow \Mor(T, M)
$$
pour tout schéma $T$, la famille ainsi obtenue étant compatible avec les changements de base. Soit $a : T \to \mathcal{X}$
un objet du membre de gauche. Nous obtenons un recouvrement \'etale
$\{T_i \to T\}$, avec $T_i = \mathcal{X}_i \times_\mathcal{X} T$,
et des morphismes $a_i : T_i \to \mathcal{X}_i$. Nous obtenons alors
$b_i = f_i \circ a_i : T_i \to M_i$. Puisque
$T_i \times_T T_j = \mathcal{X}_{ij} \times_\mathcal{X} T$,
nous obtenons de plus un morphisme $a_{ij} : T_i \times_T T_j \to \mathcal{X}_{ij}$.
En posant $b_{ij} = f_{ij} \circ a_{ij}$, nous constatons que
$b_i \times b_j$ se factorise par le monomorphisme
$M_{ij} \to M_i \times M_j$. Par conséquent, les morphismes
$$
T_i \xrightarrow{b_i} M_i \to M
$$
coïncident sur $T_i \times_T T_j$. Comme $M$ est un faisceau pour la
topologie \'etale, nous voyons que ces morphismes se recollent en un unique
morphisme $b = f(a) : T \to M$. Nous omettons de vérifier que
cette construction est compatible avec les changements de base, ainsi que de vérifier
que les diagrammes
(\ref{stacks-more-morphisms-equation-fundamental-diagram}) sont commutatifs.

\medskip\noindent
Affirmation : les diagrammes (\ref{stacks-more-morphisms-equation-fundamental-diagram}) sont cartésiens.
Pour le voir, étudions le morphisme induit
$$
h_i : \mathcal{X}_i \longrightarrow M_i \times_M \mathcal{X}
$$
C'est un morphisme de champs \'etale sur $\mathcal{X}$,
donc $h_i$ est \'etale (Morphismes de champs, lemme
\ref{stacks-morphisms-lemma-etale-permanence}).
Comme $g_i$ est séparé, nous voyons que $h_i$ est séparé
(utiliser Morphismes de champs, lemme
\ref{stacks-morphisms-lemma-compose-after-separated}, et le fait,
vu plus haut, que la diagonale de $\mathcal{X}$ est séparée).
Le morphisme $h_i$ induit des isomorphismes sur les groupes d'automorphismes
(Morphismes de champs, remarque
\ref{stacks-morphisms-remark-identify-automorphism-groups})
puisqu'il en est ainsi pour $g_i$. Pour un corps algébriquement clos $k$,
le diagramme
$$
\xymatrix{
\Mor(\Spec(k), M_i \times_M \mathcal{X}) \ar[r] \ar[d] &
\Mor(\Spec(k), \mathcal{X}) \ar[d] \\
M_i(k) \ar[r] &
M(k)
}
$$
est un diagramme cartésien de catégories, et la flèche du haut
induit des bijections sur les groupes d'automorphismes.
D'autre part, nous avons
$$
M(k) = U(k)/R(k) = U(k)/
U(k) \times_{|\Mor(\Spec(k), \mathcal{X})|} U(k) =
|\Mor(\Spec(k), \mathcal{X})|
$$
d'après ce qui précède. Ainsi, la flèche verticale de droite du
diagramme cartésien ci-dessus est bijective sur les classes d'isomorphie.
Nous concluons que
$|\Mor(\Spec(k), M_i \times_M \mathcal{X})| \to M_i(k)$ est bijective.
Récapitulons : $h_i$ est séparé et \'etale, il induit des isomorphismes sur les
groupes d'automorphismes (comme dans Morphismes de champs, remarque
\ref{stacks-morphisms-remark-identify-automorphism-groups}),
et une équivalence sur toute catégorie fibre au-dessus d'un corps algébriquement clos.
Il s'agit donc d'un isomorphisme d'après Morphismes de champs, lemme
\ref{stacks-morphisms-lemma-etale-iso}.

\medskip\noindent
L'affirmation donne en particulier ceci :
il existe un morphisme \'etale surjectif $U \to M$
tel que le changement de base de $f$ est séparé, quasi-compact
et un homéomorphisme universel. Il s'ensuit que $f$ est séparé,
quasi-compact et un homéomorphisme universel.
Voir Morphismes de champs, lemme
\ref{stacks-morphisms-lemma-check-separated-covering},
\ref{stacks-morphisms-lemma-check-quasi-compact-covering}, et
\ref{stacks-morphisms-lemma-check-universal-homeomorphism-covering}.

\medskip\noindent
Pour achever la démonstration, il reste à montrer que $f : \mathcal{X} \to M$
est un espace de modules catégorique uniforme.
Pour cela, il suffit de montrer que, pour tout morphisme plat
$M' \to M$ d'espaces algébriques, le changement de base
$$
M' \times_M \mathcal{X} \longrightarrow M'
$$
est un espace de modules catégorique. Considérons donc un morphisme
$$
\theta : M' \times_M \mathcal{X} \longrightarrow E
$$
où $E$ est un espace algébrique. Pour chaque $i$, nous savons que
$f_i$ est un espace de modules catégorique uniforme. Nous obtenons donc
$$
\xymatrix{
M' \times_M \mathcal{X}_i \ar[d] \ar[r] &
M' \times_M \mathcal{X} \ar[d]^\theta \\
M' \times_M M_i \ar[r]^{\psi_i} &
E
}
$$
Puisque $\{M' \times_M M_i \to M'\}$ est un recouvrement \'etale,
pour obtenir le morphisme souhaité $\psi : M' \to E$, il suffit
de montrer que $\psi_i$ et $\psi_j$ coïncident sur
$M' \times_M M_i \times_M M_j = M' \times_M M_{ij}$.
Cela résulte aussitôt du fait que
$f_{ij} : \mathcal{X}_{ij} =
\mathcal{X}_i \times_\mathcal{X} \mathcal{X}_j \to M_{ij}$ est un quotient catégorique
uniforme ; nous omettons les détails.
Enfin, on montre que $\psi$ s'insère dans le diagramme commutatif
$$
\xymatrix{
M' \times_M \mathcal{X} \ar[d] \ar[rd]^\theta \\
M' \ar[r]^\psi &
E
}
$$
car « $\{M' \times_M \mathcal{X}_i \to M' \times_M \mathcal{X}\}$
est un recouvrement \'etale » et que les morphismes $\psi_i$ s'insèrent dans les
diagrammes commutatifs correspondants par construction.
Ceci achève la démonstration du théorème de Keel--Mori.
\end{proof}

\noindent
Le lemme suivant souligne le caractère local pour la topologie \'etale de la construction.

\begin{lemma}
\label{stacks-more-morphisms-lemma-etale-separated-over-keel-mori}
Soit $h : \mathcal{X}' \to \mathcal{X}$ un morphisme de champs algébriques.
Supposons que
\begin{enumerate}
\item $\mathcal{I}_\mathcal{X} \to \mathcal{X}$ est fini,
\item $h$ est \'etale et séparé, et induit des isomorphismes sur les
groupes d'automorphismes (Morphismes de champs, remarque
\ref{stacks-morphisms-remark-identify-automorphism-groups}).
\end{enumerate}
Il existe alors un diagramme cartésien
$$
\xymatrix{
\mathcal{X}' \ar[d] \ar[r] &
\mathcal{X} \ar[d] \\
M' \ar[r] &
M
}
$$
où $M' \to M$ est un morphisme d'espaces algébriques séparé et \'etale, et
les flèches verticales sont les espaces de modules construits dans le
théorème \ref{stacks-more-morphisms-theorem-keel-mori}.
\end{lemma}

\begin{proof}
D'après Morphismes de champs, lemme \ref{stacks-morphisms-lemma-aut-iso-unramified},
nous voyons que
$\mathcal{I}_{\mathcal{X}'} \to
\mathcal{X}' \times_\mathcal{X} \mathcal{I}_\mathcal{X}$
est un isomorphisme. Ainsi, $\mathcal{I}_{\mathcal{X}'} \to \mathcal{X}'$
est fini comme changement de base de $\mathcal{I}_\mathcal{X} \to \mathcal{X}$.
Soient $f' : \mathcal{X}' \to M'$ et $f : \mathcal{X} \to M$ ceux du
théorème \ref{stacks-more-morphisms-theorem-keel-mori}.
Nous obtenons un diagramme commutatif comme dans le lemme, car
$f'$ est un espace de modules catégorique.
Choisissons $I$ et $g'_i : \mathcal{X}'_i \to \mathcal{X}'$ comme dans le
lemme \ref{stacks-more-morphisms-lemma-etale-local-finite-inertia}.
Remarquons que $g_i = h \circ g'_i$
est \'etale et séparé, et induit des isomorphismes sur les
groupes d'automorphismes (Morphismes de champs, remarque
\ref{stacks-morphisms-remark-identify-automorphism-groups}).
Soit $f'_i : \mathcal{X}'_i \to M'_i$ le morphisme du
lemme \ref{stacks-more-morphisms-lemma-well-nigh-affine-moduli-space}.
Dans la démonstration du théorème \ref{stacks-more-morphisms-theorem-keel-mori},
nous avons vu que les diagrammes
$$
\xymatrix{
\mathcal{X}'_i \ar[d]_{f'_i} \ar[r]_{g'_i} &
\mathcal{X}' \ar[d]^{f'} \\
M'_i \ar[r] &
M'
}
\quad\text{et}\quad
\xymatrix{
\mathcal{X}'_i \ar[d]_{f'_i} \ar[r]_{g_i} &
\mathcal{X} \ar[d]^f \\
M'_i \ar[r] &
M
}
$$
sont cartésiens et que $M'_i \to M'$ et $M'_i \to M$ sont \'etales
(cela résulte aussi directement des propriétés des morphismes
$g'_i, g_i, f', f'_i, f$ énumérées jusqu'ici, en raisonnant exactement de la même façon).
Nous en déduisons d'une part que $M' \to M$ est \'etale et d'autre part que
le diagramme du lemme est cartésien. Il reste à montrer
que $M' \to M$ est séparé. Pour cela, considérons le
diagramme
$$
\xymatrix{
\mathcal{X}' \ar[r] \ar[d] &
\mathcal{X}' \times_\mathcal{X} \mathcal{X}' \ar[d] \\
M' \ar[r] &
M' \times_M M'
}
$$
La flèche horizontale supérieure est universellement fermée, car
$\mathcal{X}' \to \mathcal{X}$ est séparé.
Les flèches verticales sont celles du théorème \ref{stacks-more-morphisms-theorem-keel-mori}
(en tant que changements de base plats de $\mathcal{X} \to M$)
et sont donc des homéomorphismes universels. La flèche horizontale inférieure
est par conséquent universellement fermée. Comme elle est en outre un monomorphisme
\'etale d'espaces algébriques, c'est l'immersion fermée voulue.
\end{proof}





\section{Propriétés des espaces de modules}
\label{stacks-more-morphisms-section-properties-moduli-spaces}

\noindent
Une fois établie l'existence d'un espace de modules,
il devient possible (et généralement aisé) d'établir
des propriétés de ces espaces de modules.

\begin{lemma}
\label{stacks-more-morphisms-lemma-keel-mori-finite-type}
Soit $p : \mathcal{X} \to Y$ un morphisme d'un champ algébrique vers un
espace algébrique. Supposons que
\begin{enumerate}
\item $\mathcal{I}_\mathcal{X} \to \mathcal{X}$ est fini,
\item $Y$ est localement noethérien, et
\item $p$ est localement de type fini.
\end{enumerate}
Soit $f : \mathcal{X} \to M$ l'espace de modules construit dans le
théorème \ref{stacks-more-morphisms-theorem-keel-mori}.
Alors $M \to Y$ est localement de type fini.
\end{lemma}

\begin{proof}
Puisque $f$ est un espace de modules catégorique uniforme, nous obtenons le
morphisme $M \to Y$. Il suffit de vérifier que $M \to Y$
est localement de type fini, localement pour la topologie \'etale sur $M$ et $Y$.
Puisque $f$ est un espace de modules catégorique uniforme, nous
pouvons d'abord remplacer $Y$ par un schéma affine \'etale sur $Y$.
Ensuite, nous pouvons choisir $I$ et $g_i : \mathcal{X}_i \to \mathcal{X}$
comme dans le lemme \ref{stacks-more-morphisms-lemma-etale-local-finite-inertia}.
Alors, par le lemme \ref{stacks-more-morphisms-lemma-etale-separated-over-keel-mori},
nous nous ramenons au cas $\mathcal{X} = \mathcal{X}_i$.
Autrement dit, nous pouvons supposer que $\mathcal{X}$ est presque affine.
Dans ce cas, nous avons $Y = \Spec(A_0)$,
$\mathcal{X} = [U/R]$ avec $U = \Spec(A)$ et
$M = \Spec(C)$, où $C \subset A$ est l'ensemble des fonctions invariantes par $R$
sur $U$. Voir les
lemmes \ref{stacks-more-morphisms-lemma-well-nigh-affine} et
\ref{stacks-more-morphisms-lemma-well-nigh-affine-moduli-space}.
Alors $A_0$ est noethérien et $A_0 \to A$ est de type fini.
De plus, $A$ est entier sur $C$ d'après
Groupoïdes, lemme \ref{groupoids-lemma-integral-over-invariants},
donc fini sur $C$
(cet anneau étant de type fini sur $A_0$).
Nous pouvons donc enfin appliquer
Algèbre, lemme \ref{algebra-lemma-Artin-Tate}
pour conclure.
\end{proof}

\begin{lemma}
\label{stacks-more-morphisms-lemma-keel-mori-diagonal}
Soit $\mathcal{X}$ un champ algébrique. Supposons que
$\mathcal{I}_\mathcal{X} \to \mathcal{X}$ est fini.
Soit $f : \mathcal{X} \to M$ l'espace de modules construit dans le
théorème \ref{stacks-more-morphisms-theorem-keel-mori}.
\begin{enumerate}
\item Si $\mathcal{X}$ est quasi-séparé, alors $M$ est quasi-séparé.
\item Si $\mathcal{X}$ est séparé, alors $M$ est séparé.
\item Ajouter ici d'autres énoncés, par exemple des versions relatives de ceux qui précèdent.
\end{enumerate}
\end{lemma}

\begin{proof}
Pour le prouver, considérons le diagramme suivant
$$
\xymatrix{
\mathcal{X} \ar[d]_f \ar[r]_{\Delta_\mathcal{X}} &
\mathcal{X} \times \mathcal{X} \ar[d]^{f \times f} \\
M \ar[r]^{\Delta_M} &
M \times M
}
$$
Puisque $f$ est un homéomorphisme universel, nous voyons que
$f \times f$ est un homéomorphisme universel.

\medskip\noindent
Si $\mathcal{X}$ est séparé, alors $\Delta_\mathcal{X}$ est propre,
donc $\Delta_\mathcal{X}$ est universellement fermé ; par conséquent,
$\Delta_M$ est universellement fermé, et $M$ est donc séparé
d'après Morphismes d'espaces, lemme
\ref{spaces-morphisms-lemma-separated-diagonal-proper}.

\medskip\noindent
Si $\mathcal{X}$ est quasi-séparé, alors $\Delta_\mathcal{X}$ est
quasi-compact ; par conséquent, $\Delta_M$ est quasi-compact, et $M$ est donc
quasi-séparé.
\end{proof}

\begin{lemma}
\label{stacks-more-morphisms-lemma-keel-mori-proper}
Soit $p : \mathcal{X} \to Y$ un morphisme d'un champ algébrique
vers un espace algébrique. Supposons que
\begin{enumerate}
\item $\mathcal{I}_\mathcal{X} \to \mathcal{X}$ est fini,
\item $p$ est propre, et
\item $Y$ est localement noethérien.
\end{enumerate}
Soit $f : \mathcal{X} \to M$ l'espace de modules construit dans le
théorème \ref{stacks-more-morphisms-theorem-keel-mori}. Alors $M \to Y$ est propre.
\end{lemma}

\begin{proof}
D'après le lemme \ref{stacks-more-morphisms-lemma-keel-mori-finite-type},
nous voyons que $M \to Y$ est localement de type fini.
D'après le lemme \ref{stacks-more-morphisms-lemma-keel-mori-diagonal}, nous voyons que
$M \to Y$ est séparé.
Bien entendu, $M \to Y$ est quasi-compact et universellement fermé,
car ces propriétés sont topologiques, le morphisme $\mathcal{X} \to Y$
les possède et $\mathcal{X} \to M$ est un
homéomorphisme universel.
\end{proof}





\section{Champs et recouvrements fpqc}
\label{stacks-more-morphisms-section-fpqc}

\noindent
Certains champs algébriques satisfont à la descente fpqc.
L'analogue de cette section pour les espaces algébriques
est Propriétés des espaces, section \ref{spaces-properties-section-fpqc}.


\begin{proposition}
\label{stacks-more-morphisms-proposition-stack-fpqc}
\begin{reference}
Proposition 3.3.6 de « Intro to Algebraic Stacks », par
Anatoly Preygel.
\end{reference}
Soit $\mathcal{X}$ un champ algébrique dont la diagonale est quasi-affine\footnote{Il suffit
de la supposer ind-quasi-affine.}. Alors
$\mathcal{X}$ satisfait à la descente pour les recouvrements fpqc.
\end{proposition}

\begin{proof}
Selon nos conventions, $\mathcal{X}$ est un champ
en groupoïdes $p : \mathcal{X} \to (\Sch/S)_{fppf}$
sur la catégorie des schémas sur un schéma de base $S$, munie
de la topologie fppf. L'énoncé signifie ce qui suit :
étant donné un recouvrement fpqc $\mathcal{U} = \{U_i \to U\}_{i \in I}$
de schémas sur $S$, le foncteur
$$
\mathcal{X}_U \longrightarrow DD(\mathcal{U})
$$
est une équivalence. À gauche se trouve la catégorie
des objets de $\mathcal{X}$ au-dessus de $U$, et à droite
la catégorie des données de descente dans $\mathcal{X}$ relativement à $\mathcal{U}$.
Voir l'explication dans Champs, section \ref{stacks-section-descent-data}.

\medskip\noindent
Pleine fidélité. Supposons que nous ayons deux objets $x, y$ de $\mathcal{X}$
au-dessus de $U$. Alors $J = \mathit{Isom}(x, y)$ est un espace algébrique sur $U$.
Par conséquent, une famille de sections de $J$ sur les $U_i$ dont les restrictions
à $U_i \times_U U_j$ coïncident provient d'une unique section sur $U$, d'après
l'analogue de la proposition pour les espaces algébriques ; voir
Propriétés des espaces, proposition
\ref{spaces-properties-proposition-sheaf-fpqc}.
Notre foncteur est donc pleinement fidèle.

\medskip\noindent
Surjectivité essentielle. On se donne ici des objets
$x_i$ au-dessus de $U_i$ et des isomorphismes
$\varphi_{ij} : \text{pr}_0^*x_i \to \text{pr}_1^*x_j$
sur $U_i \times_U U_j$ satisfaisant à la condition de cocycle
sur $U_i \times_U U_j \times_U U_k$.

\medskip\noindent
Soit $W$ un schéma affine et soit $W \to \mathcal{X}$
un morphisme. Pour chaque $i$, nous pouvons former
$$
W_i = U_i \times_{x_i, \mathcal{X}} W
$$
La projection $W_i \to U_i$ est quasi-affine, car la
diagonale de $\mathcal{X}$ est quasi-affine. Pour toute paire
$i, j \in I$, l'isomorphisme $\varphi_{ij}$ induit
un isomorphisme
$$
W_i \times_U U_j =
(U_i \times_U U_j) \times_{x_i \circ \text{pr}_0, \mathcal{X}} W
\to
(U_i \times_U U_j) \times_{x_j \circ \text{pr}_1, \mathcal{X}} W =
U_i \times_U W_j
$$
De plus, ces isomorphismes satisfont à la condition de cocycle sur
$U_i \times_U U_j \times_U U_k$. Autrement dit, ils
définissent une donnée de descente sur les schémas $W_i/U_i$ relativement à $\mathcal{U}$.
D'après Descente, lemme \ref{descent-lemma-quasi-affine},
nous voyons que cette donnée de descente est effective\footnote{Ou bien utiliser
Compléments sur les groupoïdes, lemme \ref{more-groupoids-lemma-ind-quasi-affine},
dans le cas d'une diagonale ind-quasi-affine.}.
Nous en concluons qu'il existe un morphisme quasi-affine
$W' \to U$ et un diagramme commutatif
$$
\xymatrix{
W' \ar[d] &
W_i \ar[l] \ar[d] \ar[r] &
W \ar[d] \\
U &
U_i \ar[l] \ar[r]^{x_i} &
\mathcal{X}
}
$$
dont les carrés sont cartésiens. Puisque $\{W_i \to W'\}_{i \in I}$
est le changement de base de $\mathcal{U}$ par $W' \to U$, nous concluons
qu'il s'agit d'un recouvrement fpqc. Puisque $W$ satisfait à la condition de faisceau
pour les recouvrements fpqc, nous obtenons un unique morphisme
$W' \to W$ tel que $W_i \to W' \to W$ soit le morphisme donné
$W_i \to W$. Autrement dit, nous avons les diagrammes commutatifs
$$
\xymatrix{
W_i \ar[d] \ar[r] &
W' \ar[d] \ar[r] &
W \ar[d] \\
U_i \ar[r] \ar@/_1pc/[rr]_{x_i} &
U &
\mathcal{X}
}
$$
compatibles avec les isomorphismes $\varphi_{ij}$ et dont le
carré et le rectangle sont cartésiens.

\medskip\noindent
Choisissons une famille de schémas affines $W_\alpha$, $\alpha \in A$,
et des morphismes lisses $W_\alpha \to \mathcal{X}$ tels que
$\coprod W_\alpha \to \mathcal{X}$ soit surjectif.
Par le procédé du paragraphe précédent, nous construisons
un diagramme
$$
\xymatrix{
W_{\alpha, i} \ar[d] \ar[r] &
W_\alpha' \ar[d] \ar[r] &
W_\alpha \ar[d] \\
U_i \ar[r] \ar@/_1pc/[rr]_{x_i} &
U &
\mathcal{X}
}
$$
pour chaque $\alpha$. Les morphismes $W_\alpha' \to U$ sont alors
lisses et conjointement surjectifs.

\medskip\noindent
Notons $x_\alpha$ l'objet de $\mathcal{X}$ au-dessus de $W_\alpha'$
correspondant à $W_\alpha' \to W_\alpha \to \mathcal{X}$.
Puisque $\mathcal{X}$ est un champ pour la topologie fppf et que
$\{W_\alpha' \to U\}$ est un recouvrement fppf, il suffit
de montrer qu'il existe des isomorphismes
$\text{pr}_0^*x_\alpha \to \text{pr}_1^*x_\beta$
sur $W_\alpha' \times_U W'_\beta$ satisfaisant à la
condition de cocycle. Or, après restriction à
$W_{\alpha, i}$, nous disposons bien de tels isomorphismes sur
$W_{\alpha, i} \times_{U_i} W_{\beta, i} =
U_i \times_U (W_\alpha' \times_U W'_\beta)$
car l'image inverse de $x_\alpha$ sur $W_{\alpha, i}$
est isomorphe à l'image inverse de $x_i$ sur $W_{\alpha, i}$.
Puisque $\{U_i \times_U (W_\alpha' \times_U W'_\beta) \to
W_\alpha' \times_U W'_\beta\}_{i \in I}$
est un recouvrement fpqc et que, par la compatibilité susmentionnée
des diagrammes ci-dessus avec $\varphi_{ij}$, ces
isomorphismes descendent à $W_\alpha' \times_U W'_\beta$,
la démonstration est achevée.
\end{proof}










\section{Foncteurs tensoriels}
\label{stacks-more-morphisms-section-tensor-functors}

\noindent
Soit $f : \mathcal{Y} \to \mathcal{X}$ un morphisme de champs algébriques
noethériens. Le foncteur image inverse
$$
f^* :
\textit{Coh}(\mathcal{O}_\mathcal{X})
\longrightarrow
\textit{Coh}(\mathcal{O}_\mathcal{Y})
$$
est un foncteur tensoriel exact à droite : il est additif, exact à droite
et commute aux produits tensoriels de modules cohérents. On peut
se demander dans quelle mesure tout foncteur tensoriel exact à droite
$F : \textit{Coh}(\mathcal{O}_\mathcal{X}) \to
\textit{Coh}(\mathcal{O}_\mathcal{Y})$ provient d'un
morphisme $f : \mathcal{Y} \to \mathcal{X}$.
Le lecteur pourra consulter
\cite{Hall-Rydh-coherent} pour un résultat très général
de ce type.
Le but de cette section est de donner une courte démonstration du
théorème \ref{stacks-more-morphisms-theorem-main} en guise d'introduction à ces idées.

\medskip\noindent
Commençons par quelques lemmes.

\begin{lemma}
\label{stacks-more-morphisms-lemma-extend}
Soient $\mathcal{X}$ et $\mathcal{Y}$ des champs algébriques noethériens.
Tout foncteur tensoriel exact à droite $F : \textit{Coh}(\mathcal{O}_\mathcal{X}) \to
\textit{Coh}(\mathcal{O}_\mathcal{Y})$ se prolonge de manière unique en un
foncteur tensoriel exact à droite
$F : \QCoh(\mathcal{O}_\mathcal{X}) \to \QCoh(\mathcal{O}_\mathcal{Y})$
commutant à toutes les limites inductives.
\end{lemma}

\begin{proof}
L'existence et l'unicité du prolongement sont un fait général ; voir
Catégories, lemme \ref{categories-lemma-extend-functor-by-colim}.
Pour voir que le lemme s'applique, remarquons que les modules cohérents sur des
champs algébriques localement noethériens sont, par définition, des modules
de présentation finie ; voir
Cohomologie des champs, définition \ref{stacks-cohomology-definition-coherent}.
Ainsi, un module cohérent sur $\mathcal{X}$ est un
objet compact au sens catégorique de $\QCoh(\mathcal{O}_\mathcal{X})$ d'après
Cohomologie des champs, lemme
\ref{stacks-cohomology-lemma-finite-presentation-quasi-compact-colimit}.
Enfin, tout module quasi-cohérent est une limite inductive filtrante de
ses sous-modules cohérents d'après Cohomologie des champs, lemme
\ref{stacks-cohomology-lemma-directed-colimit-coherent}.

\medskip\noindent
Puisque $F$ est additif, le prolongement de $F$ l'est aussi (détails omis).
Puisque $F$ est un foncteur tensoriel et que les limites inductives de modules
commutent à la formation des produits tensoriels, le prolongement de $F$
est lui aussi un foncteur tensoriel (détails omis).

\medskip\noindent
Dans ce paragraphe, montrons que le prolongement commute aux sommes directes
arbitraires. Si $\mathcal{F} = \bigoplus_{j \in J} \mathcal{H}_j$
avec $\mathcal{H}_j$ quasi-cohérent, alors
$\mathcal{F} = \colim_{J' \subset J\text{ fini}}
\bigoplus_{j \in J'} \mathcal{H}_j$.
En notant encore $F$ le prolongement de $F$, nous obtenons
\begin{align*}
F(\mathcal{F})
& =
\colim_{J' \subset J\text{ fini}}
F(\bigoplus\nolimits_{j \in J'} \mathcal{H}_j) \\
& =
\colim_{J' \subset J\text{ fini}}
\bigoplus\nolimits_{j \in J'} F(\mathcal{H}_j) \\
& =
\bigoplus\nolimits_{j \in J} F(\mathcal{H}_j)
\end{align*}
Ainsi, $F$ commute aux sommes directes arbitraires.

\medskip\noindent
Dans ce paragraphe, montrons que le prolongement est exact à droite.
Supposons que $0 \to \mathcal{F} \to \mathcal{F}' \to \mathcal{F}'' \to 0$
soit une suite exacte courte de $\mathcal{O}_\mathcal{X}$-modules quasi-cohérents.
Écrivons alors $\mathcal{F}' = \bigcup_{i \in I} \mathcal{F}'_i$ comme
réunion de ses sous-modules cohérents (voir la référence donnée plus haut).
Notons $\mathcal{F}''_i \subset \mathcal{F}''$ l'image de $\mathcal{F}'_i$
et posons $\mathcal{F}_i = \mathcal{F} \cap \mathcal{F}'_i =
\Ker(\mathcal{F}'_i \to \mathcal{F}''_i)$. Il est alors clair que
$\mathcal{F} = \bigcup \mathcal{F}_i$ et
$\mathcal{F}'' = \bigcup \mathcal{F}''_i$
et que nous avons des suites exactes courtes
$$
0 \to \mathcal{F}_i \to \mathcal{F}_i' \to \mathcal{F}_i'' \to 0
$$
Puisque le prolongement commute aux limites inductives filtrantes, nous avons
$F(\mathcal{F}) = \colim_{i \in I} F(\mathcal{F}_i)$,
$F(\mathcal{F}') = \colim_{i \in I} F(\mathcal{F}'_i)$, et
$F(\mathcal{F}'') = \colim_{i \in I} F(\mathcal{F}''_i)$.
Comme les limites inductives filtrantes de faisceaux de modules sont exactes, nous
concluons que le prolongement de $F$ est exact à droite.

\medskip\noindent
La démonstration est achevée, car tout foncteur exact à droite qui commute aux
coproduits commute à toutes les limites inductives ; voir
Catégories, lemme \ref{categories-lemma-colimits-coproducts-coequalizers}.
\end{proof}

\begin{lemma}
\label{stacks-more-morphisms-lemma-affine}
Soit $\mathcal{X}$ un champ algébrique à diagonale affine.
Soit $B$ un anneau. Soit $F : \QCoh(\mathcal{O}_\mathcal{X}) \to \text{Mod}_B$
un foncteur tensoriel exact à droite qui commute aux sommes directes.
Soit $g : U \to \mathcal{X}$ un morphisme, où $U = \Spec(A)$ est affine. Alors
\begin{enumerate}
\item $C = F(g_{\QCoh, *}\mathcal{O}_U)$ est une $B$-algèbre commutative, et
\item il existe un homomorphisme d'anneaux $A \to C$
\end{enumerate}
tel que $F \circ g_{\QCoh, *} : \text{Mod}_A \to \text{Mod}_B$
envoie $M$ sur $M \otimes_A C$ considéré comme $B$-module.
\end{lemma}

\begin{proof}
Remarquons que $g$ est quasi-compact et quasi-séparé ; voir
Morphismes de champs, lemme
\ref{stacks-morphisms-lemma-quasi-compact-quasi-separated-permanence}.
Dans Cohomologie des champs, proposition
\ref{stacks-cohomology-proposition-direct-image-quasi-coherent}
nous avons construit le foncteur
$g_{\QCoh, *} : \QCoh(\mathcal{O}_U) \to \QCoh(\mathcal{O}_\mathcal{X})$.
D'après Cohomologie des champs, remarques
\ref{stacks-cohomology-remark-direct-image-quasi-coherent-tensor} et
\ref{stacks-cohomology-remark-QCoh-tensor}
nous obtenons une multiplication
$$
\mu :
g_{\QCoh, *}\mathcal{O}_U
\otimes_{\mathcal{O}_\mathcal{X}}
g_{\QCoh, *}\mathcal{O}_U
\longrightarrow
g_{\QCoh, *}\mathcal{O}_U
$$
qui munit $g_{\QCoh, *}\mathcal{O}_U$ d'une structure de
$\mathcal{O}_\mathcal{X}$-algèbre commutative. Ainsi, $C = F(g_{\QCoh, *}\mathcal{O}_U)$
est une algèbre commutative dans $\text{Mod}_B$ ; autrement dit, $C$ est une
$B$-algèbre commutative. Remarquons que nous avons un homomorphisme d'anneaux
$\kappa : A \to \text{End}_{\mathcal{O}_\mathcal{X}}(g_{\QCoh, *}\mathcal{O}_U)$
tel que, pour tout $a \in A$, le diagramme
$$
\xymatrix{
g_{\QCoh, *}\mathcal{O}_U \otimes_{\mathcal{O}_\mathcal{X}}
g_{\QCoh, *}\mathcal{O}_U
\ar[d]_{\kappa(a) \otimes 1} \ar[rr]_-\mu & &
g_{\QCoh, *}\mathcal{O}_U \ar[d]^{\kappa(a)} \\
g_{\QCoh, *}\mathcal{O}_U \otimes_{\mathcal{O}_\mathcal{X}}
g_{\QCoh, *}\mathcal{O}_U
\ar[rr]^-\mu & &
g_{\QCoh, *}\mathcal{O}_U
}
$$
est commutatif. Il s'ensuit que l'on obtient un homomorphisme d'anneaux
$\kappa' = F(\kappa) : A \to \text{End}_B(C)$
tel que $\kappa'(a)(c) c' = \kappa'(a)(cc')$ ; bien entendu, cela signifie que
$a \mapsto \kappa'(a)(1)$ est un homomorphisme d'anneaux $A \to C$.

\medskip\noindent
Le morphisme $g : U \to \mathcal{X}$ est affine ; voir
Morphismes de champs, lemme \ref{stacks-morphisms-lemma-affine-permanence}.
Ainsi, $g_{\QCoh, *}$ est exact et commute aux sommes directes d'après
Cohomologie des champs, lemme
\ref{stacks-cohomology-lemma-quasi-coherent-pushforward-affine}.
Ainsi, $F \circ g_{\QCoh, *} : \text{Mod}_A \to \text{Mod}_B$
est un foncteur exact à droite qui commute aux sommes directes
et envoie $A$ sur $C$.
D'après Foncteurs et morphismes, lemme \ref{functors-lemma-functor},
nous voyons que le foncteur $F \circ g_{\QCoh, *}$ envoie un $A$-module
$M$ sur $M \otimes_A C$ considéré comme un $B$-module.
\end{proof}

\begin{lemma}
\label{stacks-more-morphisms-lemma-universally-injective}
Avec les notations du lemme \ref{stacks-more-morphisms-lemma-affine}, supposons que $\mathcal{X}$ soit
noethérien et que $g$ soit surjectif et plat.
Alors $B \to C$ est universellement injectif.
\end{lemma}

\begin{proof}
Considérons le morphisme naturel
$1 : \mathcal{O}_\mathcal{X} \to g_{\QCoh, *}\mathcal{O}_U$
dans $\QCoh(\mathcal{O}_\mathcal{X})$. En prenant son image inverse sur $U$ et
en utilisant l'adjonction, nous constatons que le composé
$$
\mathcal{O}_U =
g^*\mathcal{O}_\mathcal{X} \xrightarrow{g^*1}
g^*g_{\QCoh, *}\mathcal{O}_U \to \mathcal{O}_U
$$
est l'identité dans $\QCoh(\mathcal{O}_U)$.
Écrivons $g_{\QCoh, *}\mathcal{O}_U = \colim \mathcal{F}_i$
sous la forme d'une limite inductive filtrante de $\mathcal{O}_\mathcal{X}$-modules cohérents ; voir
Cohomologie des champs, lemme
\ref{stacks-cohomology-lemma-directed-colimit-coherent}.
Pour $i$ assez grand, le morphisme
$1 : \mathcal{O}_\mathcal{X} \to g_{\QCoh, *}\mathcal{O}_U$
se factorise par $\mathcal{F}_i$ ; voir
Cohomologie des champs, lemme
\ref{stacks-cohomology-lemma-finite-presentation-quasi-compact-colimit}.
Soit $s : \mathcal{O}_\mathcal{X} \to \mathcal{F}_i$ cette
factorisation. Alors
$$
\mathcal{O}_U \xrightarrow{g^*s} g^*\mathcal{F}_i \to
g^*g_{\QCoh, *}\mathcal{O}_U \to \mathcal{O}_U
$$
est l'identité. Autrement dit, nous voyons que $s$ devient
l'inclusion d'un facteur direct après image inverse sur $U$. Posons
$\mathcal{F}_i^\vee = hom(\mathcal{F}_i, \mathcal{O}_\mathcal{X})$
avec les notations de
Cohomologie des champs, lemme
\ref{stacks-cohomology-lemma-internal-hom-fp-into-qcoh}.
En particulier, il existe un morphisme d'évaluation
$ev : \mathcal{F}_i \otimes_{\mathcal{O}_\mathcal{X}} \mathcal{F}_i^\vee
\to \mathcal{O}_\mathcal{X}$.
L'évaluation en $s$ définit un morphisme
$s^\vee : \mathcal{F}_i^\vee \to \mathcal{O}_\mathcal{X}$.
Par dualité avec l'assertion concernant $s$, nous voyons que $g^*(s^\vee)$ est surjectif ;
voir Cohomologie des champs, section
\ref{stacks-cohomology-section-further-remarks}
pour la compatibilité de $hom$ et de $\otimes$ avec la restriction à $U$.
Puisque $g$ est surjectif et plat, nous concluons que $s^\vee$ est surjectif
(voir la référence citée). Puisque $F$ est exact à droite, nous concluons que
$F(\mathcal{F}_i^\vee) \to F(\mathcal{O}_\mathcal{X}) = B$ est surjectif.
Choisissons $\lambda \in F(\mathcal{F}_i^\vee)$
qui s'envoie sur $1 \in B$. Notons $e = F(s)(1) \in F(\mathcal{F}_i)$
l'image de $1$ par le morphisme
$F(s) : B = F(\mathcal{O}_\mathcal{X}) \to F(\mathcal{F}_i)$.
Alors le morphisme
$$
F(ev) :
F(\mathcal{F}_i) \otimes_B F(\mathcal{F}_i^\vee) =
F(\mathcal{F}_i \otimes_{\mathcal{O}_\mathcal{X}} \mathcal{F}_i^\vee)
\longrightarrow
F(\mathcal{O}_\mathcal{X}) = B
$$
envoie $e \otimes \lambda$ sur $1$ par construction. Par conséquent, le
morphisme $B \to F(\mathcal{F}_i)$, $b \mapsto be$,
est universellement injectif, car il admet l'inverse à gauche
$F(\mathcal{F}_i) \to B$, $\xi \mapsto F(ev)(\xi \otimes \lambda)$.
Comme cela vaut pour tout $i$ assez grand, nous concluons.
\end{proof}

\begin{lemma}
\label{stacks-more-morphisms-lemma-flat}
Soit $B \to C$ un homomorphisme d'anneaux. Si
\begin{enumerate}
\item les coprojections $C \to C \otimes_B C$ sont plates, et
\item $B \to C$ est universellement injectif,
\end{enumerate}
alors $B \to C$ est fidèlement plat.
\end{lemma}

\begin{proof}
Le morphisme $\Spec(C) \to \Spec(B)$ est surjectif, puisque $B \to C$ est universellement
injectif. Il suffit donc de montrer que $B \to C$ est plat,
ce qui résulte de
Descente, théorème \ref{descent-theorem-descend-module-properties}.
\end{proof}

\noindent
La version très simple de la formule de K\"unneth qui suit devrait devenir
obsolète lorsque nous aurons rédigé une section sur les théorèmes de K\"unneth pour la
cohomologie des modules quasi-cohérents sur les champs algébriques.

\begin{lemma}
\label{stacks-more-morphisms-lemma-easy-kunneth}
Soient $a : \mathcal{Y} \to \mathcal{X}$ et $b : \mathcal{Z} \to \mathcal{X}$
des morphismes représentables par des schémas, quasi-compacts, quasi-séparés et plats.
Alors $a_{\QCoh, *}\mathcal{O}_\mathcal{Y}
\otimes_{\mathcal{O}_\mathcal{X}}
b_{\QCoh, *}\mathcal{O}_\mathcal{Z} =
f_{\QCoh, *}\mathcal{O}_{\mathcal{Y} \times_\mathcal{X} \mathcal{Z}}$
où $f : \mathcal{Y} \times_\mathcal{X} \mathcal{Z} \to \mathcal{X}$
est le morphisme évident.
\end{lemma}

\begin{proof}
Pour abréger, posons $\mathcal{P} = \mathcal{Y} \times_\mathcal{X} \mathcal{Z}$.
Puisque $a \circ \text{pr}_1 = f$ et $b \circ \text{pr}_2 = f$,
nous obtenons des morphismes $a_*\mathcal{O}_\mathcal{Y} \to f_*\mathcal{O}_\mathcal{P}$
et $b_*\mathcal{O}_\mathcal{Z} \to f_*\mathcal{O}_\mathcal{P}$ (en utilisant
les morphismes d'image inverse relatifs ; voir Sites, section \ref{sites-section-pullback}).
Nous obtenons donc un cup-produit relatif
$$
\mu :
a_*\mathcal{O}_\mathcal{Y}
\otimes_{\mathcal{O}_\mathcal{X}}
b_*\mathcal{O}_\mathcal{Z} \longrightarrow
f_*\mathcal{O}_{\mathcal{Y} \times_\mathcal{X} \mathcal{Z}}
$$
En appliquant $Q$ et sa compatibilité avec les produits tensoriels
(Cohomologie des champs, remarque \ref{stacks-cohomology-remark-QCoh-tensor}),
nous obtenons une flèche
$Q(\mu) : a_{\QCoh, *}\mathcal{O}_\mathcal{Y}
\otimes_{\mathcal{O}_\mathcal{X}}
b_{\QCoh, *}\mathcal{O}_\mathcal{Z} \to
f_{\QCoh, *}\mathcal{O}_{\mathcal{Y} \times_\mathcal{X} \mathcal{Z}}$
dans $\QCoh(\mathcal{O}_\mathcal{X})$.
Choisissons ensuite un schéma $U$ et un morphisme lisse surjectif
$U \to \mathcal{X}$. Il suffit de montrer que la restriction de
$Q(\mu)$ à $U_\etale$ est un isomorphisme ; voir
Cohomologie des champs, section \ref{stacks-cohomology-section-further-remarks}.
À son tour, d'après les résultats de la même section,
il suffit de montrer que la restriction de $\mu$ à $U_\etale$
est un isomorphisme (on utilise ici que la source et le but de $\mu$
sont des modules localement quasi-cohérents vérifiant la propriété de changement de base).
En outre, nous pouvons calculer les images directes pour la topologie \'etale ; voir
Cohomologie des champs, proposition
\ref{stacks-cohomology-proposition-loc-qcoh-flat-base-change}.
Enfin, nous voyons que
$a_*\mathcal{O}_\mathcal{Y}|_{U_\etale} = (V \to U)_{small, *}\mathcal{O}_V$
où $V = U \times_\mathcal{X} \mathcal{Y}$.
De même pour $b_*$ et $f_*$. Ainsi, le résultat découle
de la formule de K\"unneth pour les morphismes plats, quasi-compacts et quasi-séparés
de schémas ; voir Catégories dérivées de schémas,
lemme \ref{perfect-lemma-kunneth}.
\end{proof}

\begin{lemma}
\label{stacks-more-morphisms-lemma-fully-faithful}
Soit $\mathcal{X}$ un champ algébrique à diagonale affine.
Soit $B$ un anneau. Soient $f_i : \Spec(B) \to \mathcal{X}$, $i = 1, 2$,
deux morphismes. Soit $t : f_1^* \to f_2^*$ un isomorphisme
des foncteurs tensoriels
$f_i^* : \QCoh(\mathcal{O}_\mathcal{X}) \to \text{Mod}_B$.
Alors il existe une $2$-flèche $f_1 \to f_2$ qui induit $t$.
\end{lemma}

\begin{proof}
Choisissons un schéma affine $U = \Spec(A)$ et un morphisme lisse surjectif
$g : U \to \mathcal{X}$ ; voir
Propriétés des champs, lemme \ref{stacks-properties-lemma-quasi-compact-stack}.
Puisque la diagonale de $\mathcal{X}$ est affine, nous voyons que
$U_i = \Spec(B) \times_{f_i, \mathcal{X}, g} U$ est affine.
Écrivons $U_i = \Spec(C_i)$. Alors $C_i$ est la $B$-algèbre munie
de l'homomorphisme d'anneaux $A \to C_i$ construit au lemme \ref{stacks-more-morphisms-lemma-affine}
au moyen du foncteur $F = f_i^*$. Par conséquent, $t$ induit un isomorphisme
$C_1 \to C_2$ de $B$-algèbres, compatible avec les homomorphismes d'anneaux
$A \to C_1$ et $A \to C_2$. Autrement dit, nous avons les diagrammes
commutatifs
$$
\xymatrix{
U_i \ar[r] \ar[d] & U \ar[d]^g \\
\Spec(B) \ar[r]^{f_i} & \mathcal{X}
}
\quad\quad
\xymatrix{
&
U_2 \ar[ld] \ar[d]^{\cong} \ar[rd] \\
\Spec(B) &
U_1 \ar[l] \ar[r] &
U
}
$$
Cela montre déjà que les objets $f_1$ et $f_2$ de $\mathcal{X}$
au-dessus de $\Spec(B)$ deviennent isomorphes après le recouvrement lisse
$\{U_1 \to \Spec(B)\}$. Pour montrer que cet isomorphisme descend en un isomorphisme
de $f_1$ et $f_2$ sur $\Spec(B)$, nous devons montrer que notre isomorphisme
(qui provient des diagrammes commutatifs ci-dessus) est compatible
avec les données de descente sur $U_1 \times_{\Spec(B)} U_1$. À cette fin,
remarquons que $U \times_\mathcal{X} U$ est lui aussi affine, que
nous disposons du morphisme $g' : U \times_\mathcal{X} U \to \mathcal{X}$, et
que
$$
U_i \times_{\Spec(B)} U_i =
\Spec(B) \times_{f_i, \mathcal{X}, g'}
(U \times_\mathcal{X} U)
$$
Il s'ensuit que l'isomorphisme $C_1 \otimes_B C_1 \to C_2 \otimes_B C_2$
induit par l'isomorphisme $C_1 \to C_2$ est compatible avec les
morphismes $U_i \times_{\Spec(B)} U_i \to U \times_\mathcal{X} U$.
Certains détails sont omis.
\end{proof}

\begin{lemma}
\label{stacks-more-morphisms-lemma-main}
Soit $\mathcal{X}$ un champ algébrique noethérien à diagonale affine.
Soit $B$ un anneau.
Soit $F : \QCoh(\mathcal{O}_\mathcal{X}) \to \text{Mod}_B$
un foncteur tensoriel exact à droite qui commute aux sommes directes.
Alors $F$ provient d'un unique morphisme $\Spec(B) \to \mathcal{X}$.
\end{lemma}

\begin{proof}
Choisissons un morphisme lisse surjectif $g : U \to \mathcal{X}$, avec $U = \Spec(A)$
affine ; voir Propriétés des champs, lemme
\ref{stacks-properties-lemma-quasi-compact-stack}. Appliquons
le lemme \ref{stacks-more-morphisms-lemma-affine} pour obtenir la
$B$-algèbre commutative de type fini $C = F(g_{\QCoh, *}\mathcal{O}_U)$
et l'homomorphisme d'anneaux $A \to C$.
D'après le lemme \ref{stacks-more-morphisms-lemma-universally-injective},
l'homomorphisme d'anneaux $B \to C$ est universellement injectif.
Considérons l'algèbre
$$
C \otimes_B C =
F(g_{\QCoh, *}\mathcal{O}_U
\otimes_{\mathcal{O}_\mathcal{X}}
g_{\QCoh, *}\mathcal{O}_U)
$$
Puisque $g$ est plat, quasi-compact et quasi-séparé,
le lemme \ref{stacks-more-morphisms-lemma-easy-kunneth} donne la première égalité de
$$
g_{\QCoh, *}\mathcal{O}_U
\otimes_{\mathcal{O}_\mathcal{X}}
g_{\QCoh, *}\mathcal{O}_U
=
f_{\QCoh, *}\mathcal{O}_{U \times_\mathcal{X} U} =
g_{\QCoh, *}(\text{pr}_{2, *}\mathcal{O}_{U \times_\mathcal{X} U})
$$
où
$f : U \times_\mathcal{X} U \to \mathcal{X}$ est le morphisme évident
et $\text{pr}_2 : U \times_\mathcal{X} U \to U$ la seconde
projection. La seconde égalité résulte de
Cohomologie des champs, lemme \ref{stacks-cohomology-lemma-relative-leray},
et de $f = g \circ \text{pr}_2$. Puisque la diagonale de $\mathcal{X}$
est affine, nous voyons que $U \times_\mathcal{X} U = \Spec(R)$ est affine.
Utilisons $\text{pr}_2 : A \to R$ pour considérer $R$ comme une $A$-algèbre.
Au total, nous obtenons
$$
C \otimes_B C =
F(g_{\QCoh, *}\mathcal{O}_U
\otimes_{\mathcal{O}_\mathcal{X}}
g_{\QCoh, *}\mathcal{O}_U) =
F(g_{\QCoh, *}(\text{pr}_{2, *}\mathcal{O}_{U \times_\mathcal{X} U})) =
R \otimes_A C
$$
où la dernière égalité résulte de la dernière assertion du
lemme \ref{stacks-more-morphisms-lemma-affine}. Puisque $A \to R$ est plat (car
$\text{pr}_2$ est plat en tant que changement de base de $U \to \mathcal{X}$),
nous concluons que $C \otimes_B C$ est plat sur $C$.
D'après le lemme \ref{stacks-more-morphisms-lemma-flat}, nous concluons que $B \to C$ est fidèlement plat.

\medskip\noindent
Nous affirmons qu'il existe un diagramme commutatif aux flèches pleines
$$
\xymatrix{
\Spec(C \otimes_B C) \ar@<1ex>[d] \ar@<-1ex>[d] \ar[r] &
U \times_\mathcal{X} U \ar@<1ex>[d] \ar@<-1ex>[d] \\
\Spec(C) \ar[d] \ar[r] &
U \ar[d] \\
\Spec(B) \ar@{..>}[r] &
\mathcal{X}
}
$$
La flèche $\Spec(C) \to U = \Spec(A)$ provient de l'homomorphisme d'anneaux
$A \to C$ de l'énoncé du lemme \ref{stacks-more-morphisms-lemma-affine}.
De même, la flèche $\Spec(C \otimes_B C) \to U \times_\mathcal{X} U$
provient de l'homomorphisme d'anneaux $R \to C \otimes_B C$.
Pour vérifier que le carré supérieur est commutatif, utilisons le lemme \ref{stacks-more-morphisms-lemma-fully-faithful} ;
les détails sont omis.
Nous en concluons que nous obtenons la flèche pointillée $\Spec(B) \to \mathcal{X}$
grâce à la proposition \ref{stacks-more-morphisms-proposition-stack-fpqc}.

\medskip\noindent
L'assertion selon laquelle $F$ est le foncteur image inverse associé
à la flèche pointillée découle elle aussi clairement de ceci et de l'assertion correspondante
du lemme \ref{stacks-more-morphisms-lemma-affine}. Les détails sont omis.
\end{proof}

\noindent
Pour un anneau $B$, notons $\text{Mod}^{fg}_B$ la catégorie des
$B$-modules de type fini (également appelés $B$-modules finis).

\begin{theorem}
\label{stacks-more-morphisms-theorem-main}
Soit $\mathcal{X}$ un champ algébrique noethérien à diagonale affine.
Soit $B$ un anneau noethérien.
Soit $F : \text{Coh}(\mathcal{O}_\mathcal{X}) \to \text{Mod}^{fg}_B$
un foncteur tensoriel exact à droite.
Alors $F$ provient d'un unique morphisme $\Spec(B) \to \mathcal{X}$.
\end{theorem}

\begin{proof}
Le lemme \ref{stacks-more-morphisms-lemma-extend}
permet de prolonger $F$ de manière unique en un foncteur tensoriel exact à droite
$F : \QCoh(\mathcal{O}_\mathcal{X}) \to \text{Mod}_B$
qui commute à toutes les sommes directes. Nous pouvons alors appliquer
le lemme \ref{stacks-more-morphisms-lemma-main}.
\end{proof}














% OK, start here.
%

\chapter{La géométrie des champs algébriques}



\phantomsection
\label{stacks-geometry-section-phantom}


\section{Introduction}
\label{stacks-geometry-section-introduction}

\noindent
Ce chapitre étudie quelques propriétés géométriques des champs algébriques.
Les versions initiales des sections \ref{stacks-geometry-section-multiplicities} et
\ref{stacks-geometry-section-dimension-of-algebraic-stacks}
ont été rédigées par Matthew Emerton et Toby Gee et se trouvent
sous leur forme originale dans \cite{Emerton-Gee-dim}.







\section{Anneaux versels}
\label{stacks-geometry-section-versal}

\noindent
Dans cette section, nous explicitons le lien entre les anneaux de déformation
et les anneaux locaux des champs algébriques de type fini sur une base
localement noethérienne.

\begin{situation}
\label{stacks-geometry-situation-versal}
Ici, $\mathcal{X}$ est un champ algébrique localement de type fini
sur un schéma localement noethérien $S$.
\end{situation}

\noindent
Voici la définition.

\begin{definition}
\label{stacks-geometry-definition-versal-ring-at-x}
Dans la Situation \ref{stacks-geometry-situation-versal}, soit $x_0 : \Spec(k) \to \mathcal{X}$
un morphisme, où $k$ est un corps de type fini sur $S$.
Un {\it anneau versel de $\mathcal{X}$ en $x_0$} est
une $S$-algèbre locale noethérienne complète $A$ de corps résiduel $k$,
pour laquelle il existe un objet formel versel
$(A, \xi_n, f_n)$ comme dans Axiomes d'Artin, Définition
\ref{artin-definition-versal-formal-object},
avec $\xi_1 \cong x_0$ (une $2$-isomorphie).
\end{definition}

\noindent
Nous voulons démontrer que les anneaux versels existent et sont uniques à
facteurs lisses près. Nous utiliserons pour cela les catégories de prédéformations de
Axiomes d'Artin, section \ref{artin-section-predeformation-categories}.
Dans notre situation, ce sont toujours des catégories de déformations.

\begin{lemma}
\label{stacks-geometry-lemma-deformation-category}
Dans la Situation \ref{stacks-geometry-situation-versal}, soit $x_0 : \Spec(k) \to \mathcal{X}$
un morphisme, où $k$ est un corps de type fini sur $S$.
Alors $\mathcal{F}_{\mathcal{X}, k, x_0}$
est une catégorie de déformations, et $T\mathcal{F}_{\mathcal{X}, k, x_0}$
ainsi que $\text{Inf}(\mathcal{F}_{\mathcal{X}, k, x_0})$
sont des espaces vectoriels sur $k$ de dimension finie.
\end{lemma}

\begin{proof}
Choisissons un ouvert affine $\Spec(\Lambda) \subset S$ tel que
$\Spec(k) \to S$ se factorise par celui-ci.
D'après Axiomes d'Artin, section \ref{artin-section-predeformation-categories},
nous obtenons une catégorie de prédéformations
$\mathcal{F}_{\mathcal{X}, k, x_0}$
sur la catégorie $\mathcal{C}_\Lambda$.
(Comme il est signalé au passage cité, cette catégorie ne dépend que
du morphisme $\Spec(k) \to S$, et non du choix de
$\Lambda$.) D'après Axiomes d'Artin, Lemmes
\ref{artin-lemma-deformation-category} et
\ref{artin-lemma-algebraic-stack-RS}
$\mathcal{F}_{\mathcal{X}, k, x_0}$ est en fait une catégorie de déformations.
D'après Axiomes d'Artin, Lemme \ref{artin-lemma-finite-dimension},
nous constatons que $T\mathcal{F}_{\mathcal{X}, k, x_0}$
et $\text{Inf}(\mathcal{F}_{\mathcal{X}, k, x_0})$
sont des espaces vectoriels sur $k$ de dimension finie.
\end{proof}

\begin{lemma}
\label{stacks-geometry-lemma-versal-ring}
Dans la Situation \ref{stacks-geometry-situation-versal}, soit $x_0 : \Spec(k) \to \mathcal{X}$
un morphisme, où $k$ est un corps de type fini sur $S$.
Il existe alors un anneau versel de $\mathcal{X}$ en $x_0$. Si
$A$, $A'$ en sont deux, alors $A \cong A'[[t_1, \ldots, t_r]]$
ou $A' \cong A[[t_1, \ldots, t_r]]$ comme $S$-algèbres
pour un certain $r$.
\end{lemma}

\begin{proof}
L'existence résulte du Lemme \ref{stacks-geometry-lemma-deformation-category} et de
Théorie formelle des déformations,
Définition \ref{formal-defos-definition-deformation-category} et
Lemmes \ref{formal-defos-lemma-RS-implies-S1-S2} et
\ref{formal-defos-lemma-versal-object-existence}.
D'après le résultat d'unicité de
Théorie formelle des déformations, Lemme \ref{formal-defos-lemma-minimal-versal},
il existe un anneau versel « minimal » $A$ de $\mathcal{X}$ en $x_0$
tel que tout autre anneau versel de $\mathcal{X}$ en $x_0$ soit
isomorphe à $A[[t_1, \ldots, t_r]]$ pour un certain $r$.
Ceci entraîne manifestement la seconde assertion.
\end{proof}

\begin{lemma}
\label{stacks-geometry-lemma-versal-ring-field-extension}
Dans la Situation \ref{stacks-geometry-situation-versal}, soit $x_0 : \Spec(k) \to \mathcal{X}$
un morphisme, où $k$ est un corps de type fini sur $S$.
Soit $l/k$ une extension finie de corps, et notons
$x_{l, 0} : \Spec(l) \to \mathcal{X}$ le morphisme induit.
Étant donné un anneau versel $A$ de $\mathcal{X}$ en $x_0$, il existe
un anneau versel $A'$ de $\mathcal{X}$ en $x_{l, 0}$ tel
qu'il existe un homomorphisme de $S$-algèbres $A \to A'$ induisant
l'extension de corps donnée $l/k$ et formellement lisse pour la topologie $\mathfrak m_{A'}$-adique
considérée.
\end{lemma}

\begin{proof}
Cela résulte immédiatement de
Axiomes d'Artin, Lemme \ref{artin-lemma-change-of-field},
et de
Théorie formelle des déformations, Lemme
\ref{formal-defos-lemma-change-of-field-versal-ring}.
(Nous utilisons aussi le fait que $\mathcal{X}$ satisfait (RS), d'après
Axiomes d'Artin, Lemme \ref{artin-lemma-algebraic-stack-RS}.)
\end{proof}

\begin{lemma}
\label{stacks-geometry-lemma-compare-versal-ring-completion}
Dans la Situation \ref{stacks-geometry-situation-versal}, soit $x : U \to \mathcal{X}$ un
morphisme, où $U$ est un schéma localement de type fini sur $S$.
Soit $u_0 \in U$ un point de type fini.
Posons $k = \kappa(u_0)$ et notons $x_0 : \Spec(k) \to \mathcal{X}$
l'application induite. Les conditions suivantes sont équivalentes :
\begin{enumerate}
\item $x$ est versel en $u_0$
(Axiomes d'Artin, Définition \ref{artin-definition-versal}),
\item $\hat x : \mathcal{F}_{U, k, u_0} \to \mathcal{F}_{\mathcal{X}, k, x_0}$
est lisse,
\item l'objet formel associé à
$x|_{\Spec(\mathcal{O}_{U, u_0}^\wedge)}$ est versel, et
\item il existe un voisinage ouvert $U' \subset U$ de $x$ tel que
$x|_{U'} : U' \to \mathcal{X}$ soit lisse.
\end{enumerate}
De plus, dans ce cas, le complété $\mathcal{O}_{U, u_0}^\wedge$
est un anneau versel de $\mathcal{X}$ en $x_0$.
\end{lemma}

\begin{proof}
Puisque $U \to S$ est localement de type fini (comme composé de tels morphismes),
nous voyons que $\Spec(k) \to S$ est de type fini (encore comme composé).
L'énoncé a donc un sens. L'équivalence de (1) et (2)
est la définition du fait que $x$ est versel en $u_0$.
L'équivalence de (1) et (3) est
Axiomes d'Artin, Lemme \ref{artin-lemma-versality-matches}.
Ainsi, (1), (2) et (3) sont équivalentes.

\medskip\noindent
Si $x|_{U'}$ est lisse, alors le foncteur
$\hat x : \mathcal{F}_{U, k, u_0} \to \mathcal{F}_{\mathcal{X}, k, x_0}$
est lisse d'après Axiomes d'Artin, Lemme
\ref{artin-lemma-formally-smooth-on-deformation-categories}.
Ainsi, (4) entraîne (1), (2) et (3).
Réciproquement, supposons que $x$ soit versel en $u_0$.
Choisissons un morphisme lisse surjectif $y : V \to \mathcal{X}$, où $V$
est un schéma. Posons $Z = V \times_\mathcal{X} U$ et choisissons un
point de type fini $z_0 \in |Z|$ au-dessus de $u_0$ (cela est possible d'après
Morphismes d'espaces, Lemme
\ref{spaces-morphisms-lemma-finite-type-points-surjective-morphism}).
D'après Axiomes d'Artin, Lemme \ref{artin-lemma-base-change-versal},
le morphisme $Z \to V$ est lisse en $z_0$.
Par définition, nous pouvons trouver un voisinage ouvert $W \subset Z$
de $z_0$ tel que $W \to V$ soit lisse. Puisque $Z \to U$ est ouvert,
soit $U' \subset U$ l'image de $W$. Nous voyons alors que
$U' \to \mathcal{X}$ est lisse par notre définition des morphismes lisses
de champs.

\medskip\noindent
La dernière assertion résulte des définitions, puisque
$\mathcal{O}_{U, u_0}^\wedge$
proreprésente $\mathcal{F}_{U, k, u_0}$.
\end{proof}

\begin{lemma}
\label{stacks-geometry-lemma-characterize-smoothness}
Dans la Situation \ref{stacks-geometry-situation-versal}. Soit $x_0 : \Spec(k) \to \mathcal{X}$
un morphisme tel que $\Spec(k) \to S$ soit de type fini, d'image $s$.
Soit $A$ un anneau versel de $\mathcal{X}$ en $x_0$. Les conditions suivantes
sont équivalentes :
\begin{enumerate}
\item $x_0$ appartient au lieu lisse de $\mathcal{X} \to S$
(Morphismes de champs, Lemme \ref{stacks-morphisms-lemma-where-smooth}),
\item $\mathcal{O}_{S, s} \to A$ est formellement lisse pour la
topologie $\mathfrak m_A$-adique, et
\item $\mathcal{F}_{\mathcal{X}, k, x_0}$ est sans obstruction.
\end{enumerate}
\end{lemma}

\begin{proof}
L'équivalence de (2) et (3) résulte immédiatement de
Théorie formelle des déformations, Lemme
\ref{formal-defos-lemma-smooth-power-series-classical}.

\medskip\noindent
Notons que $\mathcal{O}_{S, s} \to A$ est formellement lisse pour la
topologie $\mathfrak m_A$-adique si et seulement si
$\mathcal{O}_{S, s} \to A' = A[[t_1, \ldots, t_r]]$
est formellement lisse pour la topologie $\mathfrak m_{A'}$-adique.
Par conséquent, (2) ne dépend pas du choix
de l'anneau versel, d'après le Lemme \ref{stacks-geometry-lemma-versal-ring}.
Ensuite, soit $l/k$ une extension finie et choisissons
$A \to A'$ comme dans le Lemme \ref{stacks-geometry-lemma-versal-ring-field-extension}.
Si $\mathcal{O}_{S, s} \to A$ est formellement lisse pour la
topologie $\mathfrak m_A$-adique, alors $\mathcal{O}_{S, s} \to A'$
est formellement lisse pour la topologie $\mathfrak m_{A'}$-adique, voir
Compléments d'algèbre, Lemme \ref{more-algebra-lemma-compose-formally-smooth}.
Réciproquement, si $\mathcal{O}_{S, s} \to A'$
est formellement lisse pour la topologie $\mathfrak m_{A'}$-adique,
alors $\mathcal{O}_{S, s}^\wedge \to A'$ et $A \to A'$ sont réguliers
(Compléments d'algèbre, Proposition \ref{more-algebra-proposition-fs-regular}),
donc $\mathcal{O}_{S, s}^\wedge \to A$ est régulier
(Compléments d'algèbre, Lemme \ref{more-algebra-lemma-regular-permanence}),
et donc $\mathcal{O}_{S, s} \to A$ est formellement lisse pour la
topologie $\mathfrak m_A$-adique (même lemme que précédemment).
Ainsi, l'équivalence de (2) et (1) vaut pour
$k$ et $x_0$ si et seulement si elle vaut pour $l$ et $x_{0, l}$.

\medskip\noindent
Choisissons un schéma $U$ et un morphisme lisse $U \to \mathcal{X}$ tels
que $\Spec(k) \times_\mathcal{X} U$ soit non vide. Choisissons une extension
finie $l/k$ et un point $w_0 : \Spec(l) \to \Spec(k) \times_\mathcal{X} U$.
Soit $u_0 \in U$ l'image de $w_0$.
Nous pouvons appliquer ce qui précède à $l/k$ et à $l/\kappa(u_0)$
pour voir que nous pouvons nous ramener à $u_0$. Nous pouvons donc supposer
$A = \mathcal{O}_{U, u_0}^\wedge$, voir le
Lemme \ref{stacks-geometry-lemma-compare-versal-ring-completion}.
Observons que $x_0$ appartient au lieu lisse de $\mathcal{X} \to S$
si et seulement si $u_0$ appartient au lieu lisse de $U \to S$, voir par exemple
Morphismes de champs, Lemme \ref{stacks-morphisms-lemma-where-smooth}.
L'équivalence de (1) et (2) résulte donc de
Compléments d'algèbre, Lemme \ref{more-algebra-lemma-formally-smooth-finite-type}.
\end{proof}

\noindent
Rappelons une conséquence de l'approximation d'Artin.

\begin{lemma}
\label{stacks-geometry-lemma-Artin-approximation-by-smooth-morphism}
Dans la Situation \ref{stacks-geometry-situation-versal}. Soit $x_0 : \Spec(k) \to \mathcal{X}$
un morphisme tel que $\Spec(k) \to S$ soit de type fini, d'image $s$.
Soit $A$ un anneau versel de $\mathcal{X}$ en $x_0$.
Si $\mathcal{O}_{S, s}$ est un anneau G, nous pouvons trouver un morphisme lisse
$U \to \mathcal{X}$ dont la source est un schéma, et un point
$u_0 \in U$ de corps résiduel $k$, tels que
\begin{enumerate}
\item $\Spec(k) \to U \to \mathcal{X}$ coïncide avec le morphisme donné $x_0$,
\item il existe un isomorphisme $\mathcal{O}_{U, u_0}^\wedge \cong A$.
\end{enumerate}
\end{lemma}

\begin{proof}
Soit $(\xi_n, f_n)$ l'objet formel versel sur $A$.
D'après Axiomes d'Artin, Lemme \ref{artin-lemma-effective},
nous savons que $\xi = (A, \xi_n, f_n)$ est effectif.
Par hypothèse, $\mathcal{X}$ est localement de présentation finie sur $S$
(utiliser Morphismes de champs, Lemme
\ref{stacks-morphisms-lemma-noetherian-finite-type-finite-presentation}),
et préserve donc les limites d'après Limites de champs, Proposition
\ref{stacks-limits-proposition-characterize-locally-finite-presentation}.
Ainsi, l'approximation d'Artin, sous la forme de
Axiomes d'Artin, Lemme \ref{artin-lemma-approximate-versal},
montre que nous pouvons trouver un morphisme $U \to \mathcal{X}$ dont
la source est un $S$-schéma de type fini, contenant un point $u_0 \in U$
de corps résiduel $k$ satisfaisant (1) et (2), tel que $U \to \mathcal{X}$
soit versel en $u_0$. D'après le Lemme \ref{stacks-geometry-lemma-compare-versal-ring-completion},
quitte à rétrécir $U$, nous pouvons supposer que $U \to \mathcal{X}$ est lisse.
\end{proof}

\begin{remark}[Amélioration des anneaux versels]
\label{stacks-geometry-remark-upgrade}
Dans la Situation \ref{stacks-geometry-situation-versal}, soit $x_0 : \Spec(k) \to \mathcal{X}$
un morphisme, où $k$ est un corps de type fini sur $S$.
Soit $A$ un anneau versel de $\mathcal{X}$ en $x_0$. D'après Axiomes d'Artin,
Lemme \ref{artin-lemma-effective}, notre objet formel versel
provient en fait d'un morphisme
$$
\Spec(A) \longrightarrow \mathcal{X}
$$
sur $S$. De plus, chacun des résultats précédents peut être renforcé de façon
à être compatible avec ce morphisme. En voici la liste :
\begin{enumerate}
\item dans le Lemme \ref{stacks-geometry-lemma-versal-ring}, l'isomorphisme
$A \cong A'[[t_1, \ldots, t_r]]$ ou $A' \cong A[[t_1, \ldots, t_r]]$
peut être choisi compatible avec ces morphismes,
\item dans le Lemme \ref{stacks-geometry-lemma-versal-ring-field-extension},
l'homomorphisme $A \to A'$ peut être choisi compatible avec ces morphismes,
\item dans le Lemme \ref{stacks-geometry-lemma-compare-versal-ring-completion},
le morphisme $\Spec(\mathcal{O}_{U, u_0}^\wedge) \to \mathcal{X}$
est le composé de l'application canonique
$\Spec(\mathcal{O}_{U, u_0}^\wedge) \to U$ et de l'application donnée
$U \to \mathcal{X}$,
\item dans le Lemme \ref{stacks-geometry-lemma-Artin-approximation-by-smooth-morphism},
l'isomorphisme $\mathcal{O}_{U, u_0}^\wedge \cong A$ peut
être choisi de telle sorte que $\Spec(A) \to \mathcal{X}$ corresponde à l'application canonique
de l'alinéa précédent.
\end{enumerate}
Dans chaque cas, l'assertion résulte du fait que nos applications sont
compatibles avec les éléments formels versels; notons toutefois que les
diagrammes sous-entendus ne sont $2$-commutatifs qu'au choix (non canonique)
d'une $2$-flèche près. Cela signifie néanmoins que l'application $A' \to A$
ou $A \to A'$ sous-entendue dans (1) est bien définie à homotopie formelle près, voir
Théorie formelle des déformations, Lemme
\ref{formal-defos-lemma-versal-unique-up-to-homotopy}.
\end{remark}

\begin{lemma}
\label{stacks-geometry-lemma-versal-ring-flat}
Dans la Situation \ref{stacks-geometry-situation-versal}, soit $x_0 : \Spec(k) \to \mathcal{X}$
un morphisme, où $k$ est un corps de type fini sur $S$.
Soit $A$ un anneau versel de $\mathcal{X}$ en $x_0$.
Alors le morphisme $\Spec(A) \to \mathcal{X}$ de la
Remarque \ref{stacks-geometry-remark-upgrade} est plat.
\end{lemma}

\begin{proof}
Si l'anneau local de $S$ au point image est un anneau G, cela
résulte immédiatement du
Lemme \ref{stacks-geometry-lemma-Artin-approximation-by-smooth-morphism}
et du fait que l'application d'un anneau local noethérien dans son
complété est plate. Dans le cas général, nous procédons comme suit.

\medskip\noindent
Étape I. Si $A$ et $A'$ sont deux anneaux versels de $\mathcal{X}$ en $x_0$,
le résultat est vrai pour $A$ si et seulement s'il l'est pour $A'$.
En effet, quitte à échanger $A$ et $A'$, nous pouvons supposer qu'il existe
une application formellement lisse $\varphi : A \to A'$ telle que le composé
$$
\Spec(A') \to \Spec(A) \to \mathcal{X}
$$
soit le morphisme $\Spec(A') \to \mathcal{X}$, voir le
Lemme \ref{stacks-geometry-lemma-versal-ring} et la
Remarque \ref{stacks-geometry-remark-upgrade}.
Puisque $A \to A'$ est fidèlement plat, l'équivalence résulte de
Morphismes de champs, Lemmes \ref{stacks-morphisms-lemma-composition-flat} et
\ref{stacks-morphisms-lemma-flat-permanence}.

\medskip\noindent
Étape II. Soit $l/k$ une extension finie de corps. Soit
$x_{l, 0} : \Spec(l) \to \mathcal{X}$ le morphisme induit.
Soit $A$ un anneau versel de $\mathcal{X}$ en $x_0$ et
soit $A \to A'$ comme dans le Lemme \ref{stacks-geometry-lemma-versal-ring-field-extension}.
À nouveau, le composé
$$
\Spec(A') \to \Spec(A) \to \mathcal{X}
$$
est le morphisme $\Spec(A') \to \mathcal{X}$, voir la Remarque \ref{stacks-geometry-remark-upgrade}.
En raisonnant comme précédemment et en utilisant l'étape I pour voir que le choix des anneaux versels
est sans incidence, nous voyons que le lemme vaut pour $x_0$ si et seulement
s'il vaut pour $x_{l, 0}$.

\medskip\noindent
Étape III. Choisissons un schéma $U$ et un morphisme lisse surjectif
$U \to \mathcal{X}$. Nous pouvons alors choisir un point de type fini
$z_0$ de $Z = U \times_\mathcal{X} x_0$ (c'est un espace algébrique
non vide). Soit $u_0 \in U$ l'image de $z_0$ dans $U$.
Choisissons un schéma et une application étale surjective $W \to Z$
tels que $z_0$ soit l'image d'un point fermé $w_0 \in W$
(voir Morphismes d'espaces, section
\ref{spaces-morphisms-section-points-finite-type}).
Puisque $W \to \Spec(k)$ et $W \to U$ sont de type fini,
nous voyons que $\kappa(w_0)/k$ et $\kappa(w_0)/\kappa(u_0)$
sont des extensions finies de corps
(voir Morphismes, section \ref{morphisms-section-points-finite-type}).
En appliquant deux fois l'étape II, nous pouvons remplacer $x_0$ par
$u_0 \to U \to \mathcal{X}$. Nous voyons alors que notre morphisme est le composé
$$
\Spec(\mathcal{O}_{U, u_0}^\wedge) \to U \to \mathcal{X}
$$
La première flèche est plate parce que les complétions d'anneaux locaux noethériens
sont plates (Algèbre, Lemme \ref{algebra-lemma-completion-flat}),
et la seconde flèche est plate puisque
les morphismes lisses sont plats. Le composé est plat car la composition
préserve la platitude.
\end{proof}

\begin{remark}
\label{stacks-geometry-remark-groupoid-defo}
Dans la Situation \ref{stacks-geometry-situation-versal}, soit $x_0 : \Spec(k) \to \mathcal{X}$
un morphisme, où $k$ est un corps de type fini sur $S$.
D'après le Lemme \ref{stacks-geometry-lemma-deformation-category} et
Théorie formelle des déformations, Théorème
\ref{formal-defos-theorem-presentation-deformation-groupoid}
nous savons que $\mathcal{F}_{\mathcal{X}, k, x_0}$ admet une
présentation par un groupoïde lisse proreprésentable en
foncteurs sur $\mathcal{C}_\Lambda$.
En explicitant les définitions, cela signifie que nous pouvons choisir
\begin{enumerate}
\item une $\Lambda$-algèbre locale noethérienne complète $A$
de corps résiduel $k$, et un objet formel versel $\xi$
de $\mathcal{F}_{\mathcal{X}, k, x_0}$ sur $A$,
\item une $\Lambda$-algèbre locale noethérienne complète $B$
de corps résiduel $k$, et un isomorphisme
$$
\underline{B}|_{\mathcal{C}_\Lambda}
\longrightarrow
\underline{A}|_{\mathcal{C}_\Lambda}
\times_{\underline{\xi}, \mathcal{F}_{\mathcal{X}, k, x_0}, \underline{\xi}}
\underline{A}|_{\mathcal{C}_\Lambda}
$$
\end{enumerate}
Les projections correspondent à des applications formellement lisses
$t : A \to B$ et $s : A \to B$ (car $\xi$ est versel).
Il existe une application $c : B \to B \widehat{\otimes}_{s, A, t} B$
qui fait de $(A, B, s, t, c)$ un cogroupoïde dans la catégorie
des $\Lambda$-algèbres locales noethériennes complètes de corps résiduel $k$
(sur les foncteurs proreprésentables, cette application est construite dans
Théorie formelle des déformations, Lemme
\ref{formal-defos-lemma-presentation-construction}).
Enfin, le théorème cité nous dit que $\xi$ induit
une équivalence
$$
[\underline{A}|_{\mathcal{C}_\Lambda} / \underline{B}|_{\mathcal{C}_\Lambda}]
\longrightarrow
\mathcal{F}_{\mathcal{X}, k, x_0}
$$
de groupoïdes cofibrés sur $\mathcal{C}_\Lambda$. En fait, nous obtenons aussi
une équivalence
$$
[\underline{A}/\underline{B}]
\longrightarrow
\widehat{\mathcal{F}}_{\mathcal{X}, k, x_0}
$$
de groupoïdes cofibrés sur la catégorie complétée
$\widehat{\mathcal{C}}_\Lambda$ (voir la discussion dans
Théorie formelle des déformations, section
\ref{formal-defos-section-prorepresentable-groupoids-in-functors}
qui explique pourquoi cela fonctionne). Bien entendu, $A$ est un anneau versel de
$\mathcal{X}$ en $x_0$.
\end{remark}








\section{Multiplicités des composantes des champs algébriques}
\label{stacks-geometry-section-multiplicities}

\noindent
Si $X$ est un schéma localement noethérien, nous pouvons écrire $X$ (considéré
simplement comme espace topologique) comme réunion de composantes irréductibles,
disons $X = \bigcup T_i$. Chaque composante irréductible est l'adhérence d'un
unique point générique $\xi_i$, et l'anneau local $\mathcal O_{X,\xi_i}$
est un anneau artinien local. Nous pouvons définir la {\it multiplicité de $X$ le long de $T_i$}
ou la {\it multiplicité de $T_i$ dans $X$} par
$$
m_{T_i, X} = \text{length}_{\mathcal O_{X, \xi_i}} \mathcal O_{X, \xi_i}
$$
Autrement dit, c'est la longueur de l'anneau artinien local. Comparer
avec
Homologie de Chow, section \ref{chow-section-cycle-of-closed-subscheme}.

\medskip\noindent
Notre but est ici de généraliser cette définition aux champs algébriques
localement noethériens. Si $\mathcal{X}$ est un champ,
alors son espace topologique $|\mathcal{X}|$
(voir Propriétés des champs, Définition
\ref{stacks-properties-definition-topological-space})
est localement noethérien
(Morphismes de champs, Lemme \ref{stacks-morphisms-lemma-Noetherian-topology}).
Les composantes irréductibles de $|\mathcal{X}|$ sont parfois
appelées les composantes irréductibles de $\mathcal{X}$.
Si $\mathcal{X}$ est quasi-séparé, alors $|\mathcal{X}|$
est sobre (Morphismes de champs, Lemme
\ref{stacks-morphisms-lemma-sober-qs}),
mais il ne l'est pas nécessairement dans le cas
non quasi-séparé. Considérons par exemple l'espace algébrique non quasi-séparé
$X = \mathbf{A}^1_\mathbf{C}/\mathbf{Z}$.
En outre, il n'existe pas de faisceau structural
sur $|\mathcal{X}|$ dont les fibres puissent servir à définir les multiplicités.

\begin{lemma}
\label{stacks-geometry-lemma-map-of-components}
Soit $f : U \to \mathcal{X}$ un morphisme lisse d'un schéma
vers un champ algébrique localement noethérien. L'adhérence de l'image de toute
composante irréductible de $|U|$ est une composante irréductible de $|\mathcal{X}|$.
Si $U \to \mathcal{X}$ est surjectif, toutes les composantes irréductibles de
$|\mathcal{X}|$ s'obtiennent ainsi.
\end{lemma}

\begin{proof}
L'application $|U| \to |\mathcal{X}|$ est continue et ouverte d'après
Propriétés des champs, Lemme \ref{stacks-properties-lemma-topology-points}.
Soit $T \subset |U|$ une composante irréductible. Puisque $U$ est localement
noethérien, nous pouvons trouver un ouvert affine non vide $W \subset U$ contenu dans $T$.
Alors $f(T) \subset |\mathcal{X}|$ est irréductible et contient le
sous-ensemble ouvert non vide $f(W)$. Ainsi, l'adhérence de $f(T)$ est irréductible et
contient un ouvert non vide. Il s'ensuit que cette adhérence est une composante
irréductible.

\medskip\noindent
Supposons $U \to \mathcal{X}$ surjectif et soit $Z \subset |\mathcal{X}|$
une composante irréductible. Choisissons un sous-ensemble ouvert noethérien $V$
de $|\mathcal{X}|$ rencontrant $Z$. Après avoir retiré les autres composantes
irréductibles de $V$, nous pouvons supposer que $V \subset Z$.
Prenons une composante irréductible de l'ouvert non vide
$f^{-1}(V) \subset |U|$ et soit $T \subset |U|$ son adhérence.
C'est une composante irréductible de $|U|$, et l'adhérence de $f(T)$
doit coïncider avec $Z$ par notre choix de $T$.
\end{proof}

\noindent
Le lemme précédent s'applique en particulier aux morphismes lisses
entre schémas localement noethériens. Ce cas particulier est
invoqué implicitement dans l'énoncé du lemme suivant.

\begin{lemma}
\label{stacks-geometry-lemma-multiplicities}
Soit $U \to X$ un morphisme lisse de schémas localement noethériens.
Soit $T'$ une composante irréductible de $U$. Soit $T$ la
composante irréductible de $X$ obtenue comme adhérence de
l'image de $T'$. Alors $m_{T', U} = m_{T, X}$.
\end{lemma}

\begin{proof}
Notons $\xi'$ le point générique de $T'$, et $\xi$ le
point générique de $T$. Soient $A = \mathcal{O}_{X, \xi}$ et
$B = \mathcal{O}_{U, \xi'}$. Nous devons montrer que
$\text{length}_A A = \text{length}_B B$. Puisque
$A \to B$ est un homomorphisme local plat d'anneaux
(car les morphismes lisses sont plats), nous avons
$$
\text{length}_A(A) \text{length}_B(B/\mathfrak m_A B) =
\text{length}_B(B)
$$
d'après Algèbre, Lemme \ref{algebra-lemma-pullback-module}. Il suffit donc
de montrer que $\mathfrak m_A B = \mathfrak m_B$, ou, de manière équivalente, que
$B/\mathfrak m_A B$ est réduit. Puisque $U \to X$ est lisse,
son changement de base $U_{\xi} \to \Spec \kappa(\xi)$ l'est aussi. Comme $U_{\xi}$ est un
schéma lisse sur un corps, il est réduit, et son anneau local
en tout point l'est donc aussi
(Variétés, Lemme \ref{varieties-lemma-smooth-geometrically-normal}).
En particulier,
$$
B/\mathfrak m_A B =
\mathcal{O}_{U, \xi'}/\mathfrak m_{X, \xi}\mathcal{O}_{U, \xi'} =
\mathcal{O}_{U_\xi, \xi'}
$$
est réduit, comme voulu.
\end{proof}

\noindent
À l'aide de ce résultat, nous pouvons montrer qu'il existe une bonne notion
de multiplicité en nous plaçant localement pour la topologie lisse.

\begin{lemma}
\label{stacks-geometry-lemma-multiplicity}
Soient $U_1 \to \mathcal{X}$ et $U_2 \to \mathcal{X}$ deux morphismes lisses
de schémas vers un champ algébrique localement noethérien $\mathcal{X}$.
Soient $T_1'$ et $T_2'$ des composantes irréductibles de $|U_1|$
et $|U_2|$, respectivement. Supposons que les adhérences des images de
$T_1'$ et $T_2'$ soient la même composante irréductible $T$ de $|\mathcal{X}|$.
Alors $m_{T_1', U_1} = m_{T_2', U_2}$.
\end{lemma}

\begin{proof}
Soient $V_1$ et $V_2$ des sous-ensembles denses de $T_1'$ et $T'_2$, respectivement,
qui soient ouverts dans $U_1$ et $U_2$, respectivement (voir la preuve du
Lemme \ref{stacks-geometry-lemma-map-of-components}).
Les images de $|V_1|$ et $|V_2|$ dans $|\mathcal{X}|$ sont des sous-ensembles ouverts
non vides de l'ensemble irréductible $T$, et ont donc une intersection
non vide. D'après
Propriétés des champs, Lemme \ref{stacks-properties-lemma-points-cartesian},
l'application $|V_1 \times_\mathcal{X} V_2| \to |V_1| \times_{|\mathcal{X}|} |V_2|$
est surjective. Par conséquent, $V_1 \times_\mathcal{X} V_2$
est un espace algébrique non vide; nous pouvons donc choisir une
surjection étale $V \to V_1 \times_\mathcal{X} V_2$
dont la source est un schéma (non vide).
Si $T'$ est une composante irréductible quelconque de $V$,
alors le Lemme \ref{stacks-geometry-lemma-map-of-components} montre que l'adhérence de
l'image de $T'$ dans $U_1$ (respectivement dans $U_2$) est égale à $T'_1$
(respectivement à $T'_2$).

\medskip\noindent
En appliquant deux fois le Lemme \ref{stacks-geometry-lemma-multiplicities}, nous trouvons
que
$$
m_{T_1', U_1} = m_{T', V} = m_{T_2', U_2},
$$
comme voulu.
\end{proof}

\noindent
Nous en avons maintenant fait assez pour montrer que la définition
suivante a un sens.

\begin{definition}
\label{stacks-geometry-definition-multiplicity}
Soit $\mathcal{X}$ un champ algébrique localement noethérien. Soit
$T \subset |\mathcal{X}|$ une composante irréductible.
La {\it multiplicité} de $T$ dans $\mathcal{X}$ est définie par
$m_{T, \mathcal{X}} = m_{T', U}$, où $f : U \to \mathcal{X}$
est un morphisme lisse d'un schéma et $T' \subset |U|$
est une composante irréductible telle que $f(T') \subset T$.
\end{definition}

\noindent
Ceci est indépendant du choix de $f : U \to \mathcal{X}$
et de celui de la composante irréductible $T'$ envoyée dans $T$,
d'après les Lemmes \ref{stacks-geometry-lemma-map-of-components} et \ref{stacks-geometry-lemma-multiplicity}.

\medskip\noindent
Pour terminer, notons qu'il est parfois commode de considérer
une composante irréductible de $\mathcal{X}$ comme un sous-champ fermé.
À cette fin, si $T \subset |\mathcal{X}|$ est une composante irréductible,
nous pouvons considérer l'unique sous-champ fermé réduit
$\mathcal{T} \subset \mathcal{X}$ tel que $|\mathcal{T}| = T$, voir
Propriétés des champs, Définition
\ref{stacks-properties-definition-reduced-induced-stack}.
Si $\mathcal{X}$ est quasi-séparé, alors
une composante irréductible est un champ intègre; voir
Morphismes de champs, section \ref{stacks-morphisms-section-integral-stacks},
pour plus de détails.



\section{Branches formelles et multiplicités}
\label{stacks-geometry-section-formal-branches}

\noindent
Il sera commode de comparer la notion de multiplicité
d'une composante irréductible donnée par la Définition \ref{stacks-geometry-definition-multiplicity}
à la notion connexe de multiplicité des composantes irréductibles
(des spectres) des anneaux versels de $\mathcal{X}$ aux points de type fini.

\medskip\noindent
Dans la Situation \ref{stacks-geometry-situation-versal}, soit $x_0 : \Spec(k) \to \mathcal{X}$
un morphisme, où $k$ est un corps de type fini sur $S$.
Soient $A$, $A'$ des anneaux versels de $\mathcal{X}$ en $x_0$.
Quitte à échanger $A$ et $A'$, nous savons qu'il existe une application
formellement lisse\footnote{Au sens où
$A'$ devient isomorphe à un anneau de séries formelles sur $A$.}
$\varphi : A \to A'$ compatible avec les objets formels versels, voir le
Lemme \ref{stacks-geometry-lemma-versal-ring} et la
Remarque \ref{stacks-geometry-remark-upgrade}.
De plus, $\varphi$ est bien définie à homotopie formelle près, voir
Théorie formelle des déformations, Lemme
\ref{formal-defos-lemma-versal-unique-up-to-homotopy}.
En particulier, nous constatons que $\varphi(\mathfrak p)A'$ est un idéal
bien défini de $A'$, d'après
Théorie formelle des déformations, Lemme
\ref{formal-defos-lemma-homotopic-minimal-prime}.
Puisque $A \to A'$ est formellement lisse,
$\varphi(\mathfrak p)A'$ est en fait un idéal premier minimal de $A'$,
et tout idéal premier minimal de $A'$ est de cette forme pour
un unique idéal premier minimal $\mathfrak p \subset A$ (tout cela
se démontre facilement en écrivant $A'$ comme anneau de séries formelles sur $A$).
Ainsi, puisque les idéaux premiers minimaux correspondent
aux composantes irréductibles, la définition suivante a un sens.

\begin{definition}
\label{stacks-geometry-definition-formal-branches}
Soit $\mathcal{X}$ un champ algébrique localement de type fini
sur un schéma localement noethérien $S$. Soit $x_0 : \Spec(k) \to \mathcal{X}$
un morphisme, où $k$ est un corps de type fini sur $S$.
L'ensemble des {\it branches formelles de $\mathcal{X}$ passant par $x_0$}
est l'ensemble des composantes irréductibles de $\Spec(A)$
pour un choix quelconque d'anneau versel de $\mathcal{X}$ en $x_0$,
ces ensembles étant identifiés pour les différents choix de $A$ par le procédé décrit ci-dessus.
\end{definition}

\noindent
Supposons, dans la situation de la Définition \ref{stacks-geometry-definition-formal-branches},
qu'une extension finie $l/k$ soit donnée. Posons
$x_{l, 0} : \Spec(l) \to \mathcal{X}$ égal au composé de
$\Spec(l) \to \Spec(k)$ avec $x_0$. Soit
$A \to A'$ comme dans le Lemme \ref{stacks-geometry-lemma-versal-ring-field-extension}.
Puisque $A \to A'$ est fidèlement plat, le morphisme
$$
\Spec(A') \to \Spec(A)
$$
envoie les (points génériques des) composantes irréductibles sur les
(points génériques des) composantes irréductibles.
Cette application est surjective, mais
elle n'est en général pas bijective.
Autrement dit, nous obtenons une application surjective
$$
\text{branches formelles de }\mathcal{X}\text{ passant par }x_{l, 0}
\longrightarrow
\text{branches formelles de }\mathcal{X}\text{ passant par }x_0
$$
Il se trouve que, si $l/k$ est purement inséparable, alors
l'application est également injective (nous ajouterons ici un énoncé précis
et une preuve si nous en avons un jour besoin).

\begin{lemma}
\label{stacks-geometry-lemma-branches}
Dans la situation de la Définition \ref{stacks-geometry-definition-formal-branches},
il existe une surjection canonique de l'ensemble des branches formelles de
$\mathcal{X}$ passant par $x_0$ sur l'ensemble des composantes irréductibles de
$|\mathcal{X}|$ qui contiennent $x_0$ dans $|\mathcal{X}|$.
\end{lemma}

\begin{proof}
Soit $A$ comme dans la Définition \ref{stacks-geometry-definition-formal-branches}, et
soit $\Spec(A) \to \mathcal{X}$ comme dans la Remarque \ref{stacks-geometry-remark-upgrade}.
Nous affirmons que le point générique d'une composante irréductible de
$\Spec(A)$ est envoyé sur un point générique d'une composante irréductible
de $|\mathcal{X}|$. Choisissons un schéma $U$ et un morphisme
lisse surjectif $U \to \mathcal{X}$. Considérons le diagramme
$$
\xymatrix{
\Spec(A) \times_\mathcal{X} U \ar[d]_p \ar[r]_-q & U \ar[d]^f \\
\Spec(A) \ar[r]^j & \mathcal{X}
}
$$
D'après le Lemme \ref{stacks-geometry-lemma-versal-ring-flat}, $j$ est plat.
Par conséquent, $q$ est plat. D'autre part, $f$ est lisse surjectif,
donc $p$ est lisse surjectif. Cela entraîne que tout point générique
$\eta \in \Spec(A)$ d'une composante irréductible est l'image d'un
point de codimension $0$, $\eta'$, de l'espace algébrique
$\Spec(A) \times_\mathcal{X} U$ (voir
Propriétés des espaces, section \ref{spaces-properties-section-generic-points},
pour la notation et l'emploi du théorème de descente dans les anneaux locaux étales).
Puisque $q$ est plat, $q(\eta')$ est un point de codimension $0$ de $U$
(même raisonnement).
Puisque $U$ est un schéma, $q(\eta')$ est le point générique
d'une composante irréductible de $U$. Ainsi, l'adhérence de l'image
de $q(\eta')$ dans $|\mathcal{X}|$ est une composante irréductible
d'après le Lemme \ref{stacks-geometry-lemma-map-of-components}, comme affirmé.

\medskip\noindent
L'affirmation fournit manifestement un procédé pour définir l'application voulue.
Pour voir qu'elle est surjective, choisissons $u_0 \in U$ s'envoyant sur
$x_0$ dans $|\mathcal{X}|$. Choisissons un voisinage ouvert affine $U' \subset U$
de $u_0$. Quitte à rétrécir $U'$, nous pouvons supposer que toute
composante irréductible de $U'$ passe par $u_0$. Nous pouvons alors
remplacer $\mathcal{X}$ par le sous-champ ouvert correspondant à
l'image de $|U'| \to |\mathcal{X}|$. Nous pouvons donc supposer que $U$ est affine,
possède un point $u_0$ s'envoyant sur $x_0 \in |\mathcal{X}|$,
et que toute composante irréductible de $U$ passe par $u_0$.
D'après Propriétés des champs, Lemme
\ref{stacks-properties-lemma-points-cartesian}
il existe un point $t \in |\Spec(A) \times_\mathcal{X} U|$
s'envoyant sur le point fermé de $\Spec(A)$ et sur $u_0$.
En appliquant le théorème de descente aux homomorphismes locaux plats d'anneaux
$$
A \longrightarrow
\mathcal{O}_{\Spec(A) \times_\mathcal{X} U, \overline{t}}
\longleftarrow \mathcal{O}_{U, u_0}
$$
nous voyons que tout idéal premier minimal de $\mathcal{O}_{U, u_0}$
est l'image d'un idéal premier minimal de l'anneau local du milieu,
et qu'un tel idéal premier minimal s'envoie sur un idéal premier minimal de $A$.
Cela démontre la surjectivité. Certains détails sont omis.
\end{proof}

\noindent
Soit $A$ un anneau local noethérien complet. Les composantes
irréductibles de $\Spec(A)$ ont alors des multiplicités, voir
l'introduction de la section \ref{stacks-geometry-section-multiplicities}.
Si $A' = A[[t_1, \ldots, t_r]]$, alors le morphisme
$\Spec(A') \to \Spec(A)$ induit une bijection entre les composantes
irréductibles qui préserve les multiplicités (nous omettons la preuve facile).
Ceci, joint à la discussion précédant la
Définition \ref{stacks-geometry-definition-formal-branches},
montre que la définition suivante a un sens.

\begin{definition}
\label{stacks-geometry-definition-multiplicity-formal-branches}
Soit $\mathcal{X}$ un champ algébrique localement de type fini
sur un schéma localement noethérien $S$. Soit $x_0 : \Spec(k) \to \mathcal{X}$
un morphisme, où $k$ est un corps de type fini sur $S$.
La {\it multiplicité d'une branche formelle de $\mathcal{X}$ passant par $x_0$}
est la multiplicité de la composante irréductible correspondante de
$\Spec(A)$ pour un choix quelconque d'anneau versel de $\mathcal{X}$ en $x_0$
(voir la discussion ci-dessus).
\end{definition}

\begin{lemma}
\label{stacks-geometry-lemma-branches-multiplicity}
Soit $\mathcal{X}$ un champ algébrique localement de type fini
sur un schéma localement noethérien $S$. Soit $x_0 : \Spec(k) \to \mathcal{X}$
un morphisme, où $k$ est un corps de type fini sur $S$, dont
l'image est $s \in S$. Si $\mathcal{O}_{S, s}$ est un anneau G, alors
l'application du Lemme \ref{stacks-geometry-lemma-branches} préserve les multiplicités.
\end{lemma}

\begin{proof}
D'après le Lemme \ref{stacks-geometry-lemma-Artin-approximation-by-smooth-morphism}, nous
pouvons supposer qu'il existe un morphisme lisse $U \to \mathcal{X}$,
où $U$ est un schéma, et où un point à valeurs dans $k$, $u_0$, appartient à $U$,
tels que $\mathcal{O}_{U, u_0}^\wedge$ soit un anneau versel de
$\mathcal{X}$ en $x_0$. Par la construction de notre application dans
la preuve du Lemme \ref{stacks-geometry-lemma-branches} (qui se simplifie
considérablement puisque $A = \mathcal{O}_{U, u_0}^\wedge$), nous trouvons
qu'il suffit de montrer ceci : la multiplicité d'une composante irréductible
de $U$ passant par $u_0$ est égale à la
multiplicité de toute composante irréductible de
$\Spec(\mathcal{O}_{U, u_0}^\wedge)$ qui s'y envoie.

\medskip\noindent
Traduit en algèbre commutative, cela donne ce qui suit :
soit $C = \mathcal{O}_{U, u_0}$. Il est essentiellement de type fini
sur $\mathcal{O}_{S, s}$ et est donc un anneau G (Compléments d'algèbre,
Proposition \ref{more-algebra-proposition-finite-type-over-G-ring}).
Alors $A = C^\wedge$. Par conséquent, $C \to A$ est un homomorphisme régulier d'anneaux.
Soit $\mathfrak q \subset C$ un idéal premier minimal et soit $\mathfrak p \subset A$
un idéal premier minimal au-dessus de $\mathfrak q$. Alors
$$
R = C_\mathfrak p \longrightarrow A_\mathfrak p = R'
$$
est un homomorphisme régulier d'anneaux locaux artiniens. Pour un tel homomorphisme,
on a toujours
$$
\text{length}_R R = \text{length}_{R'} R'
$$
C'est précisément ce qu'il fallait montrer, car le membre de gauche
est la multiplicité de notre composante sur $U$, et le membre de droite
est la multiplicité de notre composante sur $\Spec(A)$.
Pour établir l'égalité, nous utilisons d'abord
$$
\text{length}_R(R) \text{length}_{R'}(R'/\mathfrak m_R R') =
\text{length}_{R'}(R')
$$
d'après Algèbre, Lemme \ref{algebra-lemma-pullback-module}. Il suffit donc
de montrer que $\mathfrak m_R R' = \mathfrak m_{R'}$, ce qui résulte
du fait qu'il s'agit d'un homomorphisme régulier d'anneaux locaux de dimension zéro.
\end{proof}













\section{Théorie de la dimension des champs algébriques}
\label{stacks-geometry-section-dimension-of-algebraic-stacks}

\noindent
À notre connaissance, les principaux résultats publiés sur la théorie de la
dimension des champs algébriques sont ceux de \cite{Osserman}, qui
étudie les notions de codimension et de dimension relative. Nous
examinons plus en détail la notion de dimension d'un
champ algébrique en un point, et démontrons divers résultats
reliant la dimension des fibres d'un morphisme en un point de la source
à la dimension de sa source et de son but. Nous démontrons aussi un résultat
(le Lemme \ref{stacks-geometry-lemma-dimension-formula} ci-dessous) qui
nous permet (sous des hypothèses appropriées) de calculer la dimension
d'un champ algébrique en un point au moyen d'un anneau versel.

\medskip\noindent
Sans avoir toujours cherché à optimiser nos résultats, nous avons
dans une large mesure essayé d'éviter les hypothèses superflues. Toutefois, dans certains
de nos résultats, où nous comparons certaines propriétés d'un champ
algébrique à celles d'un anneau versel de ce
champ en un point, nous nous sommes restreints
au cas des champs algébriques localement de présentation finie
sur un schéma de base localement noethérien dont tous les anneaux locaux sont
des anneaux $G$. Cela nous permet de disposer de l'approximation d'Artin
pour comparer la géométrie de l'anneau versel à celle
du champ lui-même. Il se peut cependant que cette hypothèse restrictive
ne soit pas nécessaire à la validité
de tous les énoncés que nous démontrons.
Comme elle est satisfaite dans les applications que nous avons en vue,
nous nous en sommes néanmoins accommodés lorsqu'elle est utile.

\medskip\noindent
Si $X$ est un schéma, nous définissons la dimension $\dim(X)$ de $X$
comme la dimension de Krull de l'
espace topologique sous-jacent à $X$,
tandis que, si $x$ est un point de $X$,
nous définissons la dimension $\dim_x (X)$ de $X$ en $x$ comme le
minimum des dimensions des sous-ensembles ouverts $U$ de $X$ contenant
$x$, voir
Propriétés, Définition \ref{properties-definition-dimension}.
On a la relation $\dim(X) = \sup_{x \in X} \dim_x(X)$, voir
Propriétés, Lemme \ref{properties-lemma-dimension}.
Si $X$ est localement noethérien, alors $\dim_x(X)$ coïncide avec le supremum
des dimensions en $x$
des composantes irréductibles de $X$ passant par $x$.

\medskip\noindent
Si $X$ est un espace algébrique et $x \in |X|$,
nous définissons $\dim_x X = \dim_u U,$ où $U$ est un schéma quelconque
admettant une surjection étale $U \to X$,
et où $u\in U$ est un point quelconque au-dessus de $x$, voir
Propriétés des espaces, Définition
\ref{spaces-properties-definition-dimension-at-point}.
Nous posons $\dim(X) = \sup_{x \in |X|} \dim_x(X)$, voir
Propriétés des espaces, Définition \ref{spaces-properties-definition-dimension}.

\begin{remark}
\label{stacks-geometry-remark-dimension-algebraic-space}
En général, la dimension de l'espace algébrique $X$ en un point $x$
peut ne pas coïncider avec la dimension de l'espace topologique sous-jacent
$|X|$ en $x$. Par exemple, si $k$ est un corps de caractéristique zéro et
$X =  \mathbf{A}^1_k / \mathbf{Z}$, alors $X$ est de dimension $1$ (la dimension
de $\mathbf{A}^1_k$) en chacun de ses points,
tandis que $|X|$ est muni de la topologie grossière et a donc une
dimension de Krull nulle. D'autre part, dans
Espaces algébriques, Exemple \ref{spaces-example-infinite-product},
on donne un exemple d'espace algébrique
qui est de dimension $0$ en chacun de ses points, tandis que $|X|$ est
irréductible, de dimension de Krull $1$, et admet un point générique (de sorte que la
dimension de $|X|$ en chacun de ses points est $1$); voir aussi la discussion
de cet exemple dans
Propriétés des espaces, section \ref{spaces-properties-section-dimension}.

\medskip\noindent
D'autre part, si $X$ est un espace algébrique {\it convenable}, au sens de
Espaces convenables, Définition \ref{decent-spaces-definition-very-reasonable}
(en particulier, si $X$ est quasi-séparé; voir
Espaces convenables, section \ref{decent-spaces-section-reasonable-decent}),
alors la dimension de $X$ en $x$ coïncide effectivement avec la dimension
de $|X|$ en $x$; voir
Espaces convenables, Lemme \ref{decent-spaces-lemma-dimension-decent-space}.
\end{remark}

\noindent
Pour définir la dimension d'un champ algébrique,
il sera utile de disposer d'abord de la notion de dimension relative,
en un point de la source,
d'un morphisme dont la source est un espace algébrique
et le but un champ algébrique. La définition est quelque peu
compliquée, simplement parce que (contrairement au cas des schémas) les points d'un
champ algébrique, ou d'un espace
algébrique, ne peuvent être décrits comme des morphismes depuis le spectre d'un corps,
mais seulement comme des classes d'équivalence de tels morphismes.

\begin{definition}
\label{stacks-geometry-definition-relative-dimension}
Si $f : T \to \mathcal{X}$ est un morphisme localement de type fini d'un
espace algébrique vers un champ algébrique,
et si $t \in |T|$ est un point d'image $x \in | \mathcal{X}|$, nous définissons
{\it la dimension relative} de $f$ en $t$, notée
$\dim_t(T_x),$
comme suit :
choisissons un morphisme $\Spec k \to \mathcal{X}$, de source le spectre d'un
corps, qui représente $x$, et choisissons un point
$t' \in |T \times_{\mathcal{X}} \Spec k|$
s'envoyant sur $t$ par la projection vers $|T|$
(un tel point $t'$ existe, d'après
Propriétés des champs, Lemme \ref{stacks-properties-lemma-points-cartesian});
alors
$$
\dim_t(T_x) = \dim_{t'}(T \times_{\mathcal{X}} \Spec k ).
$$
\end{definition}

\noindent
Notons que, puisque $T$ est un espace algébrique et $\mathcal{X}$ un
champ algébrique, le produit fibré $T \times_{\mathcal{X}} \Spec k$
est un espace algébrique; la quantité figurant au membre de droite de
cette définition est donc bien définie (voir la discussion ci-dessus).

\begin{remark}
\label{stacks-geometry-remark-relative-dimension}
(1)
On vérifie aisément (par exemple en utilisant l'invariance
par changement de base de la dimension relative des morphismes de schémas localement
de type fini; voir par exemple
Morphismes, Lemme \ref{morphisms-lemma-dimension-fibre-after-base-change})
que $\dim_t(T_x)$ est bien définie, indépendamment des choix
faits pour la calculer.

\medskip\noindent
(2)
Lorsque $\mathcal{X}$ est aussi un espace algébrique,
on vérifie immédiatement que cette définition coïncide avec
la définition de la dimension relative donnée dans
Morphismes d'espaces, Définition
\ref{spaces-morphisms-definition-dimension-fibre}.
\end{remark}

\noindent
Rappelons maintenant le lemme suivant, sur lequel repose notre étude de
la dimension d'un champ algébrique localement noethérien.

\begin{lemma}
\label{stacks-geometry-lemma-behaviour-of-dimensions-wrt-smooth-morphisms}
Si $f: U \to X$ est un morphisme lisse d'espaces algébriques localement
noethériens, et
si $u \in |U|$ est d'image $x \in |X|$, alors
$$
\dim_u (U) = \dim_x(X) + \dim_{u} (U_x)
$$
où $\dim_u (U_x)$ est défini par la
Définition \ref{stacks-geometry-definition-relative-dimension}.
\end{lemma}

\begin{proof}
Voir Morphismes d'espaces, Lemme
\ref{spaces-morphisms-lemma-smoothness-dimension-spaces}
en notant que la définition de $\dim_u (U_x)$ utilisée ici coïncide avec
celle qui y est employée, d'après la Remarque \ref{stacks-geometry-remark-relative-dimension} (2).
\end{proof}

\begin{lemma}
\label{stacks-geometry-lemma-dimension-for-stacks}
Soient $\mathcal{X}$ un champ algébrique localement noethérien et
$x \in |\mathcal{X}|$. Soit $U \to \mathcal{X}$ un morphisme lisse
d'un espace algébrique vers $\mathcal{X}$, et soit $u$ un point quelconque de $|U|$
s'envoyant sur $x$. Alors
$$
\dim_x(\mathcal{X}) =  \dim_u(U) - \dim_{u}(U_x)
$$
où la dimension relative $\dim_u(U_x)$ est définie
par la Définition \ref{stacks-geometry-definition-relative-dimension}, et la
dimension de $\mathcal{X}$ en $x$ comme dans
Propriétés des champs, Définition
\ref{stacks-properties-definition-dimension-at-point}.
\end{lemma}

\begin{proof}
Le Lemme \ref{stacks-geometry-lemma-behaviour-of-dimensions-wrt-smooth-morphisms}
permet de vérifier que le membre de droite
$\dim_u(U) + \dim_u(U_x)$ est indépendant du choix
du morphisme lisse $U \to \mathcal{X}$ et de $u \in |U|$.
Nous omettons les détails. En particulier, nous pouvons supposer que $U$ est un schéma.
Dans ce cas, nous pouvons calculer $\dim_u(U_x)$
en choisissant comme représentant de $x$
le composé $\Spec \kappa(u) \to U \to \mathcal{X}$, où
le premier morphisme est le morphisme canonique d'image $u \in U$.
Écrivons alors $R = U \times_{\mathcal{X}} U$ et notons
$e : U \to R$ le morphisme diagonal. L'invariance de la
dimension relative par changement de base montre que
$\dim_u(U_x) = \dim_{e(u)}(R_u)$. Nous voyons donc que
le membre de droite est égal à
$\dim_u (U) - \dim_{e(u)}(R_u) = \dim_x(\mathcal{X})$, comme voulu.
\end{proof}

\begin{remark}
\label{stacks-geometry-remark-dimension-DM}
Pour les champs de Deligne--Mumford suffisamment convenables
(par exemple quasi-séparés),
$\dim_x(\mathcal{X})$ coïncide à nouveau avec la
quantité définie
topologiquement $\dim_x |\mathcal{X}|$. Cependant, pour des champs d'Artin plus
généraux,
ce n'est généralement pas le cas. Par exemple, si
$\mathcal{X} = [\mathbf{A}^1/\mathbf{G}_m]$
(sur un certain corps, le quotient étant pris pour
l'action multiplicative usuelle de $\mathbf{G}_m$ sur $\mathbf{A}^1$),
alors $|\mathcal{X}|$ possède deux points, dont l'un est une spécialisation de l'autre
(correspondant
aux deux orbites de $\mathbf{G}_m$ sur $\mathbf{A}^1$), et est donc de
dimension $1$ comme
espace topologique; mais $\dim_x (\mathcal{X}) = 0$ pour les deux points
$x \in |\mathcal{X}|$.
(Un exemple encore plus extrême est fourni par l'espace classifiant
$[\Spec k/\mathbf{G}_m]$, dont la dimension en son unique point
est égale à $-1$.)
\end{remark}

\noindent
Nous pouvons maintenant étendre la Définition \ref{stacks-geometry-definition-relative-dimension}
au contexte des morphismes (localement de type fini)
entre champs algébriques (localement noethériens).

\begin{definition}
\label{stacks-geometry-definition-relative-dimension-for-stacks}
Si $f : \mathcal{T} \to \mathcal{X}$
est un morphisme localement de type fini entre
champs algébriques localement noethériens, et si
$t \in |\mathcal{T}|$ est un point d'image $x \in |\mathcal{X}|$, nous
définissons la {\it dimension relative} de $f$ en $t$, notée
$\dim_t(\mathcal{T}_x),$ comme suit :
choisissons un morphisme $\Spec k \to \mathcal{X}$, de source le spectre d'un
corps, qui représente $x$, et choisissons un point
$t' \in |\mathcal{T} \times_{\mathcal{X}} \Spec k|$
s'envoyant sur $t$ par la projection vers $|\mathcal{T}|$
(un tel point $t'$ existe, d'après
Propriétés des champs, Lemme
\ref{stacks-properties-lemma-points-cartesian}; alors
$$
\dim_t(\mathcal{T}_x) = \dim_{t'}(\mathcal{T} \times_{\mathcal{X}} \Spec k ).
$$
\end{definition}

\noindent
Notons que, puisque $\mathcal{T}$ est un champ algébrique et $\mathcal{X}$ un
champ algébrique, le produit fibré $\mathcal{T}\times_{\mathcal{X}} \Spec k$
est un champ algébrique, localement noethérien d'après
Morphismes de champs, Lemme
\ref{stacks-morphisms-lemma-locally-finite-type-locally-noetherian}.
Ainsi, la quantité figurant au membre de droite de cette définition
est définie par Propriétés des champs,
Définition \ref{stacks-properties-definition-dimension-at-point}.

\begin{remark}
\label{stacks-geometry-remark-dimension-tangent-space-well-defined}
Des manipulations usuelles montrent que $\dim_t(\mathcal{T}_x)$ est bien définie,
indépendamment des choix faits pour la calculer.
\end{remark}

\noindent
Établissons maintenant quelques propriétés élémentaires de la dimension relative,
qui sont des généralisations évidentes des énoncés correspondants dans le
cas des morphismes de schémas.

\begin{lemma}
\label{stacks-geometry-lemma-base-change-invariance-of-relative-dimension}
Supposons donné
un carré cartésien de morphismes de champs localement noethériens
$$
\xymatrix{
\mathcal{T}' \ar[d]\ar[r] & \mathcal{T} \ar[d] \\
\mathcal{X}' \ar[r] & \mathcal{X}
}
$$
où les morphismes verticaux sont localement de type fini.
Si $t' \in |\mathcal{T}'|$,
a pour images $t$, $x'$ et $x$ dans $|\mathcal{T}|$, $|\mathcal{X}'|$ et
$|\mathcal{X}|$,
respectivement, alors $\dim_{t'}(\mathcal{T}'_{x'}) = \dim_{t}(\mathcal{T}_x).$
\end{lemma}

\begin{proof}
Les deux membres peuvent (par définition) être calculés comme la
dimension du même produit fibré.
\end{proof}

\begin{lemma}
\label{stacks-geometry-lemma-behaviour-of-dimensions-wrt-smooth-morphisms-stacky}
Si $f: \mathcal{U} \to \mathcal{X}$ est un morphisme lisse de champs algébriques
localement noethériens, et
si $u \in |\mathcal{U}|$ est d'image $x \in |\mathcal{X}|$,
alors
$$
\dim_u (\mathcal{U}) = \dim_x(\mathcal{X}) + \dim_{u} (\mathcal{U}_x).
$$
\end{lemma}

\begin{proof}
Choisissons un morphisme lisse surjectif $V \to \mathcal{U}$ dont la source
est un schéma, et soit $v\in |V|$ un point s'envoyant sur $u$.
Alors le composé $V \to \mathcal{U} \to \mathcal{X}$ est lui aussi lisse,
et, d'après le Lemme \ref{stacks-geometry-lemma-behaviour-of-dimensions-wrt-smooth-morphisms},
nous avons $\dim_x(\mathcal{X}) = \dim_v(V) - \dim_v(V_x),$
tandis que $\dim_u(\mathcal{U}) = \dim_v(V) - \dim_v(V_u).$
Ainsi
$$
\dim_u(\mathcal{U}) - \dim_x(\mathcal{X}) = \dim_v (V_x) - \dim_v (V_u).
$$

\medskip\noindent
Choisissons un représentant $\Spec k \to \mathcal{X}$ de $x$
et un point $v' \in | V \times_{\mathcal{X}} \Spec k|$ au-dessus de
$v$, d'image $u'$ dans $|\mathcal{U}\times_{\mathcal{X}} \Spec k|$;
alors, par définition,
$\dim_u(\mathcal{U}_x) = \dim_{u'}(\mathcal{U}\times_{\mathcal{X}} \Spec k),$
et
$\dim_v(V_x) = \dim_{v'}(V\times_{\mathcal{X}} \Spec k).$

\medskip\noindent
Or $V\times_{\mathcal{X}} \Spec k \to \mathcal{U}\times_{\mathcal{X}}\Spec k$
est un morphisme lisse surjectif (comme changement de base
d'un tel morphisme) dont la source est un espace algébrique
(puisque $V$ et $\Spec k$ sont des schémas et $\mathcal{X}$
un champ algébrique). Ainsi, toujours par définition,
nous avons
\begin{align*}
\dim_{u'}(\mathcal{U}\times_{\mathcal{X}} \Spec k)
& =
\dim_{v'}(V\times_{\mathcal{X}} \Spec k) -
\dim_{v'}(V \times_{\mathcal{X}} \Spec k)_{u'}) \\
& = \dim_v(V_x) -
\dim_{v'}( (V\times_{\mathcal{X}} \Spec k)_{u'}).
\end{align*}
Or $V\times_{\mathcal{X}} \Spec k \cong
V\times_{\mathcal{U}} (\mathcal{U}\times_{\mathcal{X}} \Spec k),$
et donc le
Lemme \ref{stacks-geometry-lemma-base-change-invariance-of-relative-dimension}
montre que
$\dim_{v'}((V\times_{\mathcal{X}} \Spec k)_{u'})  = \dim_v(V_u).$
En réunissant ces égalités, nous trouvons
$$
\dim_u(\mathcal{U}) - \dim_x(\mathcal{X}) =
\dim_u(\mathcal{U}_x),
$$
comme voulu.
\end{proof}

\begin{lemma}
\label{stacks-geometry-lemma-relative-dimension-is-semi-continuous}
Soit $f: \mathcal{T} \to \mathcal{X}$ un morphisme localement de type fini de
champs algébriques.
\begin{enumerate}
\item
La fonction $t \mapsto \dim_t(\mathcal{T}_{f(t)})$ est semi-continue supérieurement
sur $|\mathcal{T}|$.
\item Si $f$ est lisse, alors
la fonction $t \mapsto \dim_t(\mathcal{T}_{f(t)})$ est localement constante
sur $|\mathcal{T}|$.
\end{enumerate}
\end{lemma}

\begin{proof}
Supposons d'abord que $\mathcal{T}$ soit un schéma $T$,
soit $U \to \mathcal{X}$ un morphisme lisse surjectif dont la source
est un schéma, et soit $T' = T \times_{\mathcal{X}} U$.
Soit $f': T' \to U$ le changement de base de $f$ à $U$,
et soit $g: T' \to T$ la projection.

\medskip\noindent
Le Lemme \ref{stacks-geometry-lemma-base-change-invariance-of-relative-dimension}
montre que $\dim_{t'}(T'_{f'(t')}) = \dim_{g(t')}(T_{f(g(t'))}),$
pour $t' \in T'$, tandis que,
puisque $g$ est lisse et surjectif (étant le changement de base
d'un morphisme lisse surjectif), l'application induite par $g$ sur les espaces
topologiques sous-jacents est continue et ouverte
(d'après
Propriétés des espaces, Lemme \ref{spaces-properties-lemma-topology-points}), et
surjective. Il suffit donc de noter que la partie (1) pour le morphisme $f'$
résulte de
Morphismes d'espaces, Lemme
\ref{spaces-morphisms-lemma-openness-bounded-dimension-fibres}, et la partie (2)
de l'un ou l'autre des lemmes suivants, Morphismes, Lemme
\ref{morphisms-lemma-flat-finite-presentation-CM-fibres-relative-dimension}
ou
Morphismes, Lemme \ref{morphisms-lemma-smooth-omega-finite-locally-free}
(chacun donne le résultat pour les schémas, d'où
les résultats analogues pour les espaces algébriques se
déduisent exactement comme dans
Morphismes d'espaces, Lemme
\ref{spaces-morphisms-lemma-openness-bounded-dimension-fibres}.

\medskip\noindent
Revenons maintenant au cas général,
et choisissons un morphisme lisse surjectif
$h:V \to \mathcal{T}$ dont la source est un schéma.
Si $v \in V$, alors, essentiellement par définition,
nous avons
$$
\dim_{h(v)}(\mathcal{T}_{f(h(v))}) =
\dim_{v}(V_{f(h(v))}) - \dim_{v}(V_{h(v)}).
$$
Puisque $V$ est un schéma, nous avons démontré que le premier
des termes du membre de droite de cette égalité
est semi-continu supérieur (et même localement
constant si $f$ est lisse), tandis que le second terme est
en fait localement constant.
Leur différence est donc semi-continue supérieurement
(et localement constante si $f$ est lisse),
et par conséquent la fonction
$\dim_{h(v)}(\mathcal{T}_{f(h(v))})$
est semi-continue supérieurement sur $|V|$ (et localement
constante si $f$ est lisse).
Puisque le morphisme $|V| \to |\mathcal{T}|$ est ouvert et surjectif,
le lemme en résulte.
\end{proof}

\noindent
Avant de poursuivre notre étude,
démontrons deux lemmes relatifs à la théorie de la dimension des schémas.

\medskip\noindent
Pour situer le premier lemme,
notons que, si $X$ est un schéma de dimension finie, puisque $\dim X$
est défini comme le supremum des dimensions $\dim_x X$,
il existe un point $x \in X$ tel que $\dim_x X = \dim X$.
Le lemme suivant montre que nous pouvons en outre choisir le point
$x$ de type fini.

\begin{lemma}
\label{stacks-geometry-lemma-dimension-achieved-by-finite-type-point}
Si $X$ est un schéma de dimension finie,
il existe un point fermé (et donc de type fini) $x \in X$
tel que $\dim_x X = \dim X$.
\end{lemma}

\begin{proof}
Soit $d = \dim X$,
et choisissons une chaîne strictement décroissante maximale
de sous-ensembles fermés irréductibles de $X$,
disons
\begin{equation}
\label{stacks-geometry-equation-maximal-chain}
Z_0 \supset Z_1 \supset \ldots \supset Z_d.
\end{equation}
Le sous-ensemble $Z_d$ est un sous-ensemble fermé irréductible minimal de $X$,
et tout point de $Z_d$ est donc un point générique de $Z_d$.
Puisque l'espace topologique sous-jacent au schéma $X$ est sobre,
nous concluons que $Z_d$ est un singleton, constitué d'un unique
point fermé $x \in X$.
Si $U$ est
un voisinage quelconque de $x$, alors
la chaîne
$$
U\cap Z_0 \supset U\cap Z_1 \supset \ldots \supset U\cap Z_d = Z_d =
\{x\}
$$
est une chaîne strictement décroissante de sous-ensembles
fermés irréductibles de $U$, ce qui montre que $\dim U \geq d$.
Nous trouvons donc $\dim_x X \geq d$. L'autre inégalité
étant évidente, le lemme est démontré.
\end{proof}

\noindent
Le lemme suivant montre que $\dim_x X$ est une fonction {\it constante}
sur un schéma irréductible satisfaisant quelques hypothèses supplémentaires modérées.

\begin{lemma}
\label{stacks-geometry-lemma-constancy-of-dimension}
Si $X$ est un schéma irréductible, jacobsonien, caténaire et localement
noethérien, de dimension finie,
alors $\dim U = \dim X$ pour tout
sous-ensemble ouvert non vide $U$ de $X$.
De manière équivalente, $\dim_x X$ est une fonction constante sur $X$.
\end{lemma}

\begin{proof}
L'équivalence des deux assertions résulte directement des
définitions. Supposons donc que $U\subset X$ soit un sous-ensemble ouvert
non vide.
Certes $\dim U \leq \dim X$, et nous devons montrer
que $\dim U \geq \dim X.$
Écrivons $d = \dim X$, et choisissons une chaîne strictement
décroissante maximale de sous-ensembles fermés irréductibles
de $X$, disons
$$
X = Z_0 \supset Z_1 \supset \ldots \supset Z_d.
$$
Puisque $X$ est jacobsonien, le sous-ensemble fermé irréductible
minimal $Z_d$ est égal à $\{x\}$ pour un certain
point fermé $x$.

\medskip\noindent
Si $x \in U,$ alors
$$
U = U \cap Z_0  \supset U\cap Z_1 \supset \ldots \supset
U\cap Z_d = \{x\}
$$
est une chaîne strictement décroissante de sous-ensembles fermés
irréductibles de $U$, et nous concluons donc que $\dim U \geq d$,
comme voulu. Nous pouvons donc supposer que $x \not\in U.$

\medskip\noindent
Considérons le morphisme plat $\Spec \mathcal{O}_{X,x} \to X$.
Le sous-ensemble ouvert non vide (et donc dense) $U$ de $X$
se relève en un sous-ensemble ouvert $V \subset \Spec \mathcal{O}_{X,x}$.
En remplaçant $U$ par un sous-ensemble ouvert non vide quasi-compact, donc
noethérien, nous pouvons supposer que l'inclusion
$U \to X$ est un morphisme quasi-compact. Puisque la
formation des images schématiques de morphismes quasi-compacts
commute au changement de base plat, d'après
Morphismes, Lemme
\ref{morphisms-lemma-flat-base-change-scheme-theoretic-image}
nous voyons que $V$ est dense dans $\Spec \mathcal{O}_{X,x}$,
et en particulier non vide,
et bien sûr $x \not\in V.$ (Ici, nous notons aussi $x$
le point fermé de $\Spec \mathcal{O}_{X,x}$, puisque son image
est égale au point donné $x \in X$.)
Or $\Spec \mathcal{O}_{X,x} \setminus \{x\}$ est jacobsonien d'après
Propriétés, Lemme
\ref{properties-lemma-complement-closed-point-Jacobson}
et $V$ contient donc un point fermé $z$
de $\Spec \mathcal{O}_{X,x} \setminus \{x\}$. L'adhérence
dans $X$ de l'image de $z$ est alors un sous-ensemble
fermé irréductible $Z$ de $X$ contenant $x$, dont l'intersection
avec $U$ est non vide, et
pour lequel il n'existe aucun sous-ensemble fermé
irréductible strictement contenu dans $Z$
et contenant strictement $\{x\}$
(car le changement de base à $\Spec \mathcal{O}_{X,x}$ induit
une bijection entre les sous-ensembles fermés irréductibles de $X$
contenant $x$ et les sous-ensembles fermés irréductibles de $\Spec
\mathcal{O}_{X,x}$).
Puisque $U \cap Z$ est un sous-ensemble fermé non vide de $U$,
il contient un point $u$ qui est fermé dans $X$ (puisque
$X$ est jacobsonien), et puisque $U\cap Z$
est un sous-ensemble ouvert non vide (donc dense) de l'ensemble irréductible $Z$
(qui contient un point n'appartenant pas à $U$, à savoir $x$),
l'inclusion $\{u\} \subset U\cap Z$ est stricte.

\medskip\noindent
Puisque $X$ est caténaire, la chaîne
$$
X = Z_0 \supset Z \supset \{x\} = Z_d
$$
peut être raffinée en une chaîne de longueur $d+1$, qui est nécessairement
de la forme
$$
X = Z_0 \supset W_1 \supset \ldots \supset W_{d-1} = Z \supset \{x\} = Z_d.
$$
Puisque $U\cap Z$ est non vide, nous trouvons alors que
$$
U = U \cap Z_0 \supset U \cap W_1\supset \ldots \supset U\cap W_{d-1}
= U\cap Z \supset \{u\}
$$
est une chaîne strictement décroissante de sous-ensembles fermés irréductibles
de $U$, de longueur $d+1$, ce qui montre que $\dim U \geq d$,
comme voulu.
\end{proof}

\noindent
Nous démontrerons un analogue pour les champs
du Lemme \ref{stacks-geometry-lemma-constancy-of-dimension}
dans le Lemme \ref{stacks-geometry-lemma-irreducible-implies-equidimensional} ci-dessous,
mais nous devons auparavant introduire une définition supplémentaire,
rendue nécessaire par le fait que la propriété d'un schéma d'être caténaire
n'est pas locale pour la topologie étale
(voir l'exemple de
Algèbre, Remarque \ref{algebra-remark-universally-catenary-does-not-descend},
ce qui rend difficile de définir ce que signifie, pour un espace algébrique
ou un champ algébrique, d'être caténaire
(voir la discussion de \cite[page 3]{Osserman}).
Pour certains aspects de la théorie de la dimension, la définition
suivante semble fournir un bon substitut à la notion manquante
de champ algébrique caténaire.

\begin{definition}
\label{stacks-geometry-definition-pseudo-catenary}
Nous disons qu'un champ algébrique localement noethérien $\mathcal{X}$
est {\it pseudo-caténaire} s'il existe un morphisme lisse
et surjectif $U \to \mathcal{X}$ dont la source est
un schéma universellement caténaire.
\end{definition}

\begin{example}
\label{stacks-geometry-example-pseudo-catenary}
Si $\mathcal{X}$ est localement de type fini sur un schéma universellement
caténaire et localement noethérien $S$,
et si $U\to \mathcal{X}$ est un morphisme lisse surjectif
dont la source est un schéma, alors le composé
$U \to \mathcal{X} \to S$ est localement de type fini,
et donc $U$ est universellement caténaire d'après
Morphismes, Lemme
\ref{morphisms-lemma-universally-catenary-local}.
Ainsi, $\mathcal{X}$ est pseudo-caténaire.
\end{example}

\noindent
Le lemme suivant montre que la propriété d'être pseudo-caténaire
se transmet par les morphismes de type fini.

\begin{lemma}
\label{stacks-geometry-lemma-catenary-covers}
Si $\mathcal{X}$ est un champ algébrique localement noethérien
pseudo-caténaire, et si $\mathcal{Y} \to \mathcal{X}$ est un morphisme
localement de type fini,
alors il existe un morphisme lisse surjectif $V \to \mathcal{Y}$
dont la source est un schéma universellement caténaire; ainsi
$\mathcal{Y}$ est encore pseudo-caténaire.
\end{lemma}

\begin{proof}
Par hypothèse, nous pouvons trouver un morphisme lisse surjectif
$U \to \mathcal{X}$ dont la source est un schéma universellement caténaire.
Le changement de base $U\times_{\mathcal{X}} \mathcal{Y}$ est alors un champ
algébrique; soit $V \to U\times_{\mathcal{X}} \mathcal{Y}$ un morphisme lisse
surjectif dont la source est un schéma.
Le composé $V \to U\times_{\mathcal{X}} \mathcal{Y} \to \mathcal{Y}$ est alors
lisse et surjectif (comme composé de morphismes lisses et
surjectifs), tandis que le morphisme $V \to U\times_{\mathcal{X}}
\mathcal{Y} \to U$ est localement de type fini (comme composé
de morphismes localement de type fini). Puisque $U$
est universellement caténaire, nous voyons que $V$ l'est aussi
(d'après Morphismes, Lemme
\ref{morphisms-lemma-universally-catenary-local}),
comme affirmé.
\end{proof}

\noindent
Étudions maintenant le comportement de la fonction $\dim_x(\mathcal{X})$ sur
$|\mathcal{X}|$
(pour un champ localement noethérien $\mathcal{X}$) par rapport aux
composantes
irréductibles de $|\mathcal{X}|$, ainsi que diverses
questions connexes.

\begin{lemma}
\label{stacks-geometry-lemma-irreducible-implies-equidimensional}
Si $\mathcal{X}$ est
un champ algébrique jacobsonien, pseudo-caténaire et localement noethérien
tel que $|\mathcal{X}|$ soit irréductible,
alors $\dim_x(\mathcal{X})$ est une fonction constante sur $|\mathcal{X}|$.
\end{lemma}

\begin{proof}
Il suffit de montrer que $\dim_x(\mathcal{X})$ est localement constante sur
$|\mathcal{X}|$,
car elle sera alors nécessairement constante (puisque $|\mathcal{X}|$ est connexe,
étant irréductible). Puisque $\mathcal{X}$ est pseudo-caténaire,
nous pouvons trouver un morphisme lisse surjectif $U \to \mathcal{X}$, avec $U$
un schéma universellement caténaire. Si $\{U_i\}$ est un
recouvrement de $U$ par des sous-schémas ouverts quasi-compacts, nous pouvons remplacer
$U$ par $\coprod U_i,$, et
il suffit de montrer que
la fonction $u \mapsto \dim_{f(u)}(\mathcal{X})$ est localement constante sur $U_i$.
Puisque nous vérifions ceci pour un seul $U_i$ à la fois, omettons désormais l'indice,
et écrivons simplement $U$ plutôt que $U_i$.
Puisque $U$ est quasi-compact, il
est la réunion d'un nombre fini de composantes irréductibles,
disons $T_1 \cup \ldots \cup T_n$. Notons que chaque $T_i$ est jacobsonien,
caténaire et localement noethérien,
car c'est un sous-schéma fermé du schéma jacobsonien, caténaire et localement noethérien
$U$.

\medskip\noindent
D'après le Lemme \ref{stacks-geometry-lemma-behaviour-of-dimensions-wrt-smooth-morphisms}, nous avons
$\dim_{f(u)}(\mathcal{X}) = \dim_{u}(U) - \dim_{u}(U_{f(u)}).$
Le Lemme \ref{stacks-geometry-lemma-relative-dimension-is-semi-continuous} (2)
montre que le second terme du membre de droite est localement
constant sur $U$, puisque $f$ est lisse;
nous devons donc montrer que $\dim_u(U)$
est localement constante sur $U$. Puisque $\dim_u(U)$ est le maximum
des dimensions $\dim_u T_i$, lorsque $T_i$ parcourt les composantes
de $U$ contenant $u$, il suffit de montrer
que, si un point $u$ appartient à deux composantes distinctes,
disons $T_i$ et $T_j$ (avec $i \neq j$),
alors $\dim_u T_i = \dim_u T_j$,
puis de noter que $t\mapsto \dim_t T$ est une fonction constante
sur un schéma irréductible et jacobsonien,
caténaire et localement noethérien, noté $T$,
(comme il résulte du Lemme \ref{stacks-geometry-lemma-constancy-of-dimension}).

\medskip\noindent
Soient $V = T_i \setminus (\bigcup_{i' \neq i} T_{i'})$
et $W = T_j \setminus (\bigcup_{i' \neq j} T_{i'})$.
Chacun de $V$ et $W$ est alors un sous-ensemble ouvert non vide de $U$,
et a donc une image ouverte non vide dans $|\mathcal{X}|$. Puisque $|\mathcal{X}|$ est
irréductible,
ces deux sous-ensembles ouverts non vides de $|\mathcal{X}|$ ont une intersection
non vide.
Soit $x$ un point de cette intersection, et soient $v \in V$ et
$w\in W$ des points s'envoyant sur $x$.
Nous trouvons alors
$$
\dim T_i = \dim V = \dim_v (U) = \dim_x (\mathcal{X}) + \dim_v (U_x)
$$
et, de même,
$$
\dim T_j = \dim W = \dim_w (U) = \dim_x (\mathcal{X}) + \dim_w (U_x).
$$
Puisque $u \mapsto \dim_u (U_{f(u)})$ est localement constante sur $U$,
et puisque $T_i \cup T_j$ est connexe (comme réunion de deux ensembles irréductibles,
donc connexes, d'intersection non vide),
nous voyons que $\dim_v (U_x) = \dim_w(U_x)$,
et donc, en comparant les deux équations précédentes,
que $\dim T_i = \dim T_j$, comme voulu.
\end{proof}

\begin{lemma}
\label{stacks-geometry-lemma-closed-immersions}
Si $\mathcal{Z} \hookrightarrow \mathcal{X}$ est une immersion fermée
de champs algébriques localement noethériens,
et si $z \in |\mathcal{Z}|$ a pour image $x \in |\mathcal{X}|$,
alors $\dim_z (\mathcal{Z}) \leq \dim_x(\mathcal{X})$.
\end{lemma}

\begin{proof}
Choisissons un morphisme lisse surjectif
$U\to \mathcal{X}$ dont la source est un schéma;
le morphisme obtenu par changement de base
$V = U\times_{\mathcal{X}} \mathcal{Z} \to \mathcal{Z}$
est alors lui aussi lisse et surjectif, et la projection
$V \to U$ est une immersion fermée.
Si $v \in |V|$ s'envoie sur $z \in |\mathcal{Z}|$, et
si nous notons $u$ l'image de $v$ dans $|U|$,
alors, manifestement,
$\dim_v(V) \leq \dim_u(U)$,
tandis que
$\dim_v (V_z) = \dim_u(U_x)$,
d'après le Lemme \ref{stacks-geometry-lemma-base-change-invariance-of-relative-dimension}.
Ainsi,
$$
\dim_z(\mathcal{Z})  = \dim_v(V) - \dim_v(V_z)
\leq \dim_u(U) - \dim_u(U_x) = \dim_x(\mathcal{X}),
$$
comme affirmé.
\end{proof}

\begin{lemma}
\label{stacks-geometry-lemma-dimension-via-components}
Si $\mathcal{X}$ est un champ algébrique localement noethérien, et si
$x \in |\mathcal{X}|$,
alors $\dim_x(\mathcal{X}) = \sup_{\mathcal{T}} \{ \dim_x(\mathcal{T}) \} $,
où $\mathcal{T}$ parcourt toutes les composantes irréductibles
de $|\mathcal{X}|$ passant par $x$ (munies de leur
structure réduite induite).
\end{lemma}

\begin{proof}
Le Lemme \ref{stacks-geometry-lemma-closed-immersions}
montre que
$\dim_x (\mathcal{T}) \leq \dim_x(\mathcal{X})$ pour chaque
pour chaque composante irréductible $\mathcal{T}$ passant par
le point $x$. Pour démontrer le lemme,
il suffit donc de montrer que
\begin{equation}
\label{stacks-geometry-equation-desired-inequality}
\dim_x(\mathcal{X}) \leq
\sup_{\mathcal{T}} \{\dim_x(\mathcal{T})\}.
\end{equation}
Soit $U\to\mathcal{X}$ un recouvrement lisse par un schéma. Si $T$ est une composante
irréductible de $U$, notons $\mathcal{T}$ l'adhérence de son image
dans $\mathcal{X}$; c'est une composante irréductible de $\mathcal{X}$. Soit
$u \in U$
un point s'envoyant sur $x$. Alors
$\dim_x(\mathcal{X})=\dim_uU-\dim_uU_x=\sup_T\dim_uT-\dim_uU_x$, où le
supremum est pris sur les composantes irréductibles de $U$ passant
par $u$. Choisissons une composante $T$ sur laquelle le supremum
est atteint, et notons que
$\dim_x(\mathcal{T})=\dim_uT-\dim_u T_x$.
L'inégalité voulue (\ref{stacks-geometry-equation-desired-inequality})
résulte maintenant de l'inégalité évidente $\dim_u T_x \leq \dim_u U_x.$
(Notons que, si $\Spec k \to \mathcal{X}$ est un représentant de $x$,
alors $T\times_{\mathcal{X}} \Spec k$ est un sous-espace fermé de
$U\times_{\mathcal{X}}
\Spec k$.)
\end{proof}

\begin{lemma}
\label{stacks-geometry-lemma-dimension-at-finite-type-point}
Si $\mathcal{X}$ est un champ algébrique localement noethérien, et si
$x \in |\mathcal{X}|$, alors,
pour tout sous-champ ouvert $\mathcal{V}$ de $\mathcal{X}$ contenant $x$,
il existe un point de type fini $x_0 \in |\mathcal{V}|$ tel que
$\dim_{x_0}(\mathcal{X}) = \dim_x(\mathcal{V})$.
\end{lemma}

\begin{proof}
Choisissons un morphisme lisse surjectif
$f:U \to \mathcal{X}$ dont la source est un schéma, et considérons la
fonction $u \mapsto \dim_{f(u)}(\mathcal{X});$
puisque le morphisme $|U| \to |\mathcal{X}|$ induit par $f$ est ouvert (car $f$
est lisse) et surjectif (par hypothèse),
et envoie les points de type fini sur des points de type fini (par la définition même
des points de type fini de $|\mathcal{X}|$),
il suffit de montrer que, pour tout $u \in U$ et tout voisinage ouvert de $u$,
il existe dans ce voisinage un point de type fini $u_0$ tel que
$\dim_{f(u_0)}(\mathcal{X}) = \dim_{f(u)}(\mathcal{X}).$
Avec cette reformulation
du problème, la surjectivité de $f$ n'étant plus requise,
nous pouvons remplacer $U$ par le voisinage ouvert du point $u$ considéré,
et nous ramener ainsi à montrer que, pour chaque $u \in U$,
il existe un point de type fini $u_0 \in U$ tel que
$\dim_{f(u_0)}(\mathcal{X}) = \dim_{f(u)}(\mathcal{X}).$
D'après le Lemme \ref{stacks-geometry-lemma-behaviour-of-dimensions-wrt-smooth-morphisms},
$\dim_{f(u)}(\mathcal{X}) = \dim_u(U) - \dim_u(U_{f(u)}),$
tandis que
$\dim_{f(u_0)}(\mathcal{X}) = \dim_{u_0}(U) - \dim_{u_0}(U_{f(u_0)}).$
Puisque $f$ est lisse, l'expression $\dim_{u_0}(U_{f(u_0)})$ est localement
constante lorsque $u_0$ varie dans $U$ (d'après le
Lemme \ref{stacks-geometry-lemma-relative-dimension-is-semi-continuous} (2));
quitte donc à rétrécir encore $U$ autour de
$u$, nous pouvons la supposer constante. Le problème
se réduit ainsi à montrer que nous pouvons trouver un point de type fini $u_0 \in U$
pour lequel $\dim_{u_0}(U) = \dim_u(U)$.
Comme, par définition, $\dim_u U$ est le minimum des dimensions
$\dim V$, lorsque $V$ parcourt les voisinages ouverts $V$ de $u$
dans $U$, nous pouvons encore rétrécir $U$ autour de $u$ de sorte que
$\dim_u U = \dim U$.
L'existence du point voulu $u_0$ résulte alors du
Lemme \ref{stacks-geometry-lemma-dimension-achieved-by-finite-type-point}.
\end{proof}

\begin{lemma}
\label{stacks-geometry-lemma-monomorphing-a-component-in-of-the-right-dimension}
Soit $\mathcal{T} \hookrightarrow \mathcal{X}$ un monomorphisme localement
de type fini de champs algébriques,
où $\mathcal{X}$ (et donc aussi $\mathcal{T}$)
est jacobsonien, pseudo-caténaire et localement noethérien.
Supposons en outre que $\mathcal{T}$ soit irréductible,
de dimension (finie) $d$, et que $\mathcal{X}$ soit réduit
et de dimension inférieure
ou égale à $d$.
Alors il existe un sous-champ ouvert non vide $\mathcal{V}$ de $\mathcal{T}$ tel
que le monomorphisme
induit $\mathcal{V} \hookrightarrow \mathcal{X}$ soit une immersion ouverte
qui identifie
$\mathcal{V}$ à un sous-ensemble ouvert d'une composante irréductible de $\mathcal{X}$.
\end{lemma}

\begin{proof}
Choisissons un morphisme lisse surjectif $f:U \to \mathcal{X}$ de source un schéma,
nécessairement réduit puisque $\mathcal{X}$ l'est,
et écrivons $U' = \mathcal{T}\times_{\mathcal{X}} U$. Le morphisme obtenu par changement de base
$U' \to U$ est un monomorphisme d'espaces algébriques, localement de type
fini, donc représentable d'après
Morphismes d'espaces, Lemmes
\ref{spaces-morphisms-lemma-locally-quasi-finite-separated-representable} et
\ref{spaces-morphisms-lemma-monomorphism-loc-finite-type-loc-quasi-finite};
puisque $U$ est un schéma, $U'$ l'est aussi.
La projection $f': U' \to \mathcal{T}$ est encore une surjection lisse.
Soit $u' \in U'$, d'image $u \in U$.
Le Lemme \ref{stacks-geometry-lemma-base-change-invariance-of-relative-dimension}
montre que $\dim_{u'}(U'_{f(u')}) = \dim_u(U_{f(u)}),$
tandis que $\dim_{f'(u')}(\mathcal{T}) =d
\geq \dim_{f(u)}(\mathcal{X})$ d'après le
Lemme \ref{stacks-geometry-lemma-irreducible-implies-equidimensional}
et nos hypothèses sur $\mathcal{T}$ et $\mathcal{X}$.
Nous voyons donc que
\begin{equation}
\label{stacks-geometry-equation-dim-inequality}
\dim_{u'} (U') = \dim_{u'} (U'_{f(u')}) + \dim_{f'(u')}(\mathcal{T})
\\
\geq \dim_u (U_{f(u)}) + \dim_{f(u)}(\mathcal{X}) = \dim_u (U).
\end{equation}
Puisque $U' \to U$ est un monomorphisme localement de type fini,
il est en particulier non ramifié;
par la structure locale étale des morphismes non ramifiés,
Morphismes étales, Lemme \ref{etale-lemma-finite-unramified-etale-local},
nous pouvons donc trouver un diagramme commutatif
$$
\xymatrix{
V' \ar[r]\ar[d] & V \ar[d] \\
U' \ar[r] & U
}
$$
où le schéma $V'$ est non vide,
les flèches verticales sont étales,
et la flèche horizontale supérieure est une immersion fermée.
En remplaçant $V$ par un sous-ensemble ouvert quasi-compact
dont l'image rencontre l'image de $U'$,
et $V'$ par l'image réciproque de $V$, nous pouvons en outre
supposer que $V$ (et donc $V'$) est quasi-compact.
Puisque $V$ est aussi localement noethérien,
il est donc noethérien, et est par conséquent la réunion d'un nombre fini de composantes
irréductibles.

\medskip\noindent
Puisque les morphismes étales préservent la dimension ponctuelle, d'après
Descente, Lemme \ref{descent-lemma-dimension-at-point-local},
nous déduisons de (\ref{stacks-geometry-equation-dim-inequality})
que, pour tout point $v' \in V'$,
d'image $v \in V$, nous avons
$\dim_{v'}( V') \geq \dim_v(V)$.
En particulier, l'image de $V'$ ne peut être contenue dans l'intersection
de deux composantes irréductibles distinctes de $V$; nous pouvons donc trouver
au moins un sous-ensemble ouvert irréductible de $V$ ayant une intersection non vide
avec $V'$; en remplaçant $V$ par cet ouvert, nous pouvons supposer $V$ intègre
(car à la fois réduit et irréductible). L'inégalité précédente
sur les dimensions permet de conclure que l'immersion fermée $V' \hookrightarrow V$
est en fait un isomorphisme.
Si nous notons $W$ l'image de $V'$
dans $U'$, alors $W$ est un sous-ensemble
ouvert non vide de $U'$ (car les morphismes étales sont ouverts),
et le monomorphisme induit $W \to U$ est étale
(puisqu'il l'est localement pour la topologie étale sur la source, c'est-à-dire après changement de base
à $V'$), et est donc une immersion ouverte (étant un monomorphisme étale).
Ainsi, si nous notons $\mathcal{V}$ l'image de $W$ dans $\mathcal{T}$,
alors $\mathcal{V}$ est un sous-champ ouvert dense (de manière équivalente, non vide) de
$\mathcal{T}$,
dont l'image est dense dans une composante irréductible de $\mathcal{X}$.
Enfin,
notons que le morphisme $\mathcal{V} \to \mathcal{X}$ est lisse
(puisque son composé
avec le morphisme lisse $W\to \mathcal{V}$ est lisse),
et est aussi un monomorphisme, donc une immersion ouverte.
\end{proof}

\begin{lemma}
\label{stacks-geometry-lemma-dims-of-images}
Soit $f: \mathcal{T} \to \mathcal{X}$ un morphisme localement de type fini
de champs algébriques jacobsoniens, pseudo-caténaires et localement
noethériens,
dont la source est irréductible et le but quasi-séparé,
et soit $\mathcal{Z} \hookrightarrow \mathcal{X}$ l'image
schématique de $\mathcal{T}$.
Alors, pour tout $t \in |T|$,
nous avons
$\dim_t( \mathcal{T}_{f(t)}) \geq \dim \mathcal{T}  - \dim \mathcal{Z}$,
et il existe un sous-ensemble ouvert non vide (de manière équivalente, dense)
de $|\mathcal{T}|$ sur lequel l'égalité est satisfaite.
\end{lemma}

\begin{proof}
En remplaçant $\mathcal{X}$ par $\mathcal{Z}$, nous pouvons supposer, et supposons, que $f$ est
schématiquement dominant,
et aussi que $\mathcal{X}$ est irréductible.
Par la semi-continuité supérieure des dimensions des fibres
(Lemme \ref{stacks-geometry-lemma-relative-dimension-is-semi-continuous} (1)),
il suffit de démontrer que l'égalité
$\dim_t( \mathcal{T}_{f(t)}) =\dim \mathcal{T}  - \dim \mathcal{Z}$
est satisfaite pour $t$ appartenant à
un certain sous-champ ouvert non vide de $\mathcal{T}$.
Pour cette raison, nous pouvons à tout moment dans le raisonnement
remplacer $\mathcal{T}$ par un sous-champ ouvert non vide.

\medskip\noindent
Soit $T' \to \mathcal{T}$ un morphisme lisse surjectif dont la source
est un schéma, et soit $T$ un sous-ensemble ouvert non vide quasi-compact
de $T'$. Puisque $\mathcal{Y}$ est quasi-séparé, nous trouvons
que $T \to  \mathcal{Y}$ est quasi-compact
(d'après Morphismes de champs, Lemme
\ref{stacks-morphisms-lemma-quasi-compact-permanence}, appliqué aux morphismes
 $T \to \mathcal{Y} \to \Spec \mathbf{Z}$ auquel on l'applique).
Ainsi, en remplaçant $\mathcal{T}$ par l'image de $T$ dans $\mathcal{T}$,
nous pouvons supposer (en invoquant
Morphismes de champs, Lemme
\ref{stacks-morphisms-lemma-surjection-from-quasi-compact}
que le morphisme $f:\mathcal{T} \to \mathcal{X}$ est quasi-compact.

\medskip\noindent
Si nous choisissons une surjection lisse $U \to \mathcal{X}$ avec $U$ un schéma,
alors le Lemme \ref{stacks-geometry-lemma-map-of-components} assure que
nous pouvons trouver un sous-ensemble ouvert irréductible $V$ de $U$ tel
que $V \to \mathcal{X}$ soit lisse et schématiquement dominant.
Puisque la dominance schématique des morphismes quasi-compacts
est préservée par changement de base plat,
le changement de base $\mathcal{T} \times_{\mathcal{X}} V \to V$
du morphisme schématiquement
dominant $f$ est encore
schématiquement dominant. Soit $Z$ un schéma
admettant une surjection lisse sur ce produit fibré;
alors $Z \to \mathcal{T} \times_{\mathcal{X}} V \to V$
est encore schématiquement dominant.
Nous pouvons donc trouver une composante
irréductible $C$ de $Z$ qui domine
schématiquement $V$.
Puisque le composé $Z \to \mathcal{T}\times_{\mathcal{X}} V \to \mathcal{T}$
est lisse,
et puisque $\mathcal{T}$ est irréductible,
le Lemme \ref{stacks-geometry-lemma-map-of-components} montre que toute composante irréductible
de la source a une image dense dans $|\mathcal{T}|$.
Remplaçons maintenant
$C$ par un sous-ensemble ouvert non vide $W$ disjoint de toute autre
composante irréductible de $Z$, puis
remplaçons $\mathcal{T}$ et $\mathcal{X}$ par les images de $W$
et de $V$
(et appliquons le Lemme \ref{stacks-geometry-lemma-irreducible-implies-equidimensional}
pour voir que cela
ne change la dimension ni de $\mathcal{T}$ ni de $\mathcal{X}$).
Si nous notons $\mathcal{W}$ l'image du morphisme
$W \to \mathcal{T}\times_{\mathcal{X}} V$,
alors $\mathcal{W}$ est ouvert dans $\mathcal{T}\times_{\mathcal{X}} V$ (puisque le
morphisme $W \to \mathcal{T}\times_{\mathcal{X}} V$ est lisse),
et est irréductible (comme image d'un schéma
irréductible). Nous aboutissons ainsi à un diagramme commutatif
$$
\xymatrix{
W \ar[dr] \ar[r]  & \mathcal{W} \ar[r] \ar[d] & V \ar[d] \\
& \mathcal{T} \ar[r] & \mathcal{X}
}
$$
où $W$ et $V$ sont des schémas,
les flèches verticales sont lisses et surjectives,
les flèches diagonales et la flèche horizontale
supérieure de gauche sont lisses,
et le morphisme induit $\mathcal{W} \to \mathcal{T}\times_{\mathcal{X}} V$ est
une immersion ouverte. À l'aide de ce diagramme et des définitions
des diverses dimensions intervenant dans
l'énoncé du lemme, nous ramènerons la vérification
du lemme au cas des schémas, où elle est connue.

\medskip\noindent
Fixons $w \in |W|$, d'image $w' \in |\mathcal{W}|$,
d'image $t \in |\mathcal{T}|$, d'image $v$ dans $|V|$,
et d'image $x$ dans $|\mathcal{X}|$.
Essentiellement par définition (en utilisant le
fait que $\mathcal{W}$ est ouvert dans $\mathcal{T}\times_{\mathcal{X}} V$, et que
la fibre d'un changement de base est le changement de base de la fibre),
nous obtenons les égalités
$$
\dim_v V_x = \dim_{w'} \mathcal{W}_t
$$
et
$$
\dim_t \mathcal{T}_x = \dim_{w'} \mathcal{W}_v.
$$
D'après le Lemme \ref{stacks-geometry-lemma-behaviour-of-dimensions-wrt-smooth-morphisms}
(la flèche diagonale et la flèche verticale de droite
de notre diagramme réalisent $W$ et $V$ comme des recouvrements lisses par des
schémas des champs $\mathcal{T}$ et $\mathcal{X}$), nous trouvons
$$
\dim_t \mathcal{T} = \dim_w W - \dim_w W_t
$$
et
$$
\dim_x \mathcal{X} = \dim_v V - \dim_v V_x.
$$
En combinant ces égalités, nous trouvons
$$
\dim_t \mathcal{T}_x - \dim_t \mathcal{T} + \dim_x \mathcal{X}
= \dim_{w'} \mathcal{W}_v - \dim_w W + \dim_w W_t + \dim_v V -
\dim_{w'} \mathcal{W}_t
$$
Puisque $W \to \mathcal{W}$ est une surjection lisse, il en va de même
après changement de base par le morphisme $\Spec \kappa(v) \to V$
(en considérant $W \to \mathcal{W}$ comme un morphisme sur $V$),
et ce morphisme lisse donne la première des deux
égalités suivantes
$$
\dim_w W_v - \dim_{w'} \mathcal{W}_v = \dim_w (W_v)_{w'} = \dim_w W_{w'};
$$
la seconde égalité résulte d'une comparaison directe des
deux fibres en jeu.
De même, si nous considérons $W \to \mathcal{W}$ comme un morphisme de schémas
sur $\mathcal{T}$, et effectuons le changement de base par un représentant du point
$t \in |\mathcal{T}|$, nous obtenons les égalités
$$
\dim_w W_t - \dim_{w'} \mathcal{W}_t = \dim_w (W_t)_{w'} = \dim_w W_{w'}.
$$
En réunissant le tout, nous trouvons
$$
\dim_t \mathcal{T}_x - \dim_t \mathcal{T} + \dim_x \mathcal{X}
=  \dim_w W_v - \dim_w W + \dim_v V.
$$
Notre but est de montrer que le membre de gauche de cette égalité
s'annule sur un sous-ensemble ouvert non vide
de $t$. Lorsque $w$ parcourt un sous-ensemble ouvert non vide de $W$,
son image $t \in |\mathcal{T}|$ parcourt un sous-ensemble ouvert non vide
de $|\mathcal{T}|$ (puisque $W \to \mathcal{T}$ est lisse).

\medskip\noindent
Nous sommes donc ramenés à montrer que, si $W\to V$ est un
morphisme schématiquement dominant de schémas irréductibles localement
noethériens, localement de type fini,
alors il existe un sous-ensemble ouvert non vide de
points $w\in W$ tels que $\dim_w W_v =\dim_w W - \dim_v V$
(où $v$ désigne l'image de $w$ dans $V$).
C'est un fait classique,
dont nous rappelons la preuve pour la commodité du lecteur.

\medskip\noindent
Nous pouvons remplacer $W$ et $V$ par leurs sous-schémas réduits sous-jacents
sans modifier la validité (ou non) de cette équation,
et pouvons donc les supposer en fait intègres.
Puisque $\dim_w W_v$ est localement constante sur $W,$ quitte à remplacer $W$
par un sous-ensemble ouvert non vide, nous pouvons supposer que $\dim_w W_v$
est constante, disons égale à $d$. En choisissant cet ouvert affine,
nous pouvons aussi supposer que le morphisme $W\to V$ est de type fini.
Quitte à remplacer $V$ par un sous-ensemble ouvert non vide
(et à changer ensuite la base de $W$ à cet ouvert; le changement de base obtenu
est non vide, puisque le changement de base plat d'un morphisme quasi-compact
et schématiquement
dominant demeure schématiquement dominant),
nous pouvons en outre supposer que $W$ est plat sur $V$.
Le morphisme $W\to V$ est donc de dimension relative $d$
au sens de
Morphismes, Définition
\ref{morphisms-definition-relative-dimension-d}
et il résulte de
Morphismes, Lemme \ref{morphisms-lemma-rel-dimension-dimension},
que $\dim_w(W) = \dim_v(V) + d,$ comme voulu.
\end{proof}

\begin{remark}
\label{stacks-geometry-remark-negative-dimension}
Notons que, dans le contexte du lemme précédent,
on n'a pas nécessairement $\dim \mathcal{T} \geq \dim \mathcal{Z}$; cela ne
contredit pas l'inégalité de l'énoncé du lemme, car
les fibres du morphisme $f$ sont encore des champs algébriques, et
peuvent donc avoir une dimension négative. On l'illustre en prenant
pour $k$ un corps et en appliquant le lemme au morphisme
$[\Spec k/\mathbf{G}_m] \to \Spec k$.

\medskip\noindent
Si le morphisme $f$ de l'énoncé du lemme est supposé
quasi-DM (au sens de
Morphismes de champs, Définition
\ref{stacks-morphisms-definition-separated}; les morphismes
représentables par des espaces algébriques sont par exemple quasi-DM),
alors les fibres du morphisme aux points du but
sont des champs algébriques quasi-DM, et ont donc une dimension
positive ou nulle. Dans ce cas, le lemme entraîne
que l'on a bien $\dim \mathcal{T} \geq \dim \mathcal{Z}$. En fait, nous obtenons
le résultat plus général suivant.
\end{remark}

\begin{lemma}
\label{stacks-geometry-lemma-dims-of-images-two}
Soit $f: \mathcal{T} \to \mathcal{X}$ un morphisme localement de type fini
de champs algébriques jacobsoniens, pseudo-caténaires et localement
noethériens,
qui est quasi-DM,
dont la source est irréductible et le but quasi-séparé,
et soit $\mathcal{Z} \hookrightarrow \mathcal{X}$ l'image
schématique de $\mathcal{T}$.
Alors $\dim \mathcal{Z} \leq \dim \mathcal{T}$,
et, de plus, exactement l'une des deux conditions suivantes est satisfaite :
\begin{enumerate}
\item pour tout point de type fini $t \in |T|,$
nous avons
$\dim_t(\mathcal{T}_{f(t)}) > 0,$ auquel
cas $\dim \mathcal{Z} < \dim \mathcal{T}$; ou
\item $\mathcal{T}$ et $\mathcal{Z}$
ont la même dimension.
\end{enumerate}
\end{lemma}

\begin{proof}
Comme nous l'avons observé dans la remarque précédente,
la dimension d'un champ quasi-DM est toujours positive ou nulle,
d'où nous concluons que $\dim_t \mathcal{T}_{f(t)} \geq 0$
pour tout $t \in |\mathcal{T}|$, l'égalité
$$
\dim_t \mathcal{T}_{f(t)} = \dim_t \mathcal{T} - \dim_{f(t)} \mathcal{Z}
$$
étant satisfaite
sur un sous-ensemble ouvert dense de points $t\in |\mathcal{T}|$.
\end{proof}





\section{La dimension de l'anneau local}
\label{stacks-geometry-section-dimension-local-ring}

\noindent
Un champ algébrique ne possède pas véritablement d'anneaux locaux au sens usuel,
mais nous pouvons définir la dimension de l'anneau local comme suit.

\begin{lemma}
\label{stacks-geometry-lemma-dimension-local-ring-pre}
Soit $\mathcal{X}$ un champ algébrique localement noethérien.
Soit $U \to \mathcal{X}$ un morphisme lisse et soit $u \in U$.
Alors
$$
\dim(\mathcal{O}_{U, \overline{u}}) -
\dim(\mathcal{O}_{R_u, e(\overline{u})}) =
2\dim(\mathcal{O}_{U, \overline{u}}) -
\dim(\mathcal{O}_{R, e(\overline{u})})
$$
Ici, $R = U \times_\mathcal{X} U$, de projections $s, t : R \to U$ et de
diagonale $e : U \to R$, et $R_u$ est la fibre de $s : R \to U$ en $u$.
\end{lemma}

\begin{proof}
Cela est vrai parce que
$s : \mathcal{O}_{U, \overline{u}} \to \mathcal{O}_{R, e(\overline{u})}$
est un homomorphisme local plat d'anneaux locaux noethériens, et donc
$$
\dim(\mathcal{O}_{R, e(\overline{u})}) =
\dim(\mathcal{O}_{U, \overline{u}}) +
\dim(\mathcal{O}_{R_u, e(\overline{u})})
$$
d'après
Algèbre, Lemme \ref{algebra-lemma-dimension-base-fibre-equals-total}.
\end{proof}

\begin{lemma}
\label{stacks-geometry-lemma-dimension-local-ring}
Soit $\mathcal{X}$ un champ algébrique localement noethérien.
Soit $x \in |\mathcal{X}|$ un point de type fini 
Morphismes de champs, Définition
\ref{stacks-morphisms-definition-finite-type-point}).
Soit $d \in \mathbf{Z}$.
Les conditions suivantes sont équivalentes :
\begin{enumerate}
\item il existe un schéma $U$, un morphisme lisse $U \to \mathcal{X}$,
et un point de type fini $u \in U$ s'envoyant sur $x$, tels que
$2\dim(\mathcal{O}_{U, \overline{u}}) -
\dim(\mathcal{O}_{R, e(\overline{u})}) = d$, et
\item pour tout schéma $U$, tout morphisme lisse $U \to \mathcal{X}$,
et tout point de type fini $u \in U$ s'envoyant sur $x$, nous avons
$2\dim(\mathcal{O}_{U, \overline{u}}) -
\dim(\mathcal{O}_{R, e(\overline{u})}) = d$.
\end{enumerate}
Ici, $R = U \times_\mathcal{X} U$, de projections $s, t : R \to U$ et de
diagonale $e : U \to R$, et $R_u$ est la fibre de $s : R \to U$ en $u$.
\end{lemma}

\begin{proof}
Supposons donnés deux voisinages lisses $(U, u)$ et $(U', u')$
de $x$, où $u$ et $u'$ sont des points de type fini. Quitte à rétrécir $U$
et $U'$, nous pouvons supposer que $u$ et $u'$ sont des points fermés
(par définition des points de type fini). Choisissons alors
un morphisme étale surjectif $W \to U \times_\mathcal{X} U'$.
Soit $W_u$ la fibre de $W \to U$ en $u$, et
soit $W_{u'}$ la fibre de $W \to U'$ en $u'$.
Puisque $u$ et $u'$ s'envoient sur le même point de $|\mathcal{X}|$,
nous voyons que $W_u \cap W_{u'}$ est non vide.
Nous pouvons donc choisir un point fermé $w \in W$ s'envoyant à la fois
sur $u$ et sur $u'$. Cela nous ramène à la discussion du
paragraphe suivant.

\medskip\noindent
Supposons que $(U', u') \to (U, u)$ soit un morphisme lisse de voisinages
lisses de $x$, avec $u$ et $u'$ des points fermés.
Le but est de démontrer que l'invariant défini pour $(U, u)$ est égal
à celui défini pour $(U', u')$.
Pour le voir, observons que $\mathcal{O}_{U, u} \to \mathcal{O}_{U', u'}$
est un homomorphisme local plat d'anneaux locaux noethériens, et donc
$$
\dim(\mathcal{O}_{U', \overline{u}'}) =
\dim(\mathcal{O}_{U, \overline{u}}) +
\dim(\mathcal{O}_{U'_u, \overline{u}'})
$$
d'après Algèbre, Lemme \ref{algebra-lemma-dimension-base-fibre-equals-total}.
(Nous omettons le détail de toutes les étapes reliant les propriétés des anneaux
locaux et de leurs hensélisés stricts; voir
Compléments d'algèbre, section \ref{more-algebra-section-permanence-henselization}).
D'autre part, nous avons
$$
R' = U' \times_{U, t} R \times_{s, U} U'
$$
Nous voyons donc que
$$
\dim(\mathcal{O}_{R', e(\overline{u}')}) =
\dim(\mathcal{O}_{R, e(\overline{u})}) +
\dim(\mathcal{O}_{U'_u \times_u U'_u, (\overline{u}', \overline{u}')})
$$
Pour démontrer le lemme, il suffit de montrer que
$$
\dim(\mathcal{O}_{U'_u \times_u U'_u, (\overline{u}', \overline{u}')}) =
2\dim(\mathcal{O}_{U'_u, \overline{u}'})
$$
Observons que ceci n'est pas toujours vrai (par exemple, si $U'_u$ est une courbe
et $u'$ le point générique de cette courbe). Nous savons toutefois
que $u'$ est un point fermé de l'espace algébrique $U'_u$, localement
de type fini sur $u$. Dans ce cas, l'égalité est satisfaite car, premièrement,
$\dim_{(u', u')}(U'_u \times_u U'_u) = 2\dim_{u'}(U'_u)$ d'après
Variétés, Lemme \ref{varieties-lemma-dimension-product-locally-algebraic},
et, deuxièmement, parce que la dimension coïncide avec celle des anneaux locaux
aux points fermés des schémas localement algébriques; voir
Variétés, Lemme \ref{varieties-lemma-dimension-locally-algebraic}.
Nous omettons la transposition de ces résultats pour les schémas dans le
langage des espaces algébriques.
\end{proof}

\begin{definition}
\label{stacks-geometry-definition-dimension-local-ring}
Soit $\mathcal{X}$ un champ algébrique localement noethérien.
Soit $x \in |\mathcal{X}|$ un point de type fini.
La {\it dimension de l'anneau local de $\mathcal{X}$ en $x$}
est $d \in \mathbf{Z}$ si les conditions équivalentes du
Lemme \ref{stacks-geometry-lemma-dimension-local-ring} sont satisfaites.
\end{definition}

\noindent
Bien entendu, ceci est motivé par le Lemme \ref{stacks-geometry-lemma-dimension-local-ring-pre}
et Propriétés des champs, Définition
\ref{stacks-properties-definition-dimension-at-point}.
Terminons cette section en établissant une formule permettant de
calculer $\dim_x(\mathcal{X})$ à partir des propriétés de l'anneau versel
de $\mathcal{X}$ en $x$.

\begin{lemma}
\label{stacks-geometry-lemma-dimension-formula}
Supposons que $\mathcal{X}$ soit un champ algébrique, localement de type fini
sur un schéma localement noethérien $S$. Soit $x_0 : \Spec(k) \to \mathcal{X}$
un morphisme, où $k$ est un corps de type fini sur $S$. Représentons
$\mathcal{F}_{\mathcal{X}, k, x_0}$ comme dans la Remarque \ref{stacks-geometry-remark-groupoid-defo}
par un cogroupoïde $(A, B, s, t, c)$ de $S$-algèbres locales noethériennes complètes
de corps résiduel $k$. Alors
$$
\text{la dimension de l'anneau local de }\mathcal{X}\text{ en }x_0 =
2\dim A - \dim B
$$
\end{lemma}

\begin{proof}
Soit $s \in S$ l'image de $x_0$. Si $\mathcal{O}_{S, s}$
est un anneau G (condition presque toujours satisfaite en pratique),
nous pouvons démontrer le lemme comme suit.
D'après le Lemme \ref{stacks-geometry-lemma-Artin-approximation-by-smooth-morphism},
nous pouvons trouver un morphisme lisse $U \to \mathcal{X}$, de source un schéma,
contenant un point $u_0 \in U$ de corps résiduel $k$, tel que le
morphisme induit $\Spec(k) \to U \to \mathcal{X}$ coïncide avec $x_0$
et que $A = \mathcal{O}_{U, u_0}^\wedge$.
Écrivons $R = U \times_\mathcal{X} U$. Nous pouvons alors identifier
$\mathcal{O}_{R, e(u_0)}^\wedge$ à $B$.
L'égalité résulte donc des définitions.

\medskip\noindent
Dans la suite de cette preuve, nous expliquons comment démontrer le lemme en
général, mais nous invitons vivement le lecteur à passer ce qui suit.

\medskip\noindent
Montrons d'abord que le membre de droite est indépendant du choix
de $(A, B, s, t, c)$. Supposons en effet que $(A', B', s', t', c')$
soit un second choix. Puisque $A$ et $A'$ sont des anneaux versels de $\mathcal{X}$
en $x_0$, nous pouvons choisir, quitte à échanger $A$ et $A'$,
une application formellement lisse $A \to A'$ compatible avec les objets formels
versels donnés $\xi$ et $\xi'$ sur $A$ et $A'$.
Rappelons que $\widehat{\mathcal{C}}_\Lambda$ possède
des coproduits, donnés par le produit tensoriel complété
sur $\Lambda$; voir Théorie formelle des déformations, Lemme
\ref{formal-defos-lemma-CLambdahat-coproducts}.
Alors $B$ proreprésente le foncteur des isomorphismes entre les
deux images directes de $\xi$ dans $A \widehat{\otimes}_\Lambda A$.
Il en va de même pour $B'$. Nous en concluons que
$$
B' =
B \otimes_{(A \widehat{\otimes}_\Lambda A)}
(A' \widehat{\otimes}_\Lambda A')
$$
On voit immédiatement que
$$
A \widehat{\otimes}_\Lambda A \longrightarrow
A \widehat{\otimes}_\Lambda A' \longrightarrow
A' \widehat{\otimes}_\Lambda A'
$$
est formellement lisse, de dimension relative égale à $2$ fois la
dimension relative de l'application formellement lisse $A \to A'$.
(Cela résulte de principes généraux, mais aussi du fait que,
dans ce cas particulier, $A'$ est un anneau de séries formelles sur $A$
en $r$ variables.) Ainsi, $B \to B'$ est formellement lisse, de
dimension relative $2(\dim(A') - \dim(A))$, comme voulu.

\medskip\noindent
Ensuite, soit $l/k$ une extension finie. Soit
$x_{l, 0} : \Spec(l) \to \mathcal{X}$
le point induit. Nous affirmons que le membre de droite de la formule
est le même pour $x_0$ que pour $x_{l, 0}$.
On le montre en choisissant $A \to A'$ comme dans le
Lemme \ref{stacks-geometry-lemma-versal-ring-field-extension}
et en raisonnant exactement comme au paragraphe précédent.
Nous omettons les détails.

\medskip\noindent
Enfin, en raisonnant comme dans la preuve du Lemme \ref{stacks-geometry-lemma-versal-ring-flat},
nous pouvons utiliser les compatibilités des deux paragraphes précédents
pour nous ramener au cas (étudié au premier paragraphe)
où $A$ est l'anneau local complet de $U$ en $u_0$, pour un certain
schéma lisse sur $\mathcal{X}$ et un point de type fini $u_0$.
Détails omis.
\end{proof}










% Bien, commençons ici.
%

\chapter{Champs de modules}


\phantomsection
\label{moduli-section-phantom}





\section{Introduction}
\label{moduli-section-introduction}

\noindent
Dans ce chapitre, nous vérifions des propriétés élémentaires d'espaces
et de champs de modules tels que
$\mathit{Hom}$, $\mathit{Isom}$, $\Cohstack_{X/B}$,
$\Quotfunctor_{\mathcal{F}/X/B}$, $\Hilbfunctor_{X/B}$,
$\Picardstack_{X/B}$, $\Picardfunctor_{X/B}$, $\mathit{Mor}_B(Z, X)$,
$\Spacesstack'_{fp, flat, proper}$, $\Polarizedstack$ et
$\Complexesstack_{X/B}$.
Nous avons déjà montré, sous des hypothèses convenables, qu'il s'agit
d'espaces algébriques ou de champs algébriques ; voir Quotients, sections
\ref{quot-section-hom},
\ref{quot-section-isom},
\ref{quot-section-stack-coherent-sheaves},
\ref{quot-section-not-flat},
\ref{quot-section-quot},
\ref{quot-section-hilb},
\ref{quot-section-picard-stack},
\ref{quot-section-picard-functor},
\ref{quot-section-relative-morphisms},
\ref{quot-section-stack-of-spaces},
\ref{quot-section-polarized} et
\ref{quot-section-moduli-complexes}.
Le champ des courbes, noté $\textit{Curves}$ et introduit dans
Quotients, section \ref{quot-section-curves}, est étudié au chapitre
sur les modules de courbes ; voir Modules des courbes, section
\ref{moduli-curves-section-stack-curves}.

\medskip\noindent
En un certain sens, ce chapitre suit les traces des exposés de
Grothendieck \cite{Gr-I},
\cite{Gr-II},
\cite{Gr-III},
\cite{Gr-IV},
\cite{Gr-V} et
\cite{Gr-VI}.






\section{Conventions et abus de langage}
\label{moduli-section-conventions}

\noindent
Nous continuons d'employer les conventions et les abus de langage
introduits dans Propriétés des champs, section
\ref{stacks-properties-section-conventions}.
Sauf mention contraire, notre schéma de base sera $\Spec(\mathbf{Z})$.






\section{Propriétés de Hom et Isom}
\label{moduli-section-hom-isom}

\noindent
Soit $f : X \to B$ un morphisme d'espaces algébriques de présentation
finie. Supposons que $\mathcal{F}$ et $\mathcal{G}$ soient des
$\mathcal{O}_X$-modules quasi-cohérents.
Si $\mathcal{G}$ est de présentation finie, plat sur $B$ et à support
propre sur $B$, alors le foncteur
$\mathit{Hom}(\mathcal{F}, \mathcal{G})$ défini par
$$
T/B \longmapsto \Hom_{\mathcal{O}_{X_T}}(\mathcal{F}_T, \mathcal{G}_T)
$$
est un espace algébrique affine sur $B$. Si $\mathcal{F}$ est de
présentation finie, alors
$\mathit{Hom}(\mathcal{F}, \mathcal{G}) \to B$
est de présentation finie. Voir Quotients, proposition
\ref{quot-proposition-hom}.

\medskip\noindent
Si $\mathcal{F}$ et $\mathcal{G}$ sont tous deux de présentation finie,
plats sur $B$ et à support propre sur $B$, alors le sous-foncteur
$$
\mathit{Isom}(\mathcal{F}, \mathcal{G}) \subset
\mathit{Hom}(\mathcal{F}, \mathcal{G})
$$
est un espace algébrique affine et de présentation finie sur $B$.
Voir Quotients, proposition \ref{quot-proposition-isom}.




\section{Propriétés du champ des faisceaux cohérents}
\label{moduli-section-stack-coherent-sheaves}

\noindent
Soit $f : X \to B$ un morphisme d'espaces algébriques séparé et de
présentation finie. Alors le champ $\Cohstack_{X/B}$ qui paramètre les
familles plates de modules cohérents à support propre est algébrique.
Voir Quotients, théorème
\ref{quot-theorem-coherent-algebraic-general}.

\begin{lemma}
\label{moduli-lemma-coherent-diagonal-affine-fp}
La diagonale de $\Cohstack_{X/B}$ sur $B$ est affine et de présentation
finie.
\end{lemma}

\begin{proof}
La représentabilité de la diagonale par des espaces algébriques a été
établie dans Quotients, lemme \ref{quot-lemma-coherent-diagonal}.
Il ressort de la démonstration qu'il nous faut montrer que
$\mathit{Isom}(\mathcal{F}, \mathcal{G}) \to T$
est affine et de présentation finie pour une paire de
$\mathcal{O}_{X_T}$-modules de présentation finie
$\mathcal{F}$, $\mathcal{G}$, plats sur $T$ et à support propre sur $T$.
C'est précisément ce qui a été étudié à la section
\ref{moduli-section-hom-isom}.
\end{proof}

\begin{lemma}
\label{moduli-lemma-coherent-qs-lfp}
Le morphisme $\Cohstack_{X/B} \to B$ est quasi-séparé et localement de
présentation finie.
\end{lemma}

\begin{proof}
Pour vérifier que $\Cohstack_{X/B} \to B$ est quasi-séparé, il faut
montrer que sa diagonale est quasi-compacte et quasi-séparée.
Cela résulte immédiatement du lemme
\ref{moduli-lemma-coherent-diagonal-affine-fp}.
Pour prouver que $\Cohstack_{X/B} \to B$ est localement de présentation
finie, il faut montrer que $\Cohstack_{X/B} \to B$ préserve les limites ;
voir Limites de champs, proposition
\ref{stacks-limits-proposition-characterize-locally-finite-presentation}.
Cela découle de Quotients, lemme \ref{quot-lemma-coherent-limits}
(un petit détail est omis).
\end{proof}

\begin{lemma}
\label{moduli-lemma-coherent-existence-part}
Supposons que $X \to B$ soit propre ainsi que de présentation finie.
Alors $\Cohstack_{X/B} \to B$ satisfait la partie d'existence du critère
valuatif (Morphismes de champs, définition
\ref{stacks-morphisms-definition-existence}).
\end{lemma}

\begin{proof}
Après changement de base, on se ramène immédiatement au problème suivant :
étant donnés un anneau de valuation $R$ de corps des fractions $K$, un espace
algébrique $X$ propre sur $R$ et un $\mathcal{O}_{X_K}$-module cohérent
$\mathcal{F}_K$, montrer qu'il existe un $\mathcal{O}_X$-module
$\mathcal{F}$ de présentation finie, plat sur $R$, dont la fibre générique
est $\mathcal{F}_K$.
Remarquons que, d'après Platitude sur les espaces, théorème
\ref{spaces-flat-theorem-finite-type-flat}, tout $\mathcal{O}_X$-module
quasi-cohérent $\mathcal{F}$ de type fini et plat sur $R$ est de
présentation finie.
Notons $j : X_K \to X$ l'immersion de la fibre générique.
Puisqu'il est obtenu par changement de base à partir du morphisme affine
$\Spec(K) \to \Spec(R)$, le morphisme $j$ est affine. Ainsi,
$j_*\mathcal{F}_K$ est quasi-cohérent. Écrivons
$$
j_*\mathcal{F}_K = \colim \mathcal{F}_i
$$
comme limite inductive filtrante de ses sous-$\mathcal{O}_X$-modules
quasi-cohérents de type fini ; voir Limites d'espaces, lemme
\ref{spaces-limits-lemma-directed-colimit-finite-type}.
Comme $j_*\mathcal{F}_K$ est un faisceau de $K$-espaces vectoriels sur
$X$, il est plat sur $\Spec(R)$. Chaque $\mathcal{F}_i$ est donc plat sur
$R$, car, sur un anneau de valuation, la platitude équivaut à l'absence de
torsion (Compléments d'algèbre, lemme
\ref{more-algebra-lemma-valuation-ring-torsion-free-flat}) et l'absence de
torsion passe aux sous-modules.
Enfin, nous devons montrer que l'application
$j^*\mathcal{F}_i \to \mathcal{F}_K$
est un isomorphisme pour un certain $i$.
Puisque $j^*j_*\mathcal{F}_K = \mathcal{F}_K$ (un petit détail est omis)
et que $j^*$ est exact, on voit que
$j^*\mathcal{F}_i \to \mathcal{F}_K$ est injective pour tout $i$.
Comme $j^*$ commute aux limites inductives, on a
$\mathcal{F}_K = j^*j_*\mathcal{F}_K = \colim j^*\mathcal{F}_i$.
Puisque $\mathcal{F}_K$ est cohérent (c'est-à-dire de présentation finie),
il existe un $i$ tel que $j^*\mathcal{F}_i$ contienne tous les générateurs,
en nombre fini, sur un recouvrement affine \'etale de $X$.
On obtient ainsi la surjectivité de
$j^*\mathcal{F}_i \to \mathcal{F}_K$ pour $i$ assez grand.
\end{proof}

\begin{lemma}
\label{moduli-lemma-coherent-functorial}
Soit $B$ un espace algébrique. Soit $\pi : X \to Y$ un morphisme
quasi-fini d'espaces algébriques séparés et de présentation finie sur $B$.
Alors $\pi_*$ induit un morphisme
$\Cohstack_{X/B} \to \Cohstack_{Y/B}$.
\end{lemma}

\begin{proof}
Soit $(T \to B, \mathcal{F})$ un objet de $\Cohstack_{X/B}$.
Nous affirmons que
\begin{enumerate}
\item[(a)] $(T \to B, \pi_{T, *}\mathcal{F})$ est un objet de
$\Cohstack_{Y/B}$, et que
\item[(b)] pour $T' \to T$, on a
$\pi_{T', *}(X_{T'} \to X_T)^*\mathcal{F} =
(Y_{T'} \to Y_T)^*\pi_{T, *}\mathcal{F}$.
\end{enumerate}
La partie (b) garantit que cette construction définit bien le foncteur
$\Cohstack_{X/B} \to \Cohstack_{Y/B}$ recherché.

\medskip\noindent
Soit $i : Z \to X_T$ le sous-espace fermé défini par le zéro-ième idéal
de Fitting de $\mathcal{F}$ (Diviseurs sur les espaces, section
\ref{spaces-divisors-section-fitting-ideals}).
Alors $Z \to B$ est propre par hypothèse (voir Catégories dérivées des
espaces, section \ref{spaces-perfect-section-proper-over-base}).
D'autre part, $i$ est de présentation finie (Diviseurs sur les espaces,
lemme \ref{spaces-divisors-lemma-fitting-ideal-of-finitely-presented}, et
Morphismes d'espaces, lemme
\ref{spaces-morphisms-lemma-closed-immersion-finite-presentation}).
Il existe un $\mathcal{O}_Z$-module quasi-cohérent $\mathcal{G}$ de type
fini tel que $i_*\mathcal{G} = \mathcal{F}$ (Diviseurs sur les espaces,
lemme \ref{spaces-divisors-lemma-on-subscheme-cut-out-by-Fit-0}).
En fait, $\mathcal{G}$ est de présentation finie comme
$\mathcal{O}_Z$-module, d'après Descente sur les espaces, lemme
\ref{spaces-descent-lemma-finite-finitely-presented-module}.
Remarquons que $\mathcal{G}$ est plat sur $B$, par exemple parce que les
germes de $\mathcal{G}$ et de $\mathcal{F}$ coïncident (Morphismes
d'espaces, lemme \ref{spaces-morphisms-lemma-stalk-push-closed}).
Remarquons que $\pi_T \circ i : Z \to Y_T$ est quasi-fini comme composée
de morphismes quasi-finis et que
$\pi_{T, *}\mathcal{F} = (\pi_T \circ i)_*\mathcal{G})$.
Puisque $i$ est affine, la formation de $i_*$ commute au changement de base
(Cohomologie des espaces, lemme
\ref{spaces-cohomology-lemma-affine-base-change}).
Nous pouvons donc remplacer $B$ par $T$, $X$ par $Z$,
$\mathcal{F}$ par $\mathcal{G}$ et $Y$ par $Y_T$ pour nous ramener au cas
étudié dans le paragraphe suivant.

\medskip\noindent
Supposons que $X \to B$ soit propre. Alors $\pi$ est propre d'après
Morphismes d'espaces, lemme
\ref{spaces-morphisms-lemma-universally-closed-permanence}, donc fini
d'après Compléments sur les morphismes d'espaces, lemme
\ref{spaces-more-morphisms-lemma-characterize-finite}.
Comme un morphisme fini est affine, la partie (b) résulte de Cohomologie
des espaces, lemme \ref{spaces-cohomology-lemma-affine-base-change}.
D'autre part, $\pi$ est de présentation finie d'après Morphismes d'espaces,
lemme \ref{spaces-morphisms-lemma-finite-presentation-permanence}.
Ainsi, $\pi_{T, *}\mathcal{F}$ est de présentation finie d'après Descente
sur les espaces, lemme
\ref{spaces-descent-lemma-finite-finitely-presented-module}.
Enfin, $\pi_{T, *}\mathcal{F} $ est plat sur $B$, comme on le voit par
exemple sur les germes en utilisant Cohomologie des espaces, lemme
\ref{spaces-cohomology-lemma-stalk-push-finite}.
\end{proof}

\begin{lemma}
\label{moduli-lemma-coherent-open}
Soit $B$ un espace algébrique. Soit $\pi : X \to Y$ une immersion ouverte
d'espaces algébriques séparés et de présentation finie sur $B$.
Alors le morphisme $\Cohstack_{X/B} \to \Cohstack_{Y/B}$ du lemme
\ref{moduli-lemma-coherent-functorial} est une immersion ouverte.
\end{lemma}

\begin{proof}
Omis. Indication : si $\mathcal{F}$ est un objet de $\Cohstack_{Y/B}$
sur $T$ et si, pour $t \in T$, on a
$\text{Supp}(\mathcal{F}_t) \subset |X_t|$, alors il en va de même pour
$t' \in T$ dans un voisinage de $t$.
\end{proof}

\begin{lemma}
\label{moduli-lemma-coherent-closed}
Soit $B$ un espace algébrique. Soit $\pi : X \to Y$ une immersion fermée
d'espaces algébriques séparés et de présentation finie sur $B$.
Alors le morphisme $\Cohstack_{X/B} \to \Cohstack_{Y/B}$ du lemme
\ref{moduli-lemma-coherent-functorial} est une immersion fermée.
\end{lemma}

\begin{proof}
Soit $\mathcal{I} \subset \mathcal{O}_Y$ le faisceau d'idéaux qui définit
$X$ comme sous-espace fermé de $Y$. Rappelons que $\pi_*$ induit une
équivalence entre la catégorie des $\mathcal{O}_X$-modules quasi-cohérents
et celle des $\mathcal{O}_Y$-modules quasi-cohérents annulés par
$\mathcal{I}$ ; voir Morphismes d'espaces, lemme
\ref{spaces-morphisms-lemma-i-star-equivalence}.
Il en va de même, mutatis mutandis, après changement de base par
$T \to B$, en remplaçant $\mathcal{I}$ par le faisceau d'idéaux
$\mathcal{I}_T = \Im((Y_T \to Y)^*\mathcal{I} \to \mathcal{O}_{Y_T})$.
L'analyse de la démonstration du lemme
\ref{moduli-lemma-coherent-functorial} montre que l'image essentielle de
$\Cohstack_{X/B} \to \Cohstack_{Y/B}$
est exactement formée des objets $\xi = (T \to B, \mathcal{F})$ pour
lesquels $\mathcal{F}$ est annulé par $\mathcal{I}_T$.
Autrement dit, $\xi$ appartient à l'image essentielle si et seulement si
l'application de multiplication
$$
\mathcal{F} \otimes_{\mathcal{O}_{Y_T}} (Y_T \to Y)^*\mathcal{I}
\longrightarrow
\mathcal{F}
$$
est nulle, de même qu'après tout changement de base ultérieur
$T' \to T$.
Notons que
$$
(Y_{T'} \to Y_T)^*(
\mathcal{F} \otimes_{\mathcal{O}_{Y_T}} (Y_T \to Y)^*\mathcal{I}) =
(Y_{T'} \to Y_T)^*\mathcal{F} \otimes_{\mathcal{O}_{Y_{T'}}}
(Y_{T'} \to Y)^*\mathcal{I})
$$
La nullité de l'application de multiplication sur $T'$ est donc
représentable par un sous-espace fermé de $T$, d'après Platitude sur les
espaces, lemme \ref{spaces-flat-lemma-F-zero-closed-proper}.
\end{proof}

\begin{situation}[Invariants numériques]
\label{moduli-situation-numerical}
Soit $f : X \to B$ comme dans l'introduction de cette section. Soit $I$
un ensemble et, pour $i \in I$, soit $E_i \in D(\mathcal{O}_X)$ parfait.
Étant donné un objet $(T \to B, \mathcal{F})$ de $\Cohstack_{X/B}$,
notons $E_{i, T}$ l'image inverse dérivée de $E_i$ sur $X_T$.
L'objet
$$
K_i = Rf_{T, *}(E_{i, T} \otimes_{\mathcal{O}_{X_T}}^\mathbf{L} \mathcal{F})
$$
de $D(\mathcal{O}_T)$ est parfait et sa formation commute au changement de
base ; voir Catégories dérivées des espaces, lemme
\ref{spaces-perfect-lemma-base-change-tensor-perfect}.
La fonction
$$
\chi_i : |T| \longrightarrow \mathbf{Z},\quad
\chi_i(t) =
\chi(X_t, E_{i, t} \otimes_{\mathcal{O}_{X_t}}^\mathbf{L} \mathcal{F}_t) =
\chi(K_i \otimes_{\mathcal{O}_T}^\mathbf{L} \kappa(t))
$$
est donc localement constante d'après Catégories dérivées des espaces,
lemme \ref{spaces-perfect-lemma-chi-locally-constant}.
Soit $P : I \to \mathbf{Z}$ une application. Considérons le sous-champ
$$
\Cohstack^P_{X/B} \subset \Cohstack_{X/B}
$$
formé des familles plates de faisceaux cohérents à support propre dont les
invariants numériques coïncident avec $P$. Plus précisément, un objet
$(T \to B, \mathcal{F})$ de $\Cohstack_{X/B}$ appartient à
$\Cohstack^P_{X/B}$ si et seulement si $\chi_i(t) = P(i)$ pour tout
$i \in I$ et tout $t \in T$.
\end{situation}

\begin{lemma}
\label{moduli-lemma-open-P}
Dans la situation \ref{moduli-situation-numerical}, le champ
$\Cohstack^P_{X/B}$ est algébrique et
$$
\Cohstack^P_{X/B} \longrightarrow \Cohstack_{X/B}
$$
est une immersion fermée plate. Si $I$ est fini ou si $B$ est localement
noethérien, alors $\Cohstack^P_{X/B}$ est un sous-champ ouvert et fermé de
$\Cohstack_{X/B}$.
\end{lemma}

\begin{proof}
C'est immédiat si $I$ est fini, puisque les fonctions
$t \mapsto \chi_i(t)$ sont localement constantes. Si $I$ est infini,
écrivons
$$
I = \bigcup\nolimits_{I' \subset I\text{ fini}} I'
$$
et notons $P' = P|_{I'}$. On a alors
$$
\Cohstack^P_{X/B} = \bigcap\nolimits_{I' \subset I\text{ fini}}
\Cohstack^{P'}_{X/B}
$$
Ainsi, $\Cohstack^P_{X/B}$ est toujours un champ algébrique et le morphisme
$\Cohstack^P_{X/B} \subset \Cohstack_{X/B}$ est toujours une immersion
fermée plate, mais il peut cesser d'être un sous-champ ouvert. (Nous
laissons au lecteur le soin de construire des exemples.) Toutefois, si
$B$ est localement noethérien, alors $\Cohstack_{X/B}$ l'est aussi, d'après
le lemme \ref{moduli-lemma-coherent-qs-lfp} et Morphismes de champs, lemme
\ref{stacks-morphisms-lemma-locally-finite-type-locally-noetherian}.
Par conséquent, si $U \to \Cohstack_{X/B}$ est un morphisme lisse
surjectif, où $U$ est un schéma localement noethérien, les images inverses
des sous-champs ouverts et fermés $\Cohstack^{P'}_{X/B}$ ont dans $U$ une
intersection ouverte (car les composantes connexes des espaces topologiques
localement noethériens sont ouvertes).
D'où le résultat dans ce cas.
\end{proof}

\begin{lemma}
\label{moduli-lemma-finite-list-perfect-objects}
Soit $f : X \to B$ comme dans l'introduction de cette section.
Soient $E_1, \ldots, E_r \in D(\mathcal{O}_X)$ parfaits.
Soit $I = \mathbf{Z}^{\oplus r}$ et considérons l'application
$$
I \longrightarrow D(\mathcal{O}_X),\quad
(n_1, \ldots, n_r) \longmapsto
E_1^{\otimes n_1}
\otimes \ldots \otimes
E_r^{\otimes n_r}
$$
Soit $P : I \to \mathbf{Z}$ une application. Alors
$\Cohstack^P_{X/B} \subset \Cohstack_{X/B}$,
tel qu'il est défini dans la situation \ref{moduli-situation-numerical}, est un
sous-champ ouvert et fermé.
\end{lemma}

\begin{proof}
Nous pouvons travailler localement pour la topologie \'etale sur $B$ et donc
supposer $B$ affine. Dans ce cas, nous pouvons effectuer une réduction
noethérienne absolue ; nous suggérons au lecteur de sauter la démonstration.
Écrivons en effet $B = \Spec(\Lambda)$.
Écrivons $\Lambda = \colim \Lambda_i$ comme limite inductive filtrante,
chaque $\Lambda_i$ étant de type fini sur $\mathbf{Z}$. Pour un certain
$i$, on peut trouver un morphisme d'espaces algébriques
$X_i \to \Spec(\Lambda_i)$ séparé et de présentation finie, dont le
changement de base à $\Lambda$ est $X$. Voir Limites d'espaces, lemmes
\ref{spaces-limits-lemma-descend-finite-presentation} et
\ref{spaces-limits-lemma-descend-separated-morphism}.
Après avoir augmenté $i$, on peut alors supposer qu'il existe des objets
parfaits $E_{1, i}, \ldots, E_{r, i}$ dans $D(\mathcal{O}_{X_i})$ dont
les images inverses dérivées sur $X$ sont respectivement isomorphes à
$E_1, \ldots, E_r$ ; voir Catégories dérivées des espaces, lemme
\ref{spaces-perfect-lemma-perfect-on-limit}.
On a manifestement un carré cartésien
$$
\xymatrix{
\Cohstack^P_{X/B} \ar[r] \ar[d] &
\Cohstack_{X/B} \ar[d] \\
\Cohstack^P_{X_i/\Spec(\Lambda_i)} \ar[r] &
\Cohstack_{X_i/\Spec(\Lambda_i)}
}
$$
et l'on peut donc appliquer le lemme \ref{moduli-lemma-open-P} pour achever la
démonstration.
\end{proof}

\begin{example}[Faisceaux cohérents à polynôme de Hilbert fixé]
\label{moduli-example-hilbert-polynomial}
Soit $f : X \to B$ comme dans l'introduction de cette section.
Soit $\mathcal{L}$ un $\mathcal{O}_X$-module inversible.
Soit $P : \mathbf{Z} \to \mathbf{Z}$ un polynôme numérique.
On peut alors considérer le sous-champ algébrique ouvert et fermé
$$
\Cohstack^P_{X/B} =
\Cohstack^{P, \mathcal{L}}_{X/B}
\subset \Cohstack_{X/B}
$$
formé des familles plates de faisceaux cohérents à support propre dont les
invariants numériques coïncident avec $P$ : un objet
$(T \to B, \mathcal{F})$ de $\Cohstack_{X/B}$ appartient à
$\Cohstack^P_{X/B}$ si et seulement si
$$
P(n) =
\chi(X_t, \mathcal{F}_t \otimes_{\mathcal{O}_{X_t}} \mathcal{L}_t^{\otimes n})
$$
pour tout $n \in \mathbf{Z}$ et tout $t \in T$. Il s'agit bien sûr d'un
cas particulier de la situation \ref{moduli-situation-numerical}, où
$I = \mathbf{Z} \to D(\mathcal{O}_X)$ est donné par
$n \mapsto \mathcal{L}^{\otimes n}$. Il résulte du lemme
\ref{moduli-lemma-finite-list-perfect-objects} qu'il s'agit d'un sous-champ ouvert
et fermé. Comme les fonctions
$n \mapsto
\chi(X_t, \mathcal{F}_t \otimes_{\mathcal{O}_{X_t}} \mathcal{L}_t^{\otimes n})$
sont toujours des polynômes numériques (Espaces sur les corps, lemme
\ref{spaces-over-fields-lemma-numerical-polynomial-from-euler}), on conclut
que
$$
\Cohstack_{X/B} = \coprod\nolimits_{P\text{ polynôme numérique}}
\Cohstack^P_{X/B}
$$
est une décomposition en union disjointe.
\end{example}







\section{Propriétés de Quot}
\label{moduli-section-quot}

\noindent
Soit $f : X \to B$ un morphisme d'espaces algébriques séparé et de
présentation finie. Soit $\mathcal{F}$ un $\mathcal{O}_X$-module
quasi-cohérent. Alors $\Quotfunctor_{\mathcal{F}/X/B}$ est un espace
algébrique. Si $\mathcal{F}$ est de présentation finie, alors
$\Quotfunctor_{\mathcal{F}/X/B} \to B$ est localement de présentation
finie. Voir Quotients, proposition \ref{quot-proposition-quot}.

\begin{lemma}
\label{moduli-lemma-quot-diagonal-closed}
La diagonale de $\Quotfunctor_{\mathcal{F}/X/B} \to B$ est une immersion
fermée. Si $\mathcal{F}$ est de type fini, la diagonale est une immersion
fermée de présentation finie.
\end{lemma}

\begin{proof}
Supposons donnés un schéma $T/B$ et deux quotients
$\mathcal{F}_T \to \mathcal{Q}_i$, $i = 1, 2$, correspondant à des points
à valeurs dans $T$ de $\Quotfunctor_{\mathcal{F}/X/B}$ au-dessus de $B$.
Notons $\mathcal{K}_1$ le noyau du premier et soit
$u : \mathcal{K}_1 \to \mathcal{Q}_2$ le morphisme composé.
D'après Platitude sur les espaces, lemme
\ref{spaces-flat-lemma-F-zero-closed-proper}, il existe un sous-espace fermé
de $T$ tel que $T' \to T$ se factorise par lui si et seulement si l'image
inverse $u_{T'}$ est nulle. Cela montre que la diagonale est une immersion
fermée. De plus, si $\mathcal{F}$ est de type fini, alors
$\mathcal{K}_1$ est de type fini (Modules sur les sites, lemme
\ref{sites-modules-lemma-kernel-surjection-finite-onto-finite-presentation})
et le même lemme montre que la diagonale est de présentation finie.
\end{proof}

\begin{lemma}
\label{moduli-lemma-quot-s-lfp}
Le morphisme $\Quotfunctor_{\mathcal{F}/X/B} \to B$ est séparé.
Si $\mathcal{F}$ est de présentation finie, il est aussi localement de
présentation finie.
\end{lemma}

\begin{proof}
Pour vérifier que $\Quotfunctor_{\mathcal{F}/X/B} \to B$ est séparé, il
faut montrer que sa diagonale est une immersion fermée. Cela résulte du
lemme \ref{moduli-lemma-quot-diagonal-closed}. La seconde assertion fait partie
de Quotients, proposition \ref{quot-proposition-quot}.
\end{proof}

\begin{lemma}
\label{moduli-lemma-quot-existence-part}
Supposons que $X \to B$ soit propre ainsi que de présentation finie et que
$\mathcal{F}$ soit quasi-cohérent de type fini.
Alors $\Quotfunctor_{\mathcal{F}/X/B} \to B$ satisfait la partie
d'existence du critère valuatif (Morphismes d'espaces, définition
\ref{spaces-morphisms-definition-valuative-criterion}).
\end{lemma}

\begin{proof}
Après changement de base, on se ramène immédiatement au problème suivant :
étant donnés un anneau de valuation $R$ de corps des fractions $K$, un
espace algébrique $X$ propre sur $R$, un $\mathcal{O}_X$-module
quasi-cohérent $\mathcal{F}$ de type fini et un quotient cohérent
$\mathcal{F}_K \to \mathcal{Q}_K$, montrer qu'il existe un quotient
$\mathcal{F} \to \mathcal{Q}$ où $\mathcal{Q}$ est un
$\mathcal{O}_X$-module de présentation finie, plat sur $R$, dont la fibre
générique est $\mathcal{Q}_K$.
Remarquons que, d'après Platitude sur les espaces, théorème
\ref{spaces-flat-theorem-finite-type-flat}, tout $\mathcal{O}_X$-module
quasi-cohérent $\mathcal{F}$ de type fini et plat sur $R$ est de
présentation finie. Nous établissons d'abord l'existence de $\mathcal{Q}$
localement sur les ouverts affines.

\medskip\noindent
Localement sur les ouverts affines, on aboutit au problème suivant :
soit $R \to A$ un homomorphisme d'anneaux de présentation finie, soit $M$
un $A$-module fini et soit $\varphi : M_K \to N_K$ un $A_K$-module
quotient. On peut alors considérer
$$
L = \{x \in M \mid \varphi(x \otimes 1) = 0 \}
$$
Le morphisme $M \to M/L$ définit un $A$-module quotient qui est sans
torsion comme $R$-module. Il est donc plat comme $R$-module (Compléments
d'algèbre, lemme
\ref{more-algebra-lemma-valuation-ring-torsion-free-flat}).
Puisque $M$ est fini comme $A$-module, il en va de même de $L$, et l'on
conclut que $L$ est de présentation finie comme $A$-module (par la
référence précédente). Manifestement, $M/L$ est l'unique quotient de ce
type tel que $(M/L)_K = N_K$.

\medskip\noindent
L'unicité dans la construction du paragraphe précédent garantit que ces
quotients se recollent et donnent le $\mathcal{Q}$ recherché. Donnons un
peu plus de détails. Choisissons un morphisme \'etale surjectif
$U \to X$, où $U$ est un schéma affine. La construction précédente donne
un quotient $\mathcal{F}|_U \to \mathcal{Q}_U$ qui est quasi-cohérent,
plat sur $R$ et redonne $\mathcal{Q}_K|U$ sur la fibre générique.
Comme $X$ est séparé, on voit que $U \times_X U$ est lui aussi un schéma
affine \'etale sur $X$.
Les quotients $\mathcal{F}|_{U \times_X U} \to
\text{pr}_1^*\mathcal{Q}_U$ et
$\mathcal{F}|_{U \times_X U} \to \text{pr}_2^*\mathcal{Q}_U$
coïncident alors par unicité de la construction. On peut donc descendre
$\mathcal{F}|_U \to \mathcal{Q}_U$ en une surjection
$\mathcal{F} \to \mathcal{Q}$, comme voulu (Propriétés des espaces,
proposition \ref{spaces-properties-proposition-quasi-coherent}).
\end{proof}

\begin{lemma}
\label{moduli-lemma-quot-functorial}
Soit $B$ un espace algébrique. Soit $\pi : X \to Y$ un morphisme
quasi-fini affine d'espaces algébriques séparés et de présentation finie sur
$B$. Soit $\mathcal{F}$ un $\mathcal{O}_X$-module quasi-cohérent.
Alors $\pi_*$ induit un morphisme
$\Quotfunctor_{\mathcal{F}/X/B} \to
\Quotfunctor_{\pi_*\mathcal{F}/Y/B}$.
\end{lemma}

\begin{proof}
Posons $\mathcal{G} = \pi_*\mathcal{F}$. Puisque $\pi$ est affine, pour
tout schéma $T$ sur $B$, on a
$\mathcal{G}_T = \pi_{T, *}\mathcal{F}_T$, d'après Cohomologie des
espaces, lemme \ref{spaces-cohomology-lemma-affine-base-change}.
De plus, $\pi_T$ est affine, donc $\pi_{T, *}$ est exact et transforme les
quotients en quotients. Remarquons qu'un quotient quasi-cohérent
$\mathcal{F}_T \to \mathcal{Q}$ définit un point de
$\Quotfunctor_{X/B}$ si et seulement si $\mathcal{Q}$ définit un objet de
$\Cohstack_{X/B}$ sur $T$ (et de même pour $\mathcal{G}$ et $Y$).
Comme nous avons vu au lemme \ref{moduli-lemma-coherent-functorial} que $\pi_*$
induit un morphisme $\Cohstack_{X/B} \to \Cohstack_{Y/B}$, on voit que si
$\mathcal{F}_T \to \mathcal{Q}$ appartient à
$\Quotfunctor_{\mathcal{F}/X/B}(T)$, alors
$\mathcal{G}_T \to \pi_{T, *}\mathcal{Q}$ appartient à
$\Quotfunctor_{\mathcal{G}/Y/B}(T)$.
\end{proof}

\begin{lemma}
\label{moduli-lemma-quot-open}
Soit $B$ un espace algébrique. Soit $\pi : X \to Y$ une immersion ouverte
affine d'espaces algébriques séparés et de présentation finie sur $B$.
Soit $\mathcal{F}$ un $\mathcal{O}_X$-module quasi-cohérent. Alors le
morphisme
$\Quotfunctor_{\mathcal{F}/X/B} \to
\Quotfunctor_{\pi_*\mathcal{F}/Y/B}$ du lemme
\ref{moduli-lemma-quot-functorial} est une immersion ouverte.
\end{lemma}

\begin{proof}
Omis. Indication : si $(\pi_*\mathcal{F})_T \to \mathcal{Q}$ est un
élément de $\Quotfunctor_{\pi_*\mathcal{F}/Y/B}(T)$ et si, pour
$t \in T$, on a $\text{Supp}(\mathcal{Q}_t) \subset |X_t|$, alors il en
va de même pour $t' \in T$ dans un voisinage de $t$.
\end{proof}

\begin{lemma}
\label{moduli-lemma-quot-better-open}
Soit $B$ un espace algébrique. Soit $j : X \to Y$ une immersion ouverte
d'espaces algébriques séparés et de présentation finie sur $B$.
Soit $\mathcal{G}$ un $\mathcal{O}_Y$-module quasi-cohérent et posons
$\mathcal{F} = j^*\mathcal{G}$. Il existe alors une immersion ouverte
$$
\Quotfunctor_{\mathcal{F}/X/B}
\longrightarrow
\Quotfunctor_{\mathcal{G}/Y/B}
$$
d'espaces algébriques sur $B$.
\end{lemma}

\begin{proof}
Si $\mathcal{F}_T \to \mathcal{Q}$ est un élément de
$\Quotfunctor_{\mathcal{F}/X/B}(T)$, on peut considérer
$\mathcal{G}_T \to j_{T, *}\mathcal{F}_T \to j_{T, *}\mathcal{Q}$.
L'examen des germes montre que ce morphisme est surjectif.
D'après le lemme \ref{moduli-lemma-coherent-functorial},
$j_{T, *}\mathcal{Q}$ est de présentation finie, plat sur $B$ et à support
propre sur $B$. On obtient ainsi un point à valeurs dans $T$ de
$\Quotfunctor_{\mathcal{G}/Y/B}$. Cela définit le morphisme du lemme.
Nous omettons la démonstration du fait qu'il s'agit d'une immersion ouverte.
Indication : si $\mathcal{G}_T \to \mathcal{Q}$ est un élément de
$\Quotfunctor_{\mathcal{G}/Y/B}(T)$ et si, pour $t \in T$, on a
$\text{Supp}(\mathcal{Q}_t) \subset |X_t|$, alors il en va de même pour
$t' \in T$ dans un voisinage de $t$.
\end{proof}

\begin{lemma}
\label{moduli-lemma-quot-closed}
Soit $B$ un espace algébrique. Soit $\pi : X \to Y$ une immersion fermée
d'espaces algébriques séparés et de présentation finie sur $B$.
Soit $\mathcal{F}$ un $\mathcal{O}_X$-module quasi-cohérent.
Alors le morphisme
$\Quotfunctor_{\mathcal{F}/X/B} \to
\Quotfunctor_{\pi_*\mathcal{F}/Y/B}$ du lemme
\ref{moduli-lemma-quot-functorial} est un isomorphisme.
\end{lemma}

\begin{proof}
Pour tout schéma $T$ sur $B$, le morphisme $\pi_T : X_T \to Y_T$ est une
immersion fermée. Le foncteur $\pi_{T, *}$ est donc une équivalence de
catégories entre $\QCoh(\mathcal{O}_{X_T})$ et la sous-catégorie pleine de
$\QCoh(\mathcal{O}_{Y_T})$ formée des modules quasi-cohérents annulés par
le faisceau d'idéaux de $X_T$ ; voir Morphismes d'espaces, lemme
\ref{spaces-morphisms-lemma-i-star-equivalence}.
Puisqu'un quotient de $(\pi_*\mathcal{F})_T$ est annulé par cet idéal, on
obtient la bijectivité de l'application
$\Quotfunctor_{\mathcal{F}/X/B}(T) \to
\Quotfunctor_{\pi_*\mathcal{F}/Y/B}(T)$
pour tout $T$, comme voulu.
\end{proof}

\begin{lemma}
\label{moduli-lemma-quot-quotient}
Soit $X \to B$ comme dans l'introduction de cette section. Soit
$\mathcal{F} \to \mathcal{G}$ une surjection de $\mathcal{O}_X$-modules
quasi-cohérents. Il existe alors une immersion fermée canonique
$\Quotfunctor_{\mathcal{G}/X/B} \to
\Quotfunctor_{\mathcal{F}/X/B}$.
\end{lemma}

\begin{proof}
Soit $\mathcal{K} = \Ker(\mathcal{F} \to \mathcal{G})$.
Par exactitude à droite des images inverses, on obtient une suite exacte
$\mathcal{K}_T \to \mathcal{F}_T \to \mathcal{G}_T \to 0$
pour tout schéma $T$ sur $B$.
En particulier, un quotient de $\mathcal{G}_T$ détermine un quotient de
$\mathcal{F}_T$, et l'on obtient notre transformation de foncteurs
$\Quotfunctor_{\mathcal{G}/X/B} \to
\Quotfunctor_{\mathcal{F}/X/B}$.
Cette transformation est une immersion fermée d'après Platitude sur les
espaces, lemme \ref{spaces-flat-lemma-F-zero-closed-proper}.
En effet, étant donné un élément $\mathcal{F}_T \to \mathcal{Q}$ de
$\Quotfunctor_{\mathcal{F}/X/B}(T)$, son image inverse à $T'/T$ appartient
à l'image de la transformation si et seulement si
$\mathcal{K}_{T'} \to \mathcal{Q}_{T'}$ est nul.
\end{proof}

\begin{remark}[Invariants numériques]
\label{moduli-remark-quot-numerical}
Soient $f : X \to B$ et $\mathcal{F}$ comme dans l'introduction de cette
section. Soit $I$ un ensemble et, pour $i \in I$, soit
$E_i \in D(\mathcal{O}_X)$ parfait. Soit $P : I \to \mathbf{Z}$ une
fonction. Rappelons que nous avons un morphisme
$$
\Quotfunctor_{\mathcal{F}/X/B} \longrightarrow \Cohstack_{X/B}
$$
qui envoie l'élément $\mathcal{F}_T \to \mathcal{Q}$ de
$\Quotfunctor_{\mathcal{F}/X/B}(T)$ sur l'objet $\mathcal{Q}$ de
$\Cohstack_{X/B}$ au-dessus de $T$ ; voir la démonstration de Quotients,
proposition \ref{quot-proposition-quot}. On peut donc former le diagramme
de produit fibré
$$
\xymatrix{
\Quotfunctor^P_{\mathcal{F}/X/B} \ar[r] \ar[d] &
\Cohstack^P_{X/B} \ar[d] \\
\Quotfunctor_{\mathcal{F}/X/B} \ar[r] &
\Cohstack_{X/B}
}
$$
Il s'agit du diagramme qui définit l'espace algébrique situé en haut à
gauche. La flèche verticale de gauche est une immersion fermée plate qui
est une immersion ouverte et fermée, par exemple si $I$ est fini, si $B$
est localement noethérien, ou si $I = \mathbf{Z}$ et
$E_i = \mathcal{L}^{\otimes i}$ pour un $\mathcal{O}_X$-module inversible
$\mathcal{L}$ (dans ce dernier cas, nous employons parfois la notation
$\Quotfunctor^{P, \mathcal{L}}_{\mathcal{F}/X/B}$).
Voir la situation \ref{moduli-situation-numerical}, les lemmes
\ref{moduli-lemma-open-P} et \ref{moduli-lemma-finite-list-perfect-objects}, ainsi que
l'exemple \ref{moduli-example-hilbert-polynomial}.
\end{remark}

\begin{lemma}
\label{moduli-lemma-quot-tensor-invertible}
Soient $f : X \to B$ et $\mathcal{F}$ comme dans l'introduction de cette
section. Soit $\mathcal{L}$ un $\mathcal{O}_X$-module inversible.
Le produit tensoriel par $\mathcal{L}$ définit alors un isomorphisme
$$
\Quotfunctor_{\mathcal{F}/X/B}
\longrightarrow
\Quotfunctor_{\mathcal{F} \otimes_{\mathcal{O}_X} \mathcal{L}/X/B}
$$
Étant donné un polynôme numérique $P(t)$, posons $P'(t) = P(t + 1)$ ;
cette application induit alors un isomorphisme
$\Quotfunctor^P_{\mathcal{F}/X/B}
\longrightarrow
\Quotfunctor^{P'}_{\mathcal{F} \otimes_{\mathcal{O}_X} \mathcal{L}/X/B}$
de sous-champs ouverts et fermés.
\end{lemma}

\begin{proof}
Posons $\mathcal{G} = \mathcal{F} \otimes_{\mathcal{O}_X} \mathcal{L}$.
Remarquons que
$\mathcal{G}_T = \mathcal{F}_T \otimes_{\mathcal{O}_{X_T}}
\mathcal{L}_T$.
Si $\mathcal{F}_T \to \mathcal{Q}$ est un élément de
$\Quotfunctor_{\mathcal{F}/X/B}(T)$, on l'envoie sur l'élément
$\mathcal{G}_T \to
\mathcal{Q} \otimes_{\mathcal{O}_{X_T}} \mathcal{L}_T$
de $\Quotfunctor_{\mathcal{F} \otimes_{\mathcal{O}_X}
\mathcal{L}/X/B}(T)$.
Cette construction est compatible aux images inverses et définit donc la
transformation de foncteurs voulue. Comme il existe une transformation
inverse évidente, c'est un isomorphisme. Nous omettons la démonstration de
la dernière assertion.
\end{proof}

\begin{lemma}
\label{moduli-lemma-quot-power-invertible}
Soient $f : X \to B$ et $\mathcal{F}$ comme dans l'introduction de cette
section. Soit $\mathcal{L}$ un $\mathcal{O}_X$-module inversible.
Alors
$$
\Quotfunctor^{P, \mathcal{L}}_{\mathcal{F}/X/B} =
\Quotfunctor^{P', \mathcal{L}^{\otimes n}}_{\mathcal{F}/X/B}
$$
où $P'(t) = P(nt)$.
\end{lemma}

\begin{proof}
Cela résulte immédiatement du développement des définitions.
\end{proof}

\section{Bornitude pour Quot}
\label{moduli-section-quot-bounded}

\noindent
Contrairement à ce qui se passe dans la situation classique, nous savons
déjà que le foncteur Quot est un espace algébrique, mais nous ignorons s'il
est jamais représenté par un espace algébrique de type fini.

\begin{lemma}
\label{moduli-lemma-quot-Pn}
Soient $n \geq 0$, $r \geq 1$ et $P \in \mathbf{Q}[t]$.
L'espace algébrique
$$
X = \Quotfunctor^P_{\mathcal{O}^{\oplus r}_{\mathbf{P}^n_\mathbf{Z}}/
\mathbf{P}^n_\mathbf{Z}/\mathbf{Z}}
$$
qui paramètre les quotients de
$\mathcal{O}_{\mathbf{P}^n_\mathbf{Z}}^{\oplus r}$ de polynôme de Hilbert
$P$ est propre sur $\Spec(\mathbf{Z})$.
\end{lemma}

\begin{proof}
Nous savons déjà que $X \to \Spec(\mathbf{Z})$ est séparé et localement de
présentation finie (lemme \ref{moduli-lemma-quot-s-lfp}).
Nous savons aussi que $X \to \Spec(\mathbf{Z})$ satisfait la partie
d'existence du critère valuatif ; voir le lemme
\ref{moduli-lemma-quot-existence-part}.
D'après le critère valuatif de propreté, il suffit de montrer que notre
espace Quot est quasi-compact ; voir Morphismes d'espaces, lemme
\ref{spaces-morphisms-lemma-characterize-proper}.
Il suffit donc de trouver un schéma quasi-compact $T$ et un morphisme
surjectif $T \to X$. Soit $m$ l'entier fourni par Variétés, lemme
\ref{varieties-lemma-bound-quotients-free}. Posons
$$
N = r{m + n \choose n} - P(m)
$$
Nous écrirons $\mathbf{P}^n$ pour
$\mathbf{P}^n_\mathbf{Z} = \text{Proj}(\mathbf{Z}[T_0, \ldots, T_n])$,
et les produits dépourvus d'indice seront pris sur $\Spec(\mathbf{Z})$.
L'idée de la démonstration consiste à construire une application
« universelle »
$$
\Psi :
\mathcal{O}_{T \times \mathbf{P}^n}(-m)^{\oplus N}
\longrightarrow
\mathcal{O}_{T \times \mathbf{P}^n}^{\oplus r}
$$
sur un schéma affine $T$, puis à montrer que tout point de $X$ correspond
au conoyau de cette application en un certain point de $T$.

\medskip\noindent
Définition de $T$ et de $\Psi$. Prenons $T = \Spec(A)$, où
$$
A = \mathbf{Z}[a_{i, j, E}]
$$
avec $i \in \{1, \ldots, r\}$, $j \in \{1, \ldots, N\}$ et où
$E = (e_0, \ldots, e_n)$ parcourt les multi-indices de degré total
$|E| = \sum_{k = 0, \ldots n} e_k = m$.
Nous définissons alors $\Psi$ comme l'application dont l'entrée de matrice
$(i, j)$ est l'application
$$
\sum\nolimits_{E = (e_0, \ldots, e_n)}
a_{i, j, E} T_0^{e_0} \ldots T_n^{e_n} :
\mathcal{O}_{T \times \mathbf{P}^n}(-m)
\longrightarrow
\mathcal{O}_{T \times \mathbf{P}^n}
$$
où la somme porte sur les $E$ ci-dessus ($i$ et $j$ étant bien entendu
fixés).

\medskip\noindent
Considérons le quotient $\mathcal{Q} = \Coker(\Psi)$ sur
$T \times \mathbf{P}^n$. D'après Compléments sur les morphismes, lemme
\ref{more-morphisms-lemma-generic-flatness-stratification}, il existe un
$t \geq 0$ et des sous-schémas fermés
$$
T = T_0 \supset T_1 \supset \ldots \supset T_t = \emptyset
$$
tels que l'image inverse $\mathcal{Q}_p$ de $\mathcal{Q}$ sur
$(T_p \setminus T_{p + 1}) \times \mathbf{P}^n$ soit plate sur
$T_p \setminus T_{p + 1}$. Remarquons que l'on a une suite exacte
$$
\mathcal{O}_{(T_p \setminus T_{p + 1}) \times \mathbf{P}^n}(-m)^{\oplus N}
\to
\mathcal{O}_{(T_p \setminus T_{p + 1}) \times \mathbf{P}^n}^{\oplus r}
\to
\mathcal{Q}_p
\to
0
$$
en prenant l'image inverse de la suite exacte qui définit
$\mathcal{Q} = \Coker(\Psi)$. Nous obtenons donc un morphisme
$$
\coprod (T_p \setminus T_{p + 1})
\longrightarrow
\Quotfunctor_{\mathcal{O}^{\oplus r}/\mathbf{P}/\mathbf{Z}}
\supset
\Quotfunctor^P_{\mathcal{O}^{\oplus r}/\mathbf{P}/\mathbf{Z}} = X
$$
Puisque le membre de gauche est un schéma noethérien et que l'inclusion du
membre de droite est ouverte, il suffit de montrer que tout point de $X$
appartient à l'image de ce morphisme.

\medskip\noindent
Soit $k$ un corps et soit $x \in X(k)$. Alors $x$ correspond à une
surjection $\mathcal{O}_{\mathbf{P}^n_k}^{\oplus r} \to \mathcal{F}$ de
$\mathcal{O}_{\mathbf{P}^n_k}$-modules cohérents telle que le polynôme de
Hilbert de $\mathcal{F}$ soit $P$. Considérons la suite exacte courte
$$
0 \to \mathcal{K} \to
\mathcal{O}_{\mathbf{P}^n_k}^{\oplus r} \to
\mathcal{F} \to 0
$$
D'après Variétés, lemme \ref{varieties-lemma-bound-quotients-free}, et le
choix de $m$, le faisceau $\mathcal{K}$ est $m$-régulier.
D'après Variétés, lemme
\ref{varieties-lemma-m-regular-globally-generated},
$\mathcal{K}(m)$ est engendré par ses sections globales.
D'après Variétés, lemme \ref{varieties-lemma-m-regular-up}, et la
définition de la $m$-régularité, on a
$H^i(\mathbf{P}^n_k, \mathcal{K}(m)) = 0$ pour $i > 0$.
Il vient donc
$$
\dim_k H^0(\mathbf{P}^n_k, \mathcal{K}(m)) =
\chi(\mathcal{K}(m)) =
\chi(\mathcal{O}_{\mathbf{P}^n_k}(m)^{\oplus r}) -
\chi(\mathcal{F}(m)) = N
$$
par le choix de $N$. Cela donne une surjection
$$
\mathcal{O}_{\mathbf{P}^n_k}^{\oplus N}
\longrightarrow
\mathcal{K}(m)
$$
Après torsion en sens inverse et en utilisant la suite exacte courte
ci-dessus, on voit que $\mathcal{F}$ est le conoyau d'une application
$$
\Psi_x :
\mathcal{O}_{\mathbf{P}^n_k}(-m)^{\oplus N}
\to
\mathcal{O}_{\mathbf{P}^n_k}^{\oplus r}
$$
Il existe un unique homomorphisme d'anneaux $\tau : A \to k$ tel que le
changement de base de $\Psi$ par le morphisme correspondant
$t = \Spec(\tau) : \Spec(k) \to T$ soit $\Psi_x$.
En effet, les entrées de la matrice $N \times r$ qui définit $\Psi_x$
sont des polynômes homogènes
$\sum \lambda_{i, j, E} T_0^{e_0} \ldots T_n^{e_n}$
de degré $m$ en $T_0, \ldots, T_n$, à coefficients
$\lambda_{i, j, E} \in k$, et l'on peut poser
$\tau(a_{i, j, E}) = \lambda_{i, j, E}$.
Alors $t \in T_p \setminus T_{p + 1}$ pour un certain $p$, et l'image de
$t$ par le morphisme ci-dessus est $x$, comme voulu.
\end{proof}

\begin{lemma}
\label{moduli-lemma-quot-Pn-over-base}
Soit $B$ un espace algébrique. Soit
$X = B \times \mathbf{P}^n_\mathbf{Z}$. Soit $\mathcal{L}$ l'image
inverse de $\mathcal{O}_{\mathbf{P}^n}(1)$ sur $X$.
Soit $\mathcal{F}$ un $\mathcal{O}_X$-module de présentation finie.
L'espace algébrique $\Quotfunctor^P_{\mathcal{F}/X/B}$ qui paramètre les
quotients de $\mathcal{F}$ de polynôme de Hilbert $P$ relativement à
$\mathcal{L}$ est propre sur $B$.
\end{lemma}

\begin{proof}
La question est locale pour la topologie \'etale sur $B$ ; voir Morphismes
d'espaces, lemme \ref{spaces-morphisms-lemma-proper-local}.
Nous pouvons donc supposer que $B$ est un schéma affine.
Dans ce cas, $\mathcal{L}$ est un module inversible ample sur $X$
(d'après Constructions, lemme \ref{constructions-lemma-ample-on-proj}, et la
définition des modules inversibles amples dans Propriétés, définition
\ref{properties-definition-ample}).
On peut donc trouver $r' \geq 0$ et $r \geq 0$, ainsi qu'une surjection
$$
\mathcal{O}_X^{\oplus r} \longrightarrow
\mathcal{F} \otimes_{\mathcal{O}_X} \mathcal{L}^{\otimes r'}
$$
d'après Propriétés, proposition
\ref{properties-proposition-characterize-ample}.
En vertu du lemme \ref{moduli-lemma-quot-tensor-invertible}, on peut remplacer
$\mathcal{F}$ par
$\mathcal{F} \otimes_{\mathcal{O}_X} \mathcal{L}^{\otimes r'}$
et $P(t)$ par $P(t + r')$.
Le lemme \ref{moduli-lemma-quot-quotient} fournit une immersion fermée
$$
\Quotfunctor^P_{\mathcal{F}/X/B}
\longrightarrow
\Quotfunctor^P_{\mathcal{O}_X^{\oplus r}/X/B}
$$
Comme nous avons montré que
$\Quotfunctor^P_{\mathcal{O}_X^{\oplus r}/X/B} \to B$ est propre au lemme
\ref{moduli-lemma-quot-Pn}, on conclut.
\end{proof}

\begin{lemma}
\label{moduli-lemma-quot-proper-over-base}
Soit $f : X \to B$ un morphisme propre et de présentation finie d'espaces
algébriques. Soit $\mathcal{F}$ un $\mathcal{O}_X$-module de présentation
finie. Soit $\mathcal{L}$ un $\mathcal{O}_X$-module inversible ample sur
$X/B$ ; voir Diviseurs sur les espaces, définition
\ref{spaces-divisors-definition-relatively-ample}.
L'espace algébrique $\Quotfunctor^P_{\mathcal{F}/X/B}$ qui paramètre les
quotients de $\mathcal{F}$ de polynôme de Hilbert $P$ relativement à
$\mathcal{L}$ est propre sur $B$.
\end{lemma}

\begin{proof}
La question est locale pour la topologie \'etale sur $B$ ; voir Morphismes
d'espaces, lemme \ref{spaces-morphisms-lemma-proper-local}.
Nous pouvons donc supposer que $B$ est un schéma affine.
Il existe alors une immersion fermée $i : X \to \mathbf{P}^n_B$ telle que
$i^*\mathcal{O}_{\mathbf{P}^n_B}(1) \cong \mathcal{L}^{\otimes d}$
pour un certain $d \geq 1$. Voir Morphismes, lemme
\ref{morphisms-lemma-quasi-projective-finite-type-over-S}.
Remplacer $\mathcal{L}$ par $\mathcal{L}^{\otimes d}$ et le polynôme
numérique $P(t)$ par $P(dt)$ ne modifie pas
$\Quotfunctor^P_{\mathcal{F}/X/B}$ ; certains détails sont omis.
Nous pouvons donc supposer
$\mathcal{L} = i^*\mathcal{O}_{\mathbf{P}^n_B}(1)$.
L'isomorphisme
$\Quotfunctor_{\mathcal{F}/X/B} \to
\Quotfunctor_{i_*\mathcal{F}/\mathbf{P}^n_B/B}$ du lemme
\ref{moduli-lemma-quot-closed} induit alors un isomorphisme
$\Quotfunctor^P_{\mathcal{F}/X/B} \cong
\Quotfunctor^P_{i_*\mathcal{F}/\mathbf{P}^n_B/B}$.
Comme $\Quotfunctor^P_{i_*\mathcal{F}/\mathbf{P}^n_B/B}$ est propre sur
$B$ d'après le lemme \ref{moduli-lemma-quot-Pn-over-base}, on conclut.
\end{proof}

\begin{lemma}
\label{moduli-lemma-quot-qc-over-base}
Soit $f : X \to B$ un morphisme séparé et de présentation finie d'espaces
algébriques. Soit $\mathcal{F}$ un $\mathcal{O}_X$-module de présentation
finie. Soit $\mathcal{L}$ un $\mathcal{O}_X$-module inversible ample sur
$X/B$ ; voir Diviseurs sur les espaces, définition
\ref{spaces-divisors-definition-relatively-ample}.
L'espace algébrique $\Quotfunctor^P_{\mathcal{F}/X/B}$ qui paramètre les
quotients de $\mathcal{F}$ de polynôme de Hilbert $P$ relativement à
$\mathcal{L}$ est séparé et de présentation finie sur $B$.
\end{lemma}

\begin{proof}
Nous avons déjà vu que $\Quotfunctor_{\mathcal{F}/X/B} \to B$ est séparé
et localement de présentation finie ; voir le lemme
\ref{moduli-lemma-quot-s-lfp}. Il suffit donc de montrer que le sous-espace ouvert
$\Quotfunctor^P_{\mathcal{F}/X/B}$ de la remarque
\ref{moduli-remark-quot-numerical} est quasi-compact sur $B$.

\medskip\noindent
La question est locale pour la topologie \'etale sur $B$
(Morphismes d'espaces, lemme
\ref{spaces-morphisms-lemma-quasi-compact-local}).
Nous pouvons donc supposer $B$ affine.

\medskip\noindent
Supposons $B = \Spec(\Lambda)$. Écrivons
$\Lambda = \colim \Lambda_i$ comme limite inductive de ses sous-algèbres
de type fini sur $\mathbf{Z}$. On peut alors trouver un $i$ et un système
$X_i, \mathcal{F}_i, \mathcal{L}_i$ comme dans le lemme sur
$B_i = \Spec(\Lambda_i)$, dont le changement de base à $B$ donne
$X, \mathcal{F}, \mathcal{L}$. Cela résulte de Limites d'espaces, lemmes
\ref{spaces-limits-lemma-descend-finite-presentation} (pour trouver $X_i$),
\ref{spaces-limits-lemma-descend-modules-finite-presentation} (pour trouver
$\mathcal{F}_i$), \ref{spaces-limits-lemma-descend-invertible-modules}
(pour trouver $\mathcal{L}_i$) et
\ref{spaces-limits-lemma-descend-separated} (pour rendre $X_i$ séparé).
Comme
$$
\Quotfunctor_{\mathcal{F}/X/B} = B \times_{B_i}
\Quotfunctor_{\mathcal{F}_i/X_i/B_i}
$$
et qu'il en va de même de $\Quotfunctor^P_{\mathcal{F}/X/B}$, on se
ramène au cas étudié dans le paragraphe suivant.

\medskip\noindent
Supposons $B$ affine et noethérien. Nous pouvons remplacer $\mathcal{L}$
par une puissance positive ; voir le lemme
\ref{moduli-lemma-quot-power-invertible}.
Nous pouvons ainsi supposer qu'il existe une immersion
$i : X \to \mathbf{P}^n_B$ telle que
$i^*\mathcal{O}_{\mathbf{P}^n}(1) = \mathcal{L}$.
D'après Morphismes, lemme
\ref{morphisms-lemma-quasi-compact-immersion}, il existe un sous-schéma
fermé $X' \subset \mathbf{P}^n_B$ tel que $i$ se factorise par une
immersion ouverte $j : X \to X'$.
D'après Propriétés, lemme
\ref{properties-lemma-lift-finite-presentation}, il existe un
$\mathcal{O}_{X'}$-module $\mathcal{G}$ de présentation finie tel que
$j^*\mathcal{G} = \mathcal{F}$. On obtient donc une immersion ouverte
$$
\Quotfunctor_{\mathcal{F}/X/B}
\longrightarrow
\Quotfunctor_{\mathcal{G}/X'/B}
$$
par le lemme \ref{moduli-lemma-quot-better-open}. Manifestement, cette immersion
ouverte envoie $\Quotfunctor^P_{\mathcal{F}/X/B}$ dans
$\Quotfunctor^P_{\mathcal{G}/X'/B}$. Or
$\Quotfunctor^P_{\mathcal{G}/X'/B}$ est propre sur $B$ d'après le lemme
\ref{moduli-lemma-quot-proper-over-base}.
Il est donc noethérien et, puisque tout ouvert d'un espace algébrique
noethérien est quasi-compact, la démonstration est achevée.
\end{proof}




\section{Propriétés du foncteur de Hilbert}
\label{moduli-section-hilb}

\noindent
Soit $f : X \to B$ un morphisme d'espaces algébriques séparé et de
présentation finie. Alors $\Hilbfunctor_{X/B}$ est un espace algébrique
localement de présentation finie sur $B$. Voir Quotients, proposition
\ref{quot-proposition-hilb}.

\begin{lemma}
\label{moduli-lemma-hilb-diagonal-closed}
La diagonale de $\Hilbfunctor_{X/B} \to B$ est une immersion fermée de
présentation finie.
\end{lemma}

\begin{proof}
Dans Quotients, lemme \ref{quot-lemma-hilb-is-quot}, nous avons vu que
$\Hilbfunctor_{X/B} = \Quotfunctor_{\mathcal{O}_X/X/B}$.
L'assertion résulte donc du lemme \ref{moduli-lemma-quot-diagonal-closed}.
\end{proof}

\begin{lemma}
\label{moduli-lemma-hilb-s-lfp}
Le morphisme $\Hilbfunctor_{X/B} \to B$ est séparé et localement de
présentation finie.
\end{lemma}

\begin{proof}
Pour vérifier que $\Hilbfunctor_{X/B} \to B$ est séparé, il faut montrer
que sa diagonale est une immersion fermée. Cela résulte du lemme
\ref{moduli-lemma-hilb-diagonal-closed}. La seconde assertion fait partie de
Quotients, proposition \ref{quot-proposition-hilb}.
\end{proof}

\begin{lemma}
\label{moduli-lemma-hilb-existence-part}
Supposons que $X \to B$ soit propre ainsi que de présentation finie.
Alors $\Hilbfunctor_{X/B} \to B$ satisfait la partie d'existence du critère
valuatif (Morphismes d'espaces, définition
\ref{spaces-morphisms-definition-valuative-criterion}).
\end{lemma}

\begin{proof}
Dans Quotients, lemme \ref{quot-lemma-hilb-is-quot}, nous avons vu que
$\Hilbfunctor_{X/B} = \Quotfunctor_{\mathcal{O}_X/X/B}$.
L'assertion résulte donc du lemme \ref{moduli-lemma-quot-existence-part}.
\end{proof}

\begin{lemma}
\label{moduli-lemma-hilb-open}
Soit $B$ un espace algébrique. Soit $\pi : X \to Y$ une immersion ouverte
d'espaces algébriques séparés et de présentation finie sur $B$.
Alors $\pi$ induit une immersion ouverte
$\Hilbfunctor_{X/B} \to \Hilbfunctor_{Y/B}$.
\end{lemma}

\begin{proof}
Omis. Indication : si $Z \subset X_T$ est un sous-schéma fermé propre sur
$T$, alors $Z$ est aussi fermé dans $Y_T$. On obtient ainsi la
transformation $\Hilbfunctor_{X/B} \to \Hilbfunctor_{Y/B}$.
Si $Z \subset Y_T$ est un élément de $\Hilbfunctor_{Y/B}(T)$ et si, pour
$t \in T$, on a $|Z_t| \subset |X_t|$, alors il en va de même pour
$t' \in T$ dans un voisinage de $t$.
\end{proof}

\begin{lemma}
\label{moduli-lemma-hilb-closed}
Soit $B$ un espace algébrique. Soit $\pi : X \to Y$ une immersion fermée
d'espaces algébriques séparés et de présentation finie sur $B$.
Alors $\pi$ induit une immersion fermée
$\Hilbfunctor_{X/B} \to \Hilbfunctor_{Y/B}$.
\end{lemma}

\begin{proof}
Puisque $\pi$ est une immersion fermée, il est immédiat qu'étant donné un
sous-schéma fermé $Z \subset X_T$, on peut voir $Z$ comme un sous-schéma
fermé de $X_T$. On obtient ainsi la transformation
$\Hilbfunctor_{X/B} \to \Hilbfunctor_{Y/B}$.
On voit immédiatement que cette transformation est un monomorphisme.
Pour prouver qu'elle est une immersion fermée, on peut appliquer le lemme
\ref{moduli-lemma-quot-quotient} à l'application
$\mathcal{O}_Y \to \mathcal{O}_X$ et aux identifications
$\Hilbfunctor_{X/B} = \Quotfunctor_{\mathcal{O}_X/X/B}$,
$\Hilbfunctor_{Y/B} = \Quotfunctor_{\mathcal{O}_Y/Y/B}$
de Quotients, lemme \ref{quot-lemma-hilb-is-quot}.
\end{proof}

\begin{remark}[Invariants numériques]
\label{moduli-remark-hilb-numerical}
Soit $f : X \to B$ comme dans l'introduction de cette section.
Soit $I$ un ensemble et, pour $i \in I$, soit
$E_i \in D(\mathcal{O}_X)$ parfait. Soit
$P : I \to \mathbf{Z}$ une fonction. Rappelons que
$\Hilbfunctor_{X/B} = \Quotfunctor_{\mathcal{O}_X/X/B}$ ; voir Quotients,
lemme \ref{quot-lemma-hilb-is-quot}. On peut donc définir
$$
\Hilbfunctor^P_{X/B} = \Quotfunctor^P_{\mathcal{O}_X/X/B}
$$
où $\Quotfunctor^P_{\mathcal{O}_X/X/B}$ est comme dans la remarque
\ref{moduli-remark-quot-numerical}. Le morphisme
$$
\Hilbfunctor^P_{X/B} \longrightarrow \Hilbfunctor_{X/B}
$$
est une immersion fermée plate qui est une immersion ouverte et fermée,
par exemple si $I$ est fini, si $B$ est localement noethérien, ou si
$I = \mathbf{Z}$ et $E_i = \mathcal{L}^{\otimes i}$ pour un
$\mathcal{O}_X$-module inversible $\mathcal{L}$. Dans ce dernier cas, nous
employons parfois la notation $\Hilbfunctor^{P, \mathcal{L}}_{X/B}$.
\end{remark}

\begin{lemma}
\label{moduli-lemma-hilb-proper-over-base}
Soit $f : X \to B$ un morphisme propre et de présentation finie d'espaces
algébriques. Soit $\mathcal{L}$ un $\mathcal{O}_X$-module inversible ample
sur $X/B$ ; voir Diviseurs sur les espaces, définition
\ref{spaces-divisors-definition-relatively-ample}.
L'espace algébrique $\Hilbfunctor^P_{X/B}$ qui paramètre les sous-schémas
fermés de polynôme de Hilbert $P$ relativement à $\mathcal{L}$ est propre
sur $B$.
\end{lemma}

\begin{proof}
Rappelons que
$\Hilbfunctor_{X/B} = \Quotfunctor_{\mathcal{O}_X/X/B}$ ; voir Quotients,
lemme \ref{quot-lemma-hilb-is-quot}.
Le présent lemme est donc une conséquence immédiate du lemme
\ref{moduli-lemma-quot-proper-over-base}.
\end{proof}

\begin{lemma}
\label{moduli-lemma-hilb-qc-over-base}
Soit $f : X \to B$ un morphisme séparé et de présentation finie d'espaces
algébriques. Soit $\mathcal{L}$ un $\mathcal{O}_X$-module inversible ample
sur $X/B$ ; voir Diviseurs sur les espaces, définition
\ref{spaces-divisors-definition-relatively-ample}.
L'espace algébrique $\Hilbfunctor^P_{X/B}$ qui paramètre les sous-schémas
fermés de polynôme de Hilbert $P$ relativement à $\mathcal{L}$ est séparé
et de présentation finie sur $B$.
\end{lemma}

\begin{proof}
Rappelons que
$\Hilbfunctor_{X/B} = \Quotfunctor_{\mathcal{O}_X/X/B}$ ; voir Quotients,
lemme \ref{quot-lemma-hilb-is-quot}.
Le présent lemme est donc une conséquence immédiate du lemme
\ref{moduli-lemma-quot-qc-over-base}.
\end{proof}








\section{Propriétés du champ de Picard}
\label{moduli-section-picard-stack}

\noindent
Soit $f : X \to B$ un morphisme d'espaces algébriques plat, propre et de
présentation finie. Alors le champ $\Picardstack_{X/B}$ qui paramètre les
faisceaux inversibles sur $X/B$ est algébrique ; voir Quotients,
proposition \ref{quot-proposition-pic}.

\begin{lemma}
\label{moduli-lemma-pic-diagonal-affine-fp}
La diagonale de $\Picardstack_{X/B}$ sur $B$ est affine et de présentation
finie.
\end{lemma}

\begin{proof}
Dans Quotients, lemme \ref{quot-lemma-picard-stack-open-in-coh}, nous avons
vu que $\Picardstack_{X/B}$ est un sous-champ ouvert de
$\Cohstack_{X/B}$. L'assertion résulte donc du lemme
\ref{moduli-lemma-coherent-diagonal-affine-fp}.
\end{proof}

\begin{lemma}
\label{moduli-lemma-pic-qs-lfp}
Le morphisme $\Picardstack_{X/B} \to B$ est quasi-séparé et localement de
présentation finie.
\end{lemma}

\begin{proof}
Dans Quotients, lemme \ref{quot-lemma-picard-stack-open-in-coh}, nous avons
vu que $\Picardstack_{X/B}$ est un sous-champ ouvert de
$\Cohstack_{X/B}$. L'assertion résulte donc du lemme
\ref{moduli-lemma-coherent-qs-lfp}.
\end{proof}

\begin{lemma}
\label{moduli-lemma-pic-existence-part}
Supposons en outre que $X \to B$ soit lisse, en plus d'être propre.
Alors $\Picardstack_{X/B} \to B$ satisfait la partie d'existence du critère
valuatif (Morphismes de champs, définition
\ref{stacks-morphisms-definition-existence}).
\end{lemma}

\begin{proof}
Après changement de base, on se ramène immédiatement au problème suivant :
étant donnés un anneau de valuation $R$ de corps des fractions $K$, un
espace algébrique $X$ propre et lisse sur $R$ et un
$\mathcal{O}_{X_K}$-module inversible $\mathcal{L}_K$, montrer qu'il existe
un $\mathcal{O}_X$-module inversible $\mathcal{L}$ dont la fibre générique
est $\mathcal{L}_K$.
Remarquons que $X_K$ est noethérien, séparé et régulier (utiliser Morphismes
d'espaces, lemme
\ref{spaces-morphisms-lemma-finite-presentation-noetherian}, et Espaces sur
les corps, lemme \ref{spaces-over-fields-lemma-smooth-regular}).
On peut donc écrire $\mathcal{L}_K$ comme la différence, dans le groupe de
Picard, de $\mathcal{O}_{X_K}(D_K)$ et
$\mathcal{O}_{X_K}(D'_K)$, pour deux diviseurs de Cartier effectifs
$D_K, D'_K$ dans $X_K$ ; voir Diviseurs sur les espaces, lemme
\ref{spaces-divisors-lemma-Noetherian-regular-separated-pic-effective-Cartier}.
Enfin, on sait que $D_K$ et $D'_K$ sont les restrictions de diviseurs de
Cartier effectifs $D, D' \subset X$ ; voir Diviseurs sur les espaces,
lemme
\ref{spaces-divisors-lemma-smooth-over-valuation-ring-effective-Cartier}.
\end{proof}

\begin{lemma}
\label{moduli-lemma-pic-inertia}
Supposons $f_{T, *}\mathcal{O}_{X_T} \cong \mathcal{O}_T$ pour tout schéma
$T$ sur $B$. Alors le champ d'inertie de $\Picardstack_{X/B}$ est égal à
$\mathbf{G}_m \times \Picardstack_{X/B}$.
\end{lemma}

\begin{proof}
Ceci est expliqué dans Exemples de champs, exemple
\ref{examples-stacks-example-inertia-stack-of-picard}.
\end{proof}

\begin{lemma}
\label{moduli-lemma-pic-curves-smooth}
Supposons en outre que $f : X \to B$ soit de dimension relative
$\leq 1$, en plus des autres hypothèses de cette section. Alors
$\Picardstack_{X/B} \to B$ est lisse.
\end{lemma}

\begin{proof}
Nous savons déjà que $\Picardstack_{X/B} \to B$ est localement de
présentation finie ; voir le lemme \ref{moduli-lemma-pic-qs-lfp}.
Il suffit donc de montrer que $\Picardstack_{X/B} \to B$ est formellement
lisse ; voir Compléments sur les morphismes de champs, lemme
\ref{stacks-more-morphisms-lemma-smooth-formally-smooth}.
Après changement de base, on se ramène immédiatement au problème suivant :
étant donné un épaississement du premier ordre $T \subset T'$ de schémas
affines, étant donné $X' \to T'$ propre, plat, de présentation finie et de
dimension relative $\leq 1$, et étant donné sur
$X = T \times_{T'} X'$ un $\mathcal{O}_X$-module inversible
$\mathcal{L}$, prouver qu'il existe un $\mathcal{O}_{X'}$-module inversible
$\mathcal{L}'$ dont la restriction à $X$ est $\mathcal{L}$.
Puisque $T \subset T'$ est un épaississement du premier ordre, il en va de
même de $X \subset X'$ ; voir Compléments sur les morphismes d'espaces,
lemme \ref{spaces-more-morphisms-lemma-base-change-thickening}.
D'après Compléments sur les morphismes d'espaces, lemme
\ref{spaces-more-morphisms-lemma-picard-group-first-order-thickening}, il
suffit de montrer que $H^2(X, \mathcal{I}) = 0$, où $\mathcal{I}$ est
l'idéal quasi-cohérent qui définit $X$ dans $X'$.
Notons $f : X \to T$ le morphisme structural.
D'après Cohomologie des espaces, lemme
\ref{spaces-cohomology-lemma-higher-direct-images-zero-above-dimension-fibre},
on a $R^pf_*\mathcal{I} = 0$ pour $p > 1$.
La nullité recherchée résulte donc de Cohomologie des espaces, lemme
\ref{spaces-cohomology-lemma-quasi-coherence-higher-direct-images-application}
(c'est ici que nous utilisons enfin le fait que $T$ est affine).
\end{proof}




\section{Propriétés du foncteur de Picard}
\label{moduli-section-picard-functor}

\noindent
Soit $f : X \to B$ un morphisme d'espaces algébriques plat, propre et de
présentation finie, tel que, de plus, pour tout $T/B$, l'application
canonique
$$
\mathcal{O}_T \longrightarrow f_{T, *}\mathcal{O}_{X_T}
$$
soit un isomorphisme. Alors le foncteur de Picard
$\Picardfunctor_{X/B}$ est un espace algébrique ; voir Quotients,
proposition \ref{quot-proposition-pic-functor}.
Il existe un lien étroit avec le champ de Picard.

\begin{lemma}
\label{moduli-lemma-pic-gerbe-over-pic-functor}
Le morphisme $\Picardstack_{X/B} \to \Picardfunctor_{X/B}$ fait du champ
de Picard une gerbe au-dessus du foncteur de Picard.
\end{lemma}

\begin{proof}
La définition du fait que
$\Picardstack_{X/B} \to \Picardfunctor_{X/B}$ soit une gerbe est donnée
dans Morphismes de champs, définition
\ref{stacks-morphisms-definition-gerbe}, qui renvoie à son tour à Champs,
définition \ref{stacks-definition-gerbe-over-stack-in-groupoids}.
Pour le prouver, nous allons vérifier les conditions (2)(a) et (2)(b) de
Champs, lemme \ref{stacks-lemma-when-gerbe}. Cela résulte immédiatement de
Quotients, lemme \ref{quot-lemma-pic-over-pic} ; voici une explication
détaillée.

\medskip\noindent
Condition (2)(a).
Supposons que $\xi \in \Picardfunctor_{X/B}(U)$ pour un schéma $U$ sur
$B$. Comme $\Picardfunctor_{X/B}$ est le faisceau associé pour la topologie
fppf à la règle $T \mapsto \Pic(X_T)$ sur les schémas au-dessus de $B$
(Quotients, situation \ref{quot-situation-pic}), il existe un recouvrement
fppf $\{U_i \to U\}$ tel que $\xi|_{U_i}$ corresponde à un module
inversible $\mathcal{L}_i$ sur $X_{U_i}$.
Alors $(U_i \to B, \mathcal{L}_i)$ est un objet de
$\Picardstack_{X/B}$ sur $U_i$ qui s'envoie sur $\xi|_{U_i}$.

\medskip\noindent
Condition (2)(b). Supposons que $U$ soit un schéma sur $B$ et que
$\mathcal{L}, \mathcal{N}$ soient des modules inversibles sur $X_U$ qui
s'envoient sur le même élément de $\Picardfunctor_{X/B}(U)$.
Il existe alors un recouvrement fppf $\{U_i \to U\}$ tel que
$\mathcal{L}|_{X_{U_i}}$ soit isomorphe à
$\mathcal{N}|_{X_{U_i}}$. On obtient ainsi les isomorphismes recherchés
entre
$(U \to B, \mathcal{L})|_{U_i} \to
(U \to B, \mathcal{N})|_{U_i}$.
\end{proof}

\begin{lemma}
\label{moduli-lemma-pic-functor-diagonal-qc-immersion}
La diagonale de $\Picardfunctor_{X/B}$ sur $B$ est une immersion
quasi-compacte.
\end{lemma}

\begin{proof}
La diagonale est une immersion d'après Quotients, lemme
\ref{quot-lemma-diagonal-pic}. Pour conclure, montrons qu'elle est
quasi-compacte. La diagonale de $\Picardstack_{X/B}$ est quasi-compacte
d'après le lemme \ref{moduli-lemma-pic-diagonal-affine-fp}, et
$\Picardstack_{X/B}$ est une gerbe au-dessus de
$\Picardfunctor_{X/B}$ d'après le lemme
\ref{moduli-lemma-pic-gerbe-over-pic-functor}.
On conclut par Morphismes de champs, lemme
\ref{stacks-morphisms-lemma-gerbe-diagonal-quasi-compact}.
\end{proof}

\begin{lemma}
\label{moduli-lemma-pic-functor-qs-lfp}
Le morphisme $\Picardfunctor_{X/B} \to B$ est quasi-séparé et localement
de présentation finie.
\end{lemma}

\begin{proof}
Pour vérifier que $\Picardfunctor_{X/B} \to B$ est quasi-séparé, il faut
montrer que sa diagonale est quasi-compacte. Cela résulte immédiatement du
lemme \ref{moduli-lemma-pic-functor-diagonal-qc-immersion}.
Comme le morphisme
$\Picardstack_{X/B} \to \Picardfunctor_{X/B}$ est surjectif, plat et
localement de présentation finie (d'après le lemme
\ref{moduli-lemma-pic-gerbe-over-pic-functor} et Morphismes de champs, lemme
\ref{stacks-morphisms-lemma-gerbe-fppf}), il suffit de prouver que
$\Picardstack_{X/B} \to B$ est localement de présentation finie ; voir
Morphismes de champs, lemme
\ref{stacks-morphisms-lemma-flat-finite-presentation-permanence}.
Cela résulte du lemme \ref{moduli-lemma-pic-qs-lfp}.
\end{proof}

\begin{lemma}
\label{moduli-lemma-pic-functor-uniqueness-part}
Supposons en outre que les fibres géométriques de $X \to B$ soient
intègres, en plus des autres hypothèses de cette section.
Alors $\Picardfunctor_{X/B} \to B$ est séparé.
\end{lemma}

\begin{proof}
Puisque $\Picardfunctor_{X/B} \to B$ est quasi-séparé, il suffit de
vérifier la partie d'unicité du critère valuatif ; voir Morphismes
d'espaces, lemme
\ref{spaces-morphisms-lemma-valuative-criterion-separatedness}.
On se ramène immédiatement au problème suivant : étant donnés
\begin{enumerate}
\item un anneau de valuation $R$ de corps des fractions $K$,
\item un espace algébrique $X$ propre et plat sur $R$, à fibre géométrique
intègre,
\item un élément $a \in \Picardfunctor_{X/R}(R)$ tel que
$a|_{\Spec(K)} = 0$,
\end{enumerate}
il faut prouver que $a = 0$. Appliquons Morphismes de champs, lemme
\ref{stacks-morphisms-lemma-lift-valuation-ring-through-flat-morphism},
au morphisme plat surjectif
$\Picardstack_{X/R} \to \Picardfunctor_{X/R}$
(surjectif et plat d'après le lemme
\ref{moduli-lemma-pic-gerbe-over-pic-functor} et Morphismes de champs, lemme
\ref{stacks-morphisms-lemma-gerbe-fppf}). Après avoir remplacé $R$ par une
extension, on peut supposer que $a$ est donné par un
$\mathcal{O}_X$-module inversible $\mathcal{L}$.
Comme $a|_{\Spec(K)} = 0$, on obtient
$\mathcal{L}_K \cong \mathcal{O}_{X_K}$ d'après Quotients, lemme
\ref{quot-lemma-flat-geometrically-connected-fibres}.

\medskip\noindent
Notons $f : X \to \Spec(R)$ le morphisme structural.
Soient $\eta, 0 \in \Spec(R)$ le point générique et le point fermé.
Considérons les complexes parfaits
$K = Rf_*\mathcal{L}$ et $M = Rf_*(\mathcal{L}^{\otimes -1})$
sur $\Spec(R)$ ; voir Catégories dérivées des espaces, lemme
\ref{spaces-perfect-lemma-flat-proper-perfect-direct-image-general}.
Considérons les fonctions
$\beta_{K, i}, \beta_{M, i} : \Spec(R) \to \mathbf{Z}$
de Catégories dérivées des espaces, lemme
\ref{spaces-perfect-lemma-jump-loci}, associées à $K$ et à $M$.
Comme la formation de $K$ et de $M$ commute au changement de base (voir le
lemme cité ci-dessus), on obtient
$\beta_{K, 0}(\eta) = \beta_{M, 0}(\eta) = 1$ d'après Espaces sur les
corps, lemme
\ref{spaces-over-fields-lemma-proper-geometrically-reduced-global-sections},
et notre hypothèse sur les fibres de $f$.
Comme ces fonctions sont semi-continues supérieurement, on obtient
$\beta_{K, 0}(0) \geq 1$ et $\beta_{M, 0}(0) \geq 1$.
D'après Espaces sur les corps, lemme
\ref{spaces-over-fields-lemma-characterize-trivial-pic-integral}, on
conclut que la restriction de $\mathcal{L}$ à la fibre spéciale $X_0$ est
triviale. Cela donne à son tour
$\beta_{K, 0}(0) = \beta_{M, 0} = 1$, comme ci-dessus.
D'après Compléments d'algèbre, lemme
\ref{more-algebra-lemma-lift-pseudo-coherent-from-residue-field}, on peut
alors représenter $K$ par un complexe de la forme
$$
\ldots \to 0 \to R \to R^{\oplus \beta_{K, 1}(0)} \to
R^{\oplus \beta_{K, 2}(0)} \to \ldots
$$
Or $R \to R^{\oplus \beta_{K, 1}(0)}$ est nul parce que
$\beta_{K, 0}(\eta) = 1$. Autrement dit,
$K = R \oplus \tau_{\geq 1}(K)$ dans $D(R)$, où
$\tau_{\geq 1}(K)$ est d'amplitude de Tor dans $[1, b]$ pour un certain
$b \in \mathbf{Z}$.
Il existe donc une section globale $s \in H^0(X, \mathcal{L})$ dont la
restriction $s_0$ à $X_0$ ne s'annule nulle part (de nouveau parce que la
formation de $K$ commute au changement de base).
Alors $s : \mathcal{O}_X \to \mathcal{L}$ est une application de faisceaux
inversibles dont la restriction à $X_0$ est un isomorphisme ; c'est donc
un isomorphisme, comme voulu.
\end{proof}

\begin{lemma}
\label{moduli-lemma-pic-functor-curves-smooth}
Supposons en outre que $f : X \to B$ soit de dimension relative
$\leq 1$, en plus des autres hypothèses de cette section. Alors
$\Picardfunctor_{X/B} \to B$ est lisse.
\end{lemma}

\begin{proof}
D'après le lemme \ref{moduli-lemma-pic-curves-smooth}, nous savons que
$\Picardstack_{X/B} \to B$ est lisse. Le morphisme
$\Picardstack_{X/B} \to \Picardfunctor_{X/B}$ est surjectif et lisse en
combinant le lemme \ref{moduli-lemma-pic-gerbe-over-pic-functor} et Morphismes de
champs, lemme \ref{stacks-morphisms-lemma-gerbe-smooth}.
Ainsi, si $U$ est un schéma et si
$U \to \Picardstack_{X/B}$ est surjectif et lisse, alors
$U \to \Picardfunctor_{X/B}$ est surjectif et lisse, et
$U \to B$ est surjectif et lisse (car ces propriétés sont stables par
composition). Le morphisme $\Picardfunctor_{X/B} \to B$ est donc lisse,
par exemple d'après Descente sur les espaces, lemme
\ref{spaces-descent-lemma-syntomic-smooth-etale-permanence}.
\end{proof}






\section{Propriétés des morphismes relatifs}
\label{moduli-section-relative-morphisms}

\noindent
Soit $B$ un espace algébrique. Soient $X$ et $Y$ des espaces algébriques
sur $B$ tels que $Y \to B$ soit plat, propre et de présentation finie, et
que $X \to B$ soit séparé et de présentation finie.
Alors le foncteur $\mathit{Mor}_B(Y, X)$ des morphismes relatifs est un
espace algébrique localement de présentation finie sur $B$.
Voir Quotients, proposition \ref{quot-proposition-Mor}.

\begin{lemma}
\label{moduli-lemma-Mor-diagonal-closed}
La diagonale de $\mathit{Mor}_B(Y, X) \to B$ est une immersion fermée de
présentation finie.
\end{lemma}

\begin{proof}
Il existe une immersion ouverte
$\mathit{Mor}_B(Y, X) \to \Hilbfunctor_{Y \times_B X/B}$ ; voir
Quotients, lemme \ref{quot-lemma-Mor-into-Hilb-open}.
Le lemme résulte donc du lemme \ref{moduli-lemma-hilb-diagonal-closed}.
\end{proof}

\begin{lemma}
\label{moduli-lemma-Mor-s-lfp}
Le morphisme $\mathit{Mor}_B(Y, X) \to B$ est séparé et localement de
présentation finie.
\end{lemma}

\begin{proof}
Pour vérifier que $\mathit{Mor}_B(Y, X) \to B$ est séparé, il faut montrer
que sa diagonale est une immersion fermée. Cela résulte du lemme
\ref{moduli-lemma-Mor-diagonal-closed}. La seconde assertion fait partie de
Quotients, proposition \ref{quot-proposition-Mor}.
\end{proof}

\begin{lemma}
\label{moduli-lemma-Isom-in-Mor}
Avec $B, X, Y$ comme dans l'introduction de cette section, supposons en
outre que $X \to B$ soit propre. Alors le sous-foncteur
$\mathit{Isom}_B(Y, X) \subset \mathit{Mor}_B(Y, X)$ des isomorphismes est
un sous-espace ouvert.
\end{lemma}

\begin{proof}
Cela résulte immédiatement de Compléments sur les morphismes d'espaces,
lemme \ref{spaces-more-morphisms-lemma-where-isomorphism}.
\end{proof}

\begin{remark}[Invariants numériques]
\label{moduli-remark-Mor-numerical}
Soient $B, X, Y$ comme dans l'introduction de cette section.
Soit $I$ un ensemble et, pour $i \in I$, soit
$E_i \in D(\mathcal{O}_{Y \times_B X})$ parfait.
Soit $P : I \to \mathbf{Z}$ une fonction. Rappelons que
$$
\mathit{Mor}_B(Y, X) \subset
\Hilbfunctor_{Y \times_B X/B}
$$
est un sous-espace ouvert ; voir Quotients, lemme
\ref{quot-lemma-Mor-into-Hilb-open}. On peut donc définir
$$
\mathit{Mor}^P_B(Y, X) =
\mathit{Mor}_B(Y, X) \cap \Hilbfunctor^P_{Y \times_B X/B}
$$
où $\Hilbfunctor^P_{Y \times_B X/B}$ est comme dans la remarque
\ref{moduli-remark-hilb-numerical}. Le morphisme
$$
\mathit{Mor}^P_B(Y, X) \longrightarrow \mathit{Mor}_B(Y, X)
$$
est une immersion fermée plate qui est une immersion ouverte et fermée,
par exemple si $I$ est fini, si $B$ est localement noethérien, ou si
$I = \mathbf{Z}$, $E_i = \mathcal{L}^{\otimes i}$ pour un
$\mathcal{O}_{Y \times_B X}$-module inversible $\mathcal{L}$.
Dans ce dernier cas, nous employons parfois la notation
$\mathit{Mor}^{P, \mathcal{L}}_B(Y, X)$.
\end{remark}

\begin{lemma}
\label{moduli-lemma-Mor-qc-over-base}
Avec $B, X, Y$ comme dans l'introduction de cette section, soit
$\mathcal{L}$ ample sur $X/B$ et soit $\mathcal{N}$ ample sur $Y/B$.
Voir Diviseurs sur les espaces, définition
\ref{spaces-divisors-definition-relatively-ample}.
Soit $P$ un polynôme numérique. Alors
$$
\mathit{Mor}^{P, \mathcal{M}}_B(Y, X) \longrightarrow B
$$
est séparé et de présentation finie, où
$\mathcal{M} = \text{pr}_1^*\mathcal{N}
\otimes_{\mathcal{O}_{Y \times_B X}} \text{pr}_2^*\mathcal{L}$.
\end{lemma}

\begin{proof}
D'après le lemme \ref{moduli-lemma-Mor-s-lfp}, le morphisme
$\mathit{Mor}_B(Y, X) \to B$ est séparé et localement de présentation
finie. Il suffit donc de montrer que le sous-espace ouvert et fermé
$\mathit{Mor}^{P, \mathcal{M}}_B(Y, X)$ de la remarque
\ref{moduli-remark-Mor-numerical} est quasi-compact sur $B$.

\medskip\noindent
La question est locale pour la topologie \'etale sur $B$
(Morphismes d'espaces, lemme
\ref{spaces-morphisms-lemma-quasi-compact-local}).
Nous pouvons donc supposer $B$ affine.

\medskip\noindent
Supposons $B = \Spec(\Lambda)$. Notons que $X$ et $Y$ sont des schémas et
que $\mathcal{L}$ et $\mathcal{N}$ sont des faisceaux inversibles amples
sur $X$ et $Y$ (cela résulte immédiatement des définitions).
Écrivons $\Lambda = \colim \Lambda_i$ comme limite inductive de ses
sous-algèbres de type fini sur $\mathbf{Z}$. On peut alors trouver un $i$
et un système $X_i, Y_i, \mathcal{L}_i, \mathcal{N}_i$ comme dans le lemme
sur $B_i = \Spec(\Lambda_i)$, dont le changement de base à $B$ donne
$X, Y, \mathcal{L}, \mathcal{N}$. Cela résulte de Limites, lemmes
\ref{limits-lemma-descend-finite-presentation} (pour trouver $X_i$, $Y_i$),
\ref{limits-lemma-descend-invertible-modules} (pour trouver
$\mathcal{L}_i$, $\mathcal{N}_i$),
\ref{limits-lemma-descend-separated-finite-presentation} (pour rendre
$X_i \to B_i$ séparé), \ref{limits-lemma-eventually-proper} (pour rendre
$Y_i \to B_i$ propre) et \ref{limits-lemma-limit-ample} (pour rendre
$\mathcal{L}_i$, $\mathcal{N}_i$ amples).
Comme
$$
\mathit{Mor}_B(Y, X) = B \times_{B_i} \mathit{Mor}_{B_i}(Y_i, X_i)
$$
et qu'il en va de même de $\mathit{Mor}^P_B(Y, X)$, on se ramène au cas
étudié dans le paragraphe suivant.

\medskip\noindent
Supposons $B$ affine et noethérien. D'après Propriétés, lemme
\ref{properties-lemma-ample-on-product}, $\mathcal{M}$ est ample.
D'après le lemme \ref{moduli-lemma-hilb-qc-over-base},
$\Hilbfunctor^{P, \mathcal{M}}_{Y \times_B X/B}$ est de présentation finie
sur $B$, donc noethérien. Par construction,
$$
\mathit{Mor}^{P, \mathcal{M}}_B(Y, X) =
\mathit{Mor}_B(Y, X) \cap
\Hilbfunctor^{P, \mathcal{M}}_{Y \times_B X/B}
$$
est un sous-espace ouvert de
$\Hilbfunctor^{P, \mathcal{M}}_{Y \times_B X/B}$, donc quasi-compact
(puisque tout ouvert d'un espace algébrique noethérien est quasi-compact).
\end{proof}






\section{Propriétés du champ des schémas propres polarisés}
\label{moduli-section-polarized}

\noindent
Dans cette section, nous étudions des propriétés du champ de modules
$$
\Polarizedstack \longrightarrow \Spec(\mathbf{Z})
$$
dont la catégorie des sections sur un schéma $S$ est la catégorie des
schémas propres, plats et de présentation finie sur $S$, munis d'un
faisceau inversible relativement ample. C'est un champ algébrique d'après
Quotients, théorème \ref{quot-theorem-polarized-algebraic}.

\begin{lemma}
\label{moduli-lemma-polarized-diagonal-separated-fp}
La diagonale de $\Polarizedstack$ est séparée et de présentation finie.
\end{lemma}

\begin{proof}
Rappelons que $\Polarizedstack$ est un champ algébrique qui préserve les
limites ; voir Quotients, lemme \ref{quot-lemma-polarized-limits}.
D'après Limites de champs, lemme
\ref{stacks-limits-lemma-limit-preserving-diagonal}, cela implique que
$\Delta : \Polarizedstack \to \Polarizedstack \times \Polarizedstack$
préserve les limites. Ainsi, $\Delta$ est localement de présentation finie
d'après Limites de champs, proposition
\ref{stacks-limits-proposition-characterize-locally-finite-presentation}.

\medskip\noindent
Montrons que $\Delta$ est séparée. Pour cela, il suffit de montrer qu'étant
donnés un schéma affine $U$ et deux objets
$\upsilon = (Y, \mathcal{N})$ et $\chi = (X, \mathcal{L})$
de $\Polarizedstack$ sur $U$, l'espace algébrique
$$
\mathit{Isom}_{\Polarizedstack}(\upsilon, \chi)
$$
est séparé. La règle qui, à un isomorphisme
$\upsilon_T \to \chi_T$, associe l'isomorphisme sous-jacent
$Y_T \to X_T$ définit un morphisme
$$
\mathit{Isom}_{\Polarizedstack}(\upsilon, \chi)
\longrightarrow
\mathit{Isom}_U(Y, X)
$$
Puisque nous avons vu aux lemmes \ref{moduli-lemma-Mor-s-lfp} et
\ref{moduli-lemma-Isom-in-Mor} que le but est un espace algébrique séparé, il
suffit de prouver que ce morphisme est séparé.
Étant donné un isomorphisme $f : Y_T \to X_T$ sur un schéma $T/U$, on a
manifestement
$$
\mathit{Isom}_{\Polarizedstack}(\upsilon, \chi)
\times_{\mathit{Isom}_U(Y, X), [f]} T
=
\mathit{Isom}(\mathcal{N}_T, f^*\mathcal{L}_T)
$$
Ici, $[f] : T \to \mathit{Isom}_U(Y, X)$ désigne le point à valeurs dans
$T$ correspondant à $f$, et
$\mathit{Isom}(\mathcal{N}_T, f^*\mathcal{L}_T)$ est l'espace algébrique
étudié à la section \ref{moduli-section-hom-isom}.
Comme cet espace algébrique est affine sur $U$, l'assertion implique que
$\Delta$ est séparée.

\medskip\noindent
Pour achever la démonstration, montrons que $\Delta$ est quasi-compacte.
Comme $\Delta$ est représentable par des espaces algébriques, il suffit de
vérifier que le changement de base de $\Delta$ par un morphisme lisse
surjectif $U \to \Polarizedstack \times \Polarizedstack$ est quasi-compact
(voir par exemple Propriétés des champs, lemme
\ref{stacks-properties-lemma-check-property-covering}).
Nous pouvons supposer que $U = \coprod U_i$ est une union disjointe
d'ouverts affines. Comme $\Polarizedstack$ préserve les limites (voir
ci-dessus), on voit que
$\Polarizedstack \to \Spec(\mathbf{Z})$ est localement de présentation
finie, donc que $U_i \to \Spec(\mathbf{Z})$ est localement de présentation
finie (Limites de champs, proposition
\ref{stacks-limits-proposition-characterize-locally-finite-presentation},
et Morphismes de champs, lemmes
\ref{stacks-morphisms-lemma-composition-finite-presentation} et
\ref{stacks-morphisms-lemma-smooth-locally-finite-presentation}).
En particulier, $U_i$ est affine noethérien. On se ramène ainsi au cas
étudié dans le paragraphe suivant.

\medskip\noindent
Dans ce paragraphe, étant donnés un schéma affine noethérien $U$ et deux
objets $\upsilon = (Y, \mathcal{N})$ et
$\chi = (X, \mathcal{L})$ de $\Polarizedstack$ sur $U$, montrons que
l'espace algébrique
$$
\mathit{Isom}_{\Polarizedstack}(\upsilon, \chi)
$$
est quasi-compact. Comme les composantes connexes de $U$ sont ouvertes et
fermées, on peut remplacer $U$ par celles-ci. Nous pouvons donc supposer,
et supposerons, que $U$ est connexe.
Soit $u \in U$ un point. Soit $P$ le polynôme de Hilbert
$n \mapsto \chi(Y_u, \mathcal{N}_u^{\otimes n})$ ; voir Variétés, lemme
\ref{varieties-lemma-numerical-polynomial-from-euler}.
Puisque $U$ est connexe et que les fonctions
$u \mapsto \chi(Y_u, \mathcal{N}_u^{\otimes n})$
sont localement constantes (voir Catégories dérivées des schémas, lemme
\ref{perfect-lemma-chi-locally-constant-geometric}), on obtient le même
polynôme de Hilbert en tout point de $U$. Posons
$\mathcal{M} = \text{pr}_1^*\mathcal{N}
\otimes_{\mathcal{O}_{Y \times_U X}} \text{pr}_2^*\mathcal{L}$
sur $Y \times_U X$. Étant donné
$(f, \varphi) \in \mathit{Isom}_{\Polarizedstack}(\upsilon, \chi)(T)$
pour un schéma $T$ sur $U$, on a pour tout $t \in T$
$$
\chi(Y_t, (\text{id} \times f)^*\mathcal{M}^{\otimes n}) =
\chi(Y_t,
\mathcal{N}_t^{\otimes n} \otimes_{\mathcal{O}_{Y_t}}
f_t^*\mathcal{L}_t^{\otimes n}) =
\chi(Y_t, \mathcal{N}_t^{\otimes 2n}) = P(2n)
$$
où, dans l'égalité du milieu, nous utilisons l'isomorphisme
$\varphi : f^*\mathcal{L}_T \to \mathcal{N}_T$.
Posant $P'(t) = P(2t)$, on voit que le morphisme
$$
\mathit{Isom}_{\Polarizedstack}(\upsilon, \chi)
\longrightarrow
\mathit{Isom}_U(Y, X)
$$
(voir plus haut) a son image contenue dans l'intersection
$$
\mathit{Isom}_U(Y, X) \cap \mathit{Mor}^{P', \mathcal{M}}_U(Y, X)
$$
Cette intersection est une intersection de sous-espaces ouverts de
$\mathit{Mor}_U(Y, X)$ (voir le lemme \ref{moduli-lemma-Isom-in-Mor} et la
remarque \ref{moduli-remark-Mor-numerical}).
Or $\mathit{Mor}^{P', \mathcal{M}}_U(Y, X)$ est un espace algébrique
noethérien, puisqu'il est de présentation finie sur $U$ d'après le lemme
\ref{moduli-lemma-Mor-qc-over-base}. L'intersection est donc elle aussi un espace
algébrique noethérien. Comme le morphisme
$$
\mathit{Isom}_{\Polarizedstack}(\upsilon, \chi)
\longrightarrow
\mathit{Isom}_U(Y, X) \cap \mathit{Mor}^{P', \mathcal{M}}_U(Y, X)
$$
est affine (voir plus haut), on conclut.
\end{proof}

\begin{lemma}
\label{moduli-lemma-polarized-qs-lfp}
Le morphisme $\Polarizedstack \to \Spec(\mathbf{Z})$ est quasi-séparé et
localement de présentation finie.
\end{lemma}

\begin{proof}
Pour vérifier que $\Polarizedstack \to \Spec(\mathbf{Z})$ est
quasi-séparé, il faut montrer que sa diagonale est quasi-compacte et
quasi-séparée. Cela résulte immédiatement du lemme
\ref{moduli-lemma-polarized-diagonal-separated-fp}.
Pour prouver que $\Polarizedstack \to \Spec(\mathbf{Z})$ est localement de
présentation finie, il suffit de montrer que $\Polarizedstack$ préserve les
limites ; voir Limites de champs, proposition
\ref{stacks-limits-proposition-characterize-locally-finite-presentation}.
C'est Quotients, lemme \ref{quot-lemma-polarized-limits}.
\end{proof}

\begin{lemma}
\label{moduli-lemma-bounded-polarized}
Soit $n \geq 1$ un entier et soit $P$ un polynôme numérique.
Soit
$$
T \subset |\Polarizedstack|
$$
un sous-ensemble possédant la propriété suivante : pour tout $\xi \in T$,
il existe un corps $k$ et un objet $(X, \mathcal{L})$ de
$\Polarizedstack$ sur $k$ représentant $\xi$, tels que
\begin{enumerate}
\item le polynôme de Hilbert de $\mathcal{L}$ sur $X$ soit $P$, et
\item il existe une immersion fermée $i : X \to \mathbf{P}^n_k$ telle que
$i^*\mathcal{O}_{\mathbf{P}^n}(1) \cong \mathcal{L}$.
\end{enumerate}
Alors $T$ est un espace topologique noethérien, en particulier
quasi-compact.
\end{lemma}

\begin{proof}
Remarquons que $|\Polarizedstack|$ est un espace topologique localement
noethérien ; voir Morphismes de champs, lemme
\ref{stacks-morphisms-lemma-Noetherian-topology}
(on utilise également le fait que $\Spec(\mathbf{Z})$ est noethérien et
que $\Polarizedstack$ est donc un champ algébrique localement noethérien
d'après le lemme \ref{moduli-lemma-polarized-qs-lfp} et Morphismes de champs,
lemme
\ref{stacks-morphisms-lemma-locally-finite-type-locally-noetherian}).
Tout sous-ensemble quasi-compact de $|\Polarizedstack|$ est donc un espace
topologique noethérien, et tout sous-ensemble d'un tel espace est lui aussi
noethérien ; voir Topologie, lemmes
\ref{topology-lemma-finite-union-Noetherian} et
\ref{topology-lemma-Noetherian}.
Il ne reste donc qu'à trouver un sous-ensemble quasi-compact contenant
$T$.

\medskip\noindent
D'après le lemme \ref{moduli-lemma-hilb-proper-over-base}, l'espace algébrique
$$
H =
\Hilbfunctor^{P, \mathcal{O}(1)}_{\mathbf{P}^n_\mathbf{Z}/\Spec(\mathbf{Z})}
$$
est propre sur $\Spec(\mathbf{Z})$. D'après Quotients, lemme
\ref{quot-lemma-extend-hilb-to-spaces}\footnote{Nous verrons plus loin
(insérer ici une référence future) que $H$ est un schéma ; l'emploi de ce
lemme et de Quotients, lemme
\ref{quot-lemma-extend-polarized-to-spaces}, n'est donc pas nécessaire.},
le morphisme identité de $H$ correspond à un sous-espace fermé
$$
Z \subset \mathbf{P}^n_H
$$
qui est propre, plat et de présentation finie sur $H$, tel que la
restriction $\mathcal{N} = \mathcal{O}(1)|_Z$ soit relativement ample sur
$Z/H$ et ait le polynôme de Hilbert $P$ sur les fibres de $Z \to H$.
En particulier, la paire $(Z \to H, \mathcal{N})$ définit un morphisme
$$
H \longrightarrow \Polarizedstack
$$
qui envoie un morphisme de schémas $U \to H$ sur le morphisme classifiant
de la famille $(Z_U \to U, \mathcal{N}_U)$ ; voir Quotients, lemme
\ref{quot-lemma-extend-polarized-to-spaces}.
Comme $H$ est un espace algébrique noethérien (puisqu'il est propre sur
$\mathbf{Z})$), on voit que $|H|$ est noethérien, donc quasi-compact.
L'application
$$
|H| \longrightarrow |\Polarizedstack|
$$
est continue ; son image est donc quasi-compacte.
Il suffit ainsi de prouver que $T$ est contenu dans l'image de
$|H| \to |\Polarizedstack|$.
Or les hypothèses (1) et (2) expriment exactement ce fait : tout choix
d'une immersion fermée $i : X \to \mathbf{P}^n_k$ telle que
$i^*\mathcal{O}_{\mathbf{P}^n}(1) \cong \mathcal{L}$ fournit un point à
valeurs dans $k$ de $H$, par l'interprétation modulaire de $H$.
Cela achève la démonstration du lemme.
\end{proof}






\section{Propriétés des modules de complexes sur un morphisme propre}
\label{moduli-section-complexes}

\noindent
Soit $f : X \to B$ un morphisme d'espaces algébriques propre, plat et de
présentation finie. Alors le champ $\Complexesstack_{X/B}$ qui paramètre
les complexes relativement parfaits dont les groupes d'auto-extensions de
degré négatif sont nuls est algébrique. Voir Quotients, théorème
\ref{quot-theorem-complexes-algebraic}.

\begin{lemma}
\label{moduli-lemma-complexes-diagonal-affine-fp}
La diagonale de $\Complexesstack_{X/B}$ sur $B$ est affine et de
présentation finie.
\end{lemma}

\begin{proof}
La représentabilité de la diagonale par des espaces algébriques a été
établie dans Quotients, lemme \ref{quot-lemma-complexes-diagonal}.
Il ressort de la démonstration qu'il faut montrer ceci : étant donnés un
schéma $T$ sur $B$ et des objets $E, E' \in D(\mathcal{O}_{X_T})$ tels que
$(T, E)$ et $(T, E')$ soient des objets de la catégorie fibre de
$\Complexesstack_{X/B}$ sur $T$, le morphisme
$\mathit{Isom}(E, E') \to T$ est affine et de présentation finie.
Ici, $\mathit{Isom}(E, E')$ est le foncteur
$$
(\Sch/T)^{opp} \to \textit{Ensembles},\quad
T' \mapsto \{\varphi : E_{T'} \to E'_{T'}
\text{ isomorphisme dans }D(\mathcal{O}_{X_{T'}})\}
$$
où $E_{T'}$ et $E'_{T'}$ sont les images inverses dérivées de $E$ et de
$E'$ sur $X_{T'}$. Considérons le foncteur
$H = \SheafHom(E, E')$ défini par la règle
$$
(\Sch/T)^{opp} \to \textit{Ensembles},\quad
T' \mapsto \Hom_{\mathcal{O}_{X_{T'}}}(E_T, E'_T)
$$
D'après Quotients, lemme
\ref{quot-lemma-complexes-open-neg-exts-vanishing}, c'est un espace
algébrique affine et de présentation finie sur $T$.
Il en va de même de $H' = \SheafHom(E', E)$,
$I = \SheafHom(E, E)$ et $I' = \SheafHom(E', E')$.
On voit donc que
$$
\mathit{Isom}(E, E') = (H' \times_T H) \times_{c, I \times_T I', \sigma} T
$$
où $c(\varphi', \varphi) =
(\varphi \circ \varphi', \varphi' \circ \varphi)$ et
$\sigma = (\text{id}, \text{id})$ (comparer à la démonstration de
Quotients, proposition \ref{quot-proposition-isom}).
Ainsi, $\mathit{Isom}(E, E')$ est affine sur $T$ comme produit fibré de
schémas affines sur $T$. De même, $\mathit{Isom}(E, E')$ est de
présentation finie sur $T$.
\end{proof}

\begin{lemma}
\label{moduli-lemma-complexes-qs-lfp}
Le morphisme $\Complexesstack_{X/B} \to B$ est quasi-séparé et localement
de présentation finie.
\end{lemma}

\begin{proof}
Pour vérifier que $\Complexesstack_{X/B} \to B$ est quasi-séparé, il faut
montrer que sa diagonale est quasi-compacte et quasi-séparée.
Cela résulte immédiatement du lemme
\ref{moduli-lemma-complexes-diagonal-affine-fp}.
Pour prouver que $\Complexesstack_{X/B} \to B$ est localement de
présentation finie, il faut montrer que $\Complexesstack_{X/B} \to B$
préserve les limites ; voir Limites de champs, proposition
\ref{stacks-limits-proposition-characterize-locally-finite-presentation}.
Cela résulte de Quotients, lemme \ref{quot-lemma-complexes-limits}
(un petit détail est omis).
\end{proof}





































% OK, start here.
%

\chapter{Modules de courbes}


\phantomsection
\label{moduli-curves-section-phantom}





\section{Introduction}
\label{moduli-curves-section-introduction}

\noindent
Dans ce chapitre, nous étudions quelques-uns des champs de modules de courbes
classiques. L'article célèbre de Deligne et Mumford constitue une référence,
voir \cite{DM}.




\section{Conventions et abus de langage}
\label{moduli-curves-section-conventions}

\noindent
Nous continuons d'employer les conventions et les abus de langage
introduits dans Propriétés des champs, section
\ref{stacks-properties-section-conventions}.
Sauf mention contraire, notre schéma de base sera $\Spec(\mathbf{Z})$.







\section{Le champ des courbes}
\label{moduli-curves-section-stack-curves}

\noindent
Cette section fait suite à Quot, section \ref{quot-section-curves}.
Soit $\Curvesstack$ le champ dont la catégorie des sections au-dessus d'un
schéma $S$ est la catégorie des familles de courbes sur $S$.
Il importe de garder à l'esprit qu'une
{\it famille de courbes} est un morphisme $f : X \to S$, où $X$
est un espace algébrique (!) et où $f$ est plat, propre, de présentation finie
et de dimension relative $\leq 1$.
Nous savons déjà que $\Curvesstack$ est un
champ algébrique sur $\mathbf{Z}$, voir Quot, théorème
\ref{quot-theorem-curves-algebraic}.
Si nous n'autorisions pas les espaces algébriques dans la définition de
notre champ, ce théorème serait faux.

\medskip\noindent
On note souvent un changement de base par un indice, mais nous ne pouvons pas
employer cette notation pour $\Curvesstack$, car $\Curvesstack_S$
désigne chez nous la catégorie fibre au-dessus de $S$.
C'est pourquoi, dans Quot, remarque \ref{quot-remark-curves-base-change},
nous avons noté $B\text{-}\Curvesstack$ le changement de base
$$
B\text{-}\Curvesstack = \Curvesstack \times B
$$
à l'espace algébrique $B$. Le produit de droite est pris au-dessus de
l'objet final, c'est-à-dire au-dessus de $\Spec(\mathbf{Z})$. L'objet de gauche
est le champ classifiant les familles de courbes sur la catégorie des schémas
sur $B$. En particulier, si $k$ est un corps, alors
$$
k\text{-}\Curvesstack = \Curvesstack \times \Spec(k)
$$
est le champ de modules classifiant les familles de courbes sur la catégorie
des schémas sur $k$.
Avant de poursuivre, effectuons une vérification de cohérence.

\begin{lemma}
\label{moduli-curves-lemma-extend-curves-to-spaces}
Soit $T \to B$ un morphisme d'espaces algébriques. La catégorie
$$
\Mor_B(T, B\text{-}\Curvesstack) = \Mor(T, \Curvesstack)
$$
est la catégorie des familles de courbes sur $T$.
\end{lemma}

\begin{proof}
Une famille de courbes sur $T$ est un morphisme $f : X \to T$ d'espaces
algébriques, qui est plat, propre, de présentation finie et de
dimension relative $\leq 1$ (Morphismes d'espaces, définition
\ref{spaces-morphisms-definition-relative-dimension}).
C'est exactement la même définition que dans Quot, situation
\ref{quot-situation-curves}, si ce n'est que la base $T$ peut être
un espace algébrique. Notre catégorie de base par défaut pour les
champs et espaces algébriques est la catégorie des schémas ; le lemme
ne découle donc pas immédiatement des définitions. Cela étant dit,
nous invitons le lecteur à sauter la démonstration.

\medskip\noindent
D'après la description de $B\text{-}\Curvesstack$ comme produit donnée
ci-dessus, il suffit de démontrer le lemme dans le cas absolu. Choisissons
un schéma $U$ et un morphisme étale surjectif $p : U \to T$.
Posons $R = U \times_T U$, avec projections $s, t : R \to U$.

\medskip\noindent
Soit $v : T \to \Curvesstack$ un morphisme. Alors $v \circ p$
correspond à une famille de courbes $X_U \to U$. Le
$2$-morphisme canonique $v \circ p \circ t \to v \circ p \circ s$
est un isomorphisme $\varphi : X_U \times_{U, s} R \to X_U \times_{U, t} R$.
Cet isomorphisme satisfait à la condition de cocycle sur
$R \times_{s, t} R$.
D'après Amorçage, lemme \ref{bootstrap-lemma-descend-algebraic-space},
nous obtenons un morphisme d'espaces algébriques $X \to T$
dont l'image inverse sur $U$ est égale à $X_U$, compatiblement avec $\varphi$.
Puisque $\{U \to T\}$ est un recouvrement étale, le morphisme
$X \to T$ est plat, propre et de présentation finie d'après
Descente pour les espaces, lemmes
\ref{spaces-descent-lemma-descending-property-flat},
\ref{spaces-descent-lemma-descending-property-proper} et
\ref{spaces-descent-lemma-descending-property-finite-presentation}.
De plus, $X \to T$ est de dimension relative $\leq 1$, car cette propriété
est locale pour la topologie étale. Ainsi, $X \to T$ est une famille de
courbes sur $T$.

\medskip\noindent
Réciproquement, soit $X \to T$ une famille de courbes. Alors le
changement de base $X_U$ détermine un morphisme $w : U \to \Curvesstack$,
et l'isomorphisme canonique $X_U \times_{U, s} R \to X_U \times_{U, t} R$
détermine une $2$-flèche $w \circ s \to w \circ t$ satisfaisant à la
condition de cocycle. On obtient ainsi un morphisme
$v : T = [U/R] \to \Curvesstack$ par la propriété universelle du
quotient $[U/R]$ ; voir
Groupoïdes dans les espaces, lemme
\ref{spaces-groupoids-lemma-quotient-stack-2-coequalizer}.
(En fait, il est ici bien plus simple de revenir à la situation antérieure
à l'introduction de notre abus de langage et de construire directement
le foncteur $\Sch/T \to \Curvesstack$ qui « est » le
morphisme $T \to \Curvesstack$.)

\medskip\noindent
Nous omettons de vérifier que les constructions ci-dessus
s'étendent aux morphismes entre objets et sont quasi-inverses l'une de l'autre.
\end{proof}






\section{Le champ des courbes polarisées}
\label{moduli-curves-section-polarized-curves}

\noindent
Dans cette section, nous développons une partie des considérations de
Quot, remarque \ref{quot-remark-alternative-approach-curves}.
Considérons le $2$-produit fibré
$$
\xymatrix{
\Curvesstack \times_{\Spacesstack'_{fp, flat, proper}}
\Polarizedstack \ar[r] \ar[d] &
\Polarizedstack \ar[d] \\
\Curvesstack \ar[r] &
\Spacesstack'_{fp, flat, proper}
}
$$
Nous notons ce $2$-produit fibré
$$
\textit{PolarizedCurves} =
\Curvesstack
\times_{\Spacesstack'_{fp, flat, proper}}
\Polarizedstack
$$
Ce produit fibré paramètre les courbes polarisées, c'est-à-dire les familles
de courbes munies d'un faisceau inversible relativement ample.
Plus précisément, un objet de
$\textit{PolarizedCurves}$
est un couple $(X \to S, \mathcal{L})$ tel que
\begin{enumerate}
\item $X \to S$ est un morphisme de schémas propre, plat,
de présentation finie et de dimension relative $\leq 1$, et
\item $\mathcal{L}$ est un $\mathcal{O}_X$-module inversible
qui est relativement ample sur $X/S$.
\end{enumerate}
Un morphisme $(X' \to S', \mathcal{L}') \to (X \to S, \mathcal{L})$
entre objets de
$\textit{PolarizedCurves}$
est donné par un triplet $(f, g, \varphi)$,
où $f : X' \to X$ et $g : S' \to S$
sont des morphismes de schémas s'inscrivant dans un diagramme commutatif
$$
\xymatrix{
X' \ar[d] \ar[r]_f & X \ar[d] \\
S' \ar[r]^g & S
}
$$
qui induit un isomorphisme $X' \to S' \times_S X$ ; autrement dit, le
diagramme est cartésien, et $\varphi : f^*\mathcal{L} \to \mathcal{L}'$
est un isomorphisme. La composition se définit de manière évidente.

\begin{lemma}
\label{moduli-curves-lemma-polarized-curves-in-polarized}
Le morphisme
$\textit{PolarizedCurves} \to
\Polarizedstack$ est une immersion ouverte et fermée.
\end{lemma}

\begin{proof}
Cela résulte de ce que le $1$-morphisme
$\Curvesstack \to \Spacesstack'_{fp, flat, proper}$
est représentable par des immersions ouvertes et fermées ; voir
Quot, lemme \ref{quot-lemma-curves-open-and-closed-in-spaces}.
\end{proof}

\begin{lemma}
\label{moduli-curves-lemma-polarized-curves-over-curves}
Le morphisme
$\textit{PolarizedCurves} \to \Curvesstack$
est lisse et surjectif.
\end{lemma}

\begin{proof}
Surjectivité. Soient $k$ un corps et $X$ un espace algébrique propre
sur $k$, de dimension $\leq 1$, c'est-à-dire un objet de $\Curvesstack$ sur $k$.
D'après Espaces sur les corps, lemme
\ref{spaces-over-fields-lemma-codim-1-point-in-schematic-locus}
l'espace algébrique $X$ est un schéma. Ainsi, $X$
est un schéma propre de dimension $\leq 1$ sur $k$.
D'après Variétés, lemme \ref{varieties-lemma-dim-1-proper-projective},
$X$ est H-projectif sur $\kappa$.
Il existe en particulier un $\mathcal{O}_X$-module inversible ample
$\mathcal{L}$ sur $X$. Alors $(X, \mathcal{L})$ est un objet
de $\textit{PolarizedCurves}$ sur
$k$ dont l'image est $X$.

\medskip\noindent
Lissité. Soit $X \to S$ un objet de $\Curvesstack$, c'est-à-dire un
morphisme $S \to \Curvesstack$. Il est clair que
$$
\textit{PolarizedCurves}
\times_{\Curvesstack} S
\subset \Picardstack_{X/S}
$$
est le sous-champ des objets $(T/S, \mathcal{L}/X_T)$ tels que
$\mathcal{L}$ soit ample sur $X_T/T$. C'est un sous-champ ouvert d'après
Descente pour les espaces, lemme \ref{spaces-descent-lemma-ample-in-neighbourhood}.
Puisque $\Picardstack_{X/S} \to S$ est lisse d'après
Champs de modules, lemme \ref{moduli-lemma-pic-curves-smooth},
la conclusion s'ensuit.
\end{proof}

\begin{lemma}
\label{moduli-curves-lemma-etale-locally-scheme}
Soit $X \to S$ une famille de courbes.
Il existe alors un recouvrement étale $\{S_i \to S\}$
tel que $X_i = X \times_S S_i$ soit un schéma. Nous pouvons même
supposer $X_i$ H-projectif sur $S_i$.
\end{lemma}

\begin{proof}
Il s'agit d'un corollaire immédiat du
lemme \ref{moduli-curves-lemma-polarized-curves-over-curves}.
En effet, l'explicitation des définitions montre que ce lemme fournit un
morphisme lisse surjectif $S' \to S$ tel que $X' = X \times_S S'$
soit muni d'un $\mathcal{O}_{X'}$-module inversible
$\mathcal{L}'$ qui est ample sur $X'/S'$.
Nous pouvons ensuite raffiner le recouvrement lisse $\{S' \to S\}$
par un recouvrement étale $\{S_i \to S\}$ ; voir
Compléments sur les morphismes, lemme \ref{more-morphisms-lemma-etale-dominates-smooth}.
Après remplacement de $S_i$ par un recouvrement ouvert convenable, nous
pouvons supposer $X_i \to S_i$ H-projectif ; voir
Morphismes, lemmes \ref{morphisms-lemma-proper-ample-locally-projective} et
\ref{morphisms-lemma-characterize-locally-projective}
(ce point est aussi étudié en détail dans
Compléments sur les morphismes, section \ref{more-morphisms-section-projective}).
\end{proof}






\section{Propriétés du champ des courbes}
\label{moduli-curves-section-properties}

\noindent
Le lemme suivant n'est pas vrai pour les modules de surfaces ; voir
la remarque \ref{moduli-curves-remark-boundedness-aut-does-not-work-surfaces}.

\begin{lemma}
\label{moduli-curves-lemma-curves-diagonal-separated-fp}
La diagonale de $\Curvesstack$ est séparée
et de présentation finie.
\end{lemma}

\begin{proof}
Rappelons que $\Curvesstack$ est un champ algébrique qui préserve les limites ;
voir Quot, lemme \ref{quot-lemma-curves-limits}.
D'après Limites de champs, lemme
\ref{stacks-limits-lemma-limit-preserving-diagonal}, cela implique que
$\Delta : \Polarizedstack \to \Polarizedstack \times \Polarizedstack$
préserve les limites. Par conséquent, $\Delta$ est localement de présentation
finie d'après Limites de champs, proposition
\ref{stacks-limits-proposition-characterize-locally-finite-presentation}.

\medskip\noindent
Montrons que $\Delta$ est séparée. Pour cela, il suffit de montrer que,
pour un schéma $U$ et deux objets $Y \to U$ et $X \to U$ de
$\Curvesstack$ au-dessus de $U$, l'espace algébrique
$$
\mathit{Isom}_U(Y, X)
$$
est séparé. Il résulte en effet de Champs de modules, lemmes
\ref{moduli-lemma-Mor-s-lfp} et
\ref{moduli-lemma-Isom-in-Mor}, que cet espace
algébrique est séparé.

\medskip\noindent
Pour achever la démonstration, montrons que $\Delta$ est quasi-compacte.
Puisque $\Delta$ est représentable par des espaces algébriques, il suffit
de vérifier que le changement de base de $\Delta$ par un morphisme lisse
surjectif $U \to \Curvesstack \times \Curvesstack$ est quasi-compact
(voir, par exemple, Propriétés des champs, lemme
\ref{stacks-properties-lemma-check-property-covering}).
Choisissons pour $U = \coprod U_i$ une réunion disjointe d'ouverts affines,
munie d'un morphisme lisse surjectif
$$
U \longrightarrow
\textit{PolarizedCurves} \times \textit{PolarizedCurves}
$$
Alors $U \to \Curvesstack \times \Curvesstack$ est surjectif
et lisse, car $\textit{PolarizedCurves} \to \Curvesstack$
est surjectif et lisse d'après le lemme
\ref{moduli-curves-lemma-polarized-curves-over-curves}.
Comme $\textit{PolarizedCurves}$ préserve les limites
(d'après Axiomes d'Artin, lemme
\ref{artin-lemma-fibre-product-limit-preserving},
et Quot, lemmes \ref{quot-lemma-curves-limits},
\ref{quot-lemma-polarized-limits} et
\ref{quot-lemma-spaces-limits}), le morphisme
$\textit{PolarizedCurves} \to \Spec(\mathbf{Z})$ est localement de
présentation finie. Il en va donc de même de $U_i \to \Spec(\mathbf{Z})$
(Limites de champs, proposition
\ref{stacks-limits-proposition-characterize-locally-finite-presentation}
et Morphismes de champs, lemmes
\ref{stacks-morphisms-lemma-composition-finite-presentation} et
\ref{stacks-morphisms-lemma-smooth-locally-finite-presentation}).
En particulier, $U_i$ est affine noethérien. Nous sommes ainsi ramenés
au cas étudié dans le paragraphe suivant.

\medskip\noindent
Dans ce paragraphe, étant donnés un schéma affine noethérien $U$ et deux objets
$(Y, \mathcal{N})$ et $(X, \mathcal{L})$
de $\textit{PolarizedCurves}$ au-dessus de $U$, nous montrons que l'espace algébrique
$$
\mathit{Isom}_U(Y, X)
$$
est quasi-compact. Comme les composantes connexes de $U$ sont ouvertes et
fermées, nous pouvons remplacer $U$ par celles-ci. Nous supposons donc
désormais $U$ connexe. Soit $u \in U$ un point. Soient $Q$ et $P$ les
polynômes de Hilbert de ces familles, c'est-à-dire
$$
Q(n) = \chi(Y_u, \mathcal{N}_u^{\otimes n})
\quad\text{et}\quad
P(n) = \chi(X_u, \mathcal{L}_u^{\otimes n})
$$
voir Variétés, lemme \ref{varieties-lemma-numerical-polynomial-from-euler}.
Comme $U$ est connexe et que les fonctions
$u \mapsto \chi(Y_u, \mathcal{N}_u^{\otimes n})$ et
$u \mapsto \chi(X_u, \mathcal{L}_u^{\otimes n})$
 sont localement constantes (voir
Catégories dérivées de schémas, lemme
\ref{perfect-lemma-chi-locally-constant-geometric})
nous obtenons le même polynôme de Hilbert en tout point de $U$.
Posons
$$
\mathcal{M} = \text{pr}_1^*\mathcal{N}
\otimes_{\mathcal{O}_{Y \times_U X}} \text{pr}_2^*\mathcal{L}
$$
sur $Y \times_U X$. Pour
$(f, \varphi) \in \mathit{Isom}_U(Y, X)(T)$, où $T$ est un schéma sur $U$,
et pour tout $t \in T$, on a
\begin{align*}
\chi(Y_t, (\text{id} \times f)^*\mathcal{M}^{\otimes n})
& =
\chi(Y_t,
\mathcal{N}_t^{\otimes n} \otimes_{\mathcal{O}_{Y_t}}
f_t^*\mathcal{L}_t^{\otimes n}) \\
& =
n\deg(\mathcal{N}_t) + n\deg(f_t^*\mathcal{L}_t) +
\chi(Y_t, \mathcal{O}_{Y_t}) \\
& =
Q(n) + n\deg(\mathcal{L}_t) \\
& =
Q(n) + P(n) - P(0)
\end{align*}
d'après Riemann--Roch pour les courbes propres, plus précisément d'après
Variétés, définition \ref{varieties-definition-degree-invertible-sheaf},
lemme \ref{varieties-lemma-degree-tensor-product},
et le fait que $f_t$ est un isomorphisme.
En posant $P'(t) = Q(t) + P(t) - P(0)$, nous obtenons
$$
\mathit{Isom}_U(Y, X) =
\mathit{Isom}_U(Y, X) \cap \mathit{Mor}^{P', \mathcal{M}}_U(Y, X)
$$
Cette intersection est une intersection de sous-espaces ouverts de
$\mathit{Mor}_U(Y, X)$ ; voir Champs de modules, lemme
\ref{moduli-lemma-Isom-in-Mor} et remarque
\ref{moduli-remark-Mor-numerical}.
Or $\mathit{Mor}^{P', \mathcal{M}}_U(Y, X)$
est un espace algébrique noethérien, puisqu'il est de présentation
finie sur $U$ d'après Champs de modules, lemme
\ref{moduli-lemma-Mor-qc-over-base}.
L'intersection est donc elle aussi un espace algébrique noethérien,
ce qui achève la démonstration.
\end{proof}

\begin{remark}
\label{moduli-curves-remark-boundedness-aut-does-not-work-surfaces}
L'argument de bornitude employé dans la démonstration du
lemme \ref{moduli-curves-lemma-curves-diagonal-separated-fp}
ne fonctionne pas pour les modules de surfaces ; de fait,
le résultat est faux, notamment parce que des surfaces K3
sur des corps peuvent avoir des groupes discrets d'automorphismes infinis.
La « raison » pour laquelle l'argument ne fonctionne pas est que, sur une
surface projective $S$ sur un corps, étant donnés des faisceaux inversibles
amples $\mathcal{N}$ et $\mathcal{L}$ de polynômes de Hilbert $Q$ et $P$,
il n'existe aucune borne a priori pour le polynôme de Hilbert de
$\mathcal{N} \otimes_{\mathcal{O}_S} \mathcal{L}$.
En termes de théorie de l'intersection, si $H_1$ et $H_2$ sont des
diviseurs de Cartier effectifs et amples sur $S$, il n'existe aucune
borne (supérieure) pour le nombre d'intersection $H_1 \cdot H_2$
en fonction de $H_1 \cdot H_1$ et de $H_2 \cdot H_2$.
\end{remark}

\begin{lemma}
\label{moduli-curves-lemma-curves-qs-lfp}
Le morphisme $\Curvesstack \to \Spec(\mathbf{Z})$ est quasi-séparé et
localement de présentation finie.
\end{lemma}

\begin{proof}
Pour vérifier que $\Curvesstack \to \Spec(\mathbf{Z})$ est quasi-séparé,
il faut montrer que sa diagonale est quasi-compacte et quasi-séparée.
Cela résulte immédiatement du lemme
\ref{moduli-curves-lemma-curves-diagonal-separated-fp}.
Pour montrer que $\Curvesstack \to \Spec(\mathbf{Z})$ est localement de
présentation finie, il suffit de montrer que $\Curvesstack$
préserve les limites ; voir Limites de champs, proposition
\ref{stacks-limits-proposition-characterize-locally-finite-presentation}.
C'est Quot, lemme \ref{quot-lemma-curves-limits}.
\end{proof}






\section{Sous-champs ouverts du champ des courbes}
\label{moduli-curves-section-open}

\noindent
Dans la suite, nous caractériserons souvent un sous-champ ouvert de
$\Curvesstack$ au moyen d'une propriété $P$ des morphismes d'espaces
algébriques. Pour voir que $P$ définit un sous-champ ouvert, il suffit
de vérifier la condition suivante :
\begin{enumerate}
\item[(o)] pour toute famille de courbes $f : X \to S$, il existe
un plus grand sous-schéma ouvert $S' \subset S$ tel que
$f|_{f^{-1}(S')} : f^{-1}(S') \to S'$ possède $P$ et tel que
la formation de $S'$ commute à tout changement de base.
\end{enumerate}
Supposons en effet que (o) soit satisfaite. Choisissons un schéma $U$ et
un morphisme lisse surjectif $m : U \to \Curvesstack$. Posons
$R = U \times_{\Curvesstack} U$ et notons $t, s : R \to U$
les projections. Rappelons que $\Curvesstack = [U/R]$ est une présentation ;
voir Champs algébriques, lemme \ref{algebraic-lemma-stack-presentation} et
définition \ref{algebraic-definition-presentation}.
Par construction de $\Curvesstack$ comme champ des courbes, le morphisme $m$
est le morphisme classifiant d'une famille de courbes $C \to U$. La
$2$-commutativité du diagramme
$$
\xymatrix{
R \ar[r]_s \ar[d]_t & U \ar[d] \\
U \ar[r] & \Curvesstack
}
$$
implique que $C \times_{U, s} R  \cong C \times_{U, t} R$
(isomorphisme de familles de courbes sur $R$). Soit $W \subset U$
le plus grand sous-schéma ouvert tel que
$f|_{f^{-1}(W)} : f^{-1}(W) \to W$ possède $P$, comme dans (o).
Puisque, d'après (o), la formation de $W$ commute aux changements de base,
l'isomorphisme ci-dessus donne $s^{-1}(W) = t^{-1}(W)$.
Ainsi, $W \subset U$ correspond à un sous-champ ouvert
$$
\Curvesstack^P \subset \Curvesstack
$$
d'après Propriétés des champs, lemme
\ref{stacks-properties-lemma-immersion-presentation}.

\medskip\noindent
En conservant les notations du paragraphe précédent, nous affirmons que
le sous-champ ouvert $\Curvesstack^P$ possède les deux propriétés
universelles suivantes :
\begin{enumerate}
\item pour une famille de courbes $X \to S$, les assertions suivantes
sont équivalentes :
\begin{enumerate}
\item le morphisme classifiant $S \to \Curvesstack$ se factorise par
$\Curvesstack^P$,
\item le morphisme $X \to S$ possède $P$ ;
\end{enumerate}
\item pour un schéma $X$ propre sur un corps $k$, de dimension $\leq 1$,
les assertions suivantes sont équivalentes :
\begin{enumerate}
\item le morphisme classifiant $\Spec(k) \to \Curvesstack$ se factorise
par $\Curvesstack^P$,
\item le morphisme $X \to \Spec(k)$ possède $P$.
\end{enumerate}
\end{enumerate}
Cela résulte de la considération du $2$-produit fibré
$$
\xymatrix{
T \ar[r]_p \ar[d]_q & U \ar[d] \\
S \ar[r] & \Curvesstack
}
$$
Observons que $T \to S$ est surjectif et lisse, car il se déduit de
$U \to \Curvesstack$ par changement de base. L'ouvert $S' \subset S$
donné par (o) est donc déterminé par son image inverse dans $T$.
Or, par l'invariance des ouverts de (o) sous changement de base et puisque
$X \times_S T \cong C \times_U T$ par $2$-commutativité, nous obtenons
$q^{-1}(S') = p^{-1}(W)$ comme ouverts de $T$.
Cela implique immédiatement (1). La partie (2) est un cas particulier de (1).

\medskip\noindent
Étant données deux propriétés $P$ et $Q$ des morphismes d'espaces
algébriques, supposons que nous ayons déjà établi que $\Curvesstack^Q$
est un sous-champ ouvert de $\Curvesstack$. La même méthode permet alors
de démontrer que
$\Curvesstack^{Q, P} \subset \Curvesstack^Q$ est ouvert.
Nous omettons les détails.



\section{Courbes à groupes d'automorphismes finis et réduits}
\label{moduli-curves-section-finite-aut}

\noindent
Soit $X$ un schéma propre sur un corps $k$, de dimension $\leq 1$,
c'est-à-dire un objet de $\Curvesstack$ sur $k$.
D'après le lemme \ref{moduli-curves-lemma-curves-diagonal-separated-fp},
l'espace algébrique des automorphismes $\mathit{Aut}(X)$
est de type fini et séparé sur $k$.
En particulier, $\mathit{Aut}(X)$ est un schéma en groupes ; voir
Compléments sur les groupoïdes dans les espaces, lemme
\ref{spaces-more-groupoids-lemma-group-space-scheme-locally-finite-type-over-k}.
Si la caractéristique de $k$ est nulle, alors $\mathit{Aut}(X)$
est réduit et même lisse sur $k$ (Groupoïdes, lemme
\ref{groupoids-lemma-group-scheme-characteristic-zero-smooth}).
Toutefois, en général, $\mathit{Aut}(X)$ n'est pas réduit, même
si $X$ est géométriquement réduit.

\begin{example}[Groupe d'automorphismes non réduit]
\label{moduli-curves-example-non-reduced}
Soit $k$ un corps algébriquement clos de caractéristique $2$.
Posons $Y = Z = \mathbf{P}^1_k$. Choisissons trois points $k$-rationnels
deux à deux distincts $a, b, c$ dans $\mathbf{A}^1_k$. En considérant
$\mathbf{A}^1_k \subset \mathbf{P}^1_k = Y = Z$ comme un sous-schéma ouvert,
nous obtenons une immersion fermée
$$
T =  \Spec(k[t]/(t - a)^2) \amalg \Spec(k[t]/(t - b)^2)
\amalg \Spec(k[t]/(t - c)^2)
\longrightarrow
\mathbf{P}^1_k
$$
Soit $X$ la somme amalgamée dans le diagramme
$$
\xymatrix{
T \ar[r] \ar[d] & Y \ar[d] \\
Z \ar[r] & X
}
$$
Soit $U \subset X$ la partie ouverte affine qui est l'image de
$\mathbf{A}^1_k \amalg \mathbf{A}^1_k$. Nous avons alors un diagramme
égalisateur
$$
\xymatrix{
\mathcal{O}_X(U) \ar[r] &
k[t] \times k[t] \ar@<1ex>[r] \ar@<-1ex>[r] &
k[t]/(t - a)^2 \times k[t]/(t - b)^2 \times k[t]/(t - c)^2
}
$$
Sur l'algèbre des nombres duaux $A = k[\epsilon]$, ce diagramme égalisateur
admet un automorphisme non trivial envoyant $t$ sur $t + \epsilon$. Nous
laissons au lecteur le soin de voir que cet automorphisme s'étend en un
automorphisme de $X$ sur $A$. D'autre part, le lecteur vérifie aisément que
le groupe d'automorphismes de $X$ sur $k$ est fini.
Ainsi, $\mathit{Aut}(X)$ est nécessairement non réduit.
\end{example}

\noindent
Soit $X$ un schéma propre sur un corps $k$, de dimension $\leq 1$,
c'est-à-dire un objet de $\Curvesstack$ sur $k$. Même si $\mathit{Aut}(X)$
est géométriquement réduit, il n'est pas nécessairement de
dimension $0$, quand bien même $X$ est lisse et géométriquement connexe.

\begin{example}[Groupe d'automorphismes lisse de dimension positive]
\label{moduli-curves-example-pos-dim}
Soit $k$ un corps algébriquement clos. Si $X$ est une courbe lisse
de genre $0$, resp.\ $1$, alors son groupe d'automorphismes est de
dimension $3$, resp.\ $1$. En effet, dans le cas du genre $0$, on a
$X \cong \mathbf{P}^1_k$ d'après Courbes algébriques, proposition
\ref{curves-proposition-projective-line}. Puisque
$$
\mathit{Aut}(\mathbf{P}^1_k) = \text{PGL}_{2, k}
$$
comme foncteurs, sa dimension est $3$. D'autre part,
si le genre de $X$ est $1$, alors l'application
$X = \underline{\Hilbfunctor}^1_{X/k} \to
\underline{\Picardfunctor}^1_{X/k}$ est un isomorphisme ; voir
Schémas de Picard des courbes, lemme \ref{pic-lemma-picard-pieces},
et Courbes algébriques, théorème
\ref{curves-theorem-curves-rational-maps}.
Ainsi, $X$ possède une structure de variété abélienne
(puisque $\underline{\Picardfunctor}^1_{X/k} \cong
\underline{\Picardfunctor}^0_{X/k}$).
En particulier, les fibrés tangent et cotangent de $X$ sont triviaux
(Groupoïdes, lemme \ref{groupoids-lemma-group-scheme-module-differentials}).
Nous en concluons que $\dim_k H^0(X, T_X) = 1$, donc que
$\dim \mathit{Aut}(X) \leq 1$. D'autre part, les translations
(en considérant $X$ comme un schéma en groupes) fournissent une partie
de dimension $1$ de $\text{Aut}(X)$ ; sa dimension vaut donc bien $1$.
\end{example}

\noindent
Il existe en fait un sous-champ ouvert de
$\Curvesstack$ qui paramètre les courbes dont le groupe d'automorphismes
est géométriquement réduit et fini.
En voici un énoncé précis.

\begin{lemma}
\label{moduli-curves-lemma-DM-curves}
Il existe un sous-champ ouvert $\Curvesstack^{DM} \subset \Curvesstack$
possédant les propriétés suivantes :
\begin{enumerate}
\item $\Curvesstack^{DM} \subset \Curvesstack$ est le plus grand
sous-champ ouvert qui soit DM ;
\item pour une famille de courbes $X \to S$, les assertions suivantes
sont équivalentes :
\begin{enumerate}
\item le morphisme classifiant $S \to \Curvesstack$ se factorise par
$\Curvesstack^{DM}$,
\item l'espace algébrique en groupes $\mathit{Aut}_S(X)$ est non ramifié
sur $S$ ;
\end{enumerate}
\item pour un schéma $X$ propre sur un corps $k$, de dimension $\leq 1$,
les assertions suivantes sont équivalentes :
\begin{enumerate}
\item le morphisme classifiant $\Spec(k) \to \Curvesstack$ se factorise
par $\Curvesstack^{DM}$ ;
\item $\mathit{Aut}(X)$ est géométriquement réduit sur $k$ et
de dimension $0$ ;
\item $\mathit{Aut}(X) \to \Spec(k)$ est non ramifié.
\end{enumerate}
\end{enumerate}
\end{lemma}

\begin{proof}
L'existence d'un sous-champ ouvert ayant la propriété (1) résulte de
Morphismes de champs, lemme \ref{stacks-morphisms-lemma-open-DM-locus}.
Les points de ce sous-champ ouvert sont caractérisés par (3)(c) d'après
Morphismes de champs, lemme \ref{stacks-morphisms-lemma-points-DM-locus}.
L'équivalence de (3)(b) et (3)(c) affirme qu'un espace algébrique $G$,
localement de type fini, géométriquement réduit et de dimension $0$ sur
un corps $k$, est non ramifié sur $k$.
Tout d'abord, $G$ est un schéma d'après Espaces sur les corps, lemme
\ref{spaces-over-fields-lemma-locally-finite-type-dim-zero}.
Nous pouvons ensuite prendre un ouvert affine de $G$, observer
qu'il est propre sur $k$, puis appliquer Variétés, lemme
\ref{varieties-lemma-proper-geometrically-reduced-global-sections}.
Nous omettons les détails mineurs.

\medskip\noindent
La partie (2) est vraie parce que (3) l'est. En effet, le morphisme
$\mathit{Aut}_S(X) \to S$ est localement de type fini. Nous pouvons donc
vérifier si $\mathit{Aut}_S(X) \to S$ est non ramifié en tout point de
$\mathit{Aut}_S(X)$ en le vérifiant sur les fibres aux points du schéma $S$ ;
voir Morphismes d'espaces, lemme
\ref{spaces-morphisms-lemma-where-unramified}.
Mais, après changement de base en un point de $S$, nous nous ramenons à
l'équivalence de (3)(a) et (3)(c).
\end{proof}

\begin{lemma}
\label{moduli-curves-lemma-in-DM-locus-vector-fields}
Soit $X$ un schéma propre sur un corps $k$, de dimension $\leq 1$.
Alors les propriétés (3)(a), (b), (c) sont également équivalentes à
$\text{Der}_k(\mathcal{O}_X, \mathcal{O}_X) = 0$.
\end{lemma}

\begin{proof}
Dans la discussion précédente, nous avons vu que $G = \mathit{Aut}(X)$
est un schéma en groupes de type fini et séparé sur $\Spec(k)$ ;
on utilise ici le lemme \ref{moduli-curves-lemma-curves-diagonal-separated-fp} et
Compléments sur les groupoïdes dans les espaces, lemme
\ref{spaces-more-groupoids-lemma-group-space-scheme-locally-finite-type-over-k}.
Alors $G$ est non ramifié sur $k$ si et seulement si $\Omega_{G/k} = 0$
(Morphismes, lemme \ref{morphisms-lemma-unramified-omega-zero}).
D'après Groupoïdes, lemme
\ref{groupoids-lemma-group-scheme-module-differentials}, cette annulation
a lieu si $T_{G/k, e} = 0$, où $T_{G/k, e}$ est l'espace tangent à $G$
en l'élément neutre $e \in G(k)$ ; voir Variétés, définition
\ref{varieties-definition-tangent-space}, ainsi que la formule de
Variétés, lemme \ref{varieties-lemma-tangent-space-cotangent-space}.
Puisque $\kappa(e) = k$, l'espace tangent se décrit au moyen des
morphismes $\alpha : \Spec(k[\epsilon]) \to G = \mathit{Aut}(X)$
dont la restriction à $\Spec(k)$ est $e$.
On est ainsi ramené à considérer les automorphismes
$$
\alpha :
X \times_{\Spec(k)} \Spec(k[\epsilon])
\longrightarrow
X \times_{\Spec(k)} \Spec(k[\epsilon])
$$
sur $\Spec(k[\epsilon])$ dont la restriction à $\Spec(k)$ est
$\text{id}_X$. De tels automorphismes sont appelés
automorphismes infinitésimaux.

\medskip\noindent
Les automorphismes infinitésimaux de $X$ correspondent bijectivement
aux dérivations de $\mathcal{O}_X$ sur $k$. Cela résulte de
Compléments sur les morphismes, lemmes
\ref{more-morphisms-lemma-difference-derivation} et
\ref{more-morphisms-lemma-action-by-derivations} (nous n'avons besoin que
du premier, puisque le sens réciproque ne nous intéresse pas ; consulter
aussi Compléments sur les morphismes, remarque
\ref{more-morphisms-remark-another-special-case}, pour une explication).
Pour un autre argument établissant cette égalité, nous renvoyons le lecteur
à Problèmes de déformation, lemme \ref{examples-defos-lemma-schemes-TI}.
\end{proof}





\section{Courbes de Cohen-Macaulay}
\label{moduli-curves-section-CM}

\noindent
Il existe un sous-champ ouvert de $\Curvesstack$ qui paramètre
les « courbes » de Cohen-Macaulay.

\begin{lemma}
\label{moduli-curves-lemma-CM-curves}
Il existe un sous-champ ouvert $\Curvesstack^{CM} \subset \Curvesstack$
tel que :
\begin{enumerate}
\item pour une famille de courbes $X \to S$, les assertions suivantes
sont équivalentes :
\begin{enumerate}
\item le morphisme classifiant $S \to \Curvesstack$ se factorise par
$\Curvesstack^{CM}$,
\item le morphisme $X \to S$ est Cohen-Macaulay ;
\end{enumerate}
\item pour un schéma $X$ propre sur un corps $k$ et tel que $\dim(X) \leq 1$,
les assertions suivantes sont équivalentes :
\begin{enumerate}
\item le morphisme classifiant $\Spec(k) \to \Curvesstack$ se factorise
par $\Curvesstack^{CM}$,
\item $X$ est Cohen-Macaulay.
\end{enumerate}
\end{enumerate}
\end{lemma}

\begin{proof}
Soit $f : X \to S$ une famille de courbes. D'après
Compléments sur les morphismes d'espaces, lemme
\ref{spaces-more-morphisms-lemma-flat-finite-presentation-CM-open},
l'ensemble
$$
W = \{x \in |X| : f \text{ est Cohen-Macaulay au point }x\}
$$
est ouvert dans $|X|$, et la formation de cet ouvert commute à tout
changement de base. Puisque $f$ est propre, le sous-ensemble
$$
S' = S \setminus f(|X| \setminus W)
$$
de $S$ est ouvert, et $X \times_S S' \to S'$ est Cohen-Macaulay.
De plus, la formation de $S'$ commute à tout changement de base,
car il en est ainsi de $W$. La méthode exposée à la section
\ref{moduli-curves-section-open} fournit donc le sous-champ ouvert possédant les
propriétés voulues.
\end{proof}

\begin{lemma}
\label{moduli-curves-lemma-CM-1-curves}
Il existe un sous-champ ouvert $\Curvesstack^{CM, 1} \subset \Curvesstack$
tel que :
\begin{enumerate}
\item pour une famille de courbes $X \to S$, les assertions suivantes
sont équivalentes :
\begin{enumerate}
\item le morphisme classifiant $S \to \Curvesstack$ se factorise par
$\Curvesstack^{CM, 1}$,
\item le morphisme $X \to S$ est Cohen-Macaulay et de dimension
relative $1$ (Morphismes d'espaces, définition
\ref{spaces-morphisms-definition-relative-dimension}) ;
\end{enumerate}
\item pour un schéma $X$ propre sur un corps $k$ et tel que $\dim(X) \leq 1$,
les assertions suivantes sont équivalentes :
\begin{enumerate}
\item le morphisme classifiant $\Spec(k) \to \Curvesstack$ se factorise
par $\Curvesstack^{CM, 1}$,
\item $X$ est Cohen-Macaulay et $X$ est équidimensionnel de dimension $1$.
\end{enumerate}
\end{enumerate}
\end{lemma}

\begin{proof}
D'après le lemme \ref{moduli-curves-lemma-CM-curves}, il est clair que l'on a
$\Curvesstack^{CM, 1} \subset \Curvesstack^{CM}$ si ce sous-champ existe.
Soit $f : X \to S$ une famille de courbes telle que $f$ soit
Cohen-Macaulay. D'après Compléments sur les morphismes d'espaces, lemme
\ref{spaces-more-morphisms-lemma-lfp-CM-relative-dimension},
nous avons une décomposition
$$
X = X_0 \amalg X_1
$$
en sous-espaces ouverts et fermés telle que $X_0 \to S$ soit de dimension
relative $0$ et $X_1 \to S$ de dimension relative $1$.
Puisque $f$ est propre, le sous-ensemble
$$
S' = S \setminus f(|X_0|)
$$
de $S$ est ouvert, et $X \times_S S' \to S'$ est Cohen-Macaulay
et de dimension relative $1$.
De plus, la formation de $S'$ commute à tout changement de base,
car il en est ainsi de la décomposition précédente (la dimension relative
se comportant bien par changement de base ; voir Morphismes d'espaces, lemme
\ref{spaces-morphisms-lemma-dimension-fibre-after-base-change}).
La méthode exposée à la section \ref{moduli-curves-section-open} fournit donc le sous-champ
ouvert possédant les propriétés voulues.
\end{proof}





\section{Courbes de genre donné}
\label{moduli-curves-section-genus}

\noindent
Dans le projet Champs, on convient que le genre $g$ d'un
schéma $X$ propre de dimension $1$ sur un corps $k$ n'est défini que
si $H^0(X, \mathcal{O}_X) = k$. Dans ce cas,
$g = \dim_k H^1(X, \mathcal{O}_X)$.
Voir Courbes algébriques, section \ref{curves-section-genus}.
Les conditions nécessaires pour définir le genre définissent un sous-champ ouvert,
qui est alors réunion disjointe de sous-champs ouverts, un pour chaque genre.

\begin{lemma}
\label{moduli-curves-lemma-pre-genus-curves}
Il existe un sous-champ ouvert $\Curvesstack^{h0, 1} \subset \Curvesstack$
tel que :
\begin{enumerate}
\item pour une famille de courbes $f : X \to S$, les assertions suivantes
sont équivalentes :
\begin{enumerate}
\item le morphisme classifiant $S \to \Curvesstack$ se factorise par
$\Curvesstack^{h0, 1}$,
\item $f_*\mathcal{O}_X = \mathcal{O}_S$, cette égalité reste vraie
après tout changement de base, et les fibres de $f$ sont de dimension $1$ ;
\end{enumerate}
\item pour un schéma $X$ propre sur un corps $k$ et tel que $\dim(X) \leq 1$,
les assertions suivantes sont équivalentes :
\begin{enumerate}
\item le morphisme classifiant $\Spec(k) \to \Curvesstack$ se factorise par
$\Curvesstack^{h0, 1}$,
\item $H^0(X, \mathcal{O}_X) = k$ et $\dim(X) = 1$.
\end{enumerate}
\end{enumerate}
\end{lemma}

\begin{proof}
Pour une famille de courbes $X \to S$, l'ensemble des $s \in S$ tels que
$\kappa(s) = H^0(X_s, \mathcal{O}_{X_s})$
est ouvert dans $S$ d'après Catégories dérivées des espaces, lemme
\ref{spaces-perfect-lemma-jump-loci-geometric}.
De même, l'ensemble des points de $S$ où la fibre est de
dimension $1$ est ouvert d'après Compléments sur les morphismes d'espaces, lemme
\ref{spaces-more-morphisms-lemma-dimension-fibres-proper-flat}.
En outre, si $f : X \to S$ est une famille de courbes dont toutes les fibres
sont de dimension $1$ (et, en particulier, $f$ est surjectif), alors
la condition (1)(b) équivaut à
$\kappa(s) = H^0(X_s, \mathcal{O}_{X_s})$ pour tout $s \in S$ ; voir
Catégories dérivées des espaces, lemme \ref{spaces-perfect-lemma-proper-flat-h0}.
Le lemme résulte donc de la discussion générale de la
section \ref{moduli-curves-section-open}.
\end{proof}

\begin{lemma}
\label{moduli-curves-lemma-pre-genus-in-CM-1}
On a $\Curvesstack^{h0, 1} \subset \Curvesstack^{CM, 1}$
en tant que sous-champs ouverts de $\Curvesstack$.
\end{lemma}

\begin{proof}
Voir Courbes algébriques, lemme \ref{curves-lemma-automatic}, ainsi que les
lemmes \ref{moduli-curves-lemma-pre-genus-curves} et \ref{moduli-curves-lemma-CM-1-curves}.
\end{proof}

\begin{lemma}
\label{moduli-curves-lemma-genus}
Soit $f : X \to S$ une famille de courbes telle que
$\kappa(s) = H^0(X_s, \mathcal{O}_{X_s})$ pour tout $s \in S$, c'est-à-dire que
le morphisme classifiant $S \to \Curvesstack$ se factorise par
$\Curvesstack^{h0, 1}$ (lemme \ref{moduli-curves-lemma-pre-genus-curves}). Alors :
\begin{enumerate}
\item $f_*\mathcal{O}_X = \mathcal{O}_S$, et cette égalité subsiste après tout changement de base ;
\item $R^1f_*\mathcal{O}_X$ est un $\mathcal{O}_S$-module fini localement libre ;
\item pour tout morphisme $h : S' \to S$, si $f' : X' \to S'$ est obtenu par
changement de base, alors $h^*(R^1f_*\mathcal{O}_X) = R^1f'_*\mathcal{O}_{X'}$.
\end{enumerate}
\end{lemma}

\begin{proof}
Appliquons Catégories dérivées des espaces, lemme
\ref{spaces-perfect-lemma-proper-flat-h0}.
Cela prouve l'assertion (1). Cela implique aussi que, localement sur $S$,
on peut écrire $Rf_*\mathcal{O}_X = \mathcal{O}_S \oplus P$,
où $P$ est parfait et de Tor-amplitude contenue dans $[1, \infty)$.
Rappelons que la formation de $Rf_*\mathcal{O}_X$ commute
à tout changement de base
(Catégories dérivées des espaces, lemme
\ref{spaces-perfect-lemma-flat-proper-perfect-direct-image-general}).
Ainsi, pour $s \in S$, on a
$$
H^i(P \otimes_{\mathcal{O}_S}^\mathbf{L} \kappa(s)) =
H^i(X_s, \mathcal{O}_{X_s})
\text{ pour }i \geq 1
$$
Ce groupe est nul sauf si $i = 1$, car $X_s$ est un schéma noethérien
de dimension $1$ ; voir
Cohomologie, proposition \ref{cohomology-proposition-vanishing-Noetherian}.
On a alors $P = H^1(P)[-1]$, et $H^1(P)$ est fini localement libre,
par exemple d'après Compléments sur l'algèbre, lemme
\ref{more-algebra-lemma-lift-perfect-from-residue-field}.
Comme tout est compatible avec le changement de base, on voit
également que l'assertion (3) est vraie.
\end{proof}

\begin{lemma}
\label{moduli-curves-lemma-pre-genus-one-piece-per-genus}
Il existe une décomposition en sous-champs ouverts et fermés
$$
\Curvesstack^{h0, 1} = \coprod\nolimits_{g \geq 0} \Curvesstack_g
$$
où chaque $\Curvesstack_g$ se caractérise comme suit :
\begin{enumerate}
\item pour une famille de courbes $f : X \to S$, les assertions suivantes
sont équivalentes :
\begin{enumerate}
\item le morphisme classifiant $S \to \Curvesstack$ se factorise par
$\Curvesstack_g$,
\item $f_*\mathcal{O}_X = \mathcal{O}_S$, cette égalité reste vraie après
tout changement de base, les fibres de $f$ sont de dimension $1$, et
$R^1f_*\mathcal{O}_X$ est un $\mathcal{O}_S$-module localement libre de rang $g$ ;
\end{enumerate}
\item pour un schéma $X$ propre sur un corps $k$ et tel que $\dim(X) \leq 1$,
les assertions suivantes sont équivalentes :
\begin{enumerate}
\item le morphisme classifiant $\Spec(k) \to \Curvesstack$ se factorise par
$\Curvesstack_g$,
\item $\dim(X) = 1$, $k = H^0(X, \mathcal{O}_X)$ et
le genre de $X$ est $g$.
\end{enumerate}
\end{enumerate}
\end{lemma}

\begin{proof}
Nous disposons déjà de $\Curvesstack^{h0, 1}$ comme sous-champ ouvert de
$\Curvesstack$, caractérisé par les conditions du lemme qui ne font intervenir
ni $R^1f_*$ ni $H^1$ ; voir le lemme \ref{moduli-curves-lemma-pre-genus-curves}.
L'existence de la décomposition en sous-champs ouverts et fermés
résulte immédiatement de la discussion de la section \ref{moduli-curves-section-open}
et du lemme \ref{moduli-curves-lemma-genus}. Cela prouve la caractérisation de (1).
La caractérisation de (2) résulte de la définition du
genre donnée dans Courbes algébriques, définition \ref{curves-definition-genus}.
\end{proof}









\section{Courbes géométriquement réduites}
\label{moduli-curves-section-geometrically-reduced}

\noindent
Il existe un sous-champ ouvert de $\Curvesstack$ qui paramètre
les « courbes » géométriquement réduites.

\begin{lemma}
\label{moduli-curves-lemma-geometrically-reduced-curves}
Il existe un sous-champ ouvert $\Curvesstack^{geomred} \subset \Curvesstack$
tel que :
\begin{enumerate}
\item pour une famille de courbes $X \to S$, les assertions suivantes
sont équivalentes :
\begin{enumerate}
\item le morphisme classifiant $S \to \Curvesstack$ se factorise par
$\Curvesstack^{geomred}$,
\item les fibres du morphisme $X \to S$ sont géométriquement réduites
(Compléments sur les morphismes d'espaces, définition
\ref{spaces-more-morphisms-definition-geometrically-reduced-fibre}) ;
\end{enumerate}
\item pour un schéma $X$ propre sur un corps $k$ et tel que $\dim(X) \leq 1$,
les assertions suivantes sont équivalentes :
\begin{enumerate}
\item le morphisme classifiant $\Spec(k) \to \Curvesstack$ se factorise par
$\Curvesstack^{geomred}$,
\item $X$ est géométriquement réduit sur $k$.
\end{enumerate}
\end{enumerate}
\end{lemma}

\begin{proof}
Soit $f : X \to S$ une famille de courbes. D'après
Compléments sur les morphismes d'espaces, lemme
\ref{spaces-more-morphisms-lemma-geometrically-reduced-open},
l'ensemble
$$
E = \{s \in S : \text{la fibre de }X \to S\text{ en }s
\text{ est géométriquement réduite}\}
$$
est ouvert dans $S$. La formation de cet ouvert commute à tout
changement de base d'après
Compléments sur les morphismes d'espaces, lemme
\ref{spaces-more-morphisms-lemma-base-change-fibres-geometrically-reduced}.
La méthode exposée à la section \ref{moduli-curves-section-open} fournit donc le sous-champ
ouvert possédant les propriétés voulues.
\end{proof}

\begin{lemma}
\label{moduli-curves-lemma-geomred-in-CM}
On a $\Curvesstack^{geomred} \subset \Curvesstack^{CM}$
en tant que sous-champs ouverts de $\Curvesstack$.
\end{lemma}

\begin{proof}
Cela tient à ce qu'un schéma noethérien réduit de
dimension $\leq 1$ est Cohen-Macaulay. Voir
Algèbre, lemme \ref{algebra-lemma-criterion-reduced}.
\end{proof}






\section{Courbes géométriquement réduites et connexes}
\label{moduli-curves-section-geometrically-reduced-connected}

\noindent
Il existe un sous-champ ouvert de $\Curvesstack$ qui paramètre
les « courbes » géométriquement réduites et connexes.
Nous éliminons d'emblée les objets de dimension $0$.

\begin{lemma}
\label{moduli-curves-lemma-geometrically-reduced-connected-1-curves}
Il existe un sous-champ ouvert $\Curvesstack^{grc, 1} \subset \Curvesstack$
tel que :
\begin{enumerate}
\item pour une famille de courbes $X \to S$, les assertions suivantes
sont équivalentes :
\begin{enumerate}
\item le morphisme classifiant $S \to \Curvesstack$ se factorise par
$\Curvesstack^{grc, 1}$,
\item les fibres géométriques du morphisme $X \to S$ sont
réduites, connexes et de dimension $1$ ;
\end{enumerate}
\item pour un schéma $X$ propre sur un corps $k$ et tel que $\dim(X) \leq 1$,
les assertions suivantes sont équivalentes :
\begin{enumerate}
\item le morphisme classifiant $\Spec(k) \to \Curvesstack$ se factorise par
$\Curvesstack^{grc, 1}$,
\item $X$ est géométriquement réduit, géométriquement connexe
et de dimension $1$.
\end{enumerate}
\end{enumerate}
\end{lemma}

\begin{proof}
D'après les lemmes \ref{moduli-curves-lemma-geometrically-reduced-curves},
\ref{moduli-curves-lemma-geomred-in-CM}, \ref{moduli-curves-lemma-CM-curves} et \ref{moduli-curves-lemma-CM-1-curves},
il est clair que l'on a
$$
\Curvesstack^{grc, 1}
\subset
\Curvesstack^{geomred} \cap \Curvesstack^{CM, 1}
$$
si ce sous-champ existe. Soit $f : X \to S$ une famille de courbes telle que $f$
soit Cohen-Macaulay, ait des fibres géométriquement réduites et
soit de dimension relative $1$. D'après
Compléments sur les morphismes d'espaces, lemme
\ref{spaces-more-morphisms-lemma-stein-factorization-etale},
dans la factorisation de Stein
$$
X \to T \to S
$$
le morphisme $T \to S$ est \'etale. Il en résulte qu'il existe
un sous-schéma ouvert et fermé $S' \subset S$ tel que
$X \times_S S' \to S'$ ait des fibres géométriquement
connexes (dans la décomposition de
Morphismes, lemme \ref{morphisms-lemma-finite-locally-free},
appliquée au morphisme fini localement libre $T \to S$,
cela correspond à $S_1$).
La formation de cet ouvert commute à tout changement de base,
car le nombre de composantes connexes des fibres géométriques
est invariant par changement de base (il est également vrai
que la factorisation de Stein commute au changement de base
dans notre cas particulier, mais nous n'en avons pas besoin pour conclure).
La méthode exposée à la section \ref{moduli-curves-section-open} fournit donc le sous-champ
ouvert possédant les propriétés voulues.
\end{proof}

\begin{lemma}
\label{moduli-curves-lemma-geomredcon-in-h0-1}
On a $\Curvesstack^{grc, 1} \subset \Curvesstack^{h0, 1}$
en tant que sous-champs ouverts de $\Curvesstack$. En particulier, pour toute
famille de courbes $f : X \to S$
dont les fibres géométriques sont réduites, connexes et de dimension $1$,
$R^1f_*\mathcal{O}_X$ est un $\mathcal{O}_S$-module fini localement libre
dont la formation commute à tout changement de base.
\end{lemma}

\begin{proof}
Cela résulte de Variétés, lemme
\ref{varieties-lemma-proper-geometrically-reduced-global-sections},
ainsi que des lemmes \ref{moduli-curves-lemma-pre-genus-curves} et
\ref{moduli-curves-lemma-geometrically-reduced-connected-1-curves}.
La dernière assertion résulte du lemme \ref{moduli-curves-lemma-genus}.
\end{proof}

\begin{lemma}
\label{moduli-curves-lemma-one-piece-per-genus}
Il existe une décomposition en sous-champs ouverts et fermés
$$
\Curvesstack^{grc, 1} = \coprod\nolimits_{g \geq 0} \Curvesstack^{grc, 1}_g
$$
où chaque $\Curvesstack^{grc, 1}_g$ se caractérise comme suit :
\begin{enumerate}
\item pour une famille de courbes $f : X \to S$, les assertions suivantes
sont équivalentes :
\begin{enumerate}
\item le morphisme classifiant $S \to \Curvesstack$ se factorise par
$\Curvesstack^{grc, 1}_g$,
\item les fibres géométriques du morphisme $f : X \to S$ sont
réduites, connexes, de dimension $1$, et
$R^1f_*\mathcal{O}_X$ est un $\mathcal{O}_S$-module localement libre
de rang $g$ ;
\end{enumerate}
\item pour un schéma $X$ propre sur un corps $k$ et tel que $\dim(X) \leq 1$,
les assertions suivantes sont équivalentes :
\begin{enumerate}
\item le morphisme classifiant $\Spec(k) \to \Curvesstack$ se factorise par
$\Curvesstack^{grc, 1}_g$,
\item $X$ est géométriquement réduit, géométriquement connexe,
de dimension $1$ et de genre $g$.
\end{enumerate}
\end{enumerate}
\end{lemma}

\begin{proof}
Première démonstration : posons
$\Curvesstack^{grc, 1}_g = \Curvesstack^{grc, 1} \cap \Curvesstack_g$
et combinons les lemmes \ref{moduli-curves-lemma-geomredcon-in-h0-1} et
\ref{moduli-curves-lemma-pre-genus-one-piece-per-genus}.
Seconde démonstration :
L'existence de la décomposition en sous-champs ouverts et fermés
résulte immédiatement de la discussion de la section \ref{moduli-curves-section-open}
et du lemme \ref{moduli-curves-lemma-geomredcon-in-h0-1}.
Cela prouve la caractérisation de (1).
La caractérisation de (2) en résulte également, car
le genre d'un schéma $X/k$ propre de dimension $1$,
géométriquement réduit et connexe, est défini
(Courbes algébriques, définition \ref{curves-definition-genus}, et
Variétés, lemme
\ref{varieties-lemma-proper-geometrically-reduced-global-sections})
et vaut $\dim_k H^1(X, \mathcal{O}_X)$.
\end{proof}


 
 
 
 
 
\section{Courbes de Gorenstein}
\label{moduli-curves-section-gorenstein}

\noindent
Il existe un sous-champ ouvert de $\Curvesstack$ qui paramètre
les « courbes » de Gorenstein.

\begin{lemma}
\label{moduli-curves-lemma-gorenstein-curves}
Il existe un sous-champ ouvert $\Curvesstack^{Gorenstein} \subset \Curvesstack$
tel que :
\begin{enumerate}
\item pour une famille de courbes $X \to S$, les assertions suivantes sont équivalentes :
\begin{enumerate}
\item le morphisme classifiant $S \to \Curvesstack$ se factorise
par $\Curvesstack^{Gorenstein}$,
\item le morphisme $X \to S$ est de Gorenstein ;
\end{enumerate}
\item pour un schéma $X$ propre sur un corps $k$ et tel que $\dim(X) \leq 1$,
les assertions suivantes sont équivalentes :
\begin{enumerate}
\item le morphisme classifiant $\Spec(k) \to \Curvesstack$ se factorise
par $\Curvesstack^{Gorenstein}$,
\item $X$ est de Gorenstein.
\end{enumerate}
\end{enumerate}
\end{lemma}

\begin{proof}
Soit $f : X \to S$ une famille de courbes. D'après
Compléments sur les morphismes d'espaces, lemme
\ref{spaces-more-morphisms-lemma-flat-finite-presentation-gorenstein-open},
l'ensemble
$$
W = \{x \in |X| : f \text{ est de Gorenstein au point }x\}
$$
est ouvert dans $|X|$, et la formation de cet ouvert commute à tout
changement de base. Puisque $f$ est propre, le sous-ensemble
$$
S' = S \setminus f(|X| \setminus W)
$$
de $S$ est ouvert, et $X \times_S S' \to S'$ est de Gorenstein.
De plus, la formation de $S'$ commute à tout changement de base,
car il en est ainsi de $W$.
La méthode exposée à la section \ref{moduli-curves-section-open} fournit donc le
sous-champ ouvert possédant les propriétés voulues.
\end{proof}

\begin{lemma}
\label{moduli-curves-lemma-gorenstein-1-curves}
Il existe un sous-champ ouvert
$\Curvesstack^{Gorenstein, 1} \subset \Curvesstack$ tel que :
\begin{enumerate}
\item pour une famille de courbes $X \to S$, les assertions suivantes sont équivalentes :
\begin{enumerate}
\item le morphisme classifiant $S \to \Curvesstack$ se factorise
par $\Curvesstack^{Gorenstein, 1}$,
\item le morphisme $X \to S$ est de Gorenstein et de
dimension relative $1$ (Morphismes d'espaces, définition
\ref{spaces-morphisms-definition-relative-dimension}) ;
\end{enumerate}
\item pour un schéma $X$ propre sur un corps $k$ et tel que $\dim(X) \leq 1$,
les assertions suivantes sont équivalentes :
\begin{enumerate}
\item le morphisme classifiant $\Spec(k) \to \Curvesstack$ se factorise
par $\Curvesstack^{Gorenstein, 1}$,
\item $X$ est de Gorenstein et $X$ est équidimensionnel de
dimension $1$.
\end{enumerate}
\end{enumerate}
\end{lemma}

\begin{proof}
Rappelons qu'un schéma de Gorenstein est Cohen-Macaulay
(Dualité pour les schémas, lemme \ref{duality-lemma-gorenstein-CM})
et rappelons
qu'un morphisme de Gorenstein est un morphisme de Cohen-Macaulay
(Dualité pour les schémas, lemme \ref{duality-lemma-gorenstein-CM-morphism}).
Nous pouvons donc prendre
$\Curvesstack^{Gorenstein, 1}$ égale à l'intersection
de $\Curvesstack^{Gorenstein}$ et de $\Curvesstack^{CM, 1}$
dans $\Curvesstack$, et appliquer les
lemmes \ref{moduli-curves-lemma-gorenstein-curves} et \ref{moduli-curves-lemma-CM-1-curves}.
\end{proof}






\section{Courbes d'intersection complète locale}
\label{moduli-curves-section-lci}

\noindent
Il existe un sous-champ ouvert de $\Curvesstack$ qui paramètre
les « courbes » qui sont des intersections complètes locales.

\begin{lemma}
\label{moduli-curves-lemma-lci-curves}
Il existe un sous-champ ouvert $\Curvesstack^{lci} \subset \Curvesstack$
tel que :
\begin{enumerate}
\item pour une famille de courbes $X \to S$, les assertions suivantes sont équivalentes :
\begin{enumerate}
\item le morphisme classifiant $S \to \Curvesstack$ se factorise par
$\Curvesstack^{lci}$,
\item $X \to S$ est un morphisme d'intersection complète locale, et
\item $X \to S$ est un morphisme syntomique.
\end{enumerate}
\item pour un schéma propre $X$ sur un corps $k$, de dimension $\leq 1$,
les assertions suivantes sont équivalentes :
\begin{enumerate}
\item le morphisme classifiant $\Spec(k) \to \Curvesstack$ se factorise
par $\Curvesstack^{lci}$,
\item $X$ est une intersection complète locale sur $k$.
\end{enumerate}
\end{enumerate}
\end{lemma}

\begin{proof}
Rappelons qu'être un morphisme syntomique revient à être plat et à être
un morphisme d'intersection complète locale ; voir
Compléments sur les morphismes d'espaces, lemme \ref{spaces-more-morphisms-lemma-flat-lci}.
Ainsi, (1)(b) est équivalente à (1)(c).
À la section \ref{moduli-curves-section-open}, nous avons vu
qu'il suffit de montrer que, pour toute famille de courbes
$f : X \to S$, il existe un sous-schéma ouvert $S' \subset S$
tel que $S' \times_S X \to S'$ soit un morphisme d'intersection complète
locale et que la formation de $S'$ commute à tout changement de base.
Cela découle du résultat plus général
Compléments sur les morphismes d'espaces, lemme \ref{spaces-more-morphisms-lemma-where-lci}.
\end{proof}




\section{Courbes à singularités isolées}
\label{moduli-curves-section-curves-isolated}

\noindent
Nous pouvons considérer le sous-champ ouvert de $\Curvesstack$
qui paramètre les « courbes » n'ayant qu'un nombre fini de points
singuliers (ceux-ci peuvent correspondre à des composantes de dimension $0$
dans notre cadre).

\begin{lemma}
\label{moduli-curves-lemma-isolated-sings-curves}
Il existe un sous-champ ouvert
$\Curvesstack^{+} \subset \Curvesstack$
tel que :
\begin{enumerate}
\item pour une famille de courbes $X \to S$, les assertions suivantes sont équivalentes :
\begin{enumerate}
\item le morphisme classifiant $S \to \Curvesstack$ se factorise par
$\Curvesstack^{+}$,
\item le lieu singulier de $X \to S$, muni
d'une quelconque/certaine structure de sous-espace fermé, est fini sur $S$.
\end{enumerate}
\item pour un schéma propre $X$ sur un corps $k$, de dimension $\leq 1$,
les assertions suivantes sont équivalentes :
\begin{enumerate}
\item le morphisme classifiant $\Spec(k) \to \Curvesstack$ se factorise
par $\Curvesstack^{+}$,
\item $X \to \Spec(k)$ est lisse sauf en un nombre fini de points.
\end{enumerate}
\end{enumerate}
\end{lemma}

\begin{proof}
Pour démontrer le lemme, il suffit de montrer que, pour toute famille de courbes
$f : X \to S$, il existe un sous-schéma ouvert $S' \subset S$
tel que la fibre de $S' \times_S X \to S'$ possède la propriété (2).
(La formation de l'ouvert commutera alors automatiquement au changement de base.)
Par définition, le lieu $T \subset |X|$ des points où $X \to S$
n'est pas lisse est fermé. Soit $Z \subset X$ le sous-espace fermé
défini en munissant $T$ de la structure réduite induite d'espace algébrique
(Propriétés des espaces, définition
\ref{spaces-properties-definition-reduced-induced-space}).
Si maintenant $s \in S$ est un point tel que $Z_s$ soit fini, il existe
un voisinage ouvert $U_s \subset S$ de $s$ tel que
$Z \cap f^{-1}(U_s) \to U_s$ soit fini ; voir
Compléments sur les morphismes d'espaces, lemme
\ref{spaces-more-morphisms-lemma-proper-finite-fibre-finite-in-neighbourhood}.
Cela démontre le lemme.
\end{proof}




\section{Le lieu lisse du champ des courbes}
\label{moduli-curves-section-smooth}

\noindent
Le morphisme
$$
\Curvesstack \longrightarrow \Spec(\mathbf{Z})
$$
est lisse au-dessus d'un sous-champ ouvert maximal ; voir
Morphismes de champs, lemme \ref{stacks-morphisms-lemma-where-smooth}.
Nous voulons donner un critère pour qu'une courbe appartienne à ce lieu.
Pour cela, nous utiliserons un peu de théorie des déformations.

\medskip\noindent
Soit $k$ un corps. Soit $X$ un schéma propre de dimension $\leq 1$ sur $k$.
Choisissons un anneau de Cohen $\Lambda$ pour $k$ ; voir
Algèbre, lemme \ref{algebra-lemma-cohen-rings-exist}.
Nous sommes alors dans la situation décrite dans
Problèmes de déformation, exemple \ref{examples-defos-example-schemes} et
lemme \ref{examples-defos-lemma-schemes-RS}.
Nous obtenons ainsi une catégorie de déformations $\Deformationcategory_X$
sur la catégorie $\mathcal{C}_\Lambda$ des $\Lambda$-algèbres artiniennes
locales de corps résiduel $k$.

\begin{lemma}
\label{moduli-curves-lemma-in-smooth-locus}
Dans la situation ci-dessus, les assertions suivantes sont équivalentes :
\begin{enumerate}
\item le morphisme classifiant $\Spec(k) \to \Curvesstack$ se factorise
par l'ouvert où $\Curvesstack \to \Spec(\mathbf{Z})$ est lisse,
\item la catégorie de déformations $\Deformationcategory_X$ est sans obstruction.
\end{enumerate}
\end{lemma}

\begin{proof}
Puisque $\Curvesstack \longrightarrow \Spec(\mathbf{Z})$ est localement
de présentation finie (lemme \ref{moduli-curves-lemma-curves-qs-lfp}),
la formation du sous-champ ouvert où
$\Curvesstack \longrightarrow \Spec(\mathbf{Z})$ est lisse commute au changement de base
plat
(Morphismes de champs, lemme \ref{stacks-morphisms-lemma-where-smooth}).
Puisque l'anneau de Cohen $\Lambda$ est plat sur $\mathbf{Z}$,
nous pouvons travailler sur $\Lambda$. Autrement dit, il s'agit de montrer que
$$
\Lambda\text{-}\Curvesstack \longrightarrow \Spec(\Lambda)
$$
est lisse dans un voisinage ouvert du point
$x_0 : \Spec(k) \to \Lambda\text{-}\Curvesstack$
défini par $X/k$ si et seulement si $\Deformationcategory_X$ est sans obstruction.

\medskip\noindent
Le lemme résulte maintenant de
Géométrie des champs, lemme \ref{stacks-geometry-lemma-characterize-smoothness},
et de l'égalité
$$
\Deformationcategory_X =
\mathcal{F}_{\Lambda\text{-}\Curvesstack, k, x_0}
$$
Cette égalité n'est pas tout à fait triviale à établir. En effet, dans le
membre de gauche, la catégorie de déformations classifie toutes les déformations plates
$Y \to \Spec(A)$ de $X$ comme schéma sur $A \in \Ob(\mathcal{C}_\Lambda)$.
Dans le membre de droite, la catégorie de déformations classifie tous les
morphismes plats $Y \to \Spec(A)$ de fibre spéciale $X$,
où $Y$ est un espace algébrique et où
$Y \to \Spec(A)$ est propre, de présentation finie et de
dimension relative $\leq 1$. Puisque $A$ est artinien, nous voyons
que $Y$ est un schéma, par exemple d'après Espaces sur les corps, lemme
\ref{spaces-over-fields-lemma-codim-1-point-in-schematic-locus}.
Il reste donc à montrer ceci : une déformation plate $Y \to \Spec(A)$ de
$X$ comme schéma sur un anneau local artinien $A$ de corps résiduel $k$
est propre, de présentation finie et de dimension relative $\leq 1$.
La dimension relative se définit en termes des fibres et cette propriété vaut donc
automatiquement pour $Y/A$ puisqu'elle vaut pour $X/k$.
Le morphisme $Y \to \Spec(A)$ est propre et localement de présentation finie,
car c'est vrai pour $X \to \Spec(k)$ ; voir
Compléments sur les morphismes, lemme \ref{more-morphisms-lemma-deform-property}.
\end{proof}

\noindent
Voici un « grand » ouvert du champ des courbes qui est contenu
dans le lieu lisse.

\begin{lemma}
\label{moduli-curves-lemma-big-smooth-part-curves}
Le sous-champ ouvert
$$
\Curvesstack^{lci+} =
\Curvesstack^{lci} \cap \Curvesstack^{+}
\subset \Curvesstack
$$
possède les propriétés suivantes :
\begin{enumerate}
\item $\Curvesstack^{lci+} \to \Spec(\mathbf{Z})$ est lisse,
\item pour une famille de courbes $X \to S$, les assertions suivantes sont équivalentes :
\begin{enumerate}
\item le morphisme classifiant $S \to \Curvesstack$ se factorise par
$\Curvesstack^{lci+}$,
\item $X \to S$ est un morphisme d'intersection complète locale et
le lieu singulier de $X \to S$, muni d'une quelconque/certaine structure
de sous-espace fermé, est fini sur $S$ ;
\end{enumerate}
\item pour un schéma propre $X$ sur un corps $k$, de dimension $\leq 1$,
les assertions suivantes sont équivalentes :
\begin{enumerate}
\item le morphisme classifiant $\Spec(k) \to \Curvesstack$ se factorise
par $\Curvesstack^{lci+}$,
\item $X$ est une intersection complète locale sur $k$ et
$X \to \Spec(k)$ est lisse sauf en un nombre fini de points.
\end{enumerate}
\end{enumerate}
\end{lemma}

\begin{proof}
Si nous pouvons montrer qu'il existe un sous-champ ouvert $\Curvesstack^{lci+}$
dont les points sont caractérisés par (2), nous voyons alors que
(1) résulte de la combinaison du lemme \ref{moduli-curves-lemma-in-smooth-locus} avec
Problèmes de déformation, lemme \ref{examples-defos-lemma-curve-isolated-lci}.
Puisque
$$
\Curvesstack^{lci+} = \Curvesstack^{lci} \cap \Curvesstack^{+}
$$
dans $\Curvesstack$, on conclut par les
lemmes \ref{moduli-curves-lemma-lci-curves} et \ref{moduli-curves-lemma-isolated-sings-curves}.
\end{proof}




\section{Courbes lisses}
\label{moduli-curves-section-smooth-curves}

\noindent
Dans cette section, nous étudions les sous-champs ouverts de $\Curvesstack$
qui paramètrent les ``courbes'' lisses.

\begin{lemma}
\label{moduli-curves-lemma-smooth-curves}
Il existe des sous-champs ouverts
$$
\Curvesstack^{smooth, 1} \subset \Curvesstack^{smooth} \subset \Curvesstack
$$
tels que
\begin{enumerate}
\item étant donnée une famille de courbes $f : X \to S$, les conditions suivantes sont équivalentes :
\begin{enumerate}
\item le morphisme classifiant $S \to \Curvesstack$ se factorise
par $\Curvesstack^{smooth}$, respectivement par $\Curvesstack^{smooth, 1}$,
\item $f$ est lisse, respectivement lisse de dimension relative $1$,
\end{enumerate}
\item étant donné un schéma $X$ propre sur un corps $k$ avec
$\dim(X) \leq 1$, les conditions suivantes sont équivalentes :
\begin{enumerate}
\item le morphisme classifiant $\Spec(k) \to \Curvesstack$ se factorise
par $\Curvesstack^{smooth}$, respectivement par $\Curvesstack^{smooth, 1}$,
\item $X$ est lisse sur $k$, respectivement $X$ est lisse sur $k$ et
$X$ est équidimensionnel de dimension $1$.
\end{enumerate}
\end{enumerate}
\end{lemma}

\begin{proof}
Pour démontrer les assertions concernant $\Curvesstack^{smooth}$,
il suffit de montrer qu'étant donnée une famille de courbes
$f : X \to S$, il existe un sous-schéma ouvert $S' \subset S$
tel que $S' \times_S X \to S'$ soit lisse et que la
formation de cet ouvert commute au changement de base.
Nous savons qu'il existe un plus grand ouvert $U \subset X$ tel
que $U \to S$ soit lisse et que la formation de $U$ commute
à tout changement de base, voir
Morphismes d'espaces, lemme \ref{spaces-morphisms-lemma-where-smooth}.
Si $T = |X| \setminus |U|$, alors $f(T)$ est fermé dans $S$ puisque $f$ est propre.
En posant $S' = S \setminus f(T)$, on obtient l'ouvert voulu.

\medskip\noindent
Soit $f : X \to S$ une famille de courbes telle que $f$ soit lisse.
Alors les fibres $X_s$ sont lisses sur $\kappa(s)$ et sont donc
de Cohen-Macaulay (on peut le voir, par exemple, en utilisant
Algèbre, lemmes \ref{algebra-lemma-smooth-over-field} et
\ref{algebra-lemma-lci-CM}). Ainsi, on voit que l'on peut poser
$$
\Curvesstack^{smooth, 1} = \Curvesstack^{smooth} \cap
\Curvesstack^{CM, 1}
$$
et les équivalences voulues résultent de ce qui a déjà été
montré pour $\Curvesstack^{smooth}$ et du lemme \ref{moduli-curves-lemma-CM-1-curves}.
\end{proof}

\begin{lemma}
\label{moduli-curves-lemma-smooth-curves-smooth}
Le morphisme $\Curvesstack^{smooth} \to \Spec(\mathbf{Z})$ est lisse.
\end{lemma}

\begin{proof}
Cela résulte immédiatement de l'observation que
$\Curvesstack^{smooth} \subset \Curvesstack^{lci+}$
et du lemme \ref{moduli-curves-lemma-big-smooth-part-curves}.
\end{proof}

\begin{lemma}
\label{moduli-curves-lemma-smooth-curves-h0}
Il existe un sous-champ ouvert
$\Curvesstack^{smooth, h0} \subset \Curvesstack$
tel que
\begin{enumerate}
\item étant donnée une famille de courbes $f : X \to S$, les conditions suivantes sont équivalentes :
\begin{enumerate}
\item le morphisme classifiant $S \to \Curvesstack$ se factorise
par $\Curvesstack^{smooth}$,
\item $f_*\mathcal{O}_X = \mathcal{O}_S$, cette égalité reste vraie après tout changement de base,
et $f$ est lisse de dimension relative $1$,
\end{enumerate}
\item étant donné un schéma $X$ propre sur un corps $k$ avec
$\dim(X) \leq 1$, les conditions suivantes sont équivalentes :
\begin{enumerate}
\item le morphisme classifiant $\Spec(k) \to \Curvesstack$ se factorise
par $\Curvesstack^{smooth, h0}$,
\item $X$ est lisse, $\dim(X) = 1$, et $k = H^0(X, \mathcal{O}_X)$,
\item $X$ est lisse, $\dim(X) = 1$, et $X$ est géométriquement connexe,
\item $X$ est lisse, $\dim(X) = 1$, et $X$ est géométriquement intègre, et
\item $X_{\overline{k}}$ est une courbe lisse.
\end{enumerate}
\end{enumerate}
\end{lemma}

\begin{proof}
Si l'on pose
$$
\Curvesstack^{smooth, h0} = \Curvesstack^{smooth} \cap
\Curvesstack^{h0, 1}
$$
on voit alors que (1) résulte des
lemmes \ref{moduli-curves-lemma-pre-genus-curves} et \ref{moduli-curves-lemma-smooth-curves}.
En fait, cela donne aussi l'équivalence de (2)(a) et (2)(b).
Pour achever la démonstration, il reste à montrer que
(2)(b) est équivalent à chacune des conditions (2)(c), (2)(d) et (2)(e).

\medskip\noindent
Un schéma lisse sur un corps est géométriquement normal
(Variétés, lemme \ref{varieties-lemma-smooth-geometrically-normal}),
la lissité est préservée par changement de base
(Morphismes, lemme \ref{morphisms-lemma-base-change-smooth}), et
la propriété d'être lisse est locale pour la topologie fpqc sur le but
(Descente, lemme \ref{descent-lemma-descending-property-smooth}).
Compte tenu de cela, l'équivalence de (2)(b), (2)(c), 2(d) et (2)(e)
résulte de Variétés, lemme \ref{varieties-lemma-geometrically-normal-stein}.
\end{proof}

\begin{definition}
\label{moduli-curves-definition-deligne-mumford-smooth}
\begin{reference}
\cite{DM}
\end{reference}
Nous notons $\mathcal{M}$ et appelons
{\it champ de modules des courbes propres et lisses}
le champ algébrique
$\Curvesstack^{smooth, h0}$ paramétrant les familles de courbes
introduit au lemme \ref{moduli-curves-lemma-smooth-curves-h0}.
Pour $g \geq 0$, nous notons $\mathcal{M}_g$ et appelons
{\it champ de modules des courbes propres et lisses de genre $g$}
le champ algébrique introduit au
lemme \ref{moduli-curves-lemma-smooth-one-piece-per-genus}.
\end{definition}

\noindent
Voici l'inévitable lemme.

\begin{lemma}
\label{moduli-curves-lemma-smooth-one-piece-per-genus}
Il existe une décomposition en sous-champs ouverts et fermés
$$
\mathcal{M} = \coprod\nolimits_{g \geq 0} \mathcal{M}_g
$$
où chaque $\mathcal{M}_g$ se caractérise comme suit :
\begin{enumerate}
\item étant donnée une famille de courbes $f : X \to S$, les conditions suivantes sont équivalentes :
\begin{enumerate}
\item le morphisme classifiant $S \to \Curvesstack$ se factorise
par $\mathcal{M}_g$,
\item $X \to S$ est lisse, $f_*\mathcal{O}_X = \mathcal{O}_S$,
cette égalité reste vraie après tout changement de base, et $R^1f_*\mathcal{O}_X$
est un $\mathcal{O}_S$-module localement libre de rang $g$,
\end{enumerate}
\item étant donné un schéma $X$ propre sur un corps $k$ avec
$\dim(X) \leq 1$, les conditions suivantes sont équivalentes :
\begin{enumerate}
\item le morphisme classifiant $\Spec(k) \to \Curvesstack$ se factorise
par $\mathcal{M}_g$,
\item $X$ est lisse, $\dim(X) = 1$, $k = H^0(X, \mathcal{O}_X)$,
et $X$ est de genre $g$,
\item $X$ est lisse, $\dim(X) = 1$, $X$ est géométriquement connexe, et
$X$ est de genre $g$,
\item $X$ est lisse, $\dim(X) = 1$, $X$ est géométriquement intègre, et
$X$ est de genre $g$, et
\item $X_{\overline{k}}$ est une courbe lisse de genre $g$.
\end{enumerate}
\end{enumerate}
\end{lemma}

\begin{proof}
Combiner les lemmes \ref{moduli-curves-lemma-smooth-curves-h0} et
\ref{moduli-curves-lemma-pre-genus-one-piece-per-genus}.
On peut aussi utiliser
le lemme \ref{moduli-curves-lemma-one-piece-per-genus}
à la place.
\end{proof}

\begin{lemma}
\label{moduli-curves-lemma-smooth-curves-h0-smooth}
Les morphismes $\mathcal{M} \to \Spec(\mathbf{Z})$ et
$\mathcal{M}_g \to \Spec(\mathbf{Z})$
sont lisses.
\end{lemma}

\begin{proof}
Comme $\mathcal{M}$ est un sous-champ ouvert de
$\Curvesstack^{lci+}$, cela résulte du
lemme \ref{moduli-curves-lemma-big-smooth-part-curves}.
\end{proof}





\section{Densité des courbes lisses}
\label{moduli-curves-section-smooth-is-dense}

\noindent
Le titre de cette section est trompeur, car nous n'affirmons pas que
$\Curvesstack^{smooth}$ est dense dans $\Curvesstack$.
En fait, cela est faux, comme Mumford l'a montré dans \cite{PathologiesIV}.
Cependant, nous verrons que les ``courbes'' lisses sont denses
dans un grand ouvert.

\begin{lemma}
\label{moduli-curves-lemma-smooth-dense}
L'inclusion
$$
|\Curvesstack^{smooth}| \subset |\Curvesstack^{lci+}|
$$
est celle d'un sous-ensemble ouvert dense.
\end{lemma}

\begin{proof}
Par la construction même de la topologie sur
$|\Curvesstack^{lci+}|$ dans
Propriétés des champs, section \ref{stacks-properties-section-points},
on constate que $|\Curvesstack^{smooth}|$
est un sous-ensemble ouvert. Soit $\xi \in |\Curvesstack^{lci+}|$ un point.
Il existe alors un corps $k$ et un schéma $X$ sur $k$
tel que $X$ soit propre sur $k$, que $\dim(X) \leq 1$,
que $X$ soit localement d'intersection complète sur $k$, et
que $X$ soit lisse sur $k$ sauf en un nombre fini de points, et
que $\xi$ soit la classe d'équivalence du
morphisme classifiant $\Spec(k) \to \Curvesstack^{lci+}$ déterminé par $X$.
Voir le lemme \ref{moduli-curves-lemma-big-smooth-part-curves}.
D'après Problèmes de déformation, lemme
\ref{examples-defos-lemma-smoothing-proper-curve-isolated-lci},
il existe un morphisme plat et projectif $Y \to \Spec(k[[t]])$
dont la fibre générique est lisse et la fibre spéciale est
isomorphe à $X$. Considérons le morphisme classifiant
$$
\Spec(k[[t]]) \longrightarrow \Curvesstack^{lci+}
$$
déterminé par $Y$. L'image du point fermé est $\xi$
et l'image du point générique appartient à $|\Curvesstack^{smooth}|$.
Comme le point générique se spécialise au point fermé dans
$|\Spec(k[[t]])|$, on conclut que $\xi$ appartient à l'adhérence
de $|\Curvesstack^{smooth}|$, comme voulu.
\end{proof}






\section{Courbes nodales}
\label{moduli-curves-section-nodal-curves}

\noindent
En géométrie algébrique, les courbes nodales jouent un rôle particulier.
Nous suggérons au lecteur de jeter un rapide coup d'œil à certaines des discussions
de Courbes algébriques, sections \ref{curves-section-nodal} et
\ref{curves-section-families-nodal},
ainsi que de Compléments sur les morphismes d'espaces, section
\ref{spaces-more-morphisms-section-families-nodal}.

\begin{lemma}
\label{moduli-curves-lemma-nodal-curves}
Il existe un sous-champ ouvert $\Curvesstack^{nodal} \subset \Curvesstack$
tel que
\begin{enumerate}
\item étant donnée une famille de courbes $f : X \to S$, les conditions suivantes sont équivalentes :
\begin{enumerate}
\item le morphisme classifiant $S \to \Curvesstack$ se factorise
par $\Curvesstack^{nodal}$,
\item $f$ est au plus nodal de dimension relative $1$,
\end{enumerate}
\item étant donné un schéma $X$ propre sur un corps $k$ avec
$\dim(X) \leq 1$, les conditions suivantes sont équivalentes :
\begin{enumerate}
\item le morphisme classifiant $\Spec(k) \to \Curvesstack$ se factorise
par $\Curvesstack^{nodal}$,
\item les singularités de $X$ sont au plus nodales et $X$
est équidimensionnel de dimension $1$.
\end{enumerate}
\end{enumerate}
\end{lemma}

\begin{proof}
En fait, il suffit de montrer qu'étant donnée une famille de courbes
$f : X \to S$, il existe un sous-schéma ouvert $S' \subset S$
tel que $S' \times_S X \to S'$ soit au plus nodal de dimension relative $1$
et que la formation de $S'$ commute à tout changement de base.
D'après Compléments sur les morphismes d'espaces, lemme
\ref{spaces-more-morphisms-lemma-locus-where-nodal},
il existe un plus grand sous-espace ouvert $X' \subset X$ tel
que $f|_{X'} : X' \to S$ soit au plus nodal de dimension relative $1$.
De plus, la formation de $X'$ commute au changement de base.
On peut donc prendre
$$
S' = S \setminus |f|(|X| \setminus |X'|)
$$
Cet ensemble est ouvert, car un morphisme propre est universellement fermé par
définition.
\end{proof}

\begin{lemma}
\label{moduli-curves-lemma-nodal-curves-smooth}
Le morphisme $\Curvesstack^{nodal} \to \Spec(\mathbf{Z})$ est lisse.
\end{lemma}

\begin{proof}
Cela résulte immédiatement de l'observation que
$\Curvesstack^{nodal} \subset \Curvesstack^{lci+}$
et du lemme \ref{moduli-curves-lemma-big-smooth-part-curves}.
\end{proof}




\section{Le faisceau dualisant relatif}
\label{moduli-curves-section-relative-dualizing}

\noindent
Cette section sert principalement
à introduire des notations dans le cas des familles de courbes.
La majeure partie du travail a déjà été faite dans le chapitre sur la dualité.

\medskip\noindent
Soit $f : X \to S$ une famille de courbes. Il existe un objet
$\omega_{X/S}^\bullet$ de $D_\QCoh(\mathcal{O}_X)$,
appelé le {\it complexe dualisant relatif}, qui possède la propriété
suivante : pour tout diagramme de changement de base
$$
\xymatrix{
X_U \ar[d]_{f'} \ar[r]_{g'} & X \ar[d]^f \\
U \ar[r]^g & S
}
$$
où $U = \Spec(A)$ est affine, le complexe
$\omega_{X_U/U}^\bullet = L(g')^*\omega_{X/S}^\bullet$
représente le foncteur
$$
D_\QCoh(\mathcal{O}_{X_U}) \longrightarrow \text{Mod}_A,\quad
K \longmapsto \Hom_U(Rf_*K, \mathcal{O}_U)
$$
Plus précisément, soit $(\omega_{X/S}^\bullet, \tau)$
le complexe dualisant relatif de la famille défini dans
Dualité pour les espaces, définition
\ref{spaces-duality-definition-relative-dualizing-proper-flat}.
Son existence est démontrée dans Dualité pour les espaces, lemme
\ref{spaces-duality-lemma-existence-relative-dualizing}.
De plus, la formation de $(\omega_{X/S}^\bullet, \tau)$ commute
à tout changement de base (essentiellement par définition ; une référence
précise est Dualité pour les espaces, lemme
\ref{spaces-duality-lemma-base-change-relative-dualizing}).
Dorénavant, nous identifierons le changement de base de
$\omega_{X/S}^\bullet$ au complexe dualisant relatif
de la famille obtenue par changement de base, sans plus le préciser.

\medskip\noindent
Soit $\{S_i \to S\}$ un recouvrement étale avec les $S_i$ affines tel que
$X_i = X \times_S S_i$ soit un schéma ; voir le lemme \ref{moduli-curves-lemma-etale-locally-scheme}.
D'après Dualité pour les espaces, lemme \ref{spaces-duality-lemma-compare},
on constate que $\omega_{X_i/S_i}^\bullet$ s'identifie
au complexe dualisant relatif du morphisme de schémas propre, plat et
de présentation finie $f_i : X_i \to S_i$ étudié dans
Dualité pour les schémas, remarque
\ref{duality-remark-relative-dualizing-complex}.
Ainsi, pour démontrer une propriété de $\omega_{X/S}^\bullet$
locale pour la topologie étale, nous pouvons supposer que $X \to S$ est un morphisme de schémas
et utiliser la théorie développée dans le chapitre sur la dualité pour les schémas.
Plus généralement, après tout changement de base de $X$ qui est un schéma,
le complexe dualisant relatif s'identifie au
complexe dualisant relatif de Dualité pour les schémas, remarque
\ref{duality-remark-relative-dualizing-complex}.
Dorénavant, nous utiliserons cette identification sans plus le préciser.

\medskip\noindent
En particulier, soit $\Spec(k) \to S$ un morphisme, où $k$ est un corps.
Notons $X_k$ le changement de base (c'est un schéma d'après
Espaces sur les corps, lemme
\ref{spaces-over-fields-lemma-codim-1-point-in-schematic-locus}).
Alors $\omega_{X_k/k}^\bullet$ est isomorphe
au complexe $\omega_{X_k}^\bullet$ de
Courbes algébriques, lemme \ref{curves-lemma-duality-dim-1}
(tous deux représentent le même foncteur, de sorte que l'on peut appliquer le lemme de Yoneda,
mais cela résulte en fait des remarques précédentes).
On en conclut que les faisceaux de cohomologie
$H^i(\omega_{X_k/k}^\bullet)$
ne sont non nuls que pour $i = 0, -1$.
Si $X_k$ est de Cohen-Macaulay et équidimensionnel de dimension $1$,
alors on n'a plus que $H^{-1}$ et, si $X_k$ est en outre de Gorenstein,
$H^{-1}(\omega_{X_k/k})$ est inversible ; voir
Courbes algébriques, lemmes \ref{curves-lemma-duality-dim-1-CM}
et \ref{curves-lemma-rr}.

\begin{lemma}
\label{moduli-curves-lemma-CM-dualizing}
Soit $X \to S$ une famille de courbes dont les fibres sont de Cohen-Macaulay
et équidimensionnelles de dimension $1$ (lemme \ref{moduli-curves-lemma-CM-1-curves}).
Alors $\omega_{X/S}^\bullet = \omega_{X/S}[1]$, où $\omega_{X/S}$
est un $\mathcal{O}_X$-module pseudo-cohérent et plat sur $S$, dont la
formation commute à tout changement de base.
\end{lemma}

\begin{proof}
Nous invitons vivement le lecteur à déduire cet énoncé directement de l'étude ci-dessus
du comportement après changement de base à un corps. Notre démonstration
emploiera une réduction quelque peu laborieuse au cas des schémas noethériens.

\medskip\noindent
Dès que nous aurons montré $\omega_{X/S}^\bullet = \omega_{X/S}[1]$ avec
$\omega_{X/S}$ plat sur $S$, l'assertion relative au changement de base
en découlera, puisque nous savons déjà que la formation de $\omega_{X/S}^\bullet$
commute à tout changement de base. De plus, la pseudo-cohérence
sera automatique, car $\omega_{X/S}^\bullet$ est pseudo-cohérent
par définition. L'annulation des autres faisceaux de cohomologie et la platitude
peuvent se vérifier localement pour la topologie étale. Nous pouvons donc supposer que $f : X \to S$
est un morphisme de schémas avec $S$ affine (voir la discussion précédente).
Écrivons $S = \lim S_i$ comme limite projective filtrante de schémas affines $S_i$
de type fini sur $\mathbf{Z}$.
Puisque $\Curvesstack^{CM, 1}$ est localement de présentation finie sur
$\mathbf{Z}$ (comme sous-champ ouvert de $\Curvesstack$ ; voir les
lemmes \ref{moduli-curves-lemma-CM-1-curves} et \ref{moduli-curves-lemma-curves-qs-lfp}),
on peut trouver un $i$ et une famille
de courbes $X_i \to S_i$ dont l'image réciproque est $X \to S$
(Limites de champs, lemme
\ref{stacks-limits-lemma-representable-by-spaces-limit-preserving}).
Quitte à augmenter $i$, on peut supposer que $X_i$ est un schéma ;
voir Limites d'espaces, lemme \ref{spaces-limits-lemma-limit-is-scheme}.
Puisque la formation de $\omega_{X/S}^\bullet$ commute à
tout changement de base, on peut remplacer $S$ par $S_i$.
Après ce remplacement, nous pouvons supposer, et nous le faisons, que $S_i$ est noethérien.
Alors $f$ est manifestement un morphisme de Cohen-Macaulay
(Compléments sur les morphismes, définition \ref{more-morphisms-definition-CM}),
en vertu de notre hypothèse sur les fibres.
De plus, $\omega_{X/S}^\bullet = f^!\mathcal{O}_S$
par la construction même de $f^!$ dans
Dualité pour les schémas, section \ref{duality-section-upper-shriek}.
Le lemme résulte donc de Dualité pour les schémas, lemme
\ref{duality-lemma-affine-flat-Noetherian-CM}.
\end{proof}

\begin{definition}
\label{moduli-curves-definition-relative-dualizing-sheaf}
Soit $f : X \to S$ une famille de courbes dont les fibres sont de Cohen-Macaulay
et équidimensionnelles de dimension $1$ (lemme \ref{moduli-curves-lemma-CM-1-curves}).
Le $\mathcal{O}_X$-module
$$
\omega_{X/S} = H^{-1}(\omega_{X/S}^\bullet)
$$
étudié dans le lemme \ref{moduli-curves-lemma-CM-dualizing}
est appelé le {\it faisceau dualisant relatif} de $f$.
\end{definition}

\noindent
Dans la situation de la définition \ref{moduli-curves-definition-relative-dualizing-sheaf},
le faisceau dualisant relatif $\omega_{X/S}$ possède la propriété suivante
(qui le caractérise en outre localement sur $S$) :
pour tout diagramme de changement de base
$$
\xymatrix{
X_U \ar[d]_{f'} \ar[r]_{g'} & X \ar[d]^f \\
U \ar[r]^g & S
}
$$
où $U = \Spec(A)$ est affine, le module $\omega_{X_U/U} = (g')^*\omega_{X/S}$
représente le foncteur
$$
\QCoh(\mathcal{O}_{X_U}) \longrightarrow \text{Mod}_A,\quad
\mathcal{F} \longmapsto \Hom_A(H^1(X, \mathcal{F}), A)
$$
Cela résulte immédiatement de la propriété correspondante du complexe
dualisant relatif donnée ci-dessus. En particulier, si $A = k$ est un corps,
on retrouve le module dualisant de $X_k$ introduit et étudié dans
Courbes algébriques, lemmes \ref{curves-lemma-duality-dim-1},
\ref{curves-lemma-duality-dim-1-CM} et \ref{curves-lemma-rr}.

\begin{lemma}
\label{moduli-curves-lemma-gorenstein-dualizing}
Soit $X \to S$ une famille de courbes dont les fibres sont de Gorenstein
et équidimensionnelles de dimension $1$ (lemme \ref{moduli-curves-lemma-gorenstein-1-curves}).
Alors le faisceau dualisant relatif $\omega_{X/S}$ est un
$\mathcal{O}_X$-module inversible dont la
formation commute à tout changement de base.
\end{lemma}

\begin{proof}
Cela résulte de ce que l'image réciproque du module dualisant relatif
sur une fibre est inversible, d'après la discussion précédente. On
peut aussi raisonner exactement comme dans la démonstration du
lemme \ref{moduli-curves-lemma-CM-dualizing} et déduire le résultat de
Dualité pour les schémas, lemme
\ref{duality-lemma-affine-flat-Noetherian-gorenstein}.
\end{proof}









\section{Courbes préstables}
\label{moduli-curves-section-prestable-curves}

\noindent
La définition suivante est équivalente à ce qui semble être la
notion généralement admise de famille préstable de courbes.

\begin{definition}
\label{moduli-curves-definition-prestable}
Soit $f : X \to S$ une famille de courbes. On dit que $f$ est une
{\it famille préstable de courbes} si
\begin{enumerate}
\item $f$ est au plus nodal de dimension relative $1$, et
\item $f_*\mathcal{O}_X = \mathcal{O}_S$, et cette égalité subsiste après
tout changement de base\footnote{En fait, il suffit d'exiger
$f_*\mathcal{O}_X = \mathcal{O}_S$, car la factorisation de Stein
de $f$ est étale dans ce cas ; voir
Compléments sur les morphismes d'espaces, lemme
\ref{spaces-more-morphisms-lemma-stein-factorization-etale}.
On peut aussi remplacer cette condition en demandant que les fibres
géométriques soient connexes ; voir le lemme \ref{moduli-curves-lemma-geomredcon-in-h0-1}.}.
\end{enumerate}
\end{definition}

\noindent
Soit $X$ un schéma propre sur un corps $k$ avec $\dim(X) \leq 1$.
Alors $X \to \Spec(k)$ est une famille de courbes ; on peut donc demander
si elle est préstable ou non\footnote{Nous ne pouvons pas employer le terme
``courbe préstable'' ici, car le mot courbe implique l'irréductibilité. Voir
la discussion dans Courbes algébriques, section \ref{curves-section-families-nodal}.}
au sens de la définition. En explicitant les définitions, on voit que
les conditions suivantes sont équivalentes :
\begin{enumerate}
\item $X$ est préstable,
\item les singularités de $X$ sont au plus nodales, $\dim(X) = 1$,
et $k = H^0(X, \mathcal{O}_X)$,
\item $X_{\overline{k}}$ est connexe et lisse sur $\overline{k}$
en dehors d'un nombre fini de nœuds
(Courbes algébriques, définition \ref{curves-definition-multicross}).
\end{enumerate}
Cela montre que notre définition concorde avec la plupart de celles que l'on trouve
dans la littérature.

\begin{lemma}
\label{moduli-curves-lemma-prestable-curves}
Il existe un sous-champ ouvert $\Curvesstack^{prestable} \subset \Curvesstack$
tel que
\begin{enumerate}
\item étant donnée une famille de courbes $f : X \to S$, les conditions suivantes sont équivalentes :
\begin{enumerate}
\item le morphisme classifiant $S \to \Curvesstack$ se factorise
par $\Curvesstack^{prestable}$,
\item $X \to S$ est une famille préstable de courbes,
\end{enumerate}
\item étant donné un schéma $X$ propre sur un corps $k$ avec
$\dim(X) \leq 1$, les conditions suivantes sont équivalentes :
\begin{enumerate}
\item le morphisme classifiant $\Spec(k) \to \Curvesstack$
se factorise par $\Curvesstack^{prestable}$,
\item les singularités de $X$ sont au plus nodales, $\dim(X) = 1$,
et $k = H^0(X, \mathcal{O}_X)$.
\end{enumerate}
\end{enumerate}
\end{lemma}

\begin{proof}
Étant donnée une famille de courbes $X \to S$, elle est préstable si
et seulement si le morphisme classifiant se factorise à la fois par
$\Curvesstack^{nodal}$ et $\Curvesstack^{h0, 1}$. On peut aussi
utiliser $\Curvesstack^{grc, 1}$ (puisqu'une courbe nodale est géométriquement
réduite et a donc son $H^0$ égal au corps de base si et seulement si
elle est connexe). En formule,
$$
\Curvesstack^{prestable} =
\Curvesstack^{nodal} \cap \Curvesstack^{h0, 1} =
\Curvesstack^{nodal} \cap \Curvesstack^{grc, 1}
$$
Le lemme résulte donc des
lemmes \ref{moduli-curves-lemma-pre-genus-curves} et \ref{moduli-curves-lemma-nodal-curves}.
\end{proof}

\noindent
Pour tout genre $g \geq 0$, nous disposons du champ algébrique classifiant
les courbes préstables de genre $g$. En fait, dorénavant, nous dirons
que $X \to S$ est une {\it famille préstable de courbes de genre $g$}
si et seulement si le morphisme classifiant $S \to \Curvesstack$ se factorise par
le sous-champ ouvert $\Curvesstack^{prestable}_g$ du
lemme \ref{moduli-curves-lemma-prestable-one-piece-per-genus}.

\begin{lemma}
\label{moduli-curves-lemma-prestable-one-piece-per-genus}
Il existe une décomposition en sous-champs ouverts et fermés
$$
\Curvesstack^{prestable} = \coprod\nolimits_{g \geq 0}
\Curvesstack^{prestable}_g
$$
où chaque $\Curvesstack^{prestable}_g$ est caractérisé de la manière suivante :
\begin{enumerate}
\item étant donnée une famille de courbes $f : X \to S$, les conditions suivantes sont équivalentes :
\begin{enumerate}
\item le morphisme classifiant $S \to \Curvesstack$ se factorise
par $\Curvesstack^{prestable}_g$,
\item $X \to S$ est une famille préstable de courbes et
$R^1f_*\mathcal{O}_X$ est un $\mathcal{O}_S$-module localement libre de rang $g$,
\end{enumerate}
\item étant donné un schéma $X$ propre sur un corps $k$ avec
$\dim(X) \leq 1$, les conditions suivantes sont équivalentes :
\begin{enumerate}
\item le morphisme classifiant $\Spec(k) \to \Curvesstack$
se factorise par $\Curvesstack^{prestable}_g$,
\item les singularités de $X$ sont au plus nodales, $\dim(X) = 1$,
$k = H^0(X, \mathcal{O}_X)$, et le genre de $X$ est $g$.
\end{enumerate}
\end{enumerate}
\end{lemma}

\begin{proof}
Puisque nous avons vu que $\Curvesstack^{prestable}$ est contenu
dans $\Curvesstack^{h0, 1}$, l'assertion résulte des
lemmes \ref{moduli-curves-lemma-prestable-curves} et
\ref{moduli-curves-lemma-pre-genus-one-piece-per-genus}.
\end{proof}

\begin{lemma}
\label{moduli-curves-lemma-prestable-curves-smooth}
Les morphismes
$\Curvesstack^{prestable} \to \Spec(\mathbf{Z})$ et
$\Curvesstack^{prestable}_g \to \Spec(\mathbf{Z})$ sont
lisses.
\end{lemma}

\begin{proof}
Puisque $\Curvesstack^{prestable}$ est un sous-champ ouvert de
$\Curvesstack^{nodal}$, cela résulte du
lemme \ref{moduli-curves-lemma-nodal-curves-smooth}.
\end{proof}




\section{Courbes semi-stables}
\label{moduli-curves-section-semistable-curves}

\noindent
Le lemme suivant nous aidera à comprendre les familles de courbes semi-stables.

\begin{lemma}
\label{moduli-curves-lemma-semistable}
Soit $f : X \to S$ une famille préstable de courbes de genre $g \geq 1$.
Soit $s \in S$ un point du schéma de base. Soit $m \geq 2$.
Les conditions suivantes sont équivalentes :
\begin{enumerate}
\item $X_s$ n'a pas de queue rationnelle
(Courbes algébriques, exemple \ref{curves-example-rational-tail}), et
\item $f^*f_*\omega_{X/S}^{\otimes m} \to \omega_{X/S}^{\otimes m}$
est surjectif sur $f^{-1}(U)$ pour un ouvert $s \in U \subset S$.
\end{enumerate}
\end{lemma}

\begin{proof}
Supposons (2). D'après les résultats de la section \ref{moduli-curves-section-relative-dualizing},
on en déduit que $\omega_{X_s}^{\otimes m}$ est
globalement engendré. Cependant, si $C \subset X_s$
est une queue rationnelle, alors $\deg(\omega_{X_s}|_C) < 0$ d'après
Courbes algébriques, lemme \ref{curves-lemma-rational-tail-negative},
d'où $H^0(C, \omega_{X_s}|_C) = 0$ d'après
Variétés, lemme \ref{varieties-lemma-check-invertible-sheaf-trivial},
ce qui contredit le fait qu'il soit globalement engendré.
Cela prouve (1).

\medskip\noindent
Supposons (1). Supposons d'abord que $g \geq 2$. L'hypothèse (1)
implique que $\omega_{X_s}^{\otimes m}$ est globalement engendré ;
voir Courbes algébriques, lemme \ref{curves-lemma-contracting-rational-tails}.
En outre, on a
$$
\Hom_{\kappa(s)}(H^1(X_s, \omega_{X_s}^{\otimes m}), \kappa(s)) =
H^0(X_s, \omega_{X_s}^{\otimes 1 - m})
$$
par dualité ; voir Courbes algébriques, lemme \ref{curves-lemma-duality-dim-1-CM}.
Puisque $\omega_{X_s}^{\otimes m}$ est globalement engendré, on voit
que sa restriction à chaque composante irréductible est de degré non négatif.
Par conséquent, la restriction de $\omega_{X_s}^{\otimes 1 - m}$ à chaque
composante irréductible est de degré non positif. Comme
$\deg(\omega_{X_s}^{\otimes 1 - m}) = (1 - m)(2g - 2) < 0$ par Riemann--Roch
(Courbes algébriques, lemme \ref{curves-lemma-rr}), on conclut que ce $H^0$
est nul d'après Variétés, lemme \ref{varieties-lemma-no-sections-dual-nef}.
Le théorème de cohomologie et changement de base montre que
$$
E = Rf_*\omega_{X/S}^{\otimes m}
$$
est un complexe parfait dont la formation commute à tout changement de base
(Catégories dérivées des espaces, lemme
\ref{spaces-perfect-lemma-flat-proper-perfect-direct-image-general}).
L'annulation démontrée ci-dessus montre que $E \otimes^\mathbf{L} \kappa(s)$
est égal à $H^0(X_s, \omega_{X_s}^{\otimes m})$ placé en degré $0$.
Quitte à restreindre $S$, on voit que $E = f_*\omega_{X/S}^{\otimes m}$
est un $\mathcal{O}_S$-module localement libre placé en degré $0$
(et que sa formation commute à tout changement de base, comme nous
l'avons déjà dit) ; voir Catégories dérivées des espaces, lemme
\ref{spaces-perfect-lemma-open-where-cohomology-in-degree-i-rank-r-geometric}.
Le morphisme $f^*f_*\omega_{X/S}^{\otimes m} \to \omega_{X/S}^{\otimes m}$
devient surjectif après restriction à $X_s$. Il est donc surjectif dans
un voisinage ouvert de $X_s$. Puisque $f$ est propre, ce voisinage
contient $f^{-1}(U)$ pour un voisinage ouvert
$U$ de $s$ dans $S$.

\medskip\noindent
Supposons (1) et $g = 1$. D'après
Courbes algébriques, lemme \ref{curves-lemma-contracting-rational-tails},
l'hypothèse (1) signifie que $\omega_{X_s}$ est isomorphe à
$\mathcal{O}_{X_s}$. S'il est possible de montrer que, quitte à restreindre $S$,
le faisceau inversible $\omega_{X/S}$ devient trivial, alors
la preuve est achevée. On peut supposer $S$ affine. En restreignant encore $S$,
on peut écrire
$$
Rf_*\mathcal{O}_X = (\mathcal{O}_S \xrightarrow{0} \mathcal{O}_S)
$$
placé en degrés $0$ et $1$,
de manière compatible à tout changement de base ultérieur ; voir le lemme \ref{moduli-curves-lemma-genus}.
Par dualité, cela signifie que
$$
Rf_*\omega_{X/S} = (\mathcal{O}_S \xrightarrow{0} \mathcal{O}_S)
$$
placé en degrés $0$ et $1$\footnote{On utilise que
$Rf_*\omega_{X/S}^\bullet =
Rf_*R\SheafHom_{\mathcal{O}_X}(\mathcal{O}_X. \omega_{X/S}^\bullet) =
R\SheafHom_{\mathcal{O}_S}(Rf_*\mathcal{O}_X, \mathcal{O}_S)$
d'après Dualité pour les espaces, lemme \ref{spaces-duality-lemma-iso-on-RSheafHom} et
remarque \ref{spaces-duality-remark-iso-on-RSheafHom},
puis que $\omega_{X/S}^\bullet = \omega_{X/S}[1]$ d'après
nos définitions de la section \ref{moduli-curves-section-relative-dualizing}.}.
En particulier, on obtient un isomorphisme $\mathcal{O}_S \to f_*\omega_{X/S}$
qui est compatible à tout changement de base, puisque la
formation de $Rf_*\omega_{X/S}$ est compatible au changement de base
(voir la référence donnée ci-dessus).
Par adjonction, on obtient une section globale $\sigma \in \Gamma(X, \omega_{X/S})$.
La restriction de cette section à la fibre $X_s$
est non nulle (c'est même un élément d'une base) et, puisque
$\omega_{X_s}$ est trivial sur les fibres,
cette section ne s'annule en aucun point de $X_s$.
Elle ne s'annule donc pas dans
un voisinage ouvert de $X_s$. Puisque $f$ est propre, ce voisinage
contient $f^{-1}(U)$ pour un voisinage ouvert
$U$ de $s$ dans $S$.
\end{proof}

\noindent
Le lemme \ref{moduli-curves-lemma-semistable} motive la définition suivante.

\begin{definition}
\label{moduli-curves-definition-semistable}
Soit $f : X \to S$ une famille de courbes.
On dit que $f$ est une {\it famille semi-stable de courbes} si
\begin{enumerate}
\item $X \to S$ est une famille préstable de courbes, et
\item $X_s$ est de genre $\geq 1$ et
n'a pas de queue rationnelle pour tout $s \in S$.
\end{enumerate}
\end{definition}

\noindent
En particulier, une famille préstable de courbes de genre $0$ n'est jamais
semi-stable.
Soit $X$ un schéma propre sur un corps $k$ avec $\dim(X) \leq 1$.
Alors $X \to \Spec(k)$ est une famille de courbes ; on peut donc demander
si elle est semi-stable ou non. En développant les définitions, on voit que
les conditions suivantes sont équivalentes :
\begin{enumerate}
\item $X$ est semi-stable,
\item $X$ est préstable, est de genre $\geq 1$ et n'a pas de queue rationnelle,
\item $X_{\overline{k}}$ est connexe, est lisse sur $\overline{k}$
en dehors d'un nombre fini de nœuds, est de genre $\geq 1$, et n'a aucune
composante irréductible isomorphe à $\mathbf{P}^1_{\overline{k}}$
qui rencontre le reste de $X_{\overline{k}}$ en un seul point.
\end{enumerate}
Pour voir l'équivalence de (2) et (3), on utilise que $X$ n'a pas de queue rationnelle
si et seulement si $X_{\overline{k}}$ n'a pas de queue rationnelle, d'après
Courbes algébriques, lemme \ref{curves-lemma-contracting-rational-tails}.
Cela montre que notre définition concorde avec la plupart de celles que l'on trouve
dans la littérature.

\begin{lemma}
\label{moduli-curves-lemma-semistable-curves}
Il existe un sous-champ ouvert $\Curvesstack^{semistable} \subset \Curvesstack$
tel que
\begin{enumerate}
\item étant donnée une famille de courbes $f : X \to S$, les conditions suivantes sont équivalentes :
\begin{enumerate}
\item le morphisme classifiant $S \to \Curvesstack$ se factorise
par $\Curvesstack^{semistable}$,
\item $X \to S$ est une famille semi-stable de courbes,
\end{enumerate}
\item étant donné un schéma $X$ propre sur un corps $k$ avec
$\dim(X) \leq 1$, les conditions suivantes sont équivalentes :
\begin{enumerate}
\item le morphisme classifiant $\Spec(k) \to \Curvesstack$
se factorise par $\Curvesstack^{semistable}$,
\item les singularités de $X$ sont au plus nodales, $\dim(X) = 1$,
$k = H^0(X, \mathcal{O}_X)$, le genre de $X$ est $\geq 1$, et
$X$ n'a pas de queue rationnelle,
\item les singularités de $X$ sont au plus nodales, $\dim(X) = 1$,
$k = H^0(X, \mathcal{O}_X)$, et $\omega_{X_s}^{\otimes m}$ est
globalement engendré pour $m \geq 2$.
\end{enumerate}
\end{enumerate}
\end{lemma}

\begin{proof}
L'équivalence de (2)(b) et (2)(c) est donnée par
Courbes algébriques, lemme \ref{curves-lemma-contracting-rational-tails}.
Dans la suite de la démonstration, nous travaillerons avec (2)(b),
conformément à la définition \ref{moduli-curves-definition-semistable}.

\medskip\noindent
D'après la discussion de la section \ref{moduli-curves-section-open},
il suffit de considérer les familles $f : X \to S$ de
courbes préstables. Le lemme \ref{moduli-curves-lemma-semistable}
donne l'ouverture souhaitée du lieu en question.
La formation de cet ouvert commute à tout changement de base,
car l'existence ou la non-existence de queues rationnelles est invariante
par extension du corps de base, d'après
Courbes algébriques, lemme \ref{curves-lemma-contracting-rational-tails}.
\end{proof}

\begin{lemma}
\label{moduli-curves-lemma-semistable-one-piece-per-genus}
Il existe une décomposition en sous-champs ouverts et fermés
$$
\Curvesstack^{semistable} = \coprod\nolimits_{g \geq 1}
\Curvesstack^{semistable}_g
$$
où chaque $\Curvesstack^{semistable}_g$ est caractérisé de la manière suivante :
\begin{enumerate}
\item étant donnée une famille de courbes $f : X \to S$, les conditions suivantes sont équivalentes :
\begin{enumerate}
\item le morphisme classifiant $S \to \Curvesstack$ se factorise
par $\Curvesstack^{semistable}_g$,
\item $X \to S$ est une famille semi-stable de courbes et
$R^1f_*\mathcal{O}_X$ est un $\mathcal{O}_S$-module localement libre de rang $g$,
\end{enumerate}
\item étant donné un schéma $X$ propre sur un corps $k$ avec
$\dim(X) \leq 1$, les conditions suivantes sont équivalentes :
\begin{enumerate}
\item le morphisme classifiant $\Spec(k) \to \Curvesstack$
se factorise par $\Curvesstack^{semistable}_g$,
\item les singularités de $X$ sont au plus nodales, $\dim(X) = 1$,
$k = H^0(X, \mathcal{O}_X)$, le genre de $X$ est $g$, et $X$
n'a pas de queue rationnelle,
\item les singularités de $X$ sont au plus nodales, $\dim(X) = 1$,
$k = H^0(X, \mathcal{O}_X)$, le genre de $X$ est $g$, et
$\omega_{X_s}^{\otimes m}$ est globalement engendré pour $m \geq 2$.
\end{enumerate}
\end{enumerate}
\end{lemma}

\begin{proof}
Combiner les lemmes \ref{moduli-curves-lemma-semistable-curves} et
\ref{moduli-curves-lemma-prestable-one-piece-per-genus}.
\end{proof}

\begin{lemma}
\label{moduli-curves-lemma-semistable-curves-smooth}
Les morphismes
$\Curvesstack^{semistable} \to \Spec(\mathbf{Z})$ et
$\Curvesstack^{semistable}_g \to \Spec(\mathbf{Z})$
sont lisses.
\end{lemma}

\begin{proof}
Puisque $\Curvesstack^{semistable}$ est un sous-champ ouvert de
$\Curvesstack^{nodal}$, cela résulte du
lemme \ref{moduli-curves-lemma-nodal-curves-smooth}.
\end{proof}







\section{Courbes stables}
\label{moduli-curves-section-stable-curves}

\noindent
Le lemme suivant nous aidera à comprendre les familles de courbes stables.

\begin{lemma}
\label{moduli-curves-lemma-stable}
Soit $f : X \to S$ une famille préstable de courbes de genre $g \geq 2$.
Soit $s \in S$ un point du schéma de base.
Les conditions suivantes sont équivalentes :
\begin{enumerate}
\item $X_s$ n'a ni queue rationnelle ni
pont rationnel
(Courbes algébriques, exemples
\ref{curves-example-rational-tail} et
\ref{curves-example-rational-bridge}) ;
\item $\omega_{X/S}$ est ample sur $f^{-1}(U)$ pour un ouvert $s \in U \subset S$.
\end{enumerate}
\end{lemma}

\begin{proof}
Supposons (2). Alors $\omega_{X_s}$ est ample sur $X_s$.
D'après Courbes algébriques, lemmes \ref{curves-lemma-rational-tail-negative} et
\ref{curves-lemma-rational-bridge-zero},
nous en concluons que (1) est satisfaite (nous utilisons aussi
la caractérisation des faisceaux inversibles amples donnée dans
Variétés, lemme
\ref{varieties-lemma-ampleness-in-terms-of-degrees-components}).

\medskip\noindent
Supposons (1). Alors $\omega_{X_s}$ est ample sur $X_s$ d'après
Courbes algébriques, lemme \ref{curves-lemma-contracting-rational-bridges}.
La conclusion résulte de Descente sur les espaces, lemme
\ref{spaces-descent-lemma-ample-in-neighbourhood}.
\end{proof}

\noindent
Le lemme \ref{moduli-curves-lemma-stable} motive la définition suivante.

\begin{definition}
\label{moduli-curves-definition-stable}
Soit $f : X \to S$ une famille de courbes.
On dit que $f$ est une {\it famille stable de courbes} si
\begin{enumerate}
\item $X \to S$ est une famille préstable de courbes ;
\item $X_s$ est de genre $\geq 2$ et ne possède ni queue rationnelle
ni pont rationnel pour tout $s \in S$.
\end{enumerate}
\end{definition}

\noindent
En particulier, une famille préstable de courbes de genre $0$ ou $1$ n'est jamais
stable.
Soit $X$ un schéma propre sur un corps $k$ avec $\dim(X) \leq 1$.
Alors $X \to \Spec(k)$ est une famille de courbes ; on peut donc demander
si elle est stable ou non. En explicitant les définitions, on voit que
les conditions suivantes sont équivalentes :
\begin{enumerate}
\item $X$ est stable ;
\item $X$ est préstable, est de genre $\geq 2$ et ne possède ni queue rationnelle
ni pont rationnel ;
\item $X$ est géométriquement connexe, est lisse sur $k$
en dehors d'un nombre fini de nœuds, et $\omega_X$ est ample.
\end{enumerate}
Pour voir l'équivalence de (2) et (3), on utilise
le lemme \ref{moduli-curves-lemma-stable} ci-dessus.
Cela montre que notre définition concorde avec la plupart de celles que l'on trouve
dans la littérature.

\begin{lemma}
\label{moduli-curves-lemma-stable-curves}
Il existe un sous-champ ouvert $\Curvesstack^{stable} \subset \Curvesstack$
tel que :
\begin{enumerate}
\item pour une famille de courbes $f : X \to S$, les conditions suivantes sont équivalentes :
\begin{enumerate}
\item le morphisme classifiant $S \to \Curvesstack$ se factorise
par $\Curvesstack^{stable}$ ;
\item $X \to S$ est une famille stable de courbes ;
\end{enumerate}
\item pour un schéma $X$ propre sur un corps $k$ avec
$\dim(X) \leq 1$, les conditions suivantes sont équivalentes :
\begin{enumerate}
\item le morphisme classifiant $\Spec(k) \to \Curvesstack$
se factorise par $\Curvesstack^{stable}$ ;
\item les singularités de $X$ sont au plus nodales, $\dim(X) = 1$,
$k = H^0(X, \mathcal{O}_X)$, le genre de $X$ est $\geq 2$, et
$X$ n'a ni queue rationnelle ni pont rationnel ;
\item les singularités de $X$ sont au plus nodales, $\dim(X) = 1$,
$k = H^0(X, \mathcal{O}_X)$, et $\omega_{X_s}$ est ample.
\end{enumerate}
\end{enumerate}
\end{lemma}

\begin{proof}
D'après la discussion de la section \ref{moduli-curves-section-open},
il suffit de considérer les familles $f : X \to S$ de
courbes préstables. Le lemme \ref{moduli-curves-lemma-stable}
donne l'ouverture souhaitée du lieu considéré.
La formation de cet ouvert commute à tout changement de base,
soit parce que l'existence ou la non-existence de queues ou de ponts rationnels
est invariante par extension du corps de base, d'après
Courbes algébriques, lemmes
\ref{curves-lemma-contracting-rational-tails} et
\ref{curves-lemma-contracting-rational-bridges},
soit parce que l'amplitude est invariante par extension du corps de base, d'après
Descente, lemme \ref{descent-lemma-descending-property-ample}.
\end{proof}

\begin{definition}
\label{moduli-curves-definition-deligne-mumford}
\begin{reference}
\cite{DM}
\end{reference}
Nous notons $\overline{\mathcal{M}}$ et appelons
{\it champ de modules des courbes stables} le champ algébrique
$\Curvesstack^{stable}$ qui paramètre les familles stables de courbes
introduit au lemme \ref{moduli-curves-lemma-stable-curves}.
Pour $g \geq 2$, nous notons $\overline{\mathcal{M}}_g$ et appelons
{\it champ de modules des courbes stables de genre $g$}
le champ algébrique introduit au lemme \ref{moduli-curves-lemma-stable-one-piece-per-genus}.
\end{definition}

\noindent
Voici le lemme incontournable.

\begin{lemma}
\label{moduli-curves-lemma-stable-one-piece-per-genus}
Il existe une décomposition en sous-champs ouverts et fermés
$$
\overline{\mathcal{M}} = \coprod\nolimits_{g \geq 2} \overline{\mathcal{M}}_g
$$
où chaque $\overline{\mathcal{M}}_g$ est caractérisé comme suit :
\begin{enumerate}
\item pour une famille de courbes $f : X \to S$, les conditions suivantes sont équivalentes :
\begin{enumerate}
\item le morphisme classifiant $S \to \Curvesstack$ se factorise
par $\overline{\mathcal{M}}_g$ ;
\item $X \to S$ est une famille stable de courbes et
$R^1f_*\mathcal{O}_X$ est un $\mathcal{O}_S$-module localement libre de rang $g$ ;
\end{enumerate}
\item pour un schéma $X$ propre sur un corps $k$ avec
$\dim(X) \leq 1$, les conditions suivantes sont équivalentes :
\begin{enumerate}
\item le morphisme classifiant $\Spec(k) \to \Curvesstack$
se factorise par $\overline{\mathcal{M}}_g$ ;
\item les singularités de $X$ sont au plus nodales, $\dim(X) = 1$,
$k = H^0(X, \mathcal{O}_X)$, le genre de $X$ est $g$, et $X$
n'a ni queue rationnelle ni pont rationnel ;
\item les singularités de $X$ sont au plus nodales, $\dim(X) = 1$,
$k = H^0(X, \mathcal{O}_X)$, le genre de $X$ est $g$, et
$\omega_{X_s}$ est ample.
\end{enumerate}
\end{enumerate}
\end{lemma}

\begin{proof}
Combiner les lemmes \ref{moduli-curves-lemma-stable-curves} et
\ref{moduli-curves-lemma-prestable-one-piece-per-genus}.
\end{proof}

\begin{lemma}
\label{moduli-curves-lemma-stable-curves-smooth}
Les morphismes
$\overline{\mathcal{M}} \to \Spec(\mathbf{Z})$ et
$\overline{\mathcal{M}}_g \to \Spec(\mathbf{Z})$
sont lisses.
\end{lemma}

\begin{proof}
Puisque $\overline{\mathcal{M}}$ est un sous-champ ouvert de
$\Curvesstack^{nodal}$, cela résulte du
lemme \ref{moduli-curves-lemma-nodal-curves-smooth}.
\end{proof}

\begin{lemma}
\label{moduli-curves-lemma-stable-curves-deligne-mumford}
Les champs $\overline{\mathcal{M}}$ et
$\overline{\mathcal{M}}_g$
sont des sous-champs ouverts de $\Curvesstack^{DM}$.
En particulier, $\overline{\mathcal{M}}$ et
$\overline{\mathcal{M}}_g$ sont DM
(Morphismes de champs, définition
\ref{stacks-morphisms-definition-absolute-separated})
ainsi que des champs de Deligne-Mumford
(Champs algébriques, définition \ref{algebraic-definition-deligne-mumford}).
\end{lemma}

\begin{proof}
Démontrons la première assertion.
Soit $X$ un schéma propre sur un corps $k$ dont les singularités
sont au plus nodales, avec $\dim(X) = 1$, $k = H^0(X, \mathcal{O}_X)$,
tel que le genre de $X$ soit $\geq 2$, et tel que $X$ n'ait
ni queue rationnelle ni pont rationnel.
Il faut montrer que le morphisme classifiant
$\Spec(k) \to \overline{\mathcal{M}} \to \Curvesstack$
se factorise par $\Curvesstack^{DM}$.
Nous pouvons d'abord remplacer $k$ par sa clôture algébrique
(puisque nous savons déjà que les champs considérés sont des sous-champs
ouverts du champ algébrique $\Curvesstack$).
D'après les lemmes \ref{moduli-curves-lemma-stable-curves}, \ref{moduli-curves-lemma-DM-curves} et
\ref{moduli-curves-lemma-in-DM-locus-vector-fields}, il suffit de montrer que
$\text{Der}_k(\mathcal{O}_X, \mathcal{O}_X) = 0$.
Cela est démontré dans
Courbes algébriques, lemme \ref{curves-lemma-stable-vector-fields}.

\medskip\noindent
Puisque $\Curvesstack^{DM}$ est le plus grand sous-champ ouvert de
$\Curvesstack$ qui soit DM, il en va de même du
sous-champ ouvert $\overline{\mathcal{M}}$ de $\Curvesstack^{DM}$.
Enfin, un champ algébrique DM est de Deligne-Mumford d'après
Morphismes de champs, théorème \ref{stacks-morphisms-theorem-DM}.
\end{proof}

\begin{lemma}
\label{moduli-curves-lemma-smooth-dense-in-stable}
Soit $g \geq 2$. L'inclusion
$$
|\mathcal{M}_g| \subset |\overline{\mathcal{M}}_g|
$$
est celle d'un sous-ensemble ouvert dense.
\end{lemma}

\begin{proof}
Puisque $\overline{\mathcal{M}}_g \subset \Curvesstack^{lci+}$
est ouvert et que
$\Curvesstack^{smooth} \cap \overline{\mathcal{M}}_g = \mathcal{M}_g$,
cela résulte immédiatement du
lemme \ref{moduli-curves-lemma-smooth-dense}.
\end{proof}



\section{Morphismes de contraction}
\label{moduli-curves-section-contracting}

\noindent
Nous invitons vivement le lecteur à se familiariser avec
Courbes algébriques, sections
\ref{curves-section-contracting-rational-tails},
\ref{curves-section-contracting-rational-bridges} et
\ref{curves-section-contracting-to-stable},
avant de poursuivre. Le résultat principal de cette section
est l'existence d'un morphisme de « stabilisation »
$$
\Curvesstack^{prestable}_g
\longrightarrow
\overline{\mathcal{M}}_g
$$
Voir le lemme \ref{moduli-curves-lemma-stabilization-morphism}.
En termes informels, ce morphisme envoie le point modulaire d'une
courbe nodale de genre $g$ sur le point modulaire de la courbe stable
associée construite dans
Courbes algébriques, lemme \ref{curves-lemma-characterize-contraction-to-stable}.

\begin{lemma}
\label{moduli-curves-lemma-contract}
Soient $S$ un schéma et $s \in S$ un point.
Soient $f : X \to S$ et $g : Y \to S$ des familles de courbes.
Soit $c : X \to Y$ un morphisme sur $S$. Si
$c_{s, *}\mathcal{O}_{X_s} = \mathcal{O}_{Y_s}$ et
$R^1c_{s, *}\mathcal{O}_{X_s} = 0$, alors,
après remplacement de $S$ par un voisinage ouvert de $s$,
on a $\mathcal{O}_Y = c_*\mathcal{O}_X$ et $R^1c_*\mathcal{O}_X = 0$,
et cela reste vrai après tout changement de base par un morphisme $S' \to S$.
\end{lemma}

\begin{proof}
Soit $(U, u) \to (S, s)$ un voisinage étale tel que
$\mathcal{O}_{Y_U} = (X_U \to Y_U)_*\mathcal{O}_{X_U}$ et
$R^1(X_U \to Y_U)_*\mathcal{O}_{X_U} = 0$, ces égalités restant vraies
après changement de base par $U' \to U$. Remplaçons alors $S$ par
l'image ouverte de $U \to S$. Étant donné $S' \to S$, posons $U' = U \times_S S'$ ;
nous obtenons les recouvrements étales $\{U' \to S'\}$ et
$\{Y_{U'} \to Y_{S'}\}$. La validité de l'énoncé pour
le changement de base de $c$ par $S' \to S$ résulte donc de sa validité
pour le changement de base de $X_U \to Y_U$ par
$U' \to U$. Autrement dit, la question est locale pour la
topologie étale sur $S$.
Ainsi, d'après le lemme \ref{moduli-curves-lemma-etale-locally-scheme}, nous pouvons supposer que
$X$ et $Y$ sont des schémas. D'après
Compléments sur les morphismes, lemme \ref{more-morphisms-lemma-h1-fibre-zero-isom},
il existe un sous-schéma ouvert $V \subset Y$ contenant $Y_s$
tel que $c_*\mathcal{O}_X|_V = \mathcal{O}_V$ et
$R^1c_*\mathcal{O}_X|_V = 0$, ces égalités restant vraies après
tout changement de base par $S' \to S$. Puisque $g : Y \to S$ est propre, on peut
trouver un voisinage ouvert $U \subset S$ de $s$ tel que
$g^{-1}(U) \subset V$. Alors $U$ convient.
\end{proof}

\begin{lemma}
\label{moduli-curves-lemma-contract-basic-uniqueness}
Soient $S$ un schéma et $s \in S$ un point.
Soient $f : X \to S$ et $g_i : Y_i \to S$, $i = 1, 2$, des familles de courbes.
Soient $c_i : X \to Y_i$ des morphismes sur $S$.
Supposons donné un isomorphisme $Y_{1, s} \cong Y_{2, s}$
des fibres, compatible avec $c_{1, s}$ et $c_{2, s}$.
Si $c_{1, s, *}\mathcal{O}_{X_s} = \mathcal{O}_{Y_{1, s}}$ et
$R^1c_{1, s, *}\mathcal{O}_{X_s} = 0$, alors il existe un
voisinage ouvert $U$ de $s$ et un isomorphisme
$Y_{1, U} \cong Y_{2, U}$ de familles de courbes sur $U$,
compatible avec l'isomorphisme donné des fibres ainsi qu'avec
$c_1$ et $c_2$.
\end{lemma}

\begin{proof}
Rappelons que $\mathcal{O}_{S, s} = \colim \mathcal{O}_S(U)$, où
la limite inductive est prise sur le système des voisinages affines $U$ de $s$.
La catégorie des espaces algébriques de présentation finie sur
l'anneau local est donc la limite inductive des catégories d'espaces algébriques
de présentation finie sur les voisinages affines de $s$.
Voir Limites d'espaces, lemme
\ref{spaces-limits-lemma-descend-finite-presentation}.
On se ramène ainsi au cas où $S$ est le spectre d'un
anneau local et où $s$ est le point fermé.

\medskip\noindent
Supposons $S = \Spec(A)$, où $A$ est un anneau local et $s$ son point fermé.
Écrivons $A = \colim A_j$, avec les $A_j$ locaux noethériens (par exemple essentiellement
de type fini sur $\mathbf{Z}$) et les homomorphismes de transition locaux.
Posons $S_j = \Spec(A_j)$, de point fermé $s_j$. On peut trouver un
$j$ et des familles de courbes $X_j \to S_j$, $Y_{j, i} \to S_j$ ;
voir le lemme \ref{moduli-curves-lemma-curves-qs-lfp} et
Limites de champs, lemme
\ref{stacks-limits-lemma-representable-by-spaces-limit-preserving}.
Quitte à augmenter $j$, on peut trouver des morphismes
$c_{j, i} : X_j \to Y_{j, i}$ dont le changement de base à $s$ est $c_i$ ; voir
Limites d'espaces, lemme
\ref{spaces-limits-lemma-descend-finite-presentation}.
Puisque $\kappa(s) = \colim \kappa(s_j)$, on peut de même
supposer donné un isomorphisme $Y_{j, 1, s_j} \cong Y_{j, 2, s_j}$
compatible avec $c_{j, 1, s_j}$ et $c_{j, 2, s_j}$.
Enfin, les hypothèses
$c_{1, s, *}\mathcal{O}_{X_s} = \mathcal{O}_{Y_{1, s}}$ et
$R^1c_{1, s, *}\mathcal{O}_{X_s} = 0$
sont héritées par $c_{j, 1, s_j}$, car
$\{s_j \to s\}$ est un recouvrement fpqc et
$c_{1, s}$ est le changement de base de $c_{j, 1, s_j}$ par ce recouvrement
(nous omettons les détails).
Le lemme est ainsi ramené au cas traité dans le paragraphe suivant.

\medskip\noindent
Supposons que $S$ soit le spectre d'un anneau local noethérien $\Lambda$ et que
$s$ soit le point fermé. Considérons l'image schématique $Z$ de
$$
(c_1, c_2) : X \longrightarrow Y_1 \times_S Y_2
$$
L'énoncé du lemme équivaut à dire que $Z$ s'envoie
isomorphiquement sur $Y_1$ et $Y_2$ par les morphismes de projection.
Puisque la formation de l'image schématique de ce morphisme commute
au changement de base plat (Morphismes d'espaces, lemme
\ref{spaces-morphisms-lemma-flat-base-change-scheme-theoretic-image}),
nous pouvons remplacer $\Lambda$ par son complété
(Compléments d'algèbre, section \ref{more-algebra-section-permanence-completion}).

\medskip\noindent
Supposons que $S$ soit le spectre d'un anneau local noethérien complet $\Lambda$.
Observons que $X$, $Y_1$, $Y_2$ sont alors des schémas
(Compléments sur les morphismes d'espaces, lemme
\ref{spaces-more-morphisms-lemma-projective-over-complete}).
Notons $X_n$, $Y_{1, n}$, $Y_{2, n}$ les changements de base de
$X$, $Y_1$, $Y_2$ à $\Spec(\Lambda/\mathfrak m^{n + 1})$.
Rappelons que la flèche
$$
\Deformationcategory_{X_s \to Y_{2, s}} \cong
\Deformationcategory_{X_s \to Y_{1, s}} \longrightarrow
\Deformationcategory_{X_s}
$$
est une équivalence ; voir Problèmes de déformation, lemme
\ref{examples-defos-lemma-schemes-morphisms-smooth-to-base}.
Il existe donc un isomorphisme d'objets formels
$(X_n \to Y_{1, n}) \cong (X_n \to Y_{2, n})$
de $\Deformationcategory_{X_s \to Y_{1, s}}$.
Enfin, le théorème d'algébrisation de Grothendieck
(Cohomologie des schémas, lemme \ref{coherent-lemma-algebraize-morphism})
fournit un isomorphisme $Y_1 \to Y_2$ compatible avec $c_1$ et $c_2$.
\end{proof}

\begin{lemma}
\label{moduli-curves-lemma-contract-basic}
Soit $f : X \to S$ une famille de courbes. Soit $s \in S$ un point.
Soit $h_0 : X_s \to Y_0$ un morphisme vers un schéma propre $Y_0$ sur $\kappa(s)$
tel que $h_{0, *}\mathcal{O}_{X_s} = \mathcal{O}_{Y_0}$ et
$R^1h_{0, *}\mathcal{O}_{X_s} = 0$. Alors il existe un voisinage
étale élémentaire $(U, u) \to (S, s)$, une famille de courbes $Y \to U$
et un morphisme $h : X_U \to Y$ sur $U$ dont la fibre en $u$
est isomorphe à $h_0$.
\end{lemma}

\begin{proof}
Commençons par quelques réductions ; le lecteur peut passer directement à la suite.
La question est locale sur $S$ ; on peut donc supposer $S$ affine.
Écrivons $S = \lim S_i$ comme limite projective filtrante de schémas affines $S_i$
de type fini sur $\mathbf{Z}$.
Pour un certain $i$, on peut trouver une famille de courbes $X_i \to S_i$
dont le changement de base est $X \to S$. Cela résulte du
lemme \ref{moduli-curves-lemma-curves-qs-lfp} et de
Limites de champs, lemme
\ref{stacks-limits-lemma-representable-by-spaces-limit-preserving}.
Soit $s_i \in S_i$ l'image de $s$. Observons que
$\kappa(s) = \colim \kappa(s_i)$ et que $X_s$ est un schéma
(Espaces sur les corps, lemme
\ref{spaces-over-fields-lemma-codim-1-point-in-schematic-locus}).
Quitte à augmenter $i$, on peut supposer qu'il existe un morphisme
$h_{i, 0} : X_{i, s_i} \to Y_i$
de schémas de type fini sur $\kappa(s_i)$ dont le changement de base à
$\kappa(s)$ est $h_0$ ; voir
Limites, lemme \ref{limits-lemma-descend-finite-presentation}.
Quitte à augmenter $i$, on peut supposer $Y_i$ propre sur $\kappa(s_i)$ ; voir
Limites, lemme \ref{limits-lemma-eventually-proper}.
Soit $g_{i, 0} : Y_0 \to Y_{i, 0}$ la projection. Observons qu'il
s'agit d'un morphisme fidèlement plat, car il se déduit par changement de base de
$\Spec(\kappa(s)) \to \Spec(\kappa(s_i))$.
Par changement de base plat, on a
$$
h_{0, *}\mathcal{O}_{X_s} = g_{i, 0}^*h_{i, 0, *}\mathcal{O}_{X_{i, s_i}}
\quad\text{et}\quad
R^1h_{0, *}\mathcal{O}_{X_s} = g_{i, 0}^*Rh_{i, 0, *}\mathcal{O}_{X_{i, s_i}}
$$
voir Cohomologie des schémas, lemme
\ref{coherent-lemma-flat-base-change-cohomology}.
Par fidèle platitude, on voit que $X_i \to S_i$, $s_i \in S_i$ et
$X_{i, s_i} \to Y_i$ satisfont à toutes les hypothèses du lemme.
On est ainsi ramené au cas traité dans le paragraphe suivant.

\medskip\noindent
Supposons $S$ affine de type fini sur $\mathbf{Z}$.
Soit $\mathcal{O}_{S, s}^h$ l'hensélisé de l'anneau local
de $S$ en $s$. Observons que $\mathcal{O}_{S, s}^h$ est un anneau G d'après
Compléments d'algèbre, lemme \ref{more-algebra-lemma-henselization-G-ring} et
proposition \ref{more-algebra-proposition-ubiquity-G-ring}.
Supposons que nous puissions construire une famille de courbes
$Y' \to \Spec(\mathcal{O}_{S, s}^h)$ et un morphisme
$$
h' : X \times_S \Spec(\mathcal{O}_{S, s}^h) \longrightarrow Y'
$$
sur $\Spec(\mathcal{O}_{S, s}^h)$ dont le changement de base au point
fermé soit $h_0$. Cela suffira. En effet, utilisons d'abord l'égalité
$$
\mathcal{O}_{S, s}^h = \colim_{(U, u)} \mathcal{O}_U(U)
$$
où la limite inductive est prise sur la catégorie filtrante des
voisinages étales élémentaires (Compléments sur les morphismes, lemme
\ref{more-morphisms-lemma-describe-henselization}).
Utilisons ensuite de nouveau le fait que $Y'$ descend en une famille
$Y \to U$ pour un certain $U$ (voir les références données ci-dessus).
Puis appliquons
Limites, lemme \ref{limits-lemma-descend-finite-presentation},
pour faire descendre $h'$ en un certain $h$. On est ainsi ramené au cas
traité dans le paragraphe suivant.

\medskip\noindent
Supposons $S = \Spec(\Lambda)$, où $(\Lambda, \mathfrak m, \kappa)$
est un anneau G local noethérien hensélien et $s$ le point fermé de $S$.
Rappelons que l'application
$$
\Deformationcategory_{X_s \to Y_0} \to \Deformationcategory_{X_s}
$$
est une équivalence ; voir Problèmes de déformation, lemme
\ref{examples-defos-lemma-schemes-morphisms-smooth-to-base}.
(C'est la seule étape importante de la démonstration ; tout le reste
relève de la technique.) Notons $\Lambda^\wedge$ le complété $\mathfrak m$-adique.
Les images réciproques $X_n$ de $X$ sur $\Lambda/\mathfrak m^{n + 1}$ définissent
un objet formel $\xi$ de $\Deformationcategory_{X_s}$ sur $\Lambda^\wedge$.
L'équivalence fournit un objet formel
$\xi'$ de $\Deformationcategory_{X_s \to Y_0}$ sur $\Lambda^\wedge$.
On obtient ainsi un immense diagramme commutatif
$$
\xymatrix{
\ldots \ar[r] &
X_n \ar[r] \ar[d] &
X_{n - 1} \ar[r] \ar[d] &
\ldots \ar[r] &
X_s \ar[d] \\
\ldots \ar[r] &
Y_n \ar[r] \ar[d] &
Y_{n - 1} \ar[r] \ar[d] &
\ldots \ar[r] &
Y_0 \ar[d] \\
\ldots \ar[r] &
\Spec(\Lambda/\mathfrak m^{n + 1}) \ar[r] &
\Spec(\Lambda/\mathfrak m^n) \ar[r] &
\ldots \ar[r] &
\Spec(\kappa)
}
$$
L'objet formel $(Y_n)$ provient d'une famille de courbes
$Y' \to \Spec(\Lambda^\wedge)$ d'après
Quot, lemme \ref{quot-lemma-curves-existence}.
D'après Compléments sur les morphismes d'espaces, lemme
\ref{spaces-more-morphisms-lemma-algebraize-morphism},
on obtient un morphisme $h' : X_{\Lambda^\wedge} \to Y'$ qui induit
les morphismes donnés $X_n \to Y_n$ pour tout $n$,
et en particulier le morphisme donné $X_s \to Y_0$.

\medskip\noindent
Pour terminer, appliquons un argument standard d'algébrisation et d'approximation.
Observons d'abord que l'on peut trouver une sous-$\Lambda$-algèbre de type fini
$\Lambda \subset A \subset \Lambda^\wedge$, une famille de courbes
$Y'' \to \Spec(A)$ et un morphisme $h'' : X_A \to Y''$ sur $A$
dont le changement de base à $\Lambda^\wedge$ est $h'$.
Cela résulte de ce que $\Lambda^\wedge$ est la limite inductive filtrante
de ces anneaux $A$ et de ce que l'on peut raisonner comme précédemment en utilisant
le fait que $\Curvesstack$ est localement de présentation finie
(ce qui fournit $Y''$ sur $A$ d'après
Limites de champs, lemme
\ref{stacks-limits-lemma-representable-by-spaces-limit-preserving})
et en appliquant
Limites d'espaces, lemme \ref{spaces-limits-lemma-descend-finite-presentation},
pour faire descendre $h'$ en un certain $h''$.
On peut alors appliquer la propriété d'approximation des anneaux G
(sous la forme du théorème de Lissage des homomorphismes d'anneaux,
\ref{smoothing-theorem-approximation-property})
afin de trouver une application $A \to \Lambda$ qui induise la même application
$A \to \kappa$ que celle obtenue de $A \to \Lambda^\wedge$.
En changeant la base de $h''$ à $\Lambda$, on achève la démonstration.
\end{proof}

\begin{lemma}
\label{moduli-curves-lemma-contract-prestable-to-stable}
Soit $f : X \to S$ une famille préstable de courbes de genre $g \geq 2$.
Il existe une factorisation $X \to Y \to S$ de $f$, où $g : Y \to S$
est une famille stable de courbes et où $c : X \to Y$ possède les propriétés
suivantes :
\begin{enumerate}
\item $\mathcal{O}_Y = c_*\mathcal{O}_X$ et $R^1c_*\mathcal{O}_X = 0$,
ces égalités restant vraies après tout changement de base par un morphisme $S' \to S$ ;
\item pour tout $s \in S$, le morphisme $c_s : X_s \to Y_s$ est la
contraction des queues et ponts rationnels étudiée dans
Courbes algébriques, section \ref{curves-section-contracting-to-stable}.
\end{enumerate}
De plus, $c : X \to Y$ est unique à isomorphisme unique près.
\end{lemma}

\begin{proof}
Soit $s \in S$. Soit $c_0 : X_s \to Y_0$ la contraction de
Courbes algébriques, section \ref{curves-section-contracting-to-stable}
(plus précisément, Courbes algébriques, lemme
\ref{curves-lemma-characterize-contraction-to-stable}).
D'après le lemme \ref{moduli-curves-lemma-contract-basic},
il existe un voisinage étale élémentaire
$(U, u)$ et un morphisme $c : X_U \to Y$
de familles de courbes sur $U$ qui redonne
$c_0$ comme fibre en $u$.
Puisque $\omega_{Y_0}$ est ample, quitte à rétrécir $U$,
on voit que $Y \to U$ est une famille stable de genre $g$
par l'ouverture contenue dans
les lemmes \ref{moduli-curves-lemma-stable-curves} et \ref{moduli-curves-lemma-stable-one-piece-per-genus}.
Quitte à rétrécir $U$ une nouvelle fois, l'assertion (1) du lemme pour
$c : X_U \to Y$ résulte du lemme \ref{moduli-curves-lemma-contract}.
De plus, la partie (2) résulte de l'unicité dans Courbes algébriques, lemme
\ref{curves-lemma-characterize-contraction-to-stable}.
On en conclut qu'un morphisme $c$ comme dans le lemme existe localement pour
la topologie étale sur $S$. Plus précisément, il existe un recouvrement étale
$\{U_i \to S\}$ et des morphismes $c_i : X_{U_i} \to Y_i$ sur $U_i$,
où $Y_i \to U_i$ est une famille stable de courbes
possédant les propriétés (1) et (2) énoncées dans le lemme.

\medskip\noindent
Pour achever la démonstration, il suffit de prouver l'unicité de $c : X \to Y$
(à isomorphisme unique près). En effet, une fois cela établi, nous
obtenons des isomorphismes
$$
\varphi_{ij} :
Y_i \times_{U_i} (U_i \times_S U_j)
\longrightarrow
Y_i \times_{U_j} (U_i \times_S U_j)
$$
satisfaisant la condition de cocycle (par unicité) sur
$U_i \times U_j \times U_k$. Puisque $\overline{\mathcal{M}_g}$
est un champ algébrique, les données de descente sont effectives,
et nous obtenons $Y \to S$. Les morphismes $c_i$ descendent en un morphisme
$c : X \to Y$ sur $S$. Enfin, les propriétés (1) et (2) pour $c$
résultent immédiatement des propriétés (1) et (2) pour $c_i$.

\medskip\noindent
Enfin, si $c_1 : X \to Y_i$, $i = 1, 2$, sont deux morphismes vers
des familles stables de courbes sur $S$ satisfaisant (1) et (2), alors
nous obtenons un morphisme $Y_1 \to Y_2$ compatible avec $c_1$ et $c_2$,
au moins localement sur $S$, grâce au lemme \ref{moduli-curves-lemma-contract-basic-uniqueness}.
Nous omettons la vérification de l'unicité de ces morphismes
(indication : elle résulte du fait que l'image schématique
de $c_1$ est $Y_1$). Ces morphismes définis localement se recollent donc,
ce qui achève la démonstration.
\end{proof}

\begin{lemma}
\label{moduli-curves-lemma-stabilization-morphism}
Soit $g \geq 2$. Il existe un morphisme de champs algébriques sur $\mathbf{Z}$
$$
stabilisation :
\Curvesstack^{prestable}_g
\longrightarrow
\overline{\mathcal{M}}_g
$$
qui envoie une famille préstable de courbes $X \to S$ de genre $g$
sur la famille stable $Y \to S$ qui lui est associée dans le
lemme \ref{moduli-curves-lemma-contract-prestable-to-stable}.
\end{lemma}

\begin{proof}
Pour le voir, il suffit de vérifier que la construction du
lemme \ref{moduli-curves-lemma-contract-prestable-to-stable}
est compatible aux changements de base (ainsi qu'aux isomorphismes, ce qui est immédiat) ;
voir l'(abus de) langage pour les champs algébriques introduit dans
Propriétés des champs, section \ref{stacks-properties-section-conventions}.
Il suffit pour cela de vérifier que les propriétés (1) et (2) du
lemme \ref{moduli-curves-lemma-contract-prestable-to-stable} sont stables
par changement de base. C'est immédiat pour (1).
Pour (2), cela résulte soit de ce que les contractions de
Courbes algébriques, lemmes \ref{curves-lemma-contracting-rational-tails} et
\ref{curves-lemma-contracting-rational-bridges},
sont stables par extension du corps de base, soit
de ce que les conditions caractérisant les morphismes sur les fibres dans
Courbes algébriques, lemme \ref{curves-lemma-characterize-contraction-to-stable},
sont préservées par extension du corps de base.
\end{proof}



\section{Théorème de réduction stable}
\label{moduli-curves-section-stable-reduction}

\noindent
Dans le chapitre sur la réduction semi-stable, nous avons démontré le célèbre
théorème de réduction semi-stable des courbes. Soit $K$ le corps des fractions
d'un anneau de valuation discrète $R$. Soit $C$ une courbe projective lisse
sur $K$ telle que $K = H^0(C, \mathcal{O}_C)$. Conformément à
Réduction semi-stable, définition \ref{models-definition-semistable},
on dit que $C$ a une {\it réduction semi-stable} si, au choix,
il existe une famille préstable de courbes sur $R$ de fibre générique $C$, ou
un (ce qui équivaut à tout) modèle régulier minimal de $C$ sur $R$ est préstable.
Nous montrons dans cette section que, pour les courbes de genre $g \geq 2$,
cela équivaut aussi à l'existence d'une réduction stable.

\begin{lemma}
\label{moduli-curves-lemma-stable-reduction}
Soit $R$ un anneau de valuation discrète de corps des fractions $K$.
Soit $C$ une courbe projective lisse sur $K$ telle que
$K = H^0(C, \mathcal{O}_C)$ et de genre $g \geq 2$.
Les conditions suivantes sont équivalentes :
\begin{enumerate}
\item $C$ a une réduction semi-stable
(Réduction semi-stable, définition \ref{models-definition-semistable}) ;
\item il existe une famille stable de courbes sur $R$ de fibre générique $C$.
\end{enumerate}
\end{lemma}

\begin{proof}
Une famille stable de courbes étant aussi préstable, il est immédiat que
(2) implique (1). Réciproquement, étant donnée une famille préstable de courbes sur $R$
de fibre générique $C$, le lemme \ref{moduli-curves-lemma-contract-prestable-to-stable}
permet de la contracter en une famille stable de courbes. Comme la fibre générique
est déjà stable, cette opération ne la modifie pas,
ce qui achève la démonstration.
\end{proof}

\noindent
Le lemme suivant affirme que la famille stable de courbes sur $R$
promise par le lemme \ref{moduli-curves-lemma-stable-reduction}
est unique à isomorphisme unique près.

\begin{lemma}
\label{moduli-curves-lemma-unique-stable-model}
Soit $R$ un anneau de valuation discrète de corps des fractions $K$.
Soit $C$ une courbe propre et lisse sur $K$
telle que $K = H^0(C, \mathcal{O}_C)$ et de genre $g$.
Si $X$ et $X'$ sont des modèles de $C$
(Réduction semi-stable, section \ref{models-section-models})
et si $X$ et $X'$ sont des familles stables de courbes de genre $g$ sur $R$,
alors il existe un unique isomorphisme de modèles $X \to X'$.
\end{lemma}

\begin{proof}
Soit $Y$ le modèle minimal de $C$. Rappelons que $Y$ existe, est
unique et est au plus nodal de dimension relative $1$ sur $R$ ; voir
Réduction semi-stable,
proposition \ref{models-proposition-exists-minimal-model} et
lemmes \ref{models-lemma-minimal-model-unique} et
\ref{models-lemma-semistable} (ce dernier s'applique puisque l'on dispose de $X$).
Il existe un morphisme de contraction
$$
Y \longrightarrow Z
$$
tel que $Z$ soit une famille stable de courbes de genre $g$ sur $R$
(lemme \ref{moduli-curves-lemma-contract-prestable-to-stable}). Nous affirmons
qu'il existe un unique isomorphisme de modèles $X \to Z$.
Par symétrie, il en va de même pour $X'$, ce qui achèvera la démonstration.

\medskip\noindent
D'après Réduction semi-stable, lemme \ref{models-lemma-blowup-at-worst-nodal},
il existe une suite
$$
X_m \to \ldots \to X_1 \to X_0 = X
$$
telle que $X_{i + 1} \to X_i$ soit l'éclatement d'un point fermé
$x_i$ où $X_i$ est singulier, que $X_i \to \Spec(R)$ soit au plus nodal
de dimension relative $1$, et que $X_m$ soit régulier.
D'après Réduction semi-stable, lemme \ref{models-lemma-pre-exists-minimal-model},
il existe une suite
$$
X_m = Y_n \to Y_{n - 1} \to \ldots \to Y_1 \to Y_0 = Y
$$
de modèles propres réguliers de $C$, dont chaque morphisme est une
contraction d'une courbe exceptionnelle de première espèce\footnote{En fait,
on a $X_m = Y$, c'est-à-dire que $X_m$ ne contient aucune courbe exceptionnelle
de première espèce. Nous invitons le lecteur à le vérifier,
car cela simplifie quelque peu la démonstration.}.
D'après Réduction semi-stable, lemme \ref{models-lemma-blowdown-at-worst-nodal},
chaque $Y_i$ est au plus nodal de dimension relative $1$ sur $R$.
Pour établir notre affirmation, il suffit de montrer qu'il existe un isomorphisme
$X \to Z$ compatible avec les morphismes $X_m \to X$
et $X_m = Y_n \to Y \to Z$. Soit $s \in \Spec(R)$ le point fermé.
En appliquant soit le
lemme \ref{moduli-curves-lemma-contract-basic-uniqueness}, soit le
lemme \ref{moduli-curves-lemma-contract-prestable-to-stable},
on se ramène à démontrer que les morphismes
$X_{m, s} \to X_s$ et
$X_{m, s} \to Z_s$
sont tous deux égaux au morphisme canonique de
Courbes algébriques, lemme \ref{curves-lemma-characterize-contraction-to-stable}.

\medskip\noindent
Pour un morphisme $c : U \to V$ de schémas sur $\kappa(s)$,
on dit que $c$ possède la propriété (*) si $\dim(U_v) \leq 1$ pour $v \in V$,
$\mathcal{O}_V = c_*\mathcal{O}_U$ et $R^1c_*\mathcal{O}_U = 0$.
Cette propriété est stable par composition.
Puisque $X_s$ et $Z_s$ sont tous deux des courbes stables de genre $g$ sur $\kappa(s)$,
il suffit de montrer que chacun des morphismes $Y_s \to Z_s$,
$X_{i + 1, s} \to X_{i, s}$ et $Y_{i + 1, s} \to Y_{i, s}$
possède la propriété (*) ; voir
Courbes algébriques, lemme \ref{curves-lemma-characterize-contraction-to-stable}.

\medskip\noindent
La propriété (*) est satisfaite par $Y_s \to Z_s$ par construction.

\medskip\noindent
Les morphismes $c : X_{i + 1, s} \to X_{i, s}$ sont construits et étudiés
dans la démonstration de
Réduction semi-stable, lemme \ref{models-lemma-blowup-at-worst-nodal}.
Il suffit de vérifier (*) localement pour la topologie étale sur $X_{i, s}$.
Il suffit donc de vérifier (*) pour le changement de base du morphisme
« $X_1 \to X_0$ » de Réduction semi-stable, exemple \ref{models-example-blowup},
à $R/\pi R$.
Nous laissons le calcul explicite au lecteur.

\medskip\noindent
Le morphisme $c : Y_{i + 1, s} \to Y_{i, s}$ est la restriction
de la contraction d'une courbe exceptionnelle $E \subset Y_{i + 1}$
de première espèce ; autrement dit, $b : Y_{i + 1} \to Y_i$ est une contraction de $E$,
c'est-à-dire que $b$ est l'éclatement d'un point régulier de la surface $Y_i$
(Résolution des surfaces, section \ref{resolve-section-minus-one}).
On a alors $\mathcal{O}_{Y_i} = b_*\mathcal{O}_{Y_{i + 1}}$ et
$R^1b_*\mathcal{O}_{Y_{i + 1}} = 0$ ; voir par exemple
Résolution des surfaces, lemme
\ref{resolve-lemma-cohomology-of-blowup}.
Nous en concluons que $\mathcal{O}_{Y_{i, s}} = c_*\mathcal{O}_{Y_{i + 1, s}}$
et $R^1c_*\mathcal{O}_{Y_{i + 1, s}} = 0$ d'après
Compléments sur les morphismes, lemmes \ref{more-morphisms-lemma-check-h1-fibre-zero},
\ref{more-morphisms-lemma-h1-fibre-zero} et
\ref{more-morphisms-lemma-h1-fibre-zero-check-h0-kappa}
(ce dernier ne donne que la surjectivité de
$\mathcal{O}_{Y_{i, s}} \to c_*\mathcal{O}_{Y_{i + 1, s}}$, mais
l'injectivité résulte aisément de ce que $Y_{i, s}$ est réduit
et que $c$ ne modifie la situation qu'au-dessus d'un point fermé).
Cela achève la démonstration.
\end{proof}

\noindent
Du lemme \ref{moduli-curves-lemma-stable-reduction} et du
théorème de Réduction semi-stable \ref{models-theorem-semistable-reduction},
on déduit immédiatement le théorème de réduction stable.

\begin{theorem}
\label{moduli-curves-theorem-stable-reduction}
\begin{reference}
\cite[corollaire 2.7]{DM}
\end{reference}
Soit $R$ un anneau de valuation discrète de corps des fractions $K$. Soit $C$ une
courbe projective lisse sur $K$ telle que $H^0(C, \mathcal{O}_C) = K$
et de genre $g \geq 2$. Alors :
\begin{enumerate}
\item il existe une extension d'anneaux de valuation discrète $R \subset R'$
induisant une extension finie séparable des corps des fractions $K'/K$, et
une famille stable de courbes $Y \to \Spec(R')$ de genre $g$ telle que
$Y_{K'} \cong C_{K'}$ sur $K'$ ;
\item il existe une extension finie séparable $L/K$ et une famille stable
de courbes $Y \to \Spec(A)$ de genre $g$, où $A \subset L$
est la clôture intégrale de $R$ dans $L$, telles que
$Y_L \cong C_L$ sur $L$.
\end{enumerate}
\end{theorem}

\begin{proof}
La partie (1) résulte immédiatement du lemme \ref{moduli-curves-lemma-stable-reduction} et du
théorème de Réduction semi-stable \ref{models-theorem-semistable-reduction}.

\medskip\noindent
Démontrons (2). Soit $L/K$ l'extension finie séparable obtenue dans la partie (3) du
théorème de Réduction semi-stable \ref{models-theorem-semistable-reduction}.
Soit $A \subset L$ la clôture intégrale de $R$.
Rappelons que $A$ est un anneau de Dedekind fini sur $R$, possédant
un nombre fini d'idéaux maximaux $\mathfrak m_1, \ldots, \mathfrak m_n$ ; voir
Compléments d'algèbre, remarque \ref{more-algebra-remark-finite-separable-extension}.
Posons $S = \Spec(A)$, $S_i = \Spec(A_{\mathfrak m_i})$,
$U = \Spec(L)$ et $U_i = S_i \setminus \{\mathfrak m_i\}$.
Observons que $U \cong U_i$ pour $i = 1, \ldots, n$.
Posons $X = C_L$, considéré comme schéma sur le sous-schéma ouvert $U$ de $S$.
Par le choix de $L$ et de $A$ et par le lemme \ref{moduli-curves-lemma-stable-reduction},
nous disposons de familles stables de courbes $X_i \to S_i$ et d'isomorphismes
$X \times_U U_i \cong X_i \times_{S_i} U_i$.
D'après Limites d'espaces, lemme
\ref{spaces-limits-lemma-glueing-near-multiple-closed-points},
on peut trouver un morphisme de présentation finie $Y \to S$
dont le changement de base à $S_i$ est isomorphe à $X_i$ pour $i = 1, \ldots, n$.
On peut aussi utiliser le fait que $S = \bigcup_{i = 1, \ldots, n} S_i$
est un recouvrement ouvert de $S$ et que $S_i \cap S_j = U$ pour $i \not = j$,
puis appliquer $n - 1$ fois
Limites d'espaces, lemme \ref{spaces-limits-lemma-relative-glueing},
afin d'obtenir $Y \to S$ dont le
changement de base à $S_i$ est isomorphe à $X_i$ pour $i = 1, \ldots, n$.
Il est clair que $Y \to S$ est la famille stable de courbes recherchée.
\end{proof}






\section{Propriétés du champ des courbes stables}
\label{moduli-curves-section-properties-stable}

\noindent
Dans cette section, nous démontrons le résultat structurel fondamental pour
$\overline{\mathcal{M}}_g$ lorsque $g \geq 2$.

\begin{lemma}
\label{moduli-curves-lemma-stable-separated}
Soit $g \geq 2$. Le champ $\overline{\mathcal{M}}_g$ est séparé.
\end{lemma}

\begin{proof}
Cet énoncé signifie que le morphisme
$\overline{\mathcal{M}}_g \to \Spec(\mathbf{Z})$ est séparé.
Nous allons le démontrer au moyen du critère valuatif noethérien raffiné
énoncé dans
Compléments sur les morphismes de champs, lemme
\ref{stacks-more-morphisms-lemma-refined-valuative-criterion-separated}.

\medskip\noindent
Puisque $\overline{\mathcal{M}}_g$ est un sous-champ ouvert de
$\Curvesstack$, le morphisme $\overline{\mathcal{M}}_g \to \Spec(\mathbf{Z})$
est quasi-séparé et
localement de présentation finie d'après le lemme \ref{moduli-curves-lemma-curves-qs-lfp}.
En particulier, le champ $\overline{\mathcal{M}}_g$ est localement
noethérien (Morphismes de champs, lemme
\ref{stacks-morphisms-lemma-locally-finite-type-locally-noetherian}).
D'après le lemme \ref{moduli-curves-lemma-smooth-dense-in-stable}, l'immersion ouverte
$\mathcal{M}_g \to \overline{\mathcal{M}}_g$
est d'image dense. De plus, $\mathcal{M}_g \to \overline{\mathcal{M}}_g$
est quasi-compact (Morphismes de champs, lemme
\ref{stacks-morphisms-lemma-locally-closed-in-noetherian}),
donc de type fini. Toutes les hypothèses préliminaires de
Compléments sur les morphismes de champs, lemme
\ref{stacks-more-morphisms-lemma-refined-valuative-criterion-separated},
sont ainsi satisfaites pour les morphismes
$$
\mathcal{M}_g \to \overline{\mathcal{M}}_g
\quad\text{et}\quad
\overline{\mathcal{M}}_g \to \Spec(\mathbf{Z})
$$
et il suffit de vérifier ce qui suit : pour tout diagramme $2$-commutatif
$$
\xymatrix{
\Spec(K) \ar[r] \ar[d] &
\mathcal{M}_g \ar[r] &
\overline{\mathcal{M}}_g \ar[d] \\
\Spec(R) \ar[rr] \ar@{..>}[rru] & & \Spec(\mathbf{Z})
}
$$
où $R$ est un anneau de valuation discrète de corps des fractions $K$,
la catégorie des flèches pointillées est soit vide, soit un sétoïde ayant exactement
une classe d'isomorphisme. (Observons qu'il n'est pas nécessaire de trop se soucier
des $2$-flèches ; voir Morphismes de champs, lemme
\ref{stacks-morphisms-lemma-cat-dotted-arrows-independent}).
En explicitant cet énoncé et en utilisant le fait que
$\mathcal{M}_g$, resp.\ $\overline{\mathcal{M}}_g$, sont
les champs algébriques qui paramètrent les familles lisses, resp.\ stables,
de courbes de genre $g$, on constate que l'assertion à démontrer
est exactement le résultat d'unicité énoncé et démontré dans le
lemme \ref{moduli-curves-lemma-unique-stable-model}.
\end{proof}

\begin{lemma}
\label{moduli-curves-lemma-stable-quasi-compact}
Soit $g \geq 2$. Le champ $\overline{\mathcal{M}}_g$ est quasi-compact.
\end{lemma}

\begin{proof}
Nous utiliserons les notations de la section \ref{moduli-curves-section-polarized-curves}.
Considérons le sous-ensemble
$$
T \subset |\textit{PolarizedCurves}|
$$
des points $\xi$ pour lesquels il existe un corps $k$ et une paire
$(X, \mathcal{L})$ sur $k$ représentant $\xi$
et possédant les deux propriétés suivantes :
\begin{enumerate}
\item $X$ est une courbe stable de genre $g$ ;
\item $\mathcal{L} = \omega_X^{\otimes 3}$.
\end{enumerate}
Manifestement, par l'application continue
$$
|\textit{PolarizedCurves}|
\longrightarrow
|\Curvesstack|
$$
l'image de l'ensemble $T$ est exactement le sous-ensemble ouvert
$$
|\overline{\mathcal{M}}_g| \subset |\Curvesstack|
$$
Il suffit donc de montrer que $T$ est quasi-compact.
D'après le lemme \ref{moduli-curves-lemma-polarized-curves-in-polarized}, on a une
immersion ouverte et fermée
$$
|\textit{PolarizedCurves}| \subset |\Polarizedstack|
$$
Il suffit donc de démontrer la quasi-compacité de $T$ comme sous-ensemble de
$|\Polarizedstack|$. Nous appliquons pour cela le critère de
Champs de modules, lemme \ref{moduli-lemma-bounded-polarized}.
Observons d'abord que, pour $(X, \mathcal{L})$
comme ci-dessus, le polynôme de Hilbert $P$ est la fonction
$P(t) = (6g - 6)t + (1 - g)$ d'après Riemann--Roch ; voir
Courbes algébriques, lemme \ref{curves-lemma-rr}.
Observons ensuite que $H^1(X, \mathcal{L}) = 0$
et que $\mathcal{L}$ est très ample d'après
Courbes algébriques, lemme \ref{curves-lemma-tricanonical}.
Cela signifie exactement que, pour $n = P(3) - 1$,
il existe une immersion fermée
$$
i : X \longrightarrow \mathbf{P}^n_k
$$
telle que $\mathcal{L} = i^*\mathcal{O}_{\mathbf{P}^1_k}(1)$,
comme souhaité.
\end{proof}

\noindent
Voici le théorème principal de cette section.

\begin{theorem}
\label{moduli-curves-theorem-stable-smooth-proper}
Soit $g \geq 2$. Le champ algébrique $\overline{\mathcal{M}}_g$ est un
champ de Deligne-Mumford, propre et lisse sur $\Spec(\mathbf{Z})$.
De plus, le lieu $\mathcal{M}_g$ qui paramètre les courbes lisses
est un sous-champ ouvert dense.
\end{theorem}

\begin{proof}
La plupart des propriétés mentionnées dans l'énoncé ont déjà été démontrées.
La lissité résulte du lemme \ref{moduli-curves-lemma-stable-curves-smooth}.
La propriété de Deligne-Mumford résulte du lemme \ref{moduli-curves-lemma-stable-curves-deligne-mumford}.
L'ouverture de $\mathcal{M}_g$ résulte du lemme \ref{moduli-curves-lemma-smooth-dense-in-stable}.
Nous savons que $\overline{\mathcal{M}}_g \to \Spec(\mathbf{Z})$
est séparé d'après le lemme \ref{moduli-curves-lemma-stable-separated} et que
$\overline{\mathcal{M}}_g$ est quasi-compact d'après le
lemme \ref{moduli-curves-lemma-stable-quasi-compact}.
Ainsi, pour montrer que $\overline{\mathcal{M}}_g \to \Spec(\mathbf{Z})$
est propre et achever la démonstration, on peut appliquer
Compléments sur les morphismes de champs, lemme
\ref{stacks-more-morphisms-lemma-refined-valuative-criterion-proper},
aux morphismes $\mathcal{M}_g \to \overline{\mathcal{M}}_g$ et
$\overline{\mathcal{M}}_g \to \Spec(\mathbf{Z})$.
Il suffit donc de vérifier ce qui suit : pour tout diagramme $2$-commutatif
$$
\xymatrix{
\Spec(K) \ar[r] \ar[d]_j &
\mathcal{M}_g \ar[r] &
\overline{\mathcal{M}}_g \ar[d] \\
\Spec(A) \ar[rr] & & \Spec(\mathbf{Z})
}
$$
où $A$ est un anneau de valuation discrète de corps des fractions $K$, il
existe une extension de corps $K'/K$ et un anneau de valuation $A' \subset K'$
dominant $A$, tels que la catégorie des flèches pointillées du
diagramme induit
$$
\xymatrix{
\Spec(K') \ar[r] \ar[d]_{j'} & \overline{\mathcal{M}}_g \ar[d] \\
\Spec(A') \ar[r] \ar@{..>}[ru] & \Spec(\mathbf{Z})
}
$$
soit non vide (Morphismes de champs, définition
\ref{stacks-morphisms-definition-fill-in-diagram}).
(Observons qu'il n'est pas nécessaire de trop se soucier des
$2$-flèches ; voir Morphismes de champs, lemme
\ref{stacks-morphisms-lemma-cat-dotted-arrows-independent}).
En explicitant cet énoncé et en utilisant le fait que
$\mathcal{M}_g$, resp.\ $\overline{\mathcal{M}}_g$, sont les champs algébriques
qui paramètrent les familles lisses, resp.\ stables, de courbes de genre $g$,
on constate que l'assertion à démontrer est exactement le résultat contenu
dans le théorème de réduction stable, c'est-à-dire le théorème
\ref{moduli-curves-theorem-stable-reduction}.
\end{proof}














% OK, start here.
%

\chapter{Exemples}



\phantomsection
\label{examples-section-phantom}


\section{Introduction}
\label{examples-section-introduction}

\noindent
Ce chapitre contiendra des exemples qui éclairent la théorie.





\section{Une limite vide}
\label{examples-section-empty-limit}

\noindent
Cet exemple est dû à Waterhouse, voir \cite{Waterhouse}.
Soit $S$ un ensemble non dénombrable. Pour toute partie finie
$T \subset S$, considérons l'ensemble $M_T$ des applications injectives $T \to \mathbf{N}$.
Si $T \subset T' \subset S$ sont finies, la restriction $M_{T'} \to M_T$
est surjective. Nous obtenons ainsi un système projectif indexé par
l'ensemble ordonné filtrant des parties finies de $S$,
dont les applications de transition sont surjectives.
Mais $\lim M_T = \emptyset$, car un élément de la limite
définirait une application injective $S \to \mathbf{N}$.



\section{Une limite nulle}
\label{examples-section-zero-limit}

\noindent
Soit $(S_i)_{i \in I}$ un système projectif d'ensembles non vides
indexé par un ensemble ordonné filtrant, à applications de transition surjectives,
et tel que $\lim S_i = \emptyset$, voir la section \ref{examples-section-empty-limit}.
Soit $K$ un corps et posons
$$
V_i = \bigoplus\nolimits_{s \in S_i} K
$$
Alors les applications de transition $V_i \to V_j$ sont surjectives pour $i \geq j$.
Cependant, $\lim V_i = 0$. En effet, si $v = (v_i)$ est un élément de la
limite, le support de $v_i$ serait une partie finie $T_i \subset S_i$
telle que $\lim T_i \not = \emptyset$, voir
Catégories, lemme \ref{categories-lemma-nonempty-limit}.

\medskip\noindent
Pour tout $i$, considérons l'unique application $K$-linéaire $V_i \to K$ qui envoie
chaque vecteur de base $s \in S_i$ sur $1$. Soit $W_i \subset V_i$
son noyau. Alors
$$
0 \to (W_i) \to (V_i) \to (K) \to 0
$$
est une suite exacte courte non scindée de systèmes projectifs d'espaces vectoriels
sur l'ensemble ordonné filtrant $I$. Ainsi,
$W_i$ forme un système projectif de $K$-espaces vectoriels
à applications de transition surjectives, de limite nulle et dont
$R^1\lim$ ne s'annule pas.



\section{Une limite projective non quasi-compacte d'espaces quasi-compacts}
\label{examples-section-lim-not-quasi-compact}

\noindent
Désignons par $\mathbf{N}$ l'ensemble des entiers naturels.
Pour tout entier $n$, désignons par $I_n$ l'ensemble des entiers naturels $> n$.
Définissons $T_n$ comme l'unique topologie sur $\mathbf{N}$ ayant pour base
$\{1\}, \ldots , \{n\}, I_n$.  Notons $X_n$ l'espace topologique
$(\mathbf{N}, T_n)$.  Pour tous $m < n$, l'application identité
$$
f_{n, m} : X_n \longrightarrow X_m
$$
est continue.  Il est clair que, pour $m < n < p$, la composée
$f_{p, n} \circ f_{n, m}$ est égale à $f_{p, m}$.  Ainsi, $((X_n), (f_{n,m}))$
est un système projectif d'espaces topologiques quasi-compacts.

\medskip\noindent
Soit $T$ la topologie discrète sur $\mathbf{N}$, et soit $X$
l'espace $(\mathbf{N}, T)$. Pour tout entier $n$, l'application identité
$$
f_n : X \longrightarrow X_n
$$
est continue. Nous affirmons qu'il s'agit de la limite projective du système
ci-dessus. Soit $(Y, S)$ un espace topologique quelconque. Pour tout entier $n$,
soit
$$
g_n : (Y, S) \longrightarrow (\mathbf{N}, T_n)
$$
une application continue. Supposons que, pour tous $m < n$, nous ayons
$f_{n,m} \circ g_n = g_m$, c'est-à-dire que le système $(g_n)$ soit compatible
avec le système projectif ci-dessus. En particulier, toutes les applications ensemblistes
$g_n$ sont égales à une même application ensembliste
$$
g : Y \longrightarrow \mathbf{N}.
$$
De plus, pour tout entier $n$, puisque $\{n\}$ est ouvert dans $X_n$,
$g^{-1}(\{n\}) = g_n^{-1}(\{n\})$ est également ouvert dans $Y$.
Par conséquent, l'application ensembliste $g$ est continue pour la topologie $S$ sur $Y$
et la topologie $T$ sur $\mathbf{N}$. Ainsi, $(X, (f_n))$ est la limite projective
du système projectif ci-dessus.

\medskip\noindent
Cependant, $X$ n'est manifestement pas quasi-compact, puisque le recouvrement ouvert
infini par les parties réduites à un point n'admet aucun sous-recouvrement fini.

\begin{lemma}
\label{examples-lemma-lim-not-quasi-compact}
Il existe un système projectif d'espaces topologiques quasi-compacts
indexé par $\mathbf{N}$ dont la limite projective n'est pas quasi-compacte.
\end{lemma}

\begin{proof}
Voir la discussion qui précède.
\end{proof}






\section{Le faisceau structural sur le produit fibré}
\label{examples-section-silly}

\noindent
Soient $X, Y, S, a, b, p, q, f$ comme dans l'introduction de
Catégories dérivées de schémas,
section \ref{perfect-section-kunneth}. Diagramme :
$$
\xymatrix{
& X \times_S Y \ar[ld]^p \ar[rd]_q \ar[dd]^f \\
X \ar[rd]_a & & Y \ar[ld]^b \\
& S
}
$$
Nous disposons alors d'un homomorphisme canonique
$$
can :
p^{-1}\mathcal{O}_X \otimes_{f^{-1}\mathcal{O}_S} q^{-1}\mathcal{O}_Y
\longrightarrow
\mathcal{O}_{X \times_S Y}
$$
qui n'est pas un isomorphisme en général.

\medskip\noindent
Par exemple, soient $S = \Spec(\mathbf{R})$, $X = \Spec(\mathbf{C})$ et
$Y = \Spec(\mathbf{C})$. Alors
$X \times_S Y = \Spec(\mathbf{C}) \amalg \Spec(\mathbf{C})$
est un espace discret à deux points,
et les faisceaux $p^{-1}\mathcal{O}_X$, $q^{-1}\mathcal{O}_Y$
et $f^{-1}\mathcal{O}_S$ sont les faisceaux constants de valeurs
$\mathbf{C}$, $\mathbf{C}$ et $\mathbf{R}$.
Ainsi, la source de $can$ est le faisceau constant
de valeur $\mathbf{C} \otimes_\mathbf{R} \mathbf{C}$
sur l'espace discret à deux points. Ses sections
globales sont donc de dimension $8$ comme $\mathbf{R}$-espace vectoriel,
tandis que les sections globales du but de $can$
forment $\mathbf{C} \times \mathbf{C}$,
qui est de dimension $4$ comme $\mathbf{R}$-espace vectoriel.

\medskip\noindent
Un autre exemple est le suivant. Soit $k$ un corps algébriquement
clos. Considérons
$S = \Spec(k)$, $X = \mathbf{A}^1_k$ et $Y = \mathbf{A}^1_k$.
Alors, pour tout ouvert non vide $U \subset X \times_S Y = \mathbf{A}^2_k$,
les images $p(U) \subset X = \mathbf{A}^1_k$ et $q(U) \subset \mathbf{A}^1_k$
sont ouvertes, et le lecteur vérifiera que
$$
\left(
p^{-1}\mathcal{O}_X \otimes_{f^{-1}\mathcal{O}_S} q^{-1}\mathcal{O}_Y
\right)(U) = \mathcal{O}_X(p(U)) \otimes_k \mathcal{O}_Y(q(U))
$$
Ce dernier n'est pas égal à $\mathcal{O}_{X \times_S Y}(U)$
si $U$ est le complémentaire d'une courbe irréductible $C$ dans
$X \times_S Y = \mathbf{A}^2_k$ telle que $p|_C$ et $q|_C$
soient toutes deux non constantes.

\medskip\noindent
Revenons au cas général et soit $z = (x, y, s, \mathfrak p)$
un point de $X \times_S Y$ comme dans
Schémas, lemme \ref{schemes-lemma-points-fibre-product}.
Alors, sur les fibres en $z$, l'application $can$ donne l'application
$$
can_z :
\mathcal{O}_{X, x} \otimes_{\mathcal{O}_{S, s}} \mathcal{O}_{Y, y}
\longrightarrow
\mathcal{O}_{X \times_S Y, z}
$$
Cet homomorphisme d'anneaux est plat, car le but est un anneau localisé de la
source (détails omis; indication : se ramener au cas où $X$, $Y$ et $S$
sont affines).
Notons que la source n'est pas en général un anneau local, ce qui fournit
une autre manière de voir que $can$ n'est pas un isomorphisme en général.

\medskip\noindent
Plus généralement, supposons donnés un $\mathcal{O}_X$-module $\mathcal{F}$ et un
$\mathcal{O}_Y$-module $\mathcal{G}$. Il existe alors un homomorphisme canonique
\begin{align*}
& p^{-1}\mathcal{F} \otimes_{f^{-1}\mathcal{O}_S} q^{-1}\mathcal{G} \\
& =
p^{-1}(\mathcal{F} \otimes_{\mathcal{O}_X} \mathcal{O}_X)
\otimes_{f^{-1}\mathcal{O}_S}
q^{-1}(\mathcal{O}_Y \otimes_{\mathcal{O}_Y} \mathcal{G}) \\
& =
p^{-1}\mathcal{F} \otimes_{p^{-1}\mathcal{O}_X} p^{-1}\mathcal{O}_X
\otimes_{f^{-1}\mathcal{O}_S}
q^{-1}\mathcal{O}_Y \otimes_{q^{-1}\mathcal{O}_Y} q^{-1}\mathcal{G} \\
& \xrightarrow{can}
p^{-1}\mathcal{F} \otimes_{q^{-1}\mathcal{O}_X}
\mathcal{O}_{X \times_S Y}
\otimes_{q^{-1}\mathcal{O}_Y} q^{-1}\mathcal{G}  \\
& =
p^{-1}\mathcal{F} \otimes_{q^{-1}\mathcal{O}_X}
\mathcal{O}_{X \times_S Y}
\otimes_{\mathcal{O}_{X \times_S Y}}
\mathcal{O}_{X \times_S Y}
\otimes_{q^{-1}\mathcal{O}_Y} q^{-1}\mathcal{G}  \\
& =
p^*\mathcal{F} \otimes_{\mathcal{O}_{X \times_S Y}} q^*\mathcal{G}
\end{align*}
qui est rarement un isomorphisme.





\section[Schéma connexe non intègre à anneaux locaux intègres]{Un schéma connexe non intègre dont les anneaux locaux sont intègres}
\label{examples-section-connected-locally-integral-not-integral}

\noindent
Nous donnons un exemple de schéma affine $X = \Spec(A)$ qui est
connexe, dont tous les anneaux locaux sont intègres, mais qui n'est pas intègre.
La connexité de $X$ signifie que $A$ ne possède aucun idempotent non trivial, voir
Algèbre, lemme \ref{algebra-lemma-disjoint-decomposition}.
Les anneaux locaux de $X$ sont intègres si, lorsque $fg = 0$ dans $A$, tout
point de $X$ possède un voisinage sur lequel $f$ ou $g$ s'annule.
Si $A$ n'est pas un anneau intègre, alors $X$ n'est pas intègre
(Propriétés, définition \ref{properties-definition-integral}).

\medskip\noindent
En termes imagés, la construction est la suivante : soit $X_0$ la croix
(la réunion des axes de coordonnées) dans le plan affine. Soit ensuite $X_1$
l'image inverse totale (réduite) de $X_0$ dans l'éclatement du plan ($X_1$
a trois composantes rationnelles formant une chaîne).  Éclatons ensuite la
surface obtenue aux deux points singuliers de $X_1$, et soit $X_2$
l'image inverse réduite de $X_1$ (qui a cinq composantes rationnelles), etc.
Prenons pour $X$ la limite projective. Le seul problème de cette construction
est que les éclatements ajoutent une droite projective, de sorte que $X_1$ n'est pas affine.
Corrigeons cela en recollant plutôt une droite affine (notre schéma sera donc une
partie ouverte de celui décrit ci-dessus).

\medskip\noindent
Voici une construction entièrement algébrique : pour tout $k \ge 0$, soit $A_k$
l'anneau suivant : ses éléments sont les familles de
polynômes $p_i \in \mathbf{C}[x]$, où $i = 0, \ldots, 2^k$, telles que
$p_i(1) = p_{i + 1}(0)$. Posons $X_k = \Spec(A_k)$. Remarquons que $X_k$ est
la réunion de $2^k + 1$ droites affines qui se coupent transversalement en chaîne.
Définissons un homomorphisme d'anneaux $A_k \to A_{k + 1}$ par
$$
(p_0, \ldots, p_{2^k})
\longmapsto
(p_0, p_0(1), p_1, p_1(1), \ldots, p_{2^k}),
$$
autrement dit, un polynôme sur deux est constant. Ceci identifie
$A_k$ à un sous-anneau de $A_{k + 1}$. Soit $A$ la limite inductive des $A_k$
(essentiellement, leur réunion). Posons $X = \Spec(A)$. Pour tout $k$, nous avons
une injection canonique $A_k \to A$, c'est-à-dire une application $X\to X_k$.
Chaque $A_k$ est connexe mais non intègre; il s'ensuit que $A$ est
connexe mais non intègre. Il reste à montrer que les anneaux locaux de
$A$ sont intègres.

\medskip\noindent
Soient $f, g \in A$ tels que $fg = 0$ et $x \in X$. Construisons un
voisinage de $x$ sur lequel l'un de $f$ et $g$ s'annule. Choisissons $k$
tel que $f, g \in A_{k - 1}$ (noter l'indice $k - 1$).
Soit $y$ l'image de $x$ dans $X_k$. Il suffit de montrer que $y$ possède
un voisinage sur lequel $f$ ou $g$ s'annule en tant que section de
$\mathcal{O}_{X_k}$.
Si $y$ est un point lisse de $X_k$, c'est-à-dire s'il n'appartient qu'à l'une des
$2^k + 1$ droites, c'est clair. Nous pouvons donc supposer que $y$ est l'un
des $2^k$ points singuliers; deux composantes de $X_k$ passent donc par
$y$. Cependant, sur l'une de ces deux composantes (celle d'indice impair),
$f$ et $g$ sont tous deux constants, puisque ce sont des images inverses de fonctions sur
$X_{k - 1}$. Puisque $fg = 0$ partout, $f$ ou $g$ (disons $f$)
s'annule sur l'autre composante.
Il s'ensuit que $f$ s'annule sur les deux composantes, comme voulu.






\section{Un complété non complet}
\label{examples-section-noncomplete-completion}

\noindent
Soit $R$ un anneau et soit $\mathfrak m$ un idéal maximal. Considérons le
complété
$$
R^\wedge = \lim R/\mathfrak m^n.
$$
Remarquons que $R^\wedge$ est un anneau local d'idéal maximal
$\mathfrak m' = \Ker(R^\wedge \to R/\mathfrak m)$.
En effet, si $x = (x_n) \in R^\wedge$ n'appartient pas à $\mathfrak m'$, alors
$y = (x_n^{-1}) \in R^\wedge$ vérifie $xy = 1$; ainsi, $R^\wedge$ est local d'après
Algèbre, lemme \ref{algebra-lemma-characterize-local-ring}. Or il est
toujours vrai que $R^\wedge$ est complet pour sa topologie de limite projective (voir la
discussion dans
Compléments d'algèbre, section \ref{more-algebra-section-topological-ring}).
Mais au-delà se posent les questions suivantes :
\begin{enumerate}
\item A-t-on $\mathfrak m R^\wedge = \mathfrak m'$ ?
\item $R^\wedge$, considéré comme un $R^\wedge$-module, est-il
complet pour la topologie $\mathfrak m'$-adique ?
\item $R^\wedge$, considéré comme un $R$-module, est-il complet pour la topologie $\mathfrak m$-adique ?
\end{enumerate}
Il se trouve que toutes ces questions ont une réponse négative.
L'exemple ci-dessous est tiré d'une note non publiée de
Bart de Smit et Hendrik Lenstra. Voir aussi
\cite[Exercise III.2.12]{Bourbaki-CA} et
\cite[exemple 1.8]{Yekutieli}

\medskip\noindent
Soient $k$ un corps, $R = k[x_1, x_2, x_3, \ldots]$ et
$\mathfrak m = (x_1, x_2, x_3, \ldots)$.
Nous considérerons un élément $f$ de $R^\wedge$ comme une somme éventuellement infinie
$$
f = \sum a_I x^I
$$
(en notation multi-indice) telle que, pour tout $d \geq 0$, il
n'y ait qu'un nombre fini de coefficients $a_I$ non nuls pour $|I| = d$. L'idéal
maximal $\mathfrak m' \subset R^\wedge$ est l'ensemble des $f$ dont le
terme constant est nul. En particulier, l'élément
$$
f = x_1 + x_2^2 + x_3^3 + \ldots
$$
appartient à $\mathfrak m'$ mais non à $\mathfrak m R^\wedge$, ce qui
montre que (1) est fausse dans cet exemple. Cependant, si (1) est
fausse, alors (3) l'est nécessairement car
$\mathfrak m' = \Ker(R^\wedge \to R/\mathfrak m)$
et nous pouvons appliquer
Algèbre, lemme \ref{algebra-lemma-hathat} avec $n = 1$.

\medskip\noindent
Pour terminer, démontrons que $R^\wedge$ n'est pas complet pour la topologie $\mathfrak m'$-adique.
Pour $n \geq 1$, posons $K_n = \Ker(R^\wedge \to R/\mathfrak m^n)$. Alors
nous avons des suites exactes courtes
$$
0 \to K_n/(\mathfrak m')^n \to R^\wedge/(\mathfrak m')^n \to
R/\mathfrak m^n \to 0
$$
L'application de projection $R^\wedge \to R/\mathfrak m^{n + 1}$ envoie
$(\mathfrak m')^n$ sur $\mathfrak m^n/\mathfrak m^{n + 1}$.
Il s'ensuit que $K_{n + 1} \to K_n/(\mathfrak m')^n$ est surjective.
Par conséquent, le système projectif $\left(K_n/(\mathfrak m')^n\right)$
a des applications de transition surjectives et,
en passant aux limites projectives, nous obtenons une suite exacte
$$
0 \to \lim K_n/(\mathfrak m')^n \to
\lim R^\wedge/(\mathfrak m')^n \to
\lim R/\mathfrak m^n \to 0
$$
d'après Algèbre, lemme \ref{algebra-lemma-Mittag-Leffler}.
Nous voyons ainsi que $R^\wedge$ est complet pour la topologie $\mathfrak m'$-adique
si et seulement si $K_n = (\mathfrak m')^n$ pour tout $n \geq 1$.

\medskip\noindent
Pour montrer que $R^\wedge$ n'est pas complet pour la topologie $\mathfrak m'$-adique
dans notre exemple, montrons que $K_2 = \Ker(R^\wedge \to R/\mathfrak m^2)$
n'est pas égal à $(\mathfrak m')^2$.
Remarquons qu'un élément de $(\mathfrak m')^2$
peut s'écrire comme une somme finie
\begin{equation}
\label{examples-equation-sum}
\sum\nolimits_{i = 1, \ldots, t} f_i g_i
\end{equation}
où $f_i, g_i \in R^\wedge$ ont un terme constant nul.
Pour obtenir un exemple, nous allons choisir un $z \in K_2$
de la forme
$$
z = z_1 + z_2 + z_3 + \ldots
$$
ayant les propriétés suivantes
\begin{enumerate}
\item il existe des suites $1 < d_1 < d_2 < d_3 < \ldots $ et
$0 < n_1 < n_2 < n_3 < \ldots$ telles que
$z_i \in k[x_{n_i}, x_{n_i + 1}, \ldots, x_{n_{i + 1} - 1}]$
soit homogène de degré $d_i$, et
\item dans l'anneau $k[[x_{n_i}, x_{n_i + 1}, \ldots, x_{n_{i + 1} - 1}]]$,
l'élément $z_i$ ne peut s'écrire comme une somme (\ref{examples-equation-sum})
avec $t \leq i$.
\end{enumerate}
Cela implique manifestement que $z$ n'appartient pas à $(\mathfrak m')^2$,
car l'image de la relation (\ref{examples-equation-sum}) dans
l'anneau $k[[x_{n_i}, x_{n_i + 1}, \ldots, x_{n_{i + 1} - 1}]]$
pour $i$ assez grand conduirait à une contradiction. Il suffit donc
de montrer que, pour tout $t > 0$, il existe un $d \gg 0$ et un entier
$n$ tels que nous puissions trouver un élément homogène
$z \in k[x_1, \ldots, x_n]$ de degré $d$ qui ne puisse s'écrire comme
une somme (\ref{examples-equation-sum}) pour ce $t$ dans $k[[x_1, \ldots, x_n]]$.
Prenons $n > 2t$ et un $d > 1$ quelconque, premier à la caractéristique de $k$, et
posons $z = \sum_{i = 1, \ldots, n} x_i^d$. Alors l'ensemble des zéros
de l'idéal
$$
(\frac{\partial z}{\partial x_1}, \ldots, \frac{\partial z}{\partial x_n})
=
(dx_1^{d - 1}, \ldots, dx_n^{d - 1})
$$
se réduit à un point. D'autre part,
$$
\frac{\partial ( \sum\nolimits_{i = 1, \ldots, t} f_i g_i ) }{\partial x_j}
\in (f_1, \ldots, f_t, g_1, \ldots, g_t)
$$
par la règle de Leibniz; par conséquent, l'ensemble des zéros de ces dérivées
contient au moins
$$
V(f_1, \ldots, f_t, g_1, \ldots, g_t) \subset
\Spec(k[[x_1, \ldots, x_n]]).
$$
On obtient ainsi une contradiction, puisque la dimension de
$V(f_1, \ldots, f_t, g_1, \ldots, g_t)$ est au moins $n - 2t \geq 1$.

\begin{lemma}
\label{examples-lemma-noncomplete-completion}
Il existe un anneau local $R$ et un idéal maximal $\mathfrak m$ tels que
le complété $R^\wedge$ de $R$ pour la topologie $\mathfrak m$-adique ait les
propriétés suivantes :
\begin{enumerate}
\item $R^\wedge$ est local, mais son idéal maximal n'est pas égal à
$\mathfrak m R^\wedge$,
\item $R^\wedge$ n'est pas un anneau local complet, et
\item $R^\wedge$ n'est pas complet pour la topologie $\mathfrak m$-adique comme $R$-module.
\end{enumerate}
\end{lemma}

\begin{proof}
Cela résulte de la discussion qui précède, puisque (avec $R = k[x_1, x_2, x_3, \ldots]$)
le complété du localisé $R_{\mathfrak m}$ est égal au
complété de $R$.
\end{proof}





\section{Un quotient non complet}
\label{examples-section-noncomplete-quotient}

\noindent
Soit $k$ un corps. Posons
$$
R = k[t, z_1, z_2, z_3, \ldots, w_1, w_2, w_3, \ldots, x]/
(z_it - x^iw_i, z_i w_j)
$$
Remarquons qu'en particulier $z_iz_jt = 0$ dans cet anneau. Tout élément $f$ de $R$
s'écrit de manière unique comme une somme finie
$$
f = \sum\nolimits_{i = 0, \ldots, d} f_i x^i
$$
où chaque $f_i \in k[t, z_i, w_j]$ ne comporte aucun terme contenant les produits
$z_it$ ou $z_iw_j$. De plus, si $f$ s'écrit de cette manière, alors
$f \in (x^n)$ si et seulement si $f_i = 0$ pour $i < n$.
Ainsi, $x$ est un non-diviseur de zéro et $\bigcap (x^n) = 0$.
Soit $R^\wedge$ le complété de $R$ pour la topologie définie par l'idéal $(x)$.
Remarquons que $R^\wedge$ est complet pour la topologie $(x)$-adique; voir
Algèbre, lemme \ref{algebra-lemma-hathat-finitely-generated}.
D'après ce qui précède, un élément de $R^\wedge$ s'écrit de manière unique
comme une somme infinie
$$
f = \sum\nolimits_{i = 0}^\infty f_i x^i
$$
où chaque $f_i \in k[t, z_i, w_j]$ ne comporte aucun terme contenant les produits
$z_it$ ou $z_iw_j$. Considérons l'élément
$$
f = \sum\nolimits_{i = 1}^\infty x^i w_i =
xw_1 + x^2w_2 + x^3w_3 + \ldots
$$
c'est-à-dire que $f_n = w_n$. Remarquons que $f \in (t , x^n)$ pour tout $n$
car $x^mw_m \in (t)$ pour tout $m$.
Nous affirmons que $f \not \in (t)$. Pour le montrer, supposons que
$tg = f$, où $g = \sum g_lx^l$ est sous la forme canonique ci-dessus.
Comme $tz_iz_j = 0$, nous pouvons supposer que les $g_l$ ne
comportent aucun terme contenant les produits $z_iz_j$. En examinant comment
$tg$ est mis sous forme canonique, nous obtenons les règles suivantes :
Considérons un terme $c m$ quelconque de $g_l$, où $c \in k$ et où $m$ est un
monôme en $t, z_i, w_j$; effectuons la distinction suivante :
\begin{enumerate}
\item si le monôme $m$ ne fait intervenir aucun $z_i$, alors $ctm$ est
un terme de $f_l$, et
\item si le monôme $m$ fait intervenir un $z_i$, alors il est égal à
$m = z_i$ et nous voyons que $cw_i$ est un terme de $f_{l + i}$.
\end{enumerate}
Comme $g_0$ est un polynôme, les variables $z_i$ qui y apparaissent sont en nombre fini.
Choisissons $n$ tel que $z_n$ n'apparaisse pas dans $g_0$.
Alors les règles ci-dessus montrent que $w_n$ n'apparaît pas dans $f_n$, ce qui est
une contradiction. Il s'ensuit que $R^\wedge/(t)$ n'est pas complet; voir
Algèbre, lemme \ref{algebra-lemma-quotient-complete}.

\begin{lemma}
\label{examples-lemma-noncomplete-quotient}
Il existe un anneau $R$ complet pour la topologie définie par un idéal principal
$I$ et un idéal principal $J$ tels que $R/J$ ne soit pas complet pour la topologie
$I$-adique.
\end{lemma}

\begin{proof}
Voir la discussion qui précède.
\end{proof}



\section{La complétion n'est pas exacte}
\label{examples-section-completion-not-exact}

\noindent
Voici un exemple simple. Supposons que $R = k[t]$. Notons $M^\wedge$
le complété d'un $R$-module pour la topologie définie par $I = (t)$. Posons
$P = K = \bigoplus_{n \in \mathbf{N}} R$ et
$M = \bigoplus_{n \in \mathbf{N}} R/(t^n)$. Il existe alors une suite exacte courte
$0 \to K \to P \to M \to 0$ dont la première application est donnée par la
multiplication par $t^n$ sur le $n$-ième facteur direct. Nous affirmons que
$0 \to K^\wedge \to P^\wedge \to M^\wedge \to 0$ n'est pas exacte au terme médian.
En effet, $\xi = (t^2, t^3, t^4, \ldots) \in P^\wedge$ a une image nulle dans
$M^\wedge$, mais n'appartient pas à l'image de $K^\wedge \to P^\wedge$, car
il devrait être l'image de $(t, t, t, \ldots)$, qui n'est pas un élément de
$K^\wedge$.

\medskip\noindent
Voici un exemple « plus petit ». Dans la situation du
lemme \ref{examples-lemma-noncomplete-quotient},
la suite exacte courte $0 \to J \to R \to R/J \to 0$ ne reste pas
exacte après complétion pour la topologie $I$-adique.
En effet, si $f \in J$ est un générateur, alors
$f : R \to J$ est surjective, donc $R \to J^\wedge$ est surjective, et donc
l'image de $J^\wedge \to R$ est $(f) = J$; mais le fait que
$R/J$ ne soit pas complet signifie que le noyau de la surjection
$R \to (R/J)^\wedge$ contient strictement $J$; voir
Algèbre, lemmes \ref{algebra-lemma-completion-generalities} et
\ref{algebra-lemma-quotient-complete}.
De même, la suite
$R \to R \to R/(f) \to 0$ ne reste pas exacte après complétion.

\begin{lemma}
\label{examples-lemma-completion-not-exact}
\begin{slogan}
En général, la complétion n'est exacte ni à gauche ni à droite.
\end{slogan}
En général, la complétion n'est pas un foncteur exact; elle n'est même pas
exacte à droite. Cela vaut encore lorsque $I$ est de type fini,
sur la catégorie des modules de présentation finie.
\end{lemma}

\begin{proof}
Voir la discussion qui précède.
\end{proof}



\section{La catégorie des modules complets n'est pas abélienne}
\label{examples-section-non-abelian}

\noindent
Soient $R$ un anneau et $I \subset R$ un idéal de type fini.
Considérons la catégorie $\mathcal{A}$ des modules complets pour la topologie $I$-adique
qui sont des $R$-modules; voir
Algèbre, définition \ref{algebra-definition-complete}.
Soit $\varphi : M \to N$ un morphisme de $\mathcal{A}$.
Le conoyau de $\varphi$ dans $\mathcal{A}$ est le complété $I$-adique
$(\Coker(\varphi))^\wedge$ du conoyau usuel
(comme $I$ est de type fini, ce complété est complet; voir
Algèbre, lemme \ref{algebra-lemma-hathat-finitely-generated}).
Posons $K = \Ker(\varphi)$; comme $\varphi$ est continue, c'est
un sous-module fermé de $M$. Le lemme \ref{examples-lemma-closed-is-complete}
ci-dessous montre que $K$ est complet pour la topologie $I$-adique et est donc le
noyau de $\varphi$ dans $\mathcal{A}$.

\begin{lemma}
\label{examples-lemma-closed-is-complete}
Soit $I \subset R$ un idéal d'un anneau. Soit $M$ un module complet
pour la topologie $I$-adique sur $R$ et soit $N \subset M$ un sous-module. Les conditions suivantes sont équivalentes :
\begin{enumerate}
\item $N$ est fermé dans $M$,
\item $N = \bigcap_{n \geq 1} (N + I^n M)$,
\item $M/N$ est complet pour la topologie $I$-adique.
\end{enumerate}
Si $I$ est de type fini, ces conditions impliquent que $N$
est complet pour la topologie $I$-adique.
\end{lemma}

\begin{proof}
L'équivalence de (1) et (2) résulte de
Espaces formels, lemme \ref{formal-spaces-lemma-closed}.
L'équivalence de (2) et (3) est le
lemme \ref{algebra-lemma-quotient-complete} d'Algèbre.
Supposons $I$ de type fini et (1) satisfaite.
Notons $N^\wedge$ le complété $I$-adique de $N$. Comme $M$
est complet pour la topologie $I$-adique, l'application d'inclusion se factorise en
$$
N \to N^\wedge \to M
$$
Comme $N^\wedge \to M$ est continue, que $N$ est dense dans $N^\wedge$
et que $N$ est fermé dans $M$, nous voyons que
$N^\wedge \to M$ se factorise par $N$. Ainsi $N^\wedge = N \oplus C$
pour un certain $R$-module $C$. Ainsi $N$ est un facteur direct du
module complet pour la topologie $I$-adique qu'est $N^\wedge$; voir
Algèbre, lemme \ref{algebra-lemma-hathat-finitely-generated}
(c'est ici que nous utilisons le fait que $I$ est de type fini),
et donc $N$ est complet pour la topologie $I$-adique.
\end{proof}

\noindent
Nous allons donner un exemple montrant qu'en général $\Im \not = \Coim$ dans la
catégorie $\mathcal{A}$. Prenons
$R = \mathbf{Z}_p = \lim_n \mathbf{Z}/p^n\mathbf{Z}$
pour l'anneau des entiers $p$-adiques et prenons $I = (p)$.
Considérons l'application
$$
\text{diag}(1, p, p^2, \ldots) :
\left(\bigoplus\nolimits_{n \geq 1} \mathbf{Z}_p\right)^\wedge
\longrightarrow
\prod\nolimits_{n \geq 1} \mathbf{Z}_p
$$
où le membre de gauche est le complété $p$-adique de la somme directe.
Un élément du membre de gauche est donc un vecteur $(x_1, x_2, x_3, \ldots)$
tel que $x_i \in \mathbf{Z}_p$ et que sa valuation $p$-adique vérifie $v_p(x_i) \to \infty$ lorsque
$i \to \infty$. Son image est $(x_1, px_2, p^2x_3, \ldots)$. Nous voyons donc
que $(1, p, p^2, \ldots)$ appartient à l'adhérence de l'image, mais pas à
l'image. D'après la description des noyaux et conoyaux donnée ci-dessus, il est
clair que $\Im \not = \Coim$ pour cette application.

\begin{lemma}
\label{examples-lemma-complete-modules-not-abelian}
Soient $R$ un anneau et $I \subset R$ un idéal de type fini.
La catégorie des modules complets pour la topologie $I$-adique qui sont des $R$-modules possède des noyaux et des
conoyaux, mais elle n'est pas abélienne en général, même lorsque $R$ est noethérien.
\end{lemma}

\begin{proof}
Voir ci-dessus.
\end{proof}





\section{La catégorie des modules complets au sens dérivé}
\label{examples-section-derived-complete-modules}

\noindent
Lire d'abord Compléments d'algèbre, section
\ref{more-algebra-section-derived-complete-modules}
avant d'aborder la présente section.

\medskip\noindent
Soient $A$ un anneau et $I$ un idéal de $A$, et notons $\mathcal{C}$
la catégorie des modules complets au sens dérivé définie dans
Compléments d'algèbre, définition \ref{more-algebra-definition-derived-complete}.

\medskip\noindent
Soient $T$ un ensemble et $M_t$, $t \in T$, une famille de modules complets
au sens dérivé. Nous affirmons qu'en général $\bigoplus M_t$ n'est pas un module complet
au sens dérivé. Pour un exemple précis, prenons $A = \mathbf{Z}_p$ et $I = (p)$ et
considérons $\bigoplus_{n \in \mathbf{N}} \mathbf{Z}_p$. L'application de
$\bigoplus_{n \in \mathbf{N}} \mathbf{Z}_p$ dans son complété $p$-adique
n'est pas surjective. Cela signifie que $\bigoplus_{n \in \mathbf{N}} \mathbf{Z}_p$
ne peut pas être complet au sens dérivé, car le contraire entraînerait sa surjectivité; voir
Compléments d'algèbre, lemme \ref{more-algebra-lemma-complete-derived-complete}.
Par conséquent, le foncteur d'inclusion $\mathcal{C} \to \text{Mod}_A$ ne commute
ni aux sommes directes ni aux limites inductives (filtrantes).

\medskip\noindent
Supposons $I$ de type fini. D'après la discussion de Compléments d'algèbre, section
\ref{more-algebra-section-derived-complete-modules}, la catégorie
$\mathcal{C}$ possède des limites inductives quelconques. Cependant, nous affirmons que les limites inductives filtrantes
ne sont pas exactes dans la catégorie $\mathcal{C}$. En effet, supposons que
$A = \mathbf{Z}_p$ et $I = (p)$. On a des inclusions
$f_n : \mathbf{Z}_p/p\mathbf{Z}_p \to \mathbf{Z}_p/p^n\mathbf{Z}_p$
de modules complets pour la topologie $p$-adique sur $A$, données par la multiplication par
$p^{n - 1}$. On a des diagrammes commutatifs
$$
\xymatrix{
\mathbf{Z}_p/p\mathbf{Z}_p \ar[r]_{f_n} \ar[d]^1 &
\mathbf{Z}_p/p^n\mathbf{Z}_p \ar[d]_p \\
\mathbf{Z}_p/p\mathbf{Z}_p \ar[r]^{f_{n + 1}} &
\mathbf{Z}_p/p^{n + 1}\mathbf{Z}_p
}
$$
Nous affirmons que la limite inductive de ces inclusions dans la catégorie $\mathcal{C}$
donne l'application $\mathbf{Z}_p/p\mathbf{Z}_p \to 0$. En effet, la limite inductive dans
$\text{Mod}_A$ du système de droite est $\mathbf{Q}_p/\mathbf{Z}_p$.
La limite inductive dans $\mathcal{C}$ est donc
$$
H^0((\mathbf{Q}_p/\mathbf{Z}_p)^\wedge) =
H^0(\mathbf{Z}_p[1]) = 0
$$
d'après Compléments d'algèbre, section
\ref{more-algebra-section-derived-complete-modules}
où ${}^\wedge$ désigne la complétion dérivée.
Cela démontre notre affirmation.

\begin{lemma}
\label{examples-lemma-derived-complete-modules}
Soient $A$ un anneau et $I \subset A$ un idéal.
La catégorie $\mathcal{C}$ des modules complets au sens dérivé
est abélienne, et le foncteur d'inclusion $F : \mathcal{C} \to \text{Mod}_A$
est exact et commute aux limites projectives quelconques.
Si $I$ est de type fini, alors $\mathcal{C}$ possède
des sommes directes quelconques et des limites inductives quelconques, mais $F$ ne commute pas à celles-ci
en général. Enfin, les limites inductives filtrantes ne sont pas exactes dans $\mathcal{C}$
en général; ainsi, $\mathcal{C}$ n'est pas une catégorie abélienne de Grothendieck.
\end{lemma}

\begin{proof}
Voir Compléments d'algèbre, lemme \ref{more-algebra-lemma-derived-complete-modules},
et la discussion qui précède.
\end{proof}





\section{Des complétés non plats}
\label{examples-section-nonflat}

\noindent
Le complété d'un anneau pour un idéal n'est pas toujours plat,
contrairement au cas noethérien. Nous avons vu deux exemples de ce
phénomène dans
Compléments d'algèbre, exemple \ref{more-algebra-example-not-glueing-pair}.
Nous donnons dans cette section deux exemples supplémentaires.

\begin{lemma}
\label{examples-lemma-countable-fg-tensor}
Soit $R$ un anneau. Soit $M$ un $R$-module dénombrable.
Alors $M$ est un $R$-module de type fini si et seulement si
$M \otimes_R R^\mathbf{N} \to M^\mathbf{N}$ est surjective.
\end{lemma}

\begin{proof}
Si $M$ est un module de type fini, l'application est surjective d'après Algèbre, proposition
\ref{algebra-proposition-fg-tensor}. Réciproquement, supposons l'application surjective.
Soit $m_1, m_2, m_3, \ldots$ une énumération des éléments de $M$.
Soit $\sum_{j = 1, \ldots, m} x_j \otimes a_j$ un élément du
produit tensoriel dont l'image est $(m_n) \in M^\mathbf{N}$. Alors
nous voyons que $x_1, \ldots, x_m$ engendrent $M$ sur $R$, comme dans la démonstration de
Algèbre, proposition \ref{algebra-proposition-fg-tensor}.
\end{proof}

\begin{lemma}
\label{examples-lemma-countable-fp-tensor}
Soit $R$ un anneau dénombrable. Soit $M$ un $R$-module dénombrable. Alors $M$
est de présentation finie si et seulement si l'homomorphisme canonique
$M \otimes_R R^\mathbf{N} \to M^\mathbf{N}$ est un isomorphisme.
\end{lemma}

\begin{proof}
Si $M$ est un module de présentation finie, l'application est un isomorphisme
d'après Algèbre, proposition \ref{algebra-proposition-fp-tensor}. Réciproquement,
supposons que l'application soit un isomorphisme. Le lemme \ref{examples-lemma-countable-fg-tensor}
montre que le module $M$ est de type fini. Choisissons une surjection $R^{\oplus m} \to M$ de
noyau $K$. Alors $K$ est dénombrable comme sous-module de $R^{\oplus m}$.
En raisonnant comme dans la démonstration d'Algèbre, proposition
\ref{algebra-proposition-fp-tensor}, nous voyons que
$K \otimes_R R^\mathbf{N} \to K^\mathbf{N}$ est surjective.
Nous concluons donc que $K$ est un $R$-module de type fini par le
lemme \ref{examples-lemma-countable-fg-tensor}.
Ainsi $M$ est de présentation finie.
\end{proof}

\begin{lemma}
\label{examples-lemma-countable-coherent}
Soit $R$ un anneau dénombrable. Alors $R$ est cohérent si et seulement si
$R^\mathbf{N}$ est un $R$-module plat.
\end{lemma}

\begin{proof}
Si $R$ est cohérent, alors $R^\mathbf{N}$ est un module plat d'après
Algèbre, proposition \ref{algebra-proposition-characterize-coherent}.
Supposons $R^\mathbf{N}$ plat. Soit $I \subset R$ un idéal de type
fini. Pour démontrer le lemme, nous montrons que $I$ est de
présentation finie comme $R$-module. En effet, l'application
$I \otimes_R R^\mathbf{N} \to R^\mathbf{N}$ est
injective puisque $R^\mathbf{N}$ est plat et que son image est
$I^\mathbf{N}$ d'après le lemme \ref{examples-lemma-countable-fg-tensor}.
Nous concluons donc par le lemme \ref{examples-lemma-countable-fp-tensor}.
\end{proof}

\noindent
Soit $R$ un anneau dénombrable. Remarquons que $R[[x]]$ est isomorphe à
$R^\mathbf{N}$ comme $R$-module. Le lemme \ref{examples-lemma-countable-coherent}
montre que $R \to R[[x]]$ est plat si et seulement si $R$ est cohérent.
Il existe de nombreux anneaux dénombrables non cohérents, par exemple
$$
R = k[y, z, a_1, b_1, a_2, b_2, a_3, b_3, \ldots]/
(a_1 y + b_1 z, a_2 y + b_2 z, a_3 y + b_3 z, \ldots)
$$
où $k$ est un corps dénombrable. Cet anneau n'est pas cohérent, car
l'idéal $(y, z)$ de $R$ n'est pas un $R$-module de présentation finie.
Notons que $R[[x]]$ est le complété de $R[x]$ pour l'idéal
principal $(x)$.

\begin{lemma}
\label{examples-lemma-completion-polynomial-ring-not-flat}
Il existe un anneau tel que le complété $R[[x]]$ de $R[x]$
en $(x)$ ne soit pas plat sur $R$ et, a fortiori, pas plat sur $R[x]$.
\end{lemma}

\begin{proof}
Voir la discussion qui précède.
\end{proof}

\noindent
Il existe en fait un anneau $R$ tel que $R[[x]]$ soit plat
sur $R$, mais que $R[[x]]$ ne soit pas plat sur $R[x]$. Voir
\href{https://math.stackexchange.com/users/164860/badam-baplan}{ce billet}
de Badam Baplan. En effet, soit
$R$ un anneau de valuation. Alors $R$ est cohérent (Algèbre, exemple
\ref{algebra-example-valuation-ring-coherent}) et donc
$R[[x]]$ est plat sur $R$ d'après
Algèbre, proposition \ref{algebra-proposition-characterize-coherent}.
En revanche, nous avons le lemme suivant.

\begin{lemma}
\label{examples-lemma-almost-integral-when-powerseries-flat}
Soit $R$ un anneau intègre de corps des fractions $K$.
Si $R[[x]]$ est plat sur $R[x]$, alors $R$ est normal si et seulement
si $R$ est complètement intégralement clos
(Algèbre, définition \ref{algebra-definition-almost-integral}).
\end{lemma}

\begin{proof}
Supposons donnés $\alpha \in K$ et un élément non nul $r \in R$ tels que
$r \alpha^n \in R$ pour tout $n \geq 1$. Considérons alors
$f = \sum r \alpha^n x^n$ dans $R[[x]]$. Écrivons $\alpha = a/b$
avec $a, b \in R$ et $b$ non nul. Nous voyons alors que $(a x - b)f = -rb$.
Il s'ensuit que $rb$ appartient à l'idéal $(ax - b)R[[x]]$.
Soit $S = \{h \in R[x] : h(0) = 1\}$. C'est une partie multiplicativement stable,
et la platitude de $R[x] \to R[[x]]$ implique que $S^{-1}R[x] \to R[[x]]$
est fidèlement plat (détails omis; indication : utiliser Algèbre, lemme
\ref{algebra-lemma-ff-rings}). Par conséquent,
$$
S^{-1}R[x]/(ax - b)S^{-1}R[x] \to R[[x]]/(ax - b)R[[x]]
$$
est injective. Nous en concluons que
$h rb = (ax - b) g$ pour certains $h \in S$ et $g \in R[x]$.
En écrivant $h = 1 + h_1 x + \ldots + h_d x^d$, nous obtenons
$$
1 + h_1 x + \ldots + h_d x^d = (1/r)(\alpha x - 1)g
$$
Cette factorisation dans $K[x]$ donne une factorisation correspondante
dans $K[x^{-1}]$, qui montre que $\alpha$ est racine d'un polynôme unitaire
à coefficients dans $R$, comme voulu.
\end{proof}

\begin{lemma}
\label{examples-lemma-completion-polynomial-ring-not-flat-bis}
Si $R$ est un anneau de valuation de dimension $> 1$, alors $R[[x]]$
est plat sur $R$, mais pas sur $R[x]$.
\end{lemma}

\begin{proof}
Les arguments ci-dessus montrent que, pour cela, il suffit d'établir que
$R$ n'est pas complètement intégralement clos (les anneaux de valuation sont
normaux; voir Algèbre, lemme \ref{algebra-lemma-valuation-ring-normal}).
Soit $\mathfrak p \subset \mathfrak m \subset R$ une chaîne d'idéaux premiers.
Choisissons $x \in \mathfrak p$ non nul et $y \in \mathfrak m \setminus \mathfrak p$.
Alors $x y^{-n} \in R$ pour tout $n \geq 1$ (sinon $y^n/x \in R$,
ce qui est absurde puisque $y \not \in \mathfrak p$). Ainsi $1/y$ est
presque entier sur $R$, mais n'appartient pas à $R$.
\end{proof}

\noindent
Construisons maintenant un exemple où le complété d'un localisé
n'est pas plat. Pour cela, considérons l'anneau
$$
R = k[y, z, a_1, a_2, a_3, \ldots]/(ya_i, a_i a_j)
$$
Notons $f \in R$ la classe de $z$. Nous affirmons que l'homomorphisme d'anneaux
\begin{equation}
\label{examples-equation-nonflat}
R[[x]] \longrightarrow R_f[[x]]
\end{equation}
n'est pas plat. Soit $I$ le noyau de $y : R[[x]] \to R[[x]]$. Un
élément $g$ de $I$ est typiquement de la forme $g = \sum g_{n, m} a_mx^n$
où $g_{n, m} \in k[z]$ et, pour $n$ donné, seul un nombre fini de
coefficients $g_{n, m}$ sont non nuls. Soit $J$ le noyau de $y : R_f[[x]] \to R_f[[x]]$.
Nous affirmons que $J \not = I R_f[[x]]$. En effet, si c'était le cas, nous
aurions
$$
\sum z^{-n} a_n x^n = \sum\nolimits_{i = 1, \ldots, m} h_i g_i
$$
pour un certain $m \geq 1$, avec $g_i \in I$ et $h_i \in R_f[[x]]$. Écrivons
$h_i = \bar h_i \bmod (y, a_1, a_2, a_3, \ldots)$
avec $\bar h_i \in k[z, 1/z][[x]]$. En considérant le coefficient de
$a_n$ et en utilisant la description ci-dessus des éléments $g_i$, nous obtiendrions
$$
z^{-n} x^n = \sum \bar h_i \bar g_{i, n}
$$
pour certains $\bar g_{i, n} \in k[z][[x]]$. Cela signifierait que
tous les $z^{-n}x^n$ appartiennent au $k[z][[x]]$-module de type fini
engendré par les éléments $\bar h_i$. Puisque $k[z][[x]]$ est noethérien,
il en résulterait que le sous-$k[z][[x]]$-module de $k[z, 1/z][[x]]$
engendré par $1, z^{-1}x, z^{-2}x^2, \ldots$ est de type fini. D'après
Algèbre, lemme \ref{algebra-lemma-characterize-integral-element},
nous conclurions que $z^{-1}x$ est entier sur $k[z][[x]]$,
ce qui est absurde. En revanche,
si (\ref{examples-equation-nonflat}) était plat, alors nous
aurions $J = IR_f[[x]]$ en tensorisant la suite exacte
$0 \to I \to R[[x]] \xrightarrow{y} R[[x]]$ par $R_f[[x]]$.

\begin{lemma}
\label{examples-lemma-nonflat-completion-localization}
Il existe un anneau $A$ complet pour la topologie définie par un idéal principal $I$
et un élément $f \in A$ tels que le complété $I$-adique
$A_f^\wedge$ de $A_f$ ne soit pas plat sur $A$.
\end{lemma}

\begin{proof}
Posons $A = R[[x]]$ et $I = (x)$, et remarquons que $R_f[[x]]$
est le complété de $R[[x]]_f$.
\end{proof}





\section{Catégorie non abélienne des modules quasi-cohérents}
\label{examples-section-nonabelian-QCoh}

\noindent
Dans Faisceaux sur les champs, section \ref{stacks-sheaves-section-quasi-coherent},
nous avons défini la catégorie des modules quasi-cohérents sur une catégorie fibrée en
groupoïdes au-dessus de $\Sch$. Bien que nous montrions dans
Faisceaux sur les champs, section
\ref{stacks-sheaves-section-quasi-coherent-algebraic-stacks}
que cette catégorie est abélienne pour les champs algébriques, nous montrons dans cette
section que ce n'est pas le cas pour les espaces algébriques formels.

\medskip\noindent
Plus précisément, considérons $\mathbf{Z}_p$ comme un anneau topologique, muni de
la topologie $p$-adique. Posons $X = \text{Spf}(\mathbf{Z}_p)$ ; voir
Espaces formels, définition
\ref{formal-spaces-definition-affine-formal-spectrum}.
Alors $X$ est un faisceau d'ensembles sur $(\Sch/\mathbf{Z})_{fppf}$
et donne lieu à un champ en catégories discrètes $\mathcal{X}$ ; voir
Champs, lemme \ref{stacks-lemma-when-stack-in-sets}.
La discussion de Faisceaux sur les champs, section
\ref{stacks-sheaves-section-quasi-coherent-algebraic-stacks}
s'applique donc.

\medskip\noindent
Soit $\mathcal{F}$ un module quasi-cohérent sur $\mathcal{X}$.
Puisque $X = \colim \Spec(\mathbf{Z}/p^n\mathbf{Z})$, il résulte aussitôt de
Faisceaux sur les champs, lemme \ref{stacks-sheaves-lemma-quasi-coherent},
que $\mathcal{F}$ est déterminé par une suite $(\mathcal{F}_n)$ vérifiant :
\begin{enumerate}
\item $\mathcal{F}_n$ est un module quasi-cohérent sur
$\Spec(\mathbf{Z}/p^n\mathbf{Z})$, et
\item les applications de transition induisent des isomorphismes
$\mathcal{F}_n = \mathcal{F}_{n + 1}/p^n\mathcal{F}_{n + 1}$.
\end{enumerate}
En passant aux modules, nous voyons que $\mathcal{F}$ correspond à un
système $(M_n)$ où chaque $M_n$ est un groupe abélien annulé
par $p^n$ et où les applications de transition induisent des isomorphismes
$M_n = M_{n + 1}/p^n M_{n + 1}$. Dans cette situation, le module
$M = \lim M_n$ est complet pour la topologie $p$-adique et $M_n = M/p^n M$ ; voir
Algèbre, lemme \ref{algebra-lemma-limit-complete}.
Nous en déduisons que la catégorie des modules quasi-cohérents sur $X$
est équivalente à la catégorie des groupes abéliens complets pour
la topologie $p$-adique. Cette catégorie n'est pas abélienne ; voir
la section \ref{examples-section-non-abelian}.

\begin{lemma}
\label{examples-lemma-quasi-coherent-not-abelian}
La catégorie des modules quasi-cohérents\footnote{Les modules quasi-cohérents
sont ici entendus au sens défini ci-dessus. Compte tenu de la manière dont le projet Champs
est organisé, il s'agit bien de la bonne définition ; comme nous l'avons vu ci-dessus,
notre définition coïncide avec celle que l'on aurait naïvement donnée des modules quasi-cohérents
sur $\text{Spf}(A)$, à savoir les $A$-modules complets.}
sur un espace algébrique formel
$X$ n'est pas abélienne en général, même lorsque $X$ est un espace algébrique
formel affine noethérien.
\end{lemma}

\begin{proof}
Voir la discussion qui précède.
\end{proof}





\section{Suite non scindée, mais localement scindée}
\label{examples-section-nonsplit-locally-split}

\noindent
Considérons la suite exacte courte
$$
0 \to M \to
\bigoplus\nolimits_{p\text{ premier}} \mathbf{Z}_{(p)} \to \mathbf{Q} \to 0
$$
où $M$ est le noyau de l'homomorphisme de droite, défini par l'inclusion
$\mathbf{Z}_{(p)} \to \mathbf{Q}$ sur chaque facteur. Cette suite de
$\mathbf{Z}$-modules n'est pas scindée, car il n'existe aucun
homomorphisme non nul $\mathbf{Q} \to \mathbf{Z}_{(p)}$. En revanche, si
nous localisons en un idéal premier quelconque $p$, la suite devient scindée.

\begin{lemma}
\label{examples-lemma-nonsplit-locally-split}
Il existe un anneau $R$ et une suite non scindée de modules
qui devient scindée localement pour la topologie de Zariski.
\end{lemma}

\begin{proof}
Voir la discussion qui précède.
\end{proof}





\section{Suites régulières et changement de base}
\label{examples-section-regular-base-change}

\noindent
Nous allons construire un anneau $R$ muni d'une suite régulière
$(x, y, z)$ telle qu'il existe un élément non nul $\delta \in R/zR$
vérifiant $x\delta = y\delta = 0$.

\medskip\noindent
Pour construire notre exemple, nous commençons par
construire un module particulier $E$ sur l'anneau $k[x, y, z]$,
où $k$ est un corps quelconque. Plus précisément, $E$ sera une somme amalgamée, comme
dans le diagramme suivant :
$$
\xymatrix{
\frac{xk[x, y, z, y^{-1}]}{xyk[x, y, z]} \ar[r] \ar[d]^{z/x} &
\frac{k[x, y, z, x^{-1}, y^{-1}]}{yk[x, y, z, x^{-1}]} \ar[r] \ar[d] &
\frac{k[x, y, z, x^{-1}, y^{-1}]}{yk[x, y, z, x^{-1}] + xk[x, y, z, y^{-1}]}
\ar[d] \\
\frac{k[x, y, z, y^{-1}]}{yzk[x, y, z]} \ar[r] &
E \ar[r] &
\frac{k[x, y, z, x^{-1}, y^{-1}]}{yk[x, y, z, x^{-1}] + xk[x, y, z, y^{-1}]}
}
$$
dont les lignes sont des suites exactes courtes (nous avons omis les zéros extérieurs pour
des raisons typographiques). On peut aussi décrire $E$ par
$$
E = \{(f, g) \mid f \in k[x, y, z, x^{-1}, y^{-1}],
g \in k[x, y, z, y^{-1}] \}/\sim
$$
où $(f, g) \sim (f', g')$ si et seulement s'il existe un
$h \in k[x, y, z, y^{-1}]$ tel que
$$
f = f' + xh \bmod yk[x, y, z, x^{-1}], \quad
g = g' - zh \bmod yzk[x, y, z]
$$
Nous affirmons que : (a) $x : E \to E$ est injective, (b)
$y : E/xE \to E/xE$ est injective, (c) $E/(x, y)E = 0$, (d) il
existe un élément non nul $\delta \in E/zE$ tel que
$x\delta = y\delta = 0$.

\medskip\noindent
Pour prouver (a), supposons que $(f, g)$ soit un couple qui représente un
élément de $E$ et que $(xf, xg) \sim 0$. Il existe alors un
$h \in k[x, y, z, y^{-1}]$ tel que $xf + xh \in yk[x, y, z, x^{-1}]$
et $xg - zh \in yzk[x, y, z]$. Nous pouvons supposer que
$h = \sum a_{i, j, k}x^iy^jz^k$ est une somme de monômes dans laquelle seuls
les $j \leq 0$ interviennent. Alors $xg - zh \in yzk[x, y, z]$ entraîne que
seuls les $i > 0$ interviennent, c'est-à-dire que $h = xh'$ pour un certain $h' \in k[x, y, z, y^{-1}]$.
Alors $(f, g) \sim (f + xh', g - zh')$ et nous voyons que nous pouvons supposer
que $g = 0$ et $h = 0$. Dans ce cas, $xf \in yk[x, y, z, x^{-1}]$
entraîne $f \in yk[x, y, z, x^{-1}]$ et nous voyons que $(f, g) \sim 0$.
Par conséquent, $x : E \to E$ est injective.

\medskip\noindent
Puisque la multiplication par $x$ est un isomorphisme de
$\frac{k[x, y, z, x^{-1}, y^{-1}]}{yk[x, y, z, x^{-1}]}$, nous voyons que
$E/xE$ est isomorphe à
$$
\frac{k[x, y, z, y^{-1}]}{
yzk[x, y, z] + xk[x, y, z, y^{-1}] + zk[x, y, z, y^{-1}]}
=
\frac{k[x, y, z, y^{-1}]}{xk[x, y, z, y^{-1}] + zk[x, y, z, y^{-1}]}
$$
et la multiplication par $y$ est donc un isomorphisme de $E/xE$. Cela
entraîne manifestement (b) et (c).

\medskip\noindent
Soit $e \in E$ la classe d'équivalence de $(1, 0)$.
Supposons que $e \in zE$. Il existe alors $f \in k[x, y, z, x^{-1}, y^{-1}]$,
$g \in k[x, y, z, y^{-1}]$ et $h \in k[x, y, z, y^{-1}]$ tels que
$$
1 + zf + xh \in yk[x, y, z, x^{-1}], \quad
0 + zg - zh \in yzk[x, y, z].
$$
C'est impossible : le monôme $1$ ne peut apparaître dans
$zf$, ni dans $xh$. En revanche, nous avons $ye = 0$ et
$xe = (x, 0) \sim (0, z) = z(0, 1)$. Par conséquent, en prenant pour $\delta$
la classe de $e$ dans $E/zE$, nous obtenons (d).

\begin{lemma}
\label{examples-lemma-strange-regular-sequence}
Il existe un anneau local $R$ et une suite régulière $x, y, z$
(dans l'idéal maximal) tels qu'il existe un élément non nul
$\delta \in R/zR$ vérifiant $x\delta = y\delta = 0$.
\end{lemma}

\begin{proof}
Soit $R = k[x, y, z] \oplus E$, où $E$ est le module ci-dessus, considéré
comme un idéal de carré nul. Il est alors clair que $x, y, z$ est une suite régulière
dans $R$ et que l'élément $\delta \in E/zE \subset R/zR$
a les propriétés voulues. Pour obtenir un exemple local,
nous pouvons localiser $R$ en l'idéal maximal $\mathfrak m = (x, y, z, E)$.
La suite $x, y, z$ reste régulière (car la localisation est
exacte), et l'élément $\delta$ reste non nul puisque son support est réduit
à $\mathfrak m$.
\end{proof}

\begin{lemma}
\label{examples-lemma-base-change-regular-sequence}
Il existe un homomorphisme local d'anneaux locaux $A \to B$
et une suite régulière $x, y$ dans l'idéal maximal de $B$ tels que
$B/(x, y)$ soit plat sur $A$, mais que les images
$\overline{x}, \overline{y}$ de $x, y$ dans $B/\mathfrak m_AB$ ne
forment pas une suite régulière, ni même une suite régulière au sens de Koszul.
\end{lemma}

\begin{proof}
Posons $A = k[z]_{(z)}$ et $B = (k[x, y, z] \oplus E)_{(x, y, z, E)}$.
Puisque $x, y, z$ est une suite régulière dans $B$ (voir la démonstration du
lemme \ref{examples-lemma-strange-regular-sequence}),
nous voyons que $x, y$ est une suite régulière dans $B$ et que
$B/(x, y)$ est un $A$-module sans torsion, donc plat.
En revanche, il existe un élément non nul
$\delta \in B/\mathfrak m_AB = B/zB$ qui est annulé
par $\overline{x}, \overline{y}$. Par conséquent,
$H_2(K_\bullet(B/\mathfrak m_AB, \overline{x}, \overline{y})) \not = 0$.
Ainsi, la suite $\overline{x}, \overline{y}$ n'est pas régulière au sens de Koszul ; en particulier,
elle n'est pas une suite régulière ; voir
Compléments d'algèbre, lemme \ref{more-algebra-lemma-regular-koszul-regular}.
\end{proof}





\section{Un anneau noethérien de dimension infinie}
\label{examples-section-Noetherian-infinite-dimension}

\noindent
Un anneau local noethérien est de dimension finie, comme nous l'avons vu dans
Algèbre, proposition \ref{algebra-proposition-dimension}.
Il existe cependant des anneaux noethériens de dimension infinie.
Voir \cite[Appendix, exemple 1]{Nagata}.

\medskip\noindent
Plus précisément, soit $k$ un corps, et considérons l'anneau
$$
R = k[x_1, x_2, x_3, \ldots ].
$$
Soit $\mathfrak p_i = (x_{2^{i - 1}}, x_{2^{i - 1} + 1}, \ldots, x_{2^i - 1})$
pour $i = 1, 2, \ldots$ ; ce sont des idéaux premiers de $R$. Soit $S$
la partie multiplicativement stable
$$
S = \bigcap\nolimits_{i \geq 1} (R \setminus \mathfrak p_i).
$$
Considérons l'anneau $A = S^{-1}R$.
Nous affirmons que
\begin{enumerate}
\item Les idéaux maximaux de l'anneau $A$ sont les idéaux
$\mathfrak m_i = \mathfrak p_iA$.
\item Nous avons $A_{\mathfrak m_i} = R_{\mathfrak p_i}$, qui est
un anneau local noethérien de dimension $2^{i - 1}$.
\item L'anneau $A$ est noethérien.
\end{enumerate}
Il est donc clair qu'il s'agit de l'exemple recherché.
Nous omettons les détails.




\section{Anneaux locaux dont le complété n'est pas réduit}
\label{examples-section-local-completion-nonreduced}

\noindent
Dans Algèbre, exemple \ref{algebra-example-bad-dvr-char-p}, nous avons donné l'exemple
d'un anneau local noethérien intègre de caractéristique $p$, noté $R$, et de dimension $1$,
dont le complété n'est pas réduit. Nous présentons ici l'exemple
de \cite[Proposition 3.1]{Ferrand-Raynaud}, qui fournit un anneau semblable
en caractéristique zéro.

\medskip\noindent
Soit $\mathbf{C}\{x\}$ l'anneau des séries entières convergentes sur
le corps $\mathbf{C}$ des nombres complexes. L'anneau de toutes les séries entières
$\mathbf{C}[[x]]$ est son complété. Soit $K = \mathbf{C}\{x\}[1/x]$
le corps des séries de Laurent convergentes. Le $K$-module
$\Omega_{K/\mathbf{C}}$ des différentielles algébriques
de $K$ sur $\mathbf{C}$ est un $K$-espace vectoriel de dimension infinie
(démonstration omise). Nous pouvons choisir $f_n \in x\mathbf{C}\{x\}$,
$n \geq 1$, de telle sorte que
$
\text{d}x, \text{d}f_1, \text{d}f_2, \ldots
$
fassent partie d'une base de $\Omega_{K/\mathbf{C}}$. Nous pouvons donc
trouver une $\mathbf{C}$-dérivation
$$
D : \mathbf{C}\{x\} \longrightarrow \mathbf{C}((x))
$$
telle que $D(x) = 0$ et $D(f_n) = x^{-n}$. Posons
$$
A = \{f \in \mathbf{C}\{x\} \mid D(f) \in \mathbf{C}[[x]]\}
$$
Nous affirmons que
\begin{enumerate}
\item $\mathbf{C}\{x\}$ est entier sur $A$,
\item $A$ est un anneau local intègre,
\item $\dim(A) = 1$,
\item l'idéal maximal de $A$ est engendré par $x$ et $xf_1$,
\item $A$ est noethérien, et
\item le complété de $A$ est égal à l'algèbre des nombres duaux
sur $\mathbf{C}[[x]]$.
\end{enumerate}
Puisque l'algèbre des nombres duaux n'est pas réduite, l'anneau $A$ fournit l'exemple recherché.

\medskip\noindent
Remarquons que, si $0 \not = f \in x\mathbf{C}\{x\}$, alors
nous pouvons écrire $D(f) = h/f^n$ pour un certain $n \geq 0$ et $h \in \mathbf{C}[[x]]$.
Par conséquent, $D(f^{n + 1}/(n + 1)) \in \mathbf{C}[[x]]$
et $D(f^{n + 2}/(n + 2)) \in \mathbf{C}[[x]]$. Nous voyons donc que
$f^{n + 1}, f^{n + 2} \in A$ !
En particulier, l'assertion (1) est vérifiée. Nous en déduisons aussi que
le corps des fractions de $A$ est égal au corps des fractions de
$\mathbf{C}\{x\}$. Il en résulte aussi immédiatement que
$A \cap x\mathbf{C}\{x\}$ est l'ensemble des éléments non inversibles de $A$ ; ainsi,
$A$ est un anneau local intègre de dimension $1$. Si nous pouvons prouver (4),
il s'ensuivra que $A$ est noethérien (démonstration omise).
Supposons que $f \in A \cap x\mathbf{C}\{x\}$. Écrivons
$D(f) = h$, avec $h \in \mathbf{C}[[x]]$. Écrivons $h = c + xh'$
avec $c \in \mathbf{C}$ et $h' \in \mathbf{C}[[x]]$. Alors
$D(f - cxf_1) = c + xh' - c = xh'$. D'autre part,
$f - cxf_1 = xg$ avec $g \in \mathbf{C}\{x\}$, mais le
calcul ci-dessus donne $D(g) = h' \in \mathbf{C}[[x]]$
et donc $g \in A$. Ainsi, $f = cxf_1 + xg \in (x, xf_1)$, comme voulu.

\medskip\noindent
Enfin, pourquoi le complété de $A$ n'est-il pas réduit ? Notons $\hat A$ le
complété de $A$. Celui-ci admet évidemment un homomorphisme surjectif vers le complété
$\mathbf{C}[[x]]$ de $\mathbf{C}\{x\}$, puisque $x \in A$. Notons
cet homomorphisme $\psi : \hat A \to \mathbf{C}[[x]]$.
Nous avons vu ci-dessus que $\mathfrak m_A = (x, xf_1)$,
et donc $D(\mathfrak m_A^n) \subset (x^{n - 1})$ par un calcul
facile. Ainsi, $D : A \to \mathbf{C}[[x]]$ est continue et
donne lieu à une dérivation continue $\hat D : \hat A \to \mathbf{C}[[x]]$
au-dessus de $\psi$. Nous obtenons donc un homomorphisme d'anneaux
$$
\psi + \epsilon \hat D :
\hat A
\longrightarrow
\mathbf{C}[[x]][\epsilon].
$$
Comme $\hat A$ est un anneau local noethérien complet de dimension un, si nous
montrons que cette flèche est surjective, il s'ensuivra que $\hat A$
n'est pas réduit. En fait, cet homomorphisme est un isomorphisme, mais nous en omettons
la vérification. Le sous-anneau $\mathbf{C}[x]_{(x)} \subset A$
donne lieu, par complétion, à un homomorphisme $i : \mathbf{C}[[x]] \to \hat A$ tel
que $\psi \circ i = \text{id}$ et que $\hat D \circ i = 0$
(puisque $D(x) = 0$ par construction). Considérons les éléments $x^nf_n \in A$.
Nous avons
$$
(\psi + \epsilon D)(x^nf_n) = x^n f_n + \epsilon
$$
pour tout $n \geq 1$. La surjectivité résulte aussitôt de ces remarques.





\section{Un autre anneau local dont le complété n'est pas réduit}
\label{examples-section-another-local-completion-nonreduced}

\noindent
Dans cette section, nous construisons un exemple d'anneau local noethérien intègre
de dimension $2$, complet pour la topologie définie par un idéal principal,
tel que le complété d'un localisé ne soit pas réduit.

\medskip\noindent
Soit $p$ un nombre premier. Soit $k$ un corps de caractéristique $p$
tel que $k$ soit de degré infini sur son sous-corps $k^p$ des puissances $p$-ièmes.
Par exemple, $k = \mathbf{F}_p(t_1, t_2, t_3, \ldots)$.
Considérons l'anneau
$$
A =
\left\{
\begin{matrix}
\sum a_{i, j} x^iy^j \in k[[x, y]] \text{ tel que, pour tout }n \geq 0
\text{ on ait } \\
[k^p(a_{n, n}, a_{n, n + 1}, a_{n + 1, n},
a_{n, n + 2}, a_{n + 2, n}, \ldots) : k^p] < \infty
\end{matrix}
\right\}
$$
Ensemblistement, on a
$$
k^p[[x, y]] \subset A  \subset k[[x, y]]
$$
Tout élément $f$ de $A$ s'écrit de manière unique sous la forme d'une série
$$
f = f_0 + f_1 xy + f_2 (xy)^2 + f_3 (xy)^3 + \ldots
$$
où
$$
f_n = a_{n, n} + a_{n, n + 1} y + a_{n + 1, n} x +
a_{n, n + 2} y^2 + a_{n + 2, n} x^2 + \ldots
$$
et la condition de la formule qui définit $A$ signifie que les
coefficients de $f_n$ engendrent une extension finie de $k^p$.
Cette présentation montre clairement que $A$ est une
sous-$k^p[[x, y]]$-algèbre de $k[[x, y]]$, complète pour la
topologie définie par l'idéal $xy$. De plus, nous avons manifestement
$$
A/xy A = C \times_k D
$$
où $k^p[[x]] \subset C \subset k[[x]]$ et
$k^p[[y]] \subset D \subset k[[y]]$ sont les sous-anneaux
de séries formelles de
Algèbre, exemple \ref{algebra-example-bad-dvr-char-p}.
Ainsi, $C$ et $D$ sont des anneaux de valuation discrète, et nous voyons que $A/ xy A$ est
noethérien. D'après Algèbre, lemme \ref{algebra-lemma-completion-Noetherian}, nous
concluons que $A$ est noethérien. Puisque $\dim(k[[x, y]]) = 2$, le
lemme \ref{algebra-lemma-integral-sub-dim-equal} d'Algèbre
donne $\dim(A) = 2$.

\medskip\noindent
Soit $f = \sum a_i x^i$ une série formelle telle que
$k^p(a_0, a_1, a_2, \ldots)$ soit de degré infini sur $k^p$.
Alors $f \not \in A$, mais $f^p \in A$. Posons
$$
B = A[f] \subset k[[x, y]]
$$
Comme $B$ est une $A$-algèbre finie, $B$ est noethérien.
De plus, $B$ est complet pour la topologie définie par l'idéal engendré par $xy$; voir
Algèbre, lemme \ref{algebra-lemma-completion-tensor}.
En fait, $B$ est libre sur $A$, de base $1, f, f^2, \ldots, f^{p - 1}$;
nous omettons la démonstration.

\medskip\noindent
Nous affirmons que l'anneau
$$
(B_y)^\wedge = (B[1/y])^\wedge = \lim B[1/y]/(xy)^n B[1/y] =
\lim B[1/y] / x^n B[1/y]
$$
n'est pas réduit. En effet, cet anneau est libre sur
$$
(A_y)^\wedge = (A[1/y])^\wedge = \lim A[1/y]/(xy)^n A[1/y] =
\lim A[1/y] / x^n A[1/y]
$$
de base $1, f, \ldots, f^{p - 1}$. Cependant, il existe un
élément $g \in (A_y)^\wedge$ tel que $f^p = g^p$. En effet,
il suffit de prendre $g = \sum a_i x^i$ (la même expression que
celle utilisée pour $f$), qui a un sens dans $(A_y)^\wedge$.
Nous voyons donc que
$$
(B_y)^\wedge = (A_y)^\wedge[f]/(f^p - g^p) \cong
(A_y)^\wedge[\tau]/(\tau^p)
$$
n'est pas réduit. En fait, cet exemple montre un peu plus.
En effet, remarquons que $(A_y)^\wedge$ est un anneau de valuation discrète
d'uniformisante $x$ et dont le corps résiduel est le corps des fractions
de l'anneau de valuation discrète $D$ défini ci-dessus. Nous voyons donc que même
$$
(B_y)^\wedge[1/(xy)] = ((B_y)^\wedge)_{xy}
$$
n'est pas réduit. Cela fournit un exemple du type suivant.

\begin{lemma}
\label{examples-lemma-nonreduced-recompletion}
Il existe un anneau local noethérien intègre de dimension $2$, noté $(B, \mathfrak m)$,
complet pour la topologie définie par un idéal principal $I = (b)$, et un
élément $f \in \mathfrak m$, $f \not \in I$ tels que
le complété $I$-adique $C = (B_f)^\wedge$ du localisé
principal $B_f$ ne soit pas réduit et que, de plus,
$C_b = C[1/b] = (B_f)^\wedge[1/b]$ ne soit pas réduit.
\end{lemma}

\begin{proof}
Voir la discussion qui précède.
\end{proof}








\section{Un anneau local noethérien non caténaire}
\label{examples-section-non-catenary-Noetherian-local}

\noindent
Bien qu'il existe une théorie satisfaisante de la dimension des anneaux locaux noethériens,
il existe des anneaux locaux noethériens non caténaires. On en trouve un exemple dans
\cite[Appendix, exemple 2]{Nagata}. Nous allons en fait présenter cet exemple
dans le cas le plus simple. Plus précisément, nous allons construire un anneau local noethérien intègre $A$
de dimension $2$ qui n'est pas universellement caténaire. (Remarquons que $A$ est
automatiquement caténaire; voir
Exercices, exercice
\ref{exercises-exercise-Noetherian-local-domain-dim-2-catenary}.)
L'existence d'un anneau local noethérien qui n'est pas universellement
caténaire entraîne celle d'un anneau local noethérien qui n'est pas
caténaire; nous l'expliciterons à la fin de cette section
dans l'exemple particulier considéré.

\medskip\noindent
Soit $k$ un corps, et considérons l'anneau de séries formelles
$k[[x]]$ en une indéterminée sur $k$. Soit
$$
z = \sum\nolimits_{i = 1}^\infty a_i x^i
$$
une série formelle. Nous supposons que $z$, considéré comme élément du corps
de séries de Laurent $k((x)) = k[[x]][1/x]$, est transcendant sur $k(x)$.
Posons
$$
z_j
=
x^{-j}(z - \sum\nolimits_{i = 1, \ldots, j - 1} a_i x^i)
=
\sum\nolimits_{i = j}^\infty a_i x^{i - j}
\in k[[x]].
$$
Remarquons que $z = xz_1$.
Soit $R$ le sous-anneau de $k[[x]]$ engendré par $x$, $z$ et tous les
$z_j$; autrement dit,
$$
R = k[x, z_1, z_2, z_3, \ldots ] \subset k[[x]].
$$
Considérons les idéaux $\mathfrak m = (x)$ et
$\mathfrak n = (x - 1, z_1, z_2 + a_1, z_3 + a_1 + a_2, \ldots)$ de $R$.

\medskip\noindent
Nous avons $xz_{j + 1} + a_j = z_j$. Ainsi, $R/\mathfrak m = k$
et $\mathfrak m$ est un idéal maximal. De plus, tout élément de $R$
qui n'appartient pas à $\mathfrak m$ a pour image un élément inversible de $k[[x]]$, et donc
$R_{\mathfrak m} \subset k[[x]]$. En fait, on en déduit facilement
que $R_{\mathfrak m}$ est un anneau de valuation discrète de corps
résiduel $k$.

\medskip\noindent
Nous affirmons que
$$
R/(x - 1) =
k[x, z_1, z_2, z_3, \ldots ]/(x - 1)
\cong
k[z].
$$
En effet, la relation ci-dessus implique que
$z_{j + 1} = z_j - a_j - (x - 1)z_{j + 1}$, et nous pouvons donc
exprimer la classe de $z_{j + 1}$ en fonction de $z_j$ dans
le quotient $R/(x - 1)$. Puisque le corps des fractions de $R$
est de degré de transcendance $2$ sur $k$ par construction, nous voyons que $z$ est
transcendant sur $k$ dans $R/(x - 1)$, d'où l'isomorphisme voulu.
Ainsi, $\mathfrak n = (x - 1, z)$ est un idéal maximal. En fait,
l'homomorphisme
$$
k[x, x^{-1}, z]_{(x - 1, z)} \longrightarrow R_{\mathfrak n}
$$
est un isomorphisme (puisque $x^{-1}$ est inversible dans $R_{\mathfrak n}$
et que $z_{j + 1} = x^{-1}z_j - a_j = \ldots = f_j(x, x^{-1}, z)$).
Cela montre que $R_{\mathfrak n}$ est un anneau local régulier
de dimension $2$ et de corps résiduel $k$.

\medskip\noindent
Soit $S$ la partie multiplicativement stable
$$
S =
(R \setminus \mathfrak m) \cap (R \setminus \mathfrak n) =
R \setminus (\mathfrak m \cup \mathfrak n)
$$
et posons $B = S^{-1}R$. Nous affirmons que
\begin{enumerate}
\item L'anneau $B$ est une $k$-algèbre.
\item Les idéaux maximaux de l'anneau $B$ sont les deux idéaux
$\mathfrak mB$ et $\mathfrak nB$.
\item Le corps résiduel en chacun de ces idéaux maximaux est $k$.
\item Nous avons $B_{\mathfrak mB} = R_{\mathfrak m}$
et $B_{\mathfrak nB} = R_{\mathfrak n}$,
qui sont des anneaux locaux noethériens réguliers de dimensions $1$ et $2$.
\item L'anneau $B$ est noethérien.
\end{enumerate}
Nous omettons les détails des vérifications.

\medskip\noindent
Étant donnée une $k$-algèbre $B$ possédant les propriétés ci-dessus, nous
obtenons un exemple comme suit. Posons $A = k + \text{rad}(B) \subset B$,
où $\text{rad}(B) = \mathfrak mB \cap \mathfrak nB$ est le radical de Jacobson.
On voit facilement que $B$ est une $A$-algèbre finie, et donc que $A$ est
noethérien par le théorème d'Eakin (voir \cite{Eakin}, ou
\cite[Appendix A1]{Nagata}, ou insérer ici une référence ultérieure).
De plus, $A$ est un anneau local intègre ayant le même corps des fractions que $B$ et
pour corps résiduel $k$. Puisque la dimension de $B$ est $2$, nous voyons que $A$
est également de dimension $2$, d'après
Algèbre, lemme \ref{algebra-lemma-integral-sub-dim-equal}.

\medskip\noindent
Si $A$ était universellement caténaire, la formule de dimension,
Algèbre, lemme \ref{algebra-lemma-dimension-formula},
donnerait $\dim(B_{\mathfrak mB}) = 2$, ce qui est contradictoire.

\medskip\noindent
Remarquons que $B$ est une $A$-algèbre engendrée par un élément.
Ainsi, $B = A[x]/\mathfrak p$ pour un idéal premier
$\mathfrak p$ de $A[x]$. Soit $\mathfrak m' \subset A[x]$
l'idéal maximal correspondant à $\mathfrak mB$. Alors, d'une
part, $\dim(A[x]_{\mathfrak m'}) = 3$ et, d'autre
part,
$$
(0)
\subset \mathfrak pA[x]_{\mathfrak m'}
\subset \mathfrak m'A[x]_{\mathfrak m'}
$$
est une chaîne maximale d'idéaux premiers. Ainsi, $A[x]_{\mathfrak m'}$ est
un exemple d'anneau local noethérien non caténaire.





\section{Existence d'anneaux locaux noethériens pathologiques}
\label{examples-section-bad}

\noindent
Soit $(A, \mathfrak m, \kappa)$ un anneau local noethérien complet.
Dans \cite{Lech}, il est démontré que $A$ est le complété d'un anneau
local noethérien intègre si $\text{profondeur}(A) \geq 1$ et si $A$ contient soit
$\mathbf{Q}$, soit $\mathbf{F}_p$ comme sous-anneau, ou bien contient $\mathbf{Z}$
comme sous-anneau et $A$ est sans torsion comme $\mathbf{Z}$-module.
Cela fournit de nombreux exemples d'anneaux locaux noethériens intègres aux
propriétés « bizarres ».

\medskip\noindent
En appliquant ceci, par exemple, à $A = \mathbf{C}[[x, y]]/(y^2)$, nous obtenons
un anneau local noethérien intègre dont le complété n'est pas réduit.
Comparer avec la
section \ref{examples-section-local-completion-nonreduced}.

\medskip\noindent
Dans \cite{LLPY}, on a trouvé des conditions qui caractérisent les cas où $A$ est
le complété d'un anneau local noethérien réduit.

\medskip\noindent
Dans \cite{Heitmann-completion-UFD}, il est démontré que $A$ est le complété
d'un anneau local noethérien factoriel $R$ si $\text{profondeur}(A) \geq 2$ et si $A$ contient
soit $\mathbf{Q}$, soit $\mathbf{F}_p$ comme sous-anneau, ou bien contient $\mathbf{Z}$
comme sous-anneau et $A$ est sans torsion comme $\mathbf{Z}$-module.
En particulier, $R$ est normal (Algèbre, lemme \ref{algebra-lemma-UFD-normal});
le hensélisé de $R$ est donc lui aussi un anneau normal intègre
(Compléments d'algèbre, lemme \ref{more-algebra-lemma-henselization-normal}).
Ainsi, $A$ comme ci-dessus est le complété d'un anneau local noethérien
hensélien, normal et intègre (puisque les complétés de $R$ et de son hensélisé coïncident;
voir Compléments d'algèbre, lemme \ref{more-algebra-lemma-henselization-noetherian}).

\medskip\noindent
Appliquons ce résultat pour obtenir un anneau local noethérien factoriel $R$ tel que
$R^\wedge \cong \mathbf{C}[[x, y, z, w]]/(wx, wy)$.
Remarquons que $\Spec(R^\wedge)$ est la
réunion d'une composante régulière de dimension $2$ et d'une composante régulière de dimension $3$.
L'anneau $R$ ne peut pas être universellement caténaire. Soit
$$
X \longrightarrow \Spec(R)
$$
l'éclatement de l'idéal maximal. Alors $X$ est un schéma intègre.
Il existe un point fermé $x \in X$ tel que $\dim(\mathcal{O}_{X, x}) = 2$;
plus précisément, au niveau de l'anneau local complet, nous choisissons $x$ sur la
transformée stricte de la composante de dimension $2$, mais non sur la transformée stricte
de la composante de dimension $3$. D'après
Morphismes, lemme \ref{morphisms-lemma-dimension-formula},
nous voyons que $R$ n'est pas universellement caténaire. Comparer avec la
section \ref{examples-section-non-catenary-Noetherian-local}.

\medskip\noindent
L'anneau ci-dessus est caténaire (car c'est un anneau local noethérien factoriel de dimension $3$).
Cependant, dans \cite{Ogoma-example}, l'auteur construit un anneau local
noethérien normal intègre $R$ tel que $R^\wedge \cong \mathbf{C}[[x, y, z, w]]/(wx, wy)$
et que $R$ ne soit pas caténaire. Voir aussi \cite{Heitmann-Ogoma} et
\cite{Lech-YAPO}.

\medskip\noindent
Dans \cite{Heitmann-isolated}, il est démontré que $A$ est le complété
d'un anneau local noethérien $R$ à singularité isolée,
pourvu que $A$ contienne soit $\mathbf{Q}$, soit $\mathbf{F}_p$ comme sous-anneau,
ou que $A$ soit de caractéristique résiduelle $p > 0$ et que $p$ ne puisse devenir un
diviseur de zéro non nul dans aucun localisé propre de $A$.
Nous disons ici qu'un anneau local noethérien $R$
est à singularité isolée si $R_\mathfrak p$ est
un anneau local régulier pour tout idéal premier non maximal $\mathfrak p \subset R$.

\medskip\noindent
Les articles \cite{Nishimura-few} et \cite{Nishimura-few-II}
contiennent de longues listes d'anneaux locaux
noethériens « pathologiques » dont le complété est prescrit. En particulier, ils construisent
un exemple d'anneau local de Nagata, normal et intègre, de dimension $2$, dont le
complété est $\mathbf{C}[[x, y, z]]/(yz)$, et un autre dont le complété
est $\mathbf{C}[[x, y, z]]/(y^2 - z^3)$.

\medskip\noindent
À titre de digression, il est démontré dans \cite{Loepp} que $A$ est le complété d'un
anneau local noethérien excellent et intègre si $A$ est réduit, équidimensionnel,
et si aucun entier de $A$ n'est un diviseur de zéro. Cela ne conduit toutefois pas
à des anneaux locaux noethériens « pathologiques », puisque ceux obtenus sont excellents !




\section{Dimension dans les anneaux noethériens de Jacobson}
\label{examples-section-noetherian-jacobson}

\noindent
Soit $k$ la clôture algébrique d'un corps fini.
Soient $A = k[x, y]$ et $X = \Spec(A)$.
Soit $C = V(x)$ l'axe des $y$ (on pourrait choisir tout autre
sous-schéma fermé intègre de dimension $1$ de $X$).
Soit $C_1, C_2, C_3, \ldots$ une énumération
des autres sous-schémas fermés intègres de $X$ de dimension $1$.
Soit $p_1, p_2, p_3, \ldots$ une
énumération des points fermés de $C$.

\medskip\noindent
Affirmation : pour tout $n$, il existe un fermé irréductible
$Z_n \subset X$ de dimension $1$ tel que
$$
\{p_n\} = Z_n \cap (C \cup C_1 \cup \ldots \cup C_n)
$$
ensemblistement. Pour cela, posons
$Y = C \cup C_1 \cup C_2 \cup \ldots \cup C_n$.
C'est un schéma algébrique affine réduit de dimension $1$ sur $k$.
Il suffit de trouver $f \in k[x, y]$ tel que $V(f) \cap  Y = \{p_n\}$
ensemblistement, car nous pouvons alors prendre pour $Z_n$ une composante
irréductible convenable de $V(f)$. Puisque l'application de restriction
$$
k[x, y] \longrightarrow \Gamma(Y, \mathcal{O}_Y)
$$
est surjective, il suffit de trouver une fonction régulière $g$ sur $Y$
dont l'ensemble des zéros est $\{p_n\}$ ensemblistement.
Pour voir que cela est possible, choisissons un diviseur de Cartier effectif
$D \subset Y$ dont le support est $p_n$ (cela est possible d'après
Variétés, lemme \ref{varieties-lemma-complement-codim-1-closed-points}).
Il suffit donc de montrer que $\mathcal{O}_Y(ND) \cong \mathcal{O}_Y$
pour un certain $N > 0$. Or le groupe de Picard d'un schéma algébrique
affine de dimension $1$ sur la clôture algébrique d'un corps fini est de torsion
(insérer ici une référence ultérieure), ce qui démontre l'affirmation.

\medskip\noindent
Choisissons $Z_n$ comme ci-dessus pour tout $n$. Puisque $k[x, y]$ est un anneau factoriel,
nous pouvons écrire $Z_n = V(f_n)$ pour un élément irréductible $f_n \in A$.
Soit $S \subset k[x, y]$ la partie multiplicative engendrée
par $f_1, f_2, f_3, \ldots$. Considérons l'anneau noethérien $B = S^{-1}A$.

\medskip\noindent
Il est clair que l'homomorphisme d'anneaux $A \to B$ identifie les anneaux locaux
correspondants et induit une injection $\Spec(B) \to \Spec(A)$.
De plus, en considérant la courbe $C_1$, nous voyons que
seuls les points de $C \cap C_1$ sont supprimés lors du passage
de $\Spec(A)$ à $\Spec(B)$.
En particulier, $\Spec(B)$ possède une infinité
d'idéaux maximaux correspondant à des idéaux
maximaux de $A$. En revanche, $xB$ est un idéal maximal, car
le spectre de $B/xB$ est constitué d'un unique idéal premier,
puisque nous avons supprimé tous les points fermés de $C = V(x)$ (mais non
le point générique).
Enfin, pour $i \geq 1$, considérons la courbe $C_i$.
Écrivons $C_i = V(g_i)$, où $g_i \in A$ est irréductible.
Si $C_i = Z_n$ pour un certain $n$, alors $g_iB$ est l'idéal unité.
Sinon, tous les points fermés de $C_i$ sauf un nombre fini
subsistent lors du passage de $A$ à $B$ :
plus précisément, seuls les points de
$(Z_1 \cup \ldots \cup Z_{i - 1} \cup C) \cap C_i$
sont supprimés de $C_i$.

\medskip\noindent
La structure du spectre de $B$ décrite ci-dessus
montre que $B$ est un anneau de Jacobson d'après
Algèbre, lemme \ref{algebra-lemma-noetherian-dim-1-Jacobson}.
Les idéaux maximaux sont ceux de $A$
qui appartiennent à $\Spec(B)$ (et ils sont en nombre infini),
ainsi que l'idéal maximal $xB$. Nous constatons donc que nous avons
des anneaux locaux de dimensions $1$ et $2$.

\begin{lemma}
\label{examples-lemma-Noetherian-Jacobson}
Il existe un anneau de Jacobson noethérien intègre et universellement caténaire $B$
possédant des idéaux maximaux $\mathfrak m_1, \mathfrak m_2$ tels que
$\dim(B_{\mathfrak m_1}) = 1$ et $\dim(B_{\mathfrak m_2}) = 2$.
\end{lemma}

\begin{proof}
La construction de $B$ est donnée ci-dessus. Signalons simplement que
$B$ est universellement caténaire d'après
Algèbre, lemme \ref{algebra-lemma-localization-catenary} et
Morphismes, lemme \ref{morphisms-lemma-ubiquity-uc}.
\end{proof}




\section{Espace sous-jacent noethérien, schéma non noethérien}
\label{examples-section-noetherian-not-noetherian}

\noindent
Nous donnons deux exemples montrant qu'un schéma dont l'espace topologique sous-jacent
est noethérien peut ne pas être un schéma noethérien.

\begin{example}
\label{examples-example-many-variables}
Soit $k$ un corps, et soit $A = k[x_1, x_2, x_3, \dots] /
(x_1^2, x_2^2, x_3^2, \dots)$. Tout idéal premier de $A$ contient les
éléments nilpotents $x_1, x_2, x_3, \dots$; ainsi $\mathfrak p = (x_1, x_2, x_3,
\dots)$ est l'unique idéal premier de $A$. Par conséquent, l'espace
topologique sous-jacent à $\operatorname{Spec} A$ est réduit à un seul point et est
en particulier noethérien. Toutefois, l'idéal $\mathfrak p$ n'est manifestement pas
de type fini.
\end{example}

\begin{example}
\label{examples-example-many-mononials}
Soit $k$ un corps, et soit $A \subseteq k[x, y]$ le
sous-anneau engendré par $k$ et les monômes $\{xy^i\}_{i \ge 0}$. Les idéaux
premiers de $A$ qui ne contiennent pas $x$ sont en correspondance biunivoque avec
les idéaux premiers de $A_x \cong k[x, x^{-1}, y]$. Si $\mathfrak p$ est un
idéal premier qui contient $x$, alors il contient tout $xy^i$, $i \ge 0$,
car $(xy^i)^2 = x(xy^{2i}) \in \mathfrak p$ et $\mathfrak p$ est
radical. Par conséquent, $\mathfrak p = (\{xy^i\}_{i \ge 0})$. Ainsi, l'espace
topologique sous-jacent à $\operatorname{Spec} A$ est noethérien, puisqu'il
est constitué des points du schéma noethérien
$\Spec(A[x, x^{-1}, y])$ et de l'idéal premier $\mathfrak p$.
Mais l'anneau $A$ n'est pas noethérien, car $\mathfrak p$
n'est pas de type fini. Notons que,
dans cet exemple, $A$ est également intègre.
\end{example}






\section{Variété non quasi-affine dont la normalisation est quasi-affine}
\label{examples-section-nonquasi-affine}

\noindent
L'existence d'un exemple de ce type est mentionnée dans
\cite[II remarque 6.6.13]{EGA}, qui renvoie au cinquième volume des
EGA pour un tel exemple; or ce cinquième volume n'a pas paru.

\medskip\noindent
Soit $k$ un corps.
Soit $Y = \mathbf{A}^2_k \setminus \{(0, 0)\}$.
Nous allons construire un morphisme fini, surjectif et birationnel
$\pi : Y \longrightarrow X$
où $X$ est une variété sur $k$ telle que $X$ n'est pas quasi-affine.
Plus précisément, considérons les courbes suivantes dans $Y$ :
$$
\begin{matrix}
C_1 & : & x = 0 \\
C_2 & : & y = 0
\end{matrix}
$$
Remarquons que $C_1 \cap C_2 = \emptyset$. Choisissons l'isomorphisme
$\varphi : C_1 \to C_2$, $(0, y) \mapsto (y^{-1}, 0)$.
Nous affirmons qu'il existe un unique morphisme $\pi : Y \to X$ comme ci-dessus
tel que
$$
\xymatrix{
C_1
\ar@<1ex>[rr]^{\text{id}} \ar@<-1ex>[rr]_{\varphi}
& &
Y \ar[r]^\pi & X
}
$$
soit un diagramme coégalisateur dans la catégorie des variétés (et même dans
la catégorie des schémas). Admettons-le provisoirement et
montrons qu'un tel $X$ ne peut être quasi-affine. En effet, il est clair
que nous aurions
$$
\Gamma(X, \mathcal{O}_X) =
\{ f \in k[x, y] \mid f(0, y) = f(y^{-1}, 0)\} =
k \oplus (xy) \subset k[x, y].
$$
En particulier, ces fonctions ne séparent pas les points $(1, 0)$
et $(-1, 0)$, dont les images dans $X$ (comme nous le verrons plus bas) sont distinctes
(si la caractéristique de $k$ n'est pas $2$).

\medskip\noindent
Pour montrer que $X$ existe, considérons l'ouvert de Zariski
$D(x + y) \subset Y$ de $Y$. C'est le spectre
de l'anneau
$k[x, y, 1/(x + y)]$
et les courbes $C_1$, $C_2$ sont entièrement contenues dans
$D(x + y)$. De plus, le morphisme
$$
C_1 \amalg C_2
\longrightarrow
D(x + y) \cap Y = \Spec(k[x, y, 1/(x + y)])
$$
est une immersion fermée. Il résulte de
Compléments d'algèbre, lemme \ref{more-algebra-lemma-fibre-product-finite-type},
que l'anneau
$$
A =
\{f \in k[x, y, 1/(x + y)] \mid f(0, y) = f(y^{-1}, 0)\}
$$
est une $k$-algèbre de type fini. D'autre part, nous avons l'ouvert
$D(xy) \subset Y$ de $Y$, qui est disjoint des courbes $C_1$
et $C_2$. C'est le spectre de l'anneau
$$
B = k[x, y, 1/xy].
$$
Remarquons que $A_{xy} \cong B_{x + y}$ (car $A$ contient manifestement
les éléments $xyP(x, y)$ pour tout polynôme $P$, ainsi que l'élément $xy/(x + y)$).
Le schéma $X$ est obtenu par recollement des schémas affines
$\Spec(A)$ et $\Spec(B)$ au moyen de l'isomorphisme
$A_{xy} \cong B_{x + y}$; il est donc manifestement de type fini sur
$k$. Pour voir qu'il est séparé, il faut montrer que
l'homomorphisme d'anneaux $A \otimes_k B \to B_{x + y}$ est surjectif. Pour cela,
utilisons le fait que $A \otimes_k B$ contient l'élément
$xy/(x + y) \otimes 1/xy$, dont l'image est $1/(x + y)$.
Le morphisme $Y \to X$ est donné par les morphismes naturels
$D(x + y) \to \Spec(A)$ et $D(xy) \to \Spec(B)$.
Comme ils sont tous deux finis, il en résulte que $Y \to X$ est fini,
comme souhaité. Nous omettons la vérification que $X$ est bien le
coégalisateur du diagramme affiché ci-dessus; voir toutefois
(insérer ici une référence ultérieure sur les sommes amalgamées dans la catégorie des schémas
). Remarquons que le morphisme $\pi : Y \to X$ envoie bien
les points $(1, 0)$
et $(-1, 0)$ sur des points distincts de $X$, puisque la
fonction $(x + y^3)/(x + y)^2 \in A$ prend les valeurs
$1/1$ et, respectivement,\ $-1/(-1)^2 = -1$, qui sont toujours distinctes
(sauf en caractéristique $2$ -- au lecteur de choisir d'autres points
en caractéristique $2$). Résumons cette discussion sous
la forme d'un lemme.

\begin{lemma}
\label{examples-lemma-quasi-affine-normalization-not-quasi-affine}
Soit $k$ un corps.
Il existe une variété $X$ dont la normalisation est quasi-affine, mais
qui n'est pas elle-même quasi-affine.
\end{lemma}

\begin{proof}
Voir la discussion ci-dessus et (insérer ici une référence ultérieure sur la normalisation).
\end{proof}



\section{Prise d'images schématiques}
\label{examples-section-scheme-theoretic-image}

\noindent
Soit $k$ un corps. Soit $t$ une indéterminée. Soient $Y = \Spec(k[t])$
et $X = \coprod_{n \geq 1} \Spec(k[t]/(t^n))$. Notons
$f : X \to Y$ le morphisme
défini au moyen de l'immersion fermée $\Spec(k[t]/(t^n)) \to \Spec(k[t])$
pour tout $n \geq 1$. Dans ce cas, nous avons
\begin{enumerate}
\item L'image schématique
(Morphismes, définition \ref{morphisms-definition-scheme-theoretic-image})
de $f$ est $Y$. En revanche, l'image de $f$ est le point fermé
$t = 0$ de $Y$. Ainsi, la partie fermée sous-jacente à
l'image schématique de $f$
n'est pas égale à l'adhérence de l'image de $f$.
\item La formation de l'image schématique ne commute pas
à la restriction au sous-schéma ouvert $V = \Spec(k[t, 1/t]) \subset Y$.
En effet, l'image réciproque de $V$ dans $X$ est vide; l'image schématique
de $f|_{f^{-1}(V)} : f^{-1}(V) \to V$ est donc le schéma vide.
Celui-ci n'est pas égal à $Y \cap V$.
\end{enumerate}




\section{Images de parties localement fermées}
\label{examples-section-non-chevalley}

\noindent
Le théorème de Chevalley affirme que l'image d'une partie constructible
par un morphisme de présentation finie entre schémas affines est constructible; voir
Algèbre, théorème \ref{algebra-theorem-chevalley}, et
Morphismes, section \ref{morphisms-section-constructible}.
Nous allons voir que l'énoncé analogue n'est pas vrai
pour les images de parties localement fermées.

\medskip\noindent
Soit $k$ un corps de caractéristique $0$. Considérons le morphisme de projection
$$
f :
X = \Spec(k[t, x_1, x_2, \ldots, y_1, y_2, \ldots])
\longrightarrow
\Spec(k[x_1, x_2, \ldots, y_1, y_2, \ldots]) = Y
$$
C'est un morphisme de présentation finie.
Soit $Z$ la partie fermée de $X$ définie par
$$
x_1(t - 1) = 0,\quad
x_2(t - 1)(t - 2) = 0,\quad
x_3(t - 1)(t - 2)(t - 3) = 0,\quad \ldots
$$
Soit $U = \bigcup_{j \geq 1} U_j$ l'ouvert de $X$ défini par
$$
U_j = \text{points où }y_j(t - 1)(t - 2) ... (t - j)\text{ est non nul}
$$
Alors nous avons
$$
f(Z \cap U_j) =
\text{points où }x_1, \ldots, x_j\text{ sont nuls et }y_j\text{ est non nul}
$$
Nous affirmons que $B = f(Z \cap U) = \bigcup_{j \geq 1} f(Z \cap U_j)$
n'est pas une réunion finie de parties localement fermées de $Y$.

\medskip\noindent
Démonstration de l'affirmation. Supposons que $B = A_1 \cup \cdots \cup A_m$
soit un recouvrement fini de $B$ par des parties localement fermées de $Y$.
Nous allons montrer par récurrence sur $n$ que $m \geq n$. Le cas initial
$n = 1$ est immédiat, puisque $B$ est non vide. Supposons $n > 1$ et
l'hypothèse de récurrence vraie pour $n - 1$. Comme l'adhérence de $B$ est
$(x_1 = 0)$, l'un des $A_i$ doit contenir une partie ouverte non vide de
$(x_1 = 0)$. Alors $A_i$ doit être ouvert dans $(x_1 = 0)$.
Mais une telle partie ouverte ne peut contenir aucun point où
$y_1 = 0$; en effet, pour les points de $B$, $y_1 = 0$ entraîne $x_2 = 0$,
ce qui montre que $B$ ne contient aucun voisinage de $(x, y)$ dans $(x_1 = 0)$.
Par conséquent, les $m - 1$ parties restantes induisent
un recouvrement constructible de $B \cap (y_1 = 0)$.
Or l'application de décalage à droite
$x_i \mapsto x_{i + 1}$, $y_i \mapsto y_{i + 1}$ identifie $B$
à $B \cap (y_1 = 0)$! L'hypothèse de récurrence donne donc
$m - 1 \geq n - 1$, d'où $m \geq n$. Ceci achève
la démonstration de l'étape de récurrence et établit l'affirmation.

\begin{lemma}
\label{examples-lemma-no-chevalley}
Il existe un morphisme $f : X \to Y$ de présentation finie
entre schémas affines et une partie localement fermée $T$ de $X$
tels que $f(T)$ ne soit pas une réunion finie de parties localement fermées de $Y$.
\end{lemma}

\begin{proof}
Voir la discussion ci-dessus.
\end{proof}










\section[Sous-schéma localement fermé non ouvert dans un fermé]{Un sous-schéma localement fermé qui n'est pas ouvert dans un fermé}
\label{examples-section-strange-immersion}

\noindent
Ceci est une copie de
Morphismes, exemple \ref{morphisms-example-thibaut}.
Voici un exemple d'immersion qui ne se compose pas d'une
immersion ouverte suivie d'une immersion fermée.
Soit $k$ un corps.
Soit $X = \Spec(k[x_1, x_2, x_3, \ldots])$.
Soit $U = \bigcup_{n = 1}^{\infty} D(x_n)$.
Alors $U \to X$ est une immersion ouverte.
Considérons les idéaux
$$
I_n =
(x_1^n, x_2^n, \ldots, x_{n - 1}^n, x_n - 1, x_{n + 1}, x_{n + 2}, \ldots)
\subset
k[x_1, x_2, x_3, \ldots][1/x_n].
$$
Remarquons que $I_n k[x_1, x_2, x_3, \ldots][1/x_nx_m] = (1)$
pour tout $m \not = n$. Par conséquent, les idéaux quasi-cohérents
$\widetilde I_n$ sur $D(x_n)$ coïncident sur $D(x_nx_m)$; plus précisément,
$\widetilde I_n|_{D(x_nx_m)} = \mathcal{O}_{D(x_n x_m)}$ si
$n \not = m$. Ces idéaux se recollent donc en un faisceau quasi-cohérent d'idéaux
$\mathcal{I} \subset \mathcal{O}_U$.
Soit $Z \subset U$ le sous-schéma fermé correspondant à
$\mathcal{I}$. Ainsi, $Z \to X$ est une immersion.

\medskip\noindent
Nous affirmons qu'il est impossible de factoriser $Z \to X$ sous la forme
$Z \to \overline{Z} \to X$, où $\overline{Z} \to X$ est une immersion fermée
et $Z \to \overline{Z}$ est une immersion ouverte. En effet, $\overline{Z}$ devrait
être défini par un idéal $I \subset k[x_1, x_2, x_3, \ldots]$
tel que $I_n = I k[x_1, x_2, x_3, \ldots][1/x_n]$.
Mais le seul élément $f \in k[x_1, x_2, x_3, \ldots]$
qui appartienne à tous les $I_n$ est $0$! Un tel $I$ n'existe donc pas.

\medskip\noindent
Le morphisme $Z \to X$ fournit aussi un exemple de comportement pathologique
des images schématiques d'immersions. En effet, les arguments ci-dessus
montrent que l'image schématique de l'immersion $Z \to X$ est $X$.
En revanche, nous constatons que
\begin{enumerate}
\item $Z$ n'est pas topologiquement dense dans $X$, et
\item l'image schématique de $Z = Z \cap U \to U$ est simplement
$Z$. Elle n'est pas égale à $U \cap X = U$; la formation
de l'image schématique, dans ce cas,
ne commute donc pas aux restrictions aux ouverts.
\end{enumerate}






\section{Non-existence d'ouverts convenables}
\label{examples-section-nonexistence-opens}

\noindent
Cette section complète les résultats de
Propriétés, section \ref{properties-section-finding-affine-opens}.

\medskip\noindent
Soit $k$ un corps et soit $A = k[z_1, z_2, z_3, \ldots]/I$, où $I$
est l'idéal engendré par tous les produits deux à deux
$z_iz_j$, $i \not = j$, $i, j \in \mathbf{N}$. Posons $S = \Spec(A)$.
Soit $s \in S$ le point fermé correspondant à l'idéal
maximal $(z_i)$. Nous affirmons qu'il n'existe aucun ouvert
quasi-compact $V \subset S \setminus \{s\}$ qui soit
dense dans $S \setminus \{s\}$. Remarquons que $S \setminus \{s\} = \bigcup D(z_i)$.
Chaque $D(z_i)$ est ouvert et irréductible, de point générique
$\eta_i$. Nous en déduisons que $\eta_i \in V$ pour tout $i$.
Cependant, tout ouvert affine principal de $S \setminus \{s\}$ est
de la forme $D(f)$, où $f \in (z_1, z_2, \ldots)$. Alors
$f \in (z_1, \ldots, z_n)$ pour un certain $n$, et $D(f)$ ne contient
qu'un nombre fini des points $\eta_i$. Ainsi, $V$ ne peut être quasi-compact.

\medskip\noindent
Soit $k$ un corps et soit $B = k[x, z_1, z_2, z_3, \ldots]/J$, où $J$
est l'idéal engendré par les produits $xz_i$, $i \in \mathbf{N}$, et par
tous les produits deux à deux $z_iz_j$, $i \not = j$, $i, j \in \mathbf{N}$.
Posons $T = \Spec(B)$. Considérons l'ouvert principal $U = D(x)$.
Nous affirmons qu'il n'existe aucun ouvert quasi-compact $V \subset T$ tel que
$V \cap U = \emptyset$ et que $V \cup U$ soit dense dans $T$.
Soit $t \in T$ le point fermé correspondant à l'idéal
maximal $(x, z_i)$. L'adhérence de $U$ dans $T$ est
$\overline{U} = U \cup \{t\}$. Ainsi, $V \subset \bigcup_i D(z_i)$
est un ouvert quasi-compact. Les arguments du paragraphe précédent
montrent que $V$ ne peut être dense dans $\bigcup D(z_i)$.

\begin{lemma}
\label{examples-lemma-complement-of-affine-does-not-contain-qc-dense-open}
Non-existence d'ouverts quasi-compacts dans des schémas affines :
\begin{enumerate}
\item Il existe un schéma affine $S$ et un ouvert affine $U \subset S$
tels qu'il n'existe aucun ouvert quasi-compact $V \subset S$ pour lequel
$U \cap V = \emptyset$ et $U \cup V$ est dense dans $S$.
\item Il existe un schéma affine $S$ et un point fermé $s \in S$ tels que
$S \setminus \{s\}$ ne contienne aucun ouvert quasi-compact dense.
\end{enumerate}
\end{lemma}

\begin{proof}
Voir la discussion ci-dessus.
\end{proof}

\noindent
Soit $X$ le recollement de deux copies du schéma affine $T$ (voir ci-dessus)
le long de l'ouvert affine $U$. Il existe donc un morphisme $\pi : X \to T$ et
$X = U_1 \cup U_2$ tels que $\pi$ envoie isomorphiquement $U_i$ sur $T$ et
$U_1 \cap U_2$ sur $U$. Remarquons que $X$ est quasi-séparé
(d'après Schémas, lemme \ref{schemes-lemma-characterize-quasi-separated})
et quasi-compact. Nous affirmons qu'il n'existe aucun ouvert séparé, dense et
quasi-compact $W \subset X$. Considérons en effet les deux points fermés
$x_1 \in U_1$, $x_2 \in U_2$ qui s'envoient sur le point fermé $t \in T$
introduit ci-dessus. Soit $\tilde \eta \in U_1 \cap U_2$ le point générique
qui s'envoie sur l'unique point générique $\eta$ de $U$.
Remarquons que $\tilde\eta \leadsto x_1$ et $\tilde\eta \leadsto x_2$
se trouvent au-dessus de la spécialisation $\eta \leadsto t$.
Comme $\pi|_W : W \to T$ est séparé, nous en concluons qu'on ne peut avoir simultanément
$x_1$ et $x_2 \in W$ (par le critère valuatif de séparation,
Schémas, lemme \ref{schemes-lemma-valuative-criterion-separatedness}).
Supposons $x_1 \not \in W$. Alors $W \cap U_1$ est un ouvert quasi-compact (car $X$
est quasi-séparé), dense dans $U_1$, qui ne contient pas $x_1$.
Observons maintenant qu'il existe un isomorphisme $(T, t) \cong (S, s)$
de schémas (envoyant $x$ sur $z_1$ et $z_i$ sur $z_{i + 1}$).
Le premier paragraphe de cette section fournit alors une contradiction.

\begin{lemma}
\label{examples-lemma-no-dense-separated-quasi-compact-open-in-qcqs}
Il existe un schéma quasi-compact et quasi-séparé $X$ qui ne contient
aucun ouvert séparé, quasi-compact et dense.
\end{lemma}

\begin{proof}
Voir la discussion ci-dessus.
\end{proof}







\section{Non-existence d'un sous-schéma ouvert quasi-compact dense}
\label{examples-section-nonexistence-qc-dense-open-subscheme}

\noindent
Soit $X$ un espace algébrique quasi-compact et quasi-séparé sur un corps
$k$. Nous savons que le lieu schématique $X' \subset X$ est un sous-espace
ouvert dense; voir
Propriétés des espaces, proposition
\ref{spaces-properties-proposition-locally-quasi-separated-open-dense-scheme}.
En fait, ce résultat vaut lorsque $X$ est raisonnable; voir
Espaces décents, proposition
\ref{decent-spaces-proposition-reasonable-open-dense-scheme}.
Il est naturel de demander si l'on peut trouver un sous-schéma ouvert
quasi-compact dense de $X$. Ce n'est en général pas possible.

\medskip\noindent
Supposons que la caractéristique de $k$ ne soit pas 2.
Soit $B = k[x, z_1, z_2, z_3, \ldots]/J$, où $J$ est l'idéal engendré par
les produits $xz_i$, $i \in \mathbf{N}$, et par tous les produits deux à deux
$z_iz_j$, $i \not = j$, $i, j \in \mathbf{N}$. Posons $U = \Spec(B)$.
Notons $0 \in U$ le point fermé dont toutes les coordonnées sont nulles.
Posons
$$
j : R = \Delta \amalg \Gamma \longrightarrow U \times_k U
$$
où $\Delta$ est l'image du morphisme diagonal de $U$ sur $k$ et
$$
\Gamma = \{((x, 0, 0, 0, \ldots), (-x, 0, 0, 0, \ldots))
\mid x \in \mathbf{A}^1_k, x \not = 0\}.
$$
Il est clair que $s, t : R \to U$ sont \'etales; ainsi,
$j$ est une relation d'équivalence \'etale. Le quotient $X = U/R$
est un espace algébrique (Espaces, théorème \ref{spaces-theorem-presentation}).
Remarquons que $j$ n'est pas une immersion, car
$(0, 0) \in \Delta$ appartient à l'adhérence de $\Gamma$.
Ainsi, $X$ n'est pas un schéma. En revanche, $X$ est quasi-séparé,
car $R$ est quasi-compact. Notons $0_X$ l'image du point $0 \in U$.
Nous affirmons que $X \setminus \{0_X\}$ est un schéma; plus précisément,
$$
X \setminus \{0_X\} =
\Spec\left(k[x^2, x^{-2}]\right) \amalg
\Spec\left(k[z_1, z_2, z_3, \ldots]/(z_iz_j)\right) \setminus \{0\}
$$
(détails omis). D'autre part, nous avons vu dans la
section \ref{examples-section-nonexistence-opens} que le schéma
figurant au membre de droite ne contient
aucun ouvert quasi-compact dense.

\begin{lemma}
\label{examples-lemma-nonexistence-qc-dense-open-subscheme}
Il existe un espace algébrique quasi-compact et quasi-séparé
qui ne contient aucun sous-schéma ouvert quasi-compact dense.
\end{lemma}

\begin{proof}
Voir la discussion ci-dessus.
\end{proof}

\noindent
En utilisant la construction de
Espaces, exemple \ref{spaces-example-non-representable-descent},
de la même manière que nous avons employé ci-dessus celle de
Espaces, exemple \ref{spaces-example-affine-line-involution},
on obtient un espace algébrique quasi-compact, quasi-séparé et
localement séparé qui ne contient aucun sous-schéma ouvert
quasi-compact dense.


\section{Schémas affines au-dessus d'espaces algébriques}
\label{examples-section-embedding-affines}

\medskip\noindent
Supposons que $f : Y \to X$ soit un morphisme de schémas, que $f$
soit localement de type fini et que $Y$ soit affine. Il existe alors une immersion
$Y \to \mathbf{A}^n_X$ de $Y$ dans l'espace affine de dimension $n$ sur $X$.
Voir l'énoncé légèrement plus général de
Morphismes, lemme \ref{morphisms-lemma-quasi-affine-finite-type-over-S}.

\medskip\noindent
Supposons maintenant que $f : Y \to X$ soit un morphisme d'espaces algébriques,
$f$ étant localement de type fini et $Y$ un schéma affine. En général,
il n'existe pas nécessairement d'immersion de $Y$ dans l'espace
affine de dimension $n$ sur $X$.

\medskip\noindent
Un premier contre-exemple (peu élégant) est donné par $Y = \Spec(k)$ et
$X = [\mathbf{A}^1_k/\mathbf{Z}]$, où $k$ est un corps de caractéristique nulle
et $\mathbf{Z}$ agit sur $\mathbf{A}^1_k$ par la translation $(n, t) \mapsto t + n$.
En effet, pour tout morphisme $Y \to \mathbf{A}^n_X$ au-dessus de $X$, le changement de base
au revêtement $\mathbf{A}^1_k$ de $X$ donne une réunion disjointe infinie de copies de
$\mathbf{A}^1_k$ s'envoyant dans $\mathbf{A}^{n + 1}_k$, morphisme qui n'est pas une
immersion.

\medskip\noindent
Un second contre-exemple est $Y = \mathbf{A}^1_k \to X = \mathbf{A}^1_k/R$,
où $R = \{(t, t)\} \amalg \{(t, -t), t \not = 0\}$. En effet, dans
ce cas, le morphisme $Y \to \mathbf{A}^n_X$ serait donné par des
fonctions régulières $f_1, \ldots, f_n$ sur $Y$; le
produit fibré de $Y$ avec le revêtement
$\mathbf{A}^{n + 1}_k \to \mathbf{A}^n_X$
serait donc le schéma
$$
\{(f_1(t), \ldots, f_n(t), t)\} \amalg
\{(f_1(t), \ldots, f_n(t), -t), t \not = 0\}
$$
muni du morphisme évident vers $\mathbf{A}^{n + 1}_k$, qui n'est pas une immersion.
Remarquons que ce contre-exemple a $X$ quasi-séparé.

\begin{lemma}
\label{examples-lemma-cannot-embed-into-affine}
Il existe un morphisme de type fini d'espaces algébriques $Y \to X$
tel que $Y$ soit affine et $X$ quasi-séparé, pour lequel il n'existe
aucune immersion $Y \to \mathbf{A}^n_X$ au-dessus de $X$.
\end{lemma}

\begin{proof}
Voir la discussion ci-dessus.
\end{proof}










\section{Image directe de modules quasi-cohérents}
\label{examples-section-push-quasi-coherent}

\noindent
Dans Schémas, lemme \ref{schemes-lemma-push-forward-quasi-coherent},
nous avons montré que $f_*$ envoie les modules quasi-cohérents sur des modules
quasi-cohérents lorsque $f$ est quasi-compact et quasi-séparé. Voici des
exemples montrant que ces deux conditions sont nécessaires.

\medskip\noindent
Supposons que $Y = \Spec(A)$ soit un schéma affine et que
$X = \coprod_{n \in \mathbf{N}} Y$. Nous affirmons que $f_*\mathcal{O}_X$
n'est pas quasi-cohérent, où $f : X \to Y$ est le morphisme évident.
En effet, pour $a \in A$, nous avons
$$
f_*\mathcal{O}_X(D(a)) = \prod\nolimits_{n \in \mathbf{N}} A_a
$$
Pour que $f_*\mathcal{O}_X$ soit quasi-cohérent, il faudrait donc avoir
$$
\prod\nolimits_{n \in \mathbf{N}} A_a
=
\left(\prod\nolimits_{n \in \mathbf{N}} A\right)_a
$$
pour tout $a \in A$. Ce n'est pas vrai en général : si, par exemple,
$A = \mathbf{Z}$ et $a = 2$, alors $(1, 1/2, 1/4, 1/8, \ldots)$
est un élément du membre de gauche qui n'appartient pas à celui de droite.
Remarquons que $f$ est un morphisme séparé non quasi-compact.

\medskip\noindent
Soit $k$ un corps. Posons
$$
A = k[t, z, x_1, x_2, x_3, \ldots]/(tx_1z, t^2x_2^2z, t^3x_3^3z, \ldots)
$$
Soit $Y = \Spec(A)$. Soit $V \subset Y$ le sous-schéma ouvert
$V = D(x_1) \cup D(x_2) \cup \ldots$. Soit $X$ le recollement de deux copies de $Y$
le long de $V$. Soit $f : X \to Y$ le morphisme évident. Nous avons alors
une suite exacte
$$
0 \to f_*\mathcal{O}_X \to
\mathcal{O}_Y \oplus \mathcal{O}_Y \xrightarrow{(1, -1)} j_*\mathcal{O}_V
$$
où $j : V \to Y$ est le morphisme d'inclusion. Comme
$$
A \longrightarrow \prod A_{x_n}
$$
est injectif (détails omis), nous obtenons $\Gamma(Y, f_*\mathcal{O}_X) = A$.
En revanche, le noyau de l'application
$$
A_t \longrightarrow \prod A_{tx_n}
$$
est non nul, car il contient l'élément $z$. Par conséquent,
$\Gamma(D(t), f_*\mathcal{O}_X)$ est strictement plus grand que
$A_t$, puisqu'il contient $(z, 0)$. Ainsi, $f_*\mathcal{O}_X$
n'est pas quasi-cohérent. Remarquons que $f$ est quasi-compact mais
non quasi-séparé.

\begin{lemma}
\label{examples-lemma-pushforward-quasi-coherent}
Schémas, lemme \ref{schemes-lemma-push-forward-quasi-coherent},
est optimal au sens où l'on ne peut supprimer ni l'hypothèse
de quasi-compacité ni celle de quasi-séparation.
\end{lemma}

\begin{proof}
Voir la discussion ci-dessus.
\end{proof}








\section{Un module non de type fini dont les germes sont libres de rang 1}
\label{examples-section-nonfree}

\noindent
Soit $R = \mathbf{Q}[x]$. Posons
$M = \sum_{n \in \mathbf{N}} \frac{1}{x - n}R$, considéré comme sous-module du
corps des fractions de $R$. Alors $M$ n'est pas de type fini, mais,
pour tout idéal premier $\mathfrak p$ de $R$, nous avons
$M_{\mathfrak p} \cong R_{\mathfrak p}$ comme
$R_{\mathfrak p}$-modules.

\medskip\noindent
Un exemple du même genre est donné par $R = \mathbf{Z}$ et
$M = \sum_{p \text{ premier}} \frac{1}{p} \mathbf{Z} \subset \mathbf{Q}$,
qui est l'ensemble des fractions $\frac{a}{b}$ où $b$ est non nul et
sans facteur carré.








\section{Un idéal non inversible mais inversible dans les germes}
\label{examples-section-locally-invertible-not-invertible}

\noindent
Soit $A$ un anneau intègre et soit $I \subset A$ un idéal non nul.
Rappelons que, lorsque nous disons que $I$ est inversible, nous entendons que
$I$ est inversible comme $A$-module.
Nous allons construire un exemple de cette situation où
$I$ n'est pas inversible, bien que $I_\mathfrak q = (f) \subset A_\mathfrak q$
soit un idéal principal non nul pour tout idéal premier $\mathfrak q \subset A$.
Dans la littérature, la propriété que $I_\mathfrak q$ soit principal
pour tout idéal premier $\mathfrak q$ s'exprime parfois en
disant que « $I$ est un idéal localement principal ». Nous ne pouvons employer
cette terminologie, car « local » signifie toujours ici
« local pour la topologie de Zariski » (ou pour toute autre topologie dans
laquelle nous travaillons).

\medskip\noindent
Soient $R = \mathbf{Q}[x]$ et $M = \sum \frac{1}{x - n}R$
le module construit dans la section \ref{examples-section-nonfree}.
Considérons l'anneau\footnote{L'anneau $A$ est un exemple
d'anneau intègre non noethérien dont les anneaux locaux sont noethériens.}
$$
A = \text{Sym}^*_R(M)
$$
et l'idéal $I = M A = \bigoplus_{d \geq 1} \text{Sym}^d_R(M)$.
Comme $M$ n'est pas de type fini comme $R$-module,
$I$ ne peut être engendré par un nombre fini d'éléments comme idéal de $A$.
Puisqu'un module inversible est de type fini, il s'ensuit que
$I$ n'est pas inversible.
D'autre part, soit $\mathfrak p \subset R$ un idéal premier.
Par construction, $M_\mathfrak p \cong R_\mathfrak p$. Ainsi,
$$
A_\mathfrak p = \text{Sym}^*_{R_\mathfrak p}(M_\mathfrak p) \cong
\text{Sym}^*_{R_\mathfrak p}(R_\mathfrak p) =
R_\mathfrak p[T]
$$
comme $R_\mathfrak p$-algèbre graduée. Il s'ensuit que
$I_\mathfrak p \subset A_\mathfrak p$ est engendré par
le non-diviseur de zéro $T$.
Ainsi, pour tout idéal premier $\mathfrak q \subset A$,
$I_\mathfrak q$ est certainement engendré par un seul élément.

\begin{lemma}
\label{examples-lemma-locally-principal-not-invertible}
Il existe un anneau intègre $A$ et un idéal non nul $I \subset A$
tels que $I_\mathfrak q \subset A_\mathfrak q$ soit un idéal
principal pour tout idéal premier $\mathfrak q \subset A$, mais que $I$ ne soit pas un
$A$-module inversible.
\end{lemma}

\begin{proof}
Voir la discussion ci-dessus.
\end{proof}









\section{Un module plat de type fini qui n'est pas projectif}
\label{examples-section-finite-flat-not-projective}

\noindent
Il s'agit d'une copie de
Algèbre, remarque \ref{algebra-remark-warning}.
Il est faux qu'un $R$-module de type fini
plat sur $R$ soit automatiquement projectif. Un contre-exemple
s'obtient en prenant pour $R = \mathcal{C}^\infty(\mathbf{R})$
l'anneau des fonctions indéfiniment différentiables sur
$\mathbf{R}$, et $M = R_{\mathfrak m} = R/I$, où
$\mathfrak m = \{f \in R \mid f(0) = 0\}$ et
$I = \{f \in R \mid \exists \epsilon, \epsilon > 0 :
f(x) = 0\ \forall x, |x| < \epsilon\}$.

\medskip\noindent
Le morphisme $\Spec(R/I) \to \Spec(R)$ est également
un exemple d'immersion fermée plate qui n'est pas ouverte.

\begin{lemma}
\label{examples-lemma-finite-flat-non-projective}
Modules plats pathologiques.
\begin{enumerate}
\item Il existe un anneau $R$ et un $R$-module plat de type fini $M$ qui n'est pas
projectif.
\item Il existe une immersion fermée plate qui n'est pas ouverte.
\end{enumerate}
\end{lemma}

\begin{proof}
Voir la discussion ci-dessus.
\end{proof}



\section{Un module projectif qui n'est pas localement libre}
\label{examples-section-projective-not-locally-free}

\noindent
Nous donnons deux exemples : l'un où le rang est compris entre $0$ et $1$,
l'autre où le rang est $\aleph_0$.

\begin{lemma}
\label{examples-lemma-ideal-generated-by-idempotents-projective}
Soit $R$ un anneau. Soit $I \subset R$ un idéal engendré par
une famille dénombrable d'idempotents. Alors $I$ est projectif
comme $R$-module.
\end{lemma}

\begin{proof}
Écrivons $I = (e_1, e_2, e_3, \ldots)$, où $e_n$ est un idempotent de $R$.
En remplaçant par récurrence $e_{n + 1}$ par $e_n + (1 - e_n)e_{n + 1}$,
nous pouvons supposer que $(e_1) \subset (e_2) \subset (e_3) \subset \ldots$,
et donc que $I = \bigcup_{n \geq 1} (e_n) = \colim_n e_nR$.
Dans ce cas,
$$
\Hom_R(I, M) = \Hom_R(\colim_n e_nR, M)
= \lim_n \Hom_R(e_nR, M) = \lim_n e_nM
$$
Remarquons que les applications de transition $e_{n + 1}M \to e_nM$ sont données
par la multiplication par $e_n$ et sont surjectives. Par conséquent, d'après
Algèbre, lemme \ref{algebra-lemma-ML-exact-sequence},
le foncteur $\Hom_R(I, M)$ est exact, c'est-à-dire que $I$ est un
$R$-module projectif.
\end{proof}

\noindent
Supposons que $P \subset Q$ soit une inclusion de $R$-modules, avec $Q$
est un $R$-module de type fini et $P$ est localement libre; voir
Algèbre, définition \ref{algebra-definition-locally-free}.
Supposons que $Q$ puisse être engendré par $N$ éléments comme $R$-module.
Il résulte alors de
Algèbre, lemme \ref{algebra-lemma-map-cannot-be-injective}
que $P$ est localement libre de type fini (ses parties libres étant de rang
au plus $N$). Dans ce cas, $P$ est un $R$-module de type fini; voir
Algèbre, lemme \ref{algebra-lemma-finite-projective}.

\medskip\noindent
En combinant ceci avec ce qui précède, nous voyons qu'un idéal qui n'est pas
de type fini mais qui est engendré par une famille dénombrable d'idempotents
est projectif sans être localement libre. Un exemple explicite est
$R = \prod_{n \in \mathbf{N}} \mathbf{F}_2$, avec
$I$ l'idéal engendré par les idempotents
$$
e_n = (1, 1, \ldots, 1, 0, \ldots )
$$
où la suite de $1$ est de longueur $n$.

\begin{lemma}
\label{examples-lemma-ideal-projective-not-locally-free}
Il existe un anneau $R$ et un idéal $I$ tels que $I$ soit projectif comme
$R$-module, mais ne soit pas localement libre comme $R$-module.
\end{lemma}

\begin{proof}
Voir ci-dessus.
\end{proof}

\begin{remark}
\label{examples-remark-projective-with-dual-finite-not-finite}
L'exemple ci-dessus où $R = \prod_{n \in \mathbf{N}} \mathbf{F}_2$ et
$I$ est l'idéal engendré par les idempotents
$e_n = (1, 1, \ldots, 1, 0, \ldots )$ fournit aussi un exemple $M = I$
de module projectif qui n'est pas de type fini, mais dont le dual est de type
fini et même libre de type fini. En reprenant la méthode de la démonstration du
lemme \ref{examples-lemma-ideal-generated-by-idempotents-projective},
le lecteur montre en effet que $\Hom_R(I, R) \cong R$.
\end{remark}

\begin{lemma}
\label{examples-lemma-chow-group-product}
Soit $K$ un corps.
Soient $C_i$, $i = 1, \ldots, n$, des courbes lisses, projectives et géométriquement
irréductibles sur $K$. Soit $P_i \in C_i(K)$ un point rationnel et
soit $Q_i \in C_i$ un point tel que $[\kappa(Q_i) : K] = 2$.
Alors $[P_1 \times \ldots \times P_n]$ est non nul dans
$\CH_0(U_1 \times_K \ldots \times_K U_n)$, où $U_i = C_i \setminus \{Q_i\}$.
\end{lemma}

\begin{proof}
Il existe une application degré
$\deg : \CH_0(C_1 \times_K \ldots \times_K C_n) \to \mathbf{Z}$
Comme chaque $Q_i$ est de degré $2$ sur $K$, nous voyons que
tout cycle de dimension zéro dont le support est contenu dans le « bord »
$$
C_1 \times_K \ldots \times_K C_n
\setminus
U_1 \times_K \ldots \times_K U_n
$$
a un degré divisible par $2$.
\end{proof}

\noindent
Nous pouvons construire, à l'aide du lemme précédent, un autre exemple de
module projectif non localement libre. Soient
$C_n$, $n = 1, 2, 3, \ldots$, des courbes lisses, projectives et géométriquement
irréductibles sur $\mathbf{Q}$, chacune munie d'une paire de points
$P_n, Q_n \in C_n$ tels que $\kappa(P_n) = \mathbf{Q}$ et que
$\kappa(Q_n)$ soit une extension quadratique de $\mathbf{Q}$.
Posons $U_n = C_n \setminus \{Q_n\}$; c'est une courbe affine.
Soit $\mathcal{L}_n$ l'inverse du faisceau d'idéaux de $P_n$
sur $U_n$. Remarquons que $c_1(\mathcal{L}_n) = [P_n]$ dans le groupe des
cycles de dimension zéro $\CH_0(U_n)$. Posons $A_n = \Gamma(U_n, \mathcal{O}_{U_n})$.
Soit $L_n = \Gamma(U_n, \mathcal{L}_n)$, qui est un module localement
libre de rang $1$ sur $A_n$. Posons
$$
B_n = A_1 \otimes_{\mathbf{Q}} A_2 \otimes_{\mathbf{Q}} \ldots
\otimes_{\mathbf{Q}} A_n
$$
de sorte que $\Spec(B_n) = U_1 \times \ldots \times U_n$, tous les produits
étant pris sur $\Spec(\mathbf{Q})$. Pour $i \leq n$, posons
$$
L_{n, i} =
A_1 \otimes_{\mathbf{Q}} \ldots \otimes_{\mathbf{Q}} L_i
\otimes_{\mathbf{Q}} \ldots \otimes_{\mathbf{Q}} A_n
$$
qui est un $B_n$-module localement libre de rang $1$. Remarquons qu'il s'agit
aussi des sections globales de $\text{pr}_i^*\mathcal{L}_i$. Posons
$$
B_\infty = \colim_n B_n
\quad\text{et}\quad
L_{\infty, i} = \colim_n L_{n, i}
$$
Enfin, posons
$$
M = \bigoplus\nolimits_{i \geq 1} L_{\infty, i}.
$$
C'est une somme directe de modules localement libres de type fini, donc un module projectif.
Nous affirmons que $M$ n'est pas localement libre. Supposons en effet que
$f \in B_\infty$ soit une fonction non nulle telle que $M_f$ soit libre
sur $(B_\infty)_f$. Soit $e_1, e_2, \ldots$ une base. Choisissons
$n \geq 1$ tel que $f \in B_n$.
Choisissons $m \geq n + 1$ tel que $e_1, \ldots, e_{n + 1}$ appartiennent à
$$
\bigoplus\nolimits_{1 \leq i \leq m} L_{m, i}.
$$
Comme les éléments $e_1, \ldots, e_{n + 1}$ font partie d'une base
après un changement de base fidèlement plat, nous en déduisons que
les classes de Chern
$$
c_i(\text{pr}_1^*\mathcal{L}_1 \oplus \ldots \oplus
\text{pr}_m^*\mathcal{L}_m), \quad i = m, m - 1, \ldots, m - n
$$
sont nulles dans le groupe de Chow de
$$
D(f) \subset U_1 \times \ldots \times U_m
$$
Comme $f$ est l'image réciproque d'une fonction sur $U_1 \times \ldots \times U_n$,
il en résulte en particulier que
$$
c_{m - n}(\mathcal{O}_W^{\oplus n} \oplus
\text{pr}_1^*\mathcal{L}_{n + 1} \oplus \ldots
\oplus \text{pr}_{m - n}^*\mathcal{L}_m) = 0.
$$
sur la variété
$$
W = (C_{n + 1} \times \ldots \times C_m)_K
$$
sur le corps $K = \mathbf{Q}(C_1 \times \ldots \times C_n)$.
Autrement dit, le cycle
$$
[(P_{n + 1} \times \ldots \times P_m)_K]
$$
est nul dans le groupe de Chow des cycles de dimension zéro sur $W$. Ceci contredit le
lemme \ref{examples-lemma-chow-group-product}
ci-dessus, car les points $Q_i$, $n + 1 \leq i \leq m$,
induisent des points correspondants $Q_i'$ sur $(C_i)_K$ et, comme $K/\mathbf{Q}$
est géométriquement irréductible, nous avons $[\kappa(Q_i') : K] = 2$.

\begin{lemma}
\label{examples-lemma-projective-not-locally-free}
Il existe un anneau dénombrable $R$ et un module projectif $M$
qui est somme directe d'une famille dénombrable de modules localement libres de rang $1$,
mais tel que $M$ ne soit pas localement libre.
\end{lemma}

\begin{proof}
Voir ci-dessus.
\end{proof}







\section{Anneau local de dimension zéro possédant un idéal plat non nul}
\label{examples-section-zero-dimensional-flat-ideal}

\noindent
On trouve dans \cite{Lazard} et \cite{Autour}
un exemple d'anneau local de dimension zéro possédant un
idéal plat non nul. En voici la construction. Soit $k$ un corps.
Soient $X_i, Y_i$, $i \geq 1$, des indéterminées. Prenons
$R = k[X_i, Y_i]/(X_i - Y_i X_{i + 1}, Y_i^2)$. Notons $x_i$, respectivement,\ $y_i$
l'image de $X_i$, respectivement,\ $Y_i$ dans cet anneau. Remarquons que
$$
x_i = y_i x_{i + 1} = y_i y_{i +1} x_{i + 2} =
y_i y_{i + 1} y_{i + 2} x_{i + 3} = \ldots
$$
dans cet anneau. L'anneau $R$ ne possède qu'un
seul idéal premier, à savoir $\mathfrak m = (x_i, y_i)$. Nous affirmons que
l'idéal $I = (x_i)$ est plat comme $R$-module.

\medskip\noindent
Remarquons que l'annulateur de $x_i$ dans $R$ est l'idéal
$(x_1, x_2, x_3, \ldots, y_i, y_{i + 1}, y_{i + 2}, \ldots)$.
Considérons le $R$-module $M$ engendré par des éléments $e_i$, $i \geq 1$, et
les relations $e_i = y_i e_{i + 1}$. Alors $M$ est plat, car c'est la
limite inductive $\colim_i R$ de copies de $R$ dont les applications de transition sont
$$
R \xrightarrow{y_1} R \xrightarrow{y_2} R \xrightarrow{y_3} \ldots
$$
Remarquons que l'annulateur de $e_i$ dans $M$ est l'idéal
$(x_1, x_2, x_3, \ldots, y_i, y_{i + 1}, y_{i + 2}, \ldots)$.
Puisque tout élément de $M$, respectivement,\ de $I$, peut s'écrire
$f e_i$, respectivement,\ $h x_i$, pour certains $f, h \in R$, l'application
$M \to I$, $e_i \to x_i$, est un isomorphisme; ainsi, $I$ est plat.

\begin{lemma}
\label{examples-lemma-zero-dimensional-flat-ideal}
\begin{slogan}
Anneau de dimension zéro possédant un idéal plat.
\end{slogan}
Il existe un anneau local $R$ possédant un unique idéal premier
et un idéal non nul $I \subset R$ qui est un $R$-module plat.
\end{lemma}

\begin{proof}
Voir la discussion ci-dessus.
\end{proof}








\section[Épimorphisme d'anneaux de dimension zéro non surjectif]{Un épimorphisme d'anneaux de dimension zéro qui n'est pas surjectif}
\label{examples-section-epimorphism-not-surjective}

\noindent
On trouve dans \cite{Lazard-deux} et \cite{Autour} l'exemple suivant.
Soit $k$ un corps. Considérons l'homomorphisme d'anneaux
$$
k[x_1, x_2, \ldots, z_1, z_2, \ldots]/
(x_i^{4^i}, z_i^{4^i})
\longrightarrow
k[x_1, x_2, \ldots, y_1, y_2, \ldots]/
(x_i^{4^i}, y_i - x_{i + 1}y_{i + 1}^2)
$$
qui envoie $x_i$ sur $x_i$ et $z_i$ sur $x_iy_i$. Remarquons que $y_i^{4^{i + 1}}$
est nul dans le membre de droite, tandis que $y_1$ ne l'est pas (détails omis).
Cette application n'est pas surjective : on peut considérer ce qui précède comme une application
d'algèbres $\mathbf{Z}$-graduées en posant $\deg(x_i) = -1$, $\deg(z_i) = 0$
et $\deg(y_i) = 1$; il est alors clair que $y_1$ n'appartient pas à l'image.
Enfin, l'application est un épimorphisme, car
$$
y_{i - 1} \otimes 1 = x_i y_i^2 \otimes 1 = y_i \otimes x_i y_i =
x_i y_i \otimes y_i = 1 \otimes x_i y_i^2.
$$
par conséquent, le produit tensoriel du but sur la source est isomorphe
au but.

\begin{lemma}
\label{examples-lemma-epi-not-surjective}
Il existe un épimorphisme d'anneaux locaux de dimension $0$
qui n'est pas une surjection.
\end{lemma}

\begin{proof}
Voir la discussion ci-dessus.
\end{proof}



\section{De type fini, non de présentation finie, plat en un idéal premier}
\label{examples-section-ft-not-fp-flat-at-prime}

\noindent
Soit $k$ un corps. Considérons l'anneau local $A_0 = k[x, y]_{(x, y)}$.
Notons $\mathfrak p_{0, n} = (y + x^n + x^{2n + 1})$. C'est un idéal premier.
Posons
$$
A = A_0[z_1, z_2, z_3, \ldots]/(z_n z_m, z_n(y + x^n + x^{2n + 1}))
$$
Remarquons que $A \to A_0$ est une surjection dont le noyau est un idéal de
carré nul. Ainsi, $A$ est également un anneau local, et les idéaux premiers de $A$
sont en correspondance biunivoque avec ceux de $A_0$.
Notons $\mathfrak p_n$ l'idéal premier de $A$ correspondant à
$\mathfrak p_{0, n}$. Observons que $\mathfrak p_n$ est l'annulateur
de $z_n$ dans $A$. Soit
$$
C = A[z]/(xz^2 + z + y)[\frac{1}{2zx + 1}]
$$
Remarquons que $A \to C$ est un homomorphisme d'anneaux \'etale; voir
Algèbre, exemple \ref{algebra-example-make-standard-smooth}.
Soit $\mathfrak q \subset C$ l'idéal maximal engendré par
$x$, $y$, $z$ et tous les $z_n$. Comme $A \to C$ est plat,
l'annulateur de $z_n$ dans $C$ est $\mathfrak p_nC$. Calculons
\begin{align*}
C/\mathfrak p_n C
& =
A_0[z]/(xz^2 + z + y, y + x^n + x^{2n + 1})[1/(2zx + 1)] \\
& =
k[x]_{(x)}[z]/(xz^2 + z - x^n - x^{2n + 1})[1/(2zx + 1)] \\
& =
k[x]_{(x)}[z]/(z - x^n) \times
k[x]_{(x)}[z]/(xz + x^{n + 1} + 1)[1/(2zx + 1)] \\
& =
k[x]_{(x)} \times k(x)
\end{align*}
car $(z - x^n)(xz + x^{n + 1} + 1) = xz^2 + z - x^n - x^{2n + 1}$.
Nous obtenons donc $\mathfrak p_nC = \mathfrak r_n \cap \mathfrak q_n$,
où $\mathfrak r_n = \mathfrak p_nC + (z - x^n)C$ et
$\mathfrak q_n = \mathfrak p_nC + (xz + x^{n + 1} + 1)C$.
Comme $\mathfrak q_n + \mathfrak r_n = C$, nous avons aussi
$\mathfrak p_nC = \mathfrak r_n \mathfrak q_n$.
Il s'ensuit que $\mathfrak q_n$ est l'annulateur de $\xi_n = (z - x^n)z_n$.
Observons que, d'une part, $\mathfrak r_n \subset \mathfrak q$ et que,
d'autre part, $\mathfrak q_n + \mathfrak q = C$. Cela résulte par exemple du fait
que $\mathfrak q_n$ est un idéal maximal de $C$ distinct de $\mathfrak q$.
De même, $\mathfrak q_n + \mathfrak q_m = C$ pour $n \not = m$.
Posons maintenant
$$
B = \Im(C \longrightarrow C_{\mathfrak q})
$$
Les éléments $\xi_n$ s'envoient sur zéro dans $B$, car $xz + x^{n + 1} + 1$
n'appartient pas à $\mathfrak q$. Notons $\mathfrak q' \subset B$ l'image de
$\mathfrak q$. Par construction, $B$ est une $A$-algèbre de type fini et
$B_{\mathfrak q'} \cong C_{\mathfrak q}$. En particulier,
$B_{\mathfrak q'}$ est plat sur $A$.

\medskip\noindent
Nous affirmons qu'il n'existe aucun élément $g' \in B$, $g' \not \in \mathfrak q'$,
tel que $B_{g'}$ soit de présentation finie sur $A$. Esquissons la démonstration de
cette affirmation. Choisissons un élément $g \in C$ dont l'image est $g' \in B$.
Considérons l'application $C_g \to B_{g'}$. D'après
Algèbre, lemme \ref{algebra-lemma-finite-presentation-independent},
$B_{g'}$ est de présentation finie sur $A$ si et seulement si le noyau
de $C_g \to B_{g'}$ est de type fini. Mais l'élément $g \in C$
n'appartient pas à $\mathfrak q$ et s'envoie donc sur un élément non nul de
$A_0[z]/(xz^2 + z + y)$. Par conséquent, $g$ ne peut appartenir qu'à un nombre fini
des idéaux premiers $\mathfrak q_n$, car les idéaux premiers
$(y + x^n + x^{2n + 1}, xz + x^{n + 1} + 1)$ forment une famille infinie
de points de codimension 1 de l'espace noethérien irréductible de dimension 2
$\Spec(k[x, y, z]/(xz^2 + z + y))$. L'application
$$
\bigoplus\nolimits_{g \not \in \mathfrak q_n} C/\mathfrak q_n
\longrightarrow
C_g, \quad
(c_n) \longrightarrow \sum c_n \xi_n
$$
est injective et son image est le noyau de $C_g \to B_{g'}$. Nous omettons la
démonstration de cet énoncé. (Indication : écrire $A = A_0 \oplus I$ comme $A_0$-module,
où $I$ est le noyau de $A \to A_0$. De même, écrire $C = C_0 \oplus IC$.
Écrire
$IC = \bigoplus Cz_n \cong \bigoplus (C/\mathfrak r_n \oplus C/\mathfrak q_n)$
et étudier l'effet de la multiplication par $g$ sur les facteurs.)
Ceci achève l'esquisse de la démonstration de l'affirmation.
Cela montre également que $B_{g'}$ n'est plat sur $A$ pour aucun $g'$ comme ci-dessus.
En effet, s'il était plat, l'annulateur de l'image de $z_n$ dans
$B_{g'}$ serait $\mathfrak p_nB_{g'}$ et ne contiendrait pas $z - x^n$.

\medskip\noindent
Nous pouvons par conséquent répondre par la négative à une question posée dans
\cite[Part I, Remarques (3.4.7) (\romannumeral 5)]{GruRay}.
En voici un énoncé précis.

\begin{lemma}
\label{examples-lemma-example-raynaud-gruson}
Il existe un anneau local $A$, un homomorphisme d'anneaux de type fini $A \to B$ et un idéal premier
$\mathfrak q$ au-dessus de $\mathfrak m_A$ tels que $B_{\mathfrak q}$ soit plat
sur $A$ et que, pour tout élément $g \in B$, $g \not \in \mathfrak q$,
l'anneau $B_g$ ne soit ni de présentation finie sur $A$, ni plat sur $A$.
\end{lemma}

\begin{proof}
Voir la discussion ci-dessus.
\end{proof}


\section{De type fini et plat, mais non de présentation finie}
\label{examples-section-finite-type-flat-not-finite-presentation}

\noindent
Dans cette section, nous donnons des exemples d'homomorphismes d'anneaux et de morphismes
qui sont de type fini et plats, mais non de présentation finie.

\medskip\noindent
Soit $R$ un anneau possédant un idéal $I$ tel que $R/I$ soit un module
plat de type fini mais non projectif; voir la
section \ref{examples-section-finite-flat-not-projective}
pour un exemple explicite. Cela signifie que $I$ n'est pas
de type fini; voir
Algèbre, lemme \ref{algebra-lemma-finitely-generated-pure-ideal}.
Remarquons que $I = I^2$; voir
Algèbre, lemme \ref{algebra-lemma-pure}.
Dans nos exemples, l'anneau de base sera
$R$ et, par conséquent, le schéma de base $S = \Spec(R)$.

\medskip\noindent
Considérons l'homomorphisme d'anneaux $R \to R \oplus R/I\epsilon$, où
$\epsilon^2 = 0$ par convention. C'est un homomorphisme fini et plat
qui n'est pas de présentation finie. Tous les anneaux des fibres
sont des intersections complètes et sont géométriquement irréductibles.

\medskip\noindent
Soit $A = R[x, y]/(xy, ay; a \in I)$. Remarquons que, comme $R$-module,
nous avons $A = \bigoplus_{i \geq 0} Ry^i \oplus \bigoplus_{j > 0} R/Ix^j$.
Ainsi, $R \to A$ est un homomorphisme d'anneaux plat de type fini qui n'est pas de
présentation finie. Chaque anneau de fibre
est isomorphe soit à $\kappa(\mathfrak p)[x, y]/(xy)$,
soit à $\kappa(\mathfrak p)[x]$.

\medskip\noindent
Nous pouvons transformer l'exemple précédent en un morphisme projectif en
prenant $B = R[X_0, X_1, X_2]/(X_1X_2, aX_2; a \in I)$.
Dans ce cas, $X = \text{Proj}(B) \to S$ est un morphisme propre et plat
qui n'est pas de présentation finie et tel que, pour tout $s \in S$,
la fibre $X_s$ soit isomorphe soit à $\mathbf{P}^1_s$, soit au
sous-schéma fermé de $\mathbf{P}^2_s$ défini par l'annulation de
$X_1X_2$ (c'est une courbe nodale projective de genre arithmétique $0$).

\medskip\noindent
Soit $M = R \oplus R \oplus R/I$. Posons $B = \text{Sym}_R(M)$, l'algèbre
symétrique de $M$. Posons $X = \text{Proj}(B)$.
Alors $X \to S$ est un morphisme propre et plat, non de présentation finie,
tel que, pour $s \in S$, la fibre géométrique soit isomorphe soit à
$\mathbf{P}^1_s$, soit à $\mathbf{P}^2_s$. En particulier, ces fibres sont
lisses et géométriquement irréductibles.

\begin{lemma}
\label{examples-lemma-finite-type-flat-not-finitely-presented}
Il existe des exemples de
\begin{enumerate}
\item homomorphisme d'anneaux plat de type fini, non de présentation finie,
dont les anneaux des fibres sont des intersections complètes géométriquement irréductibles;
\item homomorphisme d'anneaux plat de type fini, non de présentation finie,
dont les anneaux des fibres sont des intersections complètes de dimension 1,
géométriquement connexes et géométriquement réduites;
\item morphisme propre et plat de schémas $X \to S$, non de présentation finie,
dont chaque fibre est isomorphe soit à $\mathbf{P}^1_s$, soit au lieu d'annulation de
$X_1X_2$ dans $\mathbf{P}^2_s$; et
\item morphisme propre et plat de schémas $X \to S$ dont chaque
fibre est isomorphe soit à $\mathbf{P}^1_s$, soit à $\mathbf{P}^2_s$,
qui n'est pas de présentation finie.
\end{enumerate}
\end{lemma}

\begin{proof}
Voir la discussion ci-dessus.
\end{proof}



\section{Topologie d'un homomorphisme d'anneaux de type fini}
\label{examples-section-topology-finite-type}

\noindent
Soit $A \to B$ un homomorphisme local d'anneaux locaux intègres.
Si $A$ est noethérien, si $A \to B$ est essentiellement de type fini
et si $A/\mathfrak m_A \subset B/\mathfrak m_B$ est fini, alors il
existe un idéal premier $\mathfrak q \subset B$, $\mathfrak q \not = \mathfrak m_B$,
tel que $A \to B/\mathfrak q$ soit la localisation d'un homomorphisme d'anneaux
quasi-fini. Voir
Compléments sur les morphismes, lemme
\ref{more-morphisms-lemma-quasi-finite-quasi-section-meeting-nearby-open}.

\medskip\noindent
Dans cette section, nous donnons un exemple montrant que ce résultat devient faux si $A$
n'est plus noethérien. Soit en effet $k$ un corps et posons
$$
A = \{a_0 + a_1 x + a_2 x^2 + \ldots \mid a_0 \in k, a_i \in k((y))
\text{ pour }i\geq 1\}
$$
et
$$
C = \{a_0 + a_1 x + a_2 x^2 + \ldots \mid a_0 \in k[y], a_i \in k((y))
\text{ pour }i\geq 1\}.
$$
L'inclusion $A \to C$ est de type fini, car $C$ est engendré par $y$
sur $A$. Nous affirmons que $A$ est un anneau local d'idéal maximal
$\mathfrak m = \{a_1 x + a_2 x^2 + \ldots \in A\}$, sans autre idéal
premier que $(0)$ et $\mathfrak m$. En effet, un élément
$f = a_0 + a_1 x + a_2 x^2 + \ldots$ de $A$ est inversible dès que
$a_0 \not = 0$. Si $\mathfrak q \subset A$ est un idéal premier non nul
et si $f = a_i x^i + \ldots \in \mathfrak q$, les propriétés des
séries formelles montrent que, pour tout $g \in k((y))$, l'élément
$g^{i + 1} x^{i + 1} \in \mathfrak q$, c'est-à-dire $gx \in \mathfrak q$.
Cela prouve que $\mathfrak q = \mathfrak m$.

\medskip\noindent
Quant au spectre de l'anneau $C$, le même raisonnement montre
que tout idéal premier non nul contient l'idéal premier
$\mathfrak p = \{a_1 x + a_2 x^2 + \ldots \in C\}$, qui se trouve
au-dessus de $\mathfrak m$. Ainsi, le seul idéal premier de $C$ au-dessus de
$(0)$ est $(0)$. Posons $\mathfrak m_C = yC + \mathfrak p$ et
$B = C_{\mathfrak m_C}$. Alors $A \to B$ est l'exemple recherché.

\begin{lemma}
\label{examples-lemma-topology-finite-type}
Il existe un homomorphisme local $A \to B$ d'anneaux locaux intègres qui est
essentiellement de type fini et tel que $A/\mathfrak m_A \to B/\mathfrak m_B$
soit fini, avec la propriété que, pour tout idéal premier
$\mathfrak q \not = \mathfrak m_B$ de $B$, l'homomorphisme d'anneaux
$A \to B/\mathfrak q$ n'est pas la localisation d'un homomorphisme quasi-fini.
\end{lemma}

\begin{proof}
Voir la discussion ci-dessus.
\end{proof}



\section{Pur mais non universellement pur}
\label{examples-section-pure-not-universally}

\noindent
Soit $k$ un corps. Soit
$$
R = k[[x, xy, xy^2, \ldots]] \subset k[[x, y]].
$$
Autrement dit, une série formelle $f \in k[[x, y]]$ appartient à $R$ si et seulement si
$f(0, y)$ est constante. En particulier, $R[1/x] = k[[x, y]][1/x]$
et $R/xR$ est un anneau local dont l'idéal maximal est de carré
nul. Notons $R[y] \subset k[[x, y]]$ l'ensemble des séries formelles
$f \in k[[x, y]]$ telles que $f(0, y)$ soit un polynôme en $y$. Alors
$R \to R[y]$ est un homomorphisme d'anneaux de type fini mais non de présentation finie
qui induit un isomorphisme après inversion de $x$. Il existe aussi une
surjection $R[y]/xR[y] \to k[y]$ dont le noyau est de carré nul. Considérons
l'homomorphisme d'anneaux de présentation finie $R \to S = R[t]/(xt - xy)$.
De nouveau, $R[1/x] \to S[1/x]$ est un isomorphisme et, dans ce cas,
$S/xS \cong (R/xR)[t]/(xy)$ se surjecte sur $k[t]$ avec un noyau nilpotent.
Il existe une surjection $S \to R[y]$, $t \longmapsto y$, qui induit
un isomorphisme après inversion de $x$ et une surjection de noyau nilpotent
modulo $x$. Le noyau de $S \to R[y]$ est donc localement nilpotent.
En particulier, $S \to R[y]$ induit un homéomorphisme universel sur les spectres.

\medskip\noindent
Nous affirmons d'abord que $S$, considéré comme $S$-module, est relativement pur
sur $R$. Comme l'inversion de $x$ donne un isomorphisme, il suffit
de le vérifier en l'idéal maximal $\mathfrak m \subset R$. Puisque
$R$ est complet pour son idéal maximal, il est hensélien;
il suffit donc de vérifier que, pour tout idéal premier $\mathfrak p \subset R$,
$\mathfrak p \not = \mathfrak m$, l'unique idéal premier $\mathfrak q$
de $S$ au-dessus de $\mathfrak p$ satisfait
$\mathfrak mS + \mathfrak q \not = S$. Comme $\mathfrak p \not = \mathfrak m$,
il correspond à un unique idéal premier de $k[[x, y]][1/x]$. Ainsi,
soit $\mathfrak p = (0)$, soit $\mathfrak p = (f)$ pour un certain
élément irréductible $f \in k[[x, y]]$ non associé à $x$
(nous utilisons ici que $k[[x, y]]$ est factoriel -- insérer ici une référence ultérieure).
Dans le premier cas, $\mathfrak q = (0)$ et le résultat est clair. Dans le
second, nous pouvons multiplier $f$ par une unité afin que $f \in R[y]$
(préparation de Weierstrass; détails omis). On voit alors facilement que
$R[y]/fR[y] \cong k[[x, y]]/(f)$; ainsi, $f$ définit un idéal premier
de $R[y]$ et $\mathfrak mR[y] + fR[y] \not = R[y]$.
Comme $S \to R[y]$ induit un homéomorphisme universel sur les spectres, nous en déduisons
également le résultat voulu pour $S$.

\medskip\noindent
Nous affirmons ensuite que $S$ n'est pas universellement relativement pur sur $R$.
Pour le voir, il suffit de trouver un anneau de valuation
$\mathcal{O}$ et un homomorphisme local $R \to \mathcal{O}$ tels que
$\Spec(R[y] \otimes_R \mathcal{O}) \to \Spec(\mathcal{O})$
n'atteigne pas le point fermé de $\Spec(\mathcal{O})$.
De manière équivalente, il faut trouver $\varphi : R \to \mathcal{O}$ tel que
$\varphi(x) \not = 0$ et $v(\varphi(x)) > v(\varphi(xy))$, où $v$
est la valuation de $\mathcal{O}$.
(En effet, cela signifie que la valuation de $y$ est négative.)
Considérons pour cela l'homomorphisme d'anneaux
$$
R
\longrightarrow
\{a_0 + a_1 x + a_2 x^2 + \ldots \mid a_0 \in k[y^{-1}], a_i \in k((y))\}
$$
défini de manière évidente. On peut trouver un anneau de valuation $\mathcal{O}$
dominant le localisé du membre de droite en l'idéal
maximal $(y^{-1}, x)$, ce qui conclut.

\begin{lemma}
\label{examples-lemma-pure-not-universally-pure}
Il existe un morphisme de schémas affines de présentation finie
$X \to S$ et un $\mathcal{O}_X$-module $\mathcal{F}$ de présentation finie
tels que $\mathcal{F}$ soit pur relativement à $S$, mais non universellement
pur relativement à $S$.
\end{lemma}

\begin{proof}
Voir la discussion ci-dessus.
\end{proof}



\section[Homomorphisme formellement lisse non plat]{Un homomorphisme d'anneaux formellement lisse qui n'est pas plat}
\label{examples-section-formally-smooth-nonflat}

\noindent
Soit $k$ un corps. Considérons la $k$-algèbre $k[\mathbf{Q}]$.
C'est la $k$-algèbre ayant pour base les $x_\alpha, \alpha \in \mathbf{Q}$,
dont la multiplication est déterminée par $x_\alpha x_\beta = x_{\alpha + \beta}$.
(En particulier, $x_0 = 1$.)
Considérons l'homomorphisme de $k$-algèbres
$$
k[\mathbf{Q}] \longrightarrow k, \quad
x_\alpha \longmapsto 1.
$$
Il est surjectif, de noyau $J$ engendré par les éléments $x_\alpha - 1$.
Calculons $J/J^2$. Remarquons que la multiplication par $x_\alpha$ sur $J/J^2$
est l'application identité. Notons $z_\alpha$ la classe de $x_\alpha - 1$ modulo
$J^2$. Ces classes engendrent $J/J^2$. Comme
$$
(x_\alpha - 1)(x_\beta - 1) = x_{\alpha + \beta} - x_\alpha - x_\beta + 1 =
(x_{\alpha + \beta} - 1) - (x_\alpha - 1) - (x_\beta - 1)
$$
nous avons $z_{\alpha + \beta} = z_\alpha + z_\beta$ dans $J/J^2$.
Un élément général de $J/J^2$ est de la forme $\sum \lambda_\alpha z_\alpha$,
où $\lambda_\alpha \in k$ (un nombre fini seulement étant non nuls). Si la
caractéristique de $k$ est $p > 0$, alors
$$
0 = pz_{\alpha/p} = z_{\alpha/p} + \ldots + z_{\alpha/p} = z_\alpha
$$
et nous obtenons $J/J^2 = 0$. Si la caractéristique de $k$ est nulle, alors
$$
J/J^2 = \mathbf{Q} \otimes_{\mathbf{Z}} k \cong k
$$
(détails omis), qui n'est pas nul.

\medskip\noindent
Nous affirmons que $k[\mathbf{Q}] \to k$ est un homomorphisme d'anneaux formellement lisse
si la caractéristique de $k$ est positive. Supposons en effet donné un
diagramme commutatif formé des flèches pleines
$$
\xymatrix{
k \ar[r] \ar@{..>}[rd] & A \\
k[\mathbf{Q}] \ar[u] \ar[r]^\varphi & A' \ar[u]
}
$$
où $A' \to A$ est une surjection dont le noyau $I$ est de carré nul.
Pour montrer que $k[\mathbf{Q}] \to k$ est formellement lisse, il faut prouver
que $\varphi$ se factorise par $k$. Comme $\varphi(x_\alpha - 1)$ s'envoie
sur zéro dans $A$, l'application $\varphi$ induit une application
$\overline{\varphi} : J/J^2 \to I$ dont l'annulation est l'obstruction
à la factorisation voulue. Puisque $J/J^2 = 0$ lorsque la caractéristique
est $p > 0$, nous obtenons le résultat voulu : $k[\mathbf{Q}] \to k$ est
formellement lisse dans ce cas. Enfin, cet homomorphisme n'est pas plat :
par exemple, le non-diviseur de zéro $x_2 - 1$ s'envoie sur zéro.

\begin{lemma}
\label{examples-lemma-formally-smooth-nonflat}
Il existe un homomorphisme d'anneaux formellement lisse qui n'est pas plat.
\end{lemma}

\begin{proof}
Voir la discussion ci-dessus.
\end{proof}


\section{Un homomorphisme d'anneaux formellement \'etale qui n'est pas plat}
\label{examples-section-formally-etale-nonflat}

\noindent
Dans cette section, nous donnons un contre-exemple à la dernière phrase de
\cite[0, exemple 19.10.3(i)]{EGA} (ce point ne figurait pas parmi ceux relevés
dans leurs listes ultérieures d'errata). Considérons $A \to A/J$, où $A$ est un anneau local
et $J$ un idéal propre non nul tel que $J^2 = J$ (ainsi, $J$ n'est pas de type
fini); l'anneau de valuation d'un corps non archimédien algébriquement clos,
avec pour $J$ son idéal maximal, fournit un tel $(A, J)$. Ces
homomorphismes quotients non plats sont formellement \'etales. Supposons en effet donné un
diagramme commutatif
$$
\xymatrix{
A/J \ar[r] & R/I \\
A \ar[u] \ar[r]^\varphi & R \ar[u]
}
$$
où $I$ est un idéal de l'anneau $R$ tel que $I^2 = 0$. Alors $A \to R$
se factorise de manière unique par $A/J$, car
$$
\varphi(J) = \varphi(J^2) \subset (\varphi(J)R)^2 \subset I^2 = 0.
$$
Cela fournit également un contre-exemple au cas formellement \'etale
du « théorème de structure » pour les morphismes localement de type fini et formellement \'etales
de \cite[IV, théorème 18.4.6(i)]{EGA}
(mais non à la partie (ii), qui est celle que l'on emploie effectivement
en pratique). L'erreur dans la démonstration de ce dernier énoncé
consiste, à sa toute dernière étape, à invoquer l'énoncé incorrect
\cite[0, exemple 19.3.10(i)]{EGA};
c'est ainsi que le contre-exemple précédent s'y introduit.

\begin{lemma}
\label{examples-lemma-formally-etale-not-presented}
Il existe des homomorphismes d'anneaux formellement \'etales non plats.
\end{lemma}

\begin{proof}
Voir la discussion ci-dessus.
\end{proof}



\section{Un homomorphisme d'anneaux formellement \'etale à complexe cotangent non trivial}
\label{examples-section-formally-etale-nontrivial-cotangent-complex}

\noindent
Soit $k$ un corps. Considérons l'anneau
$$
R = k[\{x_n\}_{n \geq 1}, \{y_n\}_{n \geq 1}]/(
x_1y_1, x_{nm}^m - x_n, y_{nm}^m - y_n)
$$
Soit $A$ le localisé en l'idéal maximal engendré par
tous les $x_n, y_n$, et notons $J \subset A$ l'idéal maximal. Posons $B = A/J$.
Par construction, $J^2 = J$; ainsi, $A \to B$ est formellement \'etale (voir la
section \ref{examples-section-formally-etale-nonflat}).
Nous affirmons que l'élément $x_1 \otimes y_1$ est un élément
non nul du noyau de
$$
J \otimes_A J \longrightarrow J.
$$
En effet, $(A, J)$ est la limite inductive des localisés
$(A_n, J_n)$ des anneaux
$$
R_n = k[x_n, y_n]/(x_n^n y_n^n)
$$
en leurs idéaux maximaux correspondants. Alors
$x_1 \otimes y_1$ correspond à l'élément
$x_n^n \otimes y_n^n \in J_n \otimes_{A_n} J_n$, qui est non nul
(par un calcul explicite que nous omettons). Puisque $\otimes$ commute
aux limites inductives, nous concluons. D'après
\cite[III section 3.3]{cotangent}
nous voyons que $J$ n'est pas faiblement régulier. Par conséquent, d'après
\cite[III Proposition 3.3.3]{cotangent}
nous voyons que le complexe cotangent $L_{B/A}$ n'est pas nul. En fait, nous pouvons
être plus précis. Nous avons $H_0(L_{B/A}) = \Omega_{B/A}$ et
$H_1(L_{B/A}) = 0$, car $J/J^2 = 0$. Mais la suite exacte à cinq termes
de la suite spectrale fondamentale de Quillen
(voir Cotangent, remarque \ref{cotangent-remark-elucidate-ss}, ou
\cite[Corollary 8.2.6]{Reinhard})
et la non-annulation de
$\text{Tor}_2^A(B, B) = \Ker(J \otimes_A J \to J)$ montrent que
$H_2(L_{B/A})$ est non nul.

\begin{lemma}
\label{examples-lemma-formally-etale-nontrivial-cotangent-complex}
Il existe un homomorphisme d'anneaux surjectif et formellement \'etale $A \to B$
tel que $L_{B/A}$ ne soit pas nul.
\end{lemma}

\begin{proof}
Voir la discussion ci-dessus.
\end{proof}






\section[Plat et formellement non ramifié, non formellement \'etale]{Plat et formellement non ramifié n'implique pas formellement \'etale}
\label{examples-section-flat-formally-unramified-not-formally-etale}

\noindent
Dans Compléments sur les morphismes, lemme
\ref{more-morphisms-lemma-unramified-flat-formally-etale}
il est montré qu'un morphisme de schémas non ramifié et plat $X \to S$
est formellement \'etale. Le but de cette section est de donner deux exemples
montrant qu'on ne peut remplacer « non ramifié » par « formellement non ramifié ».
Le premier exploite des propriétés particulières des anneaux parfaits, tandis que le
second montre que le résultat échoue même pour des homomorphismes d'anneaux réguliers noethériens.

\begin{lemma}
\label{examples-lemma-perfect-closure-polynomial-ring}
Soit $A = \mathbb{F}_p[T]$ l'anneau de polynômes en une indéterminée sur
$\mathbb{F}_p$. Notons $A_{perf}$ la clôture parfaite de $A$.
Alors $A \rightarrow A_{perf}$ est plat et formellement non ramifié,
mais non formellement \'etale.
\end{lemma}

\begin{proof}
Remarquons que, pour le Frobenius $F_A : A \to A$, la copie d'arrivée de $A$
est un module libre sur la copie de départ, de base $\{1, T, \dots, T^{p - 1}\}$.
Ainsi, $F_A$ est fidèlement plat et, par conséquent,
$A \to A_{perf}$ l'est aussi, puisqu'il est une limite inductive d'homomorphismes fidèlement plats.
Comme $A_{perf}$ est un anneau parfait, le Frobenius relatif
$F_{A_{perf}/A}$ est surjectif. Autrement dit,
$A_{perf} = A[A_{perf}^p]$, ce qui implique immédiatement
$\Omega_{A_{perf}/A} = 0$. Alors
$A \rightarrow A_{perf}$ est formellement non ramifié d'après
Compléments sur les morphismes, lemme
\ref{more-morphisms-lemma-formally-unramified-differentials}

\medskip\noindent
Il suffit de montrer que $A \rightarrow A_{perf}$
n'est pas formellement lisse. Comme $A$ est une
$\mathbb{F}_p$-algèbre lisse, le complexe cotangent
$L_{A/\mathbb{F}_p} \simeq \Omega_{A/\mathbb{F}_p}[0]$
est concentré en degré $0$; voir
Cotangent, lemme \ref{cotangent-lemma-when-projective}. De plus,
$L_{A_{perf}/\mathbb{F}_p} = 0$ dans $D(A_{perf})$
d'après Cotangent, lemme \ref{cotangent-lemma-perfect-zero}.
Considérons le triangle distingué de complexes cotangents
$$
L_{A/\mathbb{F}_p} \otimes_A A_{perf} \to
L_{A_{perf}/\mathbb{F}_p} \to
L_{A_{perf}/A} \to
(L_{A/\mathbb{F}_p} \otimes_A A_{perf})[1]
$$
dans $D(A_{perf})$; voir Cotangent, section \ref{cotangent-section-triangle}.
Nous obtenons $L_{A_{perf}/A} = \Omega_{A/\mathbb{F}_p} \otimes_A A_{perf}[1]$,
c'est-à-dire que $L_{A_{perf}/A}$ est égal à un module libre de rang $1$ sur $A_{perf}$
placé en degré $-1$. Ainsi, $A \rightarrow A_{perf}$
n'est pas formellement lisse d'après
Compléments sur les morphismes, lemme \ref{more-morphisms-lemma-NL-formally-smooth},
et
Cotangent, lemme \ref{cotangent-lemma-relation-with-naive-cotangent-complex}.
\end{proof}

\noindent
L'exemple suivant porte également sur des anneaux de caractéristique positive, mais il est peut-être
un peu plus surprenant. Son inconvénient est qu'il requiert une meilleure connaissance des
phénomènes de caractéristique $p$ que l'exemple précédent. Rappelons qu'un
anneau $A$ de caractéristique positive est dit $F$-fini si le Frobenius de $A$
est fini.

\begin{lemma}
\label{examples-lemma-completion-etale}
Soit $(A, \mathfrak m, \kappa)$ un anneau local noethérien de caractéristique
$p > 0$ tel que $[\kappa : \kappa^p] < \infty$.
Alors l'application canonique $A \to A^\wedge$ vers le complété de $A$
est plate et formellement non ramifiée. Cependant, si $A$ est régulier mais non
excellent, cette application n'est pas formellement \'etale.
\end{lemma}

\begin{proof}
La platitude du complété est donnée par
Algèbre, lemme \ref{algebra-lemma-completion-flat}.
Pour montrer que l'application est formellement non ramifiée, il
suffit de montrer que $\Omega_{A^\wedge/A} = 0$; voir
Algèbre, lemme \ref{algebra-lemma-characterize-formally-unramified}.

\medskip\noindent
Esquissons une démonstration. Choisissons $x_1, \ldots, x_r \in A$ dont les images forment une $p$-base
$\overline{x}_1, \ldots, \overline{x}_r$ de $\kappa$, c'est-à-dire
telle que $\kappa$ soit engendré minimalement par les $\overline{x}_i$ sur $\kappa^p$.
Choisissons un système minimal de générateurs $y_1, \ldots, y_s$ de $\mathfrak m$.
Pour tout $n$, les éléments $x_1, \ldots, x_r, y_1, \ldots, y_s$ engendrent
$A/\mathfrak m^n$ sur $(A/\mathfrak m^n)^p$ par Frobenius.
Certains détails sont omis. Nous en déduisons que $F : A^\wedge \to A^\wedge$
est fini. Par conséquent, $\Omega_{A^\wedge/A}$ est un $A^\wedge$-module de type fini.
D'autre part, pour tout $a \in A^\wedge$ et tout $n$, il existe
$a_0 \in A$ tel que $a - a_0 \in \mathfrak m^nA^\wedge$.
Nous en déduisons que $\text{d}(a) \in \bigcap \mathfrak m^n \Omega_{A^\wedge/A}$,
ce qui implique que $\text{d}(a)$ est nul d'après
Algèbre, lemme \ref{algebra-lemma-intersect-powers-ideal-module-zero}.
Ainsi, $\Omega_{A^\wedge/A} = 0$.

\medskip\noindent
Supposons $A$ régulier. Alors, par le théorème de structure de Cohen,
$x_1, \ldots, x_r, y_1, \ldots, y_s$ est une $p$-base de l'anneau
$A^\wedge$, c'est-à-dire que nous avons
$$
A^\wedge = \bigoplus\nolimits_{I, J} (A^\wedge)^p
x_1^{i_1} \ldots x_r^{i_r} y_1^{j_1} \ldots y_s^{j_s}
$$
où $I = (i_1, \ldots, i_r)$, $J = (j_1, \ldots, j_s)$ et
$0 \leq i_a, j_b \leq p - 1$. Détails omis. En particulier,
$\Omega_{A^\wedge}$ est un $A^\wedge$-module libre ayant pour base
$\text{d}(x_1), \ldots, \text{d}(x_r), \text{d}(y_1), \ldots, \text{d}(y_s)$.

\medskip\noindent
Si maintenant $A \to A^\wedge$ est formellement \'etale, ou même seulement formellement lisse,
alors $\NL_{A^\wedge/A}$ a une cohomologie nulle en degrés $-1, 0$
d'après Algèbre, proposition \ref{algebra-proposition-characterize-formally-smooth}.
La suite de Jacobi-Zariski
(Algèbre, lemme \ref{algebra-lemma-exact-sequence-NL}) appliquée aux homomorphismes d'anneaux
$\mathbf{F}_p \to A \to A^\wedge$ donne alors un isomorphisme
$\Omega_A \otimes_A A^\wedge \cong \Omega_{A^\wedge}$.
Nous obtenons donc que $\Omega_A$ est libre ayant pour base
$\text{d}(x_1), \ldots, \text{d}(x_r), \text{d}(y_1), \ldots, \text{d}(y_s)$.
En considérant les corps des fractions et en utilisant la normalité de $A$,
nous concluons que $a \in A$ est une puissance $p$-ième si et seulement si son image dans
$A^\wedge$ est une puissance $p$-ième (détails omis; utiliser
Algèbre, lemme \ref{algebra-lemma-derivative-zero-pth-power}).
Une seconde conséquence est que les opérateurs $\partial/\partial x_a$ et
$\partial/\partial y_b$ sont définis sur $A$.

\medskip\noindent
Nous allons montrer que ce qui précède implique que $A$ est fini
sur $A^p$, avec pour $p$-base $x_1, \ldots, x_r, y_1, \ldots, y_s$.
Cela contredira la non-excellence de $A$ par un résultat de
Kunz; voir \cite[Corollary 2.6]{Kun76}.
Soit en effet $a \in A$ et écrivons
$$
a =  \sum\nolimits_{I, J} (a_{I, J})^p
x_1^{i_1} \ldots x_r^{i_r} y_1^{j_1} \ldots y_s^{j_s}
$$
avec $a_{I, J} \in A^\wedge$. Pour achever la démonstration, il suffit de
montrer que $a_{I, J} \in A$. En appliquant l'opérateur
$$
(\partial/\partial x_1)^{p - 1} \ldots
(\partial/\partial x_r)^{p - 1}
(\partial/\partial y_1)^{p - 1} \ldots
(\partial/\partial y_s)^{p - 1}
$$
aux deux membres, nous concluons que $a_{I, J}^p \in A$, où
$I = (p - 1, \ldots, p - 1)$ et $J = (p - 1, \ldots, p - 1)$.
D'après la remarque précédente, cela implique aussi $a_{I, J} \in A$.
Après avoir remplacé $a$ par $a' = a - a_{I, J}^p x^I y^J$,
nous pouvons employer un opérateur différentiel d'ordre abaissé de $1$ pour obtenir
qu'un autre coefficient $a_{I, J}$ appartient à $A$, et ainsi de suite.
\end{proof}

\begin{remark}
\label{examples-remark-reference-existence-regular-nonexcellent-rings}
On peut construire des anneaux réguliers non excellents dont le corps résiduel a une $p$-base finie,
même dans le corps de fonctions de $\mathbb{P}^2_k$, sur un corps de
caractéristique $p$, $k = \overline{k}$. Voir
\cite[$\mathsection 4.1$]{DS18}.
\end{remark}

\noindent
La démonstration du lemme \ref{examples-lemma-completion-etale} montre en fait un peu plus.

\begin{lemma}
\label{examples-lemma-excellent-regular-local-rings}
Soit $(A, \mathfrak m, \kappa)$ un anneau local régulier de caractéristique
$p > 0$. Supposons $[\kappa : \kappa^p] < \infty$. Alors $A$ est excellent
si et seulement si $A \to A^\wedge$ est formellement \'etale.
\end{lemma}

\begin{proof}
L'implication réciproque résulte du lemme \ref{examples-lemma-completion-etale}.
Pour l'implication directe, le
lemme \ref{examples-lemma-completion-etale}
montre déjà que $A \to A^\wedge$ est formellement non ramifié, ou, de manière équivalente,
que $\Omega_{A^\wedge/A}$ est nul.
Il suffit donc de montrer que l'application de complétion est formellement lisse lorsque
$A$ est excellent. Par la désingularisation de N\'eron-Popescu,
$A \to A^\wedge$ s'écrit comme une limite inductive filtrante
de $A$-algèbres lisses
(Lissage d'homomorphismes d'anneaux, théorème \ref{smoothing-theorem-popescu}).
Par conséquent, $\NL_{A^\wedge/A}$ a une cohomologie nulle en degré $-1$.
Ainsi, $A \to A^\wedge$ est formellement lisse d'après Algèbre, proposition
\ref{algebra-proposition-characterize-formally-smooth}.
\end{proof}










\section[Idéaux d'idempotents et localisation]{Idéaux engendrés par des familles d'idempotents et localisation}
\label{examples-section-ideal-locally-idempotents}

\noindent
Soit $R$ un anneau. Considérons l'anneau
$$
B(R) = R[x_n; n \in \mathbf{Z}]/(x_n(x_n - 1), x_nx_m; n \not = m)
$$
On montre facilement que tout idéal premier $\mathfrak q \subset B(R)$
est soit de la forme
$$
\mathfrak q = \mathfrak pB(R) + (x_n; n \in \mathbf{Z})
$$
soit de la forme
$$
\mathfrak q =
\mathfrak pB(R) + (x_n - 1) + (x_m; n \not = m, m \in \mathbf{Z}).
$$
Nous voyons donc que
$$
\Spec(B(R)) =
\Spec(R) \amalg \coprod\nolimits_{n \in \mathbf{Z}} \Spec(R)
$$
où la topologie n'est pas simplement celle de la réunion disjointe. Elle possède les
propriétés suivantes : chaque copie indexée par $n \in \mathbf{Z}$
est un sous-schéma ouvert, à savoir l'ouvert standard $D(x_n)$.
La copie « centrale » de $\Spec(R)$ appartient à l'adhérence de la réunion
de toute famille infinie des autres copies de $\Spec(R)$.
Remarquons que cette dernière copie de $\Spec(R)$ est définie par l'idéal
$(x_n, n \in \mathbf{Z})$, qui est engendré par les idempotents $x_n$.
Nous voyons ainsi que, si $\Spec(R)$ est connexe,
la décomposition précédente est exactement la décomposition de
$\Spec(B(R))$ en composantes connexes.

\medskip\noindent
Prenons ensuite $A = \mathbf{C}[x, y]/((y - x^2 + 1)(y + x^2 - 1))$.
Le spectre de $A$ possède deux composantes irréductibles
$C_1 = \Spec(A_1)$, $C_2 = \Spec(A_2)$,
où $A_1 = \mathbf{C}[x, y]/(y - x^2 + 1)$ et
$A_2 = \mathbf{C}[x, y]/(y + x^2 - 1)$. Remarquons qu'elles sont
paramétrées par $(x, y) = (t, t^2 - 1)$ et $(x, y) = (t, -t^2 + 1)$,
et se rencontrent en $P = (-1, 0)$ et $Q = (1, 0)$. Construisons une version
tordue de $B(A)$ en recollant $B(A_1)$ à $B(A_2)$ de la manière suivante :
au-dessus de $P$, faisons correspondre $x_n \in B(A_1) \otimes \kappa(P)$
à $x_n \in B(A_2) \otimes \kappa(P)$, tandis qu'au-dessus de $Q$,
nous faisons correspondre $x_n \in B(A_1) \otimes \kappa(Q)$
à $x_{n + 1} \in B(A_2) \otimes \kappa(Q)$.
Notons $B^{twist}(A)$ la $A$-algèbre ainsi obtenue.
Détails omis. Par construction,
$B^{twist}(A)$ est localement pour Zariski sur $A$ isomorphe à la version
non tordue. C'est le cas au-dessus des deux ouverts principaux
$\Spec(A) \setminus \{P\}$
et $\Spec(A) \setminus \{Q\}$.
Cependant, le recollement choisi produit assez de « monodromie » pour que
$\Spec(B^{twist}(A))$ soit connexe (détails omis).
Enfin, il existe une copie centrale
$\Spec(A) \to \Spec(B^{twist}(A))$
qui donne un sous-schéma fermé qui est, localement pour Zariski
sur $B^{twist}(A)$, défini par des idéaux engendrés par des idempotents, mais
ne l'est pas globalement (car $B^{twist}(A)$ ne possède aucun idempotent non trivial).

\begin{lemma}
\label{examples-lemma-not-generated-idempotents}
Il existe un schéma affine $X = \Spec(A)$ et un
sous-schéma fermé $T \subset X$ tel que $T$ soit, localement pour Zariski
sur $X$, défini par des idéaux engendrés par des idempotents, mais que
$T$ ne soit pas défini par un idéal engendré par des idempotents.
\end{lemma}

\begin{proof}
Voir ci-dessus.
\end{proof}



\section{Un homomorphisme d'anneaux qui induit des isomorphismes sur les anneaux locaux sans être ind-\'etale}
\label{examples-section-not-ind-etale}

\noindent
Remarquons que l'homomorphisme d'anneaux $R \to B(R)$ construit dans la
section \ref{examples-section-ideal-locally-idempotents} est
une limite inductive de produits finis de copies de $R$. Ainsi, $R \to B(R)$
est ind-Zariski; voir
Cohomologie pro-\'etale, définition \ref{proetale-definition-ind-zariski}.
Considérons ensuite l'homomorphisme d'anneaux $A \to B^{twist}(A)$
construit dans la section \ref{examples-section-ideal-locally-idempotents}.
Comme cet homomorphisme est, localement pour Zariski sur $\Spec(A)$, isomorphe à un
homomorphisme ind-Zariski $R \to B(R)$, nous concluons qu'il induit des isomorphismes sur les anneaux locaux
(voir Cohomologie pro-\'etale, lemme \ref{proetale-lemma-ind-zariski-implies}).
La discussion de la section \ref{examples-section-ideal-locally-idempotents}
montre qu'il existe une section
$B^{twist}(A) \to A$ dont le noyau n'est pas engendré par des idempotents.
Or, si $A \to B^{twist}(A)$ était ind-\'etale, c'est-à-dire si
$B^{twist}(A) = \colim A_i$ avec $A \to A_i$ \'etale,
alors le noyau de $A_i \to A$ serait engendré par un idempotent
(Algèbre, lemmes \ref{algebra-lemma-map-between-etale} et
\ref{algebra-lemma-surjective-flat-finitely-presented}).
Cela contredirait le résultat précédent.

\begin{lemma}
\label{examples-lemma-not-ind-etale}
Il existe un homomorphisme d'anneaux $A \to B$ qui induit des isomorphismes sur les anneaux locaux, mais
qui n'est pas ind-\'etale. A fortiori, il n'est pas ind-Zariski.
\end{lemma}

\begin{proof}
Voir la discussion ci-dessus.
\end{proof}




\section[Faisceau quasi-cohérent non flasque d'un module injectif]{Faisceau quasi-cohérent non flasque associé à un module injectif}
\label{examples-section-nonflasque}

\noindent
Pour d'autres exemples de ce type, voir \cite[Expos\'e II, Appendix I]{SGA6},
où Illusie expose certains exemples dus à Verdier.

\medskip\noindent
Considérons le schéma affine $X = \Spec(A)$, où
$$
A = k[x, y, z_1, z_2, \ldots]/(x^nz_n)
$$
est l'anneau de
Propriétés, exemple \ref{properties-example-does-not-work-in-general}.
Posons $I = (x) \subset A$. Considérons l'ouvert quasi-compact $U = D(x)$ de $X$.
Nous avons vu loc.\ cit.\ qu'il existe une section
$s \in \mathcal{O}_X(U)$ qui ne provient d'aucun homomorphisme de $A$-modules
$I^n \to A$, quel que soit $n \geq 0$.

\medskip\noindent
Soit $\alpha : A \to J$ un plongement de $A$ dans un $A$-module injectif.
Posons $Q = J/\alpha(A)$ et notons $\beta : J \to Q$ l'application quotient.
Nous affirmons que l'application
$$
\Gamma(X, \widetilde{J})
\longrightarrow
\Gamma(U, \widetilde{J})
$$
n'est pas surjective. Plus précisément, nous affirmons que $\alpha(s)$
n'appartient pas à son image. Pour le voir, raisonnons par l'absurde.
Supposons donc qu'il existe un élément $x \in J$ qui se restreigne à
$\alpha(s)$ sur $U$. Alors $\beta(x) \in Q$ se restreint à $0$ sur $U$. Nous savons donc que
$I^n\beta(x) = 0$ pour un certain $n$; voir
Propriétés,
lemme \ref{properties-lemma-sections-over-quasi-compact-open-in-affine}.
Nous obtenons par conséquent un morphisme
$\varphi : I^n \to A$, $h \mapsto \alpha^{-1}(hx)$. Il est immédiat que
ce morphisme $\varphi$ donne la section $s$ par l'application de
Propriétés,
lemme \ref{properties-lemma-sections-over-quasi-compact-open-in-affine},
ce qui est absurde.

\begin{lemma}
\label{examples-lemma-nonflasque}
Il existe un schéma affine $X = \Spec(A)$ et un $A$-module injectif $J$
tels que $\widetilde{J}$ ne soit pas un faisceau flasque sur $X$.
Même la restriction $\Gamma(X, \widetilde{J}) \to \Gamma(U, \widetilde{J})$,
où $U$ est un ouvert standard, n'est pas nécessairement surjective.
\end{lemma}

\begin{proof}
Voir ci-dessus.
\end{proof}

\noindent
En fait, une construction analogue fournit un exemple de module injectif
dont la cohomologie du faisceau quasi-cohérent associé sur un ouvert
quasi-compact n'est pas nulle. Plus précisément, partons de l'anneau
$$
A = k[x, y, w_1, u_1, w_2, u_2, \ldots]/(x^nw_n, y^nu_n, u_n^2, w_n^2)
$$
où $k$ est un corps. Choisissons un homomorphisme injectif $A \to I$, où
$I$ est un $A$-module injectif. Nous affirmons que l'élément $1/xy$ de
$A_{xy} \subset I_{xy}$ n'appartient pas à l'image de
$I_x \oplus I_y \to I_{xy}$. En raisonnant par l'absurde, supposons que
$$
\frac{1}{xy} = \frac{i}{x^n} + \frac{j}{y^n}
$$
pour un certain $n \geq 1$ et des $i, j \in I$. En chassant les dénominateurs, nous obtenons
$$
(xy)^{n + m - 1} = x^my^{n + m}i + x^{n + m}y^mj
$$
pour un certain $m \geq 0$. En multipliant par $u_{n + m}w_{n + m}$,
nous voyons que $u_{n + m}w_{n + m}(xy)^{n + m - 1} = 0$ dans $A$, ce
qui donne la contradiction cherchée.
Posons $U = D(x) \cup D(y) \subset X = \Spec(A)$. Pour tout $A$-module
$M$, nous avons une suite exacte
$$
0 \to H^0(U, \widetilde{M}) \to M_x \oplus M_y \to M_{xy} \to
H^1(U, \widetilde{M}) \to 0
$$
par Mayer--Vietoris. Nous en concluons que $H^1(U, \widetilde{I})$ n'est pas nul.

\begin{lemma}
\label{examples-lemma-nonvanishing}
Il existe un schéma affine $X = \Spec(A)$ dont l'espace
topologique sous-jacent est noethérien et un
$A$-module injectif $I$ tels que $\widetilde{I}$ ait un $H^1$ non nul
sur un ouvert quasi-compact $U$ de $X$.
\end{lemma}

\begin{proof}
Voir ci-dessus. Notons que $\Spec(A) = \Spec(k[x, y])$ comme espaces topologiques.
\end{proof}









\section{Un schéma en groupes plat non séparé}
\label{examples-section-non-separated-group-scheme}

\noindent
Tout schéma en groupes sur un corps est séparé; voir
Groupoïdes, lemme \ref{groupoids-lemma-group-scheme-over-field-separated}.
Ce n'est pas vrai pour les schémas en groupes sur une base.

\medskip\noindent
Soit $k$ un corps. Posons $S = \Spec(k[x]) = \mathbf{A}^1_k$.
Soit $G$ la droite affine dont l'origine $0$ est doublée (voir
Schémas, exemple \ref{schemes-example-affine-space-zero-doubled}),
considérée comme un schéma sur $S$. Ainsi, une fibre de $G \to S$ est soit
réduite à un point, soit un ensemble à deux éléments (l'un dans $U$, l'autre dans $V$).
Nous pouvons donc munir ces fibres d'une structure de groupe (en
prenant l'élément de $U$ pour zéro de la structure de groupe).
Plus précisément, $G$ possède deux ouverts $U, V$ qui se projettent isomorphiquement
sur $S$ et tels que $U \cap V$ se projette isomorphiquement sur
$S \setminus \{0\}$. Alors
$$
G \times_S G = U \times_S U \cup V \times_S U \cup
U \times_S V \cup V \times_S V
$$
où chaque morceau est isomorphe à $S$. Nous pouvons donc définir une multiplication
$m : G \times_S G \to G$ comme l'unique $S$-morphisme qui envoie le premier
et le dernier morceau dans $U$, et les deux morceaux intermédiaires dans $V$. Cela concorde avec
la description fibre par fibre donnée ci-dessus. Nous omettons la
vérification que cela définit une structure de schéma en groupes.

\begin{lemma}
\label{examples-lemma-non-separated-group-scheme}
Il existe sur la droite affine un schéma en groupes plat et de type fini
qui n'est pas séparé.
\end{lemma}

\begin{proof}
Voir la discussion ci-dessus.
\end{proof}

\begin{lemma}
\label{examples-lemma-non-quasi-separated-group-scheme}
Il existe sur l'espace affine de dimension infinie un schéma en groupes
plat et de type fini qui n'est pas quasi-séparé.
\end{lemma}

\begin{proof}
La même construction peut être effectuée avec l'espace affine de dimension infinie
$S = \mathbf{A}^\infty_k = \Spec k[x_1, x_2, \ldots]$ comme base
et avec l'origine $0 \in S$, qui correspond à l'idéal maximal
$(x_1, x_2, \ldots)$, comme point fermé doublé dans $G$.
Le schéma en groupes $G \rightarrow S$ ainsi obtenu
n'est pas quasi-séparé, comme l'explique
Schémas, exemple \ref{schemes-example-not-quasi-separated}.
\end{proof}



\section{Un schéma en groupes non plat à composante neutre plate}
\label{examples-section-non-flat-group-scheme}

\noindent
Soit $X \to S$ un monomorphisme de schémas. Posons $G = S \amalg X$.
Soit $m : G \times_S G \to G$ le $S$-morphisme
$$
G \times_S G = X \times_S X \amalg X \amalg X \amalg S
\longrightarrow G = X \amalg S
$$
qui envoie $X \times_S X$ et $S$ dans $S$, ainsi que chacun des
deux exemplaires de $X$ dans $X$ par le morphisme identité.
Cela définit une loi de groupe. Pour le voir, il faut montrer que
$m \circ (m \times \text{id}_G) = m \circ (\text{id}_G \times m)$
comme applications $G \times_S G \times_S G \to G$. En décomposant
$G \times_S G \times_S G$ en morceaux comme ci-dessus, nous constatons que
l'égalité doit être vérifiée sur chacun des $8$ morceaux.
Chaque morceau est isomorphe à $S$, à $X$, à $X \times_S X$ ou à
$X \times_S X \times_S X$. De plus, les deux applications envoient
respectivement ces morceaux dans $S$, $X$, $S$, $X$. Cela étant, le fait que
$X \to S$ soit un monomorphisme implique que $X \times_S X \cong X$
et $X \times_S X \times_S X \cong X$, et qu'il existe dans chaque cas
un unique $S$-morphisme $S \to S$ ou $X \to X$. Nous voyons donc
que $m \circ (m \times \text{id}_G) = m \circ (\text{id}_G \times m)$.
Ainsi, en prenant pour $X \to S$ n'importe quel
monomorphisme non plat de schémas (par exemple une immersion fermée),
nous obtenons un exemple de schéma en groupes
sur une base $S$ dont la composante neutre est $S$ (donc plate),
mais qui n'est pas plat.

\begin{lemma}
\label{examples-lemma-non-flat-group-scheme}
Il existe un schéma en groupes $G$ sur une base $S$ dont la composante
neutre est plate sur $S$, mais qui n'est pas plat sur $S$.
\end{lemma}

\begin{proof}
Voir la discussion ci-dessus.
\end{proof}




\section{Un espace algébrique en groupes non séparé sur un corps}
\label{examples-section-non-separated-group-space}

\noindent
Tout schéma en groupes sur un corps est séparé; voir
Groupoïdes, lemme \ref{groupoids-lemma-group-scheme-over-field-separated}.
Ce n'est pas vrai pour les espaces algébriques en groupes sur un corps
(mais voir la fin de cette section pour des résultats positifs).

\medskip\noindent
Soit $k$ un corps de caractéristique zéro.
Considérons l'espace algébrique $G = \mathbf{A}^1_k/\mathbf{Z}$ de
Espaces, exemple \ref{spaces-example-affine-line-translation}.
Par construction, $G$ est le faisceau fppf associé au préfaisceau
$$
T \longmapsto \Gamma(T, \mathcal{O}_T) / \mathbf{Z}
$$
sur la catégorie des schémas sur $k$. La loi d'addition évidente du préfaisceau
induit une addition $m : G \times G \to G$ qui fait de $G$ un espace
algébrique en groupes sur $\Spec(k)$. Notons que $G$ n'est pas séparé
(ni même quasi-séparé ou localement séparé). En revanche,
$G \to \Spec(k)$ est de type fini !

\begin{lemma}
\label{examples-lemma-non-separated-group-space}
Il existe sur un corps un espace algébrique en groupes de type fini
qui n'est pas séparé (ni même quasi-séparé ou localement séparé).
\end{lemma}

\begin{proof}
Voir la discussion ci-dessus.
\end{proof}

\noindent
Résultats positifs : si l'espace algébrique en groupes $G$ est
quasi-séparé ou localement séparé, ou plus généralement s'il est
un espace algébrique décent, alors $G$ est en fait séparé; voir
Compléments sur les groupoïdes dans les espaces, lemme
\ref{spaces-more-groupoids-lemma-group-scheme-over-field-separated}.
De plus, un espace algébrique en groupes de type fini et séparé sur un
corps est en fait un schéma d'après Compléments sur les groupoïdes dans les espaces, lemme
\ref{spaces-more-groupoids-lemma-group-space-scheme-locally-finite-type-over-k}.
L'idée de la preuve est que le lieu schématique est ouvert et dense; voir
Propriétés des espaces, proposition
\ref{spaces-properties-proposition-locally-quasi-separated-open-dense-scheme}
ou
Espaces décents, théorème \ref{decent-spaces-theorem-decent-open-dense-scheme}.
En prenant les translatés de cet ouvert, nous voyons que
tout point de $G$ possède un voisinage ouvert qui est un schéma.




\section{Spécialisations entre points d'une même fibre d'un morphisme étale}
\label{examples-section-specializations-fibre-etale}

\noindent
Si $f : X \to Y$ est un morphisme étale ou, plus généralement, localement
quasi-fini de schémas, il n'existe aucune spécialisation entre des points
d'une même fibre; voir
Morphismes, lemme \ref{morphisms-lemma-locally-quasi-finite-fibres}.
En revanche, cela est faux en général pour les morphismes d'espaces algébriques.

\medskip\noindent
Pour construire un exemple, soit $k$ un corps.
Posons
$$
P = k[u, u^{-1}, y, \{x_n\}_{n \in \mathbf{Z}}].
$$
Considérons l'action
de $\mathbf{Z}$ sur $P$ par des homomorphismes de $k$-algèbres, engendrée par l'automorphisme
$\tau$ défini par $\tau(u) = u$, $\tau(y) = uy$ et
$\tau(x_n) = x_{n + 1}$. Pour $d \geq 1$, posons
$I_d = ((1 - u^d)y, x_n - x_{n + d}, n \in \mathbf{Z})$.
Alors $V(I_d) \subset \Spec(P)$ est le lieu des points fixes de $\tau^d$.
Soit $S \subset P$ la partie multiplicativement stable engendrée par $y$ et
tous les $1 - u^d$, $d \in \mathbf{N}$. Nous
voyons alors que $\mathbf{Z}$ agit librement sur $U = \Spec(S^{-1}P)$.
Soit $X = U/\mathbf{Z}$ l'espace algébrique quotient; voir
Espaces, définition \ref{spaces-definition-quotient}.

\medskip\noindent
Considérons les idéaux premiers $\mathfrak p_n = (x_n, x_{n + 1}, \ldots)$ de
$S^{-1}P$. Notons que $\tau(\mathfrak p_n) = \mathfrak p_{n + 1}$.
Chacun d'eux définit donc un point $\xi_n \in U$ dont l'image dans $X$ est
le même point $x$ de $X$. De plus, nous avons les spécialisations
$$
\ldots \leadsto \xi_n \leadsto \xi_{n - 1} \leadsto \ldots
$$
Nous en concluons que $U \to X$ est un exemple du type annoncé.

\begin{lemma}
\label{examples-lemma-specializations-fibre-etale}
Il existe un morphisme étale d'espaces algébriques $f : X \to Y$
et une spécialisation non triviale de points $x \leadsto x'$ dans $|X|$, telle que
$f(x) = f(x')$ dans $|Y|$.
\end{lemma}

\begin{proof}
Voir la discussion ci-dessus.
\end{proof}







\section{Un torseur qui n'est pas un torseur fppf}
\label{examples-section-torsor-not-fppf}

\noindent
Dans
Groupoïdes, remarque \ref{groupoids-remark-fun-with-torsors},
nous demandons si tout torseur sous $G$ est un torseur fppf sous $G$.
Nous montrons dans cette section que ce n'est pas toujours le cas.

\medskip\noindent
Soit $k$ un corps. Dans la suite, tous les schémas et champs sont sur $k$.
Soit $G \to \Spec(k)$ le schéma en groupes
$$
G = (\mu_{2, k})^\infty =
\mu_{2, k} \times_k \mu_{2, k} \times_k \mu_{2, k} \times_k \ldots =
\lim_n (\mu_{2, k})^n
$$
où $\mu_{2, k}$ est le schéma en groupes des racines carrées de l'unité sur
$\Spec(k)$; voir
Groupoïdes, exemple \ref{groupoids-example-roots-of-unity}.
Comme $G$ est une limite projective de schémas affines, c'est un schéma en groupes
affine. C'est en fait le spectre de l'anneau
$k[t_1, t_2, t_3, \ldots]/(t_i^2 - 1)$. L'application de multiplication
$m : G \times_k G \to G$ est donnée au niveau des algèbres par
$t_i \mapsto t_i \otimes t_i$.

\medskip\noindent
Nous affirmons que tout torseur sous $G$ sur $k$ est de la forme
$$
P = \Spec(k[x_1, x_2, x_3, \ldots]/(x_i^2 - a_i))
$$
pour certains $a_i \in k^*$; l'action de $G$ sur cet espace, $G \times_k P \to P$,
est donnée par $x_i \to t_i \otimes x_i$ au niveau des algèbres.
Nous omettons la démonstration.
En fait, pour l'exemple, il nous suffit que $P$ soit un torseur sous $G$,
ce qui est clair puisque, sur $k' = k(\sqrt{a_1}, \sqrt{a_2}, \ldots)$,
le schéma $P$ devient isomorphe à $G$ de manière $G$-équivariante.
Notons que $P$ est trivial si et seulement si $k' = k$ : en effet, si
$P$ possède un point rationnel sur $k$, tous les $a_i$ sont des carrés.

\medskip\noindent
Nous affirmons que $P$ est un torseur fppf si et seulement si l'extension de corps
$k' = k(\sqrt{a_1}, \sqrt{a_2}, \ldots)/k$ est finie.
Si $k'$ est une extension finie de $k$, alors
$\{\Spec(k') \to \Spec(k)\}$
est un recouvrement fppf qui trivialise $P$, et nous voyons que $P$ est bien
un torseur fppf. Réciproquement, supposons que $P$ soit un torseur fppf sous $G$.
Cela signifie qu'il existe un recouvrement fppf
$\{S_i \to \Spec(k)\}$ tel que chaque $P_{S_i}$ soit trivial.
Choisissons un $i$ tel que $S_i$ ne soit pas vide. Soit $s \in S_i$ un point
fermé. D'après
Variétés, lemme \ref{varieties-lemma-locally-finite-type-Jacobson},
l'extension de corps $\kappa(s)/k$ est finie et, par construction,
$P_{\kappa(s)}$ possède un point rationnel sur $\kappa(s)$. Nous voyons donc que
$k \subset k' \subset \kappa(s)$ et $k'$ est une extension finie de $k$.

\medskip\noindent
Pour obtenir un exemple explicite, prenons $k = \mathbf{Q}$ et $a_i = i$,
par exemple (ou bien $a_i$ égal au $i$-ième nombre premier, si l'on préfère).

\begin{lemma}
\label{examples-lemma-torsors-principal-spaces-not-equal}
Soient $S$ un schéma et $G$ un schéma en groupes sur $S$.
Le champ $G\textit{-Principal}$ qui classifie les espaces principaux homogènes sous $G$
(voir Exemples de champs, sous-section
\ref{examples-stacks-subsection-principal-homogeneous-spaces})
et le champ $G\textit{-Torsors}$ qui classifie les torseurs fppf sous $G$
(voir Exemples de champs, sous-section
\ref{examples-stacks-subsection-fppf-torsors})
ne sont pas équivalents en général.
\end{lemma}

\begin{proof}
La discussion ci-dessus montre que le foncteur
$G\textit{-Torsors} \to G\textit{-Principal}$ n'est pas essentiellement
surjectif en général.
\end{proof}









\section[Champ non algébrique à recouvrement plat quasi compact]{Champ non algébrique admettant un recouvrement plat quasi compact}
\label{examples-section-not-algebraic-stack}

\noindent
Dans cette section, nous décrivons brièvement un exemple dû à Brian Conrad.
On peut trouver cet exemple en ligne à
\href{https://mathoverflow.net/questions/15082/fpqc-covers-of-stacks/15269#15269}{cette adresse}.
Notre exemple est légèrement différent.

\medskip\noindent
Soit $k$ un corps algébriquement clos.
Dans la suite, tous les schémas et champs sont sur $k$.
Soit $G \to \Spec(k)$ un schéma en groupes affine.
Dans Exemples de champs,
lemme \ref{examples-stacks-lemma-classifying-stacks},
nous avons donné plusieurs manières équivalentes de voir
$\mathcal{X} = [\Spec(k)/G]$ comme un champ en groupoïdes sur
$(\Sch/\Spec(k))_{fppf}$. En particulier, $\mathcal{X}$ classifie les
torseurs fppf sous $G$. Plus précisément, un $1$-morphisme $T \to \mathcal{X}$
correspond à un torseur fppf sous $G_T$, noté $P$, sur $T$, et les $2$-flèches
correspondent aux isomorphismes de torseurs. Il s'ensuit que
le $1$-morphisme diagonal
$$
\Delta :
\mathcal{X}
\longrightarrow
\mathcal{X} \times_{\Spec(k)} \mathcal{X}
$$
est représentable et affine. En effet, pour toute paire
de torseurs fppf sous $G_T$, notés $P_1, P_2$, sur un schéma $T/k$, le
schéma $\mathit{Isom}(P_1, P_2)$ est affine sur $T$.
Le torseur trivial sous $G$ sur $\Spec(k)$ définit un $1$-morphisme
$$
f : \Spec(k) \longrightarrow \mathcal{X}.
$$
Nous affirmons qu'il s'agit d'un $1$-morphisme surjectif. En effet,
par définition, pour tout $1$-morphisme $T \to \mathcal{X}$, il existe
un recouvrement fppf $\{T_i \to T\}$ tel que $P_{T_i}$ soit isomorphe
au torseur trivial sous $G_{T_i}$. Les composés
$T_i \to T \to \mathcal{X}$ se factorisent donc par $f$. Il est alors clair que
la projection $T \times_\mathcal{X} \Spec(k) \to T$
est surjective (c'est ainsi que nous définissons la surjectivité de $f$; voir
Champs algébriques,
définition \ref{algebraic-definition-relative-representable-property}).
On montre de même que $f$ est quasi compact et plat (nous omettons les détails).
Notons également que
$$
\Spec(k) \times_\mathcal{X} \Spec(k) \cong G
$$
comme schémas sur $k$.

\medskip\noindent
Supposons qu'il existe un morphisme lisse surjectif
$p : U \to \mathcal{X}$, où $U$ est un schéma.
Considérons le produit fibré
$$
\xymatrix{
W \ar[d] \ar[r] & U \ar[d] \\
\Spec(k) \ar[r] & \mathcal{X}
}
$$
Alors $W$ est un schéma lisse non vide sur $k$ et possède donc
un point rationnel sur $k$. Ainsi, $f$ se factorise par $U$.
Nous obtenons donc
$$
G \cong
\Spec(k) \times_\mathcal{X} \Spec(k) \cong
(\Spec(k) \times_k \Spec(k))
\times_{(U \times_k U)}
(U \times_\mathcal{X} U)
$$
et, puisque les projections $U \times_\mathcal{X} U \to U$ sont
supposées lisses, nous en concluons que
$U \times_\mathcal{X} U \to U \times_k U$ est
localement de type fini; voir
Morphismes, lemme \ref{morphisms-lemma-permanence-finite-type}.
Il s'ensuit que, dans ce cas, $G$ est localement de type fini sur $k$.
Nous avons ainsi démontré le lemme suivant (qui admet une
généralisation considérable).

\begin{lemma}
\label{examples-lemma-BG-algebraic}
Soit $k$ un corps. Soit $G$ un schéma en groupes affine sur $k$.
Si le champ $[\Spec(k)/G]$ possède un recouvrement lisse par un
schéma, alors $G$ est de type fini sur $k$.
\end{lemma}

\begin{proof}
Voir la discussion ci-dessus.
\end{proof}

\noindent
Pour obtenir un exemple explicite comme dans le titre de cette section, prenons
$G = (\mu_{2, k})^\infty$, le schéma en groupes de la
section \ref{examples-section-torsor-not-fppf}, qui n'est pas localement de type fini
sur $k$. La discussion ci-dessus montre que
$\mathcal{X} = [\Spec(k)/G]$ possède les propriétés (1) et (2) de
Champs algébriques, définition \ref{algebraic-definition-algebraic-stack},
mais non la propriété (3). Ainsi, $\mathcal{X}$ n'est pas un champ algébrique.
En revanche, il existe bien un schéma $U$ et un morphisme surjectif,
plat et quasi compact $U \to \mathcal{X}$, à savoir le morphisme
$f : \Spec(k) \to \mathcal{X}$ étudié ci-dessus.






\section{Commutation aux limites sur les objets, mais non aux limites}
\label{examples-section-limit-preserving}

\noindent
Soit $S$ un schéma non vide. Soit $\mathcal{G}$ un faisceau abélien injectif
sur $(\Sch/S)_{fppf}$. Nous obtenons un champ en groupoïdes
$$
\mathcal{G}\textit{-Torsors} \longrightarrow (\Sch/S)_{fppf}
$$
sur $S$; voir Exemples de champs, lemme
\ref{examples-stacks-lemma-torsors-sheaf-stack-in-groupoids}.
Ce champ commute aux limites sur les objets au-dessus de $(\Sch/S)_{fppf}$ (voir
Critères de représentabilité, section \ref{criteria-section-limit-preserving}),
car tout torseur sous $\mathcal{G}$ est trivial. En revanche,
$\mathcal{G}\textit{-Torsors}$ ne commute en général pas aux limites (voir
Axiomes d'Artin, définition \ref{artin-definition-limit-preserving}),
car le faisceau $\mathcal{G}$ ne commute pas nécessairement aux limites. Par exemple,
prenons un faisceau injectif non nul quelconque $\mathcal{I}$ et posons
$\mathcal{G} = \prod_{n \in \mathbf{Z}} \mathcal{I}$ pour obtenir un
exemple.

\begin{lemma}
\label{examples-lemma-limit-preserving-on-objects-not-limit-preserving}
Soit $S$ un schéma non vide. Il existe un champ en groupoïdes
$p : \mathcal{X} \to (\Sch/S)_{fppf}$
tel que $p$ commute aux limites sur les objets, mais que $\mathcal{X}$ ne
commute pas aux limites.
\end{lemma}

\begin{proof}
Voir la discussion ci-dessus.
\end{proof}



\section{Un champ classifiant non algébrique}
\label{examples-section-non-algebraic}

\noindent
Soit $S = \Spec(\mathbf{F}_p)$ et notons $\mu_p$ le
schéma en groupes des racines $p$-ièmes de l'unité sur $S$. Dans
Groupoïdes dans les espaces, section \ref{spaces-groupoids-section-stacks},
nous avons introduit le champ quotient $[S/\mu_p]$ et, dans
Exemples de champs, section \ref{examples-stacks-section-group-quotient-stacks},
nous avons montré que $[S/\mu_p]$ est le champ classifiant les
torseurs fppf sous $\mu_p$ : pour tout schéma $T$ sur $S$, la catégorie
$\Mor_S(T, [S/\mu_p])$ est canoniquement équivalente à la
catégorie des torseurs fppf sous $\mu_p$ sur $T$. Enfin, dans
Critères de représentabilité, théorème
\ref{criteria-theorem-flat-groupoid-gives-algebraic-stack}
nous avons vu que $[S/\mu_p]$ est un champ algébrique.

\medskip\noindent
Nous pouvons maintenant poser la question : « Qu'en est-il de la catégorie fibrée en groupoïdes
$\mathcal{S}$ qui classifie les torseurs étales sous $\mu_p$ ? » (Autrement dit,
$\mathcal{S}$ est une catégorie au-dessus de $\Sch/S$ dont la catégorie fibre au-dessus d'un
schéma $T$ est la catégorie des torseurs étales sous $\mu_p$ sur $T$.)

\medskip\noindent
La première objection est qu'il ne s'agit pas d'un champ pour la topologie fppf,
car la descente des objets n'y est pas satisfaite. Par exemple, le
torseur sous $\mu_p$ donné par $\Spec(\mathbf{F}_p(t)[x]/(x^p - t))$ sur
$T = \Spec(\mathbf{F}_p(t))$ est localement trivial pour la topologie fppf, mais
non pour la topologie étale.

\medskip\noindent
Pour remédier à ce premier problème, travaillons avec la topologie étale :
la descente des objets y est satisfaite. En effet,
$\mathcal{S}$ est un champ en groupoïdes sur $(\Sch/S)_\etale$.
De plus, la diagonale
$\Delta : \mathcal{S} \to \mathcal{S} \times \mathcal{S}$
est représentable (par des schémas). En effet, étant donnés deux
torseurs sous $\mu_p$ (qu'ils soient ou non localement triviaux pour la topologie étale),
le faisceau de leurs isomorphismes est représentable par un schéma.

\medskip\noindent
Nous pouvons donc enfin demander s'il existe un schéma $U$ et un
$1$-morphisme lisse et surjectif $U \to \mathcal{S}$. Nous montrerons de deux
manières que cela est impossible : par un argument direct (que nous conseillons
au lecteur de passer) et par un argument utilisant un résultat général.

\medskip\noindent
Argument direct (esquisse) : notons que le $1$-morphisme
$\mathcal{S} \to \Spec(\mathbf{F}_p)$ satisfait au
critère de relèvement infinitésimal de lissité formelle.
En effet, pour tout épaississement infinitésimal du premier ordre
de schémas $T \to T'$, le noyau de $\mu_p(T') \to \mu_p(T)$
est isomorphe aux sections du faisceau d'idéaux de $T$ dans $T'$, et
ainsi $H^1_\etale(T, \mu_p) = H^1_\etale(T', \mu_p)$.
De plus, le champ $\mathcal{S}$ commute aux limites. Par conséquent,
si $U \to \mathcal{S}$ est lisse, alors $U \to \Spec(\mathbf{F}_p)$
commute aux limites et satisfait au
critère de relèvement infinitésimal de lissité formelle.
Il en résulte que $U$ est lisse sur $\mathbf{F}_p$.
En particulier, $U$ est réduit, donc $H^1_\etale(U, \mu_p) = 0$.
Ainsi, $U \to \mathcal{S}$ se factorise en
$U \to \Spec(\mathbf{F}_p) \to \mathcal{S}$
et la première flèche est lisse. Par descente de la lissité,
la lissité de $U \to \mathcal{S}$ impliquerait celle de
$\Spec(\mathbf{F}_p) \to \mathcal{S}$. Or ce n'est
pas le cas, puisque
$\Spec(\mathbf{F}_p) \times_\mathcal{S} \Spec(\mathbf{F}_p)$
est $\mu_p$, qui n'est pas lisse sur $\Spec(\mathbf{F}_p)$.


\medskip\noindent
Argument structurel : dans
Critères de représentabilité, section \ref{criteria-section-stacks-etale},
nous avons vu que l'on peut considérer les champs algébriques comme les
champs en groupoïdes pour la topologie étale dont la diagonale est
représentable par des espaces algébriques admettant un recouvrement lisse.
Ainsi, s'il existe un morphisme lisse surjectif $U \to \mathcal{S}$, alors
$\mathcal{S}$ est un champ algébrique et satisfait en particulier à la
descente pour la topologie fppf. Mais nous avons vu ci-dessus que $\mathcal{S}$
ne satisfait pas à la descente pour la topologie fppf.

\medskip\noindent
De manière informelle, les arguments ci-dessus montrent que le champ classifiant,
pour la topologie étale, les torseurs localement triviaux pour cette topologie
sous un schéma en groupes $G$ sur une base $B$ est algébrique si et seulement
si $G$ est lisse sur $B$. L'un des avantages de travailler avec
la topologie fppf est qu'il suffit de supposer que $G \to B$
est plat et localement de présentation finie. En fait, le champ quotient
(pour la topologie fppf) $[B/G]$ est algébrique si et seulement
si $G \to B$ est plat et localement de présentation finie; voir
Critères de représentabilité, lemme \ref{criteria-lemma-BG-algebraic}.







\section{Faisceau muni d'un recouvrement plat quasi-compact qui n'est pas algébrique}
\label{examples-section-not-algebraic}

\noindent
Considérons le foncteur $F = (\mathbf{P}^1)^\infty$, c'est-à-dire que, pour un schéma $T$,
$F(T)$ est l'ensemble des $f = (f_1, f_2, f_3, \ldots)$ où chaque
$f_i : T \to \mathbf{P}^1$ est un morphisme de schémas. Remarquons que
$\mathbf{P}^1$ vérifie la propriété de faisceau pour les recouvrements fpqc; voir
Descente, lemme \ref{descent-lemma-fpqc-universal-effective-epimorphisms}.
Un produit de faisceaux étant un faisceau, $F$ vérifie aussi la propriété de faisceau pour
la topologie fpqc. La diagonale de $F$ est représentable : si $f : T \to F$
et $g : S \to F$ sont des morphismes, alors $T \times_F S$ est l'intersection schématique
des sous-schémas fermés $T \times_{f_i, \mathbf{P}^1, g_i} S$
dans le schéma $T \times S$. Considérons le schéma en groupes $\text{SL}_2$, muni
d'un morphisme affine, lisse et surjectif $\text{SL}_2 \to \mathbf{P}^1$.
Considérons ensuite $U = (\text{SL}_2)^\infty$ et son morphisme canonique (produit)
$U \to F$. Remarquons que $U$ est un schéma affine. Nous affirmons que le morphisme
$U \to F$ est plat, surjectif et universellement ouvert. En effet, supposons donné
un morphisme $f : T \to F$. Alors $Z = T \times_F U$ est le produit fibré infini
des schémas $Z_i = T \times_{f_i, \mathbf{P}^1} \text{SL}_2$
sur $T$. Chacun des morphismes $Z_i \to T$ est lisse, surjectif et
affine, ce qui entraîne que
$$
Z = Z_1 \times_T Z_2 \times_T Z_3 \times_T \ldots
$$
est un schéma plat et affine sur $T$. Un argument élémentaire de limite montre aussi que
$Z \to T$ est ouvert.

\medskip\noindent
En revanche, nous affirmons que $F$ n'est pas un espace algébrique.
En effet, si $F$ était un espace algébrique, ce serait un espace algébrique
quasi-compact et séparé (d'après notre description des produits fibrés sur $F$).
La cohomologie des faisceaux quasi-cohérents s'annulerait donc au-dessus d'un
certain degré (voir
Cohomologie des espaces, proposition \ref{spaces-cohomology-proposition-vanishing},
et les remarques qui la précèdent). Mais il est clair qu'en prenant l'image réciproque
de $\mathcal{O}(-2, -2, \ldots, -2)$ par la projection
$$
(\mathbf{P}^1)^\infty \longrightarrow (\mathbf{P}^1)^n
$$
(qui possède une section), on peut obtenir un faisceau quasi-cohérent dont la cohomologie
est non nulle en degré $n$. Nous obtenons ainsi une réponse à une question posée
par Anton Geraschenko sur MathOverflow.

\begin{lemma}
\label{examples-lemma-not-algebraic}
Il existe un foncteur $F : \Sch^{opp} \to \textit{Ensembles}$
qui satisfait à la condition de faisceau pour la topologie fpqc, dont la diagonale
$\Delta : F \to F \times F$ est représentable, et tel qu'il existe un morphisme
surjectif, plat, universellement ouvert et quasi-compact
$U \to F$, où $U$ est un schéma, mais tel que $F$ ne soit pas un espace algébrique.
\end{lemma}

\begin{proof}
Voir la discussion qui précède.
\end{proof}







\section[Topologie \'etale, Zariski et recouvrements finis \'etales]{La topologie \'etale face à celle de Zariski et aux recouvrements finis et \'etales}
\label{examples-section-etale-zariski-finite-etale}

\noindent
Nous donnons dans cette section un exemple trouvé par B.~Bhatt en réponse
à une question de C.~Deninger, qui montre que la topologie \'etale
des variétés algébriques sur $\mathbf{C}$ est strictement plus fine que la
topologie engendrée par les immersions ouvertes et les morphismes finis et \'etales.
De tels exemples sont certainement connus des spécialistes et remontent peut-être
à Artin et Grothendieck\footnote{Selon une anecdote, Grothendieck aurait d'abord
tenté de définir la topologie \'etale au moyen des recouvrements ouverts de Zariski et des
recouvrements finis et \'etales, avant qu'Artin ne le convainque du contraire
(probablement à partir d'exemples comme celui consigné ici).}.
Si vous connaissez un exemple publié, veuillez écrire à
\href{mailto:stacks.project@gmail.com}{stacks.project@gmail.com}.

\medskip\noindent
Soit $X$ une courbe affine lisse et connexe sur $\mathbf{C}$.
Choisissons un morphisme fini $g: Y \to X$ de degré $3$, où
$Y$ est une courbe lisse connexe, et un point $x \in X$
dont la fibre $g^{-1}(x) = \{u, y\}$ est formée d'un point non ramifié
$u$ (de degré 1) et d'un point ramifié $y$ (de degré 2). En remplaçant
$X$ par un voisinage ouvert de Zariski de $x$, nous pouvons supposer que $g$
est non ramifié partout sauf en $y$. Soit $U = Y \setminus \{y\}$.
Alors $f=g|_U: U \to X$ est un morphisme \'etale surjectif et
$f^{-1}(x) = \{u\}$. 

\begin{lemma}
\label{examples-lemma-etale-not-dominated-zariski-finite-etale}
Le recouvrement \'etale $U \to X$ n'admet aucun raffinement par une composition finie
d'immersions ouvertes et de morphismes finis et \'etales.
\end{lemma}

\begin{proof}
Raisonnons par l'absurde et supposons donnée une tour de morphismes
$$
W_m \to W_{m-1} \to \dots \to W_1 \to W_0 = X
$$
où chaque morphisme $W_i \to W_{i - 1}$ est soit une immersion ouverte, soit un
morphisme fini et \'etale, ainsi qu'un relèvement $W_m \to U$ de $W_m \to X$ le long de
$f : U \to X$ et un point $w \in W_m$ d'image $u \in U$
(et donc $x \in X$). En remplaçant chaque $W_k$ par la composante connexe
qui contient l'image $w_k \in W_k$ de $w$, nous pouvons supposer
que chaque $W_k$ est lisse et connexe.

\medskip\noindent
Pour chaque $0 \le i \le m$, considérons le produit fibré $Z_i = Y \times_X W_i$.
Comme $W_i \to X$ est \'etale et $Y$ est lisse, la courbe $Z_i$ est une réunion
disjointe de courbes lisses connexes, et les morphismes $Z_i \to W_i$ sont finis
de degré $3$ par changement de base. La fibre de $Z_i \to W_i$
au-dessus de $w_i$ est manifestement formée de deux points $z_i, v_i$ qui s'envoient
respectivement sur $y, u$ dans $Y$. Le morphisme $Z_i \to W_i$ est non ramifié en $v_i$
et ramifié de degré $2$ en $z_i$. Posons $n_i = \#\pi_0(Z_i)$. Comme la fibre de
$Z_i \to W_i$ possède $2$ points, nous avons $n_i \in \{1, 2\}$. De plus, $n_i$ ne peut
que croître avec $i$, car les morphismes de transition $Z_i \to Z_{i - 1}$ sont dominants.
Nous avons clairement $n_0 = 1$. D'autre part, $n_m \geq 2$ : il existe un
$X$-morphisme $W_m \to U \subset Y$, donc le changement de base $Z_m \to W_m$
possède une section. Il existe alors un plus petit indice $i \ge 1$ tel que $n_{i - 1} = 1$
et $n_i = 2$. Comme les $n_j$ sont déterminés par l'irréductibilité de $Z_j$,
le morphisme $W_i \to W_{i - 1}$ ne peut être une immersion ouverte;
il doit donc être fini et \'etale. Puisque $n_i = 2$, la courbe lisse $Z_i$
possède deux composantes connexes; puisque le morphisme vers $W_i$ est fini, l'une de
ces composantes s'envoie nécessairement isomorphiquement sur $W_i$.
Nous pouvons donc écrire $Z_i = W_i \sqcup Z'_i$, où $Z'_i \to W_i$ est de degré $2$.
Le morphisme induit $W_i \subset Z_i \to Z_{i - 1}$ est dominant
entre deux courbes connexes, toutes deux finies sur $W_{i - 1}$.
Un tel morphisme est nécessairement surjectif. Nous pouvons donc choisir un point $t \in W_i$
d'image $z_{i - 1}$. Nous obtenons des extensions
$$
\mathcal{O}_{W_{i - 1}, w_{i - 1}} \to
\mathcal{O}_{Z_{i - 1}, z_{i - 1}} \to
\mathcal{O}_{W_i, t}
$$
d'anneaux de valuation discrète. La première est ramifiée, tandis que
l'extension composée est non ramifiée puisque $W_i \to W_{i - 1}$ est \'etale; c'est
impossible, d'où la contradiction.
\end{proof}










\section{Faisceaux et spécialisations}
\label{examples-section-sheaves}

\noindent
Dans la suite, nous fixons un grand site \'etale $\Sch_\etale$
construit dans
Topologies, définition \ref{topologies-definition-big-etale-site}.
De plus, par « schéma », nous entendrons un objet de ce site.
Rappelons que, si $x, x'$ sont des points d'un schéma $X$, on dit que $x$ est une
{\it spécialisation} de $x'$, ou l'on écrit $x' \leadsto x$, si
$x \in \overline{\{x'\}}$. C'est notamment le cas lorsque $x = x'$.

\medskip\noindent
Considérons le foncteur $F : \Sch_\etale \to \textit{Ab}$
défini par les règles suivantes :
$$
F(X) = \prod\nolimits_{x \in X}
\prod\nolimits_{x' \in X, x' \leadsto x, x' \not = x}
\mathbf{Z}/2\mathbf{Z}
$$
Pour tout schéma $X$, nous notons $|X|$ son ensemble sous-jacent.
Un élément $a \in F(X)$ sera considéré comme une application ensembliste
$|X| \times |X| \to \mathbf{Z}/2\mathbf{Z}$, $(x, x') \mapsto a(x, x')$,
nulle si $x = x'$ ou si $x$ n'est pas une spécialisation de $x'$.
Pour un morphisme de schémas $f : X \to Y$, nous définissons
$$
F(f) : F(Y) \longrightarrow F(X)
$$
par la règle suivante : pour $b \in F(Y)$, nous posons
$$
F(f)(b)(x, x') =
\left\{
\begin{matrix}
0 & \text{si }x\text{ n'est pas une spécialisation de }x' \\
b(f(x), f(x')) & \text{sinon.}
\end{matrix}
\right.
$$
Remarquons que cela définit bien un élément de $F(X)$. Nous affirmons que, si
$f : X \to Y$ et $g : Y \to Z$ sont des morphismes composables, alors
$F(f) \circ F(g) = F(g \circ f)$. En effet, soient $c \in F(Z)$ et
$x' \leadsto x$ une spécialisation entre deux points de $X$; alors
$$
F(g \circ f)(c)(x, x') = c(g(f(x)), g(f(x'))) = F(f)(F(g)(c))(x, x')
$$
car $f(x') \leadsto f(x)$. (Cela vaut aussi si $f(x) = f(x')$.)

\medskip\noindent
Soit $G$ le faisceau associé à $F$ pour la topologie \'etale.

\medskip\noindent
Nous affirmons que, si $X$ est un schéma et $x' \leadsto x$ une spécialisation
avec $x' \not = x$, alors $G(X) \not = 0$. En effet, soit $a \in F(X)$
un élément tel que $a$, considéré comme la fonction
$|X| \times |X| \to \mathbf{Z}/2\mathbf{Z}$, est non nul en $(x, x')$.
Soit $\{f_i : U_i \to X\}$ un recouvrement \'etale de $X$. Nous pouvons choisir un
indice $i$ et un point $u_i \in U_i$ tels que $f_i(u_i) = x$. Comme les générisations
se relèvent le long des morphismes plats (voir
Morphismes, lemme \ref{morphisms-lemma-generalizations-lift-flat}),
nous pouvons trouver une spécialisation $u'_i \leadsto u_i$
telle que $f_i(u'_i) = x'$. Notre construction montre alors que
$F(f_i)(a) \not = 0$. Ainsi, $a$ détermine un élément non nul de $G(X)$.

\medskip\noindent
Remarquons que, si $X = \Spec(k)$, où $k$ est un corps (ou, plus généralement,
un anneau dont tous les idéaux premiers sont maximaux), alors $F(X) = 0$;
de plus, pour tout morphisme \'etale $U \to X$, nous avons $F(U) = 0$, car il
n'existe aucune spécialisation entre des points distincts dans les fibres d'un
morphisme \'etale. Par conséquent, $G(X) = 0$.

\medskip\noindent
Supposons que $X \subset X'$ soit un épaississement; voir
Compléments sur les morphismes, définition \ref{more-morphisms-definition-thickening}.
Alors la catégorie des schémas \'etales sur $X'$ est équivalente à la
catégorie des schémas \'etales sur $X$ par le foncteur de changement de base
$U' \mapsto U = U' \times_{X'} X$; voir
Cohomologie \'etale,
théorème \ref{etale-cohomology-theorem-topological-invariance}.
Puisque, dans cette situation, on a toujours $F(U) = F(U')$,
nous avons aussi $G(X) = G(X')$.

\medskip\noindent
Comme variante, nous pouvons considérer le préfaisceau $F_n$ qui associe
à un schéma $X$ l'ensemble des applications
$a : |X|^{n + 1} \to \mathbf{Z}/2\mathbf{Z}$ telles que $a(x_0, \ldots, x_n)$
ne soit non nul que si $x_n \leadsto \ldots \leadsto x_0$ est une suite de
spécialisations et $x_n \not = x_{n - 1} \not = \ldots \not = x_0$.
Soit $G_n$ le faisceau associé à $F_n$.
Exactement comme ci-dessus, on montre que $G_n$ est non nul
si $\dim(X) \geq n$ et nul si $\dim(X) < n$.

\begin{lemma}
\label{examples-lemma-sheaf-zero-on-low-dimension}
Il existe un faisceau de groupes abéliens $G$ sur
$\Sch_\etale$ qui possède les propriétés suivantes :
\begin{enumerate}
\item $G(X) = 0$ dès que $\dim(X) < n$,
\item $G(X)$ n'est pas nul si $\dim(X) \geq n$, et
\item si $X \subset X'$ est un épaississement, alors $G(X) = G(X')$.
\end{enumerate}
\end{lemma}

\begin{proof}
Voir la discussion qui précède.
\end{proof}

\begin{remark}
\label{examples-remark-specialization}
Voici quelques remarques :
\begin{enumerate}
\item Les préfaisceaux $F$ et $F_n$ sont séparés.
\item Il se trouve que $F$ et $F_n$ ne sont pas des faisceaux.
\item On peut montrer que $G$ et $G_n$ sont en fait des faisceaux pour la topologie fppf.
\end{enumerate}
Nous démontrerons ces résultats si nous en avons besoin.
\end{remark}


\section{Faisceaux et fonctions constructibles}
\label{examples-section-constructible-functions}

\noindent
Dans la suite, nous fixons un grand site \'etale $\Sch_\etale$ construit dans
Topologies, définition \ref{topologies-definition-big-etale-site}.
De plus, par « schéma », nous entendrons un objet de ce site.
Dans cette section, nous appelons {\it partition constructible}
d'un schéma $X$ toute décomposition localement finie en réunion disjointe
$X = \coprod_{i \in I} X_i$ telle que chaque $X_i \subset X$ soit une
partie localement constructible de $X$. Localement finie signifie que,
pour tout ouvert quasi-compact $U \subset X$, il n'existe qu'un nombre fini
d'indices $i \in I$ tels que $X_i \cap U$ soit non vide.
Remarquons que, si $f : X \to Y$ est un morphisme de schémas et
$Y = \coprod Y_j$ une partition constructible, alors
$X = \coprod f^{-1}(Y_j)$ est une partition constructible de $X$.
Étant donnés un ensemble $S$ et un schéma $X$, une {\it fonction constructible}
$f : |X| \to S$ est une application telle que $X = \coprod_{s \in S} f^{-1}(s)$
soit une partition constructible de $X$.
Si $G$ est un groupe abstrait et $a, b : |X| \to G$ des fonctions
constructibles, alors $ab : |X| \to G, x \mapsto a(x)b(x)$ est encore une fonction
constructible. En effet, deux partitions constructibles quelconques
admettent une troisième partition qui les raffine toutes deux.

\medskip\noindent
Soit $A$ un groupe abélien quelconque. Pour tout schéma $X$, nous définissons
$$
F(X) =
\frac{\{a : |X| \to A \mid a \text{ est une fonction constructible}\}}{\{\text{fonctions localement constantes }|X| \to A\}}
$$
Nous considérons simplement un élément $a$ de $F(X)$ comme une fonction définie
à une fonction localement constante près. Étant donnés un morphisme de schémas
$f : X \to Y$ et un élément $b \in F(Y)$, nous posons
$F(f)(b) = b \circ f$. Ainsi, $F$ est un préfaisceau sur $\Sch_\etale$.

\medskip\noindent
Remarquons que, si $\{f_i : U_i \to X\}$ est un recouvrement fppf et si $a \in F(X)$
est tel que $F(f_i)(a) = 0$ dans $F(U_i)$, alors $a \circ f_i$ est une fonction
localement constante pour tout $i$. Il en résulte que $a$ est elle-même localement
constante, puisque les morphismes $f_i$ sont ouverts. Donc $a = 0$ dans $F(X)$.
Ainsi, $F$ est un préfaisceau séparé (pour la topologie fppf,
donc a fortiori pour la topologie \'etale).

\medskip\noindent
Soit $G$ le faisceau associé à $F$ pour la topologie \'etale.
Comme $F$ est séparé et que, par exemple, $F(X) \not = 0$ lorsque
$X$ est le spectre d'un anneau de valuation discrète, nous voyons que $G$
n'est pas nul.

\medskip\noindent
Soit $X = \Spec(k)$, où $k$ est un corps.
Alors tout recouvrement \'etale de $X$ admet un raffinement
de la forme $\{\Spec(k') \to \Spec(k)\}$, où $k'/k$
est une extension finie séparable. Puisque $F(\Spec(k')) = 0$,
nous avons $G(X) = 0$.

\medskip\noindent
Supposons que $X \subset X'$ soit un épaississement; voir
Compléments sur les morphismes, définition \ref{more-morphisms-definition-thickening}.
Alors la catégorie des schémas \'etales sur $X'$ est équivalente à la
catégorie des schémas \'etales sur $X$ par le foncteur de changement de base
$U' \mapsto U = U' \times_{X'} X$; voir
Cohomologie \'etale,
théorème \ref{etale-cohomology-theorem-topological-invariance}.
Comme $F(U) = F(U')$ dans cette situation, nous avons aussi $G(X) = G(X')$.

\medskip\noindent
Le faisceau $G$ commute aux limites; voir
Limites d'espaces,
définition \ref{spaces-limits-definition-locally-finite-presentation}.
En effet, soit $R$ un anneau écrit comme une limite inductive filtrante
$R = \colim_i R_i$ d'anneaux. Posons $X = \Spec(R)$ et
$X_i = \Spec(R_i)$, de sorte que $X = \lim_i X_i$. Alors
$G(X) = \colim_i G(X_i)$. Pour le démontrer, on établit d'abord que
toute partition constructible de $\Spec(R)$ provient d'une partition constructible
d'un certain $\Spec(R_i)$. On obtient ainsi le résultat pour $F$. Pour
obtenir le résultat pour le faisceau associé, on utilise que tout homomorphisme \'etale d'anneaux
$R \to R'$ provient d'un homomorphisme \'etale $R_i \to R_i'$ pour un certain $i$.
Nous omettons les détails.

\begin{lemma}
\label{examples-lemma-weird-sheaf}
Il existe un faisceau de groupes abéliens $G$ sur
$\Sch_\etale$ qui possède les propriétés suivantes :
\begin{enumerate}
\item $G(\Spec(k)) = 0$ pour tout corps $k$,
\item $G$ commute aux limites,
\item si $X \subset X'$ est un épaississement, alors $G(X) = G(X')$, et
\item $G$ n'est pas nul.
\end{enumerate}
\end{lemma}

\begin{proof}
Voir la discussion qui précède.
\end{proof}




\section{Le site lisse-\'etale n'est pas fonctoriel}
\label{examples-section-lisse-etale-not-functorial}

\noindent
Le site {\it lisse-\'etale}
$X_{lisse,\etale}$ de $X$ est la catégorie des schémas lisses sur $X$
munie des recouvrements \'etales (usuels); voir
Cohomologie des champs, section \ref{stacks-cohomology-section-lisse-etale}.
Soit $f : X  \to Y$ un morphisme de schémas.
Il existe un foncteur
$$
u : Y_{lisse,\etale} \longrightarrow X_{lisse,\etale},\quad
V/Y \longmapsto V \times_Y X
$$
qui est continu. Nous obtenons donc une paire de foncteurs adjoints
$$
u^s :
\Sh(X_{lisse,\etale})
\longrightarrow
\Sh(Y_{lisse,\etale}),
\quad
u_s :
\Sh(Y_{lisse,\etale})
\longrightarrow
\Sh(X_{lisse,\etale}),
$$
voir Sites, section \ref{sites-section-continuous-functors}.
Nous affirmons qu'en général $u$ ne définit {\bf pas} un morphisme de sites; voir
Sites, définition \ref{sites-definition-morphism-sites}.
Autrement dit, nous affirmons que $u_s$ n'est pas exact à gauche en général. Notons que
les préfaisceaux représentables sont des faisceaux sur les sites lisses-\'etales. Par
Sites, lemme \ref{sites-lemma-pullback-representable-sheaf},
nous avons donc $u_sh_V = h_{V \times_Y X}$. Considérons deux morphismes
$$
\xymatrix{
V_1 \ar[rd] \ar@<1ex>[rr]^a \ar@<-1ex>[rr]_b & & V_2 \ar[ld] \\
& Y
}
$$
de schémas $V_1, V_2$ lisses sur $Y$. Si $u_s$ était exact à gauche, nous
aurions alors
$$
u_s \text{Égaliseur}(h_a, h_b : h_{V_1} \to h_{V_2})
=
\text{Égaliseur}(h_{a \times 1}, h_{b \times 1} :
h_{V_1 \times_Y X} \to h_{V_2 \times_Y X})
$$
Nous choisirons les morphismes $a, b : V_1 \to V_2$ de sorte qu'il n'existe
aucun morphisme d'un schéma lisse sur $Y$ vers $(a = b) \subset V_1$, c'est-à-dire
que le membre de gauche soit le faisceau vide, mais qu'après
changement de base à $X$ l'égalisateur soit non vide et lisse sur $X$.
Un exemple rudimentaire consiste à prendre $X = \Spec(\mathbf{F}_p)$,
$Y = \Spec(\mathbf{Z})$ et $V_1 = V_2 = \mathbf{A}^1_\mathbf{Z}$,
avec $a(x) = x$ et $b(x) = x + p$. Remarquons que l'égalisateur
de $a$ et $b$ est la fibre de $\mathbf{A}^1_\mathbf{Z}$ au-dessus de $(p)$.

\begin{lemma}
\label{examples-lemma-lisse-etale-not-functorial}
Le site lisse-\'etale n'est pas fonctoriel, même pour les morphismes de schémas.
\end{lemma}

\begin{proof}
Voir la discussion qui précède.
\end{proof}






\section{Faisceaux sur la catégorie des schémas noethériens}
\label{examples-section-sheaves-locally-Noetherian}

\noindent
Soit $S$ un schéma localement noethérien. Comme dans
Axiomes d'Artin, section \ref{artin-section-noetherian},
considérons le foncteur d'inclusion
$$
u : (\textit{Noetherian}/S)_{fppf} \longrightarrow (\Sch/S)_{fppf}
$$
du site fppf des schémas localement noethériens sur $S$ dans un
grand site fppf de $S$. Comme expliqué dans la section citée,
ce foncteur est continu. Nous obtenons donc une paire de foncteurs adjoints
$$
u^s :
\Sh((\Sch/S)_{fppf})
\longrightarrow
\Sh((\textit{Noetherian}/S)_{fppf})
$$
et
$$
u_s :
\Sh((\textit{Noetherian}/S)_{fppf})
\longrightarrow
\Sh((\Sch/S)_{fppf})
$$
voir Sites, section \ref{sites-section-continuous-functors}.
Nous affirmons toutefois qu'en général $u$ ne définit pas un
morphisme de sites; voir Sites, définition \ref{sites-definition-morphism-sites}.
Autrement dit, le foncteur $u_s$ n'est pas exact à gauche en général.

\medskip\noindent
Soit $p$ un nombre premier et posons $S = \Spec(\mathbf{F}_p)$.
Considérons le morphisme injectif de faisceaux
$$
a : \mathcal{F} \longrightarrow \mathcal{G}
$$
sur $(\textit{Noetherian}/S)_{fppf}$ défini comme suit : pour $U$
un schéma localement noethérien sur $S$, nous définissons
$$
\mathcal{G}(U) = \Gamma(U, \mathcal{O}_U)^* =
\Mor_S(U, \mathbf{G}_{m, S})
$$
et nous prenons
$$
\begin{aligned}
\mathcal{F}(U) = \{f \in \mathcal{G}(U) \mid {}&
\text{localement pour la topologie fppf, }f
\text{ admet des racines de degré une puissance de }p\\
&\text{arbitrairement grande}\}
\end{aligned}
$$
Une $\mathbf{F}_p$-algèbre noethérienne $A$ possède un nilradical nilpotent
$I \subset A$; les racines de degré une puissance de $p$ de $1$ dans $A$ sont les
éléments de la forme $1 + a$, $a \in I$, et aucune racine non triviale
de degré une puissance de $p$ de $1$ n'admet donc des racines de degré une puissance de $p$ arbitrairement grande. Nous en déduisons
que $\mathcal{F}(U)$ est un groupe abélien sans $p$-torsion pour tout schéma
localement noethérien $U$; nous omettons certains détails. Il s'ensuit que
$p : \mathcal{F} \to \mathcal{F}$
est un morphisme injectif de faisceaux abéliens sur $(\textit{Noetherian}/S)_{fppf}$.

\medskip\noindent
Pour obtenir une contradiction, supposons $u_s$ exact. Alors
$p : u_s\mathcal{F} \to u_s\mathcal{F}$ est lui aussi injectif,
et $(u_s\mathcal{F})(V)$ est un groupe abélien sans $p$-torsion
pour tout $V$ sur $S$. Puisque les
préfaisceaux représentables sont des faisceaux sur les sites fppf, il résulte de
Sites, lemme \ref{sites-lemma-pullback-representable-sheaf},
que $u_s\mathcal{G}$ est représenté par
$\mathbf{G}_{m, S}$. Comme $u_s\mathcal{F} \to u_s\mathcal{G}$
est injectif, nous obtenons un sous-groupe sans $p$-torsion
$$
(u_s\mathcal{F})(V) \subset \Gamma(V, \mathcal{O}_V)^*
$$
pour tout schéma $V$ sur $S$, avec la propriété suivante :
pour tout morphisme $V \to U$ de schémas sur $S$, où $U$ est
localement noethérien, le sous-groupe
$$
\mathcal{F}(U) \subset \Gamma(U, \mathcal{O}_U)^*
$$
s'envoie dans le sous-groupe $(u_s\mathcal{F})(V)$ par l'application de restriction
$\Gamma(U, \mathcal{O}_U)^* \to \Gamma(V, \mathcal{O}_V)^*$.

\medskip\noindent
La contradiction proprement dite s'obtient comme suit :
posons $k = \bigcup_{n \geq 0} \mathbf{F}_p(t^{1/{p^n}})$
et
$$
B = k \otimes_{\mathbf{F}_p(t)} k
$$
et soit $V = \Spec(B)$. Nous avons les deux morphismes de projection
$V \to \Spec(k)$ correspondant aux deux coprojections $k \to B$;
comme $\Spec(k)$ est noethérien, nous concluons que le sous-groupe
$$
(u_s\mathcal{F})(V) \subset B^*
$$
contient $k^* \otimes 1$ et $1 \otimes k^*$. C'est une contradiction,
car
$$
(t^{1/p} \otimes 1) \cdot (1 \otimes t^{-1/p}) =
t^{1/p} \otimes t^{-1/p}
$$
est une unité non triviale de $p$-torsion dans $B$.

\begin{lemma}
\label{examples-lemma-not-a-morphism-of-sites-noetherian-to-all}
Pour $S = \Spec(\mathbf{F}_p)$, le foncteur d'inclusion
$(\textit{Noetherian}/S)_{fppf} \to (\Sch/S)_{fppf}$
ne définit pas un morphisme de sites.
\end{lemma}

\begin{proof}
Voir la discussion qui précède.
\end{proof}





\section{Image directe dérivée des modules quasi-cohérents}
\label{examples-section-derived-push-quasi-coherent}

\noindent
Soit $k$ un corps de caractéristique $p > 0$. Soit $S = \Spec(k[x])$.
Soit $G = \mathbf{Z}/p\mathbf{Z}$, considéré soit comme groupe abstrait,
soit comme schéma en groupes constant sur $S$. Considérons le champ algébrique
$\mathcal{X} = [S/G]$, où $G$ agit trivialement sur $S$; voir
Exemples de champs, remarque \ref{examples-stacks-remark-X-mod-G-group},
et
Critères de représentabilité, lemme \ref{criteria-lemma-BG-algebraic}.
Considérons le morphisme structural
$$
f : \mathcal{X} \longrightarrow S
$$
Ce morphisme est quasi-compact et quasi-séparé. Nous obtenons donc un foncteur
$$
Rf_{\QCoh, *} :
D^{+}_\QCoh(\mathcal{O}_\mathcal{X})
\longrightarrow
D^{+}_\QCoh(\mathcal{O}_S),
$$
voir
Catégories dérivées des champs, proposition
\ref{stacks-perfect-proposition-derived-direct-image-quasi-coherent}.
Calculons $Rf_{\QCoh, *}\mathcal{O}_\mathcal{X}$.
Comme $D_\QCoh(\mathcal{O}_S)$ est équivalente à
la catégorie dérivée des $k[x]$-modules (voir
Catégories dérivées des schémas, lemme
\ref{perfect-lemma-affine-compare-bounded})
cela revient à calculer
$R\Gamma(\mathcal{X}, \mathcal{O}_\mathcal{X})$.
Nous pouvons pour cela utiliser le recouvrement $S \to \mathcal{X}$ et la suite
spectrale
$$
H^q(S \times_\mathcal{X} \ldots \times_\mathcal{X} S, O)
\Rightarrow H^{p + q}(\mathcal{X}, \mathcal{O}_\mathcal{X})
$$
voir
Cohomologie des champs, proposition
\ref{stacks-cohomology-proposition-smooth-covering-compute-cohomology}.
Remarquons que
$$
S \times_\mathcal{X} \ldots \times_\mathcal{X} S = S \times G^p
$$
qui est affine. Le complexe
$$
k[x] \to \text{Map}(G, k[x]) \to \text{Map}(G^2, k[x]) \to \ldots
$$
calcule donc $R\Gamma(\mathcal{X}, \mathcal{O}_\mathcal{X})$.
Ici, pour $\varphi \in \text{Map}(G^{p - 1}, k[x])$, sa différentielle est
l'application qui envoie $(g_1, \ldots, g_p)$ sur
$$
\varphi(g_2, \ldots, g_p) +
\sum\nolimits_{i = 1}^{p - 1}
(-1)^i\varphi(g_1, \ldots, g_i + g_{i + 1}, \ldots, g_p)
+ (-1)^p\varphi(g_1, \ldots, g_{p - 1}).
$$
C'est précisément le complexe qui calcule la cohomologie du groupe $G$ agissant
trivialement sur $k[x]$ (insérer ici une référence ultérieure). La cohomologie
du groupe cyclique $G$ à coefficients dans $k[x]$ est exactement une copie de $k[x]$ en chaque
degré cohomologique $\geq 0$ (insérer ici une référence ultérieure). Nous en concluons
que
$$
Rf_*\mathcal{O}_\mathcal{X} = \bigoplus\nolimits_{n \geq 0} \mathcal{O}_S[-n]
$$
Considérons maintenant le complexe
$$
E = \bigoplus\nolimits_{m \geq 0} \mathcal{O}_\mathcal{X}[m]
$$
C'est un objet de $D_\QCoh(\mathcal{O}_\mathcal{X})$. Interrompons la
discussion pour établir un résultat général.

\begin{lemma}
\label{examples-lemma-is-limit}
Soit $\mathcal{X}$ un champ algébrique. Soit $K$ un objet de
$D(\mathcal{O}_\mathcal{X})$ dont les faisceaux de cohomologie sont localement
quasi-cohérents (Faisceaux sur les champs, définition
\ref{stacks-sheaves-definition-locally-quasi-coherent})
et satisfont à la propriété de changement de base plat (Cohomologie des champs,
définition \ref{stacks-cohomology-definition-flat-base-change}).
Alors il existe un triangle distingué
$$
K \to
\prod\nolimits_{n \geq 0} \tau_{\geq -n} K \to
\prod\nolimits_{n \geq 0} \tau_{\geq -n} K \to K[1]
$$
dans $D(\mathcal{O}_\mathcal{X})$. Autrement dit, $K$ est la limite dérivée
de ses troncations canoniques.
\end{lemma}

\begin{proof}
Rappelons que nous travaillons sur le « grand site fppf » $\mathcal{X}_{fppf}$
de $\mathcal{X}$ (selon nos conventions
pour les faisceaux de $\mathcal{O}_\mathcal{X}$-modules dans les chapitres
Faisceaux sur les champs et Cohomologie sur les champs). Soit $\mathcal{B}$ l'ensemble
des objets $x$ de $\mathcal{X}_{fppf}$ qui sont au-dessus d'un schéma affine $U$.
En combinant
Faisceaux sur les champs, lemmes
\ref{stacks-sheaves-lemma-compare-fppf-etale},
\ref{stacks-sheaves-lemma-cohomology-restriction},
Descente, lemme \ref{descent-lemma-quasi-coherent-and-flat-base-change},
et
Cohomologie des schémas, lemme
\ref{coherent-lemma-quasi-coherent-affine-cohomology-zero}
nous voyons que $H^p(x, \mathcal{F}) = 0$ si $\mathcal{F}$ est
localement quasi-cohérent et $x \in \mathcal{B}$.
L'assertion découle alors de
Cohomologie sur les sites, lemme \ref{sites-cohomology-lemma-is-limit-dimension},
avec $d = 0$.
\end{proof}

\begin{lemma}
\label{examples-lemma-sum-is-product}
Soit $\mathcal{X}$ un champ algébrique. Si $\mathcal{F}_n$ est une famille
de faisceaux localement quasi-cohérents satisfaisant à la propriété de changement de base plat sur
$\mathcal{X}$, alors $\oplus_n \mathcal{F}_n[n] \to \prod_n \mathcal{F}_n[n]$
est un isomorphisme dans $D(\mathcal{O}_\mathcal{X})$.
\end{lemma}

\begin{proof}
En effet, le lemme \ref{examples-lemma-is-limit} montre que la somme directe
est isomorphe au produit.
\end{proof}

\noindent
Reprenons la discussion. Puisqu'un module quasi-cohérent est localement
quasi-cohérent et satisfait à la propriété de changement de base plat
(Faisceaux sur les champs, lemme \ref{stacks-sheaves-lemma-quasi-coherent}), nous obtenons
$$
E = \prod\nolimits_{m \geq 0} \mathcal{O}_\mathcal{X}[m]
$$
Puisque la cohomologie commute
aux limites, nous avons
$$
Rf_*E = \prod\nolimits_{m \geq 0}
\left(\bigoplus\nolimits_{n \geq 0} \mathcal{O}_S[m - n]\right)
$$
Remarquons que ce complexe n'est pas un objet de $D_\QCoh(\mathcal{O}_S)$,
car son faisceau de cohomologie en degré $0$ est un produit infini de copies
de $\mathcal{O}_S$, qui n'est même pas un
$\mathcal{O}_S$-module localement quasi-cohérent.

\begin{lemma}
\label{examples-lemma-push-not-OK}
Un morphisme quasi-compact et quasi-séparé
$f : \mathcal{X} \to \mathcal{Y}$ de champs algébriques
n'induit pas nécessairement un foncteur
$Rf_* : D_\QCoh(\mathcal{O}_\mathcal{X}) \to
D_\QCoh(\mathcal{O}_\mathcal{Y})$.
\end{lemma}

\begin{proof}
Voir la discussion qui précède.
\end{proof}



\section{Une grande catégorie abélienne}
\label{examples-section-big}

\noindent
Le but de cette section est de donner un exemple d'une « grande » catégorie abélienne
$\mathcal{A}$ et d'objets $M, N$ tels que la collection des classes d'isomorphisme
d'extensions $\Ext_\mathcal{A}(M, N)$ ne soit pas un ensemble.
Cet exemple est dû à Freyd; voir \cite[page 131, Exercise A]{Freyd}.

\medskip\noindent
Nous définissons $\mathcal{A}$ comme suit. Un objet de $\mathcal{A}$
est un triplet $(M, \alpha, f)$, où $M$ est un groupe abélien,
$\alpha$ un ordinal et $f : \alpha \to \text{End}(M)$ une application.
Un morphisme $(M, \alpha, f) \to (M', \alpha', f')$ est donné par un homomorphisme
de groupes abéliens $\varphi : M \to M'$ tel que, pour {\it tout} ordinal $\beta$,
on ait
$$
\varphi \circ f(\beta) = f'(\beta) \circ \varphi
$$
La convention est de poser $f(\beta) = 0$ si $\beta$ n'appartient pas à
$\alpha$ et, de même, $f'(\beta)$ est nul si $\beta$ n'est pas un
élément de $\alpha'$. Nous omettons la vérification que la catégorie ainsi définie
est abélienne.

\medskip\noindent
Considérons l'objet $Z = (\mathbf{Z}, \emptyset, f)$, c'est-à-dire que tous les
opérateurs sont nuls. L'observation est que,
calculé dans $\mathcal{A}$, le groupe $\Ext^1_\mathcal{A}(Z, Z)$ est une classe propre
et non un ensemble. En effet, pour tout ordinal $\alpha$, nous pouvons construire une extension
$(M, \alpha + 1, f)$ de $Z$ par $Z$ dont le groupe sous-jacent est
$M = \mathbf{Z} \oplus \mathbf{Z}$ et pour laquelle $f$ est
partout nulle sauf en
$$
f(\alpha) =
\left(
\begin{matrix}
0 & 1 \\
0 & 0
\end{matrix}
\right).
$$
Cela produit manifestement une classe propre de classes d'isomorphisme d'extensions.
En particulier, la catégorie dérivée de $\mathcal{A}$ possède des classes propres pour
ses collections de morphismes; voir
Catégories dérivées, lemme \ref{derived-lemma-ext-1}. Il faut donc faire preuve
de prudence dans la définition des quotients de Verdier de catégories
triangulées.

\begin{lemma}
\label{examples-lemma-big-abelian-category}
Il existe une « grande » catégorie abélienne $\mathcal{A}$ dont les
groupes $\Ext$ sont des classes propres.
\end{lemma}

\begin{proof}
Voir la discussion qui précède.
\end{proof}





\section{Points faiblement associés et densité schématique}
\label{examples-section-weak-ass-dense}

\noindent
Soit $k$ un corps. Soit $R = k[z, x_i, y_i]/(z^2, zx_iy_i)$, où
$i$ parcourt les éléments de $\mathbf{N}$. Remarquons que $R = R_0 \oplus M_0$, où
$R_0 = k[x_i, y_i]$ est un sous-anneau et $M_0$ un idéal de carré nul,
avec $M_0 \cong R_0/(x_iy_i)$ comme $R_0$-module. L'idéal premier
$\mathfrak p = (z, x_i)$ est faiblement associé à $R$ considéré comme $R$-module
(Algèbre, définition \ref{algebra-definition-weakly-associated}).
En effet, l'élément $z$ de $R_\mathfrak p$ est non nul mais annulé
par $\mathfrak pR_\mathfrak p$. Considérons d'autre part le sous-schéma
ouvert
$$
U = \bigcup D(x_i) \subset \Spec(R) = S
$$
Nous affirmons que $U \subset S$ est schématiquement dense
(Morphismes, définition \ref{morphisms-definition-scheme-theoretically-dense}).
Pour le démontrer, il suffit de voir que $\mathcal{O}_S \to j_*\mathcal{O}_U$
est injectif, où $j : U \to S$ est le morphisme d'inclusion; voir
Morphismes, lemme \ref{morphisms-lemma-characterize-scheme-theoretically-dense}.
Traduit en termes algébriques, il faut montrer que, pour tout $g \in R$,
l'application
$$
R_g \longrightarrow \prod R_{x_ig}
$$
est injective. Écrivons $g = g_0 + m_0$, avec $g_0 \in R_0$ et $m_0 \in M_0$.
Alors $R_g = R_{g_0}$ (détails omis). Nous pouvons donc supposer $g \in R_0$.
Nous pouvons aussi supposer $g$ non nul. Alors $R_g = (R_0)_g \oplus (M_0)_g$.
Puisque $R_0$ est intègre, l'application $(R_0)_g \to \prod (R_0)_{x_ig}$ est
injective. Si $g \in (x_iy_i)$, alors $(M_0)_g = 0$ et il n'y a rien
à démontrer. Si $g \not \in (x_iy_i)$, comme
$(x_iy_i)$ est un idéal radical de $R_0$, il reste à montrer que
$M_0 \to \prod (M_0)_{x_ig}$ est injective. Le noyau de
$R_0 \to M_0 \to (M_0)_{x_n}$ est $(x_iy_i, y_n)$. Comme
$(x_iy_i, y_n)$ est un idéal radical, si $g \not \in (x_iy_i, y_n)$,
alors le noyau de $R_0 \to M_0 \to (M_0)_{x_ng}$ est $(x_iy_i, y_n)$.
Puisque $g \not \in (x_iy_i, y_n)$ pour tout $n \gg 0$, nous concluons que
le noyau est contenu dans $\bigcap_{n \gg 0} (x_iy_i, y_n) = (x_iy_i)$,
comme voulu.

\medskip\noindent
Voici un second exemple, dû à Ofer Gabber. Soit $k$ un corps, et soient
$R$, resp.\ $R'$, l'anneau des fonctions $\mathbf{N} \to k$, resp.\ l'anneau
des fonctions $\mathbf{N} \to k$ constantes à partir d'un certain rang. Alors $\Spec(R)$,
resp.\ $\Spec(R')$, est le compactifié de
Stone-{\v C}ech\footnote{Tout élément $f \in R$ est de la
forme $ue$, où $u$ est une unité et $e$ un idempotent. Alors
Algèbre, lemme \ref{algebra-lemma-ring-with-only-minimal-primes},
montre que $\Spec(R)$ est séparé. D'autre part, $\mathbf{N}$ muni de la
topologie discrète peut être considéré comme un ouvert dense. Étant donnée une application d'ensembles
$\mathbf{N} \to X$ vers un espace topologique séparé et quasi-compact $X$,
nous obtenons un homomorphisme d'anneaux $\mathcal{C}^0(X; k) \to R$, où $\mathcal{C}^0(X; k)$
est la $k$-algèbre des applications localement constantes $X \to k$. Cela donne
$\Spec(R) \to \Spec(\mathcal{C}^0(X; k)) = X$ et démontre la propriété
universelle.} $\beta\mathbf{N}$, resp.\ le compactifié à
un point\footnote{On montre ici qu'il n'existe effectivement
qu'un seul idéal maximal supplémentaire dans $R'$.}
$\mathbf{N}^* = \mathbf{N} \cup \{\infty\}$.
Tous les points sont faiblement associés, puisque tous les idéaux premiers sont minimaux dans les
anneaux $R$ et $R'$.

\begin{lemma}
\label{examples-lemma-example-schematically-dense-missing-weakly-associated-point}
Il existe un schéma réduit $X$ et un ouvert schématiquement dense
$U \subset X$ tel qu'un certain point faiblement associé $x \in X$ n'appartienne pas à $U$.
\end{lemma}

\begin{proof}
Dans le premier exemple, $\mathfrak p \not \in U$ par construction.
Dans les exemples de Gabber, les schémas $\Spec(R)$ et $\Spec(R')$ sont réduits.
\end{proof}



\section{Exemple de non-additivité des traces}
\label{examples-section-non-additive}

\noindent
Soit $k$ un corps et soit $R = k[\epsilon]$ l'anneau des nombres duaux
sur $k$. Autrement dit, $R = k[x]/(x^2)$ et $\epsilon$ est la classe
de $x$ dans $R$. Considérons la suite exacte courte de complexes
$$
\xymatrix{
0 \ar[d] \ar[r] &
R \ar[d]^\epsilon \ar[r]_1 &
R \ar[d] \\
R \ar[r]^1 & R \ar[r] & 0
}
$$
Ici, les colonnes sont les complexes, la première ligne est placée en degré $0$ et
la seconde en degré $1$. Désignons le premier complexe (c'est-à-dire la
colonne de gauche) par $A^\bullet$, le deuxième par $B^\bullet$ et le troisième par
$C^\bullet$. Nous affirmons que le diagramme
\begin{equation}
\label{examples-equation-commutes-up-to-homotopy}
\vcenter{
\xymatrix{
A^\bullet \ar[d]_{1 + \epsilon} \ar[r] &
B^\bullet \ar[r] \ar[d]_1 &
C^\bullet \ar[d]_1 \\
A^\bullet \ar[r] & B^\bullet \ar[r] & C^\bullet
}
}
\end{equation}
est commutatif dans $K(R)$, c'est-à-dire qu'il s'agit d'un diagramme de complexes commutatif à homotopie près.
En effet, le carré de droite est commutatif et celui de gauche l'est modulo
l'homotopie $1 : A^1 \to B^0$. En revanche,
$$
\text{Tr}_{A^\bullet}(1 + \epsilon) + \text{Tr}_{C^\bullet}(1)
\not = \text{Tr}_{B^\bullet}(1).
$$

\begin{lemma}
\label{examples-lemma-nonadditivity-of-trace}
Il existe un anneau $R$, un triangle distingué
$(K, L, M, \alpha, \beta, \gamma)$ dans la catégorie homotopique $K(R)$,
et un endomorphisme $(a, b, c)$ de ce triangle distingué, tels que
$K$, $L$, $M$ soient des complexes parfaits et que
$\text{Tr}_K(a) + \text{Tr}_M(c) \not = \text{Tr}_L(b)$.
\end{lemma}

\begin{proof}
Considérons l'exemple ci-dessus. L'application $\gamma : C^\bullet \to A^\bullet[1]$
est donnée par multiplication par $\epsilon$ en degré $0$; voir
Catégories dérivées, définition \ref{derived-definition-distinguished-triangle}.
Il est donc également vrai que
$$
\xymatrix{
C^\bullet \ar[d] \ar[r]_\gamma & A^\bullet[1] \ar[d] \\
C^\bullet \ar[r]^\gamma & A^\bullet[1]
}
$$
est commutatif dans $K(R)$, puisque $\epsilon(1 + \epsilon) = \epsilon$.
Nous avons donc bien un morphisme de triangles distingués.
\end{proof}




\section{Cohomologie non finie du faisceau structural d'un schéma projectif}
\label{examples-section-nonfinite-h0}

\noindent
Soient $R$ un anneau non noethérien quelconque et $I$ un idéal de $R$
qui n'est pas de type fini. Soit $S = R[x, y]/xyI$, considéré
comme anneau gradué où $x$ et $y$ sont de degré $1$.
Alors $X = \text{Proj}(S)$ est un schéma projectif sur $R$.
Nous affirmons que $H^0(X, \mathcal{O}_X)$ n'est pas
un $R$-module de type fini.

\medskip\noindent
En décomposant les anneaux $S_x$, $S_y$ et $S_{xy}$ suivant le bidegré,
nous pouvons écrire
$$
S_x = \bigoplus_{\substack{(a, b) \in \mathbf{Z}^2 \\ b \geq 0}} R'_{a, b},
\quad\quad
S_y = \bigoplus_{\substack{(a, b) \in \mathbf{Z}^2 \\ a \geq 0}} R''_{a, b},
\quad\quad
S_{xy} = \bigoplus_{\substack{(a, b) \in \mathbf{Z}^2}} R/I,
$$
où les facteurs directs d'indice $(a, b)$ sont en degré $a + b$ et où
$$
R'_{a, b} =
\left\{
\begin{matrix}
R & b = 0, \\
R/I & b \ne 0,
\end{matrix}
\right.
\quad\quad
R''_{a, b} =
\left\{
\begin{matrix}
R & a = 0, \\
R/I & a \ne 0.
\end{matrix}
\right.
$$
Le $R$-module $H^0(X, \mathcal{O}_X)$ se calcule comme le noyau de l'application
$$
(S_x)_0 \times (S_y)_0 \to (S_{xy})_0
$$
définie par $(F, G) \mapsto F - G$. Ici, la notation $A_0$ désigne la
composante de degré $0$ d'un anneau $\mathbf{Z}$-gradué $A$. D'après la description explicite
ci-dessus de $S_x, S_y, S_{xy}$, ce noyau est isomorphe
au sous-module $M \subset R^2$ formé des couples $(F, G)$
tels que $F - G \in I$. Puisque l'application $M \to I$ de $R$-modules
donnée par $(F, G) \mapsto F - G$ est surjective,
$M$ n'est pas un $R$-module de type fini.

\medskip\noindent
La surjection $R[x, y] \to S$ correspond à une
immersion fermée $i : X \to \mathbf{P}^1_R$. Alors $\mathcal{F} = i_*\mathcal{O}_X$
est un $\mathcal{O}_{\mathbf{P}^1_R}$-module quasi-cohérent de type fini
dont le $H^0$ est le module $M$ non de type fini
calculé ci-dessus.

\begin{lemma}
\label{examples-lemma-nonfinite-h0}
$H^0$ non fini.
\begin{enumerate}
\item Il existe un anneau $R$ et un schéma projectif $X$
sur $R$ tels que $H^0(X, \mathcal{O}_X)$ ne soit pas
un $R$-module de type fini.
\item Il existe un anneau $R$ et un module quasi-cohérent de type fini
$\mathcal{F}$ sur $\mathbf{P}^1_R$ tels que
$H^0(\mathbf{P}^1_R, \mathcal{F})$ ne soit pas un $R$-module de type fini.
\end{enumerate}
\end{lemma}

\begin{proof}
Voir la discussion qui précède.
\end{proof}









\section{Être projectif n'est pas local sur la base}
\label{examples-section-non-descending-property-projective}

\noindent
Dans le chapitre sur la descente, nous avons vu que de nombreuses propriétés des morphismes
sont locales sur la base, même pour la topologie fpqc. Voir
Descente, sections \ref{descent-section-descending-properties-morphisms},
\ref{descent-section-descending-properties-morphisms-fpqc} et
\ref{descent-section-descending-properties-morphisms-fppf}.
Ce n'est pas le cas de la projectivité des morphismes.

\begin{lemma}
\label{examples-lemma-non-descending-property-projective}
Les propriétés
\begin{enumerate}
\item[] $\mathcal{P}(f) =$ « $f$ est projectif », et
\item[] $\mathcal{P}(f) =$ « $f$ est quasi-projectif »
\end{enumerate}
ne sont pas locales sur la base pour la topologie de Zariski. A fortiori, elles ne sont pas locales
sur la base pour la topologie fpqc.
\end{lemma}

\begin{proof}
Suivant Hironaka \cite[exemple B.3.4.1]{H},
nous construisons un morphisme propre entre variétés complexes lisses de dimension 3
$f:V_Y\to Y$
qui est localement projectif pour la topologie de Zariski, mais non projectif. Comme $f$ est propre
et non projectif, il n'est pas non plus quasi-projectif.

\medskip\noindent
Soit $Y$ l'espace projectif de dimension 3 sur les nombres complexes ${\mathbf C}$.
Soient $C$ et $D$ des coniques lisses de $Y$ telles
que le sous-schéma fermé $C\cap D$ soit réduit et consiste
en deux points complexes $P$ et $Q$. (Par exemple,
prenons $C=\{ [x,y,z,w]: xy=z^2, w=0\}$, $D=\{ [x,y,z,w]:
xy=w^2, z=0\}$, $P=[1,0,0,0]$,
et $Q=[0,1,0,0]$.)
Sur $Y-Q$, éclatons d'abord la courbe $C$, puis la
transformée stricte de la courbe $D$ (Diviseurs, définition
\ref{divisors-definition-strict-transform}). Sur $Y-P$, éclatons d'abord
la courbe $D$, puis la transformée stricte de la courbe
$C$. Au-dessus de $Y-P-Q$, les deux variétés ainsi construites sont canoniquement
isomorphes; nous pouvons donc les recoller au-dessus de $Y-P-Q$. Le résultat
est une variété lisse propre $V_Y$ de dimension 3 sur ${\mathbf C}$. Le morphisme
$f:V_Y\to Y$ est propre et localement projectif pour la topologie de Zariski (car
c'est un éclatement au-dessus de $Y-P$ et de $Y-Q$), d'après Diviseurs,
lemme \ref{divisors-lemma-blowing-up-projective}. Nous allons montrer que
$V_Y$ n'est pas projectif sur ${\mathbf C}$. Il en résultera que
$f$ n'est pas projectif.

\medskip\noindent
Pour cela, soit $L$ l'image réciproque dans $V_Y$ d'un point complexe
de $C-P-Q$, et $M$ celle d'un point complexe de $D-P-Q$.
Alors $L$ et $M$ sont isomorphes à la droite projective
${\mathbf P}^1_{{\mathbf C}}$.
Soit ensuite $E$ le sous-schéma de $V_Y$ donné par l'image réciproque de $C\cup D\subset Y$ dans $V_Y$;
ainsi, $E\rightarrow C\cup D$ est un morphisme propre, dont les fibres
sont isomorphes à ${\mathbf P}^1$ au-dessus de $(C\cup D)-\{P,Q\}$.
L'image
réciproque de $P$ dans $E$ est la réunion de deux droites $L_0$ et $M_0$, et nous avons
les équivalences rationnelles de cycles $L\sim L_0+M_0$ et $M\sim M_0$ sur $E$
(puisque $C$ et $D$ sont isomorphes à ${\mathbf P}^1$).
Remarquons l'asymétrie qui résulte de l'ordre dans lequel nous avons éclaté
les deux courbes. Au voisinage de $Q$, la situation inverse se produit. L'image réciproque
de $Q$ est donc la réunion de deux droites $L_0'$ et $M_0'$, et nous avons
les équivalences rationnelles $L\sim L_0'$ et $M\sim L_0'+M_0'$ sur $E$.
En combinant ces équivalences, nous obtenons $L_0+M_0'\sim 0$
sur $E$, donc sur $V_Y$. Si $V_Y$ était projectif sur ${\mathbf C}$,
il posséderait
un fibré inversible ample $H$, de degré $> 0$ sur toute courbe
de $V_Y$. En particulier, $H$ serait de degré positif sur $L_0+M_0'$,
ce qui contredit le fait que le degré d'un fibré inversible est bien défini
sur les 1-cycles modulo équivalence rationnelle d'un schéma propre
sur un corps (Homologie de Chow,
lemme \ref{chow-lemma-proper-pushforward-rational-equivalence}
et lemme \ref{chow-lemma-factors}).
Ainsi, $V_Y$ n'est pas projectif sur ${\mathbf C}$.
\end{proof}

\noindent
Dans une autre terminologie, la variété de dimension 3 de Hironaka $V_Y$ est une petite
résolution de l'éclatement $Y'$ de $Y$ le long du sous-schéma réduit
$C\cup D$; ici, $Y'$ possède deux singularités nodales. Si l'on définit $Z$ en éclatant
$Y$ le long de $C$, puis le long de la transformée stricte
de $D$, alors $Z$ est une variété projective lisse de dimension 3, et la variété non projective
$V_Y$ diffère de $Z$ par un « flop » au-dessus de $Y-P$.


\section[Descente non effective pour des schémas projectifs]{Données de descente non effectives pour les schémas projectifs}
\label{examples-section-non-effective-descent-projective}

\noindent
Dans le chapitre sur la descente, nous avons vu que les données de descente de schémas relativement
à un morphisme fpqc sont effectives pour plusieurs classes
de morphismes. En particulier, les morphismes affines et, plus généralement,
les morphismes quasi-affines satisfont à la descente pour les recouvrements fpqc
(Descente, lemme \ref{descent-lemma-quasi-affine}).
Ce n'est pas le cas des morphismes projectifs.

\begin{lemma}
\label{examples-lemma-non-effective-descent-projective}
Il existe un recouvrement étale de schémas $X\to S$ et une donnée de descente
$(V/X,\varphi)$ relativement à $X\to S$ telle que
$V\to X$ soit projectif,
mais que la donnée de descente ne soit pas effective dans la catégorie des schémas.
\end{lemma}

\begin{proof}
Nous imitons l'exemple de Hironaka d'un espace algébrique complexe
lisse et séparé de dimension 3
qui n'est pas un schéma \cite[exemple B.3.4.2]{H}.

\medskip\noindent
Considérons l'action du groupe $G = \mathbf{Z}/2 = \{1, g\}$
sur l'espace projectif de dimension 3
$\mathbf{P}^3$ sur les nombres complexes, donnée par
$$
g[x,y,z,w] = [y,x,w,z].
$$
L'action est libre en dehors des deux droites disjointes
$L_1=\{ [x,x,z,z]\}$ et $L_2=\{ [x,-x,z,-z]\}$ de
${\mathbf P}^3$. Soit $Y={\mathbf P}^3-(L_1\cup L_2)$. Il existe un
schéma quasi-projectif lisse $S=Y/G$ sur ${\mathbf C}$ tel que
$Y\to S$ soit un torseur sous $G$ (Groupoïdes,
définition \ref{groupoids-definition-principal-homogeneous-space}).
Explicitement, nous pouvons définir $S$ comme l'image de l'ouvert $Y$
de ${\mathbf P}^3$ par le morphisme
\begin{align*}
{\mathbf P}^3 & \to \text{Proj } {\mathbf C}[x,y,z,w]^G\\
   & = \text{Proj } {\mathbf C}[u_0,u_1,v_0,v_1,v_2]/(v_0v_1=v_2^2),
\end{align*}
où $u_0=x+y$, $u_1=z+w$, $v_0=(x-y)^2$, $v_1=(z-w)^2$
et $v_2=(x-y)(z-w)$, l'anneau étant gradué avec $u_0,u_1$
en degré 1 et $v_0,v_1,v_2$ en degré 2.

\medskip\noindent
Soient $C=\{ [x,y,z,w]: xy=z^2, w=0\}$ et $D=\{ [x,y,z,w]:
xy=w^2, z=0\}$. Ce sont des coniques lisses de ${\mathbf P}^3$, contenues
dans l'ouvert $G$-invariant $Y$, avec $g(C)=D$. De plus,
$C\cap D$ consiste en deux points $P:=[1,0,0,0]$
et $Q:=[0,1,0,0]$, que l'action
de $G$ échange.

\medskip\noindent
Soit $V_Y\to Y$ le schéma défini au-dessus de $Y-P$
en éclatant $D$, puis la transformée stricte
de $C$, et au-dessus de $Y-Q$ en éclatant $C$, puis
la transformée stricte de $D$. (C'est la même construction
que dans la démonstration du lemme \ref{examples-lemma-non-descending-property-projective},
à ceci près que $Y$ désigne ici un ouvert de ${\mathbf P}^3$
et non ${\mathbf P}^3$ tout entier.)
Alors l'action de $G$ sur $Y$ se relève en une action de $G$ sur $V_Y$
qui échange les images réciproques de $Y-P$ et de $Y-Q$. Cette action
de $G$ sur $V_Y$ fournit une donnée de descente $(V_Y/Y,\varphi_Y)$
sur $V_Y$ relativement au torseur sous $G$
$Y\to S$. Le morphisme $V_Y\to Y$
est propre mais non projectif, comme le montre la démonstration
du lemme \ref{examples-lemma-non-descending-property-projective}.

\medskip\noindent
Soit $X$ la réunion disjointe des ouverts $Y-P$ et $Y-Q$;
nous avons alors des morphismes étales surjectifs $X\to Y\to S$.
Soit $V$ l'image réciproque de $V_Y\to Y$ sur $X$; alors le morphisme
$V\to X$ est projectif, puisque $V_Y\to Y$ est un éclatement au-dessus de chacun des ouverts
$Y-P$ et $Y-Q$. De plus, la donnée de descente $(V_Y/Y,\varphi_Y)$
induit par image réciproque une donnée de descente $(V/X,\varphi)$ relativement au
recouvrement étale $X\to S$.

\medskip\noindent
Supposons cette donnée de descente effective dans la catégorie
des schémas. Il existe alors un schéma $U\to S$
dont l'image réciproque est le morphisme $V\to X$, muni
de sa donnée de descente. Alors $U$ serait le quotient
de $V_Y$ par l'action de $G$.
$$
\xymatrix{
V \ar[r]\ar[d]&  X\ar[d] \\
V_Y \ar[r]\ar[d]& Y\ar[d] \\
U \ar[r]& S
}
$$

\medskip\noindent
Soit $E$ l'image réciproque de $C\cup D\subset Y$ dans $V_Y$;
ainsi, $E\rightarrow C\cup D$ est un morphisme propre, dont les fibres
sont isomorphes à ${\mathbf P}^1$ au-dessus de $(C\cup D)-\{P,Q\}$.
L'image réciproque de $P$ dans $E$ est la réunion de deux droites $L_0$
et $M_0$. Il en résulte que l'image réciproque de $Q=g(P)$ dans $E$
est la réunion des deux droites $L_0'=g(M_0)$ et $M_0'=g(L_0)$.
Comme le montre la démonstration
du lemme \ref{examples-lemma-non-descending-property-projective},
nous avons une équivalence rationnelle $L_0+M_0'=L_0+g(L_0)\sim 0$ sur $E$.

\medskip\noindent
Par descente des sous-schémas fermés, il existe une courbe $L_1\subset U$
(isomorphe à ${\mathbf P}^1$)
dont l'image réciproque dans $V_Y$ est $L_0\cup g(L_0)$. (Utiliser Descente, lemme
\ref{descent-lemma-affine}, en remarquant qu'une immersion fermée est un morphisme
affine.)
Soit $R$ un point complexe de $L_1$. Puisque
nous avons supposé que $U$ est un schéma, nous pouvons choisir une fonction
$f$ dans l'anneau local $O_{U,R}$ qui s'annule en $R$, mais non
sur toute la courbe $L_1$. Soit $D_{\text{loc}}$ une composante irréductible
du fermé $\{f = 0\}$ de $\Spec O_{U,R}$; alors
$D_{\text{loc}}$ est de codimension 1.
L'adhérence de $D_{\text{loc}}$ dans $U$ est un diviseur irréductible $D_U$
de $U$ qui contient le point $R$, mais non toute la courbe $L_1$.
L'image réciproque de $D_U$ dans $V_Y$ est un diviseur effectif $D$
qui rencontre $L_0\cup g(L_0)$ mais ne contient aucune des deux
courbes $L_0$ et $g(L_0)$.

\medskip\noindent
Puisque la variété complexe $V_Y$ de dimension 3 est lisse, $O(D)$ est un fibré
inversible sur $V_Y$. Nous utilisons ici le fait qu'un anneau local régulier est à factorisation unique,
autrement dit factoriel; voir
Compléments sur l'algèbre, lemme \ref{more-algebra-lemma-regular-local-UFD}.
La restriction de $O(D)$ à la surface propre
$E\subset V_Y$ est un fibré inversible de degré positif sur le 1-cycle
$L_0+g(L_0)$, d'après les propriétés de $D$. Puisque
$L_0+g(L_0)\sim 0$ sur $E$, cela contredit
le fait que le degré d'un fibré inversible est bien défini
sur les 1-cycles modulo équivalence rationnelle d'un schéma propre
sur un corps (Homologie de Chow,
lemme \ref{chow-lemma-proper-pushforward-rational-equivalence}
et lemme \ref{chow-lemma-factors}). La donnée de descente
$(V/X,\varphi)$ n'est donc pas effective; autrement dit, $U$ n'existe pas
comme schéma.
\end{proof}

\noindent
Dans cet exemple, la donnée de descente {\it est }effective dans la catégorie
des espaces algébriques. Plus précisément,
$U$ existe comme espace algébrique lisse et séparé
de dimension 3 sur ${\mathbf C}$,
par exemple d'après Espaces algébriques, lemme \ref{spaces-lemma-quotient}.
La variété $U$ de Hironaka, de dimension 3, est une petite résolution de l'éclatement $S'$
de la variété quasi-projective lisse $S$ de dimension 3 le long de la courbe nodale irréductible
$(C\cup D)/G$; la variété $S'$ de dimension 3 possède une singularité nodale. L'autre
petite résolution de $S'$ (qui diffère de $U$ par un « flop »)
est encore un espace algébrique qui n'est pas un schéma.






\section{Une famille de courbes dont l'espace total n'est pas un schéma}
\label{examples-section-family-of-curves}

\noindent
Dans Quot, section \ref{quot-section-curves}, nous définissons une famille de courbes sur
un schéma $S$ comme un morphisme propre, plat et de présentation finie
de dimension relative $\leq 1$ d'un espace algébrique $X$ vers $S$.
Si $S$ est le spectre d'un anneau local noethérien complet, alors
$X$ est un schéma ; voir
Compléments sur les morphismes d'espaces, lemme
\ref{spaces-more-morphisms-lemma-projective-over-complete}.
Dans cette section, nous montrons que ce résultat est faux en général.

\medskip\noindent
Soit $k$ un corps. Partons d'un morphisme propre et plat
$$
Y \longrightarrow \mathbf{A}^1_k
$$
et d'un point $y \in Y(k)$ situé au-dessus de $0 \in \mathbf{A}^1_k(k)$
possédant les propriétés suivantes :
\begin{enumerate}
\item la fibre $Y_0$ est une courbe lisse géométriquement irréductible sur $k$,
\item pour tout sous-schéma fermé propre $T \subset Y$ dominant
$\mathbf{A}^1_k$, l'intersection $T \cap Y_0$ contient au moins un point
distinct de $y$.
\end{enumerate}
À partir d'une telle surface, nous construisons notre exemple comme suit.
$$
\xymatrix{
Y \ar[rd] & Z \ar[d] \ar[l] \ar[r] & X \ar[ld] \\
& \mathbf{A}^1_k
}
$$
Ici, $Z \to Y$ est l'éclatement de $Y$ en $y$. Soit $E \subset Z$ le
diviseur exceptionnel et soit $C \subset Z$ la transformée stricte de $Y_0$.
On a $Z_0 = E \cup C$ au sens schématique (pour le voir, utiliser que
$Y$ est lisse en $y$ et que, de plus, $Y \to \mathbf{A}^1_k$ est lisse en $y$).
D'après les résultats d'Artin (\cite{ArtinII} ; utiliser
Réduction semi-stable, lemme \ref{models-lemma-properties-form}
pour voir que le fibré normal de $C$ est négatif),
nous pouvons contracter la courbe $C$ dans $Z$ et obtenir un espace algébrique $X$
comme dans le diagramme. Soit $x \in X(k)$ l'image de $C$.

\medskip\noindent
Nous affirmons que $X$ n'est pas un schéma. En effet, s'il en était un,
il existerait un voisinage ouvert affine $U \subset X$ de $x$.
Posons $T = X \setminus U$. Alors $T$ domine $\mathbf{A}^1_k$
(car les fibres de $X \to \mathbf{A}^1_k$ sont propres de dimension $1$
alors que les fibres de $U \to \mathbf{A}^1_k$ sont affines, donc différentes).
Soit $T' \subset Z$ le sous-schéma fermé qui s'envoie isomorphiquement
sur $T$ (car $x \not \in T$). Alors l'image de $T'$ dans $X$
contredit la condition (2) ci-dessus (car $T' \cap Z_0$ est contenu dans
le diviseur exceptionnel $E$ de l'éclatement $Z \to Y$).

\medskip\noindent
Pour achever la discussion, il nous reste à construire $Y$.
Nous supposerons que la caractéristique de $k$ n'est pas $3$.
Écrivons $\mathbf{A}^1_k = \Spec(k[t])$ et prenons
$$
Y \quad : \quad T_0^3 + T_1^3 + T_2^3 - tT_0T_1T_2 = 0
$$
dans $\mathbf{P}^2_{k[t]}$. La fibre pour $t = 0$
est une courbe projective lisse de genre $1$. Sur l'ouvert affine
$V_+(T_0)$, nous obtenons l'équation affine
$$
1 + x^3 + y^3 - txy = 0
$$
qui définit une surface lisse sur $k$. Par symétrie, il en va de même sur les
autres ouverts affines, et nous voyons donc que $Y$ est une surface lisse.
Enfin, l'équation affine montre aussi que le corps des fractions
est $k(x, y)$ ; la surface $Y$ est donc rationnelle.
Or le groupe de Picard d'une surface rationnelle est de type fini
(insérer ici une référence future).
Ainsi, pour choisir $y \in Y_0(k)$ possédant la propriété (2),
il suffit de choisir $y$ de telle sorte que
\begin{equation}
\label{examples-equation-three}
\mathcal{O}_{Y_0}(ny) \not \in \Im(\Pic(Y) \to \Pic(Y_0))
\text{ pour tout }n > 0
\end{equation}
En effet, la somme des composantes irréductibles de dimension $1$ d'un $T$
contredisant (2) donnerait un diviseur de Cartier effectif dont l'intersection avec $Y_0$
serait le diviseur $ny$ pour un certain $n \geq 1$, et nous conclurions
que $\mathcal{O}_{Y_0}(ny)$ appartient à l'image de l'application de restriction.
Remarquons que, puisque $Y_0$ est de genre $\geq 1$, l'application
$$
Y_0(k) \to \Pic(Y_0),\quad y \mapsto \mathcal{O}_{Y_0}(y)
$$
est injective. Si maintenant $k$ est un corps algébriquement clos non dénombrable,
la dénombrabilité de $\Pic(Y)$ et la remarque
précédente permettent de trouver un $y \in Y_0(k)$ satisfaisant (\ref{examples-equation-three}),
et donc (2).

\begin{lemma}
\label{examples-lemma-family-of-curves-not-scheme}
Il existe un corps $k$ et une famille de courbes
$X \to \mathbf{A}^1_k$ telle que $X$ ne soit pas un schéma.
\end{lemma}

\begin{proof}
Voir la discussion ci-dessus.
\end{proof}






\section{Changement de base dérivé}
\label{examples-section-derived-base-change}

\noindent
Soit $R \to R'$ un morphisme d'anneaux. Dans
Compléments d'algèbre, section \ref{more-algebra-section-derived-base-change},
nous construisons un foncteur de changement de base dérivé
$- \otimes_R^\mathbf{L} R' : D(R) \to D(R')$.
Soit ensuite $R \to A$ un second morphisme d'anneaux. Considérons le diagramme
$$
\xymatrix{
	A \ar[r] & A \otimes_R R' \ar@{=}[r] & A' \\
R \ar[u] \ar[r] & R' \ar[u] \ar[ur]
}
$$
Étant donné un $A$-module $M$, le produit tensoriel $M \otimes_R R'$ est un
$A \otimes_R R'$-module, c'est-à-dire un $A'$-module. Au morphisme d'anneaux
$A \to A'$ est associé un foncteur dérivé
$$
- \otimes_A^\mathbf{L} A' : D(A) \longrightarrow D(A')
$$
mais, en général, ce foncteur ne coïncide pas avec $- \otimes_R^\mathbf{L} R'$.
Plus précisément, pour $K \in D(A)$, le morphisme canonique
$$
K \otimes_R^{\mathbf{L}} R' \longrightarrow
K \otimes_A^{\mathbf{L}} A'
$$
dans $D(R')$, construit dans
Compléments d'algèbre, équation (\ref{more-algebra-equation-comparison-map}),
n'est pas un isomorphisme en général. On peut donc se demander s'il existe
un « foncteur de changement de base dérivé » $T : D(A) \to D(A')$, c'est-à-dire un foncteur
tel que l'image de $T(K)$ soit $K \otimes_R^\mathbf{L} R'$ dans $D(R')$.
Dans cette section, nous montrons qu'un tel foncteur n'existe pas en général.

\medskip\noindent
Soit $k$ un corps. Posons $R = k[x, y]$, $R' = R/(xy)$ et $A = R/(x^2)$.
L'objet $A \otimes_R^\mathbf{L} R'$ de $D(R')$ est représenté par
$$
x^2 : R' \longrightarrow R'
$$
et $H^0(A \otimes_R^\mathbf{L} R') = A \otimes_R R'$. Nous affirmons qu'il
n'existe aucun objet $E$ de $D(A \otimes_R R')$ dont l'image soit
$A \otimes_R^\mathbf{L} R'$ dans $D(R')$. En effet, pour un tel $E$, le module
$H^0(E)$ serait libre ; par conséquent, $E$ se décomposerait sous la forme
$H^0(E)[0] \oplus H^{-1}(E)[1]$. Or il est facile de voir que
$A \otimes_R^\mathbf{L} R'$ n'est pas isomorphe à la somme de ses
groupes de cohomologie dans $D(R')$.

\begin{lemma}
\label{examples-lemma-no-derived-base-change}
Soient $R \to R'$ et $R \to A$ des morphismes d'anneaux. En général, il
n'existe pas de foncteur $T : D(A) \to D(A \otimes_R R')$
entre catégories triangulées tel qu'à tout $A$-module $M$ soit associé un
objet $T(M)$ de $D(A \otimes_R R')$ qui s'envoie sur
$M \otimes_R^\mathbf{L} R'$ par l'application $D(A \otimes_R R') \to D(R')$.
\end{lemma}

\begin{proof}
Voir la discussion ci-dessus.
\end{proof}




\section{Un objet compact intéressant}
\label{examples-section-interesting-compact}


\noindent
Soit $R$ un anneau et soit $(A, \text{d})$ une $R$-algèbre différentielle graduée.
Si $A = R$, nous savons que tout objet compact de $D(A, \text{d}) = D(R)$
est représenté par un complexe fini de modules projectifs de type fini. En
d'autres termes, les objets compacts sont parfaits ; voir
Compléments d'algèbre, proposition \ref{more-algebra-proposition-perfect-is-compact}.
Dans le langage des modules différentiels gradués, la question analogue
serait : « Tout objet compact de $D(A, \text{d})$ est-il représenté
par un $A$-module différentiel gradué $P$, de type fini et
projectif comme module gradué ? »

\medskip\noindent
Pour une algèbre différentielle graduée générale, ce n'est pas vrai. En effet,
soit $k$ un corps de caractéristique $2$
(de sorte que nous n'ayons pas à nous soucier des signes).
Soit $A = k[x, y]/(y^2)$
avec
\begin{enumerate}
\item $x$ est de degré $0$,
\item $y$ est de degré $-1$,
\item $\text{d}(x) = 0$, et
\item $\text{d}(y) = x^2 + x$.
\end{enumerate}
Alors $x : A \to A$ est un projecteur dans $K(A, \text{d})$.
Par conséquent, on a
$$
A = \Ker(x) \oplus \Im(1 - x)
$$
dans $K(A, \text{d})$ ; voir
Algèbre différentielle graduée, lemme \ref{dga-lemma-homotopy-direct-sums}, et
Catégories dérivées, lemme
\ref{derived-lemma-projectors-have-images-triangulated}.
Il est clair que $A$ est un objet compact de $D(A, \text{d})$.
Il s'ensuit que $\Ker(x)$ est un objet compact de $D(A, \text{d})$,
comme il résulte de
Catégories dérivées, lemme \ref{derived-lemma-compact-objects-subcategory}.

\medskip\noindent
Supposons ensuite que $M$ soit un $A$-module différentiel gradué (à droite)
représentant $\Ker(x)$ et que $M$ soit de type fini et projectif
comme $A$-module gradué. Tout module gradué projectif de type fini
sur $k[x, y]/(y^2)$ étant gradué libre, nous voyons que $M$ est
libre de type fini comme $k[x, y]/(y^2)$-module gradué (autrement dit,
lorsque l'on oublie la différentielle). Posons $N = M/M(x^2 + x)$.
Considérons la suite exacte
$$
0 \to M \xrightarrow{x^2 + x} M \to N \to 0
$$
Comme $x^2 + x$ est de degré $0$, appartient au centre de $A$ et
$\text{d}(x^2 + x) = 0$, c'est une suite exacte courte de
$A$-modules différentiels gradués. De plus, puisque $\text{d}(y) = x^2 + x$,
la différentielle de $N$ est linéaire. Les applications
$$
H^{-1}(N) \to H^0(M)
\quad\text{et}\quad
H^0(M) \to H^0(N)
$$
sont des isomorphismes, car $H^*(M) = H^0(M) = k$ puisque $M \cong \Ker(x)$
dans $D(A, \text{d})$. Un calcul de l'application de connexion montre que
$H^*(N) = k[x, y]/(x, y^2)$ comme module gradué ; nous omettons
les détails. Comme $N$ est un $k[x, y]/(y^2, x^2 + x)$-module libre,
nous disposons d'une résolution
$$
\ldots \to N[2] \xrightarrow{y} N[1] \xrightarrow{y} N \to N/Ny \to 0
$$
compatible aux différentielles. Comme $N$ est borné et que
$H^0(N) = k[x,y]/(x, y^2)$, il résulte de
Homologie, lemme \ref{homology-lemma-first-quadrant-ss},
que $H^0(N/Ny) = k[x]/(x)$. Mais $N/Ny$ est un complexe fini de
$k[x]/(x^2 + x) = k \times k$-modules libres ; sa cohomologie
doit donc être de dimension paire, ce qui est une contradiction.

\begin{lemma}
\label{examples-lemma-no-good-representatif-compact-object}
Il existe une algèbre différentielle graduée $(A, \text{d})$ et un objet
compact $E$ de $D(A, \text{d})$ tel que $E$ ne puisse pas
être représenté par un $A$-module différentiel gradué de type fini et
projectif comme module gradué.
\end{lemma}

\begin{proof}
Voir la discussion ci-dessus.
\end{proof}







\section{Deux catégories différentielles graduées}
\label{examples-section-nongraded-differential-graded}

\noindent
Dans cette section, nous construisons deux catégories différentielles graduées
satisfaisant aux axiomes (A), (B) et (C) de
Algèbre différentielle graduée, situation \ref{dga-situation-ABC},
mais dont les objets ne sont pas munis d'une $\mathbf{Z}$-graduation.

\medskip\noindent
{\bf Exemple I.} Soit $X$ un espace topologique. Notons
$\underline{\mathbf{Z}}$ le faisceau constant de valeur $\mathbf{Z}$.
Soit $A$ un $\underline{\mathbf{Z}}$-torseur. Dans ce contexte, un faisceau
de groupes abéliens $\mathcal{F}$ est dit {\it $A$-gradué} si, pour toute
section locale $a \in A(U)$, il existe un projecteur
$p_a : \mathcal{F}|_U \to \mathcal{F}|_U$ tel que, pour tout isomorphisme
local $\underline{\mathbf{Z}}|_U \to A|_U$, on ait
$\mathcal{F}|_U = \bigoplus_{n \in \mathbf{Z}} p_n(\mathcal{F})$.
Autrement dit, localement sur $X$, le faisceau abélien $\mathcal{F}$ possède
une $\mathbf{Z}$-graduation, mais, sur les intersections, les différents choix
de graduation diffèrent d'un décalage de degré donné par les fonctions de
transition du torseur $A$.
Nous disons qu'un couple $(\mathcal{F}, \text{d})$ est un
{\it complexe $A$-gradué de faisceaux abéliens} si $\mathcal{F}$
est un faisceau abélien $A$-gradué et si $\text{d} : \mathcal{F} \to \mathcal{F}$ est une différentielle,
c'est-à-dire si $\text{d}^2 = 0$ et si
$p_{a + 1} \circ \text{d} = \text{d} \circ p_a$ pour toute section locale
$a$ de $A$. Autrement dit, $\text{d}(p_a(\mathcal{F}))$
est contenu dans $p_{a + 1}(\mathcal{F})$.

\medskip\noindent
Considérons ensuite la catégorie $\mathcal{A}$ définie comme suit :
\begin{enumerate}
\item ses objets sont les complexes $A$-gradués de faisceaux abéliens ;
\item pour deux objets $(\mathcal{F}, \text{d})$, $(\mathcal{G}, \text{d})$, nous posons
$$
\Hom_\mathcal{A}((\mathcal{F}, \text{d}), (\mathcal{G}, \text{d})) =
\bigoplus \Hom^n(\mathcal{F}, \mathcal{G})
$$
où $\Hom^n(\mathcal{F}, \mathcal{G})$ est le groupe des
morphismes de faisceaux abéliens $f$ tels que
$f(p_a(\mathcal{F})) \subset p_{a + n}(\mathcal{G})$ pour toute
section locale $a$ de $A$. Nous prenons pour différentielle
$\text{d}(f) = \text{d} \circ f - (-1)^n f \circ \text{d}$ ; voir
Algèbre différentielle graduée, exemple \ref{dga-example-category-complexes}.
\end{enumerate}
Nous omettons la vérification qu'il s'agit bien d'une catégorie différentielle
graduée satisfaisant à (A), (B) et (C). Toutes ces propriétés se vérifient
localement sur $X$, où l'on retrouve simplement la catégorie différentielle
graduée des complexes de faisceaux abéliens. Nous obtenons ainsi une catégorie
triangulée $K(\mathcal{A})$.

\medskip\noindent
Catégorie dérivée tordue de $X$. Remarquons qu'à tout objet
$(\mathcal{F}, \text{d})$ de $\mathcal{A}$ est associé un faisceau de
cohomologie $A$-gradué bien défini $H(\mathcal{F}, \text{d})$.
La notion de quasi-isomorphisme dans $K(\mathcal{A})$ est donc claire.
En inversant les quasi-isomorphismes, nous obtenons la
{\it catégorie dérivée $D(\mathcal{A})$ des complexes de faisceaux
$A$-gradués}. Si $A$ est le torseur trivial, alors $D(\mathcal{A})$
est égale à $D(X)$ ; pour un torseur non trivial, on obtient en revanche
une sorte de catégorie dérivée {\it tordue} de $X$.

\medskip\noindent
{\bf Exemple II.} Soit $C$ une courbe lisse sur un corps parfait $k$ de
caractéristique $2$. Alors $\Omega_{C/k}$ possède une racine carrée canonique.
Plus précisément, on peut écrire $\Omega_{C/k} = \mathcal{L}^{\otimes 2}$
de telle sorte que, pour toute fonction locale $f$ sur $C$, la section
$\text{d}(f)$ soit égale à $s^{\otimes 2}$ pour une certaine section locale
$s$ de $\mathcal{L}$. La « raison » est que
$$
\text{d}(a_0 + a_1t + \ldots +a_dt^d) =
(\sum\nolimits_{i\text{ impair}} a_i^{1/2} t^{(i - 1)/2})^2\text{d}t
$$
(insérer ici une référence future). Cela détermine en particulier une
connexion canonique
$$
\nabla_{can} :
\Omega_{C/k}
\longrightarrow
\Omega_{C/k} \otimes_{\mathcal{O}_C} \Omega_{C/k}
$$
dont la $2$-courbure est nulle (à savoir l'unique connexion pour laquelle
les carrés ont une dérivée nulle).
Remarquons que la catégorie des fibrés vectoriels
munis d'une connexion est une catégorie tensorielle ; nous obtenons donc aussi
des connexions canoniques
$\nabla_{can}$ sur les faisceaux inversibles $\Omega_{C/k}^{\otimes n}$ pour tous $n \in \mathbf{Z}$.

\medskip\noindent
Soit $\mathcal{A}$ la catégorie définie comme suit :
\begin{enumerate}
\item ses objets sont les couples $(\mathcal{F}, \nabla)$ formés d'un
faisceau localement libre de type fini $\mathcal{F}$ muni d'une connexion
$$
\nabla :
\mathcal{F}
\longrightarrow
\mathcal{F} \otimes_{\mathcal{O}_C} \Omega_{C/k}
$$
dont la $2$-courbure est nulle ;
\item les morphismes entre $(\mathcal{F}, \nabla_\mathcal{F})$ et
$(\mathcal{G}, \nabla_\mathcal{G})$
sont donnés par
$$
\Hom_\mathcal{A}((\mathcal{F}, \nabla_\mathcal{F}),
(\mathcal{G}, \nabla_\mathcal{G})) =
\bigoplus \Hom_{\mathcal{O}_C}(\mathcal{F},
\mathcal{G} \otimes_{\mathcal{O}_C} \Omega_{C/k}^{\otimes n})
$$
Pour un élément
$f : \mathcal{F} \to \mathcal{G} \otimes \Omega_{C/k}^{\otimes n}$
de degré $n$, nous posons
$$
\text{d}(f) =
\nabla_{\mathcal{G} \otimes \Omega_{C/k}^{\otimes n}} \circ f +
f \circ \nabla_\mathcal{F}
$$
moyennant les identifications convenables.
\end{enumerate}
Nous omettons la vérification que ceci définit une catégorie différentielle
graduée satisfaisant aux propriétés (A), (B) et (C). Nous obtenons ainsi une
catégorie d'homotopie triangulée $K(\mathcal{A})$.

\medskip\noindent
Si $C = \mathbf{P}^1_k$, alors $K(\mathcal{A})$ est la catégorie nulle.
En revanche, si $C$ est une courbe lisse propre de genre $> 1$, alors
$K(\mathcal{A})$ n'est pas nulle. En effet, supposons que $\mathcal{N}$
soit un faisceau inversible
de degré $0 \leq d < g - 1$
possédant une section non nulle $\sigma$. Posons alors
$(\mathcal{F}, \nabla_\mathcal{F}) = (\mathcal{O}_C, \text{d})$
et
$(\mathcal{G}, \nabla_\mathcal{G}) = (\mathcal{N}^{\otimes 2}, \nabla_{can})$.
On voit que
$$
\Hom_\mathcal{A}^n((\mathcal{F}, \nabla_\mathcal{F}),
(\mathcal{G}, \nabla_\mathcal{G})) =
\left\{
\begin{matrix}
0 & \text{si} & n < 0 \\
\Gamma(C, \mathcal{N}^{\otimes 2}) & \text{si} & n = 0 \\
\Gamma(C, \mathcal{N}^{\otimes 2} \otimes \Omega_{C/k}) & \text{si} & n = 1
\end{matrix}
\right.
$$
Le premier $0$ apparaît
parce que le degré de
$\mathcal{N}^{\otimes 2} \otimes \Omega_{C/k}^{\otimes -1}$
est négatif, par l'hypothèse $d < g - 1$. Or la dérivée de la section
$\sigma^{\otimes 2}$ est nulle ; par conséquent, le groupe d'homomorphismes
$$
\Hom_{K(\mathcal{A})}((\mathcal{F}, \nabla_\mathcal{F}),
(\mathcal{G}, \nabla_\mathcal{G}))
$$
est non nul.






\section{Le champ des espaces algébriques propres n'est pas algébrique}
\sectionmark{Le champ des espaces propres n'est pas algébrique}
\label{examples-section-proper-spaces-not-algebraic}

\noindent
Dans Quot, section \ref{quot-section-stack-of-spaces}, nous avons
introduit et étudié le champ en groupoïdes
$$
p'_{fp, flat, proper} :
\Spacesstack'_{fp, flat, proper}
\longrightarrow
\Sch_{fppf}
$$
le champ dont la catégorie des sections sur un schéma $S$
est la catégorie des espaces algébriques plats, propres et de présentation
finie sur $S$. Nous avons démontré qu'il satisfait à de nombreux
axiomes d'Artin. Dans cette section, nous montrons que ce champ
n'est pas algébrique en établissant que l'effectivité formelle
échoue en général.

\medskip\noindent
L'exemple canonique repose sur le fait que l'espace de déformation universelle
d'une variété abélienne de dimension $g$ possède $g^2$ paramètres formels,
alors que toute déformation effective peut être définie sur un
anneau local complet de dimension $\leq g(g + 1)/2$.
Nous construirons notre exemple en explicitant une
déformation non effective convenable d'une surface K3.
Nous nous bornerons à esquisser l'argument, sans en donner tous les détails.

\medskip\noindent
Soit $k = \mathbf{C}$ le corps des nombres complexes.
Soit $X \subset \mathbf{P}^3_k$ une surface lisse de degré $4$
sur $k$. On a
$\omega_X \cong \Omega^2_{X/k} \cong \mathcal{O}_X$.
Enfin,
$\dim_k H^0(X, T_{X/k}) = 0$,
$\dim_k H^1(X, T_{X/k}) = 20$ et
$\dim_k H^2(X, T_{X/k}) = 0$.
Comme $L_{X/k} = \Omega_{X/k}$ puisque $X$ est lisse sur $k$,
comme $\Ext^i_{\mathcal{O}_X}(\Omega_{X/k}, \mathcal{O}_X) =
H^i(X, T_{X/k})$, et en vertu de Cotangent, lemme
\ref{cotangent-lemma-find-obstruction-ringed-topoi},
nous obtenons une déformation universelle de $X$
sur
$$
k[[x_1, \ldots, x_{20}]]
$$
Supposons que cette déformation universelle soit effective
(au sens de Axiomes d'Artin, section \ref{artin-section-formal-objects}).
Nous obtiendrions alors un morphisme plat et propre
$$
f : Y \longrightarrow \Spec(k[[x_1, \ldots, x_{20}]])
$$
où $Y$ est un espace algébrique qui réalise la déformation universelle.
Cela est impossible pour la raison suivante. Comme $Y$ est séparé,
nous pouvons trouver un sous-schéma ouvert affine $V \subset Y$. Puisque la
fibre spéciale $X$ de $Y$ est lisse, le morphisme $f$ est lisse. Ainsi $Y$
est régulier, puisqu'il est lisse sur une base régulière, et il s'ensuit
que le complémentaire $D$ de $V$ dans $Y$ est un diviseur de Cartier effectif.
Alors $\mathcal{O}_Y(D)$ est un élément non trivial de $\Pic(Y)$
(pour le démontrer, on montre que le complémentaire d'un ouvert affine non vide
d'un espace algébrique propre et lisse sur un corps donne toujours un élément
non trivial du groupe de Picard, puis on applique ce résultat à la fibre générique
de $f$). Enfin, afin d'obtenir une contradiction, montrons que $\Pic(Y) = 0$.
L'application $\Pic(Y) \to \Pic(X)$ est injective,
car $H^1(X, \mathcal{O}_X) = 0$ (toutes les déformations de
$\mathcal{O}_X$ à $Y \times \Spec(k[[x_i]]/\mathfrak m^n)$
sont donc triviales) et par le théorème d'existence de Grothendieck
(selon lequel deux modules cohérents induisant les mêmes faisceaux
sur les épaississements
sont isomorphes). Si $X$ est suffisamment générale, alors
$\Pic(X) = \mathbf{Z}$, engendré par $\mathcal{O}_X(1)$. Il suffit donc de montrer que
$\mathcal{O}_X(n)$, pour $n > 0$, ne se déforme pas au voisinage du
premier ordre\footnote{Cet argument fonctionne dès lors que l'application
$c_1 : \Pic(X) \to H^1(X, \Omega_{X/k})$ est injective, ce qui est vrai
pour tout corps $k$ de caractéristique zéro et toute hypersurface lisse
$X$ de degré $4$ dans $\mathbf{P}^3_k$.}.
Considérons le cup-produit
$$
H^1(X, \Omega_{X/k}) \times H^1(X, T_{X/k})
\longrightarrow H^2(X, \mathcal{O}_X)
$$
Il s'agit d'un accouplement non dégénéré par dualité cohérente.
Un calcul montre que la classe de Chern
$c_1(\mathcal{O}_X(n)) \in H^1(X, \Omega_{X/k})$
est non nulle en cohomologie de Hodge.
Il existe donc une déformation du premier ordre
dont le cup-produit avec $c_1(\mathcal{O}_X(n))$ est non nul.
Enfin, on montre que ce cup-produit est précisément la classe
d'obstruction au relèvement.

\begin{lemma}
\label{examples-lemma-proper-spaces-not-algebraic}
Le champ en groupoïdes
$$
p'_{fp, flat, proper} :
\Spacesstack'_{fp, flat, proper}
\longrightarrow
\Sch_{fppf}
$$
dont la catégorie des sections sur un schéma $S$ est la catégorie des
espaces algébriques plats, propres et de présentation finie sur $S$
(voir Quot, section \ref{quot-section-stack-of-spaces})
n'est pas un champ algébrique.
\end{lemma}

\begin{proof}
S'il s'agissait d'un champ algébrique, tout objet formel serait
effectif ; voir Axiomes d'Artin, lemme \ref{artin-lemma-effective}.
La discussion précédente montre que ce n'est pas le cas
après changement de base à $\Spec(\mathbf{C})$.
D'où la conclusion.
\end{proof}






\section{Un exemple de champ des morphismes non algébrique}
\label{examples-section-non-algebraic-hom-stack}

\noindent
Soient $\mathcal{Y}, \mathcal{Z}$ des champs algébriques sur un schéma $S$.
Le {\it champ des morphismes} $\underline{\Mor}_S(\mathcal{Y}, \mathcal{Z})$
est le champ en groupoïdes sur $S$ dont la catégorie des sections sur
un schéma $T$ est la catégorie
$$
\Mor_T(\mathcal{Y} \times_S T, \mathcal{Z} \times_S T)
$$
dont les objets sont les $1$-morphismes et les morphismes, les $2$-morphismes.
Nous omettons la preuve qu'il s'agit bien d'un champ en groupoïdes sur
$(\Sch/S)_{fppf}$ (insérer ici une référence future). Bien entendu,
le champ des morphismes n'est pas algébrique en général.
Dans cette section,
nous en donnons un exemple où $\mathcal{Y}$ est représentable par un schéma
propre et plat sur $S$ et où $\mathcal{Z}$ est lisse et propre sur $S$.

\medskip\noindent
Soit $k$ un corps algébriquement clos qui ne soit pas la clôture
algébrique d'un corps fini. Posons $S = \Spec(k[[t]])$
et $S_n = \Spec(k[t]/(t^n)) \subset S$. Soit $f : X \to S$ un morphisme
satisfaisant aux conditions suivantes :
\begin{enumerate}
\item $f$ est projectif et plat, et ses fibres sont des courbes
géométriquement connexes,
\item la fibre $X_0 = X \times_S S_0$ est une courbe nodale à composantes
irréductibles lisses, dont le graphe dual contient un lacet formé de
courbes rationnelles,
\item $X$ est un schéma régulier.
\end{enumerate}
Pour construire une telle surface $X$, nous pouvons prendre par exemple
$$
X\quad :\quad T_0T_1T_2 - t(T_0^3 + T_1^3 + T_2^3) = 0
$$
dans $\mathbf{P}^2_{k[[t]]}$. Soit $A_0$ une variété abélienne non nulle sur $k$,
par exemple une courbe elliptique. Soit $A = A_0 \times_{\Spec(k)} S$ le
schéma abélien constant\footnote{Un schéma abélien sur une base $S$ est
un schéma en groupes plat, propre et de présentation finie sur $S$, dont
toutes les fibres sont des variétés abéliennes.}
sur $S$ associé à $A_0$. Nous montrerons que le
champ $\mathcal{X} = \underline{\Mor}_S(X, [S/A])$ n'est pas algébrique.

\medskip\noindent
Rappelons que $[S/A]$ est, d'une part, le champ quotient de l'action
triviale de $A$ sur $S$ et,
d'autre part, le champ classifiant les $A$-torseurs fppf ; voir
Exemples de champs, proposition
\ref{examples-stacks-proposition-equal-quotient-stacks}.
Observons que $[S/A] = [\Spec(k)/A_0] \times_{\Spec(k)} S$. Cela nous
permet de décrire comme suit la catégorie fibre au-dessus d'un schéma $T$ :
\begin{align*}
\mathcal{X}_T
& =
\underline{\Mor}_S(X, [S/A])_T \\
& =
\Mor_T(X \times_S T, [S/A] \times_S T) \\
& =
\Mor_S(X \times_S T, [S/A]) \\
& =
\Mor_{\Spec(k)}(X \times_S T, [\Spec(k)/A_0])
\end{align*}
pour tout $S$-schéma $T$. Autrement dit, le groupoïde $\mathcal{X}_T$
est le groupoïde des $A_0$-torseurs fppf sur $X \times_S T$.
Avant d'expliquer pourquoi $\mathcal{X}$ n'est pas un champ algébrique,
il nous faut quelques lemmes.

\begin{lemma}
\label{examples-lemma-torsors-over-two-dimensional-regular}
Soit $W$ un schéma noethérien intègre régulier de dimension deux,
de corps des fonctions $K$. Soit $G \to W$ un schéma abélien.
Alors l'application
$H^1_{fppf}(W, G) \to H^1_{fppf}(\Spec(K), G)$ est injective.
\end{lemma}

\begin{proof}[Esquisse de la démonstration]
Soit $P \to W$ un $G$-torseur fppf, trivial au point générique.
Nous disposons alors d'un morphisme $\Spec(K) \to P$ au-dessus de $W$,
et pouvons considérer son image schématique $Z \subset P$. Comme $P \to W$
est propre (puisqu'il s'agit d'un torseur sous un espace algébrique en
groupes propre sur $W$), le morphisme $Z \to W$ est propre et birationnel.
D'après Espaces sur un corps, lemme \ref{spaces-over-fields-lemma-finite-in-codim-1}, le morphisme $Z \to W$
est fini en dehors d'un nombre fini de points fermés de $W$. D'après
(insérer ici une référence future sur la résolution des indéterminations
des morphismes de surfaces par transformations quadratiques),
les composantes irréductibles des fibres géométriques de $Z \to W$
sont des courbes rationnelles. Par
Compléments sur les groupoïdes en espaces, lemme
\ref{spaces-more-groupoids-lemma-no-nonconstant-morphism-from-P1-to-group}
il n'existe aucun morphisme non constant d'une courbe rationnelle
vers un schéma en groupes, ni vers un torseur sous un tel schéma.
Ainsi $Z \to W$ est fini ; par conséquent, $Z$ est un schéma et
$Z \to W$ est un isomorphisme d'après
Morphismes, lemme \ref{morphisms-lemma-finite-birational-over-normal}.
Autrement dit, le torseur $P$ est trivial.
\end{proof}

\begin{remark}
\label{examples-remark-torsors-over-regular}
On peut supprimer la condition de dimension $2$ imposée à $W$ dans le
lemme \ref{examples-lemma-torsors-over-two-dimensional-regular}. En effet, si $W$
est noethérien, régulier et intègre, alors
(1) $P(W)\to P(K)$ est bijective pour tout $G$-torseur $P$, et
(2) $H^1(W,G)\to H^1(K,G)$ est injective.
Voici une esquisse de la démonstration.
D'abord, (2) découle clairement de (1). Pour (1), l'injectivité est immédiate,
et la surjectivité revient (comme dans la preuve précédente) à montrer que
l'adhérence schématique d'une section rationnelle est une section.
Cette question est locale pour la topologie fppf sur $W$ ; après un changement
de base \'etale, nous pouvons donc supposer $P=G$ (les recouvrements \'etales
préservent la régularité). Observons ensuite que
$G=\underline{\mathrm{Pic}}^0_{G'/W}$, où $G'$ est le schéma abélien dual.
La conclusion découle du fait que $G'$ est un schéma régulier : les faisceaux
inversibles définis génériquement se prolongent. En fait, on peut remplacer
l'hypothèse « $W$ est régulier » par « les $W$-schémas lisses sont localement
factoriels ». Laurent Moret-Bailly tient cette astuce de Picard de Raynaud ;
il doit exister une référence quelque part dans les travaux de Raynaud.
\end{remark}

\begin{lemma}
\label{examples-lemma-torsors-over-field-torsion}
Soit $G$ un espace algébrique en groupes lisse et commutatif
sur un corps $K$. Alors $H^1_{fppf}(\Spec(K), G)$ est de torsion.
\end{lemma}

\begin{proof}
Tout $G$-torseur $P$ sur $\Spec(K)$ est lisse sur $K$, puisqu'il est une forme
de $G$. Par conséquent, $P$ possède un point sur une extension finie
séparable $L/K$. Posons $[L : K] = n$. Notons $[n](P)$ le $G$-torseur
dont la classe est $n$ fois celle de $P$ dans
$H^1_{fppf}(\Spec(K), G)$. Il existe un morphisme canonique
$$
P \times_{\Spec(K)} \ldots \times_{\Spec(K)} P \to [n](P)
$$
d'espaces algébriques sur $K$. Ce morphisme est symétrique puisque
$G$ est commutatif. Il se factorise donc par le quotient
$$
(P \times_{\Spec(K)} \ldots \times_{\Spec(K)} P)/S_n
$$
D'autre part, le morphisme $\Spec(L) \to P$ définit un morphisme
$$
(\Spec(L) \times_{\Spec(K)} \ldots \times_{\Spec(K)} \Spec(L))/S_n
\longrightarrow (P \times_{\Spec(K)} \ldots \times_{\Spec(K)} P)/S_n
$$
et le lecteur vérifiera que le schéma de gauche possède un point $K$-rationnel.
Il s'ensuit que $[n](P)$ est le torseur trivial.
\end{proof}

\noindent
Pour montrer que $\mathcal{X} = \underline{\Mor}_S(X, [S/A])$
n'est pas un champ algébrique, il suffit, d'après
Axiomes d'Artin, lemme \ref{artin-lemma-effective},
d'établir le résultat suivant.

\begin{lemma}
\label{examples-lemma-not-essentially-surjective}
L'application canonique $\mathcal{X}(S) \to \lim \mathcal{X}(S_n)$ n'est pas
essentiellement surjective.
\end{lemma}

\begin{proof}[Esquisse de la démonstration]
En déroulant les définitions, il suffit de vérifier que
$H^1(X, A_0) \to \lim H^1(X_n, A_0)$ n'est pas surjective.
Comme $X$ est régulier et projectif, les
lemmes \ref{examples-lemma-torsors-over-field-torsion} et
\ref{examples-lemma-torsors-over-two-dimensional-regular}
montrent que tout $A_0$-torseur sur $X$ est
de torsion. En particulier, le groupe $H^1(X, A_0)$ est de torsion.
Il suffit donc de montrer :
(a) le groupe $H^1(X_0, A_0)$ n'est pas de torsion, et
(b) les applications $H^1(X_{n + 1}, A_0) \to H^1(X_n, A_0)$ sont surjectives pour tout $n$.

\medskip\noindent
Pour (a), on construit un $A_0$-torseur $P_0$ sur la courbe nodale $X_0$
qui n'est pas de torsion, en recollant les $A_0$-torseurs triviaux sur chaque
composante de $X_0$ au moyen de points d'ordre infini de $A_0$ comme
isomorphismes au-dessus des nœuds. Plus précisément, soit $x \in X_0$ un nœud
appartenant à un lacet formé de courbes rationnelles.
Soit $X'_0 \to X_0$ la normalisation de $X_0$ dans $X_0 \setminus \{x\}$.
Soient $x', x'' \in X'_0$ les deux points au-dessus de $x$.
Nous prenons $A_0 \times_{\Spec(k)} X'_0$ et identifions
$A_0 \times \{x'\}$ à $A_0 \times \{x''\}$ par la translation
$A_0 \to A_0$ définie par un point d'ordre infini $a_0 \in A_0(k)$
(un tel point existe puisque $k$ est algébriquement clos et n'est pas la
clôture algébrique d'un corps fini --- ce fait n'est pas trivial à démontrer).
On peut montrer que le recollement est un espace algébrique (et même
un schéma) et qu'il s'agit d'un $A_0$-torseur qui n'est pas de torsion sur $X_0$.
En effet, si $[n](P_0)$ possédait une section, celle-ci fournirait un morphisme
$s : X'_0 \to A_0$ tel que
$[n](a_0) = s(x') - s(x'')$ pour la loi de groupe de $A_0(k)$. Or,
les composantes irréductibles du lacet étant rationnelles, la section $s$
y est constante (
Compléments sur les groupoïdes en espaces, lemme
\ref{spaces-more-groupoids-lemma-no-nonconstant-morphism-from-P1-to-group}).
Ainsi $s(x') = s(x'')$, ce qui donne une contradiction.

\medskip\noindent
Pour (b), la théorie des déformations montre que l'obstruction à déformer un $A_0$-torseur $P_n \to X_n$ en un $A_0$-torseur
$P_{n + 1} \to X_{n + 1}$ appartient à $H^2(X_0, \omega)$
pour un fibré vectoriel convenable $\omega$ sur $X_0$.
Ce groupe s'annule puisque $X_0$ est une courbe, ce qui prouve l'assertion.
\end{proof}

\begin{proposition}
\label{examples-proposition-nonalghomstack}
Le champ $\mathcal{X} = \underline{\Mor}_S(X, [S/A])$ n'est pas algébrique.
\end{proposition}

\begin{proof}
Voir la discussion ci-dessus.
\end{proof}

\begin{remark}
\label{examples-remark-contradict-aoki}
La proposition \ref{examples-proposition-nonalghomstack} contredit
\cite[théorème 1.1]{AokiHomStacks}. La difficulté vient de la non-effectivité
des objets formels de $\underline{\Mor}_S(X, [S/A])$.
Cette même difficulté est mentionnée dans l'erratum
\cite{AokiHomStacksErr} à \cite{AokiHomStacks}. Malheureusement,
l'erratum affirme ensuite que $\underline{\Mor}_S(\mathcal{Y}, \mathcal{Z})$
est algébrique lorsque $\mathcal{Z}$ est séparé, ce qui contredit également
la proposition \ref{examples-proposition-nonalghomstack}, puisque $[S/A]$ est séparé.
\end{remark}








\section[Champ algébrique sans effectivité formelle forte]{Un champ algébrique ne satisfaisant pas à l'effectivité formelle forte}
\label{examples-section-non-formal-effectiveness}

\noindent
Il s'agit de \cite[exemple 4.12]{Bhatt-Algebraize}.
Soit $k$ un corps algébriquement clos.
Soit $A$ une variété abélienne sur $k$.
Supposons que $A(k)$ ne soit pas de torsion (c'est toujours le cas si $k$
n'est pas la clôture algébrique d'un corps fini).
Soit $\mathcal{X} = [\Spec(k)/A]$.
Nous affirmons qu'il existe un idéal $I \subset k[x, y]$
tel que
$$
\mathcal{X}_{\Spec(k[x, y]^\wedge)}
\longrightarrow
\lim \mathcal{X}_{\Spec(k[x, y]/I^n)}
$$
ne soit pas essentiellement surjectif. En effet, soit $I$
l'idéal engendré par $xy(x + y - 1)$.
Alors $X_0 = V(I)$ est formé de trois copies de $\mathbf{A}^1_k$
recollées en triangle en trois points. Nous pouvons donc construire un torseur
$P_0$ sous $A$ sur $X_0$, d'ordre infini, en prenant le torseur trivial
sur les composantes irréductibles de $X_0$ et en les recollant
par des translations définies par des points d'ordre infini.
Exactement comme dans la preuve du lemme \ref{examples-lemma-not-essentially-surjective},
nous pouvons relever $P_0$ en un torseur $P_n$ sur $X_n = \Spec(k[x, y]/I^n)$.
Comme $k[x, y]^\wedge$ est un domaine régulier de dimension deux,
tout torseur $P$ sous $A$ sur $\Spec(k[x, y]^\wedge)$
est de torsion (lemmes \ref{examples-lemma-torsors-over-two-dimensional-regular}
et \ref{examples-lemma-torsors-over-field-torsion}). Le système de
torseurs n'appartient donc pas à l'image du foncteur affiché.

\begin{lemma}
\label{examples-lemma-non-formal-effectiveness}
Soit $k$ un corps algébriquement clos qui ne soit pas la clôture
algébrique d'un corps fini. Soit $A$ une variété abélienne sur $k$.
Soit $\mathcal{X} = [\Spec(k)/A]$.
Il existe un système projectif de $k$-algèbres $R_n$,
à morphismes de transition surjectifs dont les noyaux sont localement nilpotents,
et un système $(\xi_n)$ d'objets de $\mathcal{X}$ au-dessus du système
$(\Spec(R_n))$, tel que ce système ne soit pas effectif
au sens de Axiomes d'Artin, remarque \ref{artin-remark-strong-effectiveness}.
\end{lemma}

\begin{proof}
Voir la discussion ci-dessus.
\end{proof}








\section{Un contre-exemple au théorème d'existence de Grothendieck}
\label{examples-section-Grothendieck-existence}

\noindent
Soit $k$ un corps et soit $A = k[[t]]$. Soit $X$ le recollement de
$U = \Spec(A[x])$ et $V = \Spec(A[y])$ au moyen de l'identification
$$
U \setminus \{0_U\} \longrightarrow V \setminus \{0_V\}
$$
qui envoie $x$ sur $y$, où $0_U \in U$ et $0_V \in V$ sont les points
correspondant aux idéaux maximaux $(x, t)$ et $(y, t)$. Posons
$A_n = A/(t^n)$ et $X_n = X \times_{\Spec(A)} \Spec(A_n)$.
Soit $\mathcal{F}_n$ le faisceau cohérent sur $X_n$
correspondant au $A_n[x]$-module $A_n[x]/(x) \cong A_n$
et au $A_n[y]$-module $0$, avec le recollement évident.
Soit $\mathcal{I} \subset \mathcal{O}_X$ le faisceau
d'idéaux engendré par $t$. Alors $(\mathcal{F}_n)$ est un objet de
la catégorie $\textit{Coh}_{\text{support propre sur } A}(X, \mathcal{I})$
définie dans
Cohomologie des schémas, section \ref{coherent-section-existence-proper-support}.
D'autre part, cet objet n'appartient pas à l'image du
foncteur
de Cohomologie des schémas,
équation (\ref{coherent-equation-completion-functor-proper-over-A}).
En effet, s'il y appartenait, il existerait un $A[x]$-module $M$ de type fini,
un $A[y]$-module $N$ de type fini et un isomorphisme $M[1/t] \cong N[1/t]$
tels que $M/t^nM \cong A_n[x]/(x)$ et $N/t^nN = 0$ pour tout $n$.
Il est facile de voir que c'est impossible.

\begin{lemma}
\label{examples-lemma-counter-Grothendieck-existence}
Contre-exemples à l'algébrisation des faisceaux cohérents.
\begin{enumerate}
\item Le théorème d'existence de Grothendieck, tel qu'il est énoncé dans
Cohomologie des schémas, théorème \ref{coherent-theorem-grothendieck-existence},
est faux si l'on omet l'hypothèse que $X \to \Spec(A)$ est séparé.
\item Le champ des faisceaux cohérents $\Cohstack_{X/B}$ des
théorèmes \ref{quot-theorem-coherent-algebraic-general} et
\ref{quot-theorem-coherent-algebraic} de Quot n'est, en général,
pas algébrique si l'on omet l'hypothèse que $X \to S$ est séparé.
\item Le foncteur $\Quotfunctor_{\mathcal{F}/X/B}$ de
Quot, proposition \ref{quot-proposition-quot},
n'est pas un espace algébrique en général si l'on omet l'hypothèse
que $X \to B$ est séparé.
\end{enumerate}
\end{lemma}

\begin{proof}
Nous avons vu (1) ci-dessus. Cela montre que $\textit{Coh}_{X/A}$
ne satisfait pas à l'axiome [4] de Axiomes d'Artin,
section \ref{artin-section-axioms}. Il ne peut donc pas être un champ algébrique, d'après Axiomes d'Artin, lemme
\ref{artin-lemma-effective}. On obtient ainsi (2).
Pour (3), remarquons qu'il existe des surjections compatibles
$\mathcal{O}_{X_n} \to \mathcal{F}_n$ pour tout $n$.
Ainsi $\Quotfunctor_{\mathcal{O}_X/X/A}$ ne satisfait pas à
l'axiome [4], et (3) se démontre comme précédemment.
\end{proof}






\section{Espaces algébriques formels affines}
\label{examples-section-affine-formal-algebraic-space}

\noindent
Soit $K$ un corps et soit $(V_i)_{i \in I}$ un système projectif filtrant
d'espaces vectoriels non nuls sur $K$, à morphismes de transition surjectifs,
tel que $\lim V_i = 0$ ; voir la section \ref{examples-section-zero-limit}.
Munissons $R_i = K \oplus V_i$ de la structure de $K$-algèbre pour laquelle
$V_i$ est un idéal de carré nul. Alors les $R_i$ forment un système projectif
de $K$-algèbres, à morphismes de transition surjectifs de noyaux nilpotents,
et $\lim R_i = K$. L'espace algébrique formel affine
$X = \colim \Spec(R_i)$ est donc un exemple d'espace
algébrique formel affine qui n'est pas McQuillan.

\begin{lemma}
\label{examples-lemma-affine-not-mcquillan}
Il existe un espace algébrique formel affine qui n'est pas McQuillan.
\end{lemma}

\begin{proof}
Voir la discussion ci-dessus.
\end{proof}

\noindent
Soit $0 \to W_i \to V_i \to K \to 0$ un système de suites exactes
comme dans la section \ref{examples-section-zero-limit}. Posons
$A_i = K[V_i]/(ww'; w, w' \in W_i)$.
Il existe alors un système compatible de surjections $A_i \to K[t]$
à noyaux nilpotents, et les morphismes
de transition $A_i \to A_j$ sont eux aussi surjectifs à noyaux nilpotents.
Rappelons que $V_i$ est libre sur $K$, de base les éléments $s \in S_i$.
Si la caractéristique de $K$ est nulle, la composante de degré $d$ de
$A_i$ est libre sur $K$, de base les $s^d$, $s \in S_i$, qui s'envoient
tous sur $t^d$. Le système projectif des composantes de degré $d$ des
$A_i$ est donc isomorphe au système projectif des espaces vectoriels $V_i$.
Comme $\lim V_i = 0$, nous concluons que $\lim A_i = K$, au moins lorsque
la caractéristique de $K$ est nulle.
Cela donne un exemple d'espace algébrique formel affine
dont les fonctions régulières ne séparent pas les points.

\begin{lemma}
\label{examples-lemma-affine-formal-functions-do-not-separate-points}
Il existe un espace algébrique formel affine $X$
dont les fonctions régulières ne séparent pas les points, au sens suivant :
si l'on écrit $X = \colim X_\lambda$ comme dans
Espaces formels, définition
\ref{formal-spaces-definition-affine-formal-algebraic-space},
alors $\lim \Gamma(X_\lambda, \mathcal{O}_{X_\lambda})$
est un corps, tandis que $X_{red}$ possède une infinité de points.
\end{lemma}

\begin{proof}
Voir la discussion ci-dessus.
\end{proof}

\noindent
Soient $K$, $I$ et $(V_i)$ comme ci-dessus. Considérons les systèmes
$$
\Phi = (\Lambda, J_i \subset \Lambda, (M_i) \to (V_i))
$$
où $\Lambda$ est une $K$-algèbre augmentée,
$J_i \subset \Lambda$ est, pour $i \in I$, un idéal de carré nul,
$(M_i) \to (V_i)$ est un morphisme de systèmes projectifs de
$K$-espaces vectoriels tel que $M_i \to V_i$ soit surjectif pour tout $i$,
chaque $M_i$ possède une structure de $\Lambda$-module, les morphismes
de transition $M_i \to M_j$, pour $i > j$, sont $\Lambda$-linéaires, et
$J_j M_i \subset \Ker(M_i \to M_j)$ pour $i > j$.
Affirmation : il existe un système comme ci-dessus tel que
$M_j = M_i/J_j M_i$ pour tout $i > j$.

\medskip\noindent
Si l'affirmation est vraie, nous obtenons un morphisme représentable
$$
\colim_{i \in I} \Spec(\Lambda/J_i \oplus M_i)
\longrightarrow
\text{Spf}(\lim \Lambda/J_i)
$$
d'espaces algébriques formels affines dont la source n'est pas McQuillan,
mais dont le but l'est. Ici, $\Lambda/J_i \oplus M_i$ est muni de sa
structure usuelle de $\Lambda/J_i$-algèbre, où $M_i$ est un idéal de carré nul.
La représentabilité équivaut exactement à la condition
$M_i/J_jM_i = M_j$ pour $i > j$. La source du morphisme n'est pas
McQuillan, car les projections
$\lim_{i \in I} M_i \to M_i$ ne sont pas surjectives. Cela résulte du fait
que les applications $\lim V_i \to V_i$ ne sont pas surjectives
et que nous disposons de la surjection $M_i \to V_i$. Nous omettons certains détails.

\medskip\noindent
Démonstration de l'affirmation. Remarquons d'abord qu'il existe au moins
un système, à savoir
$$
\Phi_0 = (K, J_i = (0), (V_i) \xrightarrow{\text{id}} (V_i))
$$
Étant donné un système $\Phi$, nous montrerons qu'il existe un morphisme
de systèmes $\Phi \to \Phi'$ (les morphismes de systèmes étant définis
de manière évidente) tel que $\Ker(M_i/J_j M_i \to M_j)$ s'envoie sur
zéro dans $M'_i/J'_j M'_i$. Cela fait, nous pouvons appliquer le procédé
usuel : poser par récurrence $\Phi_n = (\Phi_{n - 1})'$ pour $n \geq 1$,
puis prendre $\Phi = \colim \Phi_n$ afin d'obtenir un système possédant
les propriétés voulues. Nous omettons les détails.

\medskip\noindent
Construction de $\Phi'$ à partir de $\Phi$. Considérons l'ensemble $U$ des triplets $u = (i, j, \xi)$ tels que $i > j$ et
$\xi \in \Ker(M_i \to M_j)$. Notons $s, t : U \to I$ les applications
$s(i, j, \xi) = i$ et $t(i, j, \xi) = j$. Posons ensuite
$\xi_u \in M_{s(u)}$ égal à la troisième composante de $u$.
Prenons
$$
\Lambda' = \Lambda[x_u; u \in U]/(x_u x_{u'}; u, u' \in U)
$$
munie de l'augmentation $\Lambda' \to K$ induite par celle de $\Lambda$
et envoyant les $x_u$ sur zéro.
Posons $J'_k = J_k \Lambda' + (x_{u,\ t(u) \geq k})$, puis
$$
M'_i = M_i \oplus \bigoplus\nolimits_{s(u) \geq i} K\epsilon_{i, u}
$$
Comme morphismes de transition $M'_i \to M'_j$ pour $i > j$, nous prenons
le morphisme donné $M_i \to M_j$ et envoyons $\epsilon_{i, u}$ sur
$\epsilon_{j, u}$. Le morphisme $M'_i \to V_i$ induit le morphisme donné
$M_i \to V_i$ et envoie $\epsilon_{i, u}$ sur zéro.
Enfin, faisons agir $\Lambda'$ sur $M'_i$ comme suit :
un élément $\lambda \in \Lambda$ agit par la structure de $\Lambda$-module
sur $M_i$ et, par l'augmentation $\Lambda \to K$, sur les $\epsilon_{i, u}$.
L'élément $x_u$ agit par $0$ sur $M_i$ pour tout $i$.
Enfin, nous posons
$$
x_u \epsilon_{i, u} = \text{image de }\xi_u\text{ dans }M_i
$$
et tous les autres produits $x_{u'} \epsilon_{i, u}$ égaux à zéro.
La formule affichée a un sens puisque $s(u) \geq i$ et
$\xi_u \in M_{s(u)}$. Les principaux points à vérifier sont que
$J'_j M'_i \subset M'_i$ s'envoie sur zéro dans $M'_j$ pour $i > j$,
et que $\Ker(M_i \to M_j)$ s'envoie sur zéro dans
$M'_i/J_j M'_i$. Ce dernier fait vient de ce que
$\xi = x_{(i, j, \xi)} \epsilon_{i, (i, j, \xi)} \in J'_j M'_i$
pour tout $\xi \in \Ker(M_i \to M_j)$.
Nous omettons les détails.

\begin{lemma}
\label{examples-lemma-representable-morphism-affine-formal-not-mcquillan-top}
Il existe un morphisme représentable $f : X \to Y$ d'espaces
algébriques formels affines tel que $Y$ soit McQuillan, mais que $X$
ne le soit pas.
\end{lemma}

\begin{proof}
Voir la discussion ci-dessus.
\end{proof}









\section[Morphismes plats et colimites de présentation finie]{Les morphismes plats ne sont pas des colimites filtrantes de morphismes plats de présentation finie}
\label{examples-section-flat-not-colimit-flat-finitely-presented}

\noindent
Le but de cette section est de donner un exemple de morphisme plat d'anneaux
qui ne soit pas une colimite filtrante de morphismes plats d'anneaux de
présentation finie. Dans \cite{gabber-nonexcellent}, il est montré que si
$A$ est un anneau local non excellent, de dimension $1$ et de caractéristique
résiduelle nulle, alors le morphisme plat $A \to A^\wedge$ vers son complété
n'est pas une colimite filtrante de morphismes plats d'anneaux de type fini.
Dans notre exemple, l'anneau source sera excellent. Nous invitons le lecteur
à proposer d'autres exemples ;
si vous en connaissez un, écrivez à
\href{mailto:stacks.project@gmail.com}{stacks.project@gmail.com}.

\medskip\noindent
Pour la construction, fixons un nombre premier $p$ et posons
$A = \mathbf{F}_p[x_1, \ldots, x_n]$.
Choisissons une clôture intégrale absolue $A^+$ de $A$, c'est-à-dire que $A^+$ est la
clôture intégrale de $A$ dans une clôture algébrique de son corps des fractions.
Il est montré dans \cite[\S 6.7]{HHBigCM} que $A \to A^+$ est plat.

\medskip\noindent
Nous affirmons que la $A$-algèbre $A^+$ n'est pas une colimite filtrante de
$A$-algèbres plates de présentation finie lorsque $n \geq 3$.

\medskip\noindent
Nous esquissons l'argument pour $n = 3$ et laissons au lecteur la généralisation
aux valeurs supérieures de $n$. Il suffit de prouver l'énoncé analogue pour
le morphisme $R \to R^+$, où $R$ est l'hensélisé strict de $A$ à l'origine
et $R^+$ sa clôture intégrale absolue. Remarquons que $R$
est un anneau local régulier hensélien dont le corps résiduel $k$
est une clôture algébrique de $\mathbf{F}_p$.

\medskip\noindent
Choisissons une surface abélienne ordinaire $X$ sur $k$ et un fibré inversible très ample $L$ sur $X$. L'anneau des sections
$\Gamma_*(X, L) = \bigoplus_n H^0(X,L^n)$ est l'anneau de coordonnées du cône affine sur $X$ relativement à $L$. Il est normal dès que $L$
est suffisamment positif. Notons $S$ l'hensélisé de $\Gamma_*(X, L)$
au sommet du cône. Alors $S$ est un domaine local hensélien noethérien
normal de dimension $3$. Nous obtenons un morphisme fini injectif
$R \to S$ en hensélisant une normalisation de Noether de la $k$-algèbre de type fini $\Gamma_*(X, L)$. Puisque $R^+$ est une
clôture intégrale absolue de $R$, nous pouvons aussi fixer un plongement
$S \to R^+$. Ainsi $R^+$ est également la clôture intégrale absolue de $S$.
Pour montrer que $R^+$ n'est pas une colimite filtrante de $R$-algèbres plates, il suffit d'établir les deux assertions suivantes :
\begin{enumerate}
\item Si $S \to P$ est un morphisme d'anneaux et si $P$ est plate et
de type fini sur $R$, il existe un morphisme d'anneaux $S \to T$ avec
$T$ finie et plate sur $R$.
\item Pour tout morphisme fini d'anneaux $S \to T$, l'anneau $T$
n'est pas plat sur $R$.
\end{enumerate}
En effet, comme $S$ est de présentation finie sur $R$, si l'on pouvait écrire
$R^+ = \colim_i P_i$ comme une colimite filtrante de $R$-algèbres plates
de présentation finie $P_i$, le morphisme $S \to R^+$ se factoriserait
par $S \to P_i \to R^+$ pour $i \gg 0$, contrairement aux deux assertions
précédentes. L'assertion (1) découle de l'hensélianité de $R$ et d'un
argument de sectionnement ; voir Compléments sur les morphismes, lemme
\ref{more-morphisms-lemma-qf-fp-flat-neighbourhood-dominates-fppf}.
L'assertion (2) a été démontrée dans \cite{BhattSmallCMMod} ; pour la
commodité du lecteur, nous rappelons l'argument.

\medskip\noindent
Soit $U \subset \Spec(S)$ le spectre épointé ; nous avons des morphismes naturels $X \leftarrow U \subset \Spec(S)$. Le premier fournit une
identification $H^1(U, \mathcal{O}_U) \simeq H^1(X, \mathcal{O}_X)$.
En passant aux vecteurs de Witt du perfectionné et en utilisant la suite
d'Artin--Schreier\footnote{On utilise ici que $S$ est un anneau local
strictement hensélien de caractéristique $p$ et que, par conséquent,
$S \to S$, $f \mapsto f^p - f$, est surjective.
De plus, $S$ est un domaine normal, donc
$\Gamma(U, \mathcal{O}_U) = S$. Ainsi $H^1_\etale(U, \mathbf{Z}/p)$
est le noyau de l'application
$H^1(U, \mathcal{O}_U) \to H^1(U, \mathcal{O}_U)$
induite par $f \mapsto f^p - f$.}, on obtient une identification
$H^1_\etale(U, \mathbf{Z}_p) \simeq H^1_\etale(X, \mathbf{Z}_p)$.
En particulier, ce groupe est un $\mathbf{Z}_p$-module libre de type fini
de rang $2$ (puisque $X$ est ordinaire). Pour obtenir une contradiction, supposons qu'il existe une $R$-algèbre plate $T$ comme dans (2).
Notons $V \subset \Spec(T)$ l'image réciproque de $U$, et
$f : V \to U$ le morphisme fini surjectif induit. Comme $U$ est normal,
il existe une application trace $f_*\mathbf{Z}_p \to \mathbf{Z}_p$
sur $U_\etale$, dont la composée avec l'image inverse
$\mathbf{Z}_p \to f_*\mathbf{Z}_p$ est la multiplication par $d = \deg(f)$.
En passant à la cohomologie et en utilisant que
$H^1_\etale(U, \mathbf{Z}_p)$ n'est pas de torsion, on en déduit que
$H^1_\etale(V, \mathbf{Z}_p)$ est non nul. Comme
$H^1_\etale(V, \mathbf{Z}_p) \simeq \lim H^1_\etale(V, \mathbf{Z}/p^n)$
et qu'aucun terme $R^1\lim$ n'intervient, le groupe
$H^1(V_\etale,\mathbf{Z}/p)$ est nécessairement non nul.
Puisque $T$ est plate sur $R$, on a $\Gamma(V, \mathcal{O}_V) = T$ ; cet anneau est strictement hensélien,
et la suite d'Artin--Schreier montre que $H^1(V, \mathcal{O}_V) \neq 0$.
Cela équivaut à $H^2_\mathfrak m(T) \neq 0$, où
$\mathfrak m \subset R$ est l'idéal maximal. On obtient donc une contradiction,
car $T$ est finie et plate (c'est-à-dire libre de rang fini) comme
$R$-module, et $H^2_\mathfrak m(R) = 0$. Cette contradiction prouve (2).

\begin{lemma}
\label{examples-lemma-weird-flat-map}
Il existe un anneau commutatif $A$ et une $A$-algèbre plate $B$
qui ne peut s'écrire comme colimite filtrante de $A$-algèbres plates
de présentation finie. En fait, on peut choisir $A$ soit comme
$\mathbf{F}_p$-algèbre de type fini, soit comme anneau local noethérien
de dimension $1$, à corps résiduel de caractéristique $0$.
\end{lemma}

\begin{proof}
Voir la discussion ci-dessus.
\end{proof}




\section{La catégorie des modules modulo les modules de torsion}
\label{examples-section-serre-quotient-modulo-torsion-modules}

\noindent
La catégorie des groupes de torsion est une sous-catégorie de Serre
(Homologie, définition \ref{homology-definition-serre-subcategory})
de la catégorie de tous les groupes abéliens. Plus généralement, pour tout
anneau $A$, la catégorie des $A$-modules de torsion est une sous-catégorie
de Serre de la catégorie de tous les $A$-modules ; voir Compléments d'algèbre, section
\ref{more-algebra-section-abelian-categories-modules}.
Si $A$ est un domaine intègre, la catégorie quotient
(Homologie, lemme \ref{homology-lemma-serre-subcategory-is-kernel})
est équivalente à la catégorie des espaces vectoriels sur le corps des fractions.
Ce fait découle de la proposition plus générale suivante.

\begin{proposition}
\label{examples-proposition-localization-and-serre-quotients}
Soit $A$ un anneau et soit $S$ une partie multiplicative de $A$.
Notons $\text{Mod}_A$ la catégorie des $A$-modules et $\mathcal{T}$ sa sous-catégorie de Serre formée des modules dont tout élément est annulé
par un certain élément de $S$.
Il existe alors une équivalence canonique
$\text{Mod}_A/\mathcal{T} \rightarrow \text{Mod}_{S^{-1}A}$.
\end{proposition}

\begin{proof}
Le foncteur $\text{Mod}_A \to \text{Mod}_{S^{-1}A}$ défini par $M
\mapsto M \otimes_A S^{-1}A$ est exact (Algèbre, proposition
\ref{algebra-proposition-localization-exact})
et envoie les modules de $\mathcal{T}$ sur zéro.
La propriété universelle donnée par Homologie, lemme \ref{homology-lemma-serre-subcategory-is-kernel},
permet donc de le factoriser en un foncteur
$\text{Mod}_A/\mathcal{T} \to \text{Mod}_{S^{-1}A}$.

\medskip\noindent
Réciproquement, tout $A$-module $M$ tel que
$M \otimes_A S^{-1}A = 0$ appartient à $\mathcal{T}$, puisque
$M \otimes_A S^{-1}A \cong S^{-1} M$
(Algèbre, lemme \ref{algebra-lemma-tensor-localization}).
Homologie, lemme
\ref{homology-lemma-quotient-by-kernel-exact-functor}, montre donc que
le foncteur $\text{Mod}_A/\mathcal{T} \to \text{Mod}_{S^{-1}A}$ est fidèle.

\medskip\noindent
De plus, ce plongement est essentiellement surjectif : un antécédent d'un
$S^{-1}A$-module $N$ est $N_A$, c'est-à-dire $N$ considéré comme $A$-module,
car l'application canonique $N_A \otimes_A S^{-1}A \to N$, qui envoie
$x \otimes a/s$ sur $(a/s) \cdot x$, est un isomorphisme de $S^{-1}A$-modules.
\end{proof}

\begin{proposition}
\label{examples-proposition-quotient-by-torsion-modules}
Soit $A$ un anneau. Notons $Q(A)$ son anneau total des fractions
(comme dans Algèbre, exemple \ref{algebra-example-localize-at-prime}).
Notons $\text{Mod}_A$ la catégorie des $A$-modules et $\mathcal{T}$ sa
sous-catégorie de Serre des modules de torsion. Notons
$\text{Mod}_{Q(A)}$ la catégorie des $Q(A)$-modules.
Il existe alors une équivalence canonique
$\text{Mod}_A/\mathcal{T} \rightarrow \text{Mod}_{Q(A)}$.
\end{proposition}

\begin{proof}
Cela résulte immédiatement de la proposition \ref{examples-proposition-localization-and-serre-quotients}, appliquée à la partie
multiplicative
$S = \{f \in A \mid f \text{ n'est pas un diviseur de zéro dans }A\}$,
puisqu'un module est de torsion si et seulement si chacun de ses éléments
est annulé par un certain élément de $S$.
\end{proof}

\begin{proposition}
\label{examples-proposition-quotient-by-finitely-generated-torsion-modules}
Soit $A$ un domaine intègre noethérien et soit $K$ son corps des fractions.
Notons $\text{Mod}_A^{fg}$ la catégorie des $A$-modules de type fini
et $\mathcal{T}^{fg}$ sa sous-catégorie de Serre des modules de torsion
de type fini. Alors $\text{Mod}_A^{fg}/\mathcal{T}^{fg}$ est canoniquement
équivalente à la catégorie des $K$-espaces vectoriels de dimension finie.
\end{proposition}

\begin{proof}
L'équivalence de la proposition
\ref{examples-proposition-quotient-by-torsion-modules} se restreint, le long du
plongement $\text{Mod}_A^{fg}/\mathcal{T}^{fg} \to \text{Mod}_A/\mathcal{T}$,
en une équivalence
$\text{Mod}_A^{fg}/\mathcal{T}^{fg} \to \text{Vect}_K^{fd}$.
L'hypothèse noethérienne garantit que $\text{Mod}_A^{fg}$ est une
catégorie abélienne (voir Compléments d'algèbre, section
\ref{more-algebra-section-abelian-categories-modules}) et que le foncteur canonique
$\text{Mod}_A^{fg}/\mathcal{T}^{fg} \to \text{Mod}_A/\mathcal{T}$
est plein (sinon, les sous-modules de torsion des modules de type fini
pourraient ne pas être des objets de $\mathcal{T}^{fg}$).
\end{proof}

\begin{proposition}
\label{examples-proposition-quotient-abelian-groups-by-torsion-groups}
Le quotient de la catégorie des groupes abéliens par sa
sous-catégorie de Serre des groupes de torsion est la catégorie des
$\mathbf{Q}$-espaces vectoriels.
\end{proposition}

\begin{proof}
L'assertion découle directement de la
proposition \ref{examples-proposition-quotient-by-torsion-modules}.
\end{proof}




\section{Différentes topologies de colimite}
\label{examples-section-colimit-topology}

\noindent
Cet exemple est \cite[exemple 1.2, page 553]{TSH}. Soit
$G_n = \mathbf{Q} \times \mathbf{R}^n$, $n \geq 1$, considéré comme groupe topologique
additif muni de la topologie usuelle (euclidienne). Considérons les
plongements fermés $G_n \to G_{n + 1}$ qui envoient
$(x_0, \ldots, x_n)$ sur $(x_0, \ldots, x_n, 0)$. Nous affirmons que
$G = \colim G_n$ muni de la topologie
$$
U \subset G\text{ ouvert} \Leftrightarrow G_n \cap U\text{ ouvert }\forall n
$$
ne définit pas un groupe topologique.

\medskip\noindent
Pour le voir, considérons l'ensemble
$$
U = \{(x_0, x_1, x_2, \ldots)\text{ tels que }
|x_j| < |\cos(jx_0)| \text{ pour } j > 0\}
$$
Comme $jx_0$ n'est jamais un multiple entier de $\pi/2$, puisque $\pi$
n'est pas rationnel, il est facile de montrer que $U \cap G_n$ est ouvert.
Puisque $0 \in U$, si la topologie ci-dessus faisait de $G$ un groupe
topologique, il existerait un voisinage ouvert $V \subset G$ de $0$
tel que $V + V \subset U$. Pour tout $j \geq 0$, il existerait alors
$\epsilon_j > 0$ tel que $(0, \ldots, 0, x_j, 0, \ldots) \in V$
dès que $|x_j| < \epsilon_j$. Puisque $V + V \subset U$, on aurait
$$
(x_0, 0, \ldots, 0, x_j, 0, \ldots) \in U
$$
pour $|x_0| < \epsilon_0$ et $|x_j| < \epsilon_j$. Or, si l'on choisit
$j$ assez grand pour que $j \epsilon_0 > \pi/2$, on peut choisir
$x_0 \in \mathbf{Q}$ de telle sorte que $|\cos(jx_0)|$ soit inférieur à
$\epsilon_j$ ; il existe alors un $x_j$ tel que
$|\cos(jx_0)| < |x_j| < \epsilon_j$. Cette contradiction prouve l'affirmation.

\begin{lemma}
\label{examples-lemma-colimit-topology}
Il existe un système $G_1 \to G_2 \to G_3 \to \ldots$ de groupes
topologiques (abéliens) tel que $\colim G_n$, pris dans la catégorie des
espaces topologiques, diffère de $\colim G_n$, pris dans la catégorie
des groupes topologiques.
\end{lemma}

\begin{proof}
Voir la discussion ci-dessus.
\end{proof}






\section{Universellement submersif sans être un recouvrement V}
\label{examples-section-universally-submersive-not-V}

\noindent
Soit $A$ un anneau de valuation. Soit $\mathfrak p \subset A$ un idéal
premier qui n'est ni l'idéal premier minimal ni l'idéal maximal.
(Un cas utile à garder à l'esprit est celui où $A$ possède trois idéaux
premiers et où $\mathfrak p$ est celui « du milieu ».)
Considérons le morphisme de schémas affines
$$
\Spec(A_\mathfrak p) \amalg \Spec(A/\mathfrak p) \longrightarrow \Spec(A)
$$
Nous affirmons que ce morphisme est universellement submersif. Pour le prouver,
soit $\Spec(B) \to \Spec(A)$ un morphisme de schémas affines défini par
le morphisme d'anneaux $A \to B$. Il faut montrer que
$$
\Spec(B_\mathfrak p) \amalg \Spec(B/\mathfrak pB) \to \Spec(B)
$$
soit submersif. Il est tout d'abord surjectif. Supposons ensuite que
$T \subset \Spec(B)$ soit une partie telle que
$T_1 = \Spec(B_\mathfrak p) \cap T$ et
$T_2 = \Spec(B/\mathfrak p B) \cap T$ soient fermées. Alors $T$ est l'image
du spectre d'une $B$-algèbre, car $T_1$ et $T_2$ sont tous deux des spectres
de $B$-algèbres. Pour montrer que $T$ est fermé, il suffit donc d'établir
que $T$ est stable par spécialisation ; voir
Algèbre, lemme \ref{algebra-lemma-image-stable-specialization-closed}.
Supposons que $p \leadsto q$ soit une spécialisation de points de
$\Spec(B)$, avec $p \in T$. Soit $A'$ un anneau de valuation et soit
$\Spec(A') \to \Spec(B)$ un morphisme tel que le point générique $\eta$
de $\Spec(A')$ s'envoie sur $p$ et le point fermé $s$ de $\Spec(A')$ sur $q$ ; voir
Schémas, lemme \ref{schemes-lemma-points-specialize}.
Observons que l'image du composé $\gamma : \Spec(A') \to \Spec(A)$
est exactement
l'ensemble des points
$\xi \in \Spec(A)$ avec $\gamma(\eta) \leadsto \xi \leadsto \gamma(s)$ (les détails sont omis).
Si $\mathfrak p \not \in \Im(\gamma)$, alors soit
$p, q \in \Spec(B_\mathfrak p)$, soit
$p, q \in \Spec(B/\mathfrak pB)$. Dans ce cas, puisque
$T_1$, resp.\ $T_2$, est fermée,
on a $q \in T_1$, resp.\ $q \in T_2$, et donc $q \in T$.
Supposons enfin que $\mathfrak p \in \Im(\gamma)$, disons
$\mathfrak p = \gamma(r)$. Nous avons alors les spécialisations
$p \leadsto r$ et $r \leadsto q$. Dans ce cas,
$p, r \in \Spec(B_\mathfrak p)$ et
$r, q \in \Spec(B/\mathfrak pB)$. On conclut d'abord que
$r \in T_1 \subset T$, puis que $r \in T_2$ puisque $r$ s'envoie sur
$\mathfrak p$, et enfin que $q \in T_2 \subset T$, comme voulu.

\medskip\noindent
D'autre part, nous affirmons que la famille réduite à un seul morphisme
$$
\{\Spec(A_\mathfrak p) \amalg \Spec(A/\mathfrak p) \longrightarrow \Spec(A)\}
$$
n'est pas un recouvrement V. Voir
Topologies, définition \ref{topologies-definition-V-covering}.
En effet, si elle était un recouvrement V, il existerait une extension
d'anneaux de valuation $A \subset B$ telle que
$\Spec(B) \to \Spec(A)$ se factorise par
$\Spec(A_\mathfrak p) \amalg \Spec(A/\mathfrak p)$.
Cela impliquerait que $\Spec(B)$ est non connexe, ce qui est absurde.

\begin{lemma}
\label{examples-lemma-universally-submersive-not-V}
Il existe un morphisme $X \to Y$ de schémas affines,
universellement submersif, tel que $\{X \to Y\}$
ne soit pas un recouvrement V.
\end{lemma}

\begin{proof}
Voir la discussion ci-dessus.
\end{proof}









\section{Le spectre des entiers n'est pas quasi-compact}
\label{examples-section-canonical}

\noindent
Bien entendu, le titre de cette section ne concerne pas le spectre des entiers
considéré comme espace topologique, car tout spectre est quasi-compact
en tant qu'espace topologique
(Algèbre, lemme \ref{algebra-lemma-quasi-compact}).
Il concerne plutôt le spectre des entiers pour la topologie canonique
sur la catégorie des schémas et la définition d'un objet quasi-compact
dans un site
(Sites, définition \ref{sites-definition-quasi-compact}).

\medskip\noindent
Soit $U$ un ultrafiltre non principal sur l'ensemble $P$ des nombres premiers.
Pour une partie $T \subset P$, notons $T^c = P \setminus T$ son complément.
Pour $A \in U$, soit $S_A \subset \mathbf{Z}$ la partie multiplicative
engendrée par les $p \in A$. Posons
$$
\mathbf{Z}_A = S_A^{-1}\mathbf{Z}
$$
Observons que $\Spec(\mathbf{Z}_A) = \{(0)\} \cup A^c \subset \Spec(\mathbf{Z})$
si l'on identifie $P$ à l'ensemble des points fermés de $\Spec(\mathbf{Z})$.
Si $A, B \in U$, alors $A \cap B \in U$ et $A \cup B \in U$
et nous avons
$$
\mathbf{Z}_{A \cap B} =
\mathbf{Z}_A \times_{\mathbf{Z}_{A \cup B}} \mathbf{Z}_B
$$
(produit fibré d'anneaux). En particulier, pour tout
entier $n$ et tous $A_1, \ldots, A_n \in U$, le morphisme
$$
\Spec(\mathbf{Z}_{A_1}) \amalg \ldots \amalg \Spec(\mathbf{Z}_{A_n})
\longrightarrow \Spec(\mathbf{Z})
$$
se factorise par $\Spec(\mathbf{Z}[1/p])$ pour un certain $p$
(en fait, pour tout $p \in A_1 \cap \ldots \cap A_n$).
Nous concluons que la famille de morphismes plats
$\{\Spec(\mathbf{Z}_A) \to \Spec(\mathbf{Z})\}_{A \in U}$
est conjointement surjective, mais qu'aucune sous-famille finie ne l'est.

\medskip\noindent
Pour un $\mathbf{Z}$-module $M$, posons
$$
M_A = S_A^{-1}M = M \otimes_{\mathbf{Z}} \mathbf{Z}_A
$$
Affirmation I : pour tout $\mathbf{Z}$-module $M$, on a
$$
M =
\text{Égaliseur}\left(
\xymatrix{
\prod\nolimits_{A \in U} M_A \ar@<1ex>[r] \ar@<-1ex>[r] &
\prod\nolimits_{A, B \in U} M_{A \cup B}
}
\right)
$$
Supposons d'abord $M$ sans torsion. Alors $M_A \subset M_P$ pour tout
$A \in U$. Il suffit donc de montrer que
$$
M = \bigcap\nolimits_{A \in U} M_A\text{ dans }M_P = M \otimes \mathbf{Q}
$$
En effet, comme $U$ n'est pas principal, pour tout nombre premier $p$, on a
$\{p\}^c \in U$. De plus, $M_{\{p\}^c} = M_{(p)}$ est la localisation
en l'idéal premier $(p)$.
L'égalité ci-dessus est donc claire, puisque $M = \bigcap_{p \in P} M_{(p)}$.
Supposons ensuite que $M$ soit de torsion. Alors
$$
M = \bigoplus\nolimits_{p \in P} M[p^\infty]
$$
et, par conséquent,
$$
M_A = \bigoplus\nolimits_{p \not \in A} M[p^\infty]
$$
puisque nous localisons aux nombres premiers appartenant à $A$.
Supposons que $(x_A) \in \prod M_A$ appartienne à l'égalisateur.
Notons $x_p = x_{\{p\}^c} \in M[p^{\infty}]$.
La propriété d'égalisation donne
$$
x_A = (x_p)_{p \not \in A}
$$
et implique en particulier que $x_p$ est nul sauf pour un nombre fini de
$p \not \in A$. Pour achever la preuve dans le cas de torsion, il suffit
de montrer que $x_p$ est nul pour presque tout nombre premier $p$.
Sinon, écrivons
$\{p \in P \mid x_p \not = 0\} = T \amalg T'$ comme réunion disjointe
de deux ensembles infinis. Alors $T \not \in U$ ou $T' \not \in U$,
puisque $U$ est un ultrafiltre (si $T, T'$ appartenaient à $U$, alors $U$ contiendrait
$T \cap T' = \emptyset$, ce qui est impossible).
Supposons $T \not \in U$. Alors $T = A^c$,
ce qui contredit la finitude établie ci-dessus.
Enfin, soit $M$ un module arbitraire. Considérons la suite exacte courte
$$
0 \to M_{tors} \to M \to M/M_{tors} \to 0
$$
et le grand diagramme suivant :
$$
\xymatrix{
M_{tors} \ar[r] \ar[d] &
\prod\nolimits_{A \in U} M_{tors, A} \ar@<1ex>[r] \ar@<-1ex>[r] \ar[d] &
\prod\nolimits_{A, B \in U} M_{tors, A \cup B} \ar[d] \\
M \ar[r] \ar[d] &
\prod\nolimits_{A \in U} M_A \ar@<1ex>[r] \ar@<-1ex>[r] \ar[d] &
\prod\nolimits_{A, B \in U} M_{A \cup B} \ar[d] \\
M/M_{tors} \ar[r] &
\prod\nolimits_{A \in U} (M/M_{tors})_A \ar@<1ex>[r] \ar@<-1ex>[r] &
\prod\nolimits_{A, B \in U} (M/M_{tors})_{A \cup B} \\
}
$$
Une chasse au diagramme, utilisant l'exactitude des colonnes ainsi que
le résultat pour le module de torsion $M_{tors}$ et pour le module
sans torsion $M/M_{tors}$, démontre l'affirmation I pour $M$.
On obtient ainsi un exemple du phénomène décrit dans le lemme suivant.

\begin{lemma}
\label{examples-lemma-non-fpqc-descent}
Il existe un anneau $A$ et une famille infinie de morphismes plats
d'anneaux $\{A \to A_i\}_{i \in I}$ telle que, pour tout $A$-module $M$,
$$
M =
\text{Égaliseur}\left(
\xymatrix{
\prod\nolimits_{i \in I} M \otimes_A A_i \ar@<1ex>[r] \ar@<-1ex>[r] &
\prod\nolimits_{i, j \in I} M \otimes_A A_i \otimes_A A_j
}
\right)
$$
mais qu'aucune sous-famille finie ne possède cette propriété.
\end{lemma}

\begin{proof}
Voir la discussion ci-dessus.
\end{proof}

\noindent
Nous continuons avec notre ultrafiltre non principal $U$ sur l'ensemble $P$
des nombres premiers.
Soit $R$ un anneau. Notons $R_A = S_A^{-1}R = R \otimes \mathbf{Z}_A$
pour $A \in U$.
Affirmation II : étant données des parties fermées
$T_A \subset \Spec(R_A)$, $A \in U$, telles que
$$
(\Spec(R_{A \cup B}) \to \Spec(R_A))^{-1}T_A =
(\Spec(R_{A \cup B}) \to \Spec(R_B))^{-1}T_B
$$
pour tous $A, B \in U$, il existe une partie fermée $T \subset \Spec(R)$ telle que
$T_A = (\Spec(R_A) \to \Spec(R))^{-1}(T)$ pour tout $A \in U$.
Soit $I_A \subset R_A$, pour $A \in U$, l'idéal radical définissant $T_A$.
La condition de recollement implique alors
$S_{A \cup B}^{-1}I_A = S_{A \cup B}^{-1}I_B$ dans $R_{A \cup B}$ pour
tous $A, B \in U$ (car la localisation préserve le caractère radical
d'un idéal). Soit $I' \subset R$ l'ensemble des éléments
dont l'image appartient à $I_P \subset R_P = R \otimes \mathbf{Q}$.
Pour $A \in U$, on a alors :
\begin{enumerate}
\item $I_A \subset I'_A = S_A^{-1}I'$, et
\item $M_A = I'_A/I_A$ est un module de torsion.
\end{enumerate}
On obtient bien sûr des identifications canoniques
$S_{A \cup B}^{-1}M_A = S_{A \cup B}^{-1}M_B$ pour $A, B \in U$.
En décomposant les modules de torsion $M_A$ en leurs composantes
$p$-primaires, le lecteur vérifie aisément qu'il existe, pour chaque $p$, des
$R$-modules $M_p$ tels que
$$
M_A = \bigoplus\nolimits_{p \not \in A} M_p
$$
de manière compatible avec les identifications canoniques ci-dessus.
En posant $M = \bigoplus_{p \in P} M_p$, on obtient des isomorphismes
canoniques $M_A = S_A^{-1}M$, compatibles avec ces identifications.
On obtient alors une application canonique
$$
I' \longrightarrow M
$$
de $R$-modules qui redonne l'application $I'_A \to M_A$ pour tout
$A \in U$. Cela résulte de toutes les compatibilités mentionnées
ci-dessus et de l'affirmation démontrée précédemment, selon laquelle
$M$ est l'égalisateur des deux applications de
$\prod_{A \in U} M_A$ vers $\prod_{A, B \in U} M_{A \cup B}$.
Soit $I = \Ker(I' \to M)$. Alors $I$ est un idéal, et
$T = V(I)$ est une partie fermée qui redonne les parties fermées $T_A$ pour tout $A \in U$. Cela démontre l'affirmation II.

\begin{lemma}
\label{examples-lemma-Z-not-quasi-compact}
Le schéma $\Spec(\mathbf{Z})$ n'est pas quasi-compact pour
la topologie canonique sur la catégorie des schémas.
\end{lemma}

\begin{proof}
Avec les notations précédentes, considérons la famille de morphismes
$$
\mathcal{W} = \{\Spec(\mathbf{Z}_A) \to \Spec(\mathbf{Z})\}_{A \in U}
$$
D'après Descente, lemme \ref{descent-lemma-universal-effective-epimorphism}, et les deux
affirmations démontrées ci-dessus, c'est un épimorphisme effectif
universel. Dans toute catégorie $\mathcal{C}$ possédant des produits fibrés, les
épimorphismes effectifs universels y définissent une
structure de site (à quelques difficultés ensemblistes près, faciles
à résoudre) qui définit la topologie canonique.
Ainsi, $\mathcal{W}$ est un recouvrement pour la topologie canonique.
D'autre part, nous avons vu ci-dessus que toute sous-famille finie
$$
\{\Spec(\mathbf{Z}_{A_i}) \to \Spec(\mathbf{Z})\}_{i = 1, \ldots, n},\quad
n \in \mathbf{N}, A_1, \ldots, A_n \in U
$$
se factorise par $\Spec(\mathbf{Z}[1/p])$ pour un certain $p$.
Cette famille finie ne peut donc pas être un épimorphisme effectif
universel et, plus généralement, aucun épimorphisme effectif universel
$\{g_j : T_j \to \Spec(\mathbf{Z})\}$ ne peut raffiner
$\{\Spec(\mathbf{Z}_{A_i}) \to \Spec(\mathbf{Z})\}_{i = 1, \ldots, n}$.
D'après Sites, définition \ref{sites-definition-quasi-compact},
cela signifie que $\Spec(\mathbf{Z})$ n'est pas quasi-compact
pour la topologie canonique. Pour voir que notre notion de
quasi-compacité concorde avec la définition usuelle
en théorie des topos, voir Sites, lemme \ref{sites-lemma-quasi-compact}.
\end{proof}











% OK, start here.
%

\chapter{Exercices}



\phantomsection
\label{exercises-section-phantom}



\section{Algèbre}
\label{exercises-section-algebra}

\noindent
Cette première section ne contient que quelques questions diverses.

\begin{exercise}
\label{exercises-exercise-isomorphism-localizations}
Soient $A$ un anneau et ${\mathfrak m}$ un idéal maximal. Dans $A[X]$,
posons $\tilde {\mathfrak m}_1 = ({\mathfrak m}, X)$ et
$\tilde {\mathfrak m}_2 = ({\mathfrak m}, X-1)$. Montrer
que
$$
A[X]_{\tilde {\mathfrak m}_1} \cong A[X]_{\tilde {\mathfrak m}_2}.
$$
\end{exercise}

\begin{exercise}
\label{exercises-exercise-coherent}
Trouver un exemple d'anneau non noethérien $R$ tel que tout idéal
de type fini de $R$ soit de présentation finie comme $R$-module.
(On dit qu'un anneau est {\it cohérent} si cette dernière propriété est vérifiée.)
\end{exercise}

\begin{exercise}
\label{exercises-exercise-flat-ideals-pid}
Supposons que $(A, {\mathfrak m}, k)$ soit un anneau local noethérien. Pour tout
$A$-module de type fini $M$, notons $r(M)$ le nombre minimal
de générateurs de $M$ comme $A$-module. D'après le lemme de Nakayama, ce nombre vaut
$\dim_k M/{\mathfrak m} M = \dim_k M \otimes_A k$.
\begin{enumerate}
\item Montrer que $r(M \otimes_A N) = r(M)r(N)$.
\item Soit $I\subset A $ un idéal tel que $r(I) > 1$. Montrer que
$r(I^2) < r(I)^2$.
\item En déduire que, si tout idéal de $A$ est un module plat, alors
$A$ est un anneau principal (ou un corps).
\end{enumerate}
\end{exercise}

\begin{exercise}
\label{exercises-exercise-not-isomorphic}
Soit $k$ un corps. Montrer que les paires suivantes de
$k$-algèbres ne sont pas isomorphes :
\begin{enumerate}
\item $k[x_1, \ldots, x_n]$ et $k[x_1, \ldots, x_{n + 1}]$ pour tout
$n\geq 1$.
\item $k[a, b, c, d, e, f]/(ab + cd + ef)$ et $k[x_1, \ldots, x_n]$
pour $n = 5$.
\item $k[a, b, c, d, e, f]/(ab + cd + ef)$ et $k[x_1, \ldots, x_n]$
pour $n = 6$.
\end{enumerate}
\end{exercise}

\begin{remark}
\label{exercises-remark-simple-geometric}
L'idée de cet exercice est bien entendu de trouver, dans chaque cas,
un argument simple plutôt que d'appliquer un « gros » théorème.
Il est néanmoins utile de se laisser guider par des principes généraux.
\end{remark}

\begin{exercise}
\label{exercises-exercise-silly}
Algèbre. (Exercice un peu futile, qui devrait être facile.)
\begin{enumerate}
\item Donner un exemple d'un anneau $A$ et d'une
suite exacte courte non scindée de $A$-modules
$$
0 \to M_1 \to M_2 \to M_3 \to 0.
$$
\item Donner un exemple d'une suite non scindée de $A$-modules
comme ci-dessus et d'un homomorphisme fidèlement plat $A \to B$ tel que
$$
0 \to M_1\otimes_A B \to M_2\otimes_A B \to M_3\otimes_A B \to 0.
$$
soit scindée comme suite de $B$-modules.
\end{enumerate}
\end{exercise}

\begin{exercise}
\label{exercises-exercise-field-kummer}
Supposons que le corps $k$ contienne une racine primitive $n$-ième
de l'unité $\zeta$. Cela signifie que $\zeta^n = 1$, mais que $\zeta^m\not = 1$ pour
$0 < m < n$.
\begin{enumerate}
\item Montrer que la caractéristique de $k$ est première à $n$.
\item Supposons que $a \in k$ soit un élément de $k$ qui ne soit une
puissance $d$-ième dans $k$ pour aucun diviseur $d$ de $n$ tel que $n \geq d > 1$. Montrer que
$k[x]/(x^n-a)$ est un corps. (Indication : considérer un corps de décomposition de
$x^n-a$ et utiliser la théorie de Galois.)
\end{enumerate}
\end{exercise}

\begin{exercise}
\label{exercises-exercise-valuation}
Soit $\nu : k[x]\setminus \{0\}  \to {\mathbf Z}$ une application
ayant les propriétés suivantes : $\nu(fg) = \nu(f) + \nu(g)$ lorsque
$f$, $g$ sont non nuls, et $\nu(f + g) \geq min(\nu(f), \nu(g))$ lorsque
$f$, $g$ et $f + g$ sont non nuls, et $\nu(c) = 0$ pour tout $c\in k^*$.
\begin{enumerate}
\item Montrer que, si $f$, $g$ et $f + g$ sont non nuls et si
$\nu(f) \not = \nu(g)$, alors $\nu(f + g) = min(\nu(f), \nu(g))$.
\item Montrer que, si $f = \sum a_i x^i$, $f\not = 0$, alors
$\nu(f) \geq min(\{i\nu(x)\}_{a_i\not = 0})$. Quand a-t-on égalité ?
\item Montrer que, si $\nu$ prend une valeur négative, alors
$\nu(f) = -n \deg(f)$ pour un certain $n\in {\mathbf N}$.
\item Supposons que $\nu(x) \geq 0$. Montrer que
$\{f \mid f = 0, \ \text{ou}\ \nu(f) > 0\}$ est un idéal premier de $k[x]$.
\item Décrire toutes les applications $\nu$ possibles.
\end{enumerate}
\end{exercise}


\noindent
Soit $A$ un anneau. Un {\it idempotent} est un élément $e \in A$
tel que $e^2 = e$. Les éléments $1$ et $0$ sont toujours idempotents.
Un {\it idempotent non trivial} est un idempotent distinct de zéro et de l'unité.
Deux idempotents $e, e' \in A$ sont dits {\it orthogonaux}
si $ee' = 0$.

\begin{exercise}
\label{exercises-exercise-product}
Soit $A$ un anneau. Montrer que $A$ est un produit de deux anneaux non nuls si
et seulement si $A$ possède un idempotent non trivial.
\end{exercise}

\begin{exercise}
\label{exercises-exercise-lift-idempotents}
Soient $A$ un anneau et $I \subset A$ un idéal localement nilpotent.
Montrer que l'application $A \to A/I$ induit une bijection sur les idempotents.
(Indication : il peut être plus facile de commencer par le cas où $I$ est nilpotent.
Utiliser ensuite la « réduction noethérienne absolue » pour se ramener à ce cas.)
\end{exercise}







% Fin de la tranche source louée.
\section{Limites inductives}
\label{exercises-section-colimits}


\begin{definition}
\label{exercises-definition-directed-poset}
Un {\it ensemble préordonné filtrant} est un ensemble non vide $I$ muni d'un préordre $\leq$
tel que, pour toute paire $i, j \in I$, il existe un $k \in I$ tel que
$i \leq k$ et $j \leq k$. Un {\it système inductif d'anneaux} indexé par $I$ est la donnée d'un
anneau $A_i$ pour chaque $i \in I$ et d'un homomorphisme d'anneaux $\varphi_{ij} : A_i \to A_j$
dès que $i \leq j$, de sorte que le composé $A_i \to A_j \to A_k$ soit égal à
$A_i \to A_k$ dès que $i \leq j \leq k$.
\end{definition}

\noindent
On définit de même les systèmes inductifs de groupes, de modules sur un anneau fixé,
d'espaces vectoriels sur un corps, etc.

\begin{exercise}
\label{exercises-exercise-directed-colimit}
Soit $I$ un ensemble préordonné filtrant et soit
$(A_i, \varphi_{ij})$ un système inductif d'anneaux indexé par $I$.
Montrer qu'il existe un anneau $A$ et des homomorphismes $\varphi_i : A_i \to A$
tels que $\varphi_j \circ \varphi_{ij} = \varphi_i$ pour tous $i \leq j$
et possédant la propriété universelle suivante : étant donnés un anneau $B$
et des homomorphismes $\psi_i : A_i \to B$ tels que
$\psi_j \circ \varphi_{ij} = \psi_i$ pour tous $i \leq j$, il
existe un unique homomorphisme d'anneaux $\psi : A \to B$ tel que
$\psi_i = \psi \circ \varphi_i$.
\end{exercise}

\begin{definition}
\label{exercises-definition-colimit}
L'anneau $A$ construit dans l'exercice \ref{exercises-exercise-directed-colimit}
est appelé la {\it limite inductive} du système. Notation : $\colim A_i$.
\end{definition}

\begin{exercise}
\label{exercises-exercise-prime-in-colimit}
Soit $(I, \geq)$ un ensemble préordonné filtrant et soit
$(A_i, \varphi_{ij})$ un système inductif d'anneaux indexé par $I$, de limite inductive
$A$. Montrer qu'il existe une bijection
$$
\Spec(A) = \{(\mathfrak p_i)_{i \in I} \mid
\mathfrak p_i \subset A_i \text{ et }
\mathfrak p_i = \varphi_{ij}^{-1}(\mathfrak p_j)\ \forall i \leq j\}
\subset \prod\nolimits_{i \in I} \Spec(A_i)
$$
L'ensemble figurant au membre de droite de l'égalité est la limite projective des ensembles
$\Spec(A_i)$. Notation : $\lim \Spec(A_i)$.
\end{exercise}

\begin{exercise}
\label{exercises-exercise-colimit-surjective}
Soit $(I, \geq)$ un ensemble préordonné filtrant et soit
$(A_i, \varphi_{ij})$ un système inductif d'anneaux indexé par $I$, de limite inductive
$A$. Supposons que l'application $\Spec(A_j) \to \Spec(A_i)$ soit
surjective pour tous $i \leq j$. Montrer que
$\Spec(A) \to \Spec(A_i)$ est surjective pour tout $i$.
(Indication : on peut essayer d'utiliser le théorème de Tychonoff, mais il existe aussi une
démonstration algébrique directe essentiellement triviale fondée sur
Algèbre, lemme \ref{algebra-lemma-in-image}.)
\end{exercise}

\begin{exercise}
\label{exercises-exercise-integral-colimit-finite}
Soit une extension entière d'anneaux $A \subset B$. Montrer que
$\Spec(B) \to \Spec(A)$ est surjective.
Utiliser les exercices ci-dessus, le fait que l'assertion est vraie pour une extension finie
d'anneaux (ce qui a été démontré dans les cours), et montrer que
$B = \colim B_i$ est une limite inductive filtrante d'extensions finies $A \subset B_i$.
\end{exercise}

\begin{exercise}
\label{exercises-exercise-colimit-tensor}
Soit $(I, \geq)$ un ensemble préordonné filtrant.
Soit $A$ un anneau et soit $(N_i, \varphi_{i, i'})$ un système inductif de
$A$-modules indexé par $I$. Soit en outre $M$ un $A$-module. Montrer
que
$$
\colim_{i\in I} M \otimes_A N_i\cong
M \otimes_A \Big( \colim_{i\in I} N_i\Big).
$$
\end{exercise}

\begin{definition}
\label{exercises-definition-finite-presentation}
Un module $M$ sur $R$ est dit de {\it présentation finie} sur
$R$ s'il est isomorphe au conoyau d'un homomorphisme de modules libres de type fini
$ R^{\oplus n} \to R^{\oplus m}$.
\end{definition}

\begin{exercise}
\label{exercises-exercise-colimit-modules}
Montrer que tout module sur un anneau quelconque est
\begin{enumerate}
\item la limite inductive de ses sous-modules de type fini, et
\item, d'une certaine manière, une limite inductive de modules de présentation finie.
\end{enumerate}
\end{exercise}





\section{Catégories additives et abéliennes}
\label{exercises-section-additive}

\begin{exercise}
\label{exercises-exercise-filtered-vector-spaces}
Soit $k$ un corps. Soit $\mathcal{C}$ la catégorie des espaces vectoriels filtrés
sur $k$ ; voir
Homologie, définition \ref{homology-definition-filtered}
pour la définition d'un objet filtré d'une catégorie quelconque.
\begin{enumerate}
\item Montrer que c'est une catégorie additive (expliquer soigneusement ce qu'est la
somme directe de deux objets).
\item Soit $f : (V, F) \to (W, F)$ un morphisme de $\mathcal{C}$.
Montrer que $f$ possède un noyau et un conoyau (expliquer précisément
ce que sont le noyau et le conoyau de $f$).
\item Donner un exemple d'un morphisme de $\mathcal{C}$ tel que
le morphisme canonique $\Coim(f) \to \Im(f)$ ne soit pas un isomorphisme.
\end{enumerate}
\end{exercise}

\begin{exercise}
\label{exercises-exercise-torsion-free}
Soit $R$ un anneau intègre noethérien. Soit $\mathcal{C}$ la catégorie des
$R$-modules sans torsion de type fini.
\begin{enumerate}
\item Montrer que c'est une catégorie additive.
\item Soit $f : N \to M$ un morphisme de $\mathcal{C}$.
Montrer que $f$ possède un noyau et un conoyau (veiller à définir précisément
ce que sont le noyau et le conoyau de $f$).
\item Donner un exemple d'un anneau intègre noethérien $R$ et d'un morphisme de
$\mathcal{C}$ tel que le morphisme canonique $\Coim(f) \to \Im(f)$
ne soit pas un isomorphisme.
\end{enumerate}
\end{exercise}

\begin{exercise}
\label{exercises-exercise-other}
Donner un exemple d'une catégorie additive qui possède des noyaux
et des conoyaux, mais qui soit autre que celles des
exercices \ref{exercises-exercise-filtered-vector-spaces} et
\ref{exercises-exercise-torsion-free}.
\end{exercise}





\section{Produit tensoriel}
\label{exercises-section-tensor-product}

\noindent
Les produits tensoriels sont introduits dans Algèbre, section
\ref{algebra-section-tensor-product}.
Soit $R$ un anneau. Soit $\text{Mod}_R$ la catégorie des $R$-modules.
Nous dirons qu'un foncteur $F : \text{Mod}_R \to \text{Mod}_R$
\begin{enumerate}
\item est additif si
$F : \Hom_R(M, N) \to \Hom_R(F(M), F(N))$
est un homomorphisme de groupes abéliens
pour tous $R$-modules $M, N$ ; voir
Homologie, définition \ref{homology-definition-preadditive} ;
\item est $R$-linéaire si $F : \Hom_R(M, N) \to \Hom_R(F(M), F(N))$ est $R$-linéaire
pour tous $R$-modules $M, N$ ;
\item est exact à droite si, pour toute suite exacte courte
$0 \to M_1 \to M_2 \to M_3 \to 0$, la suite
$F(M_1) \to F(M_2) \to F(M_3) \to 0$ est exacte ;
\item est exact à gauche si, pour toute suite exacte courte
$0 \to M_1 \to M_2 \to M_3 \to 0$, la suite
$0 \to F(M_1) \to F(M_2) \to F(M_3)$ est exacte ;
\item commute aux sommes directes si, étant donnés un ensemble
$I$ et des $R$-modules $M_i$, les homomorphismes
$F(M_i) \to F(\bigoplus M_i)$ induisent un isomorphisme
$\bigoplus F(M_i) = F(\bigoplus M_i)$.
\end{enumerate}

\begin{exercise}
\label{exercises-exercise-characterize-tensor-functor}
Soit $R$ un anneau. On conserve les notations ci-dessus.
\begin{enumerate}
\item Donner un exemple d'un anneau $R$ et d'un foncteur additif
$F : \text{Mod}_R \to \text{Mod}_R$ qui ne soit pas $R$-linéaire.
\item Soit $N$ un $R$-module. Montrer que le foncteur
$F(M) = M \otimes_R N$ est $R$-linéaire,
exact à droite et commute aux sommes directes.
\item Réciproquement, montrer que tout foncteur $F : \text{Mod}_R \to \text{Mod}_R$
qui est $R$-linéaire,
exact à droite et commute aux sommes directes est de la
forme $F(M) = M \otimes_R N$ pour un certain $R$-module $N$.
\item Montrer que, si dans (3) on supprime l'hypothèse
selon laquelle $F$ commute aux sommes directes, la conclusion n'est
plus vraie.
\end{enumerate}
\end{exercise}





\section{Homomorphismes plats d'anneaux}
\label{exercises-section-flat}

\begin{exercise}
\label{exercises-exercise-localization-flat}
Soit $S$ une partie multiplicative de l'anneau $A$.
\begin{enumerate}
\item Pour un $A$-module $M$, montrer que $S^{-1}M = S^{-1}A \otimes_A M$.
\item Montrer que $S^{-1}A$ est plat sur $A$.
\end{enumerate}
\end{exercise}

\begin{exercise}
\label{exercises-exercise-examples-not-flat}
Trouver une injection $M_1 \to M_2$ de $A$-modules telle que
$M_1\otimes N \to M_2 \otimes N$ ne soit pas injective dans les
cas suivants :
\begin{enumerate}
\item $A = k[x, y]$ et $N = (x, y) \subset A$. (Ici et ci-dessous, $k$ est un corps.)
\item $A = k[x, y]$ et $N = A/(x, y)$.
\end{enumerate}
\end{exercise}

\begin{exercise}
\label{exercises-exercise-flat-not-projective}
Donner un exemple d'un anneau $A$ et d'un $A$-module de type fini $M$
qui soit plat mais non projectif comme $A$-module.
\end{exercise}

\begin{remark}
\label{exercises-remark-flat-not-projective}
Si $M$ est de présentation finie et plat sur $A$,
alors $M$ est projectif sur $A$. L'exemple devra donc faire
intervenir un anneau $A$ qui ne soit pas noethérien. Je connais un exemple
où $A$ est l'anneau des fonctions de classe ${\mathcal C}^\infty$ sur ${\mathbf R}$.
\end{remark}

\begin{exercise}
\label{exercises-exercise-flat-not-free-dvr}
Trouver un module plat mais non libre sur ${\mathbf Z}_{(2)}$.
\end{exercise}

\begin{exercise}
\label{exercises-exercise-flat-deformations}
Déformations plates.
\begin{enumerate}
\item Supposons que $k$ soit un corps et que $k[\epsilon]$ soit l'algèbre des
nombres duaux $k[\epsilon] = k[x]/(x^2)$, avec $\epsilon = \bar x$. Montrer que, pour
toute $k$-algèbre $A$, il existe une $k[\epsilon]$-algèbre plate $B$ telle que
$A$ soit isomorphe à $B/\epsilon B$.
\item Supposons que $k = {\mathbf F}_p = {\mathbf Z}/p{\mathbf Z}$ et
$$
A = k[x_1, x_2, x_3, x_4, x_5, x_6]/
(x_1^p, x_2^p, x_3^p, x_4^p, x_5^p, x_6^p).
$$
Montrer qu'il existe une ${\mathbf Z}/p^2{\mathbf Z}$-algèbre plate $B$ telle
que $B/pB$ soit isomorphe à $A$. (Ainsi, $p$ joue ici le rôle de $\epsilon$.)
\item Posons maintenant $p = 2$ et considérons la même question pour
$k = {\mathbf F}_2 = {\mathbf Z}/2{\mathbf Z}$ et
$$
A = k[x_1, x_2, x_3, x_4, x_5, x_6]/
(x_1^2, x_2^2, x_3^2, x_4^2, x_5^2, x_6^2, x_1x_2 + x_3x_4 + x_5x_6).
$$
Cependant, montrer que dans ce cas il n'existe {\it pas} de
${\mathbf Z}/4{\mathbf Z}$-algèbre plate $B$ telle que $B/2B$ soit isomorphe à
$A$. (Trouver l'astuce ! Le même exemple fonctionne en toute caractéristique
$p > 0$, mais le calcul est plus difficile.)
\end{enumerate}
\end{exercise}

\begin{exercise}
\label{exercises-exercise-flat-given-residue-field-extension}
Soit $(A, {\mathfrak m}, k)$ un anneau local et soit $k'/k$
une extension finie de corps. Montrer qu'il existe un homomorphisme local et plat
d'anneaux locaux $A \to B$ tel que ${\mathfrak m}_B = {\mathfrak m} B$ et que
$B/{\mathfrak m} B$ soit
isomorphe à $k'$ comme $k$-algèbre. (Indication : commencer par le cas où
l'extension $k \subset k'$ est engendrée par un seul élément.)
\end{exercise}

\begin{remark}
\label{exercises-remark-flat-given-residue-field-extension-general}
Le même résultat vaut pour toute extension de corps $K/k$.
\end{remark}





% Fin de la tranche source assignée.




\section{Le spectre d'un anneau}
\label{exercises-section-spectrum-ring}

\begin{exercise}
\label{exercises-exercise-spec-Z}
Déterminer $\Spec(\mathbf{Z})$ en tant qu'ensemble et décrire sa topologie.
\end{exercise}

\begin{exercise}
\label{exercises-exercise-basis-opens-standard}
Soit $A$ un anneau quelconque. Pour $f\in A$, on définit
$D(f):= \{\mathfrak p \subset A \mid f \not \in \mathfrak p\}$.
Montrer que les ouverts $D(f)$ forment une base de la topologie de
$\Spec(A)$.
\end{exercise}

\begin{exercise}
\label{exercises-exercise-radical-ideals-closed}
Montrer que l'application $I\mapsto V(I)$
définit une bijection naturelle
$$
\{I\subset A\text{ tel que }I = \sqrt{I}\}
\longrightarrow
\{T\subset \Spec(A)\text{ fermée}\}
$$
\end{exercise}

\begin{definition}
\label{exercises-definition-quasi-compact}
Un espace topologique $X$ est dit {\it quasi-compact}
si, pour tout recouvrement ouvert $X = \bigcup_{i\in I} U_i$, il existe une
partie finie $\{i_1, \ldots, i_n\}\subset I$ telle que $X = U_{i_1}\cup\ldots
U_{i_n}$.
\end{definition}

\begin{exercise}
\label{exercises-exercise-spec-quasi-compact}
Montrer que $\Spec(A)$ est quasi-compact pour tout anneau $A$.
\end{exercise}

\begin{definition}
\label{exercises-definition-Hausdorff}
On dit qu'un espace topologique $X$ vérifie l'axiome de séparation $T_0$
si, pour tout couple de points $x, y\in X$ tels que $x\not = y$, il existe un
ouvert de $X$ qui contient l'un mais pas l'autre.
On dit que $X$ est {\it séparé} si, pour tout couple $x, y\in X$ tel que $x\not = y$,
il existe des ouverts disjoints $U, V$ tels que $x\in U$
et $y\in V$.
\end{definition}

\begin{exercise}
\label{exercises-exercise-not-hausdorff}
Montrer que $\Spec(A)$ n'est {\bf pas} séparé en général.
Montrer que $\Spec(A)$ est $T_0$. Donner un exemple d'espace topologique
$X$ qui n'est pas $T_0$.
\end{exercise}

\begin{remark}
\label{exercises-remark-not-hausdorff}
En général, on réserve le mot compact aux espaces quasi-compacts et
séparés.
\end{remark}

\begin{definition}
\label{exercises-definition-irreducible}
Un espace topologique $X$ est dit {\it irréductible} si $X$ est non vide
et si $X = Z_1\cup Z_2$, avec $Z_1, Z_2\subset X$ fermés, alors soit
$Z_1 = X$, soit $Z_2 = X$. Une partie $T\subset X$ d'un espace topologique
est dite {\it irréductible} si elle est un espace topologique
irréductible pour la topologie induite par $X$.
Cette définition implique que $T$ est irréductible si et seulement
si l'adhérence $\bar T$ de $T$ dans $X$ est irréductible.
\end{definition}

\begin{exercise}
\label{exercises-exercise-irreducible-spec}
Montrer que $\Spec(A)$ est irréductible si et seulement si
$Nil(A)$ est un idéal premier et que, dans ce cas, c'est l'unique
idéal premier minimal de $A$.
\end{exercise}

\begin{exercise}
\label{exercises-exercise-irreducible-prime}
Montrer qu'une partie fermée $T\subset \Spec(A)$
est irréductible si et seulement si elle est de la forme $T = V({\mathfrak p})$ pour
un idéal premier ${\mathfrak p}\subset A$.
\end{exercise}

\begin{definition}
\label{exercises-definition-generic-point}
Un point $x$ d'un espace topologique irréductible $X$ est appelé
un {\it point générique} de $X$ si $X$ est égal à l'adhérence de
la partie $\{x\}$.
\end{definition}

\begin{exercise}
\label{exercises-exercise-irreducible-T0-at-most-one-generic}
Montrer que, dans un espace $T_0$ $X$, toute partie fermée
irréductible admet au plus un point générique.
\end{exercise}

\begin{exercise}
\label{exercises-exercise-spec-sober}
Montrer que, dans $\Spec(A)$, toute
partie fermée irréductible {\it admet bien} un point générique.
Montrer en fait que l'application
${\mathfrak p} \mapsto \overline{\{{\mathfrak p}\}}$ est
une bijection de $\Spec(A)$ sur l'ensemble des parties fermées
irréductibles de $X$.
\end{exercise}

\begin{exercise}
\label{exercises-exercise-irreducible-subset-not-generic}
Donner un exemple montrant qu'une partie irréductible de
$\Spec(\mathbf{Z})$ ne possède pas nécessairement de point générique.
\end{exercise}

\begin{definition}
\label{exercises-definition-Noetherian-space}
Un espace topologique $X$ est dit {\it noethérien} si toute
suite décroissante $Z_1\supset Z_2 \supset Z_3\supset \ldots$
de parties fermées de $X$ est stationnaire.
(Il est dit {\it artinien} si toute suite croissante de parties
fermées est stationnaire.)
\end{definition}

\begin{exercise}
\label{exercises-exercise-Noetherian-spec}
Montrer que, si l'anneau $A$ est noethérien, alors
l'espace topologique $\Spec(A)$ est noethérien. Donner un
exemple montrant que la réciproque est fausse. (Même question dans le cas
artinien, si vous le souhaitez.)
\end{exercise}

\begin{definition}
\label{exercises-definition-irreducible-component}
Une partie irréductible maximale $T\subset X$ est appelée une
{\it composante irréductible} de l'espace $X$. Une telle composante
irréductible de $X$ est automatiquement une partie fermée de $X$.
\end{definition}

\begin{exercise}
\label{exercises-exercise-irreducible-in-irreducible}
Montrer que toute partie irréductible
de $X$ est contenue dans une composante irréductible de $X$.
\end{exercise}

\begin{exercise}
\label{exercises-exercise-Noetherian-finite-nr-irreducible}
Montrer qu'un espace topologique noethérien $X$
n'a qu'un nombre fini de composantes irréductibles, disons $X_1, \ldots, X_n$,
et que $X = X_1\cup X_2\cup\ldots\cup X_n$. (Noter que
tout $X$ est toujours la réunion de ses composantes irréductibles, mais que,
si $X = {\mathbf R}$ est par exemple muni de sa topologie usuelle, les composantes irréductibles
de $X$ sont les parties réduites à un point. Ce n'est guère
intéressant.)
\end{exercise}

\begin{exercise}
\label{exercises-exercise-irreducible-components-minimal-primes}
Montrer que les composantes irréductibles de $\Spec(A)$
correspondent aux idéaux premiers minimaux de $A$.
\end{exercise}

\begin{definition}
\label{exercises-definition-closed}
Un point $x\in X$ est dit {\it fermé} si $\overline{\{x\}} = \{ x\}$.
Soient $x, y$ des points de $X$. On dit que $x$ est une {\it spécialisation}
de $y$, ou que $y$ est une {\it générisation} de $x$, si
$x\in \overline{\{y\}}$.
\end{definition}

\begin{exercise}
\label{exercises-exercise-closed-maximal}
Montrer que les points fermés de $\Spec(A)$
correspondent aux idéaux maximaux de $A$.
\end{exercise}

\begin{exercise}
\label{exercises-exercise-generalization}
Montrer que ${\mathfrak p}$ est une générisation de ${\mathfrak q}$
dans $\Spec(A)$ si et seulement si ${\mathfrak p}\subset {\mathfrak q}$.
Caractériser les points fermés,
les idéaux maximaux, les points génériques et les idéaux premiers minimaux en termes de
générisation et de spécialisation. (On utilise ici la terminologie selon laquelle un point
d'un espace topologique éventuellement réductible $X$ est appelé point générique
s'il est un point générique de l'une des composantes irréductibles de $X$.)
\end{exercise}

\begin{exercise}
\label{exercises-exercise-disjoint-closed-spec}
Soient $I$ et $J$ des idéaux de $A$.
À quelle condition $V(I)$ et $V(J)$ sont-ils disjoints ?
\end{exercise}

\begin{definition}
\label{exercises-definition-connected-component}
Un espace topologique $X$ est dit {\it connexe} s'il est non vide et n'est pas la
réunion de deux ouverts non vides et disjoints. Une {\it composante connexe}
de $X$ est une partie connexe maximale. Tout point de $X$ est contenu
dans une composante connexe de $X$, et toute composante connexe de $X$ est
fermée dans $X$. (Mais, en général, une composante connexe n'est pas nécessairement ouverte dans $X$.)
\end{definition}

\begin{exercise}
\label{exercises-exercise-disconnected-spec}
Soit $A$ un anneau non nul.
Montrer que $\Spec(A)$ est non connexe
si et seulement si $A\cong B \times C$ pour certains anneaux non nuls $B, C$.
\end{exercise}

\begin{exercise}
\label{exercises-exercise-connected-component-stable-generalization}
Soit $T$ une composante connexe
de $\Spec(A)$. Montrer que $T$ est stable par générisation.
Montrer que $T$ est une partie ouverte de $\Spec(A)$ si $A$ est noethérien.
(Remarque : cette assertion est fausse lorsque $A$ est, par exemple, un produit infini de copies de
${\mathbf F}_2$. Le spectre de cet anneau est constitué d'une infinité
de points fermés.)
\end{exercise}

\begin{exercise}
\label{exercises-exercise-primes-kx}
Déterminer $\Spec(k[x])$, c'est-à-dire décrire
les idéaux premiers de cet anneau, décrire les spécialisations possibles et
décrire la topologie. (Traiter le cas où $k$ est algébriquement clos, mais
aussi celui où $k$ ne l'est pas.)
\end{exercise}

\begin{exercise}
\label{exercises-exercise-primes-kxy}
 Déterminer $\Spec(k[x, y])$, où $k$ est algébriquement
clos.
[Indication : utiliser le morphisme
$\varphi : \Spec(k[x, y]) \to \Spec(k[x])$ ; si
$\varphi({\mathfrak p}) = (0)$, alors localiser suivant
$S = \{f\in k[x] \mid f \not = 0\}$ et utiliser le résultat du cours
sur la localisation et $\Spec$.]
(Pourquoi, à votre avis, les géomètres algébristes appellent-ils cet espace l'espace affine de dimension 2 ?)
\end{exercise}

\begin{exercise}
\label{exercises-exercise-primes-Zy}
Déterminer $\Spec(\mathbf{Z}[y])$.
[Indication : procéder comme ci-dessus.] (Espace affine de dimension 1 sur $\mathbf{Z}$.)
\end{exercise}






\section{Localisation}
\label{exercises-section-localization}


\begin{exercise}
\label{exercises-exercise-submodule-localization}
Soit $A$ un anneau. Soit $S \subset A$ une partie multiplicative.
Soit $M$ un $A$-module. Soit $N \subset S^{-1}M$ un $S^{-1}A$-sous-module.
Montrer qu'il existe un $A$-sous-module $N' \subset M$ tel que
$N = S^{-1}N'$. (Ce résultat utile s'applique en particulier aux idéaux
de $S^{-1}A$.)
\end{exercise}

\begin{exercise}
\label{exercises-exercise-localize-zero}
Soit $A$ un anneau. Soit $M$ un $A$-module. Soit $m \in M$.
\begin{enumerate}
\item Montrer que $I = \{a \in A \mid am = 0\}$ est un idéal de $A$.
\item Pour un idéal premier $\mathfrak p$ de $A$, montrer que l'image de $m$
dans $M_\mathfrak p$ est nulle si et seulement si $I \not \subset \mathfrak p$.
\item Montrer que $m$ est nul si et seulement si l'image de $m$ est nulle
dans $M_\mathfrak p$ pour tout idéal premier $\mathfrak p$ de $A$.
\item Montrer que $m$ est nul si et seulement si l'image de $m$ est nulle
dans $M_\mathfrak m$ pour tout idéal maximal $\mathfrak m$ de $A$.
\item Montrer que $M = 0$ si et seulement si $M_{\mathfrak m}$ est nul
pour tout idéal maximal $\mathfrak m$.
\end{enumerate}
\end{exercise}

\begin{exercise}
\label{exercises-exercise-localization-is-field}
Trouver un couple $(A, f)$ où $A$ est un anneau intègre possédant au moins trois idéaux
premiers deux à deux distincts et où $f \in A$ est tel que le localisé
principal $A_f = \{1, f, f^2, \ldots \}^{-1}A$ soit un corps.
\end{exercise}

\begin{exercise}
\label{exercises-exercise-localize-finite-module-zero}
Soit $A$ un anneau. Soit $M$ un $A$-module de type fini. Soit $S \subset A$
une partie multiplicative. Supposons que $S^{-1}M = 0$. Montrer qu'il
existe $f \in S$ tel que le localisé principal
$M_f = \{1, f, f^2, \ldots \}^{-1}M$ soit nul.
\end{exercise}
% Borne de la tranche : la ligne source 750 est vide.
\begin{exercise}
\label{exercises-exercise-localization-is-quotient}
Donner un exemple de triplet $(A, I, S)$ où $A$ est un anneau,
$0 \not = I \not = A$ est un idéal propre non nul et $S \subset A$
est une partie multiplicative
telle que $A/I \cong S^{-1}A$ comme $A$-algèbres.
\end{exercise}




\section{Lemme de Nakayama}
\label{exercises-section-nakayama}

\begin{exercise}
\label{exercises-exercise-nakayama}
Soit $A$ un anneau.
Soit $I$ un idéal de $A$.
Soit $M$ un $A$-module.
Soient $x_1, \ldots, x_n \in M$.
Supposons satisfaites les conditions suivantes :
\begin{enumerate}
\item $M/IM$ est engendré par $x_1, \ldots, x_n$,
\item $M$ est un $A$-module de type fini,
\item $I$ est contenu dans tout idéal maximal de $A$.
\end{enumerate}
Montrer que $x_1, \ldots, x_n$ engendrent $M$. (Solution suggérée :
se ramener, par localisation, au cas d'un idéal maximal de $A$ à l'aide de
l'exercice \ref{exercises-exercise-localize-zero} et de l'exactitude de la localisation.
Se ramener ensuite à l'énoncé du lemme de Nakayama donné dans le cours
en considérant le quotient de $M$ par le sous-module engendré par
$x_1, \ldots, x_n$.)
\end{exercise}






\section{Longueur}
\label{exercises-section-length}

\begin{definition}
\label{exercises-definition-length}
Soient $A$ un anneau et $M$ un $A$-module. La
{\it longueur} de $M$ comme $A$-module est
$$
\text{longueur}_A(M)
=
\sup
\{
n
\mid
\exists\ 0 = M_0 \subset M_1 \subset \ldots \subset M_n = M,
\text{ }M_i \not = M_{i + 1}
\}.
$$
Autrement dit, c'est la borne supérieure des longueurs des chaînes de sous-modules.
\end{definition}

\begin{exercise}
\label{exercises-exercise-length-is-one}
Montrer qu'un module $M$ sur un anneau $A$ est de longueur $1$ si et seulement s'il
est isomorphe à $A/\mathfrak m$ pour un idéal maximal $\mathfrak m$
de $A$.
\end{exercise}

\begin{exercise}
\label{exercises-exercise-length-easy}
Calculer la longueur des modules suivants sur les anneaux indiqués.
Expliquer brièvement (!) chaque réponse. (On pourra utiliser librement l'additivité de
la fonction longueur dans les suites exactes courtes ; voir
Algèbre, lemme \ref{algebra-lemma-length-additive}.)
\begin{enumerate}
\item La longueur de $\mathbf{Z}/120\mathbf{Z}$ sur $\mathbf{Z}$.
\item La longueur de $\mathbf{C}[x]/(x^{100} + x + 1)$ sur $\mathbf{C}[x]$.
\item La longueur de $\mathbf{R}[x]/(x^4 + 2x^2 + 1)$ sur $\mathbf{R}[x]$.
\end{enumerate}
\end{exercise}

\begin{exercise}
\label{exercises-exercise-compute-length}
Soit $A = k[x, y]_{(x, y)}$ l'anneau local du plan affine en
l'origine. On pourra imposer au corps $k$ toute hypothèse que l'on souhaite. Supposons
que $f = x^3 + x^2y^2 + y^{100}$ et $g = y^3 - x^{999}$. Quelle est la longueur
de $A/(f, g)$ comme $A$-module ? (Une démarche possible consiste à considérer
l'idéal engendré par $f$ et $g$ dans les quotients de la forme $A/{\mathfrak m}_A^n=
k[x, y]/(x, y)^n$, où $n$ varie. Essayer de trouver $n$ tel que
$A/(f, g)+{\mathfrak m}_A^n \cong A/(f, g)+{\mathfrak m}_A^{n + 1}$
et utiliser NAK.)
\end{exercise}






\section{Idéaux premiers associés}
\label{exercises-section-ass}

\noindent
Les idéaux premiers associés sont étudiés dans
Algèbre, section \ref{algebra-section-ass}.

\begin{exercise}
\label{exercises-exercise-compute-ass}
Calculer l'ensemble des idéaux premiers associés à chacun des modules suivants.
\begin{enumerate}
\item $R = k[x, y]$ et $M = R/(xy(x + y))$,
\item $R = \mathbf{Z}[x]$ et $M = R/(300x + 75)$, et
\item $R = k[x, y, z]$ et $M = R/(x^3, x^2y, xz)$.
\end{enumerate}
Comme d'habitude, $k$ désigne un corps.
\end{exercise}

\begin{exercise}
\label{exercises-exercise-prime-power-not-primary}
Donner un exemple d'anneau noethérien $R$ et d'idéal premier $\mathfrak p$
tel que $\mathfrak p$ ne soit pas le seul idéal premier associé à $R/\mathfrak p^2$.
\end{exercise}

\begin{exercise}
\label{exercises-exercise-product-primes-not-only-associated}
Soit $R$ un anneau noethérien possédant deux idéaux premiers incomparables
$\mathfrak p$, $\mathfrak q$, c'est-à-dire $\mathfrak p \not \subset \mathfrak q$
et $\mathfrak q \not \subset \mathfrak p$.
\begin{enumerate}
\item Montrer que, pour $N = R/(\mathfrak p \cap \mathfrak q)$,
on a $\text{Ass}(N) = \{\mathfrak p, \mathfrak q\}$.
\item Montrer par un exemple que le module
$M = R/\mathfrak p \mathfrak q$ peut
avoir un idéal premier associé distinct de $\mathfrak p$ et de $\mathfrak q$.
\end{enumerate}
\end{exercise}





\section{Groupes d'Ext}
\label{exercises-section-ext}

\noindent
Les groupes d'Ext sont définis dans Algèbre, section \ref{algebra-section-ext}.

\begin{exercise}
\label{exercises-exercise-compute-ext-abelian-groups}
Calculer tous les groupes d'Ext $\Ext^i(M, N)$ des modules donnés
dans la catégorie des $\mathbf{Z}$-modules
(appelée aussi catégorie des groupes abéliens).
\begin{enumerate}
\item $M = \mathbf{Z}$ et $N = \mathbf{Z}$,
\item $M = \mathbf{Z}/4\mathbf{Z}$ et $N = \mathbf{Z}/8\mathbf{Z}$,
\item $M = \mathbf{Q}$ et $N = \mathbf{Z}/2\mathbf{Z}$, et
\item $M = \mathbf{Z}/2\mathbf{Z}$ et $N = \mathbf{Q}/\mathbf{Z}$.
\end{enumerate}
\end{exercise}

\begin{exercise}
\label{exercises-exercise-compute-in-regular}
Soit $R = k[x, y]$, où $k$ est un corps.
\begin{enumerate}
\item Montrer directement que le complexe de Koszul
$$
0 \to R
\xrightarrow{
\left(
\begin{matrix}
y \\
-x
\end{matrix}
\right)
}
R^{\oplus 2} \xrightarrow{(x, y)} R \xrightarrow{f \mapsto f(0, 0)} k \to 0
$$
est exact.
\item Calculer $\Ext^i_R(k, k)$, où $k = R/(x, y)$ est considéré comme un $R$-module.
\end{enumerate}
\end{exercise}

\begin{exercise}
\label{exercises-exercise-infinitely-many-nonzero-ext}
Donner un exemple d'anneau noethérien $R$ et de modules de type fini $M$, $N$
tels que $\Ext^i_R(M, N)$ soit non nul pour tout $i \geq 0$.
\end{exercise}

\begin{exercise}
\label{exercises-exercise-infinite-ext}
Donner un exemple d'anneau $R$ et d'idéal $I$ tels que
$\Ext^1_R(R/I, R/I)$ ne soit pas un $R$-module de type fini.
(On sait que cela ne peut pas se produire si $R$ est noethérien, d'après
Algèbre, lemme \ref{algebra-lemma-ext-noetherian}.)
\end{exercise}







\section{Profondeur}
\label{exercises-section-depth}

\noindent
La profondeur est définie dans Algèbre, section \ref{algebra-section-depth},
et étudiée plus avant dans
Complexes dualisants, section \ref{dualizing-section-depth}.

\begin{exercise}
\label{exercises-exercise-compute-depth}
Soient $R$ un anneau, $I \subset R$ un idéal et $M$ un $R$-module.
Calculer $\text{prof}_I(M)$ dans les cas suivants.
\begin{enumerate}
\item $R = \mathbf{Z}$, $I = (30)$, $M = \mathbf{Z}$,
\item $R = \mathbf{Z}$, $I = (30)$, $M = \mathbf{Z}/(300)$,
\item $R = \mathbf{Z}$, $I = (30)$, $M = \mathbf{Z}/(7)$,
\item $R = k[x, y, z]/(x^2 + y^2 + z^2)$, $I = (x, y, z)$, $M = R$,
\item $R = k[x, y, z, w]/(xz, xw, yz, yw)$, $I = (x, y, z, w)$,
$M = R$.
\end{enumerate}
Ici, $k$ est un corps. Dans les deux derniers cas, on pourra localiser
en l'idéal maximal $I$.
\end{exercise}

\begin{exercise}
\label{exercises-exercise-depth-not-inherited-localization}
Donner un exemple d'anneau local noethérien $(R, \mathfrak m, \kappa)$
de profondeur $\geq 1$ et d'idéal premier $\mathfrak p$ tels que
\begin{enumerate}
\item $\text{prof}_\mathfrak m(R) \geq 1$,
\item $\text{prof}_\mathfrak p(R_\mathfrak p) = 0$, et
\item $\dim(R_\mathfrak p) \geq 1$.
\end{enumerate}
Si l'on ne demande pas (3), l'exercice est trop facile. Pourquoi ?
\end{exercise}

\begin{exercise}
\label{exercises-exercise-depth-torsion-free}
Soit $(R, \mathfrak m)$ un anneau local noethérien intègre.
Soit $M$ un $R$-module de type fini.
\begin{enumerate}
\item Si $M$ est sans torsion, montrer que $M$ est de profondeur au moins $1$ sur $R$.
\item Donner un exemple où la profondeur vaut $1$.
\end{enumerate}
\end{exercise}

\begin{exercise}
\label{exercises-exercise-depth-examples}
Pour tous $m \geq n \geq 0$, donner un exemple d'anneau local
noethérien $R$ tel que $\dim(R) = m$ et $\text{prof}(R) = n$.
\end{exercise}

\begin{exercise}
\label{exercises-exercise-make-depth-1}
Soit $(R, \mathfrak m)$ un anneau local noethérien.
Soit $M$ un $R$-module de type fini.
Montrer qu'il existe une suite exacte courte canonique
$$
0 \to K \to M \to Q \to 0
$$
ayant les propriétés suivantes :
\begin{enumerate}
\item $\text{prof}(Q) \geq 1$,
\item $K$ est nul ou $\text{Supp}(K) = \{\mathfrak m\}$, et
\item $\text{longueur}_R(K) < \infty$.
\end{enumerate}
Indication : en utilisant la propriété noethérienne, montrer qu'il
existe un sous-module $K$ maximal parmi ceux qui vérifient (2),
puis que $Q = M/K$ vérifie (1) et que $K$ vérifie (3).
\end{exercise}

\begin{exercise}
\label{exercises-exercise-make-depth-2}
Soit $(R, \mathfrak m)$ un anneau local noethérien.
Soit $M$ un $R$-module de type fini et de profondeur $\geq 2$.
Soit $N \subset M$ un sous-module non nul.
\begin{enumerate}
\item Montrer que $\text{prof}(N) \geq 1$.
\item Montrer que $\text{prof}(N) = 1$ si et seulement si le
module quotient $M/N$ vérifie $\text{prof}(M/N) = 0$.
\item Montrer qu'il existe un sous-module $N' \subset M$ tel que
$N \subset N'$ et que le quotient soit de longueur finie, c'est-à-dire que
$\text{longueur}_R(N'/N) < \infty$, et tel que $N'$ soit
de profondeur $\geq 2$. Indication : appliquer l'exercice \ref{exercises-exercise-make-depth-1} à $M/N$
et prendre pour $N'$ l'image inverse de $K$.
\end{enumerate}
\end{exercise}

\begin{exercise}
\label{exercises-exercise-Hartshorne-reduced}
Soit $(R, \mathfrak m)$ un anneau local noethérien.
Supposons $R$ réduit, c'est-à-dire que $R$ ne possède aucun élément
nilpotent non nul. Supposons en outre que $R$ possède deux
idéaux premiers minimaux distincts $\mathfrak p$ et $\mathfrak q$.
\begin{enumerate}
\item Montrer que la suite de $R$-modules
$$
0 \to R \to R/\mathfrak p \oplus R/\mathfrak q \to
R/\mathfrak p + \mathfrak q \to 0
$$
est exacte (le vérifier en chaque terme). Les applications sont
$x \mapsto (x \bmod \mathfrak p, x \bmod \mathfrak q)$
et $(y \bmod \mathfrak p, z \bmod \mathfrak q) \mapsto
(y - z \bmod \mathfrak p + \mathfrak q)$.
\item Montrer que si $\text{prof}(R) \geq 2$, alors
$\dim(R/\mathfrak p + \mathfrak q) \geq 1$.
\item Montrer que si $\text{prof}(R) \geq 2$, alors
$U = \Spec(R) \setminus \{\mathfrak m\}$ est un espace topologique
connexe.
\end{enumerate}
Cela démontre un cas très particulier du théorème de connexité de Hartshorne,
selon lequel le spectre épointé $U$ d'un anneau local
noethérien de $\text{prof} \geq 2$ est connexe.
\end{exercise}

\begin{exercise}
\label{exercises-exercise-depth-2}
Soit $(R, \mathfrak m)$ un anneau local noethérien.
Supposons que $x, y \in \mathfrak m$ forment une suite régulière de longueur $2$.
Pour tout $n \geq 2$, montrer qu'il n'existe pas de $a, b \in R$ tels que
$$
x^{n - 1}y^{n - 1} = a x^n + b y^n
$$
Indication : commencer par traiter le cas $n = 2$ afin de voir comment procéder.
Remarque : ce résultat admet une vaste généralisation
appelée la conjecture des monômes.
\end{exercise}







\section{Modules et anneaux de Cohen-Macaulay}
\label{exercises-section-CM}

\noindent
Les modules de Cohen-Macaulay sont étudiés dans
Algèbre, section \ref{algebra-section-CM},
et les anneaux de Cohen-Macaulay sont étudiés dans
Algèbre, section \ref{algebra-section-CM-ring}.

\begin{exercise}
\label{exercises-exercise-examples-CM}
Dans chacun des cas suivants, répondre par oui ou par non.
Aucune explication ni démonstration n'est nécessaire.
\begin{enumerate}
\item Soit $p$ un nombre premier.
L'anneau local $\mathbf{Z}_{(p)}$ est-il un anneau local de Cohen-Macaulay ?
\item Soit $p$ un nombre premier.
L'anneau local $\mathbf{Z}_{(p)}$ est-il un anneau local régulier ?
\item Soit $k$ un corps.
L'anneau local $k[x]_{(x)}$ est-il un anneau local de Cohen-Macaulay ?
\item Soit $k$ un corps.
L'anneau local $k[x]_{(x)}$ est-il un anneau local régulier ?
\item Soit $k$ un corps.
L'anneau local $(k[x, y]/(y^2 - x^3))_{(x, y)} =
k[x, y]_{(x, y)}/(y^2 - x^3)$ est-il un anneau local de Cohen-Macaulay ?
\item Soit $k$ un corps.
L'anneau local $(k[x, y]/(y^2, xy))_{(x, y)} =
k[x, y]_{(x, y)}/(y^2, xy)$ est-il un anneau local de Cohen-Macaulay ?
\end{enumerate}
\end{exercise}













\section{Singularités}
\label{exercises-section-singularities}

\begin{exercise}
\label{exercises-exercise-singularities}
Soit $k$ un corps quelconque. Supposons que $A = k[[x, y]]/(f)$ et
$B = k[[u, v]]/(g)$, où $f = xy$ et $g = uv + \delta$ avec
$\delta \in (u, v)^3$. Montrer que $A$ et $B$ sont des anneaux isomorphes.
\end{exercise}

\begin{remark}
\label{exercises-remark-singularities}
Une singularité d'une courbe sur un corps $k$ est appelée
point double ordinaire si l'anneau local complété de la courbe au
point est de la forme $k'[[x, y]]/(f)$, où (a) $k'$ est une extension finie
séparable de $k$, (b) le terme initial de $f$ est de degré deux, c'est-à-dire qu'il
est de la forme $q = ax^2 + bxy + cy^2$ pour des $a, b, c\in k'$ non tous nuls, et
(c) $q$ est une forme quadratique non dégénérée sur $k'$ (en caractéristique 2,
cela signifie que $b$ est non nul). En général, il y a une classe d'isomorphisme
de tels anneaux pour chaque classe d'isomorphisme de couples $(k', q)$.
\end{remark}

\begin{exercise}
\label{exercises-exercise-periodic-resolution}
Soit $R$ un anneau. Soit $n \geq 1$. Soient $A$, $B$ des matrices $n \times n$ à
coefficients dans $R$ telles que $AB = f 1_{n \times n}$
pour un non-diviseur de zéro $f$ de $R$. Posons $S = R/(f)$. Montrer que
$$
\ldots \to
S^{\oplus n} \xrightarrow{B}
S^{\oplus n} \xrightarrow{A}
S^{\oplus n} \xrightarrow{B}
S^{\oplus n} \to \ldots
$$
est exacte.
\end{exercise}







\section{Ensembles constructibles}
\label{exercises-section-constructible}

\noindent
Soit $k$ un corps algébriquement clos, par exemple le corps $\mathbf{C}$
des nombres complexes. Soit $n \geq 0$. Un polynôme $f \in k[x_1, \ldots, x_n]$
définit par évaluation une fonction $f : k^n \to k$. Un sous-ensemble $Z \subset k^n$
est appelé un {\it ensemble algébrique} s'il est l'ensemble des zéros communs d'une
famille de polynômes.

\begin{exercise}
\label{exercises-exercise-finite-nr-equations}
Démontrer qu'un ensemble algébrique peut toujours s'écrire comme l'ensemble des zéros
d'un nombre fini de polynômes.
\end{exercise}

\noindent
Avec les notations ci-dessus, un sous-ensemble $E \subset k^n$ est dit
{\it constructible} s'il est une réunion finie d'ensembles de la forme
$Z \cap \{f \not = 0\}$, où $f$ est un polynôme.

\begin{exercise}
\label{exercises-exercise-constructible-classical}
Montrer les assertions suivantes :
\begin{enumerate}
\item le complémentaire d'un ensemble constructible est constructible,
\item une réunion finie d'ensembles constructibles est constructible,
\item une intersection finie d'ensembles constructibles est constructible, et
\item tout ensemble constructible $E$ peut s'écrire comme une réunion disjointe finie
$E = \coprod E_i$, chaque $E_i$ étant de la forme $Z \cap \{f \not = 0\}$,
où $Z$ est un ensemble algébrique et $f$ un polynôme.
\end{enumerate}
\end{exercise}

\begin{exercise}
\label{exercises-exercise-division-with-remainder}
Soit $R$ un anneau. Soit
$f = a_d x^d + a_{d - 1} x^{d - 1} + \ldots + a_0 \in R[x]$.
(Comme d'habitude, cette notation signifie que $a_0, \ldots, a_d \in R$.) Soit $g \in R[x]$.
Démontrer qu'il existe $N \geq 0$ et $r, q \in R[x]$ tels que
$$
a_d^N g = q f + r
$$
avec $\deg(r) < d$, c'est-à-dire que, pour certains $c_i \in R$, on a
$r = c_0 + c_1 x + \ldots + c_{d - 1}x^{d - 1}$.
\end{exercise}







\section{Nullstellensatz de Hilbert}
\label{exercises-section-Hilbert-Nullstellensatz}


\begin{exercise}
\label{exercises-exercise-uncountable}
{\it Un argument saugrenu utilisant les nombres complexes !}
Soit ${\mathbf C}$ le corps des nombres complexes. Soit $V$ un espace vectoriel
sur ${\mathbf C}$. Le spectre d'un opérateur linéaire
$T : V \to V$ est l'ensemble des nombres complexes $\lambda \in {\mathbf C}$
tels que l'opérateur $T - \lambda \text{id}_V$ ne soit pas inversible.
\begin{enumerate}
\item Montrer que $\mathbf{C}(X)$
est de dimension non dénombrable sur ${\mathbf C}$.
\item Montrer que tout opérateur linéaire sur $V$ possède un
spectre non vide si la dimension de $V$ est finie ou
dénombrable.
\item Montrer que si une ${\mathbf C}$-algèbre de type fini
$R$ est un corps, alors l'application ${\mathbf C}\to R$ est un isomorphisme.
\item Montrer que tout idéal maximal ${\mathfrak m}$ de
${\mathbf C}[x_1, \ldots, x_n]$ est de la forme
$(x_1-\alpha_1, \ldots, x_n-\alpha_n)$ pour certains $\alpha_i \in {\mathbf C}$.
\end{enumerate}
\end{exercise}

\begin{remark}
\label{exercises-remark-HNSS}
Soit $k$ un corps. Alors, pour tout entier $n\in {\mathbf N}$ et
tout idéal maximal ${\mathfrak m} \subset k[x_1, \ldots, x_n]$,
le quotient $k[x_1, \ldots, x_n]/{\mathfrak m}$ est une extension finie
de $k$. Ce résultat sera démontré plus loin dans le cours. Bien entendu
(le vérifier), il entraîne un énoncé analogue pour les idéaux maximaux
des $k$-algèbres de type fini. L'exercice précédent le démontre
dans le cas $k = {\mathbf C}$.
\end{remark}

\begin{exercise}
\label{exercises-exercise-Hilbert-Nullstellensatz}
Soit $k$ un corps. Utiliser la remarque \ref{exercises-remark-HNSS}.
\begin{enumerate}
\item Soit $R$ une $k$-algèbre. Supposons que $\dim_k R < \infty$
et que $R$ soit intègre. Montrer que $R$ est un corps.
\item Supposons que $R$ soit une $k$-algèbre de type fini et que
$f\in R$ ne soit pas nilpotent. Montrer qu'il existe un idéal maximal
${\mathfrak m} \subset R$ tel que $f\not\in {\mathfrak m}$.
\item Montrer par un exemple que cette assertion est fausse lorsque $R$
n'est pas de type fini sur un corps.
\item Montrer que tout idéal radical $I \subset {\mathbf C}[x_1, \ldots, x_n]$
est l'intersection des idéaux maximaux qui le contiennent.
\end{enumerate}
\end{exercise}

\begin{remark}
\label{exercises-remark-Hilbert-Nullstellensatz}
C'est le Nullstellensatz de Hilbert. Il affirme en effet
que les parties fermées de $\Spec(k[x_1, \ldots, x_n])$
(qui correspondent aux idéaux radicaux d'après un exercice précédent)
sont déterminées par les points fermés qu'elles contiennent.
\end{remark}

\begin{exercise}
\label{exercises-exercise-product-matrices-ring}
Soit $A =
{\mathbf C}[x_{11}, x_{12}, x_{21}, x_{22}, y_{11}, y_{12}, y_{21}, y_{22}]$.
Soit $I$ l'idéal de $A$ engendré par les coefficients de la
matrice $XY$, où
$$
X = \left(
\begin{matrix}
x_{11} & x_{12}\\
x_{21} & x_{22}
\end{matrix}
\right)
\quad\text{et}\quad
Y = \left(
\begin{matrix}
y_{11} & y_{12}\\
y_{21} & y_{22}
\end{matrix}
\right).
$$
Déterminer les composantes irréductibles du sous-ensemble fermé $V(I)$ de
$\Spec(A)$.
(Il s'agit de les décrire et de donner des équations pour chacune. Il n'est pas
nécessaire de démontrer que les équations écrites définissent des idéaux premiers.) Indications :
\begin{enumerate}
\item On peut utiliser le Nullstellensatz de Hilbert, et il suffit de déterminer des
sous-ensembles localement fermés irréductibles qui recouvrent l'ensemble des points fermés de
$V(I)$.
\item Il y a deux composantes faciles à trouver.
\item L'image d'un ensemble irréductible par une application continue est
irréductible.
\end{enumerate}
\end{exercise}




\section{Dimension}
\label{exercises-section-dimension}

\begin{exercise}
\label{exercises-exercise-dimension-bigger-one-finite-nr-primes}
Construire un anneau $A$ de dimension $> 1$ qui n'a qu'un nombre fini d'idéaux premiers.
\end{exercise}

\begin{exercise}
\label{exercises-exercise-hypersurface-in-A2-dimension-one}
Soit $f \in \mathbf{C}[x, y]$ un polynôme non constant.
Montrer que $\mathbf{C}[x, y]/(f)$ est de dimension $1$.
\end{exercise}

\begin{exercise}
\label{exercises-exercise-dimension-polynomial-ring}
Soit $(R, \mathfrak m)$ un anneau local noethérien.
Soit $n \geq 1$. Soit $\mathfrak m' = (\mathfrak m, x_1, \ldots, x_n)$
dans l'anneau de polynômes $R[x_1, \ldots, x_n]$.
Montrer que
$$
\dim(R[x_1, \ldots, x_n]_{\mathfrak m'}) = \dim(R) + n.
$$
\end{exercise}





\section{Anneaux caténaires}
\label{exercises-section-catenary}

\begin{definition}
\label{exercises-definition-catenary}
On dit qu'un anneau noethérien $A$ est {\it caténaire}
si, pour tout triplet d'idéaux premiers
${\mathfrak p}_1 \subset {\mathfrak p}_2 \subset {\mathfrak p}_3$,
on a
$$
ht({\mathfrak p}_3 / {\mathfrak p}_1) = ht({\mathfrak p}_3/{\mathfrak p}_2) +
ht({\mathfrak p}_2/{\mathfrak p}_1).
$$
Ici, $ht(\mathfrak p/\mathfrak q)$ désigne la hauteur de
$\mathfrak p/\mathfrak q$ dans l'anneau $A/\mathfrak q$.
Autrement dit,
$$
ht(\mathfrak p/\mathfrak q) =
\dim(A_\mathfrak p/\mathfrak qA_\mathfrak p) =
\dim((A/\mathfrak q)_\mathfrak p) =
\dim((A/\mathfrak q)_{\mathfrak p/\mathfrak q})
$$
Un espace topologique $X$ est {\it caténaire} si, étant donnés $T \subset T' \subset X$
avec $T$ et $T'$ fermés et irréductibles, il existe une chaîne maximale
de parties fermées irréductibles
$$
T = T_0 \subset T_1 \subset \ldots \subset T_n = T'
$$
et toutes ces chaînes ont la même longueur (finie).
\end{definition}

\begin{exercise}
\label{exercises-exercise-catenary-the-same}
Montrer que la notion d'anneau caténaire définie dans
Algèbre, définition \ref{algebra-definition-catenary}
coïncide avec celle de la définition \ref{exercises-definition-catenary}
pour les anneaux noethériens.
\end{exercise}

\begin{exercise}
\label{exercises-exercise-Noetherian-local-domain-dim-2-catenary}
Montrer qu'un anneau local noethérien intègre de dimension $2$ est caténaire.
\end{exercise}

\begin{exercise}
\label{exercises-exercise-finite-type-over-field-catenary}
Soit $k$ un corps.
Montrer qu'une $k$-algèbre de type fini est caténaire.
\end{exercise}

\begin{exercise}
\label{exercises-exercise-example-no-dim-function}
Donner un exemple d'espace topologique fini, sobre et caténaire
$X$ qui n'admette aucune fonction de dimension $\delta : X \to \mathbf{Z}$.
Ici, $\delta : X \to \mathbf{Z}$ est une fonction de dimension si,
pour $x, y \in X$, on a
\begin{enumerate}
\item $x \leadsto y$ et $x \not = y$ implique $\delta(x) > \delta(y)$,
\item $x \leadsto y$ et $\delta(x) \geq \delta(y) + 2$ implique
qu'il existe $z \in X$ tel que $x \leadsto z \leadsto y$ et
$\delta(x) > \delta(z) > \delta(y)$.
\end{enumerate}
Décrire clairement cet espace et expliquer brièvement pourquoi il
ne peut admettre de fonction de dimension.
\end{exercise}




\section{Corps des fractions}
\label{exercises-section-fraction-fields}

\begin{exercise}
\label{exercises-exercise-find-fraction-field}
Considérer l'anneau intègre
$$
{\mathbf Q}[r, s, t]/(s^2-(r-1)(r-2)(r-3), t^2-(r + 1)(r + 2)(r + 3)).
$$
Trouver un anneau intègre de la forme ${\mathbf Q}[x, y]/(f)$ dont le
corps des fractions soit isomorphe à celui de l'anneau ci-dessus.
\end{exercise}



\section{Degré de transcendance}
\label{exercises-section-transcendence}

\begin{exercise}
\label{exercises-exercise-algebraic-extension}
Soient $K'/K/k$ des extensions de corps telles que $K'$ soit
algébrique sur $K$. Montrer que $\text{deg.tr.}_k(K) = \text{deg.tr.}_k(K')$.
(Indication : montrer que si $x_1, \ldots, x_d \in K$ sont algébriquement indépendants
sur $k$ et si $d < \text{deg.tr.}_k(K')$, alors l'extension $k(x_1, \ldots, x_d) \subset K$
ne peut être algébrique.)
\end{exercise}

\begin{exercise}
\label{exercises-exercise-growth-powers-subvector-space}
Soient $k$ un corps et $K/k$ une extension de type fini
de degré de transcendance $d$.
Si $V, W \subset K$ sont des sous-espaces vectoriels sur $k$ de dimension finie,
posons
$$
VW = \{f \in K \mid f = \sum\nolimits_{i = 1, \ldots, n} v_i w_i
\text{ pour un certain }n\text{ et }v_i \in V, w_i \in W\}
$$
C'est un sous-espace vectoriel sur $k$ de dimension finie.
Posons $V^2 = VV$, $V^3 = V V^2$, etc.
\begin{enumerate}
\item Montrer qu'on peut trouver $V \subset K$ et $\epsilon > 0$ tels que
$\dim V^n \geq \epsilon n^d$ pour tout $n \geq 1$.
\item Réciproquement, montrer que, pour tout sous-espace $V \subset K$ de dimension finie,
il existe $C > 0$ tel que $\dim V^n \leq C n^d$ pour tout $n \geq 1$.
(Une démarche possible consiste à procéder comme suit :
établir d'abord le résultat pour les sous-espaces vectoriels de $k[x_1, \ldots, x_d]$,
puis pour ceux de $k(x_1, \ldots, x_d)$.
Enfin, si $K/k(x_1, \ldots, x_d)$ est une extension finie,
choisir une base de $K$ sur $k(x_1, \ldots, x_d)$ et
raisonner en développant les éléments sur cette base.)
\item En déduire que l'on peut redéfinir le degré de transcendance en fonction
de la croissance des puissances des sous-espaces vectoriels de dimension finie de $K$.
\end{enumerate}
Cela est lié à la dimension de Gelfand--Kirillov des algèbres
(non commutatives) sur $k$.
\end{exercise}








\section{Dimension des fibres}
\label{exercises-section-dimension-fibres}

\noindent
Quelques questions liées à la formule de dimension ; voir
Algèbre, section \ref{algebra-section-dimension-formula}.

\begin{exercise}
\label{exercises-exercise-nr-components-fibre}
Soit $k$ un corps algébriquement clos quelconque.
Dans la suite, $k[x]$ et $k[x, y]$ désignent les anneaux de polynômes.
\begin{enumerate}
\item Pour tout entier $n \geq 0$, trouver une extension de type fini
$k[x] \subset A$ d'anneaux intègres telle que le spectre de $A/xA$
ait exactement $n$ composantes irréductibles.
\item Construire un exemple d'extension de type fini $k[x] \subset A$
d'anneaux intègres telle que le spectre de $A/(x - \alpha)A$
soit non vide et réductible pour tout $\alpha \in k$.
\item Construire un exemple d'extension de type fini $k[x, y] \subset A$
d'anneaux intègres telle que le spectre de $A/(x - \alpha, y - \beta)A$
soit irréductible\footnote{Rappelons qu'irréductible implique non vide.}
pour tout $(\alpha, \beta) \in k^2 \setminus \{(0, 0)\}$
et que le spectre de $A/(x, y)A$ soit non vide et réductible.
\end{enumerate}
\end{exercise}

\begin{exercise}
\label{exercises-exercise-codim-1}
Soit $k$ un corps algébriquement clos quelconque.
Soit $n \geq 1$.
Soit $k[x_1, \ldots, x_n]$ l'anneau de polynômes.
Posons $\mathfrak m = (x_1, \ldots, x_n)$.
Soit $k[x_1, \ldots, x_n] \subset A$ une extension de type fini
d'anneaux intègres. Posons $d = \dim(A)$.
\begin{enumerate}
\item Montrer que $d - 1 \geq \dim(A/\mathfrak m A) \geq d - n$
si $A/\mathfrak mA \not = 0$.
\item Montrer par des exemples que toutes les valeurs peuvent se produire.
\item Montrer par un exemple que $\Spec(A/\mathfrak m A)$ peut
avoir des composantes irréductibles de dimensions différentes.
\end{enumerate}
\end{exercise}




\section{Modules localement libres de type fini}
\label{exercises-section-finite-locally-free}

\begin{definition}
\label{exercises-definition-finite-locally-free}
Soit $A$ un anneau. Rappelons qu'un $A$-module $M$ est
{\it localement libre de type fini} si, pour tout ${\mathfrak p} \in \Spec(A)$,
il existe un
$f\in A$, $f \not \in {\mathfrak p}$, tel que $M_f$ soit un $A_f$-module libre
de type fini. On dit que $M$ est un {\it module inversible} si
$M$ est localement libre de type fini et de rang $1$, c'est-à-dire si, pour tout
${\mathfrak p} \in \Spec(A)$, il existe un
$f\in A$, $f \not \in \mathfrak p$, tel que $M_f \cong A_f$
comme $A_f$-module.
\end{definition}

\begin{exercise}
\label{exercises-exercise-tensor-finite-locally-free}
Montrer que le produit tensoriel de modules localement libres de type fini
est localement libre de type fini. Montrer que le produit tensoriel de deux
modules inversibles est inversible.
\end{exercise}

\begin{definition}
\label{exercises-definition-class-group}
Soit $A$ un anneau. Le {\it groupe des classes de $A$}, parfois appelé
{\it groupe de Picard de $A$}, est l'ensemble $\Pic(A)$
des classes d'isomorphie de $A$-modules inversibles, muni
de la loi de groupe définie par le produit tensoriel (voir
l'exercice \ref{exercises-exercise-tensor-finite-locally-free}).
\end{definition}

\noindent
Notons que le groupe des classes de $A$ est trivial si et seulement si tout module
inversible est isomorphe à un module libre de rang 1.

\begin{exercise}
\label{exercises-exercise-class-group-trivial}
Montrer que les groupes des classes des anneaux suivants sont triviaux :
\begin{enumerate}
\item un anneau de polynômes $A = k[x]$, où $k$ est un corps,
\item l'anneau des entiers $A = \mathbf{Z}$,
\item un anneau de polynômes $A = k[x, y]$, où $k$ est un corps, et
\item l'anneau quotient $k[x, y]/(xy)$, où $k$ est un corps.
\end{enumerate}
\end{exercise}

\begin{exercise}
\label{exercises-exercise-class-group-not-trivial}
Montrer que le groupe des classes de l'anneau
$A = k[x, y]/(y^2 - f(x))$, où $k$ est un corps de caractéristique différente de $2$
et où $f(x) = (x - t_1) \ldots (x - t_n)$ avec $t_1, \ldots, t_n \in k$
distincts et $n \geq 3$ impair, n'est pas trivial. (Indication : montrer que
l'idéal $(y, x - t_1)$ définit un élément non trivial de $\Pic(A)$.)
\end{exercise}

\begin{exercise}
\label{exercises-exercise-trace-det}
Soit $A$ un anneau.
\begin{enumerate}
\item Supposons que $M$ soit un $A$-module localement libre de type fini et
que $\varphi : M \to M$ soit un endomorphisme. Définir/construire
la {\it trace} et le {\it déterminant} de $\varphi$, et montrer que cette
construction est « fonctorielle par rapport au triplet $(A, M, \varphi)$ ».
\item Montrer que, si $M, N$ sont des $A$-modules localement libres de type fini,
et si $\varphi : M \to N$ et $\psi : N \to M$, alors
$\text{Trace}(\varphi \circ \psi) = \text{Trace}(\psi \circ \varphi)$ et
$\det(\varphi \circ \psi) = \det(\psi \circ \varphi)$.
\item Lorsque $M$ est localement libre de type fini, montrer que
$\text{Trace}$ définit une application $A$-linéaire $\text{End}_A(M) \to A$ et que
$\det$ définit une application multiplicative $\text{End}_A(M) \to A$.
\end{enumerate}
\end{exercise}

\begin{exercise}
\label{exercises-exercise-trace-det-rings}
Supposons maintenant que $B$ soit une $A$-algèbre localement libre
de type fini comme $A$-module ; autrement dit, $B$ est une $A$-algèbre
localement libre de type fini.
\begin{enumerate}
\item Définir $\text{Trace}_{B/A}$ et $\text{Norme}_{B/A}$ au moyen de
$\text{Trace}$ et de $\det$ dans l'exercice \ref{exercises-exercise-trace-det}.
\item Soit $b\in B$ et soit $\pi : \Spec(B) \to \Spec(A)$
le morphisme induit. Montrer que $\pi(V(b)) = V(\text{Norme}_{B/A}(b))$.
(Rappelons que $V(f) = \{ {\mathfrak p} \mid f \in {\mathfrak p}\}$.)
\item (Changement de base.) Supposons que $i : A \to A'$ soit un homomorphisme d'anneaux. Posons
$B' = B \otimes_A A'$. Indiquer pourquoi $i(\text{Norme}_{B/A}(b))$ est égal à
$\text{Norme}_{B'/A'}(b \otimes 1)$.
\item Calculer $\text{Norme}_{B/A}(b)$ lorsque
$B = A \times A \times A \times \ldots \times A$
et $b = (a_1, \ldots, a_n)$.
\item Calculer la norme de $y-y^3$ pour l'homomorphisme fini et plat
${\mathbf Q}[x] \to {\mathbf Q}[y]$, $x \to y^n$. (Indication : utiliser
le « changement de base »
$A = {\mathbf Q}[x] \subset A' = {\mathbf Q}(\zeta_n)(x^{1/n})$.)
\end{enumerate}
\end{exercise}



\section{Recollement}
\label{exercises-section-glueing}

\begin{exercise}
\label{exercises-exercise-cover}
Supposons que $A$ soit un anneau et que $M$ soit un $A$-module.
Soit donnée une famille d'éléments $f_i$, $i \in I$, de $A$ telle que
$$
\Spec(A) = \bigcup D(f_i).
$$
\begin{enumerate}
\item Montrer que, si $M_{f_i}$ est un $A_{f_i}$-module de type fini,
alors $M$ est un $A$-module de type fini.
\item Montrer que, si $M_{f_i}$ est un $A_{f_i}$-module plat,
alors $M$ est un $A$-module plat.
(C'est presque tautologique lorsqu'on y réfléchit.)
\end{enumerate}
\end{exercise}

\begin{remark}
\label{exercises-remark-cover}
Dans le langage de la géométrie algébrique, cela signifie que la propriété
d'« être de type fini » ou d'« être plat » est locale pour la topologie de Zariski
(en un sens approprié). On peut aussi le montrer pour la propriété
d'« être de présentation finie ».
\end{remark}

\begin{exercise}
\label{exercises-exercise-cover-ring-map}
Supposons que $A \to B$ soit un homomorphisme d'anneaux.
Soient $f_i \in A$, $i \in I$, et $g_j \in B$, $j \in J$, deux familles
d'éléments telles que
$$
\Spec(A) = \bigcup D(f_i)
\quad\text{et}\quad
\Spec(B) = \bigcup D(g_j).
$$
Montrer que, si $A_{f_i} \to B_{f_ig_j}$ est de type fini pour tous $i, j$,
alors $A \to B$ est de type fini.
\end{exercise}




\section{Montée et descente}
\label{exercises-section-going-up}


\begin{definition}
\label{exercises-definition-GU-GD}
Soit $\phi : A \to B$ un homomorphisme d'anneaux. On dit
que le {\it théorème de montée} est valable pour $\phi$ si la
condition suivante est satisfaite :
\begin{itemize}
\item[(GU)] pour tous ${\mathfrak p}, {\mathfrak p}' \in \Spec(A)$ tels que
${\mathfrak p} \subset {\mathfrak p}'$, et pour tout $P \in \Spec(B)$ situé
au-dessus de ${\mathfrak p}$, il existe $P'\in \Spec(B)$ situé
au-dessus de ${\mathfrak p}'$ tel que $P \subset P'$.
\end{itemize}
De même, on dit que le {\it théorème de descente} est valable pour $\phi$
si la condition suivante est satisfaite :
\begin{itemize}
\item[(GD)] pour tous ${\mathfrak p}, {\mathfrak p}' \in \Spec(A)$ tels que
${\mathfrak p} \subset {\mathfrak p}'$, et pour tout
$P' \in \Spec(B)$ situé
au-dessus de ${\mathfrak p}'$, il existe $P\in \Spec(B)$ situé
au-dessus de ${\mathfrak p}$ tel que $P \subset P'$.
\end{itemize}
\end{definition}

\begin{exercise}
\label{exercises-exercise-GU-GD}
Dans chacun des cas suivants, déterminer lesquelles des conditions
(GU) et (GD) sont satisfaites, et expliquer pourquoi. (Utiliser toute proposition,
tout théorème ou tout lemme disponible, mais vérifier les hypothèses dans chaque cas.)
\begin{enumerate}
\item $k$ est un corps, $A = k$, $B = k[x]$.
\item $k$ est un corps, $A = k[x]$, $B = k[x, y]$.
\item $A = {\mathbf Z}$, $B = {\mathbf Z}[1/11]$.
\item $k$ est un corps algébriquement clos, $A = k[x, y]$,
$B = k[x, y, z]/(x^2-y, z^2-x)$.
\item $A = {\mathbf Z}$, $B = {\mathbf Z}[i, 1/(2 + i)]$.
\item $A = {\mathbf Z}$, $B = {\mathbf Z}[i, 1/(14 + 7i)]$.
\item $k$ est un corps algébriquement clos, $A = k[x]$,
$B = k[x, y, 1/(xy-1)]/(y^2-y)$.
\end{enumerate}
\end{exercise}

\begin{exercise}
\label{exercises-exercise-image}
Soit $A$ un anneau. Soit $B = A[x]$ l'algèbre polynomiale à
une indéterminée sur $A$. Soit $f = a_0 + a_1 x + \ldots + a_r x^r \in B = A[x]$.
Démontrer soigneusement que l'image de $D(f)$ dans $\Spec(A)$ est égale à
$D(a_0) \cup \ldots \cup D(a_r)$.
\end{exercise}

\begin{exercise}
\label{exercises-exercise-images}
Soit $k$ un corps algébriquement clos. Calculer l'image dans
$\Spec(k[x, y])$
des morphismes suivants :
\begin{enumerate}
\item $\Spec(k[x, yx^{-1}]) \to \Spec(k[x, y])$, où
$k[x, y] \subset k[x, yx^{-1}] \subset k[x, y, x^{-1}]$.
(Indication : pour éviter toute confusion, donner un autre nom à l'élément $yx^{-1}$.)
\item $\Spec(k[x, y, a, b]/(ax-by-1))\to \Spec(k[x, y])$.
\item $\Spec(k[t, 1/(t-1)]) \to \Spec(k[x, y])$, induit par $x
\mapsto t^2$,
et $y \mapsto t^3$.
\item $k = {\mathbf C}$ (corps des nombres complexes),
$\Spec(k[s, t]/(s^3 + t^3-1)) \to \Spec(k[x, y])$, où
$x\mapsto s^2$, $y \mapsto t^2$.
\end{enumerate}
\end{exercise}

\begin{remark}
\label{exercises-remark-elimination-theory}
La détermination de l'image comme ci-dessus se fait habituellement au moyen de la théorie de l'élimination.
\end{remark}




\section{Idéaux de Fitting}
\label{exercises-section-fitting-ideals}

\begin{exercise}
\label{exercises-exercise-fitting}
Soient $R$ un anneau et $M$ un $R$-module de type fini.
Choisir une présentation
$$
\bigoplus\nolimits_{j \in J} R \longrightarrow R^{\oplus n}
\longrightarrow M \longrightarrow 0.
$$
de $M$. Noter que l'homomorphisme $R^{\oplus n} \to M$ est défini par une famille
d'éléments $x_1, \ldots, x_n$ de $M$. Les éléments $x_i$
sont des {\it générateurs} de $M$. L'homomorphisme $\bigoplus_{j \in J} R \to R^{\oplus n}$
est défini par une matrice $n \times J$ notée $A$, à coefficients dans $R$.
Autrement dit, $A = (a_{ij})_{i = 1, \ldots, n, j \in J}$.
Les colonnes $(a_{1j}, \ldots, a_{nj})$, $j \in J$, de $A$
sont appelées les {\it relations}. Tout vecteur $(r_i) \in R^{\oplus n}$
tel que $\sum r_i x_i = 0$ est une combinaison linéaire des colonnes de $A$.
Bien entendu, tout $R$-module de type fini possède de nombreuses présentations différentes.
\begin{enumerate}
\item Montrer que l'idéal engendré par les mineurs $(n - k) \times (n - k)$ de
$A$ est indépendant du choix de la présentation.
Cet idéal est l'{\it idéal de Fitting d'indice $k$ de $M$}. On le note $Fit_k(M)$.
\item Montrer que
$Fit_0(M) \subset Fit_1(M) \subset Fit_2(M) \subset \ldots$.
(Indication : utiliser le développement d'un déterminant suivant une colonne.)
\item Montrer que les conditions suivantes sont équivalentes :
\begin{enumerate}
\item $Fit_{r - 1}(M) = (0)$ et $Fit_r(M) = R$, et
\item $M$ est localement libre de rang $r$.
\end{enumerate}
\end{enumerate}
\end{exercise}




\section{Fonctions de Hilbert}
\label{exercises-section-hilbert}

\begin{definition}
\label{exercises-definition-numerical-polynomial}
Un {\it polynôme numérique} est un polynôme $f(x) \in {\mathbf Q}[x]$
tel que $f(n) \in {\mathbf Z}$ pour tout entier $n$.
\end{definition}

\begin{definition}
\label{exercises-definition-graded-module}
Un {\it module gradué} $M$ sur un anneau $A$ est un $A$-module $M$
muni d'une décomposition en somme directe de la forme
$
\bigoplus\nolimits_{n \in {\mathbf Z}} M_n
$
en sous-$A$-modules. On dit que $M$ est {\it localement fini} si tous les
$M_n$ sont des $A$-modules de type fini. Supposons que $A$ soit un anneau noethérien et
que $\varphi$ soit une {\it fonction d'Euler-Poincar\'e} sur les $A$-modules de type fini.
Cela signifie que, pour tout $A$-module de type fini $M$, on se donne un
entier $\varphi(M) \in {\mathbf Z}$ et que, pour toute suite exacte courte
$$
0
\longrightarrow
M'
\longrightarrow
M
\longrightarrow
M''
\longrightarrow
0
$$
on a $\varphi(M) = \varphi(M') + \varphi(M'')$. La {\it fonction de Hilbert}
d'un module gradué localement fini $M$ (relativement à $\varphi$) est la
fonction $\chi_\varphi(M, n) = \varphi(M_n)$. On dit que $M$ admet un
{\it polynôme de Hilbert} s'il existe un polynôme numérique
$P_\varphi$ tel que $\chi_\varphi(M, n) = P_\varphi(n)$ pour tout entier $n$
assez grand.
\end{definition}

\begin{definition}
\label{exercises-definition-graded-algebra}
Une {\it $A$-algèbre graduée} est un $A$-module gradué
$B = \bigoplus_{n \geq 0} B_n$ muni d'une application $A$-bilinéaire
$$
B \times B \longrightarrow B, \ (b, b') \longmapsto bb'
$$
qui munit $B$ d'une structure de $A$-algèbre telle que $B_n \cdot B_m \subset B_{n + m}$.
Enfin, un {\it module gradué $M$ sur une $A$-algèbre graduée $B$} est la donnée
d'un $A$-module gradué $M$ muni d'une structure (compatible) de $B$-module
telle que $B_n \cdot M_d \subset M_{n + d}$. On peut alors définir les
{\it homomorphismes de modules/anneaux gradués}, les {\it sous-modules gradués},
les {\it idéaux gradués}, les {\it suites exactes de modules gradués}, etc.
\end{definition}

\begin{exercise}
\label{exercises-exercise-Euler-Poincare-field}
Soit le corps $A = k$. Quelles sont toutes les fonctions d'Euler-Poincar\'e
possibles sur les $A$-modules de type fini dans ce cas ?
\end{exercise}

\begin{exercise}
\label{exercises-exercise-Euler-Poincare-Z}
Soit $A ={\mathbf Z}$. Quelles sont toutes les fonctions d'Euler-Poincar\'e
possibles sur les $A$-modules de type fini dans ce cas ?
\end{exercise}

\begin{exercise}
\label{exercises-exercise-Euler-Poincare-node}
Soit $A = k[x, y]/(xy)$, où $k$ est un corps algébriquement clos. Quelles sont, dans ce cas, toutes les
fonctions d'Euler-Poincar\'e possibles sur les $A$-modules de type fini ?
\end{exercise}

\begin{exercise}
\label{exercises-exercise-kernel-locally-finite}
Supposons que $A$ soit noethérien. Montrer que le noyau d'un homomorphisme
de $A$-modules gradués localement finis est localement fini.
\end{exercise}

\begin{exercise}
\label{exercises-exercise-no-hilbert}
Soit $k$ un corps, posons $A = k$ et munissons $B = k[x, y]$ de la graduation
déterminée par $\deg(x) = 2$ et $\deg(y) = 3$. Soit $\varphi(M) = \dim_k(M)$.
Calculer la fonction de Hilbert de $B$ comme $k$-module gradué. Existe-t-il
un polynôme de Hilbert dans ce cas ?
\end{exercise}

\begin{exercise}
\label{exercises-exercise-no-hilbert-or-is-there}
Soit $k$ un corps, posons $A = k$ et munissons $B = k[x, y]/(x^2, xy)$ de la graduation
déterminée par $\deg(x) = 2$ et $\deg(y) = 3$. Soit $\varphi(M) = \dim_k(M)$.
Calculer la fonction de Hilbert de $B$ comme $k$-module gradué. Existe-t-il
un polynôme de Hilbert dans ce cas ?
\end{exercise}

\begin{exercise}
\label{exercises-exercise-hilbert-to-compute}
Soit $k$ un corps et posons $A = k$. Soit $\varphi(M) = \dim_k(M)$.
Fixons $d\in {\mathbf N}$. Considérons l'algèbre graduée sur $A$
$B = k[x, y, z]/(x^d + y^d + z^d)$, où $x, y, z$ sont tous de degré $1$.
Calculer la fonction de Hilbert de $B$. Existe-t-il un polynôme de Hilbert
dans ce cas ?
\end{exercise}



\section{Proj d'un anneau}
\label{exercises-section-proj-ring}

\begin{definition}
\label{exercises-definition-homogeneous-ideal}
Soit $R$ un anneau gradué. Un idéal {\it homogène} est simplement un idéal
$I \subset R$ qui est en outre un sous-module gradué de $R$. De manière équivalente,
c'est un idéal engendré par des éléments homogènes. De manière équivalente, si
$f \in I$ et si
$$
f = f_0 + f_1 + \ldots + f_n
$$
est la décomposition de $f$ en composantes homogènes dans $R$, alors $f_i \in I$
pour tout $i$.
\end{definition}

\begin{definition}
\label{exercises-definition-Proj-R}
On définit le {\it spectre homogène $\text{Proj}(R)$}
de l'anneau gradué $R$ comme l'ensemble des idéaux premiers homogènes
${\mathfrak p}$ de $R$ tels que $R_{+} \not \subset {\mathfrak p}$.
Noter que $\text{Proj}(R)$ est une partie de $\Spec(R)$ et porte donc naturellement
la topologie induite.
\end{definition}

\begin{definition}
\label{exercises-definition-Dplus-Vplus}
Soit $R = \oplus_{d \geq 0} R_d$ un anneau gradué, soit $f\in R_d$ et
supposons que $d \geq 1$. On définit {\it $R_{(f)}$} comme le sous-anneau de
$R_f$ formé des éléments de la forme $r/f^n$, où $r$ est homogène et
$\deg(r) = nd$. On définit en outre
$$
D_{+}(f) = \{ {\mathfrak p} \in \text{Proj}(R) | f \not\in {\mathfrak p} \}.
$$
Enfin, pour un idéal homogène $I \subset R$, on définit
$V_{+}(I) = V(I) \cap \text{Proj}(R)$.
\end{definition}

\begin{exercise}
\label{exercises-exercise-topology-proj}
Sur la topologie de $\text{Proj}(R)$. Avec les définitions et les notations
ci-dessus, démontrer les assertions suivantes.
\begin{enumerate}
\item Montrer que $D_{+}(f)$ est ouvert dans $\text{Proj}(R)$.
\item Montrer que $D_{+}(ff') = D_{+}(f) \cap D_{+}(f')$.
\item Soit $g = g_0 + \ldots + g_m$ un élément
de $R$ avec $g_i \in R_i$. Exprimer $D(g) \cap \text{Proj}(R)$
en fonction des $D_{+}(g_i)$, $i \geq 1$, et de $D(g_0) \cap \text{Proj}(R)$.
Aucune démonstration n'est nécessaire.
\item Soit $g\in R_0$ un élément homogène de degré $0$.
Exprimer $D(g) \cap \text{Proj}(R)$ en fonction des $D_{+}(f_\alpha)$
pour une famille convenable d'éléments homogènes $f_\alpha \in R$
de degré strictement positif.
\item Montrer que la famille $\{D_{+}(f)\}$ d'ouverts forme une
base de la topologie de $\text{Proj}(R)$.
\item
\label{exercises-item-bijection}
Montrer qu'il existe une bijection canonique $D_{+}(f) \to \Spec(R_{(f)})$.
(Indication : imiter la démonstration pour $\Spec$, mais, à un certain stade, faire intervenir le
radical d'un idéal.)
\item Montrer que l'application définie en (\ref{exercises-item-bijection}) est un homéomorphisme.
\item Donner un exemple d'un anneau $R$ tel que $\text{Proj}(R)$ ne soit pas
quasi-compact. Aucune démonstration n'est nécessaire.
\item Montrer que toute partie fermée $T \subset \text{Proj}(R)$ est de
la forme $V_{+}(I)$ pour un idéal homogène $I \subset R$.
\end{enumerate}
\end{exercise}

\begin{remark}
\label{exercises-remark-continuous-proj-spec}
Il existe une application continue $ \text{Proj}(R) \longrightarrow \Spec(R_0) $.
\end{remark}

\begin{exercise}
\label{exercises-exercise-iso-polynomial-ring-one-variable}
Si $R = A[X]$ avec $\deg(X) = 1$, montrer que l'application naturelle
$\text{Proj}(R) \to \Spec(A)$ est une bijection et même
un homéomorphisme.
\end{exercise}

\begin{exercise}
\label{exercises-exercise-blowing-up-I}
Éclatement : première partie.
Dans cet exercice, $R = Bl_I(A) = A \oplus I \oplus I^2 \oplus \ldots$.
Considérer l'application naturelle $b : \text{Proj}(R) \to \Spec(A)$.
Poser $U = \Spec(A) - V(I)$. Montrer que
$$
b : b^{-1}(U) \longrightarrow U
$$
est un homéomorphisme.
On peut donc identifier $U$ à une partie ouverte de $\text{Proj}(R)$.
Soit $Z \subset \Spec(A)$ un sous-schéma fermé irréductible
de point générique $\xi \in Z$. Supposer que $\xi \not\in V(I)$,
autrement dit $Z \not\subset V(I)$, autrement dit
$\xi \in U$, autrement dit $Z\cap U \not = \emptyset$. On définit
le {\it transformé pur} $Z'$ de $Z$ comme l'adhérence de l'unique
point $\xi'$ situé au-dessus de $\xi$. Autrement dit,
$Z'$ est l'adhérence dans $\text{Proj}(R)$ de la partie localement fermée
$Z\cap U \subset U \subset \text{Proj}(R)$.
\end{exercise}

\begin{exercise}
\label{exercises-exercise-blowing-up-II}
Éclatement : deuxième partie.
Soit $A = k[x, y]$, où $k$ est un corps, et soit $I = (x, y)$.
Soit $R$ l'algèbre d'éclatement associée à $A$ et $I$.
\begin{enumerate}
\item Montrer que les transformés purs de $Z_1 = V(\{x\})$ et de
$Z_2 = V(\{y\})$ sont disjoints.
\item Montrer que les transformés purs de $Z_1 = V(\{x\})$ et de
$Z_2 = V(\{x-y^2\})$ ne sont pas disjoints.
\item Trouver un idéal $J \subset A$ tel que $V(J) = V(I)$
et tel que les transformés purs de $Z_1 = V(\{x\})$ et de
$Z_2 = V(\{x-y^2\})$ dans l'éclatement le long de $J$ soient disjoints.
\end{enumerate}
\end{exercise}

\begin{exercise}
\label{exercises-exercise-proj-when-empty}
Soit $R$ un anneau gradué.
\begin{enumerate}
\item Montrer que $\text{Proj}(R)$ est vide si $R_n = (0)$ pour tout $n >> 0$.
\item Montrer que $\text{Proj}(R)$ est un espace topologique irréductible
si $R$ est intègre et si $R_{+}$ n'est pas nul. (Rappelons que l'espace
topologique vide n'est pas irréductible.)
\end{enumerate}
\end{exercise}

\begin{exercise}
\label{exercises-exercise-blowing-up-III}
Éclatement : troisième partie.
On reprend $A$, $I$, $U$ et $Z$ comme dans la définition du transformé pur.
Soit $Z = V({\mathfrak p})$ pour un idéal premier ${\mathfrak p}$. Poser $\bar A =
A/{\mathfrak p}$ et noter
$\bar I$ l'image de $I$ dans $\bar A$.
\begin{enumerate}
\item Montrer qu'il existe un homomorphisme d'anneaux surjectif
$R: = Bl_I(A) \to \bar R: = Bl_{\bar I}(\bar A)$.
\item Montrer que l'homomorphisme d'anneaux ci-dessus induit une application bijective
de $\text{Proj}(\bar R)$ sur le transformé pur $Z'$ de $Z$. (Cette question
n'est pas si facile. Indication : utiliser 5(b) ci-dessus.)
\item En déduire que le transformé pur $Z' = V_{+}(P)$, où
$P \subset R$ est l'idéal homogène défini par
$P_d = I^d \cap {\mathfrak p}$.
\item Supposer que $Z_1 = V({\mathfrak p})$ et
$Z_2 = V({\mathfrak q})$ sont des parties fermées irréductibles
définies par des idéaux premiers, que $Z_1 \not \subset Z_2$
et que $Z_2 \not \subset Z_1$. Montrer que l'éclatement de l'idéal
$I = {\mathfrak p} + {\mathfrak q}$ sépare les
transformés purs de $Z_1$ et de $Z_2$,
c'est-à-dire $Z_1' \cap Z_2' = \emptyset$. (Indication : considérer les idéaux homogènes
$P$ et $Q$ de la partie (c), puis $V(P + Q)$.)
\end{enumerate}
\end{exercise}


\section{Anneaux de Cohen-Macaulay de dimension 1}
\label{exercises-section-CM-dim-1}

\begin{definition}
\label{exercises-definition-CM}
On dit qu'un anneau local noethérien $A$ est un {\it anneau de Cohen-Macaulay}
de dimension $d$ s'il est de dimension $d$ et s'il existe un système
de paramètres $x_1, \ldots, x_d$ pour $A$ tel que $x_i$ soit un non-diviseur de zéro
dans $A/(x_1, \ldots, x_{i-1})$ pour $i = 1, \ldots, d$.
\end{definition}

\begin{exercise}
\label{exercises-exercise-CM-dim-1-I}
Anneaux de Cohen-Macaulay de dimension 1. Première partie : théorie.
\begin{enumerate}
\item Soit $(A, {\mathfrak m})$ un anneau local noethérien tel que $\dim A = 1$.
Montrer que, si $x\in {\mathfrak m}$ n'est pas un diviseur de zéro, alors
\begin{enumerate}
\item $\dim A/xA = 0$, autrement dit $A/xA$ est artinien,
autrement dit $\{x\}$ est un système de paramètres pour $A$.
\item $A$ n'a aucun idéal premier associé immergé.
\end{enumerate}
\item Réciproquement, soit $(A, {\mathfrak m})$ un anneau local noethérien de
dimension $1$. Montrer que, si $A$ n'a aucun idéal premier associé immergé, il existe
un non-diviseur de zéro dans ${\mathfrak m}$.
\end{enumerate}
\end{exercise}

\begin{exercise}
\label{exercises-exercise-CM-dim-1-II}
Anneaux de Cohen-Macaulay de dimension 1. Deuxième partie : exemples.
\begin{enumerate}
\item Soit $A$ l'anneau local en $(x, y)$ de $k[x, y]/(x^2, xy)$.
\begin{enumerate}
\item Montrer que $A$ est de dimension 1.
\item Démontrer que tout élément de ${\mathfrak m}\subset A$ est un
diviseur de zéro.
\item Trouver $z\in {\mathfrak m}$ tel que $\dim A/zA = 0$
(aucune démonstration n'est nécessaire).
\end{enumerate}
\item Soit $A$ l'anneau local en $(x, y)$ de $k[x, y]/(x^2)$.
Trouver un non-diviseur de zéro dans ${\mathfrak m}$ (aucune démonstration n'est nécessaire).
\end{enumerate}
\end{exercise}

\begin{exercise}
\label{exercises-exercise-embedding-dim-1}
Anneaux locaux de dimension de plongement $1$.
Supposer que $(A, {\mathfrak m}, k)$ soit un anneau local noethérien
de dimension de plongement $1$, c'est-à-dire
$$
\dim_k {\mathfrak m}/{\mathfrak m}^2 = 1.
$$
Montrer que la fonction $f(n) = \dim_k {\mathfrak m}^n/{\mathfrak m}^{n + 1}$
est soit constante de valeur $1$, soit telle que ses valeurs soient
$$
1, 1, \ldots, 1, 0, 0, 0, 0, 0, \ldots
$$
\end{exercise}

\begin{exercise}
\label{exercises-exercise-regular-local-dim-1}
Anneaux locaux réguliers de dimension $1$.
Supposer que $(A, {\mathfrak m}, k)$ soit un anneau local noethérien régulier de
dimension $1$. Rappelons que cela signifie que $A$ est de dimension $1$
et de dimension de plongement $1$, c'est-à-dire
$$
\dim_k {\mathfrak m}/{\mathfrak m}^2 = 1.
$$
Soit $x\in{\mathfrak m}$ un élément dont la classe dans ${\mathfrak m}/{\mathfrak
m}^2$ n'est pas nulle.
\begin{enumerate}
\item Montrer que, pour tout élément $y$
de ${\mathfrak m}$, il existe un entier $n$ tel que $y$ puisse s'écrire
$y = ux^n$, où $u\in A^\ast$ est une unité.
\item Montrer que $x$ est un non-diviseur de zéro dans $A$.
\item En déduire que $A$ est intègre.
\end{enumerate}
\end{exercise}

\begin{exercise}
\label{exercises-exercise-nonzerodivisor-graded}
Soit $(A, {\mathfrak m}, k)$ un anneau local noethérien, et soit
$Gr_{\mathfrak m}(A)$ son gradué associé.
\begin{enumerate}
\item Supposer que $x\in {\mathfrak m}^d$ a pour image un non-diviseur de zéro
$\bar x \in {\mathfrak m}^d/{\mathfrak m}^{d + 1}$ dans la composante de degré $d$ de
$Gr_{\mathfrak m}(A)$.
Montrer que $x$ est un non-diviseur de zéro.
\item Supposer que la profondeur de $A$ est au moins $1$.
Autrement dit, supposer qu'il existe un non-diviseur de zéro $y \in {\mathfrak m}$.
Dans ce cas, on peut améliorer le résultat : supposer seulement que $x\in {\mathfrak m}^d$ a pour image
l'élément $\bar x \in {\mathfrak m}^d/{\mathfrak m}^{d + 1}$ de degré $d$
de $Gr_{\mathfrak m}(A)$, lequel est un non-diviseur de zéro en degrés assez
élevés : $\exists N$ tel que, pour tout $n \geq N$, l'application
de multiplication par $\bar x$
$$
{\mathfrak m}^n/{\mathfrak m}^{n + 1} \longrightarrow
{\mathfrak m}^{n + d}/{\mathfrak m}^{n + d + 1}
$$
soit injective. Montrer alors que $x$ est un non-diviseur de zéro.
\end{enumerate}
\end{exercise}

\begin{exercise}
\label{exercises-exercise-embedding-2-dim-1}
Supposer que $(A, {\mathfrak m}, k)$ soit un anneau local noethérien de
dimension $1$. Supposer également que la dimension de plongement de $A$ soit
$2$, c'est-à-dire supposer que
$$
\dim_k {\mathfrak m}/{\mathfrak m}^2 = 2.
$$
Notation : $f(n) = \dim_k {\mathfrak m}^n/{\mathfrak m}^{n + 1}$.
Choisir des générateurs $x, y \in {\mathfrak m}$
et écrire $Gr_{\mathfrak m}(A) = k[\bar x, \bar y]/I$ pour un certain
idéal homogène $I$.
\begin{enumerate}
\item Montrer qu'il existe un élément homogène
$F\in k[\bar x, \bar y]$ tel que $I \subset (F)$, avec égalité
en tout degré assez élevé.
\item Montrer que $f(n) \leq n + 1$.
\item Montrer que, si $f(n) < n + 1$, alors $n \geq \deg(F)$.
\item Montrer que, si $f(n) < n + 1$, alors $f(n + 1) \leq f(n)$.
\item Montrer que $f(n) = \deg(F)$ pour tout $n >> 0$.
\end{enumerate}
\end{exercise}

\begin{exercise}
\label{exercises-exercise-CM-dim-1-embedding-dim-2}
Anneaux de Cohen-Macaulay de dimension 1 et de dimension de plongement 2.
Supposer que $(A, {\mathfrak m}, k)$ soit un anneau local noethérien de
Cohen-Macaulay
et de dimension $1$. Supposer également que la dimension de plongement de $A$ soit
$2$, c'est-à-dire supposer que
$$
\dim_k {\mathfrak m}/{\mathfrak m}^2 = 2.
$$
Notations : $f$, $F$, $x, y\in {\mathfrak m}$ et $I$ comme dans
l'exercice \ref{exercises-exercise-embedding-2-dim-1}. On pourra utiliser librement
tous les résultats des problèmes ci-dessus.
\begin{enumerate}
\item Supposer que $z\in {\mathfrak m}$ est un élément dont la classe
dans ${\mathfrak m}/{\mathfrak m}^2$ est la forme linéaire
$\alpha \bar x + \beta \bar y \in k[\bar x, \bar y]$
qui est première à $F$.
\begin{enumerate}
\item Montrer que $z$ est un non-diviseur de zéro dans $A$.
\item Poser $d = \deg(F)$.
Montrer que ${\mathfrak m}^n = z^{n + 1-d}{\mathfrak m}^{d-1}$
pour tout $n$ assez grand. (Indication : montrer d'abord que
$z^{n + 1-d}{\mathfrak m}^{d-1} \to {\mathfrak m}^n/{\mathfrak m}^{n + 1}$
est surjective en utilisant ce que l'on sait de $Gr_{\mathfrak m}(A)$. Puis utiliser le lemme de Nakayama.)
\end{enumerate}
\item Quelle condition sur $k$ garantit l'existence
d'un tel $z$ ? (Aucune démonstration n'est nécessaire; c'est trop facile.)

\noindent
Supposons désormais qu'il existe un $z$ comme ci-dessus. Cette hypothèse
est en fait sans restriction essentielle (au sens où l'on peut se ramener à
la situation où elle est satisfaite afin d'obtenir les résultats des
parties (d) et (e) ci-dessous).
\item Montrer maintenant que ${\mathfrak m}^\ell = z^{\ell - d + 1} {\mathfrak m}^{d-1}$
pour tout $\ell \geq d$.
\item En déduire que $I = (F)$.
\item En déduire que la fonction $f$ prend les valeurs
$$
2, 3, 4, \ldots, d-1, d, d, d, d, d, d, d, \ldots
$$
\end{enumerate}
\end{exercise}

\begin{remark}
\label{exercises-remark-CM-dim-1-embedding-dim-2}
Cela donne à penser qu'un anneau local noethérien de Cohen-Macaulay de dimension 1
et de dimension de plongement 2 est de la forme $B/FB$, où $B$ est un anneau local
régulier de dimension 2. C'est à peu près exact (sous des hypothèses
convenables de ``bonne tenue'' de l'anneau).
\end{remark}






\section{Une infinité de nombres premiers}
\label{exercises-section-many-primes}

\noindent
section regroupant des questions singulières sur les anneaux dans lesquels
les nombres premiers non inversibles sont en nombre infini.

\begin{exercise}
\label{exercises-exercise-not-in-Q}
Donner un exemple d'une ${\mathbf Z}$-algèbre $R$ de type fini
possédant les deux propriétés suivantes :
\begin{enumerate}
\item Il n'existe aucun homomorphisme d'anneaux $R \to {\mathbf Q}$.
\item Pour tout nombre premier $p$, il existe un idéal maximal
${\mathfrak m} \subset R$ tel que $R/{\mathfrak m} \cong {\mathbf F}_p$.
\end{enumerate}
\end{exercise}

\begin{exercise}
\label{exercises-exercise-strange-fp-1}
Pour $f \in {\mathbf Z}[x, u]$, on pose $f_p(x)
= f(x, x^p) \bmod p \in {\mathbf F}_p[x]$. Donner un exemple
d'un $f \in {\mathbf Z}[x, u]$ tel que les deux propriétés
suivantes soient satisfaites :
\begin{enumerate}
\item Il existe une infinité de nombres premiers $p$ tels que $f_p$
n'ait pas de racine dans ${\mathbf F}_p$.
\item Pour tout $p >> 0$, le polynôme $f_p$ possède un facteur
linéaire ou quadratique.
\end{enumerate}
\end{exercise}

\begin{exercise}
\label{exercises-exercise-strange-fp-2}
Pour $f \in {\mathbf Z}[x, y, u, v]$, on pose $f_p(x, y)
= f(x, y, x^p, y^p) \bmod p \in {\mathbf F}_p[x, y]$. Donner un exemple
« intéressant » d'un $f$ tel que $f_p$ soit réductible pour tout $p >> 0$.
Par exemple, $f = xv-yu$ avec $f_p = xy^p-x^py = xy(x^{p-1}-y^{p-1})$ est
« sans intérêt » ; tout $f$ ne dépendant que de $x, u$ est « sans intérêt »,
etc.
\end{exercise}

\begin{remark}
\label{exercises-remark-strange-fp}
Soit $h \in {\mathbf Z}[y]$ un polynôme unitaire de degré $d$.
Alors :
\begin{enumerate}
\item L'homomorphisme $A = {\mathbf Z}[x] \to B ={\mathbf Z}[y]$,
$x \mapsto h$ est fini et localement libre de rang $d$.
\item Pour tout nombre premier $p$, l'homomorphisme
$A_p = {\mathbf F}_p[x]\to B_p = {\mathbf F}_p[y]$,
$y \mapsto h(y) \bmod p$ est fini et localement libre de rang $d$.
\end{enumerate}
\end{remark}

\begin{exercise}
\label{exercises-exercise-strange-fp-3}
Soient $h, A, B, A_p, B_p$ comme dans la remarque. Pour $f \in {\mathbf Z}[x, u]$, on
pose $f_p(x) = f(x, x^p) \bmod p \in {\mathbf F}_p[x]$. Pour
$g \in {\mathbf Z}[y, v]$, on pose
$g_p(y) = g(y, y^p) \bmod p \in {\mathbf F}_p[y]$.
\begin{enumerate}
\item Donner un exemple de $h$ et de $g$ pour lesquels
il n'existe aucun $f$ ayant la propriété
$$
f_p  =  Norm_{B_p/A_p}(g_p).
$$
\item Montrer que, quels que soient $h$ et $g$ comme ci-dessus,
il existe un $f$ non nul tel que, pour tout $p$, on ait
$$
Norm_{B_p/A_p}(g_p)\quad\text{divides}\quad f_p .
$$
On peut, si on le souhaite, se restreindre au cas $h = y^n$, même avec $n = 2$,
mais l'énoncé est vrai en général.
\item Discuter le rapport avec les exercices 6 et 7 de la série
précédente.
\end{enumerate}
\end{exercise}

\begin{exercise}
\label{exercises-exercise-strange-fp-unsolved}
Problèmes non résolus. Ils sont peut-être très difficiles ou peut-être faciles.
Je l'ignore.
\begin{enumerate}
\item Existe-t-il un $f \in {\mathbf Z}[x, u]$ tel que $f_p$ soit
irréductible pour une infinité de $p$ ? (Indication : oui, c'est le cas pour
$f(x, u) = u - x - 1$ ainsi que pour $f(x, u) = u^2 - x^2 + 1$.)
\item Soit $f \in {\mathbf Z}[x, u]$ non nul, et supposons que
$\deg_x(f_p) = dp + d'$ pour tout $p$ assez grand. (Autrement dit, $\deg_u(f) = d$
et le coefficient $c$ de $u^d$ dans $f$ vérifie $\deg_x(c) = d'$.) Supposons que l'on
puisse écrire $d = d_1 + d_2$ et $d' = d'_1 + d'_2$ avec $d_1, d_2 > 0$
et $d'_1, d'_2 \geq 0$, de telle sorte que, pour tout $p$ assez grand,
il existe une factorisation
$$
f_p = f_{1, p} f_{2, p}
$$
avec $\deg_x(f_{1, p}) = d_1p + d'_1$. Est-il vrai que $f$ provient d'une
construction par la norme comme dans l'exercice 4 ? (Plus précisément, existe-t-il $h$ et
$g$ tels que $Norm_{B_p/A_p}(g_p)$ divise $f_p$ pour tout $p >> 0$ ?)
\item Question analogue à celle de (b), mais avec cette fois
$f \in {\mathbf Z}[x_1, x_2, u_1, u_2]$ irréductible et en supposant seulement que
$f_p(x_1, x_2) = f(x_1, x_2, x_1^p, x_2^p) \bmod p$ se factorise pour tout
$p >> 0$.
\end{enumerate}
\end{exercise}


















\section{Catégorie dérivée filtrée}
\label{exercises-section-filtered-derived}

\noindent
Pour traiter les exercices de cette section, lire au préalable les développements
d'Homologie, section \ref{homology-section-filtrations}. Nous dirons que
$A$ est un objet filtré de $\mathcal{A}$ pour signifier que $A$ est muni
d'une filtration $F$, que nous omettons de la notation.

\begin{exercise}
\label{exercises-exercise-split-injective}
Soit $\mathcal{A}$ une catégorie abélienne.
Soit $I$ un objet filtré de $\mathcal{A}$.
Supposons que la filtration de $I$ soit finie
et que chaque $\text{gr}^p(I)$ soit un objet injectif de $\mathcal{A}$.
Montrer qu'il existe un isomorphisme
$I \cong \bigoplus \text{gr}^p(I)$ respectant les filtrations, la filtration
$F^p(I)$ correspondant à $\bigoplus_{p' \geq p} \text{gr}^p(I)$.
\end{exercise}

\begin{exercise}
\label{exercises-exercise-filtered-injective}
Soit $\mathcal{A}$ une catégorie abélienne.
Soit $I$ un objet filtré de $\mathcal{A}$.
Supposons que la filtration de $I$ soit finie.
Montrer l'équivalence des propriétés suivantes :
\begin{enumerate}
\item Pour tout diagramme dont les flèches pleines sont données
$$
\xymatrix{
A \ar[r]_\alpha \ar[d] & B \ar@{-->}[ld] \\
I &
}
$$
d'objets filtrés, si
(\romannumeral1) les filtrations de $A$ et de $B$ sont finies,
et (\romannumeral2) $\text{gr}(\alpha)$ est injectif, alors la
flèche en pointillé existe et rend le diagramme commutatif.
\item Chaque $\text{gr}^p I$ est injectif.
\end{enumerate}
\end{exercise}

\noindent
Notons que, si $\alpha : A \to B$ est un morphisme d'objets filtrés
à filtrations finies, dire que $\text{gr}(\alpha)$ est injectif
revient à dire que $\alpha$ est un {\it monomorphisme strict}
dans la catégorie $\text{Fil}(\mathcal{A})$. En effet,
être un monomorphisme signifie que $\Ker(\alpha) = 0$, tandis qu'être strict signifie que
cela entraîne aussi $\Ker(\text{gr}(\alpha)) = 0$.
Voir Homologie, lemme \ref{homology-lemma-characterize-strict}.
(Nous n'employons le terme « injectif » pour un morphisme que dans une catégorie abélienne,
bien qu'il ait un sens dans toute catégorie additive admettant des noyaux.)
Les exercices ci-dessus justifient la définition suivante.

\begin{definition}
\label{exercises-definition-injective-filtered}
Soit $\mathcal{A}$ une catégorie abélienne.
Soit $I$ un objet filtré de $\mathcal{A}$.
Supposons que la filtration de $I$ soit finie.
Nous dirons que $I$ est {\it injectif filtré} si chaque $\text{gr}^p(I)$ est
un objet injectif de $\mathcal{A}$.
\end{definition}

\noindent
Nous adoptons la définition suivante afin de ne pas avoir à répéter
partout « muni d'une filtration finie ».

\begin{definition}
\label{exercises-definition-finite-filtration-category}
Soit $\mathcal{A}$ une catégorie abélienne.
Nous noterons {\it $\text{Fil}^f(\mathcal{A})$} la sous-catégorie pleine
de $\text{Fil}(\mathcal{A})$ formée des objets
$A \in \Ob(\text{Fil}(\mathcal{A}))$
dont la filtration est finie.
\end{definition}

\begin{exercise}
\label{exercises-exercise-inject-into-injective}
Soit $\mathcal{A}$ une catégorie abélienne.
Supposons que $\mathcal{A}$ admette suffisamment d'injectifs.
Soit $A$ un objet de $\text{Fil}^f(\mathcal{A})$.
Montrer qu'il existe un monomorphisme strict $\alpha : A \to I$
de $A$ dans un objet injectif filtré $I$ de $\text{Fil}^f(\mathcal{A})$.
\end{exercise}

\begin{definition}
\label{exercises-definition-filtered-quasi-isomorphism}
Soit $\mathcal{A}$ une catégorie abélienne.
Soit $\alpha : K^\bullet \to L^\bullet$ un morphisme de
complexes de $\text{Fil}(\mathcal{A})$. Nous dirons que
$\alpha$ est un {\it quasi-isomorphisme filtré} si,
pour tout $p \in \mathbf{Z}$, le morphisme
$\text{gr}^p(K^\bullet) \to \text{gr}^p(L^\bullet)$ est
un quasi-isomorphisme.
\end{definition}

\begin{definition}
\label{exercises-definition-filtered-acyclic}
Soit $\mathcal{A}$ une catégorie abélienne.
Soit $K^\bullet$ un complexe de $\text{Fil}^f(\mathcal{A})$.
Nous dirons que $K^\bullet$ est {\it acyclique filtré} si,
pour tout $p \in \mathbf{Z}$, le complexe $\text{gr}^p(K^\bullet)$ est
acyclique.
\end{definition}

\begin{exercise}
\label{exercises-exercise-filtered-quasi-isomorphism}
Soit $\mathcal{A}$ une catégorie abélienne.
Soit $\alpha : K^\bullet \to L^\bullet$ un morphisme de complexes bornés inférieurement
de $\text{Fil}^f(\mathcal{A})$. (Noter l'exposant $f$.)
Montrer l'équivalence des propriétés suivantes :
\begin{enumerate}
\item $\alpha$ est un quasi-isomorphisme filtré,
\item pour tout $p \in \mathbf{Z}$, l'application
$\alpha : F^pK^\bullet \to F^pL^\bullet$ est un quasi-isomorphisme,
\item pour tout $p \in \mathbf{Z}$, l'application
$\alpha : K^\bullet/F^pK^\bullet \to L^\bullet/F^pL^\bullet$
est un quasi-isomorphisme, et
\item le cône de $\alpha$ (voir
Catégories dérivées, définition \ref{derived-definition-cone})
est un complexe acyclique filtré.
\end{enumerate}
Montrer de plus que, si $\alpha$ est un quasi-isomorphisme filtré,
alors $\alpha$ est aussi un quasi-isomorphisme au sens usuel.
\end{exercise}

\begin{exercise}
\label{exercises-exercise-injective-resolution}
Soit $\mathcal{A}$ une catégorie abélienne.
Supposons que $\mathcal{A}$ admette suffisamment d'injectifs.
Soit $A$ un objet de $\text{Fil}^f(\mathcal{A})$.
Montrer qu'il existe un complexe
$I^\bullet$ de $\text{Fil}^f(\mathcal{A})$,
et un morphisme $A[0] \to I^\bullet$ tels que
\begin{enumerate}
\item chaque $I^p$ est injectif filtré,
\item $I^p = 0$ pour $p < 0$, et
\item $A[0] \to I^\bullet$ est un quasi-isomorphisme filtré.
\end{enumerate}
\end{exercise}

\begin{exercise}
\label{exercises-exercise-injective-resolution-complex}
Soit $\mathcal{A}$ une catégorie abélienne.
Supposons que $\mathcal{A}$ admette suffisamment d'injectifs.
Soit $K^\bullet$ un complexe borné inférieurement d'objets de
$\text{Fil}^f(\mathcal{A})$. Montrer qu'il existe un
quasi-isomorphisme filtré $\alpha : K^\bullet \to I^\bullet$
où $I^\bullet$ est un complexe de $\text{Fil}^f(\mathcal{A})$
dont les termes $I^n$ sont injectifs filtrés et qui est borné inférieurement.
En fait, on peut choisir $\alpha$ de sorte que chaque $\alpha^n$ soit
un monomorphisme strict.
\end{exercise}

\begin{exercise}
\label{exercises-exercise-morphisms-lift}
Soit $\mathcal{A}$ une catégorie abélienne.
Considérons un diagramme dont les flèches pleines sont données
$$
\xymatrix{
K^\bullet \ar[r]_\alpha \ar[d]_\gamma & L^\bullet \ar@{-->}[dl]^\beta \\
I^\bullet
}
$$
de complexes de $\text{Fil}^f(\mathcal{A})$. Supposons que
$K^\bullet$, $L^\bullet$ et $I^\bullet$ soient bornés inférieurement et que
chaque $I^n$ soit un objet injectif filtré.
Supposons aussi que $\alpha$ soit un quasi-isomorphisme filtré.
\begin{enumerate}
\item Il existe un morphisme de complexes $\beta$ rendant le diagramme
commutatif à homotopie près.
\item Si $\alpha$ est un monomorphisme strict en chaque degré,
on peut trouver un $\beta$ rendant le diagramme commutatif.
\end{enumerate}
\end{exercise}

\begin{exercise}
\label{exercises-exercise-acyclic-is-zero}
Soit $\mathcal{A}$ une catégorie abélienne.
Soient $K^\bullet$, $K^\bullet$ des complexes de $\text{Fil}^f(\mathcal{A})$.
Supposons que
\begin{enumerate}
\item $K^\bullet$ soit borné inférieurement et acyclique filtré, et
\item $I^\bullet$ soit borné inférieurement et formé d'objets injectifs filtrés.
\end{enumerate}
Alors tout morphisme $K^\bullet \to I^\bullet$ est homotope à zéro.
\end{exercise}

\begin{exercise}
\label{exercises-exercise-morphisms-equal-up-to-homotopy}
Soit $\mathcal{A}$ une catégorie abélienne.
Considérons un diagramme dont les flèches pleines sont données
$$
\xymatrix{
K^\bullet \ar[r]_\alpha \ar[d]_\gamma & L^\bullet \ar@{-->}[dl]^{\beta_i} \\
I^\bullet
}
$$
de complexes de $\text{Fil}^f(\mathcal{A})$.
Supposons que $K^\bullet$, $L^\bullet$ et $I^\bullet$ soient bornés inférieurement et que
chaque $I^n$ soit un objet injectif
filtré. Supposons aussi que $\alpha$ soit un quasi-isomorphisme filtré.
Deux morphismes quelconques $\beta_1, \beta_2$ rendant le diagramme commutatif
à homotopie près sont homotopes.
\end{exercise}







\section{Fonctions régulières}
\label{exercises-section-regular-functions}

\begin{exercise}
\label{exercises-exercise-extra-function}
Considérons la courbe affine $X$ définie par l'équation
$t^2 = s^5 + 8$ dans $\mathbf{C}^2$ muni des coordonnées $s, t$.
Soit $x \in X$ le point de coordonnées $(1, 3)$.
Posons $U = X \setminus \{x\}$. Montrer qu'il existe une
fonction régulière sur $U$ qui ne soit pas la restriction
d'une fonction régulière sur $\mathbf{C}^2$, autrement dit qui ne soit pas
la restriction d'un polynôme en $s$ et $t$ à $U$.
\end{exercise}

\begin{exercise}
\label{exercises-exercise-no-extra-function}
Soit $n \geq 2$. Soit $E \subset \mathbf{C}^n$ une partie finie.
Montrer que toute fonction régulière sur $\mathbf{C}^n \setminus E$
est un polynôme.
\end{exercise}

\begin{exercise}
\label{exercises-exercise-cone}
Soit $X \subset \mathbf{C}^n$ une variété affine. Disons que $X$ est un
{\it cône} si $x = (a_1, \ldots, a_n) \in X$ et $\lambda \in \mathbf{C}$
entraînent $(\lambda a_1, \ldots, \lambda a_n) \in X$.
Bien entendu, si $\mathfrak p \subset \mathbf{C}[x_1, \ldots, x_n]$
est un idéal premier engendré par des polynômes homogènes en $x_1, \ldots, x_n$,
alors la variété affine $X = V(\mathfrak p) \subset \mathbf{C}^n$ est un cône.
Montrer que, réciproquement, l'idéal premier
$\mathfrak p \subset \mathbf{C}[x_1, \ldots, x_n]$
correspondant à un cône peut être engendré par des polynômes homogènes
en $x_1, \ldots, x_n$.
\end{exercise}

\begin{exercise}
\label{exercises-exercise-extra-function-cone}
Donner un exemple d'une variété affine $X \subset \mathbf{C}^n$
qui soit un cône (voir l'exercice \ref{exercises-exercise-cone})
et d'une fonction régulière $f$ sur $U = X \setminus \{(0, \ldots, 0)\}$
qui ne soit pas la restriction d'une fonction polynomiale
sur $\mathbf{C}^n$.
\end{exercise}

\begin{exercise}
\label{exercises-exercise-regular-functions}
Dans cet exercice, nous cherchons à comprendre le comportement des fonctions régulières
sur des corps non algébriquement clos. Soit $k$ un corps.
Soit $Z \subset k^n$ une partie localement fermée pour la topologie de Zariski, c'est-à-dire qu'il
existe des idéaux $I \subset J \subset k[x_1, \ldots, x_n]$ tels que
$$
Z = \{a \in k^n \mid
f(a) = 0\ \forall\ f \in I,\ \exists\ g \in J,\ g(a) \not = 0\}.
$$
Une fonction $\varphi : Z \to k$ est dite {\it régulière} si, pour tout
$z \in Z$, il existe un voisinage ouvert de Zariski $z \in U \subset Z$
et des polynômes $f, g \in k[x_1, \ldots, x_n]$ tels que
$g(u) \not = 0$ pour tout $u \in U$ et que
$\varphi(u) = f(u)/g(u)$ pour tout $u \in U$.
\begin{enumerate}
\item Si $k = \bar k$ et $Z = k^n$, montrer que les fonctions régulières sont
données par des polynômes. (Ne traiter ceci que si vous n'avez pas déjà vu cet
argument.)
\item Si $k$ est fini, montrer que (a) toute fonction $\varphi$ est régulière,
(b) l'anneau des fonctions régulières est de dimension finie sur $k$.
(Si vous le souhaitez, vous pouvez prendre $Z = k^n$, voire $n = 1$.)
\item Si $k = \mathbf{R}$, donner un exemple d'une fonction régulière sur
$Z = \mathbf{R}$ qui ne soit pas donnée par un polynôme.
\item Si $k = \mathbf{Q}_p$, donner un exemple d'une fonction régulière sur
$Z = \mathbf{Q}_p$ qui ne soit pas donnée par un polynôme.
\end{enumerate}
\end{exercise}






\section{Faisceaux}
\label{exercises-section-sheaves}

\noindent
Un morphisme $f : X \to Y$ dans une catégorie $\mathcal{C}$ est un {\it monomorphisme}
si, pour tout couple de morphismes $a, b : W \to X$, on a
$f \circ a = f \circ b \Rightarrow a = b$. Dans la catégorie
des ensembles, un monomorphisme est une application injective.

\begin{exercise}
\label{exercises-exercise-mono-sheaves-sets}
Démontrer soigneusement qu'un morphisme de faisceaux d'ensembles est un monomorphisme
(dans la catégorie des faisceaux d'ensembles) si et seulement si les applications induites
sur toutes les fibres sont injectives.
\end{exercise}

\noindent
Un morphisme $f : X \to Y$ dans une catégorie $\mathcal{C}$ est un {\it isomorphisme}
s'il existe un morphisme $g : Y \to X$ tel que $f \circ g = \text{id}_Y$
et $g \circ f = \text{id}_X$. Dans la catégorie des ensembles,
un isomorphisme est une application bijective.

\begin{exercise}
\label{exercises-exercise-isomorphism-sheaves-sets}
Démontrer soigneusement qu'un morphisme de faisceaux d'ensembles est un isomorphisme
(dans la catégorie des faisceaux d'ensembles) si et seulement si les applications induites
sur toutes les fibres sont bijectives.
\end{exercise}

\noindent
Un morphisme $f : X \to Y$ dans une catégorie $\mathcal{C}$ est un {\it épimorphisme}
si, pour tout couple de morphismes $a, b : Y \to Z$, on a
$a \circ f = b \circ f \Rightarrow a = b$. Dans la
catégorie des ensembles, un épimorphisme est une application surjective.

\begin{exercise}
\label{exercises-exercise-epi-sheaves-sets}
Démontrer soigneusement qu'un morphisme de faisceaux d'ensembles est un épimorphisme
(dans la catégorie des faisceaux d'ensembles) si et seulement si les applications induites
sur toutes les fibres sont surjectives.
\end{exercise}

\begin{exercise}
\label{exercises-exercise-adjoint-push-pull}
Soit $f : X \to Y$ une application d'espaces topologiques.
Montrer que l'image directe $f_\ast$ et l'image réciproque $f^{-1}$ des faisceaux d'{\bf ensembles}
forment un couple de foncteurs adjoints.
\end{exercise}

\begin{exercise}
\label{exercises-exercise-j-shriek}
Soit $j : U \to X$ une immersion ouverte. Montrer que
\begin{enumerate}
\item L'image réciproque $j^{-1} : \Sh(X) \to \Sh(U)$ admet un adjoint à gauche
$j_{!} : \Sh(U) \to \Sh(X)$ appelé {\it prolongement par le vide}.
\item Caractériser les fibres de $j_{!}({\mathcal G})$ pour
$\mathcal{G} \in \Sh(U)$.
\item L'image réciproque $j^{-1} : \textit{Ab}(X) \to \textit{Ab}(U)$ admet un adjoint à gauche
$j_{!} : \textit{Ab}(U) \to \textit{Ab}(X)$ appelé {\it prolongement par zéro}.
\item Caractériser les fibres de $j_{!}({\mathcal G})$ pour
$\mathcal{G} \in \textit{Ab}(U)$.
\end{enumerate}
Observer que le prolongement par zéro diffère du prolongement par le
vide !
\end{exercise}

\begin{exercise}
\label{exercises-exercise-not-locally-generated-by-sections}
Soit $X = \mathbf{R}$ muni de la topologie usuelle.
Soit $\mathcal{O}_X = \underline{\mathbf{Z}/2\mathbf{Z}}_X$.
Soit $i : Z = \{0\} \to X$ l'inclusion et soit
$\mathcal{O}_Z = \underline{\mathbf{Z}/2\mathbf{Z}}_Z$.
Démontrer les assertions suivantes (les trois premières découlent des définitions, mais, si
celles-ci ne vous sont pas claires, il convient de les expliciter) :
\begin{enumerate}
\item $i_*\mathcal{O}_Z$ est un faisceau gratte-ciel.
\item Il existe un morphisme surjectif canonique
$\underline{\mathbf{Z}/2\mathbf{Z}}_X \to
i_*\underline{\mathbf{Z}/2\mathbf{Z}}_Z$.
Noter $\mathcal{I} \subset \mathcal{O}_X$ le noyau de ce morphisme.
\item $\mathcal{I}$ est un faisceau d'idéaux de $\mathcal{O}_X$.
\item Le faisceau $\mathcal{I}$ sur $X$ ne peut pas être localement engendré
par des sections (au sens de
Modules, définition \ref{modules-definition-locally-generated}.)
\end{enumerate}
\end{exercise}

\begin{exercise}
\label{exercises-exercise-quotient-j-shriek-Z}
Soit $X$ un espace topologique.
Soit ${\mathcal F}$ un faisceau abélien sur $X$. Montrer
que ${\mathcal F}$ est le quotient d'une somme directe (éventuellement très grande)
de faisceaux dont tous les termes sont de la forme
$$
j_{!}(\underline{{\mathbf Z}}_U)
$$
où $U \subset X$ est ouvert et où $\underline{{\mathbf Z}}_U$ désigne le
faisceau constant de valeur ${\mathbf Z}$ sur $U$.
\end{exercise}

\begin{remark}
\label{exercises-remark-direct-sum-stalk-abelian}
Soit $X$ un espace topologique.
Dans la catégorie des faisceaux abéliens, la somme directe
d'une famille de faisceaux $\{{\mathcal F}_i\}_{i\in I}$ est le faisceau associé au
préfaisceau $U \mapsto \oplus {\mathcal F}_i(U)$. Par conséquent, la fibre de
la somme directe en un point $x$ est la somme directe des fibres des
${\mathcal F}_i$ en $x$.
\end{remark}

\begin{exercise}
\label{exercises-exercise-product-over-points}
Soit $X$ un espace topologique. Supposons donnée une famille de
groupes abéliens $A_x$ indexée par $x \in X$. Montrer que la règle
$U \mapsto \prod_{x \in U} A_x$, munie des applications de restriction évidentes,
définit un faisceau $\mathcal{G}$ de groupes abéliens. Montrer, à l'aide d'un exemple,
que l'on n'a généralement pas $\mathcal{G}_x = A_x$ pour $x \in X$.
\end{exercise}

\begin{exercise}
\label{exercises-exercise-modified-product-over-points}
Soient $X$, $A_x$, $\mathcal{G}$ comme dans
l'exercice \ref{exercises-exercise-product-over-points}.
Soit $\mathcal{B}$ une base de la topologie de $X$, voir
Topologie, définition \ref{topology-definition-base}.
Pour $U \in \mathcal{B}$, soit $A_U$ un sous-groupe
$A_U \subset \mathcal{G}(U) = \prod_{x \in U} A_x$. Supposons que, si
$U \subset V$ avec $U, V \in \mathcal{B}$, les applications de restriction
envoient $A_V$ dans $A_U$. Pour tout ouvert $U \subset X$, posons
$$
\mathcal{F}(U) =
\left\{
(s_x)_{x \in U}
\middle|
\begin{matrix}
\text{ pour tout }x\text{ dans }U\text{ il existe } V \in \mathcal{B} \\
x \in V \subset U\text{ tel que } (s_y)_{y \in V} \in A_V
\end{matrix}
\right\}
$$
Montrer que $\mathcal{F}$ définit un faisceau de groupes abéliens sur $X$.
Montrer, à l'aide d'un exemple, que l'on n'a généralement pas
$\mathcal{F}(U) = A_U$ pour $U \in \mathcal{B}$.
\end{exercise}

\begin{exercise}
\label{exercises-exercise-exact-but-not-a-stalk-functor}
Donner un exemple d'un espace topologique $X$ et d'un foncteur
$$
F : \Sh(X) \longrightarrow \textit{Ensembles}
$$
qui soit exact (c'est-à-dire qui commute aux produits finis et aux égaliseurs, ainsi
qu'aux coproduits finis et aux coégaliseurs ; voir Catégories, section
\ref{categories-section-exact-functor}), mais pour lequel il n'existe aucun point $x \in X$
tel que $F$ soit isomorphe au foncteur fibre
$\mathcal{F} \mapsto \mathcal{F}_x$.
\end{exercise}




\section{Schémas}
\label{exercises-section-schemes}

\noindent
Soit $LRS$ la catégorie des espaces localement annelés.
Un schéma affine est un objet de $LRS$ isomorphe dans $LRS$ à
une paire de la forme $(\Spec(A), {\mathcal O}_{\Spec(A)})$.
Un schéma est un
objet $(X, {\mathcal O}_X)$ de $LRS$ tel que tout point $x\in X$
possède un voisinage ouvert $U \subset X$ pour lequel la paire
$(U, {\mathcal O}_X|_U)$ est un schéma affine.

\begin{exercise}
\label{exercises-exercise-one-point}
Trouver un espace localement annelé ayant $1$ point qui ne soit pas un schéma.
\end{exercise}

\begin{exercise}
\label{exercises-exercise-two-points}
Supposons que $X$ soit un schéma dont l'espace topologique
sous-jacent possède 2 points. Montrer que $X$ est un schéma affine.
\end{exercise}

\begin{exercise}
\label{exercises-exercise-discrete-finite-set-points}
Supposons que $X$ soit un schéma dont l'espace topologique sous-jacent est un
ensemble discret fini. Montrer que $X$ est un schéma affine.
\end{exercise}

\begin{exercise}
\label{exercises-exercise-three-points}
Montrer qu'il existe un schéma non affine possédant trois points.
\end{exercise}

\begin{exercise}
\label{exercises-exercise-quasi-compact-closed-point}
Supposons que $X$ soit un schéma quasi-compact non vide.
Montrer que $X$ possède un point fermé.
\end{exercise}

\begin{remark}
\label{exercises-remark-open-immersion}
Si $(X, {\mathcal O}_X)$ est un espace annelé et si $U \subset X$
est une partie ouverte, alors $(U, {\mathcal O}_X|_U)$ est un espace annelé. Notation :
${\mathcal O}_U = {\mathcal O}_X|_U$. Il existe un morphisme canonique
d'espaces annelés
$$
j : (U, {\mathcal O}_U) \longrightarrow (X, {\mathcal O}_X).
$$
Si $(X, {\mathcal O}_X)$ est un espace localement annelé, il en est de même de
$(U, {\mathcal O}_U)$, et
$j$ est un morphisme d'espaces localement annelés. Si $(X, {\mathcal O}_X)$
est un schéma,
il en est de même de $(U, {\mathcal O}_U)$, et $j$ est un morphisme de schémas. On dit
que
$(U, {\mathcal O}_U)$ est un {\it sous-schéma ouvert} de $(X, {\mathcal O}_X)$
et que
$j$ est une {\it immersion ouverte}. Plus généralement, tout morphisme
$j' : (V, {\mathcal O}_V) \to (X, {\mathcal O}_X)$ qui est {\it isomorphe}
à un
morphisme $j : (U, {\mathcal O}_U) \to (X, {\mathcal O}_X)$ comme ci-dessus est
appelé une
immersion ouverte.
\end{remark}

\begin{exercise}
\label{exercises-exercise-open-affine-not-affine}
Donner un exemple d'un schéma affine $(X, {\mathcal O}_X)$
et d'un ouvert $U \subset X$ tels que $(U, {\mathcal O}_X|U)$ ne soit pas un
schéma affine.
\end{exercise}

\begin{exercise}
\label{exercises-exercise-morphism-does-not-extend}
Donner un exemple d'une paire de schémas affines
$(X, {\mathcal O}_X)$, $(Y, {\mathcal O}_Y)$,
d'un sous-schéma ouvert $(U, {\mathcal O}_X|_U)$
de $X$ et d'un morphisme de schémas
$(U, {\mathcal O}_X|_U) \to (Y, {\mathcal O}_Y)$
qui ne se prolonge pas en un morphisme de schémas
$(X, {\mathcal O}_X) \to (Y, {\mathcal O}_Y)$.
\end{exercise}

\begin{exercise}
\label{exercises-exercise-closed-subscheme-does-not-extend}
(Cet exercice est assez difficile.)
Donner un exemple d'un schéma $X$, d'un sous-schéma ouvert $U \subset X$
et d'un sous-schéma fermé $Z \subset U$ tel que $Z$ ne se prolonge pas
en un sous-schéma fermé de $X$.
\end{exercise}

\begin{exercise}
\label{exercises-exercise-not-morphism-schemes}
Donner un exemple d'un schéma $X$, d'un corps $K$ et d'un
morphisme d'espaces annelés $\Spec(K) \to X$ qui
N'EST PAS un morphisme de schémas.
\end{exercise}

\begin{exercise}
\label{exercises-exercise-just-kidding}
Faire tous les exercices de \cite[chapitre II]{H},
sections 1 et 2...\ \ Je plaisante !
\end{exercise}

\begin{definition}
\label{exercises-definition-integral}
Un schéma $X$ est dit {\it intègre} si $X$ est non vide et si,
pour tout ouvert affine non vide $U \subset X$, l'anneau
$\Gamma(U, \mathcal{O}_X) = \mathcal{O}_X(U)$ est intègre.
\end{definition}

\begin{exercise}
\label{exercises-exercise-morphism-integral-schemes-surjective-stalks-not-closed}
Donner un exemple d'un morphisme de schémas {\it intègres}
$f : X \to Y$ tel que les homomorphismes induits ${\mathcal O}_{Y, f(x)}
\to {\mathcal O}_{X, x}$ soient surjectifs pour tout $x\in X$, mais que $f$
ne soit pas une immersion fermée.
\end{exercise}

\begin{exercise}
\label{exercises-exercise-fibre-product-affines-not-affine}
Donner un exemple d'un produit fibré $X \times_S Y$ tel que $X$ et $Y$
soient affines, mais que $X \times_S Y$ ne le soit pas.
\end{exercise}

\begin{remark}
\label{exercises-remark-separated-base-fibre-product-affines-affine}
Cela ne peut en fait pas se produire lorsque $S$ est séparé. Savez-vous pourquoi ?
\end{remark}

\begin{exercise}
\label{exercises-exercise-not-geometrically-integral}
Donner un exemple d'un schéma
$V$ qui soit un schéma intègre de dimension 1 et de type fini
sur ${\mathbf Q}$ tel que
$\Spec({\mathbf C}) \times_{\Spec({\mathbf Q})} V$
ne soit pas intègre.
\end{exercise}

\begin{exercise}
\label{exercises-exercise-not-geometrically-reduced}
Donner un exemple d'un schéma
$V$ qui soit un schéma intègre de dimension 1 et de type fini
sur un corps $k$ tel que $\Spec(k') \times_{\Spec(k)} V$
ne soit pas réduit pour une certaine extension finie de corps $k'/k$.
\end{exercise}

\begin{remark}
\label{exercises-remark-affine-dimension}
Si votre schéma est affine, sa dimension est
égale à la dimension de Krull de l'anneau sous-jacent. Vous pouvez donc
utiliser les résultats du semestre dernier pour calculer la dimension.
\end{remark}






\section{Morphismes}
\label{exercises-section-morphisms}

\noindent
Une question importante est de savoir, étant donné un morphisme $\pi : X \to S$,
si ce morphisme possède une section ou une section rationnelle.
Voici quelques exercices à titre d'exemple.

\begin{exercise}
\label{exercises-exercise-no-section}
Considérons le morphisme de schémas
$$
\pi :
X = \Spec(\mathbf{C}[x, t, 1/xt])
\longrightarrow
S = \Spec(\mathbf{C}[t]).
$$
\begin{enumerate}
\item Montrer qu'il n'existe aucun morphisme $\sigma : S \to X$
tel que $\pi \circ \sigma = \text{id}_S$.
\item Montrer qu'il existe un ouvert non vide $U \subset S$ et
un morphisme $\sigma : U \to X$ tel que $\pi \circ \sigma = \text{id}_U$.
\end{enumerate}
\end{exercise}

\begin{exercise}
\label{exercises-exercise-no-rational-section}
Considérons le morphisme de schémas
$$
\pi :
X = \Spec(\mathbf{C}[x, t]/(x^2 + t))
\longrightarrow
S = \Spec(\mathbf{C}[t]).
$$
Montrer qu'il n'existe aucun ouvert non vide $U \subset S$ et
aucun morphisme $\sigma : U \to X$ tel que $\pi \circ \sigma = \text{id}_U$.
\end{exercise}

\begin{exercise}
\label{exercises-exercise-has-rational-section}
Soient $A, B, C \in \mathbf{C}[t]$ des polynômes non nuls.
Considérons le morphisme de schémas
$$
\pi :
X = \Spec(\mathbf{C}[x, y, t]/(A + Bx^2 + Cy^2))
\longrightarrow
S = \Spec(\mathbf{C}[t]).
$$
Montrer qu'il existe un ouvert non vide $U \subset S$ et
un morphisme $\sigma : U \to X$ tel que $\pi \circ \sigma = \text{id}_U$.
(Indication : symboliquement, écrire $x = X/Z$, $y = Y/Z$ avec
$X, Y, Z \in \mathbf{C}[t]$ de degré $\leq d$, pour un certain $d$,
et déterminer la condition pour que cela résolve l'équation.
Montrer ensuite, au moyen de la théorie de la dimension, que si $d >> 0$, on peut trouver
des $X, Y, Z$ non nuls qui résolvent l'équation.)
\end{exercise}

\begin{remark}
\label{exercises-remark-tsen}
L'exercice \ref{exercises-exercise-has-rational-section}
est un cas particulier du « théorème de Tsen ».
L'exercice \ref{exercises-exercise-no-section-curve} montre que la
méthode est limitée aux équations de faible degré (aux coniques lorsque la base et
la fibre sont de dimension 1).
\end{remark}

\begin{exercise}
\label{exercises-exercise-no-section-curve}
Considérons le morphisme de schémas
$$
\pi :
X = \Spec(\mathbf{C}[x, y, t]
/(1 + t x^3 + t^2 y^3))
\longrightarrow
S = \Spec(\mathbf{C}[t])
$$
Montrer qu'il n'existe aucun ouvert non vide $U \subset S$ et
aucun morphisme $\sigma : U \to X$ tel que $\pi \circ \sigma = \text{id}_U$.
\end{exercise}

\begin{exercise}
\label{exercises-exercise-no-section-surface}
Considérons les schémas
$$
X = \Spec(\mathbf{C}[\{x_i\}_{i = 1}^{8}, s, t]
/(1 + s x_1^3 + s^2 x_2^3 +
t x_3^3 + st x_4^3 + s^2t x_5^3 +
t^2 x_6^3 + st^2 x_7^3 + s^2t^2 x_8^3))
$$
et
$$
S = \Spec(\mathbf{C}[s, t])
$$
ainsi que le morphisme de schémas
$$
\pi : X \longrightarrow S
$$
Montrer qu'il n'existe aucun ouvert non vide $U \subset S$ et
aucun morphisme $\sigma : U \to X$ tel que $\pi \circ \sigma = \text{id}_U$.
\end{exercise}

\begin{exercise}
\label{exercises-exercise-for-number-theorists}
(Pour les théoriciens des nombres.) Donner un exemple d'un sous-schéma fermé
$$
Z \subset \Spec\left({\mathbf Z}[x, \frac{1 }{ x(x-1)(2x-1)}]\right)
$$
tel que le morphisme $Z \to \Spec({\mathbf Z})$ soit fini
et surjectif.
\end{exercise}

\begin{exercise}
\label{exercises-exercise-quasi-section}
Si la théorie des nombres ne vous plaît pas, vous pouvez essayer la
variante où l'on considère
$$
\Spec\left({\mathbf F}_p[t, x, \frac{1 }{ x(x-t)(tx-1)}]\right)
\longrightarrow
\Spec({\mathbf F}_p[t])
$$
et chercher un sous-schéma fermé du schéma supérieur
dont le morphisme vers celui du bas soit fini et surjectif. (Il existe une
raison théorique de supposer ici le corps de base fini, bien que
cela ne soit peut-être pas nécessaire dans ce cas particulier.)
\end{exercise}

\begin{remark}
\label{exercises-remark-interpretation-skolem-noether}
Dans les exercices \ref{exercises-exercise-for-number-theorists} et
\ref{exercises-exercise-quasi-section}, on considère un morphisme $f : X \to S$
où $S$ est le spectre d'un anneau de Dedekind $A$, le morphisme $f$ étant plat et
surjectif, à fibre générique géométriquement irréductible. Dans les deux cas,
on conclut qu'il existe un morphisme fini et surjectif $S' \to S$
tel qu'après changement de base, $X_{S'} \to S'$ possède une section.
Ce résultat vaut lorsque $A$ est excellent, que ses corps résiduels
aux idéaux maximaux sont des extensions algébriques de corps finis,
et que $\Pic(A')$ est de torsion lorsque $A'$ est la clôture intégrale
de $A$ dans une extension finie de son corps des fractions; voir \cite{MB1}.
C'est par exemple le cas si $A = \mathbf{Z}$ ou $A = \mathbf{F}_p[x]$, ou pour une extension
finie de l'un de ces anneaux. Cependant, il
existe un anneau de Dedekind $A$ dont les corps résiduels aux
idéaux maximaux sont finis et dont le groupe de Picard possède un élément non de torsion; voir \cite{Goldman},
et le résultat est faux pour le spectre d'un tel anneau.
L'exercice \ref{exercises-exercise-no-quasi-section} donne un exemple géométrique
où cela ne fonctionne pas.
\end{remark}

\begin{exercise}
\label{exercises-exercise-no-quasi-section}
Démontrer qu'il existe un $f \in \mathbf{C}[x, t]$ qui n'est divisible
par $t - \alpha$ pour aucun $\alpha \in \mathbf{C}$ et tel qu'il
n'existe aucun $Z \subset \Spec(\mathbf{C}[x, t, 1/f])$
dont le morphisme vers $\Spec(\mathbf{C}[t])$ soit fini et surjectif.
(Je pense que $f(x, t) = (xt - 2)(x - t + 3)$ convient. Il n'est pas si facile,
me semble-t-il, de montrer que ce candidat possède la propriété requise.)
\end{exercise}

\begin{exercise}
\label{exercises-exercise-finite}
Soit $A \to B$ un homomorphisme d'anneaux de type fini. Supposons que
$\Spec(B) \to \Spec(A)$ se factorise par une immersion fermée
$\Spec(B) \to \mathbf{P}^n_A$ pour un certain $n$.
Démontrer que $A \to B$ est un homomorphisme fini, c'est-à-dire que
$B$ est un $A$-module de type fini.
Indication : si $A$ est noethérien (supposez-le simplement),
on peut utiliser le fait que $H^i(Z, \mathcal{O}_Z)$, pour $i \in \mathbf{Z}$,
est un $A$-module de type fini pour tout sous-schéma fermé
$Z \subset \mathbf{P}^n_A$.
\end{exercise}

\begin{exercise}
\label{exercises-exercise-projective-finite}
Soit $k$ un corps algébriquement clos.
Soit $f : X \to Y$ un morphisme de variétés
projectives tel que $f^{-1}(\{y\})$ soit fini pour
tout point fermé $y \in Y$.
Démontrer que $f$ est fini en tant que morphisme de schémas.
Indications : (a) la finitude est une propriété locale,
(b) essayer de se ramener à l'exercice \ref{exercises-exercise-finite}, et
(c) utiliser une immersion fermée $X \to \mathbf{P}^n_k$ pour obtenir
une immersion fermée $X \to \mathbf{P}^n_Y$ au-dessus de $Y$.
\end{exercise}






\section{Espaces tangents}
\label{exercises-section-tangent-space}

\begin{definition}
\label{exercises-definition-dual-numbers}
Pour tout anneau $R$, nous notons $R[\epsilon]$ l'anneau
des {\it nombres duaux}. Comme $R$-module, il est libre de
base $1$, $\epsilon$. La structure d'anneau s'obtient en posant
$\epsilon^2 = 0$.
\end{definition}

\begin{exercise}
\label{exercises-exercise-tangent-space-Zariski}
Soit $f : X \to S$ un morphisme de schémas.
Soit $x \in X$ un point, et soit $s = f(x)$.
Considérer le diagramme commutatif suivant, où les flèches pleines sont données
$$
\xymatrix{
\Spec(\kappa(x)) \ar[r] \ar[dr] \ar@/^1pc/[rr] &
\Spec(\kappa(x)[\epsilon]) \ar@{.>}[r] \ar[d]&
X \ar[d] \\
&
\Spec(\kappa(s)) \ar[r] &
S
}
$$
où la flèche courbe est le morphisme canonique de
$\Spec(\kappa(x))$ vers $X$.
Si $\kappa(x) = \kappa(s)$, montrer que l'ensemble des flèches
pointillées qui rendent le diagramme commutatif est en bijection
avec l'ensemble des applications linéaires
$$
\Hom_{\kappa(x)}(
\frac{\mathfrak m_x}{\mathfrak m_x^2 + \mathfrak m_s\mathcal{O}_{X, x}},
\kappa(x))
$$
Autrement dit, décrire une telle bijection.
(Cela vaut plus généralement si $\kappa(x) \supset \kappa(s)$ est une
extension algébrique séparable.)
\end{exercise}

\begin{definition}
\label{exercises-definition-tangent-space}
Soit $f : X \to S$ un morphisme de schémas.
Soit $x \in X$. Nous appelons l'ensemble des flèches pointillées
de l'exercice \ref{exercises-exercise-tangent-space-Zariski}
l'{\it espace tangent de $X$ sur $S$}
et nous le notons $T_{X/S, x}$. Un élément de cet
espace est appelé {\it vecteur tangent} à $X/S$ en $x$.
\end{definition}

\begin{exercise}
\label{exercises-exercise-simple-push-out}
Pour tout corps $K$, démontrer que le diagramme
$$
\xymatrix{
\Spec(K) \ar[r] \ar[d] & \Spec(K[\epsilon_1]) \ar[d] \\
\Spec(K[\epsilon_2]) \ar[r] &
\Spec(K[\epsilon_1, \epsilon_2]/(\epsilon_1\epsilon_2))
}
$$
est cocartésien dans la catégorie des schémas.
(Ici $\epsilon_i^2 = 0$ comme précédemment.)
\end{exercise}

\begin{exercise}
\label{exercises-exercise-tangent-space-vectors-space}
Soit $f : X \to S$ un morphisme de schémas.
Soit $x \in X$. Définir l'addition des vecteurs tangents,
à l'aide de l'exercice \ref{exercises-exercise-simple-push-out}
et d'un morphisme convenable
$$
\Spec(K[\epsilon])
\longrightarrow
\Spec(K[\epsilon_1, \epsilon_2]/(\epsilon_1\epsilon_2)).
$$
Définir de même la multiplication des vecteurs tangents par les scalaires (c'est plus simple).
Montrer que vos constructions munissent $T_{X/S, x}$ d'une structure
d'espace vectoriel sur $\kappa(x)$.
\end{exercise}

\begin{exercise}
\label{exercises-exercise-compute-TS}
Soit $k$ un corps. Considérer
le morphisme structural $f : X = \mathbf{A}^1_k \to \Spec(k) = S$.
\begin{enumerate}
\item Soit $x \in X$ un point fermé. Quelle est la dimension de
$T_{X/S, x}$ ?
\item Soit $\eta \in X$ le point générique. Quelle est la dimension
de $T_{X/S, \eta}$ ?
\item Considérer maintenant $X$ comme un schéma sur $\Spec(\mathbf{Z})$.
Quelles sont les dimensions de $T_{X/\mathbf{Z}, x}$ et de
$T_{X/\mathbf{Z}, \eta}$ ?
\end{enumerate}
\end{exercise}

\begin{remark}
\label{exercises-remark-tangent-space-relative}
L'exercice \ref{exercises-exercise-compute-TS} explique pourquoi il est nécessaire
de considérer l'espace tangent à $X$ relativement à $S$ pour obtenir une bonne notion.
\end{remark}

\begin{exercise}
\label{exercises-exercise-compute-TS-field}
Considérer le morphisme de schémas
$$
f : X = \Spec(\mathbf{F}_p(t))
\longrightarrow
\Spec(\mathbf{F}_p(t^p)) = S
$$
Calculer l'espace tangent de $X/S$ à l'unique point de $X$.
N'est-ce pas étrange ? Que se passe-t-il, à votre avis, si l'on considère le morphisme
de schémas correspondant à $\mathbf{F}_p[t^p] \to \mathbf{F}_p[t]$ ?
\end{exercise}

\begin{exercise}
\label{exercises-exercise-compute-TS-cusp}
Soit $k$ un corps.
Calculer l'espace tangent de $X/k$ au point $x = (0, 0)$
où $X = \Spec(k[x, y]/(x^2 - y^3))$.
\end{exercise}

\begin{exercise}
\label{exercises-exercise-map-tangent-spaces}
Soit $f : X \to Y$ un morphisme de schémas au-dessus de $S$.
Soit $x \in X$ un point. Poser $y = f(x)$.
Supposons que l'application naturelle $\kappa(y) \to \kappa(x)$ soit
bijective. Montrer, à partir de la définition,
que $f$ induit une application linéaire naturelle
$$
\text{d}f : T_{X/S, x} \longrightarrow T_{Y/S, y}.
$$
Identifier celle-ci à l'application obtenue sur les anneaux locaux à l'aide de
l'exercice \ref{exercises-exercise-tangent-space-Zariski} lorsque $\kappa(x) = \kappa(s)$.
\end{exercise}

\begin{exercise}
\label{exercises-exercise-Jacobian}
Soit $k$ un corps algébriquement clos.
Soit
\begin{eqnarray*}
f : \mathbf{A}_k^n & \longrightarrow & \mathbf{A}^m_k \\
(x_1, \ldots, x_n) & \longmapsto & (f_1(x_i), \ldots, f_m(x_i))
\end{eqnarray*}
un morphisme de schémas au-dessus de $k$. Il est défini par
$m$ polynômes $f_1, \ldots, f_m$ en $n$ variables.
Considérer la matrice
$$
A = \left( \frac{\partial f_j}{\partial x_i} \right)
$$
Soit $x \in \mathbf{A}^n_k$ un point fermé.
Poser $y =  f(x)$. Montrer que l'application sur les espaces tangents
$T_{\mathbf{A}^n_k/k, x} \to T_{\mathbf{A}^m_k/k, y}$
est donnée par la valeur de la matrice $A$ au point $x$.
\end{exercise}














\section{Faisceaux quasi-cohérents}
\label{exercises-section-quasi-coherent}

\begin{definition}
\label{exercises-definition-quasi-coherent}
Soit $X$ un schéma.
Un faisceau $\mathcal{F}$ de $\mathcal{O}_X$-modules est {\it quasi-cohérent}
si, pour tout ouvert affine $\Spec(R) = U \subset X$, la restriction
$\mathcal{F}|_U$ est de la forme $\widetilde M$ pour un certain $R$-module
$M$.
\end{definition}

\noindent
Il suffit de vérifier cette condition sur les membres d'un
recouvrement ouvert affine de $X$.
Voir Schémas, section \ref{schemes-section-quasi-coherent}
pour davantage de résultats.

\begin{definition}
\label{exercises-definition-specialization}
Soit $X$ un espace topologique. Soient $x, x' \in X$.
On dit que $x$ est une {\it spécialisation} de $x'$
si et seulement si $x \in \overline{\{x'\}}$.
\end{definition}

\begin{exercise}
\label{exercises-exercise-quasi-coherent-specialization-points}
Soit $X$ un schéma. Soient $x, x' \in X$. Soit $\mathcal{F}$
un faisceau quasi-cohérent de $\mathcal{O}_X$-modules.
Supposons que (a) $x$ soit une spécialisation de $x'$ et que (b)
$\mathcal{F}_{x'} \not = 0$. Montrer que $\mathcal{F}_x \not = 0$.
\end{exercise}

\begin{exercise}
\label{exercises-exercise-O-module-specialization-points}
Trouver un exemple d'un schéma $X$, de points $x, x' \in X$ et
d'un faisceau de $\mathcal{O}_X$-modules
$\mathcal{F}$ tels que (a) $x$ soit une spécialisation de $x'$ et que (b)
$\mathcal{F}_{x'} \not = 0$ et $\mathcal{F}_x = 0$.
\end{exercise}

\begin{definition}
\label{exercises-definition-Noetherian-scheme}
Un schéma $X$ est dit {\it localement noethérien} si et seulement si,
pour tout point $x \in X$, il existe un ouvert affine
$\Spec(R) = U \subset X$ tel que $R$ soit noethérien.
Un schéma est {\it noethérien} s'il est localement noethérien et quasi-compact.
\end{definition}

\noindent
Si $X$ est localement noethérien, tout ouvert affine de $X$
est le spectre d'un anneau noethérien, voir
Propriétés, lemme \ref{properties-lemma-locally-Noetherian}.

\begin{definition}
\label{exercises-definition-coherent}
Soit $X$ un schéma localement noethérien.
Soit $\mathcal{F}$ un faisceau quasi-cohérent de
$\mathcal{O}_X$-modules. On dit que $\mathcal{F}$ est {\it cohérent}
si, pour tout point $x \in X$, il existe un ouvert affine
$\Spec(R) = U \subset X$ tel que $\mathcal{F}|_U$
soit isomorphe à $\widetilde M$ pour un $R$-module $M$ de type fini.
\end{definition}

\begin{exercise}
\label{exercises-exercise-extend-quasi-coherent}
Soit $X = \Spec(R)$ un schéma affine.
\begin{enumerate}
\item Soit $f \in R$. Soit $\mathcal{G}$ un
faisceau quasi-cohérent de $\mathcal{O}_{D(f)}$-modules
sur le sous-schéma ouvert $D(f)$.
Montrer que $\mathcal{G} = \mathcal{F}|_{D(f)}$ pour un certain
faisceau quasi-cohérent de $\mathcal{O}_X$-modules
$\mathcal{F}$.
\item Soit $I \subset R$ un idéal.
Soit $i : Z \to X$ le sous-schéma fermé de $X$ correspondant
à $I$. Soit $\mathcal{G}$ un
faisceau quasi-cohérent de $\mathcal{O}_Z$-modules
sur le sous-schéma fermé $Z$.
Montrer que $\mathcal{G} = i^*\mathcal{F}$ pour un certain
faisceau quasi-cohérent de $\mathcal{O}_X$-modules $\mathcal{F}$.
(Pourquoi cela est-il trivial ?)
\item Supposons que $R$ soit noethérien.
Soit $f \in R$. Soit $\mathcal{G}$ un
faisceau cohérent de $\mathcal{O}_{D(f)}$-modules
sur le sous-schéma ouvert $D(f)$.
Montrer que $\mathcal{G} = \mathcal{F}|_{D(f)}$ pour un certain
faisceau cohérent de $\mathcal{O}_X$-modules $\mathcal{F}$.
\end{enumerate}
\end{exercise}

\begin{remark}
\label{exercises-remark-extend-off-open}
Si $U \to X$ est une immersion quasi-compacte, tout
faisceau quasi-cohérent sur $U$ est la restriction d'un
faisceau quasi-cohérent sur $X$.
Si $X$ est un schéma noethérien et si $U \subset X$ est ouvert,
tout faisceau cohérent sur $U$ est la restriction d'un
faisceau cohérent sur $X$.
L'exercice ci-dessus est bien sûr plus facile et ne devrait pas faire appel à ces résultats généraux.
\end{remark}
















\section{Proj et schémas projectifs}
\label{exercises-section-proj}

\begin{exercise}
\label{exercises-exercise-graded-ring-specified-result}
Donner des exemples d'anneaux gradués $S$ tels que
\begin{enumerate}
\item $\text{Proj}(S)$ soit affine et non vide, et
\item $\text{Proj}(S)$ soit intègre, non vide, mais non isomorphe
à ${\mathbf P}^n_A$ pour aucun $n\geq 0$ ni aucun anneau $A$.
\end{enumerate}
\end{exercise}

\begin{exercise}
\label{exercises-exercise-nonconstant-morphism-proj}
Donner un exemple de morphisme non constant
de schémas ${\mathbf P}^1_{\mathbf C} \to {\mathbf P}^5_{\mathbf C}$ au-dessus de
$\Spec({\mathbf C})$.
\end{exercise}

\begin{exercise}
\label{exercises-exercise-isomorphism-P1-conic}
Donner un exemple d'isomorphisme de schémas
$$
{\mathbf P}^1_{\mathbf C} \to
\text{Proj}({\mathbf C}[X_0, X_1, X_2]/(X_0^2 + X_1^2 + X_2^2))
$$
\end{exercise}

\begin{exercise}
\label{exercises-exercise-family-special-fibre-different}
Donner un exemple de morphisme de schémas
$f : X \to {\mathbf A}^1_{\mathbf C} = \Spec({\mathbf C}[T])$
tel que la fibre (schématique) $X_t$ de $f$ au-dessus de
$t \in {\mathbf A}^1_{\mathbf C}$ soit
(a) isomorphe à ${\mathbf P}^1_{\mathbf C}$ lorsque $t$
est un point fermé distinct de $0$, et
(b) non isomorphe à ${\mathbf P}^1_{\mathbf C}$ lorsque $t = 0$.
Nous appellerons $X_0$ la {\it fibre spéciale} du morphisme.
On peut le faire de très nombreuses façons. Essayer de donner des exemples qui
satisfassent (chacune des) contraintes supplémentaires suivantes (à moins
que cela ne soit impossible) :
\begin{enumerate}
\item Peut-on le faire avec une fibre spéciale projective ?
\item Peut-on le faire avec une fibre spéciale irréductible et projective ?
\item Peut-on le faire avec une fibre spéciale intègre et projective ?
\item Peut-on le faire avec une fibre spéciale lisse et projective ?
\item Peut-on le faire avec $f$ plat ?
Cela signifie simplement que, pour tout ouvert affine $\Spec(A) \subset X$,
l'homomorphisme d'anneaux induit $\mathbf{C}[t] \to A$ est plat, ce qui signifie ici
que tout polynôme non nul en $t$ est un non-diviseur de zéro dans $A$.
\item Peut-on le faire avec $f$ plat et projectif ?
\item Peut-on le faire avec $f$ plat et projectif et avec une fibre spéciale réduite ?
\item Peut-on le faire avec $f$ plat et projectif et avec une fibre spéciale irréductible ?
\item Peut-on le faire avec $f$ plat et projectif et avec une fibre spéciale intègre ?
\end{enumerate}
Que se passe-t-il, selon vous, lorsque l'on remplace
${\mathbf P}^1_{\mathbf C}$ par une autre variété sur ${\mathbf C}$ ?
(Cela peut devenir très difficile selon celle des variantes ci-dessus que l'on
demande.)
\end{exercise}

\begin{exercise}
\label{exercises-exercise-affine-onto-projective-space}
Soit $n \geq 1$ un entier positif quelconque.
Donner un exemple de morphisme surjectif
$X \to {\mathbf P}^n_{\mathbf C}$ avec $X$ affine.
\end{exercise}

\begin{exercise}
\label{exercises-exercise-morphism-proj}
Morphismes de $\text{Proj}$. Soient $R$ et $S$ des anneaux gradués.
Supposons donné un homomorphisme d'anneaux
$$
\psi : R \to S
$$
et un entier $e \geq 1$ tel que $\psi(R_d) \subset S_{de}$
pour tout $d \geq 0$. (Selon nos conventions, ce n'est pas un homomorphisme
d'anneaux gradués, sauf si $e = 1$.)
\begin{enumerate}
\item Pour quels éléments $\mathfrak p \in \text{Proj}(S)$ existe-t-il
un point correspondant bien défini dans $\text{Proj}(R)$ ? Autrement dit,
trouver un ouvert convenable $U \subset \text{Proj}(S)$ tel que $\psi$ définisse
une application continue $r_\psi : U \to \text{Proj}(R)$.
\item Donner un exemple où $U \not = \text{Proj}(S)$.
\item Donner un exemple où $U = \text{Proj}(S)$.
\item (Ne pas rédiger ceci.) Se convaincre que
l'application continue $U \to \text{Proj}(R)$ est canoniquement munie
d'un morphisme de faisceaux qui fait de $r_\psi$ un morphisme de schémas :
$$
\text{Proj}(S) \supset U \longrightarrow \text{Proj}(R).
$$
\item Que peut-on dire de ce morphisme si
$R = \bigoplus_{d \geq 0} S_{de}$ (considéré comme anneau gradué où
$S_e$, $S_{2e}$, etc. sont respectivement de degrés $1$, $2$, etc.) et si $\psi$
est l'inclusion ?
\end{enumerate}
\end{exercise}

\noindent
{\bf Notation.} Soit $R$ un anneau gradué comme ci-dessus et
soit $n \geq 0$ un entier. Soit $X = \text{Proj}(R)$. Il existe alors un unique
${\mathcal O}_X$-module quasi-cohérent ${\mathcal O}_X(n)$ sur $X$ tel que,
pour tout élément homogène $f \in R$ de degré positif, la restriction
${\mathcal O}_X |_{D_{+}(f)}$ soit le faisceau quasi-cohérent associé au
$R_{(f)} = (R_f)_0$-module $(R_f)_n$ ($ = $ les éléments homogènes de degré
$n$ dans $R_f = R[1/f]$). Voir \cite[page 116+]{H}. Noter qu'il existe
des morphismes naturels
$$
{\mathcal O}_X(n_1) \otimes_{{\mathcal O}_X} {\mathcal O}_X(n_2)
\longrightarrow
{\mathcal O}_X(n_1 + n_2)
$$

\begin{exercise}
\label{exercises-exercise-pathologies-proj}
Pathologies dans $\text{Proj}$.
Donner des exemples de $R$ comme ci-dessus tels que
\begin{enumerate}
\item ${\mathcal O}_X(1)$ ne soit pas un ${\mathcal O}_X$-module inversible.
\item ${\mathcal O}_X(1)$ soit inversible, mais que le
morphisme naturel ${\mathcal O}_X(1) \otimes_{{\mathcal O}_X} {\mathcal O}_X(1) \to
{\mathcal O}_X(2)$ ne soit PAS un isomorphisme.
\end{enumerate}
\end{exercise}

\begin{exercise}
\label{exercises-exercise-finitely-many-points-in-affine}
Soit $S$ un anneau gradué.
Soit $X = \text{Proj}(S)$.
Montrer que tout ensemble fini de points de $X$ est contenu dans un ouvert
affine standard.
\end{exercise}

\begin{exercise}
\label{exercises-exercise-prepare-glueing}
Soit $S$ un anneau gradué.
Soit $X = \text{Proj}(S)$.
Soient $Z, Z' \subset X$ deux sous-schémas fermés.
Soit $\varphi : Z \to Z'$ un isomorphisme.
Supposons que $Z \cap Z' = \emptyset$.
Montrer que, pour tout $z \in Z$, il existe un ouvert affine
$U \subset X$ tel que $z \in U$, $\varphi(z) \in U$ et
$\varphi(Z \cap U) = Z' \cap U$.
(Indication : utiliser l'exercice \ref{exercises-exercise-finitely-many-points-in-affine}
et un argument analogue à celui de
Schémas, lemme \ref{schemes-lemma-standard-open-two-affines}.)
\end{exercise}







\section{Morphismes de la droite projective}
\label{exercises-section-from-P1}

\noindent
Dans cette section, nous étudions les morphismes de $\mathbf{P}^1$
vers des schémas projectifs.

\begin{exercise}
\label{exercises-exercise-from-generic-point}
Soit $k$ un corps. Soit $k[t] \subset k(t)$
l'inclusion de l'anneau de polynômes dans son corps des fractions.
Soit $X$ un schéma de type fini sur $k$.
Montrer que, pour tout morphisme
$$
\varphi : \Spec(k(t)) \longrightarrow X
$$
au-dessus de $k$, il existe un élément non nul $f \in k[t]$ et un morphisme
$\psi : \Spec(k[t, 1/f]) \to X$ au-dessus de $k$
tels que $\varphi$ soit la composée
$$
\Spec(k(t)) \longrightarrow \Spec(k[t, 1/f]) \longrightarrow X
$$
\end{exercise}

\begin{exercise}
\label{exercises-exercise-from-generic-point-Pn}
Soit $k$ un corps. Soit $k[t] \subset k(t)$
l'inclusion de l'anneau de polynômes dans son corps des fractions.
Montrer que, pour tout morphisme
$$
\varphi : \Spec(k(t)) \longrightarrow \mathbf{P}^n_k
$$
au-dessus de $k$, il existe un morphisme $\psi : \Spec(k[t]) \to \mathbf{P}^n_k$
au-dessus de $k$ tel que $\varphi$ soit la composée
$$
\Spec(k(t)) \longrightarrow \Spec(k[t]) \longrightarrow \mathbf{P}^n_k
$$
Indication : l'image de $\varphi$ est contenue dans un ouvert standard $D_+(T_i)$
pour un certain $i$ ; montrer alors que l'on peut ``chasser les dénominateurs''.
\end{exercise}

\begin{exercise}
\label{exercises-exercise-from-generic-point-projective}
Soit $k$ un corps. Soit $k[t] \subset k(t)$
l'inclusion de l'anneau de polynômes dans son corps des fractions.
Soit $X$ un schéma projectif sur $k$.
Montrer que, pour tout morphisme
$$
\varphi : \Spec(k(t)) \longrightarrow X
$$
au-dessus de $k$, il existe un morphisme $\psi : \Spec(k[t]) \to X$
au-dessus de $k$ tel que $\varphi$ soit la composée
$$
\Spec(k(t)) \longrightarrow \Spec(k[t]) \longrightarrow X
$$
Indication : utiliser l'exercice \ref{exercises-exercise-from-generic-point-Pn}.
\end{exercise}

\begin{exercise}
\label{exercises-exercise-from-generic-point-P1-to-projective}
Soit $k$ un corps. Soit $X$ un schéma projectif sur $k$.
Soit $K$ le corps des fonctions rationnelles de $\mathbf{P}^1_k$ (voir
l'indication ci-dessous). Montrer que, pour tout morphisme
$$
\varphi : \Spec(K) \longrightarrow X
$$
au-dessus de $k$, il existe un morphisme $\psi : \mathbf{P}^1_k \to X$
au-dessus de $k$ tel que $\varphi$ soit la composée
$$
\Spec(k(t)) \longrightarrow \mathbf{P}^1_k \longrightarrow X
$$
Indication : utiliser l'exercice \ref{exercises-exercise-from-generic-point-projective}
pour chacun des deux ouverts du recouvrement affine
$\mathbf{P}^1_k = D_+(T_0) \cup D_+(T_1)$, et utiliser le fait que
$D_+(T_0)$ est le spectre d'un anneau de polynômes et que
$K$ est le corps des fractions de cet anneau de polynômes.
\end{exercise}







\section{Morphismes de surfaces vers des courbes}
\label{exercises-section-from-surfaces-to-curves}

\begin{exercise}
\label{exercises-exercise-points-projective-space}
Soit $R$ un anneau.
Soit $R \to k$ un homomorphisme de $R$ dans un corps.
Soit $n \geq 0$.
Montrer que
$$
\Mor_{\Spec(R)}(\Spec(k), \mathbf{P}^n_R)
=
(k^{n + 1} \setminus \{0\})/k^*
$$
où $k^*$ agit sur $k^{n + 1}$ par multiplication scalaire.
Désormais, nous noterons $(x_0 : \ldots : x_n)$ le
morphisme $\Spec(k) \to \mathbf{P}^n_k$ correspondant
à la classe d'équivalence de l'élément
$(x_0, \ldots, x_n) \in k^{n + 1} \setminus \{0\}$.
\end{exercise}

\begin{exercise}
\label{exercises-exercise-curve-projective-plane}
Soit $k$ un corps. Soit $Z \subset \mathbf{P}^2_k$ un
sous-schéma fermé irréductible et réduit.
Montrer que soit (a) $Z$ est un point fermé, soit (b) il existe
un polynôme homogène irréductible $F \in k[X_0, X_1, X_2]$ de degré $> 0$
tel que $Z = V_{+}(F)$, soit (c) $Z = \mathbf{P}^2_k$.
(Indication : se placer sur un ouvert affine standard.)
\end{exercise}

\begin{exercise}
\label{exercises-exercise-bezout}
Soit $k$ un corps. Soient $Z_1, Z_2 \subset \mathbf{P}^2_k$ des
sous-schémas fermés irréductibles de la forme $V_{+}(F)$
pour certains polynômes homogènes irréductibles $F_i \in k[X_0, X_1, X_2]$ de degré $> 0$.
Montrer que $Z_1 \cap Z_2$ est non vide.
(Indication : utiliser la théorie de la dimension pour estimer la dimension de
l'anneau local de $k[X_0, X_1, X_2]/(F_1, F_2)$ au point $0$.)
\end{exercise}

\begin{exercise}
\label{exercises-exercise-no-nonconstant-morphism-proj}
Montrer qu'il n'existe aucun morphisme non constant de schémas
$\mathbf{P}^2_{\mathbf{C}} \to \mathbf{P}^1_{\mathbf{C}}$
au-dessus de $\Spec(\mathbf{C})$. Ici, un {\it morphisme constant} est
un morphisme dont l'image est réduite à un seul point.
(Indication : si le morphisme n'est pas constant, considérer les fibres au-dessus de
$0$ et de $\infty$, puis montrer qu'elles doivent se rencontrer, ce qui donne une contradiction.)
\end{exercise}

\begin{exercise}
\label{exercises-exercise-nonconstant-morphism}
Soit $k$ un corps.
Supposons que $X \subset \mathbf{P}^3_k$ soit un sous-schéma fermé
défini par une seule équation homogène $F \in k[X_0, X_1, X_2, X_3]$.
En d'autres termes,
$$
X = \text{Proj}(k[X_0, X_1, X_2, X_3]/(F)) \subset \mathbf{P}^3_k
$$
comme cela a été expliqué dans le cours. Supposons que
$$
F = X_0 G + X_1 H
$$
pour des polynômes homogènes $G, H \in k[X_0, X_1, X_2, X_3]$
de degré positif. Montrer que, si $X_0, X_1, G, H$ n'ont aucun zéro commun,
alors il existe un morphisme non constant
$$
X \longrightarrow \mathbf{P}^1_k
$$
de schémas au-dessus de $\Spec(k)$
qui, sur les points à valeurs dans un corps (voir l'exercice \ref{exercises-exercise-points-projective-space}),
s'écrit $(x_0 : x_1 : x_2 : x_3) \mapsto (x_0 : x_1)$ lorsque
$x_0$ ou $x_1$ est non nul.
\end{exercise}





\section{Faisceaux inversibles}
\label{exercises-section-invertible-sheaves}

\begin{definition}
\label{exercises-definition-invertible-sheaf}
Soit $X$ un espace localement annelé.
Un {\it ${\mathcal O}_X$-module inversible} sur $X$
est un faisceau de ${\mathcal O}_X$-modules ${\mathcal L}$ tel que tout point
possède un voisinage ouvert $U \subset X$ tel que ${\mathcal L}|_U$
soit isomorphe à ${\mathcal O}_U$ comme ${\mathcal O}_U$-module.
On dit que ${\mathcal L}$ est trivial s'il est isomorphe à
${\mathcal O}_X$ comme ${\mathcal O}_X$-module.
\end{definition}

\begin{exercise}
\label{exercises-exercise-general-facts-invertible}
Généralités.
\begin{enumerate}
\item Montrer qu'un ${\mathcal O}_X$-module inversible sur
un schéma $X$ est quasi-cohérent.
\item Supposons que $X\to Y$ soit un morphisme d'espaces localement annelés,
et que ${\mathcal L}$ soit un ${\mathcal O}_Y$-module inversible.
Montrer que $f^\ast {\mathcal L}$ est un ${\mathcal O}_X$-module inversible.
\end{enumerate}
\end{exercise}

\begin{exercise}
\label{exercises-exercise-invertible-algebra}
Algèbre.
\begin{enumerate}
\item Montrer qu'un ${\mathcal O}_X$-module inversible sur
un schéma affine $\Spec(A)$ correspond à un $A$-module $M$ qui est
(i) de type fini, (ii) projectif, (iii) localement libre de rang 1,
et par conséquent (iv) plat et (v) de présentation finie. (On pourra
citer librement des résultats du cours du semestre dernier ou de livres d'algèbre.)
\item Supposons que $A$ soit un anneau intègre et que $M$ soit
un module comme en (a). Montrer que $M$ est isomorphe, comme $A$-module,
à un idéal $I \subset A$ tel que $IA_{\mathfrak p}$ soit principal pour
tout idéal premier ${\mathfrak p}$.
\end{enumerate}
\end{exercise}

\begin{definition}
\label{exercises-definition-invertible-module}
Soit $R$ un anneau. Un {\it module inversible $M$} est un $R$-module
$M$ tel que $\widetilde M$ soit un faisceau inversible sur le
spectre de $R$. On dit que $M$ est {\it trivial} si $M \cong R$ comme
$R$-module.
\end{definition}

\noindent
Autrement dit, $M$ est inversible si et seulement s'il
satisfait à toutes les conditions suivantes :
il est plat, de présentation finie, projectif et
localement libre de rang 1. (Il suffit bien entendu qu'il
soit localement libre de rang 1.)

\begin{exercise}
\label{exercises-exercise-simple-examples-invertible}
Exemples simples.
\begin{enumerate}
\item
\label{exercises-item-affine-line}
Soit $k$ un corps. Soit $A = k[x]$.
Montrer que $X = \Spec(A)$ ne possède que des
${\mathcal O}_X$-modules inversibles triviaux. Autrement dit, montrer que tout
$A$-module inversible est libre de rang 1.
\item
\label{exercises-item-affine-line-with-0-and-1-identified}
Soit $A$ l'anneau
$$
A = \{ f\in k[x] \mid f(0) = f(1) \}.
$$
Montrer qu'il existe un $A$-module inversible non trivial, sauf si
$k = {\mathbf F}_2$. (Indication : voir $\Spec(A)$ comme obtenu en identifiant
$0$ et $1$ dans ${\mathbf A}^1_k = \Spec(k[x])$.)
\item
\label{exercises-item-affine-line-with-cusp}
Même question qu'en (\ref{exercises-item-affine-line-with-0-and-1-identified})
pour l'anneau $A = k[x^2, x^3] \subset k[x]$
(cette fois, $k = {\mathbf F}_2$ convient également).
\end{enumerate}
\end{exercise}

\begin{exercise}
\label{exercises-exercise-higher-dimension-invertible}
Dimensions supérieures.
\begin{enumerate}
\item Prouver que tout faisceau inversible sur l'espace affine de dimension
deux est trivial. Plus précisément, soit
${\mathbf A}^2_k = \Spec(k[x, y])$ où $k$ est un corps.
Montrer que tout faisceau inversible sur ${\mathbf A}^2_k$ est trivial.
(Indication : une méthode consiste à considérer le module
$M$ correspondant, à examiner $M \otimes_{k[x, y]} k(x)[y]$, puis
à utiliser
l'exercice \ref{exercises-exercise-simple-examples-invertible} (\ref{exercises-item-affine-line})
pour lui trouver un générateur ; il reste alors encore à réfléchir.
Une autre méthode consiste à utiliser
l'exercice \ref{exercises-exercise-invertible-algebra}
et ce que l'on sait des idéaux de l'anneau
de polynômes : les idéaux premiers de hauteur un sont engendrés par un polynôme
irréductible ; il reste alors encore à réfléchir.)
\item Prouver que tout faisceau inversible sur tout sous-schéma
ouvert de l'espace affine de dimension deux est trivial. Plus précisément, soit
$U \subset {\mathbf A}^2_k$ un sous-schéma ouvert, où $k$ est un corps.
Montrer que tout faisceau inversible sur $U$ est trivial. Indication : montrer que tout
faisceau inversible sur $U$ se prolonge à ${\mathbf A}^2_k$. Ce n'est pas facile,
mais on peut le trouver dans \cite{H}.
\item Trouver un exemple de faisceau inversible non trivial
sur un cône épointé au-dessus d'un corps. Plus précisément,
soit $k$ un corps et soit $C = \Spec(k[x, y, z]/(xy-z^2))$.
Soit $U = C \setminus \{ (x, y, z) \}$. Trouver un faisceau inversible non trivial
sur $U$. Indication : il peut être plus facile de calculer le
groupe des classes d'isomorphisme de faisceaux inversibles sur $U$ que d'en
trouver simplement un. Noter que $U$ est recouvert par les ouverts
$\Spec(k[x, y, z, 1/x]/(xy-z^2))$ et
$\Spec(k[x, y, z, 1/y]/(xy-z^2))$,
qui sont « faciles » à traiter.
\end{enumerate}
\end{exercise}

\begin{definition}
\label{exercises-definition-picard-group}
Soit $X$ un espace localement annelé.
Le {\it groupe de Picard de $X$} est l'ensemble $\Pic(X)$
des classes d'isomorphisme de $\mathcal{O}_X$-modules inversibles,
la loi d'addition étant donnée par le produit tensoriel.
Voir Modules, définition \ref{modules-definition-pic}.
Pour un anneau $R$, on pose $\Pic(R) = \Pic(\Spec(R))$.
\end{definition}

\begin{exercise}
\label{exercises-exercise-traverso}
Soit $R$ un anneau.
\begin{enumerate}
\item Montrer que si $R$ est un anneau intègre normal noethérien, alors
$\Pic(R) = \Pic(R[t])$. [Indication : il existe un homomorphisme
$R[t] \to R$, $t \mapsto 0$ qui est un inverse à gauche de l'homomorphisme
$R \to R[t]$. Il suffit donc de montrer que tout
$R[t]$-module inversible $M$ tel que $M/tM \cong R$ est libre de rang $1$.
Soit $K$ le corps des fractions de $R$.
Choisir une trivialisation $K[t] \to M \otimes_{R[t]} K[t]$, ce qui est possible par
l'exercice \ref{exercises-exercise-simple-examples-invertible} (\ref{exercises-item-affine-line}).
L'ajuster pour qu'elle induise la trivialisation
ci-dessus de $M/tM$. Montrer qu'elle est en fait une trivialisation de
$M$ sur $R[t]$ (c'est ici qu'intervient la normalité).]
\item Soit $k$ un corps. Montrer que
$\Pic(k[x^2, x^3, t]) \not = \Pic(k[x^2, x^3])$.
\end{enumerate}
\end{exercise}













\section{Cohomologie de {\v C}ech}
\label{exercises-section-cech-cohomology}

\begin{exercise}
\label{exercises-exercise-cech-cohomology}
Cohomologie de {\v C}ech. Ici, $k$ est un corps.
\begin{enumerate}
\item Soit $X$ un schéma muni d'un recouvrement ouvert
${\mathcal U} : X = U_1 \cup U_2$, avec $U_1 = \Spec(k[x])$,
$U_2 =  \Spec(k[y])$,
$U_1 \cap U_2 = \Spec(k[z, 1/z])$ et les immersions ouvertes
$U_1 \cap U_2 \to U_1$ resp.\ $U_1 \cap U_2 \to U_2$ étant déterminées
par $x \mapsto z$ resp.\ $y \mapsto z$ (et c'est bien ce que je veux dire).
(Nous avons vu en cours qu'un tel $X$ existe ; c'est la droite affine
à origine dédoublée.) Calculer ${\check H}^1({\mathcal U}, {\mathcal
O})$ ;
par exemple,\ en donner une base comme espace vectoriel sur $k$.
\item Pour chaque élément de
${\check H}^1({\mathcal U}, {\mathcal O})$,
construire une suite exacte de faisceaux de ${\mathcal O}_X$-modules
$$
0 \to {\mathcal O}_X \to E \to {\mathcal O}_X \to 0
$$
telle que le cobord $\delta(1) \in {\check H}^1({\mathcal U},
{\mathcal O})$
soit l'élément donné. (Il faut notamment donner un sens à cet énoncé.
Voir aussi ci-dessous.
On peut également montrer abstraitement qu'un tel objet doit exister.)
\end{enumerate}
\end{exercise}

\begin{definition}
\label{exercises-definition-delta}
(Définition de delta.) Supposons que
$$
0 \to {\mathcal F}_1 \to {\mathcal F}_2 \to {\mathcal F}_3 \to 0
$$
soit une suite exacte courte de faisceaux abéliens sur un espace topologique $X$.
L'application cobord
$\delta : H^0(X, {\mathcal F}_3) \to {\check H}^1(X, {\mathcal F}_1)$
est définie comme suit. Soit un élément $\tau \in H^0(X, {\mathcal F}_3)$.
Choisir un recouvrement ouvert ${\mathcal U} : X = \bigcup_{i\in I} U_i$ tel
que, pour tout $i$, il existe une section $\tilde \tau_i \in {\mathcal F}_2$
relevant la restriction de $\tau$ à $U_i$. Considérer alors la correspondance
$$
(i_0, i_1) \longmapsto
\tilde \tau_{i_0}|_{U_{i_0i_1}} - \tilde \tau_{i_1}|_{U_{i_0i_1}}.
$$
C'est clairement un 1-cobord dans le complexe de {\v C}ech
${\check C}^\ast({\mathcal U}, {\mathcal F}_2)$. Mais on observe que
(en considérant ${\mathcal F}_1$ comme un sous-faisceau de ${\mathcal F}_2$) le membre de droite
est toujours une section de ${\mathcal F}_1$ sur $U_{i_0i_1}$. On voit donc
que cette correspondance définit une 1-cochaîne dans le complexe
${\check C}^\ast({\mathcal U}, {\mathcal F}_2)$. La classe de cohomologie
de cette 1-cochaîne est par définition {\it $\delta(\tau)$}.
\end{definition}



\section{Cohomologie}
\label{exercises-section-cohomology}

\begin{exercise}
\label{exercises-exercise-cohomology-not-zero}
Soit $X = \mathbf{R}$ muni de la topologie euclidienne usuelle.
En utilisant seulement les propriétés formelles de la cohomologie (fonctorialité
et suite exacte longue de cohomologie), montrer qu'il
existe un faisceau $\mathcal{F}$ sur $X$ tel que $H^1(X, \mathcal{F})$ soit non nul.
\end{exercise}

\begin{exercise}
\label{exercises-exercise-mayer-vietoris}
Soit $X = U \cup V$ un espace topologique écrit comme la
réunion de deux ouverts. On a alors une suite exacte longue
de Mayer--Vietoris
$$
0 \to
H^0(X, \mathcal{F}) \to
H^0(U, \mathcal{F}) \oplus H^0(V, \mathcal{F}) \to
H^0(U \cap V, \mathcal{F}) \to
H^1(X, \mathcal{F}) \to \ldots
$$
Quelle propriété des faisceaux injectifs est essentielle à la construction
de la suite exacte longue de Mayer--Vietoris ? Pourquoi est-elle satisfaite ?
\end{exercise}

\begin{exercise}
\label{exercises-exercise-cohomology-two-point-space}
Soit $X$ un espace topologique.
\begin{enumerate}
\item Montrer que $H^i(X, \mathcal{F})$ est nul pour $i > 0$
si $X$ possède au plus $2$ points.
\item Que se passe-t-il si $X$ possède $3$ points ?
\end{enumerate}
\end{exercise}

\begin{exercise}
\label{exercises-exercise-cohomology-spec-local-ring}
Soit $X$ le spectre d'un anneau local. Montrer que
$H^i(X, \mathcal{F})$ est nul pour $i > 0$ et pour tout
faisceau de groupes abéliens $\mathcal{F}$.
\end{exercise}

\begin{exercise}
\label{exercises-exercise-affine-morphism}
Soit $f : X \to Y$ un morphisme affine de schémas.
Prouver que $H^i(X, \mathcal{F}) = H^i(Y, f_*\mathcal{F})$
pour tout $\mathcal{O}_X$-module quasi-cohérent $\mathcal{F}$.
On pourra imposer des conditions supplémentaires à $X$ et à $Y$
et utiliser la coïncidence de la cohomologie de {\v C}ech et de la cohomologie
pour les faisceaux quasi-cohérents et les recouvrements ouverts affines des schémas séparés.
\end{exercise}

\begin{exercise}
\label{exercises-exercise-affine-morphism-to-Pn}
Soit $A$ un anneau. Soit $\mathbf{P}^n_A = \text{Proj}(A[T_0, \ldots, T_n])$
l'espace projectif sur $A$. Soit
$\mathbf{A}^{n + 1}_A = \Spec(A[T_0, \ldots, T_n])$ et soit
$$
U = \bigcup\nolimits_{i = 0, \ldots, n} D(T_i) \subset \mathbf{A}^{n + 1}_A
$$
le complémentaire de l'image de l'immersion fermée
$0 : \Spec(A) \to \mathbf{A}^{n + 1}_A$.
Construire un morphisme affine surjectif
$$
f : U \longrightarrow \mathbf{P}^n_A
$$
et prouver que $f_*\mathcal{O}_U =
\bigoplus_{d \in \mathbf{Z}} \mathcal{O}_{\mathbf{P}^n_A}(d)$.
Plus généralement, montrer que, pour un $A[T_0, \ldots, T_n]$-module gradué $M$,
on a
$$
f_*(\widetilde{M}|_U) =
\bigoplus\nolimits_{d \in \mathbf{Z}} \widetilde{M(d)}
$$
où, au membre de gauche, figure le faisceau quasi-cohérent
$\widetilde{M}$ associé à $M$ sur $\mathbf{A}^{n + 1}_A$
et, au membre de droite, les faisceaux quasi-cohérents
$\widetilde{M(d)}$ associés au module gradué $M(d)$.
\end{exercise}

\begin{exercise}
\label{exercises-exercise-compute-Pn}
Soit $A$ un anneau et soit $\mathbf{P}^n_A = \text{Proj}(A[T_0, \ldots, T_n])$
l'espace projectif sur $A$.
Calculer soigneusement la cohomologie des tordus de Serre
$\mathcal{O}_{\mathbf{P}^n_A}(d)$ du faisceau
structural sur $\mathbf{P}^n_A$. On pourra utiliser la cohomologie de {\v C}ech
et la coïncidence de la cohomologie de {\v C}ech et de la cohomologie
pour les faisceaux quasi-cohérents et les recouvrements ouverts affines des schémas séparés.
\end{exercise}

\begin{exercise}
\label{exercises-exercise-cohomology-hypersurface}
Soit $A$ un anneau et soit $\mathbf{P}^n_A = \text{Proj}(A[T_0, \ldots, T_n])$
l'espace projectif sur $A$. Soit
$F \in A[T_0, \ldots, T_n]$ homogène de degré $d$.
Soit $X \subset \mathbf{P}^n_A$ le sous-schéma fermé
correspondant à l'idéal gradué $(F)$ de $A[T_0, \ldots, T_n]$.
Que peut-on dire de $H^i(X, \mathcal{O}_X)$ ?
\end{exercise}

\begin{exercise}
\label{exercises-exercise-characterize-finite-product-fields}
Soit $R$ un anneau tel que, pour tout foncteur exact à gauche
$F : \text{Mod}_R \to \textit{Ab}$, on ait
$R^iF = 0$ pour $i > 0$. Montrer que $R$ est un produit fini
de corps.
\end{exercise}






\section{Compléments de cohomologie}
\label{exercises-section-more-cohomology}

\begin{exercise}
\label{exercises-exercise-cohomology-coordinate-axes}
Soit $k$ un corps. Soit $X \subset \mathbf{P}^n_k$
la « croix des axes de coordonnées ». Autrement dit, soit $X$
défini par les équations homogènes
$$
T_i T_j = 0\text{ pour }i > j > 0
$$
où, comme d'habitude, on écrit $\mathbf{P}^n_k = \text{Proj}(k[T_0, \ldots, T_n])$.
Autrement dit, $X$ est le sous-schéma fermé correspondant au quotient
$k[T_0, \ldots, T_n]/(T_iT_j; i > j > 0)$ de l'anneau de polynômes.
Calculer $H^i(X, \mathcal{O}_X)$ pour tout $i$. Indication : utiliser la cohomologie de {\v C}ech.
\end{exercise}

\begin{exercise}
\label{exercises-exercise-compute-cohomology-punctured}
Soit $A$ un anneau. Soit $I = (f_1, \ldots, f_t)$
un idéal de type fini de $A$.
Soit $U \subset \Spec(A)$ le complémentaire de $V(I)$.
Pour tout $A$-module $M$, écrire un complexe
de $A$-modules (en fonction de $A$, $f_1, \ldots, f_t$, $M$)
dont les groupes de cohomologie donnent $H^n(U, \widetilde{M})$.
\end{exercise}

\begin{exercise}
\label{exercises-exercise-compute-cohomology-affine-space-punctured}
Soit $k$ un corps. Soit $U \subset \mathbf{A}^d_k$ le
complémentaire du point fermé $0$ de $\mathbf{A}^d_k$.
Calculer $H^n(U, \mathcal{O}_U)$ pour tout $n$.
\end{exercise}

\begin{exercise}
\label{exercises-exercise-find-curve-genus-one-hundred}
Soit $k$ un corps. Trouver explicitement un schéma $X$ projectif sur $k$,
de dimension $1$, tel que $H^0(X, \mathcal{O}_X) = k$ et
$\dim_k H^1(X, \mathcal{O}_X) = 100$.
\end{exercise}

\begin{exercise}
\label{exercises-exercise-degree-2-cover}
Soit $f : X \to Y$ un morphisme fini localement libre de degré $2$.
Supposons que $X$ et $Y$ soient des schémas intègres et que $2$ soit inversible
dans le faisceau structural de $Y$, c'est-à-dire que $2 \in \Gamma(Y, \mathcal{O}_Y)$
soit inversible. Montrer que l'homomorphisme de $\mathcal{O}_Y$-modules
$$
f^\sharp : \mathcal{O}_Y \longrightarrow f_*\mathcal{O}_X
$$
possède un inverse à gauche, c'est-à-dire qu'il existe un homomorphisme de $\mathcal{O}_Y$-modules
$\tau : f_*\mathcal{O}_X \to \mathcal{O}_Y$
tel que $\tau \circ f^\sharp = \text{id}$.
En déduire que $H^n(Y, \mathcal{O}_Y) \to H^n(X, \mathcal{O}_X)$
est injectif\footnote{Il existe effectivement un morphisme fini localement libre
$X \to Y$ entre schémas intègres, de degré $2$, pour lequel l'homomorphisme
$H^1(Y, \mathcal{O}_Y) \to H^1(X, \mathcal{O}_X)$ n'est pas injectif.}.
\end{exercise}

\begin{exercise}
\label{exercises-exercise-pic}
Soit $X$ un schéma (ou un espace localement annelé). La règle
$U \mapsto \mathcal{O}_X(U)^*$ définit un faisceau de groupes
noté $\mathcal{O}_X^*$.
Expliquer brièvement pourquoi le groupe de Picard de $X$
(définition \ref{exercises-definition-picard-group}) est
égal à $H^1(X, \mathcal{O}_X^*)$.
\end{exercise}

\begin{exercise}
\label{exercises-exercise-pic-nontrivial}
Donner un exemple de schéma affine $X$ tel que $\Pic(X)$ soit non trivial.
En déduire, à l'aide de l'exercice \ref{exercises-exercise-pic}, que $H^1(X, -)$ n'est pas le
foncteur nul pour un tel $X$.
\end{exercise}

\begin{exercise}
\label{exercises-exercise-kill-cohomology-complement}
Soit $A$ un anneau. Soit $I = (f_1, \ldots, f_t)$ un idéal de type fini
de $A$. Soit $U \subset \Spec(A)$ le complémentaire de $V(I)$.
Étant donnés un $\mathcal{O}_{\Spec(A)}$-module quasi-cohérent $\mathcal{F}$
et $\xi \in H^p(U, \mathcal{F})$ avec $p > 0$, montrer qu'il existe
$n > 0$ tel que $f_i^n \xi = 0$ pour $i = 1, \ldots, t$.
Indication : une méthode possible consiste à utiliser le complexe
trouvé dans l'exercice \ref{exercises-exercise-compute-cohomology-punctured}.
\end{exercise}

\begin{exercise}
\label{exercises-exercise-h0-complement}
Soit $A$ un anneau. Soit $I = (f_1, \ldots, f_t)$ un
idéal de type fini de $A$. Soit $U \subset \Spec(A)$
le complémentaire de $V(I)$. Soit $M$ un $A$-module
dont la $I$-torsion est nulle, c'est-à-dire
$0 = \Ker((f_1, \ldots, f_t) : M \to M^{\oplus t})$.
Montrer qu'il existe un isomorphisme canonique
$$
H^0(U, \widetilde{M}) = \colim \Hom_A(I^n, M).
$$
Attention : ce n'est pas trivial.
\end{exercise}

\begin{exercise}
\label{exercises-exercise-Noetherian}
Soit $A$ un anneau noethérien. Soit $I$ un idéal de $A$.
Soit $M$ un $A$-module. Soit $M[I^\infty]$ l'ensemble des éléments de
torsion annulés par une puissance de $I$, défini par
$$
M[I^\infty] = \{x \in M \mid
\text{ il existe un }n \geq 1\text{ tel que }I^nx = 0\}
$$
Posons $M' = M/M[I^\infty]$. La torsion annulée par une puissance de $I$ dans $M'$ est alors nulle.
Montrer que
$$
\colim \Hom_A(I^n, M) = \colim \Hom_A(I^n, M').
$$
Attention : ce n'est pas trivial. Indications : (1) essayer de se ramener au cas où $M$ est de type fini,
(2) montrer que tout élément de $\Ext^1_A(I^n, N)$
s'envoie sur zéro dans $\Ext^1_A(I^{n + m}, N)$ pour un certain $m > 0$
si $N = M[I^\infty]$ et si $M$ est de type fini, (3) montrer la même assertion
comme en (2) pour $\Hom_A(I^n, N)$, (3) considérer la suite exacte longue
de cohomologie
$$
0 \to \Hom_A(I^n, M[I^\infty]) \to \Hom_A(I^n, M) \to \Hom_A(I^n, M')
\to \Ext^1_A(I^n, M[I^\infty])
$$
pour $M$ de type fini et la comparer à la suite relative à $I^{n + m}$ pour conclure.
\end{exercise}








\section{Retour sur la cohomologie}
\label{exercises-section-more-more-cohomology}


\begin{exercise}
\label{exercises-exercise-nonkillable}
Construire un exemple d'un corps $k$, d'une courbe $X$ sur $k$,
d'un $\mathcal{O}_X$-module inversible $\mathcal{L}$ et
d'une classe de cohomologie $\xi \in H^1(X, \mathcal{L})$ vérifiant
la propriété suivante : pour tout morphisme fini surjectif
$\pi : Y \to X$ de schémas, l'élément $\xi$
a pour image réciproque un élément non nul de $H^1(Y, \pi^*\mathcal{L})$.
Indication : construire $X$, $k$, $\mathcal{L}$ de telle sorte que
l'on ait une suite exacte courte
$0 \to \mathcal{L} \to \mathcal{O}_X \to i_*\mathcal{O}_Z \to 0$
où $Z \subset X$ soit un sous-schéma fermé constitué de
plus de $1$ point fermé. Examiner ensuite ce qui se passe par
image réciproque.
\end{exercise}

\begin{exercise}
\label{exercises-exercise-cohomology-cycle-curves}
Soit $k$ un corps algébriquement clos.
Soit $X$ un schéma projectif de dimension $1$.
Supposons que $X$ contienne un cycle de courbes, c'est-à-dire qu'il
existe un $n \geq 2$, des sous-schémas fermés intègres
de dimension $1$ deux à deux distincts
$C_1, \ldots, C_n$, ainsi que des points fermés deux à deux distincts
$x_1, \ldots, x_n \in X$, tels que $x_n \in C_n \cap C_1$
et $x_i \in C_i \cap C_{i + 1}$
pour $i = 1, \ldots, n - 1$.
Démontrer que $H^1(X, \mathcal{O}_X)$ est non nul.
Indication : soit $\mathcal{F}$ l'image du morphisme
$\mathcal{O}_X \to \bigoplus \mathcal{O}_{C_i}$,
et montrer que $H^1(X, \mathcal{F})$ est non nul
en utilisant $\kappa(x_i) = k$ et $H^0(C_i, \mathcal{O}_{C_i}) = k$.
Utiliser aussi $H^2(X, -) = 0$, d'après le théorème de Grothendieck.
\end{exercise}

\begin{exercise}
\label{exercises-exercise-surface}
Soit $X$ une surface projective sur un corps algébriquement clos $k$.
Démontrer qu'il existe un sous-schéma fermé strict $Z \subset X$ tel que
$H^1(Z, \mathcal{O}_Z)$ soit non nul. Indication : utiliser
l'exercice \ref{exercises-exercise-cohomology-cycle-curves}.
\end{exercise}

\begin{exercise}
\label{exercises-exercise-surface-better}
Soit $X$ une surface projective sur un corps algébriquement clos $k$.
Montrer que, pour tout $n \geq 0$,
il existe un sous-schéma fermé strict $Z \subset X$ tel que
$\dim_k H^1(Z, \mathcal{O}_Z) > n$.
Indiquer seulement comment procéder en modifiant les arguments
de l'exercice \ref{exercises-exercise-surface} et de
\ref{exercises-exercise-cohomology-cycle-curves} ; ne pas donner tous les détails.
\end{exercise}

\begin{exercise}
\label{exercises-exercise-surface-h2-nonzero}
Soit $X$ une surface projective sur un corps algébriquement clos $k$.
Démontrer qu'il existe un $\mathcal{O}_X$-module cohérent $\mathcal{F}$
tel que $H^2(X, \mathcal{F})$ soit non nul. Indication : utiliser le résultat
de l'exercice \ref{exercises-exercise-surface-better} ainsi qu'une
suite exacte judicieusement choisie.
\end{exercise}

\begin{exercise}
\label{exercises-exercise-kunneth}
Soient $X$ et $Y$ des schémas sur un corps $k$ (on pourra supposer
$X$ et $Y$ suffisamment raisonnables, par exemple qcqs ou projectifs sur $k$).
Posons $X \times Y = X \times_{\Spec(k)} Y$ et notons
$p : X \times Y \to X$ et $q : X \times Y \to Y$ les projections. Pour
un $\mathcal{O}_X$-module quasi-cohérent $\mathcal{F}$ et
un $\mathcal{O}_Y$-module quasi-cohérent $\mathcal{G}$, démontrer que
$$
H^n(X \times Y,
p^*\mathcal{F} \otimes_{\mathcal{O}_{X \times Y}} q^*\mathcal{G}) =
\bigoplus\nolimits_{a + b = n}
H^a(X, \mathcal{F}) \otimes_k H^b(Y, \mathcal{G})
$$
ou montrer seulement que l'égalité vaut après passage aux dimensions.
Des points supplémentaires seront accordés aux solutions « élégantes ».
\end{exercise}

\begin{exercise}
\label{exercises-exercise-Oab}
Soit $k$ un corps. Soit $X = \mathbf{P}|^1 \times \mathbf{P}^1$
le produit de la droite projective sur $k$ par elle-même,
muni des projections $p : X \to \mathbf{P}^1_k$ et $q : X \to \mathbf{P}^1_k$.
Posons
$$
\mathcal{O}(a, b) = p^*\mathcal{O}_{\mathbf{P}^1_k}(a)
\otimes_{\mathcal{O}_X} q^*\mathcal{O}_{\mathbf{P}^1_k}(b)
$$
Calculer les dimensions de $H^i(X, \mathcal{O}(a, b))$
pour tous $i, a, b$. Indication : utiliser l'exercice \ref{exercises-exercise-kunneth}.
\end{exercise}










\section{Cohomologie et polynômes de Hilbert}
\label{exercises-section-cohomology-hilbert-polynomials}

\begin{situation}
\label{exercises-situation-hilbert-polynomial}
Soit $k$ un corps. Soit $X = \mathbf{P}^n_k$
l'espace projectif de dimension $n$. Soit $\mathcal{F}$
un $\mathcal{O}_X$-module cohérent. Rappelons que
$$
\chi(X, \mathcal{F}) =
\sum\nolimits_{i = 0}^n (-1)^i \dim_k H^i(X, \mathcal{F})
$$
Rappelons que le {\it polynôme de Hilbert} de $\mathcal{F}$ est la fonction
$$
t \longmapsto \chi(X, \mathcal{F}(t))
$$
Rappelons aussi que
$\mathcal{F}(t) = \mathcal{F} \otimes_{\mathcal{O}_X} \mathcal{O}_X(t)$
où $\mathcal{O}_X(t)$ est le $t$-ième tordu du faisceau structural
comme dans Constructions, définition \ref{constructions-definition-twist}.
Dans Variétés, sous-section \ref{varieties-subsection-hilbert}, nous avons
démontré que le polynôme de Hilbert est un polynôme en $t$.
\end{situation}

\begin{exercise}
\label{exercises-exercise-hilbert-pol-easy}
Dans la situation \ref{exercises-situation-hilbert-polynomial}.
\begin{enumerate}
\item Si $P(t)$ est le polynôme de Hilbert de $\mathcal{F}$,
quel est le polynôme de Hilbert de $\mathcal{F}(-13)$ ?
\item Si $P_i$ est le polynôme de Hilbert de $\mathcal{F}_i$,
quel est le polynôme de Hilbert de $\mathcal{F}_1 \oplus \mathcal{F}_2$ ?
\item Si $P_i$ est le polynôme de Hilbert de $\mathcal{F}_i$ et si
$\mathcal{F}$ est le noyau d'un morphisme surjectif
$\mathcal{F}_1 \to \mathcal{F}_2$, quel est le polynôme de Hilbert de
$\mathcal{F}$ ?
\end{enumerate}
\end{exercise}

\begin{exercise}
\label{exercises-exercise-find-given-hilbert-pol-dim-1}
Dans la situation \ref{exercises-situation-hilbert-polynomial}, supposons $n \geq 1$.
Trouver un faisceau cohérent dont le polynôme de Hilbert est $t - 101$.
\end{exercise}

\begin{exercise}
\label{exercises-exercise-find-given-hilbert-pol-dim-2}
Dans la situation \ref{exercises-situation-hilbert-polynomial}, supposons $n \geq 2$.
Trouver un faisceau cohérent dont le polynôme de Hilbert est
$t^2/2 + t/2 - 1$. (Cet exercice est assez délicat ; il suffit
de montrer qu'un tel faisceau cohérent existe.)
\end{exercise}

\begin{exercise}
\label{exercises-exercise-bound-genus-in-degree}
Dans la situation \ref{exercises-situation-hilbert-polynomial}, supposons $n \geq 2$
et que $k$ soit algébriquement clos.
Soit $C \subset X$ un sous-schéma fermé intègre de dimension $1$.
Autrement dit, $C$ est une courbe projective.
Soit $d t + e$ le polynôme de Hilbert de $\mathcal{O}_C$
considéré comme faisceau cohérent sur $X$.
\begin{enumerate}
\item Donner une borne supérieure de $e$. (Indications : utiliser le fait que
$\mathcal{O}_C(t)$ n'a de cohomologie qu'aux degrés $0$ et $1$
et étudier $H^0(C, \mathcal{O}_C)$.)
\item Choisir une section globale $s$ de $\mathcal{O}_X(1)$
qui coupe $C$ transversalement, c'est-à-dire telle qu'il
existe des points fermés deux à deux distincts $c_1, \ldots, c_r \in C$ et
une suite exacte courte
$$
0 \to \mathcal{O}_C \xrightarrow{s} \mathcal{O}_C(1) \to
\bigoplus\nolimits_{i = 1, \ldots, r} k_{c_i} \to 0
$$
où $k_{c_i}$ est le faisceau gratte-ciel de valeur $k$ en $c_i$.
(Une telle section $s$ existe ; on l'admettra.)
Montrer que $r = d$. (Indication : tordre la suite et examiner ce que l'on obtient.)
\item En tordant la suite exacte courte, on obtient une application $k$-linéaire
$\varphi_t : \Gamma(C, \mathcal{O}_C(t)) \to \bigoplus_{i = 1, \ldots, d} k$
pour tout $t$. Montrer que cette application est surjective si $t \geq d - 1$.
\item Donner une borne inférieure de $e$ en fonction de $d$. (Indication : montrer que
$H^1(C, \mathcal{O}_C(d - 2)) = 0$ à l'aide du résultat de (3), puis utiliser
un résultat d'annulation.)
\end{enumerate}
\end{exercise}

\begin{exercise}
\label{exercises-exercise-three-quadrics-in-plane}
Dans la situation \ref{exercises-situation-hilbert-polynomial}, supposons $n = 2$.
Soient $s_1, s_2, s_3 \in \Gamma(X, \mathcal{O}_X(2))$ trois
équations quadratiques. Considérons le faisceau cohérent
$$
\mathcal{F} = \Coker\left(\mathcal{O}_X(-2)^{\oplus 3}
\xrightarrow{s_1, s_2, s_3} \mathcal{O}_X\right)
$$
Énumérer les polynômes de Hilbert possibles d'un tel $\mathcal{F}$.
(Essayer de visualiser les intersections de quadriques dans le plan projectif.)
\end{exercise}





\section{Courbes}
\label{exercises-section-curves}

\begin{exercise}
\label{exercises-exercise-closed-subset-vanshing}
Soit $k$ un corps algébriquement clos. Soit $X$ une courbe projective
sur $k$. Soit $\mathcal{L}$ un $\mathcal{O}_X$-module inversible.
Soient $s_0, \ldots, s_n \in H^0(X, \mathcal{L})$
des sections globales de $\mathcal{L}$.
Démontrer qu'il existe un sous-schéma fermé naturel
$$
Z \subset
\mathbf{P}^n \times X
$$
tel que le point fermé
$((\lambda_0 : \ldots : \lambda_n), x)$
appartienne à $Z$ si et seulement si la section
$\lambda_0 s_0 + \ldots + \lambda_n s_n$
s'annule en $x$. (Indication : décrire $Z$ localement sur des ouverts affines.)
\end{exercise}

\begin{exercise}
\label{exercises-exercise-diagonal-cartier}
Soit $k$ un corps algébriquement clos. Soit $X$ une courbe lisse sur $k$.
Soit $r \geq 1$. Montrer que le sous-ensemble fermé
$$
D \subset X \times X^r = X^{r + 1}
$$
dont les points fermés sont les uplets $(x, x_1, \ldots, x_r)$ tels que
$x = x_i$ pour un certain $i$, possède un faisceau d'idéaux inversible. Autrement dit,
montrer que
$D$ est un diviseur de Cartier effectif. Indications : se ramener au cas $r = 1$ et
utiliser le fait que $X$ est une courbe lisse pour obtenir une propriété de la diagonale
(consulter les ouvrages sur ce point).
\end{exercise}

\begin{exercise}
\label{exercises-exercise-one-divisor-in-another}
Soit $k$ un corps algébriquement clos. Soit $X$ une courbe projective lisse
sur $k$. Soit $T$ un schéma de type fini sur $k$ et soient
$$
D_1 \subset X \times T
\quad\text{et}\quad
D_2 \subset X \times T
$$
deux diviseurs de Cartier effectifs tels que, pour $t \in T$, les
fibres $D_{i, t} \subset X_t$ ne soient pas denses (c'est-à-dire ne contiennent pas
le point générique de la courbe $X_t$).
Démontrer qu'il existe un sous-schéma fermé canonique
$Z \subset T$ tel qu'un point fermé $t \in T$ appartienne à $Z$
si et seulement si, pour les fibres schématiques $D_{1, t}$, $D_{2, t}$
de $D_1$, $D_2$, on a
$$
D_{1, t} \subset D_{2, t}
$$
comme sous-schémas fermés de $X_t$.
Indications : montrer que, quitte à restreindre $T$, on peut
supposer $T = \Spec(A)$ affine et qu'il existe un ouvert affine $U \subset X$
tel que $D_i \subset U \times T$. Puis montrer que
$M_1 = \Gamma(D_1, \mathcal{O}_{D_1})$ est un $A$-module localement libre de rang fini
(il faudra ici un peu d'algèbre non triviale --- demander à vos amis).
Après avoir restreint $T$, on peut supposer que $M_1$ est un $A$-module libre.
Considérer alors
$$
\Gamma(U \times T, \mathcal{I}_{D_2}) \to M_1 = A^{\oplus N}
$$
et définir $Z$ comme le sous-schéma fermé découpé par l'idéal engendré
par les coordonnées des vecteurs de l'image de ce morphisme. Expliquer pourquoi cela fonctionne
(cela demandera peut-être encore un peu d'algèbre commutative).
\end{exercise}

\begin{exercise}
\label{exercises-exercise-closed-subset-vanshing-r}
Soit $k$ un corps algébriquement clos. Soit $X$ une courbe projective lisse
sur $k$. Soit $\mathcal{L}$ un $\mathcal{O}_X$-module inversible.
Soient $s_0, \ldots, s_n \in H^0(X, \mathcal{L})$
des sections globales de $\mathcal{L}$. Soit $r \geq 1$.
Démontrer qu'il existe un sous-schéma fermé naturel
$$
Z \subset
\mathbf{P}^n \times X \times \ldots \times X = \mathbf{P}^n \times X^r
$$
tel que le point fermé
$((\lambda_0 : \ldots : \lambda_n), x_1, \ldots, x_r)$
appartienne à $Z$ si et seulement si la section
$s_\lambda = \lambda_0 s_0 + \ldots + \lambda_n s_n$
s'annule sur le diviseur $D = x_1 + \ldots + x_r$, c'est-à-dire que
la section $s_\lambda$ appartient à $\mathcal{L}(-D)$.
Indication : expliquer comment cela résulte de la combinaison des résultats des
exercices \ref{exercises-exercise-diagonal-cartier} et
\ref{exercises-exercise-one-divisor-in-another}.
\end{exercise}

\begin{exercise}
\label{exercises-exercise-effective}
Soit $k$ un corps algébriquement clos. Soit $X$ une courbe projective lisse
sur $k$. Soit $\mathcal{L}$ un $\mathcal{O}_X$-module inversible.
Montrer qu'il existe un sous-ensemble fermé naturel
$$
Z \subset X^r
$$
tel qu'un point fermé $(x_1, \ldots, x_r)$ de $X^r$ appartienne à $Z$
si et seulement si $\mathcal{L}(-x_1 - \ldots -x_r)$ admet une section
globale non nulle. Indication : utiliser l'exercice \ref{exercises-exercise-closed-subset-vanshing-r}.
\end{exercise}

\begin{exercise}
\label{exercises-exercise-difference-effective}
Soit $k$ un corps algébriquement clos. Soit $X$ une courbe projective lisse
sur $k$. Soient $r \geq s$ des entiers. Montrer qu'il existe un sous-ensemble fermé
naturel
$$
Z \subset X^r \times X^s
$$
tel qu'un point fermé $(x_1, \ldots, x_r, y_1, \ldots, y_s)$ de
$X^r \times X^s$ appartienne à $Z$
si et seulement si $x_1 + \ldots + x_r - y_1 - \ldots - y_s$
est linéairement équivalent à un diviseur effectif.
Indication : choisir un module inversible auxiliaire $\mathcal{L}$
de degré suffisamment grand pour que $\mathcal{L}(-D)$ admette une section
non nulle pour tout diviseur effectif $D$ de degré $r$.
Utiliser ensuite deux fois le résultat de l'exercice \ref{exercises-exercise-effective}.
\end{exercise}

\begin{exercise}
\label{exercises-exercise-genus-7}
Choisir un corps algébriquement clos $k$ de son choix.
Déterminer aussi complètement que possible tous les $\mathfrak g^r_d$
susceptibles d'exister sur une courbe de genre $7$.
Au cours de cette étude, essayer aussi de
\begin{enumerate}
\item déterminer dans quels cas le $\mathfrak g^r_d$ est sans point de base, et
\item déterminer dans quels cas le $\mathfrak g^r_d$ définit une immersion fermée
dans $\mathbf{P}^r$.
\end{enumerate}
Faire de même si l'on suppose la courbe « générale » (adopter
la notion de généralité de son choix -- cela peut être plus facile que la question précédente).
Faire de même si l'on suppose la courbe hyperelliptique.
Faire de même si l'on suppose la courbe trigonale
(et non hyperelliptique). Etc.
\end{exercise}





\section{Modules}
\label{exercises-section-moduli}

\noindent
Dans cette section, nous examinons quelques approches naïves des problèmes de modules
d'objets de géométrie algébrique.

\medskip\noindent
Soit $k$ un corps algébriquement clos. Supposons que $M$ soit un
problème de modules sur $k$. Nous ne définirons pas ici avec précision
ce que cela signifie, mais le sens devrait être clair dans chaque exercice.
Pour comprendre ce qui suit, il suffit de connaître
les objets de $M$ sur $k$, les isomorphismes entre
ces objets de $M$ sur $k$, et les familles d'objets de $M$
sur une variété. On dit alors que le {\it nombre de modules de $M$}
est $d \geq 0$ si les conditions suivantes sont vérifiées :
\begin{enumerate}
\item il existe un nombre fini de familles $X_i \to V_i$, $i = 1, \ldots, n$
telles que tout objet de $M$ sur $k$ soit isomorphe à une fibre de l'une
d'entre elles et que $\max \dim(V_i) = d$, et
\item il est impossible de le faire avec une valeur plus petite de $d$.
\end{enumerate}
Il ne s'agit en réalité que d'une approximation très grossière de notions meilleures
que l'on trouve dans la littérature.

\begin{exercise}
\label{exercises-exercise-number-moduli-lines-conics}
Soit $k$ un corps algébriquement clos. Soient $d \geq 1$ et $n \geq 1$.
Convenons que le problème de modules des hypersurfaces de degré $d$ dans $P^n$ est donné par
\begin{enumerate}
\item un objet est une hypersurface $X \subset \mathbf{P}^n_k$
de degré $d$,
\item un isomorphisme entre deux objets $X \subset \mathbf{P}^n_k$
et $Y \subset \mathbf{P}^n_k$ est un élément $g \in \text{PGL}_n(k)$
tel que $g(X) = Y$, et
\item une famille d'hypersurfaces sur une variété $V$ est un sous-schéma
fermé $X \subset \mathbf{P}^n_V$ tel que, pour tout $v \in V$,
la fibre schématique $X_v$ de $X \to V$ soit
une hypersurface dans $\mathbf{P}^n_v$.
\end{enumerate}
Calculer (si possible -- les questions deviennent de plus en plus difficiles)
\begin{enumerate}
\item le nombre de modules lorsque $n = 1$ et $d$ est arbitraire,
\item le nombre de modules lorsque $n = 2$ et $d = 1$,
\item le nombre de modules lorsque $n = 2$ et $d = 2$,
\item le nombre de modules lorsque $n \geq 1$ et $d = 2$,
\item le nombre de modules lorsque $n = 2$ et $d = 3$,
\item le nombre de modules lorsque $n = 3$ et $d = 3$, et
\item le nombre de modules lorsque $n = 2$ et $d = 4$.
\end{enumerate}
\end{exercise}

\begin{exercise}
\label{exercises-exercise-number-moduli-hyperelliptic}
Soit $k$ un corps algébriquement clos. Soit $g \geq 2$.
Convenons que le problème de modules des courbes hyperelliptiques de genre $g$ est donné par
\begin{enumerate}
\item un objet est une courbe hyperelliptique projective lisse $X$
de genre $g$,
\item un isomorphisme entre deux objets $X$ et $Y$ est un
isomorphisme $X \to Y$ de schémas sur $k$, et
\item une famille de courbes hyperelliptiques de genre $g$ sur une variété $V$ est un
morphisme propre et plat\footnote{On peut supprimer cette hypothèse sans changer
la réponse à la question.} $X \to Y$ tel que toutes les fibres
schématiques de $X \to V$ soient des courbes hyperelliptiques projectives lisses
de genre $g$.
\end{enumerate}
Montrer que le nombre de modules est $2g - 1$.
\end{exercise}





\section{Ext globaux}
\label{exercises-section-global-exts}

\begin{exercise}
\label{exercises-exercise-ext-line-plane-p3}
Soit $k$ un corps. Soit $X = \mathbf{P}^3_k$. Soient $L \subset X$
et $P \subset X$ respectivement une droite et un plan, considérés comme des sous-schémas fermés
définis respectivement par $1$ et par $2$ équations linéaires. Calculer
les dimensions de
$$
\Ext^i_X(\mathcal{O}_L, \mathcal{O}_P)
$$
pour tout $i$.
Traiter à la fois le cas où $L$ est contenu dans $P$ et
celui où $L$ n'est pas contenu dans $P$.
\end{exercise}

\begin{exercise}
\label{exercises-exercise-ext-point}
Soit $k$ un corps. Soit $X = \mathbf{P}^n_k$. Soit $Z \subset X$
un point fermé rationnel sur $k$, considéré comme un sous-schéma fermé.
Par exemple, le point de coordonnées homogènes
$(1 : 0 : \ldots : 0)$. Calculer les dimensions de
$$
\Ext^i_X(\mathcal{O}_Z, \mathcal{O}_Z)
$$
pour tout $i$.
\end{exercise}

\begin{exercise}
\label{exercises-exercise-ext-pairings}
Soit $X$ un espace annelé. Définir les applications de cup-produit
$$
\Ext^i_X(\mathcal{G}, \mathcal{H})
\times
\Ext^j_X(\mathcal{F}, \mathcal{G})
\longrightarrow
\Ext^{i + j}_X(\mathcal{F}, \mathcal{H})
$$
pour des $\mathcal{O}_X$-modules $\mathcal{F}, \mathcal{G}, \mathcal{H}$.
(Indication : il s'agit d'une construction extrêmement générale.)
\end{exercise}

\begin{exercise}
\label{exercises-exercise-move-out-locally-free}
Soit $X$ un espace annelé. Soit $\mathcal{E}$ un
$\mathcal{O}_X$-module localement libre de rang fini, de dual
$\mathcal{E}^\vee = \SheafHom_{\mathcal{O}_X}(\mathcal{E}, \mathcal{O}_X)$.
Démontrer les assertions suivantes
\begin{enumerate}
\item $\SheafExt^i_{\mathcal{O}_X}(
\mathcal{F} \otimes_{\mathcal{O}_X} \mathcal{E}, \mathcal{G}) =
\SheafExt^i_{\mathcal{O}_X}(
\mathcal{F},
\mathcal{E}^\vee \otimes_{\mathcal{O}_X} \mathcal{G}) =
\SheafExt^i_{\mathcal{O}_X}(
\mathcal{F}, \mathcal{G}) \otimes_{\mathcal{O}_X} \mathcal{E}^\vee$, et
\item $\Ext^i_X(
\mathcal{F} \otimes_{\mathcal{O}_X} \mathcal{E}, \mathcal{G}) =
\Ext^i_X(\mathcal{F},
\mathcal{E}^\vee \otimes_{\mathcal{O}_X} \mathcal{G})$.
\end{enumerate}
Ici, $\mathcal{F}$ et $\mathcal{G}$ sont des $\mathcal{O}_X$-modules. En déduire
que
$$
\Ext^i_X(\mathcal{E}, \mathcal{G}) =
H^i(X, \mathcal{E}^\vee \otimes_{\mathcal{O}_X} \mathcal{G})
$$
\end{exercise}

\begin{exercise}
\label{exercises-exercise-trace-on-ext}
Soit $X$ un espace annelé. Soit $\mathcal{E}$ un
$\mathcal{O}_X$-module localement libre de rang fini. Construire un homomorphisme trace
$$
\Ext^i_X(\mathcal{E}, \mathcal{E}) \to H^i(X, \mathcal{O}_X)
$$
pour tout $i$. Le généraliser en un homomorphisme trace
$$
\Ext^i_X(\mathcal{E}, \mathcal{E} \otimes_{\mathcal{O}_X} \mathcal{F})
\to H^i(X, \mathcal{F})
$$
pour tout $\mathcal{O}_X$-module $\mathcal{F}$.
\end{exercise}

\begin{exercise}
\label{exercises-exercise-duality}
Soit $k$ un corps. Soit $X = \mathbf{P}^d_k$. Posons
$\omega_{X/k} = \mathcal{O}_X(-d - 1)$.
Démontrer que, pour des modules localement libres de rang fini
$\mathcal{E}$, $\mathcal{F}$, le cup-produit sur les Ext
composé avec l'homomorphisme trace sur les Ext
$$
\Ext^i_X(\mathcal{E}, \mathcal{F} \otimes_{\mathcal{O}_X} \omega_{X/k})
\times
\Ext^{d - i}_X(\mathcal{F}, \mathcal{E})
\to
\Ext_X^d(\mathcal{F}, \mathcal{F} \otimes_{\mathcal{O}_X} \omega_{X/k})
\to
H^d(X, \omega_{X/k}) = k
$$
définit un accouplement non dégénéré. Indication : on peut soit redémontrer la dualité
dans ce cadre, soit se ramener à la cohomologie des faisceaux et appliquer le
théorème de dualité de Serre démontré dans le cours.
\end{exercise}







\section{Diviseurs}
\label{exercises-section-divisors}

\noindent
Nous rassemblons ici, pour plus de commodité, toutes les définitions utiles.

\begin{definition}
\label{exercises-definition-divisor}
Dans toute cette définition, soit $S$ un schéma quelconque et soit
$X$ un schéma noethérien intègre.
\begin{enumerate}
\item Un {\it diviseur de Weil} sur $X$ est une combinaison linéaire formelle
$\Sigma n_i[Z_i]$ de diviseurs premiers $Z_i$ à coefficients entiers.
\item Un {\it diviseur premier} est un sous-schéma fermé $Z \subset X$,
intègre et de point générique $\xi \in Z$, tel que
${\mathcal O}_{X, \xi}$ soit de dimension $1$. Nous utiliserons la notation
${\mathcal O}_{X, Z} = {\mathcal O}_{X, \xi}$
lorsque $\xi \in Z \subset X$ est comme ci-dessus. Noter que ${\mathcal O}_{X, Z} \subset
K(X)$ est un sous-anneau du corps des fonctions de $X$.
\item Le {\it diviseur de Weil associé à une fonction rationnelle
$f \in K(X)^\ast$} est la somme $\Sigma v_Z(f)[Z]$. Ici, $v_Z(f)$ est
défini comme suit
\begin{enumerate}
\item Si $f \in {\mathcal O}_{X, Z}^\ast$, alors $v_Z(f) = 0$.
\item Si $f \in {\mathcal O}_{X, Z}$, alors
$$
v_Z(f) = \text{length}_{{\mathcal O}_{X, Z}}({\mathcal O}_{X, Z}/(f)).
$$
\item Si $f = \frac{a}{b}$ avec $a, b \in {\mathcal O}_{X, Z}$,
alors
$$
v_Z(f) = \text{length}_{{\mathcal O}_{X, Z}}({\mathcal O}_{X, Z}/(a)) -
\text{length}_{{\mathcal O}_{X, Z}}({\mathcal O}_{X, Z}/(b)).
$$
\end{enumerate}
\item Un {\it diviseur de Cartier effectif} sur un schéma $S$
est un sous-schéma fermé $D \subset S$ tel que tout point $d\in D$
possède un voisinage ouvert affine $\Spec(A) = U \subset S$ dans $S$
tel que $D \cap U = \Spec(A/(f))$, où $f \in A$ est un élément régulier.
\item Le {\it diviseur de Weil $[D]$ associé à un diviseur
de Cartier effectif $D \subset X$} de notre schéma noethérien intègre
$X$ est défini comme la somme $\Sigma v_Z(D)[Z]$, où
$v_Z(D)$ est défini comme suit
\begin{enumerate}
\item Si le point générique $\xi$ de $Z$ n'appartient pas à $D$,
alors $v_Z(D) = 0$.
\item Si le point générique $\xi$ de $Z$ appartient à $D$,
alors
$$
v_Z(D) = \text{length}_{{\mathcal O}_{X, Z}}({\mathcal O}_{X, Z}/(f))
$$
où $f \in {\mathcal O}_{X, Z} = {\mathcal O}_{X, \xi}$ est l'élément régulier
qui définit $D$ dans un voisinage affine de $\xi$ (comme en (4) ci-dessus).
\end{enumerate}
\item Soit $S$ un schéma. Le {\it faisceau des anneaux totaux de
fractions ${\mathcal K}_S$} est le faisceau de ${\mathcal O}_S$-algèbres obtenu par
le faisceau associé au préfaisceau ${\mathcal K}'$ défini comme suit.
Pour $U \subset S$ ouvert, posons ${\mathcal K}'(U) = S_U^{-1}{\mathcal O}_S(U)$,
où $S_U \subset {\mathcal O}_S(U)$ est la partie multiplicative
formée des sections $f \in {\mathcal O}_S(U)$ telles que le germe
de $f$ dans ${\mathcal O}_{S, u}$ soit régulier pour tout $u\in U$.
En particulier, tous les éléments de $S_U$ sont réguliers.
Ainsi, ${\mathcal O}_S$ est un sous-faisceau de ${\mathcal K}_S$, et l'on obtient une
suite exacte
$$
0 \to {\mathcal O}_S^\ast \to {\mathcal K}_S^\ast \to
{\mathcal K}_S^\ast/{\mathcal O}_S^\ast \to 0.
$$
\item Un {\it diviseur de Cartier} sur un schéma $S$ est une section globale
du faisceau quotient ${\mathcal K}_S^\ast/{\mathcal O}_S^\ast$.
\item Le {\it diviseur de Weil associé à un diviseur de Cartier}
$\tau \in \Gamma(X, {\mathcal K}_X^\ast/{\mathcal O}_X^\ast)$ sur notre schéma
noethérien intègre
$X$ est la somme $\Sigma v_Z(\tau)[Z]$, où $v_Z(\tau)$ est défini
par la procédure suivante
\begin{enumerate}
\item Si le germe de $\tau$ au point générique $\xi$
de $Z$ est nul -- autrement dit, si l'image de $\tau$ dans la fibre
$({\mathcal K}^\ast/{\mathcal O}^\ast)_\xi$ est ``nulle'' --, alors $v_Z(\tau) = 0$.
\item Choisir un voisinage ouvert affine $\Spec(A) = U \subset X$
tel que $\tau|_U$ soit l'image d'une section $f \in {\mathcal K}(U)$
et que, de plus, $f = a/b$ avec $a, b \in A$. On pose alors
$$
v_Z(f) = \text{length}_{{\mathcal O}_{X, Z}}({\mathcal O}_{X, Z}/(a)) -
\text{length}_{{\mathcal O}_{X, Z}}({\mathcal O}_{X, Z}/(b)).
$$
\end{enumerate}
\end{enumerate}
\end{definition}

\begin{remarks}
\label{exercises-remarks-divisors}
Voici quelques remarques immédiates.
\begin{enumerate}
\item Sur un schéma noethérien intègre $X$, le
faisceau ${\mathcal K}_X$ est constant de valeur le corps des fonctions $K(X)$.
\item Pour que les définitions ci-dessus aient un sens, il faut
montrer que
$$
\text{length}_{\mathcal O}({\mathcal O}/(ab)) =
\text{length}_{\mathcal O}({\mathcal O}/(a)) +
\text{length}_{\mathcal O}({\mathcal O}/(b))
$$
pour toute paire $(a, b)$ d'éléments non nuls d'un anneau local noethérien intègre
de dimension 1, noté ${\mathcal O}$. Cela sera démontré dans le cours.
\end{enumerate}
\end{remarks}

\begin{exercise}
\label{exercises-exercise-effective-cartier-cartier}
(Sur un schéma quelconque.)
Décrire comment associer un diviseur de Cartier à un diviseur de Cartier effectif.
\end{exercise}

\begin{exercise}
\label{exercises-exercise-rational-function-cartier}
(Sur un schéma intègre.)
Décrire comment associer un diviseur de Cartier $D$ à une fonction rationnelle
$f$ de telle sorte que les diviseurs de Weil associés à $D$ et à $f$ coïncident.
(C'est tautologique.)
\end{exercise}

\begin{exercise}
\label{exercises-exercise-weil-not-cartier}
Donner un exemple de diviseur de Weil sur une variété qui ne soit
associé à aucun diviseur de Cartier.
\end{exercise}

\begin{exercise}
\label{exercises-exercise-weil-Q-cartier}
Donner un exemple de diviseur de Weil $D$ sur une variété qui ne soit
associé à aucun diviseur de Cartier, mais
tel que $nD$ soit associé à un diviseur de Cartier
pour un certain $n > 1$.
\end{exercise}

\begin{exercise}
\label{exercises-exercise-weil-not-Q-cartier}
Donner un exemple de diviseur de Weil $D$ sur une variété qui ne soit
associé à aucun diviseur de Cartier et
tel que $nD$ NE SOIT associé à aucun diviseur de Cartier
pour tout $n > 1$.
(Indication : considérer un cône, par exemple $X : xy - zw = 0$ dans
$\mathbf{A}^4_k$. Essayer de montrer que $D = [x = 0, z = 0]$ convient.)
\end{exercise}

\begin{exercise}
\label{exercises-exercise-cartier-not-difference-effective-cartier}
Sur un schéma séparé $X$ de type fini sur un corps :
donner un exemple de diviseur de Cartier qui ne soit pas la différence de
deux diviseurs de Cartier effectifs.
Indication : trouver un $X$ qui ne possède aucun diviseur de Cartier effectif non vide,
par exemple le schéma construit dans
\cite[III, exercice 5.9]{H}. Il existe même un exemple
où $X$ est une variété : c'est la variété de
l'exercice \ref{exercises-exercise-nonprojective}.
\end{exercise}

\begin{exercise}
\label{exercises-exercise-nonprojective}
Exemple de variété propre non projective.
Soit $k$ un corps. Soit $L \subset \mathbf{P}^3_k$ une droite et
soit $C \subset \mathbf{P}^3_k$ une conique non singulière. Supposons que
$C \cap L = \emptyset$. Choisissons un
isomorphisme $\varphi : L \to C$. Soit $X$ la $k$-variété obtenue
en recollant $C$ à $L$ au moyen de $\varphi$. Autrement dit, il existe
un morphisme propre birationnel surjectif
$$
\pi : \mathbf{P}^3_k \longrightarrow X
$$
et un ouvert $U \subset X$ tel que $\pi : \pi^{-1}(U) \to U$ soit
un isomorphisme, $\pi^{-1}(U) = \mathbf{P}^3_k \setminus (L \cup C)$,
et tel que $\pi|_L = \pi|_C \circ \varphi$. (Ces conditions ne
définissent pas encore $X$ de façon unique. Pour cela, il faut préciser le
faisceau structural de $X$ au voisinage des points de $Z = X \setminus U$.)
Montrer que $X$ existe et que c'est une variété propre, mais non projective.
(Indication : pour l'existence, utiliser le résultat de
l'exercice \ref{exercises-exercise-prepare-glueing}. Pour la non-projectivité, utiliser
$\Pic(\mathbf{P}^3_k) = \mathbf{Z}$ pour montrer que $X$ ne peut pas posséder
de faisceau inversible ample.)
\end{exercise}



\section{Différentielles}
\label{exercises-section-differentials}

\noindent
{\bf Définitions et résultats.} Différentielles de K\"ahler.
\begin{enumerate}
\item Soit $R \to A$ un homomorphisme d'anneaux. Le {\it module des différentielles de K\"ahler
de $A$ sur $R$} se note $\Omega_{A/R}$.
Il est engendré par les éléments $\text{d}a$, $a \in A$,
soumis aux relations :
$$
\text{d}(a_1 + a_2) = \text{d}a_1 + \text{d}a_2,\quad
\text{d}(a_1a_2) = a_1\text{d}a_2 + a_2\text{d}a_1,\quad
\text{d}r = 0
$$
La $R$-dérivation universelle canonique $\text{d} : A \to \Omega_{A/R}$
envoie $a\mapsto \text{d}a$.
\item Considérons la suite exacte
$$
0 \to I \to A \otimes_R A \to A \to 0
$$
qui définit l'idéal $I$. Il existe une dérivation canonique
$\text{d} : A \to I/I^2$ qui envoie $a$ sur la classe de
$a \otimes 1 - 1 \otimes a$. C'est une autre présentation du
module des différentielles de $A$ sur $R$; autrement dit,
$$
(I/I^2, \text{d}) \cong (\Omega_{A/R}, \text{d}).
$$
\item Pour des parties multiplicatives $S_R \subset R$ et
$S_A \subset A$ telles que l'image de $S_R$ soit contenue dans $S_A$, on a
$$
\Omega_{S_A^{-1}A / S_R^{-1}R} =
S_A^{-1}\Omega_{A/R}.
$$
\item Si $A$ est une $R$-algèbre de présentation finie, alors
$\Omega_{A/R}$ est un $A$-module de présentation finie. Ainsi, dans
ce cas, les idéaux de {\it Fitting} de $\Omega_{A/R}$ sont définis.
\item Soit $f : X \to S$ un morphisme de schémas. Il existe
un faisceau quasi-cohérent de ${\mathcal O}_X$-modules $\Omega_{X/S}$
et une dérivation ${\mathcal O}_S$-linéaire
$$
\text{d} : {\mathcal O}_X \longrightarrow \Omega_{X/S}
$$
tels que, pour tous ouverts affines $\Spec(A) = U \subset X$,
$\Spec(R) = V \subset S$,
avec $f(U) \subset V$, on ait
$$
\Gamma(\Spec(A), \Omega_{X/S}) = \Omega_{A/R}
$$
d'une manière compatible avec $\text{d}$.
\end{enumerate}

\begin{exercise}
\label{exercises-exercise-dual-numbers}
Soit $k[\epsilon]$ l'anneau des nombres duaux
sur le corps $k$, c'est-à-dire que $\epsilon^2 = 0$.
\begin{enumerate}
\item Considérons l'homomorphisme d'anneaux
$$
R = k[\epsilon] \to A = k[x, \epsilon]/(\epsilon x)
$$
Montrer que les idéaux de Fitting de $\Omega_{A/R}$ sont (en commençant par
l'idéal de Fitting d'indice zéro)
$$
(\epsilon), A, A, \ldots
$$
\item Considérons l'homomorphisme $R = k[t] \to
A = k[x, y, t]/(x(y-t)(y-1), x(x-t))$. Montrer que les idéaux de Fitting de
$\Omega_{A/R}$ dans $A$ sont (supposer, pour simplifier, que $k$ est de caractéristique
nulle)
$$
x(2x-t)(2y-t-1)A, \ (x, y, t)\cap (x, y-1, t), \ A, \ A, \ldots
$$
Ainsi, l'idéal de Fitting d'indice $0$ est défini par un unique élément de $A$,
l'idéal de Fitting d'indice $1$ définit deux points fermés de $\Spec(A)$, et
tous les autres sont triviaux.
\item Considérons l'homomorphisme $R = k[t] \to A = k[x, y, t]/(xy-t^n)$.
Calculer les idéaux de Fitting de $\Omega_{A/R}$.
\end{enumerate}
\end{exercise}

\begin{remark}
\label{exercises-remark-fitting-omega-not-sings}
On utilise couramment le $k$-ième idéal de Fitting de $\Omega_{X/S}$
pour définir le schéma des singularités du morphisme $X \to S$ lorsque $X$ est de dimension relative
$k$ sur $S$. Toutefois, comme le montre la partie (a), cette définition demande de la prudence
lorsque la dimension des fibres de la famille n'est pas ``constante'', par exemple lorsque
la famille n'est pas plate. Comme le montre la partie (b), la platitude ne suffit pas non plus
(et oui, cette famille est plate). Dans les ``bons cas'' -- comme en (c) --, pour les
familles de courbes, on s'attend à ce que l'idéal de Fitting d'indice $0$ soit nul et
que celui d'indice $1$ définisse, schématiquement, le lieu singulier.
\end{remark}

\begin{exercise}
\label{exercises-exercise-formally-smooth}
Supposons que $R$ soit un anneau et que
$$
A = R[x_1, \ldots, x_n]/(f_1, \ldots, f_n).
$$
Noter que nous supposons $A$ présenté par autant
d'équations que de variables. La matrice des dérivées
partielles
$$
( \partial f_i / \partial x_j )
$$
est $n \times n$, donc carrée. Supposons que
son déterminant soit inversible comme élément de $A$. Noter que
c'est exactement la condition qui exprime que $\Omega_{A/R} = (0)$
dans ce cas de $n$ générateurs et de $n$ relations.
Soit $\pi : B' \to B$ une surjection de $R$-algèbres
dont le noyau $J$ est de carré nul (comme idéal de $B'$).
Soit $\varphi : A \to B$ un homomorphisme de $R$-algèbres.
Montrer qu'il existe un unique homomorphisme de $R$-algèbres
$\varphi' : A \to B'$ tel que $\varphi = \pi \circ \varphi'$.
\end{exercise}

\begin{exercise}
\label{exercises-exercise-formally-smooth-one-equation}
Généraliser le résultat de l'exercice \ref{exercises-exercise-formally-smooth}
au cas où $A = R[x, y]/(f)$.
\end{exercise}

\begin{exercise}
\label{exercises-exercise-Jacobian-criterion}
Soit $k$ un corps, soient $f_1, \ldots, f_c \in k[x_1, \ldots, x_n]$, et
soit $A = k[x_1, \ldots, x_n]/(f_1, \ldots, f_c)$.
Supposons que $f_j(0, \ldots, 0) = 0$. Cela signifie que
$\mathfrak m = (x_1, \ldots, x_n)A$ est un idéal maximal.
Démontrer que l'anneau local $A_\mathfrak m$ est régulier
si la matrice
$$
(\partial f_j/ \partial x_i)|_{(x_1, \ldots, x_n) = (0, \ldots, 0)}
$$
est de rang $c$. Quelle est alors la dimension de $A_\mathfrak m$ ?
Montrer que la réciproque est fausse en donnant un exemple où
$A_\mathfrak m$ est régulier mais où le rang est strictement inférieur à $c$;
quelle est la dimension de $A_\mathfrak m$ dans cet exemple ?
\end{exercise}




\section{Schémas, examen final, automne 2007}
\label{exercises-section-final-exam-fall-2007}

\noindent
Voici les questions de l'examen final d'un cours sur les schémas,
donné au printemps 2007 à l'Université Columbia.

\begin{exercise}[Définitions]
\label{exercises-exercise-definitions}
Donner la définition des notions suivantes.
\begin{enumerate}
\item $X$ est un {\it schéma}
\item le morphisme de schémas $f : X \to Y$ est {\it fini}
\item le morphisme de schémas $f : X \to Y$ est {\it de type fini}
\item le schéma $X$ est {\it noethérien}
\item le ${\mathcal O}_X$-module ${\mathcal L}$ sur
le schéma $X$ est {\it inversible}
\item le {\it genre} d'une courbe non singulière
projective sur un corps algébriquement clos
\end{enumerate}
\end{exercise}

\begin{exercise}
\label{exercises-exercise-kill-global-sections}
Soit $X = \Spec({\mathbf Z}[x, y])$, et soit ${\mathcal F}$ un
${\mathcal O}_X$-module quasi-cohérent.
Supposons que ${\mathcal F}$ soit nul après restriction à
cet
ouvert affine standard $D(x)$.
\begin{enumerate}
\item Montrer que toute section globale $s$ de ${\mathcal F}$ est annulée par une
puissance de $x$, c'est-à-dire $x^ns = 0$ pour un certain $n\in {\mathbf N}$.
\item Pensez-vous que le même énoncé reste vrai si l'on ne suppose pas que ${\mathcal F}$
est quasi-cohérent ?
\end{enumerate}
\end{exercise}

\begin{exercise}
\label{exercises-exercise-empty-fibre-empty}
Supposons que $X \to \Spec(R)$ soit un morphisme propre et que
$R$ soit un anneau de valuation discrète de corps résiduel $k$. Supposons que
$X \times_{\Spec(R)} \Spec(k)$ soit le schéma vide. Montrer que
$X$ est le schéma vide.
\end{exercise}

\begin{exercise}
\label{exercises-exercise-curve-p1-p1}
Considérons la
variété projective\footnote{Le plongement projectif est
$((X_0, X_1), (Y_0, Y_1))\mapsto (X_0Y_0, X_0Y_1, X_1Y_0, X_1Y_1)$
ou, autrement dit, $(x, y)\mapsto (1, y, x, xy)$.}
$$
{\mathbf P}^1 \times {\mathbf P}^1 = {\mathbf P}^1_{{\mathbf C}}
\times_{\Spec({\mathbf C})} {\mathbf P}^1_{\mathbf C}
$$
sur le corps des nombres complexes ${\mathbf C}$. Elle est recouverte par quatre morceaux
affines,
correspondant aux couples de morceaux affines standards de ${\mathbf P}^1_{\mathbf
C}$. Par exemple,
prenons les coordonnées homogènes $X_0, X_1$ sur le premier facteur et
$Y_0, Y_1$ sur le second. Posons $x = X_1/X_0$, et $y = Y_1/Y_0$. Les 4
morceaux affines sont alors les spectres des anneaux
$$
{\mathbf C}[x, y], \quad
{\mathbf C}[x^{-1}, y], \quad
{\mathbf C}[x, y^{-1}], \quad
{\mathbf C}[x^{-1}, y^{-1}].
$$
Soit $X \subset {\mathbf P}^1 \times {\mathbf P}^1$ le sous-schéma fermé
qui est
l'adhérence de la partie fermée du premier morceau affine définie par l'équation
$$
y^3(x^4 + 1) = x^4 -1.
$$
\begin{enumerate}
\item Montrer que $X$ est contenu dans la réunion du premier et
du dernier des 4 morceaux affines.
\item Montrer que $X$ est une courbe projective non singulière.
\item Considérons le morphisme $pr_2 : X \to {\mathbf P}^1$ (projection sur
le premier facteur). Sur le premier morceau affine, il est donné par $(x, y) \mapsto x$.
Expliquer brièvement pourquoi il est de degré $3$.
\item Calculer les points de ramification et les indices de ramification
du morphisme $pr_2 : X \to {\mathbf P}^1$.
\item Calculer le genre de $X$.
\end{enumerate}
\end{exercise}

\begin{exercise}
\label{exercises-exercise-finite-type-over-Z}
Soit $X \to \Spec({\mathbf Z})$ un morphisme de type fini.
Supposons qu'il existe une infinité de nombres premiers $p$ tels que
$X \times_{\Spec({\mathbf Z})} \Spec({\mathbf F}_p)$ ne soit pas vide.
\begin{enumerate}
\item Montrer que $X \times_{\Spec({\mathbf Z})}\Spec(\mathbf{Q})$
n'est pas vide.
\item Pensez-vous que le même énoncé reste vrai si l'on remplace la condition
« de type fini » par la condition « localement de type fini » ?
\end{enumerate}
\end{exercise}




\section{Schémas, examen final, printemps 2009}
\label{exercises-section-final-exam-spring-2009}

\noindent
Voici les questions de l'examen final d'un cours sur les schémas,
donné au printemps 2009 à l'Université Columbia.

\begin{exercise}
\label{exercises-exercise-Noetherian-coherent}
Soit $X$ un schéma noethérien.
Soit $\mathcal{F}$ un faisceau cohérent sur $X$.
Soit $x \in X$ un point.
Supposons que $\text{Supp}(\mathcal{F}) = \{ x \}$.
\begin{enumerate}
\item Montrer que $x$ est un point fermé de $X$.
\item Montrer que $H^0(X, \mathcal{F})$ n'est pas nul.
\item Montrer que $\mathcal{F}$ est engendré par ses sections globales.
\item Montrer que $H^p(X, \mathcal{F}) = 0$ pour $p > 0$.
\end{enumerate}
\end{exercise}

\begin{remark}
\label{exercises-remark-invertible-projective-space}
Soit $k$ un corps.
Soit $\mathbf{P}^2_k = \text{Proj}(k[X_0, X_1, X_2])$.
Tout faisceau inversible sur $\mathbf{P}^2_k$ est isomorphe à
$\mathcal{O}_{\mathbf{P}^2_k}(n)$ pour un certain $n \in \mathbf{Z}$.
Rappelons que
$$
\Gamma(\mathbf{P}^2_k, \mathcal{O}_{\mathbf{P}^2_k}(n)) =
k[X_0, X_1, X_2]_n
$$
est la composante homogène de degré $n$ de l'anneau de polynômes.
Pour un faisceau quasi-cohérent $\mathcal{F}$ sur $\mathbf{P}^2_k$, posons
$\mathcal{F}(n) =
\mathcal{F}
\otimes_{\mathcal{O}_{\mathbf{P}^2_k}}
\mathcal{O}_{\mathbf{P}^2_k}(n)$
comme d'habitude.
\end{remark}

\begin{exercise}
\label{exercises-exercise-nonsplit-vectorbundle}
Soit $k$ un corps.
Soit $\mathcal{E}$ un fibré vectoriel sur $\mathbf{P}^2_k$, c'est-à-dire un
$\mathcal{O}_{\mathbf{P}^2_k}$-module localement libre de rang fini.
On dit que $\mathcal{E}$ est {\it scindé} si $\mathcal{E}$ est isomorphe à
une somme directe de $\mathcal{O}_{\mathbf{P}^2_k}$-modules inversibles.
\begin{enumerate}
\item Montrer que $\mathcal{E}$ est scindé si et seulement si $\mathcal{E}(n)$
est scindé.
\item Montrer que, si $\mathcal{E}$ est scindé, alors
$H^1({\mathbf{P}^2_k}, \mathcal{E}(n)) = 0$
pour tout $n \in \mathbf{Z}$.
\item Soit
$$
\varphi :
\mathcal{O}_{\mathbf{P}^2_k}
\longrightarrow
\mathcal{O}_{\mathbf{P}^2_k}(1) \oplus
\mathcal{O}_{\mathbf{P}^2_k}(1) \oplus
\mathcal{O}_{\mathbf{P}^2_k}(1)
$$
défini par les formes linéaires
$L_0, L_1, L_2 \in \Gamma(\mathbf{P}^2_k, \mathcal{O}_{\mathbf{P}^2_k}(1))$.
Supposons que $L_i \not = 0$ pour un certain $i$.
Quelle condition doivent vérifier $L_0, L_1, L_2$ pour que
le conoyau de $\varphi$ soit un fibré vectoriel ?
Pourquoi ?
\item Donner un exemple d'un tel $\varphi$.
\item Montrer que $\Coker(\varphi)$ n'est pas scindé (lorsqu'il s'agit
d'un fibré vectoriel).
\end{enumerate}
\end{exercise}

\begin{remark}
\label{exercises-remark-recall-dimension-theory}
On pourra utiliser librement les faits suivants de théorie de la dimension
(et en ajouter d'autres au besoin).
\begin{enumerate}
\item La dimension d'un schéma est la borne supérieure des longueurs des chaînes
de parties fermées irréductibles.
\item La dimension d'un schéma de type fini sur un corps est le maximum
des dimensions de ses ouverts affines.
\item La dimension d'un schéma noethérien est le maximum des dimensions
de ses composantes irréductibles.
\item La dimension d'un schéma affine
coïncide avec la dimension de l'anneau correspondant.
\item Soit $k$ un corps et soit $A$ une $k$-algèbre de type fini.
Si $A$ est intègre et $x \not = 0$, alors $\dim(A) = \dim(A/xA) + 1$.
\end{enumerate}
\end{remark}

\begin{exercise}
\label{exercises-exercise-irreducible-fibres-same-dimension-irreducible}
Soit $k$ un corps.
Soit $X$ un schéma projectif réduit sur $k$.
Soit $f : X \to \mathbf{P}^1_k$ un morphisme de schémas sur $k$.
Supposons qu'il existe un entier $d \geq 0$ tel que,
pour tout point $t \in \mathbf{P}^1_k$, la fibre $X_t = f^{-1}(t)$
soit irréductible de dimension $d$. (Rappelons qu'un espace irréductible
n'est pas vide.)
\begin{enumerate}
\item Montrer que $\dim(X) = d + 1$.
\item Soit $X_0 \subset X$ une composante irréductible de $X$ de dimension
$d + 1$. Démontrer que, pour tout $t \in \mathbf{P}^1_k$, la fibre
$X_{0, t}$ est de dimension $d$.
\item Que peut-on déduire de ce qui précède au sujet de $X_t$ et de $X_{0, t}$ ?
\item Montrer que $X$ est irréductible.
\end{enumerate}
\end{exercise}

\begin{remark}
\label{exercises-remark-chi}
Étant donnés un schéma projectif $X$ sur un corps $k$ et
un faisceau cohérent $\mathcal{F}$ sur $X$, nous posons
$$
\chi(X, \mathcal{F}) =
\sum\nolimits_{i \geq 0} (-1)^i\dim_k H^i(X, \mathcal{F}).
$$
\end{remark}

\begin{exercise}
\label{exercises-exercise-complete-intersection}
Soit $k$ un corps.
Écrivons $\mathbf{P}^3_k = \text{Proj}(k[X_0, X_1, X_2, X_3])$.
Soit $C \subset \mathbf{P}^3_k$ une
{\it courbe intersection complète de type $(5, 6)$}.
Cela signifie qu'il existe $F \in k[X_0, X_1, X_2, X_3]_5$ et
$G \in k[X_0, X_1, X_2, X_3]_6$ tels que
$$
C = \text{Proj}(k[X_0, X_1, X_2, X_3]/(F, G))
$$
soit une variété de dimension $1$. (Une variété est en particulier réduite et irréductible,
mais on peut, si on le souhaite, supposer que $C$ est non singulière.)
Soit $i : C \to \mathbf{P}^3_k$ l'immersion fermée associée.
Le fait d'être une intersection complète implique aussi que
$$
\xymatrix{
0 \ar[r] &
\mathcal{O}_{\mathbf{P}^3_k}(-11)
\ar[r]^-{
\left(
\begin{matrix}
-G \\
F
\end{matrix}
\right)
} &
\mathcal{O}_{\mathbf{P}^3_k}(-5) \oplus \mathcal{O}_{\mathbf{P}^3_k}(-6)
\ar[r]^-{(F, G)} &
\mathcal{O}_{\mathbf{P}^3_k} \ar[r] &
i_*\mathcal{O}_C \ar[r] &
0
}
$$
est une suite exacte de faisceaux. Utiliser ces faits pour :
\begin{enumerate}
\item calculer $\chi(C, i^*\mathcal{O}_{\mathbf{P}^3_k}(n))$ pour
tout $n \in \mathbf{Z}$, et
\item calculer la dimension de $H^1(C, \mathcal{O}_C)$.
\end{enumerate}
\end{exercise}

\begin{exercise}
\label{exercises-exercise-glueing}
Soit $k$ un corps.
Considérons les anneaux
\begin{align*}
A & = k[x, y]/(xy) \\
B & = k[u, v]/(uv) \\
C & = k[t, t^{-1}] \times k[s, s^{-1}]
\end{align*}
et les morphismes de $k$-algèbres
$$
\begin{matrix}
A \longrightarrow C, &
x \mapsto (t, 0), &
y \mapsto (0, s) \\
B \longrightarrow C, &
u \mapsto (t^{-1}, 0), &
v \mapsto (0, s^{-1})
\end{matrix}
$$
Ces morphismes induisent effectivement des isomorphismes
$A_{x + y} \to C$ et $B_{u + v} \to C$. Par conséquent, les morphismes $A \to C$
et $B \to C$ identifient $\Spec(C)$ à des ouverts de
$\Spec(A)$ et de $\Spec(B)$. Soit $X$ le schéma obtenu
en recollant $\Spec(A)$ et $\Spec(B)$ le long de $\Spec(C)$ :
$$
X = \Spec(A) \amalg_{\Spec(C)} \Spec(B).
$$
Comme nous l'avons vu dans le cours, un tel schéma existe et possède des
ouverts affines $\Spec(A) \subset X$ et $\Spec(B) \subset X$
dont l'intersection est exactement $\Spec(C)$, identifié à un ouvert de
chacun d'eux au moyen des morphismes ci-dessus.
\begin{enumerate}
\item Pourquoi $X$ est-il séparé ?
\item Pourquoi $X$ est-il de type fini sur $k$ ?
\item Calculer $H^1(X, \mathcal{O}_X)$, ou, autrement dit, quelle est sa dimension ?
\item Quelle description plus géométrique peut-on donner de $X$ ?
\end{enumerate}
\end{exercise}





\section{Schémas, examen final, automne 2010}
\label{exercises-section-final-exam-fall-2010}

\noindent
Voici les questions de l'examen final d'un cours sur les schémas,
donné à l'automne 2010 à l'Université Columbia.

\begin{exercise}[Définitions]
\label{exercises-exercise-definitions-fall-2010}
Donner la définition des notions suivantes.
\begin{enumerate}
\item un schéma séparé,
\item un morphisme quasi-compact de schémas,
\item un morphisme affine de schémas,
\item une partie multiplicative d'un anneau,
\item un schéma noethérien,
\item une variété.
\end{enumerate}
\end{exercise}

\begin{exercise}
\label{exercises-exercise-prime-avoidance}
Évitement des idéaux premiers.
\begin{enumerate}
\item Soit $A$ un anneau. Soit $I \subset A$ un idéal et
soient $\mathfrak q_1$, $\mathfrak q_2$ des idéaux premiers tels que
$I \not \subset \mathfrak q_i$. Montrer que
$I \not \subset \mathfrak q_1 \cup \mathfrak q_2$.
\item Quelle est une interprétation géométrique de (1) ?
\item Soit $X = \text{Proj}(S)$ pour un anneau gradué $S$.
Soient $x_1, x_2 \in X$. Montrer qu'il existe un ouvert standard
$D_{+}(F)$ qui contient à la fois $x_1$ et $x_2$.
\end{enumerate}
\end{exercise}

\begin{exercise}
\label{exercises-exercise-compose-affine}
Pourquoi le composé de morphismes affines est-il affine ?
\end{exercise}

\begin{exercise}[Exemples]
\label{exercises-exercise-examples}
Donner des exemples des objets suivants :
\begin{enumerate}
\item Un schéma projectif réductible sur un corps $k$.
\item Un schéma à 100 points.
\item Un morphisme de schémas qui n'est pas affine.
\end{enumerate}
\end{exercise}

\begin{exercise}
\label{exercises-exercise-chevalley-hilbert-nullstellensatz}
Le théorème de Chevalley et le théorème des zéros de Hilbert.
\begin{enumerate}
\item Soit $\mathfrak p \subset \mathbf{Z}[x_1, \ldots, x_n]$
un idéal maximal. Que déduit-on du théorème de Chevalley au sujet de
$\mathfrak p \cap \mathbf{Z}$ ?
\item Que déduit-on alors du théorème des zéros de Hilbert au sujet de
$\kappa(\mathfrak p)$ ?
\end{enumerate}
\end{exercise}

\begin{exercise}
\label{exercises-exercise-P0}
Soit $A$ un anneau. Soit $S = A[X]$ une $A$-algèbre graduée
où $X$ est de degré $1$. Montrer que $\text{Proj}(S) \cong \Spec(A)$
en tant que schémas sur $A$.
\end{exercise}

\begin{exercise}
\label{exercises-exercise-finite-is-projective}
Soit $A \to B$ un homomorphisme fini d'anneaux. Montrer que
$\Spec(B)$ est un schéma H-projectif sur $\Spec(A)$.
\end{exercise}

\begin{exercise}
\label{exercises-exercise-not-geometrically-irreducible}
Donner un exemple d'un schéma $X$ sur un corps $k$ tel que
$X$ soit irréductible et que, pour une certaine extension finie $k'/k$,
le changement de base $X_{k'} = X \times_{\Spec(k)} \Spec(k')$
soit connexe mais réductible.
\end{exercise}





\section{Schémas, examen final, printemps 2011}
\label{exercises-section-final-exam-spring-2011}

\noindent
Voici les questions de l'examen final d'un cours sur les schémas,
donné au printemps 2011 à l'Université Columbia.

\begin{exercise}[Définitions]
\label{exercises-exercise-definitions-spring-2011}
Donner la définition des notions en italique.
\begin{enumerate}
\item un schéma {\it séparé},
\item un morphisme de schémas {\it universellement fermé},
\item {\it $A$ domine $B$} pour des anneaux locaux $A, B$ contenus dans un
même corps,
\item la {\it dimension} d'un schéma $X$,
\item la {\it codimension} d'un sous-schéma fermé irréductible $Y$
d'un schéma $X$,
\end{enumerate}
\end{exercise}

\begin{exercise}[Résultats]
\label{exercises-exercise-results-spring-2011}
Pour chacun des points suivants, énoncer un résultat formellement équivalent au fait étudié
dans le cours.
\begin{enumerate}
\item Le critère valuatif de propreté pour un morphisme
$X \to Y$ de variétés, par exemple.
\item La relation entre $\dim(X)$ et le corps des fonctions
$k(X)$ de $X$ pour une variété $X$ sur un corps $k$.
\item Compléter : La catégorie des courbes projectives non singulières sur
$k$ et des morphismes non constants est antiéquivalente à $\ldots\ldots\ldots$.
\item La normalisation de Noether.
\item Le critère jacobien.
\end{enumerate}
\end{exercise}

\begin{exercise}
\label{exercises-exercise-genus-plane-curve}
Soit $k$ un corps.
Soit $F \in k[X_0, X_1, X_2]$ une forme homogène de degré $d$.
Supposons que $C = V_{+}(F) \subset \mathbf{P}^2_k$ soit une courbe lisse sur $k$.
Notons $i : C \to \mathbf{P}^2_k$ l'immersion fermée associée.
\begin{enumerate}
\item Montrer qu'il existe une suite exacte courte
$$
0 \to \mathcal{O}_{\mathbf{P}^2_k}(-d)
\to \mathcal{O}_{\mathbf{P}^2_k} \to i_*\mathcal{O}_C \to 0
$$
de faisceaux cohérents sur $\mathbf{P}^2_k$ : indiquer quelles sont les flèches
et expliquer brièvement pourquoi cette suite est exacte.
\item En déduire que $H^0(C, \mathcal{O}_C) = k$.
\item Calculer le genre de $C$.
\item Supposons maintenant que $P = (0 : 0 : 1)$ n'appartienne pas à $C$. Démontrer que
$\pi : C \to \mathbf{P}^1_k$ défini par $(a_0 : a_1 : a_2) \mapsto (a_0 : a_1)$
est de degré $d$.
\item Supposons que $k$ soit algébriquement clos, que tous les indices de ramification
(les ``$e_i$'') soient égaux à $1$ ou $2$, et que la caractéristique de $k$
ne soit pas égale à $2$. Combien de points de ramification
le morphisme $\pi : C \to \mathbf{P}^1_k$ possède-t-il ?
\item En fonction de $F$, quelles équations proposeriez-vous pour
l'ensemble des points de ramification de $\pi$ ?
\item Pouvez-vous deviner $K_C$ ?
\end{enumerate}
\end{exercise}

\begin{exercise}
\label{exercises-exercise-Pic-triangle}
Soit $k$ un corps. Soit $X$ un ``triangle'' sur $k$, c'est-à-dire
que l'on obtient $X$ en recollant trois copies de $\mathbf{A}^1_k$ les unes aux autres en
identifiant $0$ de la première copie à $1$ de la deuxième, $0$ de la
deuxième copie à $1$ de la troisième et $0$ de la troisième à $1$ de
la première. Il se trouve que $X$ est isomorphe à
$\Spec(k[x, y]/(xy(x + y + 1)))$ ; on pourra utiliser ce fait.
Calculer le groupe de Picard de $X$.
\end{exercise}

\begin{exercise}
\label{exercises-exercise-birational-morphism-curves-ample}
Soit $k$ un corps. Soit $\pi : X \to Y$ un morphisme fini birationnel
de courbes, où $X$ est une courbe projective non singulière sur $k$.
Il résulte du cours que $Y$ est une courbe propre et
que $\pi$ est le morphisme de normalisation de $Y$.
Nous avons aussi vu dans le cours qu'il existe un ouvert dense $V \subset Y$
tel que $U = \pi^{-1}(V)$ soit un ouvert dense de $X$ et que $\pi : U \to V$
soit un isomorphisme.
\begin{enumerate}
\item Montrer qu'il existe un diviseur de Cartier effectif $D \subset X$
tel que $D \subset U$ et que $\mathcal{O}_X(D)$ soit ample sur $X$.
\item Soit $D$ comme en (1). Montrer que $E = \pi(D)$ est un diviseur de Cartier
effectif sur $Y$.
\item Indiquer brièvement pourquoi
\begin{enumerate}
\item le morphisme $\mathcal{O}_Y \to \pi_*\mathcal{O}_X$ a un conoyau cohérent
$Q$ dont le support est contenu dans $Y \setminus V$, et
\item pour tout $n$, il existe un morphisme correspondant
$\mathcal{O}_Y(nE) \to \pi_*\mathcal{O}_X(nD)$
dont le conoyau est isomorphe à $Q$.
\end{enumerate}
\item Montrer que
$\dim_k H^0(X, \mathcal{O}_X(nD)) - \dim_k H^0(Y, \mathcal{O}_Y(nE))$
est bornée (par quoi ?) et en déduire que le faisceau inversible
$\mathcal{O}_Y(nE)$ possède beaucoup de sections pour $n$ grand (pourquoi ?).
\end{enumerate}
\end{exercise}





\section{Schémas, examen final, automne 2011}
\label{exercises-section-final-exam-fall-2011}

\noindent
Ces questions ont été posées lors de l'examen final d'un cours
d'algèbre commutative, donné à l'automne 2011 à l'Université Columbia.


\begin{exercise}[Définitions]
\label{exercises-exercise-definitions-fall-2011}
Donner la définition des notions en italique.
\begin{enumerate}
\item un anneau {\it noethérien},
\item un schéma {\it noethérien},
\item un homomorphisme d'anneaux {\it fini},
\item un morphisme de schémas {\it fini},
\item la {\it dimension} d'un anneau.
\end{enumerate}
\end{exercise}

\begin{exercise}[Résultats]
\label{exercises-exercise-results-fall-2011}
Énoncer un résultat formellement équivalent à celui qui a été étudié
dans le cours.
\begin{enumerate}
\item Le théorème principal de Zariski.
\item La normalisation de Noether.
\item Le théorème chinois.
\item La montée pour les homomorphismes d'anneaux finis.
\end{enumerate}
\end{exercise}

\begin{exercise}
\label{exercises-exercise-dimension-of-ring}
Soit $(A, \mathfrak m, \kappa)$ un anneau local noethérien
dont le corps résiduel est de caractéristique différente de $2$.
Supposons que $\mathfrak m$ soit engendré par trois éléments
$x, y, z$ et que $x^2 + y^2 + z^2 = 0$ dans $A$.
\begin{enumerate}
\item Quelles sont les valeurs possibles de $\dim(A)$ ?
\item Donner un exemple montrant que chacune de ces valeurs est possible.
\item Montrer que $A$ est intègre si $\dim(A) = 2$.
(Indication : considérer
$\bigoplus_{n \geq 0} \mathfrak m^n/\mathfrak m^{n + 1}$.)
\end{enumerate}
\end{exercise}

\begin{exercise}
\label{exercises-exercise-localization}
Soit $A$ un anneau.
Soient $S \subset T \subset A$ des parties multiplicatives.
Supposons que
$$
\{\mathfrak q \mid \mathfrak q \cap S = \emptyset\} =
\{\mathfrak q \mid \mathfrak q \cap T = \emptyset\}.
$$
Montrer que $S^{-1}A \to T^{-1}A$ est un isomorphisme.
\end{exercise}

\begin{exercise}
\label{exercises-exercise-locus-of-rank-1}
Soit $k$ un corps algébriquement clos. Soit
$$
V_0 = \{ A \in \text{Mat}(3 \times 3, k) \mid \text{rank}(A) = 1\}
\subset \text{Mat}(3 \times 3, k) = k^9.
$$
\begin{enumerate}
\item Montrer que $V_0$ est l'ensemble des points fermés d'une
partie localement fermée (pour la topologie de Zariski) $V \subset \mathbf{A}^9_k$.
\item $V$ est-il irréductible ?
\item Que vaut $\dim(V)$ ?
\end{enumerate}
\end{exercise}

\begin{exercise}
\label{exercises-exercise-not-complete-intersection}
Démontrer que l'idéal $(x^2, xy, y^2)$ de $\mathbf{C}[x, y]$
ne peut pas être engendré par $2$ éléments.
\end{exercise}

\begin{exercise}
\label{exercises-exercise-finite-projection}
Soit $f \in \mathbf{C}[x, y]$ un polynôme non constant.
Montrer qu'il existe $\alpha, \beta \in \mathbf{C}$ tels que
l'homomorphisme de $\mathbf{C}$-algèbres
$$
\mathbf{C}[t] \longrightarrow \mathbf{C}[x, y]/(f),\quad
t \longmapsto \alpha x + \beta y
$$
soit fini.
\end{exercise}

\begin{exercise}
\label{exercises-exercise-union-of-two-affines}
Montrer que, pour tous points $p_1, \ldots, p_n \in \mathbf{C}^2$,
le schéma $\mathbf{A}^2_\mathbf{C} \setminus \{p_1, \ldots, p_n\}$
est réunion de deux ouverts affines.
\end{exercise}

\begin{exercise}
\label{exercises-exercise-surjection-A1-P1}
Montrer qu'il existe un morphisme surjectif de schémas
$\mathbf{A}^1_\mathbf{C} \to \mathbf{P}^1_\mathbf{C}$.
(Surjectif signifie simplement surjectif sur les ensembles de points sous-jacents.)
\end{exercise}

\begin{exercise}
\label{exercises-exercise-almost-surjective}
Soit $k$ un corps algébriquement clos. Soit $A \subset B$ une extension
d'anneaux intègres qui sont tous deux des $k$-algèbres de type fini.
Montrer que l'image de $\Spec(B) \to \Spec(A)$
contient un ouvert non vide de $\Spec(A)$ en suivant les étapes ci-dessous :
\begin{enumerate}
\item Le démontrer lorsque $A \to B$ est en outre fini.
\item Le démontrer lorsque le corps des fractions de $B$ est une extension finie
du corps des fractions de $A$.
\item Réduire l'énoncé au cas précédent.
\end{enumerate}
\end{exercise}





\section{Schémas, examen final, automne 2013}
\label{exercises-section-final-exam-fall-2013}

\noindent
Ces questions ont été posées lors de l'examen final d'un cours
d'algèbre commutative, donné à l'automne 2013 à l'Université Columbia.

\begin{exercise}[Définitions]
\label{exercises-exercise-definitions-fall-2013}
Donner la définition des notions en italique.
\begin{enumerate}
\item un {\it idéal radical} d'un anneau,
\item un homomorphisme d'anneaux {\it de type fini},
\item une {\it différentielle au sens de Weil},
\item un {\it schéma}.
\end{enumerate}
\end{exercise}

\begin{exercise}[Résultats]
\label{exercises-exercise-results-fall-2013}
Énoncer un résultat formellement équivalent à celui qui a été étudié
dans le cours.
\begin{enumerate}
\item un résultat sur les polynômes de Hilbert des modules gradués.
\item la dimension d'un anneau local noethérien $(R, \mathfrak m)$ et de
$\bigoplus_{n \geq 0} \mathfrak m^n/\mathfrak m^{n + 1}$.
\item Riemann-Roch.
\item le théorème de Clifford.
\item le théorème de Chevalley.
\end{enumerate}
\end{exercise}

\begin{exercise}
\label{exercises-exercise-surjective-after-localization}
Soit $A \to B$ un homomorphisme d'anneaux. Soit $S \subset A$ une partie multiplicative.
Supposons que $A \to B$ soit de type fini et que $S^{-1}A \to S^{-1}B$ soit
surjectif. Montrer qu'il existe $f \in S$ tel que $A_f \to B_f$
soit surjectif.
\end{exercise}

\begin{exercise}
\label{exercises-exercise-injective-local-ring-map-not-surjective}
Donner un exemple d'un homomorphisme local injectif $A \to B$
d'anneaux locaux tel que $\Spec(B) \to \Spec(A)$ ne soit pas surjectif.
\end{exercise}

\begin{situation}[Notation : courbe plane]
\label{exercises-situation-curve-in-the-plane}
Soit $k$ un corps algébriquement clos. Soit
$F(X_0, X_1, X_2) \in k[X_0, X_1, X_2]$ un polynôme irréductible
homogène de degré $d$. Posons
$$
D = V(F) \subset \mathbf{P}^2
$$
la courbe plane projective définie par l'annulation de $F$.
Posons $x = X_1/X_0$ et $y = X_2/X_0$, ainsi que
$f(x, y) = X_0^{-d}F(X_0, X_1, X_2) = F(1, x, y)$.
Notons $K$ le corps des fractions de l'anneau intègre $k[x, y]/(f)$.
Soit $C$ la courbe abstraite correspondant à $K$.
Rappelons (d'après le cours) qu'il existe un morphisme surjectif $C \to D$
qui est bijectif au-dessus du lieu non singulier de $D$ et un
isomorphisme si $D$ est non singulière.
Posons $f_x = \partial f/\partial x$ et $f_y = \partial f/\partial y$.
Enfin, notons $\omega = \text{d}x/f_y = - \text{d}y/f_x$ cet
élément de $\Omega_{K/k}$ étudié dans le cours.
Notons $K_C$ le diviseur des zéros et des pôles de $\omega$.
\end{situation}

\begin{exercise}
\label{exercises-exercise-node-in-the-plane}
Dans la situation \ref{exercises-situation-curve-in-the-plane}, supposons
$d \geq 3$ et que la courbe $D$ ait exactement un point singulier,
à savoir $P = (1 : 0 : 0)$. Supposons en outre que le développement
$$
f(x, y) = xy + h.o.t
$$
soit valable au voisinage de $P = (0, 0)$. Alors $C$ a deux points $v$ et $w$ au-dessus de
$P$, caractérisés par
$$
v(x) = 1, v(y) > 1
\quad\text{et}\quad
w(x) > 1, w(y) = 1
$$
\begin{enumerate}
\item Montrer que l'élément
$\omega = \text{d}x/f_y = - \text{d}y/f_x$ de $\Omega_{K/k}$
a un pôle simple en $v$ et en $w$. (Le comportement de
$\omega$ aux points non singuliers est celui étudié dans le cours.)
\item Nous avons montré dans le cours que $\omega$ s'annule à l'ordre
$d - 3$ le long du diviseur $X_0 = 0$ ramené sur $C$ par le morphisme
$C \to D$. En combinant ce fait avec l'information de (1), quel est le degré
du diviseur des zéros et des pôles de $\omega$ sur $C$ ?
\item Quel est le genre de la courbe $C$ ?
\end{enumerate}
\end{exercise}

\begin{exercise}
\label{exercises-exercise-smooth-plane-curve-linear-system}
Dans la situation \ref{exercises-situation-curve-in-the-plane}, supposons $d = 5$
et que la courbe $C = D$ soit non singulière. Nous avons montré dans le cours
que le genre de $C$ vaut $6$ et que le système linéaire $K_C$ est donné par
$$
L(K_C) = \{h\omega \mid h \in k[x, y],\ \deg(h) \leq 2\}
$$
où $\deg$ désigne le degré total\footnote{On obtient $\leq 2$ parce que
$d - 3 = 5 - 3 = 2$.}. Soient $P_1, P_2, P_3, P_4, P_5 \in D$
des points deux à deux distincts situés dans l'ouvert affine $X_0 \not = 0$. Notons
$\sum P_i = P_1 + P_2 + P_3 + P_4 + P_5$ le diviseur correspondant sur $C$.
\begin{enumerate}
\item Décrire $L(K_C - \sum P_i)$ en termes de polynômes.
\item Quelles sont les valeurs possibles de $l(\sum P_i)$ ?
\end{enumerate}
\end{exercise}

\begin{exercise}
\label{exercises-exercise-rational-curve-high-degree}
Donner un $F$ comme dans la situation \ref{exercises-situation-curve-in-the-plane}
avec $d = 100$ tel que le genre de $C$ soit $0$.
\end{exercise}

\begin{exercise}
\label{exercises-exercise-high-degree-curve-quadratic-equation}
Soit $k$ un corps algébriquement clos. Soit $K/k$ une extension de corps
finiment engendrée de degré de transcendance $1$.
Soit $C$ la courbe abstraite correspondant à $K$.
Soit $V \subset K$ un $g^r_d$ et soit $\Phi : C \to \mathbf{P}^r$
le morphisme correspondant. Montrer que l'image de $C$
est contenue dans une quadrique\footnote{Une quadrique est une hypersurface de degré $2$,
c'est-à-dire le lieu des zéros dans $\mathbf{P}^r$ d'un polynôme homogène
de degré $2$.} si $V$ est un système linéaire complet
et si $d$ est assez grand par rapport au genre de $C$.
(Bonus : donner une bonne borne pour le degré nécessaire.)
\end{exercise}

\begin{exercise}
\label{exercises-exercise-surjective-map-a2-p2}
Notations de la situation \ref{exercises-situation-curve-in-the-plane}.
Soit $U \subset \mathbf{P}^2_k$ le sous-schéma ouvert
dont le complémentaire est $D$. Décrire la $k$-algèbre
$A = \mathcal{O}_{\mathbf{P}^2_k}(U)$. Donner une borne supérieure du
nombre de générateurs de $A$ comme $k$-algèbre.
\end{exercise}



\section{Schémas, examen final, printemps 2014}
\label{exercises-section-final-exam-spring-2014}

\noindent
Ces questions ont été posées lors de l'examen final d'un cours sur les
schémas, donné au printemps 2014 à l'Université Columbia.

\begin{exercise}[Définitions]
\label{exercises-exercise-definitions-spring-2014}
Soit $(X, \mathcal{O}_X)$ un schéma. Donner la définition des
notions en italique.
\begin{enumerate}
\item l'{\it anneau local de $X$ en un point $x$},
\item un faisceau {\it quasi-cohérent} de $\mathcal{O}_X$-modules,
\item un faisceau {\it cohérent} de $\mathcal{O}_X$-modules (supposer
que $X$ est localement noethérien),
\item un {\it ouvert affine} de $X$,
\item un {\it morphisme fini de schémas} $X \to Y$.
\end{enumerate}
\end{exercise}

\begin{exercise}[Théorèmes]
\label{exercises-exercise-results-spring-2014}
Énoncer précisément un résultat non trivial étudié dans le cours qui se rapporte
à chacun des sujets suivants.
\begin{enumerate}
\item l'invariance birationnelle des plurigenres des variétés,
\item le fait que la propriété d'un morphisme d'être affine est locale,
\item la topologie d'une fibre schématique d'un morphisme, et
\item le critère valuatif de propreté.
\end{enumerate}
\end{exercise}

\begin{exercise}
\label{exercises-exercise-miss-curve}
Soit $X = \mathbf{A}^2_\mathbf{C}$, où $\mathbf{C}$ est le corps
des nombres complexes. Une {\it droite} désignera un
sous-schéma fermé de $X$ défini par une équation linéaire $ax + by + c = 0$ pour
certains $a, b, c \in \mathbf{C}$ tels que $(a, b) \not = (0, 0)$.
Une {\it courbe} désignera un sous-schéma fermé irréductible (donc non vide)
$C \subset X$ de dimension $1$.
Une {\it conique} désignera une courbe définie par une
équation quadratique $ax^2 + bxy + cy^2 + dx + ey + f = 0$
pour certains $a, b, c, d, e, f \in \mathbf{C}$ tels que
$(a, b, c) \not = (0, 0, 0)$.
\begin{enumerate}
\item Trouver une courbe $C$ telle que toute droite ait une intersection non vide avec $C$.
\item Trouver une courbe $C$ telle que toute droite et toute conique aient une
intersection non vide avec $C$.
\item Montrer que, pour toute courbe $C$, il existe une autre courbe
telle que $C \cap C' = \emptyset$.
\end{enumerate}
\end{exercise}

\begin{exercise}
\label{exercises-exercise-normal-bundle-exceptional-curve}
Soit $k$ un corps. Soit $b : X \to \mathbf{A}^2_k$ l'éclatement
du plan affine à l'origine. Autrement dit, si
$\mathbf{A}^2_k = \Spec(k[x, y])$, alors
$X = \text{Proj}(\bigoplus_{n \geq 0} \mathfrak m^n)$
où $\mathfrak m = (x, y) \subset k[x, y]$.
Démontrer les assertions suivantes :
\begin{enumerate}
\item la fibre schématique $E$ de $b$ au-dessus de l'origine
est isomorphe à $\mathbf{P}^1_k$,
\item $E$ est un diviseur de Cartier effectif sur $X$,
\item la restriction de $\mathcal{O}_X(-E)$
à $E$ est un faisceau inversible de degré $1$.
\end{enumerate}
(Rappelons que $\mathcal{O}_X(-E)$ est le faisceau d'idéaux de $E$ dans $X$.)
\end{exercise}

\begin{exercise}
\label{exercises-exercise-surjective-map-affine-variety-projective-variety}
Soit $k$ un corps. Soit $X$ une variété projective sur $k$.
Montrer qu'il existe une variété affine $U$ sur $k$ et un morphisme surjectif
de variétés $U \to X$.
\end{exercise}

\begin{exercise}
\label{exercises-exercise-vandermonde}
Soit $k$ un corps de caractéristique $p > 0$ différente de $2,3$.
Considérons le sous-schéma fermé $X$ de $\mathbf{P}^n_k$ défini par
$$
\sum\nolimits_{i = 0, \ldots, n} X_i = 0,\quad
\sum\nolimits_{i = 0, \ldots, n} X_i^2 = 0,\quad
\sum\nolimits_{i = 0, \ldots, n} X_i^3 = 0
$$
Pour quelles paires $(n, p)$ cette variété est-elle singulière ?
\end{exercise}






\section{Algèbre commutative, examen final, automne 2016}
\label{exercises-section-final-exam-fall-2016}

\noindent
Ces questions ont été posées lors de l'examen final d'un cours
d'algèbre commutative, donné à l'automne 2016 à l'Université Columbia.

\begin{exercise}[Définitions]
\label{exercises-exercise-definitions-fall-2016}
Soit $R$ un anneau.
Donner la définition des notions en italique.
\begin{enumerate}
\item l'{\it anneau local de $R$ en un idéal premier $\mathfrak p$},
\item un $R$-module {\it de type fini},
\item un $R$-module {\it de présentation finie},
\item $R$ est {\it régulier},
\item $R$ est {\it caténaire},
\item $R$ est {\it de Cohen-Macaulay}.
\end{enumerate}
\end{exercise}

\begin{exercise}[Théorèmes]
\label{exercises-exercise-results-fall-2016}
Énoncer précisément un résultat non trivial étudié dans le cours qui se rapporte
à chacun des sujets suivants.
\begin{enumerate}
\item les anneaux réguliers,
\item les idéaux premiers associés des modules de Cohen-Macaulay,
\item la dimension d'un anneau intègre de type fini sur un corps, et
\item le théorème de Chevalley.
\end{enumerate}
\end{exercise}

\begin{exercise}
\label{exercises-exercise-finite-injective}
Soit $A \to B$ un homomorphisme d'anneaux tel que
\begin{enumerate}
\item $A$ soit local, d'idéal maximal $\mathfrak m$,
\item $A \to B$ soit fini\footnote{Rappelons que cela signifie que $B$
est de type fini en tant que $A$-module.},
\item $A \to B$ soit injectif (nous considérons $A$ comme un sous-anneau de $B$).
\end{enumerate}
Montrer qu'il existe un idéal premier $\mathfrak q \subset B$ tel que
$\mathfrak m = A \cap \mathfrak q$.
\end{exercise}

\begin{exercise}
\label{exercises-exercise-find-components}
Soit $k$ un corps. Soit $R = k[x, y, z, w]$.
Considérons l'idéal $I = (xy, xz, xw)$.
Quelles sont les composantes irréductibles de
$V(I) \subset \Spec(R)$ et quelles sont leurs dimensions ?
\end{exercise}

\begin{exercise}
\label{exercises-exercise-no-nonconstant-morphism}
Soit $k$ un corps. Soient $A = k[x, x^{-1}]$ et $B = k[y]$.
Montrer que tout homomorphisme de $k$-algèbres $\varphi : A \to B$ envoie $x$ sur une constante.
\end{exercise}

\begin{exercise}
\label{exercises-exercise-regular-over-Z}
Considérons l'anneau $R = \mathbf{Z}[x, y]/(xy - 7)$.
Démontrer que $R$ est régulier.
\end{exercise}

\noindent
Étant donné un anneau local noethérien $(R, \mathfrak m, \kappa)$,
pour $n \geq 0$, posons
$\varphi_R(n) = \dim_\kappa(\mathfrak m^n/\mathfrak m^{n + 1})$.

\begin{exercise}
\label{exercises-exercise-hilbert-function-allowed}
Existe-t-il un anneau local noethérien $R$ tel que
$\varphi_R(n) = n + 1$ pour tout $n \geq 0$ ?
\end{exercise}

\begin{exercise}
\label{exercises-exercise-hilbert-function-prohibited}
Soit $R$ un anneau local noethérien.
Supposons que $\varphi_R(0) = 1$, $\varphi_R(1) = 3$,
$\varphi_R(2) = 5$. Montrer que $\varphi_R(3) \leq 7$.
\end{exercise}





\section{Schémas, examen final, printemps 2017}
\label{exercises-section-final-exam-spring-2017}

\noindent
Ces questions ont été posées lors de l'examen final d'un cours sur les schémas,
donné au printemps 2017 à l'Université Columbia.

\begin{exercise}[Définitions]
\label{exercises-exercise-definitions-spring-2017}
Soit $f : X \to Y$ un morphisme de schémas.
Donner une brève définition des notions en italique.
\begin{enumerate}
\item la {\it fibre schématique} de $f$ en $y \in Y$,
\item $f$ est un {\it morphisme fini},
\item un {\it faisceau quasi-cohérent} de $\mathcal{O}_X$-modules,
\item $X$ est une {\it variété},
\item $f$ est un {\it morphisme lisse},
\item $f$ est un {\it morphisme propre}.
\end{enumerate}
\end{exercise}

\begin{exercise}[Théorèmes]
\label{exercises-exercise-results-spring-2017}
Énoncer avec précision, mais brièvement, un résultat non trivial étudié dans le cours
en rapport avec chacun des sujets suivants.
\begin{enumerate}
\item l'image directe des faisceaux quasi-cohérents,
\item la cohomologie des faisceaux cohérents sur les variétés projectives,
\item la dualité de Serre pour un schéma projectif sur un corps, et
\item la formule de Riemann-Hurwitz.
\end{enumerate}
\end{exercise}

\begin{exercise}
\label{exercises-exercise-compute-degree}
Soit $k$ un corps algébriquement clos.
Soit $\ell > 100$ un nombre premier différent de la
caractéristique de $k$. Soit $X$ le modèle projectif non singulier
de la courbe affine donnée par l'équation
$$
y^\ell = x(x - 1)^3
$$
dans $\mathbf{A}^2_k$. Répondre aux questions suivantes :
\begin{enumerate}
\item Quel est le genre de $X$ ?
\item Donner une borne supérieure pour la
gonalité\footnote{La gonalité est le plus petit degré d'un morphisme non constant
de $X$ vers $\mathbf{P}^1_k$.} de $X$.
\end{enumerate}
\end{exercise}

\begin{exercise}
\label{exercises-exercise-count-components}
Soit $k$ un corps algébriquement clos. Soit $X$ un schéma réduit,
projectif sur $k$, dont toutes les composantes irréductibles
sont de même dimension $1$. Soit $\omega_{X/k}$ le
module dualisant relatif. Montrer que, si
$\dim_k H^1(X, \omega_{X/k}) > 1$, alors $X$ est non connexe.
\end{exercise}

\begin{exercise}
\label{exercises-exercise-sections-L-L-inverse}
Donner un exemple d'un schéma $X$ et d'un
$\mathcal{O}_X$-module inversible non trivial $\mathcal{L}$ tels que
$H^0(X, \mathcal{L})$ et $H^0(X, \mathcal{L}^{\otimes -1})$
soient tous deux non nuls.
\end{exercise}

\begin{exercise}
\label{exercises-exercise-in-product}
Soit $k$ un corps algébriquement clos. Soit $g \geq 3$.
Soient $X$ et $X'$ des courbes projectives lisses sur $k$, de
genres respectifs $g$ et $g + 1$. Soit $Y \subset X \times X'$
une courbe telle que les projections $Y \to X$ et $Y \to X'$
soient non constantes. Démontrer que le modèle projectif non singulier
de $Y$ est de genre $\geq 2g + 1$.
\end{exercise}

\begin{exercise}
\label{exercises-exercise-finitely-many}
Soit $k$ un corps fini. Soit $g > 1$. Esquisser une démonstration de ce qui suit :
il n'existe qu'un nombre fini de classes d'isomorphisme de
courbes projectives lisses sur $k$ et de genre $g$.
(Même une simple tentative de réponse sera prise en compte.)
\end{exercise}







\section{Algèbre commutative, examen final, automne 2017}
\label{exercises-section-final-exam-fall-2017}

\noindent
Voici les questions posées à l'examen final d'un cours d'algèbre commutative,
à l'automne 2017, à l'Université Columbia.

\begin{exercise}[Définitions]
\label{exercises-exercise-definitions-fall-2017}
Donner une brève définition de chacun des concepts en italique.
\begin{enumerate}
\item l'{\it adjoint à gauche} d'un foncteur $F : \mathcal{A} \to \mathcal{B}$,
\item le {\it degré de transcendance} d'une extension de corps $L/K$,
\item une {\it fonction régulière} sur une variété affine classique
$X \subset k^n$,
\item un {\it faisceau} sur un espace topologique,
\item un {\it anneau local}, et
\item un morphisme de schémas $f : X \to Y$ qui est {\it affine}.
\end{enumerate}
\end{exercise}

\begin{exercise}[Théorèmes]
\label{exercises-exercise-results-fall-2017}
Énoncer avec précision mais brièvement un résultat non trivial vu en cours
relatif à chacun des points suivants (s'il en existe plusieurs, en choisir
un seul).
\begin{enumerate}
\item lemme de Yoneda,
\item suite de Mayer-Vietoris,
\item dimension et cohomologie,
\item polynôme de Hilbert, et
\item dualité pour l'espace projectif.
\end{enumerate}
\end{exercise}

\begin{exercise}
\label{exercises-exercise-compute-dimension}
Soit $k$ un corps algébriquement clos. Considérons la partie fermée
$X$ de $k^5$ muni de la topologie de Zariski et des coordonnées $x_1, x_2, x_3, x_4, x_5$,
définie par les équations
$$
x_1^2 - x_4 = 0,\quad
x_2^5 - x_5 = 0,\quad
x_3^2 + x_3 + x_4 + x_5 = 0
$$
Quelle est la dimension de $X$, et pourquoi ?
\end{exercise}

\begin{exercise}
\label{exercises-exercise-can-there-be}
Soit $k$ un corps. Soit $X = \mathbf{P}^1_k$ l'espace projectif
de dimension $1$ sur $k$. Soit $\mathcal{E}$ un
$\mathcal{O}_X$-module localement libre de rang fini. Pour $d \in \mathbf{Z}$, notons
$\mathcal{E}(d) = \mathcal{E} \otimes_{\mathcal{O}_X} \mathcal{O}_X(d)$
le tordu de Serre d'indice $d$ de $\mathcal{E}$ et
$h^i(X, \mathcal{E}(d)) = \dim_k H^i(X, \mathcal{E}(d))$.
\begin{enumerate}
\item Pourquoi n'existe-t-il pas de $\mathcal{E}$ tel que
$h^0(X, \mathcal{E}) = 5$ et $h^0(X, \mathcal{E}(1)) = 4$ ?
\item Pourquoi n'existe-t-il pas de $\mathcal{E}$ tel que
$h^1(X, \mathcal{E}(1)) = 5$ et $h^1(X, \mathcal{E}) = 4$ ?
\item Pour quels $a \in \mathbf{Z}$ peut-il exister un fibré vectoriel
$\mathcal{E}$ sur $X$ tel que
$$
\begin{matrix}
h^0(X, \mathcal{E})\phantom{(1)} = 1 &
h^1(X, \mathcal{E})\phantom{(1)} = 1 \\
h^0(X, \mathcal{E}(1)) = 2 &
h^1(X, \mathcal{E}(1)) = 0 \\
h^0(X, \mathcal{E}(2)) = 4 &
h^1(X, \mathcal{E}(2)) = a
\end{matrix}
$$
\end{enumerate}
Les réponses partielles sont les bienvenues et vivement encouragées.
\end{exercise}

\begin{exercise}
\label{exercises-exercise-banana}
Soit $X$ un espace topologique qui est la réunion
$X = Y \cup Z$ de deux parties fermées
$Y$ et $Z$, dont l'intersection est notée $W = Y \cap Z$.
Notons $i : Y \to X$, $j : Z \to X$ et $k : W \to X$ les applications
d'inclusion.
\begin{enumerate}
\item Montrer qu'il existe une suite exacte courte de faisceaux
$$
0 \to \underline{\mathbf{Z}}_X \to
i_*(\underline{\mathbf{Z}}_Y) \oplus
j_*(\underline{\mathbf{Z}}_Z) \to
k_*(\underline{\mathbf{Z}}_W) \to 0
$$
où $\underline{\mathbf{Z}}_X$ désigne le faisceau constant
de valeur $\mathbf{Z}$ sur $X$, et ainsi de suite.
\item Quelle relation peut-on en déduire entre les
groupes de cohomologie de $X$, $Y$, $Z$ et $W$ à coefficients dans $\mathbf{Z}$ ?
\end{enumerate}
\end{exercise}

\begin{exercise}
\label{exercises-exercise-cohomology-infinite-punctured}
Soit $k$ un corps. Soit $A = k[x_1, x_2, x_3, \ldots]$ l'anneau de polynômes
en une infinité d'indéterminées.
Notons $\mathfrak m$ l'idéal maximal de $A$ engendré par
toutes les indéterminées. Soient $X = \Spec(A)$ et $U = X \setminus \{\mathfrak m\}$.
\begin{enumerate}
\item Montrer que $H^1(U, \mathcal{O}_U) = 0$. Indication : calcul de
cohomologie de {\v C}ech.
\item Quelle valeur conjecturez-vous pour $H^i(U, \mathcal{O}_U)$ lorsque $i \geq 1$ ?
\end{enumerate}
\end{exercise}

\begin{exercise}
\label{exercises-exercise-principal}
Soit $A$ un anneau local. Soit $a \in A$ un non diviseur de zéro.
Soient $I, J \subset A$ des idéaux tels que $IJ = (a)$. Montrer que
l'idéal $I$ est principal, c'est-à-dire engendré par un seul élément
(qui se révélera être un non diviseur de zéro).
\end{exercise}



\section{Schémas, examen final, printemps 2018}
\label{exercises-section-final-exam-spring-2018}

\noindent
Voici les questions posées à l'examen final d'un cours sur les schémas,
au printemps 2018, à l'Université Columbia.

\begin{exercise}[Définitions]
\label{exercises-exercise-definitions-spring-2018}
Donner une brève définition de chacun des concepts en italique.
Soit $k$ un corps algébriquement clos.
Soit $X$ une courbe projective sur $k$.
\begin{enumerate}
\item une algèbre {\it lisse} sur $k$,
\item le {\it degré} d'un $\mathcal{O}_X$-module inversible sur $X$,
\item le {\it genre} de $X$,
\item le {\it groupe des classes de diviseurs de Weil} de $X$,
\item $X$ est {\it hyperelliptique}, et
\item le {\it nombre d'intersection}
de deux courbes sur une surface projective lisse sur $k$.
\end{enumerate}
\end{exercise}

\begin{exercise}[Théorèmes]
\label{exercises-exercise-results-spring-2018}
Énoncer avec précision mais brièvement un résultat non trivial vu en cours
relatif à chacun des points suivants (s'il en existe plusieurs, en choisir
un seul).
\begin{enumerate}
\item théorème de Riemann-Hurwitz,
\item théorème de Clifford,
\item factorisation des morphismes entre surfaces projectives lisses,
\item théorème de l'indice de Hodge, et
\item hypothèse de Riemann pour les courbes sur les corps finis.
\end{enumerate}
\end{exercise}

\begin{exercise}
\label{exercises-exercise-hyperelliptic}
Soit $k$ un corps algébriquement clos. Soit $X \subset \mathbf{P}^3_k$
une courbe lisse de degré $d$ et de genre $\geq 2$. Supposons que $X$
ne soit contenue dans aucun plan et qu'il existe
une droite $\ell$ de $\mathbf{P}^3_k$ rencontrant $X$ en $d - 2$ points.
Montrer que $X$ est hyperelliptique.
\end{exercise}

\begin{exercise}
\label{exercises-exercise-singular}
Soit $k$ un corps algébriquement clos. Soit $X$ une courbe projective
ayant pour points singuliers deux à deux distincts $p_1, \ldots, p_n$.
Expliquer pourquoi le genre de la normalisée de $X$ est au plus
$-n + \dim_k H^1(X, \mathcal{O}_X)$.
\end{exercise}

\begin{exercise}
\label{exercises-exercise-how-many-blowups}
Soit $k$ un corps. Soit $X = \Spec(k[x, y])$ l'espace affine de dimension $2$.
Soit
$$
I = (x^3, x^2y, xy^2, y^3) \subset k[x, y].
$$
Soit $Y \subset X$ le sous-schéma fermé correspondant à $I$.
Soit $b : X' \to X$ l'éclatement de l'idéal $(x, y)$, c'est-à-dire
l'éclatement de l'espace affine en l'origine.
\begin{enumerate}
\item Montrer que l'image réciproque schématique $b^{-1}Y \subset X'$
est un diviseur de Cartier effectif.
\item Donner un exemple d'idéal $J \subset k[x, y]$
tel que $I \subset J \subset (x, y)$ et que, si $Z \subset X$
est le sous-schéma fermé correspondant à $J$, alors l'image réciproque
schématique $b^{-1}Z$ ne soit pas un diviseur de Cartier effectif.
\end{enumerate}
\end{exercise}

\begin{exercise}
\label{exercises-exercise-rational-surface}
Soit $k$ un corps algébriquement clos. Considérons les types
de surfaces suivants :
\begin{enumerate}
\item $S = C_1 \times C_2$ où $C_1$ et $C_2$ sont des courbes projectives lisses,
\item $S = C_1 \times C_2$ où $C_1$ et $C_2$ sont des courbes projectives lisses
et le genre de $C_1$ est $> 0$,
\item $S \subset \mathbf{P}^3_k$ est une hypersurface de degré $4$, et
\item $S \subset \mathbf{P}^3_k$ est une hypersurface lisse de degré $4$.
\end{enumerate}
Pour chaque type, indiquer brièvement si la classe des surfaces de
ce type contient ou non des surfaces rationnelles, et pourquoi.
\end{exercise}

\begin{exercise}
\label{exercises-exercise-self-square-line}
Soit $k$ un corps algébriquement clos. Soit $S \subset \mathbf{P}^3_k$
une hypersurface lisse de degré $d$. Supposons que $S$ contienne une
droite $\ell$. Quelle est l'auto-intersection de $\ell$ considérée comme diviseur sur $S$ ?
\end{exercise}






\section{Algèbre commutative, examen final, automne 2019}
\label{exercises-section-final-exam-fall-2019}

\noindent
Voici les questions posées à l'examen final d'un cours d'algèbre commutative,
à l'automne 2019, à l'Université Columbia.

\begin{exercise}[Définitions]
\label{exercises-exercise-definitions-fall-2019}
Donner une brève définition de chacun des concepts en italique.
\begin{enumerate}
\item une {\it partie constructible} d'un espace topologique noethérien,
\item la {\it localisation} d'un $R$-module $M$ en un idéal premier $\mathfrak p$,
\item la {\it longueur} d'un module sur un anneau local noethérien
$(A, \mathfrak m, \kappa)$,
\item un {\it module projectif} sur un anneau $R$, et
\item un {\it module de Cohen-Macaulay} sur un
anneau local noethérien $(A, \mathfrak m, \kappa)$.
\end{enumerate}
\end{exercise}

\begin{exercise}[Théorèmes]
\label{exercises-exercise-results-fall-2019}
Énoncer avec précision mais brièvement un résultat non trivial vu en cours
relatif à chacun des points suivants (s'il en existe plusieurs, en choisir
un seul).
\begin{enumerate}
\item images de parties constructibles,
\item théorème des zéros de Hilbert,
\item dimension des algèbres de type fini sur un corps,
\item lemme de normalisation de Noether, et
\item anneaux locaux réguliers.
\end{enumerate}
\end{exercise}

\noindent
Pour un anneau $R$ et un idéal $I \subset R$, rappelons que $V(I)$
désigne l'ensemble des $\mathfrak p \in \Spec(R)$ tels que $I \subset \mathfrak p$.

\begin{exercise}[Construction d'idéaux premiers]
\label{exercises-exercise-infinitely-many-primes}
Construire une infinité d'idéaux premiers distincts
$\mathfrak p \subset \mathbf{C}[x, y]$
tels que $V(\mathfrak p)$ contienne $(x, y)$ et $(x - 1, y - 1)$.
\end{exercise}

\begin{exercise}[Aucun idéal premier]
\label{exercises-exercise-no-prime}
Soit $R = \mathbf{C}[x, y, z]/(xy)$. Justifier brièvement qu'il n'existe pas
d'idéal premier $\mathfrak p \subset R$ tel que
$V(\mathfrak p)$ contienne $(x, y - 1, z - 5)$ et $(x - 1, y, z - 7)$.
\end{exercise}

\begin{exercise}[Frobenius]
\label{exercises-exercise-frobenius}
Soit $p$ un nombre premier (on pourra supposer $p = 2$ pour simplifier les formules).
Soit $R$ un anneau tel que $p = 0$ dans $R$.
\begin{enumerate}
\item Montrer que l'application $F : R \to R$, $x \mapsto x^p$ est un homomorphisme d'anneaux.
\item Montrer que $\Spec(F) : \Spec(R) \to \Spec(R)$ est l'identité.
\end{enumerate}
\end{exercise}

\noindent
Rappelons qu'une {\it spécialisation} $x \leadsto y$ de points d'un espace topologique
signifie simplement que $y$ appartient à l'adhérence de $x$. Nous disons que $x \leadsto y$ est une
{\it spécialisation immédiate} s'il n'existe aucun $z$ distinct
de $x$ et de $y$ tel que $x \leadsto z$ et $z \leadsto y$.

\begin{exercise}[Dimension]
\label{exercises-exercise-dimension}
Supposons que nous ayons un espace topologique sobre $X$ contenant $5$ points distincts
$x, y, z, u, v$ et ayant les spécialisations suivantes :
$$
\xymatrix{
x \ar[d] \ar[r] & u & v \ar[l] \ar[dl] \\
y \ar[r] & z
}
$$
Quelle est la dimension minimale que peut avoir un tel $X$ ? Si $X$ est le spectre
d'une algèbre de type fini sur un corps et si $x \leadsto u$ est une
spécialisation immédiate, que peut-on dire de la
spécialisation $v \leadsto z$ ?
\end{exercise}

\begin{exercise}[Calcul de Tor]
\label{exercises-exercise-tor-computation}
Soit $R = \mathbf{C}[x, y, z]$. Soient $M = R/(x, z)$ et $N = R/(y, z)$.
Pour quels $i \in \mathbf{Z}$ le module $\text{Tor}_i^R(M, N)$ est-il non nul ?
\end{exercise}

\begin{exercise}
\label{exercises-exercise-depth-goes-up}
Soit $A \to B$ un homomorphisme local plat d'anneaux locaux noethériens.
Montrer que, si $A$ est de profondeur $k$, alors $B$ est de profondeur au moins $k$.
\end{exercise}









\section{Géométrie algébrique, examen final, printemps 2020}
\label{exercises-section-final-exam-spring-2020}

\noindent
Voici les questions posées à l'examen final d'un cours de géométrie algébrique,
au printemps 2020, à l'Université Columbia.

\begin{exercise}[Définitions]
\label{exercises-exercise-definitions-spring-2020}
Donner une brève définition de chacun des concepts en italique.
\begin{enumerate}
\item un {\it schéma},
\item un {\it morphisme de schémas},
\item un {\it module quasi-cohérent} sur un schéma,
\item une {\it variété} sur un corps $k$,
\item une {\it courbe} sur un corps $k$,
\item un {\it morphisme fini} de schémas,
\item la {\it cohomologie} d'un faisceau de groupes abéliens $\mathcal{F}$
sur un espace topologique $X$,
\item un {\it faisceau dualisant} sur un schéma $X$ de dimension $d$
propre sur un corps $k$, et
\item une {\it application rationnelle} d'une variété $X$ dans une variété $Y$.
\end{enumerate}
\end{exercise}

\begin{exercise}[Théorèmes]
\label{exercises-exercise-results-spring-2020}
Énoncer avec précision mais brièvement un résultat non trivial vu en cours
relatif à chacun des points suivants (s'il en existe plusieurs, en choisir
un seul).
\begin{enumerate}
\item cohomologie des faisceaux abéliens sur un espace topologique noethérien
$X$ de dimension $d$,
\item faisceau des différentielles $\Omega^1_{X/k}$ d'une variété lisse
sur un corps $k$,
\item faisceau dualisant $\omega_X$ d'une variété projective lisse $X$
sur le corps $k$,
\item une courbe propre lisse
de genre $0$ sur un corps algébriquement clos $k$, et
\item le genre d'une courbe plane de degré $d$.
\end{enumerate}
\end{exercise}

\begin{exercise}
\label{exercises-exercise-coh-P1-infty-doubled}
Soit $k$ un corps. Soit $X$ un schéma sur $k$. Supposons que $X = X_1 \cup X_2$
soit un recouvrement ouvert, $X_1$ et $X_2$ étant tous deux isomorphes à $\mathbf{P}^1_k$,
et $X_1 \cap X_2$ étant isomorphe à $\mathbf{A}^1_k$. (Un tel schéma existe :
par exemple, on peut prendre $\mathbf{P}^1_k$ avec le point $\infty$ dédoublé.)
Montrer que $\dim_k H^1(X, \mathcal{O}_X)$ est infinie.
\end{exercise}

\begin{exercise}
\label{exercises-exercise-no-morphism}
Soit $k$ un corps algébriquement clos. Soit $Y$ une courbe projective lisse
de genre $10$. Trouver une bonne borne inférieure pour le
genre d'une courbe projective lisse
$X$ telle qu'il existe un morphisme non constant $f : X \to Y$
qui ne soit pas un isomorphisme.
\end{exercise}

\begin{exercise}
\label{exercises-exercise-element-order-n-in-pic}
Soit $k$ un corps algébriquement clos de caractéristique $0$.
Soit
$$
X : T_0^d + T_1^d - T_2^d = 0 \subset \mathbf{P}^2_k
$$
la courbe de Fermat de degré $d \geq 3$. Considérons les points fermés
$p = [1 : 0 : 1]$ et $q = [0 : 1 : 1]$ de $X$. Posons $D = [p] - [q]$.
\begin{enumerate}
\item Montrer que $D$ est non trivial dans le groupe des classes de diviseurs de Weil.
\item Montrer que $d D$ est trivial dans le groupe des classes de diviseurs de Weil.
(Indication : essayer de montrer que $d[p]$ et $d[q]$ sont tous deux l'intersection
de $X$ avec une droite du plan.)
\end{enumerate}
\end{exercise}

\begin{exercise}
\label{exercises-exercise-number-equations}
Soit $k$ un corps algébriquement clos. Considérons le plongement $2$-uple
$$
\varphi : \mathbf{P}^2 \longrightarrow \mathbf{P}^5
$$
Dans la terminologie et les notations du cours, il s'agit du morphisme
$$
\varphi = \varphi_{\mathcal{O}_{\mathbf{P}^2}(2)} :
\mathbf{P}^2
\longrightarrow
\mathbf{P}(\Gamma(\mathbf{P}^2, \mathcal{O}_{\mathbf{P}^2}(2)))
$$
En coordonnées homogènes, il est donné par
$$
[a_0 : a_1 : a_2] \longmapsto
[a_0^2 : a_0a_1 : a_0a_2 : a_1^2 : a_1a_2 : a_2^2]
$$
C'est une immersion fermée (on pourra simplement l'admettre).
Soit $I \subset k[T_0, \ldots, T_5]$ l'idéal homogène
de $\varphi(\mathbf{P}^2)$, c'est-à-dire que les éléments de la
composante homogène $I_d$ sont les
polynômes homogènes $F(T_0, \ldots, T_5)$ de degré $d$ dont la
restriction au sous-schéma fermé $\varphi(\mathbf{P}^2)$ est nulle.
Calculer $\dim_k I_d$ en fonction de $d$.
\end{exercise}

\begin{exercise}
\label{exercises-exercise-dualizing-sheaf-rank-1}
Soit $k$ un corps algébriquement clos. Soit $X$ un schéma propre de
dimension $d$ sur $k$, muni d'un module dualisant $\omega_X$.
On dispose des informations suivantes :
\begin{enumerate}
\item $\text{Ext}^i_X(\mathcal{F}, \omega_X) \times
H^{d - i}(X, \mathcal{F}) \to H^d(X, \omega_X) \xrightarrow{t} k$
est non dégénéré pour tout $i$ et tout $\mathcal{O}_X$-module cohérent
$\mathcal{F}$, et
\item $\omega_X$ est localement libre de rang fini, disons $r$.
\end{enumerate}
Montrer que $r = 1$. (Indication : examiner ce qui se passe en prenant pour $\mathcal{F}$
un module convenable dont le support est un point fermé.)
\end{exercise}







\section{Algèbre commutative, examen final, automne 2021}
\label{exercises-section-final-exam-fall-2021}

\noindent
Les questions suivantes ont été posées lors de l'examen final d'un cours d'algèbre commutative,
donné à l'automne 2021 à l'Université Columbia.

\begin{exercise}[Définitions]
\label{exercises-exercise-definitions-fall-2021}
Donner des définitions succinctes des notions en italique.
\begin{enumerate}
\item une {\it partie multiplicative} d'un anneau $A$,
\item un {\it anneau artinien} $A$,
\item le {\it spectre d'un anneau} $A$ en tant qu'espace topologique,
\item un {\it homomorphisme d'anneaux plat} $A \to B$,
\item la {\it hauteur} d'un idéal premier $\mathfrak p$ de $A$, et
\item les foncteurs {\it $\text{Tor}^A_i(-, -)$} sur un anneau $A$.
\end{enumerate}
\end{exercise}

\begin{exercise}[Théorèmes]
\label{exercises-exercise-results-fall-2021}
Énoncer avec précision mais brièvement un résultat non trivial vu en cours
se rapportant à chacun des points suivants (s'il y en a plusieurs, n'en choisir
qu'un seul).
\begin{enumerate}
\item les anneaux artiniens,
\item la platitude et les idéaux premiers,
\item les longueurs des $A/\mathfrak m^n$ lorsque $(A, \mathfrak m)$ est un anneau local noethérien,
\item la formule de dimension pour les anneaux noethériens universellement caténaires,
\item la complétion d'un anneau local noethérien, et
\item la dualité de Matlis pour les anneaux locaux artiniens.
\end{enumerate}
\end{exercise}

\begin{exercise}[Unités]
\label{exercises-exercise-simple-units}
Quelle est la structure du groupe des unités de $\mathbf{Z}[x, 1/x]$
en tant que groupe abélien ? Aucune justification n'est nécessaire.
\end{exercise}

\begin{exercise}[Idéaux]
\label{exercises-exercise-ideals}
Soit $A = \mathbf{F}_2[x, y]/(x^2, xy, y^2)$ et notons
$\overline{x}$ et $\overline{y}$ les images de $x$ et $y$ dans $A$.
Énumérer les idéaux de $A$. Aucune justification n'est nécessaire.
\end{exercise}

\begin{exercise}[Tor et Ext]
\label{exercises-exercise-compute-tor-ext}
Soit $(A, \mathfrak m, \kappa)$ un anneau local noethérien.
Posons $\varphi(n) = \dim_\kappa \mathfrak m^n/\mathfrak m^{n + 1}$.
\begin{enumerate}
\item Montrer que $\text{Tor}_1^A(A/\mathfrak m^n, \kappa)$
est de dimension $\varphi(n)$ comme espace vectoriel sur $\kappa$.
\item Montrer que $\text{Ext}^1_A(A/\mathfrak m^n, \kappa)$
est de dimension $\varphi(n)$ comme espace vectoriel sur $\kappa$.
\end{enumerate}
\end{exercise}

\begin{exercise}[Deux vecteurs]
\label{exercises-exercise-two-vectors}
Soit $A = \mathbf{Z}[a_1, a_2, a_3, b_1, b_2, b3]$. Posons
$a = (a_1, a_2, a_3)$ et $b = (b_1, b_2, b_3)$ dans $A^{\oplus 3}$.
Considérons l'ensemble
$$
Z = \{\mathfrak p \in \Spec(A) \mid a, b
\text{ ont des images linéairement dépendantes dans }
\kappa(\mathfrak p)^{\oplus 3}\}
$$
\begin{enumerate}
\item Prouver que $Z$ est une partie fermée de $\Spec(A)$.
\item Quelle est la dimension $\dim(Z)$ de $Z$ ?
\item Que deviendrait $\dim(Z)$ si l'on remplaçait $\mathbf{Z}$ par un corps ?
\end{enumerate}
\end{exercise}

\begin{exercise}[Injectifs]
\label{exercises-exercise-injective-artinian-local}
Soit $(A, \mathfrak m, \kappa)$ un anneau local artinien.
Supposons que $A$ soit injectif comme $A$-module. Montrer que
$\Hom_A(\kappa, A)$ est de dimension $1$ comme
espace vectoriel sur $\kappa$.
\end{exercise}



\section{Géométrie algébrique, examen final, printemps 2022}
\label{exercises-section-final-exam-spring-2022}

\noindent
Les questions suivantes ont été posées lors de l'examen final d'un cours de géométrie algébrique,
donné au printemps 2022 à l'Université Columbia.

\begin{exercise}[Définitions]
\label{exercises-exercise-definitions-spring-2022}
Donner des définitions succinctes des notions en italique.
\begin{enumerate}
\item un {\it schéma},
\item un {\it module quasi-cohérent} sur un schéma $X$,
\item un morphisme {\it plat} de schémas $X \to Y$,
\item un morphisme {\it fini} de schémas $X \to Y$,
\item un {\it schéma en groupes} $G$ au-dessus d'un schéma de base $S$,
\item une {\it famille de variétés} au-dessus d'un schéma de base $S$,
\item le {\it degré} d'un point fermé $x$ d'une variété $X$
sur le corps $k$,
\item la {\it hauteur logarithmique} usuelle d'un point
$p = (a_0 : \ldots : a_n)$ de $\mathbf{P}^n(\mathbf{Q})$, et
\item un {\it corps $C_i$}.
\end{enumerate}
\end{exercise}

\begin{exercise}[Théorèmes]
\label{exercises-exercise-results-spring-2022}
Énoncer avec précision mais brièvement un résultat non trivial vu en cours
se rapportant à chacun des points suivants (s'il y en a plusieurs, n'en choisir
qu'un seul).
\begin{enumerate}
\item les morphismes d'un schéma $X$ vers le schéma affine $\Spec(A)$,
\item la cohomologie d'un module quasi-cohérent $\mathcal{F}$ sur
un schéma affine $X$,
\item le groupe de Picard de $\mathbf{P}^1_k$, où $k$ est un corps,
\item les dimensions des fibres d'un morphisme propre et plat
$X \to S$ lorsque $S$ est noethérien,
\item les modules $\mathbf{G}_m$-équivariants sur un schéma $S$, et
\item le théorème de Bézout sur les intersections (on pourra se limiter à
un cas particulier).
\end{enumerate}
\end{exercise}

\begin{exercise}[Hypersurfaces cubiques]
\label{exercises-exercise-cubics}
Soit $F \in \mathbf{C}[T_0, \ldots, T_n]$ homogène de degré $3$.
Étant donnés $3$ vecteurs $x, y, z \in \mathbf{C}^{n + 1}$, considérons la condition
$$
(*)\quad
F(\lambda x + \mu y + \nu z) = 0 \text{ dans } \mathbf{C}[\lambda, \mu, \nu]
$$
\begin{enumerate}
\item Quelle est la dimension de l'espace de tous les choix de $x, y, z$ ?
\item Combien d'équations portant sur les coordonnées de $x$, $y$ et $z$
la condition (*) représente-t-elle ?
\item Quelle est la dimension attendue de l'espace de tous les triplets
$x, y, z$ vérifiant (*) ?
\item Quelle est la dimension de l'espace de tous les triplets tels que
$x, y, z$ soient linéairement dépendants ?
\item En déduire que, sur une hypersurface de degré $3$ dans $\mathbf{P}^n$,
on s'attend à trouver un sous-espace linéaire de dimension $2$ pourvu que
$n \geq a$, où il vous appartient de déterminer $a$.
\end{enumerate}
\end{exercise}

\begin{exercise}[Hauteurs]
\label{exercises-exercise-heights}
Soit $K$ un corps. Soit $h_n : \mathbf{P}^n(K) \to \mathbf{R}$, $n \geq 0$,
une famille de fonctions satisfaisant les $2$ axiomes étudiés
en cours. Soit $X$ une variété projective
sur $K$. Soit $\mathcal{L}$ un $\mathcal{O}_X$-module inversible
et rappelons que nous lui avons associé en cours une
fonction hauteur $h_\mathcal{L} : X(K) \to \mathbf{R}$.
Soit $\alpha : X \to X$ un automorphisme de $X$ au-dessus de $K$.
\begin{enumerate}
\item Montrer que l'application $P \mapsto h_\mathcal{L}(\alpha(P))$
diffère de la fonction $h_{\alpha^*\mathcal{L}}$
par une fonction bornée. (Indication : rappeler que s'il existe un morphisme
$\varphi : X \to \mathbf{P}^n$ tel que
$\mathcal{L} = \varphi^*\mathcal{O}_{\mathbf{P}^n}(1)$, alors, par construction,
$h_\mathcal{L}(P) = h_n(\varphi(P))$ et jouer avec cette identité.
Dans le cas général, écrire $\mathcal{L}$ comme différence de deux modules de ce type.)
\item Supposons que $h_\mathcal{L}(P) - h_\mathcal{L}(\alpha(P))$
ne soit pas bornée sur $X(K)$. Montrer que
$h_\mathcal{N}$, où
$\mathcal{N} = \mathcal{L} \otimes \alpha^*\mathcal{L}^{\otimes -1}$,
n'est pas bornée sur $X(K)$.
\item Supposons que $X$ soit une courbe elliptique et que
$\mathcal{L}$ soit un module inversible ample et symétrique sur $X$
tel que $h_\mathcal{L}$ ne soit pas bornée sur $X(K)$. Montrer
qu'il existe un module inversible $\mathcal{N}$ de degré $0$
tel que $h_\mathcal{N}$ ne soit pas bornée.
(Indications : rappeler que $X$ est une variété abélienne de dimension $1$.
Ainsi, $h_\mathcal{L}$ est quadratique à une constante près, d'après les résultats
vus en cours. Choisir un point convenable $P_0 \in X(K)$.
Soit $\alpha : X \to X$ la translation par $P_0$. Considérer
$P \mapsto h_\mathcal{L}(P) - h_\mathcal{L}(P + P_0)$.
Appliquer les résultats démontrés ci-dessus.)
\end{enumerate}
\end{exercise}

\begin{exercise}[Monomorphismes]
\label{exercises-exercise-monomorphisms}
Soit $f : X \to Y$ un monomorphisme dans la catégorie des schémas :
pour toute paire de morphismes $a, b : T \to X$ de schémas, si
$f \circ a = f \circ b$, alors $a = b$. Montrer que $f$ est injectif
sur les points. Votre argument montre-t-il autre chose ?
\end{exercise}

\begin{exercise}[Points fixes]
\label{exercises-exercise-fix-points}
Soit $k$ un corps algébriquement clos.
\begin{enumerate}
\item Si $G = \mathbf{G}_{m, k}$, montrer que, si $G$ agit sur une
variété projective $X$ sur $k$, alors l'action admet un point fixe, c'est-à-dire
montrer qu'il existe un point $x \in X(k)$ tel que $a(g, x) = x$ pour
tout $g \in G(k)$.
\item Même question pour $G = (\mathbf{G}_{m, k})^n$, produit de
$n \geq 1$ exemplaires du groupe multiplicatif.
\item Donner un exemple d'une action d'un schéma en groupes connexe $G$
sur une variété projective lisse $X$ qui n'admet aucun point fixe.
\end{enumerate}
\end{exercise}



\section{Géométrie algébrique, examen final, printemps 2025}
\label{exercises-section-final-exam-spring-2025}

\noindent
Les questions suivantes ont été posées lors de l'examen final d'un cours de géométrie algébrique,
donné au printemps 2025 à l'Université Columbia.

\begin{exercise}[Définitions]
\label{exercises-exercise-definitions-spring-2025}
Donner des définitions succinctes des notions en italique.
\begin{enumerate}
\item un {\it schéma} $X$,
\item un {\it module quasi-cohérent} sur un schéma $X$,
\item le {\it groupe de Picard} d'un schéma $X$,
\item étant donnés un morphisme $f : X \to Y$ de schémas et un
$\mathcal{O}_X$-module $\mathcal{F}$, son {\it image directe} $f_*\mathcal{F}$,
\item une {\it immersion fermée} de schémas, et
\item une {\it variété} sur un corps donné $k$.
\end{enumerate}
\end{exercise}

\begin{exercise}[Théorèmes]
\label{exercises-exercise-results-spring-2025}
Énoncer avec précision et concision un résultat non trivial vu en cours
se rapportant à chacun des points suivants (s'il y en a plusieurs, n'en choisir qu'un seul).
\begin{enumerate}
\item l'espace topologique d'un schéma $X$,
\item un critère de représentabilité pour les foncteurs sur la catégorie des schémas,
\item le foncteur de Picard d'une courbe projective lisse $X$ sur un
corps algébriquement clos $k$,
\item la torsion dans le groupe de Picard d'une courbe projective lisse $X$ sur un
corps algébriquement clos $k$,
\item un théorème sur le groupe de Picard d'un produit $X \times Y \times Z$
de variétés projectives lisses $X$, $Y$, $Z$ sur un corps algébriquement
clos $k$.
\end{enumerate}
\end{exercise}

\begin{exercise}
\label{exercises-exercise-affine-line-infinite-nr-points}
Soit $k$ un corps. Expliquer pourquoi $\Spec(k[x])$ a une infinité de points.
\end{exercise}

\begin{exercise}
\label{exercises-exercise-hypersurface-variety}
Soit $k$ un corps et soient $x_1, \ldots, x_n$ des indéterminées.
Soit $f \in k[x_1, \ldots, x_n]$ un polynôme non constant.
Soit $X = \Spec(k[x_1, \ldots, x_n]/(f))$. Quelle propriété
$f$ doit-il posséder pour que $X$ soit une variété sur $k$ ?
\end{exercise}

\begin{exercise}
\label{exercises-exercise-find-singular-point}
Trouver un point singulier (il n'est pas demandé de les trouver tous)
du spectre de $\mathbf{C}[x, y]/(x^5 - 5xy^4 + 20y - 16)$
considéré comme une variété sur $\mathbf{C}$.
\end{exercise}

\begin{exercise}
\label{exercises-exercise-kernel-restriction-pic-curve}
Soit $X$ une courbe projective lisse sur un corps algébriquement clos
$k$. Soit $x \in X(k)$ un point fermé de $X$. Soit
$U = X \setminus \{x\}$. Que peut-on dire du noyau de
l'homomorphisme de restriction $\Pic(X) \to \Pic(U)$ ?
\end{exercise}

\begin{exercise}
\label{exercises-exercise-sections-degree-2g-minus-4}
Soit $X$ une courbe projective lisse sur un corps algébriquement clos
$k$ de genre $g \geq 100$. Quelles valeurs
$h^0(\mathcal{L}) = \dim_k H^0(X, \mathcal{L})$ peut-il prendre lorsque
$\mathcal{L}$ est un $\mathcal{O}_X$-module inversible de degré $2g - 4$
sur $X$ ? Répondre de manière aussi complète que possible.
\end{exercise}














% Bon, commençons ici.
%

\chapter{Guide bibliographique}


\phantomsection
\label{guide-section-phantom}


\section{Courts articles d'introduction}
\label{guide-section-short-introductions}


\begin{itemize}
\item Barbara Fantechi: \emph{Stacks for Everybody} \cite{fantechi_stacks}
\item Dan Edidin: \emph{What is a stack?} \cite{edidin_whatis}
\item Dan Edidin: \emph{Notes on the construction of the moduli space of
curves} \cite{edidin_notes}
\item Angelo Vistoli: \emph{Intersection theory on algebraic
stacks and on their moduli spaces}, et surtout l'appendice.
\cite{vistoli_intersection}
\end{itemize}


\section{Références classiques}
\label{guide-section-classics}

\begin{itemize}
\item Mumford: \emph{Picard groups of moduli problems}
\cite{mumford_picard}
\begin{quote}
Mumford n'emploie jamais ici le terme « champ », mais la
notion est implicite dans l'article ; il calcule le groupe de Picard du champ
de modules des courbes elliptiques.
\end{quote}
\item Deligne, Mumford: \emph{The irreducibility of the space of curves of
given genus} \cite{DM}
\begin{quote}
Cet article influent introduit les « champs algébriques » au sens de ce que
l'on appelle aujourd'hui universellement des champs de Deligne--Mumford (champs à
diagonale représentable qui admettent des présentations étales par des schémas).
Il contient de nombreux résultats fondateurs \emph{sans démonstration}. L'article
utilise les champs pour donner deux démonstrations de l'irréductibilité de l'espace de modules des courbes de genre
$g$.
\end{quote}
\item Artin: \emph{Versal deformations and algebraic stacks}
\cite{ArtinVersal}
\begin{quote}
Cet article introduit les « champs algébriques » qui généralisent les champs de
Deligne--Mumford et que l'on appelle aujourd'hui couramment \emph{champs d'Artin} :
ce sont les champs à diagonale représentable qui admettent des présentations
lisses par des schémas. L'article donne un critère fondé sur la théorie des déformations,
connu sous le nom de critère d'Artin, qui permet de démontrer qu'un champ de
modules donné est un champ d'Artin sans exhiber explicitement de présentation.
\end{quote}
\end{itemize}




\section{Livres et notes en ligne}
\label{guide-section-books}


\begin{itemize}
\item Laumon, Moret-Bailly: \emph{Champs Alg\'ebriques} \cite{LM-B}
\begin{quote}
Ce livre constitue actuellement la référence la plus exhaustive sur les champs
et contient de nombreux résultats fondateurs. Il suppose que le lecteur connaît
les espaces algébriques et renvoie fréquemment au livre de Knutson \cite{Kn}. Le
chapitre 12 contient une erreur concernant la fonctorialité du site lisse-étale
d'un champ algébrique. Cette erreur ne prête plus à conséquence, car Martin
Olsson l'a corrigée (voir \cite{olsson_sheaves}) et les résultats des chapitres
suivants sont justes, moyennant éventuellement de légères modifications.
\end{quote}
\item The Stacks Project Authors: \emph{Stacks Project}
\cite{stacks-project}.
\begin{quote}
Vous êtes en train de le lire !
\end{quote}
\item Anton Geraschenko:
\emph{Lecture notes for Martin Olsson's class on stacks} \cite{olsson_stacks}
\begin{quote}
Ce cours développe systématiquement la théorie des espaces
algébriques avant d'introduire les champs algébriques (définis pour la première
fois au cours 27 !). Outre les propriétés fondamentales, il traite de
l'équivalence entre être de Deligne--Mumford et avoir une diagonale non ramifiée,
du site lisse-étale d'un champ d'Artin, de la théorie des faisceaux quasi-cohérents,
du théorème de Keel--Mori, de la descente cohomologique et des gerbes (ainsi que
de leur relation avec le groupe de Brauer). Il propose également des exercices.
\end{quote}
\item
Behrend, Conrad, Edidin, Fantechi, Fulton, G\"ottsche et Kresch :
\emph{Algebraic stacks}, notes en ligne pour un livre en cours de rédaction
\cite{stacks_book}
\begin{quote}
Ce livre vise à offrir une introduction accessible aux champs, sans exiger de
prérequis approfondis, en privilégiant les exemples et les applications.
Contrairement à \cite{LM-B}, il ne suppose pas que le lecteur maîtrise la théorie
des espaces algébriques. Les champs de Deligne--Mumford y sont plutôt introduits,
les espaces algébriques constituant un cas particulier ; l'un des objectifs est de
développer une théorie suffisante pour démontrer les assertions de \cite{DM}. La
théorie générale des champs d'Artin doit être développée dans la seconde partie.
Seule une partie du livre est actuellement disponible sur le site de Kresch.
\end{quote}
\item Olsson, Martin: \emph{Algebraic spaces and stacks}, \cite{olsson_book}
\begin{quote}
Introduction vivement recommandée à la théorie des espaces algébriques et des
champs algébriques, accessible à un lecteur qui maîtrise le livre
de Hartshorne sur la géométrie algébrique.
\end{quote}
\end{itemize}


\section{Références connexes sur les fondements de la théorie des champs}
\label{guide-section-related}


\begin{itemize}
\item Vistoli:
\emph{Notes on Grothendieck topologies, fibered categories and descent theory}
\cite{vistoli_fga}
\begin{quote}
Contient des résultats utiles sur les catégories fibrées, les champs et la
théorie de la descente pour la topologie fpqc, ainsi que des démonstrations rigoureuses.
\end{quote}
\item Knutson: \emph{Algebraic Spaces} \cite{Kn}
\begin{quote}
Ce livre, issu de sa thèse de doctorat dirigée par Michael Artin,
contient les fondements de la théorie des espaces algébriques. Le livre
\cite{LM-B} y renvoie fréquemment. Voir aussi les articles d'Artin sur les
espaces algébriques : \cite{Artin-Algebraic-Approximation},
\cite{ArtinI}, \cite{Artin-Implicit-Function},
\cite{ArtinII}, \cite{Artin-Construction-Techniques},
\cite{Artin-Algebraic-Spaces}, \cite{Artin-Theorem-Representability} et
\cite{ArtinVersal}
\end{quote}
\item Grothendieck et al., \emph{Th\'eorie des Topos et Cohomologie \'Etale des
Sch\'emas I, II, III}, également connu sous le nom de SGA 4 \cite{SGA4}
\begin{quote}
Le volume 1 contient de nombreux résultats généraux sur les univers, les sites
et les catégories fibrées. Le mot « champ » (terme français correspondant à
« stack ») apparaît dans l'exposé XVIII de Deligne.
\end{quote}
\item Jean Giraud: \emph{Cohomologie non ab\'elienne} \cite{giraud}
\begin{quote}
Le livre étudie les catégories fibrées, les champs, les torseurs et les gerbes
sur des sites généraux, mais non les champs algébriques. Par exemple, si $G$ est
un faisceau de groupes abéliens sur $X$, de même que $H^1(X, G)$ s'identifie aux
$G$-torseurs, $H^2(X, G)$ s'identifie à un ensemble convenablement défini de
$G$-gerbes. Lorsque $G$ n'est pas abélien, $H^2(X, G)$ est défini comme
l'ensemble des $G$-gerbes.
\end{quote}
\item Kelly and Street: \emph{Review of the elements of 2-categories}
\cite{kelly-street}
\begin{quote}
La catégorie des champs forme une 2-catégorie, certes d'un type simple où les
2-morphismes sont inversibles. Il s'agit d'une référence sur les 2-catégories
générales. Je ne l'ai jamais utilisée et ne puis donc en apprécier l'utilité.
Notons également que \cite{stacks-project} contient quelques notions de base sur les 2-catégories.
\end{quote}
\end{itemize}




\section{Articles de recherche}
\label{guide-section-papers}

\noindent
On trouvera ci-dessous une liste d'articles de recherche contenant des résultats
fondamentaux sur les champs et les espaces algébriques. Les résumés visent seulement
à signaler les résultats qui contribuent à la théorie des champs ; dans bien des cas,
ces résultats sont secondaires par rapport aux objectifs principaux de l'article.
Nous répartissons les articles en catégories, certains relevant de plusieurs d'entre elles.


\subsection{Théorie des déformations et champs algébriques}
\label{guide-subsection-deformation-theory}

\noindent
Les trois premiers articles d'Artin ne traitent pas des champs, mais ils
contiennent des résultats puissants, les deux premiers étant essentiels à
\cite{ArtinVersal}.
\begin{itemize}
\item Artin: \emph{Algebraic approximation of structures over
complete local rings} \cite{Artin-Algebraic-Approximation}
\begin{quote}
On démontre que, sous des hypothèses peu restrictives, toute déformation formelle
effective peut être approchée : si $F: (\Sch/S) \to (\textit{Sets})$
est un foncteur contravariant
localement de présentation finie, si $S$ est de type fini sur un corps ou sur un
anneau de valuation discrète excellent, si $s \in S$, et si
$\hat{\xi} \in F(\hat{\mathcal{O}}_{S, s})$ est une déformation formelle
effective, alors, pour tout $n > 0$, il existe un voisinage étale
$(S', s') \to (S, s)$ induisant un isomorphisme sur les corps résiduels, et
un élément $\xi' \in F(S')$, tels que $\xi'$ et
$\hat{\xi}$ coïncident à l'ordre $n$ (c'est-à-dire aient la même restriction dans
$F(\mathcal{O}_{S, s} / \mathfrak m^n)$).
\end{quote}
\item
Artin: \emph{Algebraization of formal moduli I} \cite{ArtinI}
\begin{quote}
On démontre que, sous des hypothèses peu restrictives, toute déformation formelle
verselle effective est algébrisable. Soit
$F: (\Sch/S) \to (\textit{Sets})$ un foncteur contravariant
localement de présentation finie, où $S$ est de type fini sur un corps ou sur un
anneau de valuation discrète excellent ; soit $s \in S$ un point localement
fermé, soit $\hat A$ une $\mathcal{O}_S$-algèbre locale noethérienne complète,
de corps résiduel $k'$, extension finie de $k(s)$,
et soit $\hat{\xi} \in F(\hat A)$ une déformation formelle verselle effective d'un
élément $\xi_0 \in F(k')$. Il existe alors un schéma $X$ de type fini sur $S$,
un point fermé $x \in X$ de corps résiduel $k(x) = k'$, et un élément
$\xi \in F(X)$, ainsi qu'un isomorphisme
$\hat{\mathcal{O}}_{X, x} \cong \hat{A}$
qui identifie les restrictions de $\xi$ et de $\hat{\xi}$ dans chaque
$F(\hat A / \mathfrak m^n)$. L'algébrisation
est unique si $\hat{\xi}$ est une déformation universelle. On en donne des
applications à la représentabilité des schémas de Hilbert et de Picard.
\end{quote}
\item Artin: \emph{Algebraization of formal moduli. II}
\cite{ArtinII}
\begin{quote}
En termes imprécis, il est démontré que, si l'on peut contracter formellement un
sous-ensemble fermé $Y' \subset X'$ au voisinage de $Y'$, alors il existe un
morphisme global $X' \to X$ qui contracte $Y$, où $X$ est un espace
algébrique.
\end{quote}
\item
Artin: \emph{Versal deformations and algebraic stacks} \cite{ArtinVersal}
\begin{quote}
Cet article décisif s'appuie sur les travaux d'Artin dans
\cite{Artin-Algebraic-Approximation} et \cite{ArtinI}. Il
introduit le critère d'Artin, qui permet d'établir l'algébricité d'un
champ en vérifiant des propriétés relevant de la théorie des déformations. Plus
précisément (mais sans prétendre à une précision parfaite), Artin construit une
présentation d'un champ $\mathcal{X}$ qui commute aux limites inductives,
au voisinage d'un point $x \in \mathcal{X}(k)$, de la manière suivante : si le
champ $\mathcal{X}$
satisfait au critère de Schlessinger (\cite{Sch}), il existe une déformation
formelle verselle
$\hat{\xi} \in \lim \mathcal{X}(\hat A / \mathfrak m^n)$ de $x$.
Si l'on suppose que les déformations formelles sont effectives (c'est-à-dire que
$\mathcal{X}(\hat{A}) \to \lim \mathcal{X}(\hat A / \mathfrak m^n)$
est bijective), on obtient une déformation formelle verselle effective
$\xi \in \mathcal{X}(\hat A)$. Les résultats de
\cite{ArtinI} fournissent alors un schéma $U$ de type fini et un
élément $\xi_U: U \to \mathcal{X}$ qui est formellement versel en un point
$u \in U$ situé au-dessus de $x$. Si l'on suppose en outre que $\mathcal{X}$
admet une théorie des déformations et des obstructions
qui satisfait certaines conditions (compatibilité avec la localisation étale
et avec la complétion, ainsi qu'une condition de constructibilité), Artin montre
à la section 4 que la versalité formelle est une condition ouverte ; quitte à
restreindre $U$, le morphisme $U \to \mathcal{X}$ est donc lisse.
Artin démontre également que tout champ admettant une présentation fppf par
un schéma admet une présentation lisse par un schéma ; on peut donc, en
particulier, former des champs quotients par des schémas en groupes plats,
séparés et de présentation finie.
\end{quote}
\item Conrad, de Jong: \emph{Approximation of Versal Deformations}
\cite{conrad-dejong}
\begin{quote}
Cet article propose une approche du théorème d'algébrisation d'Artin fondée sur
le puissant théorème de Popescu : si $A$ est un anneau noethérien et $B$ une
$A$-algèbre noethérienne, l'homomorphisme $A \to B$ est régulier si et seulement
si $B$ est une limite inductive d'$A$-algèbres lisses. On voit aisément que le
théorème de Popescu entraîne le théorème d'approximation d'Artin sur tout schéma
excellent (l'hypothèse d'excellence implique que, pour un anneau local $A$,
l'homomorphisme $A^{\text{h}} \to \hat A$ de l'hensélisé vers le complété est
régulier). L'article utilise le théorème de Popescu pour donner une
généralisation « aux groupoïdes » du théorème principal de \cite{ArtinI},
valable sur tout schéma de base excellent et en tout point $s \in S$.
En particulier, les résultats de \cite{ArtinVersal} valent sur une base
excellente arbitraire. Les auteurs étudient l'unicité locale pour la topologie
étale de l'algébrisation, ainsi que la question de savoir si le groupe des
automorphismes de l'objet agit naturellement sur l'hensélisé de
l'algébrisation.
\end{quote}
\item Jason Starr: \emph{Artin's axioms, composition, and moduli spaces}
\cite{starr_artin}
\begin{quote}
L'article établit que les axiomes d'Artin pour l'algébrisation sont compatibles
avec la composition des 1-morphismes.
\end{quote}
\item Martin Olsson: \emph{Deformation theory of representable
morphism of algebraic stacks} \cite{olsson_deformation}
\begin{quote}
Cet article étend aux morphismes représentables de champs algébriques, en termes
du complexe cotangent, les résultats usuels de la théorie des déformations des
morphismes de schémas. Ces résultats ne peuvent être considérés comme des
conséquences de la théorie générale d'Illusie, car le complexe cotangent d'un
morphisme représentable $X \to \mathcal{X}$ n'est pas
défini au moyen du complexe cotangent d'un morphisme de topos annelés (le
site lisse-étale n'étant pas fonctoriel).
\end{quote}
\end{itemize}

\subsection{Espaces de modules grossiers}
\label{guide-subsection-coarse-moduli-spaces}

\noindent
Articles consacrés aux espaces de modules grossiers.
\begin{itemize}
\item Keel, Mori: \emph{Quotients in Groupoids} \cite{K-M}
\begin{quote}
Il semble que l'on ait longtemps tenu pour « bien connu » que les champs de
Deligne--Mumford séparés admettaient des espaces de modules grossiers. On trouve
ici une démonstration rigoureuse, quoique concise, du théorème suivant : si
$\mathcal{X}$ est un champ d'Artin localement de
type fini sur un schéma de base noethérien et si le morphisme d'inertie
$I_\mathcal{X} \to \mathcal{X}$ est fini, alors il existe un espace de
modules grossier $\phi : \mathcal{X} \to Y$
tel que $\phi$ soit séparé et que $Y$ soit un espace algébrique localement de
type fini sur $S$. La finitude de l'inertie est précisément la bonne
condition : il existe un espace de modules grossier
$\phi : \mathcal{X} \to Y$ tel que $\phi$
soit séparé si et seulement si l'inertie est finie.
\end{quote}
\item Conrad: \emph{The Keel-Mori Theorem via Stacks} \cite{conrad}
\begin{quote}
L'article de Keel et Mori \cite{K-M} est rédigé dans le langage des groupoïdes,
que certains trouvent difficile à maîtriser. Brian Conrad en donne une version
formulée dans le langage des champs, dont la démonstration est fort transparente
bien qu'elle emploie un formalisme élaboré. Conrad supprime également
l'hypothèse noethérienne.
\end{quote}
\item Rydh: \emph{Existence of quotients by finite groups and coarse moduli
spaces} \cite{rydh_quotients}
\begin{quote}
Rydh supprime l'hypothèse, imposée dans \cite{K-M} et \cite{conrad}, selon
laquelle $\mathcal{X}$
doit être de présentation finie sur une base.
\end{quote}
\item
Abramovich, Olsson, Vistoli: \emph{Tame stacks in positive characteristic}
\cite{tame}
\begin{quote}
Les auteurs définissent un \emph{champ d'Artin modéré} comme un champ d'Artin à
inertie finie tel que, si $\phi : \mathcal{X} \to Y$ est son espace de modules
grossier, le foncteur $\phi_*$ soit exact
sur les faisceaux quasi-cohérents. Ils démontrent que, pour un champ d'Artin à
inertie finie, les conditions suivantes sont équivalentes : $\mathcal{X}$ est
modéré ; les stabilisateurs de $\mathcal{X}$ sont linéairement réductifs ;
localement pour la topologie étale sur l'espace de
modules grossier, $\mathcal{X}$ est le quotient d'un schéma affine par un schéma
en groupes linéairement réductif. L'espace de modules grossier d'un champ
d'Artin modéré possède des propriétés particulièrement favorables.
Ainsi, sa formation commute à tout changement de base, tandis que
la formation de l'espace de modules grossier d'un champ d'Artin à inertie finie
ne commute en général qu'aux changements de base plats.
\end{quote}
\item Alper: \emph{Good moduli spaces for Artin stacks} \cite{alper_good}
\begin{quote}
Pour les champs d'Artin généraux dont les groupes de stabilisateurs affines sont
infinis (ces champs étant nécessairement non séparés), les espaces de modules
grossiers n'existent souvent pas. L'exemple le plus simple est
$[\mathbf{A}^1 / \mathbf{G}_m]$. On dit ici qu'un morphisme quasi-compact
$\phi : \mathcal{X} \to Y$ définit un \emph{bon espace de modules} si
$\mathcal{O}_Y \to \phi_*
\mathcal{O}_\mathcal{X}$ est un isomorphisme et si $\phi_*$ est exact sur les
faisceaux quasi-cohérents.
Cette notion généralise celle de champ d'Artin modéré de \cite{tame} et
englobe la théorie géométrique des invariants de Mumford : si $G$ est un groupe
réductif agissant linéairement sur $X \subset \mathbf{P}^n$, alors le morphisme
du champ quotient du lieu semi-stable vers le quotient GIT,
$[X^{ss}/G] \to X//G$, définit un bon espace de modules. La notion de bon
espace de modules possède de nombreuses propriétés géométriques remarquables :
(1) $\phi$ est surjectif, universellement fermé et universellement
submersif ; (2) $\phi$ identifie les points de $Y$ aux classes d'équivalence
par adhérence des points de $\mathcal{X}$ ; (3) $\phi$ est universel parmi les
morphismes vers des espaces algébriques ; (4)
les bons espaces de modules sont stables par tout changement de base ; et (5)
un fibré vectoriel sur un champ d'Artin descend au bon espace de modules si et
seulement si les représentations correspondantes sont triviales aux points
fermés.
\end{quote}
\end{itemize}


\subsection{Théorie de l'intersection}
\label{guide-subsection-intersection-theory}

\noindent
Articles consacrés à la théorie de l'intersection sur les champs algébriques.
\begin{itemize}
\item
Vistoli: \emph{Intersection theory on algebraic stacks and on their moduli
spaces} \cite{vistoli_intersection}
\begin{quote}
Cet article établit les fondements d'une théorie de l'intersection à
coefficients rationnels pour les champs de Deligne--Mumford. Si $\mathcal{X}$
est un champ de Deligne--Mumford séparé, le groupe de Chow
$\CH_*(\mathcal{X})$ à coefficients rationnels est
défini comme le quotient du groupe abélien libre engendré par les sous-champs
fermés intègres de dimension $k$ par l'équivalence rationnelle.
On dispose d'une image réciproque plate, d'une image directe propre
et d'un homomorphisme de Gysin généralisé pour les immersions locales
régulières. Si $\phi : \mathcal{X} \to Y$ est un espace de modules
(c'est-à-dire un morphisme propre bijectif sur les
points géométriques), il existe une image directe induite
$\CH_*(\mathcal{X}) \to \CH_k(Y)$
qui est un isomorphisme.
\end{quote}
\item Edidin, Graham: \emph{Equivariant Intersection Theory}
\cite{edidin-graham}
\begin{quote}
Le but de cet article est de développer une théorie de l'intersection à
coefficients entiers pour le champ quotient $[X/G]$ associé à l'action d'un
groupe algébrique $G$ sur un espace algébrique $X$, autrement dit une théorie
de l'intersection $G$-équivariante de $X$. Définir les groupes de Chow
équivariants à l'aide des seuls cycles invariants ne donne pas une théorie
possédant de bonnes propriétés. Les auteurs généralisent plutôt la définition de
Totaro dans le cas de $BG$ et s'appuient sur le fait que, si
$V \to X$ est un fibré vectoriel, alors
$\CH_i(X) \cong \CH_i(V)$ canoniquement. Ils définissent
$\CH_i^G(X)$ comme suit :
posons $\dim(X) = n$ et $\dim(G) = g$. Pour chaque $i$, choisissons une
représentation de dimension $l$ du groupe $G$ sur un espace $V$, où $G$ agit librement dans
un ouvert $U \subset V$ dont le complémentaire est de codimension
$d > n - i$. Alors $X_G = [X \times U / G]$ est un
espace algébrique (on peut même choisir les données de sorte que ce soit un
schéma). On pose
$\CH_i^G(X) = \CH_{i + l - g}(X_G)$. Pour le champ quotient, on définit
$\CH_i( [X/G]) = \CH_{i + g}^G(X) = \CH_{i + l}(X_G)$.
En particulier, $\CH_i([X/G]) = 0$ pour $i > \dim [X/G] = n - g$,
mais ce groupe peut être non nul pour $i \ll 0$. Par exemple,
$\CH_i(B \mathbf{G}_m) = \mathbf{Z}$ pour $i \le 0$.
Les auteurs établissent que ces groupes de Chow équivariants possèdent les mêmes
propriétés fonctorielles que les groupes de Chow ordinaires. Ils montrent en
outre qu'un isomorphisme
$[X / G] \cong [Y / H]$ entraîne
$\CH_i([X/G]) = \CH_i([Y/H])$ ; la définition
ne dépend donc pas de la présentation du champ comme champ quotient.
\end{quote}
\item
Kresch: \emph{Cycle Groups for Artin Stacks} \cite{kresch_cycle}
\begin{quote}
Kresch définit les groupes de Chow des champs d'Artin arbitraires ; dans le cas
d'un champ quotient, sa définition coïncide avec celle d'Edidin et Graham dans
\cite{edidin-graham}. Pour les champs algébriques dont les groupes de
stabilisateurs sont affines, la théorie possède les propriétés usuelles.
\end{quote}
\item Behrend et Fantechi: \emph{The intrinsic normal cone}
\cite{behrend-fantechi}
\begin{quote}
Généralisant une construction due à Li et Tian, Behrend et Fantechi construisent
une classe fondamentale virtuelle pour tout champ de Deligne--Mumford.
\end{quote}
\end{itemize}


\subsection{Champs quotients}
\label{guide-subsection-quotient-stacks}

\noindent
Les champs quotients\footnote{Dans la littérature,
\emph{champ quotient} désigne souvent un champ de la
forme $[X/G]$, où $X$ est un espace algébrique et $G$ un sous-schéma en groupes
de $\text{GL}_n$, plutôt qu'un schéma en groupes plat quelconque.}
forment une sous-classe très importante des champs d'Artin, qui comprend presque
tous les champs de modules étudiés par les géomètres algébristes. La géométrie
du champ quotient $[X/G]$ est la géométrie $G$-équivariante de $X$. Il est souvent
plus facile d'établir des propriétés pour les champs quotients, et certains
résultats ne sont connus que pour eux. Les articles suivants abordent ces questions :
quand un champ algébrique est-il un champ quotient global ? Un champ algébrique
est-il « localement » un champ quotient ?
\begin{itemize}
\item Laumon, Moret-Bailly: \cite[chapitre 6]{LM-B}
\begin{quote}
Le chapitre 6 contient plusieurs résultats sur la structure locale et globale
des champs algébriques. Il y est démontré qu'un champ algébrique $\mathcal{X}$ sur $S$
est un
champ quotient $[Y/G]$, où $Y$ est un espace algébrique (resp. un schéma, resp. un
schéma affine) et $G$ un groupe fini, si et seulement s'il existe un espace algébrique
(resp. un schéma, resp. un schéma affine) $Y'$ et un morphisme fini étale $Y'
\to \mathcal{X}$. On montre que, pour tout champ de Deligne--Mumford sur $S$ et tout
$x : \Spec(K) \to \mathcal{X}$, il existe un morphisme représentable, étale et
séparé $\phi : [X/G] \to \mathcal{X}$, où $G$ est un groupe fini
agissant sur un schéma affine sur $S$, tel
que $\Spec(K) = [X/G] \times_\mathcal{X} \Spec(K)$.
L'existence de présentations
à fibres géométriquement connexes est également étudiée en détail.
\end{quote}
\item
Edidin, Hassett, Kresch, Vistoli:
\emph{Brauer Groups and Quotient stacks} \cite{ehkv}
\begin{quote}
Ils établissent d'abord quelques faits fondamentaux (quoique peu difficiles)
concernant
les conditions sous lesquelles un champ algébrique donné (toujours supposé de type
fini sur un schéma noethérien dans cet article) est un champ quotient. Pour un champ
algébrique $\mathcal{X}$ : $\mathcal{X}$ est un champ quotient si et seulement s'il
existe un fibré vectoriel $V \to \mathcal{X}$ tel que, pour tout
point géométrique, le stabilisateur agisse fidèlement sur la fibre,
si et seulement s'il existe un fibré vectoriel $V \to \mathcal{X}$ et
un sous-champ localement fermé
$V^0 \subset V$ tel que $V^0$ soit représentable et se projette surjectivement sur
$\mathcal{X}$. Ils
établissent qu'un champ algébrique est un champ quotient s'il existe un revêtement
fini et plat par un espace algébrique. Tout champ de Deligne--Mumford lisse à
stabilisateur génériquement trivial est un champ quotient.
Ils montrent qu'une gerbe en $\mathbf{G}_m$ sur un schéma noethérien $X$
correspondant à
$\beta \in H^2(X, \mathbf{G}_m)$ est un champ quotient si et seulement si $\beta$
appartient à
l'image de l'application de Brauer $\text{Br}(X) \to \text{Br}'(X)$. Ils en déduisent
un exemple de
champ de Deligne--Mumford non séparé qui n'est pas un champ quotient.
\end{quote}
\item Totaro: \emph{The resolution property for schemes and stacks}
\cite{totaro_resolution}
\begin{quote}
Un champ possède la propriété de résolution si tout faisceau cohérent est quotient
d'un fibré vectoriel. Le premier théorème principal affirme que, si $\mathcal{X}$ est
un champ algébrique normal noethérien dont les groupes stabilisateurs aux points
fermés sont affines, alors les conditions suivantes sont équivalentes : (1)
$\mathcal{X}$ possède la propriété de résolution et
(2)
$\mathcal{X} = [Y/\text{GL}_n]$, où $Y$ est quasi-affine. Lorsque
$\mathcal{X}$ est de type fini sur
un corps, (1) et (2) équivalent aussi à : (3)
$\mathcal{X} = [\Spec(A)/G]$, où $G$ est
un schéma en groupes affine de type fini sur $k$. L'implication selon laquelle les
champs quotients possèdent la propriété de résolution a été démontrée par Thomason.
Le second théorème principal affirme que, si $\mathcal{X}$ est un champ de
Deligne--Mumford lisse sur
un corps, dont le groupe stabilisateur $I_\mathcal{X}
\to \mathcal{X}$ est fini et génériquement trivial, et dont l'espace de modules
grossier est un schéma à diagonale affine, alors
$\mathcal{X}$ possède la propriété de résolution. Un autre résultat remarquable
affirme que, si $\mathcal{X}$ est
un champ algébrique noethérien possédant la propriété de résolution, alors
$\mathcal{X}$ a
une diagonale affine si et seulement si les stabilisateurs des points fermés sont affines.
\end{quote}
\item Kresch: \emph{On the Geometry of Deligne-Mumford Stacks}
\cite{kresch_geometry}
\begin{quote}
Cet article résume des résultats généraux sur la structure des champs de
Deligne--Mumford (de type fini sur un corps) et contient plusieurs résultats
intéressants sur les champs quotients. Il y est démontré que tout champ de
Deligne--Mumford lisse, séparé et génériquement modéré, dont l'espace de modules
grossier est quasi-projectif, est un champ quotient $[Y/G]$, où $Y$ est
quasi-projectif et $G$ un groupe algébrique. Si $\mathcal{X}$ est un champ de
Deligne--Mumford dont l'espace de modules grossier
est un schéma, alors $\mathcal{X}$ est localement pour la topologie de Zariski un
champ quotient si et seulement s'il
admet un recouvrement ouvert de Zariski par des champs quotients de schémas par des
groupes finis. Si $\mathcal{X}$ est un champ de Deligne--Mumford propre sur un corps
de caractéristique 0
et d'espace de modules grossier $Y$, alors les conditions suivantes sont équivalentes :
$Y$ est projectif et $\mathcal{X}$ est un champ quotient; $Y$ est projectif et
$\mathcal{X}$ possède un faisceau générateur; enfin,
$\mathcal{X}$ admet une immersion fermée dans un champ de Deligne--Mumford lisse et
propre dont l'espace de modules grossier est projectif. Cela conduit à dire qu'un
champ de Deligne--Mumford est \emph{projectif} s'il existe une immersion fermée dans
un champ de Deligne--Mumford lisse et propre dont l'espace de modules grossier est projectif.
\end{quote}
\item Kresch, Vistoli \emph{On coverings of Deligne-Mumford stacks and
surjectivity of the Brauer map} \cite{kresch-vistoli}
\begin{quote}
Il est démontré qu'en caractéristique 0 et pour un entier $n$ fixé, les deux
assertions suivantes sont équivalentes : (1) tout champ de Deligne--Mumford lisse de
dimension $n$ est un champ quotient, et (2) le groupe de Brauer d'Azumaya coïncide
avec le groupe de Brauer cohomologique pour les schémas lisses de dimension $n$.
\end{quote}
\item Kresch: \emph{Cycle Groups for Artin Stacks} \cite{kresch_cycle}
\begin{quote}
Il est démontré que tout champ d'Artin réduit de type fini sur un corps, dont les
groupes stabilisateurs sont affines, admet une stratification par champs quotients.
\end{quote}
\item Abramovich-Vistoli:
\emph{Compactifying the space of stable maps} \cite{abramovich-vistoli}
\begin{quote}
Le lemme 2.2.3 établit que tout champ de Deligne--Mumford séparé est, localement pour
la topologie étale sur son espace de modules grossier, un champ quotient $[U/G]$, où
$U$ est affine et $G$ un groupe fini. Dans cet argument,
\cite[théorème 2.12]{olsson_homstacks} montre même que $G$ est le groupe stabilisateur.
\end{quote}
\item Abramovich, Olsson, Vistoli:
\emph{Tame stacks in positive characteristic} \cite{tame}
\begin{quote}
Cet article montre qu'un champ d'Artin modéré est, localement pour la topologie étale
sur l'espace de modules grossier, le quotient d'un schéma affine par le groupe
stabilisateur.
\end{quote}
\item Alper: \emph{On the local quotient structure of Artin stacks}
\cite{alper_quotient}
\begin{quote}
On conjecture que, pour un champ d'Artin $\mathcal{X}$ et un point fermé $x
\in \mathcal{X}$
à stabilisateur linéairement réductif, il existe un morphisme étale $[V/G_x]
\to \mathcal{X}$, où $V$ est un espace algébrique. Plusieurs éléments à l'appui de
cette conjecture sont
donnés. Un argument simple de théorie des déformations (fondé sur les idées de
\cite{tame}) montre que l'assertion est vraie dans le voisinage formel. Une
démonstration, dans le langage des champs, du théorème de la tranche étale de Luna
est présentée; elle établit que, pour les champs
$\mathcal{X} = [\Spec(A)/G]$
avec $G$ linéairement réductif, $\mathcal{X}$ est, localement pour la topologie étale
sur le quotient GIT $\Spec(A^G)$, un champ quotient par le stabilisateur.
\end{quote}
\end{itemize}

\subsection{Cohomologie}
\label{guide-subsection-cohomology}

\noindent
Articles consacrés à la cohomologie des faisceaux sur les champs algébriques.
\begin{itemize}
\item Olsson: \emph{Sheaves on Artin stacks} \cite{olsson_sheaves}
\begin{quote}
Cet article développe la théorie des faisceaux quasi-cohérents et constructibles
et démontre les propriétés cohomologiques fondamentales. Il corrige une erreur de
\cite{LM-B} concernant la fonctorialité du site lisse-étale. Le complexe cotangent
y est construit. En outre, les théorèmes suivants y sont démontrés :
le théorème fondamental de Grothendieck pour les morphismes propres, le théorème
d'existence de Grothendieck, le théorème de connexité de Zariski et un théorème de
finitude pour les images directes propres de faisceaux cohérents et constructibles.
\end{quote}
\item Behrend: \emph{Derived $l$-adic categories for algebraic stacks}
\cite{behrend_derived}
\begin{quote}
Démontre la formule des traces de Lefschetz pour les champs algébriques.
\end{quote}
\item Behrend: \emph{Cohomology of stacks} \cite{behrend_cohomology}
\begin{quote}
Définit la cohomologie de de Rham des champs différentiables et la cohomologie
singulière des champs topologiques.
\end{quote}
\item Faltings:
\emph{Finiteness of coherent cohomology for proper fppf stacks}
\cite{faltings_finiteness}
\begin{quote}
Démontre la cohérence des images directes de faisceaux cohérents par des morphismes propres.
\end{quote}
\item Abramovich, Corti, Vistoli:
\emph{Twisted bundles and admissible covers} \cite{acv}
\begin{quote}
L'appendice contient le théorème de changement de base propre pour la cohomologie
étale des champs de Deligne--Mumford modérés.
\end{quote}
\end{itemize}


\subsection{Existence de revêtements finis par des schémas}
\label{guide-subsection-finite-covers}

\noindent
L'existence de revêtements finis de champs de Deligne--Mumford par des schémas
est un résultat important. Dans la théorie de l'intersection sur les champs de
Deligne--Mumford, c'est un ingrédient essentiel de la définition de l'image directe
propre pour les morphismes non représentables. Plusieurs
résultats sur $\overline{\mathcal{M}}_g$ reposent sur l'existence d'un revêtement
fini par un schéma \emph{lisse}, démontrée par Looijenga. Le premier résultat dans
cette direction est peut-être \cite[théorème 6.1]{seshadri_quotients},
qui traite le cadre équivariant.
\begin{itemize}
\item Vistoli: \emph{Intersection theory on algebraic stacks and on their
moduli spaces} \cite{vistoli_intersection}
\begin{quote}
Si $\mathcal{X}$ est un champ de Deligne--Mumford muni d'un espace de modules
(c'est-à-dire d'un morphisme propre
bijectif sur les points géométriques), alors il existe un morphisme fini
$X \to \mathcal{X}$ provenant d'un schéma $X$.
\end{quote}
\item Laumon, Moret-Bailly: \cite[chapitre 16]{LM-B}
\begin{quote}
Comme application du théorème principal de Zariski, le théorème 16.6 établit que,
si $\mathcal{X}$ est un champ de Deligne--Mumford de type fini sur un schéma
noethérien, alors
il existe un morphisme fini, surjectif et génériquement étale $Z \to
\mathcal{X}$,
où $Z$ est un schéma. Le corollaire 16.6.2 montre également que tout espace
algébrique normal noethérien est isomorphe à l'espace algébrique quotient $X'/G$
d'un schéma normal $X'$ par l'action d'un groupe fini $G$.
\end{quote}
\item Edidin, Hassett, Kresch, Vistoli:
\emph{Brauer Groups and Quotient stacks} \cite{ehkv}
\begin{quote}
Le théorème 2.7 affirme que,
si $\mathcal{X}$ est un champ algébrique de type fini sur un
schéma de base noethérien $S$, alors la diagonale
$\mathcal{X} \to \mathcal{X} \times_S \mathcal{X}$ est
quasi-finie si et seulement s'il existe un morphisme fini surjectif
$X \to \mathcal{X}$ provenant d'un schéma $X$.
\end{quote}
\item Kresch, Vistoli: \emph{On coverings of Deligne-Mumford stacks and
surjectivity of the Brauer map} \cite{kresch-vistoli}
\begin{quote}
Il est démontré que tout champ de Deligne--Mumford lisse et séparé, de type fini
sur un corps et dont l'espace de modules grossier est quasi-projectif, admet un
revêtement fini et plat par un schéma lisse quasi-projectif.
\end{quote}
\item Olsson: \emph{On proper coverings of Artin stacks} \cite{olsson_proper}
\begin{quote}
Il est démontré que, si $\mathcal{X}$ est un champ d'Artin séparé
et de type fini sur $S$, alors
il existe un morphisme propre surjectif $X \to \mathcal{X}$ provenant d'un schéma
$X$ quasi-projectif sur $S$. Comme application, Olsson démontre la cohérence et
la constructibilité des faisceaux images directes par les morphismes propres. Il
en déduit également le théorème d'existence de Grothendieck pour les champs
d'Artin propres.
\end{quote}
\item Rydh: \emph{Noetherian approximation of algebraic spaces and stacks}
\cite{rydh_approx}
\begin{quote}
Le théorème B de cet article s'énonce comme suit.
Soit $X$ un champ algébrique quasi-compact à diagonale quasi-finie et séparée
(resp.\ un champ de Deligne--Mumford quasi-compact à diagonale quasi-compacte
et séparée). Il existe alors un schéma $Z$ et un morphisme $Z \to X$ fini,
de présentation finie et surjectif, qui est plat (resp.\ étale) au-dessus d'un
sous-champ ouvert dense quasi-compact $U \subset X$.
\end{quote}
\end{itemize}


\subsection{Rigidification}
\label{guide-subsection-rigidification}

\noindent
La rigidification est un procédé qui retire un sous-groupe plat de l'inertie.
Par exemple, si $X$ est une variété projective, le morphisme du champ de Picard
vers le schéma de Picard est la rigidification par le groupe d'automorphismes
$\mathbf{G}_m$.
\begin{itemize}
\item Abramovich, Corti, Vistoli :
\emph{Twisted bundles and admissible covers} \cite{acv}
\begin{quote}
Soient $\mathcal{X}$ un champ algébrique sur $S$ et $H$ un schéma en groupes
plat, séparé et de présentation finie sur $S$. Supposons que, pour tout objet
$\xi \in \mathcal{X}(T)$, il existe un plongement
$H(T) \hookrightarrow \text{Aut}_{\mathcal{X}(T)}(\xi)$ compatible aux
changements de base, au sens où, pour toute flèche
$\phi : \xi \rightarrow \xi'$ au-dessus de $f: T \rightarrow T'$ et tout
$g \in H(T')$, on a $g \circ \phi = \phi \circ f^*g$.
Il existe alors un champ algébrique $\mathcal{X}/H$ et un morphisme
$\rho : \mathcal{X} \rightarrow \mathcal{X}/H$ qui est une gerbe fppf et tel
que, pour tout $\xi \in \mathcal{X}(T)$, le morphisme
$\text{Aut}_{\mathcal{X}(T)} (\xi)
\rightarrow \text{Aut}_{\mathcal{X}/H (T)} (\xi) $
est surjectif de noyau $H(T)$.
\end{quote}
\item Romagny : \emph{Group actions on stacks and applications}
\cite{romagny_actions}
\begin{quote}
Étudie le comportement des actions de groupes vis-à-vis des rigidifications.
\end{quote}
\item Abramovich, Graber, Vistoli :
\emph{Gromov-Witten theory for Deligne-Mumford stacks} \cite{agv}
\begin{quote}
L'appendice donne une synthèse de la rigidification au sens de \cite{acv},
ainsi que deux interprétations alternatives. Cet article contient aussi des
constructions permettant de recoller des champs algébriques le long de
sous-champs fermés et de prendre des racines de fibrés en droites.
\end{quote}
\item
Abramovich, Olsson, Vistoli : \emph{Tame stacks in positive characteristic}
(\cite{tame})
\begin{quote}
L'appendice traite la situation plus complexe où le sous-champ en groupes plat
de l'inertie $H \subset I_\mathcal{X}$ est normal, mais pas nécessairement
central.
\end{quote}
\end{itemize}

\subsection{Courbes stacky}
\label{guide-subsection-stacky-curves}

\noindent
Articles consacrés aux courbes stacky.
\begin{itemize}
\item Abramovich, Vistoli : \emph{Compactifying the space of stable maps}
\cite{abramovich-vistoli}
\begin{quote}
Cet article introduit les \emph{courbes tordues}. L'espace de modules des
applications stables de courbes stables dans un champ algébrique n'est en
général pas compact. En utilisant des applications définies sur des courbes
tordues, les auteurs construisent un champ de modules qui est propre lorsque
le but est un champ de Deligne--Mumford modéré dont l'espace de modules
grossier est projectif.
\end{quote}
\item Behrend, Noohi : \emph{Uniformization of Deligne-Mumford curves}
\cite{behrend-noohi}
\begin{quote}
Démontre un théorème d'uniformisation des courbes analytiques de
Deligne--Mumford.
\end{quote}
\end{itemize}

\subsection{Champs de Hilbert, de Quot, de Hom et de variétés branchées}
\label{guide-subsection-hilbert-quot-hom}

\noindent
Articles consacrés aux schémas de Hilbert et aux constructions analogues.
\begin{itemize}
\item Vistoli : \emph{The Hilbert stack and the theory of moduli of families}
\cite{vistoli_hilbert}
\begin{quote}
Si $\mathcal{X}$ est un champ algébrique séparé et localement de type fini sur
un espace algébrique $S$ localement noethérien et localement séparé, Vistoli
définit le champ de Hilbert $\mathcal{H}\text{ilb}(\mathcal{X} / S)$ qui
paramètre les morphismes finis et non ramifiés provenant de schémas propres.
Il est affirmé sans démonstration que
$\mathcal{H}\text{ilb}(\mathcal{X} / S)$ est un champ algébrique. On en déduit
que, pour $\mathcal{X}$ comme ci-dessus, le champ de Hom
$\mathcal{H} \text{om}_S(T, \mathcal{X})$ est un champ algébrique si $T$ est
propre et plat sur $S$.
\end{quote}
\item Olsson, Starr : \emph{Quot functors for Deligne-Mumford stacks}
\cite{olsson-starr}
\begin{quote}
Si $\mathcal{X}$ est un champ de Deligne--Mumford séparé et localement de
présentation finie sur un espace algébrique $S$, et si $\mathcal{F}$ est un
$\mathcal{O}_\mathcal{X}$-module localement de présentation finie, le foncteur
Quot $\text{Quot}(\mathcal{F} / \mathcal{X} / S)$ est représenté par un espace
algébrique séparé et localement de présentation finie sur $S$. Cet article
définit aussi les faisceaux générateurs et démontre l'existence d'un tel
faisceau pour les champs de Deligne--Mumford modérés et séparés qui sont des
champs quotients globaux d'un schéma par un groupe fini.
\end{quote}
\item Olsson : \emph{Hom-stacks and Restrictions of Scalars}
\cite{olsson_homstacks}
\begin{quote}
Supposons que $\mathcal{X}$ et $\mathcal{Y}$ soient des champs d'Artin
localement de présentation finie sur un espace algébrique $S$, à diagonale
finie, que $\mathcal{X}$ soit propre et plat sur $S$, et que, localement pour
la topologie fppf sur $S$, $\mathcal{X}$ admette un revêtement fini, plat et de
présentation finie par un espace algébrique (par exemple si $\mathcal{X}$ est
de Deligne--Mumford ou est un champ d'Artin modéré). Alors
$\Hom_S(\mathcal{X}, \mathcal{Y})$ est un champ d'Artin localement de
présentation finie sur $S$.
\end{quote}
\item Alexeev et Knutson : \emph{Complete moduli spaces of branchvarieties}
(\cite{alexeev-knutson})
\begin{quote}
Ils définissent une variété branchée de $\mathbf{P}^n$ comme un morphisme fini
$X \rightarrow \mathbf{P}^n$ provenant d'un schéma \emph{réduit} $X$. Ils
démontrent que le champ de modules des variétés branchées de polynôme de
Hilbert fixé et de degrés totaux fixés pour les composantes de dimension $i$
est un champ d'Artin propre à stabilisateur fini. Ils comparent le champ des
variétés branchées au schéma de Hilbert, au schéma de Chow et à l'espace de
modules des applications stables.
\end{quote}
\item Lieblich : \emph{Remarks on the stack of coherent algebras}
\cite{lieblich_remarks}
\begin{quote}
Cet article construit une généralisation du champ des variétés branchées
d'Alexeev et Knutson au-dessus d'un schéma $Y$, en le réalisant comme un champ
d'algèbres sur le faisceau structural de $Y$. Il donne des démonstrations de
l'existence des espaces $\text{Quot}$ et $\Hom$.
\end{quote}
\item Starr : \emph{Artin's axioms, composition, and moduli spaces}
\cite{starr_artin}
\begin{quote}
Comme application du résultat principal, l'article fournit une généralisation
commune du champ de Hilbert de Vistoli \cite{vistoli_hilbert} et du champ des
variétés branchées d'Alexeev et Knutson \cite{alexeev-knutson}. Si
$\mathcal{X}$ est un champ algébrique localement de type fini sur un schéma
excellent $S$, à diagonale finie, alors le champ $\mathcal{H}$ qui paramètre
les morphismes $g: T \rightarrow \mathcal{X}$ provenant d'un espace algébrique
propre $T$ muni d'un fibré en droites $g$-ample $L$ est un champ d'Artin
localement de type fini sur $S$.
\end{quote}
\item Lundkvist et Skjelnes :
\emph{Non-effective deformations of Grothendieck's Hilbert functor}
\cite{lundkvist-skjelnes}
\begin{quote}
Montre que le foncteur de Hilbert d'un schéma non séparé n'est pas représenté,
car il existe des déformations non effectives.
\end{quote}
\item Halpern-Leistner et Preygel :
\emph{Mapping stacks and categorical notions of properness}
\cite{HL-P}
\begin{quote}
Cet article démontre que le champ de Hom est algébrique sous des hypothèses sur
la source et le but qui sont plus générales que celles de l'article d'Olsson,
ou du moins différentes.
\end{quote}
\end{itemize}



\subsection{Champs toriques}
\label{guide-subsection-toric}

\noindent
Les champs toriques fournissent une vaste classe d'exemples et un banc d'essai
naturel pour les conjectures, grâce au dictionnaire entre la géométrie d'un
champ torique et la combinatoire de son éventail stacky, tout comme les
variétés toriques fournissent des exemples et des contre-exemples en théorie
des schémas.
\begin{itemize}
\item Borisov, Chen et Smith : \emph{The orbifold Chow ring of toric
Deligne-Mumford stacks} \cite{bcs}
\begin{quote}
Inspiré par la construction de Cox pour les variétés toriques, cet article
définit les champs toriques de Deligne--Mumford lisses comme des champs
quotients explicites associés à un objet combinatoire appelé
\emph{éventail stacky}.
\end{quote}
\item Iwanari : \emph{The category of toric stacks} \cite{iwanari_toric}
\begin{quote}
Cet article définit un \emph{triplet torique} comme un champ de
Deligne--Mumford lisse $\mathcal{X}$ muni d'une immersion ouverte
$\mathbf{G}_m \hookrightarrow \mathcal{X}$ d'image dense (de sorte que
$\mathcal{X}$ est un orbifold) et d'une action
$\mathcal{X} \times \mathbf{G}_m \rightarrow \mathcal{X}$. Il est démontré
qu'il existe une équivalence entre la 2-catégorie des triplets toriques et la
1-catégorie des éventails stacky. L'article étudie la relation entre les
triplets toriques et la définition des champs toriques de Deligne--Mumford
lisses donnée dans \cite{bcs}.
\end{quote}
\item Iwanari : \emph{Integral Chow rings for toric stacks}
\cite{iwanari_chow}
\begin{quote}
Généralise les $\Delta$-collections de Cox pour les variétés toriques aux
orbifolds toriques.
\end{quote}
\item Perroni : \emph{A note on toric Deligne-Mumford stacks}
\cite{perroni}
\begin{quote}
Généralise les $\Delta$-collections de Cox et l'article d'Iwanari
\cite{iwanari_chow} aux champs toriques de Deligne--Mumford lisses généraux.
\end{quote}
\item Fantechi, Mann et Nironi : \emph{Smooth toric DM stacks}
\cite{fmn}
\begin{quote}
Cet article définit un champ torique de Deligne--Mumford lisse comme un champ
de Deligne--Mumford lisse $\mathcal{X}$ muni de l'action d'un tore de
Deligne--Mumford $\mathcal{T}$ (c'est-à-dire d'un champ de Picard isomorphe à
$T \times BG$, où $G$ est fini) qui possède une orbite ouverte dense isomorphe
à $\mathcal{T}$. Les auteurs en donnent une « description ascendante » et
démontrent une équivalence entre les champs toriques de Deligne--Mumford
lisses et les éventails stacky.
\end{quote}
\item Geraschenko et Satriano : \emph{Toric Stacks I and II}
\cite{gs_toric1} et \cite{gs_toric2}
\begin{quote}
Ces articles définissent un champ torique comme le champ quotient d'une
variété torique par un sous-groupe de son tore. Un champ torique génériquement
stacky est défini comme un sous-champ invariant par le tore d'un champ
torique. Cette définition englobe et étend les définitions antérieures des
champs toriques. Le premier article établit un dictionnaire entre la
combinatoire des éventails stacky et la géométrie des champs
correspondants. Il donne aussi une interprétation modulaire des champs
toriques lisses, généralisant celle de \cite{perroni}. Le second article
démontre une caractérisation intrinsèque des champs toriques.
\end{quote}
\end{itemize}



\subsection{Théorème des fonctions formelles et théorème d'existence de Grothendieck}
\label{guide-subsection-theorem-formal-functions}

\noindent
Ces articles donnent des généralisations du théorème des fonctions formelles
\cite[III.4.1.5]{EGA} (parfois appelé théorème fondamental de Grothendieck
pour les morphismes propres) et du théorème d'existence de Grothendieck
\cite[III.5.1.4]{EGA}.
\begin{itemize}
\item Knutson : \emph{Algebraic spaces} \cite[chapitre V]{Kn}
\begin{quote}
Généralise ces théorèmes aux espaces algébriques.
\end{quote}
\item Abramovich--Vistoli : \emph{Compactifying the space of stable maps}
\cite[A.1.1]{abramovich-vistoli}
\begin{quote}
Généralise ces théorèmes aux champs de Deligne--Mumford modérés.
\end{quote}
\item Olsson et Starr : \emph{Quot functors for Deligne-Mumford stacks}
\cite{olsson-starr}
\begin{quote}
Généralise ces théorèmes aux champs de Deligne--Mumford séparés.
\end{quote}
\item Olsson : \emph{On proper coverings of Artin stacks}
\cite{olsson_proper}
\begin{quote}
En fournit une généralisation aux champs d'Artin propres.
\end{quote}
\item Conrad : \emph{Formal GAGA on Artin stacks} \cite{conrad_gaga}
\begin{quote}
En fournit une généralisation aux champs d'Artin propres et démontre un
théorème de GAGA formel.
\end{quote}
\item Olsson : \emph{Sheaves on Artin stacks} \cite{olsson_sheaves}
\begin{quote}
Donne une autre démonstration de la généralisation aux champs d'Artin propres.
\end{quote}
\end{itemize}


\subsection{Actions de groupes sur les champs}
\label{guide-subsection-group-actions}

\noindent
Les actions de groupes sur les champs algébriques apparaissent naturellement.
Par exemple, le groupe symétrique $S_n$ agit sur
$\overline{\mathcal{M}}_{g, n}$ et, lorsqu'un groupe $G$ agit sur un schéma
$X$, le normalisateur de $G$ dans $\text{Aut}(X)$ agit sur $[X/G]$. En outre,
les actions de tores sur les champs apparaissent souvent en théorie de
Gromov--Witten.
\begin{itemize}
\item Romagny : \emph{Group actions on stacks and applications}
\cite{romagny_actions}
\begin{quote}
Cet article précise ce que signifie l'action d'un groupe sur un champ
algébrique et démontre l'existence de points fixes ainsi que celle de
quotients pour les actions de schémas en groupes sur les champs algébriques.
Voir aussi la note antérieure de Romagny \cite{romagny_notes}.
\end{quote}
\end{itemize}

\subsection{Prise de racines de fibrés en droites}
\label{guide-subsection-root-stacks}

\noindent
Cette construction utile a été découverte indépendamment par Cadman et par
Abramovich, Graber et Vistoli. Étant donné un schéma $X$ muni d'un diviseur de
Cartier effectif $D$, le champ des racines $r$-ièmes est un champ d'Artin
ramifié au-dessus de $X$ le long de $D$, avec un stabilisateur $\mu_r$ au-dessus
de $D$, et il est schématique hors de $D$.
\begin{itemize}
\item Charles Cadman
\emph{Using Stacks to Impose Tangency Conditions on Curves}
\cite{cadman}
\item Abramovich, Graber, Vistoli : \emph{Gromov-Witten theory for
Deligne-Mumford stacks} \cite{agv}
\end{itemize}

\subsection{Autres articles}
\label{guide-subsection-other}

\noindent
Un assortiment d'autres articles.
\begin{itemize}
\item Lieblich : \emph{Moduli of twisted sheaves} \cite{lieblich_twisted}
\begin{quote}
Cet article contient une synthèse sur les gerbes et les faisceaux tordus.
Si $\mathcal{X} \rightarrow X$ est une $\mu_n$-gerbe, où $X$ est une surface
relative projective à fibres géométriques lisses et connexes, il est démontré
que le champ des faisceaux $\mathcal{X}$-tordus semi-stables est un champ
d'Artin localement de présentation finie sur $S$. L'article développe aussi
la théorie des points associés et de la pureté des faisceaux sur les champs
d'Artin.
\end{quote}
\item Lieblich, Osserman :
\emph{Functorial reconstruction theorem for stacks}
\cite{lieblich-osserman}
\begin{quote}
Démontre plusieurs résultats surprenants et intéressants sur les conditions
permettant de reconstruire un champ algébrique à partir du foncteur qui lui
est associé.
\end{quote}
\item David Rydh :
\emph{Noetherian approximation of algebraic spaces and stacks}
\cite{rydh_approx}
\begin{quote}
Cet article montre que tout champ algébrique quasi-compact à diagonale
quasi-finie peut être obtenu par approximation par un champ noethérien. Cela
permet notamment de supprimer l'hypothèse noethérienne dans des résultats de
Chevalley, Serre, Zariski et Chow.
\end{quote}
\end{itemize}



\section{Les champs dans d'autres domaines}
\label{guide-section-stacks-areas}

\begin{itemize}
\item Behrend et Noohi : \emph{Uniformization of Deligne-Mumford curves}
\cite{behrend-noohi}
\begin{quote}
Donne une vue d'ensemble et une comparaison des champs topologiques,
analytiques et algébriques.
\end{quote}
\item Behrang Noohi : \emph{Foundations of topological stacks I} \cite{noohi}
\item David Metzler : \emph{Topological and smooth stacks} \cite{metzler}
\end{itemize}


\section{Champs supérieurs}
\label{guide-section-higher-stacks}

\begin{itemize}
\item Lurie : \emph{Higher topos theory} \cite{lurie_topos}
\item Lurie : \emph{Derived Algebraic Geometry I - V}
\cite{dag1}, \cite{dag2}, \cite{dag3}, \cite{dag4}, \cite{dag5}
\item To\"en : \emph{Higher and derived stacks: a global overview}
\cite{toen_higher}
\item To\"en et Vezzosi : \emph{Homotopical algebraic geometry I, II}
\cite{hag1}, \cite{hag2}
\end{itemize}






% OK, commencer ici.
%

\chapter{Desiderata}


\phantomsection
\label{desirables-section-phantom}



\section{Introduction}
\label{desirables-section-introduction}

\noindent
Il s'agit essentiellement d'une liste de sujets que nous souhaitons ajouter au
projet Champs. Comme nous enrichissons continuellement le projet Champs, cette liste reste
toujours quelque peu en retard sur l'état actuel du projet Champs. En fait, c'était peut-être
une erreur de tenter d'énumérer les sujets à ajouter, car il semble
impossible de la tenir à jour.

\medskip\noindent
Dernière mise à jour : jeudi 31 août 2017.


\section{Conventions}
\label{desirables-section-conventions}

\noindent
Nous devrions consacrer un chapitre à une courte liste des conventions employées dans ce document.
Ce chapitre existe déjà, voir
Conventions, section \ref{conventions-section-comments},
mais on pourrait y ajouter bien davantage. Il serait particulièrement utile de repérer
les conventions ``cachées'' et les hypothèses tacites, puis de les y consigner.


\section{Sites et topos}
\label{desirables-section-sites}

\noindent
Nous disposons d'un chapitre consacré aux sites et aux faisceaux, voir
Sites, section \ref{sites-section-introduction}.
Nous disposons d'un chapitre consacré aux sites annelés (et aux topos) et aux modules sur ceux-ci, voir
Modules sur les sites, section \ref{sites-modules-section-introduction}.
Nous disposons d'un chapitre consacré à la cohomologie dans ce cadre, voir
Cohomologie sur les sites, section \ref{sites-cohomology-section-introduction}.
Mais on pourrait y ajouter bien davantage, surtout dans le chapitre consacré à la cohomologie.


\section{Champs}
\label{desirables-section-stacks}

\noindent
Nous disposons d'un chapitre consacré aux champs (abstraits), voir
Champs, section \ref{stacks-section-introduction}.
Il serait souhaitable :
\begin{enumerate}
\item d'améliorer la présentation de la ``champification'',
\item de donner des exemples de champification,
\item de donner davantage d'exemples en général,
\item d'améliorer la présentation des gerbes.
\end{enumerate}
Voici un exemple de résultat qui n'a pas encore été ajouté : étant donné un faisceau de groupes
abéliens $\mathcal{F}$
sur $\mathcal{C}$, l'ensemble des classes d'équivalence de gerbes de lien
$\mathcal{F}$ est en bijection avec $H^2(\mathcal{C}, \mathcal{F})$.


\section{Méthodes simpliciales}
\label{desirables-section-simplicial}

\noindent
Nous disposons d'un chapitre consacré aux méthodes simpliciales, voir
Objets simpliciaux, section \ref{simplicial-section-introduction}.
Il doit être revu et amélioré. La présentation des
rapports entre l'homotopie simpliciale (également appelée
homotopie combinatoire) et les complexes de Kan devrait être améliorée.
Il existe un chapitre consacré aux espaces simpliciaux, voir
Espaces simpliciaux, section \ref{spaces-simplicial-section-introduction}.
Ce chapitre traite brièvement des
espaces topologiques simpliciaux, des sites simpliciaux et des topos simpliciaux.
Nous pouvons développer davantage la ``géométrie algébrique simpliciale'' afin d'étudier
les schémas simpliciaux (ou les espaces algébriques simpliciaux, ou
les champs algébriques simpliciaux) et les questions géométriques, leur cohomologie,
etc.


\section{Cohomologie des schémas}
\label{desirables-section-schemes-cohomology}

\noindent
Il existe déjà un chapitre consacré à la cohomologie des faisceaux quasi-cohérents, voir
Cohomologie des schémas, section \ref{coherent-section-introduction}.
Nous disposons d'un chapitre qui traite de la catégorie dérivée des
faisceaux quasi-cohérents sur un schéma, voir
Catégories dérivées des schémas, section \ref{perfect-section-introduction}.
Nous disposons d'un chapitre consacré à la dualité pour les schémas noethériens
et à la dualité relative pour les morphismes de schémas, voir
Dualité pour les schémas, section \ref{duality-section-introduction}.
Nous avons également des chapitres sur la cohomologie \'etale des schémas et sur la
cohomologie cristalline des schémas. Mais la majeure partie du contenu de ces
chapitres est très élémentaire, et on pourrait/devrait y ajouter beaucoup.


\section{Théorie des déformations \`a la Schlessinger}
\label{desirables-section-deformation-schlessinger}

\noindent
Nous disposons d'un chapitre consacré à cette matière, voir
Théorie formelle des déformations, section \ref{formal-defos-section-introduction}.
Nous disposons d'un chapitre qui présente des exemples de la théorie générale, voir
Problèmes de déformation, section \ref{examples-defos-section-introduction}.
Nous disposons d'un chapitre, voir
Théorie des déformations, section \ref{defos-section-introduction},
qui traite des déformations d'anneaux (et de modules),
des déformations d'espaces annelés (et de faisceaux de modules),
des déformations de topos annelés (et de faisceaux de modules).
Dans ce chapitre, nous employons le complexe cotangent naïf
pour décrire les obstructions, les déformations du premier ordre et
les automorphismes infinitésimaux. Cette matière a trouvé quelques
applications à l'algébricité des champs de modules dans des chapitres ultérieurs.
Il existe également un chapitre consacré au complexe cotangent complet, voir
Complexe cotangent, section \ref{cotangent-section-introduction}.


\section{Définition des champs algébriques}
\label{desirables-section-definition-algebraic-stacks}

\noindent
Un champ algébrique est un champ en groupoïdes sur la catégorie des schémas
munie de la topologie fppf, dont la diagonale est représentable par des
espaces algébriques et qui est la cible d'un morphisme lisse surjectif provenant d'un schéma.
Voir Champs algébriques, section \ref{algebraic-section-algebraic-stacks}.
Un ``champ de Deligne--Mumford'' est un champ algébrique pour lequel il existe
un schéma et un morphisme \'etale surjectif de ce schéma vers le champ,
comme dans l'article \cite{DM} de Deligne et Mumford, voir
Champs algébriques, Définition \ref{algebraic-definition-deligne-mumford}.
Nous réserverons l'expression ``champ d'Artin'' aux champs considérés dans les articles
d'Artin, voir \cite{ArtinI}, \cite{ArtinII} et \cite{ArtinVersal}.
Une définition possible consiste à demander qu'un champ d'Artin soit un champ algébrique
$\mathcal{X}$ au-dessus d'un schéma localement noethérien $S$ tel que
$\mathcal{X} \to S$ soit
localement de type fini\footnote{Ce sont précisément les champs algébriques
sur $S$ qui satisfont les axiomes [-1], [0], [1], [2], [3], [4], [5]
des axiomes d'Artin, section \ref{artin-section-axioms}.}.


\section{Exemples de schémas, d'espaces algébriques et de champs algébriques}
\label{desirables-section-examples-stacks}

\noindent
Le projet Champs contient actuellement deux chapitres consacrés aux
champs de modules et à leurs propriétés, voir
Champs de modules, section \ref{moduli-section-introduction} et
Modules des courbes, section \ref{moduli-curves-section-introduction}.
À terme, nous comptons en ajouter d'autres, par exemple :
\begin{enumerate}
\item $\mathcal{A}_g$,
c'est-à-dire les schémas abéliens principalement polarisés de genre $g$,
\item $\mathcal{A}_1 = \mathcal{M}_{1, 1}$, c'est-à-dire
les courbes projectives lisses de genre $1$ à $1$ point marqué,
\item $\mathcal{M}_{g, n}$, c'est-à-dire les courbes projectives lisses de genre $g$
munies de $n$ points marqués deux à deux distincts,
\item $\overline{\mathcal{M}}_{g, n}$, c'est-à-dire les
courbes projectives nodales stables à $n$ points marqués, de genre $g$,
\item $\SheafHom_S(\mathcal{X}, \mathcal{Y})$, le champ de modules des morphismes
(sous des conditions appropriées sur les champs $\mathcal{X}$, $\mathcal{Y}$
et le schéma de base $S$),
\item $\textit{Bun}_G(X) = \SheafHom_S(X, BG)$, le champ des torseurs sous $G$
du programme de Langlands géométrique (sous des conditions appropriées sur le schéma
$X$, le schéma en groupes $G$ et le schéma de base $S$),
\item $\Picardstack_{\mathcal{X}/S}$, c'est-à-dire le champ de Picard associé
à un champ algébrique sur un schéma (ou espace) de base.
\end{enumerate}
Plus généralement, le projet Champs manque quelque peu
d'exemples significatifs du point de vue géométrique.


\section{Propriétés des champs algébriques}
\label{desirables-section-stacks-properties}

\noindent
C'est peut-être l'un des projets les plus faciles à entreprendre, puisque l'essentiel de la
théorie élémentaire est maintenant en place. Bien entendu, il s'agit en réalité de propriétés
de morphismes de champs. Nous pouvons définir des singularités (à des facteurs lisses près),
etc. Démontrer qu'un champ normal connexe est irréductible, etc.


\section{Site lisse-\'etale d'un champ algébrique}
\label{desirables-section-lisse-etale}

\noindent
Il a été introduit dans
Cohomologie des champs, section \ref{stacks-cohomology-section-lisse-etale}.
Un exemple montrant qu'il n'est pas fonctoriel par rapport aux $1$-morphismes
de champs algébriques est étudié dans
Exemples, section \ref{examples-section-lisse-etale-not-functorial}.
On pourrait bien sûr en dire beaucoup plus, mais il s'avère
très utile de démontrer les résultats à l'aide du ``grand'' site \'etale
dans toute la mesure du possible.



\section{Ce que vous avez toujours voulu savoir sans jamais oser le demander}
\label{desirables-section-stacks-fun-lemmas}

\noindent
Il y aura beaucoup de lemmes utiles que l'on emploie sans cesse,
mais qui ne sont pas vraiment énoncés explicitement dans la littérature,
ou pour lesquels il n'est pas aisé de trouver des références. Une boîte à outils.

\medskip\noindent
Exemple : étant donnés deux groupoïdes en schémas $R\Rightarrow U$ et
$R' \Rightarrow U'$, comment exprimer un $1$-morphisme
$[U/R] \to [U'/R']$ uniquement en termes de groupoïdes en schémas ?



\section{Faisceaux quasi-cohérents sur les champs}
\label{desirables-section-quasi-coherent}

\noindent
Ils sont définis et étudiés dans le chapitre
Cohomologie des champs, section \ref{stacks-cohomology-section-introduction}.
Les catégories dérivées de modules sont étudiées dans le chapitre
Catégories dérivées des champs, section \ref{stacks-perfect-section-introduction}.
On pourrait enrichir considérablement ces chapitres.



\section{Plat et lisse}
\label{desirables-section-flat-smooth}

\noindent
Le théorème d'Artin, selon lequel l'existence d'un morphisme plat surjectif dont la source est un schéma
peut se substituer à la condition d'existence d'un morphisme lisse surjectif, est désormais établi dans
Critères de représentabilité, Théorème \ref{criteria-theorem-bootstrap}.


\section{Théorème de représentabilité d'Artin}
\label{desirables-section-representability}

\noindent
Ce sujet est traité dans le chapitre
Axiomes d'Artin, section \ref{artin-section-introduction}.
Nous disposons également d'une application, voir
Quot, Théorème \ref{quot-theorem-coherent-algebraic}.
Il devrait y avoir bien davantage d'applications, et le chapitre
lui-même doit également être remanié.


\section{Les champs DM admettent des morphismes finis surjectifs provenant de schémas}
\label{desirables-section-dm-finite-cover}

\noindent
Nous disposons déjà du résultat correspondant pour les espaces algébriques, voir
Limites d'espaces, section \ref{spaces-limits-section-finite-cover}.
Il manque le résultat pour les champs DM et quasi-DM.


\section{Article de Martin Olsson sur la propreté}
\label{desirables-section-proper-parametrization}

\noindent
Cet article démontre que deux notions de propreté coïncident. La première partie de ce résultat
est désormais disponible sous la forme du lemme de Chow pour les champs algébriques, voir
Compléments sur les morphismes de champs, Théorème
\ref{stacks-more-morphisms-theorem-chow-finite-type}.
Il en résulte que, dans certains cas, il suffit de recourir aux anneaux de valuation discrète
pour vérifier le critère valuatif de propreté
des champs algébriques; voir
Compléments sur les morphismes de champs, section
\ref{stacks-more-morphisms-section-Noetherian-valuative-criterion}.


\section{Image directe de faisceaux cohérents par un morphisme propre}
\label{desirables-section-proper-pushforward}

\noindent
Nous pouvons maintenant commencer à étudier cette question, puisque nous disposons du lemme de Chow pour
les champs algébriques; voir la section précédente.


\section{Keel et Mori}
\label{desirables-section-keel-mori}

\noindent
Voir \cite{K-M}. Leur résultat a été ajouté dans
Compléments sur les morphismes de champs, section
\ref{stacks-more-morphisms-section-Keel-Mori}.


\section{Ajouter d'autres sujets ici}
\label{desirables-section-add-more}

\noindent
En réalité, non : nous n'aurions jamais dû commencer cette liste dans le
projet Champs lui-même ! Il existe ailleurs une liste de tâches
qui est bien plus facile à mettre à jour.







% OK, start here.
%

\chapter{Conventions de codage}



\phantomsection
\label{coding-section-phantom}


\section{Liste des conventions de style}
\label{coding-section-style}

\noindent
Ces conventions évolueront avec le temps, mais le fait d'en fixer dès
maintenant quelques-unes devrait favoriser un style LaTeX uniforme.
Nous appellerons « code\footnote{Tout est la faute de Knuth. Voir
\cite{Knuth}.} » le contenu des fichiers sources.

\begin{enumerate}
\item Dans tous les fichiers tex, limitez chaque ligne à 80 caractères.
\item N'indentez pas les fichiers tex. Servez-vous plutôt de la coloration
syntaxique de votre éditeur pour repérer les environnements, entre autres.
\item Utilisez
\begin{verbatim}
\medskip\noindent
\end{verbatim}
pour commencer un nouveau paragraphe, et
\begin{verbatim}
\noindent
\end{verbatim}
pour commencer un nouveau paragraphe juste après un environnement.
\item Évitez, si possible, de répartir le code des formules mathématiques sur
plusieurs lignes. Si le code complet, signes dollar englobants compris, ne
tient pas sur une ligne, placez le premier signe dollar au tout début de la
ligne suivante. S'il ne tient toujours pas, coupez le code à un endroit
mathématiquement naturel.
\item Codez comme suit les formules mathématiques hors texte :
\begin{verbatim}
$$
...
...
$$
\end{verbatim}
Autrement dit, ouvrez par deux signes dollar seuls sur une ligne, et fermez
de la même manière.
\item {\it N'utilisez aucune macro}. En voici la raison : le fichier tex est
ainsi plus facile à lire, et l'on peut commencer à en modifier n'importe
quelle partie sans devoir apprendre d'innombrables macros.
Cela ne rend pas non plus la saisie plus difficile ni plus longue.
L'inconvénient est bien sûr qu'un même objet mathématique peut être composé
différemment en TeX à différents endroits du texte, mais cela devrait se
repérer facilement.
\item Les environnements d'énoncés que nous utilisons sont :
« theorem », « proposition », « lemma » (style plain),
« definition », « example », « exercise », « situation »
(style definition), « remark », « remarks » (style remark).
Il existe bien sûr aussi un environnement « proof ».
\item Codez comme suit un environnement « foo » :
\begin{verbatim}
\begin{foo}
...
...
\end{foo}
\end{verbatim}
de même que pour les formules mathématiques hors texte.
\item Au lieu d'un environnement « corollary », employez simplement
l'environnement « lemma », puisque le résultat servira vraisemblablement de
toute façon à démontrer le théorème plus important qui suit.
\item Chaque lemme, proposition ou théorème doit être immédiatement suivi de
sa démonstration. Évitez les démonstrations imbriquées.
\item Les fichiers preamble.tex, chapters.tex et fdl.tex sont des fichiers tex
particuliers. Mis à part ceux-ci, chaque fichier tex a la structure suivante :
\begin{verbatim}
\input{preamble}
\begin{document}
\title{Title}
\maketitle
\tableofcontents
...
...
\begin{multicols}{2}[\section{Autres chapitres}]
\noindent
Préliminaires
\begin{enumerate}
\item \hyperref[introduction-section-phantom]{Introduction}
\item \hyperref[conventions-section-phantom]{Conventions}
\item \hyperref[sets-section-phantom]{Théorie des ensembles}
\item \hyperref[categories-section-phantom]{Catégories}
\item \hyperref[topology-section-phantom]{Topologie}
\item \hyperref[sheaves-section-phantom]{Faisceaux sur les espaces}
\item \hyperref[sites-section-phantom]{Sites et faisceaux}
\item \hyperref[stacks-section-phantom]{Champs}
\item \hyperref[fields-section-phantom]{Corps}
\item \hyperref[algebra-section-phantom]{Algèbre commutative}
\item \hyperref[brauer-section-phantom]{Groupes de Brauer}
\item \hyperref[homology-section-phantom]{Algèbre homologique}
\item \hyperref[derived-section-phantom]{Catégories dérivées}
\item \hyperref[simplicial-section-phantom]{Méthodes simpliciales}
\item \hyperref[more-algebra-section-phantom]{Compléments d'algèbre}
\item \hyperref[smoothing-section-phantom]{Lissage des morphismes d'anneaux}
\item \hyperref[modules-section-phantom]{Faisceaux de modules}
\item \hyperref[sites-modules-section-phantom]{Modules sur les sites}
\item \hyperref[injectives-section-phantom]{Objets injectifs}
\item \hyperref[cohomology-section-phantom]{Cohomologie des faisceaux}
\item \hyperref[sites-cohomology-section-phantom]{Cohomologie sur les sites}
\item \hyperref[dga-section-phantom]{Algèbre différentielle graduée}
\item \hyperref[dpa-section-phantom]{Algèbre à puissances divisées}
\item \hyperref[sdga-section-phantom]{Faisceaux différentiels gradués}
\item \hyperref[hypercovering-section-phantom]{Hyperrecouvrements}
\end{enumerate}
Schémas
\begin{enumerate}
\setcounter{enumi}{25}
\item \hyperref[schemes-section-phantom]{Schémas}
\item \hyperref[constructions-section-phantom]{Constructions de schémas}
\item \hyperref[properties-section-phantom]{Propriétés des schémas}
\item \hyperref[morphisms-section-phantom]{Morphismes de schémas}
\item \hyperref[coherent-section-phantom]{Cohomologie des schémas}
\item \hyperref[divisors-section-phantom]{Diviseurs}
\item \hyperref[limits-section-phantom]{Limites de schémas}
\item \hyperref[varieties-section-phantom]{Variétés}
\item \hyperref[topologies-section-phantom]{Topologies sur les schémas}
\item \hyperref[descent-section-phantom]{Descente}
\item \hyperref[perfect-section-phantom]{Catégories dérivées de schémas}
\item \hyperref[more-morphisms-section-phantom]{Compléments sur les morphismes}
\item \hyperref[flat-section-phantom]{Compléments sur la platitude}
\item \hyperref[groupoids-section-phantom]{Schémas en groupoïdes}
\item \hyperref[more-groupoids-section-phantom]{Compléments sur les schémas en groupoïdes}
\item \hyperref[etale-section-phantom]{Morphismes étales de schémas}
\end{enumerate}
Thèmes de théorie des schémas
\begin{enumerate}
\setcounter{enumi}{41}
\item \hyperref[chow-section-phantom]{Homologie de Chow}
\item \hyperref[intersection-section-phantom]{Théorie de l'intersection}
\item \hyperref[pic-section-phantom]{Schémas de Picard des courbes}
\item \hyperref[weil-section-phantom]{Théories cohomologiques de Weil}
\item \hyperref[adequate-section-phantom]{Modules adéquats}
\item \hyperref[dualizing-section-phantom]{Complexes dualisants}
\item \hyperref[duality-section-phantom]{Dualité pour les schémas}
\item \hyperref[discriminant-section-phantom]{Discriminants et différentes}
\item \hyperref[derham-section-phantom]{Cohomologie de de Rham}
\item \hyperref[local-cohomology-section-phantom]{Cohomologie locale}
\item \hyperref[algebraization-section-phantom]{Géométrie algébrique et formelle}
\item \hyperref[curves-section-phantom]{Courbes algébriques}
\item \hyperref[resolve-section-phantom]{Résolution des surfaces}
\item \hyperref[models-section-phantom]{Réduction semi-stable}
\item \hyperref[functors-section-phantom]{Foncteurs et morphismes}
\item \hyperref[equiv-section-phantom]{Catégories dérivées de variétés}
\item \hyperref[pione-section-phantom]{Groupes fondamentaux de schémas}
\item \hyperref[etale-cohomology-section-phantom]{Cohomologie étale}
\item \hyperref[crystalline-section-phantom]{Cohomologie cristalline}
\item \hyperref[proetale-section-phantom]{Cohomologie pro-étale}
\item \hyperref[relative-cycles-section-phantom]{Cycles relatifs}
\item \hyperref[more-etale-section-phantom]{Compléments de cohomologie étale}
\item \hyperref[trace-section-phantom]{La formule des traces}
\end{enumerate}
Espaces algébriques
\begin{enumerate}
\setcounter{enumi}{64}
\item \hyperref[spaces-section-phantom]{Espaces algébriques}
\item \hyperref[spaces-properties-section-phantom]{Propriétés des espaces algébriques}
\item \hyperref[spaces-morphisms-section-phantom]{Morphismes d'espaces algébriques}
\item \hyperref[decent-spaces-section-phantom]{Espaces algébriques décents}
\item \hyperref[spaces-cohomology-section-phantom]{Cohomologie des espaces algébriques}
\item \hyperref[spaces-limits-section-phantom]{Limites d'espaces algébriques}
\item \hyperref[spaces-divisors-section-phantom]{Diviseurs sur les espaces algébriques}
\item \hyperref[spaces-over-fields-section-phantom]{Espaces algébriques sur des corps}
\item \hyperref[spaces-topologies-section-phantom]{Topologies sur les espaces algébriques}
\item \hyperref[spaces-descent-section-phantom]{Descente et espaces algébriques}
\item \hyperref[spaces-perfect-section-phantom]{Catégories dérivées d'espaces}
\item \hyperref[spaces-more-morphisms-section-phantom]{Compléments sur les morphismes d'espaces}
\item \hyperref[spaces-flat-section-phantom]{Platitude sur les espaces algébriques}
\item \hyperref[spaces-groupoids-section-phantom]{Groupoïdes dans les espaces algébriques}
\item \hyperref[spaces-more-groupoids-section-phantom]{Compléments sur les groupoïdes dans les espaces}
\item \hyperref[bootstrap-section-phantom]{Amorçage}
\item \hyperref[spaces-pushouts-section-phantom]{Sommes amalgamées d'espaces algébriques}
\end{enumerate}
Thèmes de géométrie
\begin{enumerate}
\setcounter{enumi}{81}
\item \hyperref[spaces-chow-section-phantom]{Groupes de Chow des espaces}
\item \hyperref[groupoids-quotients-section-phantom]{Quotients de groupoïdes}
\item \hyperref[spaces-more-cohomology-section-phantom]{Compléments sur la cohomologie des espaces}
\item \hyperref[spaces-simplicial-section-phantom]{Espaces simpliciaux}
\item \hyperref[spaces-duality-section-phantom]{Dualité pour les espaces}
\item \hyperref[formal-spaces-section-phantom]{Espaces algébriques formels}
\item \hyperref[restricted-section-phantom]{Algébrisation des espaces formels}
\item \hyperref[spaces-resolve-section-phantom]{Résolution des surfaces revisitée}
\end{enumerate}
Théorie des déformations
\begin{enumerate}
\setcounter{enumi}{89}
\item \hyperref[formal-defos-section-phantom]{Théorie formelle des déformations}
\item \hyperref[defos-section-phantom]{Théorie des déformations}
\item \hyperref[cotangent-section-phantom]{Le complexe cotangent}
\item \hyperref[examples-defos-section-phantom]{Problèmes de déformation}
\end{enumerate}
Champs algébriques
\begin{enumerate}
\setcounter{enumi}{93}
\item \hyperref[algebraic-section-phantom]{Champs algébriques}
\item \hyperref[examples-stacks-section-phantom]{Exemples de champs}
\item \hyperref[stacks-sheaves-section-phantom]{Faisceaux sur les champs algébriques}
\item \hyperref[criteria-section-phantom]{Critères de représentabilité}
\item \hyperref[artin-section-phantom]{Axiomes d'Artin}
\item \hyperref[quot-section-phantom]{Espaces de Quot et de Hilbert}
\item \hyperref[stacks-properties-section-phantom]{Propriétés des champs algébriques}
\item \hyperref[stacks-morphisms-section-phantom]{Morphismes de champs algébriques}
\item \hyperref[stacks-limits-section-phantom]{Limites de champs algébriques}
\item \hyperref[stacks-cohomology-section-phantom]{Cohomologie des champs algébriques}
\item \hyperref[stacks-perfect-section-phantom]{Catégories dérivées de champs}
\item \hyperref[stacks-introduction-section-phantom]{Introduction aux champs algébriques}
\item \hyperref[stacks-more-morphisms-section-phantom]{Compléments sur les morphismes de champs}
\item \hyperref[stacks-geometry-section-phantom]{La géométrie des champs}
\end{enumerate}
Thèmes de théorie des espaces de modules
\begin{enumerate}
\setcounter{enumi}{107}
\item \hyperref[moduli-section-phantom]{Champs de modules}
\item \hyperref[moduli-curves-section-phantom]{Espaces de modules de courbes}
\end{enumerate}
Divers
\begin{enumerate}
\setcounter{enumi}{109}
\item \hyperref[examples-section-phantom]{Exemples}
\item \hyperref[exercises-section-phantom]{Exercices}
\item \hyperref[guide-section-phantom]{Guide bibliographique}
\item \hyperref[desirables-section-phantom]{Desiderata}
\item \hyperref[coding-section-phantom]{Style de programmation}
\item \hyperref[obsolete-section-phantom]{Obsolète}
\item \hyperref[fdl-section-phantom]{Licence de documentation libre GNU}
\item \hyperref[index-section-phantom]{Index généré automatiquement}
\end{enumerate}
\end{multicols}

\bibliography{my}
\bibliographystyle{amsalpha}
\end{document}
\end{verbatim}
\item Essayez d'ajouter des étiquettes aux lemmes, propositions et théorèmes,
ainsi qu'aux remarques, exercices et autres environnements. Pour étiqueter un
lemme, utilisez par exemple
\begin{verbatim}
\begin{lemma}
\label{lemma-bar}
...
\end{lemma}
\end{verbatim}
et procédez de même pour tous les autres environnements. Autrement dit,
l'étiquette d'un environnement nommé « foo » commence par « foo- ».
En outre, toutes les étiquettes doivent être constituées uniquement de lettres
minuscules, de chiffres et du symbole « - ».
\item Ne renvoyez jamais au « lemme ci-dessus » (ni à la proposition
ci-dessus, etc.). Utilisez plutôt :
\begin{verbatim}
Lemma \ref{lemma-bar} above
\end{verbatim}
Ainsi, déplacer ultérieurement des lemmes ne pose pour ainsi dire aucun
problème.
\item Renvois entre fichiers. Pour renvoyer à un lemme d'étiquette
« lemma-bar » dans le fichier foo.tex, dont le titre est « Foo », utilisez le
code suivant :
\begin{verbatim}
Foo, Lemma \ref{foo-lemma-bar}
\end{verbatim}
Si cela ne fonctionne pas, consultez le fichier preamble.tex afin d'y trouver
la syntaxe correcte. Le fichier de sortie contiendra alors
« Foo, lemme $<$link$>$ » : il sera donc clair que le lien pointe vers un autre
fichier.
\item Dans les environnements de démonstration, évitez autant que possible les
renvois à des résultats ultérieurs. (On devrait pouvoir écrire un test
automatique à cette fin.)
\item Ne commencez aucune phrase par un symbole mathématique.
\item N'écrivez pas une phrase du type « Cela résulte de ce qui suit » juste
avant un lemme, une proposition ou un théorème. Chaque phrase se termine par
un point.
\item Énoncez toutes les hypothèses dans chaque lemme, proposition et théorème.
Le lecteur peut ainsi déterminer plus facilement si un lemme, une proposition
ou un théorème donné s'applique à son problème particulier.
\item Faites des démonstrations courtes : moins d'une page dans le PDF ou le
DVI. On peut toujours y parvenir en scindant la démonstration en lemmes, entre
autres.
\item Pour une propriété définie appelée foobar, utilisez
\begin{verbatim}
{\it foobar}
\end{verbatim}
dans le code à l'intérieur de l'environnement de définition.
Procédez de même si la définition figure dans le corps du document.
Le lecteur repérera ainsi plus facilement ce qui est défini.
\item Placez dans son propre environnement de définition toute définition qui
sera utilisée hors de la section où elle apparaît. Les définitions temporaires
peuvent être données dans le texte. Les constructions mathématiques constituent
un cas délicat : elles tiennent souvent lieu de définitions faisant intervenir
plusieurs lemmes liés. Une bonne solution consiste peut-être à leur consacrer
une courte section, afin que les utilisateurs puissent renvoyer à celle-ci
plutôt qu'à une définition.
\item Ne numérotez les équations que si elles font effectivement l'objet d'un
renvoi dans le texte. Nous pourrons toujours ajouter des étiquettes plus tard.
\item Dans les énoncés des lemmes, propositions et théorèmes, ainsi que dans les
démonstrations, faites des phrases courtes. Par exemple, au lieu d'écrire
« Soit $R$ un anneau et soit $M$ un $R$-module. », écrivez « Soit $R$ un
anneau. Soit $M$ un $R$-module. ». En voici la raison : les parties plus
délicates des démonstrations et des énoncés sont ainsi plus faciles à analyser.
\item Utilisez la commande
\begin{verbatim}
\section
\end{verbatim}
pour créer des sections, mais essayez d'éviter les sous-sections et les
sous-sous-sections.
\item Évitez les constructions LaTeX complexes.
\end{enumerate}










% OK, start here.
%

\chapter{Obsolète}


\phantomsection
\label{obsolete-section-phantom}





\section{Introduction}
\label{obsolete-section-introduction}

\noindent
Dans ce chapitre, nous rassemblons certains lemmes devenus «~obsolètes~»
(voir \cite{Miller}).


\section{Préliminaires}
\label{obsolete-section-preliminaries}

\begin{remark}
\label{obsolete-remark-composition-of-adjoints-isomorphic-to-identity}
Les informations qui figuraient auparavant dans cette remarque sont désormais
contenues dans la combinaison des lemmes \ref{categories-lemma-adjoint-fully-faithful}
et \ref{categories-lemma-left-adjoint-composed-fully-faithful} du chapitre
Catégories.
\end{remark}


\section{Algèbre homologique}
\label{obsolete-section-homological-algebra}

\begin{remark}
\label{obsolete-remark-weak-serre-subcategory}
Les remarques suivantes sont obsolètes, car leur contenu est désormais couvert
par les lemmes \ref{homology-lemma-biregular-ss-converges} et
\ref{homology-lemma-first-quadrant-ss} du chapitre Homologie.
Soit $\mathcal{A}$ une catégorie abélienne.
Soit $\mathcal{C} \subset \mathcal{A}$
une sous-catégorie de Serre faible (voir la définition
\ref{homology-definition-serre-subcategory} du chapitre Homologie).
Supposons que $K^{\bullet, \bullet}$ soit un bicomplexe auquel
le lemme \ref{homology-lemma-first-quadrant-ss} du chapitre Homologie
s'applique et que, pour un certain $r \geq 0$, tous les objets
${}'E_r^{p, q}$ appartiennent à $\mathcal{C}$. Alors tous les objets de cohomologie
$H^n(sK^\bullet)$ appartiennent à $\mathcal{C}$. En effet, les hypothèses entraînent
que les noyaux et les images des ${}'d_r^{p, q}$ appartiennent à $\mathcal{C}$.
On en déduit que chaque ${}'E_{r + 1}^{p, q}$ appartient à $\mathcal{C}$.
Par récurrence, chaque ${}'E_\infty^{p, q}$ appartient à $\mathcal{C}$.
Par conséquent, chaque $H^n(sK^\bullet)$ possède une filtration finie dont les
sous-quotients appartiennent à $\mathcal{C}$. Comme $\mathcal{C}$ est stable par
extensions, on conclut que $H^n(sK^\bullet)$ appartient à $\mathcal{C}$, comme annoncé.
Le même résultat vaut pour la seconde suite spectrale associée
à $K^{\bullet, \bullet}$. De même, si $(K^\bullet, F)$ est un complexe filtré auquel
le lemme \ref{homology-lemma-biregular-ss-converges} du chapitre Homologie
s'applique et si, pour un certain $r \geq 0$, tous les objets $E_r^{p, q}$
appartiennent à $\mathcal{C}$, alors chaque $H^n(K^\bullet)$ est
un objet de $\mathcal{C}$.
\end{remark}




\section{Lemmes d'algèbre obsolètes}
\label{obsolete-section-algebra}

\begin{lemma}
\label{obsolete-lemma-finite-presentation-module-independent}
Soit $M$ un $R$-module de présentation finie.
Pour toute surjection $\alpha : R^{\oplus n} \to M$, le
noyau de $\alpha$ est un $R$-module de type fini.
\end{lemma}

\begin{proof}
C'est un cas particulier du lemme \ref{algebra-lemma-extension} du chapitre Algèbre.
\end{proof}

\begin{lemma}
\label{obsolete-lemma-p-ring-map}
Soit $\varphi : R \to S$ un homomorphisme d'anneaux. Si
\begin{enumerate}
\item pour tout $x \in S$, il existe $n > 0$ tel que
$x^n$ appartienne à l'image de $\varphi$, et
\item pour tout $x \in \Ker(\varphi)$, il existe $n > 0$
tel que $x^n = 0$,
\end{enumerate}
alors $\varphi$ induit un homéomorphisme sur les spectres. Soit
$p$ un nombre premier tel que
\begin{enumerate}
\item[(a)] $S$ soit engendrée comme $R$-algèbre par des éléments $x$ pour
lesquels il existe $n > 0$ tel que $x^{p^n} \in \varphi(R)$ et
$p^nx \in \varphi(R)$, et
\item[(b)] le noyau de $\varphi$ soit engendré par des éléments nilpotents.
\end{enumerate}
Alors (1) et (2) sont satisfaites et, pour tout homomorphisme d'anneaux $R \to R'$,
l'homomorphisme $R' \to R' \otimes_R S$ satisfait aussi (a), (b), (1) et (2) ;
en particulier, il induit un homéomorphisme sur les spectres.
\end{lemma}

\begin{proof}
Cela résulte de la combinaison des lemmes
\ref{algebra-lemma-powers} et \ref{algebra-lemma-p-ring-map} du chapitre Algèbre.
\end{proof}

\noindent
Le lemme technique suivant affirme que l'on peut relever à la situation globale,
en un sens convenable, toute suite de relations donnée sur une fibre d'un
homomorphisme d'anneaux essentiellement de type fini.

\begin{lemma}
\label{obsolete-lemma-lift-elements-ideal}
Soit $R \to S$ un homomorphisme d'anneaux.
Soit $\mathfrak p \subset R$ un idéal premier.
Soit $\mathfrak q \subset S$ un idéal premier au-dessus de $\mathfrak p$.
Supposons que $S_{\mathfrak q}$ soit essentiellement de type fini sur $R_\mathfrak p$.
Supposons donnés
\begin{enumerate}
\item un entier $n \geq 0$,
\item un idéal premier $\mathfrak a \subset \kappa(\mathfrak p)[x_1, \ldots, x_n]$,
\item un homomorphisme surjectif de $\kappa(\mathfrak p)$-algèbres
$$
\psi : (\kappa(\mathfrak p)[x_1, \ldots, x_n])_{\mathfrak a}
\longrightarrow
S_{\mathfrak q}/\mathfrak p S_{\mathfrak q},
$$
et
\item des éléments $\overline{f}_1, \ldots, \overline{f}_e$ de $\Ker(\psi)$.
\end{enumerate}
Il existe alors
\begin{enumerate}
\item un entier $m \geq 0$,
\item un élément $g \in S$, $g \not\in \mathfrak q$,
\item un homomorphisme
$$
\Psi :
R[x_1, \ldots, x_n, x_{n + 1}, \ldots, x_{n + m}]
\longrightarrow
S_g,
$$
et
\item des éléments $f_1, \ldots, f_e, f_{e + 1}, \ldots, f_{e + m}$
de $\Ker(\Psi)$
\end{enumerate}
tels que
\begin{enumerate}
\item le diagramme suivant soit commutatif :
$$
\xymatrix{
R[x_1, \ldots, x_{n + m}] \ar[d]_\Psi
\ar[rr]_-{x_{n + j} \mapsto 0} & &
(\kappa(\mathfrak p)[x_1, \ldots, x_n])_{\mathfrak a} \ar[d]^\psi \\
S_g \ar[rr] & &
S_{\mathfrak q}/\mathfrak p S_{\mathfrak q}
},
$$
\item pour $i \leq e$, l'élément $f_i$ ait pour image un multiple de
$\overline{f}_i$ par une unité dans l'anneau local
$$
(\kappa(\mathfrak p)[x_1, \ldots, x_{n + m}])_{
(\mathfrak a, x_{n + 1}, \ldots, x_{n + m})},
$$
\item l'élément $f_{e + j}$ ait pour image un multiple de
$x_{n + j}$ par une unité dans le même anneau local, et
\item l'homomorphisme induit $R[x_1, \ldots, x_{n + m}]_{\mathfrak b}
\to S_{\mathfrak q}$ soit surjectif, où
$\mathfrak b = \Psi^{-1}(\mathfrak qS_g)$.
\end{enumerate}
\end{lemma}

\begin{proof}
Nous affirmons qu'il suffit de démontrer le lemme lorsque $R$
et $S$ sont locaux, d'idéaux maximaux respectifs $\mathfrak p$ et $\mathfrak q$.
Supposons en effet avoir construit
$$
\Psi' : R_{\mathfrak p}[x_1, \ldots, x_{n + m}]
\longrightarrow
S_{\mathfrak q}
$$
et $f_1', \ldots, f_{e + m}' \in R_{\mathfrak p}[x_1, \ldots, x_{n + m}]$
possédant toutes les propriétés requises. Il existe alors un élément
$f \in R$, $f \not \in \mathfrak p$, tel que chacun des
$ff_k'$ provienne d'un élément $f_k \in R[x_1, \ldots, x_{n + m}]$.
De plus, pour un $g \in S$, $g \not \in \mathfrak q$, convenable,
les éléments $\Psi'(x_i)$ sont les images d'éléments
$y_i \in S_g$. Soit $\Psi$ l'homomorphisme de $R$-algèbres défini
par la règle $\Psi(x_i) = y_i$. Puisque $\Psi(f_i)$ est nul
dans le localisé $S_{\mathfrak q}$, nous pouvons, quitte à remplacer
$g$, supposer que $\Psi(f_i) = 0$. L'affirmation est démontrée.

\medskip\noindent
Nous pouvons donc supposer $R$ et $S$ locaux,
d'idéaux maximaux $\mathfrak p$ et $\mathfrak q$.
Choisissons $y_1, \ldots, y_n \in S$ tels que
$y_i \bmod \mathfrak pS = \psi(x_i)$.
Soient $y_{n + 1}, \ldots, y_{n + m} \in S$ des éléments qui engendrent
une $R$-algèbre dont $S$ est une localisation.
Ils existent puisque $S$ est essentiellement de type fini sur $R$.
Comme $\psi$ est surjectif, on peut écrire
$y_{n + j} \bmod \mathfrak pS = \psi(h_j)$ pour
un certain $h_j \in \kappa(\mathfrak p)[x_1, \ldots, x_n]_{\mathfrak a}$.
Écrivons $h_j = g_j/d$, avec $g_j \in \kappa(\mathfrak p)[x_1, \ldots, x_n]$
et un dénominateur commun $d \in \kappa(\mathfrak p)[x_1, \ldots, x_n]$,
$d \not \in \mathfrak a$. Choisissons des relèvements
$G_j, D \in R[x_1, \ldots, x_n]$ de $g_j$ et $d$. Posons
$y_{n + j}' = D(y_1, \ldots, y_n) y_{n + j} - G_j(y_1, \ldots, y_n)$.
Par construction, $y_{n + j}' \in \mathfrak p S$.
Il est clair que $y_1, \ldots, y_n, y_{n + 1}', \ldots, y_{n + m}'$
engendrent une $R$-sous-algèbre de $S$ dont $S$ est une localisation.
On définit
$$
\Psi : R[x_1, \ldots, x_{n + m}] \to S
$$
comme l'homomorphisme qui envoie $x_i$ sur $y_i$ pour $i = 1, \ldots, n$
et $x_{n + j}$ sur $y'_{n + j}$ pour $j = 1, \ldots, m$. Les propriétés
(1) et (4) sont immédiates par construction. De plus, l'idéal
$\mathfrak b$ s'envoie sur l'idéal
$(\mathfrak a, x_{n + 1}, \ldots, x_{n + m})$
de l'anneau de polynômes $\kappa(\mathfrak p)[x_1, \ldots, x_{n + m}]$.

\medskip\noindent
Notons $J = \Ker(\Psi)$. On a une suite exacte courte
$$
0 \to J_{\mathfrak b}
\to R[x_1, \ldots, x_{n + m}]_{\mathfrak b}
\to S_{\mathfrak q}
\to 0.
$$
La surjectivité résulte du choix fait ci-dessus de
$y_1, \ldots, y_n, y_{n + 1}', \ldots, y_{n + m}'$.
Il en résulte que la suite
$$
J_{\mathfrak b}/ \mathfrak pJ_{\mathfrak b}
\to \kappa(\mathfrak p)[x_1, \ldots, x_{n + m}]_{
(\mathfrak a, x_{n + 1}, \ldots, x_{n + m})}
\to S_{\mathfrak q}/\mathfrak pS_{\mathfrak q}
\to 0
$$
est exacte. Par construction, $x_i$ s'envoie sur $\psi(x_i)$ et
$x_{n + j}$ sur zéro par le dernier homomorphisme.
On choisit alors aisément les $f_i$ comme en
(2) et (3) du lemme.
\end{proof}

\begin{lemma}
\label{obsolete-lemma-smooth-independent-presentation}
Soit $R \to S$ un homomorphisme d'anneaux de présentation finie.
Si, pour une présentation $\alpha$ de $S$ sur $R$, le
complexe cotangent naïf $\NL(\alpha)$ est quasi-isomorphe
à un $S$-module projectif de type fini placé en degré $0$, il en va
de même pour toute présentation.
\end{lemma}

\begin{proof}
Cela résulte immédiatement du lemme \ref{algebra-lemma-NL-homotopy} du chapitre Algèbre.
\end{proof}

\begin{remark}[Résolutions projectives]
\label{obsolete-remark-projective-resolution}
Soit $R$ un anneau.
Pour tout ensemble $S$, notons $F(S)$ le $R$-module libre de base $S$.
Tout $R$-module à gauche possède alors la résolution en deux étapes suivante :
$$
F(M \times M) \oplus F(R \times M) \to F(M) \to M \to 0.
$$
Le premier homomorphisme est donné par la règle
$$
[m_1, m_2] \oplus [r, m] \mapsto [m_1 + m_2] - [m_1] - [m_2] + [rm] - r[m].
$$
\end{remark}

\begin{lemma}
\label{obsolete-lemma-spec-localization-first}
Soit $S$ une partie multiplicative de $A$. Alors l'homomorphisme
$$
f: \Spec(S^{-1}A)\longrightarrow \Spec(A)
$$
induit par l'homomorphisme canonique d'anneaux
$A \to S^{-1}A$ est un homéomorphisme sur son image et
$\Im(f) = \{ \mathfrak p \in \Spec(A) : \mathfrak p\cap S = \emptyset \}$.
\end{lemma}

\begin{proof}
C'est un doublon du lemme \ref{algebra-lemma-spec-localization} du chapitre Algèbre.
\end{proof}

\begin{lemma}
\label{obsolete-lemma-finite-type-flat-over-integral-algebra}
Soit $A \to B$ un homomorphisme d'anneaux plat de type fini, où $A$ est
intègre. Alors $B$ est une $A$-algèbre de présentation finie.
\end{lemma}

\begin{proof}
C'est un cas particulier de la proposition
\ref{flat-proposition-flat-finite-type-finite-presentation-domain}
du chapitre Compléments sur la platitude.
\end{proof}

\begin{lemma}
\label{obsolete-lemma-helper-finite-type-flat-finite-presentation}
Soit $R$ un anneau intègre de corps des fractions $K$.
Soit $S = R[x_1, \ldots, x_n]$ un anneau de polynômes sur $R$.
Soit $M$ un $S$-module de type fini. Supposons que $M$ soit plat sur $R$.
Si, pour tout sous-anneau $R \subset R' \subset K$, $R \not = R'$,
le module $M \otimes_R R'$ est de présentation finie
sur $S \otimes_R R'$, alors $M$ est de présentation finie sur $S$.
\end{lemma}

\begin{proof}
Ce lemme est vrai parce que $M$ est de présentation finie même sans
l'hypothèse selon laquelle $M \otimes_R R'$ est de présentation finie pour tout
$R'$ comme dans l'énoncé. Cela résulte de la proposition
\ref{flat-proposition-flat-finite-type-finite-presentation-domain}
du chapitre Compléments sur la platitude.
À l'origine, ce lemme comportait une preuve erronée (nous remercions Ofer Gabber
d'en avoir repéré la lacune) et servait dans une autre preuve de la proposition
citée. Pour rétablir ce lemme, il nous faut un argument correct lorsque $R$ est
un anneau intègre local normal, n'utilisant que des résultats des chapitres
d'algèbre commutative ; si vous connaissez un tel argument, veuillez écrire à
\href{mailto:stacks.project@gmail.com}{stacks.project@gmail.com}.
\end{proof}

\begin{lemma}
\label{obsolete-lemma-relative-effective-cartier-algebra}
Soit $A \to B$ un homomorphisme d'anneaux. Soit $f \in B$. Supposons que
\begin{enumerate}
\item $A \to B$ soit plat,
\item $f$ soit un non-diviseur de zéro, et
\item $A \to B/fB$ soit plat.
\end{enumerate}
Alors, pour tout idéal $I \subset A$, l'homomorphisme
$f : B/IB \to B/IB$ est injectif.
\end{lemma}

\begin{proof}
Remarquons que $IB = I \otimes_A B$ et $I(B/fB) = I \otimes_A B/fB$
par platitude de $B$ et de $B/fB$ sur $A$.
En particulier, $IB/fIB \cong I \otimes_A B/fB$ s'injecte
dans $B/fB$. Le résultat découle donc du lemme du serpent appliqué
au diagramme
$$
\xymatrix{
0 \ar[r] &
I \otimes_A B \ar[r] \ar[d]^f &
B \ar[r] \ar[d]^f &
B/IB \ar[r] \ar[d]^f &
0 \\
0 \ar[r] &
I \otimes_A B \ar[r] &
B \ar[r] &
B/IB \ar[r] &
0
}
$$
dont les lignes sont exactes.
\end{proof}

\begin{lemma}
\label{obsolete-lemma-faithfully-flat-injective}
Si $R \to S$ est un homomorphisme d'anneaux fidèlement plat, alors, pour tout
$R$-module $M$, l'homomorphisme $M \to S \otimes_R M$, $x \mapsto 1 \otimes x$,
est injectif.
\end{lemma}

\begin{proof}
Ce lemme est un doublon du lemme
\ref{algebra-lemma-faithfully-flat-universally-injective} du chapitre Algèbre.
\end{proof}

\begin{remark}
\label{obsolete-remark-section-colimits}
Cette référence et cette étiquette renvoyaient autrefois à une section du
chapitre Lissage des homomorphismes d'anneaux, mais son contenu a depuis été
repris dans la section \ref{algebra-section-colimits-flat} du chapitre Algèbre.
\end{remark}

\begin{lemma}
\label{obsolete-lemma-not-domain}
Soit $(R, \mathfrak m)$ un anneau local noethérien réduit de dimension $1$,
et soit $x \in \mathfrak m$ un non-diviseur de zéro. Soient
$\mathfrak q_1, \ldots, \mathfrak q_r$ les idéaux premiers minimaux de $R$.
Alors
$$
\text{length}_R(R/(x)) = \sum\nolimits_i \text{ord}_{R/\mathfrak q_i}(x)
$$
\end{lemma}

\begin{proof}
C'est un cas particulier (très simple) du lemme
\ref{chow-lemma-additivity-divisors-restricted} du chapitre Homologie de Chow.
\end{proof}

\begin{lemma}
\label{obsolete-lemma-bound-primes}
Soit $A$ un anneau intègre local noethérien normal de dimension $2$.
Pour $f \in \mathfrak m$ non nul, notons
$\text{div}(f) = \sum n_i (\mathfrak p_i)$
le diviseur associé à $f$ sur le spectre épointé de $A$.
Posons $|f| = \sum n_i$. Il existe des entiers $N$ et $M$
tels que $|f + g| \leq M$ pour tout $g \in \mathfrak m^N$.
\end{lemma}

\begin{proof}
Choisissons $h \in \mathfrak m$ tel que $f, h$ soit une suite régulière dans $A$
(cela résulte des lemmes \ref{algebra-lemma-criterion-normal} et
\ref{algebra-lemma-depth-drops-by-one} du chapitre Algèbre).
Nous allons démontrer le lemme avec $M = \text{length}_A(A/(f, h))$ et avec
n'importe quel entier $N$ tel que $\mathfrak m^N \subset (f, h)$. Un tel
entier $N$ existe puisque $\sqrt{(f, h)} = \mathfrak m$. Remarquons que
$M = \text{length}_A(A/(f + g, h))$ pour tout $g \in \mathfrak m^N$
car $(f, h) = (f + g, h)$. Cela entraîne en outre que $f + g, h$
est encore une suite régulière dans $A$ ; voir le lemme
\ref{algebra-lemma-reformulate-CM} du chapitre Algèbre.
Supposons maintenant que $\text{div}(f + g ) = \sum m_j (\mathfrak q_j)$.
Considérons alors l'homomorphisme
$$
c : A/(f + g) \longrightarrow \prod A/\mathfrak q_j^{(m_j)},
$$
où $\mathfrak q_j^{(m_j)}$ est la puissance symbolique ; voir la section
\ref{algebra-section-symbolic-power} du chapitre Algèbre.
Comme $A$ est normal, $A_{\mathfrak q_i}$ est un anneau de valuation discrète et
$$
A_{\mathfrak q_i}/(f + g) =
A_{\mathfrak q_i}/\mathfrak q_i^{m_i} A_{\mathfrak q_i} =
(A/\mathfrak q_i^{(m_i)})_{\mathfrak q_i}
$$
Puisque $V(f + g, h) = \{\mathfrak m\}$, il en résulte que $c$ devient
un isomorphisme après inversion de $h$ (nous omettons un détail mineur).
Comme $h$ est un non-diviseur de zéro sur $A/(f + g)$, la longueur de
$A/(f + g, h)$ est égale au quotient de Herbrand
$e_A(A/(f + g), 0, h)$ défini dans la section
\ref{chow-section-periodic-complexes} du chapitre Homologie de Chow.
De même, la longueur de $A/(h, \mathfrak q_j^{(m_j)})$ est égale à
$e_A(A/\mathfrak q_j^{(m_j)}, 0, h)$. On a alors
\begin{align*}
M & = \text{length}_A(A/(f + g, h)) \\
& =
e_A(A/(f + g), 0, h) \\
& =
\sum\nolimits_j e_A(A/\mathfrak q_j^{(m_j)}, 0, h) \\
& =
\sum\nolimits_j \sum\nolimits_{m = 0, \ldots, m_j - 1}
e_A(\mathfrak q_j^{(m)}/\mathfrak q_j^{(m + 1)}, 0, h).
\end{align*}
Les égalités résultent des lemmes
\ref{chow-lemma-additivity-periodic-length} et
\ref{chow-lemma-finite-periodic-length} du chapitre Homologie de Chow,
en utilisant notamment le fait que
le conoyau de $c$ est de longueur finie, comme expliqué ci-dessus.
Le lemme de Nakayama permet de vérifier directement que
$e_A(\mathfrak q^{(m)}/\mathfrak q^{(m + 1)}, 0, h)$
est au moins égal à $1$. Cela achève la démonstration du lemme.
\end{proof}

\begin{lemma}
\label{obsolete-lemma-flat-over-gorenstein-gorenstein-fibre}
Soit $A \to B$ un homomorphisme local plat d'anneaux locaux noethériens.
Si $A$ et $B/\mathfrak m_A B$ sont de Gorenstein, alors $B$ est de Gorenstein.
\end{lemma}

\begin{proof}
Cela résulte immédiatement du lemme
\ref{dualizing-lemma-flat-under-gorenstein} du chapitre Complexes dualisants.
\end{proof}

\begin{lemma}
\label{obsolete-lemma-kill-local}
Soit $(A, \mathfrak m)$ un anneau local noethérien.
Soit $I \subset A$ un idéal. Soit $M$ un $A$-module de type fini.
Soit $s$ un entier. Supposons que
\begin{enumerate}
\item $A$ possède un complexe dualisant,
\item si $\mathfrak p \not \in V(I)$ et
$V(\mathfrak p) \cap V(I) \not = \{\mathfrak m\}$, alors
$\text{depth}_{A_\mathfrak p}(M_\mathfrak p) + \dim(A/\mathfrak p) > s$.
\end{enumerate}
Il existe alors un entier $n > 0$ et un idéal $J \subset A$
tels que $V(J) \cap V(I) = \{\mathfrak m\}$ et que $JI^n$ annule
$H^i_\mathfrak m(M)$ pour $i \leq s$.
\end{lemma}

\begin{proof}
D'après le lemme \ref{local-cohomology-lemma-sitting-in-degrees}
du chapitre Cohomologie locale, il suffit de le démontrer pour le $A$-module de type fini
$E^i = \text{Ext}^{-i}_A(M, \omega_A^\bullet)$
pour $i \leq s$. Le support $Z$ de $E^0 \oplus \ldots \oplus E^s$
est fermé dans $\Spec(A)$ et ne contient aucun idéal premier comme en (2).
Il est donc contenu dans $V(JI^n)$ pour un certain $J$ comme dans
l'énoncé du lemme.
\end{proof}

\begin{lemma}
\label{obsolete-lemma-algebraize-local-cohomology-bis-bis}
Soit $(A, \mathfrak m)$ un anneau local noethérien.
Soit $I \subset A$ un idéal. Soit $M$ un $A$-module de type fini.
Soient $s$ et $d$ des entiers. Supposons que
\begin{enumerate}
\item[(a)] $A$ possède un complexe dualisant,
\item[(b)] $\text{cd}(A, I) \leq d$,
\item[(c)] si $\mathfrak p \not \in V(I)$, alors
$\text{depth}_{A_\mathfrak p}(M_\mathfrak p) > s$ ou
$\text{depth}_{A_\mathfrak p}(M_\mathfrak p) + \dim(A/\mathfrak p) > d + s$.
\end{enumerate}
Alors les hypothèses du lemme
\ref{algebraization-lemma-algebraize-local-cohomology-bis}
du chapitre Géométrie algébrique et formelle sont satisfaites
pour $A, I, \mathfrak m, M$, et l'homomorphisme
$H^i_\mathfrak m(M) \to \lim H^i_\mathfrak m(M/I^nM)$
est un isomorphisme pour $i \leq s$ ; ces modules sont en outre
annulés par une puissance de $I$.
\end{lemma}

\begin{proof}
Les hypothèses du lemme
\ref{algebraization-lemma-algebraize-local-cohomology-bis}
du chapitre Géométrie algébrique et formelle découlent du lemme plus général
\ref{algebraization-lemma-bootstrap-bis-bis} du même chapitre.
La conclusion du lemme
\ref{algebraization-lemma-algebraize-local-cohomology-bis}
donne alors la seconde assertion.
\end{proof}

\begin{lemma}
\label{obsolete-lemma-combine-one}
Dans la situation \ref{algebraization-situation-bootstrap} du chapitre
Géométrie algébrique et formelle, on a
$H^s_\mathfrak a(M) = \lim H^s_\mathfrak a(M/I^nM)$.
\end{lemma}

\begin{proof}
Cela résulte immédiatement du théorème
\ref{algebraization-theorem-final-bootstrap} du chapitre Géométrie algébrique et formelle.
La version originale de ce lemme, qui comportait des hypothèses supplémentaires,
a été remplacée par ce théorème.
\end{proof}

\begin{lemma}
\label{obsolete-lemma-fully-faithful-simple-one}
Soit $A$ un anneau noethérien. Soit $f \in \mathfrak a$ un élément d'un
idéal de $A$. Soit $U = \Spec(A) \setminus V(\mathfrak a)$. Supposons que
\begin{enumerate}
\item $A$ possède un complexe dualisant et soit complet pour la topologie $f$-adique,
\item $A_f$ soit $(S_2)$ et que, pour tout idéal premier minimal
$\mathfrak p \subset A$, on ait $f \not \in \mathfrak p$ et, pour tout
$\mathfrak q \in V(\mathfrak p) \cap V(\mathfrak a)$,
$\dim((A/\mathfrak p)_\mathfrak q) \geq 3$.
\end{enumerate}
Alors le foncteur de complétion
$$
\textit{Coh}(\mathcal{O}_U)
\longrightarrow
\textit{Coh}(U, I\mathcal{O}_U),
\quad
\mathcal{F} \longmapsto \mathcal{F}^\wedge
$$
est pleinement fidèle sur la sous-catégorie pleine des objets localement libres
de type fini.
\end{lemma}

\begin{proof}
Ce lemme est un cas particulier du lemme
\ref{algebraization-lemma-fully-faithful-simple-two}
du chapitre Géométrie algébrique et formelle.
\end{proof}



\section{Lemmes relatifs au th\'eor\`eme principal de Zariski}
\label{obsolete-section-ZMT}

\noindent
Les lemmes de cette section servaient initialement dans la d\'emonstration de
la version alg\'ebrique du th\'eor\`eme principal de Zariski,
Alg\`ebre, Th\'eor\`eme \ref{algebra-theorem-main-theorem}.

\begin{lemma}
\label{obsolete-lemma-change-equation-multiply}
Soient $R$ un anneau et $\varphi : R[x] \to S$ un homomorphisme
d'anneaux. Soit $t \in S$. Si $t$ est entier sur
$R[x]$, alors il existe un entier $\ell \geq 0$ tel que,
pour tout $a \in R$, l'\'el\'ement $\varphi(a)^\ell t$
soit entier relativement \`a $\varphi_a : R[y] \to S$, d\'efini par
$y \mapsto \varphi(ax)$ et $r \mapsto \varphi(r)$
pour $r\in R$.
\end{lemma}

\begin{proof}
\'Ecrivons $t^d + \sum_{i < d} \varphi(f_i)t^i = 0$
avec $f_i \in R[x]$. Soit $\ell$ le maximum des degr\'es
en $x$ de tous les $f_i$. En multipliant l'\'equation
par $\varphi(a)^\ell$, on obtient
$\varphi(a)^\ell t^d + \sum_{i < d} \varphi(a^\ell f_i)t^i = 0$.
Observons que chaque $\varphi(a^\ell f_i)$ appartient \`a l'image de
$\varphi_a$. Le r\'esultat d\'ecoule alors de
Alg\`ebre, Lemme \ref{algebra-lemma-make-integral-trivial}.
\end{proof}

\begin{lemma}
\label{obsolete-lemma-make-integral-less-trivial}
Soit $\varphi : R \to S$ un homomorphisme d'anneaux.
Supposons que $t \in S$ v\'erifie la
relation $\varphi(a_0) + \varphi(a_1)t + \ldots + \varphi(a_n) t^n = 0$.
Posons $u_n = \varphi(a_n)$, $u_{n-1} = u_n t + \varphi(a_{n-1})$,
et ainsi de suite jusqu'\`a $u_1 = u_2 t + \varphi(a_1)$.
Alors les \'el\'ements $u_n, u_{n-1}, \ldots, u_1$ ainsi que
$u_nt, u_{n-1}t, \ldots, u_1t$ sont entiers sur $R$,
et les id\'eaux $(\varphi(a_0), \ldots, \varphi(a_n))$ et
$(u_n, \ldots, u_1)$ de $S$ sont \'egaux.
\end{lemma}

\begin{proof}
Nous raisonnons par r\'ecurrence sur $n$. Comme $u_n = \varphi(a_n)$,
il r\'esulte d'Alg\`ebre, Lemme \ref{algebra-lemma-make-integral-trivial},
que $u_nt$ est entier sur $R$. Bien entendu,
$u_n = \varphi(a_n)$ est entier sur $R$. Par cons\'equent,
$u_{n - 1} = u_n t  + \varphi(a_{n - 1})$ est entier sur $R$ (voir
Alg\`ebre, Lemme \ref{algebra-lemma-integral-closure-is-ring})
et l'on a
$$
\varphi(a_0) + \varphi(a_1)t + \ldots + \varphi(a_{n - 1})t^{n - 1} +
u_{n - 1}t^{n - 1} = 0.
$$
L'hypoth\`ese de r\'ecurrence, appliqu\'ee \`a l'homomorphisme
$S' \to S$, o\`u $S'$ est la cl\^oture int\'egrale de $R$ dans $S$,
et \`a l'\'equation affich\'ee, montre donc que
$u_{n-1}, \ldots, u_1$ et $u_{n-1}t, \ldots, u_1t$
appartiennent eux aussi tous \`a $S'$. L'assertion concernant les id\'eaux
d\'ecoule imm\'ediatement de la forme des \'el\'ements et du fait que
$u_1t + \varphi(a_0) = 0$.
\end{proof}

\begin{lemma}
\label{obsolete-lemma-make-integral-not-in-ideal}
Soit $\varphi : R \to S$ un homomorphisme d'anneaux.
Supposons que $t \in S$ v\'erifie la
relation $\varphi(a_0) + \varphi(a_1)t + \ldots + \varphi(a_n) t^n = 0$.
Soit $J \subset S$ un id\'eal tel qu'il existe
au moins un $i$ pour lequel $\varphi(a_i) \not \in J$.
Alors il existe un \'el\'ement $u \in S$, $u \not\in J$, tel
que $u$ et $ut$ soient tous deux entiers sur $R$.
\end{lemma}

\begin{proof}
Cela d\'ecoule imm\'ediatement du Lemme \ref{obsolete-lemma-make-integral-less-trivial},
car l'un des \'el\'ements $u_i$ n'appartient pas \`a $J$.
\end{proof}

\noindent
Les deux lemmes suivants permettent de d\'ecrire les sous-sch\'emas ferm\'es de
$\mathbf{P}^1_R$ d\'efinis par une seule \'equation (non d\'eg\'en\'er\'ee).

\begin{lemma}
\label{obsolete-lemma-P1}
Soit $R$ un anneau.
Soit $F(X, Y) \in R[X, Y]$ un polyn\^ome homog\`ene de degr\'e
$d$. Supposons que, pour tout id\'eal premier $\mathfrak p$ de $R$,
au moins un coefficient de $F$ n'appartienne pas \`a $\mathfrak p$.
Munissons $S = R[X, Y]/(F)$ de sa structure d'anneau gradu\'e.
Alors, pour tout $n \geq d$, le $R$-module $S_n$
est localement libre de type fini, de rang $d$.
\end{lemma}

\begin{proof}
Le $R$-module $S_n$ admet une pr\'esentation
$$
R[X, Y]_{n-d} \to R[X, Y]_n \to S_n \to 0.
$$
D'apr\`es Alg\`ebre, Lemme \ref{algebra-lemma-cokernel-flat},
il suffit donc de montrer que la multiplication
par $F$ induit une application injective
$\kappa(\mathfrak p)[X, Y]
\to \kappa(\mathfrak p)[X, Y]$
pour tout id\'eal premier $\mathfrak p$.
Cela r\'esulte aussit\^ot de l'hypoth\`ese selon laquelle
l'image de $F$ modulo $\mathfrak p$ n'est pas le polyn\^ome nul.
L'assertion concernant les rangs en d\'ecoule elle aussi.
\end{proof}

\begin{lemma}
\label{obsolete-lemma-rel-prime-pols}
Soit $k$ un corps. Soient $F, G \in k[X, Y]$ des polyn\^omes homog\`enes
de degr\'es respectifs $d, e$. Supposons que $F, G$ soient premiers entre eux.
Alors la multiplication par $G$ est injective sur $S = k[X, Y]/(F)$.
\end{lemma}

\begin{proof}
C'est une mani\`ere de d\'efinir l'expression \og premiers entre eux \fg. Si
l'on adopte une autre d\'efinition, on peut montrer qu'elle lui est \'equivalente.
\end{proof}

\begin{lemma}
\label{obsolete-lemma-P1-localize}
Soit $R$ un anneau. Soit $F(X, Y) \in R[X, Y]$ homog\`ene de degr\'e
$d$. Munissons $S = R[X, Y]/(F)$ de sa structure d'anneau gradu\'e.
Soit $\mathfrak p \subset R$ un id\'eal premier tel qu'un
coefficient de $F$ n'appartienne pas \`a $\mathfrak p$.
Il existe un \'el\'ement $f \in R$, $f \not\in \mathfrak p$,
un entier $e$ et un \'el\'ement $G \in R[X, Y]_e$
tels que la multiplication par $G$ induise des isomorphismes
$(S_n)_f \to (S_{n + e})_f$ pour tout $n \geq d$.
\end{lemma}

\begin{proof}
Au cours de la d\'emonstration, nous pourrons remplacer $R$ par $R_f$
pour un $f\in R$, $f\not\in \mathfrak p$ (un nombre fini de fois).
Dans un premier temps, effectuons un tel remplacement de sorte
qu'un coefficient de $F$ soit inversible dans $R$.
En particulier, les modules $S_n$ sont d\`es lors localement
libres de rang $d$ pour $n \geq d$, d'apr\`es le Lemme \ref{obsolete-lemma-P1}.
Choisissons un \'el\'ement $G \in R[X, Y]_e$ tel que les images de
$G$ dans $\kappa(\mathfrak p)[X, Y]$ et de
$F(X, Y)$ soient premi\`eres entre elles (un tel choix est possible pour un certain $e$).
Appliquons Alg\`ebre, Lemme \ref{algebra-lemma-cokernel-flat}, \`a l'application
induite par la multiplication par $G$ de $S_d \to S_{d + e}$.
Le choix de $G$ et le Lemme \ref{obsolete-lemma-rel-prime-pols}
montrent que
$S_d \otimes \kappa(\mathfrak p) \to S_{d + e} \otimes \kappa(\mathfrak p)$
est bijective. Ainsi, apr\`es avoir remplac\'e $R$ par $R_f$ pour un
$f$ convenable, nous pouvons supposer que $G : S_d \to S_{d + e}$
est bijective. Il s'ensuit que l'image
de $G$ dans $\kappa(\mathfrak p')[X, Y]$ et celle de
$F$ sont premi\`eres entre elles, pour tout id\'eal premier $\mathfrak p'$
de $R$. En appliquant de nouveau Alg\`ebre, Lemme
\ref{algebra-lemma-cokernel-flat}, on voit alors que toutes les applications
$G : S_n \to S_{n + e}$, $n \geq d$, sont des isomorphismes.
\end{proof}

\begin{remark}
\label{obsolete-remark-algebra}
Soit $R$ un anneau. Supposons donn\'es $F \in R[X, Y]_d$
et $G \in R[X, Y]_e$ tels que, si l'on pose $S = R[X, Y]/(F)$,
on ait (1) $S_n$ localement libre de type fini, de rang $d$, pour
tout $n \geq d$, et (2) que la multiplication par $G$ d\'efinisse des
isomorphismes $S_n \to S_{n + e}$ pour tout $n \geq d$. Dans ce
cas, on peut d\'efinir une $R$-alg\`ebre finie et localement libre
$A$ de la mani\`ere suivante :
\begin{enumerate}
\item comme $R$-module, $A = S_{ed}$ ;
\item la multiplication $A \times A \to A$ est donn\'ee par
la r\`egle suivante : $H_1 H_2 = H_3$ si et seulement si $G^d H_3 = H_1 H_2$
dans $S_{2ed}$.
\end{enumerate}
Cette d\'efinition est licite, car la multiplication par $G^d$
induit une application bijective $S_{de} \to S_{2de}$.
On v\'erifie ais\'ement qu'elle d\'efinit une structure d'anneau.
Notons le fait d\'econcertant que l'\'el\'ement $G^d$
d\'efinit l'\'el\'ement unit\'e de l'anneau $A$.
\end{remark}

\begin{lemma}
\label{obsolete-lemma-finite-after-localization}
Soient $R$ un anneau et $f \in R$.
Supposons donn\'es $S$, $S'$ et les fl\`eches pleines
formant le diagramme commutatif d'anneaux suivant :
$$
\xymatrix{
& S'' \ar@{-->}[rd] \ar@{-->}[dd] &
\\
R \ar[rr] \ar@{-->}[ru] \ar[d] &  & S \ar[d]
\\
R_f \ar[r] & S' \ar[r] & S_f
}
$$
Supposons que $R_f \to S'$ soit fini. Alors on peut trouver
un homomorphisme d'anneaux fini $R \to S''$ et des fl\`eches pointill\'ees comme
dans le diagramme, de sorte que $S' = (S'')_f$.
\end{lemma}

\begin{proof}
En effet, supposons que $S'$ soit engendr\'ee par
les $x_i$ sur $R_f$, $i = 1, \ldots, w$. Soit $P_i(t) \in R_f[t]$
un polyn\^ome unitaire tel que $P_i(x_i) = 0$.
Le polyn\^ome $P_i$ est de degr\'e $d_i > 0$. \'Ecrivons
$P_i(t) = t^{d_i} + \sum_{j < d_i} (a_{ij}/f^n) t^j$
pour un m\^eme entier $n$. \'Ecrivons aussi
l'image de $x_i$ dans $S_f$ sous la forme $g_i / f^n$
avec un \'el\'ement convenable $g_i \in S$. On sait alors
que l'\'el\'ement
$\xi_i = f^{nd_i} g_i^{d_i} + \sum_{j < d_i} f^{n(d_i - j)} a_{ij} g_i^j$
de $S$ est annul\'e par une puissance de $f$.
En rempla\c cant $n$ par un entier $n'$ plus grand, ce qui revient \`a remplacer
$g_i$ par $f^{n' - n}g_i$, on peut donc supposer que $\xi_i = 0$.
L'alg\`ebre $S'$ est alors engendr\'ee par les \'el\'ements
$f^n x_i$, chacun \'etant un z\'ero du
polyn\^ome unitaire $Q_i(t) = t^{d_i} +
\sum_{j < d_i} f^{n(d_i - j)} a_{ij} t^j$
\`a coefficients dans $R$. Par construction, on a aussi
$Q_i(f^ng_i) = 0$ dans $S$. On obtient ainsi une $R$-alg\`ebre finie
$S'' = R[z_1, \ldots, z_w]/(Q_1(z_1), \ldots, Q_w(z_w))$
qui s'ins\`ere dans un diagramme commutatif comme ci-dessus.
L'homomorphisme $\alpha : S'' \to S$ envoie $z_i$ sur $f^ng_i$, et
l'homomorphisme $\beta : S'' \to S'$ envoie $z_i$ sur $f^nx_i$.
Il se peut que $\beta$ n'induise pas encore un
isomorphisme $(S'')_f \cong S'$.
Pour l'instant, on sait seulement que cette application
est surjective. Le probl\`eme est qu'il pourrait exister des
\'el\'ements $h/f^n \in (S'')_f$ qui s'envoient sur z\'ero
dans $S'$ sans \^etre eux-m\^emes nuls. Dans ce cas, $\alpha(h)$
est un \'el\'ement de $S$ tel que $f^N \alpha(h) = 0$
pour un certain $N$. Ainsi, $f^N h$ appartient \`a l'id\'eal
$J = \{h \in S'' \mid \alpha(h) = 0 \text{ et }
\beta(h) = 0\}$ de $S''$. Il est alors facile de voir que
$S''/J$ convient.
\end{proof}
\section{Homomorphismes formellement lisses}
\label{obsolete-section-formally-smooth}

\begin{lemma}
\label{obsolete-lemma-formally-smooth-smooth}
Soit $R$ un anneau. Soit $S$ une $R$-alg\`ebre.
Si $S$ est de pr\'esentation finie et formellement lisse sur $R$,
alors $S$ est lisse sur $R$.
\end{lemma}

\begin{proof}
Voir Alg\`ebre, Proposition \ref{algebra-proposition-smooth-formally-smooth}.
\end{proof}

\begin{remark}
\label{obsolete-remark-equation-derivatives}
Cette \`etiquette renvoyait \`a une \`equation de la d\'emonstration de
la Proposition \ref{restricted-proposition-approximate}
d'Alg\'ebrisation des espaces formels.
Apr\`es r\'eorganisation du contenu, cette \`etiquette n'est plus utilis\'ee.
\end{remark}

\begin{remark}
\label{obsolete-remark-equation-ci}
Cette \`etiquette renvoyait \`a une \`equation de la d\'emonstration de
la Proposition \ref{restricted-proposition-approximate}
d'Alg\'ebrisation des espaces formels.
Apr\`es r\'eorganisation du contenu, cette \`etiquette n'est plus utilis\'ee.
\end{remark}

\begin{remark}
\label{obsolete-remark-equation-in-ideal}
Cette \`etiquette renvoyait \`a une \`equation de la d\'emonstration de
la Proposition \ref{restricted-proposition-approximate}
d'Alg\'ebrisation des espaces formels.
Apr\`es r\'eorganisation du contenu, cette \`etiquette n'est plus utilis\'ee.
\end{remark}

\begin{remark}
\label{obsolete-remark-equation-derivatives-analogue}
Cette \`etiquette renvoyait \`a une \`equation de la d\'emonstration de
la Proposition \ref{restricted-proposition-approximate}
d'Alg\'ebrisation des espaces formels.
Apr\`es r\'eorganisation du contenu, cette \`etiquette n'est plus utilis\'ee.
\end{remark}

\begin{remark}
\label{obsolete-remark-equation-go-down}
Cette \`etiquette renvoyait \`a une \`equation de la d\'emonstration du
Lemme \ref{restricted-lemma-lift-approximation}
d'Alg\'ebrisation des espaces formels.
Apr\`es r\'eorganisation du contenu, cette \`etiquette n'est plus utilis\'ee.
\end{remark}

\begin{lemma}
\label{obsolete-lemma-get-morphism-general}
Soit $A$ un anneau noeth\'erien. Soit $I \subset A$ un id\'eal.
Soit $t$ le nombre minimum de g\'en\'erateurs de $I$.
Soit $C$ une $A$-alg\`ebre noeth\'erienne $I$-adiquement compl\`ete.
Il existe un entier $d \geq 0$ ne d\'ependant que de
$I \subset A \to C$ et poss\'edant la propri\'et\'e suivante : \`etant donn\'es
\begin{enumerate}
\item un entier $c \geq 0$ et l'alg\`ebre $B$ de l'\`Equation
(\ref{restricted-equation-C-prime}) d'Alg\'ebrisation des espaces formels,
tels que, pour tout $a \in I^c$,
la multiplication par $a$ sur $\NL_{B/A}^\wedge$ soit nulle dans $D(B)$,
\item un entier $n > 2t\max(c, d)$,
\item un homomorphisme de $A/I^n$-alg\`ebres
$\psi_n : B/I^nB \to C/I^nC$,
\end{enumerate}
il existe un homomorphisme de $A$-alg\`ebres $\varphi : B \to C$ tel
que $\psi_n \bmod I^{m - c} = \varphi \bmod I^{m - c}$,
o\`u $m = \lfloor \frac{n}{t} \rfloor$.
\end{lemma}

\begin{proof}
Ce lemme est devenu obsol\`ete en raison du r\'esultat plus fort
d'Alg\'ebrisation des espaces formels,
Lemme \ref{restricted-lemma-get-morphism-general-better}.
Nous allons en fait l'en d\'eduire.

\medskip\noindent
Soit $I \subset A \to C$ comme dans l'\'enonc\'e ci-dessus.
Notons $d(\text{Gr}_I(C))$ et $q(\text{Gr}_I(C))$ les entiers obtenus dans
Cohomologie locale, section \ref{local-cohomology-section-uniform}.
Remarquons que $t$ majore le nombre minimum de g\'en\'erateurs de $IC$;
on a donc $d(\text{Gr}_I(C)) + 1 \leq t$, voir la discussion de
Cohomologie locale, section \ref{local-cohomology-section-uniform}.
Nous pouvons supposer, et nous supposerons, que $t \geq 1$, car sinon le lemme
n'affirme rien. Nous affirmons que le lemme est vrai en prenant
$$
d = q(\text{Gr}_I(C))
$$
En effet, supposons que $c$, $B$, $n$ et $\psi_n$ soient comme dans l'\'enonc\'e
ci-dessus. On a alors
$$
n > 2t\max(c, d) \Rightarrow n \geq 2tc + 1 \Rightarrow
n \geq 2(d(\text{Gr}_I(C)) + 1)c + 1
$$
D'autre part,
$$
n > 2t\max(c, d) \Rightarrow n > t(c + d) \Rightarrow
n \geq q(C) + tc \geq q(\text{Gr}_I(C)) + (d(\text{Gr}_I(C)) + 1)c
$$
Les hypoth\`eses du Lemme \ref{restricted-lemma-get-morphism-general-better}
d'Alg\'ebrisation des espaces formels sont donc satisfaites,
et nous obtenons un homomorphisme de $A$-alg\`ebres
$\varphi : B \to C$ congru \`a $\psi_n$
modulo $I^{n - (d(\text{Gr}_I(C)) + 1)c}C$.
Comme
\begin{align*}
n - (d(\text{Gr}_I(C)) + 1)c
& = \frac{n}{t} + \frac{(t - 1)n}{t} - (d(\text{Gr}_I(C)) + 1)c \\
& \geq \frac{n}{t} + \frac{d(\text{Gr}_I(C))n}{t} - (d(\text{Gr}_I(C)) + 1)c \\
& > \frac{n}{t} + \frac{d(\text{Gr}_I(C))2tc}{t} - (d(\text{Gr}_I(C)) + 1)c \\
& = \frac{n}{t} + 2d(\text{Gr}_I(C))c - (d(\text{Gr}_I(C)) + 1)c \\
& = \frac{n}{t} + d(\text{Gr}_I(C))c - c \\
& \geq m - c
\end{align*}
on voit que $\varphi$ et $\psi_n$ sont congrus modulo $I^{m - c}C$, comme voulu.
\end{proof}
\section{Sites et faisceaux}
\label{obsolete-section-sites}

\begin{remark}[Absence de morphisme du prolongement par zéro vers l'image directe]
\label{obsolete-remark-from-shriek-to-star}
Soit $U$ un objet d'un site $\mathcal{C}$. Pour tout faisceau abélien
$\mathcal{G}$ sur $\mathcal{C}/U$, on peut se demander s'il existe
un morphisme canonique
$$
c : j_{U!}\mathcal{G} \longrightarrow j_{U*}\mathcal{G}
$$
Construire un tel morphisme revient à construire un morphisme
$j_U^{-1}j_{U!}\mathcal{G} \to \mathcal{G}$.
Remarquons que la restriction commute au passage au faisceau associé.
Nous pouvons donc utiliser le préfaisceau du lemme
Modules sur les sites, lemme \ref{sites-modules-lemma-extension-by-zero}.
Il suffit ainsi de définir, pour $V/U$, un morphisme
$$
\bigoplus\nolimits_{\varphi \in \Mor_\mathcal{C}(V, U)}
\mathcal{G}(V \xrightarrow{\varphi} U)
\longrightarrow
\mathcal{G}(V/U)
$$
compatible aux restrictions. Il semble que l'on puisse prendre le morphisme
nul sur tous les facteurs directs, sauf sur celui pour lequel $\varphi$
est le morphisme structural $\varphi_0 : V \to U$, où l'on prend $1$.
Cependant, ce morphisme n'est pas compatible aux applications de restriction.
En effet, si $\alpha : V' \to V$ est un morphisme de $\mathcal{C}$, notons
$V'/U$ l'objet de $\mathcal{C}/U$ dont le morphisme structural est
$\varphi'_0 = \varphi_0 \circ \alpha$.
Il nous faut vérifier que le diagramme
$$
\xymatrix{
\bigoplus\nolimits_{\varphi \in \Mor_\mathcal{C}(V, U)}
\mathcal{G}(V \xrightarrow{\varphi} U)
\ar[d] \ar[r] &
\mathcal{G}(V/U) \ar[d] \\
\bigoplus\nolimits_{\varphi' \in \Mor_\mathcal{C}(V', U)}
\mathcal{G}(V' \xrightarrow{\varphi'} U)
\ar[r] &
\mathcal{G}(V'/U)
}
$$
est commutatif. Le problème est qu'il peut exister un morphisme
$\varphi : V \to U$, distinct de $\varphi_0$, tel que
$\varphi \circ \alpha = \varphi'_0$.
La flèche verticale de gauche envoie alors le facteur direct correspondant
à $\varphi$ dans celui sur lequel la flèche horizontale inférieure est
égale à $1$, et il est presque certain que le diagramme ne commute pas.
\end{remark}




\section{Cohomologie}
\label{obsolete-section-cohomology}

\noindent
Lemmes devenus obsolètes concernant la cohomologie.

\begin{lemma}
\label{obsolete-lemma-ML-general}
Soit $I$ un idéal d'un anneau $A$. Soit $X$ un schéma sur $\Spec(A)$. Soit
$$
\ldots \to \mathcal{F}_3 \to \mathcal{F}_2 \to \mathcal{F}_1
$$
un système projectif de $\mathcal{O}_X$-modules tel que
$\mathcal{F}_n = \mathcal{F}_{n + 1}/I^n\mathcal{F}_{n + 1}$.
Supposons que
$$
\bigoplus\nolimits_{n \geq 0} H^1(X, I^n\mathcal{F}_{n + 1})
$$
satisfasse la condition de chaîne croissante en tant que
$\bigoplus_{n \geq 0} I^n/I^{n + 1}$-module gradué.
Alors le système projectif $M_n = \Gamma(X, \mathcal{F}_n)$ satisfait
la condition de Mittag-Leffler.
\end{lemma}

\begin{proof}
C'est un cas particulier du lemme plus général
Cohomologie, lemme \ref{cohomology-lemma-ML-general}.
\end{proof}

\begin{lemma}
\label{obsolete-lemma-ML-general-better}
Soit $I$ un idéal d'un anneau $A$. Soit $X$ un schéma sur $\Spec(A)$. Soit
$$
\ldots \to \mathcal{F}_3 \to \mathcal{F}_2 \to \mathcal{F}_1
$$
un système projectif de $\mathcal{O}_X$-modules tel que
$\mathcal{F}_n = \mathcal{F}_{n + 1}/I^n\mathcal{F}_{n + 1}$.
Pour tout $n$, posons
$$
H^1_n =
\bigcap\nolimits_{m \geq n}
\Im\left(
H^1(X, I^n\mathcal{F}_{m + 1}) \to H^1(X, I^n\mathcal{F}_{n + 1})
\right)
$$
Si $\bigoplus H^1_n$ satisfait la condition de chaîne croissante en tant que
$\bigoplus_{n \geq 0} I^n/I^{n + 1}$-module gradué, alors le système
projectif $M_n = \Gamma(X, \mathcal{F}_n)$ satisfait la condition
de Mittag-Leffler.
\end{lemma}

\begin{proof}
C'est un cas particulier du lemme plus général
Cohomologie, lemme \ref{cohomology-lemma-ML-general-better}.
\end{proof}

\begin{lemma}
\label{obsolete-lemma-topology-I-adic-general}
Soit $I$ un idéal de type fini d'un anneau $A$.
Soit $X$ un schéma sur $\Spec(A)$. Soit
$$
\ldots \to \mathcal{F}_3 \to \mathcal{F}_2 \to \mathcal{F}_1
$$
un système projectif de $\mathcal{O}_X$-modules tel que
$\mathcal{F}_n = \mathcal{F}_{n + 1}/I^n\mathcal{F}_{n + 1}$.
Supposons que
$$
\bigoplus\nolimits_{n \geq 0} H^0(X, I^n\mathcal{F}_{n + 1})
$$
satisfasse la condition de chaîne croissante en tant que
$\bigoplus_{n \geq 0} I^n/I^{n + 1}$-module gradué.
Alors la topologie limite sur $M = \lim \Gamma(X, \mathcal{F}_n)$
est la topologie $I$-adique.
\end{lemma}

\begin{proof}
C'est un cas particulier du lemme plus général
Cohomologie, lemme \ref{cohomology-lemma-topology-I-adic-general}.
\end{proof}

\begin{lemma}
\label{obsolete-lemma-pullback-K-flat}
Soit $(\Sh(\mathcal{C}), \mathcal{O}_\mathcal{C})$ un topos annelé.
Pour tout complexe de $\mathcal{O}_\mathcal{C}$-modules
$\mathcal{G}^\bullet$, il existe un quasi-isomorphisme
$\mathcal{K}^\bullet \to \mathcal{G}^\bullet$ tel que
$f^*\mathcal{K}^\bullet$ soit un complexe K-plat de
$\mathcal{O}_\mathcal{D}$-modules pour tout morphisme de topos annelés
$f : (\Sh(\mathcal{D}), \mathcal{O}_\mathcal{D}) \to
(\Sh(\mathcal{C}), \mathcal{O}_\mathcal{C})$.
\end{lemma}

\begin{proof}
Cela résulte de Cohomologie sur les sites, lemmes
\ref{sites-cohomology-lemma-K-flat-resolution} et
\ref{sites-cohomology-lemma-pullback-K-flat}.
\end{proof}

\begin{remark}
\label{obsolete-remark-pullback-K-flat}
Cette remarque examinait autrefois ce que l'on savait du caractère K-plat,
ou non, des images réciproques de complexes K-plats; elle est désormais
rendue obsolète par Cohomologie sur les sites, lemme
\ref{sites-cohomology-lemma-pullback-K-flat}.
\end{remark}

\noindent
Le lemme suivant calcule, dans un cas particulier, les faisceaux de
cohomologie de la limite dérivée.

\begin{lemma}
\label{obsolete-lemma-Rlim-of-system}
Soit $(\mathcal{C}, \mathcal{O})$ un site annelé. Soit $(K_n)$
un système projectif d'objets de $D(\mathcal{O})$.
Soit $\mathcal{B} \subset \Ob(\mathcal{C})$ un sous-ensemble.
Soit $d \in \mathbf{N}$. Supposons que
\begin{enumerate}
\item $K_n$ soit un objet de $D^+(\mathcal{O})$ pour tout $n$,
\item pour tout $q \in \mathbf{Z}$, il existe
$n(q)$ tel que $H^q(K_{n + 1}) \to H^q(K_n)$ soit un isomorphisme pour
$n \geq n(q)$,
\item tout objet de $\mathcal{C}$ admette un recouvrement dont les membres
appartiennent à $\mathcal{B}$,
\item pour tout $U \in \mathcal{B}$, on ait $H^p(U, H^q(K_n)) = 0$
pour $p > d$ et tout $q$.
\end{enumerate}
Alors $H^m(R\lim K_n) = \lim H^m(K_n)$ pour tout $m \in \mathbf{Z}$.
\end{lemma}

\begin{proof}
Posons $K = R\lim K_n$. Soit $U \in \mathcal{B}$. Pour chaque $n$, il existe
une suite spectrale
$$
H^p(U, H^q(K_n)) \Rightarrow H^{p + q}(U, K_n)
$$
qui converge puisque $K_n$ est borné inférieurement; voir
Catégories dérivées, lemme \ref{derived-lemma-two-ss-complex-functor}.
Fixons $m \in \mathbf{Z}$. L'hypothèse (4) montre que seuls les groupes
$H^p(U, H^q(K_n))$ contribuent à $H^m(U, K_n)$ pour
$0 \leq p \leq d$ et $m - d \leq q \leq m$.
L'hypothèse (2) implique donc que
$H^m(U, K_{n + 1}) \to H^m(U, K_n)$ est un isomorphisme dès que
$n \geq \max{n(m), \ldots, n(m - d)}$. Le foncteur $R\Gamma(U, -)$
commute aux limites dérivées d'après
Injectifs, lemme \ref{injectives-lemma-RF-commutes-with-Rlim}.
Ainsi,
$$
H^m(U, K) = H^m(R\lim R\Gamma(U, K_n))
$$
D'autre part, nous venons de voir que les groupes de cohomologie des
complexes $R\Gamma(U, K_n)$ sont stationnaires à partir d'un certain rang.
Il résulte donc de Compléments d'algèbre, remarque
\ref{more-algebra-remark-compare-derived-limit}, que $H^m(U, K)$ est égal
à $H^m(U, K_n)$ pour tout $n \gg 0$, avec une borne indépendante de
$U \in \mathcal{B}$. Fixons un tel $n$. Enfin, rappelons que $H^m(K)$
est le faisceau associé au préfaisceau $U \mapsto H^m(U, K)$, tandis que
$H^m(K_n)$ est le faisceau associé au préfaisceau
$U \mapsto H^m(U, K_n)$. Sur les éléments de $\mathcal{B}$, ces
préfaisceaux prennent les mêmes valeurs. L'hypothèse (3) garantit donc
que les faisceaux associés sont eux aussi identiques. Le lemme en résulte.
\end{proof}

\begin{lemma}
\label{obsolete-lemma-trivialities-cohomological-descent-abelian}
Dans la situation de Espaces simpliciaux,
\ref{spaces-simplicial-situation-simplicial-site}, soit $a_0$ une
augmentation vers un site $\mathcal{D}$, comme dans Espaces simpliciaux,
remarque \ref{spaces-simplicial-remark-augmentation-site}.
Supposons données des sous-catégories de Serre faibles strictement pleines
$$
\mathcal{A} \subset \textit{Ab}(\mathcal{D}),\quad
\mathcal{A}_n \subset \textit{Ab}(\mathcal{C}_n)
$$
Alors
\begin{enumerate}
\item[(1)]
la collection des faisceaux abéliens $\mathcal{F}$ sur
$\mathcal{C}_{total}$ dont la restriction à $\mathcal{C}_n$ appartient à
$\mathcal{A}_n$ pour tout $n$ est une sous-catégorie de Serre faible
strictement pleine
$\mathcal{A}_{total} \subset \textit{Ab}(\mathcal{C}_{total})$.
\end{enumerate}
Si $a_n^{-1}$ envoie $\mathcal{A}$ dans $\mathcal{A}_n$ pour tout $n$, alors
\begin{enumerate}
\item[(2)] $a^{-1}$ envoie $\mathcal{A}$ dans $\mathcal{A}_{total}$, et
\item[(3)] $a^{-1}$ envoie $D_\mathcal{A}(\mathcal{D})$ dans
$D_{\mathcal{A}_{total}}(\mathcal{C}_{total})$.
\end{enumerate}
Si $R^qa_{n, *}$ envoie $\mathcal{A}_n$ dans $\mathcal{A}$
pour tous $n, q$, alors
\begin{enumerate}
\item[(4)] $R^qa_*$ envoie $\mathcal{A}_{total}$ dans $\mathcal{A}$ pour tout
$q$, et
\item[(5)] $Ra_*$ envoie
$D_{\mathcal{A}_{total}}^+(\mathcal{C}_{total})$
dans $D_\mathcal{A}^+(\mathcal{D})$.
\end{enumerate}
\end{lemma}

\begin{proof}
Les seules assertions qui demandent une explication sont (4) et (5).
L'assertion (4) résulte de la suite spectrale de Espaces simpliciaux,
lemme \ref{spaces-simplicial-lemma-augmentation-spectral-sequence},
et de Homologie, lemme \ref{homology-lemma-biregular-ss-converges}.
L'assertion (5) s'en déduit en considérant la suite spectrale associée à
la filtration canonique par les troncations d'un objet
$K$ de $D_{\mathcal{A}_{total}}^+(\mathcal{C}_{total})$.
Nous omettons les détails.
\end{proof}

\begin{remark}
\label{obsolete-remark-cohomology-topics}
Ce tag renvoyait autrefois à une section du chapitre consacré à la
cohomologie qui énumérait les sujets restant à traiter.
\end{remark}

\begin{remark}
\label{obsolete-remark-sites-cohomology-topics}
Ce tag renvoyait autrefois à une section du chapitre consacré à la
cohomologie qui énumérait les sujets restant à traiter.
\end{remark}

\begin{remark}
\label{obsolete-remark-V-implies-C}
Ce tag renvoyait autrefois au cas particulier de
Cohomologie sur les sites, lemme
\ref{sites-cohomology-lemma-V-implies-C-general}, dans la situation décrite
par Cohomologie sur les sites, lemme
\ref{sites-cohomology-lemma-compare-qc-zar}.
\end{remark}

\begin{remark}
\label{obsolete-remark-V-implies-cohomology}
Ce tag renvoyait autrefois au cas particulier de
Cohomologie sur les sites, lemme
\ref{sites-cohomology-lemma-V-implies-cohomology-general}, dans la situation
décrite par Cohomologie sur les sites, lemme
\ref{sites-cohomology-lemma-compare-qc-zar}.
\end{remark}

\begin{remark}
\label{obsolete-remark-induction-step-V-C}
Ce tag renvoyait autrefois au cas particulier de
Cohomologie sur les sites, lemme
\ref{sites-cohomology-lemma-induction-step-V-C-general}, dans la situation
décrite par Cohomologie sur les sites, lemme
\ref{sites-cohomology-lemma-compare-qc-zar}.
\end{remark}

\begin{remark}
\label{obsolete-remark-V-implies-C-etale-fppf}
Ce tag renvoyait autrefois au cas particulier de
Cohomologie sur les sites, lemme
\ref{sites-cohomology-lemma-V-implies-C-general}, dans la situation décrite
par Cohomologie étale, lemme
\ref{etale-cohomology-lemma-compare-fppf-etale}.
\end{remark}

\begin{remark}
\label{obsolete-remark-V-implies-cohomology-etale-fppf}
Ce tag renvoyait autrefois au cas particulier de
Cohomologie sur les sites, lemme
\ref{sites-cohomology-lemma-V-implies-cohomology-general}, dans la situation
décrite par Cohomologie étale, lemme
\ref{etale-cohomology-lemma-compare-fppf-etale}.
\end{remark}

\begin{remark}
\label{obsolete-remark-induction-step-V-C-etale-fppf}
Ce tag renvoyait autrefois au cas particulier de
Cohomologie sur les sites, lemme
\ref{sites-cohomology-lemma-induction-step-V-C-general}, dans la situation
décrite par Cohomologie étale, lemme
\ref{etale-cohomology-lemma-compare-fppf-etale}.
\end{remark}

\begin{remark}
\label{obsolete-remark-V-implies-C-etale-ph}
Ce tag renvoyait autrefois au cas particulier de
Cohomologie sur les sites, lemme
\ref{sites-cohomology-lemma-V-implies-C-general}, dans la situation décrite
par Cohomologie étale, lemme
\ref{etale-cohomology-lemma-compare-ph-etale}.
\end{remark}

\begin{remark}
\label{obsolete-remark-V-implies-cohomology-etale-ph}
Ce tag renvoyait autrefois au cas particulier de
Cohomologie sur les sites, lemme
\ref{sites-cohomology-lemma-V-implies-cohomology-general}, dans la situation
décrite par Cohomologie étale, lemme
\ref{etale-cohomology-lemma-compare-ph-etale}.
\end{remark}

\begin{remark}
\label{obsolete-remark-V-implies-cohomology-etale-ph-extra}
Ce tag renvoyait autrefois au cas particulier de
Cohomologie sur les sites, lemme
\ref{sites-cohomology-lemma-V-implies-cohomology-extra-general}, dans la
situation décrite par Cohomologie étale, lemme
\ref{etale-cohomology-lemma-compare-ph-etale}.
\end{remark}

\begin{remark}
\label{obsolete-remark-make-class-zero}
Ce tag renvoyait autrefois au cas particulier de
Cohomologie sur les sites, lemme
\ref{sites-cohomology-lemma-make-class-zero-general}, dans la situation
décrite par Cohomologie étale, lemme
\ref{etale-cohomology-lemma-compare-ph-etale}.
\end{remark}

\begin{remark}
\label{obsolete-remark-induction-step-V-C-etale-ph}
Ce tag renvoyait autrefois au cas particulier de
Cohomologie sur les sites, lemme
\ref{sites-cohomology-lemma-induction-step-V-C-general}, dans la situation
décrite par Cohomologie étale, lemme
\ref{etale-cohomology-lemma-compare-ph-etale}.
\end{remark}

\begin{remark}
\label{obsolete-remark-V-implies-C-etale-h}
Ce tag renvoyait autrefois au cas particulier de
Cohomologie sur les sites, lemme
\ref{sites-cohomology-lemma-V-implies-C-general}, dans la situation décrite
par Cohomologie étale, lemme
\ref{etale-cohomology-lemma-compare-h-etale}.
\end{remark}

\begin{remark}
\label{obsolete-remark-V-implies-cohomology-etale-h}
Ce tag renvoyait autrefois au cas particulier de
Cohomologie sur les sites, lemme
\ref{sites-cohomology-lemma-V-implies-cohomology-general}, dans la situation
décrite par Cohomologie étale, lemme
\ref{etale-cohomology-lemma-compare-h-etale}.
\end{remark}

\begin{remark}
\label{obsolete-remark-V-implies-cohomology-etale-h-extra}
Ce tag renvoyait autrefois au cas particulier de
Cohomologie sur les sites, lemme
\ref{sites-cohomology-lemma-V-implies-cohomology-extra-general}, dans la
situation décrite par Cohomologie étale, lemme
\ref{etale-cohomology-lemma-compare-h-etale}.
\end{remark}

\begin{remark}
\label{obsolete-remark-induction-step-V-C-etale-h}
Ce tag renvoyait autrefois au cas particulier de
Cohomologie sur les sites, lemme
\ref{sites-cohomology-lemma-induction-step-V-C-general}, dans la situation
décrite par Cohomologie étale, lemme
\ref{etale-cohomology-lemma-compare-h-etale}.
\end{remark}

\begin{remark}
\label{obsolete-remark-how-used}
Ce tag se trouvait autrefois dans le chapitre consacré à la cohomologie
étale, mais une réorganisation fait qu'il n'y a désormais plus sa place.
Son contenu était le suivant :
Cohomologie étale, lemme \ref{etale-cohomology-lemma-when-ctf},
permet de prouver que, si $f : X \to Y$ est un morphisme de schémas séparé
et de type fini et si $Y$ est noethérien, alors $Rf_!$ induit un foncteur
$D_{ctf}(X_\etale, \Lambda) \to D_{ctf}(Y_\etale, \Lambda)$.
Un exemple de cet argument, lorsque $Y$ est le spectre d'un corps et que
$X$ est une courbe, est donné dans La formule des traces,
proposition \ref{trace-proposition-projective-curve-constructible-cohomology}.
\end{remark}

\begin{lemma}
\label{obsolete-lemma-glue-f-upper-shriek}
Soit $f : X \to Y$ un morphisme de schémas localement quasi-fini.
Il existe un unique foncteur
$f^! : \textit{Ab}(Y_\etale) \to \textit{Ab}(X_\etale)$ tel que
\begin{enumerate}
\item pour toute immersion ouverte $j : U \to X$ telle que $f \circ j$
soit séparé, il existe un isomorphisme canonique
$j^! \circ f^! = (f \circ j)^!$, et
\item ces isomorphismes, pour $U \subset U' \subset X$, soient compatibles
avec ceux de Compléments de cohomologie étale, lemme
\ref{more-etale-lemma-upper-shriek-restriction}.
\end{enumerate}
\end{lemma}

\begin{proof}
C'est une conséquence immédiate de Compléments de cohomologie étale, lemmes
\ref{more-etale-lemma-lqf-f-upper-shriek} et
\ref{more-etale-lemma-upper-shriek-restriction}.
\end{proof}

\begin{proposition}
\label{obsolete-proposition-lqf-shriek}
Soit $f : X \to Y$ un morphisme localement quasi-fini. Il existe des
foncteurs adjoints $f_! : \textit{Ab}(X_\etale) \to \textit{Ab}(Y_\etale)$
et $f^! : \textit{Ab}(Y_\etale) \to \textit{Ab}(X_\etale)$
ayant les propriétés suivantes :
\begin{enumerate}
\item le foncteur $f^!$ est celui construit dans Compléments de cohomologie
étale, lemme \ref{more-etale-lemma-lqf-f-upper-shriek},
\item pour toute immersion ouverte $j : U \to X$ telle que $f \circ j$
soit séparé, il existe un isomorphisme canonique
$f_! \circ j_! = (f \circ j)_!$, et
\item ces isomorphismes, pour $U \subset U' \subset X$, sont compatibles
avec ceux de Compléments de cohomologie étale, lemme
\ref{more-etale-lemma-f-shriek-composition}.
\end{enumerate}
\end{proposition}

\begin{proof}
Voir Compléments de cohomologie étale, sections
\ref{more-etale-section-finite-support} et
\ref{more-etale-section-duality-locally-quasi-finite}.
\end{proof}

\begin{lemma}
\label{obsolete-lemma-lqf-f-upper-shriek-stalk}
Soit $f : X \to Y$ un morphisme de schémas localement quasi-fini.
Pour un groupe abélien $A$ et un point géométrique
$\overline{y} : \Spec(k) \to Y$, on a
$f^!(\overline{y}_*A) = \prod\nolimits_{f(\overline{x}) = \overline{y}}
\overline{x}_*A$.
\end{lemma}

\begin{proof}
Cela résulte de l'assertion correspondante de Compléments de cohomologie
étale, lemme \ref{more-etale-lemma-lqf-f-upper-shriek}.
\end{proof}

\begin{lemma}
\label{obsolete-lemma-lqf-f-shriek-composition}
Soient $f : X \to Y$ et $g : Y \to Z$ des morphismes de schémas
localement quasi-finis et composables. Alors
$g_! \circ f_! = (g \circ f)_!$ et
$f^! \circ g^! = (g \circ f)^!$.
\end{lemma}

\begin{proof}
Cela résulte de la combinaison des lemmes de Compléments de cohomologie étale
\ref{more-etale-lemma-lqf-shriek-composition} et
\ref{more-etale-lemma-upper-shriek-restriction}.
\end{proof}
\section{Algèbre différentielle graduée}
\label{obsolete-section-dga}


\begin{lemma}
\label{obsolete-lemma-P-not-preserved-base-change}
Soient $(A, \text{d})$ et $(B, \text{d})$ des algèbres différentielles graduées.
Soit $N$ un $(A, B)$-bimodule différentiel gradué possédant la propriété
(P). Soit $M$ un $A$-module différentiel gradué possédant la propriété (P).
Alors $Q = M \otimes_A N$ est un $B$-module différentiel gradué qui représente
$M \otimes_A^\mathbf{L} N$ dans $D(B)$ et qui admet une filtration
$$
0 = F_{-1}Q \subset F_0Q \subset F_1Q \subset \ldots \subset Q
$$
par des sous-modules différentiels gradués telle que $Q = \bigcup F_pQ$,
que les inclusions $F_iQ \to F_{i + 1}Q$ soient des monomorphismes admissibles,
et que les quotients $F_{i + 1}Q/F_iQ$ soient isomorphes, comme
$B$-modules différentiels gradués, à une somme directe de modules de la forme
$(A \otimes_R B)[k]$.
\end{lemma}

\begin{proof}
Choisissons des filtrations $F_\bullet$ sur $M$ et $N$. Considérons alors la filtration
de $Q = M \otimes_A N$ définie par
$$
F_n(Q) = \sum\nolimits_{i + j = n}
F_i(M) \otimes_A F_j(N)
$$
Il s'agit manifestement d'un $B$-sous-module différentiel gradué. On a
$$
F_n(Q)/F_{n - 1}(Q) =
\bigoplus\nolimits_{i + j = n}
F_i(M)/F_{i - 1}(M) \otimes_A F_j(N)/F_{j - 1}(N)
$$
par exemple parce que la filtration de $M$ est scindée dans la catégorie
des $A$-modules gradués. Par hypothèse, les quotients du membre de droite
sont isomorphes à des sommes directes de décalés de $A$ et de
$A \otimes_R B$; comme
$A \otimes_A (A \otimes_R B) = A \otimes_R B$,
on en conclut que le membre de gauche est une somme directe de décalés
de $A \otimes_R B$ comme $B$-module différentiel gradué.
(Attention : $Q$ n'est pas muni d'une structure de $(A, B)$-bimodule.)
Cela établit l'assertion relative à la filtration.
L'assertion relative au produit tensoriel dérivé résulte immédiatement
de la définition du foncteur dans
Algèbre différentielle graduée, lemme \ref{dga-lemma-derived-bc}.
\end{proof}













\section{Méthodes simpliciales}
\label{obsolete-section-simplicial}


\begin{lemma}
\label{obsolete-lemma-equiv}
Sous les hypothèses et avec les notations de
Méthodes simpliciales, lemme \ref{simplicial-lemma-section},
il existe une section $g : U \to V$ du morphisme $f$, et
le composé $g \circ f$ est homotope à l'identité de $V$.
En particulier, le morphisme $f$ est une équivalence d'homotopie.
\end{lemma}

\begin{proof}
Cela résulte immédiatement de Méthodes simpliciales, lemmes
\ref{simplicial-lemma-section} et
\ref{simplicial-lemma-trivial-kan-homotopy}.
\end{proof}

\begin{lemma}
\label{obsolete-lemma-cosk-hom-deltak}
Soit $\mathcal{C}$ une catégorie admettant des sommes finies
et des limites finies. Soit $X$ un objet de $\mathcal{C}$.
Soit $k \geq 0$. Le morphisme canonique
$$
\Hom(\Delta[k], X)
\longrightarrow
\text{cosk}_1 \text{sk}_1 \Hom(\Delta[k], X)
$$
est un isomorphisme.
\end{lemma}

\begin{proof}
Pour tout objet simplicial $V$, on a
\begin{eqnarray*}
\Mor(V, \text{cosk}_1 \text{sk}_1 \Hom(\Delta[k], X))
& = &
\Mor(\text{sk}_1 V, \text{sk}_1 \Hom(\Delta[k], X)) \\
& = &
\Mor(i_{1!} \text{sk}_1 V, \Hom(\Delta[k], X)) \\
& = &
\Mor(i_{1!} \text{sk}_1 V \times \Delta[k], X)
\end{eqnarray*}
La première égalité résulte de l'adjonction entre $\text{sk}$ et $\text{cosk}$,
la deuxième de l'adjonction entre $i_{1!}$ et $\text{sk}_1$, et
la troisième de
Méthodes simpliciales, définition \ref{simplicial-definition-hom-from-simplicial-set},
où le dernier $X$ désigne l'objet simplicial constant de valeur $X$.
D'après Méthodes simpliciales, lemme \ref{simplicial-lemma-augmentation-howto},
un élément de cet ensemble ne dépend que des composantes de degrés $0$ et $1$
de $i_{1!} \text{sk}_1 V \times \Delta[k]$. Celles-ci
coïncident avec les composantes de degrés $0$ et $1$ de
$V \times \Delta[k]$; voir
Méthodes simpliciales, lemme \ref{simplicial-lemma-recovering-U-for-real}.
L'ensemble ci-dessus s'identifie donc à
$\Mor(V \times \Delta[k], X) = \Mor(V, \Hom(\Delta[k], X))$.
\end{proof}

\begin{lemma}
\label{obsolete-lemma-cosk0-hom-deltak}
Soit $\mathcal{C}$ une catégorie. Soit $X$ un objet de $\mathcal{C}$
dont les puissances cartésiennes $X \times \ldots \times X$ existent.
Soit $k \geq 0$, et soit $C[k]$ l'objet défini dans
Méthodes simpliciales, exemple \ref{simplicial-example-simplex-cosimplicial-set}.
Avec les notations de
Méthodes simpliciales, lemme \ref{simplicial-lemma-morphism-into-product},
le morphisme canonique
$$
\Hom(C[k], X)_1
\longrightarrow
(\text{cosk}_0 \text{sk}_0 \Hom(C[k], X))_1
$$
s'identifie à l'application
$$
\prod\nolimits_{\alpha : [k] \to [1]} X
\longrightarrow
X \times X
$$
qui est la projection sur les facteurs correspondant aux applications
constantes $\alpha$.
\end{lemma}

\begin{proof}
C'est établi dans la démonstration de
Hyperrecouvrements, lemme \ref{hypercovering-lemma-covering}.
\end{proof}



\section{Résultats sur les schémas}
\label{obsolete-section-devissage}

\noindent
Lemmes qui paraissent superflus.

\begin{lemma}
\label{obsolete-lemma-stein-projective}
Soit $(R, \mathfrak m, \kappa)$ un anneau local.
Soit $X \subset \mathbf{P}^n_R$ un sous-schéma fermé.
Supposons que $R = \Gamma(X, \mathcal{O}_X)$. Alors la fibre spéciale
$X_\kappa$ est géométriquement connexe.
\end{lemma}

\begin{proof}
C'est un cas particulier de
Compléments sur les morphismes, théorème
\ref{more-morphisms-theorem-stein-factorization-general}.
\end{proof}

\begin{lemma}
\label{obsolete-lemma-property-irreducible-higher-rank}
Soit $X$ un schéma noethérien.
Soit $Z_0 \subset X$ une partie fermée irréductible de point générique $\xi$.
Soit $\mathcal{P}$ une propriété des faisceaux cohérents sur $X$ telle que
\begin{enumerate}
\item Pour toute suite exacte courte de faisceaux cohérents, si deux
des trois termes possèdent la propriété $\mathcal{P}$, alors le troisième
la possède aussi.
\item Si une somme directe de deux faisceaux cohérents possède la propriété
$\mathcal{P}$, alors chacun des deux la possède.
\item Pour tout sous-schéma fermé intègre $Z \subset Z_0 \subset X$,
$Z \not = Z_0$, et tout faisceau quasi-cohérent d'idéaux
$\mathcal{I} \subset \mathcal{O}_Z$, la propriété $\mathcal{P}$ est satisfaite
par $(Z \to X)_*\mathcal{I}$.
\item Il existe un faisceau cohérent $\mathcal{G}$ sur $X$ tel que
\begin{enumerate}
\item $\text{Supp}(\mathcal{G}) = Z_0$,
\item $\mathcal{G}_\xi$ soit annulé par $\mathfrak m_\xi$, et
\item la propriété $\mathcal{P}$ soit satisfaite par $\mathcal{G}$.
\end{enumerate}
\end{enumerate}
Alors la propriété $\mathcal{P}$ est satisfaite par tout faisceau cohérent
$\mathcal{F}$ sur $X$ dont le support est contenu dans $Z_0$.
\end{lemma}

\begin{proof}
La démonstration est une variante de celle de
Cohomologie des schémas, lemme \ref{coherent-lemma-property-irreducible}.
Exactement comme dans cette démonstration, on voit que
tout faisceau cohérent dont le support est strictement contenu dans $Z_0$
possède la propriété $\mathcal{P}$.

\medskip\noindent
Considérons un faisceau cohérent $\mathcal{G}$ comme en (4).
D'après Cohomologie des schémas, lemme \ref{coherent-lemma-prepare-filter-irreducible},
il existe un faisceau d'idéaux $\mathcal{I}$ sur $Z_0$ et
une suite exacte courte
$$
0 \to
\left((Z_0 \to X)_*\mathcal{I}\right)^{\oplus r} \to
\mathcal{G} \to
\mathcal{Q} \to 0
$$
où le support de $\mathcal{Q}$ est strictement contenu dans $Z_0$.
En particulier, $r > 0$ et $\mathcal{I}$ est non nul,
car le support de $\mathcal{G}$ est égal à $Z_0$.
Puisque $\mathcal{Q}$ possède la propriété $\mathcal{P}$, on en conclut que
$\left((Z_0 \to X)_*\mathcal{I}\right)^{\oplus r}$
possède également la propriété $\mathcal{P}$.
D'après (2), on en déduit que la propriété $\mathcal{P}$ est satisfaite par
$(Z_0 \to X)_*\mathcal{I}$. En insérant ce fait au point voulu dans la démonstration de
Cohomologie des schémas, lemme \ref{coherent-lemma-property-irreducible},
on obtient le lemme.
Nous omettons quelques détails.
\end{proof}

\begin{lemma}
\label{obsolete-lemma-property-higher-rank}
Soit $X$ un schéma noethérien.
Soit $\mathcal{P}$ une propriété des faisceaux cohérents sur $X$ telle que
\begin{enumerate}
\item Pour toute suite exacte courte de faisceaux cohérents, si deux
des trois termes possèdent la propriété $\mathcal{P}$, alors le troisième
la possède aussi.
\item Si une somme directe de deux faisceaux cohérents possède la propriété
$\mathcal{P}$, alors chacun des deux la possède.
\item Pour tout sous-schéma fermé intègre $Z \subset X$
de point générique $\xi$, il existe
un faisceau cohérent $\mathcal{G}$ tel que
\begin{enumerate}
\item $\text{Supp}(\mathcal{G}) = Z$,
\item $\mathcal{G}_\xi$ soit annulé par $\mathfrak m_\xi$, et
\item la propriété $\mathcal{P}$ soit satisfaite par $\mathcal{G}$.
\end{enumerate}
\end{enumerate}
Alors la propriété $\mathcal{P}$ est satisfaite par tout faisceau cohérent sur $X$.
\end{lemma}

\begin{proof}
Cela résulte du lemme \ref{obsolete-lemma-property-irreducible-higher-rank}
exactement de la même manière que
Cohomologie des schémas, lemme \ref{coherent-lemma-property} résulte de
Cohomologie des schémas, lemme \ref{coherent-lemma-property-irreducible}.
\end{proof}

\begin{lemma}
\label{obsolete-lemma-section-maps-back-into}
Soit $X$ un schéma.
Soit $\mathcal{L}$ un $\mathcal{O}_X$-module inversible.
Soit $s \in \Gamma(X, \mathcal{L})$ une section.
Soient $\mathcal{F}' \subset \mathcal{F}$ des
$\mathcal{O}_X$-modules quasi-cohérents. Supposons que
\begin{enumerate}
\item $X$ soit quasi-compact,
\item $\mathcal{F}$ soit de type fini, et
\item $\mathcal{F}'|_{X_s} = \mathcal{F}|_{X_s}$.
\end{enumerate}
Alors il existe un entier $n \geq 0$ tel que
la multiplication par $s^n$ sur $\mathcal{F}$ se factorise
par $\mathcal{F}'$.
\end{lemma}

\begin{proof}
Autrement dit, nous affirmons que
$s^n\mathcal{F} \subset
\mathcal{F}' \otimes_{\mathcal{O}_X} \mathcal{L}^{\otimes n}$
pour un certain $n \geq 0$. Autrement dit, nous affirmons que le morphisme
quotient $\mathcal{F} \to \mathcal{F}/\mathcal{F}'$ devient
nul après multiplication par une puissance de $s$.
Cela résulte de Propriétés, lemme
\ref{properties-lemma-section-maps-backwards}.
\end{proof}

\begin{lemma}
\label{obsolete-lemma-bound-degree-in-nbhd-generic-point}
Soit $f : X \to Y$ un morphisme de schémas. Supposons que
\begin{enumerate}
\item $X$ et $Y$ soient des schémas intègres,
\item $f$ soit localement de type fini et dominant,
\item $f$ soit quasi-compact ou séparé,
\item $f$ soit génériquement fini, c'est-à-dire que l'une des conditions
(1) -- (5) de Morphismes, lemme \ref{morphisms-lemma-finite-degree}, soit satisfaite.
\end{enumerate}
Alors il existe un ouvert non vide $V \subset Y$ tel que
$f^{-1}(V) \to V$ soit fini localement libre de degré $\deg(X/Y)$.
En particulier, les degrés des fibres de $f^{-1}(V) \to V$
sont majorés par $\deg(X/Y)$.
\end{lemma}

\begin{proof}
On peut choisir $V$ de sorte que $f^{-1}(V) \to V$ soit fini.
Quitte à restreindre $V$, on peut alors supposer que $f^{-1}(V) \to V$
est plat et de présentation finie, d'après le théorème de platitude générique
(Morphismes, proposition \ref{morphisms-proposition-generic-flatness}).
Le morphisme est alors fini localement libre d'après
Morphismes, lemme \ref{morphisms-lemma-finite-flat}.
Comme $V$ est irréductible, le morphisme est de degré constant.
La dernière assertion en résulte avec
Morphismes, lemme \ref{morphisms-lemma-finite-locally-free-universally-bounded}.
\end{proof}





\section{Catégories dérivées des variétés}
\label{obsolete-section-equiv}

\noindent
Un lemme qui figurait initialement dans le chapitre consacré aux catégories
dérivées des variétés, mais qui n'est désormais plus nécessaire.

\begin{lemma}
\label{obsolete-lemma-preserves-Coh}
Soit $k$ un corps. Soit $X$ un schéma séparé de type fini sur $k$ qui est
régulier. Soit $F : D_{perf}(\mathcal{O}_X) \to D_{perf}(\mathcal{O}_X)$
un foncteur exact $k$-linéaire. Supposons que, pour tout
$\mathcal{O}_X$-module cohérent $\mathcal{F}$ tel que
$\dim(\text{Supp}(\mathcal{F})) = 0$, il existe un isomorphisme d'espaces
vectoriels sur $k$
$$
\Hom_X(\mathcal{F}, M) = \Hom_X(\mathcal{F}, F(M))
$$
fonctoriel en $M$ dans $D_{perf}(\mathcal{O}_X)$. Alors il existe un
automorphisme $f : X \to X$ sur $k$ qui induit l'identité sur l'espace
sous-jacent\footnote{Cela force souvent $f$ à être l'identité ;
voir Variétés, Lemme \ref{varieties-lemma-automorphism}.}
et un $\mathcal{O}_X$-module inversible $\mathcal{L}$ tels que $F$ et
$F'(M) = f^*M \otimes_{\mathcal{O}_X}^\mathbf{L} \mathcal{L}$ soient
apparentés.
\end{lemma}

\begin{proof}
Le Lemme \ref{equiv-lemma-duality-at-point} de Catégories dérivées des
variétés montre que, pour tout $\mathcal{O}_X$-module cohérent
$\mathcal{F}$ dont le support est un point fermé, il existe des
isomorphismes
$$
H^0(X, M \otimes^\mathbf{L}_{\mathcal{O}_X} \mathcal{F}) =
H^0(X, F(M) \otimes^\mathbf{L}_{\mathcal{O}_X} \mathcal{F})
$$
fonctoriels en $M$.

\medskip\noindent
Soit $x \in X$ un point fermé et appliquons ce qui précède à
$\mathcal{F} = \mathcal{O}_x$, le faisceau gratte-ciel de fibre
$\kappa(x)$ en $x$. On obtient
$$
\dim_{\kappa(x)}
\text{Tor}^{\mathcal{O}_{X, x}}_p(M_x, \kappa(x)) =
\dim_{\kappa(x)}
\text{Tor}^{\mathcal{O}_{X, x}}_p(F(M)_x, \kappa(x))
$$
pour tout $p \in \mathbf{Z}$. En particulier, si $H^i(M) = 0$ pour
$i > 0$, alors $H^i(F(M)) = 0$ pour $i > 0$, d'après le Lemme
\ref{equiv-lemma-orthogonal-point-sheaf} de Catégories dérivées des
variétés.

\medskip\noindent
Si $\mathcal{E}$ est localement libre de rang $r$, alors
$F(\mathcal{E})$ est localement libre de rang $r$. En effet, un complexe
parfait $K$ sur $\mathcal{O}_{X, x}$ tel que
$$
\dim_{\kappa(x)} \text{Tor}^{\mathcal{O}_{X, x}}_i(K, \kappa(x)) =
\left\{
\begin{matrix}
r & \text{si} & i = 0 \\
0 & \text{si} & i \not = 0
\end{matrix}
\right.
$$
est isomorphe à un module libre de rang $r$ placé en degré $0$. Voir, par
exemple, Algèbre, compléments, Lemme
\ref{more-algebra-lemma-lift-perfect-from-residue-field}.

\medskip\noindent
Si $M$ est supporté sur un sous-schéma fermé $Z \subset X$, alors $F(M)$
est également supporté sur $Z$. En effet,
$M \otimes_{\mathcal{O}_X}^\mathbf{L} \mathcal{O}_x = 0$ pour
$x \not \in Z$ ; il en va donc de même pour $F(M)$, et la conclusion
résulte du Lemme \ref{equiv-lemma-orthogonal-point-sheaf} de Catégories
dérivées des variétés.

\medskip\noindent
En particulier, $F(\mathcal{O}_x)$ est supporté sur $\{x\}$. Soit
$i \in \mathbf{Z}$ le plus petit entier tel que
$H^i(F(\mathcal{O}_x)) \not = 0$. On sait que $i \leq 0$. Si $i < 0$,
il existe alors un morphisme
$\mathcal{O}_x[-i] \to F(\mathcal{O}_x)$, ce qui contredit le fait que
tous les morphismes $\mathcal{O}_x[-i] \to \mathcal{O}_x$ sont nuls.
Ainsi $F(\mathcal{O}_x) = \mathcal{H}[0]$, où $\mathcal{H}$ est un
faisceau gratte-ciel en $x$.

\medskip\noindent
Soit $\mathcal{G}$ un $\mathcal{O}_X$-module cohérent tel que
$\dim(\text{Supp}(\mathcal{G})) = 0$. Il existe alors une filtration
$$
0 = \mathcal{G}_0 \subset \mathcal{G}_1 \subset \ldots \subset
\mathcal{G}_n = \mathcal{G}
$$
telle que, pour $n \geq i \geq 1$, le quotient
$\mathcal{G}_i/\mathcal{G}_{i - 1}$ soit isomorphe à
$\mathcal{O}_{x_i}$ pour un certain point fermé $x_i \in X$. On obtient
alors des triangles distingués
$$
F(\mathcal{G}_{i - 1}) \to F(\mathcal{G}_i) \to F(\mathcal{O}_{x_i})
$$
et une récurrence montre que $F(\mathcal{G}_i)$ est un faisceau cohérent
placé en degré $0$.

\medskip\noindent
Soit $\mathcal{G}$ un $\mathcal{O}_X$-module cohérent. Nous savons que
$H^i(F(\mathcal{G})) = 0$ pour $i > 0$. Supposons, en vue d'une
contradiction, que $H^i(F(\mathcal{G}))$ soit non nul pour un certain
$i < 0$. Choisissons $i$ minimal avec cette propriété, de sorte que nous
disposons d'un morphisme
$H^i(F(\mathcal{G}))[-i] \to F(\mathcal{G})$ dans
$D_{perf}(\mathcal{O}_X)$. Choisissons un point fermé $x \in X$ dans le
support de $H^i(F(\mathcal{G}))$. D'après Algèbre, compléments, Lemme
\ref{more-algebra-lemma-kollar-kovacs-pseudo-coherent}, il existe un
$n > 0$ tel que
$$
H^i(F(\mathcal{G}))_x \otimes_{\mathcal{O}_{X, x}}
\mathcal{O}_{X, x}/\mathfrak m_x^n
\longrightarrow
\text{Tor}^{\mathcal{O}_{X, x}}_{-i}(F(\mathcal{G})_x,
\mathcal{O}_{X, x}/\mathfrak m_x^n)
$$
soit non nul. Prenons ensuite $m \geq 1$ et considérons la suite exacte
courte
$$
0 \to \mathfrak m_x^m \mathcal{G} \to \mathcal{G} \to
\mathcal{G}/\mathfrak m_x^m\mathcal{G} \to 0
$$
D'après ce qui précède,
$F(\mathcal{G}/\mathfrak m_x^m\mathcal{G})$ est un faisceau placé en
degré $0$. Par conséquent,
$H^i(F(\mathfrak m_x^m \mathcal{G})) \to H^i(F(\mathcal{G}))$ est un
isomorphisme. Considérons le diagramme commutatif
$$
\xymatrix{
H^i(F(\mathfrak m_x^m\mathcal{G}))_x \otimes_{\mathcal{O}_{X, x}}
\mathcal{O}_{X, x}/\mathfrak m_x^n \ar[r] \ar[d] &
\text{Tor}^{\mathcal{O}_{X, x}}_{-i}(F(\mathfrak m_x^m\mathcal{G})_x,
\mathcal{O}_{X, x}/\mathfrak m_x^n) \ar[d] \\
H^i(F(\mathcal{G}))_x \otimes_{\mathcal{O}_{X, x}}
\mathcal{O}_{X, x}/\mathfrak m_x^n \ar[r] &
\text{Tor}^{\mathcal{O}_{X, x}}_{-i}(F(\mathcal{G})_x,
\mathcal{O}_{X, x}/\mathfrak m_x^n)
}
$$
Puisque la flèche verticale de gauche est un isomorphisme et que la flèche
horizontale du bas est non nulle, la flèche verticale de droite est non
nulle pour tout $m \geq 1$. D'autre part, le premier paragraphe de la
démonstration montre que cette flèche s'identifie à la flèche
$$
\text{Tor}^{\mathcal{O}_{X, x}}_{-i}(\mathfrak m_x^m\mathcal{G}_x,
\mathcal{O}_{X, x}/\mathfrak m_x^n)
\longrightarrow
\text{Tor}^{\mathcal{O}_{X, x}}_{-i}(\mathcal{G}_x,
\mathcal{O}_{X, x}/\mathfrak m_x^n)
$$
Cependant, cette flèche est nulle pour $m \gg n$ d'après Algèbre,
compléments, Lemme \ref{more-algebra-lemma-tor-annihilated}, ce qui donne
la contradiction cherchée.

\medskip\noindent
Nous savons donc que $F$ préserve les modules cohérents. D'après le Lemme
\ref{equiv-lemma-exact-functor-preserving-Coh} de Catégories dérivées des
variétés, $F$ est apparenté au foncteur de Fourier--Mukai $F'$ défini par
un $\mathcal{O}_{X \times X}$-module cohérent $\mathcal{K}$, plat sur
$X$ via $\text{pr}_1$ et fini sur $X$ via $\text{pr}_2$. Comme
$F(\mathcal{O}_X)$ est un $\mathcal{O}_X$-module inversible
$\mathcal{L}$ placé en degré $0$, on a
$$
\mathcal{L} \cong F(\mathcal{O}_X) \cong F'(\mathcal{O}_X) \cong
\text{pr}_{2, *}\mathcal{K}
$$
Le Lemme \ref{functors-lemma-pushforward-invertible-pre} de Foncteurs et
morphismes fournit donc un morphisme $s : X \to X \times X$ tel que
$\text{pr}_2 \circ s = \text{id}_X$ et $\mathcal{K} = s_*\mathcal{L}$.
Posons $f = \text{pr}_1 \circ s$. On a alors
\begin{align*}
F'(M)
& =
R\text{pr}_{2, *}(L\text{pr}_1^*M \otimes \mathcal{K}) \\
& =
R\text{pr}_{2, *}(L\text{pr}_1^*M \otimes s_*\mathcal{L}) \\
& =
R\text{pr}_{2, *}(Rs_*(Lf^*M \otimes \mathcal{L})) \\
& =
Lf^*M \otimes \mathcal{L}
\end{align*}
Dans la troisième étape, nous avons utilisé Catégories dérivées des
schémas, Lemme \ref{perfect-lemma-cohomology-base-change}.
Comme, pour tout point fermé $x \in X$, le module
$F(\mathcal{O}_x)$ est supporté en $x$, on voit que $f$ induit l'identité
sur l'espace sous-jacent de $X$. Il reste à montrer que $f$
est un isomorphisme ; c'est l'objet du paragraphe suivant.

\medskip\noindent
Soit $x \in X$ un point fermé. Pour $n \geq 1$, notons
$\mathcal{O}_{x, n}$ le faisceau gratte-ciel en $x$ de fibre
$\mathcal{O}_{X, x}/\mathfrak m_x^n$. Nous avons
$$
\Hom_X(\mathcal{O}_{x, m}, \mathcal{O}_{x, n}) \cong
\Hom_X(\mathcal{O}_{x, m}, F(\mathcal{O}_{x, n})) \cong
\Hom_X(\mathcal{O}_{x, m}, f^*\mathcal{O}_{x, n} \otimes \mathcal{L})
$$
fonctoriellement par rapport aux homomorphismes de $\mathcal{O}_X$-modules
entre les $\mathcal{O}_{x, n}$. (Le premier isomorphisme existe par
hypothèse, et le second parce que $F$ et $F'$ sont apparentés.) Pour
$m \geq n$, l'action sur $\mathcal{O}_{x, n}$ donne
$\mathcal{O}_{X, x}/\mathfrak m_x^n =
\Hom_X(\mathcal{O}_{x, m}, \mathcal{O}_{x, n})$ ; on en déduit que
$f^\sharp : \mathcal{O}_{X, x}/\mathfrak m_x^n \to
\mathcal{O}_{X, x}/\mathfrak m_x^n$ est bijectif pour tout $n$. Ainsi,
$f$ induit des isomorphismes sur les anneaux locaux complétés aux points
fermés et est donc étale (Morphismes étales, Lemme
\ref{etale-lemma-characterize-etale-completions}). En considérant les
points fermés, on voit que
$\Delta_f : X \to X \times_{f, X, f} X$ (qui est une immersion ouverte
puisque $f$ est étale) est bijective, donc est un isomorphisme. Par suite,
$f$ est un monomorphisme. Enfin, Descente, Lemme
\ref{descent-lemma-flat-surjective-quasi-compact-monomorphism-isomorphism}
montre que $f$ est un isomorphisme.
\end{proof}








\section{Représentabilité dans le cas propre et régulier}
\label{obsolete-section-regular-proper}

\noindent
Cette section est obsolète : le Théorème
\ref{equiv-theorem-bondal-van-den-bergh} de Catégories dérivées des variétés
a été amélioré de sorte qu'il s'applique à tous les schémas propres sur un
corps (alors qu'auparavant nous ne l'avions démontré que pour les schémas
projectifs sur un corps).

\begin{lemma}
\label{obsolete-lemma-trace-map}
\begin{reference}
La démonstration donnée ici reprend l'argument de
\cite[remarque 3.4]{MS}.
\end{reference}
Soit $f : X' \to X$ un morphisme propre et birationnel entre schémas
noethériens intègres, avec $X$ régulier. Le morphisme
$\mathcal{O}_X \to Rf_*\mathcal{O}_{X'}$ admet canoniquement une rétraction
dans $D(\mathcal{O}_X)$.
\end{lemma}

\begin{proof}
Posons $E = Rf_*\mathcal{O}_{X'}$ dans $D(\mathcal{O}_X)$.
Remarquons que $E$ appartient à
$D^b_{\textit{Coh}}(\mathcal{O}_X)$ d'après
Catégories dérivées des schémas, Lemme
\ref{perfect-lemma-direct-image-coherent}.
D'après Catégories dérivées des schémas, Lemme
\ref{perfect-lemma-perfect-on-regular}, $E$ est donc un objet parfait de
$D(\mathcal{O}_X)$. Comme $\mathcal{O}_{X'}$ est un faisceau d'algèbres,
le cup-produit relatif fournit le morphisme
$\mu : E \otimes_{\mathcal{O}_X}^\mathbf{L} E \to E$ de Cohomologie,
Remarque \ref{cohomology-remark-cup-product}.
Soit $\sigma : E \otimes E^\vee \to E^\vee \otimes E$ la contrainte de
commutativité de la catégorie monoïdale symétrique $D(\mathcal{O}_X)$
(Cohomologie, Lemme \ref{cohomology-lemma-symmetric-monoidal-derived}).
Notons $\eta : \mathcal{O}_X \to E \otimes E^\vee$ et
$\epsilon : E^\vee \otimes E \to \mathcal{O}_X$ les morphismes
construits dans Cohomologie, Exemple \ref{cohomology-example-dual-derived}.
Nous pouvons alors considérer le morphisme
$$
E \xrightarrow{\eta \otimes 1} E \otimes E^\vee \otimes E
\xrightarrow{\sigma \otimes 1} E^\vee \otimes E \otimes E
\xrightarrow{1 \otimes \mu} E^\vee \otimes E
\xrightarrow{\epsilon} \mathcal{O}_X
$$
Nous affirmons que ce morphisme est un inverse à gauche du morphisme figurant
dans l'énoncé du lemme. Pour le vérifier, il suffit de montrer que le composé
$\mathcal{O}_X \to \mathcal{O}_X$ est l'identité. Il suffit de le vérifier au
point générique de $X$ (ou sur un sous-schéma ouvert de $X$ au-dessus duquel
$f$ est un isomorphisme). Dans ce cas, $E = \mathcal{O}_X$ et $\mu$ est le
morphisme de multiplication usuel, d'où l'assertion.
\end{proof}

\begin{lemma}
\label{obsolete-lemma-characterize-dbcoh-proper-regular}
Soit $X$ un schéma régulier, propre sur un corps $k$. Soit
$K \in \Ob(D_\QCoh(\mathcal{O}_X))$. Les assertions suivantes sont
équivalentes :
\begin{enumerate}
\item $K \in D^b_{\textit{Coh}}(\mathcal{O}_X) = D_{perf}(\mathcal{O}_X)$ ;
\item $\sum_{i \in \mathbf{Z}} \dim_k \Ext^i_X(E, K) < \infty$
pour tout objet parfait $E$ de $D(\mathcal{O}_X)$.
\end{enumerate}
\end{lemma}

\begin{proof}
L'égalité de (1) résulte de Catégories dérivées des schémas,
Lemme \ref{perfect-lemma-perfect-on-regular}.
L'implication (1) $\Rightarrow$ (2) résulte de
Catégories dérivées des variétés, Lemme \ref{equiv-lemma-finiteness}.
L'implication (2) $\Rightarrow$ (1) résulte de
Compléments sur les morphismes, Lemme
\ref{more-morphisms-lemma-characterize-relatively-perfect}.
\end{proof}

\begin{lemma}
\label{obsolete-lemma-bondal-van-den-bergh}
Soit $X$ un schéma régulier, propre sur un corps $k$.
\begin{enumerate}
\item Soit $F : D_{perf}(\mathcal{O}_X)^{opp} \to \text{Vect}_k$
un foncteur cohomologique $k$-linéaire tel que
$$
\sum\nolimits_{n \in \mathbf{Z}} \dim_k F(E[n]) < \infty
$$
pour tout $E \in D_{perf}(\mathcal{O}_X)$. Alors $F$ est isomorphe à un
foncteur de la forme $E \mapsto \Hom_X(E, K)$, où
$K \in D_{perf}(\mathcal{O}_X)$.
\item Soit $G : D_{perf}(\mathcal{O}_X) \to \text{Vect}_k$
un foncteur homologique $k$-linéaire tel que
$$
\sum\nolimits_{n \in \mathbf{Z}} \dim_k G(E[n]) < \infty
$$
pour tout $E \in D_{perf}(\mathcal{O}_X)$. Alors $G$ est isomorphe à un
foncteur de la forme $E \mapsto \Hom_X(K, E)$, où
$K \in D_{perf}(\mathcal{O}_X)$.
\end{enumerate}
\end{lemma}

\begin{proof}
Cela résulte de Catégories dérivées des variétés, Théorème
\ref{equiv-theorem-bondal-van-den-bergh}, et du
Lemme \ref{equiv-lemma-homological-representable}.
Nous donnons également une autre démonstration ci-dessous.

\medskip\noindent
Démonstration de (1). La catégorie dérivée $D_\QCoh(\mathcal{O}_X)$ admet les
sommes directes, est compactement engendrée, et
$D_{perf}(\mathcal{O}_X)$ en est la sous-catégorie pleine des objets compacts ;
voir Catégories dérivées des schémas, Lemme
\ref{perfect-lemma-quasi-coherence-direct-sums},
Théorème \ref{perfect-theorem-bondal-van-den-Bergh}, et
Proposition \ref{perfect-proposition-compact-is-perfect}.
Le lemme de Catégories dérivées des variétés
\ref{equiv-lemma-van-den-bergh} permet de supposer que
$F(E) = \Hom_X(E, K)$ pour un certain
$K \in \Ob(D_\QCoh(\mathcal{O}_X))$.
$K$ appartient alors à $D^b_{\textit{Coh}}(\mathcal{O}_X)$ d'après le
Lemme \ref{obsolete-lemma-characterize-dbcoh-proper-regular}.

\medskip\noindent
Démonstration de (2). Considérons le foncteur contravariant
$E \mapsto E^\vee$ sur $D_{perf}(\mathcal{O}_X)$ ; voir
Cohomologie, Lemme \ref{cohomology-lemma-dual-perfect-complex}.
C'est une anti-équivalence exacte de $D_{perf}(\mathcal{O}_X)$ sur elle-même.
On peut donc appliquer (1) au foncteur $F(E) = G(E^\vee)$ : il existe
$K \in D_{perf}(\mathcal{O}_X)$ tel que
$G(E^\vee) = \Hom_X(E, K)$.
Il s'ensuit que $G(E) = \Hom_X(E^\vee, K) = \Hom_X(K^\vee, E)$ ;
ainsi, $K^\vee$ convient.
\end{proof}








\section{Foncteur des quotients}
\label{obsolete-section-quotients}

\begin{lemma}
\label{obsolete-lemma-factors-through-quotient}
Soit $S = \Spec(R)$ un sch\'ema affine. Soit $X$ un espace alg\'ebrique sur
$S$. Soient $q_i : \mathcal{F} \to \mathcal{Q}_i$, $i = 1, 2$,
des morphismes surjectifs de $\mathcal{O}_X$-modules quasi-coh\'erents.
Supposons $\mathcal{Q}_1$ plat sur $S$. Soit $T \to S$ un morphisme quasi-compact
de sch\'emas tel qu'il existe une factorisation
$$
\xymatrix{
& \mathcal{F}_T \ar[rd]^{q_{2, T}} \ar[ld]_{q_{1, T}} \\
\mathcal{Q}_{1, T} & & \mathcal{Q}_{2, T} \ar@{..>}[ll]
}
$$
Il existe alors un sous-sch\'ema ferm\'e $Z \subset S$ tel que
(a) $T \to S$ se factorise par $Z$ et (b)
$q_{1, Z}$ se factorise par $q_{2, Z}$.
Si $\Ker(q_2)$ est un $\mathcal{O}_X$-module de type fini et si $X$
est quasi-compact, on peut choisir $Z \to S$ de pr\'esentation finie.
\end{lemma}

\begin{proof}
On applique Platitude sur les espaces, Lemme \ref{spaces-flat-lemma-F-zero-somewhat-closed},
au morphisme $\Ker(q_2) \to \mathcal{Q}_1$.
\end{proof}







\section{Espaces et recouvrements fpqc}
\label{obsolete-section-fpqc}

\noindent
Le contenu de cette section a \'et\'e rendu obsol\`ete par l'argument de Gabber, qui montre que
les espaces alg\'ebriques v\'erifient la condition de faisceau pour les
recouvrements fpqc. Voir
Propri\'et\'es des espaces, section \ref{spaces-properties-section-fpqc}.

\begin{lemma}
\label{obsolete-lemma-separated-fpqc}
Soit $S$ un sch\'ema. Soit $X$ un espace alg\'ebrique sur $S$.
Soit $\{f_i : T_i \to T\}_{i \in I}$ un recouvrement fpqc de sch\'emas sur $S$.
Alors l'application
$$
\Mor_S(T, X)
\longrightarrow
\prod\nolimits_{i \in I} \Mor_S(T_i, X)
$$
est injective.
\end{lemma}

\begin{proof}
C'est une cons\'equence imm\'ediate de la Proposition
\ref{spaces-properties-proposition-sheaf-fpqc}
de Propri\'et\'es des espaces.
\end{proof}

\begin{lemma}
\label{obsolete-lemma-sheaf-fpqc-open-covering}
Soit $S$ un sch\'ema. Soit $X$ un espace alg\'ebrique sur $S$.
Soit $X = \bigcup_{j \in J} X_j$ un recouvrement de Zariski; voir
Espaces, D\'efinition \ref{spaces-definition-Zariski-open-covering}.
Si chacun des $X_j$ v\'erifie la condition de faisceau pour la topologie fpqc,
alors $X$ v\'erifie la condition de faisceau pour la topologie fpqc.
\end{lemma}

\begin{proof}
Cela r\'esulte du fait que tous les espaces alg\'ebriques v\'erifient la condition
de faisceau pour la topologie fpqc; voir
Propri\'et\'es des espaces, Proposition
\ref{spaces-properties-proposition-sheaf-fpqc}.
\end{proof}

\begin{lemma}
\label{obsolete-lemma-sheaf-fpqc-quasi-separated}
Soit $S$ un sch\'ema. Soit $X$ un espace alg\'ebrique sur $S$.
Si $X$ est localement quasi-s\'epar\'e sur $S$ pour la topologie de Zariski,
alors $X$ v\'erifie la condition de faisceau pour la topologie fpqc.
\end{lemma}

\begin{proof}
C'est une cons\'equence imm\'ediate de la Proposition g\'en\'erale
\ref{spaces-properties-proposition-sheaf-fpqc}
de Propri\'et\'es des espaces.
\end{proof}

\begin{remark}
\label{obsolete-remark-proof-works-when}
Cette remarque examinait autrefois dans quelle mesure la d\'emonstration originelle du
Lemme \ref{obsolete-lemma-sheaf-fpqc-quasi-separated} (dat\'ee du 18 d\'ecembre 2009)
se g\'en\'eralise.
\end{remark}





\section{Espaces algébriques très raisonnables}
\label{obsolete-section-very-reasonable}

\noindent
Éléments quelque peu obsolètes.

\begin{lemma}
\label{obsolete-lemma-reasonable-kolmogorov}
Soit $S$ un schéma.
Soit $X$ un espace algébrique raisonnable au-dessus de $S$.
Alors $|X|$ est un espace de Kolmogorov (voir
Topologie, Définition \ref{topology-definition-generic-point}).
\end{lemma}

\begin{proof}
Cela résulte des définitions et du
Lemme \ref{decent-spaces-lemma-kolmogorov} d'Espaces décents.
\end{proof}

\noindent
Dans la suite de cette section, nous formulons quelques remarques sur les
espaces algébriques très raisonnables. S'il existe un schéma $U$ et un
morphisme surjectif, étale et quasi-compact
$U \to X$, alors $X$ est très raisonnable ; voir
Lemme \ref{decent-spaces-lemma-characterize-very-reasonable} d'Espaces décents.

\begin{lemma}
\label{obsolete-lemma-scheme-very-reasonable}
Tout schéma est très raisonnable.
\end{lemma}

\begin{proof}
Cela tient à ce que le morphisme identité est étale, quasi-compact
et surjectif.
\end{proof}

\begin{lemma}
\label{obsolete-lemma-very-reasonable-Zariski-local}
Soit $S$ un schéma.
Soit $X$ un espace algébrique au-dessus de $S$.
S'il existe un recouvrement ouvert de Zariski $X = \bigcup X_i$ tel que
chaque $X_i$ soit très raisonnable, alors $X$ est très raisonnable.
\end{lemma}

\begin{proof}
C'est le cas $(\epsilon)$ du
Lemme \ref{decent-spaces-lemma-properties-local} d'Espaces décents.
\end{proof}

\begin{lemma}
\label{obsolete-lemma-quasi-separated-very-reasonable}
Tout espace algébrique localement quasi-séparé pour la topologie de Zariski est très raisonnable.
En particulier, tout espace algébrique quasi-séparé est très raisonnable.
\end{lemma}

\begin{proof}
C'est l'une des implications du
Lemme \ref{decent-spaces-lemma-bounded-fibres} d'Espaces décents.
\end{proof}

\begin{lemma}
\label{obsolete-lemma-representable-very-reasonable}
Soit $S$ un schéma.
Soient $X$ et $Y$ des espaces algébriques au-dessus de $S$.
Soit $Y \to X$ un morphisme représentable.
Si $X$ est très raisonnable, alors $Y$ l'est aussi.
\end{lemma}

\begin{proof}
C'est le cas $(\epsilon)$ du
Lemme \ref{decent-spaces-lemma-representable-properties} d'Espaces décents.
\end{proof}

\begin{remark}
\label{obsolete-remark-very-reasonable-Zariski-locally-quasi-separated}
Les espaces algébriques très raisonnables forment une classe strictement plus
grande que celle des espaces algébriques localement quasi-séparés pour la
topologie de Zariski. Considérons un espace algébrique de la forme
$X = [U/G]$ (voir
Espaces, Définition \ref{spaces-definition-quotient}),
où $G$ est un groupe fini agissant sans points fixes sur un schéma
$U$ non quasi-séparé. En effet, dans ce cas,
$U \times_X U = U \times G$, et les deux projections vers $U$ sont
manifestement quasi-compactes ; ainsi, $X$ est très raisonnable. En revanche,
la diagonale $U \times_X U \to U \times U$ n'est pas quasi-compacte ; cet
espace algébrique n'est donc pas quasi-séparé. Prenons maintenant pour $U$
l'espace affine de dimension infinie sur un corps $k$ de caractéristique
$\not = 2$, dont l'origine est dédoublée ; voir
Schémas, Exemple \ref{schemes-example-not-quasi-separated}.
Notons $0_1, 0_2$ les deux origines de $U$. Soit $G = \{+1, -1\}$ et
faisons agir $-1$ par multiplication par $-1$ sur toutes les coordonnées,
tout en échangeant $0_1$ et $0_2$. Alors $[U/G]$ est très raisonnable, mais
n'est pas localement quasi-séparé pour la topologie de Zariski (détails omis).
\end{remark}

\noindent
Avertissement : le lemme suivant doit être utilisé avec prudence, car les
schémas $U_i$ qui y interviennent ne sont pas nécessairement séparés, ni même
quasi-séparés.

\begin{lemma}
\label{obsolete-lemma-very-reasonable-quasi-compact-pieces}
Soit $S$ un schéma.
Soit $X$ un espace algébrique très raisonnable au-dessus de $S$.
Il existe une famille de schémas
$U_i$ et des morphismes $U_i \to X$ tels que
\begin{enumerate}
\item chaque $U_i$ soit un schéma quasi-compact,
\item chaque $U_i \to X$ soit étale,
\item les deux projections $U_i \times_X U_i \to U_i$ soient quasi-compactes, et
\item le morphisme $\coprod U_i \to X$ soit surjectif (et étale).
\end{enumerate}
\end{lemma}

\begin{proof}
La Définition \ref{decent-spaces-definition-very-reasonable} d'Espaces décents
assure l'existence de $U_i \to X$ satisfaisant (2), (3) et (4).
Fixons $i$, posons $R_i = U_i \times_X U_i$ et notons $s, t : R_i \to U_i$
les projections.
Pour tout ouvert affine $W \subset U_i$, l'ouvert $W' = t(s^{-1}(W)) \subset U_i$
est quasi-compact et invariant par $R_i$ (voir
Groupoïdes, Lemme \ref{groupoids-lemma-constructing-invariant-opens}).
Ainsi, $W'$ est un schéma quasi-compact, $W' \to X$ est étale et
$W' \times_X W' = s^{-1}(W') = t^{-1}(W')$ ; les deux projections
$W' \times_X W' \to W'$ sont donc quasi-compactes. La famille des
$W' \to X$, où $W \subset U_i$ parcourt les membres de recouvrements ouverts
affines des $U_i$, possède par conséquent les propriétés voulues.
\end{proof}




\section{Lemmes obsolètes sur les espaces algébriques}
\label{obsolete-section-obsolete-on-spaces}

\noindent
Lemmes qui semblent superflus ou ne sont plus utilisés dans le texte.

\begin{lemma}
\label{obsolete-lemma-vanishing-surjective}
Dans la Situation \ref{spaces-cohomology-situation-vanishing} de Cohomologie des espaces,
le morphisme $p : X \to \Spec(A)$ est surjectif.
\end{lemma}

\begin{proof}
Ce lemme intervenait à l'origine dans la démonstration de la
Proposition \ref{spaces-cohomology-proposition-vanishing-affine} de
Cohomologie des espaces, mais il en est désormais une conséquence.
\end{proof}

\begin{lemma}
\label{obsolete-lemma-vanishing-universally-closed}
Dans la Situation \ref{spaces-cohomology-situation-vanishing} de Cohomologie des espaces,
le morphisme $p : X \to \Spec(A)$ est universellement fermé.
\end{lemma}

\begin{proof}
Ce lemme intervenait à l'origine dans la démonstration de la
Proposition \ref{spaces-cohomology-proposition-vanishing-affine} de
Cohomologie des espaces, mais il en est désormais une conséquence.
\end{proof}

\begin{remark}
\label{obsolete-remark-equation-first}
Ce tag renvoyait autrefois à une équation de la démonstration du
Lemme \ref{formal-spaces-lemma-iff-adic-star} d'Espaces formels.
\end{remark}

\begin{remark}
\label{obsolete-remark-equation-second}
Ce tag renvoyait autrefois à une équation de la démonstration du
Lemme \ref{formal-spaces-lemma-iff-adic-star} d'Espaces formels.
\end{remark}



\section{Lemmes obsolètes sur les champs algébriques}
\label{obsolete-section-obsolete-on-stacks}

\noindent
Lemmes qui semblent superflus ou ne sont plus utilisés dans le texte.

\begin{lemma}
\label{obsolete-lemma-infinite-sequence}
Soit $S$ un schéma localement noethérien. Soit $\mathcal{X}$ une catégorie
fibrée en groupoïdes au-dessus de $(\Sch/S)_{fppf}$ vérifiant (RS*).
Soit $x$ un objet de
$\mathcal{X}$ au-dessus d'un schéma affine $U$ de type fini sur $S$.
Soient $u_n \in U$, $n \geq 1$, des points de type fini deux à deux distincts
tels que $x$ ne soit versel en $u_n$ pour tout $n$. Après avoir remplacé
$u_n$ par une sous-suite, il existe des morphismes
$$
x \to x_1 \to x_2 \to \ldots
\quad\text{dans }\mathcal{X}\text{, au-dessus de}\quad
U \to U_1 \to U_2 \to \ldots
$$
sur $S$ tels que
\begin{enumerate}
\item pour tout $n$, le morphisme $U \to U_n$ est un épaississement
du premier ordre,
\item pour tout $n$, on a une suite exacte courte
$$
0 \to \kappa(u_n) \to \mathcal{O}_{U_n} \to \mathcal{O}_{U_{n - 1}} \to 0
$$
où $U_0 = U$ lorsque $n = 1$,
\item pour tout $n$, il n'existe {\bf aucun} couple $(W, \alpha)$
formé d'un voisinage ouvert $W \subset U_n$ de $u_n$
et d'un morphisme $\alpha : x_n|_W \to x$
tel que le composé
$$
x|_{U \cap W} \xrightarrow{\text{restriction de }x \to x_n}
x_n|_W \xrightarrow{\alpha} x
$$
soit le morphisme canonique $x|_{U \cap W} \to x$.
\end{enumerate}
\end{lemma}

\begin{proof}
À l'origine, ce lemme intervenait dans la preuve d'un critère
d'ouverture de la versalité
(Axiomes d'Artin, lemme \ref{artin-lemma-SGE-implies-openness-versality}),
mais il a été remplacé par Axiomes d'Artin,
lemme \ref{artin-lemma-infinite-sequence-pre}, dont il découle immédiatement.
En effet, quitte à remplacer $u_n$, $n \geq 1$, par une sous-suite,
nous pouvons supposer qu'il n'existe aucune relation de spécialisation
entre ces points ; voir Properties,
lemme \ref{properties-lemma-thin-infinite-sequence}.
Il suffit alors d'appliquer Axiomes d'Artin,
lemme \ref{artin-lemma-infinite-sequence-pre} pour achever la preuve.
\end{proof}




\section{Variantes du complexe cotangent pour les schémas}
\label{obsolete-section-cotangent-schemes-variant}

\noindent
Cette section donne une autre construction du complexe cotangent
d'un morphisme de schémas. Elle figure actuellement dans le chapitre
obsolète, car, pour les applications, la version plus élémentaire étudiée dans
Complexe cotangent, section \ref{cotangent-section-cotangent-schemes-variant}
suffit.

\medskip\noindent
Soit $f : X \to Y$ un morphisme de schémas. Soit $\mathcal{C}_{X/Y}$ la
catégorie dont les objets sont les diagrammes commutatifs
\begin{equation}
\label{obsolete-equation-object}
\vcenter{
\xymatrix{
X \ar[d] & U \ar[l] \ar[d] \ar[r]_i & A \ar[ld] \\
Y & V \ar[l]
}
}
\end{equation}
de schémas tels que
\begin{enumerate}
\item $U$ soit un sous-schéma ouvert de $X$,
\item $V$ soit un sous-schéma ouvert de $Y$, et
\item il existe un isomorphisme $A = V \times \Spec(P)$ au-dessus de $V$,
où $P$ est une algèbre de polynômes sur $\mathbf{Z}$ (en un ensemble
quelconque de variables).
\end{enumerate}
Autrement dit, $A$ est un espace affine (de dimension éventuellement infinie)
sur $V$. Les morphismes sont donnés par des diagrammes commutatifs.

\medskip\noindent
{\bf Notation.} Un objet de $\mathcal{C}_{X/Y}$, c'est-à-dire un diagramme
(\ref{obsolete-equation-object}), sera souvent noté $U \to A$, étant entendu que
(a) $U$ est un sous-schéma ouvert de $X$, (b) $U \to A$ est un morphisme
au-dessus de $Y$, (c) l'image du morphisme structural $A \to Y$ est un ouvert
$V \subset Y$, et (d) le morphisme $A \to V$ présente sa source comme un
espace affine sur sa base.
Nous écrirons $U \to A/V$ pour indiquer que $V \subset Y$ est l'image de
$A \to Y$.
Rappelons que $X_{Zar}$ désigne le petit site de Zariski de $X$.
Il existe des foncteurs d'oubli
$$
\mathcal{C}_{X/Y} \to X_{Zar},\ (U \to A) \mapsto U
\quad\text{et}\quad
\mathcal{C}_{X/Y} \to Y_{Zar},\ (U \to A/V) \mapsto V.
$$

\begin{lemma}
\label{obsolete-lemma-category-fibred}
Soit $X \to Y$ un morphisme de schémas.
\begin{enumerate}
\item La catégorie $\mathcal{C}_{X/Y}$ est fibrée sur $X_{Zar}$.
\item La catégorie $\mathcal{C}_{X/Y}$ est fibrée sur $Y_{Zar}$.
\item La catégorie $\mathcal{C}_{X/Y}$ est fibrée sur la
catégorie des couples $(U, V)$, où $U \subset X$ et $V \subset Y$ sont
ouverts et où $f(U) \subset V$.
\end{enumerate}
\end{lemma}

\begin{proof}
Pour (1). Étant donnés un objet $U \to A$ de $\mathcal{C}_{X/Y}$ et un
morphisme $U' \to U$ de $X_{Zar}$, considérons l'objet $i' : U' \to A$ de
$\mathcal{C}_{X/Y}$, où $i'$ est le composé de $i$ et de $U' \to U$.
Le morphisme $(U' \to A) \to (U \to A)$ de $\mathcal{C}_{X/Y}$
est fortement cartésien au-dessus de $X_{Zar}$.

\medskip\noindent
Pour (2). Étant donnés un objet $U \to A/V$ et un morphisme $V' \to V$,
posons $U' = U \cap f^{-1}(V')$ et $A' = V' \times_V A$; on obtient ainsi
un morphisme fortement cartésien $(U' \to A') \to (U \to A)$ au-dessus de
$V' \to V$.

\medskip\noindent
Pour (3). Notons $(X/Y)_{Zar}$ la catégorie de (3). Étant donnés
$U \to A/V$ et un morphisme $(U', V') \to (U, V)$ dans $(X/Y)_{Zar}$,
considérons $A' = V' \times_V A$. Alors le morphisme
$(U' \to A'/V') \to (U \to A/V)$ est fortement cartésien dans
$\mathcal{C}_{X/Y}$ au-dessus de $(X/Y)_{Zar}$.
\end{proof}

\noindent
Nous obtenons une topologie $\tau_X$ sur $\mathcal{C}_{X/Y}$ en
prenant la topologie induite par celle de $X_{Zar}$ (voir
Champs, section \ref{stacks-section-topology}). Sauf mention contraire,
c'est cette topologie sur $\mathcal{C}_{X/Y}$ que nous considérerons.
Précisément, une famille de morphismes
$\{(U_i \to A_i) \to (U \to A)\}$ est couvrante dans
$\mathcal{C}_{X/Y}$ si et seulement si
\begin{enumerate}
\item $U = \bigcup U_i$, et
\item $A_i = A$ pour tout $i$.
\end{enumerate}
On obtient la même catégorie de faisceaux si l'on autorise
$A_i \cong A$ dans (2). Le foncteur d'oubli vers $X_{Zar}$ définit un
morphisme de topos
$\pi : \Sh(\mathcal{C}_{X/Y}) \to \Sh(X_{Zar})$.

\medskip\noindent
Le site $\mathcal{C}_{X/Y}$ est muni de plusieurs faisceaux d'anneaux.
\begin{enumerate}
\item Le faisceau $\mathcal{O}$ défini par la règle
$(U \to A) \mapsto \mathcal{O}(A)$.
\item Le faisceau $\underline{\mathcal{O}}_X = \pi^{-1}\mathcal{O}_X$ défini par
la règle $(U \to A) \mapsto \mathcal{O}(U)$.
\item Le faisceau $\underline{\mathcal{O}}_Y$ défini par la règle
$(U \to A/V) \mapsto \mathcal{O}(V)$.
\end{enumerate}
Nous en déduisons des morphismes de topos annelés
\begin{equation}
\label{obsolete-equation-pi-schemes}
\vcenter{
\xymatrix{
(\Sh(\mathcal{C}_{X/Y}), \underline{\mathcal{O}}_X) \ar[r]_i \ar[d]_\pi &
(\Sh(\mathcal{C}_{X/Y}), \mathcal{O}) \\
(\Sh(X_{Zar}), \mathcal{O}_X)
}
}
\end{equation}
Le morphisme $i$ est l'identité sur les topos sous-jacents, et
$i^\sharp : \mathcal{O} \to \underline{\mathcal{O}}_X$
est le morphisme évident.
Le morphisme $\pi$ est un cas particulier de Cohomologie sur les sites, Situation
\ref{sites-cohomology-situation-fibred-category}.
Les foncteurs dérivés
$
Li^* : D(\mathcal{O}) \longrightarrow D(\underline{\mathcal{O}}_X)
$,
adjoint à gauche de
$Ri_* = i_* : D(\underline{\mathcal{O}}_X) \to D(\mathcal{O})$,
et
$
L\pi_! : D(\underline{\mathcal{O}}_X) \longrightarrow D(\mathcal{O}_X)
$,
adjoint à gauche de
$\pi^* = \pi^{-1} : D(\mathcal{O}_X) \to D(\underline{\mathcal{O}}_X)$,
joueront un rôle important dans la suite.

\begin{remark}
\label{obsolete-remark-different-topologies}
Nous obtenons une deuxième topologie $\tau_Y$ sur $\mathcal{C}_{X/Y}$
en prenant la topologie induite par celle de $Y_{Zar}$.
Il existe une troisième topologie $\tau_{X \to Y}$, pour laquelle une famille
de morphismes $\{(U_i \to A_i) \to (U \to A)\}$ est couvrante si et seulement
si $U = \bigcup U_i$, $V = \bigcup V_i$ et $A_i \cong V_i \times_V A$.
C'est la topologie induite par celle du site $(X/Y)_{Zar}$, dont la catégorie
sous-jacente est la catégorie des couples $(U, V)$ du
Lemme \ref{obsolete-lemma-category-fibred}, partie (3). Les familles couvrantes de
$(X/Y)_{Zar}$ sont les familles $\{(U_i, V_i) \to (U, V)\}$ telles que
$U = \bigcup U_i$ et $V = \bigcup V_i$. Il existe des morphismes de topos
$$
\xymatrix{
\Sh(\mathcal{C}_{X/Y})
= \Sh(\mathcal{C}_{X/Y}, \tau_X) &
\Sh(\mathcal{C}_{X/Y}, \tau_{X \to Y}) \ar[l] \ar[r] &
\Sh(\mathcal{C}_{X/Y}, \tau_Y)
}
$$
(rappelons que $\tau_X$ est notre topologie «~par défaut~»). Les foncteurs
d'image inverse associés à ces flèches sont donnés par le passage au faisceau
associé, tandis que l'image directe est l'identité sur les préfaisceaux
sous-jacents. Le diagramme de topos
$$
\xymatrix{
\Sh(X_{Zar}) \ar[d]^f & \Sh(\mathcal{C}_{X/Y}) \ar[l]^\pi &
\Sh(\mathcal{C}_{X/Y}, \tau_{X \to Y}) \ar[l] \ar[d] \\
\Sh(Y_{Zar}) & & \Sh(\mathcal{C}_{X/Y}, \tau_Y) \ar[ll]
}
$$
n'est {\bf pas} commutatif. En effet, l'image inverse d'un faisceau abélien
non nul sur $Y$ est un faisceau abélien non nul sur
$(\mathcal{C}_{X/Y}, \tau_{X \to Y})$, mais il existe certainement des
exemples où son image inverse sur $X$ est nulle. Notons que tout préfaisceau
$\mathcal{F}$ sur $Y_{Zar}$ définit un faisceau $\underline{\mathcal{F}}$ sur
$\mathcal{C}_{X/Y}$ par la règle qui associe à $(U \to A/V)$ l'ensemble
$\mathcal{F}(V)$. Même lorsque $\mathcal{F}$ est un faisceau, on n'a pas en
général $\underline{\mathcal{F}} = \pi^{-1}f^{-1}\mathcal{F}$.
Cela tient à la non-commutativité du diagramme ci-dessus, car
$\underline{\mathcal{F}}$ est l'image directe de l'image inverse de
$\mathcal{F}$ sur $\Sh(\mathcal{C}_{X/Y}, \tau_{X \to Y})$ par les flèches
horizontale inférieure et verticale droite. Le faisceau
$\underline{\mathcal{O}}_Y$ en est un exemple.
En revanche, il existe un morphisme
$\underline{\mathcal{F}} \to \pi^{-1}f^{-1}\mathcal{F}$
qui, comme nous le verrons plus loin, devient un isomorphisme après application
de $\pi_!$.
\end{remark}





\section{D\'eformations et obstructions des modules plats}
\label{obsolete-section-modules}

\noindent
Dans cette section, nous esquissons la construction d'une th\'eorie des
d\'eformations du champ des faisceaux coh\'erents sur un espace alg\'ebrique
$X$ quelconque au-dessus d'un anneau $\Lambda$. Ce d\'eveloppement est
obsol\`ete en raison de l'expos\'e am\'elior\'e de
Quot, section \ref{quot-section-not-flat}.

\medskip\noindent
Nous nous placerons dans le cadre suivant. On se donne
\begin{enumerate}
\item un anneau $\Lambda$,
\item un espace alg\'ebrique $X$ sur $\Lambda$,
\item une $\Lambda$-alg\`ebre $A$; posons
$X_A = X \times_{\Spec(\Lambda)} \Spec(A)$, et
\item un $\mathcal{O}_{X_A}$-module $\mathcal{F}$ de pr\'esentation finie,
plat sur $A$.
\end{enumerate}
Dans cette situation, nous consid\'ererons toutes les surjections possibles
$$
0 \to I \to A' \to A \to 0
$$
o\`u $A'$ est une $\Lambda$-alg\`ebre dont le noyau $I$ est un id\'eal de carr\'e
nul dans $A'$. Pour tout $A'$, on obtient un \'epaississement du premier
ordre $X_A \to X_{A'}$ d'espaces alg\'ebriques au-dessus de
$\Spec(\Lambda)$. Pour chacun de ces \'epaississements, nous consid\'erons le
probl\`eme du rel\`evement de $\mathcal{F}$ en un module $\mathcal{F}'$ de
pr\'esentation finie sur $X_{A'}$, plat sur $A'$. Nous voudrions retrouver
dans ce cadre les r\'esultats de Th\'eorie des d\'eformations, Lemme
\ref{defos-lemma-flat-ringed-topoi}.

\medskip\noindent
Plus pr\'ecis\'ement, notons $\textit{Lift}(\mathcal{F}, A')$ la cat\'egorie
des couples $(\mathcal{F}', \alpha)$, o\`u $\mathcal{F}'$ est un module de
pr\'esentation finie sur $X_{A'}$, plat sur $A'$, et o\`u
$\alpha : \mathcal{F}'|_{X_A} \to \mathcal{F}$ est un isomorphisme.
Les morphismes $(\mathcal{F}'_1, \alpha_1) \to (\mathcal{F}'_2, \alpha_2)$
sont les isomorphismes $\mathcal{F}'_1 \to \mathcal{F}'_2$ compatibles avec
$\alpha_1$ et $\alpha_2$.
L'ensemble des classes d'isomorphisme de $\textit{Lift}(\mathcal{F}, A')$
est not\'e $\text{Lift}(\mathcal{F}, A')$.

\medskip\noindent
Soit $\mathcal{G}$ un faisceau de $\mathcal{O}_X \otimes_\Lambda A$-modules
sur $X_\etale$, plat sur $A$. Nous introduisons la cat\'egorie
$\textit{Lift}(\mathcal{G}, A')$ des couples
$(\mathcal{G}', \beta)$, o\`u $\mathcal{G}'$ est un faisceau de
$\mathcal{O}_X \otimes_\Lambda A'$-modules plat sur $A'$ et $\beta$
est un isomorphisme $\mathcal{G}' \otimes_{A'} A \to \mathcal{G}$.

\begin{lemma}
\label{obsolete-lemma-equivalence}
Notations et hypoth\`eses comme ci-dessus. Soit $p : X_A \to X$ la projection.
Pour tout $A'$, notons $p' : X_{A'} \to X$ la projection. Le foncteur $p'_*$
induit une \'equivalence de cat\'egories entre
\begin{enumerate}
\item la cat\'egorie $\textit{Lift}(\mathcal{F}, A')$, et
\item la cat\'egorie $\textit{Lift}(p_*\mathcal{F}, A')$.
\end{enumerate}
\end{lemma}

\begin{proof}
FIXME.
\end{proof}

\noindent
Soit $\mathcal{H}$ un faisceau de $\mathcal{O} \otimes_\Lambda A$-modules
sur $\mathcal{C}_{X/\Lambda}$, plat sur $A$. Nous introduisons la cat\'egorie
$\textit{Lift}_\mathcal{O}(\mathcal{H}, A')$
dont les objets sont les couples $(\mathcal{H}', \gamma)$, o\`u $\mathcal{H}'$
est un faisceau de $\mathcal{O} \otimes_\Lambda A'$-modules plat sur $A'$,
et o\`u $\gamma : \mathcal{H}' \otimes_{A'} A \to \mathcal{H}$
est un isomorphisme de $\mathcal{O} \otimes_\Lambda A$-modules.

\medskip\noindent
Soit $\mathcal{G}$ un faisceau de $\mathcal{O}_X \otimes_\Lambda A$-modules
sur $X_\etale$, plat sur $A$.
Consid\'erons les morphismes $i$ et $\pi$ de
Complexe cotangent, \'Equation (\ref{cotangent-equation-pi-spaces}).
Notons $\underline{\mathcal{G}} = \pi^{-1}(\mathcal{G})$. Il est
simplement donn\'e par la r\`egle $(U \to \mathbf{A}) \mapsto \mathcal{G}(U)$;
c'est donc un faisceau de
$\underline{\mathcal{O}}_X \otimes_\Lambda A$-modules.
Notons $i_*\underline{\mathcal{G}}$ le m\^eme faisceau, consid\'er\'e comme un
faisceau de $\mathcal{O} \otimes_\Lambda A$-modules.

\begin{lemma}
\label{obsolete-lemma-second-equivalence}
Notations et hypoth\`eses comme ci-dessus.
Le foncteur $\pi_!$ induit une \'equivalence de cat\'egories entre
\begin{enumerate}
\item la cat\'egorie
$\textit{Lift}_\mathcal{O}(i_*\underline{\mathcal{G}}, A')$, et
\item la cat\'egorie $\textit{Lift}(\mathcal{G}, A')$.
\end{enumerate}
\end{lemma}

\begin{proof}
FIXME.
\end{proof}

\begin{lemma}
\label{obsolete-lemma-second-equivalence-obs}
Notations et hypoth\`eses du Lemme \ref{obsolete-lemma-second-equivalence}.
Consid\'erons l'objet
$$
L = L(\Lambda, X, A, \mathcal{G}) = L\pi_!(Li^*(i_*(\underline{\mathcal{G}})))
$$
de $D(\mathcal{O}_X \otimes_\Lambda A)$. Pour toute surjection $A' \to A$
de $\Lambda$-alg\`ebres dont le noyau $I$ est de carr\'e nul, on a
\begin{enumerate}
\item La cat\'egorie $\textit{Lift}(\mathcal{G}, A')$ est non vide si et
seulement si une certaine classe
$\xi \in \Ext^2_{\mathcal{O}_X \otimes A}(L, \mathcal{G} \otimes_A I)$
est nulle.
\item Si $\textit{Lift}(\mathcal{G}, A')$ est non vide, alors
$\text{Lift}(\mathcal{G}, A')$ est un ensemble principal homog\`ene sous
$\Ext^1_{\mathcal{O}_X \otimes A}(L, \mathcal{G} \otimes_A I)$.
\item \'Etant donn\'e un rel\`evement $\mathcal{G}'$, l'ensemble des
automorphismes de $\mathcal{G}'$ dont l'image inverse est
$\text{id}_\mathcal{G}$ est
canoniquement isomorphe \`a
$\Ext^0_{\mathcal{O}_X \otimes A}(L, \mathcal{G} \otimes_A I)$.
\end{enumerate}
\end{lemma}

\begin{proof}
FIXME.
\end{proof}

\noindent
Enfin, rassemblons ces r\'esultats comme suit.

\begin{proposition}
\label{obsolete-proposition-conclusion}
Pour $\Lambda$, $X$, $A$ et $\mathcal{F}$ comme ci-dessus, il existe un objet
canonique $L = L(\Lambda, X, A, \mathcal{F})$ de $D(X_A)$ tel que, pour toute
surjection $A' \to A$ de $\Lambda$-alg\`ebres dont le noyau $I$ est de carr\'e
nul, on ait
\begin{enumerate}
\item La cat\'egorie $\textit{Lift}(\mathcal{F}, A')$ est non vide si et
seulement si une certaine classe $\xi \in \Ext^2_{X_A}(L,
\mathcal{F} \otimes_A I)$ est nulle.
\item Si $\textit{Lift}(\mathcal{F}, A')$ est non vide, alors
$\text{Lift}(\mathcal{F}, A')$ est un ensemble principal homog\`ene sous
$\Ext^1_{X_A}(L, \mathcal{F} \otimes_A I)$.
\item \'Etant donn\'e un rel\`evement $\mathcal{F}'$, l'ensemble des
automorphismes de $\mathcal{F}'$ dont l'image inverse est
$\text{id}_\mathcal{F}$ est
canoniquement isomorphe \`a
$\Ext^0_{X_A}(L, \mathcal{F} \otimes_A I)$.
\end{enumerate}
\end{proposition}

\begin{proof}
FIXME.
\end{proof}

\begin{lemma}
\label{obsolete-lemma-pseudo-coherent}
Dans la situation de la Proposition \ref{obsolete-proposition-conclusion}, si
$X \to \Spec(\Lambda)$ est localement de type fini et si $\Lambda$ est
noeth\'erien, alors $L$ est pseudo-coh\'erent.
\end{lemma}

\begin{proof}
FIXME.
\end{proof}


\section{Le champ des faisceaux cohérents dans le cas non plat}
\label{obsolete-section-not-flat}

\noindent
Dans Quot, théorème \ref{quot-theorem-coherent-algebraic},
l'hypothèse selon laquelle $f : X \to B$ est plat n'est pas nécessaire.
Dans cette section, nous modifions la méthode de démonstration, en nous
appuyant sur des idées issues de la géométrie algébrique dérivée, afin de nous
affranchir de l'hypothèse de platitude. Une méthode entièrement différente est
employée dans Quot, section \ref{quot-section-not-flat} pour obtenir exactement
le même résultat ; c'est pourquoi la méthode de la présente section est
obsolète.

\medskip\noindent
La seule étape de la démonstration de
Quot, théorème \ref{quot-theorem-coherent-algebraic}
qui utilise la platitude est l'application de
Quot, lemme \ref{quot-lemma-coherent-defo-thy}.
Ce lemme sert à construire une théorie de l'obstruction comme dans
Axiomes d'Artin, section \ref{artin-section-dual}.
Sa démonstration repose sur
Théorie des déformations, lemmes \ref{defos-lemma-flat-ringed-topoi} et
\ref{defos-lemma-verify-iv-ringed-topoi} de
Théorie des déformations, section \ref{defos-section-flat-ringed-topoi}.
C'est ainsi qu'apparaît l'hypothèse de platitude de $f$.
Avant de poursuivre, remarquons que les résultats (2) et (3) de
Théorie des déformations, lemmes \ref{defos-lemma-flat-ringed-topoi}
restent valables sans supposer $f$ plat, puisqu'ils reposent sur
Théorie des déformations, lemmes \ref{defos-lemma-inf-ext-rel-ringed-topoi}
et \ref{defos-lemma-inf-map-rel-ringed-topoi},
qui ne comportent aucune hypothèse de platitude.

\medskip\noindent
Avant d'entrer dans les détails, motivons la construction provenant de la
géométrie algébrique dérivée, car elle éclairera, nous semble-t-il, la suite.
Soit $A$ une algèbre de type fini sur la base localement noethérienne $S$.
Notons $X \otimes^\mathbf{L} A$ un « changement de base dérivé » de $X$ à
$A$, et notons $i : X_A \to X \otimes^\mathbf{L} A$ le morphisme d'inclusion
canonique. L'objet $X \otimes^\mathbf{L} A$ n'est pas (encore) défini dans le
Projet Champs ; on peut le concevoir comme l'espace algébrique $X_A$ muni d'un
faisceau simplicial d'anneaux
$\mathcal{O}_{X \otimes^\mathbf{L} A}$ dont les faisceaux d'homologie sont
$$
H_i(\mathcal{O}_{X \otimes^\mathbf{L} A}) =
\text{Tor}^{\mathcal{O}_S}_i(\mathcal{O}_X, A).
$$
Le morphisme $X \otimes^\mathbf{L} A \to \Spec(A)$ est plat
(les termes du faisceau simplicial d'anneaux étant plats sur $A$),
de sorte que les résultats usuels sur les déformations de modules plats lui
sont applicables. On obtient ainsi une théorie de l'obstruction faisant
intervenir les groupes
$$
\Ext^i_{X \otimes^\mathbf{L} A}(i_*\mathcal{F},
i_*\mathcal{F} \otimes_A M)
$$
où $i = 0, 1, 2$ correspondent respectivement aux automorphismes
infinitésimaux, aux déformations infinitésimales et aux obstructions.
Remarquons que toute déformation plate de $i_*\mathcal{F}$ sur
$X \otimes^\mathbf{L} A'$ est nécessairement de la forme
$i'_*\mathcal{F}'$, où $\mathcal{F}'$ est une déformation plate de
$\mathcal{F}$. Par adjonction entre les foncteurs $Li^*$ et
$i_* = Ri_*$, ces groupes d'Ext sont égaux à
$$
\Ext^i_{X_A}(Li^*(i_*\mathcal{F}), \mathcal{F} \otimes_A M)
$$
Nous obtenons donc des groupes d'obstruction exactement de la même forme que
dans la démonstration de Quot, lemme \ref{quot-lemma-coherent-defo-thy},
à ceci près que la première occurrence de $\mathcal{F}$ est remplacée par le
complexe $Li^*(i_*\mathcal{F})$.

\medskip\noindent
Nous démontrons ci-dessous la version non plate du lemme en construisant
« directement » $E(\mathcal{F}) = Li^*(i_*\mathcal{F})$, puis en établissant
directement son lien avec la théorie des déformations de $\mathcal{F}$.
En fait, il suffit de construire $\tau_{\geq -2}E(\mathcal{F})$, puisque seuls
nous intéressent les groupes d'Ext
$\Ext^i_{X_A}(Li^*(i_*\mathcal{F}), \mathcal{F} \otimes_A M)$
pour $i = 0, 1, 2$. On peut même déterminer les faisceaux de cohomologie
$$
H^i(E(\mathcal{F})) =
\left\{
\begin{matrix}
0 & \text{si }i > 0 \\
\mathcal{F} & \text{si } i = 0 \\
0 & \text{si } i = -1 \\
\text{Tor}_1^{\mathcal{O}_S}(\mathcal{O}_X, A)
\otimes_{\mathcal{O}_X} \mathcal{F} &
\text{si } i = -2
\end{matrix}
\right.
$$
Cette observation guidera notre construction de $E(\mathcal{F})$ dans les
remarques qui suivent.

\begin{remark}[Construction directe]
\label{obsolete-remark-construction-E}
Soit $S$ un schéma. Soit $f : X \to B$ un morphisme d'espaces algébriques
sur $S$. Soit $U$ un autre espace algébrique sur $B$. Notons
$q : X \times_B U \to U$ la seconde projection. Considérons le triangle
distingué
$$
Lq^*L_{U/B} \to L_{X \times_B U/B} \to E \to Lq^*L_{U/B}[1]
$$
de Cotangent, section \ref{cotangent-section-fibre-product}.
Pour tout faisceau $\mathcal{F}$ de
$\mathcal{O}_{X \times_B U}$-modules, on dispose de la classe d'Atiyah
$$
\mathcal{F} \to
L_{X \times_B U/B}
\otimes_{\mathcal{O}_{X \times_B U}}^\mathbf{L} \mathcal{F}[1]
$$
voir Cotangent, section \ref{cotangent-section-atiyah-general}.
Composons-la avec le morphisme vers $E$, puis choisissons un triangle distingué
$$
E(\mathcal{F}) \to \mathcal{F} \to
\mathcal{F} \otimes_{\mathcal{O}_{X \times_B U}}^\mathbf{L} E[1] \to
E(\mathcal{F})[1]
$$
dans $D(\mathcal{O}_{X \times_B U})$.
Par construction, la classe d'Atiyah se relève en un morphisme
$$
e_\mathcal{F} :
E(\mathcal{F})
\longrightarrow
Lq^*L_{U/B} \otimes_{\mathcal{O}_{X \times_B U}}^\mathbf{L} \mathcal{F}[1]
$$
qui s'inscrit dans un morphisme de triangles distingués
$$
\xymatrix{
\mathcal{F} \otimes^\mathbf{L} Lq^*L_{U/B}[1] \ar[r] &
\mathcal{F} \otimes^\mathbf{L} L_{X \times_B U/B}[1] \ar[r] &
\mathcal{F} \otimes^\mathbf{L} E[1] \\
E(\mathcal{F}) \ar[r] \ar[u]^{e_\mathcal{F}} &
\mathcal{F} \ar[r] \ar[u]^{Atiyah} &
\mathcal{F} \otimes^\mathbf{L} E[1] \ar[u]^{=}
}
$$
Une fois donnés $S, B, X, f, U, \mathcal{F}$, nous fixons un choix de
$E(\mathcal{F})$ et de $e_\mathcal{F}$.
\end{remark}

\begin{remark}[Construction de la classe d'obstruction]
\label{obsolete-remark-construction-ob}
Avec les notations de la remarque \ref{obsolete-remark-construction-E}, soit
$i : U \to U'$ un épaississement du premier ordre de $U$ sur $B$. Soit
$\mathcal{I} \subset \mathcal{O}_{U'}$ le faisceau quasi-cohérent d'idéaux
définissant $U$ dans $U'$. Le triangle fondamental
$$
Li^*L_{U'/B} \to L_{U/B} \to L_{U/U'} \to Li^*L_{U'/B}[1]
$$
et le morphisme $L_{U/U'} \to \mathcal{I}[1]$ définissent ensemble un
morphisme $e_{U'} : L_{U/B} \to \mathcal{I}[1]$. En le combinant au morphisme
$e_\mathcal{F}$ de la remarque précédente, on obtient
$$
(\text{id}_\mathcal{F} \otimes Lq^*e_{U'}) \cup e_\mathcal{F} :
E(\mathcal{F})
\longrightarrow
\mathcal{F} \otimes_{\mathcal{O}_{X \times_B U}} q^*\mathcal{I}[2]
$$
(on a en outre composé avec le morphisme du produit tensoriel dérivé vers le
produit tensoriel usuel). Autrement dit, on obtient un élément
$$
\xi_{U'} \in
\Ext^2_{\mathcal{O}_{X \times_B U}}(
E(\mathcal{F}),
\mathcal{F} \otimes_{\mathcal{O}_{X \times_B U}} q^*\mathcal{I})
$$
\end{remark}

\begin{lemma}
\label{obsolete-lemma-ob-is-obstruction}
Dans la situation de la remarque \ref{obsolete-remark-construction-ob}, supposons que
$\mathcal{F}$ soit plat sur $U$. Alors l'annulation de la classe $\xi_{U'}$
est une condition nécessaire et suffisante pour qu'il existe un
$\mathcal{O}_{X \times_B U'}$-module $\mathcal{F}'$, plat sur $U'$, tel que
$i^*\mathcal{F}' \cong \mathcal{F}$.
\end{lemma}

\begin{proof}[Démonstration (esquisse)]
Nous utiliserons le critère de
Théorie des déformations, lemme \ref{defos-lemma-inf-obs-ext-rel-ringed-topoi}.
Pour abréger, nous écrirons
$\mathcal{O} = \mathcal{O}_{X \times_B U}$ et
$\mathcal{O}' = \mathcal{O}_{X \times_B U'}$.
Considérons la suite exacte courte
$$
0 \to \mathcal{I} \to \mathcal{O}_{U'} \to \mathcal{O}_U \to 0.
$$
Soit $\mathcal{J} \subset \mathcal{O}'$ le faisceau quasi-cohérent d'idéaux
définissant $X \times_B U$. D'après ce qui précède, on obtient une suite exacte
$$
\text{Tor}_1^{\mathcal{O}_B}(\mathcal{O}_X, \mathcal{O}_U) \to
q^*\mathcal{I} \to \mathcal{J} \to 0
$$
où $\text{Tor}_1^{\mathcal{O}_B}(\mathcal{O}_X, \mathcal{O}_U)$ est une
abréviation de
$$
\text{Tor}_1^{h^{-1}\mathcal{O}_B}(p^{-1}\mathcal{O}_X, q^{-1}\mathcal{O}_U)
\otimes_{(p^{-1}\mathcal{O}_X\otimes_{h^{-1}\mathcal{O}_B}q^{-1}\mathcal{O}_U)}
\mathcal{O}.
$$
En tensorisant par $\mathcal{F}$, on obtient la suite exacte
$$
\mathcal{F} \otimes_\mathcal{O}
\text{Tor}_1^{\mathcal{O}_B}(\mathcal{O}_X, \mathcal{O}_U) \to
\mathcal{F} \otimes_\mathcal{O}
q^*\mathcal{I} \to
\mathcal{F} \otimes_\mathcal{O} \mathcal{J} \to 0
$$
(Remarquons que les rôles des lettres $\mathcal{I}$ et $\mathcal{J}$ sont
intervertis par rapport aux notations de
Théorie des déformations, lemme \ref{defos-lemma-inf-obs-ext-rel-ringed-topoi}.)
La condition (1) du lemme affirme que le dernier morphisme ci-dessus est un
isomorphisme, autrement dit que le premier morphisme est nul.
L'annulation de ce morphisme peut se vérifier sur les fibres en des points géométriques
$\overline{z} = (\overline{x}, \overline{u}) : \Spec(k) \to X \times_B U$.
Posons $R = \mathcal{O}_{B, \overline{b}}$,
$A = \mathcal{O}_{X, \overline{x}}$,
$B = \mathcal{O}_{U, \overline{u}}$ et
$C = \mathcal{O}_{\overline{z}}$.
D'après Cotangent, lemme \ref{cotangent-lemma-fibre-product}
et le triangle qui définit $E(\mathcal{F})$, on a
$$
H^{-2}(E(\mathcal{F}))_{\overline{z}} =
\mathcal{F}_{\overline{z}} \otimes \text{Tor}_1^R(A, B)
$$
Le morphisme $\xi_{U'}$ induit donc un morphisme
$$
\mathcal{F}_{\overline{z}} \otimes \text{Tor}_1^R(A, B)
\longrightarrow
\mathcal{F}_{\overline{z}} \otimes_B \mathcal{I}_{\overline{u}}
$$
Nous affirmons que ce morphisme coïncide avec la fibre en ce point du
morphisme décrit ci-dessus (démonstration omise : il s'agit d'un énoncé
purement algébrique). La condition (1) de
Théorie des déformations, lemme \ref{defos-lemma-inf-obs-ext-rel-ringed-topoi}
équivaut donc à l'annulation de
$H^{-2}(\xi_{U'}) :
H^{-2}(E(\mathcal{F})) \to \mathcal{F} \otimes \mathcal{I}$.

\medskip\noindent
Pour achever la démonstration, montrons que, sous l'hypothèse que la condition
(1) est satisfaite, la condition (2) équivaut à l'annulation de $\xi_{U'}$.
Dans la suite de la démonstration, nous écrivons
$\mathcal{F} \otimes \mathcal{I}$ pour désigner
$\mathcal{F} \otimes_\mathcal{O} q^*\mathcal{I} =
\mathcal{F} \otimes_\mathcal{O} \mathcal{J}$. L'examen de la suite spectrale
$$
\Ext^i(H^{-j}(E(\mathcal{F})), \mathcal{F} \otimes \mathcal{I})
\Rightarrow
\Ext^{i + j}(E(\mathcal{F}), \mathcal{F} \otimes \mathcal{I})
$$
et les égalités $H^0(E(\mathcal{F})) = \mathcal{F}$ et
$H^{-1}(E(\mathcal{F})) = 0$ montrent que l'on a une suite exacte
$$
0 \to
\Ext^2(\mathcal{F}, \mathcal{F} \otimes \mathcal{I}) \to
\Ext^2(E(\mathcal{F}), \mathcal{F} \otimes \mathcal{I}) \to
\Hom(H^{-2}(E(\mathcal{F})), \mathcal{F} \otimes \mathcal{I})
$$
Ainsi, notre élément $\xi_{U'}$ appartient à
$\Ext^2(\mathcal{F}, \mathcal{F} \otimes \mathcal{I})$.
Il reste à montrer qu'il coïncide avec l'élément de
Théorie des déformations, lemme \ref{defos-lemma-inf-obs-ext-rel-ringed-topoi} ;
nous omettons cette vérification.
\end{proof}

\begin{lemma}
\label{obsolete-lemma-coherent-defo-thy-general}
Dans Quot, situation \ref{quot-situation-coherent}, supposons que $S$ soit un
schéma localement noethérien et que $S = B$. Posons
$\mathcal{X} = \textit{Coh}_{X/B}$. Alors la versalité est ouverte pour
$\mathcal{X}$ (voir Axiomes d'Artin, définition
\ref{artin-definition-openness-versality}).
\end{lemma}

\begin{proof}[Démonstration (esquisse)]
Soit $U \to S$ un morphisme de schémas de type fini, soit $x$ un objet de
$\mathcal{X}$ au-dessus de $U$, et soit $u_0 \in U$ un point de type fini tel
que $x$ soit versel en $u_0$. En rétrécissant $U$, on peut supposer que $u_0$
est un point fermé (Morphismes, lemme
\ref{morphisms-lemma-point-finite-type}) et que $U = \Spec(A)$, le morphisme
$U \to S$ se factorisant par un ouvert affine $\Spec(\Lambda)$ de $S$.
Nous appliquerons Axiomes d'Artin, lemme \ref{artin-lemma-dual-openness} pour
démontrer le lemme. Soit $\mathcal{F}$ le module cohérent sur
$X_A = \Spec(A) \times_S X$, plat sur $A$, qui correspond à l'objet $x$.

\medskip\noindent
Choisissons $E(\mathcal{F})$ et $e_\mathcal{F}$ comme dans
la remarque \ref{obsolete-remark-construction-E}.
La description des faisceaux de cohomologie de $E(\mathcal{F})$ montre que
$$
\Ext^1(E(\mathcal{F}), \mathcal{F} \otimes_A M) =
\Ext^1(\mathcal{F}, \mathcal{F} \otimes_A M)
$$
pour tout $A$-module $M$. D'après cette égalité et
Théorie des déformations, lemme \ref{defos-lemma-inf-ext-rel-ringed-topoi},
on dispose d'un isomorphisme de foncteurs
$$
T_x(M) = \Ext^1_{X_A}(E(\mathcal{F}), \mathcal{F} \otimes_A M)
$$
D'après le lemme \ref{obsolete-lemma-ob-is-obstruction}, toute surjection $A' \to A$
de $\Lambda$-algèbres dont le noyau $I$ est de carré nul fournit une classe
d'obstruction
$$
\xi_{A'} \in \Ext^2_{X_A}(E(\mathcal{F}), \mathcal{F} \otimes_A I)
$$
Appliquons Catégories dérivées des espaces, lemme
\ref{spaces-perfect-lemma-compute-ext} au calcul des groupes d'Ext
$\Ext^i_{X_A}(E(\mathcal{F}), \mathcal{F} \otimes_A M)$
pour $i \leq m$, où $m = 2$. Nous omettons la vérification que
$E(\mathcal{F})$ appartient à $D^-_{\textit{Coh}}$ ; indication : appliquer
Cotangent, lemme \ref{cotangent-lemma-cotangent-finite}.
On obtient un objet parfait $K \in D(A)$ et des isomorphismes fonctoriels
$$
H^i(K \otimes_A^\mathbf{L} M)
\longrightarrow
\Ext^i_{X_A}(E(\mathcal{F}), \mathcal{F} \otimes_A M)
$$
pour $i \leq m$, compatibles aux morphismes de bord. Cet objet $K$, muni des
identifications affichées ci-dessus, fournit une donnée comme dans
Axiomes d'Artin, situation \ref{artin-situation-dual}.
Enfin, la condition (iv) de
Axiomes d'Artin, lemme \ref{artin-lemma-dual-obstruction}
résulte d'une variante de
Théorie des déformations, lemme \ref{defos-lemma-verify-iv-ringed-topoi},
dont nous omettons l'énoncé et la démonstration.
Ainsi, Axiomes d'Artin, lemme \ref{artin-lemma-dual-openness}
s'applique, ce qui démontre le lemme.
\end{proof}

\begin{theorem}
\label{obsolete-theorem-coherent-algebraic-general}
Soit $S$ un schéma. Soit $f : X \to B$ un morphisme d'espaces algébriques
sur $S$. Supposons que $f$ soit de présentation finie et séparé. Alors
$\textit{Coh}_{X/B}$ est un champ algébrique sur $S$.
\end{theorem}

\begin{proof}
Ce théorème reproduit Quot, théorème
\ref{quot-theorem-coherent-algebraic-general}.
Nous le reprenons ici parce que le contenu de la présente section en fournit
une seconde démonstration (comme nous l'avons expliqué au début de la section).
En effet, nous raisonnons exactement comme dans la démonstration de
Quot, théorème \ref{quot-theorem-coherent-algebraic},
si ce n'est que nous y substituons
le lemme \ref{obsolete-lemma-coherent-defo-thy-general} à
Quot, lemme \ref{quot-lemma-coherent-defo-thy}.
\end{proof}




\section{Modifications}
\label{obsolete-section-modifications}

\noindent
Voici un résultat désormais obsolète concernant la catégorie de
Algébrisation des espaces formels, équation
(\ref{restricted-equation-modification}).
On trouvera le traitement actuel dans Algébrisation des espaces formels,
section \ref{restricted-section-modifications}.

\begin{lemma}
\label{obsolete-lemma-henselian}
Soit $(A, \mathfrak m, \kappa)$ un anneau local noethérien.
La catégorie de
Algébrisation des espaces formels, équation
(\ref{restricted-equation-modification})
pour $A$ est équivalente à la catégorie donnée par
Algébrisation des espaces formels, équation
(\ref{restricted-equation-modification})
pour l'hensélisé $A^h$ de $A$.
\end{lemma}

\begin{proof}
C'est un cas particulier d'Algébrisation des espaces formels, lemme
\ref{restricted-lemma-equivalence-to-completion}.
\end{proof}

\noindent
Le lemme suivant sur les singularités rationnelles n'est plus nécessaire
dans le chapitre consacré à la résolution des singularités de surfaces.

\begin{lemma}
\label{obsolete-lemma-double-dual-rational}
Plaçons-nous dans Résolution des surfaces, situation
\ref{resolve-situation-rational}.
Soit $M$ un $A$-module réflexif fini. Notons
$M \otimes_A \mathcal{O}_X$ l'image réciproque du
$\mathcal{O}_S$-module associé. Alors le morphisme canonique de
$M \otimes_A \mathcal{O}_X$ vers son bidual est surjectif.
\end{lemma}

\begin{proof}
Soit $\mathcal{F} = (M \otimes_A \mathcal{O}_X)^{**}$ le bidual, et soit
$\mathcal{F}' \subset \mathcal{F}$ l'image du morphisme d'évaluation
$M \otimes_A \mathcal{O}_X \to \mathcal{F}$. On dispose alors d'une suite
exacte
$$
0 \to \mathcal{F}' \to \mathcal{F} \to \mathcal{Q} \to 0
$$
Puisque $X$ est normal, les anneaux locaux $\mathcal{O}_{X, x}$ sont des
anneaux de valuation discrète aux points de codimension $1$ (voir
Propriétés, lemme \ref{properties-lemma-criterion-normal}).
Par conséquent, $\mathcal{Q}_x = 0$ en ces points, d'après
Compléments d'algèbre, lemme
\ref{more-algebra-lemma-cokernel-map-double-dual-dvr}.
Ainsi, $\mathcal{Q}$ est supporté par un nombre fini de points fermés et est
engendré par ses sections globales, d'après Cohomologie des schémas, lemme
\ref{coherent-lemma-coherent-support-dimension-0}.
On obtient la suite exacte
$$
0 \to H^0(X, \mathcal{F}') \to H^0(X, \mathcal{F}) \to H^0(X, \mathcal{Q}) \to 0
$$
parce que $\mathcal{F}'$ est engendré par ses sections globales
(Résolution des surfaces, lemme
\ref{resolve-lemma-globally-generated}).
Comme $X \to \Spec(A)$ est un isomorphisme au-dessus du complémentaire du
point fermé et que $M$ est réflexif, les morphismes
$$
M \to H^0(X, \mathcal{F}') \to H^0(X, \mathcal{F})
$$
induisent des isomorphismes après localisation en tout idéal premier non
maximal de $A$. Ces morphismes sont donc des isomorphismes, d'après
Compléments d'algèbre, lemme
\ref{more-algebra-lemma-check-isomorphism-via-depth-and-ass},
et le fait que les modules réflexifs sur les anneaux normaux vérifient la
condition $(S_2)$ (Compléments d'algèbre, lemme
\ref{more-algebra-lemma-reflexive-over-normal}).
On conclut donc, comme voulu, que $\mathcal{Q} = 0$.
\end{proof}


\section{Théorie de l'intersection}
\label{obsolete-section-intersection-theory}

\begin{lemma}
\label{obsolete-lemma-good-blowing-up}
Soit $b : X' \to X$ l'éclatement d'un schéma projectif lisse
sur un corps $k$ le long d'un sous-schéma fermé lisse $Z \subset X$.
Considérons le diagramme
$$
\xymatrix{
E \ar[r]_j \ar[d]_\pi & X' \ar[d]^b \\
Z \ar[r]^i & X
}
$$
Supposons qu'il existe un élément de $K_0(X)$ dont la restriction à
$Z$ est égale à la classe de $\mathcal{C}_{Z/X}$ dans $K_0(Z)$.
Alors $[Lb^*\mathcal{O}_Z] = [\mathcal{O}_E] \cdot \alpha''$
dans $K_0(X')$ pour un certain $\alpha'' \in K_0(X')$.
\end{lemma}

\begin{proof}
Les schémas $X$, $X'$, $E$, $Z$ sont projectifs et lisses sur
$k$ et nous avons donc $K'_0(X) = K_0(X) = K_0(\textit{Vect}(X)) =
K_0(D^b_{\textit{Coh}}(X)))$
et de même pour les $3$ autres. Voir Catégories dérivées des schémas, Lemmes
\ref{perfect-lemma-Noetherian-Kprime},
\ref{perfect-lemma-Kprime-K} et \ref{perfect-lemma-K-is-old-K}.
Nous passerons librement d'une de ces versions à l'autre dans cette démonstration.
Considérons la suite exacte
$$
0 \to \mathcal{F} \to \pi^*\mathcal{C}_{Z/X} \to \mathcal{C}_{E/X'} \to 0
$$
de $\mathcal{O}_E$-modules localement libres de type fini qui définit $\mathcal{F}$.
Remarquons que $\mathcal{C}_{E/X'} = \mathcal{O}_{X'}(-E)|_E$
est la restriction du $\mathcal{O}_X$-module inversible
$\mathcal{O}_{X'}(-E)$.
Soit $\alpha \in K_0(X)$ un élément tel que
$i^*\alpha = [\mathcal{C}_{Z/X}]$ dans $K_0(Z)$.
Posons $\alpha' = b^*\alpha - [\mathcal{O}_{X'}(-E)]$.
Alors $j^*\alpha' = [\mathcal{F}]$. On en déduit que
$j^*\lambda^i(\alpha') = [\wedge^i(\mathcal{F})]$ d'après
Théories cohomologiques de Weil, Lemme \ref{weil-lemma-lambda-operations}.
Cela signifie que $[\mathcal{O}_E] \cdot \alpha' = [\wedge^i\mathcal{F}]$
dans $K_0(X)$; voir
Catégories dérivées des schémas, Lemme \ref{perfect-lemma-projection-formula}.
Soit $r$ la codimension maximale d'une composante irréductible de $Z$
dans $X$. Un calcul que nous omettons montre que
$H^{-i}(Lb^*\mathcal{O}_Z) = \wedge^i\mathcal{F}$
pour $i \geq 0, 1, \ldots, r - 1$ et est nul dans les autres degrés.
Il s'ensuit que, dans $K_0(X)$, on a
\begin{align*}
[Lb^*\mathcal{O}_Z] & =
\sum\nolimits_{i = 0, \ldots, r - 1} (-1)^i[\wedge^i\mathcal{F}] \\
& =
\sum\nolimits_{i = 0, \ldots, r - 1} (-1)^i[\mathcal{O}_E] \lambda^i(\alpha') \\
& =
[\mathcal{O}_E] \left(\sum\nolimits_{i = 0, \ldots, r - 1}
(-1)^i \lambda^i(\alpha')\right)
\end{align*}
Cela démontre le lemme en prenant
$\alpha'' = \sum_{i = 0, \ldots, r - 1} (-1)^i \lambda^i(\alpha')$.
\end{proof}



\begin{lemma}
\label{obsolete-lemma-gysin-factors-principal}
Soit $(S, \delta)$ comme dans Homologie de Chow, Situation \ref{chow-situation-setup}.
Soit $X$ localement de type fini sur $S$.
Supposons $X$ intègre et $n = \dim_\delta(X)$.
Soit $a \in \Gamma(X, \mathcal{O}_X)$ une fonction non nulle.
Soit $i : D = Z(a) \to X$ l'immersion fermée du sous-schéma des zéros de $a$.
Soit $f \in R(X)^*$.
Dans ce cas, $i^*\text{div}_X(f) = 0$ dans $A_{n - 2}(D)$.
\end{lemma}

\begin{proof}
Cas particulier d'Homologie de Chow, Lemme \ref{chow-lemma-gysin-factors-general}.
\end{proof}

\begin{remark}
\label{obsolete-remark-not-true-not-quasi-compact}
Cette remarque disait autrefois que l'on ne savait pas si les flèches de
Homologie de Chow, Lemme \ref{chow-lemma-cycles-k-group}, étaient des isomorphismes
en général. Nous avons cependant trouvé depuis une démonstration de ce fait.
\end{remark}







\subsection{Lemmes d'éclatement}
\label{obsolete-section-blowing-up-lemmas}

\noindent
Dans cette section, nous démontrons quelques lemmes permettant de représenter
des diviseurs de Cartier par des diviseurs de Cartier effectifs convenables
sur des éclatements. Ces lemmes se trouvent dans \cite[section 2.4]{F}.
Nous avons adapté leur énoncé afin qu'ils restent valables
sans hypothèse de type fini. Il se peut que le morphisme $b$
du Lemme \ref{obsolete-lemma-blowing-up-intersections} soit un composé
d'une infinité d'éclatements, mais, au-dessus de tout ouvert quasi-compact
$W \subset X$ donné, un nombre fini d'éclatements suffit
(et c'est le résultat de loc.\ cit.).

\begin{lemma}
\label{obsolete-lemma-push-pull-effective-Cartier}
Soit $(S, \delta)$ comme dans Homologie de Chow, Situation \ref{chow-situation-setup}.
Soient $X$, $Y$ localement de type fini sur $S$.
Soit $f : X \to Y$ un morphisme propre.
Soit $D \subset Y$ un diviseur de Cartier effectif.
Supposons $X$, $Y$ intègres, $n = \dim_\delta(X) = \dim_\delta(Y)$ et
que $f$ soit dominant. Alors
$$
f_*[f^{-1}(D)]_{n - 1} = [R(X) : R(Y)] [D]_{n - 1}.
$$
En particulier, si $f$ est birationnel, alors $f_*[f^{-1}(D)]_{n - 1} = [D]_{n - 1}$.
\end{lemma}

\begin{proof}
Cela résulte immédiatement d'Homologie de Chow, Lemme \ref{chow-lemma-equal-c1-as-cycles},
et du fait que $D$ est le sous-schéma des zéros
de la section canonique $1_D$ de $\mathcal{O}_X(D)$.
\end{proof}

\begin{lemma}
\label{obsolete-lemma-blowing-up-denominators}
Soit $(S, \delta)$ comme dans Homologie de Chow, Situation \ref{chow-situation-setup}.
Soit $X$ localement de type fini sur $S$.
Supposons $X$ intègre avec $\dim_\delta(X) = n$.
Soit $\mathcal{L}$ un $\mathcal{O}_X$-module inversible.
Soit $s$ une section méromorphe non nulle de $\mathcal{L}$.
Soit $U \subset X$ le plus grand sous-schéma ouvert tel que
$s$ corresponde à une section de $\mathcal{L}$ sur $U$.
Il existe un morphisme projectif
$$
\pi : X' \longrightarrow X
$$
tel que
\begin{enumerate}
\item $X'$ est intègre,
\item $\pi|_{\pi^{-1}(U)} : \pi^{-1}(U) \to U$ est un isomorphisme,
\item il existe des diviseurs de Cartier effectifs $D, E \subset X'$
tels que
$$
\pi^*\mathcal{L} = \mathcal{O}_{X'}(D - E),
$$
\item la section méromorphe $s$ correspond, par l'isomorphisme ci-dessus,
à la section méromorphe $1_D \otimes (1_E)^{-1}$ (voir
Diviseurs, Définition
\ref{divisors-definition-invertible-sheaf-effective-Cartier-divisor}),
\item on a
$$
\pi_*([D]_{n - 1} - [E]_{n - 1}) = \text{div}_\mathcal{L}(s)
$$
dans $Z_{n - 1}(X)$.
\end{enumerate}
\end{lemma}

\begin{proof}
Soit $\mathcal{I} \subset \mathcal{O}_X$ le faisceau quasi-cohérent d'idéaux
des dénominateurs de $s$; voir Diviseurs, Définition
\ref{divisors-definition-regular-meromorphic-ideal-denominators}.
D'après Diviseurs, Lemme \ref{divisors-lemma-blowing-up-denominators},
on obtient (2), (3) et (4).
D'après Diviseurs, Lemme \ref{divisors-lemma-blow-up-integral-scheme},
on obtient (1). D'après Diviseurs, Lemme \ref{divisors-lemma-blowing-up-projective},
le morphisme $\pi$ est projectif.
Il reste à démontrer (5).
D'après Homologie de Chow, Lemme \ref{chow-lemma-equal-c1-as-cycles}, on a
$$
\pi_*(\text{div}_{\mathcal{L}'}(s')) = \text{div}_\mathcal{L}(s).
$$
Il suffit donc de montrer que
$\text{div}_{\mathcal{L}'}(s') = [D]_{n - 1} - [E]_{n - 1}$.
Cela résulte de l'égalité
$s' = 1_D \otimes 1_E^{-1}$ et de l'additivité; voir
Diviseurs, Lemme \ref{divisors-lemma-c1-additive}.
\end{proof}

\begin{definition}
\label{obsolete-definition-epsilon}
Soit $(S, \delta)$ comme dans Homologie de Chow, Situation \ref{chow-situation-setup}.
Soit $X$ localement de type fini sur $S$.
Supposons $X$ intègre et $\dim_\delta(X) = n$.
Soient $D_1, D_2$ deux diviseurs de Cartier effectifs sur $X$.
Soit $Z \subset X$ un sous-schéma fermé intègre
avec $\dim_\delta(Z) = n - 1$. L'{\it invariant $\epsilon$}
de cette situation est
$$
\epsilon_Z(D_1, D_2) = n_Z \cdot m_Z
$$
où $n_Z$, resp.\ $m_Z$, est le coefficient de
$Z$ dans le $(n - 1)$-cycle $[D_1]_{n - 1}$, resp.\ $[D_2]_{n - 1}$.
\end{definition}

\begin{lemma}
\label{obsolete-lemma-two-divisors}
Soit $(S, \delta)$ comme dans Homologie de Chow, Situation \ref{chow-situation-setup}.
Soit $X$ localement de type fini sur $S$.
Supposons $X$ intègre et $\dim_\delta(X) = n$.
Soient $D_1, D_2$ deux diviseurs de Cartier effectifs sur $X$.
Soit $Z$ un sous-schéma ouvert et fermé du schéma $D_1 \cap D_2$.
Supposons $\dim_\delta(D_1 \cap D_2 \setminus Z) \leq n - 2$.
Alors il existe un morphisme
$b : X' \to X$ et des diviseurs de Cartier
$D_1', D_2', E$ sur $X'$ ayant les propriétés suivantes :
\begin{enumerate}
\item $X'$ est intègre,
\item $b$ est projectif,
\item $b$ est l'éclatement de $X$ le long du sous-schéma fermé $Z$,
\item $E = b^{-1}(Z)$,
\item $b^{-1}(D_1) = D'_1 + E$, et $b^{-1}D_2 = D_2' + E$,
\item $\dim_\delta(D'_1 \cap D'_2) \leq n - 2$, et si
$Z = D_1 \cap D_2$ alors $D'_1 \cap D'_2 = \emptyset$,
\item pour tout sous-schéma fermé intègre $W'$
avec $\dim_\delta(W') = n - 1$, on a
\begin{enumerate}
\item si $\epsilon_{W'}(D'_1, E) > 0$, alors, en posant
$W = b(W')$, on a
$\dim_\delta(W) = n - 1$ et
$$
\epsilon_{W'}(D'_1, E) < \epsilon_W(D_1, D_2),
$$
\item si $\epsilon_{W'}(D'_2, E) > 0$, alors, en posant
$W = b(W')$, on a
$\dim_\delta(W) = n - 1$ et
$$
\epsilon_{W'}(D'_2, E) < \epsilon_W(D_1, D_2),
$$
\end{enumerate}
\end{enumerate}
\end{lemma}

\begin{proof}
Remarquons que le faisceau quasi-cohérent d'idéaux
$\mathcal{I} = \mathcal{I}_{D_1} + \mathcal{I}_{D_2}$
définit l'intersection schématique $D_1 \cap D_2 \subset X$.
Comme $Z$ est une réunion de composantes connexes de $D_1 \cap D_2$,
on voit que, pour tout $z \in Z$, le noyau de
$\mathcal{O}_{X, z} \to \mathcal{O}_{Z, z}$ est égal à $\mathcal{I}_z$.
Soit $b : X' \to X$ l'éclatement de $X$ le long de $Z$. (Ainsi, localement pour la topologie de Zariski
au voisinage de $Z$, c'est l'éclatement de $X$ le long de $\mathcal{I}$.)
Notons $E = b^{-1}(Z)$ le diviseur de Cartier effectif correspondant; voir
Diviseurs,
Lemme \ref{divisors-lemma-blowing-up-gives-effective-Cartier-divisor}.
Comme $Z \subset D_1$, on a $E \subset f^{-1}(D_1)$ et donc
$D_1 = D_1' + E$ pour un certain diviseur de Cartier effectif $D'_1 \subset X'$;
voir Diviseurs, Lemme \ref{divisors-lemma-difference-effective-Cartier-divisors}.
De même, $D_2 = D_2' + E$. Cela établit les assertions (1) -- (5).

\medskip\noindent
Remarquons que, si $W'$ est comme en (7) (a) ou (7) (b), alors l'image $W$
de $W'$ est contenue dans $D_1 \cap D_2$. Si $W$ n'est pas contenue dans
$Z$, alors $b$ est un isomorphisme au point générique de $W$ et
on voit que $\dim_\delta(W) = \dim_\delta(W') = n - 1$, ce qui
contredit l'hypothèse
$\dim_\delta(D_1 \cap D_2 \setminus Z) \leq n - 2$.
Ainsi, $W \subset Z$. Cela signifie que,
pour démontrer (6) et (7), on peut travailler localement au voisinage de $Z$ sur $X$.

\medskip\noindent
On peut donc supposer que $X = \Spec(A)$, où
$A$ est un anneau intègre noethérien, et que $D_1 = \Spec(A/a)$,
$D_2 = \Spec(A/b)$ et $Z = D_1 \cap D_2$.
Posons $I = (a, b)$. Comme $A$ est intègre et $a, b \not = 0$, on peut
recouvrir l'éclatement par deux cartes affines, à savoir
$U = \Spec(A[s]/(as - b))$ et $V = \Spec(A[t]/(bt -a))$.
Ces cartes sont recollées au moyen de l'isomorphisme
$A[s, s^{-1}]/(as - b) \cong A[t, t^{-1}]/(bt - a)$
qui envoie $s$ sur $t^{-1}$.
Le diviseur de Cartier effectif $E$ est décrit par
$\Spec(A[s]/(as - b, a)) \subset U$ et
$\Spec(A[t]/(bt - a, b)) \subset V$.
Le sous-schéma fermé $D'_1$ correspond à
$\Spec(A[t]/(bt - a, t)) \subset U$.
Le sous-schéma fermé $D'_2$ correspond à
$\Spec(A[s]/(as -b, s)) \subset V$.
Puisque ``$ts = 1$'', on a $D'_1 \cap D'_2 = \emptyset$.

\medskip\noindent
Supposons donné un idéal premier $\mathfrak q \subset A[s]/(as - b)$
de hauteur un tel que $s, a \in \mathfrak q$.
Soit $\mathfrak p \subset A$ l'idéal premier correspondant de $A$.
Remarquons que $a, b \in \mathfrak p$.
D'après la formule des dimensions, on a aussi $\dim(A_{\mathfrak p}) = 1$.
La dernière assertion à démontrer est
$$
\text{ord}_{A_{\mathfrak p}}(a)
\text{ord}_{A_{\mathfrak p}}(b)
>
\text{ord}_{B_{\mathfrak q}}(a)
\text{ord}_{B_{\mathfrak q}}(s)
$$
où $B = A[s]/(as - b)$. D'après
Algèbre, Lemme \ref{algebra-lemma-quasi-finite-extension-dim-1},
on a $\text{ord}_{A_{\mathfrak p}}(x) \geq \text{ord}_{B_{\mathfrak q}}(x)$
pour $x = a, b$. Puisque $\text{ord}_{B_{\mathfrak q}}(s) > 0$, on conclut
par additivité de la fonction $\text{ord}$ et par le fait que
$as = b$.
\end{proof}

\begin{definition}
\label{obsolete-definition-locally-finite-sum-effective-Cartier-divisors}
Soit $X$ un schéma.
Soit $\{D_i\}_{i \in I}$ une famille localement finie
de diviseurs de Cartier effectifs sur $X$.
Supposons donnée une application
$I \to \mathbf{Z}_{\geq 0}$, $i \mapsto n_i$.
La {\it somme des diviseurs de Cartier effectifs}
$D = \sum n_i D_i$ est l'unique diviseur de Cartier effectif
$D \subset X$ tel que, pour tout ouvert quasi-compact $U \subset X$,
on ait $D|_U = \sum_{D_i \cap U \not = \emptyset} n_iD_i|_U$,
la somme du membre de droite étant celle de Diviseurs,
Définition \ref{divisors-definition-sum-effective-Cartier-divisors}.
\end{definition}

\begin{lemma}
\label{obsolete-lemma-sum-divisors-associated-Weil}
Soit $(S, \delta)$ comme dans Homologie de Chow, situation \ref{chow-situation-setup}.
Soit $X$ localement de type fini sur $S$.
Supposons $X$ intègre et $\dim_\delta(X) = n$.
Soit $\{D_i\}_{i \in I}$ une famille localement finie
de diviseurs de Cartier effectifs sur $X$.
Donnons-nous des $n_i \geq 0$ pour $i \in I$.
Alors
$$
[D]_{n - 1} = \sum\nolimits_i n_i[D_i]_{n - 1}
$$
dans $Z_{n - 1}(X)$.
\end{lemma}

\begin{proof}
Puisque nous démontrons une égalité de cycles, nous pouvons raisonner localement sur $X$.
On se ramène ainsi à une somme finie, puis, par récurrence, à une somme de
deux diviseurs de Cartier effectifs $D = D_1 + D_2$.
D'après Homologie de Chow, lemme \ref{chow-lemma-compute-c1}, on a
$D_1 = \text{div}_{\mathcal{O}_X(D_1)}(1_{D_1})$, où
$1_{D_1}$ désigne la section canonique de $\mathcal{O}_X(D_1)$.
La même assertion vaut bien entendu pour $D_2$ et $D$.
Comme $1_D = 1_{D_1} \otimes 1_{D_2}$ sous l'identification
$\mathcal{O}_X(D) = \mathcal{O}_X(D_1) \otimes \mathcal{O}_X(D_2)$,
le résultat découle de Diviseurs, lemme \ref{divisors-lemma-c1-additive}.
\end{proof}

\begin{lemma}
\label{obsolete-lemma-blowing-up-intersections}
Soit $(S, \delta)$ comme dans Homologie de Chow, situation \ref{chow-situation-setup}.
Soit $X$ localement de type fini sur $S$.
Supposons $X$ intègre et $\dim_\delta(X) = d$.
Soit $\{D_i\}_{i \in I}$ une famille localement finie de
diviseurs de Cartier effectifs sur $X$.
Supposons que, pour tout $\{i, j, k\} \subset I$ tel que $\#\{i, j, k\} = 3$,
on ait $D_i \cap D_j \cap D_k = \emptyset$.
Il existe alors
\begin{enumerate}
\item un sous-schéma ouvert $U \subset X$ tel que
$\dim_\delta(X \setminus U) \leq d - 3$,
\item un morphisme $b : U' \to U$, et
\item des diviseurs de Cartier effectifs $\{D'_j\}_{j \in J}$ sur $U'$,
\end{enumerate}
vérifiant les propriétés suivantes :
\begin{enumerate}
\item $b$ est un morphisme propre $b : U' \to U$,
\item $U'$ est intègre,
\item $b$ est un isomorphisme au-dessus du complémentaire de la réunion des
intersections deux à deux des $D_i|_U$,
\item $\{D'_j\}_{j \in J}$ est une famille localement finie de diviseurs de
Cartier effectifs sur $U'$,
\item $\dim_\delta(D'_j \cap D'_{j'}) \leq d - 2$ si $j \not = j'$, et
\item $b^{-1}(D_i|_U) = \sum n_{ij} D'_j$ pour certains $n_{ij} \geq 0$.
\end{enumerate}
De plus, si $X$ est quasi-compact, on peut supposer $U = X$ ci-dessus.
\end{lemma}

\begin{proof}
Commençons par le cas quasi-compact, qui est peut-être
le plus intéressant. Dans ce cas, nous construisons par récurrence une suite
d'éclatements
$$
X = X_0 \xleftarrow{b_0} X_1 \xleftarrow{b_1} X_2 \leftarrow \ldots
$$
et des ensembles finis de diviseurs de Cartier effectifs $\{D_{n, i}\}_{i \in I_n}$.
À chaque étape, ceux-ci auront la propriété que toute intersection triple
$D_{n, i} \cap D_{n, j} \cap D_{n, k}$ est vide.
En outre, pour tout $n \geq 0$, nous aurons
$I_{n + 1} = I_n \amalg P(I_n)$, où $P(I_n)$ désigne
l'ensemble des paires d'éléments de $I_n$. Enfin, nous aurons
$$
b_n^{-1}(D_{n, i}) = D_{n + 1, i} +
\sum\nolimits_{i' \in I_n, i' \not = i} D_{n + 1, \{i, i'\}}
$$
On en déduit que, pour tout $n \geq 0$,
$(b_0 \circ \ldots \circ b_n)^{-1}(D_i)$ est une combinaison linéaire
à coefficients entiers positifs ou nuls des diviseurs $D_{n + 1, j}$, $j \in I_{n + 1}$.

\medskip\noindent
Pour amorcer la récurrence, posons $X_0 = X$,
$I_0 = I$ et $D_{0, i} = D_i$.

\medskip\noindent
Étant donné $(X_n, \{D_{n, i}\}_{i \in I_n})$, soit $X_{n + 1}$
l'éclatement de $X_n$ le long du sous-schéma fermé
$Z_n = \bigcup_{\{i, i'\} \in P(I_n)} D_{n, i} \cap D_{n, i'}$.
Remarquons que les sous-schémas fermés $D_{n, i} \cap D_{n, i'}$ sont deux à deux
disjoints en raison de notre hypothèse sur les intersections triples.
Autrement dit, on peut écrire
$Z_n = \coprod_{\{i, i'\} \in P(I_n)} D_{n, i} \cap D_{n, i'}$.
De plus, dans un voisinage de Zariski de $D_{n, i} \cap D_{n, i'}$, le
morphisme $b_n$ est l'éclatement du schéma $X_n$
le long du sous-schéma fermé $D_{n, i} \cap D_{n, i'}$, et les conclusions
du lemme \ref{obsolete-lemma-two-divisors} s'appliquent.
En posant donc $D_{n + 1, \{i, i'\}} = b_n^{-1}(D_{n, i} \cap D_{n, i'})$,
on obtient un diviseur de Cartier effectif.
Les diviseurs de Cartier $D_{n + 1, \{i, i'\}}$ sont deux à deux disjoints.
Manifestement, on a
$b_n^{-1}(D_{n, i}) \supset D_{n + 1, \{i, i'\}}$ pour
tout $i' \in I_n$, $i' \not = i$. Par conséquent, en appliquant
Diviseurs, lemme \ref{divisors-lemma-difference-effective-Cartier-divisors},
on voit que l'on a bien $b_n^{-1}(D_{n, i}) = D_{n + 1, i} +
\sum\nolimits_{i' \in I_n, i' \not = i} D_{n + 1, \{i, i'\}}$
pour un certain diviseur de Cartier effectif $D_{n + 1, i}$ sur $X_{n + 1}$.
Au voisinage de $D_{n + 1, \{i, i'\}}$, ces diviseurs
$D_{n + 1, i}$ jouent le rôle des diviseurs primés du
lemme \ref{obsolete-lemma-two-divisors}. En particulier, on en déduit que
$D_{n + 1, i} \cap D_{n + 1, i'} = \emptyset$ si $i \not = i'$,
$i, i' \in I_n$, d'après le point (6) du lemme \ref{obsolete-lemma-two-divisors}.
Cela implique déjà que les intersections triples
des diviseurs $D_{n + 1, i}$ sont vides.

\medskip\noindent
À ce stade, nous pouvons utiliser la quasi-compacité de $X$
pour conclure que l'invariant
\begin{equation}
\label{obsolete-equation-invariant}
\epsilon(X, \{D_i\}_{i \in I}) =
\max\{\epsilon_Z(D_i, D_{i'}) \mid
Z \subset X,
\dim_\delta(Z) = d - 1,
\{i, i'\} \in P(I)\}
\end{equation}
est fini, car chaque $D_i$ n'a en définitive qu'un nombre fini de composantes
irréductibles. Nous affirmons que, pour un certain $n$, l'invariant
$\epsilon(X_n, \{D_{n, i}\}_{i \in I_n})$ est nul. En effet, sinon,
le lemme \ref{obsolete-lemma-two-divisors} donnerait une suite strictement décroissante
$$
\epsilon(X, \{D_i\}_{i \in I})
=
\epsilon(X_0, \{D_{0, i}\}_{i \in I_0})
>
\epsilon(X_1, \{D_{1, i}\}_{i \in I_1})
>
\ldots
$$
d'entiers strictement positifs, ce qui est contradictoire. Choisissons $n$ tel que
l'invariant $\epsilon(X_n, \{D_{n, i}\}_{i \in I_n})$ soit nul.
Cela signifie qu'il n'existe ni sous-schéma fermé intègre $Z \subset X_n$,
ni paire d'indices $i, i' \in I_n$
tels que $\epsilon_Z(D_{n, i}, D_{n, i'}) > 0$.
Autrement dit, $\dim_\delta(D_{n, i} \cap D_{n, i'}) \leq d - 2$ pour
toute paire $\{i, i'\} \in P(I_n)$, comme souhaité.

\medskip\noindent
Passons maintenant au cas général où l'on ne suppose plus
le schéma $X$ quasi-compact. Le problème que pose la méthode de
la première partie de la démonstration est que l'on peut obtenir une suite infinie
d'éclatements dont les centres dominent un point fixé de $X$. Pour
éviter cela, nous retirons à chaque étape des sous-ensembles fermés convenables de codimension $\geq 3$.
Plus précisément, nous construirons par récurrence
une suite de morphismes de la forme suivante :
$$
\xymatrix{
X = X_0 \\
U_0 \ar[u]^{j_0} & X_1 \ar[l]_{b_0} \\
 & U_1 \ar[u]^{j_1} & X_2 \ar[l]_{b_1} \\
 & & U_2 \ar[u]^{j_2} & X_3 \ar[l]_{b_2}
}
$$
Chacun des morphismes $j_n : U_n \to X_n$ sera une immersion ouverte.
Chacun des morphismes $b_n : X_{n + 1} \to U_n$ sera un morphisme propre
birationnel entre schémas intègres. Comme dans le cas quasi-compact, nous aurons
des diviseurs de Cartier effectifs $\{D_{n, i}\}_{i \in I_n}$ sur $X_n$.
À chaque étape, ceux-ci auront la propriété que toute intersection triple
$D_{n, i} \cap D_{n, j} \cap D_{n, k}$ est vide.
En outre, pour tout $n \geq 0$, nous aurons
$I_{n + 1} = I_n \amalg P(I_n)$, où $P(I_n)$ désigne
l'ensemble des paires d'éléments de $I_n$.
Enfin, nous ferons en sorte que
$$
b_n^{-1}(D_{n, i}|_{U_n}) = D_{n + 1, i} +
\sum\nolimits_{i' \in I_n, i' \not = i} D_{n + 1, \{i, i'\}}
$$

\medskip\noindent
Amorçons la récurrence en posant $X_0 = X$,
$I_0 = I$ et $D_{0, i} = D_i$.

\medskip\noindent
Étant donné $(X_n, \{D_{n, i}\})$, construisons le sous-schéma ouvert
$U_n$ comme suit. Pour toute paire $\{i, i'\} \in P(I_n)$, considérons
le sous-schéma fermé $D_{n, i} \cap D_{n, i'}$. Ses composantes irréductibles
« bonnes » sont celles de $\delta$-dimension $d - 2$, et ses composantes irréductibles
« mauvaises » celles de $\delta$-dimension $d - 1$.
Posons
$$
\text{Bad}(i, i')
=
\bigcup\nolimits_{W \subset D_{n, i} \cap D_{n, i'}
\text{ comp.\ irr.\ avec }\dim_\delta(W) = d - 1} W
$$
et, de même,
$$
\text{Good}(i, i')
=
\bigcup\nolimits_{W \subset D_{n, i} \cap D_{n, i'}
\text{ comp.\ irr.\ avec }\dim_\delta(W) = d - 2} W.
$$
Alors $D_{n, i} \cap D_{n, i'} = \text{Bad}(i, i') \cup \text{Good}(i, i')$
et, de plus, on a
$\dim_\delta(\text{Bad}(i, i') \cap \text{Good}(i, i')) \leq d - 3$.
Voici le choix de $U_n$ :
$$
U_n
=
X_n
\setminus
\bigcup\nolimits_{\{i, i'\} \in P(I_n)}
\text{Bad}(i, i') \cap \text{Good}(i, i').
$$
Notre condition sur les intersections triples des diviseurs $D_{n, i}$
montre que la réunion est en fait disjointe. De plus,
on voit que, comme schéma,
$$
D_{n, i}|_{U_n} \cap D_{n, i'}|_{U_n}
=
Z_{n, i, i'} \amalg G_{n, i, i'}
$$
où $Z_{n, i, i'}$ est $\delta$-équidimensionnel de dimension $d - 1$
et $G_{n, i, i'}$ est $\delta$-équidimensionnel de dimension $d - 2$.
(Ainsi, topologiquement, $Z_{n, i, i'}$ est la réunion des composantes mauvaises,
privée de ses intersections avec les composantes bonnes.) Enfin, posons
$$
Z_n =
\bigcup\nolimits_{\{i, i'\} \in P(I_n)} Z_{n, i, i'} =
\coprod\nolimits_{\{i, i'\} \in P(I_n)} Z_{n, i, i'},
$$
et soit $b_n : X_{n + 1} \to U_n$ l'éclatement de centre $Z_n$.
Remarquons que le lemme \ref{obsolete-lemma-two-divisors}
s'applique au morphisme $b_n : X_{n + 1} \to U_n$, localement au voisinage
de chacun des lieux $D_{n, i}|_{U_n} \cap D_{n, i'}|_{U_n}$. Ainsi,
exactement comme dans la première partie de la démonstration, nous obtenons des diviseurs de
Cartier effectifs $D_{n + 1, \{i, i'\}}$ pour $\{i, i'\} \in P(I_n)$
et des diviseurs de Cartier effectifs $D_{n + 1, i}$ pour $i \in I_n$
tels que
$b_n^{-1}(D_{n, i}|_{U_n}) = D_{n + 1, i} +
\sum\nolimits_{i' \in I_n, i' \not = i} D_{n + 1, \{i, i'\}}$.
Pour tout $n$, notons $\pi_n : X_n \to X$ le morphisme obtenu
comme le composé $j_0 \circ b_0 \circ \ldots \circ j_{n - 1} \circ b_{n - 1}$.

\medskip\noindent
{\bf Assertion :} étant donné un ouvert quasi-compact $V \subset X$,
pour tout $n$ assez grand, les morphismes
$$
\pi_n^{-1}(V) \leftarrow \pi_{n + 1}^{-1}(V) \leftarrow \ldots
$$
sont tous des isomorphismes. En effet, si le morphisme
$\pi_n^{-1}(V) \leftarrow \pi_{n + 1}^{-1}(V)$ n'est pas un isomorphisme,
alors $Z_{n, i, i'} \cap \pi_n^{-1}(V) \not = \emptyset$ pour une certaine
paire $\{i, i'\} \in P(I_n)$. Il existe donc une composante irréductible
$W \subset D_{n, i} \cap D_{n, i'}$ telle que $\dim_\delta(W) = d - 1$.
En particulier, on a $\epsilon_W(D_{n, i}, D_{n, i'}) > 0$.
En appliquant à plusieurs reprises le lemme \ref{obsolete-lemma-two-divisors}, on obtient
$$
\epsilon_W(D_{n, i}, D_{n, i'})
<
\epsilon(V, \{D_i|_V\}) - n
$$
où $\epsilon(V, \{D_i|_V\})$ est défini comme dans (\ref{obsolete-equation-invariant}).
Puisque $V$ est quasi-compact, on a
$\epsilon(V, \{D_i|_V\}) < \infty$ et, en choisissant $n > \epsilon(V, \{D_i|_V\})$,
on obtient le résultat.

\medskip\noindent
Notons que, par construction, la différence $X_n \setminus U_n$
vérifie $\dim_\delta(X_n \setminus U_n) \leq d - 3$.
Soit $T_n = \pi_n(X_n \setminus U_n)$ son image dans $X$.
En remontant le diagramme de morphismes ci-dessus et en utilisant le fait que chaque $b_n$ est fermé,
on obtient que $T_0 \cup \ldots \cup T_n$ est un sous-ensemble fermé de $X$
pour tout $n$. Par construction, tout $t \in T_n$ vérifie $\delta(t) \leq d - 3$.
Ainsi, $\overline{T_n} \subset X$ est un sous-ensemble fermé
tel que $\dim_\delta(T_n) \leq d - 3$. D'après l'assertion ci-dessus,
pour tout ouvert quasi-compact $V \subset X$, on a
$T_n \cap V \not = \emptyset$ pour au plus un nombre fini de $n$.
Ainsi, $\{\overline{T_n}\}_{n \geq 0}$ est une famille localement finie
de sous-ensembles fermés, et l'on peut poser
$U = X \setminus \bigcup \overline{T_n}$. Ce sera le
$U$ du lemme.

\medskip\noindent
Par construction de $U$, on a $U_n \cap \pi_n^{-1}(U) = \pi_n^{-1}(U)$.
Tous les morphismes
$$
b_n : \pi_{n + 1}^{-1}(U) \longrightarrow \pi_n^{-1}(U)
$$
sont donc propres. De plus, d'après l'assertion, ils deviennent des isomorphismes
à partir d'un certain rang au-dessus de tout ouvert quasi-compact de $X$. On peut donc définir
$$
U' = \lim_n \pi_n^{-1}(U).
$$
Le morphisme induit $b : U' \to U$ est propre, car la propreté est locale
sur $U$ et, au-dessus de tout ouvert quasi-compact, la limite est stationnaire. De même,
posons $J = \bigcup_{n \geq 0} I_n$ au moyen des inclusions
$I_n \to I_{n + 1}$ issues de la construction. Pour $j \in J$, choisissons
un $n_0$ tel que $j$ corresponde à $i \in I_{n_0}$ et définissons
$D'_j = \lim_{n \geq n_0} D_{n, i}$. Cette définition a encore un sens,
car, localement sur $X$, les morphismes se stabilisent.
Les autres assertions du lemme se vérifient comme dans le cas où
$X$ est quasi-compact.
\end{proof}









\section{Commutativité de l'intersection des diviseurs}
\label{obsolete-section-commutativity}

\noindent
Les résultats de cette section ont été initialement utilisés pour donner une autre
démonstration des lemmes de Homologie de Chow, section \ref{chow-section-commutativity},
ainsi qu'une version faible de
Homologie de Chow, Lemme \ref{chow-lemma-gysin-commutes-gysin}.

\begin{lemma}
\label{obsolete-lemma-improved-additivity}
Soit $(S, \delta)$ comme dans Homologie de Chow, Situation \ref{chow-situation-setup}.
Soit $X$ localement de type fini sur $S$.
Soit $\{i_j : D_j \to X \}_{j \in J}$ une famille localement finie
de diviseurs de Cartier effectifs sur $X$. Soient $n_j > 0$, $j\in J$.
Posons $D = \sum_{j \in J} n_j D_j$, et notons $i : D \to X$ le
morphisme d'inclusion. Soit $\alpha \in Z_{k + 1}(X)$. Alors
$$
p : \coprod\nolimits_{j \in J} D_j \longrightarrow D
$$
est propre et
$$
i^*\alpha = p_*\left(\sum n_j i_j^*\alpha\right)
$$
dans $\CH_k(D)$.
\end{lemma}

\begin{proof}
La démonstration de ce lemme est un peu plus longue que prévu
en raison d'une subtilité concernant les sommes infinies d'équivalences rationnelles.
Dans le cas quasi-compact, la famille $D_j$ est finie et le résultat
est tout à fait élémentaire : il découle directement de
Homologie de Chow, Lemme \ref{chow-lemma-compute-c1}, de
Diviseurs, Lemme \ref{divisors-lemma-c1-additive}, et des définitions.

\medskip\noindent
Le morphisme $p$ est propre puisque la famille $\{D_j\}_{j \in J}$
est localement finie. Écrivons $\alpha = \sum_{a \in A} m_a [W_a]$
avec $W_a \subset X$ un sous-schéma intègre fermé de
$\delta$-dimension $k + 1$.
Notons $i_a : W_a \to X$ l'immersion fermée.
Nous supposons que $m_a \not = 0$ pour tout $a \in A$ et que
$\{W_a\}_{a \in A}$ est localement finie sur $X$.

\medskip\noindent
Observons que, d'après
Homologie de Chow, Définition \ref{chow-definition-gysin-homomorphism},
la classe $i^*\alpha$ est celle d'un cycle
$\sum m_a\beta_a$ pour certains $\beta_a \in Z_k(W_a \cap D)$.
Plus précisément, si $W_a \not \subset D$, alors $\beta_a = [D \cap W_a]_k$,
et si $W_a \subset D$, alors $\beta_a$ est un cycle
représentant $c_1(\mathcal{O}_X(D)) \cap [W_a]$.

\medskip\noindent
Pour chaque $a \in A$, écrivons $J = J_{a, 1} \amalg J_{a, 2} \amalg J_{a, 3}$,
où
\begin{enumerate}
\item $j \in J_{a, 1}$ si et seulement si $W_a \cap D_j = \emptyset$,
\item $j \in J_{a, 2}$ si et seulement si
$W_a \not = W_a \cap D_1 \not = \emptyset$, et
\item $j \in J_{a, 3}$ si et seulement si $W_a \subset D_j$.
\end{enumerate}
Puisque la famille $\{D_j\}$ est localement finie, nous voyons que
$J_{a, 3}$ est un ensemble fini. Pour tous $a \in A$ et $j \in J$,
nous choisissons un cycle $\beta_{a, j} \in Z_k(W_a \cap D_j)$ comme suit :
\begin{enumerate}
\item si $j \in J_{a, 1}$, nous posons $\beta_{a, j} = 0$,
\item si $j \in J_{a, 2}$, nous posons $\beta_{a, j} = [D_j \cap W_a]_k$, et
\item si $j \in J_{a, 3}$, nous choisissons $\beta_{a, j} \in Z_k(W_a)$
représentant $c_1(i_a^*\mathcal{O}_X(D_j)) \cap [W_j]$.
\end{enumerate}
Nous affirmons que
$$
\beta_a \sim_{rat}
\sum\nolimits_{j \in J} n_j \beta_{a, j}
$$
dans $\CH_k(W_a \cap D)$.

\medskip\noindent
Cas I : $W_a \not \subset D$. Dans ce cas, $J_{a, 3} = \emptyset$.
Il suffit donc de montrer que
$[D \cap W_a]_k = \sum n_j [D_j \cap W_a]_k$ comme cycles.
C'est le Lemme \ref{obsolete-lemma-sum-divisors-associated-Weil}.

\medskip\noindent
Cas II : $W_a \subset D$. Dans ce cas, $\beta_a$ est un cycle représentant
$c_1(i_a^*\mathcal{O}_X(D)) \cap [W_a]$.
Écrivons $D = D_{a, 1} + D_{a, 2} + D_{a, 3}$ avec
$D_{a, s} = \sum_{j \in J_{a, s}} n_jD_j$.
D'après Diviseurs, Lemme \ref{divisors-lemma-c1-additive}, nous avons
\begin{eqnarray*}
c_1(i_a^*\mathcal{O}_X(D)) \cap [W_a] & = &
c_1(i_a^*\mathcal{O}_X(D_{a, 1})) \cap [W_a] +
c_1(i_a^*\mathcal{O}_X(D_{a, 2})) \cap [W_a] \\
& &
 + c_1(i_a^*\mathcal{O}_X(D_{a, 3})) \cap [W_a].
\end{eqnarray*}
Il est clair que le premier terme de la somme est nul.
Comme $J_{a, 3}$ est fini, nous voyons que le dernier terme coïncide
avec $\sum\nolimits_{j \in J_{a, 3}} n_jc_1(i_a^*\mathcal{L}_j) \cap [W_a]$ ;
voir Diviseurs, Lemme \ref{divisors-lemma-c1-additive}.
Ce cycle est représenté par $\sum_{j \in J_{a, 3}} n_j \beta_{a, j}$.
Enfin, le Cas I montre que le terme intermédiaire est représenté par le cycle
$\sum\nolimits_{j \in J_{a, 2}} n_j[D_j \cap W_a]_k =
\sum_{j \in J_{a, 2}} n_j\beta_{a, j}$.
D'où l'affirmation dans ce cas.

\medskip\noindent
Nous sommes maintenant en mesure d'achever la démonstration du lemme.
En effet, nous avons $i^*D \sim_{rat} \sum m_a\beta_a$ par notre
choix de $\beta_a$. Pour chaque $a$, nous avons
$\beta_a \sim_{rat} \sum_j \beta_{a, j}$, l'équivalence rationnelle
ayant lieu sur $D \cap W_a$.
Comme la famille de sous-schémas fermés $D \cap W_a$
est localement finie sur $D$, nous voyons aussi que
$\sum m_a \beta_a \sim_{rat} \sum_{a, j} m_a\beta_{a, j}$
sur $D$ (voir
Homologie de Chow, Remarque \ref{chow-remark-infinite-sums-rational-equivalences}).
Enfin, il est clair que $\sum_a m_a\beta_{a, j}$ (considéré
comme un cycle sur $D_j$) représente $i_j^*\alpha$, et donc
$\sum_{a, j} m_a\beta_{a, j}$ représente $p_* \sum_j i_j^*\alpha$,
ce qui achève la démonstration.
\end{proof}

\begin{lemma}
\label{obsolete-lemma-commutativity-effective-Cartier-proper-intersection}
Soit $(S, \delta)$ comme dans Homologie de Chow, Situation \ref{chow-situation-setup}.
Soit $X$ localement de type fini sur $S$.
Supposons $X$ intègre et $\dim_\delta(X) = n$.
Soient $D$ et $D'$ des diviseurs de Cartier effectifs sur $X$.
Supposons $\dim_\delta(D \cap D') = n - 2$. Soient $i : D \to X$,
resp.\ $i' : D' \to X$, les immersions fermées correspondantes.
Alors
\begin{enumerate}
\item il existe un cycle $\alpha \in Z_{n - 2}(D \cap D')$
dont l'image directe sur $D$ représente
$i^*[D']_{n - 1} \in \CH_{n - 2}(D)$
et dont l'image directe sur $D'$ représente
$(i')^*[D]_{n - 1} \in \CH_{n - 2}(D')$, et
\item nous avons
$$
D \cdot [D']_{n - 1}
=
D' \cdot [D]_{n - 1}
$$
dans $\CH_{n - 2}(X)$.
\end{enumerate}
\end{lemma}

\begin{proof}
La partie (2) est une conséquence immédiate de la partie (1).
Écrivons $[D]_{n - 1} = \sum n_a[Z_a]$ et
$[D']_{n - 1} = \sum m_b[Z_b]$, où les $Z_a$ sont les composantes irréductibles
de $D$ et les $[Z_b]$ les composantes irréductibles
de $D'$. D'après
Homologie de Chow, Définition \ref{chow-definition-gysin-homomorphism},
nous avons $i^*D' = \sum m_b i^*[Z_b]$ et $(i')^*D = \sum n_a(i')^*[Z_a]$.
Par hypothèse, aucune des composantes irréductibles $Z_b$
n'est contenue dans $D$, et donc $i^*[Z_b] = [Z_b\cap D]_{n - 2}$
par définition. De même, $(i')^*[Z_a] = [Z_a \cap D']_{n - 2}$.
Nous cherchons donc à démontrer l'égalité de cycles
$$
\sum n_a[Z_a \cap D']_{n - 2} = \sum m_b[Z_b \cap D]_{n - 2}
$$
qui sont bien supportés sur $D \cap D'$.
Soit $W \subset X$ un sous-schéma intègre fermé
tel que $\dim_\delta(W) = n - 2$. Soit $\xi \in W$ son point générique.
Posons $R = \mathcal{O}_{X, \xi}$. C'est un anneau local noethérien intègre.
Notons que $\dim(R) = 2$. Soit $f \in R$, resp.\ $f' \in R$,
un élément définissant l'idéal de $D$, resp.\ de $D'$.
Par hypothèse, $\dim(R/(f, f')) = 0$. Soient
$\mathfrak q'_1, \ldots, \mathfrak q'_t \subset R$ les idéaux premiers minimaux
contenant $(f')$, et $\mathfrak q_1, \ldots, \mathfrak q_s \subset R$ les idéaux premiers minimaux
contenant $(f)$.
L'égalité ci-dessus se ramène à l'égalité
$$
\sum_{i = 1, \ldots, s}
\text{length}_{R_{\mathfrak q_i}}(R_{\mathfrak q_i}/(f))
\text{ord}_{R/\mathfrak q_i}(f')
=
\sum_{j = 1, \ldots, t}
\text{length}_{R_{\mathfrak q'_j}}(R_{\mathfrak q'_j}/(f'))
\text{ord}_{R/\mathfrak q'_j}(f).
$$
D'après Homologie de Chow, Lemme \ref{chow-lemma-length-multiplication}, 
appliqué avec $M = R/(f)$, le membre de gauche de
cette égalité est égal à
$$
\text{length}_R(R/(f, f'))
-
\text{length}_R(\Ker(f' : R/(f) \to R/(f)))
$$
Or nous observons que
$\Ker(f' : R/(f) \to R/(f))$ est canoniquement isomorphe
à $((f) \cap (f'))/(ff')$ par l'application $x \bmod (f) \mapsto
f'x \bmod (ff')$. Le membre de gauche vaut donc
$$
\text{length}_R(R/(f, f'))
-
\text{length}_R((f) \cap (f')/(ff'))
$$
Comme cette expression est symétrique en $f$ et $f'$, la démonstration est achevée.
\end{proof}

\begin{lemma}
\label{obsolete-lemma-commutativity-effective-Cartier-proper-intersection-infinite}
Soit $(S, \delta)$ comme dans Homologie de Chow, Situation \ref{chow-situation-setup}.
Soit $X$ localement de type fini sur $S$.
Supposons $X$ intègre et $\dim_\delta(X) = n$.
Soit $\{D_j\}_{j \in J}$ une famille localement finie de
diviseurs de Cartier effectifs sur $X$. Soient $n_j, m_j \geq 0$ des
familles d'entiers positifs ou nuls. Posons
$D = \sum n_j D_j$ et $D' = \sum m_j D_j$.
Supposons que $\dim_\delta(D_j \cap D_{j'}) = n - 2$ pour tous
$j \not = j'$. Alors $D \cdot [D']_{n - 1} = D' \cdot [D]_{n - 1}$ dans
$\CH_{n - 2}(X)$.
\end{lemma}

\begin{proof}
Ce lemme découle immédiatement des
Lemmes \ref{obsolete-lemma-sum-divisors-associated-Weil} et
\ref{obsolete-lemma-commutativity-effective-Cartier-proper-intersection}
lorsque les sommes sont finies, par exemple si $X$ est quasi-compact.
Le lecteur pourra donc omettre la démonstration.

\medskip\noindent
Voici la démonstration dans le cas général.
Soient $i_j : D_j \to X$ les immersions fermées
Soit $p : \coprod D_j \to X$ le coproduit des morphismes $i_j$.
Soit $\{Z_a\}_{a \in A}$ la famille des composantes irréductibles de
$\bigcup D_j$. Pour chaque $j$, écrivons
$$
[D_j]_{n - 1} = \sum d_{j, a}[Z_a].
$$
D'après le Lemme \ref{obsolete-lemma-sum-divisors-associated-Weil}, nous avons
$$
[D]_{n - 1} = \sum n_j d_{j, a} [Z_a],
\quad
[D']_{n - 1} = \sum m_j d_{j, a} [Z_a].
$$
D'après le Lemme \ref{obsolete-lemma-improved-additivity},
nous avons
$$
D \cdot [D']_{n - 1} = p_*\left(\sum n_j i_j^*[D']_{n - 1} \right),
\quad
D' \cdot [D]_{n - 1} = p_*\left(\sum m_{j'} i_{j'}^*[D]_{n - 1} \right).
$$
Comme dans la définition des homomorphismes de Gysin (voir
Homologie de Chow, Définition \ref{chow-definition-gysin-homomorphism}),
nous choisissons des cycles $\beta_{a, j}$ sur $D_j \cap Z_a$ représentant
$i_j^*[Z_a]$. (Notons qu'en fait $\beta_{a, j} = [D_j \cap Z_a]_{n - 2}$
si $Z_a$ n'est pas contenu dans $D_j$, c'est-à-dire qu'il n'y a aucun choix dans ce cas.)
Puisque la restriction de $p$ à chacun des $D_j$ est une immersion fermée,
nous pouvons (et nous allons) considérer $\beta_{a, j}$ comme un cycle sur $X$.
En substituant les formules pour $[D]_{n - 1}$ et $[D']_{n - 1}$ obtenues
ci-dessus, on voit que
$$
D \cdot [D']_{n - 1} =
\sum\nolimits_{j, j', a} n_j m_{j'} d_{j', a} \beta_{a, j},
\quad
D' \cdot [D]_{n - 1} =
\sum\nolimits_{j, j', a} m_{j'} n_j d_{j, a} \beta_{a, j'}.
$$
De plus, avec les mêmes conventions, on a aussi
$$
D_j \cdot [D_{j'}]_{n - 1} = \sum d_{j', a} \beta_{a, j}.
$$
Avec ces notations, le
lemme \ref{obsolete-lemma-commutativity-effective-Cartier-proper-intersection}
(voir aussi sa démonstration)
affirme que, pour $j \not = j'$, les cycles
$\sum d_{j', a} \beta_{a, j}$ et $\sum d_{j, a} \beta_{a, j'}$
sont égaux en tant que cycles ! On en déduit que
\begin{eqnarray*}
D \cdot [D']_{n - 1}
& = &
\sum\nolimits_{j, j', a} n_j m_{j'} d_{j', a} \beta_{a, j} \\
& = &
\sum\nolimits_{j \not = j'} n_j m_{j'}
\left(\sum\nolimits_a d_{j', a} \beta_{a, j}\right) +
\sum\nolimits_{j, a} n_j m_j d_{j, a} \beta_{a, j} \\
& = &
\sum\nolimits_{j \not = j'} n_j m_{j'}
\left(\sum\nolimits_a d_{j, a} \beta_{a, j'}\right) +
\sum\nolimits_{j, a} n_j m_j d_{j, a} \beta_{a, j} \\
& = &
\sum\nolimits_{j, j', a} m_{j'} n_j d_{j, a} \beta_{a, j'} \\
& = &
D' \cdot [D]_{n - 1}
\end{eqnarray*}
ce qui achève la démonstration.
\end{proof}

\begin{lemma}
\label{obsolete-lemma-commutativity-effective-Cartier}
Soit $(S, \delta)$ comme dans Homologie de Chow, situation \ref{chow-situation-setup}.
Soit $X$ localement de type fini sur $S$.
Supposons que $X$ soit intègre et que $\dim_\delta(X) = n$.
Soient $D$ et $D'$ des diviseurs de Cartier effectifs sur $X$.
Alors
$$
D \cdot [D']_{n - 1} = D' \cdot [D]_{n - 1}
$$
dans $\CH_{n - 2}(X)$.
\end{lemma}

\begin{proof}[Première démonstration du lemme \ref{obsolete-lemma-commutativity-effective-Cartier}]
Montrons d'abord l'assertion lorsque $X$ est quasi-compact.
Dans ce cas, appliquons
le lemme \ref{obsolete-lemma-blowing-up-intersections} à $X$ et à
l'ensemble à deux éléments $\{D, D'\}$ de diviseurs de Cartier effectifs.
On obtient ainsi un morphisme propre $b : X' \to X$,
une famille finie de diviseurs de Cartier
effectifs $D'_j \subset X'$ dont les intersections deux à deux sont de codimension $\geq 2$,
avec $b^{-1}(D) = \sum n_j D'_j$, et $b^{-1}(D') = \sum m_j D'_j$.
Notons que $b_*[b^{-1}(D)]_{n - 1} = [D]_{n - 1}$ dans $Z_{n - 1}(X)$
et de même pour $D'$,
voir le lemme \ref{obsolete-lemma-push-pull-effective-Cartier}.
Par conséquent, d'après Homologie de Chow, lemme \ref{chow-lemma-pushforward-cap-c1}, on a
$$
D \cdot [D']_{n - 1} = b_*\left(b^{-1}(D) \cdot [b^{-1}(D')]_{n - 1}\right)
$$
dans $\CH_{n - 2}(X)$ et de même pour l'autre terme. Le
lemme résulte donc de l'égalité
$b^{-1}(D) \cdot [b^{-1}(D')]_{n - 1} = b^{-1}(D') \cdot [b^{-1}(D)]_{n - 1}$
dans $\CH_{n - 2}(X')$ donnée par le lemme
\ref{obsolete-lemma-commutativity-effective-Cartier-proper-intersection-infinite}.

\medskip\noindent
Notons que, dans la démonstration ci-dessus, chacun des lemmes cités s'applique aussi
au cas général (où $X$ n'est pas supposé quasi-compact). La
seule modification mineure dans le cas général est que le morphisme
$b : U' \to U$ obtenu en appliquant
le lemme \ref{obsolete-lemma-blowing-up-intersections}
a pour but
un ouvert $U \subset X$ dont le complémentaire est de codimension $\geq 3$.
D'après Homologie de Chow, lemme \ref{chow-lemma-restrict-to-open}, on voit donc que
$\CH_{n - 2}(U) = \CH_{n - 2}(X)$
et, après avoir remplacé $X$ par $U$, le reste de la démonstration s'applique
sans changement.
\end{proof}

\begin{proof}[Seconde démonstration du lemme \ref{obsolete-lemma-commutativity-effective-Cartier}]
Soient $\mathcal{I} = \mathcal{O}_X(-D)$ et
$\mathcal{I}' = \mathcal{O}_X(-D')$ les faisceaux
d'idéaux inversibles de $D$ et de $D'$.
Posons
$\mathcal{I}_{D'} = \mathcal{I} \otimes_{\mathcal{O}_X} \mathcal{O}_{D'}$
et
$\mathcal{I}'_D = \mathcal{I}' \otimes_{\mathcal{O}_X} \mathcal{O}_D$.
On peut restreindre le morphisme d'inclusion $\mathcal{I} \to \mathcal{O}_X$
à $D'$ pour obtenir un morphisme
$$
\varphi : \mathcal{I}_{D'} \longrightarrow \mathcal{O}_{D'}
$$
et, de même,
$$
\psi : \mathcal{I}'_D \longrightarrow \mathcal{O}_D
$$
Il est clair que
$$
\Coker(\varphi)
\cong
\mathcal{O}_{D \cap D'}
\cong
\Coker(\psi)
$$
et que
$$
\Ker(\varphi)
\cong
\frac{\mathcal{I} \cap \mathcal{I}'}{\mathcal{I}\mathcal{I}'}
\cong
\Ker(\psi).
$$
On voit donc que
$$
\gamma =
[\mathcal{I}_{D'}] - [\mathcal{O}_{D'}]
=
[\mathcal{I}'_D] - [\mathcal{O}_D]
$$
dans $K_0(\textit{Coh}_{\leq n - 1}(X))$. D'autre part, il est clair que
$$
[\mathcal{I}'_D]_{n - 1} = [D]_{n - 1}, \quad
[\mathcal{I}_{D'}]_{n - 1} = [D']_{n - 1}.
$$
On a également
$$
\mathcal{O}_X(D') \otimes \mathcal{I}'_D = \mathcal{O}_D, \quad
\mathcal{O}_X(D) \otimes \mathcal{I}_{D'} = \mathcal{O}_{D'}.
$$
D'après Homologie de Chow, lemme \ref{chow-lemma-coherent-sheaf-cap-c1}
(appliqué deux fois),
cela signifie que l'élément $\gamma$ appartient à $B_{n - 2}(X)$ et
a pour image à la fois $c_1(\mathcal{O}_X(D')) \cap [D]_{n - 1}$ et
$c_1(\mathcal{O}_X(D)) \cap [D']_{n - 1}$, ce qui achève la démonstration (puisque
l'application $B_{n - 2}(X) \to \CH_{n - 2}(X)$ est bien définie -- ce qui constitue
le point clef de cette démonstration).
\end{proof}
      
 
    
     
  
   
 




\section{Modules dualisants sur les modèles propres réguliers}
\label{obsolete-section-dualizing}

\noindent
Dans Réduction semi-stable, situation
\ref{models-situation-regular-model}, on pose
$\omega_{X/R}^\bullet = f^!\mathcal{O}_{\Spec(R)}$,
le complexe dualisant relatif de $f : X \to \Spec(R)$ introduit dans
Dualité pour les schémas, remarque
\ref{duality-remark-relative-dualizing-complex}.
Comme $f$ est de Gorenstein de dimension relative $1$, d'après
Réduction semi-stable, lemme \ref{models-lemma-gorenstein}, les lemmes
\ref{duality-lemma-affine-flat-Noetherian-gorenstein},
\ref{duality-lemma-CM-shriek} et
\ref{duality-lemma-gorenstein-CM-morphism}
de Dualité pour les schémas donnent
$$
\omega_{X/R}^\bullet = \omega_X[1]
$$
pour un certain $\mathcal{O}_X$-module inversible $\omega_X$.
Ce module inversible est souvent appelé le
{\it module dualisant relatif de $X$ sur $R$}.
Comme $R$ est régulier (donc de Gorenstein) de dimension $1$, le complexe
$\omega_R^\bullet = R[1]$ est un complexe dualisant normalisé pour $R$.
Ainsi, $\omega_X = H^{-2}(f^!\omega_R^\bullet)$ et $\omega_X$ n'est pas
seulement un module dualisant relatif, mais aussi un module dualisant; voir
Dualité pour les schémas, exemple
\ref{duality-example-proper-over-local}.
Le module $\omega_X$ représente donc le foncteur
$$
\textit{Coh}(\mathcal{O}_X) \to \textit{Ens},\quad
\mathcal{F} \mapsto \Hom_R(H^1(X, \mathcal{F}), R)
$$
d'après Dualité pour les schémas, lemme
\ref{duality-lemma-dualizing-module-proper-over-A}.
On obtient ainsi une autre définition du module dualisant relatif dans
Réduction semi-stable, situation \ref{models-situation-regular-model}.
La formation de $\omega_X$ commute à tout changement de base
(pour tout morphisme propre de Gorenstein de dimension relative donnée);
cela résulte du fait correspondant pour le complexe dualisant relatif,
exposé dans Dualité pour les schémas, remarque
\ref{duality-remark-relative-dualizing-complex},
qui remonte à Dualité pour les schémas, lemme
\ref{duality-lemma-proper-flat-base-change}.
L'image réciproque de $\omega_X$ est donc le module dualisant $\omega_C$ de
$C$ sur $K$ étudié dans Courbes algébriques, lemme
\ref{curves-lemma-duality-dim-1-CM}.
Notons que $\omega_C$ est isomorphe à $\Omega_{C/K}$ d'après
Courbes algébriques, lemme \ref{curves-lemma-duality-dim-1}.
De même, $\omega_X|_{X_k}$ est le module dualisant $\omega_{X_k}$ de $X_k$
sur $k$.

\begin{lemma}
\label{obsolete-lemma-dualizing-components}
Dans Réduction semi-stable, situation
\ref{models-situation-regular-model}, le module dualisant de $C_i$ sur $k$
est
$$
\omega_{C_i} = \omega_X(C_i)|_{C_i}
$$
où $\omega_X$ est défini ci-dessus.
\end{lemma}

\begin{proof}
Soit $t : C_i \to X$ l'immersion fermée. Comme $t$ est l'inclusion d'un
diviseur de Cartier effectif, les lemmes
\ref{duality-lemma-twisted-inverse-image-closed} et
\ref{duality-lemma-sheaf-with-exact-support-effective-Cartier}
de Dualité pour les schémas donnent
$t^!(\mathcal{L}) = \mathcal{L}(C_i)|_{C_i}$ pour tout
$\mathcal{O}_X$-module inversible $\mathcal{L}$.
Considérons le diagramme commutatif
$$
\xymatrix{
C_i \ar[r]_t \ar[d]_g & X \ar[d]^f \\
\Spec(k) \ar[r]^s & \Spec(R)
}
$$
La courbe $C_i$ est de Gorenstein (Réduction semi-stable, lemme
\ref{models-lemma-gorenstein}) et son module dualisant inversible
$\omega_{C_i}$ est caractérisé par la propriété
$\omega_{C_i}[0] = g^!\mathcal{O}_{\Spec(k)}$. Voir Courbes algébriques,
lemme \ref{curves-lemma-duality-dim-1}, sa démonstration, ainsi que les lemmes
\ref{curves-lemma-duality-dim-1-CM} et \ref{curves-lemma-rr}
de Courbes algébriques.
D'autre part, $s^!(R[1]) = k$, donc
$$
\omega_{C_i}[0] =
g^! s^!(R[1]) = t^!f^!(R[1]) = t^!\omega_X
$$
En combinant ces observations, on obtient l'énoncé du lemme.
\end{proof}






\section{Renvois en double ou scindés}
\label{obsolete-section-duplicates}

\noindent
Cette section rassemble les doublons ainsi que les renvois qui énonçaient
autrefois plusieurs faits simultanément, mais qui sont désormais scindés en
leurs différentes composantes.

\begin{lemma}
\label{obsolete-lemma-characterize-epi-mono}
Soit $X$ un espace topologique.
\begin{enumerate}
\item Les épimorphismes (resp.\ les monomorphismes) dans la catégorie des
préfaisceaux sont exactement les morphismes surjectifs (resp.\ injectifs) de
préfaisceaux.
\item Les épimorphismes (resp.\ les monomorphismes) dans la catégorie des
faisceaux sont exactement les morphismes surjectifs (resp.\ injectifs) de
faisceaux, et exactement ceux qui sont surjectifs (resp.\ injectifs) sur
toutes les fibres.
\item Le morphisme de faisceaux associé à un morphisme surjectif
(resp.\ injectif) de préfaisceaux d'ensembles est surjectif
(resp.\ injectif).
\end{enumerate}
\end{lemma}

\begin{proof}
On combine les lemmes \ref{sheaves-lemma-characterize-epi-mono} et
\ref{sheaves-lemma-sheafify-epi-mono} de Faisceaux.
\end{proof}

\begin{lemma}
\label{obsolete-lemma-directed-colimit-finite-type}
Soit $X$ un schéma quasi-compact et quasi-séparé.
Soit $\mathcal{F}$ un $\mathcal{O}_X$-module quasi-cohérent.
Alors $\mathcal{F}$ est la limite inductive filtrante de ses sous-modules
quasi-cohérents de type fini.
\end{lemma}

\begin{proof}
C'est un doublon de Propriétés, lemme
\ref{properties-lemma-quasi-coherent-colimit-finite-type}.
\end{proof}

\begin{lemma}
\label{obsolete-lemma-points-monomorphism}
Soit $S$ un schéma et soit $X$ un espace algébrique sur $S$.
L'application $\{\Spec(k) \to X \text{ monomorphisme}\} \to |X|$ est injective.
\end{lemma}

\begin{proof}
C'est un doublon de Propriétés des espaces, lemme
\ref{spaces-properties-lemma-points-monomorphism}.
\end{proof}

\begin{theorem}
\label{obsolete-theorem-equivalence-sheaves-point}
Soit $S = \Spec(K)$, où $K$ est un corps.
Soit $\overline{s}$ un point géométrique de $S$.
Notons $G = \text{Gal}_{\kappa(s)}$ le groupe de Galois absolu.
Il existe alors une équivalence de catégories
$\Sh(S_\etale) \to G\textit{-Ens}$,
$\mathcal{F} \mapsto \mathcal{F}_{\overline{s}}$.
\end{theorem}

\begin{proof}
C'est un doublon de Cohomologie \'etale, théorème
\ref{etale-cohomology-theorem-equivalence-sheaves-point}.
\end{proof}

\begin{remark}
\label{obsolete-remark-tangent-spaces}
Ce renvoi mène ici en raison d'une étiquette dupliquée. Pour le contenu
proprement dit, voir Théorie formelle des déformations, section
\ref{formal-defos-section-tangent-spaces}.
\end{remark}

\begin{lemma}
\label{obsolete-lemma-locally-ringed-space-direct-summand-free}
Soit $X$ un espace localement annelé. Tout facteur direct d'un
$\mathcal{O}_X$-module libre de type fini est localement libre de type fini.
\end{lemma}

\begin{proof}
C'est un doublon de Modules, lemme
\ref{modules-lemma-direct-summand-of-locally-free-is-locally-free}.
\end{proof}

\begin{lemma}
\label{obsolete-lemma-characterize-injective}
Soit $R$ un anneau et soit $E$ un $R$-module. Les conditions suivantes sont
équivalentes :
\begin{enumerate}
\item $E$ est un $R$-module injectif;
\item pour tout idéal $I \subset R$ et tout morphisme de modules
$\varphi : I \to E$, il existe un prolongement de $\varphi$ en un morphisme
de $R$-modules $R \to E$.
\end{enumerate}
\end{lemma}

\begin{proof}
C'est le critère de Baer; voir Injectifs, lemme
\ref{injectives-lemma-criterion-baer}.
\end{proof}

\begin{lemma}
\label{obsolete-lemma-periodic-length}
Soit $R$ un anneau local.
\begin{enumerate}
\item Si $(M, N, \varphi, \psi)$ est un complexe $2$-périodique tel que
$M$ et $N$ soient de longueur finie, alors
$e_R(M, N, \varphi, \psi) = \text{longueur}_R(M) - \text{longueur}_R(N)$.
\item Si $(M, \varphi, \psi)$ est un complexe $(2, 1)$-périodique tel que
$M$ soit de longueur finie, alors
$e_R(M, \varphi, \psi) = 0$.
\item Supposons donnée une suite exacte de complexes $2$-périodiques
$$
0 \to (M_1, N_1, \varphi_1, \psi_1)
\to (M_2, N_2, \varphi_2, \psi_2)
\to (M_3, N_3, \varphi_3, \psi_3)
\to 0
$$
Si deux de ces trois complexes ont des modules de cohomologie de longueur
finie, il en va de même du troisième, et l'on a
$$
e_R(M_2, N_2, \varphi_2, \psi_2) =
e_R(M_1, N_1, \varphi_1, \psi_1) +
e_R(M_3, N_3, \varphi_3, \psi_3).
$$
\end{enumerate}
\end{lemma}

\begin{proof}
Cela résulte des lemmes \ref{chow-lemma-additivity-periodic-length} et
\ref{chow-lemma-finite-periodic-length} d'Homologie de Chow.
\end{proof}




\begin{lemma}
\label{obsolete-lemma-extensions-of-rings}
Soient $A$ un anneau et $I$ un $A$-module.
\begin{enumerate}
\item L'ensemble des extensions d'anneaux $0 \to I \to A' \to A \to 0$
où $I$ est un idéal de carré nul est en bijection canonique avec
$\Ext^1_A(\NL_{A/\mathbf{Z}}, I)$.
\item Étant donnés un homomorphisme d'anneaux $A \to B$, un $B$-module $N$,
un homomorphisme de $A$-modules $c : I \to N$, ainsi que les extensions
d'anneaux suivantes, à noyaux de carré nul :
\begin{enumerate}
\item[(a)] $0 \to I \to A' \to A \to 0$ correspondant à
$\alpha \in \Ext^1_A(\NL_{A/\mathbf{Z}}, I)$, et
\item[(b)] $0 \to N \to B' \to B \to 0$ correspondant à
$\beta \in \Ext^1_B(\NL_{B/\mathbf{Z}}, N)$,
\end{enumerate}
alors il existe un homomorphisme $A' \to B'$ s'inscrivant dans
le diagramme de Théorie des déformations, équation (\ref{defos-equation-to-solve}),
si et seulement si $\beta$ et $\alpha$
ont la même image dans
$\Ext^1_A(\NL_{A/\mathbf{Z}}, N)$.
\end{enumerate}
\end{lemma}

\begin{proof}
Cela résulte de Théorie des déformations, lemmes
\ref{defos-lemma-extensions-of-algebras} et
\ref{defos-lemma-extensions-of-algebras-functorial}.
\end{proof}

\begin{lemma}
\label{obsolete-lemma-extensions-of-ringed-spaces}
Soient $(S, \mathcal{O}_S)$ un espace annelé et $\mathcal{J}$
un $\mathcal{O}_S$-module.
\begin{enumerate}
\item L'ensemble des extensions de faisceaux d'anneaux
$0 \to \mathcal{J} \to \mathcal{O}_{S'} \to \mathcal{O}_S \to 0$
où $\mathcal{J}$ est un idéal de carré nul est en bijection canonique avec
$\Ext^1_{\mathcal{O}_S}(\NL_{S/\mathbf{Z}}, \mathcal{J})$.
\item Étant donnés un morphisme d'espaces annelés
$f : (X, \mathcal{O}_X) \to (S, \mathcal{O}_S)$, un $\mathcal{O}_X$-module
$\mathcal{G}$, un $f$-morphisme $c : \mathcal{J} \to \mathcal{G}$, ainsi que
les extensions suivantes de faisceaux d'anneaux, à noyaux de carré nul :
\begin{enumerate}
\item[(a)] $0 \to \mathcal{J} \to \mathcal{O}_{S'} \to \mathcal{O}_S \to 0$
correspondant à
$\alpha \in \Ext^1_{\mathcal{O}_S}(\NL_{S/\mathbf{Z}}, \mathcal{J})$,
\item[(b)] $0 \to \mathcal{G} \to \mathcal{O}_{X'} \to \mathcal{O}_X \to 0$
correspondant à
$\beta \in \Ext^1_{\mathcal{O}_X}(\NL_{X/\mathbf{Z}}, \mathcal{G})$,
\end{enumerate}
alors il existe un morphisme $X' \to S'$ s'inscrivant dans
le diagramme de Théorie des déformations, équation
(\ref{defos-equation-to-solve-ringed-spaces}),
si et seulement si $\beta$ et $\alpha$
ont la même image dans
$\Ext^1_{\mathcal{O}_X}(Lf^*\NL_{S/\mathbf{Z}}, \mathcal{G})$.
\end{enumerate}
\end{lemma}

\begin{proof}
Cela résulte de Théorie des déformations, lemmes
\ref{defos-lemma-extensions-of-relative-ringed-spaces} et
\ref{defos-lemma-extensions-of-relative-ringed-spaces-functorial}.
\end{proof}

\begin{lemma}
\label{obsolete-lemma-extensions-of-ringed-topoi}
Soient $(\Sh(\mathcal{B}), \mathcal{O}_\mathcal{B})$ un topos annelé
et $\mathcal{J}$ un $\mathcal{O}_\mathcal{B}$-module.
\begin{enumerate}
\item L'ensemble des extensions de faisceaux d'anneaux
$0 \to \mathcal{J} \to \mathcal{O}_{\mathcal{B}'} \to
\mathcal{O}_\mathcal{B} \to 0$
où $\mathcal{J}$ est un idéal de carré nul est en bijection canonique avec
$\Ext^1_{\mathcal{O}_\mathcal{B}}(
\NL_{\mathcal{O}_\mathcal{B}/\mathbf{Z}}, \mathcal{J})$.
\item Étant donnés un morphisme de topos annelés
$f : (\Sh(\mathcal{C}), \mathcal{O}) \to
(\Sh(\mathcal{B}), \mathcal{O}_\mathcal{B})$, un $\mathcal{O}$-module
$\mathcal{G}$, un homomorphisme de $f^{-1}\mathcal{O}_\mathcal{B}$-modules
$c : f^{-1}\mathcal{J} \to \mathcal{G}$, ainsi que les extensions suivantes
de faisceaux d'anneaux, à noyaux de carré nul :
\begin{enumerate}
\item[(a)] $0 \to \mathcal{J} \to \mathcal{O}_{\mathcal{B}'} \to
\mathcal{O}_\mathcal{B} \to 0$ correspondant à
$\alpha \in \Ext^1_{\mathcal{O}_\mathcal{B}}(
\NL_{\mathcal{O}_\mathcal{B}/\mathbf{Z}}, \mathcal{J})$,
\item[(b)] $0 \to \mathcal{G} \to \mathcal{O}' \to \mathcal{O} \to 0$
correspondant à
$\beta \in \Ext^1_\mathcal{O}(\NL_{\mathcal{O}/\mathbf{Z}}, \mathcal{G})$,
\end{enumerate}
alors il existe un morphisme $(\Sh(\mathcal{C}), \mathcal{O}') \to
(\Sh(\mathcal{B}), \mathcal{O}_{\mathcal{B}'})$ s'inscrivant dans
le diagramme de Théorie des déformations, équation (\ref{defos-equation-to-solve-ringed-topoi}),
si et seulement si $\beta$ et $\alpha$ ont la même image dans
$\Ext^1_\mathcal{O}(
Lf^*\NL_{\mathcal{O}_\mathcal{B}/\mathbf{Z}}, \mathcal{G})$.
\end{enumerate}
\end{lemma}

\begin{proof}
Cela résulte de Théorie des déformations, lemmes
\ref{defos-lemma-extensions-of-relative-ringed-topoi} et
\ref{defos-lemma-extensions-of-relative-ringed-topoi-functorial}.
\end{proof}

\begin{remark}
\label{obsolete-remark-examples-formal-defos}
Cette étiquette renvoyait autrefois à une section décrivant plusieurs exemples
de problèmes de déformation. Ces exemples ont désormais chacun leur propre section.
Voir Problèmes de déformation, sections
\ref{examples-defos-section-finite-projective-modules},
\ref{examples-defos-section-representations},
\ref{examples-defos-section-continuous-representations}, et
\ref{examples-defos-section-graded-algebras}.
\end{remark}

\begin{lemma}
\label{obsolete-lemma-examples-have-RS}
Les exemples de Problèmes de déformation
\ref{examples-defos-example-finite-projective-modules},
\ref{examples-defos-example-representations},
\ref{examples-defos-example-continuous-representations}, et
\ref{examples-defos-example-graded-algebras}
satisfont à la condition de Rim--Schlessinger (RS).
\end{lemma}

\begin{proof}
Cela résulte de Problèmes de déformation, lemmes
\ref{examples-defos-lemma-finite-projective-modules-RS},
\ref{examples-defos-lemma-representations-RS},
\ref{examples-defos-lemma-continuous-representations-RS}, et
\ref{examples-defos-lemma-graded-algebras-RS}.
\end{proof}

\begin{lemma}
\label{obsolete-lemma-tangent-and-inf}
On a les identifications canoniques suivantes d'espaces vectoriels sur $k$ :
\begin{enumerate}
\item Dans Problèmes de déformation, exemple
\ref{examples-defos-example-finite-projective-modules},
si $x_0 = (k, V)$, alors $T_{x_0}\mathcal{F} = (0)$
et $\text{Inf}_{x_0}(\mathcal{F}) = \text{End}_k(V)$
sont de dimension finie.
\item Dans Problèmes de déformation, exemple
\ref{examples-defos-example-representations},
si $x_0 = (k, V, \rho_0)$, alors
$T_{x_0}\mathcal{F} = \Ext^1_{k[\Gamma]}(V, V) = H^1(\Gamma, \text{End}_k(V))$
et $\text{Inf}_{x_0}(\mathcal{F}) = H^0(\Gamma, \text{End}_k(V))$
sont de dimension finie si $\Gamma$ est engendré par un nombre fini d'éléments.
\item Dans Problèmes de déformation, exemple
\ref{examples-defos-example-continuous-representations},
si $x_0 = (k, V, \rho_0)$, alors
$T_{x_0}\mathcal{F} = H^1_{cont}(\Gamma, \text{End}_k(V))$
et
$\text{Inf}_{x_0}(\mathcal{F}) = H^0_{cont}(\Gamma, \text{End}_k(V))$
sont de dimension finie si $\Gamma$ est topologiquement de type fini.
\item Dans Problèmes de déformation, exemple
\ref{examples-defos-example-graded-algebras},
si $x_0 = (k, P)$, alors
$T_{x_0}\mathcal{F}$ et $\text{Inf}_{x_0}(\mathcal{F}) = \text{Der}_k(P, P)$
sont de dimension finie si $P$ est une $k$-algèbre de type fini.
\end{enumerate}
\end{lemma}

\begin{proof}
Cela résulte de Problèmes de déformation, lemmes
\ref{examples-defos-lemma-finite-projective-modules-TI},
\ref{examples-defos-lemma-representations-TI},
\ref{examples-defos-lemma-continuous-representations-TI}, et
\ref{examples-defos-lemma-graded-algebras-TI}.
\end{proof}










% OK, start here.
%
%---------------------------------------------------------------------
%\tableofcontents

\chapter{GNU Free Documentation License}

\phantomsection
\label{fdl-section-phantom}

\begin{center}

      Version 1.2, November 2002


Copyright \copyright 2000, 2001, 2002  Free Software Foundation, Inc.

\bigskip

    51 Franklin St, Fifth Floor, Boston, MA  02110-1301  USA

\bigskip

Everyone is permitted to copy and distribute verbatim copies
of this license document, but changing it is not allowed.
\end{center}


\begin{center}
{\bf\large Preamble}
\end{center}

The purpose of this License is to make a manual, textbook, or other
functional and useful document "free" in the sense of freedom: to
assure everyone the effective freedom to copy and redistribute it,
with or without modifying it, either commercially or noncommercially.
Secondarily, this License preserves for the author and publisher a way
to get credit for their work, while not being considered responsible
for modifications made by others.

This License is a kind of "copyleft", which means that derivative
works of the document must themselves be free in the same sense.  It
complements the GNU General Public License, which is a copyleft
license designed for free software.

We have designed this License in order to use it for manuals for free
software, because free software needs free documentation: a free
program should come with manuals providing the same freedoms that the
software does.  But this License is not limited to software manuals;
it can be used for any textual work, regardless of subject matter or
whether it is published as a printed book.  We recommend this License
principally for works whose purpose is instruction or reference.


\section{APPLICABILITY AND DEFINITIONS}
\label{fdl-section-applicability-and-definitions}

This License applies to any manual or other work, in any medium, that
contains a notice placed by the copyright holder saying it can be
distributed under the terms of this License.  Such a notice grants a
world-wide, royalty-free license, unlimited in duration, to use that
work under the conditions stated herein.  The \textbf{"Document"}, below,
refers to any such manual or work.  Any member of the public is a
licensee, and is addressed as \textbf{"you"}.  You accept the license if you
copy, modify or distribute the work in a way requiring permission
under copyright law.

A \textbf{"Modified Version"} of the Document means any work containing the
Document or a portion of it, either copied verbatim, or with
modifications and/or translated into another language.

A \textbf{"Secondary Section"} is a named appendix or a front-matter section of
the Document that deals exclusively with the relationship of the
publishers or authors of the Document to the Document's overall subject
(or to related matters) and contains nothing that could fall directly
within that overall subject.  (Thus, if the Document is in part a
textbook of mathematics, a Secondary Section may not explain any
mathematics.)  The relationship could be a matter of historical
connection with the subject or with related matters, or of legal,
commercial, philosophical, ethical or political position regarding
them.

The \textbf{"Invariant Sections"} are certain Secondary Sections whose titles
are designated, as being those of Invariant Sections, in the notice
that says that the Document is released under this License.  If a
section does not fit the above definition of Secondary then it is not
allowed to be designated as Invariant.  The Document may contain zero
Invariant Sections.  If the Document does not identify any Invariant
Sections then there are none.

The \textbf{"Cover Texts"} are certain short passages of text that are listed,
as Front-Cover Texts or Back-Cover Texts, in the notice that says that
the Document is released under this License.  A Front-Cover Text may
be at most 5 words, and a Back-Cover Text may be at most 25 words.

A \textbf{"Transparent"} copy of the Document means a machine-readable copy,
represented in a format whose specification is available to the
general public, that is suitable for revising the document
straightforwardly with generic text editors or (for images composed of
pixels) generic paint programs or (for drawings) some widely available
drawing editor, and that is suitable for input to text formatters or
for automatic translation to a variety of formats suitable for input
to text formatters.  A copy made in an otherwise Transparent file
format whose markup, or absence of markup, has been arranged to thwart
or discourage subsequent modification by readers is not Transparent.
An image format is not Transparent if used for any substantial amount
of text.  A copy that is not "Transparent" is called \textbf{"Opaque"}.

Examples of suitable formats for Transparent copies include plain
ASCII without markup, Texinfo input format, LaTeX input format, SGML
or XML using a publicly available DTD, and standard-conforming simple
HTML, PostScript or PDF designed for human modification.  Examples of
transparent image formats include PNG, XCF and JPG.  Opaque formats
include proprietary formats that can be read and edited only by
proprietary word processors, SGML or XML for which the DTD and/or
processing tools are not generally available, and the
machine-generated HTML, PostScript or PDF produced by some word
processors for output purposes only.

The \textbf{"Title Page"} means, for a printed book, the title page itself,
plus such following pages as are needed to hold, legibly, the material
this License requires to appear in the title page.  For works in
formats which do not have any title page as such, "Title Page" means
the text near the most prominent appearance of the work's title,
preceding the beginning of the body of the text.

A section \textbf{"Entitled XYZ"} means a named subunit of the Document whose
title either is precisely XYZ or contains XYZ in parentheses following
text that translates XYZ in another language.  (Here XYZ stands for a
specific section name mentioned below, such as \textbf{"Acknowledgements"},
\textbf{"Dedications"}, \textbf{"Endorsements"}, or \textbf{"History"}.)
To \textbf{"Preserve the Title"}
of such a section when you modify the Document means that it remains a
section "Entitled XYZ" according to this definition.

The Document may include Warranty Disclaimers next to the notice which
states that this License applies to the Document.  These Warranty
Disclaimers are considered to be included by reference in this
License, but only as regards disclaiming warranties: any other
implication that these Warranty Disclaimers may have is void and has
no effect on the meaning of this License.


\section{VERBATIM COPYING}
\label{fdl-section-verbatim-copying}

You may copy and distribute the Document in any medium, either
commercially or noncommercially, provided that this License, the
copyright notices, and the license notice saying this License applies
to the Document are reproduced in all copies, and that you add no other
conditions whatsoever to those of this License.  You may not use
technical measures to obstruct or control the reading or further
copying of the copies you make or distribute.  However, you may accept
compensation in exchange for copies.  If you distribute a large enough
number of copies you must also follow the conditions in section 3.

You may also lend copies, under the same conditions stated above, and
you may publicly display copies.


\section{COPYING IN QUANTITY}
\label{fdl-section-copying-in-quantity}


If you publish printed copies (or copies in media that commonly have
printed covers) of the Document, numbering more than 100, and the
Document's license notice requires Cover Texts, you must enclose the
copies in covers that carry, clearly and legibly, all these Cover
Texts: Front-Cover Texts on the front cover, and Back-Cover Texts on
the back cover.  Both covers must also clearly and legibly identify
you as the publisher of these copies.  The front cover must present
the full title with all words of the title equally prominent and
visible.  You may add other material on the covers in addition.
Copying with changes limited to the covers, as long as they preserve
the title of the Document and satisfy these conditions, can be treated
as verbatim copying in other respects.

If the required texts for either cover are too voluminous to fit
legibly, you should put the first ones listed (as many as fit
reasonably) on the actual cover, and continue the rest onto adjacent
pages.

If you publish or distribute Opaque copies of the Document numbering
more than 100, you must either include a machine-readable Transparent
copy along with each Opaque copy, or state in or with each Opaque copy
a computer-network location from which the general network-using
public has access to download using public-standard network protocols
a complete Transparent copy of the Document, free of added material.
If you use the latter option, you must take reasonably prudent steps,
when you begin distribution of Opaque copies in quantity, to ensure
that this Transparent copy will remain thus accessible at the stated
location until at least one year after the last time you distribute an
Opaque copy (directly or through your agents or retailers) of that
edition to the public.

It is requested, but not required, that you contact the authors of the
Document well before redistributing any large number of copies, to give
them a chance to provide you with an updated version of the Document.


\section{MODIFICATIONS}
\label{fdl-section-modifications}

You may copy and distribute a Modified Version of the Document under
the conditions of sections 2 and 3 above, provided that you release
the Modified Version under precisely this License, with the Modified
Version filling the role of the Document, thus licensing distribution
and modification of the Modified Version to whoever possesses a copy
of it.  In addition, you must do these things in the Modified Version:

\begin{itemize}
\item[A.]
   Use in the Title Page (and on the covers, if any) a title distinct
   from that of the Document, and from those of previous versions
   (which should, if there were any, be listed in the History section
   of the Document).  You may use the same title as a previous version
   if the original publisher of that version gives permission.

\item[B.]
   List on the Title Page, as authors, one or more persons or entities
   responsible for authorship of the modifications in the Modified
   Version, together with at least five of the principal authors of the
   Document (all of its principal authors, if it has fewer than five),
   unless they release you from this requirement.

\item[C.]
   State on the Title page the name of the publisher of the
   Modified Version, as the publisher.

\item[D.]
   Preserve all the copyright notices of the Document.

\item[E.]
   Add an appropriate copyright notice for your modifications
   adjacent to the other copyright notices.

\item[F.]
   Include, immediately after the copyright notices, a license notice
   giving the public permission to use the Modified Version under the
   terms of this License, in the form shown in the Addendum below.

\item[G.]
   Preserve in that license notice the full lists of Invariant Sections
   and required Cover Texts given in the Document's license notice.

\item[H.]
   Include an unaltered copy of this License.

\item[I.]
   Preserve the section Entitled "History", Preserve its Title, and add
   to it an item stating at least the title, year, new authors, and
   publisher of the Modified Version as given on the Title Page.  If
   there is no section Entitled "History" in the Document, create one
   stating the title, year, authors, and publisher of the Document as
   given on its Title Page, then add an item describing the Modified
   Version as stated in the previous sentence.

\item[J.]
   Preserve the network location, if any, given in the Document for
   public access to a Transparent copy of the Document, and likewise
   the network locations given in the Document for previous versions
   it was based on.  These may be placed in the "History" section.
   You may omit a network location for a work that was published at
   least four years before the Document itself, or if the original
   publisher of the version it refers to gives permission.

\item[K.]
   For any section Entitled "Acknowledgements" or "Dedications",
   Preserve the Title of the section, and preserve in the section all
   the substance and tone of each of the contributor acknowledgements
   and/or dedications given therein.

\item[L.]
   Preserve all the Invariant Sections of the Document,
   unaltered in their text and in their titles.  Section numbers
   or the equivalent are not considered part of the section titles.

\item[M.]
   Delete any section Entitled "Endorsements".  Such a section
   may not be included in the Modified Version.

\item[N.]
   Do not retitle any existing section to be Entitled "Endorsements"
   or to conflict in title with any Invariant Section.

\item[O.]
   Preserve any Warranty Disclaimers.
\end{itemize}

If the Modified Version includes new front-matter sections or
appendices that qualify as Secondary Sections and contain no material
copied from the Document, you may at your option designate some or all
of these sections as invariant.  To do this, add their titles to the
list of Invariant Sections in the Modified Version's license notice.
These titles must be distinct from any other section titles.

You may add a section Entitled "Endorsements", provided it contains
nothing but endorsements of your Modified Version by various
parties--for example, statements of peer review or that the text has
been approved by an organization as the authoritative definition of a
standard.

You may add a passage of up to five words as a Front-Cover Text, and a
passage of up to 25 words as a Back-Cover Text, to the end of the list
of Cover Texts in the Modified Version.  Only one passage of
Front-Cover Text and one of Back-Cover Text may be added by (or
through arrangements made by) any one entity.  If the Document already
includes a cover text for the same cover, previously added by you or
by arrangement made by the same entity you are acting on behalf of,
you may not add another; but you may replace the old one, on explicit
permission from the previous publisher that added the old one.

The author(s) and publisher(s) of the Document do not by this License
give permission to use their names for publicity for or to assert or
imply endorsement of any Modified Version.


\section{COMBINING DOCUMENTS}
\label{fdl-section-combining-documents}


You may combine the Document with other documents released under this
License, under the terms defined in section 4 above for modified
versions, provided that you include in the combination all of the
Invariant Sections of all of the original documents, unmodified, and
list them all as Invariant Sections of your combined work in its
license notice, and that you preserve all their Warranty Disclaimers.

The combined work need only contain one copy of this License, and
multiple identical Invariant Sections may be replaced with a single
copy.  If there are multiple Invariant Sections with the same name but
different contents, make the title of each such section unique by
adding at the end of it, in parentheses, the name of the original
author or publisher of that section if known, or else a unique number.
Make the same adjustment to the section titles in the list of
Invariant Sections in the license notice of the combined work.

In the combination, you must combine any sections Entitled "History"
in the various original documents, forming one section Entitled
"History"; likewise combine any sections Entitled "Acknowledgements",
and any sections Entitled "Dedications".  You must delete all sections
Entitled "Endorsements".

\section{COLLECTIONS OF DOCUMENTS}
\label{fdl-section-collections-of-documents}

You may make a collection consisting of the Document and other documents
released under this License, and replace the individual copies of this
License in the various documents with a single copy that is included in
the collection, provided that you follow the rules of this License for
verbatim copying of each of the documents in all other respects.

You may extract a single document from such a collection, and distribute
it individually under this License, provided you insert a copy of this
License into the extracted document, and follow this License in all
other respects regarding verbatim copying of that document.


\section{AGGREGATION WITH INDEPENDENT WORKS}
\label{fdl-section-aggregation-with-independent-works}


A compilation of the Document or its derivatives with other separate
and independent documents or works, in or on a volume of a storage or
distribution medium, is called an "aggregate" if the copyright
resulting from the compilation is not used to limit the legal rights
of the compilation's users beyond what the individual works permit.
When the Document is included in an aggregate, this License does not
apply to the other works in the aggregate which are not themselves
derivative works of the Document.

If the Cover Text requirement of section 3 is applicable to these
copies of the Document, then if the Document is less than one half of
the entire aggregate, the Document's Cover Texts may be placed on
covers that bracket the Document within the aggregate, or the
electronic equivalent of covers if the Document is in electronic form.
Otherwise they must appear on printed covers that bracket the whole
aggregate.


\section{TRANSLATION}
\label{fdl-section-translation}


Translation is considered a kind of modification, so you may
distribute translations of the Document under the terms of section 4.
Replacing Invariant Sections with translations requires special
permission from their copyright holders, but you may include
translations of some or all Invariant Sections in addition to the
original versions of these Invariant Sections.  You may include a
translation of this License, and all the license notices in the
Document, and any Warranty Disclaimers, provided that you also include
the original English version of this License and the original versions
of those notices and disclaimers.  In case of a disagreement between
the translation and the original version of this License or a notice
or disclaimer, the original version will prevail.

If a section in the Document is Entitled "Acknowledgements",
"Dedications", or "History", the requirement (section 4) to Preserve
its Title (section 1) will typically require changing the actual
title.


\section{TERMINATION}
\label{fdl-section-termination}


You may not copy, modify, sublicense, or distribute the Document except
as expressly provided for under this License.  Any other attempt to
copy, modify, sublicense or distribute the Document is void, and will
automatically terminate your rights under this License.  However,
parties who have received copies, or rights, from you under this
License will not have their licenses terminated so long as such
parties remain in full compliance.


\section{FUTURE REVISIONS OF THIS LICENSE}
\label{fdl-section-future-revisions-of-this-license}


The Free Software Foundation may publish new, revised versions
of the GNU Free Documentation License from time to time.  Such new
versions will be similar in spirit to the present version, but may
differ in detail to address new problems or concerns.  See
https://www.gnu.org/copyleft/.

Each version of the License is given a distinguishing version number.
If the Document specifies that a particular numbered version of this
License "or any later version" applies to it, you have the option of
following the terms and conditions either of that specified version or
of any later version that has been published (not as a draft) by the
Free Software Foundation.  If the Document does not specify a version
number of this License, you may choose any version ever published (not
as a draft) by the Free Software Foundation.


\section{ADDENDUM: How to use this License for your documents}
\label{fdl-section-addendum-how-to-use-this-license-for-your-documents}

To use this License in a document you have written, include a copy of
the License in the document and put the following copyright and
license notices just after the title page:

\bigskip
\begin{quote}
    Copyright \copyright  YEAR  YOUR NAME.
    Permission is granted to copy, distribute and/or modify this document
    under the terms of the GNU Free Documentation License, Version 1.2
    or any later version published by the Free Software Foundation;
    with no Invariant Sections, no Front-Cover Texts, and no Back-Cover Texts.
    A copy of the license is included in the section entitled "GNU
    Free Documentation License".
\end{quote}
\bigskip

If you have Invariant Sections, Front-Cover Texts and Back-Cover Texts,
replace the "with...Texts." line with this:

\bigskip
\begin{quote}
    with the Invariant Sections being LIST THEIR TITLES, with the
    Front-Cover Texts being LIST, and with the Back-Cover Texts being LIST.
\end{quote}
\bigskip

If you have Invariant Sections without Cover Texts, or some other
combination of the three, merge those two alternatives to suit the
situation.

If your document contains nontrivial examples of program code, we
recommend releasing these examples in parallel under your choice of
free software license, such as the GNU General Public License,
to permit their use in free software.

%---------------------------------------------------------------------

\bibliography{my}
\bibliographystyle{amsalpha}
\end{document}
